% =====================================================================
%  Analog Game Design, lezione 2 — I casi studio: una panoramica a volo
%  d'aquila dei giochi pervasivi, dei giochi di ruolo dal vivo, dei
%  giochi di ruolo da tavolo e dei giochi da tavolo
%  Fonte: trascrizione della seconda lezione (AGD2) del corso breve di
%  Analog Game Design, seconda edizione: una lezione online su Zoom di
%  circa tre ore, di venerdì sera, tenuta a più voci dai tre docenti
%  "cuore" del corso (Ruggero, Daniele e Beatrice), con una pausa di
%  pochi minuti e moltissime slide di esempi.
%  Pacchetti richiesti nel preambolo del libro (già presenti nel
%  main.tex): \usepackage{booktabs}, \usepackage{tikz} con
%    \usetikzlibrary{shapes.geometric,shapes.symbols,shapes.misc,
%      arrows.meta,positioning,calc,fit}
%  (le librerie vanno caricate nel preambolo, non nel capitolo).
%  Consigliati, come per gli altri capitoli: \usepackage[italian]{babel}
%  e \usepackage[T1]{fontenc} (il capitolo compila anche senza).
%
%  Collocazione nel libro: parte "Analog Game Design", secondo capitolo
%  (numerazione propria 2.2). Nel file principale del libro, subito dopo
%  il capitolo 2.1:
%
%    % =====================================================================
%  Analog Game Design, lezione 1 — Introduzione al corso: docenti,
%  metodo, project work, figure dell'industria e primo glossario
%  Fonte: trascrizione della prima lezione (AGD1) del corso breve di
%  Analog Game Design, seconda edizione: una lezione online su Zoom di
%  circa tre ore, di pomeriggio, tenuta insieme dai tre docenti "cuore"
%  del corso (Daniele, Ruggero e Beatrice), con nove iscritti, una
%  pausa di una decina di minuti, molte slide e una breve dimostrazione
%  dal vivo di Tabletop Simulator.
%  Pacchetti richiesti nel preambolo del libro (già presenti nel
%  main.tex): \usepackage{booktabs}, \usepackage{tikz} con
%    \usetikzlibrary{shapes.geometric,shapes.symbols,shapes.misc,
%      arrows.meta,positioning,calc,fit}
%  (le librerie vanno caricate nel preambolo, non nel capitolo).
%  Consigliati, come per gli altri capitoli: \usepackage[italian]{babel}
%  e \usepackage[T1]{fontenc} (il capitolo compila anche senza).
%
%  Collocazione nel libro: parte "Analog Game Design", primo capitolo
%  (numerazione propria 2.1). Nel file principale del libro:
%
%    \part{Analog Game Design}
%    \renewcommand{\thechapter}{2.\arabic{chapter}}
%    \setcounter{chapter}{0}
%    % =====================================================================
%  Analog Game Design, lezione 1 — Introduzione al corso: docenti,
%  metodo, project work, figure dell'industria e primo glossario
%  Fonte: trascrizione della prima lezione (AGD1) del corso breve di
%  Analog Game Design, seconda edizione: una lezione online su Zoom di
%  circa tre ore, di pomeriggio, tenuta insieme dai tre docenti "cuore"
%  del corso (Daniele, Ruggero e Beatrice), con nove iscritti, una
%  pausa di una decina di minuti, molte slide e una breve dimostrazione
%  dal vivo di Tabletop Simulator.
%  Pacchetti richiesti nel preambolo del libro (già presenti nel
%  main.tex): \usepackage{booktabs}, \usepackage{tikz} con
%    \usetikzlibrary{shapes.geometric,shapes.symbols,shapes.misc,
%      arrows.meta,positioning,calc,fit}
%  (le librerie vanno caricate nel preambolo, non nel capitolo).
%  Consigliati, come per gli altri capitoli: \usepackage[italian]{babel}
%  e \usepackage[T1]{fontenc} (il capitolo compila anche senza).
%
%  Collocazione nel libro: parte "Analog Game Design", primo capitolo
%  (numerazione propria 2.1). Nel file principale del libro:
%
%    \part{Analog Game Design}
%    \renewcommand{\thechapter}{2.\arabic{chapter}}
%    \setcounter{chapter}{0}
%    % =====================================================================
%  Analog Game Design, lezione 1 — Introduzione al corso: docenti,
%  metodo, project work, figure dell'industria e primo glossario
%  Fonte: trascrizione della prima lezione (AGD1) del corso breve di
%  Analog Game Design, seconda edizione: una lezione online su Zoom di
%  circa tre ore, di pomeriggio, tenuta insieme dai tre docenti "cuore"
%  del corso (Daniele, Ruggero e Beatrice), con nove iscritti, una
%  pausa di una decina di minuti, molte slide e una breve dimostrazione
%  dal vivo di Tabletop Simulator.
%  Pacchetti richiesti nel preambolo del libro (già presenti nel
%  main.tex): \usepackage{booktabs}, \usepackage{tikz} con
%    \usetikzlibrary{shapes.geometric,shapes.symbols,shapes.misc,
%      arrows.meta,positioning,calc,fit}
%  (le librerie vanno caricate nel preambolo, non nel capitolo).
%  Consigliati, come per gli altri capitoli: \usepackage[italian]{babel}
%  e \usepackage[T1]{fontenc} (il capitolo compila anche senza).
%
%  Collocazione nel libro: parte "Analog Game Design", primo capitolo
%  (numerazione propria 2.1). Nel file principale del libro:
%
%    \part{Analog Game Design}
%    \renewcommand{\thechapter}{2.\arabic{chapter}}
%    \setcounter{chapter}{0}
%    \input{AGD/AGD1_introduzione_all_analog_game_design}
%
%  (Con hyperref conviene aggiungere anche
%  \renewcommand{\theHchapter}{agd.\arabic{chapter}} nella parte, per
%  evitare ancore duplicate con i capitoli 1, 2, ... di Game Design.)
%  Tutte le etichette di questo capitolo hanno il prefisso "a1:"
%  (a1:cap:..., a1:sec:..., a1:fig:..., a1:tab:...). Il capitolo è
%  autonomo: non contiene \ref verso altri capitoli.
%  Le figure sono ricostruzioni schematiche di quanto i docenti spiegano
%  a voce e con le slide, non copie delle slide.
% =====================================================================

\chapter[Introduzione all'Analog Game Design]{Introduzione all'Analog Game
Design: il corso, le figure dell'industria e un primo linguaggio comune}
\label{a1:cap:intro}

Questa è la \textbf{prima lezione} del corso di \textbf{Analog Game Design},
cioè di progettazione di giochi \emph{analogici}: giochi da tavolo, giochi di
carte, giochi di ruolo da tavolo, giochi di ruolo dal vivo (LARP), giochi
\emph{pervasivi} ed esperienze a evento. È un corso \textbf{breve}, di circa tre
mesi (``da qui a dicembre''), separato dal corso di Game Design, e siamo alla sua
\textbf{seconda edizione}. Come avvertono subito i docenti, è ``una giornata di
introduzione'': si spiega come funziona il corso, chi sono i docenti, con quali
metodi e con quale calendario si lavorerà, e ci si presenta a vicenda. ``C'è poco
gioco analogico'', si scusano. Ma nella seconda parte la lezione comincia
comunque a costruire un \textbf{linguaggio comune}: le figure professionali
dell'industria del gioco analogico, la differenza tra \emph{game design} e
\emph{game development}, la situazione del mercato italiano rispetto a quello
estero, e un primo glossario di design (\emph{game} e \emph{play}, il framework
MDA, le parole ``vietate'', il playtest e l'iterazione). Si chiude con lo
strumento di settore per i prototipi digitali dei giochi da tavolo,
\textbf{Tabletop Simulator}, e con una lunga sessione di domande.

La lezione si svolge \textbf{online}, su Zoom, ed è condotta a più voci dai tre
docenti principali del corso: \textbf{Daniele}, narrative designer, che guida la
parte organizzativa; \textbf{Ruggero}, game designer, che guida la parte sui
fondamenti e le domande finali; \textbf{Beatrice}, progettista di giochi
pervasivi, che interviene soprattutto nella prima metà. Gli studenti iscritti
sono \textbf{nove}; uno è assente e un altro arriva a lezione iniziata. Ci sono
una pausa di una decina di minuti (si riprende alle 18:15) e, dopo le 19:15,
un'ultima mezz'ora di domande con il solo Ruggero.

Tre avvertenze di lettura. La prima: la trascrizione è automatica e storpia
molti nomi propri, compresi quelli dei docenti e dell'accademia (``Dive'',
``IVE'', ``Accademia Live'', ``Ai''); dove la ricostruzione è sicura, o molto
probabile, il capitolo usa la forma corretta e lo segnala in nota; dove non lo
è, riporta la forma della trascrizione e lo dice. La seconda: durante la pausa la
trascrizione registra, ripetuta in loop per decine di righe, una frase senza
senso che non ha nulla a che fare con la lezione, e che il capitolo ignora. La
terza: in molti passaggi più docenti parlano uno dopo l'altro senza che la
trascrizione segni il cambio di voce; il capitolo attribuisce le parole a chi le
pronuncia quando il contesto lo rende chiaro (chi parla di sé, chi viene
chiamato per nome, chi ``passa la parola''), e resta generico altrimenti.

Quanto alla data, la lezione si colloca con ogni probabilità \textbf{nella prima
metà di settembre del 2026}\footnote{La data non è mai detta. La si deduce da
alcuni indizi: il form iniziale va compilato ``entro e non oltre il 15
settembre''; i gruppi verranno annunciati il 25 settembre; le date del
calendario (sabato 3 ottobre, lunedì 12 ottobre, sabato 7 e 14 novembre,
venerdì 4 dicembre) cadono nei giorni della settimana indicati soltanto nel
2026; un docente parla di un argomento ``molto complesso nel 2026''; e fa
ancora molto caldo (``le temperature sono ancora indecenti''). È una
deduzione, non un'informazione esplicita.}: i riferimenti a ``oggi'',
``quest'estate'' e ``l'anno scorso'' vanno letti rispetto a quel momento.

Il percorso del capitolo segue quello della lezione:
\begin{enumerate}
  \item l'avvio, i \textbf{docenti} e le \textbf{presentazioni} dei partecipanti
        (sezioni da~\ref{a1:sec:avvio} a~\ref{a1:sec:partecipanti});
  \item i \textbf{materiali} e gli strumenti del corso: il form, la cartella
        condivisa, Discord (sezione~\ref{a1:sec:materiali});
  \item lo \textbf{scopo} e il \textbf{metodo} del corso, le regole di
        comportamento e il carico di lavoro (sezioni~\ref{a1:sec:scopo}
        e~\ref{a1:sec:regole});
  \item il \textbf{project work} e il \textbf{calendario} (sezione~\ref{a1:sec:project});
  \item la regola sull'\textbf{intelligenza artificiale} e la lunga discussione
        che ne segue (sezione~\ref{a1:sec:ia});
  \item le \textbf{figure professionali} dell'industria del gioco analogico,
        con un approfondimento su marketing e sales
        (sezioni~\ref{a1:sec:figure} e~\ref{a1:sec:marketing});
  \item la differenza tra \textbf{game design} e \textbf{game development}
        (sezione~\ref{a1:sec:design-development});
  \item il \textbf{mercato}: Italia ed estero (sezione~\ref{a1:sec:mercato});
  \item i \textbf{fondamenti di design}: perché serve una terminologia comune,
        \emph{game} e \emph{play}, il framework \textbf{MDA} (e il DPE), le
        \textbf{parole vietate}, chi crea, chi gioca e chi facilita, il
        \textbf{playtest} e l'\textbf{iterazione} (sezioni
        da~\ref{a1:sec:terminologia} a~\ref{a1:sec:playtest});
  \item \textbf{Tabletop Simulator} e il primo task (sezione~\ref{a1:sec:tts});
  \item le \textbf{domande finali}: il tema del project work, il bilanciamento,
        i ruoli nei team, la pubblicazione, il gioco di ruolo in Italia, le
        \emph{house rule} di Monopoly, i giochi ``comunicati male''
        (sezioni~\ref{a1:sec:domande-daniele} e~\ref{a1:sec:domande}).
\end{enumerate}

% ---------------------------------------------------------------------
\section{L'avvio della lezione}
\label{a1:sec:avvio}

La lezione comincia con qualche minuto di attesa. ``Siamo in nove in classe'', e
si aspetta che arrivino tutti; uno dei docenti scherza sul fatto che l'unica
persona che arriverà ``sempre in ritardo di venerdì'' è Gianmarco. Mancano
Giordano e Tiziana; Tiziana arriva, mentre la segreteria fa sapere poco dopo che
Giordano non potrà esserci e vedrà la registrazione. Si aspetta anche Beatrice,
la terza docente, che arriva a sua volta dopo pochi minuti.

Nell'attesa, Daniele anticipa il programma della giornata e qualche
informazione pratica:
\begin{itemize}
  \item oggi è una giornata di \textbf{introduzione}: come funziona il corso, chi
        sono i docenti, le presentazioni degli studenti, le \textbf{metodologie}
        e il \textbf{calendario}. Gli studenti avranno ``un fracco di domande'', e
        alla fine della lezione c'è ``proprio uno slot apposta'' per le domande;
  \item le lezioni sono \textbf{registrate}: di solito la registrazione viene
        consegnata dopo circa \textbf{12--24 ore} sul profilo personale della
        piattaforma dell'accademia, che va attivato con la segreteria (Tiziana,
        già studentessa, probabilmente ce l'ha già);
  \item i docenti ``cuore'' del corso sono tre (Daniele, Ruggero e Beatrice),
        ``soprattutto Ruggero''; a loro si aggiungeranno tre docenti ospiti, per
        meno ore, probabilmente una lezione a testa, ``tranne Marta, che avrà due
        lezioni quest'anno''.
\end{itemize}
E una battuta sul caldo: ``spero che voi siate carichi, anche se le temperature
sono ancora indecenti''. ``Carichi moderatamente, per quello che potete''.

Poi Daniele condivide lo schermo (``non solo le slide, ma anche un po' di altre
cose'') e apre ufficialmente il corso: ``benvenuti e benvenute nella seconda
edizione del corso di Analog Game Design. Questa è la prima lezione''. Il
percorso è di circa \textbf{tre mesi}, fino a dicembre, e durante il corso i
docenti cercheranno di far provare agli studenti ``cosa significa fare Analog
Game Design'' con un \textbf{project work}, un progetto di gruppo di cui si
spiegheranno più avanti la struttura e le richieste.

% ---------------------------------------------------------------------
\section{I docenti}
\label{a1:sec:docenti}

``Anzitutto partiamo da delle presentazioni doverose''. Ognuno dei tre docenti
principali si presenta con una slide che riporta i contatti e una selezione dei
lavori fatti.

\subsection{Daniele: narrativa, fiction e gioco di ruolo}

\textbf{Daniele}\footnote{La trascrizione riporta ``Daniel Cusetto''; nel
resto della lezione tutti lo chiamano ``Daniele''. Il cognome non è
ricostruibile con certezza.} ha già conosciuto quasi tutti gli studenti, perché
ha fatto loro il colloquio di ammissione. È \textbf{narrative designer}, o meglio
\emph{game narrative designer}; insegna nell'accademia che ospita il
corso\footnote{La trascrizione storpia il nome dell'accademia in molti modi
(``Dive'', ``IVE'', ``Accademia Live'', ``Ai'', ``AIVA''). Con ogni probabilità
si tratta dell'\textbf{AIV}, l'Accademia Italiana Videogiochi, che ha sede a
Roma: lo fanno pensare i corsi triennali di game design e narrative design citati
dai docenti e il fatto che una studentessa racconti di essersi trasferita a Roma.
È una ricostruzione, non un'informazione esplicita; nel seguito il capitolo
parla semplicemente dell'``accademia''.} ma lavora soprattutto come \textbf{brand
advisor}. Nel corso si occuperà della parte di \textbf{gioco di ruolo},
\textbf{narrazione} e \textbf{fiction}.

I suoi contatti, ``in quest'ordine'', sono email, Telegram e Discord. Ci tiene a
dare alcune indicazioni precise:
\begin{itemize}
  \item se si vuole avere la certezza di una risposta, bisogna scrivergli su
        \textbf{Telegram}, perché lì ha tutti i suoi clienti;
  \item Discord lo guarda di meno; l'\textbf{email} è per le cose ufficiali;
  \item non bisogna contattarlo altrove: non su Facebook, non su LinkedIn, non su
        Instagram, ``perché li vedo molto di meno''.
\end{itemize}
La slide mostra anche alcuni giochi, attività e brand su cui ha lavorato negli
anni, ma non si dilunga: ``nel caso ne parleremo''.

\subsection{Ruggero: game design e ricerca e sviluppo}

\textbf{Ruggero} (spesso chiamato ``Ruger'') ha fatto il \textbf{game designer}
per varie aziende e oggi è game designer e si occupa di \textbf{ricerca e
sviluppo} per un'azienda del settore\footnote{La trascrizione dice ``per MS'';
il nome dell'azienda non è ricostruibile con certezza. Più avanti Ruggero cita
tra i giochi su cui ha lavorato il gioco da tavolo ufficiale di
\emph{Cyberpunk}, \emph{Marvel United} e \emph{Status}, un gioco sulla
geopolitica.}. Ha un background ``misto'': nel suo portfolio ci sono giochi da
tavolo, giochi di ruolo, giochi analogici, LARP, videogiochi, giochi in realtà
virtuale. Si occupa anche di \textbf{divulgazione} sul gioco e di analisi di
\textbf{game studies} e game design. Più avanti, nelle domande finali, dirà di
avere ``50 e più giochi prodotti e progettati'' e di insegnare da vari anni.

Lo si può contattare liberamente via Gmail e Telegram, e cerca di rispondere
\textbf{entro 24 ore}, salvo imprevisti o eventi: ``sentitevi in libertà di
scrivermi se avete dubbi, senza problemi''. Quando arriva a parlare dei propri
giochi, la slide mostra ``un mix di varie cose, dalle più famose alle meno
famose'', ma Ruggero si imbarazza a presentarli, tra le risate dei colleghi.

\subsection{Beatrice: i giochi pervasivi}

\textbf{Beatrice} (``Bea'') si occupa di \textbf{giochi pervasivi}
(\emph{pervasive games}). ``Se non li avete mai sentiti nominare, non
preoccupatevi'': questo corso è ``uno dei pochi, forse l'unico in Italia'' che
li racconta esplicitamente. Li si vedrà nelle prossime lezioni, ma la definizione
di base è questa: i giochi pervasivi \textbf{utilizzano gli spazi e gli oggetti
della vita quotidiana e li trasformano in spazi e oggetti di gioco}.

Beatrice lavora come \textbf{freelance} e si occupa di diversi tipi di
progettazione: lavora sia con le \textbf{aziende} sia con gli \textbf{eventi dal
vivo}. La sua attività diretta si trova sotto l'etichetta \textbf{Minime
Differenze}\footnote{La trascrizione riporta il nome dell'etichetta come
``Minime Differenze'' e l'indirizzo del sito come ``mininedifferenze.com''; la
grafia esatta del sito non è verificabile dalla trascrizione.}, mentre i lavori
per le aziende sono ``sparpagliati'' sui siti delle varie aziende. Molte delle
cose mostrate nella sua slide sono attività proposte al pubblico ``con un
biglietto, con un gioco in formato evento'': un tema che si toccherà nel corso,
per capire come \textbf{comunicare} i propri giochi in modo che siano
comprensibili, accessibili per il pubblico e \textbf{vendibili}.

Ammette un ``grosso limite'' con i messaggi: è un po' lenta a rispondere, ``ma
rispondo sempre a tutto''. La si può contattare per email (cognome punto nome) o
su Telegram. È ``molto curiosa di vedere cosa produrrete''.

\subsection{I docenti ospiti}

Daniele presenta poi i tre docenti che il corso avrà ``solo per alcune
lezioni'', tutte verso \textbf{novembre}:
\begin{itemize}
  \item \textbf{Diego Ricchiuti}, direttore del corso di Game Design
        dell'accademia, farà una lezione sull'\textbf{RNG} (\emph{random number
        generation}, la casualità nei giochi): ``vi parleremo di RNG quando sarà
        necessario. Comunque avrete una lezione con lui, sappiatelo'';
  \item \textbf{Marta}\footnote{La trascrizione riporta il cognome come
        ``Cercasassi''; la grafia non è verificabile.} farà \textbf{due lezioni
        sul development}: quest'anno ``abbiamo deciso di fare doppietta'',
        perché l'anno scorso il suo intervento è stato ``molto, molto utile'';
  \item \textbf{Flavio}\footnote{La trascrizione riporta il cognome come
        ``Mortarino'' e, più avanti, ``Mortalino'', e il nome della sua casa
        editrice come ``Grampibert''; nessuno dei due è ricostruibile con
        certezza.} farà la lezione di \textbf{game production}; è un editore, e
        di solito partecipa anche alla tavola rotonda con gli editori della
        giornata finale.
\end{itemize}
Le slide sono già online e gli studenti avranno accesso alla cartella condivisa:
``non vi preoccupate di queste cose, ne parleremo a breve''. Ma prima di parlare
del corso, ``vogliamo conoscere voi''.

% ---------------------------------------------------------------------
\section{Le presentazioni dei partecipanti}
\label{a1:sec:partecipanti}

Le domande per la presentazione sono semplici, da affrontare ``a turno, senza
fretta, con calma'':
\begin{itemize}
  \item nome e occupazione (``se volete dirla'');
  \item perché si vogliono fare giochi analogici, e a che livello: che cosa ha
        portato lo studente al corso;
  \item che tipo di esperienze o conoscenze si hanno già nell'ambito, se si è già
        fatto qualcosa oppure no.
\end{itemize}
``Chi vuole partire senza nessun tipo di ritegno e di vergogne?''. Per far
parlare gli studenti guardandoli, Daniele interrompe la condivisione dello
schermo e ricorda lui le domande quando serve. Le presentazioni sono interessanti
anche per i commenti dei docenti, che anticipano diversi principi del corso.

\subsection{Tiziana}

Tiziana fa l'\textbf{infermiera}, ma nella vita ha fatto molti lavori:
impiegata, insegnante e ``tanti altri lavoretti''. Gioca ``da sempre'': il gioco
fa parte anche della sua vita familiare, perché suo fratello fa il programmatore
di videogiochi. Ha già \textbf{organizzato} giochi: una caccia al tesoro per far
conoscere il territorio, due scene con un delitto (giochi con omicidio da
risolvere), giochi per intrattenere le persone. Ha riscoperto di recente i giochi
da tavolo, arrivando a \emph{Carcassonne} con famiglia e amici; poiché è
difficile trovare persone con voglia di giocare nella propria cerchia ristretta,
ha scoperto il mondo delle \textbf{associazioni} ludiche, ``e lì è stato
un'esplosione'': prima in Veneto (è di Vicenza), poi a Roma, dove si è
trasferita.

Ha in mente un progetto ambizioso: far conoscere una certa realtà con la modalità
del \textbf{gioco di ruolo}, o comunque proporre un gioco suo.

\paragraph{Il commento: gli strumenti, non il gioco dei sogni.} Ruggero coglie
l'occasione per un'avvertenza che vale per tutti, visto che molti studenti
arrivano con idee già pronte. Nel corso si lavorerà a un progetto, ma
l'obiettivo dei docenti è dare \textbf{strumenti} da applicare anche
\emph{fuori} dal progetto, ai giochi che gli studenti vorranno fare un domani:
``che non siano esercizi per il corso, ma per voi in primis''. Quindi: ``non
rimaneteci male se non riuscirete a lavorare al vostro gioco dei sogni in questo
corso''. Anche perché ``sarebbe impossibile, in due mesi, fare il gioco dei
sogni''. Il corso darà invece la possibilità di fare un'esperienza ludica
insieme e di imparare a replicarla sulle proprie idee: ``le idee non buttatele,
anzi, il contrario: tenetele''.

\subsection{Nicola}

Nicola è uno \textbf{studente}. Sei mesi prima ha fatto un colloquio con Daniele
perché stava valutando di passare all'accademia, e ha deciso, ``nel dubbio'', di
cominciare da questo corso breve per capire come funziona la formazione;
eventualmente si iscriverà al corso di game design vero e proprio. Gli interessa
in particolare il lato dell'esperienza del giocatore, su cui è ``molto critico''
(qui usa la parola ``divertimento'', e Daniele lo ferma con un ``ne parleremo'':
è una delle parole che il corso vieta, come si vedrà nella
sezione~\ref{a1:sec:parole-vietate}).

Ha 21 anni ed è \textbf{giocatore di ruolo da circa otto}. Prima ancora di
conoscere i giochi di ruolo si era creato un sistema tutto suo, ``pessimo,
disgustoso, terribile'', un sistema a percentuale con il d100 ``terrificante'';
poi ha scoperto i giochi di ruolo da tavolo, i classici (\emph{Dungeons \&
Dragons}, \emph{Daggerheart}, \emph{Vampire}\footnote{La trascrizione riporta
``DD'', ``Daggerart'' e ``Vampire'': quasi certamente \emph{Dungeons \&
Dragons}, \emph{Daggerheart} e \emph{Vampire: The Masquerade}.}). Gioca anche a
\textbf{card game} e ad altri giochi sociali, pur avendo sempre avuto difficoltà
a trovare gruppi, ``avendo vissuto in un buco dimenticato dal mondo''. Come
giocatore è ``l'eterno narratore'' del gruppo, quello che prepara la storia e si
costruisce le regole. Come autore ha fatto soprattutto giochi su carta, ``perché
sono incapace a programmare, per adesso''.

Daniele osserva che il tema del giocatore e del suo rapporto con il gioco è ``un
argomento importante'', anzi ``un vaso di Pandora enorme'' (espressione che
tornerà molte volte nella lezione), e che il corso tratterà anche che cos'è
l'analog game design in generale e quanti tipi di prodotti esistono.

\subsection{Rita}

Rita\footnote{Daniele la chiama prima ``Nino'' (probabilmente un errore della
trascrizione) e poi la ringrazia come Rita.} ha 30 anni e si collega da una
regione del Sud, dove ``ancora si muore di caldo''. Ha studiato
all'\textbf{Accademia di Belle Arti di Napoli}, \textbf{design della
comunicazione visiva}: si occupa dell'aspetto più imprenditoriale della
progettazione artistica. Non volendo progettare per un'impresa (``fare loghi,
robe del genere''), vorrebbe applicare la propria creatività alla realizzazione
di giochi analogici.

Non è una giocatrice di ruolo e non lo è mai stata: è una giocatrice
\textbf{tradizionale}, cresciuta con Monopoly, la dama e giochi tradizionali
cinesi. Ha già \textbf{idee, prototipi e progetti} realizzati, e vuole capire
come renderli \textbf{prodotti per la vendita} e se questo è il campo adatto a
lei; magari anche confrontarsi con i colleghi. La sua progettazione è orientata
al \textbf{sociale} e all'\textbf{educativo}: ``responsabilizzare un po' le nuove
generazioni su come sta andando il mondo e su come potrebbe andare, visto che è
tutto nelle nostre mani''.

\paragraph{Il commento: siate curiosi.} ``Pane per i tuoi denti'', le risponde
Daniele, soprattutto per i LARP e i giochi pervasivi: ``vi apriremo un po' la
testa da questo punto di vista''. E aggiunge un consiglio per tutti: ``siate
curiosi'', perché imparare da altri ambiti dà moltissimo. Ognuno è lì per fare
qualcosa di preciso, ``però siate aperti anche ad altre esperienze''.

\subsection{Tancredi}

Tancredi ha 21 anni e non ha ancora grande esperienza lavorativa: qualche
``piccolo progettino'' che non è andato come voleva. Gli interessa il mondo dei
giochi, digitali e non, con un approccio \textbf{educativo} e \textbf{sociale}.
È arrivato soprattutto dai videogiochi, a cui gioca da quando era piccolo; ai
giochi da tavolo si è avvicinato solo nell'ultimo anno, grazie a una realtà
ludica, perché in famiglia ``non siamo molto da giochi''. Ha realizzato qualche
piccolo videogioco con vari engine: prima quelli senza un vero linguaggio di
programmazione, ``più a nodi'', poi qualcosa di più complicato, e ora sta
lavorando a un piccolo prototipo.

\paragraph{Il commento: niente di quello che vedrete è sprecato.} ``Non butterai
via niente di analog game design'', gli risponde Daniele, ``perché il 99 per cento
delle cose che vedremo si ritrovano quasi uno a uno nel mondo dei videogiochi''.
Tanto che nei corsi triennali dell'accademia (anche in quello di narrative
design, dove insegna Ruggero) il primo anno si parte dai \textbf{giochi da tavolo
e di ruolo}, proprio per spiegare meccaniche più complesse che servono poi per i
videogiochi. Ruggero aggiunge che vale anche il contrario: alcuni dei
\textbf{framework} di game design che si vedranno nelle lezioni tecniche vengono
dal videogioco e sono stati riapplicati all'analogico. ``Evitare la
crossmedialità qui è impossibile''. Anzi, tutte le conoscenze ``mediatiche
adiacenti'' (musica, letteratura, cinema) saranno utili.

\subsection{Lorenzo Vergata}

I Lorenzo in classe sono due; il primo a presentarsi è \textbf{Lorenzo
Vergata}, 26 anni, \textbf{dottorando al Politecnico di Milano}. Gioca ai
videogiochi da quando aveva due anni e mezzo, ma ha giocato al suo primo gioco da
tavolo ``credo nove mesi fa''. Gli fa piacere sentire che molti in classe sono
interessati ai giochi \textbf{sociali} ed \textbf{educativi}, un settore che
interessa anche a lui e che ``può essere esplorato molto'', anche con qualche
aiuto dell'intelligenza artificiale per esperienze più ricercate.

Di esperienza nell'analogico dice di avere fatto ``il minimo sindacale''. Ma
racconta che diciotto anni prima, da bambino, era ``più game designer di
adesso'': prendeva un gioco e ne cambiava le regole, anche al biliardino (``no,
ora introduciamo questa regola''). Gli piace sperimentare, creare regole e
vedere che cosa accade in sé e negli altri. Soprattutto vorrebbe ``una chiave di
lettura per poter rimanipolare la realtà'': tre settimane prima è stato a fare la
\textbf{vendemmia in Francia} e ha visto tantissime dinamiche che potrebbero
diventare un gioco da tavolo o un videogioco. Per esempio: finire contro un rovo
pieno di semi che si attaccano ai vestiti lo rallentava nella raccolta
dell'uva, mentre in un altro filare era più rapido. ``Si può fare pure una
dinamica'', e vorrebbe gli strumenti per farne giochi.

Arriva subito la battuta: una ``vendemmia gamificata'' da proporre in Oltrepò
Pavese (``io sono dell'Oltrepò Pavese, quindi devo per forza bere vino'')\footnote{La
trascrizione dice ``nell'altro papavese'' e non chiarisce chi pronunci la
battuta.}. Daniele conferma che quest'anno il corso avrà ``tante cose sul
sociale'': ``ci piace parecchio''.

\subsection{L'altro Lorenzo}

Il secondo Lorenzo (``secondo, ma non meno importante'') è uno \textbf{scrittore
di racconti brevi} e sta scrivendo il suo primo romanzo. Gli è sempre piaciuto
raccontare storie, e da più grande scriverle; gioca ai videogiochi ed è ``un
amante delle storie e della cultura pop''. Negli ultimi anni è entrato nel mondo
del \textbf{gioco di ruolo da tavolo} e vorrebbe unire le due passioni, le
storie e il gioco di ruolo, magari facendone un lavoro; gli interessa anche lo
\emph{strategic writing}. Daniele gli ricorda che ci sarà una lezione proprio
sulla \textbf{fiction}.

\subsection{Angelo}

Angelo si collega senza videocamera dal lavoro: fa il \textbf{birraio} (in un
piccolo birrificio in Brianza, a Bernareggio, come racconterà dopo la pausa) e,
con il caldo, ``vi risparmio l'orrore''. Per molto tempo era convinto che i
giochi da tavolo fossero solo i classici (Monopoly, Risiko) o i party game; poi,
quattro o cinque anni fa, ha ``sbattuto la faccia'' contro giochi davvero
appassionanti, con una specie di narrazione dietro che li rendeva ancora più
coinvolgenti, e non ha più smesso. È uno che ``ha sempre un botto di idee, ma poi
non mette in pratica assolutamente nulla'': cercava un corso che gli desse le
basi e soprattutto lo \textbf{motivasse} a provare a fare qualcosa.

\paragraph{Il commento: bisogna buttarsi.} Daniele racconta che anche l'anno
prima un suo caro amico si era iscritto proprio per questo motivo. Il corso è
breve, ma cerca di portare gli studenti ``da 0 a 100'' e di far realizzare un
\textbf{prodotto} da presentare alla fine. Il punto è che i giochi ``bisogna
farli, bisogna buttarsi''. E, venendo dalla narrazione, Daniele lo paragona allo
scrivere: si scrive e si riscrive, e allo stesso modo si fa \textbf{playtest}, si
rifà, si aggiorna. ``Non è che si fanno così''.

\subsection{Gianmarco (dopo le slide)}

Gianmarco arriva a lezione iniziata e si presenta più tardi, alla fine della
prima parte, dopo che i docenti hanno chiuso le slide organizzative e la
discussione sull'intelligenza artificiale (``Gianmarco, ti giuro, ti facciamo
presentare''). Lavora nel \textbf{team di innovazione} di un'azienda italiana.
Vuole fare giochi perché è una sua passione, che lo gratifica, e in passato ci ha
provato ``in maniera scomposta''. Cerca un \textbf{metodo} e un confronto con
persone più esperte: ``studiare solo dei manuali è molto sterile'', ``stare
chiusi nel pertugio non aiuta'', ed ``essere autodidatta ha estremamente i suoi
limiti''. Come esperienza è un giocatore e ha letto ``un po' di materiale, di
blog, di manuali'', ``niente di rilevante''.

Manca solo Giordano, che si presenterà alla lezione successiva.

\paragraph{Uno sguardo d'insieme.} La classe è molto varia: un'infermiera,
una designer della comunicazione visiva, un birraio, un dottorando, uno
scrittore, un impiegato nell'innovazione, due studenti ventunenni. Alcuni
vengono dal gioco di ruolo, altri dai giochi tradizionali, altri dai
videogiochi; quasi tutti hanno scoperto i giochi da tavolo moderni da pochi anni
o da pochi mesi. E almeno tre studenti dichiarano un interesse per i giochi
\textbf{sociali} ed \textbf{educativi}. Come si vedrà nella
sezione~\ref{a1:sec:scopo}, i docenti considerano questa varietà una risorsa.

% ---------------------------------------------------------------------
\section{I materiali e gli strumenti del corso}
\label{a1:sec:materiali}

Prima delle slide sul corso, Daniele sbriga un po' di ``clean up casalingo''.

\paragraph{Il form iniziale.} C'è un \textbf{form} da compilare ``entro e non
oltre il 15 settembre'': il QR code è nelle slide, e la segreteria lo ricorderà
su Discord. Richiede circa 10--15 minuti; si può fare subito o in giornata, ``e
ve lo togliete di torno''. È un form standardizzato per tutti i corsi (``non
siamo io e Ruggero a farvi passare la tortura del test''), diviso in due parti:
\begin{itemize}
  \item una parte di \textbf{conoscenza del candidato}: quanto tempo si ha a
        disposizione per il corso, come ci si organizza con il lavoro, che cosa
        piace giocare. È la parte che interessa ai docenti, perché li aiuta a
        organizzare le lezioni;
  \item una parte di \textbf{test logico}, che i docenti non hanno scritto (``io
        non sono responsabile del form'', ``no, no, nemmeno io'').
\end{itemize}

\paragraph{La cartella condivisa.} Tutti gli studenti vengono aggiunti alla
cartella del corso su \textbf{Google Drive}, che si trova tra i file
``condivisi con me''. Gli studenti ne sono \textbf{gestori dei contenuti} e
possono caricarci quello che vogliono. Dentro ci sono:
\begin{itemize}
  \item le \textbf{attività in classe};
  \item le \textbf{dispense};
  \item i \textbf{materiali didattici}, con le slide (quelle della lezione sono
        già lì);
  \item le cartelle dei \textbf{project work} futuri;
  \item i \textbf{task}, che ``saranno pochi, ma ci saranno'';
  \item il \textbf{calendario studenti}, ``importantissimo'', con due schede: il
        calendario delle lezioni e quello del project work.
\end{itemize}

\paragraph{Discord.} Il corso ha anche un \textbf{server Discord}, a cui tutti
sono stati invitati; ci sono i docenti (Daniele, Ruggero, Diego e altri), che si
possono taggare. Di solito è lì che la segreteria pubblica le informazioni, per
esempio se una lezione ``dovesse per caso essere cancellata''. Nelle domande
finali Ruggero aggiungerà che c'è un canale apposito per il corso, che gli
studenti possono usare liberamente per organizzarsi tra loro.

\paragraph{Le reaction di Zoom.} Durante le lezioni su Zoom si possono usare le
\textbf{reaction} (il pollice in alto, il cuoricino, la mano alzata): un modo
rapido per far capire ai docenti se si è capito, se si vuole intervenire. Nel
frattempo arrivano i primi pollici.

% ---------------------------------------------------------------------
\section{Lo scopo e il metodo del corso}
\label{a1:sec:scopo}

La prima serie di slide (che Daniele, per un lapsus, presenta come ``cos'è il
corso di pulire il pavimento di casa'') spiega che cosa il corso vuole fare e
come.

\subsection{Diventare autori analogici}

Lo \textbf{scopo del corso} è fornire le competenze per \textbf{diventare autori
analogici} e stimolare la creatività e l'\textbf{autorialità} di ciascuno: ``voi
imparerete a diventare autori di giochi analogici''. La differenza tra
\emph{game designer} e \emph{game developer} verrà chiarita più avanti (sezione~\ref{a1:sec:design-development}),
ma la parola chiave qui è \emph{autorialità}: ``vogliamo voi, vogliamo le vostre
idee''. Magari non si potrà usare l'idea del gioco della propria vita, ma si
potranno portare nelle lezioni le proprie idee e le proprie sensibilità.

Gli obiettivi elencati nella slide sono questi:
\begin{itemize}
  \item dare \textbf{maggiore consapevolezza nel design e nello sviluppo}. In
        Italia, purtroppo, ``si fa design a volte in maniera non consapevole''.
        Consapevole vuol dire avere almeno un'introduzione alle basi del design e
        dello sviluppo, e fare le cose ``con metodo, con raziocinio, con
        consapevolezza'';
  \item fornire gli \textbf{strumenti essenziali} per un design consapevole;
  \item stimolare la \textbf{sperimentazione} e la \textbf{trasversalità
        mediatica}. Non si chiede di fare ``il \emph{Root} della situazione o il
        \emph{Carcassonne} della situazione''\footnote{La trascrizione dice ``il
        lumaven della situazione o il carcassone della situazione''; il primo
        titolo non è ricostruibile con certezza. Il senso è chiaro: non si
        chiede di rifare un grande classico del gioco da tavolo.}, ma un
        project work \textbf{calibrato sulla durata del corso}, e però
        \textbf{sperimentale} e ``assolutamente transmediale'';
  \item portare \textbf{numerosi esempi} di tipologie di gioco, dall'analogico al
        tavolo, al pervasivo ``e oltre''. Lo si vedrà nella lezione sui
        \textbf{casi studio} del venerdì successivo: ``voi non avete idea di
        quante cose esistono, bellissime, sconosciutissime e favolose''. Daniele
        confessa che preparando e aggiornando le slide i colleghi gli hanno
        mostrato cose che lui stesso non sapeva nemmeno che esistessero;
  \item dare gli strumenti per \textbf{analizzare} i giochi altrui. Ruggero lo
        sottolinea: se il corso riesce nel suo intento, gli studenti sapranno
        non solo fare i propri giochi ma anche analizzare giochi già fatti, che è
        ``una delle skill importanti anche nella parte di ricerca e sviluppo''.
        Si vedranno casi studio e giochi pubblicati, sempre con l'obiettivo di
        capire \textbf{che cosa fanno e perché lo fanno}, per poi fare meglio il
        proprio lavoro;
  \item imparare a \textbf{fare giochi insieme}. Esistono gli autori che lavorano
        da soli (i \emph{solo developer}), ``ma sappiate che anche l'autore che
        lavora da solo poi si deve interfacciare con il pubblico, con un editore,
        magari con un editor, con un publisher''. ``Non lavorerete mai da soli'':
        bisogna imparare a lavorare con gli altri.
\end{itemize}

\subsection{Un corso intersezionale}

Il corso è \textbf{intersezionale}: \textbf{non ci sono lezioni dedicate a una
singola tipologia di gioco}. Non ci sarà ``la lezione sul gioco da tavolo'',
``la lezione sul gioco di ruolo'', ``la lezione sul gioco pervasivo''. Ci saranno
lezioni su temi diversi, ciascuna pensata per essere utile a \textbf{tutte e tre
le sfaccettature} del corso; alcune saranno più importanti per il gioco da
tavolo, altre per il pervasivo, altre per il gioco di ruolo, ma nessuna sarà
esclusiva. I docenti vogliono ``contaminazioni e applicazioni assolutamente
trasversali''.

\subsection{Un corso cross-generazionale}

È anche un corso \textbf{cross-generazionale}: i gradi di esperienza di gioco
quest'anno sono abbastanza allineati, ma le esperienze di vita, le età, le
estrazioni sociali sono varie. ``Ed è un pro, non è un contro. Più diversificata
è la classe, più imparate meglio anche voi''. La slide ha una piccola nota per
fasce d'età:
\begin{description}
  \item[Under 30] Concentratevi sulla vostra \textbf{crescita}. Evitate di
        confrontarvi con le persone più grandi che magari ``ce la fanno meglio''
        in alcune cose.
  \item[Over 30] Concentratevi sul \textbf{consolidamento}, ma siate
        \textbf{umili}: i docenti diranno cose che faranno ``storcere il naso''.
        ``Bisogna un po' reimparare per imparare'': bisogna sapersi adattare.
\end{description}
Ruggero contestualizza: lo dicono perché in passato hanno visto studenti che,
essendo più adulti o avendo giocato di più, pensavano che l'esperienza
\emph{di gioco} si trasformasse automaticamente in competenza \emph{di
creazione} di giochi. Non è così. E non bisogna sottovalutare chi gioca meno:
proprio perché gioca cose diverse, per risolvere un problema di design specifico
``è una ricchezza incredibile avere qualcun altro che riesce a uscire dalla
scatola''.

\subsection{Un corso pratico e adattabile}

Il corso è \textbf{pratico}: gli studenti si concentreranno su un progetto di
gruppo, il \textbf{project work}, lavorando in team. Ed è \textbf{adattabile}: il
programma è definito, ma ``si aggiorna sui vostri stimoli, priorità e
necessità''. Se durante il corso i docenti vedono che serve più una cosa
dell'altra, sono ``prontissimi'' a dare più spazio a un argomento invece che a
un altro.

\subsection{Inquiry-based learning: domandate}

Il metodo è quello dell'\textbf{inquiry-based learning}, l'apprendimento basato
sulle domande. Di solito le persone dicono ``ho una domanda, magari sarà
stupida''. ``Non esistono domande stupide. Pensate che ci sia una domanda
stupida? Fatela, alzate la mano''. Vanno bene anche le
\textbf{provocazioni}: se qualcosa non torna, è meglio chiedere. La condivisione
è fondamentale.

I docenti danno \textbf{metodi, dati e conoscenze}, ma lasciano trovare le
risposte agli studenti: ``vi lasciamo un pochettino annegare nel mare prima di
darvi la soluzione, in modo tale che la possiate trovare da soli''. Di
conseguenza \textbf{non ci saranno interrogazioni sulle definizioni}: interessa
vedere \textbf{come} gli studenti applicano ciò che imparano. Se una
definizione non viene in mente, non è un problema; l'importante è saperla usare.

\subsection{Feedback, non valutazione}

``Sì al feedback, no alla valutazione''. I docenti danno \textbf{feedback
verbali}, anche sulle attività da fare a casa (i task, che non saranno tanti).
Il corso è calibrato per \textbf{studenti lavoratori}: c'è un po' da lavorare
fuori dalle lezioni, ma in modo calibrato.

Il principio è \textbf{fallire oggi per imparare per il futuro}: \emph{fail
forward}, \emph{fail to learn}. I docenti lasciano provare e fallire almeno una
volta prima di dare la soluzione: ``il classico metodo Montessori, quello
originale, non quello travisato''.

\subsection{Le soft skill e i criteri di valutazione}

Grande attenzione va alle \textbf{soft skill}: i lavori di gruppo, lo scambio
propositivo, l'atteggiamento positivo, la propositività, la proattività.
``Lavoro di gruppo'' non significa stare ogni giorno in sala studio con gli
altri: va benissimo che un gruppo si organizzi dandosi priorità e tempi in base
alla disponibilità di ciascuno. L'importante è \textbf{collaborare}, essere
\textbf{professionali}, rispettarsi a vicenda, rispettare le idee proprie e
altrui e lasciare spazio agli altri, perché i progetti ``non saranno solo
vostri, saranno di tutte le persone del vostro gruppo''.

La valutazione del project work non riguarda solo la \textbf{documentazione}
prodotta, ma anche:
\begin{itemize}
  \item la \textbf{presentazione} finale;
  \item l'\textbf{autovalutazione};
  \item la \textbf{crescita personale};
  \item l'\textbf{atteggiamento} tenuto nel corso.
\end{itemize}
Daniele cita: ``da ognuno secondo le sue capacità, a ognuno secondo i suoi
bisogni''. Lavorare in gruppo non vuol dire che, in tre, ciascuno lavori
esattamente il 33\% del tempo. Ai docenti serve piuttosto capire \textbf{chi fa
che cosa} nel gruppo, e che il processo interno sia \textbf{trasparente}.

\paragraph{Lavorare in cloud.} Per lo stesso motivo il consiglio è di
\textbf{non lavorare in locale}, ma direttamente nelle cartelle condivise del
project work e dei task. Si è liberi di fare altrimenti, ma lavorando in cloud
si dà ai docenti la possibilità di vedere \emph{come} si lavora, il proprio
``processo mentale''.

% ---------------------------------------------------------------------
\section{Le regole di comportamento e il carico di lavoro}
\label{a1:sec:regole}

Alcune slide sono di quelle che si ripetono in ogni corso (``\emph{repetita
iuvant}'', storpiato per scherzo in ``ripetita Juventus'').

\paragraph{Comportamento.} Prestare attenzione durante le lezioni; interagire;
portare rispetto; essere proattivi; alzare la mano e fare domande.
\textbf{Tolleranza zero} per pregiudizi, maleducazione, mancanza di rispetto,
\emph{cheat} e scorciatoie. Daniele ammette che ``cheat e scorciatoie'' viene da
altri corsi dove il tema è più importante, ma Ruggero lo ferma: ``non
sottovalutate questa cosa'', perché implica anche \textbf{non far fare a ChatGPT
i task}. Daniele ha una slide apposta, ``perché è un discorso molto complesso nel
2026'' (sezione~\ref{a1:sec:ia}).

\paragraph{Le lezioni sono intense.} Le lezioni possono essere \textbf{dense},
``soprattutto le mie, sappiatelo: io sono logorroico'', avverte Daniele. Bisogna
chiedere spiegazioni se serve, prendere appunti, rivedere le registrazioni. Le
lezioni sono ``delle scintille'': modi per far esplorare i temi oltre le slide.
``Andate oltre, se riuscite''.

\paragraph{Il carico di lavoro.} Le regole sui task sono queste:
\begin{itemize}
  \item i task \textbf{sono per voi, non per noi};
  \item se si ha poco tempo, è meglio farli rapidamente piuttosto che non farli:
        ``provateci'';
  \item i docenti cercano di non chiedere più di \textbf{un'ora o due a
        settimana} oltre alle ore di lezione;
  \item non tutte le lezioni avranno un task, e il \textbf{project work} è il più
        importante di tutti.
\end{itemize}

\paragraph{I feedback dei docenti.} I feedback ``si guadagnano'' man mano che si
reagisce bene: non bisogna aspettarsi ``dei papiri'' sui task. Se i docenti
hanno voglia di dare feedback più ampi lo fanno, soprattutto se vedono che lo
studente se li merita. Ad Analog di solito non ci sono task \emph{rejected}; ma
se un docente rifiuta un task, vuol dire che si è fatto \textbf{meno del livello
minimo} richiesto dal corso. I feedback possono non essere immediati: ``siate un
pochettino impietositi dal fatto che siamo dei docenti''. E se un feedback non
arriva, o se ne vogliono di più, bisogna scrivere, ``perché siamo tre esseri
umani e quindi possiamo anche dimenticarci''.

% ---------------------------------------------------------------------
\section{Il project work e il calendario}
\label{a1:sec:project}

\subsection{Come funziona il project work}

Una prima slide con le date del project work è da aggiornare (``perdonatemi il
refuso''); Daniele la corregge al volo e rimanda alla slide successiva, più
precisa. Il funzionamento è questo:
\begin{itemize}
  \item i \textbf{gruppi li formano i docenti}. Gli studenti sono nove, quindi si
        faranno \textbf{tre gruppi da tre}: una dimensione pensata anche per il
        fatto che molti sono studenti lavoratori, così da non ritrovarsi ``in due,
        a esplodere'';
  \item ogni gruppo deve prima preparare un \textbf{pitch} del proprio gioco, da
        presentare in anticipo: un'idea, una bozza del gioco;
  \item poi si lavora al gioco \textbf{nelle pause del corso}: da qui a dicembre
        ci saranno due pause, quindi non ci sarà lezione tutti i lunedì e tutti i
        venerdì;
  \item ci saranno \textbf{review intermedie} con i \textbf{tutoraggi}, per lo più
        il \textbf{sabato mattina}, dalle 11 alle 12, con Ruggero su Google Meet.
        Servono ad aiutare i gruppi a portare avanti il project work.
\end{itemize}

\subsection{Il calendario delle lezioni}

Il \textbf{calendario} è nel calendario studenti della cartella condivisa.
Attenzione: le date \textbf{possono variare}. ``Il 90\% rimarranno così'', ma
può capitare che, per motivi di lavoro, qualcuno stia male o non possa, o che
``scompaia l'elettricità da tutte le città d'Italia''. E l'ordine delle lezioni
può cambiare se i docenti vedono che agli studenti serve prima una lezione
invece di un'altra.

Ci sono \textbf{due pause importantissime}:
\begin{itemize}
  \item dal \textbf{13 ottobre al 5 novembre}: è il periodo delle grandi fiere
        (\textbf{Essen} e \textbf{Lucca}\footnote{La trascrizione dice ``Essen,
        Luca, Cose'': si tratta con ogni probabilità della fiera Spiel di Essen,
        in Germania, e di Lucca Comics \& Games, che si tengono entrambe in
        ottobre.}), ``un periodo un po' frenetico'', e cominciano anche i corsi
        accademici, quindi la pausa permette ai docenti di respirare. Per gli
        studenti è la pausa più importante: quasi un mese in cui sviluppare la
        parte importante del gioco;
  \item dal \textbf{24 novembre all'11 dicembre}, dopo qualche lezione a novembre.
\end{itemize}

E ci sono \textbf{due date da evidenziare}:
\begin{itemize}
  \item il \textbf{25 settembre}: nei primi 5--10 minuti della lezione (che per il
        resto sarà tenuta da Beatrice) i docenti annunceranno i tre gruppi;
  \item il \textbf{12 ottobre}: il \textbf{pitch} dei giochi, con l'idea
        embrionale. Tra il 25 settembre e il 12 ottobre ogni gruppo prepara un
        breve pitch in formato slide; quel giorno i docenti lo rivedono e danno
        feedback: se il gioco va bene, se non va bene, come svilupparlo.
\end{itemize}
Dal 12 ottobre fino all'inizio di dicembre si sviluppa il gioco, con i
tutoraggi: ``sfruttate questi momenti di pausa''.

\subsection{Il calendario del project work}

La slide successiva, ``reale e aggiornata'', riporta il calendario del project
work (tabella~\ref{a1:tab:calendario} e figura~\ref{a1:fig:calendario}).
``Per ora non pensate al project work: ci penserete dal 25 settembre in poi''.

\begin{table}[htbp]
\centering
\begin{tabular}{@{}lp{9.2cm}@{}}
\toprule
\textbf{Data} & \textbf{Che cosa succede} \\
\midrule
ven 25 settembre\footnotemark & lezione: annuncio dei gruppi e del
\emph{mandate}, cioè di che cosa si dovrà fare \\
sab 3 ottobre, ore 11 & prima review e tutoraggio con Ruggero (dopo una
decina di giorni di lavoro) \\
lun 12 ottobre & pitch del gioco durante la lezione \\
13 ottobre -- 5 novembre & prima pausa: sviluppo del gioco \\
sab 7 novembre & tutoraggio \\
sab 14 novembre & tutoraggio \\
24 novembre -- 11 dicembre & seconda pausa: sviluppo del gioco \\
ven 4 dicembre, entro mezzanotte & consegna del gioco sul Drive (non c'è
lezione) \\
ven 11 dicembre & i giochi vengono visti a lezione \\
14 dicembre & giornata finale, con la tavola rotonda con gli editori \\
\bottomrule
\end{tabular}
\caption{Il calendario del project work come lo presenta Daniele. Le date
possono variare.}
\label{a1:tab:calendario}
\end{table}
\footnotetext{Daniele dice ``lunedì 25 settembre'', ma nel 2026 il 25 settembre
è un venerdì; il giorno della settimana nella tabella è corretto di
conseguenza.}

\paragraph{Perché la consegna è il 4 dicembre.} Il 4 dicembre non c'è lezione:
si deve solo consegnare il gioco, entro mezzanotte, sul Drive. La presentazione
a lezione è l'11 dicembre, ma i tre docenti hanno bisogno di \textbf{una
settimana} per rivedere il gioco in privato e valutare anche la
\textbf{documentazione}; per questo la consegna non coincide con la
presentazione. Che cosa voglia dire esattamente ``consegnare il gioco'' verrà
spiegato più avanti.

\paragraph{La giornata finale.} Il 14 dicembre c'è la \textbf{giornata finale},
con una \textbf{tavola rotonda con gli editori}: di solito viene invitato Flavio,
che è anche docente di game production, più qualche altro editore, ``vediamo
quest'anno chi riusciamo a portarvi''.

\begin{figure}[htbp]
\centering
\begin{tikzpicture}[
    font=\footnotesize,
    x=0.1cm,
    evento/.style={draw, thick, rounded corners=2pt, fill=white,
                   text width=1.8cm, align=center, inner sep=2pt},
    pausa/.style={fill=gray!20, draw=gray!60},
    tacca/.style={thick}
  ]
  % asse: giorno 0 = 25 settembre; 1 unità = 1 giorno
  \draw[pausa] (18,-0.35) rectangle (41,0.35);
  \draw[pausa] (60,-0.35) rectangle (77,0.35);
  \node[font=\scriptsize\itshape] at (29.5,0.18) {pausa};
  \node[font=\scriptsize\itshape] at (65,0.18) {pausa};
  \draw[-{Stealth[length=2.5mm]}, thick] (-3,0) -- (85,0);
  % eventi sopra
  \draw[tacca] (0,-0.15) -- (0,0.9);
  \node[evento, above] at (0,0.9) {25 set\\ gruppi e \emph{mandate}};
  \draw[tacca] (17,-0.15) -- (17,2.2);
  \node[evento, above, fill=gray!8] at (17,2.2) {12 ott\\ \textbf{pitch}};
  \draw[tacca] (70,-0.15) -- (70,0.9);
  \node[evento, above, fill=gray!8] at (70,0.9) {4 dic\\ \textbf{consegna}};
  % eventi sotto
  \draw[tacca] (8,0.15) -- (8,-0.9);
  \node[evento, below] at (8,-0.9) {3 ott\\ tutoraggio};
  \draw[tacca] (43,0.15) -- (43,-0.7);
  \draw[tacca] (50,0.15) -- (50,-0.7);
  \draw (43,-0.7) -- (50,-0.7);
  \draw[tacca] (46.5,-0.7) -- (46.5,-0.9);
  \node[evento, below] at (46.5,-0.9) {7 e 14 nov\\ tutoraggi};
  \draw[tacca] (77,0.15) -- (77,-0.9);
  \node[evento, below] at (77,-0.9) {11 dic\\ giochi a lezione};
  \draw[tacca] (80,-0.15) -- (80,2.2);
  \node[evento, above] at (80,2.2) {14 dic\\ giornata finale};
\end{tikzpicture}
\caption{Una linea del tempo del project work, ricostruita dalle slide (scala
approssimativa: dal 25 settembre al 14 dicembre). Le bande grigie sono le due
pause senza lezioni, in cui i gruppi sviluppano il gioco.}
\label{a1:fig:calendario}
\end{figure}

Finita questa parte, Daniele apre alle domande: ``Domande sul calendario? Sulle
slide che ho fatto? Considerazioni? Non abbiate timore, io non ci credo''. E la
domanda arriva, sull'unico tema che aveva lasciato da parte.

% ---------------------------------------------------------------------
\section{L'intelligenza artificiale nel corso}
\label{a1:sec:ia}

\subsection{La slide di Daniele}

Daniele aveva già presentato una slide che definisce ``importantissima''. ``Lo so
che è paradossale, me ne rendo conto. Lo so che magari mi darete del luddista,
dell'anti-tecnologia: non è così. Uso l'intelligenza artificiale per lavoro. Ma
in questo corso, vi prego, \textbf{non usate l'intelligenza artificiale}''. La
richiesta vale per le attività, per i task e \textbf{anche per la
prototipazione}. Le ragioni:
\begin{itemize}
  \item ``voglio che siate voi a fare il gioco, voglio che siate voi a imparare a
        farlo''. Daniele parla da utilizzatore: prima di ChatGPT, che ha smesso
        di usare, oggi soprattutto di Claude, di cui ``si sottovaluta quanto
        possa essere incerto, insicuro, impreciso, non all'altezza'';
  \item nella game industry, \textbf{i publisher rifiutano automaticamente} i
        pitch e i prototipi fatti con l'intelligenza artificiale. Daniele lo ha
        visto di persona in ``tantissime pitch review'' quell'estate, anche al
        \textbf{First Playable} di Firenze\footnote{La trascrizione dice ``First
        Playboy a Firenze'': si tratta del First Playable, l'evento di business
        dell'industria videoludica che si tiene a Firenze.}: anche nel mondo dei
        videogiochi i publisher che vedono pitch o prototipi fatti con l'IA
        ``rigettano automaticamente la build'', perché hanno ancora timore
        dell'IA e delle sue possibili conseguenze;
  \item ``non so come sarà il futuro'': non sa se bisognerà abituarsi a usarla
        tutto il giorno. Non è ``totalmente contrario'', purché sia usata ``con
        raziocinio''; secondo lui oggi, negli Stati Uniti, non lo è;
  \item ma in un corso che si chiama \emph{Analog Game Design} gli studenti devono
        saper fare i giochi da soli. In futuro, se vorranno farsi aiutare
        dall'IA per essere più efficienti, saranno liberissimi di farlo; durante
        il corso no, perché i docenti vogliono valutare \textbf{le loro
        competenze e capacità}.
\end{itemize}

\subsection{La domanda: che cosa vuol dire usare l'IA in un gioco analogico?}

La domanda di uno studente è precisa: che cosa si intende per ``task fatti con
l'IA''? Qual è l'uso dell'IA in un gioco analogico?

\paragraph{Ruggero: l'esperienza di chi valuta i prototipi.} Ruggero risponde
dalla sua esperienza nel settore, dove riceve prototipi da valutare per
un'eventuale pubblicazione. Distingue due usi:
\begin{enumerate}
  \item le \textbf{immagini generate}: illustrazioni, graphic design e così via.
        Per molti publisher è ``un \emph{reject} immediato, proprio immediato'';
  \item l'IA usata per il \textbf{regolamento}, per scrivere le regole e le
        meccaniche. È più difficile da individuare, ma molti publisher se ne
        accorgono comunque, perché le IA ``non sono ancora brave abbastanza da
        creare un regolamento coeso meccanicamente e scritto con un senso per il
        vostro gioco''. Quello che fanno è prendere parti di altri regolamenti,
        elaborare le strutture che vedono nei regolamenti pubblicati e
        riapplicarle. Ma molti giochi si prestano poco a un regolamento
        strutturato come quello di altri, perché \textbf{ogni gioco è un
        ecosistema meccanico a sé}. Il risultato sono regolamenti con
        riferimenti non immediati, e soprattutto con \textbf{la stessa meccanica
        descritta in modo diverso in pagine diverse}: due difetti che fanno
        perdere tempo all'editore, e in questo caso ai docenti.
\end{enumerate}
Poi aggiunge un punto ``personale, soggettivo'': i docenti sono interessati a
vedere il lavoro, la creatività, ``il frutto della vostra creazione''. E
viceversa gli studenti vogliono i commenti dei docenti: ``immaginatevi se voi
fate un gioco e noi vi diamo i commenti scritti con ChatGPT''. Altrimenti, dice
ridendo, gli studenti fanno il gioco con ChatGPT, i docenti lo danno a ChatGPT
per i commenti, ``e facciamo molto alla svelta''.

\paragraph{Daniele: perché gli studenti usano l'IA.} Daniele racconta la sua
esperienza con gli studenti, che secondo lui ricorrono all'IA in due momenti:
\begin{enumerate}
  \item quando vanno in \textbf{crisi}: in ansia da prestazione o per i tempi.
        Cercano di risolvere il problema della propria disorganizzazione
        facendolo risolvere all'IA, senza accorgersi che l'output è ``molto,
        molto scarso'';
  \item quando hanno l'impressione che l'intelligenza dell'IA sia
        \textbf{migliore} di quello che scrivono loro, o di quello che c'è in
        giro. Per Daniele è una ``mancanza di palato critico'': l'output dell'IA
        sembra buono in superficie, ma quando si va nel dettaglio e ci si fa
        delle domande ``di solito fa un po' di \emph{crumbling}'', si sbriciola.
\end{enumerate}
Lui usa Claude ogni giorno e ne conosce i limiti: ``non c'è nemmeno una frase
che io copio direttamente''. Gli fa fare il lavoro più noioso, più ripetitivo e
soprattutto \textbf{strutturato}, ``replicare delle cose come se fosse una
macro''. Lo usa per il testo e per il codice, mai per le immagini: per motivi di
\textbf{diritto d'autore}, e perché ha molti amici artisti ed è stufo di vedere
il loro lavoro ``dilaniato'' dall'IA.

Ammette che i docenti non possono controllare gli studenti; per questo
``ve lo consiglio''. E riprende quella che definisce la cosa più bella detta da
Ruggero (``ti ruberò la frase per il corso di narrative''): i docenti i feedback
li danno di persona. Daniele a volte manda 30 minuti di audio a uno studente, o
fa intere call per rivedere script, pitch e giochi: ``è il mio lavoro, mi piace
farlo e voglio farlo io, quindi mi aspetto che lo studente lo faccia lui''. Come
freelance, poi, è ``stufo marcio'' dei clienti che arrivano con qualcosa fatto
con l'IA (``guarda che bello, pubblica questo''), perché il più delle volte il
risultato è pessimo: non sanno come usarla. ``È anche una questione di
competenza''.

\paragraph{Il prototipo non deve essere bello.} Questo vale anche per il
\textbf{prototipo}, su cui si tornerà nelle lezioni di prototipazione: in un
prototipo non servono ``graficoni allucinanti''. Le grafiche bellissime fatte con
l'IA spesso \textbf{peccano di user experience}: sembrano bellissime, ma quando
il giocatore deve leggere la carta non capisce qual è il costo, qual è il
potere, qual è l'effetto.

\paragraph{Beatrice: si perde l'occasione di imparare.} Beatrice aggiunge due
motivi che valgono ``al di là dei nostri credo''. Il primo: chi usa l'IA in
sostituzione (per esempio per scrivere un testo, per accelerare i tempi o per
trovare strategie) \textbf{perde la parte in cui può imparare}. Nel game design
bisogna ``tenere insieme i pezzi'' di un regolamento, ma anche dell'esperienza
dell'utente e di come trasmetterla; i passaggi necessari (creare il regolamento,
creare le meccaniche, sperimentare, raccogliere i feedback) ``si capiscono solo
vivendoli''. E si sta pagando un corso per imparare. Il secondo, sul piano
editoriale: chi presenta a un editore un prodotto che ha come base l'IA rischia
da un lato di non avere imparato abbastanza (e di ritrovarsi ``a piedi'' sul
campo senza IA), dall'altro che l'editore se ne accorga, ed è un tema
\textbf{reputazionale}. Ci sono editori ``più illuminati di altri'', quindi
anche qui c'è un mercato da capire; ma i docenti preferiscono dire che cosa è
più sensato per il mercato come lo conoscono.

\paragraph{Ruggero: non abbiate paura del giudizio.} Ruggero vuole dare uno
``spotlight'' a un punto detto da Daniele e rimasto ``semi-ignorato'': a volte
l'uso dell'IA nasce dalla \textbf{paura del giudizio}, dal voler consegnare
qualcosa di più pulito, più strutturato possibile. ``Vi giuro, non ve ne
preoccupate''. Ruggero detesta i voti e non ragiona per voti; non valuterà
nessuno per grafiche fatte a matita. Insegna da anni, ha avuto studenti con
dislessia e discalculia: refusi ed errori non gli interessano, ``non è un corso
di editing o di proofreading''. Gli interessa che le \textbf{meccaniche siano
coerenti con lo scopo di design}. Immagini fatte a matita o con Paint ``non sono
un problema, anzi, le preferisco'' rispetto all'``ennesima immagine fatta con
ChatGPT o Midjourney''.

\paragraph{Daniele: il pensiero critico, costo e qualità.} L'ultimo intervento
è di Daniele, sul \textbf{pensiero critico}. Ragionando per assurdo (``un
pensiero di Belzebù''), immagina che un domani l'IA venga usata ovunque. I suoi
limiti, secondo lui, resterebbero due:
\begin{description}
  \item[Il costo] Tutti dicono che sta diminuendo e che diminuirà, ma non se ne
        può essere sicuri al 100\%.
  \item[La qualità] Cioè la capacità di leggere l'output. Se si dà un regolamento
        a Claude o a ChatGPT chiedendo quali problemi ha, il modello risponde con
        i problemi che, in percentuale, è più probabile che gli esseri umani
        abbiano già individuato nei dati su cui è stato addestrato. Ne derivano
        due problemi: se il problema del proprio gioco non è mai stato
        individuato prima, il modello non lo vedrà; e viceversa può segnalare
        problemi che il gioco non ha.
\end{description}
Fa l'esempio di Grammarly, che usa nella versione premium per rivedere i testi
in inglese: spesso non accetta i suoi consigli, perché se vuole scrivere un
dialogo di videogioco in \textbf{dialetto Cockney} londinese, lo deve scrivere
in Cockney. Se Grammarly dice che è sbagliato, ``sei tu che non hai capito il
motivo per cui voglio scrivere questo testo in Cockney''. La conclusione: prima
di usare un'IA per un'attività bisogna sapere \textbf{che cosa si vuole
ottenere}, qual è l'obiettivo, e quali \textbf{competenze} servono per valutare
in modo critico se l'output di un qualunque LLM è davvero valido. Daniele non
vuole dilungarsi oltre: è un argomento ``molto caldo'' per lui, che è anche
scrittore e si sente dire da tutti ``ti sostituiranno''.

Subito dopo si presenta finalmente Gianmarco (lo si è visto nella
sezione~\ref{a1:sec:partecipanti}), e poi si fa una pausa di una decina di
minuti. Daniele avverte che nella terza parte lui non ci sarà (alle sette e un
quarto ha un impegno di lavoro), e lo stesso vale per Beatrice: la parte finale
di domande e risposte sarà con il solo Ruggero, ``in ottime mani''.

% ---------------------------------------------------------------------
\section{Le figure professionali dell'industria analogica}
\label{a1:sec:figure}

Dopo la pausa (in cui si scopre che Angelo lavora in un piccolo birrificio in
Brianza e si scherza sul rischio di un \emph{drinking game}, tra la birra e la
vendemmia) la lezione riprende con le slide di Ruggero. ``Oggi è una lezione
introduttiva'': gli argomenti si vedranno ``molto dall'alto'' e verranno
approfonditi nelle lezioni successive. Ma servono già ora, per creare
\textbf{l'inizio di un linguaggio condiviso}: ci sono termini che torneranno
spesso, a cominciare da quelli del mercato e dell'industria del gioco
analogico, un'industria ``molto varia, molto caotica, anche molto mutevole'',
con tanti tipi diversi di giochi e altrettanti tipi diversi di entità che ci
lavorano.

\subsection{Una decina di etichette}

I docenti hanno raccolto ``una decina di etichette'', cioè di \textbf{ruoli} che
nell'industria sono o possono essere presenti. Due avvertenze:
\begin{itemize}
  \item nelle aziende più piccole \textbf{la stessa persona copre spesso più
        ruoli}, per forza di cose. Sarà anche il caso degli studenti, che
        lavoreranno in gruppi di tre e non di venti: è facile che ciascuno finisca
        a fare più cose. I docenti ne terranno conto, e qui entreranno in gioco
        anche competenze diverse dal game design (programmi di grafica,
        illustrazione, impaginazione, layout, \emph{proofreading}\ldots);
  \item Daniele precisa che \textbf{non si chiede agli studenti di diventare}
        illustratori o altro. Ruggero concorda, ma aggiunge che è utile sapere
        anche \emph{che cosa non è il proprio lavoro}. Se in un gruppo nessuno sa
        illustrare, ai docenti va benissimo un prototipo che mostri dove vanno
        le immagini: \textbf{placeholder}, concept, anche semplici box di testo
        con scritto ``qui vorremmo un'immagine di\ldots''. ``È importante anche
        sapere che cosa sapete fare e che cosa non sapete fare''.
\end{itemize}
La tabella~\ref{a1:tab:figure} riassume le figure; le sottosezioni seguenti le
riprendono una per una, con i commenti dei docenti.

\begin{table}[htbp]
\centering
\small
\begin{tabular}{@{}p{3.3cm}p{9.6cm}@{}}
\toprule
\textbf{Figura} & \textbf{Di che cosa si occupa} \\
\midrule
Game designer & crea le regole alla base dell'esperienza: le meccaniche, come
interagiscono tra loro e come chi gioca interagisce con esse \\
Game developer & fa da ponte tra chi fa game design e il publisher: bilancia il
prototipo, cura i modelli matematici, adatta il gioco (e il paragame) alle
necessità del publisher \\
Editor (proofreader) & revisiona e adatta i testi a partire dalla prima stesura
del designer (correzione, linguaggio inclusivo\ldots) \\
Publisher (editore) & pubblica e produce fisicamente il gioco; il suo marchio è
sulla scatola o sulla copertina \\
Sales e marketing & il lato commerciale: vendite (anche B2B) e comunicazione del
prodotto \\
Graphic designer & icone meccaniche, wireframe, UX: il lato funzionale degli
elementi visivi \\
Illustratore & le illustrazioni: il \emph{look and feel}, il world building, il
``vibe'' del gioco \\
Sensitivity reader (cultural consultant) & verifica l'accuratezza e la
rappresentazione di un tema sensibile che il team non conosce \\
Producer & coordina il progetto (quasi project management), parla con le
tipografie, verifica che i materiali siano sostenibili economicamente \\
\bottomrule
\end{tabular}
\caption{Le figure professionali dell'industria del gioco analogico presentate
nella slide di Ruggero, in sintesi.}
\label{a1:tab:figure}
\end{table}

\subsection{Game designer}

Il \textbf{game designer} (o la game designer) è chi si occupa, ``sembra banale
a dirlo'', di \textbf{creare le regole alla base dell'esperienza} che si vuole
far arrivare. Vale per tutti i tipi di giochi, non solo per quelli analogici, e
che cosa faccia concretamente in un progetto o in una serie specifica ``può
mutare tantissimo''. Nell'analogico non verrà chiesto di programmare o di
prototipare in digitale; verrà chiesto di \textbf{creare le meccaniche di
gioco}, di capire \textbf{come le meccaniche si interfacciano tra loro} e
\textbf{come le persone che giocano si interfacceranno con le meccaniche}.

\subsection{Game developer}

Il \textbf{game developer} (o la game developer) sarà l'argomento delle due
lezioni di Marta. La parola \emph{development} è nota soprattutto
dall'ambito videoludico, ma nell'analogico ha un significato molto specifico
``che non c'entra niente con la programmazione'': il game developer è
\textbf{un ponte tra chi fa game design e il publisher}. Di solito si occupa di
\textbf{bilanciare il prototipo}, di curare i \textbf{modelli matematici}, di
adattare il gioco e il \textbf{paragame}\footnote{La trascrizione dice ``il
Wix, il paragame''. Il \emph{paragame} è, in termini generali, tutto ciò che
accompagna il gioco senza esserne il sistema di regole (regolamento, scatola,
materiali, comunicazione); i docenti lo tratteranno in una lezione dedicata, e
qui il termine è solo citato.} alle necessità del publisher. La distinzione
viene ripresa e approfondita nella sezione~\ref{a1:sec:design-development}.

\subsection{Editor}

L'\textbf{editor}, a volte chiamato \textbf{proofreader} a seconda dell'azienda,
si occupa della revisione dei \textbf{testi}: non della stesura iniziale, perché
di norma, data la scala dell'industria, è chi fa game design a scrivere un
\textbf{primo draft}, e l'editor ci passa sopra. Non solo per correggere i
refusi: spesso il suo lavoro è anche \textbf{adattare} i testi, per esempio
renderli \textbf{gender neutral}. Se ne vedranno casi studio più avanti.

\subsection{Publisher}

Il \textbf{publisher}, o \textbf{editore}, è l'entità che nell'ambito
commerciale \textbf{pubblica} giochi, giochi di ruolo o giochi da tavolo. È una
figura ``più solidificata'' nel gioco da tavolo e nel gioco di ruolo da tavolo;
\textbf{nell'ambito pervasivo non esiste proprio in questa forma}, come
mostrerà Beatrice parlando dei casi studio. Se si compra un gioco da tavolo o un
gioco di ruolo in un negozio, il marchio del publisher è di solito sul retro
della copertina o sul fronte della scatola: è l'entità che, di solito,
\textbf{produce fisicamente} il gioco.

\subsection{Sales e marketing}

L'ambito \textbf{sales e marketing}, a volte affidato a un'unica persona e a
volte separato, riguarda il lato commerciale. In alcune aziende è interno al
publisher, e si occupa di parlare con le fabbriche e le tipografie, di ``far
quadrare i conti''. In altre realtà è completamente diverso, specialmente nel
\textbf{B2B} o nei giochi \textbf{didattici} e \textbf{formativi}, di cui
Daniele e Beatrice ``sapranno parlare meglio''. Il tema viene ripreso più avanti
in dettaglio (sezione~\ref{a1:sec:marketing}).

\subsection{Graphic designer e illustratore}

Ruggero presenta insieme \textbf{graphic designer} e \textbf{illustratore}
proprio perché sono ruoli molto diversi che si tende a far coincidere:
\begin{itemize}
  \item chi fa \textbf{graphic design} si occupa delle \textbf{icone meccaniche},
        del \textbf{wireframing}, della \textbf{UX}: del lato \textbf{meccanico e
        funzionale} degli elementi visivi, cioè di tutto ciò che serve
        meccanicamente nel gioco, o anche alla sua ``qualità della vita'';
  \item chi fa \textbf{illustrazione} si occupa dell'\emph{art}, delle
        illustrazioni: copre il \textbf{look and feel}, la parte di
        \textbf{world building} e il ``vibe'' che si dà al gioco.
\end{itemize}

\paragraph{Daniele: che cos'è la user experience.} Daniele interviene sulla
\textbf{user experience}, ``importantissima per me'': ha un background da
\textbf{layout artist} (impaginava magazine), un mestiere in cui la UX conta
moltissimo. Per user experience non si intende solo l'\textbf{interfaccia
grafica}, la UI, che anche nei videogiochi è disegnata da un artista. La UX è
qualcosa che può fare anche il game designer: \textbf{disegnare l'esperienza del
giocatore}, cioè ``il gesto, l'intenzione, l'intuito'' che stanno dietro
l'aspetto visivo dell'interfaccia. Per fare UX non serve saper disegnare.
L'esempio è quello della \textbf{porta}: se si progetta una porta, bisogna
indicare all'utente come si apre, come si prende la maniglia; ``quella è user
experience''. Il design della maniglia in sé, invece, è
\textbf{interfaccia}.

\paragraph{Un aneddoto sulle icone.} Ruggero aggiunge che molte persone che
fanno \textbf{graphic design ben fatto non sanno illustrare}, e spesso vale il
viceversa. Lo si nota quando le icone di un gioco sono bellissime ma non
funzionali: andando a guardare i crediti, di solito si scopre che le icone le ha
fatte la persona che ha fatto le illustrazioni. ``Comunque non andiamo adesso in
questo, perché è un altro vaso di Pandora''.

\subsection{Sensitivity reader}

Il \textbf{sensitivity reader}, o \textbf{cultural consultant} a seconda del
progetto, fa un lavoro di \emph{sensitivity reading} del gioco. È una figura
molto usata nell'ambito \textbf{formativo e didattico}. L'esempio di Ruggero: se
si fa un gioco per un ospedale su una malattia specifica, chi lo progetta non è
un medico, ma il gioco deve essere \textbf{accurato} e \textbf{rappresentativo}
delle dinamiche che quella malattia comporta. Serve una persona che conosca quel
tema, anche se non sa che cosa sia il game design. Sono figure chiamate di
solito \textbf{a progetto}, per giochi specifici, che compensano competenze che
il team non ha ``per forza di cose, perché noi non sappiamo tutto''; e a volte,
per lavoro, si devono fare giochi su temi che non si conoscono personalmente.

\paragraph{Daniele: non è un censore.} Daniele segnala un problema italiano:
alcune persone, anche online, hanno confuso il sensitivity reader con un
\textbf{censore}, ``una persona che fa censura, soprattutto in materia di
\emph{woke}''. ``Assolutamente non è così''. La parola \emph{sensitive} non
significa che il pubblico sia sensibile a un tema, ma che \textbf{l'argomento di
cui tratta il gioco è sensibile}. In un gioco sulla medicina, realistico o no,
non si possono inserire imprecisioni. Per questo esistono i sensitivity reader
``più disparati'': da quelli per la cultura indiana o afroamericana fino a quello
su ``come si allevano gli gnu in Africa''. È una persona competente su un
argomento che aiuta a essere \textbf{veritieri}. Certo, uno degli ambiti
principali in cui si usano è quello culturale, ma ``non fate l'errore di pensare
che siano dei censori''. Daniele ha lavorato con ``centinaia'' di sensitivity
reader e nessuno gli ha mai detto ``non puoi scrivere questa cosa''. Di solito
sono come degli \textbf{editor}: danno opzioni, mostrano i punti problematici,
anche quelli di cui non ci si era nemmeno accorti.

Ruggero conferma con un esempio di quando lavorava in \textbf{Danimarca} con
l'\textbf{OMS}: giochi sul \textbf{diabete di tipo 2} nei bambini, per un target
di ragazzi delle scuole medie. C'era un team dell'OMS specializzato nella
comunicazione sul diabete, per garantire che il gioco fosse accurato e ne desse
una rappresentazione adatta al target. Il sensitivity reader aiuta a rendere il
gioco più accurato e ``più tematicamente coeso''.

\subsection{Producer}

Il \textbf{producer} è una figura a volte interna al publisher. In un progetto
grande è la figura di riferimento, ``quasi'' un project manager. In un progetto
piccolo si occupa spesso anche di parlare con le \textbf{tipografie} e di fare la
\textbf{lista dei materiali}, per assicurarsi che i materiali non siano solo
funzionali al gioco ma anche \textbf{economicamente sostenibili} da produrre:
``magari è un gioco bellissimo, ma poi costa più di quanto potete permettervi, e
nessuno vi comprerà mai un card game a 200 euro''.

\subsection{Perché conoscere queste figure}

Ruggero chiude la slide ``tornando a ombrello'': perché parlarne, se gli studenti
sono lì per fare game design? Non solo per saper dare un \textbf{nome} alle
figure. Non c'è niente di male a creare un gioco solo per sé; ma chi vuole fare
giochi a livello lavorativo, perché li giochino altre persone, ``magari più
persone possibile'', \textbf{si interfaccerà con queste figure}, e sapere che
cosa fanno e come parlarci è parte importante del percorso. Nel corso se ne
toccheranno diverse: il development con Marta, la produzione con Flavio, e
lezioni sulla UX, sul paragame, sul graphic design. ``Abituatevi mentalmente:
non siete in solitaria contro il mondo''.

% ---------------------------------------------------------------------
\section{Marketing e sales}
\label{a1:sec:marketing}

Arrivato all'angolo della slide con marketing e sales, Ruggero passa la parola
ai colleghi: ``io è meglio che non mi fate parlare di marketing e sales'', anche
per esperienze personali.

\subsection{Beatrice: i giochi trattati come prodotti}

Beatrice premette che le due voci sono in parte autoesplicative. Nel mondo
capitalista in cui viviamo, i giochi sono ``strumenti molto potenti per fare
\emph{n} cose'', ma sono trattati come \textbf{prodotti}, e il corso ragionerà
anche in quest'ottica. Nel \textbf{marketing} delle industrie ludiche e
videoludiche ci sono varie professionalità che gli studenti probabilmente già
conoscono: dal \emph{copy} alla grafica, tutto ciò che riguarda il
\textbf{vendere} il gioco o l'attività di gioco.

Il lato \textbf{sales} è forse la distinzione meno nota. Come in tante aziende, è
la parte che si occupa delle vendite \textbf{B2B}, cioè tra aziende: fare
accordi per vendere prodotti ludici a diverse realtà, per esempio ai
\textbf{distributori}, o a una casa editrice che vende alle librerie, che poi
vendono al consumatore finale. Chi si occupa di sales nel mondo ludico capisce
come vendere \textbf{alle aziende}. I due mondi si parlano spesso, e soprattutto
nelle realtà piccole si trovano figure che fanno ``un po' sales, un po'
marketing''.

\subsection{Daniele: che cos'è (e che cosa non è) il marketing}

Daniele il marketing si è trovato ``purtroppo'' a farlo, e lo fa anche per
l'accademia. ``Purtroppo'' perché in Italia è mal compreso: molti lo
identificano con l'\emph{advertisement} online. Non è così: il marketing è una
\textbf{disciplina molto più antica}, e in Italia ``abbiamo molto dimenticato''
come si fa il marketing tradizionale. Alcune cose da ricordare:
\begin{enumerate}
  \item \textbf{Il marketing non dice che cosa pubblicare, ma come pubblicarlo}.
        Esistono, nei videogiochi e nei giochi da tavolo, realtà che fanno
        prodotti \textbf{marketing driven}: l'esempio sono gli
        \textbf{hyper-casual game} per smartphone, quelli con il personaggio che
        va a destra e a sinistra e l'esercito che si allarga. Possono anche
        guadagnare, ma ``sono un po' la morte della creatività''. Il marketing
        ideale \textbf{prende un prodotto pensato da altri e gli dà le maggiori
        chance di essere venduto}: cerca il pubblico ideale per il prodotto che
        si ha, senza doverlo cambiare. Anzi, il mondo videoludico, e ludico in
        generale, ``è lastricato'' di errori di dipartimenti marketing che hanno
        cambiato completamente non il gioco, ma il modo in cui l'hanno
        pubblicizzato, e il gioco è fallito.
  \item \textbf{Usare strumenti del marketing non significa fare marketing}. Game
        design e game development a volte usano strumenti che vengono dal
        marketing (gli studenti ne vedranno con Beatrice, con Daniele e con
        Marta).
  \item \textbf{Il marketing è un gerundio}: si fa. È un'azione che si compie per
        comunicare; la comunicazione è attiva, e il marketing è l'insieme delle
        attività con cui si comunica il proprio gioco, dal \emph{copy} ai social.
        Può includere il digital marketing, ``ma non è solo quello''.
  \item \textbf{Il marketing esiste anche se si crede che non esista}. In Italia
        molti non vogliono ``fare marketing'' perché lo associano al lucro. Ma
        il marketing è \textbf{lo studio di come gli esseri umani comprano e
        vendono}, ed è innato in noi: ci sono logiche di cui siamo vittime, e
        ``anche se sei il bastian contrario, l'anticapitalista'', magari si
        compra comunque su Amazon perché fa comodo. Sono meccanismi che
        esistono, ``affascinanti anche, a volte''.
\end{enumerate}
Gli studenti non dovranno fare marketing, ma dovranno ricordare alcuni di questi
concetti, soprattutto nelle lezioni con Marta, che tratta argomenti legati al
marketing. E, chiude Daniele, come game designer e game developer ``per lo più
speriamo che il marketing non ci rovini i giochi'', soprattutto nel videoludico.

% ---------------------------------------------------------------------
\section{Game design e game development}
\label{a1:sec:design-development}

Ruggero torna sulla distinzione più importante della slide, ``proseguendo con
questa strada'':
\begin{description}
  \item[Game design] è chi \textbf{progetta l'esperienza} tramite la
        \textbf{creazione delle meccaniche} di gioco.
  \item[Game development] è chi fa \textbf{mediazione} tra la persona che scrive
        il game design e la \textbf{casa editrice}.
\end{description}

\subsection{I red flag dei designer}

Ruggero lo ripete perché lo vede ogni giorno nel suo lavoro: tra giochi
pubblicati e prototipi, ne valuta ``quasi uno al giorno, se faccio un conto
annuale'', e solo una minima percentuale arriva alla pubblicazione. Molte volte
i \textbf{red flag} vengono proprio da designer che non hanno presente questa
definizione, e che cominciano a mettere \textbf{clausole sulla non modifica del
gioco}: il gioco deve restare com'è. ``Questa cosa è impossibile con un
publisher, spesso''. Il publisher, per forza di cose, potrà chiedere di cambiare
la grafica, di reimpaginare il regolamento, di cambiare i materiali perché
troppo costosi: chiederà dei \textbf{compromessi}. E la persona che, all'interno
del publisher, fa questo lavoro è il \textbf{developer}.

Sta al designer capire \textbf{su che cosa del proprio gioco si può trovare un
compromesso e su che cosa no}. Publisher diversi lavorano in modo diverso, e il
designer ha sempre l'autorità di dire ``no, non voglio pubblicare questo gioco
se queste sono le modifiche''. Ma deve sapere chi fa il lavoro di development,
perché se il suo obiettivo è la \textbf{pubblicazione tradizionale} si
interfaccerà con questa figura.

\subsection{Daniele: una distinzione che manca al videogioco}

Daniele sottolinea che questa distinzione è \textbf{tipica del mondo analogico}
ed è ``molto interessante''. In Italia la parola \emph{developer} è diventata
negli anni quasi un sinonimo di programmatore di videogiochi, o di chiunque
lavori nel videoludico. Secondo lui c'entra anche il fatto che qualche anno fa
il sito \textbf{Gamasutra} ha cambiato nome in \textbf{Game Developer}, e la
parola è diventata un ombrello per tutte le professioni del videogioco. ``Mi fa
un po' ribrezzo'': il game developer dovrebbe essere un ruolo di design, un
designer che, come nell'analogico, fa \textbf{mediazione}. Nel videogioco i
publisher hanno in effetti i loro developer interni, che adattano il gioco, per
esempio alle console e ai loro kit di sviluppo. Ma la distinzione dovrebbe
esistere anche nel videogioco, perché è ``difficilissimo'' avere la mentalità
per fare le due cose insieme: chi fa design ha un focus che raramente si può
sovrapporre a quello di chi fa mediazione.

A Daniele piace la parola \emph{develop}, \textbf{sviluppare}, nel senso di
prendere qualcosa che è stato progettato e \textbf{renderlo possibile nel mondo
reale}. Ci sono progetti che hanno bisogno di uno sviluppo che ``polishi''
l'esperienza: la semplifichi, la migliori, la renda comprensibile e la concluda,
ma soprattutto la renda \textbf{stampabile} e \textbf{fruibile} per l'utente. La
metafora crossmediale è quella della \textbf{fotografia}: chi sviluppa le
fotografie non è chi scatta la foto, ma chi la rende visibile, anche
fisicamente. ``È un po' retro come metafora, ma ci capiamo''. ``È vero, però è
giusto'', conferma Ruggero.

% ---------------------------------------------------------------------
\section{Il mercato: Italia ed estero}
\label{a1:sec:mercato}

Il corso è in italiano e si svolge in Italia, ma bisogna conoscere anche il
\textbf{mercato estero}: molti giochi progettati da autori italiani vengono
pubblicati all'estero, e conoscere il mercato estero aiuta a capire come il
proprio gioco può arrivarci, ``o viceversa, come non può arrivarci''. La
tabella~\ref{a1:tab:mercato} riassume il confronto della slide.

\begin{table}[htbp]
\centering
\small
\begin{tabular}{@{}p{6.4cm}p{6.4cm}@{}}
\toprule
\textbf{Italia} & \textbf{Estero (soprattutto Europa)} \\
\midrule
il gioco è ancora visto come un'attività \textbf{infantile}, anche se la
normalizzazione è in crescita &
forte tradizione di gioco anche in un contesto \textbf{familiare adulto},
soprattutto ``dalla Germania in su'' \\
risposta \textbf{conservativa} sul tema ludico &
\textbf{premi di prestigio} riconosciuti in tutto il mondo (Spiel des Jahres,
As d'Or) \\
forte focus sulla \textbf{localizzazione} di giochi stranieri &
si cercano soprattutto giochi \textbf{originali} \\
investimento \textbf{molto basso} sulla creazione di giochi originali interni;
pochissimi designer assunti &
\textbf{developer interni} pagati con un contratto (Francia, Germania,
Danimarca\ldots) \\
\bottomrule
\end{tabular}
\caption{Il mercato del gioco analogico in Italia e all'estero, secondo la
slide di Ruggero e i commenti di Daniele.}
\label{a1:tab:mercato}
\end{table}

\subsection{L'Italia}

\begin{itemize}
  \item In Italia il gioco è ancora visto come un'\textbf{attività infantile},
        anche se la \textbf{normalizzazione} è in crescita. C'è ancora una forte
        risposta \textbf{conservativa} sul tema ludico.
  \item C'è un forte focus sulla \textbf{localizzazione} dei giochi stranieri: in
        una qualunque ludoteca o negozio, molti dei giochi da tavolo e di ruolo
        sono di publisher o designer esteri, pubblicati in italiano da publisher
        italiani. (La localizzazione, cioè portare un titolo dall'estero in
        Italia, esiste come funzione all'interno del publisher; Ruggero non l'ha
        messa nella slide delle figure perché non è il focus del corso.)
  \item C'è, ``purtroppo, preparatevi psicologicamente'', un \textbf{investimento
        molto basso} sulla creazione di giochi originali interni alla casa
        editrice. Pochissime case editrici italiane hanno designer interni pagati
        a tempo pieno con contratto a tempo indeterminato. Daniele lo dice senza
        mezzi termini: ``il mercato è molto piccolo, conosciamo praticamente
        tutto il settore'', e le persone che in Italia fanno game design
        analogico a tempo pieno con un contratto a tempo indeterminato ``sono
        meno delle dita delle mie due mani''. Di solito la formula è quella del
        \textbf{freelance}, del contratto, delle \textbf{royalty}. Chi va da un
        publisher aspettandosi di essere assunto a tempo pieno come game
        designer deve sapere che ``non siamo purtroppo ancora in quel tipo di
        situazione''.
  \item C'è però un contraltare: in Italia, soprattutto nel gioco di ruolo ma
        anche nel gioco da tavolo, c'è stata molta \textbf{imprenditoria}, anche
        grazie a \textbf{Kickstarter}, che ha aperto le porte ad altri sistemi
        di \textbf{crowdfunding}. È un altro argomento, perché richiede un team,
        un investimento e così via.
  \item E gli autori non mancano, anzi: le occasioni di incontro tra autori
        vedono ogni anno un aumento dei possibili autori\footnote{La trascrizione
        dice ``i DAG a Parma''; il riferimento è probabilmente a IDEAG, il
        circuito di incontri tra autori di giochi italiani, ma non è
        ricostruibile con certezza.}. Bisogna ``sgomitare'' tra tanti autori,
        che magari arrivano da esperienze molto diverse.
\end{itemize}
Ruggero precisa: l'investimento basso è ``dall'alto''. Le persone che fanno e
vogliono fare giochi sono tantissime e in aumento; scarseggia l'investimento,
sia dei publisher disposti a pagare stipendi sia dei \textbf{bandi} pubblici. I
bandi esistono (Beatrice ne parlerà per i giochi pervasivi, che spesso si
muovono in quella direzione, e Ruggero cita qualcuno che sta facendo un card
game con un bando italiano, ``una cosa assurda, ma lo sta facendo''), ma sono
una rarità, soprattutto rispetto all'estero.

\subsection{L'estero}

\begin{itemize}
  \item L'estero (dipende dalla nazione, ma in generale il resto d'Europa) ha una
        forte \textbf{tradizione di gioco} anche in un contesto \textbf{familiare
        adulto}, soprattutto ``dalla Germania in su''.
  \item Ha \textbf{premi di prestigio} riconosciuti in tutto il mondo, come lo
        \textbf{Spiel des Jahres} tedesco e l'\textbf{As d'Or} francese\footnote{La
        trascrizione dice ``lo Spield AR, la SDOR''.}.
  \item Cerca soprattutto \textbf{giochi originali}. Lo si vede dal fatto che
        molti autori italiani di giochi da tavolo pubblicano in Germania o in
        Francia e \textbf{non vengono pubblicati in Italia}, ``che è tutto
        dire''.
  \item Ha \textbf{figure di development interne}: non è raro vedere publisher
        francesi, tedeschi, danesi con developer pagati internamente, a
        contratto.
\end{itemize}
``Poi la situazione è ancora più drammatica'', ammette Ruggero, ``ma per ora
vogliamo tenere il morale un po' alto. È l'introduzione. Ci sarà tempo di andare
nella tristezza''.

Prima di proseguire, Ruggero insiste sulle domande: con lui si può interrompere
anche mentre parla, alzando la mano con le reaction o semplicemente togliendo il
microfono e chiedendo. ``Non mi offenderò mai, anzi, sarò felice di vedervelo
fare e mi ricorderò positivamente di voi''.

% ---------------------------------------------------------------------
\section{Perché serve una terminologia comune}
\label{a1:sec:terminologia}

La serie di slide successiva è ``una piccolissima parte di fondamenti di
design'' (``non spaventatevi, vi vedo che vi spaventate un po'\,''). Anche qui
molti temi verranno ripresi in dettaglio, ma l'idea è arrivare alla prossima
lezione con almeno un \textbf{glossario}, un \textbf{linguaggio condiviso}: se
qualcuno dice \emph{game developer}, tutti sanno di che cosa sta parlando.

Ruggero parte da una domanda, e precisa che non è retorica: \textbf{perché ci
serve una terminologia comune?} Gli studenti rispondono:
\begin{itemize}
  \item Gianmarco, ``molto pragmatico'': perché bisogna lavorare in team e
        capirsi. Vero, nel corso e anche fuori;
  \item un'altra studentessa: per la \textbf{comunicazione}, tra di noi e con
        tutti quelli che sono interessati nel settore. Vero anche questo:
        conoscere i termini comuni permette di parlare con chi fa game design o
        development nel dettaglio, nel tecnico, nello specifico.
\end{itemize}
Ruggero aggiunge altre due ragioni.

\paragraph{Saper comunicare il proprio gioco.} La prima riguarda la
\textbf{comunicazione del gioco} verso l'esterno. ``E lo dico non tirando acqua
al mio mulino, ma anzi togliendola'': \textbf{non è spesso il buon game design
a far diffondere un gioco}. ``Vorrei che lo fosse, sarebbe bellissimo, ma
purtroppo non è così''. Ci sono giochi bellissimi che non sanno comunicare, che
non sanno dire che cosa fanno, e che quindi è difficile, da fuori, trovare
invitanti; e ci sono giochi mediocri dal punto di vista del design che si
comunicano molto bene, e che quindi conoscono in molti. Saper parlare di quello
che si fa è utile \textbf{a sé stessi}, in primis.

\paragraph{Documentare per sé stessi.} La seconda è la \textbf{documentazione}:
documenti di design, di brainstorming e così via. ``Sembrerà una battuta, ma
provate a recuperare un documento di brainstorming una settimana dopo, se non
avete scritto bene le cose che avete scritto''. Chi l'ha fatto sa già di che
cosa si parla; ``chi non l'ha ancora vissuto, forse lo vivrà''.

% ---------------------------------------------------------------------
\section{Game e play}
\label{a1:sec:game-play}

La \textbf{prima, importantissima divisione terminologica} è quella tra
\emph{game} e \emph{play}. In italiano c'è un problema: \textbf{gioco} e
\textbf{giocare} hanno la stessa radice, mentre in altre lingue, tra cui
l'inglese, sono due parole strutturalmente diverse:
\begin{description}
  \item[Game] il \textbf{gioco come artefatto}: il gioco come ecosistema di
        meccaniche, come struttura di regole, come prodotto (``tra
        virgolette''), come oggetto fisico.
  \item[Play] l'\textbf{atto}, il verbo: le persone che giocano, l'approccio
        fisico all'attività del gioco.
\end{description}
È importante perché gli studenti faranno \textbf{game design}; esistono anche
ambiti di \emph{play design}, ma qui ``creerete giochi, non giocattoli''.

\subsection{Giochi e giocattoli}

Qual è la differenza? Uno studente risponde: \textbf{la presenza di regole}.
Esatto: un gioco ha delle regole, un'attività di \emph{play} non
necessariamente. Con un giocattolo si può fare \emph{play}: c'è chi gioca con i
Lego, chi gioca a guardia e ladri o a nascondino al parco, e magari alcune di
quelle regole sono ``più o meno semi-solidificate''. Ma il \emph{play} permette
spesso un approccio non strettamente legato alle regole.

Viceversa, \textbf{un game può esistere anche se nessuno lo sta giocando}. Se si
prepara su un tavolo una scacchiera con i pezzi degli scacchi, chi passa e la
vede riconosce il gioco degli scacchi, anche se nessuno sta giocando. Se invece
si vedono dei bambini correre in un parco giochi, quei bambini \emph{stanno
giocando}, ma non stanno giocando \emph{a un gioco} specifico. ``È importante
separare questa cosa''; ci si tornerà, ``ma per ora abituatevi a questa
differenza''.

\paragraph{Approfondimento: giocattoli, giochi e regole.} La stessa distinzione
è il punto di partenza del capitolo che Jesse Schell, in \emph{The Art of Game
Design}, dedica alla definizione di gioco (\emph{The Experience Rises Out of a
Game}). Schell nota che usiamo perfino un linguaggio diverso: ``un gioco è
qualcosa a cui si gioca'' (\emph{a game is something you play}), ``un giocattolo
è qualcosa con cui si gioca'' (\emph{a toy is something you play with}); e
arriva a definire il giocattolo come ``un oggetto con cui si gioca''. Analizzando
poi la definizione di gioco di Avedon e Sutton-Smith, ne ricava che ``i giochi
hanno regole, i giocattoli no'', e che le regole sono uno degli aspetti
\emph{definitori} del gioco, insieme alla volontarietà, agli obiettivi, al
conflitto e alla possibilità di vincere o perdere. Schell osserva anche che
molti giochi sono costruiti \textbf{sopra} un giocattolo (``una palla è un
giocattolo, ma il baseball è un gioco''), e tra i suoi consigli di
prototipazione c'è proprio \emph{Build the Toy First}: costruire prima il
giocattolo, cioè qualcosa che sia già piacevole da manipolare, e poi il gioco
intorno. È una lettura utile della distinzione di Ruggero: il \emph{play} può
fare a meno delle regole, il \emph{game} no.

% ---------------------------------------------------------------------
\section{Il framework MDA (e il DPE)}
\label{a1:sec:mda}

\subsection{Meccaniche, dinamiche, estetiche}

Un'altra terminologia fondamentale è quella del framework \textbf{MDA}
(\emph{Mechanics, Dynamics, Aesthetics}; in italiano a volte ``MDE''). Viene da
un \textbf{paper} nato nell'ambito del videogioco\footnote{È il paper di Robin
Hunicke, Marc LeBlanc e Robert Zubek, \emph{MDA: A Formal Approach to Game
Design and Game Research}, del 2004. Daniele aggiunge il paper alla cartella
condivisa durante la lezione.}, che ha ``solidificato'' queste terminologie, e
divide un gioco in tre aree:
\begin{description}
  \item[Meccaniche] le \textbf{regole per come sono scritte}. ``Ogni volta che
        pescate un 2 fate dei punti'': quella è una meccanica.
  \item[Dinamiche] l'\textbf{applicazione in tempo reale} delle meccaniche. La
        dinamica è il giocatore che, pescando una carta, \textbf{spera} che esca
        un 2. Quel processo non è una regola: ``il gioco non vi dice che dovete
        sperare che esca un 2''.
  \item[Estetiche] l'insieme di quello che il gioco vuole comunicare,
        l'\textbf{esperienziale}.
\end{description}

L'esempio di Ruggero mette insieme le tre aree. In un gioco in cui pescare un 2
fa fare punti (meccanica) e in cui quindi si spera che esca un 2 (dinamica),
\textbf{che estetica si ottiene se nel mazzo di 100 carte c'è un solo 2? E se ce
ne sono 50?} L'emozione di aver pescato \emph{l'unico} 2 del mazzo è
completamente diversa da quella di pescarne uno sapendo che ce ne sono
cinquanta e che è molto probabile trovarlo. Le regole sono quasi identiche; cambia
l'esperienza.

\subsection{Il falso amico: ``estetica''}

Daniele avverte di non sottovalutare la parola \textbf{estetica}, che in italiano
è un \textbf{falso amico}: sembra indicare la parte grafica, e invece non lo è.
Si incontreranno designer che hanno studiato MDA e altri che non l'hanno
studiato, e lo si capirà proprio da questo: i secondi diranno che le
\emph{aesthetics} sono la grafica. Le \emph{aesthetics} sono la parte
\textbf{esperienziale}: come nell'esempio di Ruggero, quello che lascia al
giocatore aver pescato l'unico 2 del mazzo.

\subsection{Il DPE: design, play, experience}

Per gli italiani, aggiunge Daniele, è più facile capire l'MDA attraverso un
altro paper, quello del \textbf{DPE}\footnote{È il framework \emph{Design, Play,
Experience} di Brian Winn, proposto come estensione dell'MDA per i giochi
educativi (\emph{serious games}). Anche questo paper viene aggiunto da Daniele
alla cartella condivisa.}, che di fatto fa un parallelismo assegnando
\textbf{design, play, experience}:
\begin{itemize}
  \item le \textbf{meccaniche} sono il lato \textbf{design};
  \item le \textbf{dinamiche} sono il lato \textbf{play};
  \item le \textbf{estetiche} sono l'\textbf{experience}, l'esperienza che riceve
        chi gioca.
\end{itemize}
``Se vivete le \emph{aesthetics} come l'esperienza che vive chi sta giocando,
state a postissimo''. In questo momento Beatrice lascia la lezione (``ciao,
Bea''). I paper di MDA e DPE sono in inglese; chi vuole può dargli una lettura.

\begin{figure}[htbp]
\centering
\begin{tikzpicture}[
    font=\small,
    blocco/.style={draw, thick, rounded corners=3pt, fill=gray!12,
                   minimum width=3.2cm, minimum height=1.0cm, align=center},
    dpe/.style={font=\footnotesize\itshape, align=center},
    esempio/.style={draw, rounded corners=2pt, text width=3.3cm,
                    align=left, font=\footnotesize, inner sep=4pt},
    freccia/.style={-{Stealth[length=2.5mm]}, thick}
  ]
  \node[blocco] (m) at (0,0)   {\textbf{Meccaniche}\\ \footnotesize le regole scritte};
  \node[blocco] (d) at (4.3,0) {\textbf{Dinamiche}\\ \footnotesize le regole in azione};
  \node[blocco] (a) at (8.6,0) {\textbf{Estetiche}\\ \footnotesize l'esperienza};
  \draw[freccia] (m) -- (d);
  \draw[freccia] (d) -- (a);
  \node[dpe, above=0.15cm of m] {DPE: \textbf{design}};
  \node[dpe, above=0.15cm of d] {DPE: \textbf{play}};
  \node[dpe, above=0.15cm of a] {DPE: \textbf{experience}};
  \node[esempio, below=0.45cm of m] (em)
    {``Ogni volta che peschi un 2 fai punti''};
  \node[esempio, below=0.45cm of d] (ed)
    {pescando, il giocatore \emph{spera} che esca un 2 (nessuna regola glielo
     chiede)};
  \node[esempio, below=0.45cm of a] (ea)
    {un solo 2 su 100 carte: pescarlo è un evento;\\
     50 su 100: è quasi scontato};
  \draw[dashed] (m) -- (em);
  \draw[dashed] (d) -- (ed);
  \draw[dashed] (a) -- (ea);
  % prospettive
  \draw[freccia, gray] ([yshift=-0.3cm]em.south west) --
    node[below, font=\footnotesize\itshape, text=black]
    {il designer agisce sulle meccaniche}
    ([yshift=-0.3cm]ea.south east);
  \draw[freccia, gray] ([yshift=-1.2cm]ea.south east) --
    node[below, font=\footnotesize\itshape, text=black]
    {chi gioca incontra prima l'esperienza}
    ([yshift=-1.2cm]em.south west);
\end{tikzpicture}
\caption{Il framework MDA con l'esempio del ``2'' usato da Ruggero, e la
corrispondenza con il DPE suggerita da Daniele. Le due frecce in basso
riassumono la lettura classica del paper: il designer costruisce le meccaniche
per ottenere un'esperienza, mentre il giocatore parte dall'esperienza.}
\label{a1:fig:mda}
\end{figure}

\paragraph{Approfondimento: le due direzioni dell'MDA e le otto estetiche.} Le
frecce della figura~\ref{a1:fig:mda} riprendono un'idea centrale del paper
originale, che la lezione non esplicita: designer e giocatore guardano il gioco
da lati opposti. In \emph{Challenges for Game Designers}, Brenda Brathwaite e
Ian Schreiber la riassumono così: i designer si chiedono quale estetica
vogliono ottenere, definiscono le dinamiche che porterebbero a quella
sensazione, e poi creano le meccaniche che producono quelle dinamiche. Lo stesso
libro conferma l'avvertenza di Daniele sul falso amico: ``le \emph{aesthetics}
non si riferiscono all'aspetto del gioco, ma alla risposta emotiva'' che il
designer vuole suscitare. E offre esempi di dinamiche nel senso di Ruggero: in
uno sparatutto i punti di \emph{spawn} sono una meccanica, mentre appostarsi
accanto a uno di essi per colpire chi riappare (lo \emph{spawn camping}) è una
dinamica; negli scacchi i movimenti dei pezzi sono meccaniche, mentre le
aperture codificate sono dinamiche. Per dare un nome alle estetiche, il paper
MDA propone una tassonomia, poi detta degli ``otto tipi di \emph{fun}'', che
Schell riprende in \emph{The Art of Game Design} come \emph{LeBlanc's Taxonomy of
Game Pleasures}: \textbf{sensazione}, \textbf{fantasia}, \textbf{narrativa},
\textbf{sfida}, \textbf{compagnia} (\emph{fellowship}), \textbf{scoperta},
\textbf{espressione} e \textbf{sottomissione} (il piacere di entrare nel
``cerchio magico'' e lasciarsi andare a un passatempo). È un primo esempio di
quello che i docenti chiedono con le parole vietate della prossima sezione:
invece di dire che un gioco è ``divertente'', dire \emph{quale} esperienza
offre.

% ---------------------------------------------------------------------
\section{Le parole vietate}
\label{a1:sec:parole-vietate}

``Attenzione, adesso qua triggeriamo delle persone''. Ruggero annuncia una regola
che ``tornerà'': nel corso ci sono delle \textbf{parole vietate}. ``Non perché
siamo dei censori''; il perché si vedrà in una lezione apposita, perché
spiegarlo adesso richiederebbe un'ora. Ma Daniele vuole spiegarne ``un
minimo''.

\subsection{La regola}

In \textbf{nessun materiale scritto} consegnato ai docenti, task compresi, sono
accettate queste parole, né in inglese né in italiano:
\begin{center}
\begin{tabular}{@{}ll@{}}
\toprule
\textbf{Inglese} & \textbf{Italiano} \\
\midrule
\emph{fun} & divertimento \\
\emph{immersion} & immersione \\
\emph{entertainment} & intrattenimento \\
\bottomrule
\end{tabular}
\end{center}
``Non li voglio neanche vedere''. Daniele lo ripete con un paragone: ``se mi fate
i disegni con Paint, benissimo, top, fantastico, siamo ancora best friend. Se mi
scrivete \emph{divertimento}, io sbrocco''. È una \textbf{regola dell'accademia}:
negli stessi corsi triennali di narrative design e game design le parole sono
vietate agli studenti per tre anni; ``per voi solo per tre mesi''. Daniele
aggiungerebbe volentieri un quarto termine, \textbf{gamification}.

\subsection{Perché}

Il motivo approfondito lo darà Ruggero nella lezione apposita. Daniele dà il
suo. Ruggero e Diego Ricchiuti sono molto contrari a \emph{fun} e a
\emph{immersion}; lui ha aggiunto \emph{intrattenimento}. Sono \textbf{parole
troppo passe-partout e troppo ambigue}:
\begin{itemize}
  \item \textbf{intrattenimento} in Italia è stata usata per anni come un modo
        per \textbf{sminuire}: ``ma sì, tanto è un gioco solo di
        intrattenimento, è un romanzo solo di intrattenimento''. Ma anche se un
        prodotto è ``solo di intrattenimento'', non vuol dire che non ci si debba
        lavorare con attenzione. E poi intrattenimento ``non significa niente'':
        ci sono tanti modi in cui una cosa può intrattenere;
  \item lo stesso vale per \textbf{fun}, e per \textbf{immersion}, che ha un
        problema in più: è una parola passe-partout per cui si può dire che un
        gioco è ``immersivo'' nel senso generale del termine, ma poi non si sa
        davvero che cosa voglia dire, \emph{immersivo in che senso}.
\end{itemize}

L'obiettivo didattico è che gli studenti \textbf{vadano oltre queste parole}:
``un po' come quando alle medie vi dicono smettetela di dire che tutto è bello,
perché ci sono vari gradi di bello, ci sono altri modi di dire bello''. Bisogna
\textbf{decostruirle}: quando si sta per usarle, chiedersi \textbf{che cosa si
intende} con quella parola e cercarne una \textbf{più precisa}.

Ruggero chiede, in uno dei ``rari casi'', di \textbf{non fare domande adesso} su
questa regola: ogni anno arrivano domande e commenti, e ci saranno lezioni
apposite di tre ore, con strumenti per capire il perché. Ma da subito gli
studenti devono abituarsi a evitare questi termini, sia perché si spera che
abbiano termini più precisi, sia perché queste parole ``portano con sé dei
problemi agganciati'' che si vedranno. ``Non spaventatevi, non verrete sparati in
un vicolo cieco di notte''.

\paragraph{Una piccola gamification.} Il divieto va visto come uno
\textbf{stimolo}, ``un po' come se stessimo gamificando la vostra esperienza'',
una sorta di user experience delle lezioni. Imponendosi di non usare queste
parole, si noterà come migliora la capacità di \textbf{descrivere i giochi} e le
proprie esperienze ludiche, e di essere più precisi. È un invito a ``obbligarsi a
essere più precisi e ad avere una mentalità più critica''. La lezione di
Ruggero sull'\textbf{involvement}, promette Daniele, sarà ``fenomenale'' e
``vi aprirà la testa''.

\paragraph{Approfondimento: anche i libri faticano a definire il ``fun''.} La
difficoltà che il divieto vuole far emergere è visibile anche nei testi di game
design. Schell, in \emph{The Art of Game Design}, prova a definire il \emph{fun}
e propone ``piacere con sorprese'' (\emph{pleasure with surprises}), ammettendo
subito che è strano poter usare una parola per tutta la vita, sapendo con
certezza che cosa significa, senza riuscire a esprimerla chiaramente quando
qualcuno lo chiede; e invita a mettere alla prova la definizione con
controesempi. La sua lente sul \emph{fun} (\emph{The Lens of Fun}) si limita,
non a caso, a chiedere quali parti del gioco siano \emph{fun} e \emph{perché}.
Anche \emph{immersion} ha molte letture: in \emph{21st Century Game Design},
Bateman e Boon la descrivono, tra l'altro, come la capacità di ``perdersi''
nel mondo artificiale presentato dal gioco, legata alla sospensione
dell'incredulità e ai giochi di mimesi. Sono significati diversi sotto la stessa
parola: proprio il motivo per cui, in un documento di design, conviene dire
quale si intende.

% ---------------------------------------------------------------------
\section{Chi crea, chi gioca, chi facilita}
\label{a1:sec:ruoli-esperienza}

Un'altra distinzione importante è quella tra \textbf{chi crea}, \textbf{chi
gioca} e \textbf{chi facilita}. Spesso, quando si fa un gioco, lo si fa provare
di persona a chi gioca. Ma ``la norma dei giochi pubblicati non funziona così'':
quando si pubblica un gioco non si vede chi lo compra e chi lo gioca, e non si
sa se qualcuno, a casa propria con i propri amici, segue le regole alla lettera e
interpreta correttamente quello che il designer aveva descritto. Quindi già in
fase di progettazione bisogna tenere a mente \textbf{chi sta facendo game design
e chi sta giocando}.

E nei giochi a evento, nei giochi pervasivi, nei giochi di ruolo da tavolo e dal
vivo (LARP) c'è un terzo pezzo del puzzle, ``importantissimo'': \textbf{chi
facilita}. La domanda da farsi è: \textbf{il regolamento è scritto per chi gioca
o per chi deve facilitare il gioco?} Bisogna abituarsi a pensare a tutte le
persone coinvolte nell'esperienza.

% ---------------------------------------------------------------------
\section{Playtest e iterazione}
\label{a1:sec:playtest}

\subsection{Nessun gioco nasce perfetto}

L'ultima voce del glossario è ``un'altra cosa importantissima'':
\textbf{playtest} e \textbf{iterazione}. Il corso è \textbf{pragmatico}:
``specialmente da parte mia'', dice Ruggero, non si vedranno teorie ``messe lì
per fare accademia sul game design nelle nostre torri d'avorio''. Ci saranno
teorie e framework, ma tutti estremamente pragmatici, perché l'obiettivo non è
solo capire i giochi, ma \textbf{saperli fare e iterare}.

Playtest e iterazione sono \textbf{inevitabili} nel game design, perché
\textbf{nessun gioco nasce perfetto}, soprattutto perché le persone sono
diverse: un gioco perfetto per una persona può essere incomprensibile per
un'altra. Chi progetta giochi per altre persone deve capire, dopo averli fatti
provare, \textbf{che cosa cambiare, che cosa sistemare, che cosa aggiungere, che
cosa togliere}. Ci saranno lezioni su playtest, prototipazione, raccolta dei
feedback, metodi etnografici. Ma bisogna abituarsi già ora all'idea che
\textbf{non esiste la consegna definitiva}. A un certo punto i task vanno
consegnati per forza di cose, ma, come diceva Daniele all'inizio, alla fine del
corso gli studenti avranno un \textbf{prototipo}, ``che sarà un prototipo'', e
gli strumenti per portare avanti il lavoro anche dopo.

\subsection{L'aiuto dopo il corso}

Non solo gli strumenti: anche l'aiuto dei docenti. Daniele racconta di aver
portato a un editore, a Play\footnote{Probabilmente \emph{Play -- Festival del
Gioco} di Modena; la trascrizione non permette di esserne certi.}, uno dei
progetti realizzati l'anno precedente da due studenti, un gioco di ruolo in
solitario (``un \emph{solo game} RPG molto, molto bello''). Se gli studenti
vorranno trasformare il proprio lavoro in un prototipo vero e proprio, possono
contattarlo: lo fa ``per lavoro, ma anche per passione''. Gestisce anche una
community per aiutare i designer ``un po' depressi'' a rimettersi a progettare.

Quindi: ``non sentitevi persi se durante la fase di iterazione vi sentite
persi, è normale''.

\subsection{Il dizionario comune e il production dictionary}

Daniele riprende un punto detto ``benissimo'' da Ruggero e che secondo lui non è
stato sottolineato abbastanza: il \textbf{dizionario comune} serve anche a
\textbf{non perdersi} mentre si itera. Se da un'iterazione all'altra i
documenti usano termini completamente diversi, e non c'è un dizionario di base,
ci si perde.

Nel videogioco, una delle primissime cose che un designer fa quando entra in
un'azienda è il \textbf{production dictionary}: si annota, in un file privato,
tutte le parole che usano i colleghi, per allinearsi al dizionario comune
dell'azienda, perché lì su un prototipo possono lavorare anche cento persone.
Nel gioco da tavolo non si fa molto, perché il mondo dell'analogico è
``talmente più piccolo''; ma serve comunque imparare un \textbf{dizionario
minimo}, dai docenti e dai professionisti. E il corso ne darà una parte,
entrando nel dettaglio di alcuni termini e framework.

\paragraph{L'esempio del feedback.} Perché serva, lo mostra un esempio. Si
testa un gioco e ci si annota come feedback: ``questo gioco deve essere più
divertente''. Che cosa vuol dire? Come lo si migliora? Quella frase non dà
elementi utili. Se invece si scrive che deve essere \textbf{più adrenalinico},
\textbf{più frenetico}, \textbf{più rilassato}, \textbf{più cozy}, o
\textbf{più bilanciato matematicamente}, si sa dove andare a lavorare.

\paragraph{Approfondimento: la regola del ciclo.} Il principio ``nessun gioco
nasce perfetto'' ha un nome preciso in \emph{The Art of Game Design}: è la
\textbf{Rule of the Loop}, la regola del ciclo. ``Più volte provi e migliori il
tuo design, migliore sarà il tuo gioco''. Schell precisa che non è una lente,
cioè un punto di vista, ma ``una verità assoluta'', senza eccezioni: ci saranno
momenti in cui si cercherà di convincersi che ``questa volta il design è così
buono che non serve testarlo'', e ogni volta se ne pagherà il prezzo. Ed è
interessante la ragione per cui la regola pesa di più sul digitale: secondo
Schell il tempo e il denaro necessari per provare e correggere un sistema sono
molto maggiori nei videogiochi che nei giochi tradizionali. Per un gioco di
carte o da tavolo, con il tempo per provarlo e cambiarlo più e più volte, il
consiglio è semplicemente di farlo. Il gioco analogico è, da questo punto di
vista, la palestra ideale dell'iterazione.

% ---------------------------------------------------------------------
\section{Tabletop Simulator}
\label{a1:sec:tts}

\subsection{Il software di settore}

``Last but not least'': il corso usa un software, \textbf{Tabletop Simulator},
che si compra su \textbf{Steam} e costa tra i \textbf{10 e i 20 euro} (``aspettate
che vada in sconto, va in sconto ogni due per tre''). Daniele spiega perché:
\begin{itemize}
  \item è \textbf{il software di settore} del gioco da tavolo, quello che i
        publisher usano per i prodotti digitali, interni ed esterni;
  \item essendo su Steam, permette di fare \textbf{playtest} e di giocare con gli
        amici di Steam. Daniele stesso sta testando lì il proprio gioco;
  \item funziona anche su Mac: Windows o Mac, nessun problema.
\end{itemize}
Ci sarà una lezione apposita, di Ruggero, sulla creazione del proprio
\textbf{modulo} di gioco.

\paragraph{Che cosa non è.} Daniele chiarisce subito un ``fraintendimento
gigantesco'', soprattutto per chi ha già esperienza: Tabletop Simulator
\textbf{non richiede programmazione} e \textbf{non è un software per
trasformare i giochi da tavolo in videogiochi}. È un software che
\textbf{simula dei componenti fisici su un tavolo}: serve per giocare i board
game in digitale. Simula anche la fisica: c'è perfino il comando per
\textbf{ribaltare il tavolo} (``non usatelo ora'').

\paragraph{La partnership mancata.} Da anni Daniele prova a fare una
partnership con gli sviluppatori: due anni prima ha contattato l'azienda per
acquistare delle chiavi per gli studenti, dato che il software si usa nei corsi.
L'azienda ha risposto che non fa accordi commerciali. ``Credo sia la prima
azienda nella mia vita che dice che non vuole che io le dia dei soldi''. Gli
studenti dovranno quindi provare ad acquistarlo da sé; se i docenti riusciranno a
darglielo, bene.

\paragraph{La domanda di Gianmarco.} Gianmarco segue le lezioni da un computer
aziendale ``super chiuso'', su cui non può installare Steam: deve usare Tabletop
Simulator durante le lezioni? Ruggero risponde che per la lezione dedicata è
meglio poter seguire e rifare in diretta. Ma il software non sarà richiesto
necessariamente per il project work: un gruppo potrebbe fare un gioco pervasivo
o un gioco di ruolo, non un gioco da tavolo. Essendo un software di settore, è
comunque bene conoscerlo. Ruggero aggiunge che le versioni non ufficiali che
circolano non permettono di giocare in molti, e quindi non servono per
playtestare o giocare insieme.

\subsection{Una dimostrazione}

Daniele mostra il software aperto da Steam:
\begin{itemize}
  \item si può creare una partita \textbf{single player} o \textbf{multiplayer};
        per giocare tra compagni si crea un server multiplayer e ci si aggiunge
        come amici su Steam;
  \item si dà al server un nome, un numero di giocatori, eventualmente una
        password; creato il server, si apre la scelta dei giochi;
  \item ci sono giochi gratuiti, giochi da acquistare (DLC, non obbligatori) e
        soprattutto un \textbf{Workshop} ``ampissimo'', con decine di migliaia di
        giochi caricati dagli utenti: ``se c'è un board game che conoscete,
        quel board game è su Tabletop Simulator, giocabile dal Workshop''.
        L'esempio mostrato è l'open beta ufficiale del gioco da tavolo di
        \emph{Slay the Spire}\footnote{La trascrizione dice ``Slade Spire''.}:
        molti publisher, per farsi pubblicità, mettono il proprio gioco in demo o
        giocabile in digitale;
  \item c'è un \textbf{tutorial} iniziale che spiega come muoversi sul tavolo;
        con \emph{Esc} e \emph{Help} si trovano i comandi da tastiera.
\end{itemize}

\subsection{Il primo task}

Il primo (e per ora unico) task è questo: \textbf{entro il 10 ottobre}, giorno
della lezione di Ruggero su Tabletop Simulator, giocare \textbf{tre o quattro
volte in totale}, online, tra compagni di corso o con amici. Non si tratta di
imparare a creare moduli, ma di prendere la mano \textbf{lato giocatore}: aprire
il software, impararne i comandi, giocare. ``Per me potete anche andare a
prendere domino o mahjong'' (``mahjong è bellissimo tra l'altro'', aggiunge
Daniele, che sta facendo un videogioco sul mahjong), o il pachisi o il go. Gli
anni precedenti si chiedeva a ciascuno di proporre ai compagni un gioco da
tavolo che apprezzava. È anche un modo per conoscersi. Ruggero spiega il
perché: usarlo per giocare è facile e intuitivo, usarlo per creare moduli è più
complesso; ma chi arriva al 10 ottobre avendo già aperto il software e giocato,
troverà più intuitivo anche creare un modulo. E l'avvertimento: ``non riducetevi
all'ultimo, non fate una partita il 9 sera''.

\subsection{Tabletopia e altri vantaggi}

Tancredi racconta di aver provato a caricare un suo gioco su
\textbf{Tabletopia}. ``Questo è un vaso di Pandora'', risponde Daniele.
Tabletopia è molto usato, ma secondo lui il suo strumento di creazione dei moduli
è pessimo, ed è il motivo per cui la maggioranza dei moduli ufficiali dei
publisher è su Tabletop Simulator. Tabletopia ha un vantaggio vero, è
\textbf{gratuito}, ma ha scelte di ``qualità della vita'' poco felici.

Tabletop Simulator ha invece un vantaggio ``fenomenale'' che si tende a
sottovalutare: essendo su Steam, \textbf{il modulo si pubblica direttamente sul
Workshop di Steam}; si possono invitare gli amici a iscriversi e giocare con
loro, purché abbiano un account Steam. Rispetto ad altri software da imparare,
costa ``tra i 10 e i 20 euro, ragazzi: si parla proprio di spiccioli''. ``Io con
Tabletop Simulator ci ho playtestato la vita''.

Altre caratteristiche utili:
\begin{itemize}
  \item si possono creare \textbf{moduli non elencati} (\emph{unlisted}),
        perfetti da passare a un publisher prima che il gioco sia pubblicato:
        non sono pubblici, e nessuno può rubare i materiali;
  \item il software ha una grafica datata (``\emph{World Art 1998}'', scherza
        uno dei docenti), ma lo stanno portando su \textbf{Unity} con molti investimenti
        e aggiornamenti; ascolta la comunità e introduce funzioni utili anche per
        chi fa design;
  \item il suo server Discord è molto frequentato, e c'è un canale YouTube
        dedicato solo al design di giochi di carte\footnote{La trascrizione ne
        riporta il nome come ``Distraction Maker''; non è verificabile.} con un
        server Discord dove si fanno playtest: il 99\% è su Tabletop Simulator;
  \item molti editori, quando si propone loro un gioco da tavolo,
        \textbf{chiedono se c'è un modulo per Tabletop Simulator}.
\end{itemize}

% ---------------------------------------------------------------------
\section{Le domande a Daniele: il tema del project work}
\label{a1:sec:domande-daniele}

Daniele ha ``un minuto'' per rispondere alle domande sul corso prima di dover
andare. Uno studente chiede, ``giusto per curiosità'', se il project work avrà
un tema o sarà libero.

\begin{itemize}
  \item C'è un \textbf{mandate}, che i docenti propongono ogni anno e
        aggiornano leggermente; verrà presentato il \textbf{25 settembre}.
        ``Non facciamo spoiler''. Conterrà quello che si chiede di fare e quello
        che non si chiede: dei \textbf{must}, dei \textbf{must not} e dei
        \textbf{nice to have}.
  \item Verrà chiesto di lavorare su un'\textbf{idea originale} che non si è
        portata al corso: originale nel senso di \textbf{creata apposta}, non
        nel senso di ``mai vista prima''.
  \item L'idea va \textbf{partorita insieme}. Altri anni alcuni studenti hanno
        provato a imporre la propria idea al gruppo, ``e l'esperienza è stata
        pessima per tutti''. Quindi: \textbf{non imponete le vostre idee, fate
        brainstorming con gli altri}. Daniele, che fa ``molto spesso
        brainstorming'', darà strumenti per farlo: nessuno arriverà impreparato
        a quella lezione.
\end{itemize}

Ruggero aggiunge che il mandate è \textbf{abbastanza flessibile} perché tutti
possano fare qualcosa che li interessi: ``l'ultima cosa che vogliamo è forzarvi a
lavorare su qualcosa che non vi piace''. Bisogna lavorare su ciò che piace e per
cui si ha uno stimolo personale, anche nella \textbf{tipologia} di gioco: nessuno
sarà obbligato a fare un gioco di ruolo o un board game, e la lezione sui casi
studio mostrerà quante forme diverse di gioco analogico il corso accetta.

L'unica cosa da \textbf{non fare}, ``perché è contraddittorio per questo corso'',
sono i \textbf{giochi grandi}: ``lo chiediamo per voi''. Niente progetti
mastodontici\footnote{La trascrizione dice ``non fate delle cose da gnome'':
probabilmente il nome di un gioco da tavolo molto grande e complesso, non
ricostruibile con certezza.}, niente ``cose alla \emph{Dungeons \& Dragons}''.
Si può fare un gioco di ruolo, ma con un'esperienza ``meno corposa'': ci sono
tanti giochi di ruolo da 15 o 20 pagine, ``bellissimi, fichissimi'', che si
possono leggere online. Il consiglio è farsi una cultura anche su
\textbf{itch.io}, ``uno dei portali più importanti del gioco di ruolo
internazionale'', di cui si parlerà anche nella parte di ricerca e sviluppo.

Poi Daniele saluta (``noi ci vediamo venerdì a lezione'') e lascia la
conduzione a Ruggero.

% ---------------------------------------------------------------------
\section{Le domande finali con Ruggero}
\label{a1:sec:domande}

\subsection{Le regole delle domande}

Ruggero spiega che cosa vuol dire ``stare con lui''. Alla fine di ogni sua
lezione ci saranno \textbf{almeno 20--30 minuti di domande} libere, \emph{free
form}. Vuole che fare domande sia \textbf{normalizzato}: ``io gestisco malissimo
l'\emph{horror vacui} e divento super cringe molto alla svelta se non ci sono
domande in giro''. Sono benvenute le domande sulla lezione appena fatta, ma non
solo: il corso è crossmediale, e anche le domande che vengono da
\textbf{divagazioni}, da parti laterali, ``fatele''. Quello che sembra scontato o
banale a uno studente può non esserlo per il docente, e viceversa; e può essere
utile anche agli altri. ``Se vi sembrano stupide, fatele il doppio''. ``Magari
sembro incazzato, ma vi giuro sono rilassatissimo. Non mordo, magari abbaio un
po'\,''.

\subsection{Come approcciarsi a Steam e Tabletop Simulator}

Rita non ha mai usato né Steam né Tabletop Simulator e chiede come
approcciarsi, e se c'è modo di mettersi d'accordo tra compagni su Discord, magari
con un sondaggio sui giochi, per formare gruppi di gioco. Ruggero conferma: sul
Discord dell'accademia c'è un \textbf{canale apposito per il corso}, ed è degli
studenti. Lo possono usare per organizzarsi a giocare, per chiedere a un
compagno di prestare un gioco, per domande tecniche (come si scarica, come si
installa). ``Sarebbe proprio ideale che vi parlaste e vi organizzaste''.

Anche perché, aggiunge con un ``piccolo spoiler'' (``nessuno vi ha detto
niente''), se fossero i docenti a organizzare i gruppi di gioco, il task lo
starebbero facendo loro. Bisogna \textbf{abituarsi a organizzare le sessioni di
gioco}, perché arriveranno i task sul \textbf{playtest} e su come organizzarlo,
e a quel punto si sarà ``meno disperati''. Sembra scontato, ma saper
\textbf{organizzarsi per giocare in autonomia} ``vi sarà utilissimo come skill in
questo settore, ve lo giuro''.

\subsection{Il bilanciamento}

Angelo confessa ``una paranoia grossissima'': ci saranno lezioni sul
\textbf{bilanciamento}? ``Ma certo''. Sono proprio le lezioni di Marta sul
development: il bilanciamento è di solito un tema che tocca \textbf{i developer
più dei designer}. L'anno precedente c'era una sola lezione con Marta, e i
docenti si sono resi conto che serviva andare ancora più a fondo: per questo
quest'anno le lezioni sono due. ``Non è una paranoia, è una legittima domanda'':
Ruggero insegna da anni, vede altri corsi, parla con altri docenti e studenti, e
``mi incazzo spesso'' perché molti corsi di game design \textbf{non fanno
bilanciamento}. I docenti volevano invece un corso che coprisse tutte le parti
del game design: non solo l'idea del gioco, ma il bilanciamento, la produzione,
i playtest, l'iterazione, la scrittura dei regolamenti.

\subsection{Si può giocare con i docenti?}

Una studentessa chiede se si può giocare con i docenti. Per motivi etici,
durante il corso questo è limitato alle parti di \textbf{valutazione} dei giochi
e di \textbf{tutoraggio}. Dopo il corso, ``nessun problema, almeno con me''.

\subsection{Dove si trovano i giochi su Tabletop Simulator}

Ruggero mostra dove si trovano i giochi: nella propria libreria di Steam, nella
pagina di Tabletop Simulator, alla voce \textbf{Workshop}, dove ci sono ``più o
meno 36.000'' giochi. Una studentessa scopre intanto che l'app di Steam per
telefono è solo una \emph{companion app}: per registrarsi e giocare serve il
PC.

Per provare, chiede agli studenti un titolo:
\begin{itemize}
  \item \emph{Keep the Heroes Out}: non compare, ``non tutti i giochi ci sono,
        perché molti publisher non li mettono''; forse, essendoci già una
        versione del gioco su Steam, il modulo è stato tolto;
  \item \emph{Magic: The Gathering}: ``almeno 50 moduli'', anzi ``123 pagine'';
  \item \emph{Bang!}: tre moduli.
\end{itemize}
Bisogna fare attenzione: alcuni moduli sono in lingue specifiche (``magari il
primo è in russo''), alcuni hanno già l'espansione, altri sono solo il gioco
base. Conviene sfogliare e non fermarsi al primo risultato. E nel Workshop si
trovano anche cose molto particolari: giochi fuori produzione, giochi vecchi,
giochi ancora in prototipo.

\subsection{I ruoli nei team e le dinamiche tra le persone}

Uno studente torna sulla slide delle figure: nelle realtà piccole molti ruoli
hanno ``il doppio o triplo cappello''. Quali ruoli restano distinti, e come si
creano le sinergie tra persone con obiettivi diversi?

Ruggero premette che l'ideale ``platonico'' è un ambiente di lavoro costruttivo,
basato sulla fiducia e ``il più \emph{safe} possibile''; ma nell'industria non
sempre è così, e le strutture variano moltissimo: alcuni publisher sono molto
verticali, altri hanno ruoli già solidificati (se un publisher ha già i suoi
developer e non sta crescendo, entrarci come developer è praticamente
impossibile). ``Se prendiamo due aziende in Italia, vediamo due strutture
interne completamente diverse''. Non esiste una scorciatoia per farsi assumere
da un editore.

Poi prende ``l'esempio più comune di tutti'', la media dei giochi da tavolo
pubblicati (figura~\ref{a1:fig:relazioni}):
\begin{itemize}
  \item il \textbf{game designer} è una figura \textbf{esterna}, con un
        \textbf{contratto a royalty} con il publisher;
  \item la sua figura di riferimento è il \textbf{publisher} stesso, e la
        conversazione designer--publisher è ``molto calda'';
  \item se il publisher ha dei \textbf{developer}, il developer entra in questa
        conversazione: designer e developer collaborano molto, e internamente
        developer e publisher collaborano molto. Spesso, nei publisher piccoli,
        è il publisher a fare le veci del developer;
  \item \textbf{graphic design} e \textbf{illustrazione} tendono a parlare molto
        con il publisher e \textbf{molto poco con il designer}. A meno che non si
        tratti di uno studio indie molto piccolo, è difficile che il designer
        sappia anche solo chi è l'illustratore prima che il gioco esca.
\end{itemize}

\begin{figure}[htbp]
\centering
\begin{tikzpicture}[
    font=\small,
    ruolo/.style={draw, thick, rounded corners=3pt, fill=gray!12,
                  minimum width=2.8cm, minimum height=0.9cm, align=center},
    esterno/.style={draw, thick, dashed, rounded corners=3pt,
                    minimum width=2.8cm, minimum height=0.9cm, align=center},
    calda/.style={{Stealth[length=2.5mm]}-{Stealth[length=2.5mm]}, very thick},
    tiepida/.style={{Stealth[length=2mm]}-{Stealth[length=2mm]}, thin, dashed}
  ]
  \node[esterno] (gd) at (0,0) {Game designer\\ \footnotesize esterno, a royalty};
  \node[ruolo] (dev) at (5,1.4) {Developer\\ \footnotesize se c'è};
  \node[ruolo] (pub) at (5,-1.4) {Publisher};
  \node[ruolo] (art) at (10,-1.4) {Graphic designer\\ e illustratore};
  \draw[calda] (gd) -- (dev);
  \draw[calda] (gd) -- (pub);
  \draw[calda] (dev) -- (pub);
  \draw[calda] (pub) -- (art);
  \draw[tiepida] (gd.south) to[bend right=25]
     node[below, font=\footnotesize\itshape] {quasi nessun contatto} (art.south);
  \node[draw, dotted, rounded corners=4pt, fit=(dev)(pub)(art), inner sep=8pt,
        label={[font=\footnotesize\itshape]above right:interno al publisher (di solito)}] {};
\end{tikzpicture}
\caption{Le relazioni tipiche intorno a un gioco da tavolo pubblicato, secondo la
risposta di Ruggero. Le frecce spesse indicano una collaborazione stretta e
frequente; quella tratteggiata, un contatto quasi assente. Nei publisher piccoli
il publisher fa spesso anche da developer.}
\label{a1:fig:relazioni}
\end{figure}

\subsection{Fare da developer di sé stessi?}

Rita chiede in chat se un designer può fare anche da developer del proprio
gioco, o se Ruggero lo sconsiglia. In generale \textbf{lo sconsiglia}, per lo
stesso motivo per cui si sconsiglia in altri media: i \textbf{darling}. Il
designer ha delle meccaniche a cui è affezionato, e il developer a volte dovrà
\textbf{tagliarle}; è molto difficile, se si è ``gestori di sé stessi'', tagliare
cose che si vogliono nel gioco. ``È un po' come lasciare fare il lavoro di editing
allo scrittore che ha già scritto il libro'': può rivedersi, certo, ma difficilmente
cambierà cose che farebbe bene cambiare. Nelle lezioni con Marta si vedrà perché,
e alla fine del corso ``non avrete dubbi nemmeno voi'' nel separare designer e
developer.

Ma tutto dipende dallo \textbf{scopo del progetto}: se si fa un gioco in tre,
qualcuno quel lavoro deve farlo per forza, e magari il designer fa anche il
developer, compensando con \textbf{più playtest}. ``Il mondo attuale e il mondo
ideale sono due cose molto diverse, e lo vediamo non solo dai giochi ma dalla
politica'', ma ``non entriamo in questo vaso di Pandora''.

\subsection{Il ruolo indispensabile e la pubblicazione}

Lo studente di prima ne trae una conclusione: senza game designer il gioco non
esiste. ``Quello sicuramente''. Ma Ruggero coglie l'occasione per approfondire
un punto che forse non era passato: \textbf{fare design di un gioco per
pubblicarlo non è l'unica opzione possibile}. Un gioco non è ``meno gioco''
se non è pubblicato da un editore. È legittimo mettere un gioco su itch.io in
\textbf{print and play}, fare un gioco a evento senza venderlo, fare un
\emph{murder party}, una escape room, un board game per una persona amica:
``quell'esperienza lì non è meno gioco''. Dipende dall'obiettivo del progetto.

Ruggero ha lavorato molto anche con le \textbf{scuole}, dove il publisher non
esiste: è un gioco per le classi, e la figura più vicina è l'editoria scolastica,
``che però non fa il publisher''. Ci sono compromessi diversi sulla materialità
e su come il gioco si presenta, che si vedranno nella prototipazione e nei casi
studio. Quindi l'unica figura che deve esserci sempre è \textbf{chi fa game
design}, anche se magari non sa di chiamarsi game designer e non ha mai
sentito parlare di game design: ``questo non la rende meno game designer''.

\paragraph{Un libro consigliato.} Ruggero fa anche mentoring e gestisce gruppi
di settore, e spesso gli chiedono che cosa leggere per imparare il game design.
Non è il tipo che obbliga a leggere libri (``non vi obbligherò a leggere
niente''), ma ce n'è uno che consiglia sempre, ed è il primo che ha scelto di
leggere: \textbf{\emph{Rise of the Videogame Zinesters}} di \textbf{Anna
Anthropy}\footnote{La trascrizione dice ``Anna Antrofi'' e ``Rise of the
Videogames in Esters''.}. È un libro un po' datato, ma ancora reperibile; parla
soprattutto di videogiochi, ma ``per l'80\%'' il contenuto si applica al gioco e
al game design in generale. È un libro sulla \textbf{democratizzazione del gioco}:
il gioco come strumento per parlare di sé e come \textbf{medium legittimo}. Come
esistono film drammatici, tristi, intensi, possono esistere giochi su tematiche
delicate, e se ne vedranno nei casi studio. Molti giochi fatti con questa
direttiva non hanno un publisher, che tende a pubblicare ciò che dà un ritorno
commerciale.

\paragraph{Commerciale e personale.} ``Io spero che voi non farete giochi solo
per scopo commerciale''. Certo, servono un tetto sulla testa, il cibo in
tavola, un posto in cui dormire, e così si finisce per fare giochi che non sono
il proprio gioco dei sogni. Spesso è un mix: Ruggero ha fatto giochi commerciali
``per pagarmi l'affitto'' e giochi che voleva fare e su cui ha perso soldi,
``sapendolo proprio che sarebbero stati un fallimento''. I giochi, come i libri e
le canzoni, si fanno per molti motivi, e quelli commerciali di successo non sono
più importanti degli altri. Tra i giochi della sua slide c'erano il gioco da
tavolo ufficiale di \emph{Cyberpunk}\footnote{La trascrizione dice ``Cyberpunk
del 1977'': quasi certamente un gioco legato a \emph{Cyberpunk 2077}.} e
\emph{Marvel United}, su IP grosse e famose, ma anche giochi fatti di ``una
candela e una busta su un tavolo''. E se gli si chiede quale preferisce in
termini di contenuto, non sono quelli sui supereroi.

\subsection{Il gioco di ruolo in Italia}

Nicola, che è interessato soprattutto al gioco di ruolo e conosce già persone
del settore, ha ``una leggera paura'': che chi studia game design finisca per
fare il \textbf{freelance} a vita, ``perché non ci sarà mai un'occupazione per
bene'', o il \textbf{copywriter}. Conosce un ragazzo che ha lavorato in quel
mondo e che ora fa altro.

Ruggero risponde che la maggioranza dei giochi nel suo portfolio sono giochi di
ruolo, quindi il tema lo tocca da vicino, e che in Italia quello che dice Nicola
``è semplicemente vero''. Molti publisher piccoli hanno una persona di
riferimento che copre più ruoli: impagina, fa i playtest, fa un po' di
development. È ``un settore economicamente completamente disastrato'', per
motivi culturali che non apre: oggi in Italia è praticamente impossibile
\textbf{vivere di giochi di ruolo}. Chi ci lavora lo fa nel tempo libero, al di
fuori del lavoro principale, oppure coprendo molte altre cose: molti publisher
italiani di giochi di ruolo pubblicano anche titoli esteri e localizzano, non
fanno solo cose originali. ``È un grosso calderone''. Poi si ferma: ``questo è
un \emph{rant}'', e ``i \emph{rant} tornano sempre'', ma non vuole che gli
studenti escano dalla prima lezione ``con la disperazione nel cuore''. Ci sarà
spazio per parlarne quando avranno gli strumenti per capirlo, ``per non farvi
abbattere''.

\subsection{Che cosa bisogna sapere già}

Uno studente chiede se ci sono nozioni che dovrebbero già avere e che possono
recuperare: per esempio, sente parlare di giochi ``German'' o ``American'' e non
sa che cosa siano. Ruggero spiega l'organizzazione del corso:
\begin{itemize}
  \item le \textbf{prime due lezioni} sono di fondamenta. Questa era ``dall'alto'':
        come è strutturato il corso, le tempistiche;
  \item la \textbf{prossima lezione} è di tre ore di soli \textbf{casi studio},
        ``a manetta, a raffica'', divisi per categorie e tipologie, con
        ragionamenti su categorie e generi di gioco: lì si vedranno German,
        American ``e anche altre più utili'';
  \item nel resto del corso le terminologie specifiche sono inserite nelle
        lezioni dedicate: il \textbf{paragame} nelle lezioni sul paragame, il
        \textbf{modello matematico} nelle lezioni di Marta sul development.
\end{itemize}
Per ora agli studenti si chiede di sapere quello che si è visto oggi:
meccaniche, dinamiche ed estetiche; le figure professionali, designer e
developer; e soprattutto \textbf{game e play}.

\subsection{Game e play: i quattro quadranti}

``Specialmente \emph{game} e \emph{play}'', insiste Ruggero: è più difficile di
quello che sembra, perché l'italiano non abitua a separare le due cose, ed è
possibile averle usate per anni come sinonimi. Nella sua esperienza con chi si
avvicina al design è ``uno dei punti più critici''. Spesso i commenti che si
fanno sui giochi giocati sono commenti sul \emph{play} applicati al
\emph{game}: ``il gioco non mi è piaciuto perché, giocando con Tizio, è successo
questo''; ``il gioco mi sembra sbilanciato perché Tizio al tavolo ha fatto
questo''.

Separare \emph{game} e \emph{play} serve a capire quali giochi sono progettati
bene e quali male, e soprattutto a riconoscere i \textbf{quattro possibili
incroci} (figura~\ref{a1:fig:quadranti}): un gioco progettato bene che piace
giocare; uno progettato male che non piace; uno progettato male che piace; uno
progettato bene che non piace. ``Non è importantissimo in generale nella vita,
ci sono cose più importanti'', ma per chi fa game design lo è: ``a volte vi
chiederanno di lavorare su un gioco che non vi piace giocare, e quel gioco deve
essere fatto bene lo stesso''.

\begin{figure}[htbp]
\centering
\begin{tikzpicture}[
    font=\small,
    cella/.style={draw, thick, minimum width=4.4cm, minimum height=1.9cm,
                  text width=4.0cm, align=center}
  ]
  \node[cella, fill=gray!20] (a) at (0,0)     {\textbf{progettato bene}\\ e \textbf{mi piace} giocarlo};
  \node[cella]               (b) at (4.4,0)   {\textbf{progettato bene}\\ ma \textbf{non mi piace} giocarlo};
  \node[cella]               (c) at (0,-1.9)  {\textbf{progettato male}\\ ma \textbf{mi piace} giocarlo};
  \node[cella, fill=gray!8]  (d) at (4.4,-1.9){\textbf{progettato male}\\ e \textbf{non mi piace} giocarlo};
  \node[font=\footnotesize\itshape, above=0.15cm of a] {mi piace (\emph{play})};
  \node[font=\footnotesize\itshape, above=0.15cm of b] {non mi piace (\emph{play})};
  \node[font=\footnotesize\itshape, rotate=90, left=0.15cm of a, anchor=south]
       {bene (\emph{game})};
  \node[font=\footnotesize\itshape, rotate=90, left=0.15cm of c, anchor=south]
       {male (\emph{game})};
\end{tikzpicture}
\caption{I quattro incroci tra la qualità del \emph{game} (come è progettato) e
l'esperienza di \emph{play} (se piace giocarlo), descritti da Ruggero. Per un
designer i quadranti ``scomodi'' sono i due fuori diagonale: un gioco può
piacere pur essendo progettato male, e va saputo progettare bene anche un gioco
che non piace.}
\label{a1:fig:quadranti}
\end{figure}

\paragraph{La domanda di Tancredi.} Tancredi chiede conferma: il \emph{play}
non è il ``giocare senza regole'', ma può essere anche il giocare a un gioco? Sì:
il \emph{play} è l'\textbf{atto}, e può essere fatto con un \emph{game} o senza.
E un \emph{game} può esistere senza che nessuno lo giochi: gli scacchi non sono
meno scacchi perché in quel momento nessuno ci gioca. E se, giocando, una
persona applica le regole in modo diverso? Allora \textbf{non si sta giocando
quel gioco}.

\subsection{House rule: il caso di Monopoly}

Qui Ruggero fa un po' di \emph{foreshadowing} su una lezione di Daniele: il
discorso degli \textbf{hack}, delle \textbf{house rule}, delle \textbf{mod},
cioè di come modificare un gioco esistente. Il punto è che \textbf{un gioco
modificato da una house rule è un altro gioco}: cambiando le regole cambiano
necessariamente anche le \textbf{dinamiche} che ne derivano, e quindi
l'esperienza. Si immagini \emph{Monopoly} in cui ogni turno si possa, volendo,
prendere gratis dei soldi dalla banca: tutto il resto è identico, eppure il gioco
cambia.

E questo è talmente interiorizzato che ci sono regole che si applicano ``senza
neanche rendercene conto''. Ruggero chiede quanto duri in media una partita a
\emph{Monopoly}. Gianmarco: ``quattro ore''; altri: ``un pomeriggio''; un altro:
``finché uno dei giocatori non aggiunge la regola: ora vince chi\ldots''. In
realtà, \textbf{seguendo le regole scritte}, una partita non dura mai più di
\textbf{un'ora e venti, un'ora e mezza}. La domanda che fa sempre, e la risposta
è sempre la stessa: perché si pensa che \emph{Monopoly} duri tantissimo? Qual è
la regola che nessuno rispetta?

Uno studente indovina: \textbf{le aste}. Secondo il regolamento, se chi finisce
su un terreno non lo compra, il terreno \textbf{va all'asta} tra i giocatori, e
chiunque può aggiudicarselo. Molti invece giocano come se finire su un terreno
desse solo la \emph{facoltà} di comprarlo o di metterlo all'asta\footnote{Più
precisamente, nel regolamento ufficiale, se il giocatore che finisce su una
proprietà libera decide di non comprarla, la banca la mette all'asta e tutti i
giocatori possono fare offerte, compreso chi l'ha rifiutata. L'asta quindi non è
una scelta del giocatore, ma la conseguenza obbligata del suo rifiuto.}. Con la
regola applicata, di aste se ne fanno moltissime: all'ottavo turno ci sono già
otto terreni assegnati, al ventesimo venti, la gente va in bancarotta
velocissimamente e la partita dura un'ora e mezza. Anche le partite competitive
di \emph{Monopoly}, nota Ruggero, durano un'ora e venti, un'ora e mezza.

Uno studente obietta che giocava con le aste e durava comunque quattro ore.
``Avete dei nemici con l'analisi paralisi'', ride Ruggero; e cita un'altra house
rule tipica dell'infanzia: genitori che danno soldi per non far andare in
bancarotta. Ma \textbf{la bancarotta è parte del gioco}: se la si toglie,
\emph{Monopoly} non regge. Sul tema che \emph{Monopoly} vuole trattare, e sul
suo antenato (\emph{The Landlord's Game}), ``non apriamo adesso un altro vaso di
Pandora''.

\paragraph{Approfondimento: Monopoly, la durata e le house rule nei libri.}
L'esempio di Ruggero si ritrova quasi identico in \emph{The Art of Game Design}.
Parlando del bilanciamento tra partite brevi e lunghe (\emph{Short vs. Long}),
Schell scrive che \emph{Monopoly}, giocato secondo le regole ufficiali,
``finisce spesso in circa novanta minuti'', ma molti giocatori trovano quelle
regole troppo dure e le modificano, distribuendo premi in denaro e allentando i
vincoli sull'acquisto delle proprietà: con l'effetto collaterale di far durare la
partita molto di più, ``tipicamente tre ore, o anche di più''. Schell classifica
le \textbf{house rule} come regole create dai giocatori in risposta a un
difetto percepito dopo qualche partita. Anche \emph{Challenges for Game
Designers}, nel capitolo sull'aggiungere e togliere meccaniche, parte dallo
stesso gioco: il \emph{Parcheggio gratuito}, che per le regole originali non fa
succedere nulla, mentre quasi tutti giocano con una house rule che dà un bonus a
chi ci finisce sopra. Le due fonti confermano la lezione: una piccola regola
cambiata cambia le dinamiche, e le dinamiche cambiano l'esperienza.

\subsection{Il fine del game è il play?}

Un altro studente commenta: va bene la distinzione, ma il fine del \emph{game} è
il \emph{play}; quello che si vede è sempre l'effetto. Ruggero è d'accordo: di
solito, se si fa un gioco, lo si fa per farlo giocare, come si scrive un libro
per farlo leggere (anche se esistono, ed ``è già stato esplorato'', libri scritti
per non essere letti e giochi per zero giocatori). Ma la distinzione serve a una
cosa precisa: \textbf{separare, nel playtest, i feedback sul gioco dai feedback
sul gruppo di gioco} e sulle sue dinamiche. Se una persona esce frustratissima da
un playtest, non è necessariamente colpa del gioco; se una persona dice ``questo
è il gioco migliore del mondo'', potrebbe non essere merito del gioco ma del
gruppo. Capire a che cosa si riferisce un feedback serve a capire \textbf{dove
iterare} il prototipo.

Lo studente osserva che l'area è sfumata: alcuni giochi si basano proprio
sull'avere le persone giuste al tavolo. Un gioco di \textbf{negoziazione} con
giocatori abituati ai giochi ``multi-solitario'' diventa ingiocabile; e nei giochi
a \textbf{economia chiusa} (cita \emph{Container}\footnote{La trascrizione cita
anche un secondo titolo non ricostruibile (``Tokyo su Vigimarket'').}) un gruppo
economicamente conservativo, che spende poco e non fa girare l'economia, fa
durare la partita il doppio. Ruggero rilancia: e se il gioco dura il doppio,
questa dinamica cambia la \textbf{percezione di peso} del gioco, le energie spese,
quanto si è stanchi a fine partita. ``Tutto è collegato. \textbf{I giochi sono
ecosistemi di meccaniche}, e cambiandone una, le altre cambiano di
conseguenza''.

\paragraph{La parola vietata, alla prova.} Lo studente aggiunge che ci sono giochi
che ha trovato ``molto sbagliati, ma che mi divertono''. ``Attenzione: non si può
dire \emph{divertimento}!'', lo ferma Ruggero; nelle presentazioni aveva lasciato
correre, ``ma lì non vi avevo ancora letto la regola''. Lo studente ci prova:
``giochi che poi voglio \textbf{rigiocare}'', anche se ``sono evidentemente
sbagliati, rotti''. Ruggero apprezza, ma non vuole anticipare cose che meritano il
loro tempo: se ne parlerà nelle lezioni sui framework e soprattutto in quella di
\textbf{emotion design}, ``forse la mia lezione preferita del corso''. (Lo
studente ammette che ``forse alla fine si gioca per quello, per l'emozione''.)

\subsection{I giochi ``comunicati male''}

L'ultima domanda corposa riprende un'affermazione di inizio lezione: ci sono
giochi belli ma comunicati male. Qualche esempio? ``Mi basta girarmi per
vederne 48'', risponde Ruggero indicando gli scaffali alle sue spalle. Più che
titoli, elenca le \textbf{ragioni}:
\begin{itemize}
  \item \textbf{l'Italia localizza lentamente}. Basta guardare i premi del gioco
        dell'anno in Italia: spesso vanno a giochi usciti due o tre anni prima in
        Germania, Francia, Stati Uniti o Regno Unito, perché localizzare richiede i
        suoi tempi. Il mercato italiano è quindi molto dipendente dall'estero,
        ``molto impaurito'', e spesso aspetta che un gioco vada bene all'estero
        prima di portarlo in Italia; e alcuni giochi ``perdono il treno'';
  \item \textbf{il publisher originale chiude}: come si pubblica un gioco se non
        c'è più l'azienda con cui fare un contratto? A Ruggero è successo più
        volte;
  \item \textbf{il dipartimento commerciale italiano non sa comunicare il gioco}:
        si fa male un'edizione, la stampa vende pochissimo e non viene più
        ristampata;
  \item \textbf{la scala del mercato}: in Italia un gioco da tavolo di peso medio
        (un \emph{eurogame}, un \emph{German}) che vende 1.500 copie ``è tanto'';
        all'estero se ne vendono magari 10.000. Con 10.000 copie in inglese ci
        sono i soldi per una campagna di comunicazione: un trailer, i
        recensori, i tutorial online. In Italia tutto questo è ``ridotto all'osso'',
        e la cosa si riflette su stampa e divulgazione: anche i canali per
        informarsi su un gioco prima di comprarlo sono pochi (se ne parlerà nella
        lezione su ricerca e sviluppo). Un gioco spesso non arriva al mercato non
        perché sia brutto, ma perché l'editore non lo sa comunicare, ne fa poche
        copie e va fuori produzione;
  \item \textbf{il tempismo}: party game progettati nel 2019 e usciti nel 2020, in
        pieno Covid, quando non si poteva giocare al tavolo, sono falliti tutti, e
        gli editori ci hanno perso. ``Non potevano immaginarselo'';
  \item \textbf{i temi caldi}: dopo l'invasione russa dell'Ucraina e la guerra in
        Palestina, alcuni editori non pubblicano o non localizzano giochi
        strategici e di guerra. Ruggero ha lavorato a \textbf{\emph{Status}}, un
        gioco sulla geopolitica che tocca, tra gli altri temi, l'invasione russa
        dell'Ucraina e il genocidio in Palestina, e alcuni editori che di
        solito localizzano i loro giochi (per esempio il partner di riferimento in
        Russia) non lo hanno localizzato per motivi politici.
\end{itemize}
Lo studente voleva capire che cosa voglia dire comunicare male: non metterci
abbastanza energia, o comunicare una cosa per un'altra, ``come il trailer
sbagliato, che poi vai al cinema e ti trovi un altro film''. ``È un mix'', e la
domanda ``era complessa e meritava più spazio''.

\subsection{I giochi di carte collezionabili e i saluti}

L'ultima domanda (``posso rompere le balle un'ultima volta?'') riguarda i
\textbf{giochi di carte collezionabili}: funzionano con le stesse regole dei
giochi da tavolo? Se ne parlerà alla prossima lezione: è proprio il motivo per
cui la lezione sui casi studio dura tre ore, perché ogni genere di gioco ha le
sue dinamiche e le sue regole, e anche il modo di pubblicare cambia molto tra un
TCG, un gioco di ruolo da tavolo, un LARP e un gioco pervasivo.

Ruggero saluta: le email dei tre docenti sono all'inizio delle slide, la
registrazione arriverà, ``ci si vede questo venerdì''.

% ---------------------------------------------------------------------
\section*{Punti chiave del capitolo}
\addcontentsline{toc}{section}{Punti chiave del capitolo}

\begin{itemize}
  \item Il corso di \textbf{Analog Game Design} (seconda edizione) dura circa tre
        mesi, fino a dicembre, ed è tenuto da tre docenti principali:
        \textbf{Daniele} (narrative design, gioco di ruolo, fiction),
        \textbf{Ruggero} (game design, ricerca e sviluppo) e \textbf{Beatrice}
        (giochi pervasivi), con tre ospiti su RNG, development e game production.
  \item Lo scopo è formare \textbf{autori analogici} consapevoli: strumenti da
        usare anche fuori dal corso, non il ``gioco dei sogni''. Il corso è
        \textbf{intersezionale} (nessuna lezione dedicata a un solo tipo di
        gioco), \textbf{cross-generazionale}, \textbf{pratico e adattabile},
        basato sulle \textbf{domande} (nessuna domanda è stupida), sul
        \textbf{feedback} invece che sul voto e sul \emph{fail forward}.
  \item Il \textbf{project work} si fa in tre gruppi da tre, formati dai docenti;
        pitch il 12 ottobre, sviluppo nelle due pause, tutoraggi il sabato,
        consegna il 4 dicembre, presentazione l'11 e giornata finale con gli
        editori il 14. Si valutano documentazione, presentazione,
        autovalutazione, crescita e atteggiamento.
  \item \textbf{Niente intelligenza artificiale} nel corso: i publisher
        rifiutano pitch e prototipi con immagini generate, i regolamenti scritti
        dall'IA descrivono la stessa meccanica in modi diversi, e soprattutto chi
        la usa perde l'occasione di imparare. Un prototipo può essere fatto a
        matita.
  \item Le \textbf{figure dell'industria}: game designer, game developer,
        editor, publisher, sales e marketing, graphic designer, illustratore,
        sensitivity reader (che non è un censore) e producer. Nelle realtà
        piccole una persona copre più ruoli.
  \item Nell'analogico il \textbf{game designer} progetta l'esperienza creando
        le meccaniche; il \textbf{game developer} fa da \textbf{mediatore} con il
        publisher (bilanciamento, modelli matematici, adattamenti). Pretendere
        che il proprio gioco non venga modificato è un \emph{red flag}.
  \item Il \textbf{marketing} non decide che cosa pubblicare ma come; è
        un'azione, ed esiste anche se lo si ignora.
  \item In \textbf{Italia} il gioco è ancora visto come infantile, il mercato
        punta sulla localizzazione e investe poco su giochi e designer interni;
        all'estero (soprattutto dalla Germania in su) ci sono tradizione adulta,
        premi di prestigio e developer interni.
  \item \textbf{Game} è l'artefatto (le regole, il sistema, l'oggetto);
        \textbf{play} è l'atto di giocare. Si fanno giochi, non giocattoli.
        Separarli serve a distinguere i feedback sul gioco da quelli sul gruppo
        e a riconoscere i quattro incroci tra ``progettato bene/male'' e ``mi
        piace/non mi piace''.
  \item \textbf{MDA}: meccaniche (regole scritte), dinamiche (regole in azione),
        estetiche (esperienza, non grafica); nel \textbf{DPE}: design, play,
        experience.
  \item \textbf{Parole vietate} in ogni materiale scritto: \emph{fun},
        \emph{immersion}, \emph{entertainment} (divertimento, immersione,
        intrattenimento). Bisogna trovare parole più precise.
  \item Distinguere \textbf{chi crea, chi gioca e chi facilita}; nessun gioco
        nasce perfetto: \textbf{playtest e iterazione} sono inevitabili, e un
        \textbf{dizionario comune} aiuta a non perdersi.
  \item \textbf{Tabletop Simulator} è il software di settore per giocare (non
        creare videogiochi) i giochi da tavolo in digitale. Primo task: giocarci
        tre o quattro volte entro il 10 ottobre.
  \item Una \textbf{house rule} crea un altro gioco: \emph{Monopoly} con le aste
        applicate come da regolamento dura circa un'ora e mezza.
\end{itemize}

%
%  (Con hyperref conviene aggiungere anche
%  \renewcommand{\theHchapter}{agd.\arabic{chapter}} nella parte, per
%  evitare ancore duplicate con i capitoli 1, 2, ... di Game Design.)
%  Tutte le etichette di questo capitolo hanno il prefisso "a1:"
%  (a1:cap:..., a1:sec:..., a1:fig:..., a1:tab:...). Il capitolo è
%  autonomo: non contiene \ref verso altri capitoli.
%  Le figure sono ricostruzioni schematiche di quanto i docenti spiegano
%  a voce e con le slide, non copie delle slide.
% =====================================================================

\chapter[Introduzione all'Analog Game Design]{Introduzione all'Analog Game
Design: il corso, le figure dell'industria e un primo linguaggio comune}
\label{a1:cap:intro}

Questa è la \textbf{prima lezione} del corso di \textbf{Analog Game Design},
cioè di progettazione di giochi \emph{analogici}: giochi da tavolo, giochi di
carte, giochi di ruolo da tavolo, giochi di ruolo dal vivo (LARP), giochi
\emph{pervasivi} ed esperienze a evento. È un corso \textbf{breve}, di circa tre
mesi (``da qui a dicembre''), separato dal corso di Game Design, e siamo alla sua
\textbf{seconda edizione}. Come avvertono subito i docenti, è ``una giornata di
introduzione'': si spiega come funziona il corso, chi sono i docenti, con quali
metodi e con quale calendario si lavorerà, e ci si presenta a vicenda. ``C'è poco
gioco analogico'', si scusano. Ma nella seconda parte la lezione comincia
comunque a costruire un \textbf{linguaggio comune}: le figure professionali
dell'industria del gioco analogico, la differenza tra \emph{game design} e
\emph{game development}, la situazione del mercato italiano rispetto a quello
estero, e un primo glossario di design (\emph{game} e \emph{play}, il framework
MDA, le parole ``vietate'', il playtest e l'iterazione). Si chiude con lo
strumento di settore per i prototipi digitali dei giochi da tavolo,
\textbf{Tabletop Simulator}, e con una lunga sessione di domande.

La lezione si svolge \textbf{online}, su Zoom, ed è condotta a più voci dai tre
docenti principali del corso: \textbf{Daniele}, narrative designer, che guida la
parte organizzativa; \textbf{Ruggero}, game designer, che guida la parte sui
fondamenti e le domande finali; \textbf{Beatrice}, progettista di giochi
pervasivi, che interviene soprattutto nella prima metà. Gli studenti iscritti
sono \textbf{nove}; uno è assente e un altro arriva a lezione iniziata. Ci sono
una pausa di una decina di minuti (si riprende alle 18:15) e, dopo le 19:15,
un'ultima mezz'ora di domande con il solo Ruggero.

Tre avvertenze di lettura. La prima: la trascrizione è automatica e storpia
molti nomi propri, compresi quelli dei docenti e dell'accademia (``Dive'',
``IVE'', ``Accademia Live'', ``Ai''); dove la ricostruzione è sicura, o molto
probabile, il capitolo usa la forma corretta e lo segnala in nota; dove non lo
è, riporta la forma della trascrizione e lo dice. La seconda: durante la pausa la
trascrizione registra, ripetuta in loop per decine di righe, una frase senza
senso che non ha nulla a che fare con la lezione, e che il capitolo ignora. La
terza: in molti passaggi più docenti parlano uno dopo l'altro senza che la
trascrizione segni il cambio di voce; il capitolo attribuisce le parole a chi le
pronuncia quando il contesto lo rende chiaro (chi parla di sé, chi viene
chiamato per nome, chi ``passa la parola''), e resta generico altrimenti.

Quanto alla data, la lezione si colloca con ogni probabilità \textbf{nella prima
metà di settembre del 2026}\footnote{La data non è mai detta. La si deduce da
alcuni indizi: il form iniziale va compilato ``entro e non oltre il 15
settembre''; i gruppi verranno annunciati il 25 settembre; le date del
calendario (sabato 3 ottobre, lunedì 12 ottobre, sabato 7 e 14 novembre,
venerdì 4 dicembre) cadono nei giorni della settimana indicati soltanto nel
2026; un docente parla di un argomento ``molto complesso nel 2026''; e fa
ancora molto caldo (``le temperature sono ancora indecenti''). È una
deduzione, non un'informazione esplicita.}: i riferimenti a ``oggi'',
``quest'estate'' e ``l'anno scorso'' vanno letti rispetto a quel momento.

Il percorso del capitolo segue quello della lezione:
\begin{enumerate}
  \item l'avvio, i \textbf{docenti} e le \textbf{presentazioni} dei partecipanti
        (sezioni da~\ref{a1:sec:avvio} a~\ref{a1:sec:partecipanti});
  \item i \textbf{materiali} e gli strumenti del corso: il form, la cartella
        condivisa, Discord (sezione~\ref{a1:sec:materiali});
  \item lo \textbf{scopo} e il \textbf{metodo} del corso, le regole di
        comportamento e il carico di lavoro (sezioni~\ref{a1:sec:scopo}
        e~\ref{a1:sec:regole});
  \item il \textbf{project work} e il \textbf{calendario} (sezione~\ref{a1:sec:project});
  \item la regola sull'\textbf{intelligenza artificiale} e la lunga discussione
        che ne segue (sezione~\ref{a1:sec:ia});
  \item le \textbf{figure professionali} dell'industria del gioco analogico,
        con un approfondimento su marketing e sales
        (sezioni~\ref{a1:sec:figure} e~\ref{a1:sec:marketing});
  \item la differenza tra \textbf{game design} e \textbf{game development}
        (sezione~\ref{a1:sec:design-development});
  \item il \textbf{mercato}: Italia ed estero (sezione~\ref{a1:sec:mercato});
  \item i \textbf{fondamenti di design}: perché serve una terminologia comune,
        \emph{game} e \emph{play}, il framework \textbf{MDA} (e il DPE), le
        \textbf{parole vietate}, chi crea, chi gioca e chi facilita, il
        \textbf{playtest} e l'\textbf{iterazione} (sezioni
        da~\ref{a1:sec:terminologia} a~\ref{a1:sec:playtest});
  \item \textbf{Tabletop Simulator} e il primo task (sezione~\ref{a1:sec:tts});
  \item le \textbf{domande finali}: il tema del project work, il bilanciamento,
        i ruoli nei team, la pubblicazione, il gioco di ruolo in Italia, le
        \emph{house rule} di Monopoly, i giochi ``comunicati male''
        (sezioni~\ref{a1:sec:domande-daniele} e~\ref{a1:sec:domande}).
\end{enumerate}

% ---------------------------------------------------------------------
\section{L'avvio della lezione}
\label{a1:sec:avvio}

La lezione comincia con qualche minuto di attesa. ``Siamo in nove in classe'', e
si aspetta che arrivino tutti; uno dei docenti scherza sul fatto che l'unica
persona che arriverà ``sempre in ritardo di venerdì'' è Gianmarco. Mancano
Giordano e Tiziana; Tiziana arriva, mentre la segreteria fa sapere poco dopo che
Giordano non potrà esserci e vedrà la registrazione. Si aspetta anche Beatrice,
la terza docente, che arriva a sua volta dopo pochi minuti.

Nell'attesa, Daniele anticipa il programma della giornata e qualche
informazione pratica:
\begin{itemize}
  \item oggi è una giornata di \textbf{introduzione}: come funziona il corso, chi
        sono i docenti, le presentazioni degli studenti, le \textbf{metodologie}
        e il \textbf{calendario}. Gli studenti avranno ``un fracco di domande'', e
        alla fine della lezione c'è ``proprio uno slot apposta'' per le domande;
  \item le lezioni sono \textbf{registrate}: di solito la registrazione viene
        consegnata dopo circa \textbf{12--24 ore} sul profilo personale della
        piattaforma dell'accademia, che va attivato con la segreteria (Tiziana,
        già studentessa, probabilmente ce l'ha già);
  \item i docenti ``cuore'' del corso sono tre (Daniele, Ruggero e Beatrice),
        ``soprattutto Ruggero''; a loro si aggiungeranno tre docenti ospiti, per
        meno ore, probabilmente una lezione a testa, ``tranne Marta, che avrà due
        lezioni quest'anno''.
\end{itemize}
E una battuta sul caldo: ``spero che voi siate carichi, anche se le temperature
sono ancora indecenti''. ``Carichi moderatamente, per quello che potete''.

Poi Daniele condivide lo schermo (``non solo le slide, ma anche un po' di altre
cose'') e apre ufficialmente il corso: ``benvenuti e benvenute nella seconda
edizione del corso di Analog Game Design. Questa è la prima lezione''. Il
percorso è di circa \textbf{tre mesi}, fino a dicembre, e durante il corso i
docenti cercheranno di far provare agli studenti ``cosa significa fare Analog
Game Design'' con un \textbf{project work}, un progetto di gruppo di cui si
spiegheranno più avanti la struttura e le richieste.

% ---------------------------------------------------------------------
\section{I docenti}
\label{a1:sec:docenti}

``Anzitutto partiamo da delle presentazioni doverose''. Ognuno dei tre docenti
principali si presenta con una slide che riporta i contatti e una selezione dei
lavori fatti.

\subsection{Daniele: narrativa, fiction e gioco di ruolo}

\textbf{Daniele}\footnote{La trascrizione riporta ``Daniel Cusetto''; nel
resto della lezione tutti lo chiamano ``Daniele''. Il cognome non è
ricostruibile con certezza.} ha già conosciuto quasi tutti gli studenti, perché
ha fatto loro il colloquio di ammissione. È \textbf{narrative designer}, o meglio
\emph{game narrative designer}; insegna nell'accademia che ospita il
corso\footnote{La trascrizione storpia il nome dell'accademia in molti modi
(``Dive'', ``IVE'', ``Accademia Live'', ``Ai'', ``AIVA''). Con ogni probabilità
si tratta dell'\textbf{AIV}, l'Accademia Italiana Videogiochi, che ha sede a
Roma: lo fanno pensare i corsi triennali di game design e narrative design citati
dai docenti e il fatto che una studentessa racconti di essersi trasferita a Roma.
È una ricostruzione, non un'informazione esplicita; nel seguito il capitolo
parla semplicemente dell'``accademia''.} ma lavora soprattutto come \textbf{brand
advisor}. Nel corso si occuperà della parte di \textbf{gioco di ruolo},
\textbf{narrazione} e \textbf{fiction}.

I suoi contatti, ``in quest'ordine'', sono email, Telegram e Discord. Ci tiene a
dare alcune indicazioni precise:
\begin{itemize}
  \item se si vuole avere la certezza di una risposta, bisogna scrivergli su
        \textbf{Telegram}, perché lì ha tutti i suoi clienti;
  \item Discord lo guarda di meno; l'\textbf{email} è per le cose ufficiali;
  \item non bisogna contattarlo altrove: non su Facebook, non su LinkedIn, non su
        Instagram, ``perché li vedo molto di meno''.
\end{itemize}
La slide mostra anche alcuni giochi, attività e brand su cui ha lavorato negli
anni, ma non si dilunga: ``nel caso ne parleremo''.

\subsection{Ruggero: game design e ricerca e sviluppo}

\textbf{Ruggero} (spesso chiamato ``Ruger'') ha fatto il \textbf{game designer}
per varie aziende e oggi è game designer e si occupa di \textbf{ricerca e
sviluppo} per un'azienda del settore\footnote{La trascrizione dice ``per MS'';
il nome dell'azienda non è ricostruibile con certezza. Più avanti Ruggero cita
tra i giochi su cui ha lavorato il gioco da tavolo ufficiale di
\emph{Cyberpunk}, \emph{Marvel United} e \emph{Status}, un gioco sulla
geopolitica.}. Ha un background ``misto'': nel suo portfolio ci sono giochi da
tavolo, giochi di ruolo, giochi analogici, LARP, videogiochi, giochi in realtà
virtuale. Si occupa anche di \textbf{divulgazione} sul gioco e di analisi di
\textbf{game studies} e game design. Più avanti, nelle domande finali, dirà di
avere ``50 e più giochi prodotti e progettati'' e di insegnare da vari anni.

Lo si può contattare liberamente via Gmail e Telegram, e cerca di rispondere
\textbf{entro 24 ore}, salvo imprevisti o eventi: ``sentitevi in libertà di
scrivermi se avete dubbi, senza problemi''. Quando arriva a parlare dei propri
giochi, la slide mostra ``un mix di varie cose, dalle più famose alle meno
famose'', ma Ruggero si imbarazza a presentarli, tra le risate dei colleghi.

\subsection{Beatrice: i giochi pervasivi}

\textbf{Beatrice} (``Bea'') si occupa di \textbf{giochi pervasivi}
(\emph{pervasive games}). ``Se non li avete mai sentiti nominare, non
preoccupatevi'': questo corso è ``uno dei pochi, forse l'unico in Italia'' che
li racconta esplicitamente. Li si vedrà nelle prossime lezioni, ma la definizione
di base è questa: i giochi pervasivi \textbf{utilizzano gli spazi e gli oggetti
della vita quotidiana e li trasformano in spazi e oggetti di gioco}.

Beatrice lavora come \textbf{freelance} e si occupa di diversi tipi di
progettazione: lavora sia con le \textbf{aziende} sia con gli \textbf{eventi dal
vivo}. La sua attività diretta si trova sotto l'etichetta \textbf{Minime
Differenze}\footnote{La trascrizione riporta il nome dell'etichetta come
``Minime Differenze'' e l'indirizzo del sito come ``mininedifferenze.com''; la
grafia esatta del sito non è verificabile dalla trascrizione.}, mentre i lavori
per le aziende sono ``sparpagliati'' sui siti delle varie aziende. Molte delle
cose mostrate nella sua slide sono attività proposte al pubblico ``con un
biglietto, con un gioco in formato evento'': un tema che si toccherà nel corso,
per capire come \textbf{comunicare} i propri giochi in modo che siano
comprensibili, accessibili per il pubblico e \textbf{vendibili}.

Ammette un ``grosso limite'' con i messaggi: è un po' lenta a rispondere, ``ma
rispondo sempre a tutto''. La si può contattare per email (cognome punto nome) o
su Telegram. È ``molto curiosa di vedere cosa produrrete''.

\subsection{I docenti ospiti}

Daniele presenta poi i tre docenti che il corso avrà ``solo per alcune
lezioni'', tutte verso \textbf{novembre}:
\begin{itemize}
  \item \textbf{Diego Ricchiuti}, direttore del corso di Game Design
        dell'accademia, farà una lezione sull'\textbf{RNG} (\emph{random number
        generation}, la casualità nei giochi): ``vi parleremo di RNG quando sarà
        necessario. Comunque avrete una lezione con lui, sappiatelo'';
  \item \textbf{Marta}\footnote{La trascrizione riporta il cognome come
        ``Cercasassi''; la grafia non è verificabile.} farà \textbf{due lezioni
        sul development}: quest'anno ``abbiamo deciso di fare doppietta'',
        perché l'anno scorso il suo intervento è stato ``molto, molto utile'';
  \item \textbf{Flavio}\footnote{La trascrizione riporta il cognome come
        ``Mortarino'' e, più avanti, ``Mortalino'', e il nome della sua casa
        editrice come ``Grampibert''; nessuno dei due è ricostruibile con
        certezza.} farà la lezione di \textbf{game production}; è un editore, e
        di solito partecipa anche alla tavola rotonda con gli editori della
        giornata finale.
\end{itemize}
Le slide sono già online e gli studenti avranno accesso alla cartella condivisa:
``non vi preoccupate di queste cose, ne parleremo a breve''. Ma prima di parlare
del corso, ``vogliamo conoscere voi''.

% ---------------------------------------------------------------------
\section{Le presentazioni dei partecipanti}
\label{a1:sec:partecipanti}

Le domande per la presentazione sono semplici, da affrontare ``a turno, senza
fretta, con calma'':
\begin{itemize}
  \item nome e occupazione (``se volete dirla'');
  \item perché si vogliono fare giochi analogici, e a che livello: che cosa ha
        portato lo studente al corso;
  \item che tipo di esperienze o conoscenze si hanno già nell'ambito, se si è già
        fatto qualcosa oppure no.
\end{itemize}
``Chi vuole partire senza nessun tipo di ritegno e di vergogne?''. Per far
parlare gli studenti guardandoli, Daniele interrompe la condivisione dello
schermo e ricorda lui le domande quando serve. Le presentazioni sono interessanti
anche per i commenti dei docenti, che anticipano diversi principi del corso.

\subsection{Tiziana}

Tiziana fa l'\textbf{infermiera}, ma nella vita ha fatto molti lavori:
impiegata, insegnante e ``tanti altri lavoretti''. Gioca ``da sempre'': il gioco
fa parte anche della sua vita familiare, perché suo fratello fa il programmatore
di videogiochi. Ha già \textbf{organizzato} giochi: una caccia al tesoro per far
conoscere il territorio, due scene con un delitto (giochi con omicidio da
risolvere), giochi per intrattenere le persone. Ha riscoperto di recente i giochi
da tavolo, arrivando a \emph{Carcassonne} con famiglia e amici; poiché è
difficile trovare persone con voglia di giocare nella propria cerchia ristretta,
ha scoperto il mondo delle \textbf{associazioni} ludiche, ``e lì è stato
un'esplosione'': prima in Veneto (è di Vicenza), poi a Roma, dove si è
trasferita.

Ha in mente un progetto ambizioso: far conoscere una certa realtà con la modalità
del \textbf{gioco di ruolo}, o comunque proporre un gioco suo.

\paragraph{Il commento: gli strumenti, non il gioco dei sogni.} Ruggero coglie
l'occasione per un'avvertenza che vale per tutti, visto che molti studenti
arrivano con idee già pronte. Nel corso si lavorerà a un progetto, ma
l'obiettivo dei docenti è dare \textbf{strumenti} da applicare anche
\emph{fuori} dal progetto, ai giochi che gli studenti vorranno fare un domani:
``che non siano esercizi per il corso, ma per voi in primis''. Quindi: ``non
rimaneteci male se non riuscirete a lavorare al vostro gioco dei sogni in questo
corso''. Anche perché ``sarebbe impossibile, in due mesi, fare il gioco dei
sogni''. Il corso darà invece la possibilità di fare un'esperienza ludica
insieme e di imparare a replicarla sulle proprie idee: ``le idee non buttatele,
anzi, il contrario: tenetele''.

\subsection{Nicola}

Nicola è uno \textbf{studente}. Sei mesi prima ha fatto un colloquio con Daniele
perché stava valutando di passare all'accademia, e ha deciso, ``nel dubbio'', di
cominciare da questo corso breve per capire come funziona la formazione;
eventualmente si iscriverà al corso di game design vero e proprio. Gli interessa
in particolare il lato dell'esperienza del giocatore, su cui è ``molto critico''
(qui usa la parola ``divertimento'', e Daniele lo ferma con un ``ne parleremo'':
è una delle parole che il corso vieta, come si vedrà nella
sezione~\ref{a1:sec:parole-vietate}).

Ha 21 anni ed è \textbf{giocatore di ruolo da circa otto}. Prima ancora di
conoscere i giochi di ruolo si era creato un sistema tutto suo, ``pessimo,
disgustoso, terribile'', un sistema a percentuale con il d100 ``terrificante'';
poi ha scoperto i giochi di ruolo da tavolo, i classici (\emph{Dungeons \&
Dragons}, \emph{Daggerheart}, \emph{Vampire}\footnote{La trascrizione riporta
``DD'', ``Daggerart'' e ``Vampire'': quasi certamente \emph{Dungeons \&
Dragons}, \emph{Daggerheart} e \emph{Vampire: The Masquerade}.}). Gioca anche a
\textbf{card game} e ad altri giochi sociali, pur avendo sempre avuto difficoltà
a trovare gruppi, ``avendo vissuto in un buco dimenticato dal mondo''. Come
giocatore è ``l'eterno narratore'' del gruppo, quello che prepara la storia e si
costruisce le regole. Come autore ha fatto soprattutto giochi su carta, ``perché
sono incapace a programmare, per adesso''.

Daniele osserva che il tema del giocatore e del suo rapporto con il gioco è ``un
argomento importante'', anzi ``un vaso di Pandora enorme'' (espressione che
tornerà molte volte nella lezione), e che il corso tratterà anche che cos'è
l'analog game design in generale e quanti tipi di prodotti esistono.

\subsection{Rita}

Rita\footnote{Daniele la chiama prima ``Nino'' (probabilmente un errore della
trascrizione) e poi la ringrazia come Rita.} ha 30 anni e si collega da una
regione del Sud, dove ``ancora si muore di caldo''. Ha studiato
all'\textbf{Accademia di Belle Arti di Napoli}, \textbf{design della
comunicazione visiva}: si occupa dell'aspetto più imprenditoriale della
progettazione artistica. Non volendo progettare per un'impresa (``fare loghi,
robe del genere''), vorrebbe applicare la propria creatività alla realizzazione
di giochi analogici.

Non è una giocatrice di ruolo e non lo è mai stata: è una giocatrice
\textbf{tradizionale}, cresciuta con Monopoly, la dama e giochi tradizionali
cinesi. Ha già \textbf{idee, prototipi e progetti} realizzati, e vuole capire
come renderli \textbf{prodotti per la vendita} e se questo è il campo adatto a
lei; magari anche confrontarsi con i colleghi. La sua progettazione è orientata
al \textbf{sociale} e all'\textbf{educativo}: ``responsabilizzare un po' le nuove
generazioni su come sta andando il mondo e su come potrebbe andare, visto che è
tutto nelle nostre mani''.

\paragraph{Il commento: siate curiosi.} ``Pane per i tuoi denti'', le risponde
Daniele, soprattutto per i LARP e i giochi pervasivi: ``vi apriremo un po' la
testa da questo punto di vista''. E aggiunge un consiglio per tutti: ``siate
curiosi'', perché imparare da altri ambiti dà moltissimo. Ognuno è lì per fare
qualcosa di preciso, ``però siate aperti anche ad altre esperienze''.

\subsection{Tancredi}

Tancredi ha 21 anni e non ha ancora grande esperienza lavorativa: qualche
``piccolo progettino'' che non è andato come voleva. Gli interessa il mondo dei
giochi, digitali e non, con un approccio \textbf{educativo} e \textbf{sociale}.
È arrivato soprattutto dai videogiochi, a cui gioca da quando era piccolo; ai
giochi da tavolo si è avvicinato solo nell'ultimo anno, grazie a una realtà
ludica, perché in famiglia ``non siamo molto da giochi''. Ha realizzato qualche
piccolo videogioco con vari engine: prima quelli senza un vero linguaggio di
programmazione, ``più a nodi'', poi qualcosa di più complicato, e ora sta
lavorando a un piccolo prototipo.

\paragraph{Il commento: niente di quello che vedrete è sprecato.} ``Non butterai
via niente di analog game design'', gli risponde Daniele, ``perché il 99 per cento
delle cose che vedremo si ritrovano quasi uno a uno nel mondo dei videogiochi''.
Tanto che nei corsi triennali dell'accademia (anche in quello di narrative
design, dove insegna Ruggero) il primo anno si parte dai \textbf{giochi da tavolo
e di ruolo}, proprio per spiegare meccaniche più complesse che servono poi per i
videogiochi. Ruggero aggiunge che vale anche il contrario: alcuni dei
\textbf{framework} di game design che si vedranno nelle lezioni tecniche vengono
dal videogioco e sono stati riapplicati all'analogico. ``Evitare la
crossmedialità qui è impossibile''. Anzi, tutte le conoscenze ``mediatiche
adiacenti'' (musica, letteratura, cinema) saranno utili.

\subsection{Lorenzo Vergata}

I Lorenzo in classe sono due; il primo a presentarsi è \textbf{Lorenzo
Vergata}, 26 anni, \textbf{dottorando al Politecnico di Milano}. Gioca ai
videogiochi da quando aveva due anni e mezzo, ma ha giocato al suo primo gioco da
tavolo ``credo nove mesi fa''. Gli fa piacere sentire che molti in classe sono
interessati ai giochi \textbf{sociali} ed \textbf{educativi}, un settore che
interessa anche a lui e che ``può essere esplorato molto'', anche con qualche
aiuto dell'intelligenza artificiale per esperienze più ricercate.

Di esperienza nell'analogico dice di avere fatto ``il minimo sindacale''. Ma
racconta che diciotto anni prima, da bambino, era ``più game designer di
adesso'': prendeva un gioco e ne cambiava le regole, anche al biliardino (``no,
ora introduciamo questa regola''). Gli piace sperimentare, creare regole e
vedere che cosa accade in sé e negli altri. Soprattutto vorrebbe ``una chiave di
lettura per poter rimanipolare la realtà'': tre settimane prima è stato a fare la
\textbf{vendemmia in Francia} e ha visto tantissime dinamiche che potrebbero
diventare un gioco da tavolo o un videogioco. Per esempio: finire contro un rovo
pieno di semi che si attaccano ai vestiti lo rallentava nella raccolta
dell'uva, mentre in un altro filare era più rapido. ``Si può fare pure una
dinamica'', e vorrebbe gli strumenti per farne giochi.

Arriva subito la battuta: una ``vendemmia gamificata'' da proporre in Oltrepò
Pavese (``io sono dell'Oltrepò Pavese, quindi devo per forza bere vino'')\footnote{La
trascrizione dice ``nell'altro papavese'' e non chiarisce chi pronunci la
battuta.}. Daniele conferma che quest'anno il corso avrà ``tante cose sul
sociale'': ``ci piace parecchio''.

\subsection{L'altro Lorenzo}

Il secondo Lorenzo (``secondo, ma non meno importante'') è uno \textbf{scrittore
di racconti brevi} e sta scrivendo il suo primo romanzo. Gli è sempre piaciuto
raccontare storie, e da più grande scriverle; gioca ai videogiochi ed è ``un
amante delle storie e della cultura pop''. Negli ultimi anni è entrato nel mondo
del \textbf{gioco di ruolo da tavolo} e vorrebbe unire le due passioni, le
storie e il gioco di ruolo, magari facendone un lavoro; gli interessa anche lo
\emph{strategic writing}. Daniele gli ricorda che ci sarà una lezione proprio
sulla \textbf{fiction}.

\subsection{Angelo}

Angelo si collega senza videocamera dal lavoro: fa il \textbf{birraio} (in un
piccolo birrificio in Brianza, a Bernareggio, come racconterà dopo la pausa) e,
con il caldo, ``vi risparmio l'orrore''. Per molto tempo era convinto che i
giochi da tavolo fossero solo i classici (Monopoly, Risiko) o i party game; poi,
quattro o cinque anni fa, ha ``sbattuto la faccia'' contro giochi davvero
appassionanti, con una specie di narrazione dietro che li rendeva ancora più
coinvolgenti, e non ha più smesso. È uno che ``ha sempre un botto di idee, ma poi
non mette in pratica assolutamente nulla'': cercava un corso che gli desse le
basi e soprattutto lo \textbf{motivasse} a provare a fare qualcosa.

\paragraph{Il commento: bisogna buttarsi.} Daniele racconta che anche l'anno
prima un suo caro amico si era iscritto proprio per questo motivo. Il corso è
breve, ma cerca di portare gli studenti ``da 0 a 100'' e di far realizzare un
\textbf{prodotto} da presentare alla fine. Il punto è che i giochi ``bisogna
farli, bisogna buttarsi''. E, venendo dalla narrazione, Daniele lo paragona allo
scrivere: si scrive e si riscrive, e allo stesso modo si fa \textbf{playtest}, si
rifà, si aggiorna. ``Non è che si fanno così''.

\subsection{Gianmarco (dopo le slide)}

Gianmarco arriva a lezione iniziata e si presenta più tardi, alla fine della
prima parte, dopo che i docenti hanno chiuso le slide organizzative e la
discussione sull'intelligenza artificiale (``Gianmarco, ti giuro, ti facciamo
presentare''). Lavora nel \textbf{team di innovazione} di un'azienda italiana.
Vuole fare giochi perché è una sua passione, che lo gratifica, e in passato ci ha
provato ``in maniera scomposta''. Cerca un \textbf{metodo} e un confronto con
persone più esperte: ``studiare solo dei manuali è molto sterile'', ``stare
chiusi nel pertugio non aiuta'', ed ``essere autodidatta ha estremamente i suoi
limiti''. Come esperienza è un giocatore e ha letto ``un po' di materiale, di
blog, di manuali'', ``niente di rilevante''.

Manca solo Giordano, che si presenterà alla lezione successiva.

\paragraph{Uno sguardo d'insieme.} La classe è molto varia: un'infermiera,
una designer della comunicazione visiva, un birraio, un dottorando, uno
scrittore, un impiegato nell'innovazione, due studenti ventunenni. Alcuni
vengono dal gioco di ruolo, altri dai giochi tradizionali, altri dai
videogiochi; quasi tutti hanno scoperto i giochi da tavolo moderni da pochi anni
o da pochi mesi. E almeno tre studenti dichiarano un interesse per i giochi
\textbf{sociali} ed \textbf{educativi}. Come si vedrà nella
sezione~\ref{a1:sec:scopo}, i docenti considerano questa varietà una risorsa.

% ---------------------------------------------------------------------
\section{I materiali e gli strumenti del corso}
\label{a1:sec:materiali}

Prima delle slide sul corso, Daniele sbriga un po' di ``clean up casalingo''.

\paragraph{Il form iniziale.} C'è un \textbf{form} da compilare ``entro e non
oltre il 15 settembre'': il QR code è nelle slide, e la segreteria lo ricorderà
su Discord. Richiede circa 10--15 minuti; si può fare subito o in giornata, ``e
ve lo togliete di torno''. È un form standardizzato per tutti i corsi (``non
siamo io e Ruggero a farvi passare la tortura del test''), diviso in due parti:
\begin{itemize}
  \item una parte di \textbf{conoscenza del candidato}: quanto tempo si ha a
        disposizione per il corso, come ci si organizza con il lavoro, che cosa
        piace giocare. È la parte che interessa ai docenti, perché li aiuta a
        organizzare le lezioni;
  \item una parte di \textbf{test logico}, che i docenti non hanno scritto (``io
        non sono responsabile del form'', ``no, no, nemmeno io'').
\end{itemize}

\paragraph{La cartella condivisa.} Tutti gli studenti vengono aggiunti alla
cartella del corso su \textbf{Google Drive}, che si trova tra i file
``condivisi con me''. Gli studenti ne sono \textbf{gestori dei contenuti} e
possono caricarci quello che vogliono. Dentro ci sono:
\begin{itemize}
  \item le \textbf{attività in classe};
  \item le \textbf{dispense};
  \item i \textbf{materiali didattici}, con le slide (quelle della lezione sono
        già lì);
  \item le cartelle dei \textbf{project work} futuri;
  \item i \textbf{task}, che ``saranno pochi, ma ci saranno'';
  \item il \textbf{calendario studenti}, ``importantissimo'', con due schede: il
        calendario delle lezioni e quello del project work.
\end{itemize}

\paragraph{Discord.} Il corso ha anche un \textbf{server Discord}, a cui tutti
sono stati invitati; ci sono i docenti (Daniele, Ruggero, Diego e altri), che si
possono taggare. Di solito è lì che la segreteria pubblica le informazioni, per
esempio se una lezione ``dovesse per caso essere cancellata''. Nelle domande
finali Ruggero aggiungerà che c'è un canale apposito per il corso, che gli
studenti possono usare liberamente per organizzarsi tra loro.

\paragraph{Le reaction di Zoom.} Durante le lezioni su Zoom si possono usare le
\textbf{reaction} (il pollice in alto, il cuoricino, la mano alzata): un modo
rapido per far capire ai docenti se si è capito, se si vuole intervenire. Nel
frattempo arrivano i primi pollici.

% ---------------------------------------------------------------------
\section{Lo scopo e il metodo del corso}
\label{a1:sec:scopo}

La prima serie di slide (che Daniele, per un lapsus, presenta come ``cos'è il
corso di pulire il pavimento di casa'') spiega che cosa il corso vuole fare e
come.

\subsection{Diventare autori analogici}

Lo \textbf{scopo del corso} è fornire le competenze per \textbf{diventare autori
analogici} e stimolare la creatività e l'\textbf{autorialità} di ciascuno: ``voi
imparerete a diventare autori di giochi analogici''. La differenza tra
\emph{game designer} e \emph{game developer} verrà chiarita più avanti (sezione~\ref{a1:sec:design-development}),
ma la parola chiave qui è \emph{autorialità}: ``vogliamo voi, vogliamo le vostre
idee''. Magari non si potrà usare l'idea del gioco della propria vita, ma si
potranno portare nelle lezioni le proprie idee e le proprie sensibilità.

Gli obiettivi elencati nella slide sono questi:
\begin{itemize}
  \item dare \textbf{maggiore consapevolezza nel design e nello sviluppo}. In
        Italia, purtroppo, ``si fa design a volte in maniera non consapevole''.
        Consapevole vuol dire avere almeno un'introduzione alle basi del design e
        dello sviluppo, e fare le cose ``con metodo, con raziocinio, con
        consapevolezza'';
  \item fornire gli \textbf{strumenti essenziali} per un design consapevole;
  \item stimolare la \textbf{sperimentazione} e la \textbf{trasversalità
        mediatica}. Non si chiede di fare ``il \emph{Root} della situazione o il
        \emph{Carcassonne} della situazione''\footnote{La trascrizione dice ``il
        lumaven della situazione o il carcassone della situazione''; il primo
        titolo non è ricostruibile con certezza. Il senso è chiaro: non si
        chiede di rifare un grande classico del gioco da tavolo.}, ma un
        project work \textbf{calibrato sulla durata del corso}, e però
        \textbf{sperimentale} e ``assolutamente transmediale'';
  \item portare \textbf{numerosi esempi} di tipologie di gioco, dall'analogico al
        tavolo, al pervasivo ``e oltre''. Lo si vedrà nella lezione sui
        \textbf{casi studio} del venerdì successivo: ``voi non avete idea di
        quante cose esistono, bellissime, sconosciutissime e favolose''. Daniele
        confessa che preparando e aggiornando le slide i colleghi gli hanno
        mostrato cose che lui stesso non sapeva nemmeno che esistessero;
  \item dare gli strumenti per \textbf{analizzare} i giochi altrui. Ruggero lo
        sottolinea: se il corso riesce nel suo intento, gli studenti sapranno
        non solo fare i propri giochi ma anche analizzare giochi già fatti, che è
        ``una delle skill importanti anche nella parte di ricerca e sviluppo''.
        Si vedranno casi studio e giochi pubblicati, sempre con l'obiettivo di
        capire \textbf{che cosa fanno e perché lo fanno}, per poi fare meglio il
        proprio lavoro;
  \item imparare a \textbf{fare giochi insieme}. Esistono gli autori che lavorano
        da soli (i \emph{solo developer}), ``ma sappiate che anche l'autore che
        lavora da solo poi si deve interfacciare con il pubblico, con un editore,
        magari con un editor, con un publisher''. ``Non lavorerete mai da soli'':
        bisogna imparare a lavorare con gli altri.
\end{itemize}

\subsection{Un corso intersezionale}

Il corso è \textbf{intersezionale}: \textbf{non ci sono lezioni dedicate a una
singola tipologia di gioco}. Non ci sarà ``la lezione sul gioco da tavolo'',
``la lezione sul gioco di ruolo'', ``la lezione sul gioco pervasivo''. Ci saranno
lezioni su temi diversi, ciascuna pensata per essere utile a \textbf{tutte e tre
le sfaccettature} del corso; alcune saranno più importanti per il gioco da
tavolo, altre per il pervasivo, altre per il gioco di ruolo, ma nessuna sarà
esclusiva. I docenti vogliono ``contaminazioni e applicazioni assolutamente
trasversali''.

\subsection{Un corso cross-generazionale}

È anche un corso \textbf{cross-generazionale}: i gradi di esperienza di gioco
quest'anno sono abbastanza allineati, ma le esperienze di vita, le età, le
estrazioni sociali sono varie. ``Ed è un pro, non è un contro. Più diversificata
è la classe, più imparate meglio anche voi''. La slide ha una piccola nota per
fasce d'età:
\begin{description}
  \item[Under 30] Concentratevi sulla vostra \textbf{crescita}. Evitate di
        confrontarvi con le persone più grandi che magari ``ce la fanno meglio''
        in alcune cose.
  \item[Over 30] Concentratevi sul \textbf{consolidamento}, ma siate
        \textbf{umili}: i docenti diranno cose che faranno ``storcere il naso''.
        ``Bisogna un po' reimparare per imparare'': bisogna sapersi adattare.
\end{description}
Ruggero contestualizza: lo dicono perché in passato hanno visto studenti che,
essendo più adulti o avendo giocato di più, pensavano che l'esperienza
\emph{di gioco} si trasformasse automaticamente in competenza \emph{di
creazione} di giochi. Non è così. E non bisogna sottovalutare chi gioca meno:
proprio perché gioca cose diverse, per risolvere un problema di design specifico
``è una ricchezza incredibile avere qualcun altro che riesce a uscire dalla
scatola''.

\subsection{Un corso pratico e adattabile}

Il corso è \textbf{pratico}: gli studenti si concentreranno su un progetto di
gruppo, il \textbf{project work}, lavorando in team. Ed è \textbf{adattabile}: il
programma è definito, ma ``si aggiorna sui vostri stimoli, priorità e
necessità''. Se durante il corso i docenti vedono che serve più una cosa
dell'altra, sono ``prontissimi'' a dare più spazio a un argomento invece che a
un altro.

\subsection{Inquiry-based learning: domandate}

Il metodo è quello dell'\textbf{inquiry-based learning}, l'apprendimento basato
sulle domande. Di solito le persone dicono ``ho una domanda, magari sarà
stupida''. ``Non esistono domande stupide. Pensate che ci sia una domanda
stupida? Fatela, alzate la mano''. Vanno bene anche le
\textbf{provocazioni}: se qualcosa non torna, è meglio chiedere. La condivisione
è fondamentale.

I docenti danno \textbf{metodi, dati e conoscenze}, ma lasciano trovare le
risposte agli studenti: ``vi lasciamo un pochettino annegare nel mare prima di
darvi la soluzione, in modo tale che la possiate trovare da soli''. Di
conseguenza \textbf{non ci saranno interrogazioni sulle definizioni}: interessa
vedere \textbf{come} gli studenti applicano ciò che imparano. Se una
definizione non viene in mente, non è un problema; l'importante è saperla usare.

\subsection{Feedback, non valutazione}

``Sì al feedback, no alla valutazione''. I docenti danno \textbf{feedback
verbali}, anche sulle attività da fare a casa (i task, che non saranno tanti).
Il corso è calibrato per \textbf{studenti lavoratori}: c'è un po' da lavorare
fuori dalle lezioni, ma in modo calibrato.

Il principio è \textbf{fallire oggi per imparare per il futuro}: \emph{fail
forward}, \emph{fail to learn}. I docenti lasciano provare e fallire almeno una
volta prima di dare la soluzione: ``il classico metodo Montessori, quello
originale, non quello travisato''.

\subsection{Le soft skill e i criteri di valutazione}

Grande attenzione va alle \textbf{soft skill}: i lavori di gruppo, lo scambio
propositivo, l'atteggiamento positivo, la propositività, la proattività.
``Lavoro di gruppo'' non significa stare ogni giorno in sala studio con gli
altri: va benissimo che un gruppo si organizzi dandosi priorità e tempi in base
alla disponibilità di ciascuno. L'importante è \textbf{collaborare}, essere
\textbf{professionali}, rispettarsi a vicenda, rispettare le idee proprie e
altrui e lasciare spazio agli altri, perché i progetti ``non saranno solo
vostri, saranno di tutte le persone del vostro gruppo''.

La valutazione del project work non riguarda solo la \textbf{documentazione}
prodotta, ma anche:
\begin{itemize}
  \item la \textbf{presentazione} finale;
  \item l'\textbf{autovalutazione};
  \item la \textbf{crescita personale};
  \item l'\textbf{atteggiamento} tenuto nel corso.
\end{itemize}
Daniele cita: ``da ognuno secondo le sue capacità, a ognuno secondo i suoi
bisogni''. Lavorare in gruppo non vuol dire che, in tre, ciascuno lavori
esattamente il 33\% del tempo. Ai docenti serve piuttosto capire \textbf{chi fa
che cosa} nel gruppo, e che il processo interno sia \textbf{trasparente}.

\paragraph{Lavorare in cloud.} Per lo stesso motivo il consiglio è di
\textbf{non lavorare in locale}, ma direttamente nelle cartelle condivise del
project work e dei task. Si è liberi di fare altrimenti, ma lavorando in cloud
si dà ai docenti la possibilità di vedere \emph{come} si lavora, il proprio
``processo mentale''.

% ---------------------------------------------------------------------
\section{Le regole di comportamento e il carico di lavoro}
\label{a1:sec:regole}

Alcune slide sono di quelle che si ripetono in ogni corso (``\emph{repetita
iuvant}'', storpiato per scherzo in ``ripetita Juventus'').

\paragraph{Comportamento.} Prestare attenzione durante le lezioni; interagire;
portare rispetto; essere proattivi; alzare la mano e fare domande.
\textbf{Tolleranza zero} per pregiudizi, maleducazione, mancanza di rispetto,
\emph{cheat} e scorciatoie. Daniele ammette che ``cheat e scorciatoie'' viene da
altri corsi dove il tema è più importante, ma Ruggero lo ferma: ``non
sottovalutate questa cosa'', perché implica anche \textbf{non far fare a ChatGPT
i task}. Daniele ha una slide apposta, ``perché è un discorso molto complesso nel
2026'' (sezione~\ref{a1:sec:ia}).

\paragraph{Le lezioni sono intense.} Le lezioni possono essere \textbf{dense},
``soprattutto le mie, sappiatelo: io sono logorroico'', avverte Daniele. Bisogna
chiedere spiegazioni se serve, prendere appunti, rivedere le registrazioni. Le
lezioni sono ``delle scintille'': modi per far esplorare i temi oltre le slide.
``Andate oltre, se riuscite''.

\paragraph{Il carico di lavoro.} Le regole sui task sono queste:
\begin{itemize}
  \item i task \textbf{sono per voi, non per noi};
  \item se si ha poco tempo, è meglio farli rapidamente piuttosto che non farli:
        ``provateci'';
  \item i docenti cercano di non chiedere più di \textbf{un'ora o due a
        settimana} oltre alle ore di lezione;
  \item non tutte le lezioni avranno un task, e il \textbf{project work} è il più
        importante di tutti.
\end{itemize}

\paragraph{I feedback dei docenti.} I feedback ``si guadagnano'' man mano che si
reagisce bene: non bisogna aspettarsi ``dei papiri'' sui task. Se i docenti
hanno voglia di dare feedback più ampi lo fanno, soprattutto se vedono che lo
studente se li merita. Ad Analog di solito non ci sono task \emph{rejected}; ma
se un docente rifiuta un task, vuol dire che si è fatto \textbf{meno del livello
minimo} richiesto dal corso. I feedback possono non essere immediati: ``siate un
pochettino impietositi dal fatto che siamo dei docenti''. E se un feedback non
arriva, o se ne vogliono di più, bisogna scrivere, ``perché siamo tre esseri
umani e quindi possiamo anche dimenticarci''.

% ---------------------------------------------------------------------
\section{Il project work e il calendario}
\label{a1:sec:project}

\subsection{Come funziona il project work}

Una prima slide con le date del project work è da aggiornare (``perdonatemi il
refuso''); Daniele la corregge al volo e rimanda alla slide successiva, più
precisa. Il funzionamento è questo:
\begin{itemize}
  \item i \textbf{gruppi li formano i docenti}. Gli studenti sono nove, quindi si
        faranno \textbf{tre gruppi da tre}: una dimensione pensata anche per il
        fatto che molti sono studenti lavoratori, così da non ritrovarsi ``in due,
        a esplodere'';
  \item ogni gruppo deve prima preparare un \textbf{pitch} del proprio gioco, da
        presentare in anticipo: un'idea, una bozza del gioco;
  \item poi si lavora al gioco \textbf{nelle pause del corso}: da qui a dicembre
        ci saranno due pause, quindi non ci sarà lezione tutti i lunedì e tutti i
        venerdì;
  \item ci saranno \textbf{review intermedie} con i \textbf{tutoraggi}, per lo più
        il \textbf{sabato mattina}, dalle 11 alle 12, con Ruggero su Google Meet.
        Servono ad aiutare i gruppi a portare avanti il project work.
\end{itemize}

\subsection{Il calendario delle lezioni}

Il \textbf{calendario} è nel calendario studenti della cartella condivisa.
Attenzione: le date \textbf{possono variare}. ``Il 90\% rimarranno così'', ma
può capitare che, per motivi di lavoro, qualcuno stia male o non possa, o che
``scompaia l'elettricità da tutte le città d'Italia''. E l'ordine delle lezioni
può cambiare se i docenti vedono che agli studenti serve prima una lezione
invece di un'altra.

Ci sono \textbf{due pause importantissime}:
\begin{itemize}
  \item dal \textbf{13 ottobre al 5 novembre}: è il periodo delle grandi fiere
        (\textbf{Essen} e \textbf{Lucca}\footnote{La trascrizione dice ``Essen,
        Luca, Cose'': si tratta con ogni probabilità della fiera Spiel di Essen,
        in Germania, e di Lucca Comics \& Games, che si tengono entrambe in
        ottobre.}), ``un periodo un po' frenetico'', e cominciano anche i corsi
        accademici, quindi la pausa permette ai docenti di respirare. Per gli
        studenti è la pausa più importante: quasi un mese in cui sviluppare la
        parte importante del gioco;
  \item dal \textbf{24 novembre all'11 dicembre}, dopo qualche lezione a novembre.
\end{itemize}

E ci sono \textbf{due date da evidenziare}:
\begin{itemize}
  \item il \textbf{25 settembre}: nei primi 5--10 minuti della lezione (che per il
        resto sarà tenuta da Beatrice) i docenti annunceranno i tre gruppi;
  \item il \textbf{12 ottobre}: il \textbf{pitch} dei giochi, con l'idea
        embrionale. Tra il 25 settembre e il 12 ottobre ogni gruppo prepara un
        breve pitch in formato slide; quel giorno i docenti lo rivedono e danno
        feedback: se il gioco va bene, se non va bene, come svilupparlo.
\end{itemize}
Dal 12 ottobre fino all'inizio di dicembre si sviluppa il gioco, con i
tutoraggi: ``sfruttate questi momenti di pausa''.

\subsection{Il calendario del project work}

La slide successiva, ``reale e aggiornata'', riporta il calendario del project
work (tabella~\ref{a1:tab:calendario} e figura~\ref{a1:fig:calendario}).
``Per ora non pensate al project work: ci penserete dal 25 settembre in poi''.

\begin{table}[htbp]
\centering
\begin{tabular}{@{}lp{9.2cm}@{}}
\toprule
\textbf{Data} & \textbf{Che cosa succede} \\
\midrule
ven 25 settembre\footnotemark & lezione: annuncio dei gruppi e del
\emph{mandate}, cioè di che cosa si dovrà fare \\
sab 3 ottobre, ore 11 & prima review e tutoraggio con Ruggero (dopo una
decina di giorni di lavoro) \\
lun 12 ottobre & pitch del gioco durante la lezione \\
13 ottobre -- 5 novembre & prima pausa: sviluppo del gioco \\
sab 7 novembre & tutoraggio \\
sab 14 novembre & tutoraggio \\
24 novembre -- 11 dicembre & seconda pausa: sviluppo del gioco \\
ven 4 dicembre, entro mezzanotte & consegna del gioco sul Drive (non c'è
lezione) \\
ven 11 dicembre & i giochi vengono visti a lezione \\
14 dicembre & giornata finale, con la tavola rotonda con gli editori \\
\bottomrule
\end{tabular}
\caption{Il calendario del project work come lo presenta Daniele. Le date
possono variare.}
\label{a1:tab:calendario}
\end{table}
\footnotetext{Daniele dice ``lunedì 25 settembre'', ma nel 2026 il 25 settembre
è un venerdì; il giorno della settimana nella tabella è corretto di
conseguenza.}

\paragraph{Perché la consegna è il 4 dicembre.} Il 4 dicembre non c'è lezione:
si deve solo consegnare il gioco, entro mezzanotte, sul Drive. La presentazione
a lezione è l'11 dicembre, ma i tre docenti hanno bisogno di \textbf{una
settimana} per rivedere il gioco in privato e valutare anche la
\textbf{documentazione}; per questo la consegna non coincide con la
presentazione. Che cosa voglia dire esattamente ``consegnare il gioco'' verrà
spiegato più avanti.

\paragraph{La giornata finale.} Il 14 dicembre c'è la \textbf{giornata finale},
con una \textbf{tavola rotonda con gli editori}: di solito viene invitato Flavio,
che è anche docente di game production, più qualche altro editore, ``vediamo
quest'anno chi riusciamo a portarvi''.

\begin{figure}[htbp]
\centering
\begin{tikzpicture}[
    font=\footnotesize,
    x=0.1cm,
    evento/.style={draw, thick, rounded corners=2pt, fill=white,
                   text width=1.8cm, align=center, inner sep=2pt},
    pausa/.style={fill=gray!20, draw=gray!60},
    tacca/.style={thick}
  ]
  % asse: giorno 0 = 25 settembre; 1 unità = 1 giorno
  \draw[pausa] (18,-0.35) rectangle (41,0.35);
  \draw[pausa] (60,-0.35) rectangle (77,0.35);
  \node[font=\scriptsize\itshape] at (29.5,0.18) {pausa};
  \node[font=\scriptsize\itshape] at (65,0.18) {pausa};
  \draw[-{Stealth[length=2.5mm]}, thick] (-3,0) -- (85,0);
  % eventi sopra
  \draw[tacca] (0,-0.15) -- (0,0.9);
  \node[evento, above] at (0,0.9) {25 set\\ gruppi e \emph{mandate}};
  \draw[tacca] (17,-0.15) -- (17,2.2);
  \node[evento, above, fill=gray!8] at (17,2.2) {12 ott\\ \textbf{pitch}};
  \draw[tacca] (70,-0.15) -- (70,0.9);
  \node[evento, above, fill=gray!8] at (70,0.9) {4 dic\\ \textbf{consegna}};
  % eventi sotto
  \draw[tacca] (8,0.15) -- (8,-0.9);
  \node[evento, below] at (8,-0.9) {3 ott\\ tutoraggio};
  \draw[tacca] (43,0.15) -- (43,-0.7);
  \draw[tacca] (50,0.15) -- (50,-0.7);
  \draw (43,-0.7) -- (50,-0.7);
  \draw[tacca] (46.5,-0.7) -- (46.5,-0.9);
  \node[evento, below] at (46.5,-0.9) {7 e 14 nov\\ tutoraggi};
  \draw[tacca] (77,0.15) -- (77,-0.9);
  \node[evento, below] at (77,-0.9) {11 dic\\ giochi a lezione};
  \draw[tacca] (80,-0.15) -- (80,2.2);
  \node[evento, above] at (80,2.2) {14 dic\\ giornata finale};
\end{tikzpicture}
\caption{Una linea del tempo del project work, ricostruita dalle slide (scala
approssimativa: dal 25 settembre al 14 dicembre). Le bande grigie sono le due
pause senza lezioni, in cui i gruppi sviluppano il gioco.}
\label{a1:fig:calendario}
\end{figure}

Finita questa parte, Daniele apre alle domande: ``Domande sul calendario? Sulle
slide che ho fatto? Considerazioni? Non abbiate timore, io non ci credo''. E la
domanda arriva, sull'unico tema che aveva lasciato da parte.

% ---------------------------------------------------------------------
\section{L'intelligenza artificiale nel corso}
\label{a1:sec:ia}

\subsection{La slide di Daniele}

Daniele aveva già presentato una slide che definisce ``importantissima''. ``Lo so
che è paradossale, me ne rendo conto. Lo so che magari mi darete del luddista,
dell'anti-tecnologia: non è così. Uso l'intelligenza artificiale per lavoro. Ma
in questo corso, vi prego, \textbf{non usate l'intelligenza artificiale}''. La
richiesta vale per le attività, per i task e \textbf{anche per la
prototipazione}. Le ragioni:
\begin{itemize}
  \item ``voglio che siate voi a fare il gioco, voglio che siate voi a imparare a
        farlo''. Daniele parla da utilizzatore: prima di ChatGPT, che ha smesso
        di usare, oggi soprattutto di Claude, di cui ``si sottovaluta quanto
        possa essere incerto, insicuro, impreciso, non all'altezza'';
  \item nella game industry, \textbf{i publisher rifiutano automaticamente} i
        pitch e i prototipi fatti con l'intelligenza artificiale. Daniele lo ha
        visto di persona in ``tantissime pitch review'' quell'estate, anche al
        \textbf{First Playable} di Firenze\footnote{La trascrizione dice ``First
        Playboy a Firenze'': si tratta del First Playable, l'evento di business
        dell'industria videoludica che si tiene a Firenze.}: anche nel mondo dei
        videogiochi i publisher che vedono pitch o prototipi fatti con l'IA
        ``rigettano automaticamente la build'', perché hanno ancora timore
        dell'IA e delle sue possibili conseguenze;
  \item ``non so come sarà il futuro'': non sa se bisognerà abituarsi a usarla
        tutto il giorno. Non è ``totalmente contrario'', purché sia usata ``con
        raziocinio''; secondo lui oggi, negli Stati Uniti, non lo è;
  \item ma in un corso che si chiama \emph{Analog Game Design} gli studenti devono
        saper fare i giochi da soli. In futuro, se vorranno farsi aiutare
        dall'IA per essere più efficienti, saranno liberissimi di farlo; durante
        il corso no, perché i docenti vogliono valutare \textbf{le loro
        competenze e capacità}.
\end{itemize}

\subsection{La domanda: che cosa vuol dire usare l'IA in un gioco analogico?}

La domanda di uno studente è precisa: che cosa si intende per ``task fatti con
l'IA''? Qual è l'uso dell'IA in un gioco analogico?

\paragraph{Ruggero: l'esperienza di chi valuta i prototipi.} Ruggero risponde
dalla sua esperienza nel settore, dove riceve prototipi da valutare per
un'eventuale pubblicazione. Distingue due usi:
\begin{enumerate}
  \item le \textbf{immagini generate}: illustrazioni, graphic design e così via.
        Per molti publisher è ``un \emph{reject} immediato, proprio immediato'';
  \item l'IA usata per il \textbf{regolamento}, per scrivere le regole e le
        meccaniche. È più difficile da individuare, ma molti publisher se ne
        accorgono comunque, perché le IA ``non sono ancora brave abbastanza da
        creare un regolamento coeso meccanicamente e scritto con un senso per il
        vostro gioco''. Quello che fanno è prendere parti di altri regolamenti,
        elaborare le strutture che vedono nei regolamenti pubblicati e
        riapplicarle. Ma molti giochi si prestano poco a un regolamento
        strutturato come quello di altri, perché \textbf{ogni gioco è un
        ecosistema meccanico a sé}. Il risultato sono regolamenti con
        riferimenti non immediati, e soprattutto con \textbf{la stessa meccanica
        descritta in modo diverso in pagine diverse}: due difetti che fanno
        perdere tempo all'editore, e in questo caso ai docenti.
\end{enumerate}
Poi aggiunge un punto ``personale, soggettivo'': i docenti sono interessati a
vedere il lavoro, la creatività, ``il frutto della vostra creazione''. E
viceversa gli studenti vogliono i commenti dei docenti: ``immaginatevi se voi
fate un gioco e noi vi diamo i commenti scritti con ChatGPT''. Altrimenti, dice
ridendo, gli studenti fanno il gioco con ChatGPT, i docenti lo danno a ChatGPT
per i commenti, ``e facciamo molto alla svelta''.

\paragraph{Daniele: perché gli studenti usano l'IA.} Daniele racconta la sua
esperienza con gli studenti, che secondo lui ricorrono all'IA in due momenti:
\begin{enumerate}
  \item quando vanno in \textbf{crisi}: in ansia da prestazione o per i tempi.
        Cercano di risolvere il problema della propria disorganizzazione
        facendolo risolvere all'IA, senza accorgersi che l'output è ``molto,
        molto scarso'';
  \item quando hanno l'impressione che l'intelligenza dell'IA sia
        \textbf{migliore} di quello che scrivono loro, o di quello che c'è in
        giro. Per Daniele è una ``mancanza di palato critico'': l'output dell'IA
        sembra buono in superficie, ma quando si va nel dettaglio e ci si fa
        delle domande ``di solito fa un po' di \emph{crumbling}'', si sbriciola.
\end{enumerate}
Lui usa Claude ogni giorno e ne conosce i limiti: ``non c'è nemmeno una frase
che io copio direttamente''. Gli fa fare il lavoro più noioso, più ripetitivo e
soprattutto \textbf{strutturato}, ``replicare delle cose come se fosse una
macro''. Lo usa per il testo e per il codice, mai per le immagini: per motivi di
\textbf{diritto d'autore}, e perché ha molti amici artisti ed è stufo di vedere
il loro lavoro ``dilaniato'' dall'IA.

Ammette che i docenti non possono controllare gli studenti; per questo
``ve lo consiglio''. E riprende quella che definisce la cosa più bella detta da
Ruggero (``ti ruberò la frase per il corso di narrative''): i docenti i feedback
li danno di persona. Daniele a volte manda 30 minuti di audio a uno studente, o
fa intere call per rivedere script, pitch e giochi: ``è il mio lavoro, mi piace
farlo e voglio farlo io, quindi mi aspetto che lo studente lo faccia lui''. Come
freelance, poi, è ``stufo marcio'' dei clienti che arrivano con qualcosa fatto
con l'IA (``guarda che bello, pubblica questo''), perché il più delle volte il
risultato è pessimo: non sanno come usarla. ``È anche una questione di
competenza''.

\paragraph{Il prototipo non deve essere bello.} Questo vale anche per il
\textbf{prototipo}, su cui si tornerà nelle lezioni di prototipazione: in un
prototipo non servono ``graficoni allucinanti''. Le grafiche bellissime fatte con
l'IA spesso \textbf{peccano di user experience}: sembrano bellissime, ma quando
il giocatore deve leggere la carta non capisce qual è il costo, qual è il
potere, qual è l'effetto.

\paragraph{Beatrice: si perde l'occasione di imparare.} Beatrice aggiunge due
motivi che valgono ``al di là dei nostri credo''. Il primo: chi usa l'IA in
sostituzione (per esempio per scrivere un testo, per accelerare i tempi o per
trovare strategie) \textbf{perde la parte in cui può imparare}. Nel game design
bisogna ``tenere insieme i pezzi'' di un regolamento, ma anche dell'esperienza
dell'utente e di come trasmetterla; i passaggi necessari (creare il regolamento,
creare le meccaniche, sperimentare, raccogliere i feedback) ``si capiscono solo
vivendoli''. E si sta pagando un corso per imparare. Il secondo, sul piano
editoriale: chi presenta a un editore un prodotto che ha come base l'IA rischia
da un lato di non avere imparato abbastanza (e di ritrovarsi ``a piedi'' sul
campo senza IA), dall'altro che l'editore se ne accorga, ed è un tema
\textbf{reputazionale}. Ci sono editori ``più illuminati di altri'', quindi
anche qui c'è un mercato da capire; ma i docenti preferiscono dire che cosa è
più sensato per il mercato come lo conoscono.

\paragraph{Ruggero: non abbiate paura del giudizio.} Ruggero vuole dare uno
``spotlight'' a un punto detto da Daniele e rimasto ``semi-ignorato'': a volte
l'uso dell'IA nasce dalla \textbf{paura del giudizio}, dal voler consegnare
qualcosa di più pulito, più strutturato possibile. ``Vi giuro, non ve ne
preoccupate''. Ruggero detesta i voti e non ragiona per voti; non valuterà
nessuno per grafiche fatte a matita. Insegna da anni, ha avuto studenti con
dislessia e discalculia: refusi ed errori non gli interessano, ``non è un corso
di editing o di proofreading''. Gli interessa che le \textbf{meccaniche siano
coerenti con lo scopo di design}. Immagini fatte a matita o con Paint ``non sono
un problema, anzi, le preferisco'' rispetto all'``ennesima immagine fatta con
ChatGPT o Midjourney''.

\paragraph{Daniele: il pensiero critico, costo e qualità.} L'ultimo intervento
è di Daniele, sul \textbf{pensiero critico}. Ragionando per assurdo (``un
pensiero di Belzebù''), immagina che un domani l'IA venga usata ovunque. I suoi
limiti, secondo lui, resterebbero due:
\begin{description}
  \item[Il costo] Tutti dicono che sta diminuendo e che diminuirà, ma non se ne
        può essere sicuri al 100\%.
  \item[La qualità] Cioè la capacità di leggere l'output. Se si dà un regolamento
        a Claude o a ChatGPT chiedendo quali problemi ha, il modello risponde con
        i problemi che, in percentuale, è più probabile che gli esseri umani
        abbiano già individuato nei dati su cui è stato addestrato. Ne derivano
        due problemi: se il problema del proprio gioco non è mai stato
        individuato prima, il modello non lo vedrà; e viceversa può segnalare
        problemi che il gioco non ha.
\end{description}
Fa l'esempio di Grammarly, che usa nella versione premium per rivedere i testi
in inglese: spesso non accetta i suoi consigli, perché se vuole scrivere un
dialogo di videogioco in \textbf{dialetto Cockney} londinese, lo deve scrivere
in Cockney. Se Grammarly dice che è sbagliato, ``sei tu che non hai capito il
motivo per cui voglio scrivere questo testo in Cockney''. La conclusione: prima
di usare un'IA per un'attività bisogna sapere \textbf{che cosa si vuole
ottenere}, qual è l'obiettivo, e quali \textbf{competenze} servono per valutare
in modo critico se l'output di un qualunque LLM è davvero valido. Daniele non
vuole dilungarsi oltre: è un argomento ``molto caldo'' per lui, che è anche
scrittore e si sente dire da tutti ``ti sostituiranno''.

Subito dopo si presenta finalmente Gianmarco (lo si è visto nella
sezione~\ref{a1:sec:partecipanti}), e poi si fa una pausa di una decina di
minuti. Daniele avverte che nella terza parte lui non ci sarà (alle sette e un
quarto ha un impegno di lavoro), e lo stesso vale per Beatrice: la parte finale
di domande e risposte sarà con il solo Ruggero, ``in ottime mani''.

% ---------------------------------------------------------------------
\section{Le figure professionali dell'industria analogica}
\label{a1:sec:figure}

Dopo la pausa (in cui si scopre che Angelo lavora in un piccolo birrificio in
Brianza e si scherza sul rischio di un \emph{drinking game}, tra la birra e la
vendemmia) la lezione riprende con le slide di Ruggero. ``Oggi è una lezione
introduttiva'': gli argomenti si vedranno ``molto dall'alto'' e verranno
approfonditi nelle lezioni successive. Ma servono già ora, per creare
\textbf{l'inizio di un linguaggio condiviso}: ci sono termini che torneranno
spesso, a cominciare da quelli del mercato e dell'industria del gioco
analogico, un'industria ``molto varia, molto caotica, anche molto mutevole'',
con tanti tipi diversi di giochi e altrettanti tipi diversi di entità che ci
lavorano.

\subsection{Una decina di etichette}

I docenti hanno raccolto ``una decina di etichette'', cioè di \textbf{ruoli} che
nell'industria sono o possono essere presenti. Due avvertenze:
\begin{itemize}
  \item nelle aziende più piccole \textbf{la stessa persona copre spesso più
        ruoli}, per forza di cose. Sarà anche il caso degli studenti, che
        lavoreranno in gruppi di tre e non di venti: è facile che ciascuno finisca
        a fare più cose. I docenti ne terranno conto, e qui entreranno in gioco
        anche competenze diverse dal game design (programmi di grafica,
        illustrazione, impaginazione, layout, \emph{proofreading}\ldots);
  \item Daniele precisa che \textbf{non si chiede agli studenti di diventare}
        illustratori o altro. Ruggero concorda, ma aggiunge che è utile sapere
        anche \emph{che cosa non è il proprio lavoro}. Se in un gruppo nessuno sa
        illustrare, ai docenti va benissimo un prototipo che mostri dove vanno
        le immagini: \textbf{placeholder}, concept, anche semplici box di testo
        con scritto ``qui vorremmo un'immagine di\ldots''. ``È importante anche
        sapere che cosa sapete fare e che cosa non sapete fare''.
\end{itemize}
La tabella~\ref{a1:tab:figure} riassume le figure; le sottosezioni seguenti le
riprendono una per una, con i commenti dei docenti.

\begin{table}[htbp]
\centering
\small
\begin{tabular}{@{}p{3.3cm}p{9.6cm}@{}}
\toprule
\textbf{Figura} & \textbf{Di che cosa si occupa} \\
\midrule
Game designer & crea le regole alla base dell'esperienza: le meccaniche, come
interagiscono tra loro e come chi gioca interagisce con esse \\
Game developer & fa da ponte tra chi fa game design e il publisher: bilancia il
prototipo, cura i modelli matematici, adatta il gioco (e il paragame) alle
necessità del publisher \\
Editor (proofreader) & revisiona e adatta i testi a partire dalla prima stesura
del designer (correzione, linguaggio inclusivo\ldots) \\
Publisher (editore) & pubblica e produce fisicamente il gioco; il suo marchio è
sulla scatola o sulla copertina \\
Sales e marketing & il lato commerciale: vendite (anche B2B) e comunicazione del
prodotto \\
Graphic designer & icone meccaniche, wireframe, UX: il lato funzionale degli
elementi visivi \\
Illustratore & le illustrazioni: il \emph{look and feel}, il world building, il
``vibe'' del gioco \\
Sensitivity reader (cultural consultant) & verifica l'accuratezza e la
rappresentazione di un tema sensibile che il team non conosce \\
Producer & coordina il progetto (quasi project management), parla con le
tipografie, verifica che i materiali siano sostenibili economicamente \\
\bottomrule
\end{tabular}
\caption{Le figure professionali dell'industria del gioco analogico presentate
nella slide di Ruggero, in sintesi.}
\label{a1:tab:figure}
\end{table}

\subsection{Game designer}

Il \textbf{game designer} (o la game designer) è chi si occupa, ``sembra banale
a dirlo'', di \textbf{creare le regole alla base dell'esperienza} che si vuole
far arrivare. Vale per tutti i tipi di giochi, non solo per quelli analogici, e
che cosa faccia concretamente in un progetto o in una serie specifica ``può
mutare tantissimo''. Nell'analogico non verrà chiesto di programmare o di
prototipare in digitale; verrà chiesto di \textbf{creare le meccaniche di
gioco}, di capire \textbf{come le meccaniche si interfacciano tra loro} e
\textbf{come le persone che giocano si interfacceranno con le meccaniche}.

\subsection{Game developer}

Il \textbf{game developer} (o la game developer) sarà l'argomento delle due
lezioni di Marta. La parola \emph{development} è nota soprattutto
dall'ambito videoludico, ma nell'analogico ha un significato molto specifico
``che non c'entra niente con la programmazione'': il game developer è
\textbf{un ponte tra chi fa game design e il publisher}. Di solito si occupa di
\textbf{bilanciare il prototipo}, di curare i \textbf{modelli matematici}, di
adattare il gioco e il \textbf{paragame}\footnote{La trascrizione dice ``il
Wix, il paragame''. Il \emph{paragame} è, in termini generali, tutto ciò che
accompagna il gioco senza esserne il sistema di regole (regolamento, scatola,
materiali, comunicazione); i docenti lo tratteranno in una lezione dedicata, e
qui il termine è solo citato.} alle necessità del publisher. La distinzione
viene ripresa e approfondita nella sezione~\ref{a1:sec:design-development}.

\subsection{Editor}

L'\textbf{editor}, a volte chiamato \textbf{proofreader} a seconda dell'azienda,
si occupa della revisione dei \textbf{testi}: non della stesura iniziale, perché
di norma, data la scala dell'industria, è chi fa game design a scrivere un
\textbf{primo draft}, e l'editor ci passa sopra. Non solo per correggere i
refusi: spesso il suo lavoro è anche \textbf{adattare} i testi, per esempio
renderli \textbf{gender neutral}. Se ne vedranno casi studio più avanti.

\subsection{Publisher}

Il \textbf{publisher}, o \textbf{editore}, è l'entità che nell'ambito
commerciale \textbf{pubblica} giochi, giochi di ruolo o giochi da tavolo. È una
figura ``più solidificata'' nel gioco da tavolo e nel gioco di ruolo da tavolo;
\textbf{nell'ambito pervasivo non esiste proprio in questa forma}, come
mostrerà Beatrice parlando dei casi studio. Se si compra un gioco da tavolo o un
gioco di ruolo in un negozio, il marchio del publisher è di solito sul retro
della copertina o sul fronte della scatola: è l'entità che, di solito,
\textbf{produce fisicamente} il gioco.

\subsection{Sales e marketing}

L'ambito \textbf{sales e marketing}, a volte affidato a un'unica persona e a
volte separato, riguarda il lato commerciale. In alcune aziende è interno al
publisher, e si occupa di parlare con le fabbriche e le tipografie, di ``far
quadrare i conti''. In altre realtà è completamente diverso, specialmente nel
\textbf{B2B} o nei giochi \textbf{didattici} e \textbf{formativi}, di cui
Daniele e Beatrice ``sapranno parlare meglio''. Il tema viene ripreso più avanti
in dettaglio (sezione~\ref{a1:sec:marketing}).

\subsection{Graphic designer e illustratore}

Ruggero presenta insieme \textbf{graphic designer} e \textbf{illustratore}
proprio perché sono ruoli molto diversi che si tende a far coincidere:
\begin{itemize}
  \item chi fa \textbf{graphic design} si occupa delle \textbf{icone meccaniche},
        del \textbf{wireframing}, della \textbf{UX}: del lato \textbf{meccanico e
        funzionale} degli elementi visivi, cioè di tutto ciò che serve
        meccanicamente nel gioco, o anche alla sua ``qualità della vita'';
  \item chi fa \textbf{illustrazione} si occupa dell'\emph{art}, delle
        illustrazioni: copre il \textbf{look and feel}, la parte di
        \textbf{world building} e il ``vibe'' che si dà al gioco.
\end{itemize}

\paragraph{Daniele: che cos'è la user experience.} Daniele interviene sulla
\textbf{user experience}, ``importantissima per me'': ha un background da
\textbf{layout artist} (impaginava magazine), un mestiere in cui la UX conta
moltissimo. Per user experience non si intende solo l'\textbf{interfaccia
grafica}, la UI, che anche nei videogiochi è disegnata da un artista. La UX è
qualcosa che può fare anche il game designer: \textbf{disegnare l'esperienza del
giocatore}, cioè ``il gesto, l'intenzione, l'intuito'' che stanno dietro
l'aspetto visivo dell'interfaccia. Per fare UX non serve saper disegnare.
L'esempio è quello della \textbf{porta}: se si progetta una porta, bisogna
indicare all'utente come si apre, come si prende la maniglia; ``quella è user
experience''. Il design della maniglia in sé, invece, è
\textbf{interfaccia}.

\paragraph{Un aneddoto sulle icone.} Ruggero aggiunge che molte persone che
fanno \textbf{graphic design ben fatto non sanno illustrare}, e spesso vale il
viceversa. Lo si nota quando le icone di un gioco sono bellissime ma non
funzionali: andando a guardare i crediti, di solito si scopre che le icone le ha
fatte la persona che ha fatto le illustrazioni. ``Comunque non andiamo adesso in
questo, perché è un altro vaso di Pandora''.

\subsection{Sensitivity reader}

Il \textbf{sensitivity reader}, o \textbf{cultural consultant} a seconda del
progetto, fa un lavoro di \emph{sensitivity reading} del gioco. È una figura
molto usata nell'ambito \textbf{formativo e didattico}. L'esempio di Ruggero: se
si fa un gioco per un ospedale su una malattia specifica, chi lo progetta non è
un medico, ma il gioco deve essere \textbf{accurato} e \textbf{rappresentativo}
delle dinamiche che quella malattia comporta. Serve una persona che conosca quel
tema, anche se non sa che cosa sia il game design. Sono figure chiamate di
solito \textbf{a progetto}, per giochi specifici, che compensano competenze che
il team non ha ``per forza di cose, perché noi non sappiamo tutto''; e a volte,
per lavoro, si devono fare giochi su temi che non si conoscono personalmente.

\paragraph{Daniele: non è un censore.} Daniele segnala un problema italiano:
alcune persone, anche online, hanno confuso il sensitivity reader con un
\textbf{censore}, ``una persona che fa censura, soprattutto in materia di
\emph{woke}''. ``Assolutamente non è così''. La parola \emph{sensitive} non
significa che il pubblico sia sensibile a un tema, ma che \textbf{l'argomento di
cui tratta il gioco è sensibile}. In un gioco sulla medicina, realistico o no,
non si possono inserire imprecisioni. Per questo esistono i sensitivity reader
``più disparati'': da quelli per la cultura indiana o afroamericana fino a quello
su ``come si allevano gli gnu in Africa''. È una persona competente su un
argomento che aiuta a essere \textbf{veritieri}. Certo, uno degli ambiti
principali in cui si usano è quello culturale, ma ``non fate l'errore di pensare
che siano dei censori''. Daniele ha lavorato con ``centinaia'' di sensitivity
reader e nessuno gli ha mai detto ``non puoi scrivere questa cosa''. Di solito
sono come degli \textbf{editor}: danno opzioni, mostrano i punti problematici,
anche quelli di cui non ci si era nemmeno accorti.

Ruggero conferma con un esempio di quando lavorava in \textbf{Danimarca} con
l'\textbf{OMS}: giochi sul \textbf{diabete di tipo 2} nei bambini, per un target
di ragazzi delle scuole medie. C'era un team dell'OMS specializzato nella
comunicazione sul diabete, per garantire che il gioco fosse accurato e ne desse
una rappresentazione adatta al target. Il sensitivity reader aiuta a rendere il
gioco più accurato e ``più tematicamente coeso''.

\subsection{Producer}

Il \textbf{producer} è una figura a volte interna al publisher. In un progetto
grande è la figura di riferimento, ``quasi'' un project manager. In un progetto
piccolo si occupa spesso anche di parlare con le \textbf{tipografie} e di fare la
\textbf{lista dei materiali}, per assicurarsi che i materiali non siano solo
funzionali al gioco ma anche \textbf{economicamente sostenibili} da produrre:
``magari è un gioco bellissimo, ma poi costa più di quanto potete permettervi, e
nessuno vi comprerà mai un card game a 200 euro''.

\subsection{Perché conoscere queste figure}

Ruggero chiude la slide ``tornando a ombrello'': perché parlarne, se gli studenti
sono lì per fare game design? Non solo per saper dare un \textbf{nome} alle
figure. Non c'è niente di male a creare un gioco solo per sé; ma chi vuole fare
giochi a livello lavorativo, perché li giochino altre persone, ``magari più
persone possibile'', \textbf{si interfaccerà con queste figure}, e sapere che
cosa fanno e come parlarci è parte importante del percorso. Nel corso se ne
toccheranno diverse: il development con Marta, la produzione con Flavio, e
lezioni sulla UX, sul paragame, sul graphic design. ``Abituatevi mentalmente:
non siete in solitaria contro il mondo''.

% ---------------------------------------------------------------------
\section{Marketing e sales}
\label{a1:sec:marketing}

Arrivato all'angolo della slide con marketing e sales, Ruggero passa la parola
ai colleghi: ``io è meglio che non mi fate parlare di marketing e sales'', anche
per esperienze personali.

\subsection{Beatrice: i giochi trattati come prodotti}

Beatrice premette che le due voci sono in parte autoesplicative. Nel mondo
capitalista in cui viviamo, i giochi sono ``strumenti molto potenti per fare
\emph{n} cose'', ma sono trattati come \textbf{prodotti}, e il corso ragionerà
anche in quest'ottica. Nel \textbf{marketing} delle industrie ludiche e
videoludiche ci sono varie professionalità che gli studenti probabilmente già
conoscono: dal \emph{copy} alla grafica, tutto ciò che riguarda il
\textbf{vendere} il gioco o l'attività di gioco.

Il lato \textbf{sales} è forse la distinzione meno nota. Come in tante aziende, è
la parte che si occupa delle vendite \textbf{B2B}, cioè tra aziende: fare
accordi per vendere prodotti ludici a diverse realtà, per esempio ai
\textbf{distributori}, o a una casa editrice che vende alle librerie, che poi
vendono al consumatore finale. Chi si occupa di sales nel mondo ludico capisce
come vendere \textbf{alle aziende}. I due mondi si parlano spesso, e soprattutto
nelle realtà piccole si trovano figure che fanno ``un po' sales, un po'
marketing''.

\subsection{Daniele: che cos'è (e che cosa non è) il marketing}

Daniele il marketing si è trovato ``purtroppo'' a farlo, e lo fa anche per
l'accademia. ``Purtroppo'' perché in Italia è mal compreso: molti lo
identificano con l'\emph{advertisement} online. Non è così: il marketing è una
\textbf{disciplina molto più antica}, e in Italia ``abbiamo molto dimenticato''
come si fa il marketing tradizionale. Alcune cose da ricordare:
\begin{enumerate}
  \item \textbf{Il marketing non dice che cosa pubblicare, ma come pubblicarlo}.
        Esistono, nei videogiochi e nei giochi da tavolo, realtà che fanno
        prodotti \textbf{marketing driven}: l'esempio sono gli
        \textbf{hyper-casual game} per smartphone, quelli con il personaggio che
        va a destra e a sinistra e l'esercito che si allarga. Possono anche
        guadagnare, ma ``sono un po' la morte della creatività''. Il marketing
        ideale \textbf{prende un prodotto pensato da altri e gli dà le maggiori
        chance di essere venduto}: cerca il pubblico ideale per il prodotto che
        si ha, senza doverlo cambiare. Anzi, il mondo videoludico, e ludico in
        generale, ``è lastricato'' di errori di dipartimenti marketing che hanno
        cambiato completamente non il gioco, ma il modo in cui l'hanno
        pubblicizzato, e il gioco è fallito.
  \item \textbf{Usare strumenti del marketing non significa fare marketing}. Game
        design e game development a volte usano strumenti che vengono dal
        marketing (gli studenti ne vedranno con Beatrice, con Daniele e con
        Marta).
  \item \textbf{Il marketing è un gerundio}: si fa. È un'azione che si compie per
        comunicare; la comunicazione è attiva, e il marketing è l'insieme delle
        attività con cui si comunica il proprio gioco, dal \emph{copy} ai social.
        Può includere il digital marketing, ``ma non è solo quello''.
  \item \textbf{Il marketing esiste anche se si crede che non esista}. In Italia
        molti non vogliono ``fare marketing'' perché lo associano al lucro. Ma
        il marketing è \textbf{lo studio di come gli esseri umani comprano e
        vendono}, ed è innato in noi: ci sono logiche di cui siamo vittime, e
        ``anche se sei il bastian contrario, l'anticapitalista'', magari si
        compra comunque su Amazon perché fa comodo. Sono meccanismi che
        esistono, ``affascinanti anche, a volte''.
\end{enumerate}
Gli studenti non dovranno fare marketing, ma dovranno ricordare alcuni di questi
concetti, soprattutto nelle lezioni con Marta, che tratta argomenti legati al
marketing. E, chiude Daniele, come game designer e game developer ``per lo più
speriamo che il marketing non ci rovini i giochi'', soprattutto nel videoludico.

% ---------------------------------------------------------------------
\section{Game design e game development}
\label{a1:sec:design-development}

Ruggero torna sulla distinzione più importante della slide, ``proseguendo con
questa strada'':
\begin{description}
  \item[Game design] è chi \textbf{progetta l'esperienza} tramite la
        \textbf{creazione delle meccaniche} di gioco.
  \item[Game development] è chi fa \textbf{mediazione} tra la persona che scrive
        il game design e la \textbf{casa editrice}.
\end{description}

\subsection{I red flag dei designer}

Ruggero lo ripete perché lo vede ogni giorno nel suo lavoro: tra giochi
pubblicati e prototipi, ne valuta ``quasi uno al giorno, se faccio un conto
annuale'', e solo una minima percentuale arriva alla pubblicazione. Molte volte
i \textbf{red flag} vengono proprio da designer che non hanno presente questa
definizione, e che cominciano a mettere \textbf{clausole sulla non modifica del
gioco}: il gioco deve restare com'è. ``Questa cosa è impossibile con un
publisher, spesso''. Il publisher, per forza di cose, potrà chiedere di cambiare
la grafica, di reimpaginare il regolamento, di cambiare i materiali perché
troppo costosi: chiederà dei \textbf{compromessi}. E la persona che, all'interno
del publisher, fa questo lavoro è il \textbf{developer}.

Sta al designer capire \textbf{su che cosa del proprio gioco si può trovare un
compromesso e su che cosa no}. Publisher diversi lavorano in modo diverso, e il
designer ha sempre l'autorità di dire ``no, non voglio pubblicare questo gioco
se queste sono le modifiche''. Ma deve sapere chi fa il lavoro di development,
perché se il suo obiettivo è la \textbf{pubblicazione tradizionale} si
interfaccerà con questa figura.

\subsection{Daniele: una distinzione che manca al videogioco}

Daniele sottolinea che questa distinzione è \textbf{tipica del mondo analogico}
ed è ``molto interessante''. In Italia la parola \emph{developer} è diventata
negli anni quasi un sinonimo di programmatore di videogiochi, o di chiunque
lavori nel videoludico. Secondo lui c'entra anche il fatto che qualche anno fa
il sito \textbf{Gamasutra} ha cambiato nome in \textbf{Game Developer}, e la
parola è diventata un ombrello per tutte le professioni del videogioco. ``Mi fa
un po' ribrezzo'': il game developer dovrebbe essere un ruolo di design, un
designer che, come nell'analogico, fa \textbf{mediazione}. Nel videogioco i
publisher hanno in effetti i loro developer interni, che adattano il gioco, per
esempio alle console e ai loro kit di sviluppo. Ma la distinzione dovrebbe
esistere anche nel videogioco, perché è ``difficilissimo'' avere la mentalità
per fare le due cose insieme: chi fa design ha un focus che raramente si può
sovrapporre a quello di chi fa mediazione.

A Daniele piace la parola \emph{develop}, \textbf{sviluppare}, nel senso di
prendere qualcosa che è stato progettato e \textbf{renderlo possibile nel mondo
reale}. Ci sono progetti che hanno bisogno di uno sviluppo che ``polishi''
l'esperienza: la semplifichi, la migliori, la renda comprensibile e la concluda,
ma soprattutto la renda \textbf{stampabile} e \textbf{fruibile} per l'utente. La
metafora crossmediale è quella della \textbf{fotografia}: chi sviluppa le
fotografie non è chi scatta la foto, ma chi la rende visibile, anche
fisicamente. ``È un po' retro come metafora, ma ci capiamo''. ``È vero, però è
giusto'', conferma Ruggero.

% ---------------------------------------------------------------------
\section{Il mercato: Italia ed estero}
\label{a1:sec:mercato}

Il corso è in italiano e si svolge in Italia, ma bisogna conoscere anche il
\textbf{mercato estero}: molti giochi progettati da autori italiani vengono
pubblicati all'estero, e conoscere il mercato estero aiuta a capire come il
proprio gioco può arrivarci, ``o viceversa, come non può arrivarci''. La
tabella~\ref{a1:tab:mercato} riassume il confronto della slide.

\begin{table}[htbp]
\centering
\small
\begin{tabular}{@{}p{6.4cm}p{6.4cm}@{}}
\toprule
\textbf{Italia} & \textbf{Estero (soprattutto Europa)} \\
\midrule
il gioco è ancora visto come un'attività \textbf{infantile}, anche se la
normalizzazione è in crescita &
forte tradizione di gioco anche in un contesto \textbf{familiare adulto},
soprattutto ``dalla Germania in su'' \\
risposta \textbf{conservativa} sul tema ludico &
\textbf{premi di prestigio} riconosciuti in tutto il mondo (Spiel des Jahres,
As d'Or) \\
forte focus sulla \textbf{localizzazione} di giochi stranieri &
si cercano soprattutto giochi \textbf{originali} \\
investimento \textbf{molto basso} sulla creazione di giochi originali interni;
pochissimi designer assunti &
\textbf{developer interni} pagati con un contratto (Francia, Germania,
Danimarca\ldots) \\
\bottomrule
\end{tabular}
\caption{Il mercato del gioco analogico in Italia e all'estero, secondo la
slide di Ruggero e i commenti di Daniele.}
\label{a1:tab:mercato}
\end{table}

\subsection{L'Italia}

\begin{itemize}
  \item In Italia il gioco è ancora visto come un'\textbf{attività infantile},
        anche se la \textbf{normalizzazione} è in crescita. C'è ancora una forte
        risposta \textbf{conservativa} sul tema ludico.
  \item C'è un forte focus sulla \textbf{localizzazione} dei giochi stranieri: in
        una qualunque ludoteca o negozio, molti dei giochi da tavolo e di ruolo
        sono di publisher o designer esteri, pubblicati in italiano da publisher
        italiani. (La localizzazione, cioè portare un titolo dall'estero in
        Italia, esiste come funzione all'interno del publisher; Ruggero non l'ha
        messa nella slide delle figure perché non è il focus del corso.)
  \item C'è, ``purtroppo, preparatevi psicologicamente'', un \textbf{investimento
        molto basso} sulla creazione di giochi originali interni alla casa
        editrice. Pochissime case editrici italiane hanno designer interni pagati
        a tempo pieno con contratto a tempo indeterminato. Daniele lo dice senza
        mezzi termini: ``il mercato è molto piccolo, conosciamo praticamente
        tutto il settore'', e le persone che in Italia fanno game design
        analogico a tempo pieno con un contratto a tempo indeterminato ``sono
        meno delle dita delle mie due mani''. Di solito la formula è quella del
        \textbf{freelance}, del contratto, delle \textbf{royalty}. Chi va da un
        publisher aspettandosi di essere assunto a tempo pieno come game
        designer deve sapere che ``non siamo purtroppo ancora in quel tipo di
        situazione''.
  \item C'è però un contraltare: in Italia, soprattutto nel gioco di ruolo ma
        anche nel gioco da tavolo, c'è stata molta \textbf{imprenditoria}, anche
        grazie a \textbf{Kickstarter}, che ha aperto le porte ad altri sistemi
        di \textbf{crowdfunding}. È un altro argomento, perché richiede un team,
        un investimento e così via.
  \item E gli autori non mancano, anzi: le occasioni di incontro tra autori
        vedono ogni anno un aumento dei possibili autori\footnote{La trascrizione
        dice ``i DAG a Parma''; il riferimento è probabilmente a IDEAG, il
        circuito di incontri tra autori di giochi italiani, ma non è
        ricostruibile con certezza.}. Bisogna ``sgomitare'' tra tanti autori,
        che magari arrivano da esperienze molto diverse.
\end{itemize}
Ruggero precisa: l'investimento basso è ``dall'alto''. Le persone che fanno e
vogliono fare giochi sono tantissime e in aumento; scarseggia l'investimento,
sia dei publisher disposti a pagare stipendi sia dei \textbf{bandi} pubblici. I
bandi esistono (Beatrice ne parlerà per i giochi pervasivi, che spesso si
muovono in quella direzione, e Ruggero cita qualcuno che sta facendo un card
game con un bando italiano, ``una cosa assurda, ma lo sta facendo''), ma sono
una rarità, soprattutto rispetto all'estero.

\subsection{L'estero}

\begin{itemize}
  \item L'estero (dipende dalla nazione, ma in generale il resto d'Europa) ha una
        forte \textbf{tradizione di gioco} anche in un contesto \textbf{familiare
        adulto}, soprattutto ``dalla Germania in su''.
  \item Ha \textbf{premi di prestigio} riconosciuti in tutto il mondo, come lo
        \textbf{Spiel des Jahres} tedesco e l'\textbf{As d'Or} francese\footnote{La
        trascrizione dice ``lo Spield AR, la SDOR''.}.
  \item Cerca soprattutto \textbf{giochi originali}. Lo si vede dal fatto che
        molti autori italiani di giochi da tavolo pubblicano in Germania o in
        Francia e \textbf{non vengono pubblicati in Italia}, ``che è tutto
        dire''.
  \item Ha \textbf{figure di development interne}: non è raro vedere publisher
        francesi, tedeschi, danesi con developer pagati internamente, a
        contratto.
\end{itemize}
``Poi la situazione è ancora più drammatica'', ammette Ruggero, ``ma per ora
vogliamo tenere il morale un po' alto. È l'introduzione. Ci sarà tempo di andare
nella tristezza''.

Prima di proseguire, Ruggero insiste sulle domande: con lui si può interrompere
anche mentre parla, alzando la mano con le reaction o semplicemente togliendo il
microfono e chiedendo. ``Non mi offenderò mai, anzi, sarò felice di vedervelo
fare e mi ricorderò positivamente di voi''.

% ---------------------------------------------------------------------
\section{Perché serve una terminologia comune}
\label{a1:sec:terminologia}

La serie di slide successiva è ``una piccolissima parte di fondamenti di
design'' (``non spaventatevi, vi vedo che vi spaventate un po'\,''). Anche qui
molti temi verranno ripresi in dettaglio, ma l'idea è arrivare alla prossima
lezione con almeno un \textbf{glossario}, un \textbf{linguaggio condiviso}: se
qualcuno dice \emph{game developer}, tutti sanno di che cosa sta parlando.

Ruggero parte da una domanda, e precisa che non è retorica: \textbf{perché ci
serve una terminologia comune?} Gli studenti rispondono:
\begin{itemize}
  \item Gianmarco, ``molto pragmatico'': perché bisogna lavorare in team e
        capirsi. Vero, nel corso e anche fuori;
  \item un'altra studentessa: per la \textbf{comunicazione}, tra di noi e con
        tutti quelli che sono interessati nel settore. Vero anche questo:
        conoscere i termini comuni permette di parlare con chi fa game design o
        development nel dettaglio, nel tecnico, nello specifico.
\end{itemize}
Ruggero aggiunge altre due ragioni.

\paragraph{Saper comunicare il proprio gioco.} La prima riguarda la
\textbf{comunicazione del gioco} verso l'esterno. ``E lo dico non tirando acqua
al mio mulino, ma anzi togliendola'': \textbf{non è spesso il buon game design
a far diffondere un gioco}. ``Vorrei che lo fosse, sarebbe bellissimo, ma
purtroppo non è così''. Ci sono giochi bellissimi che non sanno comunicare, che
non sanno dire che cosa fanno, e che quindi è difficile, da fuori, trovare
invitanti; e ci sono giochi mediocri dal punto di vista del design che si
comunicano molto bene, e che quindi conoscono in molti. Saper parlare di quello
che si fa è utile \textbf{a sé stessi}, in primis.

\paragraph{Documentare per sé stessi.} La seconda è la \textbf{documentazione}:
documenti di design, di brainstorming e così via. ``Sembrerà una battuta, ma
provate a recuperare un documento di brainstorming una settimana dopo, se non
avete scritto bene le cose che avete scritto''. Chi l'ha fatto sa già di che
cosa si parla; ``chi non l'ha ancora vissuto, forse lo vivrà''.

% ---------------------------------------------------------------------
\section{Game e play}
\label{a1:sec:game-play}

La \textbf{prima, importantissima divisione terminologica} è quella tra
\emph{game} e \emph{play}. In italiano c'è un problema: \textbf{gioco} e
\textbf{giocare} hanno la stessa radice, mentre in altre lingue, tra cui
l'inglese, sono due parole strutturalmente diverse:
\begin{description}
  \item[Game] il \textbf{gioco come artefatto}: il gioco come ecosistema di
        meccaniche, come struttura di regole, come prodotto (``tra
        virgolette''), come oggetto fisico.
  \item[Play] l'\textbf{atto}, il verbo: le persone che giocano, l'approccio
        fisico all'attività del gioco.
\end{description}
È importante perché gli studenti faranno \textbf{game design}; esistono anche
ambiti di \emph{play design}, ma qui ``creerete giochi, non giocattoli''.

\subsection{Giochi e giocattoli}

Qual è la differenza? Uno studente risponde: \textbf{la presenza di regole}.
Esatto: un gioco ha delle regole, un'attività di \emph{play} non
necessariamente. Con un giocattolo si può fare \emph{play}: c'è chi gioca con i
Lego, chi gioca a guardia e ladri o a nascondino al parco, e magari alcune di
quelle regole sono ``più o meno semi-solidificate''. Ma il \emph{play} permette
spesso un approccio non strettamente legato alle regole.

Viceversa, \textbf{un game può esistere anche se nessuno lo sta giocando}. Se si
prepara su un tavolo una scacchiera con i pezzi degli scacchi, chi passa e la
vede riconosce il gioco degli scacchi, anche se nessuno sta giocando. Se invece
si vedono dei bambini correre in un parco giochi, quei bambini \emph{stanno
giocando}, ma non stanno giocando \emph{a un gioco} specifico. ``È importante
separare questa cosa''; ci si tornerà, ``ma per ora abituatevi a questa
differenza''.

\paragraph{Approfondimento: giocattoli, giochi e regole.} La stessa distinzione
è il punto di partenza del capitolo che Jesse Schell, in \emph{The Art of Game
Design}, dedica alla definizione di gioco (\emph{The Experience Rises Out of a
Game}). Schell nota che usiamo perfino un linguaggio diverso: ``un gioco è
qualcosa a cui si gioca'' (\emph{a game is something you play}), ``un giocattolo
è qualcosa con cui si gioca'' (\emph{a toy is something you play with}); e
arriva a definire il giocattolo come ``un oggetto con cui si gioca''. Analizzando
poi la definizione di gioco di Avedon e Sutton-Smith, ne ricava che ``i giochi
hanno regole, i giocattoli no'', e che le regole sono uno degli aspetti
\emph{definitori} del gioco, insieme alla volontarietà, agli obiettivi, al
conflitto e alla possibilità di vincere o perdere. Schell osserva anche che
molti giochi sono costruiti \textbf{sopra} un giocattolo (``una palla è un
giocattolo, ma il baseball è un gioco''), e tra i suoi consigli di
prototipazione c'è proprio \emph{Build the Toy First}: costruire prima il
giocattolo, cioè qualcosa che sia già piacevole da manipolare, e poi il gioco
intorno. È una lettura utile della distinzione di Ruggero: il \emph{play} può
fare a meno delle regole, il \emph{game} no.

% ---------------------------------------------------------------------
\section{Il framework MDA (e il DPE)}
\label{a1:sec:mda}

\subsection{Meccaniche, dinamiche, estetiche}

Un'altra terminologia fondamentale è quella del framework \textbf{MDA}
(\emph{Mechanics, Dynamics, Aesthetics}; in italiano a volte ``MDE''). Viene da
un \textbf{paper} nato nell'ambito del videogioco\footnote{È il paper di Robin
Hunicke, Marc LeBlanc e Robert Zubek, \emph{MDA: A Formal Approach to Game
Design and Game Research}, del 2004. Daniele aggiunge il paper alla cartella
condivisa durante la lezione.}, che ha ``solidificato'' queste terminologie, e
divide un gioco in tre aree:
\begin{description}
  \item[Meccaniche] le \textbf{regole per come sono scritte}. ``Ogni volta che
        pescate un 2 fate dei punti'': quella è una meccanica.
  \item[Dinamiche] l'\textbf{applicazione in tempo reale} delle meccaniche. La
        dinamica è il giocatore che, pescando una carta, \textbf{spera} che esca
        un 2. Quel processo non è una regola: ``il gioco non vi dice che dovete
        sperare che esca un 2''.
  \item[Estetiche] l'insieme di quello che il gioco vuole comunicare,
        l'\textbf{esperienziale}.
\end{description}

L'esempio di Ruggero mette insieme le tre aree. In un gioco in cui pescare un 2
fa fare punti (meccanica) e in cui quindi si spera che esca un 2 (dinamica),
\textbf{che estetica si ottiene se nel mazzo di 100 carte c'è un solo 2? E se ce
ne sono 50?} L'emozione di aver pescato \emph{l'unico} 2 del mazzo è
completamente diversa da quella di pescarne uno sapendo che ce ne sono
cinquanta e che è molto probabile trovarlo. Le regole sono quasi identiche; cambia
l'esperienza.

\subsection{Il falso amico: ``estetica''}

Daniele avverte di non sottovalutare la parola \textbf{estetica}, che in italiano
è un \textbf{falso amico}: sembra indicare la parte grafica, e invece non lo è.
Si incontreranno designer che hanno studiato MDA e altri che non l'hanno
studiato, e lo si capirà proprio da questo: i secondi diranno che le
\emph{aesthetics} sono la grafica. Le \emph{aesthetics} sono la parte
\textbf{esperienziale}: come nell'esempio di Ruggero, quello che lascia al
giocatore aver pescato l'unico 2 del mazzo.

\subsection{Il DPE: design, play, experience}

Per gli italiani, aggiunge Daniele, è più facile capire l'MDA attraverso un
altro paper, quello del \textbf{DPE}\footnote{È il framework \emph{Design, Play,
Experience} di Brian Winn, proposto come estensione dell'MDA per i giochi
educativi (\emph{serious games}). Anche questo paper viene aggiunto da Daniele
alla cartella condivisa.}, che di fatto fa un parallelismo assegnando
\textbf{design, play, experience}:
\begin{itemize}
  \item le \textbf{meccaniche} sono il lato \textbf{design};
  \item le \textbf{dinamiche} sono il lato \textbf{play};
  \item le \textbf{estetiche} sono l'\textbf{experience}, l'esperienza che riceve
        chi gioca.
\end{itemize}
``Se vivete le \emph{aesthetics} come l'esperienza che vive chi sta giocando,
state a postissimo''. In questo momento Beatrice lascia la lezione (``ciao,
Bea''). I paper di MDA e DPE sono in inglese; chi vuole può dargli una lettura.

\begin{figure}[htbp]
\centering
\begin{tikzpicture}[
    font=\small,
    blocco/.style={draw, thick, rounded corners=3pt, fill=gray!12,
                   minimum width=3.2cm, minimum height=1.0cm, align=center},
    dpe/.style={font=\footnotesize\itshape, align=center},
    esempio/.style={draw, rounded corners=2pt, text width=3.3cm,
                    align=left, font=\footnotesize, inner sep=4pt},
    freccia/.style={-{Stealth[length=2.5mm]}, thick}
  ]
  \node[blocco] (m) at (0,0)   {\textbf{Meccaniche}\\ \footnotesize le regole scritte};
  \node[blocco] (d) at (4.3,0) {\textbf{Dinamiche}\\ \footnotesize le regole in azione};
  \node[blocco] (a) at (8.6,0) {\textbf{Estetiche}\\ \footnotesize l'esperienza};
  \draw[freccia] (m) -- (d);
  \draw[freccia] (d) -- (a);
  \node[dpe, above=0.15cm of m] {DPE: \textbf{design}};
  \node[dpe, above=0.15cm of d] {DPE: \textbf{play}};
  \node[dpe, above=0.15cm of a] {DPE: \textbf{experience}};
  \node[esempio, below=0.45cm of m] (em)
    {``Ogni volta che peschi un 2 fai punti''};
  \node[esempio, below=0.45cm of d] (ed)
    {pescando, il giocatore \emph{spera} che esca un 2 (nessuna regola glielo
     chiede)};
  \node[esempio, below=0.45cm of a] (ea)
    {un solo 2 su 100 carte: pescarlo è un evento;\\
     50 su 100: è quasi scontato};
  \draw[dashed] (m) -- (em);
  \draw[dashed] (d) -- (ed);
  \draw[dashed] (a) -- (ea);
  % prospettive
  \draw[freccia, gray] ([yshift=-0.3cm]em.south west) --
    node[below, font=\footnotesize\itshape, text=black]
    {il designer agisce sulle meccaniche}
    ([yshift=-0.3cm]ea.south east);
  \draw[freccia, gray] ([yshift=-1.2cm]ea.south east) --
    node[below, font=\footnotesize\itshape, text=black]
    {chi gioca incontra prima l'esperienza}
    ([yshift=-1.2cm]em.south west);
\end{tikzpicture}
\caption{Il framework MDA con l'esempio del ``2'' usato da Ruggero, e la
corrispondenza con il DPE suggerita da Daniele. Le due frecce in basso
riassumono la lettura classica del paper: il designer costruisce le meccaniche
per ottenere un'esperienza, mentre il giocatore parte dall'esperienza.}
\label{a1:fig:mda}
\end{figure}

\paragraph{Approfondimento: le due direzioni dell'MDA e le otto estetiche.} Le
frecce della figura~\ref{a1:fig:mda} riprendono un'idea centrale del paper
originale, che la lezione non esplicita: designer e giocatore guardano il gioco
da lati opposti. In \emph{Challenges for Game Designers}, Brenda Brathwaite e
Ian Schreiber la riassumono così: i designer si chiedono quale estetica
vogliono ottenere, definiscono le dinamiche che porterebbero a quella
sensazione, e poi creano le meccaniche che producono quelle dinamiche. Lo stesso
libro conferma l'avvertenza di Daniele sul falso amico: ``le \emph{aesthetics}
non si riferiscono all'aspetto del gioco, ma alla risposta emotiva'' che il
designer vuole suscitare. E offre esempi di dinamiche nel senso di Ruggero: in
uno sparatutto i punti di \emph{spawn} sono una meccanica, mentre appostarsi
accanto a uno di essi per colpire chi riappare (lo \emph{spawn camping}) è una
dinamica; negli scacchi i movimenti dei pezzi sono meccaniche, mentre le
aperture codificate sono dinamiche. Per dare un nome alle estetiche, il paper
MDA propone una tassonomia, poi detta degli ``otto tipi di \emph{fun}'', che
Schell riprende in \emph{The Art of Game Design} come \emph{LeBlanc's Taxonomy of
Game Pleasures}: \textbf{sensazione}, \textbf{fantasia}, \textbf{narrativa},
\textbf{sfida}, \textbf{compagnia} (\emph{fellowship}), \textbf{scoperta},
\textbf{espressione} e \textbf{sottomissione} (il piacere di entrare nel
``cerchio magico'' e lasciarsi andare a un passatempo). È un primo esempio di
quello che i docenti chiedono con le parole vietate della prossima sezione:
invece di dire che un gioco è ``divertente'', dire \emph{quale} esperienza
offre.

% ---------------------------------------------------------------------
\section{Le parole vietate}
\label{a1:sec:parole-vietate}

``Attenzione, adesso qua triggeriamo delle persone''. Ruggero annuncia una regola
che ``tornerà'': nel corso ci sono delle \textbf{parole vietate}. ``Non perché
siamo dei censori''; il perché si vedrà in una lezione apposita, perché
spiegarlo adesso richiederebbe un'ora. Ma Daniele vuole spiegarne ``un
minimo''.

\subsection{La regola}

In \textbf{nessun materiale scritto} consegnato ai docenti, task compresi, sono
accettate queste parole, né in inglese né in italiano:
\begin{center}
\begin{tabular}{@{}ll@{}}
\toprule
\textbf{Inglese} & \textbf{Italiano} \\
\midrule
\emph{fun} & divertimento \\
\emph{immersion} & immersione \\
\emph{entertainment} & intrattenimento \\
\bottomrule
\end{tabular}
\end{center}
``Non li voglio neanche vedere''. Daniele lo ripete con un paragone: ``se mi fate
i disegni con Paint, benissimo, top, fantastico, siamo ancora best friend. Se mi
scrivete \emph{divertimento}, io sbrocco''. È una \textbf{regola dell'accademia}:
negli stessi corsi triennali di narrative design e game design le parole sono
vietate agli studenti per tre anni; ``per voi solo per tre mesi''. Daniele
aggiungerebbe volentieri un quarto termine, \textbf{gamification}.

\subsection{Perché}

Il motivo approfondito lo darà Ruggero nella lezione apposita. Daniele dà il
suo. Ruggero e Diego Ricchiuti sono molto contrari a \emph{fun} e a
\emph{immersion}; lui ha aggiunto \emph{intrattenimento}. Sono \textbf{parole
troppo passe-partout e troppo ambigue}:
\begin{itemize}
  \item \textbf{intrattenimento} in Italia è stata usata per anni come un modo
        per \textbf{sminuire}: ``ma sì, tanto è un gioco solo di
        intrattenimento, è un romanzo solo di intrattenimento''. Ma anche se un
        prodotto è ``solo di intrattenimento'', non vuol dire che non ci si debba
        lavorare con attenzione. E poi intrattenimento ``non significa niente'':
        ci sono tanti modi in cui una cosa può intrattenere;
  \item lo stesso vale per \textbf{fun}, e per \textbf{immersion}, che ha un
        problema in più: è una parola passe-partout per cui si può dire che un
        gioco è ``immersivo'' nel senso generale del termine, ma poi non si sa
        davvero che cosa voglia dire, \emph{immersivo in che senso}.
\end{itemize}

L'obiettivo didattico è che gli studenti \textbf{vadano oltre queste parole}:
``un po' come quando alle medie vi dicono smettetela di dire che tutto è bello,
perché ci sono vari gradi di bello, ci sono altri modi di dire bello''. Bisogna
\textbf{decostruirle}: quando si sta per usarle, chiedersi \textbf{che cosa si
intende} con quella parola e cercarne una \textbf{più precisa}.

Ruggero chiede, in uno dei ``rari casi'', di \textbf{non fare domande adesso} su
questa regola: ogni anno arrivano domande e commenti, e ci saranno lezioni
apposite di tre ore, con strumenti per capire il perché. Ma da subito gli
studenti devono abituarsi a evitare questi termini, sia perché si spera che
abbiano termini più precisi, sia perché queste parole ``portano con sé dei
problemi agganciati'' che si vedranno. ``Non spaventatevi, non verrete sparati in
un vicolo cieco di notte''.

\paragraph{Una piccola gamification.} Il divieto va visto come uno
\textbf{stimolo}, ``un po' come se stessimo gamificando la vostra esperienza'',
una sorta di user experience delle lezioni. Imponendosi di non usare queste
parole, si noterà come migliora la capacità di \textbf{descrivere i giochi} e le
proprie esperienze ludiche, e di essere più precisi. È un invito a ``obbligarsi a
essere più precisi e ad avere una mentalità più critica''. La lezione di
Ruggero sull'\textbf{involvement}, promette Daniele, sarà ``fenomenale'' e
``vi aprirà la testa''.

\paragraph{Approfondimento: anche i libri faticano a definire il ``fun''.} La
difficoltà che il divieto vuole far emergere è visibile anche nei testi di game
design. Schell, in \emph{The Art of Game Design}, prova a definire il \emph{fun}
e propone ``piacere con sorprese'' (\emph{pleasure with surprises}), ammettendo
subito che è strano poter usare una parola per tutta la vita, sapendo con
certezza che cosa significa, senza riuscire a esprimerla chiaramente quando
qualcuno lo chiede; e invita a mettere alla prova la definizione con
controesempi. La sua lente sul \emph{fun} (\emph{The Lens of Fun}) si limita,
non a caso, a chiedere quali parti del gioco siano \emph{fun} e \emph{perché}.
Anche \emph{immersion} ha molte letture: in \emph{21st Century Game Design},
Bateman e Boon la descrivono, tra l'altro, come la capacità di ``perdersi''
nel mondo artificiale presentato dal gioco, legata alla sospensione
dell'incredulità e ai giochi di mimesi. Sono significati diversi sotto la stessa
parola: proprio il motivo per cui, in un documento di design, conviene dire
quale si intende.

% ---------------------------------------------------------------------
\section{Chi crea, chi gioca, chi facilita}
\label{a1:sec:ruoli-esperienza}

Un'altra distinzione importante è quella tra \textbf{chi crea}, \textbf{chi
gioca} e \textbf{chi facilita}. Spesso, quando si fa un gioco, lo si fa provare
di persona a chi gioca. Ma ``la norma dei giochi pubblicati non funziona così'':
quando si pubblica un gioco non si vede chi lo compra e chi lo gioca, e non si
sa se qualcuno, a casa propria con i propri amici, segue le regole alla lettera e
interpreta correttamente quello che il designer aveva descritto. Quindi già in
fase di progettazione bisogna tenere a mente \textbf{chi sta facendo game design
e chi sta giocando}.

E nei giochi a evento, nei giochi pervasivi, nei giochi di ruolo da tavolo e dal
vivo (LARP) c'è un terzo pezzo del puzzle, ``importantissimo'': \textbf{chi
facilita}. La domanda da farsi è: \textbf{il regolamento è scritto per chi gioca
o per chi deve facilitare il gioco?} Bisogna abituarsi a pensare a tutte le
persone coinvolte nell'esperienza.

% ---------------------------------------------------------------------
\section{Playtest e iterazione}
\label{a1:sec:playtest}

\subsection{Nessun gioco nasce perfetto}

L'ultima voce del glossario è ``un'altra cosa importantissima'':
\textbf{playtest} e \textbf{iterazione}. Il corso è \textbf{pragmatico}:
``specialmente da parte mia'', dice Ruggero, non si vedranno teorie ``messe lì
per fare accademia sul game design nelle nostre torri d'avorio''. Ci saranno
teorie e framework, ma tutti estremamente pragmatici, perché l'obiettivo non è
solo capire i giochi, ma \textbf{saperli fare e iterare}.

Playtest e iterazione sono \textbf{inevitabili} nel game design, perché
\textbf{nessun gioco nasce perfetto}, soprattutto perché le persone sono
diverse: un gioco perfetto per una persona può essere incomprensibile per
un'altra. Chi progetta giochi per altre persone deve capire, dopo averli fatti
provare, \textbf{che cosa cambiare, che cosa sistemare, che cosa aggiungere, che
cosa togliere}. Ci saranno lezioni su playtest, prototipazione, raccolta dei
feedback, metodi etnografici. Ma bisogna abituarsi già ora all'idea che
\textbf{non esiste la consegna definitiva}. A un certo punto i task vanno
consegnati per forza di cose, ma, come diceva Daniele all'inizio, alla fine del
corso gli studenti avranno un \textbf{prototipo}, ``che sarà un prototipo'', e
gli strumenti per portare avanti il lavoro anche dopo.

\subsection{L'aiuto dopo il corso}

Non solo gli strumenti: anche l'aiuto dei docenti. Daniele racconta di aver
portato a un editore, a Play\footnote{Probabilmente \emph{Play -- Festival del
Gioco} di Modena; la trascrizione non permette di esserne certi.}, uno dei
progetti realizzati l'anno precedente da due studenti, un gioco di ruolo in
solitario (``un \emph{solo game} RPG molto, molto bello''). Se gli studenti
vorranno trasformare il proprio lavoro in un prototipo vero e proprio, possono
contattarlo: lo fa ``per lavoro, ma anche per passione''. Gestisce anche una
community per aiutare i designer ``un po' depressi'' a rimettersi a progettare.

Quindi: ``non sentitevi persi se durante la fase di iterazione vi sentite
persi, è normale''.

\subsection{Il dizionario comune e il production dictionary}

Daniele riprende un punto detto ``benissimo'' da Ruggero e che secondo lui non è
stato sottolineato abbastanza: il \textbf{dizionario comune} serve anche a
\textbf{non perdersi} mentre si itera. Se da un'iterazione all'altra i
documenti usano termini completamente diversi, e non c'è un dizionario di base,
ci si perde.

Nel videogioco, una delle primissime cose che un designer fa quando entra in
un'azienda è il \textbf{production dictionary}: si annota, in un file privato,
tutte le parole che usano i colleghi, per allinearsi al dizionario comune
dell'azienda, perché lì su un prototipo possono lavorare anche cento persone.
Nel gioco da tavolo non si fa molto, perché il mondo dell'analogico è
``talmente più piccolo''; ma serve comunque imparare un \textbf{dizionario
minimo}, dai docenti e dai professionisti. E il corso ne darà una parte,
entrando nel dettaglio di alcuni termini e framework.

\paragraph{L'esempio del feedback.} Perché serva, lo mostra un esempio. Si
testa un gioco e ci si annota come feedback: ``questo gioco deve essere più
divertente''. Che cosa vuol dire? Come lo si migliora? Quella frase non dà
elementi utili. Se invece si scrive che deve essere \textbf{più adrenalinico},
\textbf{più frenetico}, \textbf{più rilassato}, \textbf{più cozy}, o
\textbf{più bilanciato matematicamente}, si sa dove andare a lavorare.

\paragraph{Approfondimento: la regola del ciclo.} Il principio ``nessun gioco
nasce perfetto'' ha un nome preciso in \emph{The Art of Game Design}: è la
\textbf{Rule of the Loop}, la regola del ciclo. ``Più volte provi e migliori il
tuo design, migliore sarà il tuo gioco''. Schell precisa che non è una lente,
cioè un punto di vista, ma ``una verità assoluta'', senza eccezioni: ci saranno
momenti in cui si cercherà di convincersi che ``questa volta il design è così
buono che non serve testarlo'', e ogni volta se ne pagherà il prezzo. Ed è
interessante la ragione per cui la regola pesa di più sul digitale: secondo
Schell il tempo e il denaro necessari per provare e correggere un sistema sono
molto maggiori nei videogiochi che nei giochi tradizionali. Per un gioco di
carte o da tavolo, con il tempo per provarlo e cambiarlo più e più volte, il
consiglio è semplicemente di farlo. Il gioco analogico è, da questo punto di
vista, la palestra ideale dell'iterazione.

% ---------------------------------------------------------------------
\section{Tabletop Simulator}
\label{a1:sec:tts}

\subsection{Il software di settore}

``Last but not least'': il corso usa un software, \textbf{Tabletop Simulator},
che si compra su \textbf{Steam} e costa tra i \textbf{10 e i 20 euro} (``aspettate
che vada in sconto, va in sconto ogni due per tre''). Daniele spiega perché:
\begin{itemize}
  \item è \textbf{il software di settore} del gioco da tavolo, quello che i
        publisher usano per i prodotti digitali, interni ed esterni;
  \item essendo su Steam, permette di fare \textbf{playtest} e di giocare con gli
        amici di Steam. Daniele stesso sta testando lì il proprio gioco;
  \item funziona anche su Mac: Windows o Mac, nessun problema.
\end{itemize}
Ci sarà una lezione apposita, di Ruggero, sulla creazione del proprio
\textbf{modulo} di gioco.

\paragraph{Che cosa non è.} Daniele chiarisce subito un ``fraintendimento
gigantesco'', soprattutto per chi ha già esperienza: Tabletop Simulator
\textbf{non richiede programmazione} e \textbf{non è un software per
trasformare i giochi da tavolo in videogiochi}. È un software che
\textbf{simula dei componenti fisici su un tavolo}: serve per giocare i board
game in digitale. Simula anche la fisica: c'è perfino il comando per
\textbf{ribaltare il tavolo} (``non usatelo ora'').

\paragraph{La partnership mancata.} Da anni Daniele prova a fare una
partnership con gli sviluppatori: due anni prima ha contattato l'azienda per
acquistare delle chiavi per gli studenti, dato che il software si usa nei corsi.
L'azienda ha risposto che non fa accordi commerciali. ``Credo sia la prima
azienda nella mia vita che dice che non vuole che io le dia dei soldi''. Gli
studenti dovranno quindi provare ad acquistarlo da sé; se i docenti riusciranno a
darglielo, bene.

\paragraph{La domanda di Gianmarco.} Gianmarco segue le lezioni da un computer
aziendale ``super chiuso'', su cui non può installare Steam: deve usare Tabletop
Simulator durante le lezioni? Ruggero risponde che per la lezione dedicata è
meglio poter seguire e rifare in diretta. Ma il software non sarà richiesto
necessariamente per il project work: un gruppo potrebbe fare un gioco pervasivo
o un gioco di ruolo, non un gioco da tavolo. Essendo un software di settore, è
comunque bene conoscerlo. Ruggero aggiunge che le versioni non ufficiali che
circolano non permettono di giocare in molti, e quindi non servono per
playtestare o giocare insieme.

\subsection{Una dimostrazione}

Daniele mostra il software aperto da Steam:
\begin{itemize}
  \item si può creare una partita \textbf{single player} o \textbf{multiplayer};
        per giocare tra compagni si crea un server multiplayer e ci si aggiunge
        come amici su Steam;
  \item si dà al server un nome, un numero di giocatori, eventualmente una
        password; creato il server, si apre la scelta dei giochi;
  \item ci sono giochi gratuiti, giochi da acquistare (DLC, non obbligatori) e
        soprattutto un \textbf{Workshop} ``ampissimo'', con decine di migliaia di
        giochi caricati dagli utenti: ``se c'è un board game che conoscete,
        quel board game è su Tabletop Simulator, giocabile dal Workshop''.
        L'esempio mostrato è l'open beta ufficiale del gioco da tavolo di
        \emph{Slay the Spire}\footnote{La trascrizione dice ``Slade Spire''.}:
        molti publisher, per farsi pubblicità, mettono il proprio gioco in demo o
        giocabile in digitale;
  \item c'è un \textbf{tutorial} iniziale che spiega come muoversi sul tavolo;
        con \emph{Esc} e \emph{Help} si trovano i comandi da tastiera.
\end{itemize}

\subsection{Il primo task}

Il primo (e per ora unico) task è questo: \textbf{entro il 10 ottobre}, giorno
della lezione di Ruggero su Tabletop Simulator, giocare \textbf{tre o quattro
volte in totale}, online, tra compagni di corso o con amici. Non si tratta di
imparare a creare moduli, ma di prendere la mano \textbf{lato giocatore}: aprire
il software, impararne i comandi, giocare. ``Per me potete anche andare a
prendere domino o mahjong'' (``mahjong è bellissimo tra l'altro'', aggiunge
Daniele, che sta facendo un videogioco sul mahjong), o il pachisi o il go. Gli
anni precedenti si chiedeva a ciascuno di proporre ai compagni un gioco da
tavolo che apprezzava. È anche un modo per conoscersi. Ruggero spiega il
perché: usarlo per giocare è facile e intuitivo, usarlo per creare moduli è più
complesso; ma chi arriva al 10 ottobre avendo già aperto il software e giocato,
troverà più intuitivo anche creare un modulo. E l'avvertimento: ``non riducetevi
all'ultimo, non fate una partita il 9 sera''.

\subsection{Tabletopia e altri vantaggi}

Tancredi racconta di aver provato a caricare un suo gioco su
\textbf{Tabletopia}. ``Questo è un vaso di Pandora'', risponde Daniele.
Tabletopia è molto usato, ma secondo lui il suo strumento di creazione dei moduli
è pessimo, ed è il motivo per cui la maggioranza dei moduli ufficiali dei
publisher è su Tabletop Simulator. Tabletopia ha un vantaggio vero, è
\textbf{gratuito}, ma ha scelte di ``qualità della vita'' poco felici.

Tabletop Simulator ha invece un vantaggio ``fenomenale'' che si tende a
sottovalutare: essendo su Steam, \textbf{il modulo si pubblica direttamente sul
Workshop di Steam}; si possono invitare gli amici a iscriversi e giocare con
loro, purché abbiano un account Steam. Rispetto ad altri software da imparare,
costa ``tra i 10 e i 20 euro, ragazzi: si parla proprio di spiccioli''. ``Io con
Tabletop Simulator ci ho playtestato la vita''.

Altre caratteristiche utili:
\begin{itemize}
  \item si possono creare \textbf{moduli non elencati} (\emph{unlisted}),
        perfetti da passare a un publisher prima che il gioco sia pubblicato:
        non sono pubblici, e nessuno può rubare i materiali;
  \item il software ha una grafica datata (``\emph{World Art 1998}'', scherza
        uno dei docenti), ma lo stanno portando su \textbf{Unity} con molti investimenti
        e aggiornamenti; ascolta la comunità e introduce funzioni utili anche per
        chi fa design;
  \item il suo server Discord è molto frequentato, e c'è un canale YouTube
        dedicato solo al design di giochi di carte\footnote{La trascrizione ne
        riporta il nome come ``Distraction Maker''; non è verificabile.} con un
        server Discord dove si fanno playtest: il 99\% è su Tabletop Simulator;
  \item molti editori, quando si propone loro un gioco da tavolo,
        \textbf{chiedono se c'è un modulo per Tabletop Simulator}.
\end{itemize}

% ---------------------------------------------------------------------
\section{Le domande a Daniele: il tema del project work}
\label{a1:sec:domande-daniele}

Daniele ha ``un minuto'' per rispondere alle domande sul corso prima di dover
andare. Uno studente chiede, ``giusto per curiosità'', se il project work avrà
un tema o sarà libero.

\begin{itemize}
  \item C'è un \textbf{mandate}, che i docenti propongono ogni anno e
        aggiornano leggermente; verrà presentato il \textbf{25 settembre}.
        ``Non facciamo spoiler''. Conterrà quello che si chiede di fare e quello
        che non si chiede: dei \textbf{must}, dei \textbf{must not} e dei
        \textbf{nice to have}.
  \item Verrà chiesto di lavorare su un'\textbf{idea originale} che non si è
        portata al corso: originale nel senso di \textbf{creata apposta}, non
        nel senso di ``mai vista prima''.
  \item L'idea va \textbf{partorita insieme}. Altri anni alcuni studenti hanno
        provato a imporre la propria idea al gruppo, ``e l'esperienza è stata
        pessima per tutti''. Quindi: \textbf{non imponete le vostre idee, fate
        brainstorming con gli altri}. Daniele, che fa ``molto spesso
        brainstorming'', darà strumenti per farlo: nessuno arriverà impreparato
        a quella lezione.
\end{itemize}

Ruggero aggiunge che il mandate è \textbf{abbastanza flessibile} perché tutti
possano fare qualcosa che li interessi: ``l'ultima cosa che vogliamo è forzarvi a
lavorare su qualcosa che non vi piace''. Bisogna lavorare su ciò che piace e per
cui si ha uno stimolo personale, anche nella \textbf{tipologia} di gioco: nessuno
sarà obbligato a fare un gioco di ruolo o un board game, e la lezione sui casi
studio mostrerà quante forme diverse di gioco analogico il corso accetta.

L'unica cosa da \textbf{non fare}, ``perché è contraddittorio per questo corso'',
sono i \textbf{giochi grandi}: ``lo chiediamo per voi''. Niente progetti
mastodontici\footnote{La trascrizione dice ``non fate delle cose da gnome'':
probabilmente il nome di un gioco da tavolo molto grande e complesso, non
ricostruibile con certezza.}, niente ``cose alla \emph{Dungeons \& Dragons}''.
Si può fare un gioco di ruolo, ma con un'esperienza ``meno corposa'': ci sono
tanti giochi di ruolo da 15 o 20 pagine, ``bellissimi, fichissimi'', che si
possono leggere online. Il consiglio è farsi una cultura anche su
\textbf{itch.io}, ``uno dei portali più importanti del gioco di ruolo
internazionale'', di cui si parlerà anche nella parte di ricerca e sviluppo.

Poi Daniele saluta (``noi ci vediamo venerdì a lezione'') e lascia la
conduzione a Ruggero.

% ---------------------------------------------------------------------
\section{Le domande finali con Ruggero}
\label{a1:sec:domande}

\subsection{Le regole delle domande}

Ruggero spiega che cosa vuol dire ``stare con lui''. Alla fine di ogni sua
lezione ci saranno \textbf{almeno 20--30 minuti di domande} libere, \emph{free
form}. Vuole che fare domande sia \textbf{normalizzato}: ``io gestisco malissimo
l'\emph{horror vacui} e divento super cringe molto alla svelta se non ci sono
domande in giro''. Sono benvenute le domande sulla lezione appena fatta, ma non
solo: il corso è crossmediale, e anche le domande che vengono da
\textbf{divagazioni}, da parti laterali, ``fatele''. Quello che sembra scontato o
banale a uno studente può non esserlo per il docente, e viceversa; e può essere
utile anche agli altri. ``Se vi sembrano stupide, fatele il doppio''. ``Magari
sembro incazzato, ma vi giuro sono rilassatissimo. Non mordo, magari abbaio un
po'\,''.

\subsection{Come approcciarsi a Steam e Tabletop Simulator}

Rita non ha mai usato né Steam né Tabletop Simulator e chiede come
approcciarsi, e se c'è modo di mettersi d'accordo tra compagni su Discord, magari
con un sondaggio sui giochi, per formare gruppi di gioco. Ruggero conferma: sul
Discord dell'accademia c'è un \textbf{canale apposito per il corso}, ed è degli
studenti. Lo possono usare per organizzarsi a giocare, per chiedere a un
compagno di prestare un gioco, per domande tecniche (come si scarica, come si
installa). ``Sarebbe proprio ideale che vi parlaste e vi organizzaste''.

Anche perché, aggiunge con un ``piccolo spoiler'' (``nessuno vi ha detto
niente''), se fossero i docenti a organizzare i gruppi di gioco, il task lo
starebbero facendo loro. Bisogna \textbf{abituarsi a organizzare le sessioni di
gioco}, perché arriveranno i task sul \textbf{playtest} e su come organizzarlo,
e a quel punto si sarà ``meno disperati''. Sembra scontato, ma saper
\textbf{organizzarsi per giocare in autonomia} ``vi sarà utilissimo come skill in
questo settore, ve lo giuro''.

\subsection{Il bilanciamento}

Angelo confessa ``una paranoia grossissima'': ci saranno lezioni sul
\textbf{bilanciamento}? ``Ma certo''. Sono proprio le lezioni di Marta sul
development: il bilanciamento è di solito un tema che tocca \textbf{i developer
più dei designer}. L'anno precedente c'era una sola lezione con Marta, e i
docenti si sono resi conto che serviva andare ancora più a fondo: per questo
quest'anno le lezioni sono due. ``Non è una paranoia, è una legittima domanda'':
Ruggero insegna da anni, vede altri corsi, parla con altri docenti e studenti, e
``mi incazzo spesso'' perché molti corsi di game design \textbf{non fanno
bilanciamento}. I docenti volevano invece un corso che coprisse tutte le parti
del game design: non solo l'idea del gioco, ma il bilanciamento, la produzione,
i playtest, l'iterazione, la scrittura dei regolamenti.

\subsection{Si può giocare con i docenti?}

Una studentessa chiede se si può giocare con i docenti. Per motivi etici,
durante il corso questo è limitato alle parti di \textbf{valutazione} dei giochi
e di \textbf{tutoraggio}. Dopo il corso, ``nessun problema, almeno con me''.

\subsection{Dove si trovano i giochi su Tabletop Simulator}

Ruggero mostra dove si trovano i giochi: nella propria libreria di Steam, nella
pagina di Tabletop Simulator, alla voce \textbf{Workshop}, dove ci sono ``più o
meno 36.000'' giochi. Una studentessa scopre intanto che l'app di Steam per
telefono è solo una \emph{companion app}: per registrarsi e giocare serve il
PC.

Per provare, chiede agli studenti un titolo:
\begin{itemize}
  \item \emph{Keep the Heroes Out}: non compare, ``non tutti i giochi ci sono,
        perché molti publisher non li mettono''; forse, essendoci già una
        versione del gioco su Steam, il modulo è stato tolto;
  \item \emph{Magic: The Gathering}: ``almeno 50 moduli'', anzi ``123 pagine'';
  \item \emph{Bang!}: tre moduli.
\end{itemize}
Bisogna fare attenzione: alcuni moduli sono in lingue specifiche (``magari il
primo è in russo''), alcuni hanno già l'espansione, altri sono solo il gioco
base. Conviene sfogliare e non fermarsi al primo risultato. E nel Workshop si
trovano anche cose molto particolari: giochi fuori produzione, giochi vecchi,
giochi ancora in prototipo.

\subsection{I ruoli nei team e le dinamiche tra le persone}

Uno studente torna sulla slide delle figure: nelle realtà piccole molti ruoli
hanno ``il doppio o triplo cappello''. Quali ruoli restano distinti, e come si
creano le sinergie tra persone con obiettivi diversi?

Ruggero premette che l'ideale ``platonico'' è un ambiente di lavoro costruttivo,
basato sulla fiducia e ``il più \emph{safe} possibile''; ma nell'industria non
sempre è così, e le strutture variano moltissimo: alcuni publisher sono molto
verticali, altri hanno ruoli già solidificati (se un publisher ha già i suoi
developer e non sta crescendo, entrarci come developer è praticamente
impossibile). ``Se prendiamo due aziende in Italia, vediamo due strutture
interne completamente diverse''. Non esiste una scorciatoia per farsi assumere
da un editore.

Poi prende ``l'esempio più comune di tutti'', la media dei giochi da tavolo
pubblicati (figura~\ref{a1:fig:relazioni}):
\begin{itemize}
  \item il \textbf{game designer} è una figura \textbf{esterna}, con un
        \textbf{contratto a royalty} con il publisher;
  \item la sua figura di riferimento è il \textbf{publisher} stesso, e la
        conversazione designer--publisher è ``molto calda'';
  \item se il publisher ha dei \textbf{developer}, il developer entra in questa
        conversazione: designer e developer collaborano molto, e internamente
        developer e publisher collaborano molto. Spesso, nei publisher piccoli,
        è il publisher a fare le veci del developer;
  \item \textbf{graphic design} e \textbf{illustrazione} tendono a parlare molto
        con il publisher e \textbf{molto poco con il designer}. A meno che non si
        tratti di uno studio indie molto piccolo, è difficile che il designer
        sappia anche solo chi è l'illustratore prima che il gioco esca.
\end{itemize}

\begin{figure}[htbp]
\centering
\begin{tikzpicture}[
    font=\small,
    ruolo/.style={draw, thick, rounded corners=3pt, fill=gray!12,
                  minimum width=2.8cm, minimum height=0.9cm, align=center},
    esterno/.style={draw, thick, dashed, rounded corners=3pt,
                    minimum width=2.8cm, minimum height=0.9cm, align=center},
    calda/.style={{Stealth[length=2.5mm]}-{Stealth[length=2.5mm]}, very thick},
    tiepida/.style={{Stealth[length=2mm]}-{Stealth[length=2mm]}, thin, dashed}
  ]
  \node[esterno] (gd) at (0,0) {Game designer\\ \footnotesize esterno, a royalty};
  \node[ruolo] (dev) at (5,1.4) {Developer\\ \footnotesize se c'è};
  \node[ruolo] (pub) at (5,-1.4) {Publisher};
  \node[ruolo] (art) at (10,-1.4) {Graphic designer\\ e illustratore};
  \draw[calda] (gd) -- (dev);
  \draw[calda] (gd) -- (pub);
  \draw[calda] (dev) -- (pub);
  \draw[calda] (pub) -- (art);
  \draw[tiepida] (gd.south) to[bend right=25]
     node[below, font=\footnotesize\itshape] {quasi nessun contatto} (art.south);
  \node[draw, dotted, rounded corners=4pt, fit=(dev)(pub)(art), inner sep=8pt,
        label={[font=\footnotesize\itshape]above right:interno al publisher (di solito)}] {};
\end{tikzpicture}
\caption{Le relazioni tipiche intorno a un gioco da tavolo pubblicato, secondo la
risposta di Ruggero. Le frecce spesse indicano una collaborazione stretta e
frequente; quella tratteggiata, un contatto quasi assente. Nei publisher piccoli
il publisher fa spesso anche da developer.}
\label{a1:fig:relazioni}
\end{figure}

\subsection{Fare da developer di sé stessi?}

Rita chiede in chat se un designer può fare anche da developer del proprio
gioco, o se Ruggero lo sconsiglia. In generale \textbf{lo sconsiglia}, per lo
stesso motivo per cui si sconsiglia in altri media: i \textbf{darling}. Il
designer ha delle meccaniche a cui è affezionato, e il developer a volte dovrà
\textbf{tagliarle}; è molto difficile, se si è ``gestori di sé stessi'', tagliare
cose che si vogliono nel gioco. ``È un po' come lasciare fare il lavoro di editing
allo scrittore che ha già scritto il libro'': può rivedersi, certo, ma difficilmente
cambierà cose che farebbe bene cambiare. Nelle lezioni con Marta si vedrà perché,
e alla fine del corso ``non avrete dubbi nemmeno voi'' nel separare designer e
developer.

Ma tutto dipende dallo \textbf{scopo del progetto}: se si fa un gioco in tre,
qualcuno quel lavoro deve farlo per forza, e magari il designer fa anche il
developer, compensando con \textbf{più playtest}. ``Il mondo attuale e il mondo
ideale sono due cose molto diverse, e lo vediamo non solo dai giochi ma dalla
politica'', ma ``non entriamo in questo vaso di Pandora''.

\subsection{Il ruolo indispensabile e la pubblicazione}

Lo studente di prima ne trae una conclusione: senza game designer il gioco non
esiste. ``Quello sicuramente''. Ma Ruggero coglie l'occasione per approfondire
un punto che forse non era passato: \textbf{fare design di un gioco per
pubblicarlo non è l'unica opzione possibile}. Un gioco non è ``meno gioco''
se non è pubblicato da un editore. È legittimo mettere un gioco su itch.io in
\textbf{print and play}, fare un gioco a evento senza venderlo, fare un
\emph{murder party}, una escape room, un board game per una persona amica:
``quell'esperienza lì non è meno gioco''. Dipende dall'obiettivo del progetto.

Ruggero ha lavorato molto anche con le \textbf{scuole}, dove il publisher non
esiste: è un gioco per le classi, e la figura più vicina è l'editoria scolastica,
``che però non fa il publisher''. Ci sono compromessi diversi sulla materialità
e su come il gioco si presenta, che si vedranno nella prototipazione e nei casi
studio. Quindi l'unica figura che deve esserci sempre è \textbf{chi fa game
design}, anche se magari non sa di chiamarsi game designer e non ha mai
sentito parlare di game design: ``questo non la rende meno game designer''.

\paragraph{Un libro consigliato.} Ruggero fa anche mentoring e gestisce gruppi
di settore, e spesso gli chiedono che cosa leggere per imparare il game design.
Non è il tipo che obbliga a leggere libri (``non vi obbligherò a leggere
niente''), ma ce n'è uno che consiglia sempre, ed è il primo che ha scelto di
leggere: \textbf{\emph{Rise of the Videogame Zinesters}} di \textbf{Anna
Anthropy}\footnote{La trascrizione dice ``Anna Antrofi'' e ``Rise of the
Videogames in Esters''.}. È un libro un po' datato, ma ancora reperibile; parla
soprattutto di videogiochi, ma ``per l'80\%'' il contenuto si applica al gioco e
al game design in generale. È un libro sulla \textbf{democratizzazione del gioco}:
il gioco come strumento per parlare di sé e come \textbf{medium legittimo}. Come
esistono film drammatici, tristi, intensi, possono esistere giochi su tematiche
delicate, e se ne vedranno nei casi studio. Molti giochi fatti con questa
direttiva non hanno un publisher, che tende a pubblicare ciò che dà un ritorno
commerciale.

\paragraph{Commerciale e personale.} ``Io spero che voi non farete giochi solo
per scopo commerciale''. Certo, servono un tetto sulla testa, il cibo in
tavola, un posto in cui dormire, e così si finisce per fare giochi che non sono
il proprio gioco dei sogni. Spesso è un mix: Ruggero ha fatto giochi commerciali
``per pagarmi l'affitto'' e giochi che voleva fare e su cui ha perso soldi,
``sapendolo proprio che sarebbero stati un fallimento''. I giochi, come i libri e
le canzoni, si fanno per molti motivi, e quelli commerciali di successo non sono
più importanti degli altri. Tra i giochi della sua slide c'erano il gioco da
tavolo ufficiale di \emph{Cyberpunk}\footnote{La trascrizione dice ``Cyberpunk
del 1977'': quasi certamente un gioco legato a \emph{Cyberpunk 2077}.} e
\emph{Marvel United}, su IP grosse e famose, ma anche giochi fatti di ``una
candela e una busta su un tavolo''. E se gli si chiede quale preferisce in
termini di contenuto, non sono quelli sui supereroi.

\subsection{Il gioco di ruolo in Italia}

Nicola, che è interessato soprattutto al gioco di ruolo e conosce già persone
del settore, ha ``una leggera paura'': che chi studia game design finisca per
fare il \textbf{freelance} a vita, ``perché non ci sarà mai un'occupazione per
bene'', o il \textbf{copywriter}. Conosce un ragazzo che ha lavorato in quel
mondo e che ora fa altro.

Ruggero risponde che la maggioranza dei giochi nel suo portfolio sono giochi di
ruolo, quindi il tema lo tocca da vicino, e che in Italia quello che dice Nicola
``è semplicemente vero''. Molti publisher piccoli hanno una persona di
riferimento che copre più ruoli: impagina, fa i playtest, fa un po' di
development. È ``un settore economicamente completamente disastrato'', per
motivi culturali che non apre: oggi in Italia è praticamente impossibile
\textbf{vivere di giochi di ruolo}. Chi ci lavora lo fa nel tempo libero, al di
fuori del lavoro principale, oppure coprendo molte altre cose: molti publisher
italiani di giochi di ruolo pubblicano anche titoli esteri e localizzano, non
fanno solo cose originali. ``È un grosso calderone''. Poi si ferma: ``questo è
un \emph{rant}'', e ``i \emph{rant} tornano sempre'', ma non vuole che gli
studenti escano dalla prima lezione ``con la disperazione nel cuore''. Ci sarà
spazio per parlarne quando avranno gli strumenti per capirlo, ``per non farvi
abbattere''.

\subsection{Che cosa bisogna sapere già}

Uno studente chiede se ci sono nozioni che dovrebbero già avere e che possono
recuperare: per esempio, sente parlare di giochi ``German'' o ``American'' e non
sa che cosa siano. Ruggero spiega l'organizzazione del corso:
\begin{itemize}
  \item le \textbf{prime due lezioni} sono di fondamenta. Questa era ``dall'alto'':
        come è strutturato il corso, le tempistiche;
  \item la \textbf{prossima lezione} è di tre ore di soli \textbf{casi studio},
        ``a manetta, a raffica'', divisi per categorie e tipologie, con
        ragionamenti su categorie e generi di gioco: lì si vedranno German,
        American ``e anche altre più utili'';
  \item nel resto del corso le terminologie specifiche sono inserite nelle
        lezioni dedicate: il \textbf{paragame} nelle lezioni sul paragame, il
        \textbf{modello matematico} nelle lezioni di Marta sul development.
\end{itemize}
Per ora agli studenti si chiede di sapere quello che si è visto oggi:
meccaniche, dinamiche ed estetiche; le figure professionali, designer e
developer; e soprattutto \textbf{game e play}.

\subsection{Game e play: i quattro quadranti}

``Specialmente \emph{game} e \emph{play}'', insiste Ruggero: è più difficile di
quello che sembra, perché l'italiano non abitua a separare le due cose, ed è
possibile averle usate per anni come sinonimi. Nella sua esperienza con chi si
avvicina al design è ``uno dei punti più critici''. Spesso i commenti che si
fanno sui giochi giocati sono commenti sul \emph{play} applicati al
\emph{game}: ``il gioco non mi è piaciuto perché, giocando con Tizio, è successo
questo''; ``il gioco mi sembra sbilanciato perché Tizio al tavolo ha fatto
questo''.

Separare \emph{game} e \emph{play} serve a capire quali giochi sono progettati
bene e quali male, e soprattutto a riconoscere i \textbf{quattro possibili
incroci} (figura~\ref{a1:fig:quadranti}): un gioco progettato bene che piace
giocare; uno progettato male che non piace; uno progettato male che piace; uno
progettato bene che non piace. ``Non è importantissimo in generale nella vita,
ci sono cose più importanti'', ma per chi fa game design lo è: ``a volte vi
chiederanno di lavorare su un gioco che non vi piace giocare, e quel gioco deve
essere fatto bene lo stesso''.

\begin{figure}[htbp]
\centering
\begin{tikzpicture}[
    font=\small,
    cella/.style={draw, thick, minimum width=4.4cm, minimum height=1.9cm,
                  text width=4.0cm, align=center}
  ]
  \node[cella, fill=gray!20] (a) at (0,0)     {\textbf{progettato bene}\\ e \textbf{mi piace} giocarlo};
  \node[cella]               (b) at (4.4,0)   {\textbf{progettato bene}\\ ma \textbf{non mi piace} giocarlo};
  \node[cella]               (c) at (0,-1.9)  {\textbf{progettato male}\\ ma \textbf{mi piace} giocarlo};
  \node[cella, fill=gray!8]  (d) at (4.4,-1.9){\textbf{progettato male}\\ e \textbf{non mi piace} giocarlo};
  \node[font=\footnotesize\itshape, above=0.15cm of a] {mi piace (\emph{play})};
  \node[font=\footnotesize\itshape, above=0.15cm of b] {non mi piace (\emph{play})};
  \node[font=\footnotesize\itshape, rotate=90, left=0.15cm of a, anchor=south]
       {bene (\emph{game})};
  \node[font=\footnotesize\itshape, rotate=90, left=0.15cm of c, anchor=south]
       {male (\emph{game})};
\end{tikzpicture}
\caption{I quattro incroci tra la qualità del \emph{game} (come è progettato) e
l'esperienza di \emph{play} (se piace giocarlo), descritti da Ruggero. Per un
designer i quadranti ``scomodi'' sono i due fuori diagonale: un gioco può
piacere pur essendo progettato male, e va saputo progettare bene anche un gioco
che non piace.}
\label{a1:fig:quadranti}
\end{figure}

\paragraph{La domanda di Tancredi.} Tancredi chiede conferma: il \emph{play}
non è il ``giocare senza regole'', ma può essere anche il giocare a un gioco? Sì:
il \emph{play} è l'\textbf{atto}, e può essere fatto con un \emph{game} o senza.
E un \emph{game} può esistere senza che nessuno lo giochi: gli scacchi non sono
meno scacchi perché in quel momento nessuno ci gioca. E se, giocando, una
persona applica le regole in modo diverso? Allora \textbf{non si sta giocando
quel gioco}.

\subsection{House rule: il caso di Monopoly}

Qui Ruggero fa un po' di \emph{foreshadowing} su una lezione di Daniele: il
discorso degli \textbf{hack}, delle \textbf{house rule}, delle \textbf{mod},
cioè di come modificare un gioco esistente. Il punto è che \textbf{un gioco
modificato da una house rule è un altro gioco}: cambiando le regole cambiano
necessariamente anche le \textbf{dinamiche} che ne derivano, e quindi
l'esperienza. Si immagini \emph{Monopoly} in cui ogni turno si possa, volendo,
prendere gratis dei soldi dalla banca: tutto il resto è identico, eppure il gioco
cambia.

E questo è talmente interiorizzato che ci sono regole che si applicano ``senza
neanche rendercene conto''. Ruggero chiede quanto duri in media una partita a
\emph{Monopoly}. Gianmarco: ``quattro ore''; altri: ``un pomeriggio''; un altro:
``finché uno dei giocatori non aggiunge la regola: ora vince chi\ldots''. In
realtà, \textbf{seguendo le regole scritte}, una partita non dura mai più di
\textbf{un'ora e venti, un'ora e mezza}. La domanda che fa sempre, e la risposta
è sempre la stessa: perché si pensa che \emph{Monopoly} duri tantissimo? Qual è
la regola che nessuno rispetta?

Uno studente indovina: \textbf{le aste}. Secondo il regolamento, se chi finisce
su un terreno non lo compra, il terreno \textbf{va all'asta} tra i giocatori, e
chiunque può aggiudicarselo. Molti invece giocano come se finire su un terreno
desse solo la \emph{facoltà} di comprarlo o di metterlo all'asta\footnote{Più
precisamente, nel regolamento ufficiale, se il giocatore che finisce su una
proprietà libera decide di non comprarla, la banca la mette all'asta e tutti i
giocatori possono fare offerte, compreso chi l'ha rifiutata. L'asta quindi non è
una scelta del giocatore, ma la conseguenza obbligata del suo rifiuto.}. Con la
regola applicata, di aste se ne fanno moltissime: all'ottavo turno ci sono già
otto terreni assegnati, al ventesimo venti, la gente va in bancarotta
velocissimamente e la partita dura un'ora e mezza. Anche le partite competitive
di \emph{Monopoly}, nota Ruggero, durano un'ora e venti, un'ora e mezza.

Uno studente obietta che giocava con le aste e durava comunque quattro ore.
``Avete dei nemici con l'analisi paralisi'', ride Ruggero; e cita un'altra house
rule tipica dell'infanzia: genitori che danno soldi per non far andare in
bancarotta. Ma \textbf{la bancarotta è parte del gioco}: se la si toglie,
\emph{Monopoly} non regge. Sul tema che \emph{Monopoly} vuole trattare, e sul
suo antenato (\emph{The Landlord's Game}), ``non apriamo adesso un altro vaso di
Pandora''.

\paragraph{Approfondimento: Monopoly, la durata e le house rule nei libri.}
L'esempio di Ruggero si ritrova quasi identico in \emph{The Art of Game Design}.
Parlando del bilanciamento tra partite brevi e lunghe (\emph{Short vs. Long}),
Schell scrive che \emph{Monopoly}, giocato secondo le regole ufficiali,
``finisce spesso in circa novanta minuti'', ma molti giocatori trovano quelle
regole troppo dure e le modificano, distribuendo premi in denaro e allentando i
vincoli sull'acquisto delle proprietà: con l'effetto collaterale di far durare la
partita molto di più, ``tipicamente tre ore, o anche di più''. Schell classifica
le \textbf{house rule} come regole create dai giocatori in risposta a un
difetto percepito dopo qualche partita. Anche \emph{Challenges for Game
Designers}, nel capitolo sull'aggiungere e togliere meccaniche, parte dallo
stesso gioco: il \emph{Parcheggio gratuito}, che per le regole originali non fa
succedere nulla, mentre quasi tutti giocano con una house rule che dà un bonus a
chi ci finisce sopra. Le due fonti confermano la lezione: una piccola regola
cambiata cambia le dinamiche, e le dinamiche cambiano l'esperienza.

\subsection{Il fine del game è il play?}

Un altro studente commenta: va bene la distinzione, ma il fine del \emph{game} è
il \emph{play}; quello che si vede è sempre l'effetto. Ruggero è d'accordo: di
solito, se si fa un gioco, lo si fa per farlo giocare, come si scrive un libro
per farlo leggere (anche se esistono, ed ``è già stato esplorato'', libri scritti
per non essere letti e giochi per zero giocatori). Ma la distinzione serve a una
cosa precisa: \textbf{separare, nel playtest, i feedback sul gioco dai feedback
sul gruppo di gioco} e sulle sue dinamiche. Se una persona esce frustratissima da
un playtest, non è necessariamente colpa del gioco; se una persona dice ``questo
è il gioco migliore del mondo'', potrebbe non essere merito del gioco ma del
gruppo. Capire a che cosa si riferisce un feedback serve a capire \textbf{dove
iterare} il prototipo.

Lo studente osserva che l'area è sfumata: alcuni giochi si basano proprio
sull'avere le persone giuste al tavolo. Un gioco di \textbf{negoziazione} con
giocatori abituati ai giochi ``multi-solitario'' diventa ingiocabile; e nei giochi
a \textbf{economia chiusa} (cita \emph{Container}\footnote{La trascrizione cita
anche un secondo titolo non ricostruibile (``Tokyo su Vigimarket'').}) un gruppo
economicamente conservativo, che spende poco e non fa girare l'economia, fa
durare la partita il doppio. Ruggero rilancia: e se il gioco dura il doppio,
questa dinamica cambia la \textbf{percezione di peso} del gioco, le energie spese,
quanto si è stanchi a fine partita. ``Tutto è collegato. \textbf{I giochi sono
ecosistemi di meccaniche}, e cambiandone una, le altre cambiano di
conseguenza''.

\paragraph{La parola vietata, alla prova.} Lo studente aggiunge che ci sono giochi
che ha trovato ``molto sbagliati, ma che mi divertono''. ``Attenzione: non si può
dire \emph{divertimento}!'', lo ferma Ruggero; nelle presentazioni aveva lasciato
correre, ``ma lì non vi avevo ancora letto la regola''. Lo studente ci prova:
``giochi che poi voglio \textbf{rigiocare}'', anche se ``sono evidentemente
sbagliati, rotti''. Ruggero apprezza, ma non vuole anticipare cose che meritano il
loro tempo: se ne parlerà nelle lezioni sui framework e soprattutto in quella di
\textbf{emotion design}, ``forse la mia lezione preferita del corso''. (Lo
studente ammette che ``forse alla fine si gioca per quello, per l'emozione''.)

\subsection{I giochi ``comunicati male''}

L'ultima domanda corposa riprende un'affermazione di inizio lezione: ci sono
giochi belli ma comunicati male. Qualche esempio? ``Mi basta girarmi per
vederne 48'', risponde Ruggero indicando gli scaffali alle sue spalle. Più che
titoli, elenca le \textbf{ragioni}:
\begin{itemize}
  \item \textbf{l'Italia localizza lentamente}. Basta guardare i premi del gioco
        dell'anno in Italia: spesso vanno a giochi usciti due o tre anni prima in
        Germania, Francia, Stati Uniti o Regno Unito, perché localizzare richiede i
        suoi tempi. Il mercato italiano è quindi molto dipendente dall'estero,
        ``molto impaurito'', e spesso aspetta che un gioco vada bene all'estero
        prima di portarlo in Italia; e alcuni giochi ``perdono il treno'';
  \item \textbf{il publisher originale chiude}: come si pubblica un gioco se non
        c'è più l'azienda con cui fare un contratto? A Ruggero è successo più
        volte;
  \item \textbf{il dipartimento commerciale italiano non sa comunicare il gioco}:
        si fa male un'edizione, la stampa vende pochissimo e non viene più
        ristampata;
  \item \textbf{la scala del mercato}: in Italia un gioco da tavolo di peso medio
        (un \emph{eurogame}, un \emph{German}) che vende 1.500 copie ``è tanto'';
        all'estero se ne vendono magari 10.000. Con 10.000 copie in inglese ci
        sono i soldi per una campagna di comunicazione: un trailer, i
        recensori, i tutorial online. In Italia tutto questo è ``ridotto all'osso'',
        e la cosa si riflette su stampa e divulgazione: anche i canali per
        informarsi su un gioco prima di comprarlo sono pochi (se ne parlerà nella
        lezione su ricerca e sviluppo). Un gioco spesso non arriva al mercato non
        perché sia brutto, ma perché l'editore non lo sa comunicare, ne fa poche
        copie e va fuori produzione;
  \item \textbf{il tempismo}: party game progettati nel 2019 e usciti nel 2020, in
        pieno Covid, quando non si poteva giocare al tavolo, sono falliti tutti, e
        gli editori ci hanno perso. ``Non potevano immaginarselo'';
  \item \textbf{i temi caldi}: dopo l'invasione russa dell'Ucraina e la guerra in
        Palestina, alcuni editori non pubblicano o non localizzano giochi
        strategici e di guerra. Ruggero ha lavorato a \textbf{\emph{Status}}, un
        gioco sulla geopolitica che tocca, tra gli altri temi, l'invasione russa
        dell'Ucraina e il genocidio in Palestina, e alcuni editori che di
        solito localizzano i loro giochi (per esempio il partner di riferimento in
        Russia) non lo hanno localizzato per motivi politici.
\end{itemize}
Lo studente voleva capire che cosa voglia dire comunicare male: non metterci
abbastanza energia, o comunicare una cosa per un'altra, ``come il trailer
sbagliato, che poi vai al cinema e ti trovi un altro film''. ``È un mix'', e la
domanda ``era complessa e meritava più spazio''.

\subsection{I giochi di carte collezionabili e i saluti}

L'ultima domanda (``posso rompere le balle un'ultima volta?'') riguarda i
\textbf{giochi di carte collezionabili}: funzionano con le stesse regole dei
giochi da tavolo? Se ne parlerà alla prossima lezione: è proprio il motivo per
cui la lezione sui casi studio dura tre ore, perché ogni genere di gioco ha le
sue dinamiche e le sue regole, e anche il modo di pubblicare cambia molto tra un
TCG, un gioco di ruolo da tavolo, un LARP e un gioco pervasivo.

Ruggero saluta: le email dei tre docenti sono all'inizio delle slide, la
registrazione arriverà, ``ci si vede questo venerdì''.

% ---------------------------------------------------------------------
\section*{Punti chiave del capitolo}
\addcontentsline{toc}{section}{Punti chiave del capitolo}

\begin{itemize}
  \item Il corso di \textbf{Analog Game Design} (seconda edizione) dura circa tre
        mesi, fino a dicembre, ed è tenuto da tre docenti principali:
        \textbf{Daniele} (narrative design, gioco di ruolo, fiction),
        \textbf{Ruggero} (game design, ricerca e sviluppo) e \textbf{Beatrice}
        (giochi pervasivi), con tre ospiti su RNG, development e game production.
  \item Lo scopo è formare \textbf{autori analogici} consapevoli: strumenti da
        usare anche fuori dal corso, non il ``gioco dei sogni''. Il corso è
        \textbf{intersezionale} (nessuna lezione dedicata a un solo tipo di
        gioco), \textbf{cross-generazionale}, \textbf{pratico e adattabile},
        basato sulle \textbf{domande} (nessuna domanda è stupida), sul
        \textbf{feedback} invece che sul voto e sul \emph{fail forward}.
  \item Il \textbf{project work} si fa in tre gruppi da tre, formati dai docenti;
        pitch il 12 ottobre, sviluppo nelle due pause, tutoraggi il sabato,
        consegna il 4 dicembre, presentazione l'11 e giornata finale con gli
        editori il 14. Si valutano documentazione, presentazione,
        autovalutazione, crescita e atteggiamento.
  \item \textbf{Niente intelligenza artificiale} nel corso: i publisher
        rifiutano pitch e prototipi con immagini generate, i regolamenti scritti
        dall'IA descrivono la stessa meccanica in modi diversi, e soprattutto chi
        la usa perde l'occasione di imparare. Un prototipo può essere fatto a
        matita.
  \item Le \textbf{figure dell'industria}: game designer, game developer,
        editor, publisher, sales e marketing, graphic designer, illustratore,
        sensitivity reader (che non è un censore) e producer. Nelle realtà
        piccole una persona copre più ruoli.
  \item Nell'analogico il \textbf{game designer} progetta l'esperienza creando
        le meccaniche; il \textbf{game developer} fa da \textbf{mediatore} con il
        publisher (bilanciamento, modelli matematici, adattamenti). Pretendere
        che il proprio gioco non venga modificato è un \emph{red flag}.
  \item Il \textbf{marketing} non decide che cosa pubblicare ma come; è
        un'azione, ed esiste anche se lo si ignora.
  \item In \textbf{Italia} il gioco è ancora visto come infantile, il mercato
        punta sulla localizzazione e investe poco su giochi e designer interni;
        all'estero (soprattutto dalla Germania in su) ci sono tradizione adulta,
        premi di prestigio e developer interni.
  \item \textbf{Game} è l'artefatto (le regole, il sistema, l'oggetto);
        \textbf{play} è l'atto di giocare. Si fanno giochi, non giocattoli.
        Separarli serve a distinguere i feedback sul gioco da quelli sul gruppo
        e a riconoscere i quattro incroci tra ``progettato bene/male'' e ``mi
        piace/non mi piace''.
  \item \textbf{MDA}: meccaniche (regole scritte), dinamiche (regole in azione),
        estetiche (esperienza, non grafica); nel \textbf{DPE}: design, play,
        experience.
  \item \textbf{Parole vietate} in ogni materiale scritto: \emph{fun},
        \emph{immersion}, \emph{entertainment} (divertimento, immersione,
        intrattenimento). Bisogna trovare parole più precise.
  \item Distinguere \textbf{chi crea, chi gioca e chi facilita}; nessun gioco
        nasce perfetto: \textbf{playtest e iterazione} sono inevitabili, e un
        \textbf{dizionario comune} aiuta a non perdersi.
  \item \textbf{Tabletop Simulator} è il software di settore per giocare (non
        creare videogiochi) i giochi da tavolo in digitale. Primo task: giocarci
        tre o quattro volte entro il 10 ottobre.
  \item Una \textbf{house rule} crea un altro gioco: \emph{Monopoly} con le aste
        applicate come da regolamento dura circa un'ora e mezza.
\end{itemize}

%
%  (Con hyperref conviene aggiungere anche
%  \renewcommand{\theHchapter}{agd.\arabic{chapter}} nella parte, per
%  evitare ancore duplicate con i capitoli 1, 2, ... di Game Design.)
%  Tutte le etichette di questo capitolo hanno il prefisso "a1:"
%  (a1:cap:..., a1:sec:..., a1:fig:..., a1:tab:...). Il capitolo è
%  autonomo: non contiene \ref verso altri capitoli.
%  Le figure sono ricostruzioni schematiche di quanto i docenti spiegano
%  a voce e con le slide, non copie delle slide.
% =====================================================================

\chapter[Introduzione all'Analog Game Design]{Introduzione all'Analog Game
Design: il corso, le figure dell'industria e un primo linguaggio comune}
\label{a1:cap:intro}

Questa è la \textbf{prima lezione} del corso di \textbf{Analog Game Design},
cioè di progettazione di giochi \emph{analogici}: giochi da tavolo, giochi di
carte, giochi di ruolo da tavolo, giochi di ruolo dal vivo (LARP), giochi
\emph{pervasivi} ed esperienze a evento. È un corso \textbf{breve}, di circa tre
mesi (``da qui a dicembre''), separato dal corso di Game Design, e siamo alla sua
\textbf{seconda edizione}. Come avvertono subito i docenti, è ``una giornata di
introduzione'': si spiega come funziona il corso, chi sono i docenti, con quali
metodi e con quale calendario si lavorerà, e ci si presenta a vicenda. ``C'è poco
gioco analogico'', si scusano. Ma nella seconda parte la lezione comincia
comunque a costruire un \textbf{linguaggio comune}: le figure professionali
dell'industria del gioco analogico, la differenza tra \emph{game design} e
\emph{game development}, la situazione del mercato italiano rispetto a quello
estero, e un primo glossario di design (\emph{game} e \emph{play}, il framework
MDA, le parole ``vietate'', il playtest e l'iterazione). Si chiude con lo
strumento di settore per i prototipi digitali dei giochi da tavolo,
\textbf{Tabletop Simulator}, e con una lunga sessione di domande.

La lezione si svolge \textbf{online}, su Zoom, ed è condotta a più voci dai tre
docenti principali del corso: \textbf{Daniele}, narrative designer, che guida la
parte organizzativa; \textbf{Ruggero}, game designer, che guida la parte sui
fondamenti e le domande finali; \textbf{Beatrice}, progettista di giochi
pervasivi, che interviene soprattutto nella prima metà. Gli studenti iscritti
sono \textbf{nove}; uno è assente e un altro arriva a lezione iniziata. Ci sono
una pausa di una decina di minuti (si riprende alle 18:15) e, dopo le 19:15,
un'ultima mezz'ora di domande con il solo Ruggero.

Tre avvertenze di lettura. La prima: la trascrizione è automatica e storpia
molti nomi propri, compresi quelli dei docenti e dell'accademia (``Dive'',
``IVE'', ``Accademia Live'', ``Ai''); dove la ricostruzione è sicura, o molto
probabile, il capitolo usa la forma corretta e lo segnala in nota; dove non lo
è, riporta la forma della trascrizione e lo dice. La seconda: durante la pausa la
trascrizione registra, ripetuta in loop per decine di righe, una frase senza
senso che non ha nulla a che fare con la lezione, e che il capitolo ignora. La
terza: in molti passaggi più docenti parlano uno dopo l'altro senza che la
trascrizione segni il cambio di voce; il capitolo attribuisce le parole a chi le
pronuncia quando il contesto lo rende chiaro (chi parla di sé, chi viene
chiamato per nome, chi ``passa la parola''), e resta generico altrimenti.

Quanto alla data, la lezione si colloca con ogni probabilità \textbf{nella prima
metà di settembre del 2026}\footnote{La data non è mai detta. La si deduce da
alcuni indizi: il form iniziale va compilato ``entro e non oltre il 15
settembre''; i gruppi verranno annunciati il 25 settembre; le date del
calendario (sabato 3 ottobre, lunedì 12 ottobre, sabato 7 e 14 novembre,
venerdì 4 dicembre) cadono nei giorni della settimana indicati soltanto nel
2026; un docente parla di un argomento ``molto complesso nel 2026''; e fa
ancora molto caldo (``le temperature sono ancora indecenti''). È una
deduzione, non un'informazione esplicita.}: i riferimenti a ``oggi'',
``quest'estate'' e ``l'anno scorso'' vanno letti rispetto a quel momento.

Il percorso del capitolo segue quello della lezione:
\begin{enumerate}
  \item l'avvio, i \textbf{docenti} e le \textbf{presentazioni} dei partecipanti
        (sezioni da~\ref{a1:sec:avvio} a~\ref{a1:sec:partecipanti});
  \item i \textbf{materiali} e gli strumenti del corso: il form, la cartella
        condivisa, Discord (sezione~\ref{a1:sec:materiali});
  \item lo \textbf{scopo} e il \textbf{metodo} del corso, le regole di
        comportamento e il carico di lavoro (sezioni~\ref{a1:sec:scopo}
        e~\ref{a1:sec:regole});
  \item il \textbf{project work} e il \textbf{calendario} (sezione~\ref{a1:sec:project});
  \item la regola sull'\textbf{intelligenza artificiale} e la lunga discussione
        che ne segue (sezione~\ref{a1:sec:ia});
  \item le \textbf{figure professionali} dell'industria del gioco analogico,
        con un approfondimento su marketing e sales
        (sezioni~\ref{a1:sec:figure} e~\ref{a1:sec:marketing});
  \item la differenza tra \textbf{game design} e \textbf{game development}
        (sezione~\ref{a1:sec:design-development});
  \item il \textbf{mercato}: Italia ed estero (sezione~\ref{a1:sec:mercato});
  \item i \textbf{fondamenti di design}: perché serve una terminologia comune,
        \emph{game} e \emph{play}, il framework \textbf{MDA} (e il DPE), le
        \textbf{parole vietate}, chi crea, chi gioca e chi facilita, il
        \textbf{playtest} e l'\textbf{iterazione} (sezioni
        da~\ref{a1:sec:terminologia} a~\ref{a1:sec:playtest});
  \item \textbf{Tabletop Simulator} e il primo task (sezione~\ref{a1:sec:tts});
  \item le \textbf{domande finali}: il tema del project work, il bilanciamento,
        i ruoli nei team, la pubblicazione, il gioco di ruolo in Italia, le
        \emph{house rule} di Monopoly, i giochi ``comunicati male''
        (sezioni~\ref{a1:sec:domande-daniele} e~\ref{a1:sec:domande}).
\end{enumerate}

% ---------------------------------------------------------------------
\section{L'avvio della lezione}
\label{a1:sec:avvio}

La lezione comincia con qualche minuto di attesa. ``Siamo in nove in classe'', e
si aspetta che arrivino tutti; uno dei docenti scherza sul fatto che l'unica
persona che arriverà ``sempre in ritardo di venerdì'' è Gianmarco. Mancano
Giordano e Tiziana; Tiziana arriva, mentre la segreteria fa sapere poco dopo che
Giordano non potrà esserci e vedrà la registrazione. Si aspetta anche Beatrice,
la terza docente, che arriva a sua volta dopo pochi minuti.

Nell'attesa, Daniele anticipa il programma della giornata e qualche
informazione pratica:
\begin{itemize}
  \item oggi è una giornata di \textbf{introduzione}: come funziona il corso, chi
        sono i docenti, le presentazioni degli studenti, le \textbf{metodologie}
        e il \textbf{calendario}. Gli studenti avranno ``un fracco di domande'', e
        alla fine della lezione c'è ``proprio uno slot apposta'' per le domande;
  \item le lezioni sono \textbf{registrate}: di solito la registrazione viene
        consegnata dopo circa \textbf{12--24 ore} sul profilo personale della
        piattaforma dell'accademia, che va attivato con la segreteria (Tiziana,
        già studentessa, probabilmente ce l'ha già);
  \item i docenti ``cuore'' del corso sono tre (Daniele, Ruggero e Beatrice),
        ``soprattutto Ruggero''; a loro si aggiungeranno tre docenti ospiti, per
        meno ore, probabilmente una lezione a testa, ``tranne Marta, che avrà due
        lezioni quest'anno''.
\end{itemize}
E una battuta sul caldo: ``spero che voi siate carichi, anche se le temperature
sono ancora indecenti''. ``Carichi moderatamente, per quello che potete''.

Poi Daniele condivide lo schermo (``non solo le slide, ma anche un po' di altre
cose'') e apre ufficialmente il corso: ``benvenuti e benvenute nella seconda
edizione del corso di Analog Game Design. Questa è la prima lezione''. Il
percorso è di circa \textbf{tre mesi}, fino a dicembre, e durante il corso i
docenti cercheranno di far provare agli studenti ``cosa significa fare Analog
Game Design'' con un \textbf{project work}, un progetto di gruppo di cui si
spiegheranno più avanti la struttura e le richieste.

% ---------------------------------------------------------------------
\section{I docenti}
\label{a1:sec:docenti}

``Anzitutto partiamo da delle presentazioni doverose''. Ognuno dei tre docenti
principali si presenta con una slide che riporta i contatti e una selezione dei
lavori fatti.

\subsection{Daniele: narrativa, fiction e gioco di ruolo}

\textbf{Daniele}\footnote{La trascrizione riporta ``Daniel Cusetto''; nel
resto della lezione tutti lo chiamano ``Daniele''. Il cognome non è
ricostruibile con certezza.} ha già conosciuto quasi tutti gli studenti, perché
ha fatto loro il colloquio di ammissione. È \textbf{narrative designer}, o meglio
\emph{game narrative designer}; insegna nell'accademia che ospita il
corso\footnote{La trascrizione storpia il nome dell'accademia in molti modi
(``Dive'', ``IVE'', ``Accademia Live'', ``Ai'', ``AIVA''). Con ogni probabilità
si tratta dell'\textbf{AIV}, l'Accademia Italiana Videogiochi, che ha sede a
Roma: lo fanno pensare i corsi triennali di game design e narrative design citati
dai docenti e il fatto che una studentessa racconti di essersi trasferita a Roma.
È una ricostruzione, non un'informazione esplicita; nel seguito il capitolo
parla semplicemente dell'``accademia''.} ma lavora soprattutto come \textbf{brand
advisor}. Nel corso si occuperà della parte di \textbf{gioco di ruolo},
\textbf{narrazione} e \textbf{fiction}.

I suoi contatti, ``in quest'ordine'', sono email, Telegram e Discord. Ci tiene a
dare alcune indicazioni precise:
\begin{itemize}
  \item se si vuole avere la certezza di una risposta, bisogna scrivergli su
        \textbf{Telegram}, perché lì ha tutti i suoi clienti;
  \item Discord lo guarda di meno; l'\textbf{email} è per le cose ufficiali;
  \item non bisogna contattarlo altrove: non su Facebook, non su LinkedIn, non su
        Instagram, ``perché li vedo molto di meno''.
\end{itemize}
La slide mostra anche alcuni giochi, attività e brand su cui ha lavorato negli
anni, ma non si dilunga: ``nel caso ne parleremo''.

\subsection{Ruggero: game design e ricerca e sviluppo}

\textbf{Ruggero} (spesso chiamato ``Ruger'') ha fatto il \textbf{game designer}
per varie aziende e oggi è game designer e si occupa di \textbf{ricerca e
sviluppo} per un'azienda del settore\footnote{La trascrizione dice ``per MS'';
il nome dell'azienda non è ricostruibile con certezza. Più avanti Ruggero cita
tra i giochi su cui ha lavorato il gioco da tavolo ufficiale di
\emph{Cyberpunk}, \emph{Marvel United} e \emph{Status}, un gioco sulla
geopolitica.}. Ha un background ``misto'': nel suo portfolio ci sono giochi da
tavolo, giochi di ruolo, giochi analogici, LARP, videogiochi, giochi in realtà
virtuale. Si occupa anche di \textbf{divulgazione} sul gioco e di analisi di
\textbf{game studies} e game design. Più avanti, nelle domande finali, dirà di
avere ``50 e più giochi prodotti e progettati'' e di insegnare da vari anni.

Lo si può contattare liberamente via Gmail e Telegram, e cerca di rispondere
\textbf{entro 24 ore}, salvo imprevisti o eventi: ``sentitevi in libertà di
scrivermi se avete dubbi, senza problemi''. Quando arriva a parlare dei propri
giochi, la slide mostra ``un mix di varie cose, dalle più famose alle meno
famose'', ma Ruggero si imbarazza a presentarli, tra le risate dei colleghi.

\subsection{Beatrice: i giochi pervasivi}

\textbf{Beatrice} (``Bea'') si occupa di \textbf{giochi pervasivi}
(\emph{pervasive games}). ``Se non li avete mai sentiti nominare, non
preoccupatevi'': questo corso è ``uno dei pochi, forse l'unico in Italia'' che
li racconta esplicitamente. Li si vedrà nelle prossime lezioni, ma la definizione
di base è questa: i giochi pervasivi \textbf{utilizzano gli spazi e gli oggetti
della vita quotidiana e li trasformano in spazi e oggetti di gioco}.

Beatrice lavora come \textbf{freelance} e si occupa di diversi tipi di
progettazione: lavora sia con le \textbf{aziende} sia con gli \textbf{eventi dal
vivo}. La sua attività diretta si trova sotto l'etichetta \textbf{Minime
Differenze}\footnote{La trascrizione riporta il nome dell'etichetta come
``Minime Differenze'' e l'indirizzo del sito come ``mininedifferenze.com''; la
grafia esatta del sito non è verificabile dalla trascrizione.}, mentre i lavori
per le aziende sono ``sparpagliati'' sui siti delle varie aziende. Molte delle
cose mostrate nella sua slide sono attività proposte al pubblico ``con un
biglietto, con un gioco in formato evento'': un tema che si toccherà nel corso,
per capire come \textbf{comunicare} i propri giochi in modo che siano
comprensibili, accessibili per il pubblico e \textbf{vendibili}.

Ammette un ``grosso limite'' con i messaggi: è un po' lenta a rispondere, ``ma
rispondo sempre a tutto''. La si può contattare per email (cognome punto nome) o
su Telegram. È ``molto curiosa di vedere cosa produrrete''.

\subsection{I docenti ospiti}

Daniele presenta poi i tre docenti che il corso avrà ``solo per alcune
lezioni'', tutte verso \textbf{novembre}:
\begin{itemize}
  \item \textbf{Diego Ricchiuti}, direttore del corso di Game Design
        dell'accademia, farà una lezione sull'\textbf{RNG} (\emph{random number
        generation}, la casualità nei giochi): ``vi parleremo di RNG quando sarà
        necessario. Comunque avrete una lezione con lui, sappiatelo'';
  \item \textbf{Marta}\footnote{La trascrizione riporta il cognome come
        ``Cercasassi''; la grafia non è verificabile.} farà \textbf{due lezioni
        sul development}: quest'anno ``abbiamo deciso di fare doppietta'',
        perché l'anno scorso il suo intervento è stato ``molto, molto utile'';
  \item \textbf{Flavio}\footnote{La trascrizione riporta il cognome come
        ``Mortarino'' e, più avanti, ``Mortalino'', e il nome della sua casa
        editrice come ``Grampibert''; nessuno dei due è ricostruibile con
        certezza.} farà la lezione di \textbf{game production}; è un editore, e
        di solito partecipa anche alla tavola rotonda con gli editori della
        giornata finale.
\end{itemize}
Le slide sono già online e gli studenti avranno accesso alla cartella condivisa:
``non vi preoccupate di queste cose, ne parleremo a breve''. Ma prima di parlare
del corso, ``vogliamo conoscere voi''.

% ---------------------------------------------------------------------
\section{Le presentazioni dei partecipanti}
\label{a1:sec:partecipanti}

Le domande per la presentazione sono semplici, da affrontare ``a turno, senza
fretta, con calma'':
\begin{itemize}
  \item nome e occupazione (``se volete dirla'');
  \item perché si vogliono fare giochi analogici, e a che livello: che cosa ha
        portato lo studente al corso;
  \item che tipo di esperienze o conoscenze si hanno già nell'ambito, se si è già
        fatto qualcosa oppure no.
\end{itemize}
``Chi vuole partire senza nessun tipo di ritegno e di vergogne?''. Per far
parlare gli studenti guardandoli, Daniele interrompe la condivisione dello
schermo e ricorda lui le domande quando serve. Le presentazioni sono interessanti
anche per i commenti dei docenti, che anticipano diversi principi del corso.

\subsection{Tiziana}

Tiziana fa l'\textbf{infermiera}, ma nella vita ha fatto molti lavori:
impiegata, insegnante e ``tanti altri lavoretti''. Gioca ``da sempre'': il gioco
fa parte anche della sua vita familiare, perché suo fratello fa il programmatore
di videogiochi. Ha già \textbf{organizzato} giochi: una caccia al tesoro per far
conoscere il territorio, due scene con un delitto (giochi con omicidio da
risolvere), giochi per intrattenere le persone. Ha riscoperto di recente i giochi
da tavolo, arrivando a \emph{Carcassonne} con famiglia e amici; poiché è
difficile trovare persone con voglia di giocare nella propria cerchia ristretta,
ha scoperto il mondo delle \textbf{associazioni} ludiche, ``e lì è stato
un'esplosione'': prima in Veneto (è di Vicenza), poi a Roma, dove si è
trasferita.

Ha in mente un progetto ambizioso: far conoscere una certa realtà con la modalità
del \textbf{gioco di ruolo}, o comunque proporre un gioco suo.

\paragraph{Il commento: gli strumenti, non il gioco dei sogni.} Ruggero coglie
l'occasione per un'avvertenza che vale per tutti, visto che molti studenti
arrivano con idee già pronte. Nel corso si lavorerà a un progetto, ma
l'obiettivo dei docenti è dare \textbf{strumenti} da applicare anche
\emph{fuori} dal progetto, ai giochi che gli studenti vorranno fare un domani:
``che non siano esercizi per il corso, ma per voi in primis''. Quindi: ``non
rimaneteci male se non riuscirete a lavorare al vostro gioco dei sogni in questo
corso''. Anche perché ``sarebbe impossibile, in due mesi, fare il gioco dei
sogni''. Il corso darà invece la possibilità di fare un'esperienza ludica
insieme e di imparare a replicarla sulle proprie idee: ``le idee non buttatele,
anzi, il contrario: tenetele''.

\subsection{Nicola}

Nicola è uno \textbf{studente}. Sei mesi prima ha fatto un colloquio con Daniele
perché stava valutando di passare all'accademia, e ha deciso, ``nel dubbio'', di
cominciare da questo corso breve per capire come funziona la formazione;
eventualmente si iscriverà al corso di game design vero e proprio. Gli interessa
in particolare il lato dell'esperienza del giocatore, su cui è ``molto critico''
(qui usa la parola ``divertimento'', e Daniele lo ferma con un ``ne parleremo'':
è una delle parole che il corso vieta, come si vedrà nella
sezione~\ref{a1:sec:parole-vietate}).

Ha 21 anni ed è \textbf{giocatore di ruolo da circa otto}. Prima ancora di
conoscere i giochi di ruolo si era creato un sistema tutto suo, ``pessimo,
disgustoso, terribile'', un sistema a percentuale con il d100 ``terrificante'';
poi ha scoperto i giochi di ruolo da tavolo, i classici (\emph{Dungeons \&
Dragons}, \emph{Daggerheart}, \emph{Vampire}\footnote{La trascrizione riporta
``DD'', ``Daggerart'' e ``Vampire'': quasi certamente \emph{Dungeons \&
Dragons}, \emph{Daggerheart} e \emph{Vampire: The Masquerade}.}). Gioca anche a
\textbf{card game} e ad altri giochi sociali, pur avendo sempre avuto difficoltà
a trovare gruppi, ``avendo vissuto in un buco dimenticato dal mondo''. Come
giocatore è ``l'eterno narratore'' del gruppo, quello che prepara la storia e si
costruisce le regole. Come autore ha fatto soprattutto giochi su carta, ``perché
sono incapace a programmare, per adesso''.

Daniele osserva che il tema del giocatore e del suo rapporto con il gioco è ``un
argomento importante'', anzi ``un vaso di Pandora enorme'' (espressione che
tornerà molte volte nella lezione), e che il corso tratterà anche che cos'è
l'analog game design in generale e quanti tipi di prodotti esistono.

\subsection{Rita}

Rita\footnote{Daniele la chiama prima ``Nino'' (probabilmente un errore della
trascrizione) e poi la ringrazia come Rita.} ha 30 anni e si collega da una
regione del Sud, dove ``ancora si muore di caldo''. Ha studiato
all'\textbf{Accademia di Belle Arti di Napoli}, \textbf{design della
comunicazione visiva}: si occupa dell'aspetto più imprenditoriale della
progettazione artistica. Non volendo progettare per un'impresa (``fare loghi,
robe del genere''), vorrebbe applicare la propria creatività alla realizzazione
di giochi analogici.

Non è una giocatrice di ruolo e non lo è mai stata: è una giocatrice
\textbf{tradizionale}, cresciuta con Monopoly, la dama e giochi tradizionali
cinesi. Ha già \textbf{idee, prototipi e progetti} realizzati, e vuole capire
come renderli \textbf{prodotti per la vendita} e se questo è il campo adatto a
lei; magari anche confrontarsi con i colleghi. La sua progettazione è orientata
al \textbf{sociale} e all'\textbf{educativo}: ``responsabilizzare un po' le nuove
generazioni su come sta andando il mondo e su come potrebbe andare, visto che è
tutto nelle nostre mani''.

\paragraph{Il commento: siate curiosi.} ``Pane per i tuoi denti'', le risponde
Daniele, soprattutto per i LARP e i giochi pervasivi: ``vi apriremo un po' la
testa da questo punto di vista''. E aggiunge un consiglio per tutti: ``siate
curiosi'', perché imparare da altri ambiti dà moltissimo. Ognuno è lì per fare
qualcosa di preciso, ``però siate aperti anche ad altre esperienze''.

\subsection{Tancredi}

Tancredi ha 21 anni e non ha ancora grande esperienza lavorativa: qualche
``piccolo progettino'' che non è andato come voleva. Gli interessa il mondo dei
giochi, digitali e non, con un approccio \textbf{educativo} e \textbf{sociale}.
È arrivato soprattutto dai videogiochi, a cui gioca da quando era piccolo; ai
giochi da tavolo si è avvicinato solo nell'ultimo anno, grazie a una realtà
ludica, perché in famiglia ``non siamo molto da giochi''. Ha realizzato qualche
piccolo videogioco con vari engine: prima quelli senza un vero linguaggio di
programmazione, ``più a nodi'', poi qualcosa di più complicato, e ora sta
lavorando a un piccolo prototipo.

\paragraph{Il commento: niente di quello che vedrete è sprecato.} ``Non butterai
via niente di analog game design'', gli risponde Daniele, ``perché il 99 per cento
delle cose che vedremo si ritrovano quasi uno a uno nel mondo dei videogiochi''.
Tanto che nei corsi triennali dell'accademia (anche in quello di narrative
design, dove insegna Ruggero) il primo anno si parte dai \textbf{giochi da tavolo
e di ruolo}, proprio per spiegare meccaniche più complesse che servono poi per i
videogiochi. Ruggero aggiunge che vale anche il contrario: alcuni dei
\textbf{framework} di game design che si vedranno nelle lezioni tecniche vengono
dal videogioco e sono stati riapplicati all'analogico. ``Evitare la
crossmedialità qui è impossibile''. Anzi, tutte le conoscenze ``mediatiche
adiacenti'' (musica, letteratura, cinema) saranno utili.

\subsection{Lorenzo Vergata}

I Lorenzo in classe sono due; il primo a presentarsi è \textbf{Lorenzo
Vergata}, 26 anni, \textbf{dottorando al Politecnico di Milano}. Gioca ai
videogiochi da quando aveva due anni e mezzo, ma ha giocato al suo primo gioco da
tavolo ``credo nove mesi fa''. Gli fa piacere sentire che molti in classe sono
interessati ai giochi \textbf{sociali} ed \textbf{educativi}, un settore che
interessa anche a lui e che ``può essere esplorato molto'', anche con qualche
aiuto dell'intelligenza artificiale per esperienze più ricercate.

Di esperienza nell'analogico dice di avere fatto ``il minimo sindacale''. Ma
racconta che diciotto anni prima, da bambino, era ``più game designer di
adesso'': prendeva un gioco e ne cambiava le regole, anche al biliardino (``no,
ora introduciamo questa regola''). Gli piace sperimentare, creare regole e
vedere che cosa accade in sé e negli altri. Soprattutto vorrebbe ``una chiave di
lettura per poter rimanipolare la realtà'': tre settimane prima è stato a fare la
\textbf{vendemmia in Francia} e ha visto tantissime dinamiche che potrebbero
diventare un gioco da tavolo o un videogioco. Per esempio: finire contro un rovo
pieno di semi che si attaccano ai vestiti lo rallentava nella raccolta
dell'uva, mentre in un altro filare era più rapido. ``Si può fare pure una
dinamica'', e vorrebbe gli strumenti per farne giochi.

Arriva subito la battuta: una ``vendemmia gamificata'' da proporre in Oltrepò
Pavese (``io sono dell'Oltrepò Pavese, quindi devo per forza bere vino'')\footnote{La
trascrizione dice ``nell'altro papavese'' e non chiarisce chi pronunci la
battuta.}. Daniele conferma che quest'anno il corso avrà ``tante cose sul
sociale'': ``ci piace parecchio''.

\subsection{L'altro Lorenzo}

Il secondo Lorenzo (``secondo, ma non meno importante'') è uno \textbf{scrittore
di racconti brevi} e sta scrivendo il suo primo romanzo. Gli è sempre piaciuto
raccontare storie, e da più grande scriverle; gioca ai videogiochi ed è ``un
amante delle storie e della cultura pop''. Negli ultimi anni è entrato nel mondo
del \textbf{gioco di ruolo da tavolo} e vorrebbe unire le due passioni, le
storie e il gioco di ruolo, magari facendone un lavoro; gli interessa anche lo
\emph{strategic writing}. Daniele gli ricorda che ci sarà una lezione proprio
sulla \textbf{fiction}.

\subsection{Angelo}

Angelo si collega senza videocamera dal lavoro: fa il \textbf{birraio} (in un
piccolo birrificio in Brianza, a Bernareggio, come racconterà dopo la pausa) e,
con il caldo, ``vi risparmio l'orrore''. Per molto tempo era convinto che i
giochi da tavolo fossero solo i classici (Monopoly, Risiko) o i party game; poi,
quattro o cinque anni fa, ha ``sbattuto la faccia'' contro giochi davvero
appassionanti, con una specie di narrazione dietro che li rendeva ancora più
coinvolgenti, e non ha più smesso. È uno che ``ha sempre un botto di idee, ma poi
non mette in pratica assolutamente nulla'': cercava un corso che gli desse le
basi e soprattutto lo \textbf{motivasse} a provare a fare qualcosa.

\paragraph{Il commento: bisogna buttarsi.} Daniele racconta che anche l'anno
prima un suo caro amico si era iscritto proprio per questo motivo. Il corso è
breve, ma cerca di portare gli studenti ``da 0 a 100'' e di far realizzare un
\textbf{prodotto} da presentare alla fine. Il punto è che i giochi ``bisogna
farli, bisogna buttarsi''. E, venendo dalla narrazione, Daniele lo paragona allo
scrivere: si scrive e si riscrive, e allo stesso modo si fa \textbf{playtest}, si
rifà, si aggiorna. ``Non è che si fanno così''.

\subsection{Gianmarco (dopo le slide)}

Gianmarco arriva a lezione iniziata e si presenta più tardi, alla fine della
prima parte, dopo che i docenti hanno chiuso le slide organizzative e la
discussione sull'intelligenza artificiale (``Gianmarco, ti giuro, ti facciamo
presentare''). Lavora nel \textbf{team di innovazione} di un'azienda italiana.
Vuole fare giochi perché è una sua passione, che lo gratifica, e in passato ci ha
provato ``in maniera scomposta''. Cerca un \textbf{metodo} e un confronto con
persone più esperte: ``studiare solo dei manuali è molto sterile'', ``stare
chiusi nel pertugio non aiuta'', ed ``essere autodidatta ha estremamente i suoi
limiti''. Come esperienza è un giocatore e ha letto ``un po' di materiale, di
blog, di manuali'', ``niente di rilevante''.

Manca solo Giordano, che si presenterà alla lezione successiva.

\paragraph{Uno sguardo d'insieme.} La classe è molto varia: un'infermiera,
una designer della comunicazione visiva, un birraio, un dottorando, uno
scrittore, un impiegato nell'innovazione, due studenti ventunenni. Alcuni
vengono dal gioco di ruolo, altri dai giochi tradizionali, altri dai
videogiochi; quasi tutti hanno scoperto i giochi da tavolo moderni da pochi anni
o da pochi mesi. E almeno tre studenti dichiarano un interesse per i giochi
\textbf{sociali} ed \textbf{educativi}. Come si vedrà nella
sezione~\ref{a1:sec:scopo}, i docenti considerano questa varietà una risorsa.

% ---------------------------------------------------------------------
\section{I materiali e gli strumenti del corso}
\label{a1:sec:materiali}

Prima delle slide sul corso, Daniele sbriga un po' di ``clean up casalingo''.

\paragraph{Il form iniziale.} C'è un \textbf{form} da compilare ``entro e non
oltre il 15 settembre'': il QR code è nelle slide, e la segreteria lo ricorderà
su Discord. Richiede circa 10--15 minuti; si può fare subito o in giornata, ``e
ve lo togliete di torno''. È un form standardizzato per tutti i corsi (``non
siamo io e Ruggero a farvi passare la tortura del test''), diviso in due parti:
\begin{itemize}
  \item una parte di \textbf{conoscenza del candidato}: quanto tempo si ha a
        disposizione per il corso, come ci si organizza con il lavoro, che cosa
        piace giocare. È la parte che interessa ai docenti, perché li aiuta a
        organizzare le lezioni;
  \item una parte di \textbf{test logico}, che i docenti non hanno scritto (``io
        non sono responsabile del form'', ``no, no, nemmeno io'').
\end{itemize}

\paragraph{La cartella condivisa.} Tutti gli studenti vengono aggiunti alla
cartella del corso su \textbf{Google Drive}, che si trova tra i file
``condivisi con me''. Gli studenti ne sono \textbf{gestori dei contenuti} e
possono caricarci quello che vogliono. Dentro ci sono:
\begin{itemize}
  \item le \textbf{attività in classe};
  \item le \textbf{dispense};
  \item i \textbf{materiali didattici}, con le slide (quelle della lezione sono
        già lì);
  \item le cartelle dei \textbf{project work} futuri;
  \item i \textbf{task}, che ``saranno pochi, ma ci saranno'';
  \item il \textbf{calendario studenti}, ``importantissimo'', con due schede: il
        calendario delle lezioni e quello del project work.
\end{itemize}

\paragraph{Discord.} Il corso ha anche un \textbf{server Discord}, a cui tutti
sono stati invitati; ci sono i docenti (Daniele, Ruggero, Diego e altri), che si
possono taggare. Di solito è lì che la segreteria pubblica le informazioni, per
esempio se una lezione ``dovesse per caso essere cancellata''. Nelle domande
finali Ruggero aggiungerà che c'è un canale apposito per il corso, che gli
studenti possono usare liberamente per organizzarsi tra loro.

\paragraph{Le reaction di Zoom.} Durante le lezioni su Zoom si possono usare le
\textbf{reaction} (il pollice in alto, il cuoricino, la mano alzata): un modo
rapido per far capire ai docenti se si è capito, se si vuole intervenire. Nel
frattempo arrivano i primi pollici.

% ---------------------------------------------------------------------
\section{Lo scopo e il metodo del corso}
\label{a1:sec:scopo}

La prima serie di slide (che Daniele, per un lapsus, presenta come ``cos'è il
corso di pulire il pavimento di casa'') spiega che cosa il corso vuole fare e
come.

\subsection{Diventare autori analogici}

Lo \textbf{scopo del corso} è fornire le competenze per \textbf{diventare autori
analogici} e stimolare la creatività e l'\textbf{autorialità} di ciascuno: ``voi
imparerete a diventare autori di giochi analogici''. La differenza tra
\emph{game designer} e \emph{game developer} verrà chiarita più avanti (sezione~\ref{a1:sec:design-development}),
ma la parola chiave qui è \emph{autorialità}: ``vogliamo voi, vogliamo le vostre
idee''. Magari non si potrà usare l'idea del gioco della propria vita, ma si
potranno portare nelle lezioni le proprie idee e le proprie sensibilità.

Gli obiettivi elencati nella slide sono questi:
\begin{itemize}
  \item dare \textbf{maggiore consapevolezza nel design e nello sviluppo}. In
        Italia, purtroppo, ``si fa design a volte in maniera non consapevole''.
        Consapevole vuol dire avere almeno un'introduzione alle basi del design e
        dello sviluppo, e fare le cose ``con metodo, con raziocinio, con
        consapevolezza'';
  \item fornire gli \textbf{strumenti essenziali} per un design consapevole;
  \item stimolare la \textbf{sperimentazione} e la \textbf{trasversalità
        mediatica}. Non si chiede di fare ``il \emph{Root} della situazione o il
        \emph{Carcassonne} della situazione''\footnote{La trascrizione dice ``il
        lumaven della situazione o il carcassone della situazione''; il primo
        titolo non è ricostruibile con certezza. Il senso è chiaro: non si
        chiede di rifare un grande classico del gioco da tavolo.}, ma un
        project work \textbf{calibrato sulla durata del corso}, e però
        \textbf{sperimentale} e ``assolutamente transmediale'';
  \item portare \textbf{numerosi esempi} di tipologie di gioco, dall'analogico al
        tavolo, al pervasivo ``e oltre''. Lo si vedrà nella lezione sui
        \textbf{casi studio} del venerdì successivo: ``voi non avete idea di
        quante cose esistono, bellissime, sconosciutissime e favolose''. Daniele
        confessa che preparando e aggiornando le slide i colleghi gli hanno
        mostrato cose che lui stesso non sapeva nemmeno che esistessero;
  \item dare gli strumenti per \textbf{analizzare} i giochi altrui. Ruggero lo
        sottolinea: se il corso riesce nel suo intento, gli studenti sapranno
        non solo fare i propri giochi ma anche analizzare giochi già fatti, che è
        ``una delle skill importanti anche nella parte di ricerca e sviluppo''.
        Si vedranno casi studio e giochi pubblicati, sempre con l'obiettivo di
        capire \textbf{che cosa fanno e perché lo fanno}, per poi fare meglio il
        proprio lavoro;
  \item imparare a \textbf{fare giochi insieme}. Esistono gli autori che lavorano
        da soli (i \emph{solo developer}), ``ma sappiate che anche l'autore che
        lavora da solo poi si deve interfacciare con il pubblico, con un editore,
        magari con un editor, con un publisher''. ``Non lavorerete mai da soli'':
        bisogna imparare a lavorare con gli altri.
\end{itemize}

\subsection{Un corso intersezionale}

Il corso è \textbf{intersezionale}: \textbf{non ci sono lezioni dedicate a una
singola tipologia di gioco}. Non ci sarà ``la lezione sul gioco da tavolo'',
``la lezione sul gioco di ruolo'', ``la lezione sul gioco pervasivo''. Ci saranno
lezioni su temi diversi, ciascuna pensata per essere utile a \textbf{tutte e tre
le sfaccettature} del corso; alcune saranno più importanti per il gioco da
tavolo, altre per il pervasivo, altre per il gioco di ruolo, ma nessuna sarà
esclusiva. I docenti vogliono ``contaminazioni e applicazioni assolutamente
trasversali''.

\subsection{Un corso cross-generazionale}

È anche un corso \textbf{cross-generazionale}: i gradi di esperienza di gioco
quest'anno sono abbastanza allineati, ma le esperienze di vita, le età, le
estrazioni sociali sono varie. ``Ed è un pro, non è un contro. Più diversificata
è la classe, più imparate meglio anche voi''. La slide ha una piccola nota per
fasce d'età:
\begin{description}
  \item[Under 30] Concentratevi sulla vostra \textbf{crescita}. Evitate di
        confrontarvi con le persone più grandi che magari ``ce la fanno meglio''
        in alcune cose.
  \item[Over 30] Concentratevi sul \textbf{consolidamento}, ma siate
        \textbf{umili}: i docenti diranno cose che faranno ``storcere il naso''.
        ``Bisogna un po' reimparare per imparare'': bisogna sapersi adattare.
\end{description}
Ruggero contestualizza: lo dicono perché in passato hanno visto studenti che,
essendo più adulti o avendo giocato di più, pensavano che l'esperienza
\emph{di gioco} si trasformasse automaticamente in competenza \emph{di
creazione} di giochi. Non è così. E non bisogna sottovalutare chi gioca meno:
proprio perché gioca cose diverse, per risolvere un problema di design specifico
``è una ricchezza incredibile avere qualcun altro che riesce a uscire dalla
scatola''.

\subsection{Un corso pratico e adattabile}

Il corso è \textbf{pratico}: gli studenti si concentreranno su un progetto di
gruppo, il \textbf{project work}, lavorando in team. Ed è \textbf{adattabile}: il
programma è definito, ma ``si aggiorna sui vostri stimoli, priorità e
necessità''. Se durante il corso i docenti vedono che serve più una cosa
dell'altra, sono ``prontissimi'' a dare più spazio a un argomento invece che a
un altro.

\subsection{Inquiry-based learning: domandate}

Il metodo è quello dell'\textbf{inquiry-based learning}, l'apprendimento basato
sulle domande. Di solito le persone dicono ``ho una domanda, magari sarà
stupida''. ``Non esistono domande stupide. Pensate che ci sia una domanda
stupida? Fatela, alzate la mano''. Vanno bene anche le
\textbf{provocazioni}: se qualcosa non torna, è meglio chiedere. La condivisione
è fondamentale.

I docenti danno \textbf{metodi, dati e conoscenze}, ma lasciano trovare le
risposte agli studenti: ``vi lasciamo un pochettino annegare nel mare prima di
darvi la soluzione, in modo tale che la possiate trovare da soli''. Di
conseguenza \textbf{non ci saranno interrogazioni sulle definizioni}: interessa
vedere \textbf{come} gli studenti applicano ciò che imparano. Se una
definizione non viene in mente, non è un problema; l'importante è saperla usare.

\subsection{Feedback, non valutazione}

``Sì al feedback, no alla valutazione''. I docenti danno \textbf{feedback
verbali}, anche sulle attività da fare a casa (i task, che non saranno tanti).
Il corso è calibrato per \textbf{studenti lavoratori}: c'è un po' da lavorare
fuori dalle lezioni, ma in modo calibrato.

Il principio è \textbf{fallire oggi per imparare per il futuro}: \emph{fail
forward}, \emph{fail to learn}. I docenti lasciano provare e fallire almeno una
volta prima di dare la soluzione: ``il classico metodo Montessori, quello
originale, non quello travisato''.

\subsection{Le soft skill e i criteri di valutazione}

Grande attenzione va alle \textbf{soft skill}: i lavori di gruppo, lo scambio
propositivo, l'atteggiamento positivo, la propositività, la proattività.
``Lavoro di gruppo'' non significa stare ogni giorno in sala studio con gli
altri: va benissimo che un gruppo si organizzi dandosi priorità e tempi in base
alla disponibilità di ciascuno. L'importante è \textbf{collaborare}, essere
\textbf{professionali}, rispettarsi a vicenda, rispettare le idee proprie e
altrui e lasciare spazio agli altri, perché i progetti ``non saranno solo
vostri, saranno di tutte le persone del vostro gruppo''.

La valutazione del project work non riguarda solo la \textbf{documentazione}
prodotta, ma anche:
\begin{itemize}
  \item la \textbf{presentazione} finale;
  \item l'\textbf{autovalutazione};
  \item la \textbf{crescita personale};
  \item l'\textbf{atteggiamento} tenuto nel corso.
\end{itemize}
Daniele cita: ``da ognuno secondo le sue capacità, a ognuno secondo i suoi
bisogni''. Lavorare in gruppo non vuol dire che, in tre, ciascuno lavori
esattamente il 33\% del tempo. Ai docenti serve piuttosto capire \textbf{chi fa
che cosa} nel gruppo, e che il processo interno sia \textbf{trasparente}.

\paragraph{Lavorare in cloud.} Per lo stesso motivo il consiglio è di
\textbf{non lavorare in locale}, ma direttamente nelle cartelle condivise del
project work e dei task. Si è liberi di fare altrimenti, ma lavorando in cloud
si dà ai docenti la possibilità di vedere \emph{come} si lavora, il proprio
``processo mentale''.

% ---------------------------------------------------------------------
\section{Le regole di comportamento e il carico di lavoro}
\label{a1:sec:regole}

Alcune slide sono di quelle che si ripetono in ogni corso (``\emph{repetita
iuvant}'', storpiato per scherzo in ``ripetita Juventus'').

\paragraph{Comportamento.} Prestare attenzione durante le lezioni; interagire;
portare rispetto; essere proattivi; alzare la mano e fare domande.
\textbf{Tolleranza zero} per pregiudizi, maleducazione, mancanza di rispetto,
\emph{cheat} e scorciatoie. Daniele ammette che ``cheat e scorciatoie'' viene da
altri corsi dove il tema è più importante, ma Ruggero lo ferma: ``non
sottovalutate questa cosa'', perché implica anche \textbf{non far fare a ChatGPT
i task}. Daniele ha una slide apposta, ``perché è un discorso molto complesso nel
2026'' (sezione~\ref{a1:sec:ia}).

\paragraph{Le lezioni sono intense.} Le lezioni possono essere \textbf{dense},
``soprattutto le mie, sappiatelo: io sono logorroico'', avverte Daniele. Bisogna
chiedere spiegazioni se serve, prendere appunti, rivedere le registrazioni. Le
lezioni sono ``delle scintille'': modi per far esplorare i temi oltre le slide.
``Andate oltre, se riuscite''.

\paragraph{Il carico di lavoro.} Le regole sui task sono queste:
\begin{itemize}
  \item i task \textbf{sono per voi, non per noi};
  \item se si ha poco tempo, è meglio farli rapidamente piuttosto che non farli:
        ``provateci'';
  \item i docenti cercano di non chiedere più di \textbf{un'ora o due a
        settimana} oltre alle ore di lezione;
  \item non tutte le lezioni avranno un task, e il \textbf{project work} è il più
        importante di tutti.
\end{itemize}

\paragraph{I feedback dei docenti.} I feedback ``si guadagnano'' man mano che si
reagisce bene: non bisogna aspettarsi ``dei papiri'' sui task. Se i docenti
hanno voglia di dare feedback più ampi lo fanno, soprattutto se vedono che lo
studente se li merita. Ad Analog di solito non ci sono task \emph{rejected}; ma
se un docente rifiuta un task, vuol dire che si è fatto \textbf{meno del livello
minimo} richiesto dal corso. I feedback possono non essere immediati: ``siate un
pochettino impietositi dal fatto che siamo dei docenti''. E se un feedback non
arriva, o se ne vogliono di più, bisogna scrivere, ``perché siamo tre esseri
umani e quindi possiamo anche dimenticarci''.

% ---------------------------------------------------------------------
\section{Il project work e il calendario}
\label{a1:sec:project}

\subsection{Come funziona il project work}

Una prima slide con le date del project work è da aggiornare (``perdonatemi il
refuso''); Daniele la corregge al volo e rimanda alla slide successiva, più
precisa. Il funzionamento è questo:
\begin{itemize}
  \item i \textbf{gruppi li formano i docenti}. Gli studenti sono nove, quindi si
        faranno \textbf{tre gruppi da tre}: una dimensione pensata anche per il
        fatto che molti sono studenti lavoratori, così da non ritrovarsi ``in due,
        a esplodere'';
  \item ogni gruppo deve prima preparare un \textbf{pitch} del proprio gioco, da
        presentare in anticipo: un'idea, una bozza del gioco;
  \item poi si lavora al gioco \textbf{nelle pause del corso}: da qui a dicembre
        ci saranno due pause, quindi non ci sarà lezione tutti i lunedì e tutti i
        venerdì;
  \item ci saranno \textbf{review intermedie} con i \textbf{tutoraggi}, per lo più
        il \textbf{sabato mattina}, dalle 11 alle 12, con Ruggero su Google Meet.
        Servono ad aiutare i gruppi a portare avanti il project work.
\end{itemize}

\subsection{Il calendario delle lezioni}

Il \textbf{calendario} è nel calendario studenti della cartella condivisa.
Attenzione: le date \textbf{possono variare}. ``Il 90\% rimarranno così'', ma
può capitare che, per motivi di lavoro, qualcuno stia male o non possa, o che
``scompaia l'elettricità da tutte le città d'Italia''. E l'ordine delle lezioni
può cambiare se i docenti vedono che agli studenti serve prima una lezione
invece di un'altra.

Ci sono \textbf{due pause importantissime}:
\begin{itemize}
  \item dal \textbf{13 ottobre al 5 novembre}: è il periodo delle grandi fiere
        (\textbf{Essen} e \textbf{Lucca}\footnote{La trascrizione dice ``Essen,
        Luca, Cose'': si tratta con ogni probabilità della fiera Spiel di Essen,
        in Germania, e di Lucca Comics \& Games, che si tengono entrambe in
        ottobre.}), ``un periodo un po' frenetico'', e cominciano anche i corsi
        accademici, quindi la pausa permette ai docenti di respirare. Per gli
        studenti è la pausa più importante: quasi un mese in cui sviluppare la
        parte importante del gioco;
  \item dal \textbf{24 novembre all'11 dicembre}, dopo qualche lezione a novembre.
\end{itemize}

E ci sono \textbf{due date da evidenziare}:
\begin{itemize}
  \item il \textbf{25 settembre}: nei primi 5--10 minuti della lezione (che per il
        resto sarà tenuta da Beatrice) i docenti annunceranno i tre gruppi;
  \item il \textbf{12 ottobre}: il \textbf{pitch} dei giochi, con l'idea
        embrionale. Tra il 25 settembre e il 12 ottobre ogni gruppo prepara un
        breve pitch in formato slide; quel giorno i docenti lo rivedono e danno
        feedback: se il gioco va bene, se non va bene, come svilupparlo.
\end{itemize}
Dal 12 ottobre fino all'inizio di dicembre si sviluppa il gioco, con i
tutoraggi: ``sfruttate questi momenti di pausa''.

\subsection{Il calendario del project work}

La slide successiva, ``reale e aggiornata'', riporta il calendario del project
work (tabella~\ref{a1:tab:calendario} e figura~\ref{a1:fig:calendario}).
``Per ora non pensate al project work: ci penserete dal 25 settembre in poi''.

\begin{table}[htbp]
\centering
\begin{tabular}{@{}lp{9.2cm}@{}}
\toprule
\textbf{Data} & \textbf{Che cosa succede} \\
\midrule
ven 25 settembre\footnotemark & lezione: annuncio dei gruppi e del
\emph{mandate}, cioè di che cosa si dovrà fare \\
sab 3 ottobre, ore 11 & prima review e tutoraggio con Ruggero (dopo una
decina di giorni di lavoro) \\
lun 12 ottobre & pitch del gioco durante la lezione \\
13 ottobre -- 5 novembre & prima pausa: sviluppo del gioco \\
sab 7 novembre & tutoraggio \\
sab 14 novembre & tutoraggio \\
24 novembre -- 11 dicembre & seconda pausa: sviluppo del gioco \\
ven 4 dicembre, entro mezzanotte & consegna del gioco sul Drive (non c'è
lezione) \\
ven 11 dicembre & i giochi vengono visti a lezione \\
14 dicembre & giornata finale, con la tavola rotonda con gli editori \\
\bottomrule
\end{tabular}
\caption{Il calendario del project work come lo presenta Daniele. Le date
possono variare.}
\label{a1:tab:calendario}
\end{table}
\footnotetext{Daniele dice ``lunedì 25 settembre'', ma nel 2026 il 25 settembre
è un venerdì; il giorno della settimana nella tabella è corretto di
conseguenza.}

\paragraph{Perché la consegna è il 4 dicembre.} Il 4 dicembre non c'è lezione:
si deve solo consegnare il gioco, entro mezzanotte, sul Drive. La presentazione
a lezione è l'11 dicembre, ma i tre docenti hanno bisogno di \textbf{una
settimana} per rivedere il gioco in privato e valutare anche la
\textbf{documentazione}; per questo la consegna non coincide con la
presentazione. Che cosa voglia dire esattamente ``consegnare il gioco'' verrà
spiegato più avanti.

\paragraph{La giornata finale.} Il 14 dicembre c'è la \textbf{giornata finale},
con una \textbf{tavola rotonda con gli editori}: di solito viene invitato Flavio,
che è anche docente di game production, più qualche altro editore, ``vediamo
quest'anno chi riusciamo a portarvi''.

\begin{figure}[htbp]
\centering
\begin{tikzpicture}[
    font=\footnotesize,
    x=0.1cm,
    evento/.style={draw, thick, rounded corners=2pt, fill=white,
                   text width=1.8cm, align=center, inner sep=2pt},
    pausa/.style={fill=gray!20, draw=gray!60},
    tacca/.style={thick}
  ]
  % asse: giorno 0 = 25 settembre; 1 unità = 1 giorno
  \draw[pausa] (18,-0.35) rectangle (41,0.35);
  \draw[pausa] (60,-0.35) rectangle (77,0.35);
  \node[font=\scriptsize\itshape] at (29.5,0.18) {pausa};
  \node[font=\scriptsize\itshape] at (65,0.18) {pausa};
  \draw[-{Stealth[length=2.5mm]}, thick] (-3,0) -- (85,0);
  % eventi sopra
  \draw[tacca] (0,-0.15) -- (0,0.9);
  \node[evento, above] at (0,0.9) {25 set\\ gruppi e \emph{mandate}};
  \draw[tacca] (17,-0.15) -- (17,2.2);
  \node[evento, above, fill=gray!8] at (17,2.2) {12 ott\\ \textbf{pitch}};
  \draw[tacca] (70,-0.15) -- (70,0.9);
  \node[evento, above, fill=gray!8] at (70,0.9) {4 dic\\ \textbf{consegna}};
  % eventi sotto
  \draw[tacca] (8,0.15) -- (8,-0.9);
  \node[evento, below] at (8,-0.9) {3 ott\\ tutoraggio};
  \draw[tacca] (43,0.15) -- (43,-0.7);
  \draw[tacca] (50,0.15) -- (50,-0.7);
  \draw (43,-0.7) -- (50,-0.7);
  \draw[tacca] (46.5,-0.7) -- (46.5,-0.9);
  \node[evento, below] at (46.5,-0.9) {7 e 14 nov\\ tutoraggi};
  \draw[tacca] (77,0.15) -- (77,-0.9);
  \node[evento, below] at (77,-0.9) {11 dic\\ giochi a lezione};
  \draw[tacca] (80,-0.15) -- (80,2.2);
  \node[evento, above] at (80,2.2) {14 dic\\ giornata finale};
\end{tikzpicture}
\caption{Una linea del tempo del project work, ricostruita dalle slide (scala
approssimativa: dal 25 settembre al 14 dicembre). Le bande grigie sono le due
pause senza lezioni, in cui i gruppi sviluppano il gioco.}
\label{a1:fig:calendario}
\end{figure}

Finita questa parte, Daniele apre alle domande: ``Domande sul calendario? Sulle
slide che ho fatto? Considerazioni? Non abbiate timore, io non ci credo''. E la
domanda arriva, sull'unico tema che aveva lasciato da parte.

% ---------------------------------------------------------------------
\section{L'intelligenza artificiale nel corso}
\label{a1:sec:ia}

\subsection{La slide di Daniele}

Daniele aveva già presentato una slide che definisce ``importantissima''. ``Lo so
che è paradossale, me ne rendo conto. Lo so che magari mi darete del luddista,
dell'anti-tecnologia: non è così. Uso l'intelligenza artificiale per lavoro. Ma
in questo corso, vi prego, \textbf{non usate l'intelligenza artificiale}''. La
richiesta vale per le attività, per i task e \textbf{anche per la
prototipazione}. Le ragioni:
\begin{itemize}
  \item ``voglio che siate voi a fare il gioco, voglio che siate voi a imparare a
        farlo''. Daniele parla da utilizzatore: prima di ChatGPT, che ha smesso
        di usare, oggi soprattutto di Claude, di cui ``si sottovaluta quanto
        possa essere incerto, insicuro, impreciso, non all'altezza'';
  \item nella game industry, \textbf{i publisher rifiutano automaticamente} i
        pitch e i prototipi fatti con l'intelligenza artificiale. Daniele lo ha
        visto di persona in ``tantissime pitch review'' quell'estate, anche al
        \textbf{First Playable} di Firenze\footnote{La trascrizione dice ``First
        Playboy a Firenze'': si tratta del First Playable, l'evento di business
        dell'industria videoludica che si tiene a Firenze.}: anche nel mondo dei
        videogiochi i publisher che vedono pitch o prototipi fatti con l'IA
        ``rigettano automaticamente la build'', perché hanno ancora timore
        dell'IA e delle sue possibili conseguenze;
  \item ``non so come sarà il futuro'': non sa se bisognerà abituarsi a usarla
        tutto il giorno. Non è ``totalmente contrario'', purché sia usata ``con
        raziocinio''; secondo lui oggi, negli Stati Uniti, non lo è;
  \item ma in un corso che si chiama \emph{Analog Game Design} gli studenti devono
        saper fare i giochi da soli. In futuro, se vorranno farsi aiutare
        dall'IA per essere più efficienti, saranno liberissimi di farlo; durante
        il corso no, perché i docenti vogliono valutare \textbf{le loro
        competenze e capacità}.
\end{itemize}

\subsection{La domanda: che cosa vuol dire usare l'IA in un gioco analogico?}

La domanda di uno studente è precisa: che cosa si intende per ``task fatti con
l'IA''? Qual è l'uso dell'IA in un gioco analogico?

\paragraph{Ruggero: l'esperienza di chi valuta i prototipi.} Ruggero risponde
dalla sua esperienza nel settore, dove riceve prototipi da valutare per
un'eventuale pubblicazione. Distingue due usi:
\begin{enumerate}
  \item le \textbf{immagini generate}: illustrazioni, graphic design e così via.
        Per molti publisher è ``un \emph{reject} immediato, proprio immediato'';
  \item l'IA usata per il \textbf{regolamento}, per scrivere le regole e le
        meccaniche. È più difficile da individuare, ma molti publisher se ne
        accorgono comunque, perché le IA ``non sono ancora brave abbastanza da
        creare un regolamento coeso meccanicamente e scritto con un senso per il
        vostro gioco''. Quello che fanno è prendere parti di altri regolamenti,
        elaborare le strutture che vedono nei regolamenti pubblicati e
        riapplicarle. Ma molti giochi si prestano poco a un regolamento
        strutturato come quello di altri, perché \textbf{ogni gioco è un
        ecosistema meccanico a sé}. Il risultato sono regolamenti con
        riferimenti non immediati, e soprattutto con \textbf{la stessa meccanica
        descritta in modo diverso in pagine diverse}: due difetti che fanno
        perdere tempo all'editore, e in questo caso ai docenti.
\end{enumerate}
Poi aggiunge un punto ``personale, soggettivo'': i docenti sono interessati a
vedere il lavoro, la creatività, ``il frutto della vostra creazione''. E
viceversa gli studenti vogliono i commenti dei docenti: ``immaginatevi se voi
fate un gioco e noi vi diamo i commenti scritti con ChatGPT''. Altrimenti, dice
ridendo, gli studenti fanno il gioco con ChatGPT, i docenti lo danno a ChatGPT
per i commenti, ``e facciamo molto alla svelta''.

\paragraph{Daniele: perché gli studenti usano l'IA.} Daniele racconta la sua
esperienza con gli studenti, che secondo lui ricorrono all'IA in due momenti:
\begin{enumerate}
  \item quando vanno in \textbf{crisi}: in ansia da prestazione o per i tempi.
        Cercano di risolvere il problema della propria disorganizzazione
        facendolo risolvere all'IA, senza accorgersi che l'output è ``molto,
        molto scarso'';
  \item quando hanno l'impressione che l'intelligenza dell'IA sia
        \textbf{migliore} di quello che scrivono loro, o di quello che c'è in
        giro. Per Daniele è una ``mancanza di palato critico'': l'output dell'IA
        sembra buono in superficie, ma quando si va nel dettaglio e ci si fa
        delle domande ``di solito fa un po' di \emph{crumbling}'', si sbriciola.
\end{enumerate}
Lui usa Claude ogni giorno e ne conosce i limiti: ``non c'è nemmeno una frase
che io copio direttamente''. Gli fa fare il lavoro più noioso, più ripetitivo e
soprattutto \textbf{strutturato}, ``replicare delle cose come se fosse una
macro''. Lo usa per il testo e per il codice, mai per le immagini: per motivi di
\textbf{diritto d'autore}, e perché ha molti amici artisti ed è stufo di vedere
il loro lavoro ``dilaniato'' dall'IA.

Ammette che i docenti non possono controllare gli studenti; per questo
``ve lo consiglio''. E riprende quella che definisce la cosa più bella detta da
Ruggero (``ti ruberò la frase per il corso di narrative''): i docenti i feedback
li danno di persona. Daniele a volte manda 30 minuti di audio a uno studente, o
fa intere call per rivedere script, pitch e giochi: ``è il mio lavoro, mi piace
farlo e voglio farlo io, quindi mi aspetto che lo studente lo faccia lui''. Come
freelance, poi, è ``stufo marcio'' dei clienti che arrivano con qualcosa fatto
con l'IA (``guarda che bello, pubblica questo''), perché il più delle volte il
risultato è pessimo: non sanno come usarla. ``È anche una questione di
competenza''.

\paragraph{Il prototipo non deve essere bello.} Questo vale anche per il
\textbf{prototipo}, su cui si tornerà nelle lezioni di prototipazione: in un
prototipo non servono ``graficoni allucinanti''. Le grafiche bellissime fatte con
l'IA spesso \textbf{peccano di user experience}: sembrano bellissime, ma quando
il giocatore deve leggere la carta non capisce qual è il costo, qual è il
potere, qual è l'effetto.

\paragraph{Beatrice: si perde l'occasione di imparare.} Beatrice aggiunge due
motivi che valgono ``al di là dei nostri credo''. Il primo: chi usa l'IA in
sostituzione (per esempio per scrivere un testo, per accelerare i tempi o per
trovare strategie) \textbf{perde la parte in cui può imparare}. Nel game design
bisogna ``tenere insieme i pezzi'' di un regolamento, ma anche dell'esperienza
dell'utente e di come trasmetterla; i passaggi necessari (creare il regolamento,
creare le meccaniche, sperimentare, raccogliere i feedback) ``si capiscono solo
vivendoli''. E si sta pagando un corso per imparare. Il secondo, sul piano
editoriale: chi presenta a un editore un prodotto che ha come base l'IA rischia
da un lato di non avere imparato abbastanza (e di ritrovarsi ``a piedi'' sul
campo senza IA), dall'altro che l'editore se ne accorga, ed è un tema
\textbf{reputazionale}. Ci sono editori ``più illuminati di altri'', quindi
anche qui c'è un mercato da capire; ma i docenti preferiscono dire che cosa è
più sensato per il mercato come lo conoscono.

\paragraph{Ruggero: non abbiate paura del giudizio.} Ruggero vuole dare uno
``spotlight'' a un punto detto da Daniele e rimasto ``semi-ignorato'': a volte
l'uso dell'IA nasce dalla \textbf{paura del giudizio}, dal voler consegnare
qualcosa di più pulito, più strutturato possibile. ``Vi giuro, non ve ne
preoccupate''. Ruggero detesta i voti e non ragiona per voti; non valuterà
nessuno per grafiche fatte a matita. Insegna da anni, ha avuto studenti con
dislessia e discalculia: refusi ed errori non gli interessano, ``non è un corso
di editing o di proofreading''. Gli interessa che le \textbf{meccaniche siano
coerenti con lo scopo di design}. Immagini fatte a matita o con Paint ``non sono
un problema, anzi, le preferisco'' rispetto all'``ennesima immagine fatta con
ChatGPT o Midjourney''.

\paragraph{Daniele: il pensiero critico, costo e qualità.} L'ultimo intervento
è di Daniele, sul \textbf{pensiero critico}. Ragionando per assurdo (``un
pensiero di Belzebù''), immagina che un domani l'IA venga usata ovunque. I suoi
limiti, secondo lui, resterebbero due:
\begin{description}
  \item[Il costo] Tutti dicono che sta diminuendo e che diminuirà, ma non se ne
        può essere sicuri al 100\%.
  \item[La qualità] Cioè la capacità di leggere l'output. Se si dà un regolamento
        a Claude o a ChatGPT chiedendo quali problemi ha, il modello risponde con
        i problemi che, in percentuale, è più probabile che gli esseri umani
        abbiano già individuato nei dati su cui è stato addestrato. Ne derivano
        due problemi: se il problema del proprio gioco non è mai stato
        individuato prima, il modello non lo vedrà; e viceversa può segnalare
        problemi che il gioco non ha.
\end{description}
Fa l'esempio di Grammarly, che usa nella versione premium per rivedere i testi
in inglese: spesso non accetta i suoi consigli, perché se vuole scrivere un
dialogo di videogioco in \textbf{dialetto Cockney} londinese, lo deve scrivere
in Cockney. Se Grammarly dice che è sbagliato, ``sei tu che non hai capito il
motivo per cui voglio scrivere questo testo in Cockney''. La conclusione: prima
di usare un'IA per un'attività bisogna sapere \textbf{che cosa si vuole
ottenere}, qual è l'obiettivo, e quali \textbf{competenze} servono per valutare
in modo critico se l'output di un qualunque LLM è davvero valido. Daniele non
vuole dilungarsi oltre: è un argomento ``molto caldo'' per lui, che è anche
scrittore e si sente dire da tutti ``ti sostituiranno''.

Subito dopo si presenta finalmente Gianmarco (lo si è visto nella
sezione~\ref{a1:sec:partecipanti}), e poi si fa una pausa di una decina di
minuti. Daniele avverte che nella terza parte lui non ci sarà (alle sette e un
quarto ha un impegno di lavoro), e lo stesso vale per Beatrice: la parte finale
di domande e risposte sarà con il solo Ruggero, ``in ottime mani''.

% ---------------------------------------------------------------------
\section{Le figure professionali dell'industria analogica}
\label{a1:sec:figure}

Dopo la pausa (in cui si scopre che Angelo lavora in un piccolo birrificio in
Brianza e si scherza sul rischio di un \emph{drinking game}, tra la birra e la
vendemmia) la lezione riprende con le slide di Ruggero. ``Oggi è una lezione
introduttiva'': gli argomenti si vedranno ``molto dall'alto'' e verranno
approfonditi nelle lezioni successive. Ma servono già ora, per creare
\textbf{l'inizio di un linguaggio condiviso}: ci sono termini che torneranno
spesso, a cominciare da quelli del mercato e dell'industria del gioco
analogico, un'industria ``molto varia, molto caotica, anche molto mutevole'',
con tanti tipi diversi di giochi e altrettanti tipi diversi di entità che ci
lavorano.

\subsection{Una decina di etichette}

I docenti hanno raccolto ``una decina di etichette'', cioè di \textbf{ruoli} che
nell'industria sono o possono essere presenti. Due avvertenze:
\begin{itemize}
  \item nelle aziende più piccole \textbf{la stessa persona copre spesso più
        ruoli}, per forza di cose. Sarà anche il caso degli studenti, che
        lavoreranno in gruppi di tre e non di venti: è facile che ciascuno finisca
        a fare più cose. I docenti ne terranno conto, e qui entreranno in gioco
        anche competenze diverse dal game design (programmi di grafica,
        illustrazione, impaginazione, layout, \emph{proofreading}\ldots);
  \item Daniele precisa che \textbf{non si chiede agli studenti di diventare}
        illustratori o altro. Ruggero concorda, ma aggiunge che è utile sapere
        anche \emph{che cosa non è il proprio lavoro}. Se in un gruppo nessuno sa
        illustrare, ai docenti va benissimo un prototipo che mostri dove vanno
        le immagini: \textbf{placeholder}, concept, anche semplici box di testo
        con scritto ``qui vorremmo un'immagine di\ldots''. ``È importante anche
        sapere che cosa sapete fare e che cosa non sapete fare''.
\end{itemize}
La tabella~\ref{a1:tab:figure} riassume le figure; le sottosezioni seguenti le
riprendono una per una, con i commenti dei docenti.

\begin{table}[htbp]
\centering
\small
\begin{tabular}{@{}p{3.3cm}p{9.6cm}@{}}
\toprule
\textbf{Figura} & \textbf{Di che cosa si occupa} \\
\midrule
Game designer & crea le regole alla base dell'esperienza: le meccaniche, come
interagiscono tra loro e come chi gioca interagisce con esse \\
Game developer & fa da ponte tra chi fa game design e il publisher: bilancia il
prototipo, cura i modelli matematici, adatta il gioco (e il paragame) alle
necessità del publisher \\
Editor (proofreader) & revisiona e adatta i testi a partire dalla prima stesura
del designer (correzione, linguaggio inclusivo\ldots) \\
Publisher (editore) & pubblica e produce fisicamente il gioco; il suo marchio è
sulla scatola o sulla copertina \\
Sales e marketing & il lato commerciale: vendite (anche B2B) e comunicazione del
prodotto \\
Graphic designer & icone meccaniche, wireframe, UX: il lato funzionale degli
elementi visivi \\
Illustratore & le illustrazioni: il \emph{look and feel}, il world building, il
``vibe'' del gioco \\
Sensitivity reader (cultural consultant) & verifica l'accuratezza e la
rappresentazione di un tema sensibile che il team non conosce \\
Producer & coordina il progetto (quasi project management), parla con le
tipografie, verifica che i materiali siano sostenibili economicamente \\
\bottomrule
\end{tabular}
\caption{Le figure professionali dell'industria del gioco analogico presentate
nella slide di Ruggero, in sintesi.}
\label{a1:tab:figure}
\end{table}

\subsection{Game designer}

Il \textbf{game designer} (o la game designer) è chi si occupa, ``sembra banale
a dirlo'', di \textbf{creare le regole alla base dell'esperienza} che si vuole
far arrivare. Vale per tutti i tipi di giochi, non solo per quelli analogici, e
che cosa faccia concretamente in un progetto o in una serie specifica ``può
mutare tantissimo''. Nell'analogico non verrà chiesto di programmare o di
prototipare in digitale; verrà chiesto di \textbf{creare le meccaniche di
gioco}, di capire \textbf{come le meccaniche si interfacciano tra loro} e
\textbf{come le persone che giocano si interfacceranno con le meccaniche}.

\subsection{Game developer}

Il \textbf{game developer} (o la game developer) sarà l'argomento delle due
lezioni di Marta. La parola \emph{development} è nota soprattutto
dall'ambito videoludico, ma nell'analogico ha un significato molto specifico
``che non c'entra niente con la programmazione'': il game developer è
\textbf{un ponte tra chi fa game design e il publisher}. Di solito si occupa di
\textbf{bilanciare il prototipo}, di curare i \textbf{modelli matematici}, di
adattare il gioco e il \textbf{paragame}\footnote{La trascrizione dice ``il
Wix, il paragame''. Il \emph{paragame} è, in termini generali, tutto ciò che
accompagna il gioco senza esserne il sistema di regole (regolamento, scatola,
materiali, comunicazione); i docenti lo tratteranno in una lezione dedicata, e
qui il termine è solo citato.} alle necessità del publisher. La distinzione
viene ripresa e approfondita nella sezione~\ref{a1:sec:design-development}.

\subsection{Editor}

L'\textbf{editor}, a volte chiamato \textbf{proofreader} a seconda dell'azienda,
si occupa della revisione dei \textbf{testi}: non della stesura iniziale, perché
di norma, data la scala dell'industria, è chi fa game design a scrivere un
\textbf{primo draft}, e l'editor ci passa sopra. Non solo per correggere i
refusi: spesso il suo lavoro è anche \textbf{adattare} i testi, per esempio
renderli \textbf{gender neutral}. Se ne vedranno casi studio più avanti.

\subsection{Publisher}

Il \textbf{publisher}, o \textbf{editore}, è l'entità che nell'ambito
commerciale \textbf{pubblica} giochi, giochi di ruolo o giochi da tavolo. È una
figura ``più solidificata'' nel gioco da tavolo e nel gioco di ruolo da tavolo;
\textbf{nell'ambito pervasivo non esiste proprio in questa forma}, come
mostrerà Beatrice parlando dei casi studio. Se si compra un gioco da tavolo o un
gioco di ruolo in un negozio, il marchio del publisher è di solito sul retro
della copertina o sul fronte della scatola: è l'entità che, di solito,
\textbf{produce fisicamente} il gioco.

\subsection{Sales e marketing}

L'ambito \textbf{sales e marketing}, a volte affidato a un'unica persona e a
volte separato, riguarda il lato commerciale. In alcune aziende è interno al
publisher, e si occupa di parlare con le fabbriche e le tipografie, di ``far
quadrare i conti''. In altre realtà è completamente diverso, specialmente nel
\textbf{B2B} o nei giochi \textbf{didattici} e \textbf{formativi}, di cui
Daniele e Beatrice ``sapranno parlare meglio''. Il tema viene ripreso più avanti
in dettaglio (sezione~\ref{a1:sec:marketing}).

\subsection{Graphic designer e illustratore}

Ruggero presenta insieme \textbf{graphic designer} e \textbf{illustratore}
proprio perché sono ruoli molto diversi che si tende a far coincidere:
\begin{itemize}
  \item chi fa \textbf{graphic design} si occupa delle \textbf{icone meccaniche},
        del \textbf{wireframing}, della \textbf{UX}: del lato \textbf{meccanico e
        funzionale} degli elementi visivi, cioè di tutto ciò che serve
        meccanicamente nel gioco, o anche alla sua ``qualità della vita'';
  \item chi fa \textbf{illustrazione} si occupa dell'\emph{art}, delle
        illustrazioni: copre il \textbf{look and feel}, la parte di
        \textbf{world building} e il ``vibe'' che si dà al gioco.
\end{itemize}

\paragraph{Daniele: che cos'è la user experience.} Daniele interviene sulla
\textbf{user experience}, ``importantissima per me'': ha un background da
\textbf{layout artist} (impaginava magazine), un mestiere in cui la UX conta
moltissimo. Per user experience non si intende solo l'\textbf{interfaccia
grafica}, la UI, che anche nei videogiochi è disegnata da un artista. La UX è
qualcosa che può fare anche il game designer: \textbf{disegnare l'esperienza del
giocatore}, cioè ``il gesto, l'intenzione, l'intuito'' che stanno dietro
l'aspetto visivo dell'interfaccia. Per fare UX non serve saper disegnare.
L'esempio è quello della \textbf{porta}: se si progetta una porta, bisogna
indicare all'utente come si apre, come si prende la maniglia; ``quella è user
experience''. Il design della maniglia in sé, invece, è
\textbf{interfaccia}.

\paragraph{Un aneddoto sulle icone.} Ruggero aggiunge che molte persone che
fanno \textbf{graphic design ben fatto non sanno illustrare}, e spesso vale il
viceversa. Lo si nota quando le icone di un gioco sono bellissime ma non
funzionali: andando a guardare i crediti, di solito si scopre che le icone le ha
fatte la persona che ha fatto le illustrazioni. ``Comunque non andiamo adesso in
questo, perché è un altro vaso di Pandora''.

\subsection{Sensitivity reader}

Il \textbf{sensitivity reader}, o \textbf{cultural consultant} a seconda del
progetto, fa un lavoro di \emph{sensitivity reading} del gioco. È una figura
molto usata nell'ambito \textbf{formativo e didattico}. L'esempio di Ruggero: se
si fa un gioco per un ospedale su una malattia specifica, chi lo progetta non è
un medico, ma il gioco deve essere \textbf{accurato} e \textbf{rappresentativo}
delle dinamiche che quella malattia comporta. Serve una persona che conosca quel
tema, anche se non sa che cosa sia il game design. Sono figure chiamate di
solito \textbf{a progetto}, per giochi specifici, che compensano competenze che
il team non ha ``per forza di cose, perché noi non sappiamo tutto''; e a volte,
per lavoro, si devono fare giochi su temi che non si conoscono personalmente.

\paragraph{Daniele: non è un censore.} Daniele segnala un problema italiano:
alcune persone, anche online, hanno confuso il sensitivity reader con un
\textbf{censore}, ``una persona che fa censura, soprattutto in materia di
\emph{woke}''. ``Assolutamente non è così''. La parola \emph{sensitive} non
significa che il pubblico sia sensibile a un tema, ma che \textbf{l'argomento di
cui tratta il gioco è sensibile}. In un gioco sulla medicina, realistico o no,
non si possono inserire imprecisioni. Per questo esistono i sensitivity reader
``più disparati'': da quelli per la cultura indiana o afroamericana fino a quello
su ``come si allevano gli gnu in Africa''. È una persona competente su un
argomento che aiuta a essere \textbf{veritieri}. Certo, uno degli ambiti
principali in cui si usano è quello culturale, ma ``non fate l'errore di pensare
che siano dei censori''. Daniele ha lavorato con ``centinaia'' di sensitivity
reader e nessuno gli ha mai detto ``non puoi scrivere questa cosa''. Di solito
sono come degli \textbf{editor}: danno opzioni, mostrano i punti problematici,
anche quelli di cui non ci si era nemmeno accorti.

Ruggero conferma con un esempio di quando lavorava in \textbf{Danimarca} con
l'\textbf{OMS}: giochi sul \textbf{diabete di tipo 2} nei bambini, per un target
di ragazzi delle scuole medie. C'era un team dell'OMS specializzato nella
comunicazione sul diabete, per garantire che il gioco fosse accurato e ne desse
una rappresentazione adatta al target. Il sensitivity reader aiuta a rendere il
gioco più accurato e ``più tematicamente coeso''.

\subsection{Producer}

Il \textbf{producer} è una figura a volte interna al publisher. In un progetto
grande è la figura di riferimento, ``quasi'' un project manager. In un progetto
piccolo si occupa spesso anche di parlare con le \textbf{tipografie} e di fare la
\textbf{lista dei materiali}, per assicurarsi che i materiali non siano solo
funzionali al gioco ma anche \textbf{economicamente sostenibili} da produrre:
``magari è un gioco bellissimo, ma poi costa più di quanto potete permettervi, e
nessuno vi comprerà mai un card game a 200 euro''.

\subsection{Perché conoscere queste figure}

Ruggero chiude la slide ``tornando a ombrello'': perché parlarne, se gli studenti
sono lì per fare game design? Non solo per saper dare un \textbf{nome} alle
figure. Non c'è niente di male a creare un gioco solo per sé; ma chi vuole fare
giochi a livello lavorativo, perché li giochino altre persone, ``magari più
persone possibile'', \textbf{si interfaccerà con queste figure}, e sapere che
cosa fanno e come parlarci è parte importante del percorso. Nel corso se ne
toccheranno diverse: il development con Marta, la produzione con Flavio, e
lezioni sulla UX, sul paragame, sul graphic design. ``Abituatevi mentalmente:
non siete in solitaria contro il mondo''.

% ---------------------------------------------------------------------
\section{Marketing e sales}
\label{a1:sec:marketing}

Arrivato all'angolo della slide con marketing e sales, Ruggero passa la parola
ai colleghi: ``io è meglio che non mi fate parlare di marketing e sales'', anche
per esperienze personali.

\subsection{Beatrice: i giochi trattati come prodotti}

Beatrice premette che le due voci sono in parte autoesplicative. Nel mondo
capitalista in cui viviamo, i giochi sono ``strumenti molto potenti per fare
\emph{n} cose'', ma sono trattati come \textbf{prodotti}, e il corso ragionerà
anche in quest'ottica. Nel \textbf{marketing} delle industrie ludiche e
videoludiche ci sono varie professionalità che gli studenti probabilmente già
conoscono: dal \emph{copy} alla grafica, tutto ciò che riguarda il
\textbf{vendere} il gioco o l'attività di gioco.

Il lato \textbf{sales} è forse la distinzione meno nota. Come in tante aziende, è
la parte che si occupa delle vendite \textbf{B2B}, cioè tra aziende: fare
accordi per vendere prodotti ludici a diverse realtà, per esempio ai
\textbf{distributori}, o a una casa editrice che vende alle librerie, che poi
vendono al consumatore finale. Chi si occupa di sales nel mondo ludico capisce
come vendere \textbf{alle aziende}. I due mondi si parlano spesso, e soprattutto
nelle realtà piccole si trovano figure che fanno ``un po' sales, un po'
marketing''.

\subsection{Daniele: che cos'è (e che cosa non è) il marketing}

Daniele il marketing si è trovato ``purtroppo'' a farlo, e lo fa anche per
l'accademia. ``Purtroppo'' perché in Italia è mal compreso: molti lo
identificano con l'\emph{advertisement} online. Non è così: il marketing è una
\textbf{disciplina molto più antica}, e in Italia ``abbiamo molto dimenticato''
come si fa il marketing tradizionale. Alcune cose da ricordare:
\begin{enumerate}
  \item \textbf{Il marketing non dice che cosa pubblicare, ma come pubblicarlo}.
        Esistono, nei videogiochi e nei giochi da tavolo, realtà che fanno
        prodotti \textbf{marketing driven}: l'esempio sono gli
        \textbf{hyper-casual game} per smartphone, quelli con il personaggio che
        va a destra e a sinistra e l'esercito che si allarga. Possono anche
        guadagnare, ma ``sono un po' la morte della creatività''. Il marketing
        ideale \textbf{prende un prodotto pensato da altri e gli dà le maggiori
        chance di essere venduto}: cerca il pubblico ideale per il prodotto che
        si ha, senza doverlo cambiare. Anzi, il mondo videoludico, e ludico in
        generale, ``è lastricato'' di errori di dipartimenti marketing che hanno
        cambiato completamente non il gioco, ma il modo in cui l'hanno
        pubblicizzato, e il gioco è fallito.
  \item \textbf{Usare strumenti del marketing non significa fare marketing}. Game
        design e game development a volte usano strumenti che vengono dal
        marketing (gli studenti ne vedranno con Beatrice, con Daniele e con
        Marta).
  \item \textbf{Il marketing è un gerundio}: si fa. È un'azione che si compie per
        comunicare; la comunicazione è attiva, e il marketing è l'insieme delle
        attività con cui si comunica il proprio gioco, dal \emph{copy} ai social.
        Può includere il digital marketing, ``ma non è solo quello''.
  \item \textbf{Il marketing esiste anche se si crede che non esista}. In Italia
        molti non vogliono ``fare marketing'' perché lo associano al lucro. Ma
        il marketing è \textbf{lo studio di come gli esseri umani comprano e
        vendono}, ed è innato in noi: ci sono logiche di cui siamo vittime, e
        ``anche se sei il bastian contrario, l'anticapitalista'', magari si
        compra comunque su Amazon perché fa comodo. Sono meccanismi che
        esistono, ``affascinanti anche, a volte''.
\end{enumerate}
Gli studenti non dovranno fare marketing, ma dovranno ricordare alcuni di questi
concetti, soprattutto nelle lezioni con Marta, che tratta argomenti legati al
marketing. E, chiude Daniele, come game designer e game developer ``per lo più
speriamo che il marketing non ci rovini i giochi'', soprattutto nel videoludico.

% ---------------------------------------------------------------------
\section{Game design e game development}
\label{a1:sec:design-development}

Ruggero torna sulla distinzione più importante della slide, ``proseguendo con
questa strada'':
\begin{description}
  \item[Game design] è chi \textbf{progetta l'esperienza} tramite la
        \textbf{creazione delle meccaniche} di gioco.
  \item[Game development] è chi fa \textbf{mediazione} tra la persona che scrive
        il game design e la \textbf{casa editrice}.
\end{description}

\subsection{I red flag dei designer}

Ruggero lo ripete perché lo vede ogni giorno nel suo lavoro: tra giochi
pubblicati e prototipi, ne valuta ``quasi uno al giorno, se faccio un conto
annuale'', e solo una minima percentuale arriva alla pubblicazione. Molte volte
i \textbf{red flag} vengono proprio da designer che non hanno presente questa
definizione, e che cominciano a mettere \textbf{clausole sulla non modifica del
gioco}: il gioco deve restare com'è. ``Questa cosa è impossibile con un
publisher, spesso''. Il publisher, per forza di cose, potrà chiedere di cambiare
la grafica, di reimpaginare il regolamento, di cambiare i materiali perché
troppo costosi: chiederà dei \textbf{compromessi}. E la persona che, all'interno
del publisher, fa questo lavoro è il \textbf{developer}.

Sta al designer capire \textbf{su che cosa del proprio gioco si può trovare un
compromesso e su che cosa no}. Publisher diversi lavorano in modo diverso, e il
designer ha sempre l'autorità di dire ``no, non voglio pubblicare questo gioco
se queste sono le modifiche''. Ma deve sapere chi fa il lavoro di development,
perché se il suo obiettivo è la \textbf{pubblicazione tradizionale} si
interfaccerà con questa figura.

\subsection{Daniele: una distinzione che manca al videogioco}

Daniele sottolinea che questa distinzione è \textbf{tipica del mondo analogico}
ed è ``molto interessante''. In Italia la parola \emph{developer} è diventata
negli anni quasi un sinonimo di programmatore di videogiochi, o di chiunque
lavori nel videoludico. Secondo lui c'entra anche il fatto che qualche anno fa
il sito \textbf{Gamasutra} ha cambiato nome in \textbf{Game Developer}, e la
parola è diventata un ombrello per tutte le professioni del videogioco. ``Mi fa
un po' ribrezzo'': il game developer dovrebbe essere un ruolo di design, un
designer che, come nell'analogico, fa \textbf{mediazione}. Nel videogioco i
publisher hanno in effetti i loro developer interni, che adattano il gioco, per
esempio alle console e ai loro kit di sviluppo. Ma la distinzione dovrebbe
esistere anche nel videogioco, perché è ``difficilissimo'' avere la mentalità
per fare le due cose insieme: chi fa design ha un focus che raramente si può
sovrapporre a quello di chi fa mediazione.

A Daniele piace la parola \emph{develop}, \textbf{sviluppare}, nel senso di
prendere qualcosa che è stato progettato e \textbf{renderlo possibile nel mondo
reale}. Ci sono progetti che hanno bisogno di uno sviluppo che ``polishi''
l'esperienza: la semplifichi, la migliori, la renda comprensibile e la concluda,
ma soprattutto la renda \textbf{stampabile} e \textbf{fruibile} per l'utente. La
metafora crossmediale è quella della \textbf{fotografia}: chi sviluppa le
fotografie non è chi scatta la foto, ma chi la rende visibile, anche
fisicamente. ``È un po' retro come metafora, ma ci capiamo''. ``È vero, però è
giusto'', conferma Ruggero.

% ---------------------------------------------------------------------
\section{Il mercato: Italia ed estero}
\label{a1:sec:mercato}

Il corso è in italiano e si svolge in Italia, ma bisogna conoscere anche il
\textbf{mercato estero}: molti giochi progettati da autori italiani vengono
pubblicati all'estero, e conoscere il mercato estero aiuta a capire come il
proprio gioco può arrivarci, ``o viceversa, come non può arrivarci''. La
tabella~\ref{a1:tab:mercato} riassume il confronto della slide.

\begin{table}[htbp]
\centering
\small
\begin{tabular}{@{}p{6.4cm}p{6.4cm}@{}}
\toprule
\textbf{Italia} & \textbf{Estero (soprattutto Europa)} \\
\midrule
il gioco è ancora visto come un'attività \textbf{infantile}, anche se la
normalizzazione è in crescita &
forte tradizione di gioco anche in un contesto \textbf{familiare adulto},
soprattutto ``dalla Germania in su'' \\
risposta \textbf{conservativa} sul tema ludico &
\textbf{premi di prestigio} riconosciuti in tutto il mondo (Spiel des Jahres,
As d'Or) \\
forte focus sulla \textbf{localizzazione} di giochi stranieri &
si cercano soprattutto giochi \textbf{originali} \\
investimento \textbf{molto basso} sulla creazione di giochi originali interni;
pochissimi designer assunti &
\textbf{developer interni} pagati con un contratto (Francia, Germania,
Danimarca\ldots) \\
\bottomrule
\end{tabular}
\caption{Il mercato del gioco analogico in Italia e all'estero, secondo la
slide di Ruggero e i commenti di Daniele.}
\label{a1:tab:mercato}
\end{table}

\subsection{L'Italia}

\begin{itemize}
  \item In Italia il gioco è ancora visto come un'\textbf{attività infantile},
        anche se la \textbf{normalizzazione} è in crescita. C'è ancora una forte
        risposta \textbf{conservativa} sul tema ludico.
  \item C'è un forte focus sulla \textbf{localizzazione} dei giochi stranieri: in
        una qualunque ludoteca o negozio, molti dei giochi da tavolo e di ruolo
        sono di publisher o designer esteri, pubblicati in italiano da publisher
        italiani. (La localizzazione, cioè portare un titolo dall'estero in
        Italia, esiste come funzione all'interno del publisher; Ruggero non l'ha
        messa nella slide delle figure perché non è il focus del corso.)
  \item C'è, ``purtroppo, preparatevi psicologicamente'', un \textbf{investimento
        molto basso} sulla creazione di giochi originali interni alla casa
        editrice. Pochissime case editrici italiane hanno designer interni pagati
        a tempo pieno con contratto a tempo indeterminato. Daniele lo dice senza
        mezzi termini: ``il mercato è molto piccolo, conosciamo praticamente
        tutto il settore'', e le persone che in Italia fanno game design
        analogico a tempo pieno con un contratto a tempo indeterminato ``sono
        meno delle dita delle mie due mani''. Di solito la formula è quella del
        \textbf{freelance}, del contratto, delle \textbf{royalty}. Chi va da un
        publisher aspettandosi di essere assunto a tempo pieno come game
        designer deve sapere che ``non siamo purtroppo ancora in quel tipo di
        situazione''.
  \item C'è però un contraltare: in Italia, soprattutto nel gioco di ruolo ma
        anche nel gioco da tavolo, c'è stata molta \textbf{imprenditoria}, anche
        grazie a \textbf{Kickstarter}, che ha aperto le porte ad altri sistemi
        di \textbf{crowdfunding}. È un altro argomento, perché richiede un team,
        un investimento e così via.
  \item E gli autori non mancano, anzi: le occasioni di incontro tra autori
        vedono ogni anno un aumento dei possibili autori\footnote{La trascrizione
        dice ``i DAG a Parma''; il riferimento è probabilmente a IDEAG, il
        circuito di incontri tra autori di giochi italiani, ma non è
        ricostruibile con certezza.}. Bisogna ``sgomitare'' tra tanti autori,
        che magari arrivano da esperienze molto diverse.
\end{itemize}
Ruggero precisa: l'investimento basso è ``dall'alto''. Le persone che fanno e
vogliono fare giochi sono tantissime e in aumento; scarseggia l'investimento,
sia dei publisher disposti a pagare stipendi sia dei \textbf{bandi} pubblici. I
bandi esistono (Beatrice ne parlerà per i giochi pervasivi, che spesso si
muovono in quella direzione, e Ruggero cita qualcuno che sta facendo un card
game con un bando italiano, ``una cosa assurda, ma lo sta facendo''), ma sono
una rarità, soprattutto rispetto all'estero.

\subsection{L'estero}

\begin{itemize}
  \item L'estero (dipende dalla nazione, ma in generale il resto d'Europa) ha una
        forte \textbf{tradizione di gioco} anche in un contesto \textbf{familiare
        adulto}, soprattutto ``dalla Germania in su''.
  \item Ha \textbf{premi di prestigio} riconosciuti in tutto il mondo, come lo
        \textbf{Spiel des Jahres} tedesco e l'\textbf{As d'Or} francese\footnote{La
        trascrizione dice ``lo Spield AR, la SDOR''.}.
  \item Cerca soprattutto \textbf{giochi originali}. Lo si vede dal fatto che
        molti autori italiani di giochi da tavolo pubblicano in Germania o in
        Francia e \textbf{non vengono pubblicati in Italia}, ``che è tutto
        dire''.
  \item Ha \textbf{figure di development interne}: non è raro vedere publisher
        francesi, tedeschi, danesi con developer pagati internamente, a
        contratto.
\end{itemize}
``Poi la situazione è ancora più drammatica'', ammette Ruggero, ``ma per ora
vogliamo tenere il morale un po' alto. È l'introduzione. Ci sarà tempo di andare
nella tristezza''.

Prima di proseguire, Ruggero insiste sulle domande: con lui si può interrompere
anche mentre parla, alzando la mano con le reaction o semplicemente togliendo il
microfono e chiedendo. ``Non mi offenderò mai, anzi, sarò felice di vedervelo
fare e mi ricorderò positivamente di voi''.

% ---------------------------------------------------------------------
\section{Perché serve una terminologia comune}
\label{a1:sec:terminologia}

La serie di slide successiva è ``una piccolissima parte di fondamenti di
design'' (``non spaventatevi, vi vedo che vi spaventate un po'\,''). Anche qui
molti temi verranno ripresi in dettaglio, ma l'idea è arrivare alla prossima
lezione con almeno un \textbf{glossario}, un \textbf{linguaggio condiviso}: se
qualcuno dice \emph{game developer}, tutti sanno di che cosa sta parlando.

Ruggero parte da una domanda, e precisa che non è retorica: \textbf{perché ci
serve una terminologia comune?} Gli studenti rispondono:
\begin{itemize}
  \item Gianmarco, ``molto pragmatico'': perché bisogna lavorare in team e
        capirsi. Vero, nel corso e anche fuori;
  \item un'altra studentessa: per la \textbf{comunicazione}, tra di noi e con
        tutti quelli che sono interessati nel settore. Vero anche questo:
        conoscere i termini comuni permette di parlare con chi fa game design o
        development nel dettaglio, nel tecnico, nello specifico.
\end{itemize}
Ruggero aggiunge altre due ragioni.

\paragraph{Saper comunicare il proprio gioco.} La prima riguarda la
\textbf{comunicazione del gioco} verso l'esterno. ``E lo dico non tirando acqua
al mio mulino, ma anzi togliendola'': \textbf{non è spesso il buon game design
a far diffondere un gioco}. ``Vorrei che lo fosse, sarebbe bellissimo, ma
purtroppo non è così''. Ci sono giochi bellissimi che non sanno comunicare, che
non sanno dire che cosa fanno, e che quindi è difficile, da fuori, trovare
invitanti; e ci sono giochi mediocri dal punto di vista del design che si
comunicano molto bene, e che quindi conoscono in molti. Saper parlare di quello
che si fa è utile \textbf{a sé stessi}, in primis.

\paragraph{Documentare per sé stessi.} La seconda è la \textbf{documentazione}:
documenti di design, di brainstorming e così via. ``Sembrerà una battuta, ma
provate a recuperare un documento di brainstorming una settimana dopo, se non
avete scritto bene le cose che avete scritto''. Chi l'ha fatto sa già di che
cosa si parla; ``chi non l'ha ancora vissuto, forse lo vivrà''.

% ---------------------------------------------------------------------
\section{Game e play}
\label{a1:sec:game-play}

La \textbf{prima, importantissima divisione terminologica} è quella tra
\emph{game} e \emph{play}. In italiano c'è un problema: \textbf{gioco} e
\textbf{giocare} hanno la stessa radice, mentre in altre lingue, tra cui
l'inglese, sono due parole strutturalmente diverse:
\begin{description}
  \item[Game] il \textbf{gioco come artefatto}: il gioco come ecosistema di
        meccaniche, come struttura di regole, come prodotto (``tra
        virgolette''), come oggetto fisico.
  \item[Play] l'\textbf{atto}, il verbo: le persone che giocano, l'approccio
        fisico all'attività del gioco.
\end{description}
È importante perché gli studenti faranno \textbf{game design}; esistono anche
ambiti di \emph{play design}, ma qui ``creerete giochi, non giocattoli''.

\subsection{Giochi e giocattoli}

Qual è la differenza? Uno studente risponde: \textbf{la presenza di regole}.
Esatto: un gioco ha delle regole, un'attività di \emph{play} non
necessariamente. Con un giocattolo si può fare \emph{play}: c'è chi gioca con i
Lego, chi gioca a guardia e ladri o a nascondino al parco, e magari alcune di
quelle regole sono ``più o meno semi-solidificate''. Ma il \emph{play} permette
spesso un approccio non strettamente legato alle regole.

Viceversa, \textbf{un game può esistere anche se nessuno lo sta giocando}. Se si
prepara su un tavolo una scacchiera con i pezzi degli scacchi, chi passa e la
vede riconosce il gioco degli scacchi, anche se nessuno sta giocando. Se invece
si vedono dei bambini correre in un parco giochi, quei bambini \emph{stanno
giocando}, ma non stanno giocando \emph{a un gioco} specifico. ``È importante
separare questa cosa''; ci si tornerà, ``ma per ora abituatevi a questa
differenza''.

\paragraph{Approfondimento: giocattoli, giochi e regole.} La stessa distinzione
è il punto di partenza del capitolo che Jesse Schell, in \emph{The Art of Game
Design}, dedica alla definizione di gioco (\emph{The Experience Rises Out of a
Game}). Schell nota che usiamo perfino un linguaggio diverso: ``un gioco è
qualcosa a cui si gioca'' (\emph{a game is something you play}), ``un giocattolo
è qualcosa con cui si gioca'' (\emph{a toy is something you play with}); e
arriva a definire il giocattolo come ``un oggetto con cui si gioca''. Analizzando
poi la definizione di gioco di Avedon e Sutton-Smith, ne ricava che ``i giochi
hanno regole, i giocattoli no'', e che le regole sono uno degli aspetti
\emph{definitori} del gioco, insieme alla volontarietà, agli obiettivi, al
conflitto e alla possibilità di vincere o perdere. Schell osserva anche che
molti giochi sono costruiti \textbf{sopra} un giocattolo (``una palla è un
giocattolo, ma il baseball è un gioco''), e tra i suoi consigli di
prototipazione c'è proprio \emph{Build the Toy First}: costruire prima il
giocattolo, cioè qualcosa che sia già piacevole da manipolare, e poi il gioco
intorno. È una lettura utile della distinzione di Ruggero: il \emph{play} può
fare a meno delle regole, il \emph{game} no.

% ---------------------------------------------------------------------
\section{Il framework MDA (e il DPE)}
\label{a1:sec:mda}

\subsection{Meccaniche, dinamiche, estetiche}

Un'altra terminologia fondamentale è quella del framework \textbf{MDA}
(\emph{Mechanics, Dynamics, Aesthetics}; in italiano a volte ``MDE''). Viene da
un \textbf{paper} nato nell'ambito del videogioco\footnote{È il paper di Robin
Hunicke, Marc LeBlanc e Robert Zubek, \emph{MDA: A Formal Approach to Game
Design and Game Research}, del 2004. Daniele aggiunge il paper alla cartella
condivisa durante la lezione.}, che ha ``solidificato'' queste terminologie, e
divide un gioco in tre aree:
\begin{description}
  \item[Meccaniche] le \textbf{regole per come sono scritte}. ``Ogni volta che
        pescate un 2 fate dei punti'': quella è una meccanica.
  \item[Dinamiche] l'\textbf{applicazione in tempo reale} delle meccaniche. La
        dinamica è il giocatore che, pescando una carta, \textbf{spera} che esca
        un 2. Quel processo non è una regola: ``il gioco non vi dice che dovete
        sperare che esca un 2''.
  \item[Estetiche] l'insieme di quello che il gioco vuole comunicare,
        l'\textbf{esperienziale}.
\end{description}

L'esempio di Ruggero mette insieme le tre aree. In un gioco in cui pescare un 2
fa fare punti (meccanica) e in cui quindi si spera che esca un 2 (dinamica),
\textbf{che estetica si ottiene se nel mazzo di 100 carte c'è un solo 2? E se ce
ne sono 50?} L'emozione di aver pescato \emph{l'unico} 2 del mazzo è
completamente diversa da quella di pescarne uno sapendo che ce ne sono
cinquanta e che è molto probabile trovarlo. Le regole sono quasi identiche; cambia
l'esperienza.

\subsection{Il falso amico: ``estetica''}

Daniele avverte di non sottovalutare la parola \textbf{estetica}, che in italiano
è un \textbf{falso amico}: sembra indicare la parte grafica, e invece non lo è.
Si incontreranno designer che hanno studiato MDA e altri che non l'hanno
studiato, e lo si capirà proprio da questo: i secondi diranno che le
\emph{aesthetics} sono la grafica. Le \emph{aesthetics} sono la parte
\textbf{esperienziale}: come nell'esempio di Ruggero, quello che lascia al
giocatore aver pescato l'unico 2 del mazzo.

\subsection{Il DPE: design, play, experience}

Per gli italiani, aggiunge Daniele, è più facile capire l'MDA attraverso un
altro paper, quello del \textbf{DPE}\footnote{È il framework \emph{Design, Play,
Experience} di Brian Winn, proposto come estensione dell'MDA per i giochi
educativi (\emph{serious games}). Anche questo paper viene aggiunto da Daniele
alla cartella condivisa.}, che di fatto fa un parallelismo assegnando
\textbf{design, play, experience}:
\begin{itemize}
  \item le \textbf{meccaniche} sono il lato \textbf{design};
  \item le \textbf{dinamiche} sono il lato \textbf{play};
  \item le \textbf{estetiche} sono l'\textbf{experience}, l'esperienza che riceve
        chi gioca.
\end{itemize}
``Se vivete le \emph{aesthetics} come l'esperienza che vive chi sta giocando,
state a postissimo''. In questo momento Beatrice lascia la lezione (``ciao,
Bea''). I paper di MDA e DPE sono in inglese; chi vuole può dargli una lettura.

\begin{figure}[htbp]
\centering
\begin{tikzpicture}[
    font=\small,
    blocco/.style={draw, thick, rounded corners=3pt, fill=gray!12,
                   minimum width=3.2cm, minimum height=1.0cm, align=center},
    dpe/.style={font=\footnotesize\itshape, align=center},
    esempio/.style={draw, rounded corners=2pt, text width=3.3cm,
                    align=left, font=\footnotesize, inner sep=4pt},
    freccia/.style={-{Stealth[length=2.5mm]}, thick}
  ]
  \node[blocco] (m) at (0,0)   {\textbf{Meccaniche}\\ \footnotesize le regole scritte};
  \node[blocco] (d) at (4.3,0) {\textbf{Dinamiche}\\ \footnotesize le regole in azione};
  \node[blocco] (a) at (8.6,0) {\textbf{Estetiche}\\ \footnotesize l'esperienza};
  \draw[freccia] (m) -- (d);
  \draw[freccia] (d) -- (a);
  \node[dpe, above=0.15cm of m] {DPE: \textbf{design}};
  \node[dpe, above=0.15cm of d] {DPE: \textbf{play}};
  \node[dpe, above=0.15cm of a] {DPE: \textbf{experience}};
  \node[esempio, below=0.45cm of m] (em)
    {``Ogni volta che peschi un 2 fai punti''};
  \node[esempio, below=0.45cm of d] (ed)
    {pescando, il giocatore \emph{spera} che esca un 2 (nessuna regola glielo
     chiede)};
  \node[esempio, below=0.45cm of a] (ea)
    {un solo 2 su 100 carte: pescarlo è un evento;\\
     50 su 100: è quasi scontato};
  \draw[dashed] (m) -- (em);
  \draw[dashed] (d) -- (ed);
  \draw[dashed] (a) -- (ea);
  % prospettive
  \draw[freccia, gray] ([yshift=-0.3cm]em.south west) --
    node[below, font=\footnotesize\itshape, text=black]
    {il designer agisce sulle meccaniche}
    ([yshift=-0.3cm]ea.south east);
  \draw[freccia, gray] ([yshift=-1.2cm]ea.south east) --
    node[below, font=\footnotesize\itshape, text=black]
    {chi gioca incontra prima l'esperienza}
    ([yshift=-1.2cm]em.south west);
\end{tikzpicture}
\caption{Il framework MDA con l'esempio del ``2'' usato da Ruggero, e la
corrispondenza con il DPE suggerita da Daniele. Le due frecce in basso
riassumono la lettura classica del paper: il designer costruisce le meccaniche
per ottenere un'esperienza, mentre il giocatore parte dall'esperienza.}
\label{a1:fig:mda}
\end{figure}

\paragraph{Approfondimento: le due direzioni dell'MDA e le otto estetiche.} Le
frecce della figura~\ref{a1:fig:mda} riprendono un'idea centrale del paper
originale, che la lezione non esplicita: designer e giocatore guardano il gioco
da lati opposti. In \emph{Challenges for Game Designers}, Brenda Brathwaite e
Ian Schreiber la riassumono così: i designer si chiedono quale estetica
vogliono ottenere, definiscono le dinamiche che porterebbero a quella
sensazione, e poi creano le meccaniche che producono quelle dinamiche. Lo stesso
libro conferma l'avvertenza di Daniele sul falso amico: ``le \emph{aesthetics}
non si riferiscono all'aspetto del gioco, ma alla risposta emotiva'' che il
designer vuole suscitare. E offre esempi di dinamiche nel senso di Ruggero: in
uno sparatutto i punti di \emph{spawn} sono una meccanica, mentre appostarsi
accanto a uno di essi per colpire chi riappare (lo \emph{spawn camping}) è una
dinamica; negli scacchi i movimenti dei pezzi sono meccaniche, mentre le
aperture codificate sono dinamiche. Per dare un nome alle estetiche, il paper
MDA propone una tassonomia, poi detta degli ``otto tipi di \emph{fun}'', che
Schell riprende in \emph{The Art of Game Design} come \emph{LeBlanc's Taxonomy of
Game Pleasures}: \textbf{sensazione}, \textbf{fantasia}, \textbf{narrativa},
\textbf{sfida}, \textbf{compagnia} (\emph{fellowship}), \textbf{scoperta},
\textbf{espressione} e \textbf{sottomissione} (il piacere di entrare nel
``cerchio magico'' e lasciarsi andare a un passatempo). È un primo esempio di
quello che i docenti chiedono con le parole vietate della prossima sezione:
invece di dire che un gioco è ``divertente'', dire \emph{quale} esperienza
offre.

% ---------------------------------------------------------------------
\section{Le parole vietate}
\label{a1:sec:parole-vietate}

``Attenzione, adesso qua triggeriamo delle persone''. Ruggero annuncia una regola
che ``tornerà'': nel corso ci sono delle \textbf{parole vietate}. ``Non perché
siamo dei censori''; il perché si vedrà in una lezione apposita, perché
spiegarlo adesso richiederebbe un'ora. Ma Daniele vuole spiegarne ``un
minimo''.

\subsection{La regola}

In \textbf{nessun materiale scritto} consegnato ai docenti, task compresi, sono
accettate queste parole, né in inglese né in italiano:
\begin{center}
\begin{tabular}{@{}ll@{}}
\toprule
\textbf{Inglese} & \textbf{Italiano} \\
\midrule
\emph{fun} & divertimento \\
\emph{immersion} & immersione \\
\emph{entertainment} & intrattenimento \\
\bottomrule
\end{tabular}
\end{center}
``Non li voglio neanche vedere''. Daniele lo ripete con un paragone: ``se mi fate
i disegni con Paint, benissimo, top, fantastico, siamo ancora best friend. Se mi
scrivete \emph{divertimento}, io sbrocco''. È una \textbf{regola dell'accademia}:
negli stessi corsi triennali di narrative design e game design le parole sono
vietate agli studenti per tre anni; ``per voi solo per tre mesi''. Daniele
aggiungerebbe volentieri un quarto termine, \textbf{gamification}.

\subsection{Perché}

Il motivo approfondito lo darà Ruggero nella lezione apposita. Daniele dà il
suo. Ruggero e Diego Ricchiuti sono molto contrari a \emph{fun} e a
\emph{immersion}; lui ha aggiunto \emph{intrattenimento}. Sono \textbf{parole
troppo passe-partout e troppo ambigue}:
\begin{itemize}
  \item \textbf{intrattenimento} in Italia è stata usata per anni come un modo
        per \textbf{sminuire}: ``ma sì, tanto è un gioco solo di
        intrattenimento, è un romanzo solo di intrattenimento''. Ma anche se un
        prodotto è ``solo di intrattenimento'', non vuol dire che non ci si debba
        lavorare con attenzione. E poi intrattenimento ``non significa niente'':
        ci sono tanti modi in cui una cosa può intrattenere;
  \item lo stesso vale per \textbf{fun}, e per \textbf{immersion}, che ha un
        problema in più: è una parola passe-partout per cui si può dire che un
        gioco è ``immersivo'' nel senso generale del termine, ma poi non si sa
        davvero che cosa voglia dire, \emph{immersivo in che senso}.
\end{itemize}

L'obiettivo didattico è che gli studenti \textbf{vadano oltre queste parole}:
``un po' come quando alle medie vi dicono smettetela di dire che tutto è bello,
perché ci sono vari gradi di bello, ci sono altri modi di dire bello''. Bisogna
\textbf{decostruirle}: quando si sta per usarle, chiedersi \textbf{che cosa si
intende} con quella parola e cercarne una \textbf{più precisa}.

Ruggero chiede, in uno dei ``rari casi'', di \textbf{non fare domande adesso} su
questa regola: ogni anno arrivano domande e commenti, e ci saranno lezioni
apposite di tre ore, con strumenti per capire il perché. Ma da subito gli
studenti devono abituarsi a evitare questi termini, sia perché si spera che
abbiano termini più precisi, sia perché queste parole ``portano con sé dei
problemi agganciati'' che si vedranno. ``Non spaventatevi, non verrete sparati in
un vicolo cieco di notte''.

\paragraph{Una piccola gamification.} Il divieto va visto come uno
\textbf{stimolo}, ``un po' come se stessimo gamificando la vostra esperienza'',
una sorta di user experience delle lezioni. Imponendosi di non usare queste
parole, si noterà come migliora la capacità di \textbf{descrivere i giochi} e le
proprie esperienze ludiche, e di essere più precisi. È un invito a ``obbligarsi a
essere più precisi e ad avere una mentalità più critica''. La lezione di
Ruggero sull'\textbf{involvement}, promette Daniele, sarà ``fenomenale'' e
``vi aprirà la testa''.

\paragraph{Approfondimento: anche i libri faticano a definire il ``fun''.} La
difficoltà che il divieto vuole far emergere è visibile anche nei testi di game
design. Schell, in \emph{The Art of Game Design}, prova a definire il \emph{fun}
e propone ``piacere con sorprese'' (\emph{pleasure with surprises}), ammettendo
subito che è strano poter usare una parola per tutta la vita, sapendo con
certezza che cosa significa, senza riuscire a esprimerla chiaramente quando
qualcuno lo chiede; e invita a mettere alla prova la definizione con
controesempi. La sua lente sul \emph{fun} (\emph{The Lens of Fun}) si limita,
non a caso, a chiedere quali parti del gioco siano \emph{fun} e \emph{perché}.
Anche \emph{immersion} ha molte letture: in \emph{21st Century Game Design},
Bateman e Boon la descrivono, tra l'altro, come la capacità di ``perdersi''
nel mondo artificiale presentato dal gioco, legata alla sospensione
dell'incredulità e ai giochi di mimesi. Sono significati diversi sotto la stessa
parola: proprio il motivo per cui, in un documento di design, conviene dire
quale si intende.

% ---------------------------------------------------------------------
\section{Chi crea, chi gioca, chi facilita}
\label{a1:sec:ruoli-esperienza}

Un'altra distinzione importante è quella tra \textbf{chi crea}, \textbf{chi
gioca} e \textbf{chi facilita}. Spesso, quando si fa un gioco, lo si fa provare
di persona a chi gioca. Ma ``la norma dei giochi pubblicati non funziona così'':
quando si pubblica un gioco non si vede chi lo compra e chi lo gioca, e non si
sa se qualcuno, a casa propria con i propri amici, segue le regole alla lettera e
interpreta correttamente quello che il designer aveva descritto. Quindi già in
fase di progettazione bisogna tenere a mente \textbf{chi sta facendo game design
e chi sta giocando}.

E nei giochi a evento, nei giochi pervasivi, nei giochi di ruolo da tavolo e dal
vivo (LARP) c'è un terzo pezzo del puzzle, ``importantissimo'': \textbf{chi
facilita}. La domanda da farsi è: \textbf{il regolamento è scritto per chi gioca
o per chi deve facilitare il gioco?} Bisogna abituarsi a pensare a tutte le
persone coinvolte nell'esperienza.

% ---------------------------------------------------------------------
\section{Playtest e iterazione}
\label{a1:sec:playtest}

\subsection{Nessun gioco nasce perfetto}

L'ultima voce del glossario è ``un'altra cosa importantissima'':
\textbf{playtest} e \textbf{iterazione}. Il corso è \textbf{pragmatico}:
``specialmente da parte mia'', dice Ruggero, non si vedranno teorie ``messe lì
per fare accademia sul game design nelle nostre torri d'avorio''. Ci saranno
teorie e framework, ma tutti estremamente pragmatici, perché l'obiettivo non è
solo capire i giochi, ma \textbf{saperli fare e iterare}.

Playtest e iterazione sono \textbf{inevitabili} nel game design, perché
\textbf{nessun gioco nasce perfetto}, soprattutto perché le persone sono
diverse: un gioco perfetto per una persona può essere incomprensibile per
un'altra. Chi progetta giochi per altre persone deve capire, dopo averli fatti
provare, \textbf{che cosa cambiare, che cosa sistemare, che cosa aggiungere, che
cosa togliere}. Ci saranno lezioni su playtest, prototipazione, raccolta dei
feedback, metodi etnografici. Ma bisogna abituarsi già ora all'idea che
\textbf{non esiste la consegna definitiva}. A un certo punto i task vanno
consegnati per forza di cose, ma, come diceva Daniele all'inizio, alla fine del
corso gli studenti avranno un \textbf{prototipo}, ``che sarà un prototipo'', e
gli strumenti per portare avanti il lavoro anche dopo.

\subsection{L'aiuto dopo il corso}

Non solo gli strumenti: anche l'aiuto dei docenti. Daniele racconta di aver
portato a un editore, a Play\footnote{Probabilmente \emph{Play -- Festival del
Gioco} di Modena; la trascrizione non permette di esserne certi.}, uno dei
progetti realizzati l'anno precedente da due studenti, un gioco di ruolo in
solitario (``un \emph{solo game} RPG molto, molto bello''). Se gli studenti
vorranno trasformare il proprio lavoro in un prototipo vero e proprio, possono
contattarlo: lo fa ``per lavoro, ma anche per passione''. Gestisce anche una
community per aiutare i designer ``un po' depressi'' a rimettersi a progettare.

Quindi: ``non sentitevi persi se durante la fase di iterazione vi sentite
persi, è normale''.

\subsection{Il dizionario comune e il production dictionary}

Daniele riprende un punto detto ``benissimo'' da Ruggero e che secondo lui non è
stato sottolineato abbastanza: il \textbf{dizionario comune} serve anche a
\textbf{non perdersi} mentre si itera. Se da un'iterazione all'altra i
documenti usano termini completamente diversi, e non c'è un dizionario di base,
ci si perde.

Nel videogioco, una delle primissime cose che un designer fa quando entra in
un'azienda è il \textbf{production dictionary}: si annota, in un file privato,
tutte le parole che usano i colleghi, per allinearsi al dizionario comune
dell'azienda, perché lì su un prototipo possono lavorare anche cento persone.
Nel gioco da tavolo non si fa molto, perché il mondo dell'analogico è
``talmente più piccolo''; ma serve comunque imparare un \textbf{dizionario
minimo}, dai docenti e dai professionisti. E il corso ne darà una parte,
entrando nel dettaglio di alcuni termini e framework.

\paragraph{L'esempio del feedback.} Perché serva, lo mostra un esempio. Si
testa un gioco e ci si annota come feedback: ``questo gioco deve essere più
divertente''. Che cosa vuol dire? Come lo si migliora? Quella frase non dà
elementi utili. Se invece si scrive che deve essere \textbf{più adrenalinico},
\textbf{più frenetico}, \textbf{più rilassato}, \textbf{più cozy}, o
\textbf{più bilanciato matematicamente}, si sa dove andare a lavorare.

\paragraph{Approfondimento: la regola del ciclo.} Il principio ``nessun gioco
nasce perfetto'' ha un nome preciso in \emph{The Art of Game Design}: è la
\textbf{Rule of the Loop}, la regola del ciclo. ``Più volte provi e migliori il
tuo design, migliore sarà il tuo gioco''. Schell precisa che non è una lente,
cioè un punto di vista, ma ``una verità assoluta'', senza eccezioni: ci saranno
momenti in cui si cercherà di convincersi che ``questa volta il design è così
buono che non serve testarlo'', e ogni volta se ne pagherà il prezzo. Ed è
interessante la ragione per cui la regola pesa di più sul digitale: secondo
Schell il tempo e il denaro necessari per provare e correggere un sistema sono
molto maggiori nei videogiochi che nei giochi tradizionali. Per un gioco di
carte o da tavolo, con il tempo per provarlo e cambiarlo più e più volte, il
consiglio è semplicemente di farlo. Il gioco analogico è, da questo punto di
vista, la palestra ideale dell'iterazione.

% ---------------------------------------------------------------------
\section{Tabletop Simulator}
\label{a1:sec:tts}

\subsection{Il software di settore}

``Last but not least'': il corso usa un software, \textbf{Tabletop Simulator},
che si compra su \textbf{Steam} e costa tra i \textbf{10 e i 20 euro} (``aspettate
che vada in sconto, va in sconto ogni due per tre''). Daniele spiega perché:
\begin{itemize}
  \item è \textbf{il software di settore} del gioco da tavolo, quello che i
        publisher usano per i prodotti digitali, interni ed esterni;
  \item essendo su Steam, permette di fare \textbf{playtest} e di giocare con gli
        amici di Steam. Daniele stesso sta testando lì il proprio gioco;
  \item funziona anche su Mac: Windows o Mac, nessun problema.
\end{itemize}
Ci sarà una lezione apposita, di Ruggero, sulla creazione del proprio
\textbf{modulo} di gioco.

\paragraph{Che cosa non è.} Daniele chiarisce subito un ``fraintendimento
gigantesco'', soprattutto per chi ha già esperienza: Tabletop Simulator
\textbf{non richiede programmazione} e \textbf{non è un software per
trasformare i giochi da tavolo in videogiochi}. È un software che
\textbf{simula dei componenti fisici su un tavolo}: serve per giocare i board
game in digitale. Simula anche la fisica: c'è perfino il comando per
\textbf{ribaltare il tavolo} (``non usatelo ora'').

\paragraph{La partnership mancata.} Da anni Daniele prova a fare una
partnership con gli sviluppatori: due anni prima ha contattato l'azienda per
acquistare delle chiavi per gli studenti, dato che il software si usa nei corsi.
L'azienda ha risposto che non fa accordi commerciali. ``Credo sia la prima
azienda nella mia vita che dice che non vuole che io le dia dei soldi''. Gli
studenti dovranno quindi provare ad acquistarlo da sé; se i docenti riusciranno a
darglielo, bene.

\paragraph{La domanda di Gianmarco.} Gianmarco segue le lezioni da un computer
aziendale ``super chiuso'', su cui non può installare Steam: deve usare Tabletop
Simulator durante le lezioni? Ruggero risponde che per la lezione dedicata è
meglio poter seguire e rifare in diretta. Ma il software non sarà richiesto
necessariamente per il project work: un gruppo potrebbe fare un gioco pervasivo
o un gioco di ruolo, non un gioco da tavolo. Essendo un software di settore, è
comunque bene conoscerlo. Ruggero aggiunge che le versioni non ufficiali che
circolano non permettono di giocare in molti, e quindi non servono per
playtestare o giocare insieme.

\subsection{Una dimostrazione}

Daniele mostra il software aperto da Steam:
\begin{itemize}
  \item si può creare una partita \textbf{single player} o \textbf{multiplayer};
        per giocare tra compagni si crea un server multiplayer e ci si aggiunge
        come amici su Steam;
  \item si dà al server un nome, un numero di giocatori, eventualmente una
        password; creato il server, si apre la scelta dei giochi;
  \item ci sono giochi gratuiti, giochi da acquistare (DLC, non obbligatori) e
        soprattutto un \textbf{Workshop} ``ampissimo'', con decine di migliaia di
        giochi caricati dagli utenti: ``se c'è un board game che conoscete,
        quel board game è su Tabletop Simulator, giocabile dal Workshop''.
        L'esempio mostrato è l'open beta ufficiale del gioco da tavolo di
        \emph{Slay the Spire}\footnote{La trascrizione dice ``Slade Spire''.}:
        molti publisher, per farsi pubblicità, mettono il proprio gioco in demo o
        giocabile in digitale;
  \item c'è un \textbf{tutorial} iniziale che spiega come muoversi sul tavolo;
        con \emph{Esc} e \emph{Help} si trovano i comandi da tastiera.
\end{itemize}

\subsection{Il primo task}

Il primo (e per ora unico) task è questo: \textbf{entro il 10 ottobre}, giorno
della lezione di Ruggero su Tabletop Simulator, giocare \textbf{tre o quattro
volte in totale}, online, tra compagni di corso o con amici. Non si tratta di
imparare a creare moduli, ma di prendere la mano \textbf{lato giocatore}: aprire
il software, impararne i comandi, giocare. ``Per me potete anche andare a
prendere domino o mahjong'' (``mahjong è bellissimo tra l'altro'', aggiunge
Daniele, che sta facendo un videogioco sul mahjong), o il pachisi o il go. Gli
anni precedenti si chiedeva a ciascuno di proporre ai compagni un gioco da
tavolo che apprezzava. È anche un modo per conoscersi. Ruggero spiega il
perché: usarlo per giocare è facile e intuitivo, usarlo per creare moduli è più
complesso; ma chi arriva al 10 ottobre avendo già aperto il software e giocato,
troverà più intuitivo anche creare un modulo. E l'avvertimento: ``non riducetevi
all'ultimo, non fate una partita il 9 sera''.

\subsection{Tabletopia e altri vantaggi}

Tancredi racconta di aver provato a caricare un suo gioco su
\textbf{Tabletopia}. ``Questo è un vaso di Pandora'', risponde Daniele.
Tabletopia è molto usato, ma secondo lui il suo strumento di creazione dei moduli
è pessimo, ed è il motivo per cui la maggioranza dei moduli ufficiali dei
publisher è su Tabletop Simulator. Tabletopia ha un vantaggio vero, è
\textbf{gratuito}, ma ha scelte di ``qualità della vita'' poco felici.

Tabletop Simulator ha invece un vantaggio ``fenomenale'' che si tende a
sottovalutare: essendo su Steam, \textbf{il modulo si pubblica direttamente sul
Workshop di Steam}; si possono invitare gli amici a iscriversi e giocare con
loro, purché abbiano un account Steam. Rispetto ad altri software da imparare,
costa ``tra i 10 e i 20 euro, ragazzi: si parla proprio di spiccioli''. ``Io con
Tabletop Simulator ci ho playtestato la vita''.

Altre caratteristiche utili:
\begin{itemize}
  \item si possono creare \textbf{moduli non elencati} (\emph{unlisted}),
        perfetti da passare a un publisher prima che il gioco sia pubblicato:
        non sono pubblici, e nessuno può rubare i materiali;
  \item il software ha una grafica datata (``\emph{World Art 1998}'', scherza
        uno dei docenti), ma lo stanno portando su \textbf{Unity} con molti investimenti
        e aggiornamenti; ascolta la comunità e introduce funzioni utili anche per
        chi fa design;
  \item il suo server Discord è molto frequentato, e c'è un canale YouTube
        dedicato solo al design di giochi di carte\footnote{La trascrizione ne
        riporta il nome come ``Distraction Maker''; non è verificabile.} con un
        server Discord dove si fanno playtest: il 99\% è su Tabletop Simulator;
  \item molti editori, quando si propone loro un gioco da tavolo,
        \textbf{chiedono se c'è un modulo per Tabletop Simulator}.
\end{itemize}

% ---------------------------------------------------------------------
\section{Le domande a Daniele: il tema del project work}
\label{a1:sec:domande-daniele}

Daniele ha ``un minuto'' per rispondere alle domande sul corso prima di dover
andare. Uno studente chiede, ``giusto per curiosità'', se il project work avrà
un tema o sarà libero.

\begin{itemize}
  \item C'è un \textbf{mandate}, che i docenti propongono ogni anno e
        aggiornano leggermente; verrà presentato il \textbf{25 settembre}.
        ``Non facciamo spoiler''. Conterrà quello che si chiede di fare e quello
        che non si chiede: dei \textbf{must}, dei \textbf{must not} e dei
        \textbf{nice to have}.
  \item Verrà chiesto di lavorare su un'\textbf{idea originale} che non si è
        portata al corso: originale nel senso di \textbf{creata apposta}, non
        nel senso di ``mai vista prima''.
  \item L'idea va \textbf{partorita insieme}. Altri anni alcuni studenti hanno
        provato a imporre la propria idea al gruppo, ``e l'esperienza è stata
        pessima per tutti''. Quindi: \textbf{non imponete le vostre idee, fate
        brainstorming con gli altri}. Daniele, che fa ``molto spesso
        brainstorming'', darà strumenti per farlo: nessuno arriverà impreparato
        a quella lezione.
\end{itemize}

Ruggero aggiunge che il mandate è \textbf{abbastanza flessibile} perché tutti
possano fare qualcosa che li interessi: ``l'ultima cosa che vogliamo è forzarvi a
lavorare su qualcosa che non vi piace''. Bisogna lavorare su ciò che piace e per
cui si ha uno stimolo personale, anche nella \textbf{tipologia} di gioco: nessuno
sarà obbligato a fare un gioco di ruolo o un board game, e la lezione sui casi
studio mostrerà quante forme diverse di gioco analogico il corso accetta.

L'unica cosa da \textbf{non fare}, ``perché è contraddittorio per questo corso'',
sono i \textbf{giochi grandi}: ``lo chiediamo per voi''. Niente progetti
mastodontici\footnote{La trascrizione dice ``non fate delle cose da gnome'':
probabilmente il nome di un gioco da tavolo molto grande e complesso, non
ricostruibile con certezza.}, niente ``cose alla \emph{Dungeons \& Dragons}''.
Si può fare un gioco di ruolo, ma con un'esperienza ``meno corposa'': ci sono
tanti giochi di ruolo da 15 o 20 pagine, ``bellissimi, fichissimi'', che si
possono leggere online. Il consiglio è farsi una cultura anche su
\textbf{itch.io}, ``uno dei portali più importanti del gioco di ruolo
internazionale'', di cui si parlerà anche nella parte di ricerca e sviluppo.

Poi Daniele saluta (``noi ci vediamo venerdì a lezione'') e lascia la
conduzione a Ruggero.

% ---------------------------------------------------------------------
\section{Le domande finali con Ruggero}
\label{a1:sec:domande}

\subsection{Le regole delle domande}

Ruggero spiega che cosa vuol dire ``stare con lui''. Alla fine di ogni sua
lezione ci saranno \textbf{almeno 20--30 minuti di domande} libere, \emph{free
form}. Vuole che fare domande sia \textbf{normalizzato}: ``io gestisco malissimo
l'\emph{horror vacui} e divento super cringe molto alla svelta se non ci sono
domande in giro''. Sono benvenute le domande sulla lezione appena fatta, ma non
solo: il corso è crossmediale, e anche le domande che vengono da
\textbf{divagazioni}, da parti laterali, ``fatele''. Quello che sembra scontato o
banale a uno studente può non esserlo per il docente, e viceversa; e può essere
utile anche agli altri. ``Se vi sembrano stupide, fatele il doppio''. ``Magari
sembro incazzato, ma vi giuro sono rilassatissimo. Non mordo, magari abbaio un
po'\,''.

\subsection{Come approcciarsi a Steam e Tabletop Simulator}

Rita non ha mai usato né Steam né Tabletop Simulator e chiede come
approcciarsi, e se c'è modo di mettersi d'accordo tra compagni su Discord, magari
con un sondaggio sui giochi, per formare gruppi di gioco. Ruggero conferma: sul
Discord dell'accademia c'è un \textbf{canale apposito per il corso}, ed è degli
studenti. Lo possono usare per organizzarsi a giocare, per chiedere a un
compagno di prestare un gioco, per domande tecniche (come si scarica, come si
installa). ``Sarebbe proprio ideale che vi parlaste e vi organizzaste''.

Anche perché, aggiunge con un ``piccolo spoiler'' (``nessuno vi ha detto
niente''), se fossero i docenti a organizzare i gruppi di gioco, il task lo
starebbero facendo loro. Bisogna \textbf{abituarsi a organizzare le sessioni di
gioco}, perché arriveranno i task sul \textbf{playtest} e su come organizzarlo,
e a quel punto si sarà ``meno disperati''. Sembra scontato, ma saper
\textbf{organizzarsi per giocare in autonomia} ``vi sarà utilissimo come skill in
questo settore, ve lo giuro''.

\subsection{Il bilanciamento}

Angelo confessa ``una paranoia grossissima'': ci saranno lezioni sul
\textbf{bilanciamento}? ``Ma certo''. Sono proprio le lezioni di Marta sul
development: il bilanciamento è di solito un tema che tocca \textbf{i developer
più dei designer}. L'anno precedente c'era una sola lezione con Marta, e i
docenti si sono resi conto che serviva andare ancora più a fondo: per questo
quest'anno le lezioni sono due. ``Non è una paranoia, è una legittima domanda'':
Ruggero insegna da anni, vede altri corsi, parla con altri docenti e studenti, e
``mi incazzo spesso'' perché molti corsi di game design \textbf{non fanno
bilanciamento}. I docenti volevano invece un corso che coprisse tutte le parti
del game design: non solo l'idea del gioco, ma il bilanciamento, la produzione,
i playtest, l'iterazione, la scrittura dei regolamenti.

\subsection{Si può giocare con i docenti?}

Una studentessa chiede se si può giocare con i docenti. Per motivi etici,
durante il corso questo è limitato alle parti di \textbf{valutazione} dei giochi
e di \textbf{tutoraggio}. Dopo il corso, ``nessun problema, almeno con me''.

\subsection{Dove si trovano i giochi su Tabletop Simulator}

Ruggero mostra dove si trovano i giochi: nella propria libreria di Steam, nella
pagina di Tabletop Simulator, alla voce \textbf{Workshop}, dove ci sono ``più o
meno 36.000'' giochi. Una studentessa scopre intanto che l'app di Steam per
telefono è solo una \emph{companion app}: per registrarsi e giocare serve il
PC.

Per provare, chiede agli studenti un titolo:
\begin{itemize}
  \item \emph{Keep the Heroes Out}: non compare, ``non tutti i giochi ci sono,
        perché molti publisher non li mettono''; forse, essendoci già una
        versione del gioco su Steam, il modulo è stato tolto;
  \item \emph{Magic: The Gathering}: ``almeno 50 moduli'', anzi ``123 pagine'';
  \item \emph{Bang!}: tre moduli.
\end{itemize}
Bisogna fare attenzione: alcuni moduli sono in lingue specifiche (``magari il
primo è in russo''), alcuni hanno già l'espansione, altri sono solo il gioco
base. Conviene sfogliare e non fermarsi al primo risultato. E nel Workshop si
trovano anche cose molto particolari: giochi fuori produzione, giochi vecchi,
giochi ancora in prototipo.

\subsection{I ruoli nei team e le dinamiche tra le persone}

Uno studente torna sulla slide delle figure: nelle realtà piccole molti ruoli
hanno ``il doppio o triplo cappello''. Quali ruoli restano distinti, e come si
creano le sinergie tra persone con obiettivi diversi?

Ruggero premette che l'ideale ``platonico'' è un ambiente di lavoro costruttivo,
basato sulla fiducia e ``il più \emph{safe} possibile''; ma nell'industria non
sempre è così, e le strutture variano moltissimo: alcuni publisher sono molto
verticali, altri hanno ruoli già solidificati (se un publisher ha già i suoi
developer e non sta crescendo, entrarci come developer è praticamente
impossibile). ``Se prendiamo due aziende in Italia, vediamo due strutture
interne completamente diverse''. Non esiste una scorciatoia per farsi assumere
da un editore.

Poi prende ``l'esempio più comune di tutti'', la media dei giochi da tavolo
pubblicati (figura~\ref{a1:fig:relazioni}):
\begin{itemize}
  \item il \textbf{game designer} è una figura \textbf{esterna}, con un
        \textbf{contratto a royalty} con il publisher;
  \item la sua figura di riferimento è il \textbf{publisher} stesso, e la
        conversazione designer--publisher è ``molto calda'';
  \item se il publisher ha dei \textbf{developer}, il developer entra in questa
        conversazione: designer e developer collaborano molto, e internamente
        developer e publisher collaborano molto. Spesso, nei publisher piccoli,
        è il publisher a fare le veci del developer;
  \item \textbf{graphic design} e \textbf{illustrazione} tendono a parlare molto
        con il publisher e \textbf{molto poco con il designer}. A meno che non si
        tratti di uno studio indie molto piccolo, è difficile che il designer
        sappia anche solo chi è l'illustratore prima che il gioco esca.
\end{itemize}

\begin{figure}[htbp]
\centering
\begin{tikzpicture}[
    font=\small,
    ruolo/.style={draw, thick, rounded corners=3pt, fill=gray!12,
                  minimum width=2.8cm, minimum height=0.9cm, align=center},
    esterno/.style={draw, thick, dashed, rounded corners=3pt,
                    minimum width=2.8cm, minimum height=0.9cm, align=center},
    calda/.style={{Stealth[length=2.5mm]}-{Stealth[length=2.5mm]}, very thick},
    tiepida/.style={{Stealth[length=2mm]}-{Stealth[length=2mm]}, thin, dashed}
  ]
  \node[esterno] (gd) at (0,0) {Game designer\\ \footnotesize esterno, a royalty};
  \node[ruolo] (dev) at (5,1.4) {Developer\\ \footnotesize se c'è};
  \node[ruolo] (pub) at (5,-1.4) {Publisher};
  \node[ruolo] (art) at (10,-1.4) {Graphic designer\\ e illustratore};
  \draw[calda] (gd) -- (dev);
  \draw[calda] (gd) -- (pub);
  \draw[calda] (dev) -- (pub);
  \draw[calda] (pub) -- (art);
  \draw[tiepida] (gd.south) to[bend right=25]
     node[below, font=\footnotesize\itshape] {quasi nessun contatto} (art.south);
  \node[draw, dotted, rounded corners=4pt, fit=(dev)(pub)(art), inner sep=8pt,
        label={[font=\footnotesize\itshape]above right:interno al publisher (di solito)}] {};
\end{tikzpicture}
\caption{Le relazioni tipiche intorno a un gioco da tavolo pubblicato, secondo la
risposta di Ruggero. Le frecce spesse indicano una collaborazione stretta e
frequente; quella tratteggiata, un contatto quasi assente. Nei publisher piccoli
il publisher fa spesso anche da developer.}
\label{a1:fig:relazioni}
\end{figure}

\subsection{Fare da developer di sé stessi?}

Rita chiede in chat se un designer può fare anche da developer del proprio
gioco, o se Ruggero lo sconsiglia. In generale \textbf{lo sconsiglia}, per lo
stesso motivo per cui si sconsiglia in altri media: i \textbf{darling}. Il
designer ha delle meccaniche a cui è affezionato, e il developer a volte dovrà
\textbf{tagliarle}; è molto difficile, se si è ``gestori di sé stessi'', tagliare
cose che si vogliono nel gioco. ``È un po' come lasciare fare il lavoro di editing
allo scrittore che ha già scritto il libro'': può rivedersi, certo, ma difficilmente
cambierà cose che farebbe bene cambiare. Nelle lezioni con Marta si vedrà perché,
e alla fine del corso ``non avrete dubbi nemmeno voi'' nel separare designer e
developer.

Ma tutto dipende dallo \textbf{scopo del progetto}: se si fa un gioco in tre,
qualcuno quel lavoro deve farlo per forza, e magari il designer fa anche il
developer, compensando con \textbf{più playtest}. ``Il mondo attuale e il mondo
ideale sono due cose molto diverse, e lo vediamo non solo dai giochi ma dalla
politica'', ma ``non entriamo in questo vaso di Pandora''.

\subsection{Il ruolo indispensabile e la pubblicazione}

Lo studente di prima ne trae una conclusione: senza game designer il gioco non
esiste. ``Quello sicuramente''. Ma Ruggero coglie l'occasione per approfondire
un punto che forse non era passato: \textbf{fare design di un gioco per
pubblicarlo non è l'unica opzione possibile}. Un gioco non è ``meno gioco''
se non è pubblicato da un editore. È legittimo mettere un gioco su itch.io in
\textbf{print and play}, fare un gioco a evento senza venderlo, fare un
\emph{murder party}, una escape room, un board game per una persona amica:
``quell'esperienza lì non è meno gioco''. Dipende dall'obiettivo del progetto.

Ruggero ha lavorato molto anche con le \textbf{scuole}, dove il publisher non
esiste: è un gioco per le classi, e la figura più vicina è l'editoria scolastica,
``che però non fa il publisher''. Ci sono compromessi diversi sulla materialità
e su come il gioco si presenta, che si vedranno nella prototipazione e nei casi
studio. Quindi l'unica figura che deve esserci sempre è \textbf{chi fa game
design}, anche se magari non sa di chiamarsi game designer e non ha mai
sentito parlare di game design: ``questo non la rende meno game designer''.

\paragraph{Un libro consigliato.} Ruggero fa anche mentoring e gestisce gruppi
di settore, e spesso gli chiedono che cosa leggere per imparare il game design.
Non è il tipo che obbliga a leggere libri (``non vi obbligherò a leggere
niente''), ma ce n'è uno che consiglia sempre, ed è il primo che ha scelto di
leggere: \textbf{\emph{Rise of the Videogame Zinesters}} di \textbf{Anna
Anthropy}\footnote{La trascrizione dice ``Anna Antrofi'' e ``Rise of the
Videogames in Esters''.}. È un libro un po' datato, ma ancora reperibile; parla
soprattutto di videogiochi, ma ``per l'80\%'' il contenuto si applica al gioco e
al game design in generale. È un libro sulla \textbf{democratizzazione del gioco}:
il gioco come strumento per parlare di sé e come \textbf{medium legittimo}. Come
esistono film drammatici, tristi, intensi, possono esistere giochi su tematiche
delicate, e se ne vedranno nei casi studio. Molti giochi fatti con questa
direttiva non hanno un publisher, che tende a pubblicare ciò che dà un ritorno
commerciale.

\paragraph{Commerciale e personale.} ``Io spero che voi non farete giochi solo
per scopo commerciale''. Certo, servono un tetto sulla testa, il cibo in
tavola, un posto in cui dormire, e così si finisce per fare giochi che non sono
il proprio gioco dei sogni. Spesso è un mix: Ruggero ha fatto giochi commerciali
``per pagarmi l'affitto'' e giochi che voleva fare e su cui ha perso soldi,
``sapendolo proprio che sarebbero stati un fallimento''. I giochi, come i libri e
le canzoni, si fanno per molti motivi, e quelli commerciali di successo non sono
più importanti degli altri. Tra i giochi della sua slide c'erano il gioco da
tavolo ufficiale di \emph{Cyberpunk}\footnote{La trascrizione dice ``Cyberpunk
del 1977'': quasi certamente un gioco legato a \emph{Cyberpunk 2077}.} e
\emph{Marvel United}, su IP grosse e famose, ma anche giochi fatti di ``una
candela e una busta su un tavolo''. E se gli si chiede quale preferisce in
termini di contenuto, non sono quelli sui supereroi.

\subsection{Il gioco di ruolo in Italia}

Nicola, che è interessato soprattutto al gioco di ruolo e conosce già persone
del settore, ha ``una leggera paura'': che chi studia game design finisca per
fare il \textbf{freelance} a vita, ``perché non ci sarà mai un'occupazione per
bene'', o il \textbf{copywriter}. Conosce un ragazzo che ha lavorato in quel
mondo e che ora fa altro.

Ruggero risponde che la maggioranza dei giochi nel suo portfolio sono giochi di
ruolo, quindi il tema lo tocca da vicino, e che in Italia quello che dice Nicola
``è semplicemente vero''. Molti publisher piccoli hanno una persona di
riferimento che copre più ruoli: impagina, fa i playtest, fa un po' di
development. È ``un settore economicamente completamente disastrato'', per
motivi culturali che non apre: oggi in Italia è praticamente impossibile
\textbf{vivere di giochi di ruolo}. Chi ci lavora lo fa nel tempo libero, al di
fuori del lavoro principale, oppure coprendo molte altre cose: molti publisher
italiani di giochi di ruolo pubblicano anche titoli esteri e localizzano, non
fanno solo cose originali. ``È un grosso calderone''. Poi si ferma: ``questo è
un \emph{rant}'', e ``i \emph{rant} tornano sempre'', ma non vuole che gli
studenti escano dalla prima lezione ``con la disperazione nel cuore''. Ci sarà
spazio per parlarne quando avranno gli strumenti per capirlo, ``per non farvi
abbattere''.

\subsection{Che cosa bisogna sapere già}

Uno studente chiede se ci sono nozioni che dovrebbero già avere e che possono
recuperare: per esempio, sente parlare di giochi ``German'' o ``American'' e non
sa che cosa siano. Ruggero spiega l'organizzazione del corso:
\begin{itemize}
  \item le \textbf{prime due lezioni} sono di fondamenta. Questa era ``dall'alto'':
        come è strutturato il corso, le tempistiche;
  \item la \textbf{prossima lezione} è di tre ore di soli \textbf{casi studio},
        ``a manetta, a raffica'', divisi per categorie e tipologie, con
        ragionamenti su categorie e generi di gioco: lì si vedranno German,
        American ``e anche altre più utili'';
  \item nel resto del corso le terminologie specifiche sono inserite nelle
        lezioni dedicate: il \textbf{paragame} nelle lezioni sul paragame, il
        \textbf{modello matematico} nelle lezioni di Marta sul development.
\end{itemize}
Per ora agli studenti si chiede di sapere quello che si è visto oggi:
meccaniche, dinamiche ed estetiche; le figure professionali, designer e
developer; e soprattutto \textbf{game e play}.

\subsection{Game e play: i quattro quadranti}

``Specialmente \emph{game} e \emph{play}'', insiste Ruggero: è più difficile di
quello che sembra, perché l'italiano non abitua a separare le due cose, ed è
possibile averle usate per anni come sinonimi. Nella sua esperienza con chi si
avvicina al design è ``uno dei punti più critici''. Spesso i commenti che si
fanno sui giochi giocati sono commenti sul \emph{play} applicati al
\emph{game}: ``il gioco non mi è piaciuto perché, giocando con Tizio, è successo
questo''; ``il gioco mi sembra sbilanciato perché Tizio al tavolo ha fatto
questo''.

Separare \emph{game} e \emph{play} serve a capire quali giochi sono progettati
bene e quali male, e soprattutto a riconoscere i \textbf{quattro possibili
incroci} (figura~\ref{a1:fig:quadranti}): un gioco progettato bene che piace
giocare; uno progettato male che non piace; uno progettato male che piace; uno
progettato bene che non piace. ``Non è importantissimo in generale nella vita,
ci sono cose più importanti'', ma per chi fa game design lo è: ``a volte vi
chiederanno di lavorare su un gioco che non vi piace giocare, e quel gioco deve
essere fatto bene lo stesso''.

\begin{figure}[htbp]
\centering
\begin{tikzpicture}[
    font=\small,
    cella/.style={draw, thick, minimum width=4.4cm, minimum height=1.9cm,
                  text width=4.0cm, align=center}
  ]
  \node[cella, fill=gray!20] (a) at (0,0)     {\textbf{progettato bene}\\ e \textbf{mi piace} giocarlo};
  \node[cella]               (b) at (4.4,0)   {\textbf{progettato bene}\\ ma \textbf{non mi piace} giocarlo};
  \node[cella]               (c) at (0,-1.9)  {\textbf{progettato male}\\ ma \textbf{mi piace} giocarlo};
  \node[cella, fill=gray!8]  (d) at (4.4,-1.9){\textbf{progettato male}\\ e \textbf{non mi piace} giocarlo};
  \node[font=\footnotesize\itshape, above=0.15cm of a] {mi piace (\emph{play})};
  \node[font=\footnotesize\itshape, above=0.15cm of b] {non mi piace (\emph{play})};
  \node[font=\footnotesize\itshape, rotate=90, left=0.15cm of a, anchor=south]
       {bene (\emph{game})};
  \node[font=\footnotesize\itshape, rotate=90, left=0.15cm of c, anchor=south]
       {male (\emph{game})};
\end{tikzpicture}
\caption{I quattro incroci tra la qualità del \emph{game} (come è progettato) e
l'esperienza di \emph{play} (se piace giocarlo), descritti da Ruggero. Per un
designer i quadranti ``scomodi'' sono i due fuori diagonale: un gioco può
piacere pur essendo progettato male, e va saputo progettare bene anche un gioco
che non piace.}
\label{a1:fig:quadranti}
\end{figure}

\paragraph{La domanda di Tancredi.} Tancredi chiede conferma: il \emph{play}
non è il ``giocare senza regole'', ma può essere anche il giocare a un gioco? Sì:
il \emph{play} è l'\textbf{atto}, e può essere fatto con un \emph{game} o senza.
E un \emph{game} può esistere senza che nessuno lo giochi: gli scacchi non sono
meno scacchi perché in quel momento nessuno ci gioca. E se, giocando, una
persona applica le regole in modo diverso? Allora \textbf{non si sta giocando
quel gioco}.

\subsection{House rule: il caso di Monopoly}

Qui Ruggero fa un po' di \emph{foreshadowing} su una lezione di Daniele: il
discorso degli \textbf{hack}, delle \textbf{house rule}, delle \textbf{mod},
cioè di come modificare un gioco esistente. Il punto è che \textbf{un gioco
modificato da una house rule è un altro gioco}: cambiando le regole cambiano
necessariamente anche le \textbf{dinamiche} che ne derivano, e quindi
l'esperienza. Si immagini \emph{Monopoly} in cui ogni turno si possa, volendo,
prendere gratis dei soldi dalla banca: tutto il resto è identico, eppure il gioco
cambia.

E questo è talmente interiorizzato che ci sono regole che si applicano ``senza
neanche rendercene conto''. Ruggero chiede quanto duri in media una partita a
\emph{Monopoly}. Gianmarco: ``quattro ore''; altri: ``un pomeriggio''; un altro:
``finché uno dei giocatori non aggiunge la regola: ora vince chi\ldots''. In
realtà, \textbf{seguendo le regole scritte}, una partita non dura mai più di
\textbf{un'ora e venti, un'ora e mezza}. La domanda che fa sempre, e la risposta
è sempre la stessa: perché si pensa che \emph{Monopoly} duri tantissimo? Qual è
la regola che nessuno rispetta?

Uno studente indovina: \textbf{le aste}. Secondo il regolamento, se chi finisce
su un terreno non lo compra, il terreno \textbf{va all'asta} tra i giocatori, e
chiunque può aggiudicarselo. Molti invece giocano come se finire su un terreno
desse solo la \emph{facoltà} di comprarlo o di metterlo all'asta\footnote{Più
precisamente, nel regolamento ufficiale, se il giocatore che finisce su una
proprietà libera decide di non comprarla, la banca la mette all'asta e tutti i
giocatori possono fare offerte, compreso chi l'ha rifiutata. L'asta quindi non è
una scelta del giocatore, ma la conseguenza obbligata del suo rifiuto.}. Con la
regola applicata, di aste se ne fanno moltissime: all'ottavo turno ci sono già
otto terreni assegnati, al ventesimo venti, la gente va in bancarotta
velocissimamente e la partita dura un'ora e mezza. Anche le partite competitive
di \emph{Monopoly}, nota Ruggero, durano un'ora e venti, un'ora e mezza.

Uno studente obietta che giocava con le aste e durava comunque quattro ore.
``Avete dei nemici con l'analisi paralisi'', ride Ruggero; e cita un'altra house
rule tipica dell'infanzia: genitori che danno soldi per non far andare in
bancarotta. Ma \textbf{la bancarotta è parte del gioco}: se la si toglie,
\emph{Monopoly} non regge. Sul tema che \emph{Monopoly} vuole trattare, e sul
suo antenato (\emph{The Landlord's Game}), ``non apriamo adesso un altro vaso di
Pandora''.

\paragraph{Approfondimento: Monopoly, la durata e le house rule nei libri.}
L'esempio di Ruggero si ritrova quasi identico in \emph{The Art of Game Design}.
Parlando del bilanciamento tra partite brevi e lunghe (\emph{Short vs. Long}),
Schell scrive che \emph{Monopoly}, giocato secondo le regole ufficiali,
``finisce spesso in circa novanta minuti'', ma molti giocatori trovano quelle
regole troppo dure e le modificano, distribuendo premi in denaro e allentando i
vincoli sull'acquisto delle proprietà: con l'effetto collaterale di far durare la
partita molto di più, ``tipicamente tre ore, o anche di più''. Schell classifica
le \textbf{house rule} come regole create dai giocatori in risposta a un
difetto percepito dopo qualche partita. Anche \emph{Challenges for Game
Designers}, nel capitolo sull'aggiungere e togliere meccaniche, parte dallo
stesso gioco: il \emph{Parcheggio gratuito}, che per le regole originali non fa
succedere nulla, mentre quasi tutti giocano con una house rule che dà un bonus a
chi ci finisce sopra. Le due fonti confermano la lezione: una piccola regola
cambiata cambia le dinamiche, e le dinamiche cambiano l'esperienza.

\subsection{Il fine del game è il play?}

Un altro studente commenta: va bene la distinzione, ma il fine del \emph{game} è
il \emph{play}; quello che si vede è sempre l'effetto. Ruggero è d'accordo: di
solito, se si fa un gioco, lo si fa per farlo giocare, come si scrive un libro
per farlo leggere (anche se esistono, ed ``è già stato esplorato'', libri scritti
per non essere letti e giochi per zero giocatori). Ma la distinzione serve a una
cosa precisa: \textbf{separare, nel playtest, i feedback sul gioco dai feedback
sul gruppo di gioco} e sulle sue dinamiche. Se una persona esce frustratissima da
un playtest, non è necessariamente colpa del gioco; se una persona dice ``questo
è il gioco migliore del mondo'', potrebbe non essere merito del gioco ma del
gruppo. Capire a che cosa si riferisce un feedback serve a capire \textbf{dove
iterare} il prototipo.

Lo studente osserva che l'area è sfumata: alcuni giochi si basano proprio
sull'avere le persone giuste al tavolo. Un gioco di \textbf{negoziazione} con
giocatori abituati ai giochi ``multi-solitario'' diventa ingiocabile; e nei giochi
a \textbf{economia chiusa} (cita \emph{Container}\footnote{La trascrizione cita
anche un secondo titolo non ricostruibile (``Tokyo su Vigimarket'').}) un gruppo
economicamente conservativo, che spende poco e non fa girare l'economia, fa
durare la partita il doppio. Ruggero rilancia: e se il gioco dura il doppio,
questa dinamica cambia la \textbf{percezione di peso} del gioco, le energie spese,
quanto si è stanchi a fine partita. ``Tutto è collegato. \textbf{I giochi sono
ecosistemi di meccaniche}, e cambiandone una, le altre cambiano di
conseguenza''.

\paragraph{La parola vietata, alla prova.} Lo studente aggiunge che ci sono giochi
che ha trovato ``molto sbagliati, ma che mi divertono''. ``Attenzione: non si può
dire \emph{divertimento}!'', lo ferma Ruggero; nelle presentazioni aveva lasciato
correre, ``ma lì non vi avevo ancora letto la regola''. Lo studente ci prova:
``giochi che poi voglio \textbf{rigiocare}'', anche se ``sono evidentemente
sbagliati, rotti''. Ruggero apprezza, ma non vuole anticipare cose che meritano il
loro tempo: se ne parlerà nelle lezioni sui framework e soprattutto in quella di
\textbf{emotion design}, ``forse la mia lezione preferita del corso''. (Lo
studente ammette che ``forse alla fine si gioca per quello, per l'emozione''.)

\subsection{I giochi ``comunicati male''}

L'ultima domanda corposa riprende un'affermazione di inizio lezione: ci sono
giochi belli ma comunicati male. Qualche esempio? ``Mi basta girarmi per
vederne 48'', risponde Ruggero indicando gli scaffali alle sue spalle. Più che
titoli, elenca le \textbf{ragioni}:
\begin{itemize}
  \item \textbf{l'Italia localizza lentamente}. Basta guardare i premi del gioco
        dell'anno in Italia: spesso vanno a giochi usciti due o tre anni prima in
        Germania, Francia, Stati Uniti o Regno Unito, perché localizzare richiede i
        suoi tempi. Il mercato italiano è quindi molto dipendente dall'estero,
        ``molto impaurito'', e spesso aspetta che un gioco vada bene all'estero
        prima di portarlo in Italia; e alcuni giochi ``perdono il treno'';
  \item \textbf{il publisher originale chiude}: come si pubblica un gioco se non
        c'è più l'azienda con cui fare un contratto? A Ruggero è successo più
        volte;
  \item \textbf{il dipartimento commerciale italiano non sa comunicare il gioco}:
        si fa male un'edizione, la stampa vende pochissimo e non viene più
        ristampata;
  \item \textbf{la scala del mercato}: in Italia un gioco da tavolo di peso medio
        (un \emph{eurogame}, un \emph{German}) che vende 1.500 copie ``è tanto'';
        all'estero se ne vendono magari 10.000. Con 10.000 copie in inglese ci
        sono i soldi per una campagna di comunicazione: un trailer, i
        recensori, i tutorial online. In Italia tutto questo è ``ridotto all'osso'',
        e la cosa si riflette su stampa e divulgazione: anche i canali per
        informarsi su un gioco prima di comprarlo sono pochi (se ne parlerà nella
        lezione su ricerca e sviluppo). Un gioco spesso non arriva al mercato non
        perché sia brutto, ma perché l'editore non lo sa comunicare, ne fa poche
        copie e va fuori produzione;
  \item \textbf{il tempismo}: party game progettati nel 2019 e usciti nel 2020, in
        pieno Covid, quando non si poteva giocare al tavolo, sono falliti tutti, e
        gli editori ci hanno perso. ``Non potevano immaginarselo'';
  \item \textbf{i temi caldi}: dopo l'invasione russa dell'Ucraina e la guerra in
        Palestina, alcuni editori non pubblicano o non localizzano giochi
        strategici e di guerra. Ruggero ha lavorato a \textbf{\emph{Status}}, un
        gioco sulla geopolitica che tocca, tra gli altri temi, l'invasione russa
        dell'Ucraina e il genocidio in Palestina, e alcuni editori che di
        solito localizzano i loro giochi (per esempio il partner di riferimento in
        Russia) non lo hanno localizzato per motivi politici.
\end{itemize}
Lo studente voleva capire che cosa voglia dire comunicare male: non metterci
abbastanza energia, o comunicare una cosa per un'altra, ``come il trailer
sbagliato, che poi vai al cinema e ti trovi un altro film''. ``È un mix'', e la
domanda ``era complessa e meritava più spazio''.

\subsection{I giochi di carte collezionabili e i saluti}

L'ultima domanda (``posso rompere le balle un'ultima volta?'') riguarda i
\textbf{giochi di carte collezionabili}: funzionano con le stesse regole dei
giochi da tavolo? Se ne parlerà alla prossima lezione: è proprio il motivo per
cui la lezione sui casi studio dura tre ore, perché ogni genere di gioco ha le
sue dinamiche e le sue regole, e anche il modo di pubblicare cambia molto tra un
TCG, un gioco di ruolo da tavolo, un LARP e un gioco pervasivo.

Ruggero saluta: le email dei tre docenti sono all'inizio delle slide, la
registrazione arriverà, ``ci si vede questo venerdì''.

% ---------------------------------------------------------------------
\section*{Punti chiave del capitolo}
\addcontentsline{toc}{section}{Punti chiave del capitolo}

\begin{itemize}
  \item Il corso di \textbf{Analog Game Design} (seconda edizione) dura circa tre
        mesi, fino a dicembre, ed è tenuto da tre docenti principali:
        \textbf{Daniele} (narrative design, gioco di ruolo, fiction),
        \textbf{Ruggero} (game design, ricerca e sviluppo) e \textbf{Beatrice}
        (giochi pervasivi), con tre ospiti su RNG, development e game production.
  \item Lo scopo è formare \textbf{autori analogici} consapevoli: strumenti da
        usare anche fuori dal corso, non il ``gioco dei sogni''. Il corso è
        \textbf{intersezionale} (nessuna lezione dedicata a un solo tipo di
        gioco), \textbf{cross-generazionale}, \textbf{pratico e adattabile},
        basato sulle \textbf{domande} (nessuna domanda è stupida), sul
        \textbf{feedback} invece che sul voto e sul \emph{fail forward}.
  \item Il \textbf{project work} si fa in tre gruppi da tre, formati dai docenti;
        pitch il 12 ottobre, sviluppo nelle due pause, tutoraggi il sabato,
        consegna il 4 dicembre, presentazione l'11 e giornata finale con gli
        editori il 14. Si valutano documentazione, presentazione,
        autovalutazione, crescita e atteggiamento.
  \item \textbf{Niente intelligenza artificiale} nel corso: i publisher
        rifiutano pitch e prototipi con immagini generate, i regolamenti scritti
        dall'IA descrivono la stessa meccanica in modi diversi, e soprattutto chi
        la usa perde l'occasione di imparare. Un prototipo può essere fatto a
        matita.
  \item Le \textbf{figure dell'industria}: game designer, game developer,
        editor, publisher, sales e marketing, graphic designer, illustratore,
        sensitivity reader (che non è un censore) e producer. Nelle realtà
        piccole una persona copre più ruoli.
  \item Nell'analogico il \textbf{game designer} progetta l'esperienza creando
        le meccaniche; il \textbf{game developer} fa da \textbf{mediatore} con il
        publisher (bilanciamento, modelli matematici, adattamenti). Pretendere
        che il proprio gioco non venga modificato è un \emph{red flag}.
  \item Il \textbf{marketing} non decide che cosa pubblicare ma come; è
        un'azione, ed esiste anche se lo si ignora.
  \item In \textbf{Italia} il gioco è ancora visto come infantile, il mercato
        punta sulla localizzazione e investe poco su giochi e designer interni;
        all'estero (soprattutto dalla Germania in su) ci sono tradizione adulta,
        premi di prestigio e developer interni.
  \item \textbf{Game} è l'artefatto (le regole, il sistema, l'oggetto);
        \textbf{play} è l'atto di giocare. Si fanno giochi, non giocattoli.
        Separarli serve a distinguere i feedback sul gioco da quelli sul gruppo
        e a riconoscere i quattro incroci tra ``progettato bene/male'' e ``mi
        piace/non mi piace''.
  \item \textbf{MDA}: meccaniche (regole scritte), dinamiche (regole in azione),
        estetiche (esperienza, non grafica); nel \textbf{DPE}: design, play,
        experience.
  \item \textbf{Parole vietate} in ogni materiale scritto: \emph{fun},
        \emph{immersion}, \emph{entertainment} (divertimento, immersione,
        intrattenimento). Bisogna trovare parole più precise.
  \item Distinguere \textbf{chi crea, chi gioca e chi facilita}; nessun gioco
        nasce perfetto: \textbf{playtest e iterazione} sono inevitabili, e un
        \textbf{dizionario comune} aiuta a non perdersi.
  \item \textbf{Tabletop Simulator} è il software di settore per giocare (non
        creare videogiochi) i giochi da tavolo in digitale. Primo task: giocarci
        tre o quattro volte entro il 10 ottobre.
  \item Una \textbf{house rule} crea un altro gioco: \emph{Monopoly} con le aste
        applicate come da regolamento dura circa un'ora e mezza.
\end{itemize}

%    % =====================================================================
%  Analog Game Design, lezione 2 — I casi studio: una panoramica a volo
%  d'aquila dei giochi pervasivi, dei giochi di ruolo dal vivo, dei
%  giochi di ruolo da tavolo e dei giochi da tavolo
%  Fonte: trascrizione della seconda lezione (AGD2) del corso breve di
%  Analog Game Design, seconda edizione: una lezione online su Zoom di
%  circa tre ore, di venerdì sera, tenuta a più voci dai tre docenti
%  "cuore" del corso (Ruggero, Daniele e Beatrice), con una pausa di
%  pochi minuti e moltissime slide di esempi.
%  Pacchetti richiesti nel preambolo del libro (già presenti nel
%  main.tex): \usepackage{booktabs}, \usepackage{tikz} con
%    \usetikzlibrary{shapes.geometric,shapes.symbols,shapes.misc,
%      arrows.meta,positioning,calc,fit}
%  (le librerie vanno caricate nel preambolo, non nel capitolo).
%  Consigliati, come per gli altri capitoli: \usepackage[italian]{babel}
%  e \usepackage[T1]{fontenc} (il capitolo compila anche senza).
%
%  Collocazione nel libro: parte "Analog Game Design", secondo capitolo
%  (numerazione propria 2.2). Nel file principale del libro, subito dopo
%  il capitolo 2.1:
%
%    % =====================================================================
%  Analog Game Design, lezione 1 — Introduzione al corso: docenti,
%  metodo, project work, figure dell'industria e primo glossario
%  Fonte: trascrizione della prima lezione (AGD1) del corso breve di
%  Analog Game Design, seconda edizione: una lezione online su Zoom di
%  circa tre ore, di pomeriggio, tenuta insieme dai tre docenti "cuore"
%  del corso (Daniele, Ruggero e Beatrice), con nove iscritti, una
%  pausa di una decina di minuti, molte slide e una breve dimostrazione
%  dal vivo di Tabletop Simulator.
%  Pacchetti richiesti nel preambolo del libro (già presenti nel
%  main.tex): \usepackage{booktabs}, \usepackage{tikz} con
%    \usetikzlibrary{shapes.geometric,shapes.symbols,shapes.misc,
%      arrows.meta,positioning,calc,fit}
%  (le librerie vanno caricate nel preambolo, non nel capitolo).
%  Consigliati, come per gli altri capitoli: \usepackage[italian]{babel}
%  e \usepackage[T1]{fontenc} (il capitolo compila anche senza).
%
%  Collocazione nel libro: parte "Analog Game Design", primo capitolo
%  (numerazione propria 2.1). Nel file principale del libro:
%
%    \part{Analog Game Design}
%    \renewcommand{\thechapter}{2.\arabic{chapter}}
%    \setcounter{chapter}{0}
%    % =====================================================================
%  Analog Game Design, lezione 1 — Introduzione al corso: docenti,
%  metodo, project work, figure dell'industria e primo glossario
%  Fonte: trascrizione della prima lezione (AGD1) del corso breve di
%  Analog Game Design, seconda edizione: una lezione online su Zoom di
%  circa tre ore, di pomeriggio, tenuta insieme dai tre docenti "cuore"
%  del corso (Daniele, Ruggero e Beatrice), con nove iscritti, una
%  pausa di una decina di minuti, molte slide e una breve dimostrazione
%  dal vivo di Tabletop Simulator.
%  Pacchetti richiesti nel preambolo del libro (già presenti nel
%  main.tex): \usepackage{booktabs}, \usepackage{tikz} con
%    \usetikzlibrary{shapes.geometric,shapes.symbols,shapes.misc,
%      arrows.meta,positioning,calc,fit}
%  (le librerie vanno caricate nel preambolo, non nel capitolo).
%  Consigliati, come per gli altri capitoli: \usepackage[italian]{babel}
%  e \usepackage[T1]{fontenc} (il capitolo compila anche senza).
%
%  Collocazione nel libro: parte "Analog Game Design", primo capitolo
%  (numerazione propria 2.1). Nel file principale del libro:
%
%    \part{Analog Game Design}
%    \renewcommand{\thechapter}{2.\arabic{chapter}}
%    \setcounter{chapter}{0}
%    \input{AGD/AGD1_introduzione_all_analog_game_design}
%
%  (Con hyperref conviene aggiungere anche
%  \renewcommand{\theHchapter}{agd.\arabic{chapter}} nella parte, per
%  evitare ancore duplicate con i capitoli 1, 2, ... di Game Design.)
%  Tutte le etichette di questo capitolo hanno il prefisso "a1:"
%  (a1:cap:..., a1:sec:..., a1:fig:..., a1:tab:...). Il capitolo è
%  autonomo: non contiene \ref verso altri capitoli.
%  Le figure sono ricostruzioni schematiche di quanto i docenti spiegano
%  a voce e con le slide, non copie delle slide.
% =====================================================================

\chapter[Introduzione all'Analog Game Design]{Introduzione all'Analog Game
Design: il corso, le figure dell'industria e un primo linguaggio comune}
\label{a1:cap:intro}

Questa è la \textbf{prima lezione} del corso di \textbf{Analog Game Design},
cioè di progettazione di giochi \emph{analogici}: giochi da tavolo, giochi di
carte, giochi di ruolo da tavolo, giochi di ruolo dal vivo (LARP), giochi
\emph{pervasivi} ed esperienze a evento. È un corso \textbf{breve}, di circa tre
mesi (``da qui a dicembre''), separato dal corso di Game Design, e siamo alla sua
\textbf{seconda edizione}. Come avvertono subito i docenti, è ``una giornata di
introduzione'': si spiega come funziona il corso, chi sono i docenti, con quali
metodi e con quale calendario si lavorerà, e ci si presenta a vicenda. ``C'è poco
gioco analogico'', si scusano. Ma nella seconda parte la lezione comincia
comunque a costruire un \textbf{linguaggio comune}: le figure professionali
dell'industria del gioco analogico, la differenza tra \emph{game design} e
\emph{game development}, la situazione del mercato italiano rispetto a quello
estero, e un primo glossario di design (\emph{game} e \emph{play}, il framework
MDA, le parole ``vietate'', il playtest e l'iterazione). Si chiude con lo
strumento di settore per i prototipi digitali dei giochi da tavolo,
\textbf{Tabletop Simulator}, e con una lunga sessione di domande.

La lezione si svolge \textbf{online}, su Zoom, ed è condotta a più voci dai tre
docenti principali del corso: \textbf{Daniele}, narrative designer, che guida la
parte organizzativa; \textbf{Ruggero}, game designer, che guida la parte sui
fondamenti e le domande finali; \textbf{Beatrice}, progettista di giochi
pervasivi, che interviene soprattutto nella prima metà. Gli studenti iscritti
sono \textbf{nove}; uno è assente e un altro arriva a lezione iniziata. Ci sono
una pausa di una decina di minuti (si riprende alle 18:15) e, dopo le 19:15,
un'ultima mezz'ora di domande con il solo Ruggero.

Tre avvertenze di lettura. La prima: la trascrizione è automatica e storpia
molti nomi propri, compresi quelli dei docenti e dell'accademia (``Dive'',
``IVE'', ``Accademia Live'', ``Ai''); dove la ricostruzione è sicura, o molto
probabile, il capitolo usa la forma corretta e lo segnala in nota; dove non lo
è, riporta la forma della trascrizione e lo dice. La seconda: durante la pausa la
trascrizione registra, ripetuta in loop per decine di righe, una frase senza
senso che non ha nulla a che fare con la lezione, e che il capitolo ignora. La
terza: in molti passaggi più docenti parlano uno dopo l'altro senza che la
trascrizione segni il cambio di voce; il capitolo attribuisce le parole a chi le
pronuncia quando il contesto lo rende chiaro (chi parla di sé, chi viene
chiamato per nome, chi ``passa la parola''), e resta generico altrimenti.

Quanto alla data, la lezione si colloca con ogni probabilità \textbf{nella prima
metà di settembre del 2026}\footnote{La data non è mai detta. La si deduce da
alcuni indizi: il form iniziale va compilato ``entro e non oltre il 15
settembre''; i gruppi verranno annunciati il 25 settembre; le date del
calendario (sabato 3 ottobre, lunedì 12 ottobre, sabato 7 e 14 novembre,
venerdì 4 dicembre) cadono nei giorni della settimana indicati soltanto nel
2026; un docente parla di un argomento ``molto complesso nel 2026''; e fa
ancora molto caldo (``le temperature sono ancora indecenti''). È una
deduzione, non un'informazione esplicita.}: i riferimenti a ``oggi'',
``quest'estate'' e ``l'anno scorso'' vanno letti rispetto a quel momento.

Il percorso del capitolo segue quello della lezione:
\begin{enumerate}
  \item l'avvio, i \textbf{docenti} e le \textbf{presentazioni} dei partecipanti
        (sezioni da~\ref{a1:sec:avvio} a~\ref{a1:sec:partecipanti});
  \item i \textbf{materiali} e gli strumenti del corso: il form, la cartella
        condivisa, Discord (sezione~\ref{a1:sec:materiali});
  \item lo \textbf{scopo} e il \textbf{metodo} del corso, le regole di
        comportamento e il carico di lavoro (sezioni~\ref{a1:sec:scopo}
        e~\ref{a1:sec:regole});
  \item il \textbf{project work} e il \textbf{calendario} (sezione~\ref{a1:sec:project});
  \item la regola sull'\textbf{intelligenza artificiale} e la lunga discussione
        che ne segue (sezione~\ref{a1:sec:ia});
  \item le \textbf{figure professionali} dell'industria del gioco analogico,
        con un approfondimento su marketing e sales
        (sezioni~\ref{a1:sec:figure} e~\ref{a1:sec:marketing});
  \item la differenza tra \textbf{game design} e \textbf{game development}
        (sezione~\ref{a1:sec:design-development});
  \item il \textbf{mercato}: Italia ed estero (sezione~\ref{a1:sec:mercato});
  \item i \textbf{fondamenti di design}: perché serve una terminologia comune,
        \emph{game} e \emph{play}, il framework \textbf{MDA} (e il DPE), le
        \textbf{parole vietate}, chi crea, chi gioca e chi facilita, il
        \textbf{playtest} e l'\textbf{iterazione} (sezioni
        da~\ref{a1:sec:terminologia} a~\ref{a1:sec:playtest});
  \item \textbf{Tabletop Simulator} e il primo task (sezione~\ref{a1:sec:tts});
  \item le \textbf{domande finali}: il tema del project work, il bilanciamento,
        i ruoli nei team, la pubblicazione, il gioco di ruolo in Italia, le
        \emph{house rule} di Monopoly, i giochi ``comunicati male''
        (sezioni~\ref{a1:sec:domande-daniele} e~\ref{a1:sec:domande}).
\end{enumerate}

% ---------------------------------------------------------------------
\section{L'avvio della lezione}
\label{a1:sec:avvio}

La lezione comincia con qualche minuto di attesa. ``Siamo in nove in classe'', e
si aspetta che arrivino tutti; uno dei docenti scherza sul fatto che l'unica
persona che arriverà ``sempre in ritardo di venerdì'' è Gianmarco. Mancano
Giordano e Tiziana; Tiziana arriva, mentre la segreteria fa sapere poco dopo che
Giordano non potrà esserci e vedrà la registrazione. Si aspetta anche Beatrice,
la terza docente, che arriva a sua volta dopo pochi minuti.

Nell'attesa, Daniele anticipa il programma della giornata e qualche
informazione pratica:
\begin{itemize}
  \item oggi è una giornata di \textbf{introduzione}: come funziona il corso, chi
        sono i docenti, le presentazioni degli studenti, le \textbf{metodologie}
        e il \textbf{calendario}. Gli studenti avranno ``un fracco di domande'', e
        alla fine della lezione c'è ``proprio uno slot apposta'' per le domande;
  \item le lezioni sono \textbf{registrate}: di solito la registrazione viene
        consegnata dopo circa \textbf{12--24 ore} sul profilo personale della
        piattaforma dell'accademia, che va attivato con la segreteria (Tiziana,
        già studentessa, probabilmente ce l'ha già);
  \item i docenti ``cuore'' del corso sono tre (Daniele, Ruggero e Beatrice),
        ``soprattutto Ruggero''; a loro si aggiungeranno tre docenti ospiti, per
        meno ore, probabilmente una lezione a testa, ``tranne Marta, che avrà due
        lezioni quest'anno''.
\end{itemize}
E una battuta sul caldo: ``spero che voi siate carichi, anche se le temperature
sono ancora indecenti''. ``Carichi moderatamente, per quello che potete''.

Poi Daniele condivide lo schermo (``non solo le slide, ma anche un po' di altre
cose'') e apre ufficialmente il corso: ``benvenuti e benvenute nella seconda
edizione del corso di Analog Game Design. Questa è la prima lezione''. Il
percorso è di circa \textbf{tre mesi}, fino a dicembre, e durante il corso i
docenti cercheranno di far provare agli studenti ``cosa significa fare Analog
Game Design'' con un \textbf{project work}, un progetto di gruppo di cui si
spiegheranno più avanti la struttura e le richieste.

% ---------------------------------------------------------------------
\section{I docenti}
\label{a1:sec:docenti}

``Anzitutto partiamo da delle presentazioni doverose''. Ognuno dei tre docenti
principali si presenta con una slide che riporta i contatti e una selezione dei
lavori fatti.

\subsection{Daniele: narrativa, fiction e gioco di ruolo}

\textbf{Daniele}\footnote{La trascrizione riporta ``Daniel Cusetto''; nel
resto della lezione tutti lo chiamano ``Daniele''. Il cognome non è
ricostruibile con certezza.} ha già conosciuto quasi tutti gli studenti, perché
ha fatto loro il colloquio di ammissione. È \textbf{narrative designer}, o meglio
\emph{game narrative designer}; insegna nell'accademia che ospita il
corso\footnote{La trascrizione storpia il nome dell'accademia in molti modi
(``Dive'', ``IVE'', ``Accademia Live'', ``Ai'', ``AIVA''). Con ogni probabilità
si tratta dell'\textbf{AIV}, l'Accademia Italiana Videogiochi, che ha sede a
Roma: lo fanno pensare i corsi triennali di game design e narrative design citati
dai docenti e il fatto che una studentessa racconti di essersi trasferita a Roma.
È una ricostruzione, non un'informazione esplicita; nel seguito il capitolo
parla semplicemente dell'``accademia''.} ma lavora soprattutto come \textbf{brand
advisor}. Nel corso si occuperà della parte di \textbf{gioco di ruolo},
\textbf{narrazione} e \textbf{fiction}.

I suoi contatti, ``in quest'ordine'', sono email, Telegram e Discord. Ci tiene a
dare alcune indicazioni precise:
\begin{itemize}
  \item se si vuole avere la certezza di una risposta, bisogna scrivergli su
        \textbf{Telegram}, perché lì ha tutti i suoi clienti;
  \item Discord lo guarda di meno; l'\textbf{email} è per le cose ufficiali;
  \item non bisogna contattarlo altrove: non su Facebook, non su LinkedIn, non su
        Instagram, ``perché li vedo molto di meno''.
\end{itemize}
La slide mostra anche alcuni giochi, attività e brand su cui ha lavorato negli
anni, ma non si dilunga: ``nel caso ne parleremo''.

\subsection{Ruggero: game design e ricerca e sviluppo}

\textbf{Ruggero} (spesso chiamato ``Ruger'') ha fatto il \textbf{game designer}
per varie aziende e oggi è game designer e si occupa di \textbf{ricerca e
sviluppo} per un'azienda del settore\footnote{La trascrizione dice ``per MS'';
il nome dell'azienda non è ricostruibile con certezza. Più avanti Ruggero cita
tra i giochi su cui ha lavorato il gioco da tavolo ufficiale di
\emph{Cyberpunk}, \emph{Marvel United} e \emph{Status}, un gioco sulla
geopolitica.}. Ha un background ``misto'': nel suo portfolio ci sono giochi da
tavolo, giochi di ruolo, giochi analogici, LARP, videogiochi, giochi in realtà
virtuale. Si occupa anche di \textbf{divulgazione} sul gioco e di analisi di
\textbf{game studies} e game design. Più avanti, nelle domande finali, dirà di
avere ``50 e più giochi prodotti e progettati'' e di insegnare da vari anni.

Lo si può contattare liberamente via Gmail e Telegram, e cerca di rispondere
\textbf{entro 24 ore}, salvo imprevisti o eventi: ``sentitevi in libertà di
scrivermi se avete dubbi, senza problemi''. Quando arriva a parlare dei propri
giochi, la slide mostra ``un mix di varie cose, dalle più famose alle meno
famose'', ma Ruggero si imbarazza a presentarli, tra le risate dei colleghi.

\subsection{Beatrice: i giochi pervasivi}

\textbf{Beatrice} (``Bea'') si occupa di \textbf{giochi pervasivi}
(\emph{pervasive games}). ``Se non li avete mai sentiti nominare, non
preoccupatevi'': questo corso è ``uno dei pochi, forse l'unico in Italia'' che
li racconta esplicitamente. Li si vedrà nelle prossime lezioni, ma la definizione
di base è questa: i giochi pervasivi \textbf{utilizzano gli spazi e gli oggetti
della vita quotidiana e li trasformano in spazi e oggetti di gioco}.

Beatrice lavora come \textbf{freelance} e si occupa di diversi tipi di
progettazione: lavora sia con le \textbf{aziende} sia con gli \textbf{eventi dal
vivo}. La sua attività diretta si trova sotto l'etichetta \textbf{Minime
Differenze}\footnote{La trascrizione riporta il nome dell'etichetta come
``Minime Differenze'' e l'indirizzo del sito come ``mininedifferenze.com''; la
grafia esatta del sito non è verificabile dalla trascrizione.}, mentre i lavori
per le aziende sono ``sparpagliati'' sui siti delle varie aziende. Molte delle
cose mostrate nella sua slide sono attività proposte al pubblico ``con un
biglietto, con un gioco in formato evento'': un tema che si toccherà nel corso,
per capire come \textbf{comunicare} i propri giochi in modo che siano
comprensibili, accessibili per il pubblico e \textbf{vendibili}.

Ammette un ``grosso limite'' con i messaggi: è un po' lenta a rispondere, ``ma
rispondo sempre a tutto''. La si può contattare per email (cognome punto nome) o
su Telegram. È ``molto curiosa di vedere cosa produrrete''.

\subsection{I docenti ospiti}

Daniele presenta poi i tre docenti che il corso avrà ``solo per alcune
lezioni'', tutte verso \textbf{novembre}:
\begin{itemize}
  \item \textbf{Diego Ricchiuti}, direttore del corso di Game Design
        dell'accademia, farà una lezione sull'\textbf{RNG} (\emph{random number
        generation}, la casualità nei giochi): ``vi parleremo di RNG quando sarà
        necessario. Comunque avrete una lezione con lui, sappiatelo'';
  \item \textbf{Marta}\footnote{La trascrizione riporta il cognome come
        ``Cercasassi''; la grafia non è verificabile.} farà \textbf{due lezioni
        sul development}: quest'anno ``abbiamo deciso di fare doppietta'',
        perché l'anno scorso il suo intervento è stato ``molto, molto utile'';
  \item \textbf{Flavio}\footnote{La trascrizione riporta il cognome come
        ``Mortarino'' e, più avanti, ``Mortalino'', e il nome della sua casa
        editrice come ``Grampibert''; nessuno dei due è ricostruibile con
        certezza.} farà la lezione di \textbf{game production}; è un editore, e
        di solito partecipa anche alla tavola rotonda con gli editori della
        giornata finale.
\end{itemize}
Le slide sono già online e gli studenti avranno accesso alla cartella condivisa:
``non vi preoccupate di queste cose, ne parleremo a breve''. Ma prima di parlare
del corso, ``vogliamo conoscere voi''.

% ---------------------------------------------------------------------
\section{Le presentazioni dei partecipanti}
\label{a1:sec:partecipanti}

Le domande per la presentazione sono semplici, da affrontare ``a turno, senza
fretta, con calma'':
\begin{itemize}
  \item nome e occupazione (``se volete dirla'');
  \item perché si vogliono fare giochi analogici, e a che livello: che cosa ha
        portato lo studente al corso;
  \item che tipo di esperienze o conoscenze si hanno già nell'ambito, se si è già
        fatto qualcosa oppure no.
\end{itemize}
``Chi vuole partire senza nessun tipo di ritegno e di vergogne?''. Per far
parlare gli studenti guardandoli, Daniele interrompe la condivisione dello
schermo e ricorda lui le domande quando serve. Le presentazioni sono interessanti
anche per i commenti dei docenti, che anticipano diversi principi del corso.

\subsection{Tiziana}

Tiziana fa l'\textbf{infermiera}, ma nella vita ha fatto molti lavori:
impiegata, insegnante e ``tanti altri lavoretti''. Gioca ``da sempre'': il gioco
fa parte anche della sua vita familiare, perché suo fratello fa il programmatore
di videogiochi. Ha già \textbf{organizzato} giochi: una caccia al tesoro per far
conoscere il territorio, due scene con un delitto (giochi con omicidio da
risolvere), giochi per intrattenere le persone. Ha riscoperto di recente i giochi
da tavolo, arrivando a \emph{Carcassonne} con famiglia e amici; poiché è
difficile trovare persone con voglia di giocare nella propria cerchia ristretta,
ha scoperto il mondo delle \textbf{associazioni} ludiche, ``e lì è stato
un'esplosione'': prima in Veneto (è di Vicenza), poi a Roma, dove si è
trasferita.

Ha in mente un progetto ambizioso: far conoscere una certa realtà con la modalità
del \textbf{gioco di ruolo}, o comunque proporre un gioco suo.

\paragraph{Il commento: gli strumenti, non il gioco dei sogni.} Ruggero coglie
l'occasione per un'avvertenza che vale per tutti, visto che molti studenti
arrivano con idee già pronte. Nel corso si lavorerà a un progetto, ma
l'obiettivo dei docenti è dare \textbf{strumenti} da applicare anche
\emph{fuori} dal progetto, ai giochi che gli studenti vorranno fare un domani:
``che non siano esercizi per il corso, ma per voi in primis''. Quindi: ``non
rimaneteci male se non riuscirete a lavorare al vostro gioco dei sogni in questo
corso''. Anche perché ``sarebbe impossibile, in due mesi, fare il gioco dei
sogni''. Il corso darà invece la possibilità di fare un'esperienza ludica
insieme e di imparare a replicarla sulle proprie idee: ``le idee non buttatele,
anzi, il contrario: tenetele''.

\subsection{Nicola}

Nicola è uno \textbf{studente}. Sei mesi prima ha fatto un colloquio con Daniele
perché stava valutando di passare all'accademia, e ha deciso, ``nel dubbio'', di
cominciare da questo corso breve per capire come funziona la formazione;
eventualmente si iscriverà al corso di game design vero e proprio. Gli interessa
in particolare il lato dell'esperienza del giocatore, su cui è ``molto critico''
(qui usa la parola ``divertimento'', e Daniele lo ferma con un ``ne parleremo'':
è una delle parole che il corso vieta, come si vedrà nella
sezione~\ref{a1:sec:parole-vietate}).

Ha 21 anni ed è \textbf{giocatore di ruolo da circa otto}. Prima ancora di
conoscere i giochi di ruolo si era creato un sistema tutto suo, ``pessimo,
disgustoso, terribile'', un sistema a percentuale con il d100 ``terrificante'';
poi ha scoperto i giochi di ruolo da tavolo, i classici (\emph{Dungeons \&
Dragons}, \emph{Daggerheart}, \emph{Vampire}\footnote{La trascrizione riporta
``DD'', ``Daggerart'' e ``Vampire'': quasi certamente \emph{Dungeons \&
Dragons}, \emph{Daggerheart} e \emph{Vampire: The Masquerade}.}). Gioca anche a
\textbf{card game} e ad altri giochi sociali, pur avendo sempre avuto difficoltà
a trovare gruppi, ``avendo vissuto in un buco dimenticato dal mondo''. Come
giocatore è ``l'eterno narratore'' del gruppo, quello che prepara la storia e si
costruisce le regole. Come autore ha fatto soprattutto giochi su carta, ``perché
sono incapace a programmare, per adesso''.

Daniele osserva che il tema del giocatore e del suo rapporto con il gioco è ``un
argomento importante'', anzi ``un vaso di Pandora enorme'' (espressione che
tornerà molte volte nella lezione), e che il corso tratterà anche che cos'è
l'analog game design in generale e quanti tipi di prodotti esistono.

\subsection{Rita}

Rita\footnote{Daniele la chiama prima ``Nino'' (probabilmente un errore della
trascrizione) e poi la ringrazia come Rita.} ha 30 anni e si collega da una
regione del Sud, dove ``ancora si muore di caldo''. Ha studiato
all'\textbf{Accademia di Belle Arti di Napoli}, \textbf{design della
comunicazione visiva}: si occupa dell'aspetto più imprenditoriale della
progettazione artistica. Non volendo progettare per un'impresa (``fare loghi,
robe del genere''), vorrebbe applicare la propria creatività alla realizzazione
di giochi analogici.

Non è una giocatrice di ruolo e non lo è mai stata: è una giocatrice
\textbf{tradizionale}, cresciuta con Monopoly, la dama e giochi tradizionali
cinesi. Ha già \textbf{idee, prototipi e progetti} realizzati, e vuole capire
come renderli \textbf{prodotti per la vendita} e se questo è il campo adatto a
lei; magari anche confrontarsi con i colleghi. La sua progettazione è orientata
al \textbf{sociale} e all'\textbf{educativo}: ``responsabilizzare un po' le nuove
generazioni su come sta andando il mondo e su come potrebbe andare, visto che è
tutto nelle nostre mani''.

\paragraph{Il commento: siate curiosi.} ``Pane per i tuoi denti'', le risponde
Daniele, soprattutto per i LARP e i giochi pervasivi: ``vi apriremo un po' la
testa da questo punto di vista''. E aggiunge un consiglio per tutti: ``siate
curiosi'', perché imparare da altri ambiti dà moltissimo. Ognuno è lì per fare
qualcosa di preciso, ``però siate aperti anche ad altre esperienze''.

\subsection{Tancredi}

Tancredi ha 21 anni e non ha ancora grande esperienza lavorativa: qualche
``piccolo progettino'' che non è andato come voleva. Gli interessa il mondo dei
giochi, digitali e non, con un approccio \textbf{educativo} e \textbf{sociale}.
È arrivato soprattutto dai videogiochi, a cui gioca da quando era piccolo; ai
giochi da tavolo si è avvicinato solo nell'ultimo anno, grazie a una realtà
ludica, perché in famiglia ``non siamo molto da giochi''. Ha realizzato qualche
piccolo videogioco con vari engine: prima quelli senza un vero linguaggio di
programmazione, ``più a nodi'', poi qualcosa di più complicato, e ora sta
lavorando a un piccolo prototipo.

\paragraph{Il commento: niente di quello che vedrete è sprecato.} ``Non butterai
via niente di analog game design'', gli risponde Daniele, ``perché il 99 per cento
delle cose che vedremo si ritrovano quasi uno a uno nel mondo dei videogiochi''.
Tanto che nei corsi triennali dell'accademia (anche in quello di narrative
design, dove insegna Ruggero) il primo anno si parte dai \textbf{giochi da tavolo
e di ruolo}, proprio per spiegare meccaniche più complesse che servono poi per i
videogiochi. Ruggero aggiunge che vale anche il contrario: alcuni dei
\textbf{framework} di game design che si vedranno nelle lezioni tecniche vengono
dal videogioco e sono stati riapplicati all'analogico. ``Evitare la
crossmedialità qui è impossibile''. Anzi, tutte le conoscenze ``mediatiche
adiacenti'' (musica, letteratura, cinema) saranno utili.

\subsection{Lorenzo Vergata}

I Lorenzo in classe sono due; il primo a presentarsi è \textbf{Lorenzo
Vergata}, 26 anni, \textbf{dottorando al Politecnico di Milano}. Gioca ai
videogiochi da quando aveva due anni e mezzo, ma ha giocato al suo primo gioco da
tavolo ``credo nove mesi fa''. Gli fa piacere sentire che molti in classe sono
interessati ai giochi \textbf{sociali} ed \textbf{educativi}, un settore che
interessa anche a lui e che ``può essere esplorato molto'', anche con qualche
aiuto dell'intelligenza artificiale per esperienze più ricercate.

Di esperienza nell'analogico dice di avere fatto ``il minimo sindacale''. Ma
racconta che diciotto anni prima, da bambino, era ``più game designer di
adesso'': prendeva un gioco e ne cambiava le regole, anche al biliardino (``no,
ora introduciamo questa regola''). Gli piace sperimentare, creare regole e
vedere che cosa accade in sé e negli altri. Soprattutto vorrebbe ``una chiave di
lettura per poter rimanipolare la realtà'': tre settimane prima è stato a fare la
\textbf{vendemmia in Francia} e ha visto tantissime dinamiche che potrebbero
diventare un gioco da tavolo o un videogioco. Per esempio: finire contro un rovo
pieno di semi che si attaccano ai vestiti lo rallentava nella raccolta
dell'uva, mentre in un altro filare era più rapido. ``Si può fare pure una
dinamica'', e vorrebbe gli strumenti per farne giochi.

Arriva subito la battuta: una ``vendemmia gamificata'' da proporre in Oltrepò
Pavese (``io sono dell'Oltrepò Pavese, quindi devo per forza bere vino'')\footnote{La
trascrizione dice ``nell'altro papavese'' e non chiarisce chi pronunci la
battuta.}. Daniele conferma che quest'anno il corso avrà ``tante cose sul
sociale'': ``ci piace parecchio''.

\subsection{L'altro Lorenzo}

Il secondo Lorenzo (``secondo, ma non meno importante'') è uno \textbf{scrittore
di racconti brevi} e sta scrivendo il suo primo romanzo. Gli è sempre piaciuto
raccontare storie, e da più grande scriverle; gioca ai videogiochi ed è ``un
amante delle storie e della cultura pop''. Negli ultimi anni è entrato nel mondo
del \textbf{gioco di ruolo da tavolo} e vorrebbe unire le due passioni, le
storie e il gioco di ruolo, magari facendone un lavoro; gli interessa anche lo
\emph{strategic writing}. Daniele gli ricorda che ci sarà una lezione proprio
sulla \textbf{fiction}.

\subsection{Angelo}

Angelo si collega senza videocamera dal lavoro: fa il \textbf{birraio} (in un
piccolo birrificio in Brianza, a Bernareggio, come racconterà dopo la pausa) e,
con il caldo, ``vi risparmio l'orrore''. Per molto tempo era convinto che i
giochi da tavolo fossero solo i classici (Monopoly, Risiko) o i party game; poi,
quattro o cinque anni fa, ha ``sbattuto la faccia'' contro giochi davvero
appassionanti, con una specie di narrazione dietro che li rendeva ancora più
coinvolgenti, e non ha più smesso. È uno che ``ha sempre un botto di idee, ma poi
non mette in pratica assolutamente nulla'': cercava un corso che gli desse le
basi e soprattutto lo \textbf{motivasse} a provare a fare qualcosa.

\paragraph{Il commento: bisogna buttarsi.} Daniele racconta che anche l'anno
prima un suo caro amico si era iscritto proprio per questo motivo. Il corso è
breve, ma cerca di portare gli studenti ``da 0 a 100'' e di far realizzare un
\textbf{prodotto} da presentare alla fine. Il punto è che i giochi ``bisogna
farli, bisogna buttarsi''. E, venendo dalla narrazione, Daniele lo paragona allo
scrivere: si scrive e si riscrive, e allo stesso modo si fa \textbf{playtest}, si
rifà, si aggiorna. ``Non è che si fanno così''.

\subsection{Gianmarco (dopo le slide)}

Gianmarco arriva a lezione iniziata e si presenta più tardi, alla fine della
prima parte, dopo che i docenti hanno chiuso le slide organizzative e la
discussione sull'intelligenza artificiale (``Gianmarco, ti giuro, ti facciamo
presentare''). Lavora nel \textbf{team di innovazione} di un'azienda italiana.
Vuole fare giochi perché è una sua passione, che lo gratifica, e in passato ci ha
provato ``in maniera scomposta''. Cerca un \textbf{metodo} e un confronto con
persone più esperte: ``studiare solo dei manuali è molto sterile'', ``stare
chiusi nel pertugio non aiuta'', ed ``essere autodidatta ha estremamente i suoi
limiti''. Come esperienza è un giocatore e ha letto ``un po' di materiale, di
blog, di manuali'', ``niente di rilevante''.

Manca solo Giordano, che si presenterà alla lezione successiva.

\paragraph{Uno sguardo d'insieme.} La classe è molto varia: un'infermiera,
una designer della comunicazione visiva, un birraio, un dottorando, uno
scrittore, un impiegato nell'innovazione, due studenti ventunenni. Alcuni
vengono dal gioco di ruolo, altri dai giochi tradizionali, altri dai
videogiochi; quasi tutti hanno scoperto i giochi da tavolo moderni da pochi anni
o da pochi mesi. E almeno tre studenti dichiarano un interesse per i giochi
\textbf{sociali} ed \textbf{educativi}. Come si vedrà nella
sezione~\ref{a1:sec:scopo}, i docenti considerano questa varietà una risorsa.

% ---------------------------------------------------------------------
\section{I materiali e gli strumenti del corso}
\label{a1:sec:materiali}

Prima delle slide sul corso, Daniele sbriga un po' di ``clean up casalingo''.

\paragraph{Il form iniziale.} C'è un \textbf{form} da compilare ``entro e non
oltre il 15 settembre'': il QR code è nelle slide, e la segreteria lo ricorderà
su Discord. Richiede circa 10--15 minuti; si può fare subito o in giornata, ``e
ve lo togliete di torno''. È un form standardizzato per tutti i corsi (``non
siamo io e Ruggero a farvi passare la tortura del test''), diviso in due parti:
\begin{itemize}
  \item una parte di \textbf{conoscenza del candidato}: quanto tempo si ha a
        disposizione per il corso, come ci si organizza con il lavoro, che cosa
        piace giocare. È la parte che interessa ai docenti, perché li aiuta a
        organizzare le lezioni;
  \item una parte di \textbf{test logico}, che i docenti non hanno scritto (``io
        non sono responsabile del form'', ``no, no, nemmeno io'').
\end{itemize}

\paragraph{La cartella condivisa.} Tutti gli studenti vengono aggiunti alla
cartella del corso su \textbf{Google Drive}, che si trova tra i file
``condivisi con me''. Gli studenti ne sono \textbf{gestori dei contenuti} e
possono caricarci quello che vogliono. Dentro ci sono:
\begin{itemize}
  \item le \textbf{attività in classe};
  \item le \textbf{dispense};
  \item i \textbf{materiali didattici}, con le slide (quelle della lezione sono
        già lì);
  \item le cartelle dei \textbf{project work} futuri;
  \item i \textbf{task}, che ``saranno pochi, ma ci saranno'';
  \item il \textbf{calendario studenti}, ``importantissimo'', con due schede: il
        calendario delle lezioni e quello del project work.
\end{itemize}

\paragraph{Discord.} Il corso ha anche un \textbf{server Discord}, a cui tutti
sono stati invitati; ci sono i docenti (Daniele, Ruggero, Diego e altri), che si
possono taggare. Di solito è lì che la segreteria pubblica le informazioni, per
esempio se una lezione ``dovesse per caso essere cancellata''. Nelle domande
finali Ruggero aggiungerà che c'è un canale apposito per il corso, che gli
studenti possono usare liberamente per organizzarsi tra loro.

\paragraph{Le reaction di Zoom.} Durante le lezioni su Zoom si possono usare le
\textbf{reaction} (il pollice in alto, il cuoricino, la mano alzata): un modo
rapido per far capire ai docenti se si è capito, se si vuole intervenire. Nel
frattempo arrivano i primi pollici.

% ---------------------------------------------------------------------
\section{Lo scopo e il metodo del corso}
\label{a1:sec:scopo}

La prima serie di slide (che Daniele, per un lapsus, presenta come ``cos'è il
corso di pulire il pavimento di casa'') spiega che cosa il corso vuole fare e
come.

\subsection{Diventare autori analogici}

Lo \textbf{scopo del corso} è fornire le competenze per \textbf{diventare autori
analogici} e stimolare la creatività e l'\textbf{autorialità} di ciascuno: ``voi
imparerete a diventare autori di giochi analogici''. La differenza tra
\emph{game designer} e \emph{game developer} verrà chiarita più avanti (sezione~\ref{a1:sec:design-development}),
ma la parola chiave qui è \emph{autorialità}: ``vogliamo voi, vogliamo le vostre
idee''. Magari non si potrà usare l'idea del gioco della propria vita, ma si
potranno portare nelle lezioni le proprie idee e le proprie sensibilità.

Gli obiettivi elencati nella slide sono questi:
\begin{itemize}
  \item dare \textbf{maggiore consapevolezza nel design e nello sviluppo}. In
        Italia, purtroppo, ``si fa design a volte in maniera non consapevole''.
        Consapevole vuol dire avere almeno un'introduzione alle basi del design e
        dello sviluppo, e fare le cose ``con metodo, con raziocinio, con
        consapevolezza'';
  \item fornire gli \textbf{strumenti essenziali} per un design consapevole;
  \item stimolare la \textbf{sperimentazione} e la \textbf{trasversalità
        mediatica}. Non si chiede di fare ``il \emph{Root} della situazione o il
        \emph{Carcassonne} della situazione''\footnote{La trascrizione dice ``il
        lumaven della situazione o il carcassone della situazione''; il primo
        titolo non è ricostruibile con certezza. Il senso è chiaro: non si
        chiede di rifare un grande classico del gioco da tavolo.}, ma un
        project work \textbf{calibrato sulla durata del corso}, e però
        \textbf{sperimentale} e ``assolutamente transmediale'';
  \item portare \textbf{numerosi esempi} di tipologie di gioco, dall'analogico al
        tavolo, al pervasivo ``e oltre''. Lo si vedrà nella lezione sui
        \textbf{casi studio} del venerdì successivo: ``voi non avete idea di
        quante cose esistono, bellissime, sconosciutissime e favolose''. Daniele
        confessa che preparando e aggiornando le slide i colleghi gli hanno
        mostrato cose che lui stesso non sapeva nemmeno che esistessero;
  \item dare gli strumenti per \textbf{analizzare} i giochi altrui. Ruggero lo
        sottolinea: se il corso riesce nel suo intento, gli studenti sapranno
        non solo fare i propri giochi ma anche analizzare giochi già fatti, che è
        ``una delle skill importanti anche nella parte di ricerca e sviluppo''.
        Si vedranno casi studio e giochi pubblicati, sempre con l'obiettivo di
        capire \textbf{che cosa fanno e perché lo fanno}, per poi fare meglio il
        proprio lavoro;
  \item imparare a \textbf{fare giochi insieme}. Esistono gli autori che lavorano
        da soli (i \emph{solo developer}), ``ma sappiate che anche l'autore che
        lavora da solo poi si deve interfacciare con il pubblico, con un editore,
        magari con un editor, con un publisher''. ``Non lavorerete mai da soli'':
        bisogna imparare a lavorare con gli altri.
\end{itemize}

\subsection{Un corso intersezionale}

Il corso è \textbf{intersezionale}: \textbf{non ci sono lezioni dedicate a una
singola tipologia di gioco}. Non ci sarà ``la lezione sul gioco da tavolo'',
``la lezione sul gioco di ruolo'', ``la lezione sul gioco pervasivo''. Ci saranno
lezioni su temi diversi, ciascuna pensata per essere utile a \textbf{tutte e tre
le sfaccettature} del corso; alcune saranno più importanti per il gioco da
tavolo, altre per il pervasivo, altre per il gioco di ruolo, ma nessuna sarà
esclusiva. I docenti vogliono ``contaminazioni e applicazioni assolutamente
trasversali''.

\subsection{Un corso cross-generazionale}

È anche un corso \textbf{cross-generazionale}: i gradi di esperienza di gioco
quest'anno sono abbastanza allineati, ma le esperienze di vita, le età, le
estrazioni sociali sono varie. ``Ed è un pro, non è un contro. Più diversificata
è la classe, più imparate meglio anche voi''. La slide ha una piccola nota per
fasce d'età:
\begin{description}
  \item[Under 30] Concentratevi sulla vostra \textbf{crescita}. Evitate di
        confrontarvi con le persone più grandi che magari ``ce la fanno meglio''
        in alcune cose.
  \item[Over 30] Concentratevi sul \textbf{consolidamento}, ma siate
        \textbf{umili}: i docenti diranno cose che faranno ``storcere il naso''.
        ``Bisogna un po' reimparare per imparare'': bisogna sapersi adattare.
\end{description}
Ruggero contestualizza: lo dicono perché in passato hanno visto studenti che,
essendo più adulti o avendo giocato di più, pensavano che l'esperienza
\emph{di gioco} si trasformasse automaticamente in competenza \emph{di
creazione} di giochi. Non è così. E non bisogna sottovalutare chi gioca meno:
proprio perché gioca cose diverse, per risolvere un problema di design specifico
``è una ricchezza incredibile avere qualcun altro che riesce a uscire dalla
scatola''.

\subsection{Un corso pratico e adattabile}

Il corso è \textbf{pratico}: gli studenti si concentreranno su un progetto di
gruppo, il \textbf{project work}, lavorando in team. Ed è \textbf{adattabile}: il
programma è definito, ma ``si aggiorna sui vostri stimoli, priorità e
necessità''. Se durante il corso i docenti vedono che serve più una cosa
dell'altra, sono ``prontissimi'' a dare più spazio a un argomento invece che a
un altro.

\subsection{Inquiry-based learning: domandate}

Il metodo è quello dell'\textbf{inquiry-based learning}, l'apprendimento basato
sulle domande. Di solito le persone dicono ``ho una domanda, magari sarà
stupida''. ``Non esistono domande stupide. Pensate che ci sia una domanda
stupida? Fatela, alzate la mano''. Vanno bene anche le
\textbf{provocazioni}: se qualcosa non torna, è meglio chiedere. La condivisione
è fondamentale.

I docenti danno \textbf{metodi, dati e conoscenze}, ma lasciano trovare le
risposte agli studenti: ``vi lasciamo un pochettino annegare nel mare prima di
darvi la soluzione, in modo tale che la possiate trovare da soli''. Di
conseguenza \textbf{non ci saranno interrogazioni sulle definizioni}: interessa
vedere \textbf{come} gli studenti applicano ciò che imparano. Se una
definizione non viene in mente, non è un problema; l'importante è saperla usare.

\subsection{Feedback, non valutazione}

``Sì al feedback, no alla valutazione''. I docenti danno \textbf{feedback
verbali}, anche sulle attività da fare a casa (i task, che non saranno tanti).
Il corso è calibrato per \textbf{studenti lavoratori}: c'è un po' da lavorare
fuori dalle lezioni, ma in modo calibrato.

Il principio è \textbf{fallire oggi per imparare per il futuro}: \emph{fail
forward}, \emph{fail to learn}. I docenti lasciano provare e fallire almeno una
volta prima di dare la soluzione: ``il classico metodo Montessori, quello
originale, non quello travisato''.

\subsection{Le soft skill e i criteri di valutazione}

Grande attenzione va alle \textbf{soft skill}: i lavori di gruppo, lo scambio
propositivo, l'atteggiamento positivo, la propositività, la proattività.
``Lavoro di gruppo'' non significa stare ogni giorno in sala studio con gli
altri: va benissimo che un gruppo si organizzi dandosi priorità e tempi in base
alla disponibilità di ciascuno. L'importante è \textbf{collaborare}, essere
\textbf{professionali}, rispettarsi a vicenda, rispettare le idee proprie e
altrui e lasciare spazio agli altri, perché i progetti ``non saranno solo
vostri, saranno di tutte le persone del vostro gruppo''.

La valutazione del project work non riguarda solo la \textbf{documentazione}
prodotta, ma anche:
\begin{itemize}
  \item la \textbf{presentazione} finale;
  \item l'\textbf{autovalutazione};
  \item la \textbf{crescita personale};
  \item l'\textbf{atteggiamento} tenuto nel corso.
\end{itemize}
Daniele cita: ``da ognuno secondo le sue capacità, a ognuno secondo i suoi
bisogni''. Lavorare in gruppo non vuol dire che, in tre, ciascuno lavori
esattamente il 33\% del tempo. Ai docenti serve piuttosto capire \textbf{chi fa
che cosa} nel gruppo, e che il processo interno sia \textbf{trasparente}.

\paragraph{Lavorare in cloud.} Per lo stesso motivo il consiglio è di
\textbf{non lavorare in locale}, ma direttamente nelle cartelle condivise del
project work e dei task. Si è liberi di fare altrimenti, ma lavorando in cloud
si dà ai docenti la possibilità di vedere \emph{come} si lavora, il proprio
``processo mentale''.

% ---------------------------------------------------------------------
\section{Le regole di comportamento e il carico di lavoro}
\label{a1:sec:regole}

Alcune slide sono di quelle che si ripetono in ogni corso (``\emph{repetita
iuvant}'', storpiato per scherzo in ``ripetita Juventus'').

\paragraph{Comportamento.} Prestare attenzione durante le lezioni; interagire;
portare rispetto; essere proattivi; alzare la mano e fare domande.
\textbf{Tolleranza zero} per pregiudizi, maleducazione, mancanza di rispetto,
\emph{cheat} e scorciatoie. Daniele ammette che ``cheat e scorciatoie'' viene da
altri corsi dove il tema è più importante, ma Ruggero lo ferma: ``non
sottovalutate questa cosa'', perché implica anche \textbf{non far fare a ChatGPT
i task}. Daniele ha una slide apposta, ``perché è un discorso molto complesso nel
2026'' (sezione~\ref{a1:sec:ia}).

\paragraph{Le lezioni sono intense.} Le lezioni possono essere \textbf{dense},
``soprattutto le mie, sappiatelo: io sono logorroico'', avverte Daniele. Bisogna
chiedere spiegazioni se serve, prendere appunti, rivedere le registrazioni. Le
lezioni sono ``delle scintille'': modi per far esplorare i temi oltre le slide.
``Andate oltre, se riuscite''.

\paragraph{Il carico di lavoro.} Le regole sui task sono queste:
\begin{itemize}
  \item i task \textbf{sono per voi, non per noi};
  \item se si ha poco tempo, è meglio farli rapidamente piuttosto che non farli:
        ``provateci'';
  \item i docenti cercano di non chiedere più di \textbf{un'ora o due a
        settimana} oltre alle ore di lezione;
  \item non tutte le lezioni avranno un task, e il \textbf{project work} è il più
        importante di tutti.
\end{itemize}

\paragraph{I feedback dei docenti.} I feedback ``si guadagnano'' man mano che si
reagisce bene: non bisogna aspettarsi ``dei papiri'' sui task. Se i docenti
hanno voglia di dare feedback più ampi lo fanno, soprattutto se vedono che lo
studente se li merita. Ad Analog di solito non ci sono task \emph{rejected}; ma
se un docente rifiuta un task, vuol dire che si è fatto \textbf{meno del livello
minimo} richiesto dal corso. I feedback possono non essere immediati: ``siate un
pochettino impietositi dal fatto che siamo dei docenti''. E se un feedback non
arriva, o se ne vogliono di più, bisogna scrivere, ``perché siamo tre esseri
umani e quindi possiamo anche dimenticarci''.

% ---------------------------------------------------------------------
\section{Il project work e il calendario}
\label{a1:sec:project}

\subsection{Come funziona il project work}

Una prima slide con le date del project work è da aggiornare (``perdonatemi il
refuso''); Daniele la corregge al volo e rimanda alla slide successiva, più
precisa. Il funzionamento è questo:
\begin{itemize}
  \item i \textbf{gruppi li formano i docenti}. Gli studenti sono nove, quindi si
        faranno \textbf{tre gruppi da tre}: una dimensione pensata anche per il
        fatto che molti sono studenti lavoratori, così da non ritrovarsi ``in due,
        a esplodere'';
  \item ogni gruppo deve prima preparare un \textbf{pitch} del proprio gioco, da
        presentare in anticipo: un'idea, una bozza del gioco;
  \item poi si lavora al gioco \textbf{nelle pause del corso}: da qui a dicembre
        ci saranno due pause, quindi non ci sarà lezione tutti i lunedì e tutti i
        venerdì;
  \item ci saranno \textbf{review intermedie} con i \textbf{tutoraggi}, per lo più
        il \textbf{sabato mattina}, dalle 11 alle 12, con Ruggero su Google Meet.
        Servono ad aiutare i gruppi a portare avanti il project work.
\end{itemize}

\subsection{Il calendario delle lezioni}

Il \textbf{calendario} è nel calendario studenti della cartella condivisa.
Attenzione: le date \textbf{possono variare}. ``Il 90\% rimarranno così'', ma
può capitare che, per motivi di lavoro, qualcuno stia male o non possa, o che
``scompaia l'elettricità da tutte le città d'Italia''. E l'ordine delle lezioni
può cambiare se i docenti vedono che agli studenti serve prima una lezione
invece di un'altra.

Ci sono \textbf{due pause importantissime}:
\begin{itemize}
  \item dal \textbf{13 ottobre al 5 novembre}: è il periodo delle grandi fiere
        (\textbf{Essen} e \textbf{Lucca}\footnote{La trascrizione dice ``Essen,
        Luca, Cose'': si tratta con ogni probabilità della fiera Spiel di Essen,
        in Germania, e di Lucca Comics \& Games, che si tengono entrambe in
        ottobre.}), ``un periodo un po' frenetico'', e cominciano anche i corsi
        accademici, quindi la pausa permette ai docenti di respirare. Per gli
        studenti è la pausa più importante: quasi un mese in cui sviluppare la
        parte importante del gioco;
  \item dal \textbf{24 novembre all'11 dicembre}, dopo qualche lezione a novembre.
\end{itemize}

E ci sono \textbf{due date da evidenziare}:
\begin{itemize}
  \item il \textbf{25 settembre}: nei primi 5--10 minuti della lezione (che per il
        resto sarà tenuta da Beatrice) i docenti annunceranno i tre gruppi;
  \item il \textbf{12 ottobre}: il \textbf{pitch} dei giochi, con l'idea
        embrionale. Tra il 25 settembre e il 12 ottobre ogni gruppo prepara un
        breve pitch in formato slide; quel giorno i docenti lo rivedono e danno
        feedback: se il gioco va bene, se non va bene, come svilupparlo.
\end{itemize}
Dal 12 ottobre fino all'inizio di dicembre si sviluppa il gioco, con i
tutoraggi: ``sfruttate questi momenti di pausa''.

\subsection{Il calendario del project work}

La slide successiva, ``reale e aggiornata'', riporta il calendario del project
work (tabella~\ref{a1:tab:calendario} e figura~\ref{a1:fig:calendario}).
``Per ora non pensate al project work: ci penserete dal 25 settembre in poi''.

\begin{table}[htbp]
\centering
\begin{tabular}{@{}lp{9.2cm}@{}}
\toprule
\textbf{Data} & \textbf{Che cosa succede} \\
\midrule
ven 25 settembre\footnotemark & lezione: annuncio dei gruppi e del
\emph{mandate}, cioè di che cosa si dovrà fare \\
sab 3 ottobre, ore 11 & prima review e tutoraggio con Ruggero (dopo una
decina di giorni di lavoro) \\
lun 12 ottobre & pitch del gioco durante la lezione \\
13 ottobre -- 5 novembre & prima pausa: sviluppo del gioco \\
sab 7 novembre & tutoraggio \\
sab 14 novembre & tutoraggio \\
24 novembre -- 11 dicembre & seconda pausa: sviluppo del gioco \\
ven 4 dicembre, entro mezzanotte & consegna del gioco sul Drive (non c'è
lezione) \\
ven 11 dicembre & i giochi vengono visti a lezione \\
14 dicembre & giornata finale, con la tavola rotonda con gli editori \\
\bottomrule
\end{tabular}
\caption{Il calendario del project work come lo presenta Daniele. Le date
possono variare.}
\label{a1:tab:calendario}
\end{table}
\footnotetext{Daniele dice ``lunedì 25 settembre'', ma nel 2026 il 25 settembre
è un venerdì; il giorno della settimana nella tabella è corretto di
conseguenza.}

\paragraph{Perché la consegna è il 4 dicembre.} Il 4 dicembre non c'è lezione:
si deve solo consegnare il gioco, entro mezzanotte, sul Drive. La presentazione
a lezione è l'11 dicembre, ma i tre docenti hanno bisogno di \textbf{una
settimana} per rivedere il gioco in privato e valutare anche la
\textbf{documentazione}; per questo la consegna non coincide con la
presentazione. Che cosa voglia dire esattamente ``consegnare il gioco'' verrà
spiegato più avanti.

\paragraph{La giornata finale.} Il 14 dicembre c'è la \textbf{giornata finale},
con una \textbf{tavola rotonda con gli editori}: di solito viene invitato Flavio,
che è anche docente di game production, più qualche altro editore, ``vediamo
quest'anno chi riusciamo a portarvi''.

\begin{figure}[htbp]
\centering
\begin{tikzpicture}[
    font=\footnotesize,
    x=0.1cm,
    evento/.style={draw, thick, rounded corners=2pt, fill=white,
                   text width=1.8cm, align=center, inner sep=2pt},
    pausa/.style={fill=gray!20, draw=gray!60},
    tacca/.style={thick}
  ]
  % asse: giorno 0 = 25 settembre; 1 unità = 1 giorno
  \draw[pausa] (18,-0.35) rectangle (41,0.35);
  \draw[pausa] (60,-0.35) rectangle (77,0.35);
  \node[font=\scriptsize\itshape] at (29.5,0.18) {pausa};
  \node[font=\scriptsize\itshape] at (65,0.18) {pausa};
  \draw[-{Stealth[length=2.5mm]}, thick] (-3,0) -- (85,0);
  % eventi sopra
  \draw[tacca] (0,-0.15) -- (0,0.9);
  \node[evento, above] at (0,0.9) {25 set\\ gruppi e \emph{mandate}};
  \draw[tacca] (17,-0.15) -- (17,2.2);
  \node[evento, above, fill=gray!8] at (17,2.2) {12 ott\\ \textbf{pitch}};
  \draw[tacca] (70,-0.15) -- (70,0.9);
  \node[evento, above, fill=gray!8] at (70,0.9) {4 dic\\ \textbf{consegna}};
  % eventi sotto
  \draw[tacca] (8,0.15) -- (8,-0.9);
  \node[evento, below] at (8,-0.9) {3 ott\\ tutoraggio};
  \draw[tacca] (43,0.15) -- (43,-0.7);
  \draw[tacca] (50,0.15) -- (50,-0.7);
  \draw (43,-0.7) -- (50,-0.7);
  \draw[tacca] (46.5,-0.7) -- (46.5,-0.9);
  \node[evento, below] at (46.5,-0.9) {7 e 14 nov\\ tutoraggi};
  \draw[tacca] (77,0.15) -- (77,-0.9);
  \node[evento, below] at (77,-0.9) {11 dic\\ giochi a lezione};
  \draw[tacca] (80,-0.15) -- (80,2.2);
  \node[evento, above] at (80,2.2) {14 dic\\ giornata finale};
\end{tikzpicture}
\caption{Una linea del tempo del project work, ricostruita dalle slide (scala
approssimativa: dal 25 settembre al 14 dicembre). Le bande grigie sono le due
pause senza lezioni, in cui i gruppi sviluppano il gioco.}
\label{a1:fig:calendario}
\end{figure}

Finita questa parte, Daniele apre alle domande: ``Domande sul calendario? Sulle
slide che ho fatto? Considerazioni? Non abbiate timore, io non ci credo''. E la
domanda arriva, sull'unico tema che aveva lasciato da parte.

% ---------------------------------------------------------------------
\section{L'intelligenza artificiale nel corso}
\label{a1:sec:ia}

\subsection{La slide di Daniele}

Daniele aveva già presentato una slide che definisce ``importantissima''. ``Lo so
che è paradossale, me ne rendo conto. Lo so che magari mi darete del luddista,
dell'anti-tecnologia: non è così. Uso l'intelligenza artificiale per lavoro. Ma
in questo corso, vi prego, \textbf{non usate l'intelligenza artificiale}''. La
richiesta vale per le attività, per i task e \textbf{anche per la
prototipazione}. Le ragioni:
\begin{itemize}
  \item ``voglio che siate voi a fare il gioco, voglio che siate voi a imparare a
        farlo''. Daniele parla da utilizzatore: prima di ChatGPT, che ha smesso
        di usare, oggi soprattutto di Claude, di cui ``si sottovaluta quanto
        possa essere incerto, insicuro, impreciso, non all'altezza'';
  \item nella game industry, \textbf{i publisher rifiutano automaticamente} i
        pitch e i prototipi fatti con l'intelligenza artificiale. Daniele lo ha
        visto di persona in ``tantissime pitch review'' quell'estate, anche al
        \textbf{First Playable} di Firenze\footnote{La trascrizione dice ``First
        Playboy a Firenze'': si tratta del First Playable, l'evento di business
        dell'industria videoludica che si tiene a Firenze.}: anche nel mondo dei
        videogiochi i publisher che vedono pitch o prototipi fatti con l'IA
        ``rigettano automaticamente la build'', perché hanno ancora timore
        dell'IA e delle sue possibili conseguenze;
  \item ``non so come sarà il futuro'': non sa se bisognerà abituarsi a usarla
        tutto il giorno. Non è ``totalmente contrario'', purché sia usata ``con
        raziocinio''; secondo lui oggi, negli Stati Uniti, non lo è;
  \item ma in un corso che si chiama \emph{Analog Game Design} gli studenti devono
        saper fare i giochi da soli. In futuro, se vorranno farsi aiutare
        dall'IA per essere più efficienti, saranno liberissimi di farlo; durante
        il corso no, perché i docenti vogliono valutare \textbf{le loro
        competenze e capacità}.
\end{itemize}

\subsection{La domanda: che cosa vuol dire usare l'IA in un gioco analogico?}

La domanda di uno studente è precisa: che cosa si intende per ``task fatti con
l'IA''? Qual è l'uso dell'IA in un gioco analogico?

\paragraph{Ruggero: l'esperienza di chi valuta i prototipi.} Ruggero risponde
dalla sua esperienza nel settore, dove riceve prototipi da valutare per
un'eventuale pubblicazione. Distingue due usi:
\begin{enumerate}
  \item le \textbf{immagini generate}: illustrazioni, graphic design e così via.
        Per molti publisher è ``un \emph{reject} immediato, proprio immediato'';
  \item l'IA usata per il \textbf{regolamento}, per scrivere le regole e le
        meccaniche. È più difficile da individuare, ma molti publisher se ne
        accorgono comunque, perché le IA ``non sono ancora brave abbastanza da
        creare un regolamento coeso meccanicamente e scritto con un senso per il
        vostro gioco''. Quello che fanno è prendere parti di altri regolamenti,
        elaborare le strutture che vedono nei regolamenti pubblicati e
        riapplicarle. Ma molti giochi si prestano poco a un regolamento
        strutturato come quello di altri, perché \textbf{ogni gioco è un
        ecosistema meccanico a sé}. Il risultato sono regolamenti con
        riferimenti non immediati, e soprattutto con \textbf{la stessa meccanica
        descritta in modo diverso in pagine diverse}: due difetti che fanno
        perdere tempo all'editore, e in questo caso ai docenti.
\end{enumerate}
Poi aggiunge un punto ``personale, soggettivo'': i docenti sono interessati a
vedere il lavoro, la creatività, ``il frutto della vostra creazione''. E
viceversa gli studenti vogliono i commenti dei docenti: ``immaginatevi se voi
fate un gioco e noi vi diamo i commenti scritti con ChatGPT''. Altrimenti, dice
ridendo, gli studenti fanno il gioco con ChatGPT, i docenti lo danno a ChatGPT
per i commenti, ``e facciamo molto alla svelta''.

\paragraph{Daniele: perché gli studenti usano l'IA.} Daniele racconta la sua
esperienza con gli studenti, che secondo lui ricorrono all'IA in due momenti:
\begin{enumerate}
  \item quando vanno in \textbf{crisi}: in ansia da prestazione o per i tempi.
        Cercano di risolvere il problema della propria disorganizzazione
        facendolo risolvere all'IA, senza accorgersi che l'output è ``molto,
        molto scarso'';
  \item quando hanno l'impressione che l'intelligenza dell'IA sia
        \textbf{migliore} di quello che scrivono loro, o di quello che c'è in
        giro. Per Daniele è una ``mancanza di palato critico'': l'output dell'IA
        sembra buono in superficie, ma quando si va nel dettaglio e ci si fa
        delle domande ``di solito fa un po' di \emph{crumbling}'', si sbriciola.
\end{enumerate}
Lui usa Claude ogni giorno e ne conosce i limiti: ``non c'è nemmeno una frase
che io copio direttamente''. Gli fa fare il lavoro più noioso, più ripetitivo e
soprattutto \textbf{strutturato}, ``replicare delle cose come se fosse una
macro''. Lo usa per il testo e per il codice, mai per le immagini: per motivi di
\textbf{diritto d'autore}, e perché ha molti amici artisti ed è stufo di vedere
il loro lavoro ``dilaniato'' dall'IA.

Ammette che i docenti non possono controllare gli studenti; per questo
``ve lo consiglio''. E riprende quella che definisce la cosa più bella detta da
Ruggero (``ti ruberò la frase per il corso di narrative''): i docenti i feedback
li danno di persona. Daniele a volte manda 30 minuti di audio a uno studente, o
fa intere call per rivedere script, pitch e giochi: ``è il mio lavoro, mi piace
farlo e voglio farlo io, quindi mi aspetto che lo studente lo faccia lui''. Come
freelance, poi, è ``stufo marcio'' dei clienti che arrivano con qualcosa fatto
con l'IA (``guarda che bello, pubblica questo''), perché il più delle volte il
risultato è pessimo: non sanno come usarla. ``È anche una questione di
competenza''.

\paragraph{Il prototipo non deve essere bello.} Questo vale anche per il
\textbf{prototipo}, su cui si tornerà nelle lezioni di prototipazione: in un
prototipo non servono ``graficoni allucinanti''. Le grafiche bellissime fatte con
l'IA spesso \textbf{peccano di user experience}: sembrano bellissime, ma quando
il giocatore deve leggere la carta non capisce qual è il costo, qual è il
potere, qual è l'effetto.

\paragraph{Beatrice: si perde l'occasione di imparare.} Beatrice aggiunge due
motivi che valgono ``al di là dei nostri credo''. Il primo: chi usa l'IA in
sostituzione (per esempio per scrivere un testo, per accelerare i tempi o per
trovare strategie) \textbf{perde la parte in cui può imparare}. Nel game design
bisogna ``tenere insieme i pezzi'' di un regolamento, ma anche dell'esperienza
dell'utente e di come trasmetterla; i passaggi necessari (creare il regolamento,
creare le meccaniche, sperimentare, raccogliere i feedback) ``si capiscono solo
vivendoli''. E si sta pagando un corso per imparare. Il secondo, sul piano
editoriale: chi presenta a un editore un prodotto che ha come base l'IA rischia
da un lato di non avere imparato abbastanza (e di ritrovarsi ``a piedi'' sul
campo senza IA), dall'altro che l'editore se ne accorga, ed è un tema
\textbf{reputazionale}. Ci sono editori ``più illuminati di altri'', quindi
anche qui c'è un mercato da capire; ma i docenti preferiscono dire che cosa è
più sensato per il mercato come lo conoscono.

\paragraph{Ruggero: non abbiate paura del giudizio.} Ruggero vuole dare uno
``spotlight'' a un punto detto da Daniele e rimasto ``semi-ignorato'': a volte
l'uso dell'IA nasce dalla \textbf{paura del giudizio}, dal voler consegnare
qualcosa di più pulito, più strutturato possibile. ``Vi giuro, non ve ne
preoccupate''. Ruggero detesta i voti e non ragiona per voti; non valuterà
nessuno per grafiche fatte a matita. Insegna da anni, ha avuto studenti con
dislessia e discalculia: refusi ed errori non gli interessano, ``non è un corso
di editing o di proofreading''. Gli interessa che le \textbf{meccaniche siano
coerenti con lo scopo di design}. Immagini fatte a matita o con Paint ``non sono
un problema, anzi, le preferisco'' rispetto all'``ennesima immagine fatta con
ChatGPT o Midjourney''.

\paragraph{Daniele: il pensiero critico, costo e qualità.} L'ultimo intervento
è di Daniele, sul \textbf{pensiero critico}. Ragionando per assurdo (``un
pensiero di Belzebù''), immagina che un domani l'IA venga usata ovunque. I suoi
limiti, secondo lui, resterebbero due:
\begin{description}
  \item[Il costo] Tutti dicono che sta diminuendo e che diminuirà, ma non se ne
        può essere sicuri al 100\%.
  \item[La qualità] Cioè la capacità di leggere l'output. Se si dà un regolamento
        a Claude o a ChatGPT chiedendo quali problemi ha, il modello risponde con
        i problemi che, in percentuale, è più probabile che gli esseri umani
        abbiano già individuato nei dati su cui è stato addestrato. Ne derivano
        due problemi: se il problema del proprio gioco non è mai stato
        individuato prima, il modello non lo vedrà; e viceversa può segnalare
        problemi che il gioco non ha.
\end{description}
Fa l'esempio di Grammarly, che usa nella versione premium per rivedere i testi
in inglese: spesso non accetta i suoi consigli, perché se vuole scrivere un
dialogo di videogioco in \textbf{dialetto Cockney} londinese, lo deve scrivere
in Cockney. Se Grammarly dice che è sbagliato, ``sei tu che non hai capito il
motivo per cui voglio scrivere questo testo in Cockney''. La conclusione: prima
di usare un'IA per un'attività bisogna sapere \textbf{che cosa si vuole
ottenere}, qual è l'obiettivo, e quali \textbf{competenze} servono per valutare
in modo critico se l'output di un qualunque LLM è davvero valido. Daniele non
vuole dilungarsi oltre: è un argomento ``molto caldo'' per lui, che è anche
scrittore e si sente dire da tutti ``ti sostituiranno''.

Subito dopo si presenta finalmente Gianmarco (lo si è visto nella
sezione~\ref{a1:sec:partecipanti}), e poi si fa una pausa di una decina di
minuti. Daniele avverte che nella terza parte lui non ci sarà (alle sette e un
quarto ha un impegno di lavoro), e lo stesso vale per Beatrice: la parte finale
di domande e risposte sarà con il solo Ruggero, ``in ottime mani''.

% ---------------------------------------------------------------------
\section{Le figure professionali dell'industria analogica}
\label{a1:sec:figure}

Dopo la pausa (in cui si scopre che Angelo lavora in un piccolo birrificio in
Brianza e si scherza sul rischio di un \emph{drinking game}, tra la birra e la
vendemmia) la lezione riprende con le slide di Ruggero. ``Oggi è una lezione
introduttiva'': gli argomenti si vedranno ``molto dall'alto'' e verranno
approfonditi nelle lezioni successive. Ma servono già ora, per creare
\textbf{l'inizio di un linguaggio condiviso}: ci sono termini che torneranno
spesso, a cominciare da quelli del mercato e dell'industria del gioco
analogico, un'industria ``molto varia, molto caotica, anche molto mutevole'',
con tanti tipi diversi di giochi e altrettanti tipi diversi di entità che ci
lavorano.

\subsection{Una decina di etichette}

I docenti hanno raccolto ``una decina di etichette'', cioè di \textbf{ruoli} che
nell'industria sono o possono essere presenti. Due avvertenze:
\begin{itemize}
  \item nelle aziende più piccole \textbf{la stessa persona copre spesso più
        ruoli}, per forza di cose. Sarà anche il caso degli studenti, che
        lavoreranno in gruppi di tre e non di venti: è facile che ciascuno finisca
        a fare più cose. I docenti ne terranno conto, e qui entreranno in gioco
        anche competenze diverse dal game design (programmi di grafica,
        illustrazione, impaginazione, layout, \emph{proofreading}\ldots);
  \item Daniele precisa che \textbf{non si chiede agli studenti di diventare}
        illustratori o altro. Ruggero concorda, ma aggiunge che è utile sapere
        anche \emph{che cosa non è il proprio lavoro}. Se in un gruppo nessuno sa
        illustrare, ai docenti va benissimo un prototipo che mostri dove vanno
        le immagini: \textbf{placeholder}, concept, anche semplici box di testo
        con scritto ``qui vorremmo un'immagine di\ldots''. ``È importante anche
        sapere che cosa sapete fare e che cosa non sapete fare''.
\end{itemize}
La tabella~\ref{a1:tab:figure} riassume le figure; le sottosezioni seguenti le
riprendono una per una, con i commenti dei docenti.

\begin{table}[htbp]
\centering
\small
\begin{tabular}{@{}p{3.3cm}p{9.6cm}@{}}
\toprule
\textbf{Figura} & \textbf{Di che cosa si occupa} \\
\midrule
Game designer & crea le regole alla base dell'esperienza: le meccaniche, come
interagiscono tra loro e come chi gioca interagisce con esse \\
Game developer & fa da ponte tra chi fa game design e il publisher: bilancia il
prototipo, cura i modelli matematici, adatta il gioco (e il paragame) alle
necessità del publisher \\
Editor (proofreader) & revisiona e adatta i testi a partire dalla prima stesura
del designer (correzione, linguaggio inclusivo\ldots) \\
Publisher (editore) & pubblica e produce fisicamente il gioco; il suo marchio è
sulla scatola o sulla copertina \\
Sales e marketing & il lato commerciale: vendite (anche B2B) e comunicazione del
prodotto \\
Graphic designer & icone meccaniche, wireframe, UX: il lato funzionale degli
elementi visivi \\
Illustratore & le illustrazioni: il \emph{look and feel}, il world building, il
``vibe'' del gioco \\
Sensitivity reader (cultural consultant) & verifica l'accuratezza e la
rappresentazione di un tema sensibile che il team non conosce \\
Producer & coordina il progetto (quasi project management), parla con le
tipografie, verifica che i materiali siano sostenibili economicamente \\
\bottomrule
\end{tabular}
\caption{Le figure professionali dell'industria del gioco analogico presentate
nella slide di Ruggero, in sintesi.}
\label{a1:tab:figure}
\end{table}

\subsection{Game designer}

Il \textbf{game designer} (o la game designer) è chi si occupa, ``sembra banale
a dirlo'', di \textbf{creare le regole alla base dell'esperienza} che si vuole
far arrivare. Vale per tutti i tipi di giochi, non solo per quelli analogici, e
che cosa faccia concretamente in un progetto o in una serie specifica ``può
mutare tantissimo''. Nell'analogico non verrà chiesto di programmare o di
prototipare in digitale; verrà chiesto di \textbf{creare le meccaniche di
gioco}, di capire \textbf{come le meccaniche si interfacciano tra loro} e
\textbf{come le persone che giocano si interfacceranno con le meccaniche}.

\subsection{Game developer}

Il \textbf{game developer} (o la game developer) sarà l'argomento delle due
lezioni di Marta. La parola \emph{development} è nota soprattutto
dall'ambito videoludico, ma nell'analogico ha un significato molto specifico
``che non c'entra niente con la programmazione'': il game developer è
\textbf{un ponte tra chi fa game design e il publisher}. Di solito si occupa di
\textbf{bilanciare il prototipo}, di curare i \textbf{modelli matematici}, di
adattare il gioco e il \textbf{paragame}\footnote{La trascrizione dice ``il
Wix, il paragame''. Il \emph{paragame} è, in termini generali, tutto ciò che
accompagna il gioco senza esserne il sistema di regole (regolamento, scatola,
materiali, comunicazione); i docenti lo tratteranno in una lezione dedicata, e
qui il termine è solo citato.} alle necessità del publisher. La distinzione
viene ripresa e approfondita nella sezione~\ref{a1:sec:design-development}.

\subsection{Editor}

L'\textbf{editor}, a volte chiamato \textbf{proofreader} a seconda dell'azienda,
si occupa della revisione dei \textbf{testi}: non della stesura iniziale, perché
di norma, data la scala dell'industria, è chi fa game design a scrivere un
\textbf{primo draft}, e l'editor ci passa sopra. Non solo per correggere i
refusi: spesso il suo lavoro è anche \textbf{adattare} i testi, per esempio
renderli \textbf{gender neutral}. Se ne vedranno casi studio più avanti.

\subsection{Publisher}

Il \textbf{publisher}, o \textbf{editore}, è l'entità che nell'ambito
commerciale \textbf{pubblica} giochi, giochi di ruolo o giochi da tavolo. È una
figura ``più solidificata'' nel gioco da tavolo e nel gioco di ruolo da tavolo;
\textbf{nell'ambito pervasivo non esiste proprio in questa forma}, come
mostrerà Beatrice parlando dei casi studio. Se si compra un gioco da tavolo o un
gioco di ruolo in un negozio, il marchio del publisher è di solito sul retro
della copertina o sul fronte della scatola: è l'entità che, di solito,
\textbf{produce fisicamente} il gioco.

\subsection{Sales e marketing}

L'ambito \textbf{sales e marketing}, a volte affidato a un'unica persona e a
volte separato, riguarda il lato commerciale. In alcune aziende è interno al
publisher, e si occupa di parlare con le fabbriche e le tipografie, di ``far
quadrare i conti''. In altre realtà è completamente diverso, specialmente nel
\textbf{B2B} o nei giochi \textbf{didattici} e \textbf{formativi}, di cui
Daniele e Beatrice ``sapranno parlare meglio''. Il tema viene ripreso più avanti
in dettaglio (sezione~\ref{a1:sec:marketing}).

\subsection{Graphic designer e illustratore}

Ruggero presenta insieme \textbf{graphic designer} e \textbf{illustratore}
proprio perché sono ruoli molto diversi che si tende a far coincidere:
\begin{itemize}
  \item chi fa \textbf{graphic design} si occupa delle \textbf{icone meccaniche},
        del \textbf{wireframing}, della \textbf{UX}: del lato \textbf{meccanico e
        funzionale} degli elementi visivi, cioè di tutto ciò che serve
        meccanicamente nel gioco, o anche alla sua ``qualità della vita'';
  \item chi fa \textbf{illustrazione} si occupa dell'\emph{art}, delle
        illustrazioni: copre il \textbf{look and feel}, la parte di
        \textbf{world building} e il ``vibe'' che si dà al gioco.
\end{itemize}

\paragraph{Daniele: che cos'è la user experience.} Daniele interviene sulla
\textbf{user experience}, ``importantissima per me'': ha un background da
\textbf{layout artist} (impaginava magazine), un mestiere in cui la UX conta
moltissimo. Per user experience non si intende solo l'\textbf{interfaccia
grafica}, la UI, che anche nei videogiochi è disegnata da un artista. La UX è
qualcosa che può fare anche il game designer: \textbf{disegnare l'esperienza del
giocatore}, cioè ``il gesto, l'intenzione, l'intuito'' che stanno dietro
l'aspetto visivo dell'interfaccia. Per fare UX non serve saper disegnare.
L'esempio è quello della \textbf{porta}: se si progetta una porta, bisogna
indicare all'utente come si apre, come si prende la maniglia; ``quella è user
experience''. Il design della maniglia in sé, invece, è
\textbf{interfaccia}.

\paragraph{Un aneddoto sulle icone.} Ruggero aggiunge che molte persone che
fanno \textbf{graphic design ben fatto non sanno illustrare}, e spesso vale il
viceversa. Lo si nota quando le icone di un gioco sono bellissime ma non
funzionali: andando a guardare i crediti, di solito si scopre che le icone le ha
fatte la persona che ha fatto le illustrazioni. ``Comunque non andiamo adesso in
questo, perché è un altro vaso di Pandora''.

\subsection{Sensitivity reader}

Il \textbf{sensitivity reader}, o \textbf{cultural consultant} a seconda del
progetto, fa un lavoro di \emph{sensitivity reading} del gioco. È una figura
molto usata nell'ambito \textbf{formativo e didattico}. L'esempio di Ruggero: se
si fa un gioco per un ospedale su una malattia specifica, chi lo progetta non è
un medico, ma il gioco deve essere \textbf{accurato} e \textbf{rappresentativo}
delle dinamiche che quella malattia comporta. Serve una persona che conosca quel
tema, anche se non sa che cosa sia il game design. Sono figure chiamate di
solito \textbf{a progetto}, per giochi specifici, che compensano competenze che
il team non ha ``per forza di cose, perché noi non sappiamo tutto''; e a volte,
per lavoro, si devono fare giochi su temi che non si conoscono personalmente.

\paragraph{Daniele: non è un censore.} Daniele segnala un problema italiano:
alcune persone, anche online, hanno confuso il sensitivity reader con un
\textbf{censore}, ``una persona che fa censura, soprattutto in materia di
\emph{woke}''. ``Assolutamente non è così''. La parola \emph{sensitive} non
significa che il pubblico sia sensibile a un tema, ma che \textbf{l'argomento di
cui tratta il gioco è sensibile}. In un gioco sulla medicina, realistico o no,
non si possono inserire imprecisioni. Per questo esistono i sensitivity reader
``più disparati'': da quelli per la cultura indiana o afroamericana fino a quello
su ``come si allevano gli gnu in Africa''. È una persona competente su un
argomento che aiuta a essere \textbf{veritieri}. Certo, uno degli ambiti
principali in cui si usano è quello culturale, ma ``non fate l'errore di pensare
che siano dei censori''. Daniele ha lavorato con ``centinaia'' di sensitivity
reader e nessuno gli ha mai detto ``non puoi scrivere questa cosa''. Di solito
sono come degli \textbf{editor}: danno opzioni, mostrano i punti problematici,
anche quelli di cui non ci si era nemmeno accorti.

Ruggero conferma con un esempio di quando lavorava in \textbf{Danimarca} con
l'\textbf{OMS}: giochi sul \textbf{diabete di tipo 2} nei bambini, per un target
di ragazzi delle scuole medie. C'era un team dell'OMS specializzato nella
comunicazione sul diabete, per garantire che il gioco fosse accurato e ne desse
una rappresentazione adatta al target. Il sensitivity reader aiuta a rendere il
gioco più accurato e ``più tematicamente coeso''.

\subsection{Producer}

Il \textbf{producer} è una figura a volte interna al publisher. In un progetto
grande è la figura di riferimento, ``quasi'' un project manager. In un progetto
piccolo si occupa spesso anche di parlare con le \textbf{tipografie} e di fare la
\textbf{lista dei materiali}, per assicurarsi che i materiali non siano solo
funzionali al gioco ma anche \textbf{economicamente sostenibili} da produrre:
``magari è un gioco bellissimo, ma poi costa più di quanto potete permettervi, e
nessuno vi comprerà mai un card game a 200 euro''.

\subsection{Perché conoscere queste figure}

Ruggero chiude la slide ``tornando a ombrello'': perché parlarne, se gli studenti
sono lì per fare game design? Non solo per saper dare un \textbf{nome} alle
figure. Non c'è niente di male a creare un gioco solo per sé; ma chi vuole fare
giochi a livello lavorativo, perché li giochino altre persone, ``magari più
persone possibile'', \textbf{si interfaccerà con queste figure}, e sapere che
cosa fanno e come parlarci è parte importante del percorso. Nel corso se ne
toccheranno diverse: il development con Marta, la produzione con Flavio, e
lezioni sulla UX, sul paragame, sul graphic design. ``Abituatevi mentalmente:
non siete in solitaria contro il mondo''.

% ---------------------------------------------------------------------
\section{Marketing e sales}
\label{a1:sec:marketing}

Arrivato all'angolo della slide con marketing e sales, Ruggero passa la parola
ai colleghi: ``io è meglio che non mi fate parlare di marketing e sales'', anche
per esperienze personali.

\subsection{Beatrice: i giochi trattati come prodotti}

Beatrice premette che le due voci sono in parte autoesplicative. Nel mondo
capitalista in cui viviamo, i giochi sono ``strumenti molto potenti per fare
\emph{n} cose'', ma sono trattati come \textbf{prodotti}, e il corso ragionerà
anche in quest'ottica. Nel \textbf{marketing} delle industrie ludiche e
videoludiche ci sono varie professionalità che gli studenti probabilmente già
conoscono: dal \emph{copy} alla grafica, tutto ciò che riguarda il
\textbf{vendere} il gioco o l'attività di gioco.

Il lato \textbf{sales} è forse la distinzione meno nota. Come in tante aziende, è
la parte che si occupa delle vendite \textbf{B2B}, cioè tra aziende: fare
accordi per vendere prodotti ludici a diverse realtà, per esempio ai
\textbf{distributori}, o a una casa editrice che vende alle librerie, che poi
vendono al consumatore finale. Chi si occupa di sales nel mondo ludico capisce
come vendere \textbf{alle aziende}. I due mondi si parlano spesso, e soprattutto
nelle realtà piccole si trovano figure che fanno ``un po' sales, un po'
marketing''.

\subsection{Daniele: che cos'è (e che cosa non è) il marketing}

Daniele il marketing si è trovato ``purtroppo'' a farlo, e lo fa anche per
l'accademia. ``Purtroppo'' perché in Italia è mal compreso: molti lo
identificano con l'\emph{advertisement} online. Non è così: il marketing è una
\textbf{disciplina molto più antica}, e in Italia ``abbiamo molto dimenticato''
come si fa il marketing tradizionale. Alcune cose da ricordare:
\begin{enumerate}
  \item \textbf{Il marketing non dice che cosa pubblicare, ma come pubblicarlo}.
        Esistono, nei videogiochi e nei giochi da tavolo, realtà che fanno
        prodotti \textbf{marketing driven}: l'esempio sono gli
        \textbf{hyper-casual game} per smartphone, quelli con il personaggio che
        va a destra e a sinistra e l'esercito che si allarga. Possono anche
        guadagnare, ma ``sono un po' la morte della creatività''. Il marketing
        ideale \textbf{prende un prodotto pensato da altri e gli dà le maggiori
        chance di essere venduto}: cerca il pubblico ideale per il prodotto che
        si ha, senza doverlo cambiare. Anzi, il mondo videoludico, e ludico in
        generale, ``è lastricato'' di errori di dipartimenti marketing che hanno
        cambiato completamente non il gioco, ma il modo in cui l'hanno
        pubblicizzato, e il gioco è fallito.
  \item \textbf{Usare strumenti del marketing non significa fare marketing}. Game
        design e game development a volte usano strumenti che vengono dal
        marketing (gli studenti ne vedranno con Beatrice, con Daniele e con
        Marta).
  \item \textbf{Il marketing è un gerundio}: si fa. È un'azione che si compie per
        comunicare; la comunicazione è attiva, e il marketing è l'insieme delle
        attività con cui si comunica il proprio gioco, dal \emph{copy} ai social.
        Può includere il digital marketing, ``ma non è solo quello''.
  \item \textbf{Il marketing esiste anche se si crede che non esista}. In Italia
        molti non vogliono ``fare marketing'' perché lo associano al lucro. Ma
        il marketing è \textbf{lo studio di come gli esseri umani comprano e
        vendono}, ed è innato in noi: ci sono logiche di cui siamo vittime, e
        ``anche se sei il bastian contrario, l'anticapitalista'', magari si
        compra comunque su Amazon perché fa comodo. Sono meccanismi che
        esistono, ``affascinanti anche, a volte''.
\end{enumerate}
Gli studenti non dovranno fare marketing, ma dovranno ricordare alcuni di questi
concetti, soprattutto nelle lezioni con Marta, che tratta argomenti legati al
marketing. E, chiude Daniele, come game designer e game developer ``per lo più
speriamo che il marketing non ci rovini i giochi'', soprattutto nel videoludico.

% ---------------------------------------------------------------------
\section{Game design e game development}
\label{a1:sec:design-development}

Ruggero torna sulla distinzione più importante della slide, ``proseguendo con
questa strada'':
\begin{description}
  \item[Game design] è chi \textbf{progetta l'esperienza} tramite la
        \textbf{creazione delle meccaniche} di gioco.
  \item[Game development] è chi fa \textbf{mediazione} tra la persona che scrive
        il game design e la \textbf{casa editrice}.
\end{description}

\subsection{I red flag dei designer}

Ruggero lo ripete perché lo vede ogni giorno nel suo lavoro: tra giochi
pubblicati e prototipi, ne valuta ``quasi uno al giorno, se faccio un conto
annuale'', e solo una minima percentuale arriva alla pubblicazione. Molte volte
i \textbf{red flag} vengono proprio da designer che non hanno presente questa
definizione, e che cominciano a mettere \textbf{clausole sulla non modifica del
gioco}: il gioco deve restare com'è. ``Questa cosa è impossibile con un
publisher, spesso''. Il publisher, per forza di cose, potrà chiedere di cambiare
la grafica, di reimpaginare il regolamento, di cambiare i materiali perché
troppo costosi: chiederà dei \textbf{compromessi}. E la persona che, all'interno
del publisher, fa questo lavoro è il \textbf{developer}.

Sta al designer capire \textbf{su che cosa del proprio gioco si può trovare un
compromesso e su che cosa no}. Publisher diversi lavorano in modo diverso, e il
designer ha sempre l'autorità di dire ``no, non voglio pubblicare questo gioco
se queste sono le modifiche''. Ma deve sapere chi fa il lavoro di development,
perché se il suo obiettivo è la \textbf{pubblicazione tradizionale} si
interfaccerà con questa figura.

\subsection{Daniele: una distinzione che manca al videogioco}

Daniele sottolinea che questa distinzione è \textbf{tipica del mondo analogico}
ed è ``molto interessante''. In Italia la parola \emph{developer} è diventata
negli anni quasi un sinonimo di programmatore di videogiochi, o di chiunque
lavori nel videoludico. Secondo lui c'entra anche il fatto che qualche anno fa
il sito \textbf{Gamasutra} ha cambiato nome in \textbf{Game Developer}, e la
parola è diventata un ombrello per tutte le professioni del videogioco. ``Mi fa
un po' ribrezzo'': il game developer dovrebbe essere un ruolo di design, un
designer che, come nell'analogico, fa \textbf{mediazione}. Nel videogioco i
publisher hanno in effetti i loro developer interni, che adattano il gioco, per
esempio alle console e ai loro kit di sviluppo. Ma la distinzione dovrebbe
esistere anche nel videogioco, perché è ``difficilissimo'' avere la mentalità
per fare le due cose insieme: chi fa design ha un focus che raramente si può
sovrapporre a quello di chi fa mediazione.

A Daniele piace la parola \emph{develop}, \textbf{sviluppare}, nel senso di
prendere qualcosa che è stato progettato e \textbf{renderlo possibile nel mondo
reale}. Ci sono progetti che hanno bisogno di uno sviluppo che ``polishi''
l'esperienza: la semplifichi, la migliori, la renda comprensibile e la concluda,
ma soprattutto la renda \textbf{stampabile} e \textbf{fruibile} per l'utente. La
metafora crossmediale è quella della \textbf{fotografia}: chi sviluppa le
fotografie non è chi scatta la foto, ma chi la rende visibile, anche
fisicamente. ``È un po' retro come metafora, ma ci capiamo''. ``È vero, però è
giusto'', conferma Ruggero.

% ---------------------------------------------------------------------
\section{Il mercato: Italia ed estero}
\label{a1:sec:mercato}

Il corso è in italiano e si svolge in Italia, ma bisogna conoscere anche il
\textbf{mercato estero}: molti giochi progettati da autori italiani vengono
pubblicati all'estero, e conoscere il mercato estero aiuta a capire come il
proprio gioco può arrivarci, ``o viceversa, come non può arrivarci''. La
tabella~\ref{a1:tab:mercato} riassume il confronto della slide.

\begin{table}[htbp]
\centering
\small
\begin{tabular}{@{}p{6.4cm}p{6.4cm}@{}}
\toprule
\textbf{Italia} & \textbf{Estero (soprattutto Europa)} \\
\midrule
il gioco è ancora visto come un'attività \textbf{infantile}, anche se la
normalizzazione è in crescita &
forte tradizione di gioco anche in un contesto \textbf{familiare adulto},
soprattutto ``dalla Germania in su'' \\
risposta \textbf{conservativa} sul tema ludico &
\textbf{premi di prestigio} riconosciuti in tutto il mondo (Spiel des Jahres,
As d'Or) \\
forte focus sulla \textbf{localizzazione} di giochi stranieri &
si cercano soprattutto giochi \textbf{originali} \\
investimento \textbf{molto basso} sulla creazione di giochi originali interni;
pochissimi designer assunti &
\textbf{developer interni} pagati con un contratto (Francia, Germania,
Danimarca\ldots) \\
\bottomrule
\end{tabular}
\caption{Il mercato del gioco analogico in Italia e all'estero, secondo la
slide di Ruggero e i commenti di Daniele.}
\label{a1:tab:mercato}
\end{table}

\subsection{L'Italia}

\begin{itemize}
  \item In Italia il gioco è ancora visto come un'\textbf{attività infantile},
        anche se la \textbf{normalizzazione} è in crescita. C'è ancora una forte
        risposta \textbf{conservativa} sul tema ludico.
  \item C'è un forte focus sulla \textbf{localizzazione} dei giochi stranieri: in
        una qualunque ludoteca o negozio, molti dei giochi da tavolo e di ruolo
        sono di publisher o designer esteri, pubblicati in italiano da publisher
        italiani. (La localizzazione, cioè portare un titolo dall'estero in
        Italia, esiste come funzione all'interno del publisher; Ruggero non l'ha
        messa nella slide delle figure perché non è il focus del corso.)
  \item C'è, ``purtroppo, preparatevi psicologicamente'', un \textbf{investimento
        molto basso} sulla creazione di giochi originali interni alla casa
        editrice. Pochissime case editrici italiane hanno designer interni pagati
        a tempo pieno con contratto a tempo indeterminato. Daniele lo dice senza
        mezzi termini: ``il mercato è molto piccolo, conosciamo praticamente
        tutto il settore'', e le persone che in Italia fanno game design
        analogico a tempo pieno con un contratto a tempo indeterminato ``sono
        meno delle dita delle mie due mani''. Di solito la formula è quella del
        \textbf{freelance}, del contratto, delle \textbf{royalty}. Chi va da un
        publisher aspettandosi di essere assunto a tempo pieno come game
        designer deve sapere che ``non siamo purtroppo ancora in quel tipo di
        situazione''.
  \item C'è però un contraltare: in Italia, soprattutto nel gioco di ruolo ma
        anche nel gioco da tavolo, c'è stata molta \textbf{imprenditoria}, anche
        grazie a \textbf{Kickstarter}, che ha aperto le porte ad altri sistemi
        di \textbf{crowdfunding}. È un altro argomento, perché richiede un team,
        un investimento e così via.
  \item E gli autori non mancano, anzi: le occasioni di incontro tra autori
        vedono ogni anno un aumento dei possibili autori\footnote{La trascrizione
        dice ``i DAG a Parma''; il riferimento è probabilmente a IDEAG, il
        circuito di incontri tra autori di giochi italiani, ma non è
        ricostruibile con certezza.}. Bisogna ``sgomitare'' tra tanti autori,
        che magari arrivano da esperienze molto diverse.
\end{itemize}
Ruggero precisa: l'investimento basso è ``dall'alto''. Le persone che fanno e
vogliono fare giochi sono tantissime e in aumento; scarseggia l'investimento,
sia dei publisher disposti a pagare stipendi sia dei \textbf{bandi} pubblici. I
bandi esistono (Beatrice ne parlerà per i giochi pervasivi, che spesso si
muovono in quella direzione, e Ruggero cita qualcuno che sta facendo un card
game con un bando italiano, ``una cosa assurda, ma lo sta facendo''), ma sono
una rarità, soprattutto rispetto all'estero.

\subsection{L'estero}

\begin{itemize}
  \item L'estero (dipende dalla nazione, ma in generale il resto d'Europa) ha una
        forte \textbf{tradizione di gioco} anche in un contesto \textbf{familiare
        adulto}, soprattutto ``dalla Germania in su''.
  \item Ha \textbf{premi di prestigio} riconosciuti in tutto il mondo, come lo
        \textbf{Spiel des Jahres} tedesco e l'\textbf{As d'Or} francese\footnote{La
        trascrizione dice ``lo Spield AR, la SDOR''.}.
  \item Cerca soprattutto \textbf{giochi originali}. Lo si vede dal fatto che
        molti autori italiani di giochi da tavolo pubblicano in Germania o in
        Francia e \textbf{non vengono pubblicati in Italia}, ``che è tutto
        dire''.
  \item Ha \textbf{figure di development interne}: non è raro vedere publisher
        francesi, tedeschi, danesi con developer pagati internamente, a
        contratto.
\end{itemize}
``Poi la situazione è ancora più drammatica'', ammette Ruggero, ``ma per ora
vogliamo tenere il morale un po' alto. È l'introduzione. Ci sarà tempo di andare
nella tristezza''.

Prima di proseguire, Ruggero insiste sulle domande: con lui si può interrompere
anche mentre parla, alzando la mano con le reaction o semplicemente togliendo il
microfono e chiedendo. ``Non mi offenderò mai, anzi, sarò felice di vedervelo
fare e mi ricorderò positivamente di voi''.

% ---------------------------------------------------------------------
\section{Perché serve una terminologia comune}
\label{a1:sec:terminologia}

La serie di slide successiva è ``una piccolissima parte di fondamenti di
design'' (``non spaventatevi, vi vedo che vi spaventate un po'\,''). Anche qui
molti temi verranno ripresi in dettaglio, ma l'idea è arrivare alla prossima
lezione con almeno un \textbf{glossario}, un \textbf{linguaggio condiviso}: se
qualcuno dice \emph{game developer}, tutti sanno di che cosa sta parlando.

Ruggero parte da una domanda, e precisa che non è retorica: \textbf{perché ci
serve una terminologia comune?} Gli studenti rispondono:
\begin{itemize}
  \item Gianmarco, ``molto pragmatico'': perché bisogna lavorare in team e
        capirsi. Vero, nel corso e anche fuori;
  \item un'altra studentessa: per la \textbf{comunicazione}, tra di noi e con
        tutti quelli che sono interessati nel settore. Vero anche questo:
        conoscere i termini comuni permette di parlare con chi fa game design o
        development nel dettaglio, nel tecnico, nello specifico.
\end{itemize}
Ruggero aggiunge altre due ragioni.

\paragraph{Saper comunicare il proprio gioco.} La prima riguarda la
\textbf{comunicazione del gioco} verso l'esterno. ``E lo dico non tirando acqua
al mio mulino, ma anzi togliendola'': \textbf{non è spesso il buon game design
a far diffondere un gioco}. ``Vorrei che lo fosse, sarebbe bellissimo, ma
purtroppo non è così''. Ci sono giochi bellissimi che non sanno comunicare, che
non sanno dire che cosa fanno, e che quindi è difficile, da fuori, trovare
invitanti; e ci sono giochi mediocri dal punto di vista del design che si
comunicano molto bene, e che quindi conoscono in molti. Saper parlare di quello
che si fa è utile \textbf{a sé stessi}, in primis.

\paragraph{Documentare per sé stessi.} La seconda è la \textbf{documentazione}:
documenti di design, di brainstorming e così via. ``Sembrerà una battuta, ma
provate a recuperare un documento di brainstorming una settimana dopo, se non
avete scritto bene le cose che avete scritto''. Chi l'ha fatto sa già di che
cosa si parla; ``chi non l'ha ancora vissuto, forse lo vivrà''.

% ---------------------------------------------------------------------
\section{Game e play}
\label{a1:sec:game-play}

La \textbf{prima, importantissima divisione terminologica} è quella tra
\emph{game} e \emph{play}. In italiano c'è un problema: \textbf{gioco} e
\textbf{giocare} hanno la stessa radice, mentre in altre lingue, tra cui
l'inglese, sono due parole strutturalmente diverse:
\begin{description}
  \item[Game] il \textbf{gioco come artefatto}: il gioco come ecosistema di
        meccaniche, come struttura di regole, come prodotto (``tra
        virgolette''), come oggetto fisico.
  \item[Play] l'\textbf{atto}, il verbo: le persone che giocano, l'approccio
        fisico all'attività del gioco.
\end{description}
È importante perché gli studenti faranno \textbf{game design}; esistono anche
ambiti di \emph{play design}, ma qui ``creerete giochi, non giocattoli''.

\subsection{Giochi e giocattoli}

Qual è la differenza? Uno studente risponde: \textbf{la presenza di regole}.
Esatto: un gioco ha delle regole, un'attività di \emph{play} non
necessariamente. Con un giocattolo si può fare \emph{play}: c'è chi gioca con i
Lego, chi gioca a guardia e ladri o a nascondino al parco, e magari alcune di
quelle regole sono ``più o meno semi-solidificate''. Ma il \emph{play} permette
spesso un approccio non strettamente legato alle regole.

Viceversa, \textbf{un game può esistere anche se nessuno lo sta giocando}. Se si
prepara su un tavolo una scacchiera con i pezzi degli scacchi, chi passa e la
vede riconosce il gioco degli scacchi, anche se nessuno sta giocando. Se invece
si vedono dei bambini correre in un parco giochi, quei bambini \emph{stanno
giocando}, ma non stanno giocando \emph{a un gioco} specifico. ``È importante
separare questa cosa''; ci si tornerà, ``ma per ora abituatevi a questa
differenza''.

\paragraph{Approfondimento: giocattoli, giochi e regole.} La stessa distinzione
è il punto di partenza del capitolo che Jesse Schell, in \emph{The Art of Game
Design}, dedica alla definizione di gioco (\emph{The Experience Rises Out of a
Game}). Schell nota che usiamo perfino un linguaggio diverso: ``un gioco è
qualcosa a cui si gioca'' (\emph{a game is something you play}), ``un giocattolo
è qualcosa con cui si gioca'' (\emph{a toy is something you play with}); e
arriva a definire il giocattolo come ``un oggetto con cui si gioca''. Analizzando
poi la definizione di gioco di Avedon e Sutton-Smith, ne ricava che ``i giochi
hanno regole, i giocattoli no'', e che le regole sono uno degli aspetti
\emph{definitori} del gioco, insieme alla volontarietà, agli obiettivi, al
conflitto e alla possibilità di vincere o perdere. Schell osserva anche che
molti giochi sono costruiti \textbf{sopra} un giocattolo (``una palla è un
giocattolo, ma il baseball è un gioco''), e tra i suoi consigli di
prototipazione c'è proprio \emph{Build the Toy First}: costruire prima il
giocattolo, cioè qualcosa che sia già piacevole da manipolare, e poi il gioco
intorno. È una lettura utile della distinzione di Ruggero: il \emph{play} può
fare a meno delle regole, il \emph{game} no.

% ---------------------------------------------------------------------
\section{Il framework MDA (e il DPE)}
\label{a1:sec:mda}

\subsection{Meccaniche, dinamiche, estetiche}

Un'altra terminologia fondamentale è quella del framework \textbf{MDA}
(\emph{Mechanics, Dynamics, Aesthetics}; in italiano a volte ``MDE''). Viene da
un \textbf{paper} nato nell'ambito del videogioco\footnote{È il paper di Robin
Hunicke, Marc LeBlanc e Robert Zubek, \emph{MDA: A Formal Approach to Game
Design and Game Research}, del 2004. Daniele aggiunge il paper alla cartella
condivisa durante la lezione.}, che ha ``solidificato'' queste terminologie, e
divide un gioco in tre aree:
\begin{description}
  \item[Meccaniche] le \textbf{regole per come sono scritte}. ``Ogni volta che
        pescate un 2 fate dei punti'': quella è una meccanica.
  \item[Dinamiche] l'\textbf{applicazione in tempo reale} delle meccaniche. La
        dinamica è il giocatore che, pescando una carta, \textbf{spera} che esca
        un 2. Quel processo non è una regola: ``il gioco non vi dice che dovete
        sperare che esca un 2''.
  \item[Estetiche] l'insieme di quello che il gioco vuole comunicare,
        l'\textbf{esperienziale}.
\end{description}

L'esempio di Ruggero mette insieme le tre aree. In un gioco in cui pescare un 2
fa fare punti (meccanica) e in cui quindi si spera che esca un 2 (dinamica),
\textbf{che estetica si ottiene se nel mazzo di 100 carte c'è un solo 2? E se ce
ne sono 50?} L'emozione di aver pescato \emph{l'unico} 2 del mazzo è
completamente diversa da quella di pescarne uno sapendo che ce ne sono
cinquanta e che è molto probabile trovarlo. Le regole sono quasi identiche; cambia
l'esperienza.

\subsection{Il falso amico: ``estetica''}

Daniele avverte di non sottovalutare la parola \textbf{estetica}, che in italiano
è un \textbf{falso amico}: sembra indicare la parte grafica, e invece non lo è.
Si incontreranno designer che hanno studiato MDA e altri che non l'hanno
studiato, e lo si capirà proprio da questo: i secondi diranno che le
\emph{aesthetics} sono la grafica. Le \emph{aesthetics} sono la parte
\textbf{esperienziale}: come nell'esempio di Ruggero, quello che lascia al
giocatore aver pescato l'unico 2 del mazzo.

\subsection{Il DPE: design, play, experience}

Per gli italiani, aggiunge Daniele, è più facile capire l'MDA attraverso un
altro paper, quello del \textbf{DPE}\footnote{È il framework \emph{Design, Play,
Experience} di Brian Winn, proposto come estensione dell'MDA per i giochi
educativi (\emph{serious games}). Anche questo paper viene aggiunto da Daniele
alla cartella condivisa.}, che di fatto fa un parallelismo assegnando
\textbf{design, play, experience}:
\begin{itemize}
  \item le \textbf{meccaniche} sono il lato \textbf{design};
  \item le \textbf{dinamiche} sono il lato \textbf{play};
  \item le \textbf{estetiche} sono l'\textbf{experience}, l'esperienza che riceve
        chi gioca.
\end{itemize}
``Se vivete le \emph{aesthetics} come l'esperienza che vive chi sta giocando,
state a postissimo''. In questo momento Beatrice lascia la lezione (``ciao,
Bea''). I paper di MDA e DPE sono in inglese; chi vuole può dargli una lettura.

\begin{figure}[htbp]
\centering
\begin{tikzpicture}[
    font=\small,
    blocco/.style={draw, thick, rounded corners=3pt, fill=gray!12,
                   minimum width=3.2cm, minimum height=1.0cm, align=center},
    dpe/.style={font=\footnotesize\itshape, align=center},
    esempio/.style={draw, rounded corners=2pt, text width=3.3cm,
                    align=left, font=\footnotesize, inner sep=4pt},
    freccia/.style={-{Stealth[length=2.5mm]}, thick}
  ]
  \node[blocco] (m) at (0,0)   {\textbf{Meccaniche}\\ \footnotesize le regole scritte};
  \node[blocco] (d) at (4.3,0) {\textbf{Dinamiche}\\ \footnotesize le regole in azione};
  \node[blocco] (a) at (8.6,0) {\textbf{Estetiche}\\ \footnotesize l'esperienza};
  \draw[freccia] (m) -- (d);
  \draw[freccia] (d) -- (a);
  \node[dpe, above=0.15cm of m] {DPE: \textbf{design}};
  \node[dpe, above=0.15cm of d] {DPE: \textbf{play}};
  \node[dpe, above=0.15cm of a] {DPE: \textbf{experience}};
  \node[esempio, below=0.45cm of m] (em)
    {``Ogni volta che peschi un 2 fai punti''};
  \node[esempio, below=0.45cm of d] (ed)
    {pescando, il giocatore \emph{spera} che esca un 2 (nessuna regola glielo
     chiede)};
  \node[esempio, below=0.45cm of a] (ea)
    {un solo 2 su 100 carte: pescarlo è un evento;\\
     50 su 100: è quasi scontato};
  \draw[dashed] (m) -- (em);
  \draw[dashed] (d) -- (ed);
  \draw[dashed] (a) -- (ea);
  % prospettive
  \draw[freccia, gray] ([yshift=-0.3cm]em.south west) --
    node[below, font=\footnotesize\itshape, text=black]
    {il designer agisce sulle meccaniche}
    ([yshift=-0.3cm]ea.south east);
  \draw[freccia, gray] ([yshift=-1.2cm]ea.south east) --
    node[below, font=\footnotesize\itshape, text=black]
    {chi gioca incontra prima l'esperienza}
    ([yshift=-1.2cm]em.south west);
\end{tikzpicture}
\caption{Il framework MDA con l'esempio del ``2'' usato da Ruggero, e la
corrispondenza con il DPE suggerita da Daniele. Le due frecce in basso
riassumono la lettura classica del paper: il designer costruisce le meccaniche
per ottenere un'esperienza, mentre il giocatore parte dall'esperienza.}
\label{a1:fig:mda}
\end{figure}

\paragraph{Approfondimento: le due direzioni dell'MDA e le otto estetiche.} Le
frecce della figura~\ref{a1:fig:mda} riprendono un'idea centrale del paper
originale, che la lezione non esplicita: designer e giocatore guardano il gioco
da lati opposti. In \emph{Challenges for Game Designers}, Brenda Brathwaite e
Ian Schreiber la riassumono così: i designer si chiedono quale estetica
vogliono ottenere, definiscono le dinamiche che porterebbero a quella
sensazione, e poi creano le meccaniche che producono quelle dinamiche. Lo stesso
libro conferma l'avvertenza di Daniele sul falso amico: ``le \emph{aesthetics}
non si riferiscono all'aspetto del gioco, ma alla risposta emotiva'' che il
designer vuole suscitare. E offre esempi di dinamiche nel senso di Ruggero: in
uno sparatutto i punti di \emph{spawn} sono una meccanica, mentre appostarsi
accanto a uno di essi per colpire chi riappare (lo \emph{spawn camping}) è una
dinamica; negli scacchi i movimenti dei pezzi sono meccaniche, mentre le
aperture codificate sono dinamiche. Per dare un nome alle estetiche, il paper
MDA propone una tassonomia, poi detta degli ``otto tipi di \emph{fun}'', che
Schell riprende in \emph{The Art of Game Design} come \emph{LeBlanc's Taxonomy of
Game Pleasures}: \textbf{sensazione}, \textbf{fantasia}, \textbf{narrativa},
\textbf{sfida}, \textbf{compagnia} (\emph{fellowship}), \textbf{scoperta},
\textbf{espressione} e \textbf{sottomissione} (il piacere di entrare nel
``cerchio magico'' e lasciarsi andare a un passatempo). È un primo esempio di
quello che i docenti chiedono con le parole vietate della prossima sezione:
invece di dire che un gioco è ``divertente'', dire \emph{quale} esperienza
offre.

% ---------------------------------------------------------------------
\section{Le parole vietate}
\label{a1:sec:parole-vietate}

``Attenzione, adesso qua triggeriamo delle persone''. Ruggero annuncia una regola
che ``tornerà'': nel corso ci sono delle \textbf{parole vietate}. ``Non perché
siamo dei censori''; il perché si vedrà in una lezione apposita, perché
spiegarlo adesso richiederebbe un'ora. Ma Daniele vuole spiegarne ``un
minimo''.

\subsection{La regola}

In \textbf{nessun materiale scritto} consegnato ai docenti, task compresi, sono
accettate queste parole, né in inglese né in italiano:
\begin{center}
\begin{tabular}{@{}ll@{}}
\toprule
\textbf{Inglese} & \textbf{Italiano} \\
\midrule
\emph{fun} & divertimento \\
\emph{immersion} & immersione \\
\emph{entertainment} & intrattenimento \\
\bottomrule
\end{tabular}
\end{center}
``Non li voglio neanche vedere''. Daniele lo ripete con un paragone: ``se mi fate
i disegni con Paint, benissimo, top, fantastico, siamo ancora best friend. Se mi
scrivete \emph{divertimento}, io sbrocco''. È una \textbf{regola dell'accademia}:
negli stessi corsi triennali di narrative design e game design le parole sono
vietate agli studenti per tre anni; ``per voi solo per tre mesi''. Daniele
aggiungerebbe volentieri un quarto termine, \textbf{gamification}.

\subsection{Perché}

Il motivo approfondito lo darà Ruggero nella lezione apposita. Daniele dà il
suo. Ruggero e Diego Ricchiuti sono molto contrari a \emph{fun} e a
\emph{immersion}; lui ha aggiunto \emph{intrattenimento}. Sono \textbf{parole
troppo passe-partout e troppo ambigue}:
\begin{itemize}
  \item \textbf{intrattenimento} in Italia è stata usata per anni come un modo
        per \textbf{sminuire}: ``ma sì, tanto è un gioco solo di
        intrattenimento, è un romanzo solo di intrattenimento''. Ma anche se un
        prodotto è ``solo di intrattenimento'', non vuol dire che non ci si debba
        lavorare con attenzione. E poi intrattenimento ``non significa niente'':
        ci sono tanti modi in cui una cosa può intrattenere;
  \item lo stesso vale per \textbf{fun}, e per \textbf{immersion}, che ha un
        problema in più: è una parola passe-partout per cui si può dire che un
        gioco è ``immersivo'' nel senso generale del termine, ma poi non si sa
        davvero che cosa voglia dire, \emph{immersivo in che senso}.
\end{itemize}

L'obiettivo didattico è che gli studenti \textbf{vadano oltre queste parole}:
``un po' come quando alle medie vi dicono smettetela di dire che tutto è bello,
perché ci sono vari gradi di bello, ci sono altri modi di dire bello''. Bisogna
\textbf{decostruirle}: quando si sta per usarle, chiedersi \textbf{che cosa si
intende} con quella parola e cercarne una \textbf{più precisa}.

Ruggero chiede, in uno dei ``rari casi'', di \textbf{non fare domande adesso} su
questa regola: ogni anno arrivano domande e commenti, e ci saranno lezioni
apposite di tre ore, con strumenti per capire il perché. Ma da subito gli
studenti devono abituarsi a evitare questi termini, sia perché si spera che
abbiano termini più precisi, sia perché queste parole ``portano con sé dei
problemi agganciati'' che si vedranno. ``Non spaventatevi, non verrete sparati in
un vicolo cieco di notte''.

\paragraph{Una piccola gamification.} Il divieto va visto come uno
\textbf{stimolo}, ``un po' come se stessimo gamificando la vostra esperienza'',
una sorta di user experience delle lezioni. Imponendosi di non usare queste
parole, si noterà come migliora la capacità di \textbf{descrivere i giochi} e le
proprie esperienze ludiche, e di essere più precisi. È un invito a ``obbligarsi a
essere più precisi e ad avere una mentalità più critica''. La lezione di
Ruggero sull'\textbf{involvement}, promette Daniele, sarà ``fenomenale'' e
``vi aprirà la testa''.

\paragraph{Approfondimento: anche i libri faticano a definire il ``fun''.} La
difficoltà che il divieto vuole far emergere è visibile anche nei testi di game
design. Schell, in \emph{The Art of Game Design}, prova a definire il \emph{fun}
e propone ``piacere con sorprese'' (\emph{pleasure with surprises}), ammettendo
subito che è strano poter usare una parola per tutta la vita, sapendo con
certezza che cosa significa, senza riuscire a esprimerla chiaramente quando
qualcuno lo chiede; e invita a mettere alla prova la definizione con
controesempi. La sua lente sul \emph{fun} (\emph{The Lens of Fun}) si limita,
non a caso, a chiedere quali parti del gioco siano \emph{fun} e \emph{perché}.
Anche \emph{immersion} ha molte letture: in \emph{21st Century Game Design},
Bateman e Boon la descrivono, tra l'altro, come la capacità di ``perdersi''
nel mondo artificiale presentato dal gioco, legata alla sospensione
dell'incredulità e ai giochi di mimesi. Sono significati diversi sotto la stessa
parola: proprio il motivo per cui, in un documento di design, conviene dire
quale si intende.

% ---------------------------------------------------------------------
\section{Chi crea, chi gioca, chi facilita}
\label{a1:sec:ruoli-esperienza}

Un'altra distinzione importante è quella tra \textbf{chi crea}, \textbf{chi
gioca} e \textbf{chi facilita}. Spesso, quando si fa un gioco, lo si fa provare
di persona a chi gioca. Ma ``la norma dei giochi pubblicati non funziona così'':
quando si pubblica un gioco non si vede chi lo compra e chi lo gioca, e non si
sa se qualcuno, a casa propria con i propri amici, segue le regole alla lettera e
interpreta correttamente quello che il designer aveva descritto. Quindi già in
fase di progettazione bisogna tenere a mente \textbf{chi sta facendo game design
e chi sta giocando}.

E nei giochi a evento, nei giochi pervasivi, nei giochi di ruolo da tavolo e dal
vivo (LARP) c'è un terzo pezzo del puzzle, ``importantissimo'': \textbf{chi
facilita}. La domanda da farsi è: \textbf{il regolamento è scritto per chi gioca
o per chi deve facilitare il gioco?} Bisogna abituarsi a pensare a tutte le
persone coinvolte nell'esperienza.

% ---------------------------------------------------------------------
\section{Playtest e iterazione}
\label{a1:sec:playtest}

\subsection{Nessun gioco nasce perfetto}

L'ultima voce del glossario è ``un'altra cosa importantissima'':
\textbf{playtest} e \textbf{iterazione}. Il corso è \textbf{pragmatico}:
``specialmente da parte mia'', dice Ruggero, non si vedranno teorie ``messe lì
per fare accademia sul game design nelle nostre torri d'avorio''. Ci saranno
teorie e framework, ma tutti estremamente pragmatici, perché l'obiettivo non è
solo capire i giochi, ma \textbf{saperli fare e iterare}.

Playtest e iterazione sono \textbf{inevitabili} nel game design, perché
\textbf{nessun gioco nasce perfetto}, soprattutto perché le persone sono
diverse: un gioco perfetto per una persona può essere incomprensibile per
un'altra. Chi progetta giochi per altre persone deve capire, dopo averli fatti
provare, \textbf{che cosa cambiare, che cosa sistemare, che cosa aggiungere, che
cosa togliere}. Ci saranno lezioni su playtest, prototipazione, raccolta dei
feedback, metodi etnografici. Ma bisogna abituarsi già ora all'idea che
\textbf{non esiste la consegna definitiva}. A un certo punto i task vanno
consegnati per forza di cose, ma, come diceva Daniele all'inizio, alla fine del
corso gli studenti avranno un \textbf{prototipo}, ``che sarà un prototipo'', e
gli strumenti per portare avanti il lavoro anche dopo.

\subsection{L'aiuto dopo il corso}

Non solo gli strumenti: anche l'aiuto dei docenti. Daniele racconta di aver
portato a un editore, a Play\footnote{Probabilmente \emph{Play -- Festival del
Gioco} di Modena; la trascrizione non permette di esserne certi.}, uno dei
progetti realizzati l'anno precedente da due studenti, un gioco di ruolo in
solitario (``un \emph{solo game} RPG molto, molto bello''). Se gli studenti
vorranno trasformare il proprio lavoro in un prototipo vero e proprio, possono
contattarlo: lo fa ``per lavoro, ma anche per passione''. Gestisce anche una
community per aiutare i designer ``un po' depressi'' a rimettersi a progettare.

Quindi: ``non sentitevi persi se durante la fase di iterazione vi sentite
persi, è normale''.

\subsection{Il dizionario comune e il production dictionary}

Daniele riprende un punto detto ``benissimo'' da Ruggero e che secondo lui non è
stato sottolineato abbastanza: il \textbf{dizionario comune} serve anche a
\textbf{non perdersi} mentre si itera. Se da un'iterazione all'altra i
documenti usano termini completamente diversi, e non c'è un dizionario di base,
ci si perde.

Nel videogioco, una delle primissime cose che un designer fa quando entra in
un'azienda è il \textbf{production dictionary}: si annota, in un file privato,
tutte le parole che usano i colleghi, per allinearsi al dizionario comune
dell'azienda, perché lì su un prototipo possono lavorare anche cento persone.
Nel gioco da tavolo non si fa molto, perché il mondo dell'analogico è
``talmente più piccolo''; ma serve comunque imparare un \textbf{dizionario
minimo}, dai docenti e dai professionisti. E il corso ne darà una parte,
entrando nel dettaglio di alcuni termini e framework.

\paragraph{L'esempio del feedback.} Perché serva, lo mostra un esempio. Si
testa un gioco e ci si annota come feedback: ``questo gioco deve essere più
divertente''. Che cosa vuol dire? Come lo si migliora? Quella frase non dà
elementi utili. Se invece si scrive che deve essere \textbf{più adrenalinico},
\textbf{più frenetico}, \textbf{più rilassato}, \textbf{più cozy}, o
\textbf{più bilanciato matematicamente}, si sa dove andare a lavorare.

\paragraph{Approfondimento: la regola del ciclo.} Il principio ``nessun gioco
nasce perfetto'' ha un nome preciso in \emph{The Art of Game Design}: è la
\textbf{Rule of the Loop}, la regola del ciclo. ``Più volte provi e migliori il
tuo design, migliore sarà il tuo gioco''. Schell precisa che non è una lente,
cioè un punto di vista, ma ``una verità assoluta'', senza eccezioni: ci saranno
momenti in cui si cercherà di convincersi che ``questa volta il design è così
buono che non serve testarlo'', e ogni volta se ne pagherà il prezzo. Ed è
interessante la ragione per cui la regola pesa di più sul digitale: secondo
Schell il tempo e il denaro necessari per provare e correggere un sistema sono
molto maggiori nei videogiochi che nei giochi tradizionali. Per un gioco di
carte o da tavolo, con il tempo per provarlo e cambiarlo più e più volte, il
consiglio è semplicemente di farlo. Il gioco analogico è, da questo punto di
vista, la palestra ideale dell'iterazione.

% ---------------------------------------------------------------------
\section{Tabletop Simulator}
\label{a1:sec:tts}

\subsection{Il software di settore}

``Last but not least'': il corso usa un software, \textbf{Tabletop Simulator},
che si compra su \textbf{Steam} e costa tra i \textbf{10 e i 20 euro} (``aspettate
che vada in sconto, va in sconto ogni due per tre''). Daniele spiega perché:
\begin{itemize}
  \item è \textbf{il software di settore} del gioco da tavolo, quello che i
        publisher usano per i prodotti digitali, interni ed esterni;
  \item essendo su Steam, permette di fare \textbf{playtest} e di giocare con gli
        amici di Steam. Daniele stesso sta testando lì il proprio gioco;
  \item funziona anche su Mac: Windows o Mac, nessun problema.
\end{itemize}
Ci sarà una lezione apposita, di Ruggero, sulla creazione del proprio
\textbf{modulo} di gioco.

\paragraph{Che cosa non è.} Daniele chiarisce subito un ``fraintendimento
gigantesco'', soprattutto per chi ha già esperienza: Tabletop Simulator
\textbf{non richiede programmazione} e \textbf{non è un software per
trasformare i giochi da tavolo in videogiochi}. È un software che
\textbf{simula dei componenti fisici su un tavolo}: serve per giocare i board
game in digitale. Simula anche la fisica: c'è perfino il comando per
\textbf{ribaltare il tavolo} (``non usatelo ora'').

\paragraph{La partnership mancata.} Da anni Daniele prova a fare una
partnership con gli sviluppatori: due anni prima ha contattato l'azienda per
acquistare delle chiavi per gli studenti, dato che il software si usa nei corsi.
L'azienda ha risposto che non fa accordi commerciali. ``Credo sia la prima
azienda nella mia vita che dice che non vuole che io le dia dei soldi''. Gli
studenti dovranno quindi provare ad acquistarlo da sé; se i docenti riusciranno a
darglielo, bene.

\paragraph{La domanda di Gianmarco.} Gianmarco segue le lezioni da un computer
aziendale ``super chiuso'', su cui non può installare Steam: deve usare Tabletop
Simulator durante le lezioni? Ruggero risponde che per la lezione dedicata è
meglio poter seguire e rifare in diretta. Ma il software non sarà richiesto
necessariamente per il project work: un gruppo potrebbe fare un gioco pervasivo
o un gioco di ruolo, non un gioco da tavolo. Essendo un software di settore, è
comunque bene conoscerlo. Ruggero aggiunge che le versioni non ufficiali che
circolano non permettono di giocare in molti, e quindi non servono per
playtestare o giocare insieme.

\subsection{Una dimostrazione}

Daniele mostra il software aperto da Steam:
\begin{itemize}
  \item si può creare una partita \textbf{single player} o \textbf{multiplayer};
        per giocare tra compagni si crea un server multiplayer e ci si aggiunge
        come amici su Steam;
  \item si dà al server un nome, un numero di giocatori, eventualmente una
        password; creato il server, si apre la scelta dei giochi;
  \item ci sono giochi gratuiti, giochi da acquistare (DLC, non obbligatori) e
        soprattutto un \textbf{Workshop} ``ampissimo'', con decine di migliaia di
        giochi caricati dagli utenti: ``se c'è un board game che conoscete,
        quel board game è su Tabletop Simulator, giocabile dal Workshop''.
        L'esempio mostrato è l'open beta ufficiale del gioco da tavolo di
        \emph{Slay the Spire}\footnote{La trascrizione dice ``Slade Spire''.}:
        molti publisher, per farsi pubblicità, mettono il proprio gioco in demo o
        giocabile in digitale;
  \item c'è un \textbf{tutorial} iniziale che spiega come muoversi sul tavolo;
        con \emph{Esc} e \emph{Help} si trovano i comandi da tastiera.
\end{itemize}

\subsection{Il primo task}

Il primo (e per ora unico) task è questo: \textbf{entro il 10 ottobre}, giorno
della lezione di Ruggero su Tabletop Simulator, giocare \textbf{tre o quattro
volte in totale}, online, tra compagni di corso o con amici. Non si tratta di
imparare a creare moduli, ma di prendere la mano \textbf{lato giocatore}: aprire
il software, impararne i comandi, giocare. ``Per me potete anche andare a
prendere domino o mahjong'' (``mahjong è bellissimo tra l'altro'', aggiunge
Daniele, che sta facendo un videogioco sul mahjong), o il pachisi o il go. Gli
anni precedenti si chiedeva a ciascuno di proporre ai compagni un gioco da
tavolo che apprezzava. È anche un modo per conoscersi. Ruggero spiega il
perché: usarlo per giocare è facile e intuitivo, usarlo per creare moduli è più
complesso; ma chi arriva al 10 ottobre avendo già aperto il software e giocato,
troverà più intuitivo anche creare un modulo. E l'avvertimento: ``non riducetevi
all'ultimo, non fate una partita il 9 sera''.

\subsection{Tabletopia e altri vantaggi}

Tancredi racconta di aver provato a caricare un suo gioco su
\textbf{Tabletopia}. ``Questo è un vaso di Pandora'', risponde Daniele.
Tabletopia è molto usato, ma secondo lui il suo strumento di creazione dei moduli
è pessimo, ed è il motivo per cui la maggioranza dei moduli ufficiali dei
publisher è su Tabletop Simulator. Tabletopia ha un vantaggio vero, è
\textbf{gratuito}, ma ha scelte di ``qualità della vita'' poco felici.

Tabletop Simulator ha invece un vantaggio ``fenomenale'' che si tende a
sottovalutare: essendo su Steam, \textbf{il modulo si pubblica direttamente sul
Workshop di Steam}; si possono invitare gli amici a iscriversi e giocare con
loro, purché abbiano un account Steam. Rispetto ad altri software da imparare,
costa ``tra i 10 e i 20 euro, ragazzi: si parla proprio di spiccioli''. ``Io con
Tabletop Simulator ci ho playtestato la vita''.

Altre caratteristiche utili:
\begin{itemize}
  \item si possono creare \textbf{moduli non elencati} (\emph{unlisted}),
        perfetti da passare a un publisher prima che il gioco sia pubblicato:
        non sono pubblici, e nessuno può rubare i materiali;
  \item il software ha una grafica datata (``\emph{World Art 1998}'', scherza
        uno dei docenti), ma lo stanno portando su \textbf{Unity} con molti investimenti
        e aggiornamenti; ascolta la comunità e introduce funzioni utili anche per
        chi fa design;
  \item il suo server Discord è molto frequentato, e c'è un canale YouTube
        dedicato solo al design di giochi di carte\footnote{La trascrizione ne
        riporta il nome come ``Distraction Maker''; non è verificabile.} con un
        server Discord dove si fanno playtest: il 99\% è su Tabletop Simulator;
  \item molti editori, quando si propone loro un gioco da tavolo,
        \textbf{chiedono se c'è un modulo per Tabletop Simulator}.
\end{itemize}

% ---------------------------------------------------------------------
\section{Le domande a Daniele: il tema del project work}
\label{a1:sec:domande-daniele}

Daniele ha ``un minuto'' per rispondere alle domande sul corso prima di dover
andare. Uno studente chiede, ``giusto per curiosità'', se il project work avrà
un tema o sarà libero.

\begin{itemize}
  \item C'è un \textbf{mandate}, che i docenti propongono ogni anno e
        aggiornano leggermente; verrà presentato il \textbf{25 settembre}.
        ``Non facciamo spoiler''. Conterrà quello che si chiede di fare e quello
        che non si chiede: dei \textbf{must}, dei \textbf{must not} e dei
        \textbf{nice to have}.
  \item Verrà chiesto di lavorare su un'\textbf{idea originale} che non si è
        portata al corso: originale nel senso di \textbf{creata apposta}, non
        nel senso di ``mai vista prima''.
  \item L'idea va \textbf{partorita insieme}. Altri anni alcuni studenti hanno
        provato a imporre la propria idea al gruppo, ``e l'esperienza è stata
        pessima per tutti''. Quindi: \textbf{non imponete le vostre idee, fate
        brainstorming con gli altri}. Daniele, che fa ``molto spesso
        brainstorming'', darà strumenti per farlo: nessuno arriverà impreparato
        a quella lezione.
\end{itemize}

Ruggero aggiunge che il mandate è \textbf{abbastanza flessibile} perché tutti
possano fare qualcosa che li interessi: ``l'ultima cosa che vogliamo è forzarvi a
lavorare su qualcosa che non vi piace''. Bisogna lavorare su ciò che piace e per
cui si ha uno stimolo personale, anche nella \textbf{tipologia} di gioco: nessuno
sarà obbligato a fare un gioco di ruolo o un board game, e la lezione sui casi
studio mostrerà quante forme diverse di gioco analogico il corso accetta.

L'unica cosa da \textbf{non fare}, ``perché è contraddittorio per questo corso'',
sono i \textbf{giochi grandi}: ``lo chiediamo per voi''. Niente progetti
mastodontici\footnote{La trascrizione dice ``non fate delle cose da gnome'':
probabilmente il nome di un gioco da tavolo molto grande e complesso, non
ricostruibile con certezza.}, niente ``cose alla \emph{Dungeons \& Dragons}''.
Si può fare un gioco di ruolo, ma con un'esperienza ``meno corposa'': ci sono
tanti giochi di ruolo da 15 o 20 pagine, ``bellissimi, fichissimi'', che si
possono leggere online. Il consiglio è farsi una cultura anche su
\textbf{itch.io}, ``uno dei portali più importanti del gioco di ruolo
internazionale'', di cui si parlerà anche nella parte di ricerca e sviluppo.

Poi Daniele saluta (``noi ci vediamo venerdì a lezione'') e lascia la
conduzione a Ruggero.

% ---------------------------------------------------------------------
\section{Le domande finali con Ruggero}
\label{a1:sec:domande}

\subsection{Le regole delle domande}

Ruggero spiega che cosa vuol dire ``stare con lui''. Alla fine di ogni sua
lezione ci saranno \textbf{almeno 20--30 minuti di domande} libere, \emph{free
form}. Vuole che fare domande sia \textbf{normalizzato}: ``io gestisco malissimo
l'\emph{horror vacui} e divento super cringe molto alla svelta se non ci sono
domande in giro''. Sono benvenute le domande sulla lezione appena fatta, ma non
solo: il corso è crossmediale, e anche le domande che vengono da
\textbf{divagazioni}, da parti laterali, ``fatele''. Quello che sembra scontato o
banale a uno studente può non esserlo per il docente, e viceversa; e può essere
utile anche agli altri. ``Se vi sembrano stupide, fatele il doppio''. ``Magari
sembro incazzato, ma vi giuro sono rilassatissimo. Non mordo, magari abbaio un
po'\,''.

\subsection{Come approcciarsi a Steam e Tabletop Simulator}

Rita non ha mai usato né Steam né Tabletop Simulator e chiede come
approcciarsi, e se c'è modo di mettersi d'accordo tra compagni su Discord, magari
con un sondaggio sui giochi, per formare gruppi di gioco. Ruggero conferma: sul
Discord dell'accademia c'è un \textbf{canale apposito per il corso}, ed è degli
studenti. Lo possono usare per organizzarsi a giocare, per chiedere a un
compagno di prestare un gioco, per domande tecniche (come si scarica, come si
installa). ``Sarebbe proprio ideale che vi parlaste e vi organizzaste''.

Anche perché, aggiunge con un ``piccolo spoiler'' (``nessuno vi ha detto
niente''), se fossero i docenti a organizzare i gruppi di gioco, il task lo
starebbero facendo loro. Bisogna \textbf{abituarsi a organizzare le sessioni di
gioco}, perché arriveranno i task sul \textbf{playtest} e su come organizzarlo,
e a quel punto si sarà ``meno disperati''. Sembra scontato, ma saper
\textbf{organizzarsi per giocare in autonomia} ``vi sarà utilissimo come skill in
questo settore, ve lo giuro''.

\subsection{Il bilanciamento}

Angelo confessa ``una paranoia grossissima'': ci saranno lezioni sul
\textbf{bilanciamento}? ``Ma certo''. Sono proprio le lezioni di Marta sul
development: il bilanciamento è di solito un tema che tocca \textbf{i developer
più dei designer}. L'anno precedente c'era una sola lezione con Marta, e i
docenti si sono resi conto che serviva andare ancora più a fondo: per questo
quest'anno le lezioni sono due. ``Non è una paranoia, è una legittima domanda'':
Ruggero insegna da anni, vede altri corsi, parla con altri docenti e studenti, e
``mi incazzo spesso'' perché molti corsi di game design \textbf{non fanno
bilanciamento}. I docenti volevano invece un corso che coprisse tutte le parti
del game design: non solo l'idea del gioco, ma il bilanciamento, la produzione,
i playtest, l'iterazione, la scrittura dei regolamenti.

\subsection{Si può giocare con i docenti?}

Una studentessa chiede se si può giocare con i docenti. Per motivi etici,
durante il corso questo è limitato alle parti di \textbf{valutazione} dei giochi
e di \textbf{tutoraggio}. Dopo il corso, ``nessun problema, almeno con me''.

\subsection{Dove si trovano i giochi su Tabletop Simulator}

Ruggero mostra dove si trovano i giochi: nella propria libreria di Steam, nella
pagina di Tabletop Simulator, alla voce \textbf{Workshop}, dove ci sono ``più o
meno 36.000'' giochi. Una studentessa scopre intanto che l'app di Steam per
telefono è solo una \emph{companion app}: per registrarsi e giocare serve il
PC.

Per provare, chiede agli studenti un titolo:
\begin{itemize}
  \item \emph{Keep the Heroes Out}: non compare, ``non tutti i giochi ci sono,
        perché molti publisher non li mettono''; forse, essendoci già una
        versione del gioco su Steam, il modulo è stato tolto;
  \item \emph{Magic: The Gathering}: ``almeno 50 moduli'', anzi ``123 pagine'';
  \item \emph{Bang!}: tre moduli.
\end{itemize}
Bisogna fare attenzione: alcuni moduli sono in lingue specifiche (``magari il
primo è in russo''), alcuni hanno già l'espansione, altri sono solo il gioco
base. Conviene sfogliare e non fermarsi al primo risultato. E nel Workshop si
trovano anche cose molto particolari: giochi fuori produzione, giochi vecchi,
giochi ancora in prototipo.

\subsection{I ruoli nei team e le dinamiche tra le persone}

Uno studente torna sulla slide delle figure: nelle realtà piccole molti ruoli
hanno ``il doppio o triplo cappello''. Quali ruoli restano distinti, e come si
creano le sinergie tra persone con obiettivi diversi?

Ruggero premette che l'ideale ``platonico'' è un ambiente di lavoro costruttivo,
basato sulla fiducia e ``il più \emph{safe} possibile''; ma nell'industria non
sempre è così, e le strutture variano moltissimo: alcuni publisher sono molto
verticali, altri hanno ruoli già solidificati (se un publisher ha già i suoi
developer e non sta crescendo, entrarci come developer è praticamente
impossibile). ``Se prendiamo due aziende in Italia, vediamo due strutture
interne completamente diverse''. Non esiste una scorciatoia per farsi assumere
da un editore.

Poi prende ``l'esempio più comune di tutti'', la media dei giochi da tavolo
pubblicati (figura~\ref{a1:fig:relazioni}):
\begin{itemize}
  \item il \textbf{game designer} è una figura \textbf{esterna}, con un
        \textbf{contratto a royalty} con il publisher;
  \item la sua figura di riferimento è il \textbf{publisher} stesso, e la
        conversazione designer--publisher è ``molto calda'';
  \item se il publisher ha dei \textbf{developer}, il developer entra in questa
        conversazione: designer e developer collaborano molto, e internamente
        developer e publisher collaborano molto. Spesso, nei publisher piccoli,
        è il publisher a fare le veci del developer;
  \item \textbf{graphic design} e \textbf{illustrazione} tendono a parlare molto
        con il publisher e \textbf{molto poco con il designer}. A meno che non si
        tratti di uno studio indie molto piccolo, è difficile che il designer
        sappia anche solo chi è l'illustratore prima che il gioco esca.
\end{itemize}

\begin{figure}[htbp]
\centering
\begin{tikzpicture}[
    font=\small,
    ruolo/.style={draw, thick, rounded corners=3pt, fill=gray!12,
                  minimum width=2.8cm, minimum height=0.9cm, align=center},
    esterno/.style={draw, thick, dashed, rounded corners=3pt,
                    minimum width=2.8cm, minimum height=0.9cm, align=center},
    calda/.style={{Stealth[length=2.5mm]}-{Stealth[length=2.5mm]}, very thick},
    tiepida/.style={{Stealth[length=2mm]}-{Stealth[length=2mm]}, thin, dashed}
  ]
  \node[esterno] (gd) at (0,0) {Game designer\\ \footnotesize esterno, a royalty};
  \node[ruolo] (dev) at (5,1.4) {Developer\\ \footnotesize se c'è};
  \node[ruolo] (pub) at (5,-1.4) {Publisher};
  \node[ruolo] (art) at (10,-1.4) {Graphic designer\\ e illustratore};
  \draw[calda] (gd) -- (dev);
  \draw[calda] (gd) -- (pub);
  \draw[calda] (dev) -- (pub);
  \draw[calda] (pub) -- (art);
  \draw[tiepida] (gd.south) to[bend right=25]
     node[below, font=\footnotesize\itshape] {quasi nessun contatto} (art.south);
  \node[draw, dotted, rounded corners=4pt, fit=(dev)(pub)(art), inner sep=8pt,
        label={[font=\footnotesize\itshape]above right:interno al publisher (di solito)}] {};
\end{tikzpicture}
\caption{Le relazioni tipiche intorno a un gioco da tavolo pubblicato, secondo la
risposta di Ruggero. Le frecce spesse indicano una collaborazione stretta e
frequente; quella tratteggiata, un contatto quasi assente. Nei publisher piccoli
il publisher fa spesso anche da developer.}
\label{a1:fig:relazioni}
\end{figure}

\subsection{Fare da developer di sé stessi?}

Rita chiede in chat se un designer può fare anche da developer del proprio
gioco, o se Ruggero lo sconsiglia. In generale \textbf{lo sconsiglia}, per lo
stesso motivo per cui si sconsiglia in altri media: i \textbf{darling}. Il
designer ha delle meccaniche a cui è affezionato, e il developer a volte dovrà
\textbf{tagliarle}; è molto difficile, se si è ``gestori di sé stessi'', tagliare
cose che si vogliono nel gioco. ``È un po' come lasciare fare il lavoro di editing
allo scrittore che ha già scritto il libro'': può rivedersi, certo, ma difficilmente
cambierà cose che farebbe bene cambiare. Nelle lezioni con Marta si vedrà perché,
e alla fine del corso ``non avrete dubbi nemmeno voi'' nel separare designer e
developer.

Ma tutto dipende dallo \textbf{scopo del progetto}: se si fa un gioco in tre,
qualcuno quel lavoro deve farlo per forza, e magari il designer fa anche il
developer, compensando con \textbf{più playtest}. ``Il mondo attuale e il mondo
ideale sono due cose molto diverse, e lo vediamo non solo dai giochi ma dalla
politica'', ma ``non entriamo in questo vaso di Pandora''.

\subsection{Il ruolo indispensabile e la pubblicazione}

Lo studente di prima ne trae una conclusione: senza game designer il gioco non
esiste. ``Quello sicuramente''. Ma Ruggero coglie l'occasione per approfondire
un punto che forse non era passato: \textbf{fare design di un gioco per
pubblicarlo non è l'unica opzione possibile}. Un gioco non è ``meno gioco''
se non è pubblicato da un editore. È legittimo mettere un gioco su itch.io in
\textbf{print and play}, fare un gioco a evento senza venderlo, fare un
\emph{murder party}, una escape room, un board game per una persona amica:
``quell'esperienza lì non è meno gioco''. Dipende dall'obiettivo del progetto.

Ruggero ha lavorato molto anche con le \textbf{scuole}, dove il publisher non
esiste: è un gioco per le classi, e la figura più vicina è l'editoria scolastica,
``che però non fa il publisher''. Ci sono compromessi diversi sulla materialità
e su come il gioco si presenta, che si vedranno nella prototipazione e nei casi
studio. Quindi l'unica figura che deve esserci sempre è \textbf{chi fa game
design}, anche se magari non sa di chiamarsi game designer e non ha mai
sentito parlare di game design: ``questo non la rende meno game designer''.

\paragraph{Un libro consigliato.} Ruggero fa anche mentoring e gestisce gruppi
di settore, e spesso gli chiedono che cosa leggere per imparare il game design.
Non è il tipo che obbliga a leggere libri (``non vi obbligherò a leggere
niente''), ma ce n'è uno che consiglia sempre, ed è il primo che ha scelto di
leggere: \textbf{\emph{Rise of the Videogame Zinesters}} di \textbf{Anna
Anthropy}\footnote{La trascrizione dice ``Anna Antrofi'' e ``Rise of the
Videogames in Esters''.}. È un libro un po' datato, ma ancora reperibile; parla
soprattutto di videogiochi, ma ``per l'80\%'' il contenuto si applica al gioco e
al game design in generale. È un libro sulla \textbf{democratizzazione del gioco}:
il gioco come strumento per parlare di sé e come \textbf{medium legittimo}. Come
esistono film drammatici, tristi, intensi, possono esistere giochi su tematiche
delicate, e se ne vedranno nei casi studio. Molti giochi fatti con questa
direttiva non hanno un publisher, che tende a pubblicare ciò che dà un ritorno
commerciale.

\paragraph{Commerciale e personale.} ``Io spero che voi non farete giochi solo
per scopo commerciale''. Certo, servono un tetto sulla testa, il cibo in
tavola, un posto in cui dormire, e così si finisce per fare giochi che non sono
il proprio gioco dei sogni. Spesso è un mix: Ruggero ha fatto giochi commerciali
``per pagarmi l'affitto'' e giochi che voleva fare e su cui ha perso soldi,
``sapendolo proprio che sarebbero stati un fallimento''. I giochi, come i libri e
le canzoni, si fanno per molti motivi, e quelli commerciali di successo non sono
più importanti degli altri. Tra i giochi della sua slide c'erano il gioco da
tavolo ufficiale di \emph{Cyberpunk}\footnote{La trascrizione dice ``Cyberpunk
del 1977'': quasi certamente un gioco legato a \emph{Cyberpunk 2077}.} e
\emph{Marvel United}, su IP grosse e famose, ma anche giochi fatti di ``una
candela e una busta su un tavolo''. E se gli si chiede quale preferisce in
termini di contenuto, non sono quelli sui supereroi.

\subsection{Il gioco di ruolo in Italia}

Nicola, che è interessato soprattutto al gioco di ruolo e conosce già persone
del settore, ha ``una leggera paura'': che chi studia game design finisca per
fare il \textbf{freelance} a vita, ``perché non ci sarà mai un'occupazione per
bene'', o il \textbf{copywriter}. Conosce un ragazzo che ha lavorato in quel
mondo e che ora fa altro.

Ruggero risponde che la maggioranza dei giochi nel suo portfolio sono giochi di
ruolo, quindi il tema lo tocca da vicino, e che in Italia quello che dice Nicola
``è semplicemente vero''. Molti publisher piccoli hanno una persona di
riferimento che copre più ruoli: impagina, fa i playtest, fa un po' di
development. È ``un settore economicamente completamente disastrato'', per
motivi culturali che non apre: oggi in Italia è praticamente impossibile
\textbf{vivere di giochi di ruolo}. Chi ci lavora lo fa nel tempo libero, al di
fuori del lavoro principale, oppure coprendo molte altre cose: molti publisher
italiani di giochi di ruolo pubblicano anche titoli esteri e localizzano, non
fanno solo cose originali. ``È un grosso calderone''. Poi si ferma: ``questo è
un \emph{rant}'', e ``i \emph{rant} tornano sempre'', ma non vuole che gli
studenti escano dalla prima lezione ``con la disperazione nel cuore''. Ci sarà
spazio per parlarne quando avranno gli strumenti per capirlo, ``per non farvi
abbattere''.

\subsection{Che cosa bisogna sapere già}

Uno studente chiede se ci sono nozioni che dovrebbero già avere e che possono
recuperare: per esempio, sente parlare di giochi ``German'' o ``American'' e non
sa che cosa siano. Ruggero spiega l'organizzazione del corso:
\begin{itemize}
  \item le \textbf{prime due lezioni} sono di fondamenta. Questa era ``dall'alto'':
        come è strutturato il corso, le tempistiche;
  \item la \textbf{prossima lezione} è di tre ore di soli \textbf{casi studio},
        ``a manetta, a raffica'', divisi per categorie e tipologie, con
        ragionamenti su categorie e generi di gioco: lì si vedranno German,
        American ``e anche altre più utili'';
  \item nel resto del corso le terminologie specifiche sono inserite nelle
        lezioni dedicate: il \textbf{paragame} nelle lezioni sul paragame, il
        \textbf{modello matematico} nelle lezioni di Marta sul development.
\end{itemize}
Per ora agli studenti si chiede di sapere quello che si è visto oggi:
meccaniche, dinamiche ed estetiche; le figure professionali, designer e
developer; e soprattutto \textbf{game e play}.

\subsection{Game e play: i quattro quadranti}

``Specialmente \emph{game} e \emph{play}'', insiste Ruggero: è più difficile di
quello che sembra, perché l'italiano non abitua a separare le due cose, ed è
possibile averle usate per anni come sinonimi. Nella sua esperienza con chi si
avvicina al design è ``uno dei punti più critici''. Spesso i commenti che si
fanno sui giochi giocati sono commenti sul \emph{play} applicati al
\emph{game}: ``il gioco non mi è piaciuto perché, giocando con Tizio, è successo
questo''; ``il gioco mi sembra sbilanciato perché Tizio al tavolo ha fatto
questo''.

Separare \emph{game} e \emph{play} serve a capire quali giochi sono progettati
bene e quali male, e soprattutto a riconoscere i \textbf{quattro possibili
incroci} (figura~\ref{a1:fig:quadranti}): un gioco progettato bene che piace
giocare; uno progettato male che non piace; uno progettato male che piace; uno
progettato bene che non piace. ``Non è importantissimo in generale nella vita,
ci sono cose più importanti'', ma per chi fa game design lo è: ``a volte vi
chiederanno di lavorare su un gioco che non vi piace giocare, e quel gioco deve
essere fatto bene lo stesso''.

\begin{figure}[htbp]
\centering
\begin{tikzpicture}[
    font=\small,
    cella/.style={draw, thick, minimum width=4.4cm, minimum height=1.9cm,
                  text width=4.0cm, align=center}
  ]
  \node[cella, fill=gray!20] (a) at (0,0)     {\textbf{progettato bene}\\ e \textbf{mi piace} giocarlo};
  \node[cella]               (b) at (4.4,0)   {\textbf{progettato bene}\\ ma \textbf{non mi piace} giocarlo};
  \node[cella]               (c) at (0,-1.9)  {\textbf{progettato male}\\ ma \textbf{mi piace} giocarlo};
  \node[cella, fill=gray!8]  (d) at (4.4,-1.9){\textbf{progettato male}\\ e \textbf{non mi piace} giocarlo};
  \node[font=\footnotesize\itshape, above=0.15cm of a] {mi piace (\emph{play})};
  \node[font=\footnotesize\itshape, above=0.15cm of b] {non mi piace (\emph{play})};
  \node[font=\footnotesize\itshape, rotate=90, left=0.15cm of a, anchor=south]
       {bene (\emph{game})};
  \node[font=\footnotesize\itshape, rotate=90, left=0.15cm of c, anchor=south]
       {male (\emph{game})};
\end{tikzpicture}
\caption{I quattro incroci tra la qualità del \emph{game} (come è progettato) e
l'esperienza di \emph{play} (se piace giocarlo), descritti da Ruggero. Per un
designer i quadranti ``scomodi'' sono i due fuori diagonale: un gioco può
piacere pur essendo progettato male, e va saputo progettare bene anche un gioco
che non piace.}
\label{a1:fig:quadranti}
\end{figure}

\paragraph{La domanda di Tancredi.} Tancredi chiede conferma: il \emph{play}
non è il ``giocare senza regole'', ma può essere anche il giocare a un gioco? Sì:
il \emph{play} è l'\textbf{atto}, e può essere fatto con un \emph{game} o senza.
E un \emph{game} può esistere senza che nessuno lo giochi: gli scacchi non sono
meno scacchi perché in quel momento nessuno ci gioca. E se, giocando, una
persona applica le regole in modo diverso? Allora \textbf{non si sta giocando
quel gioco}.

\subsection{House rule: il caso di Monopoly}

Qui Ruggero fa un po' di \emph{foreshadowing} su una lezione di Daniele: il
discorso degli \textbf{hack}, delle \textbf{house rule}, delle \textbf{mod},
cioè di come modificare un gioco esistente. Il punto è che \textbf{un gioco
modificato da una house rule è un altro gioco}: cambiando le regole cambiano
necessariamente anche le \textbf{dinamiche} che ne derivano, e quindi
l'esperienza. Si immagini \emph{Monopoly} in cui ogni turno si possa, volendo,
prendere gratis dei soldi dalla banca: tutto il resto è identico, eppure il gioco
cambia.

E questo è talmente interiorizzato che ci sono regole che si applicano ``senza
neanche rendercene conto''. Ruggero chiede quanto duri in media una partita a
\emph{Monopoly}. Gianmarco: ``quattro ore''; altri: ``un pomeriggio''; un altro:
``finché uno dei giocatori non aggiunge la regola: ora vince chi\ldots''. In
realtà, \textbf{seguendo le regole scritte}, una partita non dura mai più di
\textbf{un'ora e venti, un'ora e mezza}. La domanda che fa sempre, e la risposta
è sempre la stessa: perché si pensa che \emph{Monopoly} duri tantissimo? Qual è
la regola che nessuno rispetta?

Uno studente indovina: \textbf{le aste}. Secondo il regolamento, se chi finisce
su un terreno non lo compra, il terreno \textbf{va all'asta} tra i giocatori, e
chiunque può aggiudicarselo. Molti invece giocano come se finire su un terreno
desse solo la \emph{facoltà} di comprarlo o di metterlo all'asta\footnote{Più
precisamente, nel regolamento ufficiale, se il giocatore che finisce su una
proprietà libera decide di non comprarla, la banca la mette all'asta e tutti i
giocatori possono fare offerte, compreso chi l'ha rifiutata. L'asta quindi non è
una scelta del giocatore, ma la conseguenza obbligata del suo rifiuto.}. Con la
regola applicata, di aste se ne fanno moltissime: all'ottavo turno ci sono già
otto terreni assegnati, al ventesimo venti, la gente va in bancarotta
velocissimamente e la partita dura un'ora e mezza. Anche le partite competitive
di \emph{Monopoly}, nota Ruggero, durano un'ora e venti, un'ora e mezza.

Uno studente obietta che giocava con le aste e durava comunque quattro ore.
``Avete dei nemici con l'analisi paralisi'', ride Ruggero; e cita un'altra house
rule tipica dell'infanzia: genitori che danno soldi per non far andare in
bancarotta. Ma \textbf{la bancarotta è parte del gioco}: se la si toglie,
\emph{Monopoly} non regge. Sul tema che \emph{Monopoly} vuole trattare, e sul
suo antenato (\emph{The Landlord's Game}), ``non apriamo adesso un altro vaso di
Pandora''.

\paragraph{Approfondimento: Monopoly, la durata e le house rule nei libri.}
L'esempio di Ruggero si ritrova quasi identico in \emph{The Art of Game Design}.
Parlando del bilanciamento tra partite brevi e lunghe (\emph{Short vs. Long}),
Schell scrive che \emph{Monopoly}, giocato secondo le regole ufficiali,
``finisce spesso in circa novanta minuti'', ma molti giocatori trovano quelle
regole troppo dure e le modificano, distribuendo premi in denaro e allentando i
vincoli sull'acquisto delle proprietà: con l'effetto collaterale di far durare la
partita molto di più, ``tipicamente tre ore, o anche di più''. Schell classifica
le \textbf{house rule} come regole create dai giocatori in risposta a un
difetto percepito dopo qualche partita. Anche \emph{Challenges for Game
Designers}, nel capitolo sull'aggiungere e togliere meccaniche, parte dallo
stesso gioco: il \emph{Parcheggio gratuito}, che per le regole originali non fa
succedere nulla, mentre quasi tutti giocano con una house rule che dà un bonus a
chi ci finisce sopra. Le due fonti confermano la lezione: una piccola regola
cambiata cambia le dinamiche, e le dinamiche cambiano l'esperienza.

\subsection{Il fine del game è il play?}

Un altro studente commenta: va bene la distinzione, ma il fine del \emph{game} è
il \emph{play}; quello che si vede è sempre l'effetto. Ruggero è d'accordo: di
solito, se si fa un gioco, lo si fa per farlo giocare, come si scrive un libro
per farlo leggere (anche se esistono, ed ``è già stato esplorato'', libri scritti
per non essere letti e giochi per zero giocatori). Ma la distinzione serve a una
cosa precisa: \textbf{separare, nel playtest, i feedback sul gioco dai feedback
sul gruppo di gioco} e sulle sue dinamiche. Se una persona esce frustratissima da
un playtest, non è necessariamente colpa del gioco; se una persona dice ``questo
è il gioco migliore del mondo'', potrebbe non essere merito del gioco ma del
gruppo. Capire a che cosa si riferisce un feedback serve a capire \textbf{dove
iterare} il prototipo.

Lo studente osserva che l'area è sfumata: alcuni giochi si basano proprio
sull'avere le persone giuste al tavolo. Un gioco di \textbf{negoziazione} con
giocatori abituati ai giochi ``multi-solitario'' diventa ingiocabile; e nei giochi
a \textbf{economia chiusa} (cita \emph{Container}\footnote{La trascrizione cita
anche un secondo titolo non ricostruibile (``Tokyo su Vigimarket'').}) un gruppo
economicamente conservativo, che spende poco e non fa girare l'economia, fa
durare la partita il doppio. Ruggero rilancia: e se il gioco dura il doppio,
questa dinamica cambia la \textbf{percezione di peso} del gioco, le energie spese,
quanto si è stanchi a fine partita. ``Tutto è collegato. \textbf{I giochi sono
ecosistemi di meccaniche}, e cambiandone una, le altre cambiano di
conseguenza''.

\paragraph{La parola vietata, alla prova.} Lo studente aggiunge che ci sono giochi
che ha trovato ``molto sbagliati, ma che mi divertono''. ``Attenzione: non si può
dire \emph{divertimento}!'', lo ferma Ruggero; nelle presentazioni aveva lasciato
correre, ``ma lì non vi avevo ancora letto la regola''. Lo studente ci prova:
``giochi che poi voglio \textbf{rigiocare}'', anche se ``sono evidentemente
sbagliati, rotti''. Ruggero apprezza, ma non vuole anticipare cose che meritano il
loro tempo: se ne parlerà nelle lezioni sui framework e soprattutto in quella di
\textbf{emotion design}, ``forse la mia lezione preferita del corso''. (Lo
studente ammette che ``forse alla fine si gioca per quello, per l'emozione''.)

\subsection{I giochi ``comunicati male''}

L'ultima domanda corposa riprende un'affermazione di inizio lezione: ci sono
giochi belli ma comunicati male. Qualche esempio? ``Mi basta girarmi per
vederne 48'', risponde Ruggero indicando gli scaffali alle sue spalle. Più che
titoli, elenca le \textbf{ragioni}:
\begin{itemize}
  \item \textbf{l'Italia localizza lentamente}. Basta guardare i premi del gioco
        dell'anno in Italia: spesso vanno a giochi usciti due o tre anni prima in
        Germania, Francia, Stati Uniti o Regno Unito, perché localizzare richiede i
        suoi tempi. Il mercato italiano è quindi molto dipendente dall'estero,
        ``molto impaurito'', e spesso aspetta che un gioco vada bene all'estero
        prima di portarlo in Italia; e alcuni giochi ``perdono il treno'';
  \item \textbf{il publisher originale chiude}: come si pubblica un gioco se non
        c'è più l'azienda con cui fare un contratto? A Ruggero è successo più
        volte;
  \item \textbf{il dipartimento commerciale italiano non sa comunicare il gioco}:
        si fa male un'edizione, la stampa vende pochissimo e non viene più
        ristampata;
  \item \textbf{la scala del mercato}: in Italia un gioco da tavolo di peso medio
        (un \emph{eurogame}, un \emph{German}) che vende 1.500 copie ``è tanto'';
        all'estero se ne vendono magari 10.000. Con 10.000 copie in inglese ci
        sono i soldi per una campagna di comunicazione: un trailer, i
        recensori, i tutorial online. In Italia tutto questo è ``ridotto all'osso'',
        e la cosa si riflette su stampa e divulgazione: anche i canali per
        informarsi su un gioco prima di comprarlo sono pochi (se ne parlerà nella
        lezione su ricerca e sviluppo). Un gioco spesso non arriva al mercato non
        perché sia brutto, ma perché l'editore non lo sa comunicare, ne fa poche
        copie e va fuori produzione;
  \item \textbf{il tempismo}: party game progettati nel 2019 e usciti nel 2020, in
        pieno Covid, quando non si poteva giocare al tavolo, sono falliti tutti, e
        gli editori ci hanno perso. ``Non potevano immaginarselo'';
  \item \textbf{i temi caldi}: dopo l'invasione russa dell'Ucraina e la guerra in
        Palestina, alcuni editori non pubblicano o non localizzano giochi
        strategici e di guerra. Ruggero ha lavorato a \textbf{\emph{Status}}, un
        gioco sulla geopolitica che tocca, tra gli altri temi, l'invasione russa
        dell'Ucraina e il genocidio in Palestina, e alcuni editori che di
        solito localizzano i loro giochi (per esempio il partner di riferimento in
        Russia) non lo hanno localizzato per motivi politici.
\end{itemize}
Lo studente voleva capire che cosa voglia dire comunicare male: non metterci
abbastanza energia, o comunicare una cosa per un'altra, ``come il trailer
sbagliato, che poi vai al cinema e ti trovi un altro film''. ``È un mix'', e la
domanda ``era complessa e meritava più spazio''.

\subsection{I giochi di carte collezionabili e i saluti}

L'ultima domanda (``posso rompere le balle un'ultima volta?'') riguarda i
\textbf{giochi di carte collezionabili}: funzionano con le stesse regole dei
giochi da tavolo? Se ne parlerà alla prossima lezione: è proprio il motivo per
cui la lezione sui casi studio dura tre ore, perché ogni genere di gioco ha le
sue dinamiche e le sue regole, e anche il modo di pubblicare cambia molto tra un
TCG, un gioco di ruolo da tavolo, un LARP e un gioco pervasivo.

Ruggero saluta: le email dei tre docenti sono all'inizio delle slide, la
registrazione arriverà, ``ci si vede questo venerdì''.

% ---------------------------------------------------------------------
\section*{Punti chiave del capitolo}
\addcontentsline{toc}{section}{Punti chiave del capitolo}

\begin{itemize}
  \item Il corso di \textbf{Analog Game Design} (seconda edizione) dura circa tre
        mesi, fino a dicembre, ed è tenuto da tre docenti principali:
        \textbf{Daniele} (narrative design, gioco di ruolo, fiction),
        \textbf{Ruggero} (game design, ricerca e sviluppo) e \textbf{Beatrice}
        (giochi pervasivi), con tre ospiti su RNG, development e game production.
  \item Lo scopo è formare \textbf{autori analogici} consapevoli: strumenti da
        usare anche fuori dal corso, non il ``gioco dei sogni''. Il corso è
        \textbf{intersezionale} (nessuna lezione dedicata a un solo tipo di
        gioco), \textbf{cross-generazionale}, \textbf{pratico e adattabile},
        basato sulle \textbf{domande} (nessuna domanda è stupida), sul
        \textbf{feedback} invece che sul voto e sul \emph{fail forward}.
  \item Il \textbf{project work} si fa in tre gruppi da tre, formati dai docenti;
        pitch il 12 ottobre, sviluppo nelle due pause, tutoraggi il sabato,
        consegna il 4 dicembre, presentazione l'11 e giornata finale con gli
        editori il 14. Si valutano documentazione, presentazione,
        autovalutazione, crescita e atteggiamento.
  \item \textbf{Niente intelligenza artificiale} nel corso: i publisher
        rifiutano pitch e prototipi con immagini generate, i regolamenti scritti
        dall'IA descrivono la stessa meccanica in modi diversi, e soprattutto chi
        la usa perde l'occasione di imparare. Un prototipo può essere fatto a
        matita.
  \item Le \textbf{figure dell'industria}: game designer, game developer,
        editor, publisher, sales e marketing, graphic designer, illustratore,
        sensitivity reader (che non è un censore) e producer. Nelle realtà
        piccole una persona copre più ruoli.
  \item Nell'analogico il \textbf{game designer} progetta l'esperienza creando
        le meccaniche; il \textbf{game developer} fa da \textbf{mediatore} con il
        publisher (bilanciamento, modelli matematici, adattamenti). Pretendere
        che il proprio gioco non venga modificato è un \emph{red flag}.
  \item Il \textbf{marketing} non decide che cosa pubblicare ma come; è
        un'azione, ed esiste anche se lo si ignora.
  \item In \textbf{Italia} il gioco è ancora visto come infantile, il mercato
        punta sulla localizzazione e investe poco su giochi e designer interni;
        all'estero (soprattutto dalla Germania in su) ci sono tradizione adulta,
        premi di prestigio e developer interni.
  \item \textbf{Game} è l'artefatto (le regole, il sistema, l'oggetto);
        \textbf{play} è l'atto di giocare. Si fanno giochi, non giocattoli.
        Separarli serve a distinguere i feedback sul gioco da quelli sul gruppo
        e a riconoscere i quattro incroci tra ``progettato bene/male'' e ``mi
        piace/non mi piace''.
  \item \textbf{MDA}: meccaniche (regole scritte), dinamiche (regole in azione),
        estetiche (esperienza, non grafica); nel \textbf{DPE}: design, play,
        experience.
  \item \textbf{Parole vietate} in ogni materiale scritto: \emph{fun},
        \emph{immersion}, \emph{entertainment} (divertimento, immersione,
        intrattenimento). Bisogna trovare parole più precise.
  \item Distinguere \textbf{chi crea, chi gioca e chi facilita}; nessun gioco
        nasce perfetto: \textbf{playtest e iterazione} sono inevitabili, e un
        \textbf{dizionario comune} aiuta a non perdersi.
  \item \textbf{Tabletop Simulator} è il software di settore per giocare (non
        creare videogiochi) i giochi da tavolo in digitale. Primo task: giocarci
        tre o quattro volte entro il 10 ottobre.
  \item Una \textbf{house rule} crea un altro gioco: \emph{Monopoly} con le aste
        applicate come da regolamento dura circa un'ora e mezza.
\end{itemize}

%
%  (Con hyperref conviene aggiungere anche
%  \renewcommand{\theHchapter}{agd.\arabic{chapter}} nella parte, per
%  evitare ancore duplicate con i capitoli 1, 2, ... di Game Design.)
%  Tutte le etichette di questo capitolo hanno il prefisso "a1:"
%  (a1:cap:..., a1:sec:..., a1:fig:..., a1:tab:...). Il capitolo è
%  autonomo: non contiene \ref verso altri capitoli.
%  Le figure sono ricostruzioni schematiche di quanto i docenti spiegano
%  a voce e con le slide, non copie delle slide.
% =====================================================================

\chapter[Introduzione all'Analog Game Design]{Introduzione all'Analog Game
Design: il corso, le figure dell'industria e un primo linguaggio comune}
\label{a1:cap:intro}

Questa è la \textbf{prima lezione} del corso di \textbf{Analog Game Design},
cioè di progettazione di giochi \emph{analogici}: giochi da tavolo, giochi di
carte, giochi di ruolo da tavolo, giochi di ruolo dal vivo (LARP), giochi
\emph{pervasivi} ed esperienze a evento. È un corso \textbf{breve}, di circa tre
mesi (``da qui a dicembre''), separato dal corso di Game Design, e siamo alla sua
\textbf{seconda edizione}. Come avvertono subito i docenti, è ``una giornata di
introduzione'': si spiega come funziona il corso, chi sono i docenti, con quali
metodi e con quale calendario si lavorerà, e ci si presenta a vicenda. ``C'è poco
gioco analogico'', si scusano. Ma nella seconda parte la lezione comincia
comunque a costruire un \textbf{linguaggio comune}: le figure professionali
dell'industria del gioco analogico, la differenza tra \emph{game design} e
\emph{game development}, la situazione del mercato italiano rispetto a quello
estero, e un primo glossario di design (\emph{game} e \emph{play}, il framework
MDA, le parole ``vietate'', il playtest e l'iterazione). Si chiude con lo
strumento di settore per i prototipi digitali dei giochi da tavolo,
\textbf{Tabletop Simulator}, e con una lunga sessione di domande.

La lezione si svolge \textbf{online}, su Zoom, ed è condotta a più voci dai tre
docenti principali del corso: \textbf{Daniele}, narrative designer, che guida la
parte organizzativa; \textbf{Ruggero}, game designer, che guida la parte sui
fondamenti e le domande finali; \textbf{Beatrice}, progettista di giochi
pervasivi, che interviene soprattutto nella prima metà. Gli studenti iscritti
sono \textbf{nove}; uno è assente e un altro arriva a lezione iniziata. Ci sono
una pausa di una decina di minuti (si riprende alle 18:15) e, dopo le 19:15,
un'ultima mezz'ora di domande con il solo Ruggero.

Tre avvertenze di lettura. La prima: la trascrizione è automatica e storpia
molti nomi propri, compresi quelli dei docenti e dell'accademia (``Dive'',
``IVE'', ``Accademia Live'', ``Ai''); dove la ricostruzione è sicura, o molto
probabile, il capitolo usa la forma corretta e lo segnala in nota; dove non lo
è, riporta la forma della trascrizione e lo dice. La seconda: durante la pausa la
trascrizione registra, ripetuta in loop per decine di righe, una frase senza
senso che non ha nulla a che fare con la lezione, e che il capitolo ignora. La
terza: in molti passaggi più docenti parlano uno dopo l'altro senza che la
trascrizione segni il cambio di voce; il capitolo attribuisce le parole a chi le
pronuncia quando il contesto lo rende chiaro (chi parla di sé, chi viene
chiamato per nome, chi ``passa la parola''), e resta generico altrimenti.

Quanto alla data, la lezione si colloca con ogni probabilità \textbf{nella prima
metà di settembre del 2026}\footnote{La data non è mai detta. La si deduce da
alcuni indizi: il form iniziale va compilato ``entro e non oltre il 15
settembre''; i gruppi verranno annunciati il 25 settembre; le date del
calendario (sabato 3 ottobre, lunedì 12 ottobre, sabato 7 e 14 novembre,
venerdì 4 dicembre) cadono nei giorni della settimana indicati soltanto nel
2026; un docente parla di un argomento ``molto complesso nel 2026''; e fa
ancora molto caldo (``le temperature sono ancora indecenti''). È una
deduzione, non un'informazione esplicita.}: i riferimenti a ``oggi'',
``quest'estate'' e ``l'anno scorso'' vanno letti rispetto a quel momento.

Il percorso del capitolo segue quello della lezione:
\begin{enumerate}
  \item l'avvio, i \textbf{docenti} e le \textbf{presentazioni} dei partecipanti
        (sezioni da~\ref{a1:sec:avvio} a~\ref{a1:sec:partecipanti});
  \item i \textbf{materiali} e gli strumenti del corso: il form, la cartella
        condivisa, Discord (sezione~\ref{a1:sec:materiali});
  \item lo \textbf{scopo} e il \textbf{metodo} del corso, le regole di
        comportamento e il carico di lavoro (sezioni~\ref{a1:sec:scopo}
        e~\ref{a1:sec:regole});
  \item il \textbf{project work} e il \textbf{calendario} (sezione~\ref{a1:sec:project});
  \item la regola sull'\textbf{intelligenza artificiale} e la lunga discussione
        che ne segue (sezione~\ref{a1:sec:ia});
  \item le \textbf{figure professionali} dell'industria del gioco analogico,
        con un approfondimento su marketing e sales
        (sezioni~\ref{a1:sec:figure} e~\ref{a1:sec:marketing});
  \item la differenza tra \textbf{game design} e \textbf{game development}
        (sezione~\ref{a1:sec:design-development});
  \item il \textbf{mercato}: Italia ed estero (sezione~\ref{a1:sec:mercato});
  \item i \textbf{fondamenti di design}: perché serve una terminologia comune,
        \emph{game} e \emph{play}, il framework \textbf{MDA} (e il DPE), le
        \textbf{parole vietate}, chi crea, chi gioca e chi facilita, il
        \textbf{playtest} e l'\textbf{iterazione} (sezioni
        da~\ref{a1:sec:terminologia} a~\ref{a1:sec:playtest});
  \item \textbf{Tabletop Simulator} e il primo task (sezione~\ref{a1:sec:tts});
  \item le \textbf{domande finali}: il tema del project work, il bilanciamento,
        i ruoli nei team, la pubblicazione, il gioco di ruolo in Italia, le
        \emph{house rule} di Monopoly, i giochi ``comunicati male''
        (sezioni~\ref{a1:sec:domande-daniele} e~\ref{a1:sec:domande}).
\end{enumerate}

% ---------------------------------------------------------------------
\section{L'avvio della lezione}
\label{a1:sec:avvio}

La lezione comincia con qualche minuto di attesa. ``Siamo in nove in classe'', e
si aspetta che arrivino tutti; uno dei docenti scherza sul fatto che l'unica
persona che arriverà ``sempre in ritardo di venerdì'' è Gianmarco. Mancano
Giordano e Tiziana; Tiziana arriva, mentre la segreteria fa sapere poco dopo che
Giordano non potrà esserci e vedrà la registrazione. Si aspetta anche Beatrice,
la terza docente, che arriva a sua volta dopo pochi minuti.

Nell'attesa, Daniele anticipa il programma della giornata e qualche
informazione pratica:
\begin{itemize}
  \item oggi è una giornata di \textbf{introduzione}: come funziona il corso, chi
        sono i docenti, le presentazioni degli studenti, le \textbf{metodologie}
        e il \textbf{calendario}. Gli studenti avranno ``un fracco di domande'', e
        alla fine della lezione c'è ``proprio uno slot apposta'' per le domande;
  \item le lezioni sono \textbf{registrate}: di solito la registrazione viene
        consegnata dopo circa \textbf{12--24 ore} sul profilo personale della
        piattaforma dell'accademia, che va attivato con la segreteria (Tiziana,
        già studentessa, probabilmente ce l'ha già);
  \item i docenti ``cuore'' del corso sono tre (Daniele, Ruggero e Beatrice),
        ``soprattutto Ruggero''; a loro si aggiungeranno tre docenti ospiti, per
        meno ore, probabilmente una lezione a testa, ``tranne Marta, che avrà due
        lezioni quest'anno''.
\end{itemize}
E una battuta sul caldo: ``spero che voi siate carichi, anche se le temperature
sono ancora indecenti''. ``Carichi moderatamente, per quello che potete''.

Poi Daniele condivide lo schermo (``non solo le slide, ma anche un po' di altre
cose'') e apre ufficialmente il corso: ``benvenuti e benvenute nella seconda
edizione del corso di Analog Game Design. Questa è la prima lezione''. Il
percorso è di circa \textbf{tre mesi}, fino a dicembre, e durante il corso i
docenti cercheranno di far provare agli studenti ``cosa significa fare Analog
Game Design'' con un \textbf{project work}, un progetto di gruppo di cui si
spiegheranno più avanti la struttura e le richieste.

% ---------------------------------------------------------------------
\section{I docenti}
\label{a1:sec:docenti}

``Anzitutto partiamo da delle presentazioni doverose''. Ognuno dei tre docenti
principali si presenta con una slide che riporta i contatti e una selezione dei
lavori fatti.

\subsection{Daniele: narrativa, fiction e gioco di ruolo}

\textbf{Daniele}\footnote{La trascrizione riporta ``Daniel Cusetto''; nel
resto della lezione tutti lo chiamano ``Daniele''. Il cognome non è
ricostruibile con certezza.} ha già conosciuto quasi tutti gli studenti, perché
ha fatto loro il colloquio di ammissione. È \textbf{narrative designer}, o meglio
\emph{game narrative designer}; insegna nell'accademia che ospita il
corso\footnote{La trascrizione storpia il nome dell'accademia in molti modi
(``Dive'', ``IVE'', ``Accademia Live'', ``Ai'', ``AIVA''). Con ogni probabilità
si tratta dell'\textbf{AIV}, l'Accademia Italiana Videogiochi, che ha sede a
Roma: lo fanno pensare i corsi triennali di game design e narrative design citati
dai docenti e il fatto che una studentessa racconti di essersi trasferita a Roma.
È una ricostruzione, non un'informazione esplicita; nel seguito il capitolo
parla semplicemente dell'``accademia''.} ma lavora soprattutto come \textbf{brand
advisor}. Nel corso si occuperà della parte di \textbf{gioco di ruolo},
\textbf{narrazione} e \textbf{fiction}.

I suoi contatti, ``in quest'ordine'', sono email, Telegram e Discord. Ci tiene a
dare alcune indicazioni precise:
\begin{itemize}
  \item se si vuole avere la certezza di una risposta, bisogna scrivergli su
        \textbf{Telegram}, perché lì ha tutti i suoi clienti;
  \item Discord lo guarda di meno; l'\textbf{email} è per le cose ufficiali;
  \item non bisogna contattarlo altrove: non su Facebook, non su LinkedIn, non su
        Instagram, ``perché li vedo molto di meno''.
\end{itemize}
La slide mostra anche alcuni giochi, attività e brand su cui ha lavorato negli
anni, ma non si dilunga: ``nel caso ne parleremo''.

\subsection{Ruggero: game design e ricerca e sviluppo}

\textbf{Ruggero} (spesso chiamato ``Ruger'') ha fatto il \textbf{game designer}
per varie aziende e oggi è game designer e si occupa di \textbf{ricerca e
sviluppo} per un'azienda del settore\footnote{La trascrizione dice ``per MS'';
il nome dell'azienda non è ricostruibile con certezza. Più avanti Ruggero cita
tra i giochi su cui ha lavorato il gioco da tavolo ufficiale di
\emph{Cyberpunk}, \emph{Marvel United} e \emph{Status}, un gioco sulla
geopolitica.}. Ha un background ``misto'': nel suo portfolio ci sono giochi da
tavolo, giochi di ruolo, giochi analogici, LARP, videogiochi, giochi in realtà
virtuale. Si occupa anche di \textbf{divulgazione} sul gioco e di analisi di
\textbf{game studies} e game design. Più avanti, nelle domande finali, dirà di
avere ``50 e più giochi prodotti e progettati'' e di insegnare da vari anni.

Lo si può contattare liberamente via Gmail e Telegram, e cerca di rispondere
\textbf{entro 24 ore}, salvo imprevisti o eventi: ``sentitevi in libertà di
scrivermi se avete dubbi, senza problemi''. Quando arriva a parlare dei propri
giochi, la slide mostra ``un mix di varie cose, dalle più famose alle meno
famose'', ma Ruggero si imbarazza a presentarli, tra le risate dei colleghi.

\subsection{Beatrice: i giochi pervasivi}

\textbf{Beatrice} (``Bea'') si occupa di \textbf{giochi pervasivi}
(\emph{pervasive games}). ``Se non li avete mai sentiti nominare, non
preoccupatevi'': questo corso è ``uno dei pochi, forse l'unico in Italia'' che
li racconta esplicitamente. Li si vedrà nelle prossime lezioni, ma la definizione
di base è questa: i giochi pervasivi \textbf{utilizzano gli spazi e gli oggetti
della vita quotidiana e li trasformano in spazi e oggetti di gioco}.

Beatrice lavora come \textbf{freelance} e si occupa di diversi tipi di
progettazione: lavora sia con le \textbf{aziende} sia con gli \textbf{eventi dal
vivo}. La sua attività diretta si trova sotto l'etichetta \textbf{Minime
Differenze}\footnote{La trascrizione riporta il nome dell'etichetta come
``Minime Differenze'' e l'indirizzo del sito come ``mininedifferenze.com''; la
grafia esatta del sito non è verificabile dalla trascrizione.}, mentre i lavori
per le aziende sono ``sparpagliati'' sui siti delle varie aziende. Molte delle
cose mostrate nella sua slide sono attività proposte al pubblico ``con un
biglietto, con un gioco in formato evento'': un tema che si toccherà nel corso,
per capire come \textbf{comunicare} i propri giochi in modo che siano
comprensibili, accessibili per il pubblico e \textbf{vendibili}.

Ammette un ``grosso limite'' con i messaggi: è un po' lenta a rispondere, ``ma
rispondo sempre a tutto''. La si può contattare per email (cognome punto nome) o
su Telegram. È ``molto curiosa di vedere cosa produrrete''.

\subsection{I docenti ospiti}

Daniele presenta poi i tre docenti che il corso avrà ``solo per alcune
lezioni'', tutte verso \textbf{novembre}:
\begin{itemize}
  \item \textbf{Diego Ricchiuti}, direttore del corso di Game Design
        dell'accademia, farà una lezione sull'\textbf{RNG} (\emph{random number
        generation}, la casualità nei giochi): ``vi parleremo di RNG quando sarà
        necessario. Comunque avrete una lezione con lui, sappiatelo'';
  \item \textbf{Marta}\footnote{La trascrizione riporta il cognome come
        ``Cercasassi''; la grafia non è verificabile.} farà \textbf{due lezioni
        sul development}: quest'anno ``abbiamo deciso di fare doppietta'',
        perché l'anno scorso il suo intervento è stato ``molto, molto utile'';
  \item \textbf{Flavio}\footnote{La trascrizione riporta il cognome come
        ``Mortarino'' e, più avanti, ``Mortalino'', e il nome della sua casa
        editrice come ``Grampibert''; nessuno dei due è ricostruibile con
        certezza.} farà la lezione di \textbf{game production}; è un editore, e
        di solito partecipa anche alla tavola rotonda con gli editori della
        giornata finale.
\end{itemize}
Le slide sono già online e gli studenti avranno accesso alla cartella condivisa:
``non vi preoccupate di queste cose, ne parleremo a breve''. Ma prima di parlare
del corso, ``vogliamo conoscere voi''.

% ---------------------------------------------------------------------
\section{Le presentazioni dei partecipanti}
\label{a1:sec:partecipanti}

Le domande per la presentazione sono semplici, da affrontare ``a turno, senza
fretta, con calma'':
\begin{itemize}
  \item nome e occupazione (``se volete dirla'');
  \item perché si vogliono fare giochi analogici, e a che livello: che cosa ha
        portato lo studente al corso;
  \item che tipo di esperienze o conoscenze si hanno già nell'ambito, se si è già
        fatto qualcosa oppure no.
\end{itemize}
``Chi vuole partire senza nessun tipo di ritegno e di vergogne?''. Per far
parlare gli studenti guardandoli, Daniele interrompe la condivisione dello
schermo e ricorda lui le domande quando serve. Le presentazioni sono interessanti
anche per i commenti dei docenti, che anticipano diversi principi del corso.

\subsection{Tiziana}

Tiziana fa l'\textbf{infermiera}, ma nella vita ha fatto molti lavori:
impiegata, insegnante e ``tanti altri lavoretti''. Gioca ``da sempre'': il gioco
fa parte anche della sua vita familiare, perché suo fratello fa il programmatore
di videogiochi. Ha già \textbf{organizzato} giochi: una caccia al tesoro per far
conoscere il territorio, due scene con un delitto (giochi con omicidio da
risolvere), giochi per intrattenere le persone. Ha riscoperto di recente i giochi
da tavolo, arrivando a \emph{Carcassonne} con famiglia e amici; poiché è
difficile trovare persone con voglia di giocare nella propria cerchia ristretta,
ha scoperto il mondo delle \textbf{associazioni} ludiche, ``e lì è stato
un'esplosione'': prima in Veneto (è di Vicenza), poi a Roma, dove si è
trasferita.

Ha in mente un progetto ambizioso: far conoscere una certa realtà con la modalità
del \textbf{gioco di ruolo}, o comunque proporre un gioco suo.

\paragraph{Il commento: gli strumenti, non il gioco dei sogni.} Ruggero coglie
l'occasione per un'avvertenza che vale per tutti, visto che molti studenti
arrivano con idee già pronte. Nel corso si lavorerà a un progetto, ma
l'obiettivo dei docenti è dare \textbf{strumenti} da applicare anche
\emph{fuori} dal progetto, ai giochi che gli studenti vorranno fare un domani:
``che non siano esercizi per il corso, ma per voi in primis''. Quindi: ``non
rimaneteci male se non riuscirete a lavorare al vostro gioco dei sogni in questo
corso''. Anche perché ``sarebbe impossibile, in due mesi, fare il gioco dei
sogni''. Il corso darà invece la possibilità di fare un'esperienza ludica
insieme e di imparare a replicarla sulle proprie idee: ``le idee non buttatele,
anzi, il contrario: tenetele''.

\subsection{Nicola}

Nicola è uno \textbf{studente}. Sei mesi prima ha fatto un colloquio con Daniele
perché stava valutando di passare all'accademia, e ha deciso, ``nel dubbio'', di
cominciare da questo corso breve per capire come funziona la formazione;
eventualmente si iscriverà al corso di game design vero e proprio. Gli interessa
in particolare il lato dell'esperienza del giocatore, su cui è ``molto critico''
(qui usa la parola ``divertimento'', e Daniele lo ferma con un ``ne parleremo'':
è una delle parole che il corso vieta, come si vedrà nella
sezione~\ref{a1:sec:parole-vietate}).

Ha 21 anni ed è \textbf{giocatore di ruolo da circa otto}. Prima ancora di
conoscere i giochi di ruolo si era creato un sistema tutto suo, ``pessimo,
disgustoso, terribile'', un sistema a percentuale con il d100 ``terrificante'';
poi ha scoperto i giochi di ruolo da tavolo, i classici (\emph{Dungeons \&
Dragons}, \emph{Daggerheart}, \emph{Vampire}\footnote{La trascrizione riporta
``DD'', ``Daggerart'' e ``Vampire'': quasi certamente \emph{Dungeons \&
Dragons}, \emph{Daggerheart} e \emph{Vampire: The Masquerade}.}). Gioca anche a
\textbf{card game} e ad altri giochi sociali, pur avendo sempre avuto difficoltà
a trovare gruppi, ``avendo vissuto in un buco dimenticato dal mondo''. Come
giocatore è ``l'eterno narratore'' del gruppo, quello che prepara la storia e si
costruisce le regole. Come autore ha fatto soprattutto giochi su carta, ``perché
sono incapace a programmare, per adesso''.

Daniele osserva che il tema del giocatore e del suo rapporto con il gioco è ``un
argomento importante'', anzi ``un vaso di Pandora enorme'' (espressione che
tornerà molte volte nella lezione), e che il corso tratterà anche che cos'è
l'analog game design in generale e quanti tipi di prodotti esistono.

\subsection{Rita}

Rita\footnote{Daniele la chiama prima ``Nino'' (probabilmente un errore della
trascrizione) e poi la ringrazia come Rita.} ha 30 anni e si collega da una
regione del Sud, dove ``ancora si muore di caldo''. Ha studiato
all'\textbf{Accademia di Belle Arti di Napoli}, \textbf{design della
comunicazione visiva}: si occupa dell'aspetto più imprenditoriale della
progettazione artistica. Non volendo progettare per un'impresa (``fare loghi,
robe del genere''), vorrebbe applicare la propria creatività alla realizzazione
di giochi analogici.

Non è una giocatrice di ruolo e non lo è mai stata: è una giocatrice
\textbf{tradizionale}, cresciuta con Monopoly, la dama e giochi tradizionali
cinesi. Ha già \textbf{idee, prototipi e progetti} realizzati, e vuole capire
come renderli \textbf{prodotti per la vendita} e se questo è il campo adatto a
lei; magari anche confrontarsi con i colleghi. La sua progettazione è orientata
al \textbf{sociale} e all'\textbf{educativo}: ``responsabilizzare un po' le nuove
generazioni su come sta andando il mondo e su come potrebbe andare, visto che è
tutto nelle nostre mani''.

\paragraph{Il commento: siate curiosi.} ``Pane per i tuoi denti'', le risponde
Daniele, soprattutto per i LARP e i giochi pervasivi: ``vi apriremo un po' la
testa da questo punto di vista''. E aggiunge un consiglio per tutti: ``siate
curiosi'', perché imparare da altri ambiti dà moltissimo. Ognuno è lì per fare
qualcosa di preciso, ``però siate aperti anche ad altre esperienze''.

\subsection{Tancredi}

Tancredi ha 21 anni e non ha ancora grande esperienza lavorativa: qualche
``piccolo progettino'' che non è andato come voleva. Gli interessa il mondo dei
giochi, digitali e non, con un approccio \textbf{educativo} e \textbf{sociale}.
È arrivato soprattutto dai videogiochi, a cui gioca da quando era piccolo; ai
giochi da tavolo si è avvicinato solo nell'ultimo anno, grazie a una realtà
ludica, perché in famiglia ``non siamo molto da giochi''. Ha realizzato qualche
piccolo videogioco con vari engine: prima quelli senza un vero linguaggio di
programmazione, ``più a nodi'', poi qualcosa di più complicato, e ora sta
lavorando a un piccolo prototipo.

\paragraph{Il commento: niente di quello che vedrete è sprecato.} ``Non butterai
via niente di analog game design'', gli risponde Daniele, ``perché il 99 per cento
delle cose che vedremo si ritrovano quasi uno a uno nel mondo dei videogiochi''.
Tanto che nei corsi triennali dell'accademia (anche in quello di narrative
design, dove insegna Ruggero) il primo anno si parte dai \textbf{giochi da tavolo
e di ruolo}, proprio per spiegare meccaniche più complesse che servono poi per i
videogiochi. Ruggero aggiunge che vale anche il contrario: alcuni dei
\textbf{framework} di game design che si vedranno nelle lezioni tecniche vengono
dal videogioco e sono stati riapplicati all'analogico. ``Evitare la
crossmedialità qui è impossibile''. Anzi, tutte le conoscenze ``mediatiche
adiacenti'' (musica, letteratura, cinema) saranno utili.

\subsection{Lorenzo Vergata}

I Lorenzo in classe sono due; il primo a presentarsi è \textbf{Lorenzo
Vergata}, 26 anni, \textbf{dottorando al Politecnico di Milano}. Gioca ai
videogiochi da quando aveva due anni e mezzo, ma ha giocato al suo primo gioco da
tavolo ``credo nove mesi fa''. Gli fa piacere sentire che molti in classe sono
interessati ai giochi \textbf{sociali} ed \textbf{educativi}, un settore che
interessa anche a lui e che ``può essere esplorato molto'', anche con qualche
aiuto dell'intelligenza artificiale per esperienze più ricercate.

Di esperienza nell'analogico dice di avere fatto ``il minimo sindacale''. Ma
racconta che diciotto anni prima, da bambino, era ``più game designer di
adesso'': prendeva un gioco e ne cambiava le regole, anche al biliardino (``no,
ora introduciamo questa regola''). Gli piace sperimentare, creare regole e
vedere che cosa accade in sé e negli altri. Soprattutto vorrebbe ``una chiave di
lettura per poter rimanipolare la realtà'': tre settimane prima è stato a fare la
\textbf{vendemmia in Francia} e ha visto tantissime dinamiche che potrebbero
diventare un gioco da tavolo o un videogioco. Per esempio: finire contro un rovo
pieno di semi che si attaccano ai vestiti lo rallentava nella raccolta
dell'uva, mentre in un altro filare era più rapido. ``Si può fare pure una
dinamica'', e vorrebbe gli strumenti per farne giochi.

Arriva subito la battuta: una ``vendemmia gamificata'' da proporre in Oltrepò
Pavese (``io sono dell'Oltrepò Pavese, quindi devo per forza bere vino'')\footnote{La
trascrizione dice ``nell'altro papavese'' e non chiarisce chi pronunci la
battuta.}. Daniele conferma che quest'anno il corso avrà ``tante cose sul
sociale'': ``ci piace parecchio''.

\subsection{L'altro Lorenzo}

Il secondo Lorenzo (``secondo, ma non meno importante'') è uno \textbf{scrittore
di racconti brevi} e sta scrivendo il suo primo romanzo. Gli è sempre piaciuto
raccontare storie, e da più grande scriverle; gioca ai videogiochi ed è ``un
amante delle storie e della cultura pop''. Negli ultimi anni è entrato nel mondo
del \textbf{gioco di ruolo da tavolo} e vorrebbe unire le due passioni, le
storie e il gioco di ruolo, magari facendone un lavoro; gli interessa anche lo
\emph{strategic writing}. Daniele gli ricorda che ci sarà una lezione proprio
sulla \textbf{fiction}.

\subsection{Angelo}

Angelo si collega senza videocamera dal lavoro: fa il \textbf{birraio} (in un
piccolo birrificio in Brianza, a Bernareggio, come racconterà dopo la pausa) e,
con il caldo, ``vi risparmio l'orrore''. Per molto tempo era convinto che i
giochi da tavolo fossero solo i classici (Monopoly, Risiko) o i party game; poi,
quattro o cinque anni fa, ha ``sbattuto la faccia'' contro giochi davvero
appassionanti, con una specie di narrazione dietro che li rendeva ancora più
coinvolgenti, e non ha più smesso. È uno che ``ha sempre un botto di idee, ma poi
non mette in pratica assolutamente nulla'': cercava un corso che gli desse le
basi e soprattutto lo \textbf{motivasse} a provare a fare qualcosa.

\paragraph{Il commento: bisogna buttarsi.} Daniele racconta che anche l'anno
prima un suo caro amico si era iscritto proprio per questo motivo. Il corso è
breve, ma cerca di portare gli studenti ``da 0 a 100'' e di far realizzare un
\textbf{prodotto} da presentare alla fine. Il punto è che i giochi ``bisogna
farli, bisogna buttarsi''. E, venendo dalla narrazione, Daniele lo paragona allo
scrivere: si scrive e si riscrive, e allo stesso modo si fa \textbf{playtest}, si
rifà, si aggiorna. ``Non è che si fanno così''.

\subsection{Gianmarco (dopo le slide)}

Gianmarco arriva a lezione iniziata e si presenta più tardi, alla fine della
prima parte, dopo che i docenti hanno chiuso le slide organizzative e la
discussione sull'intelligenza artificiale (``Gianmarco, ti giuro, ti facciamo
presentare''). Lavora nel \textbf{team di innovazione} di un'azienda italiana.
Vuole fare giochi perché è una sua passione, che lo gratifica, e in passato ci ha
provato ``in maniera scomposta''. Cerca un \textbf{metodo} e un confronto con
persone più esperte: ``studiare solo dei manuali è molto sterile'', ``stare
chiusi nel pertugio non aiuta'', ed ``essere autodidatta ha estremamente i suoi
limiti''. Come esperienza è un giocatore e ha letto ``un po' di materiale, di
blog, di manuali'', ``niente di rilevante''.

Manca solo Giordano, che si presenterà alla lezione successiva.

\paragraph{Uno sguardo d'insieme.} La classe è molto varia: un'infermiera,
una designer della comunicazione visiva, un birraio, un dottorando, uno
scrittore, un impiegato nell'innovazione, due studenti ventunenni. Alcuni
vengono dal gioco di ruolo, altri dai giochi tradizionali, altri dai
videogiochi; quasi tutti hanno scoperto i giochi da tavolo moderni da pochi anni
o da pochi mesi. E almeno tre studenti dichiarano un interesse per i giochi
\textbf{sociali} ed \textbf{educativi}. Come si vedrà nella
sezione~\ref{a1:sec:scopo}, i docenti considerano questa varietà una risorsa.

% ---------------------------------------------------------------------
\section{I materiali e gli strumenti del corso}
\label{a1:sec:materiali}

Prima delle slide sul corso, Daniele sbriga un po' di ``clean up casalingo''.

\paragraph{Il form iniziale.} C'è un \textbf{form} da compilare ``entro e non
oltre il 15 settembre'': il QR code è nelle slide, e la segreteria lo ricorderà
su Discord. Richiede circa 10--15 minuti; si può fare subito o in giornata, ``e
ve lo togliete di torno''. È un form standardizzato per tutti i corsi (``non
siamo io e Ruggero a farvi passare la tortura del test''), diviso in due parti:
\begin{itemize}
  \item una parte di \textbf{conoscenza del candidato}: quanto tempo si ha a
        disposizione per il corso, come ci si organizza con il lavoro, che cosa
        piace giocare. È la parte che interessa ai docenti, perché li aiuta a
        organizzare le lezioni;
  \item una parte di \textbf{test logico}, che i docenti non hanno scritto (``io
        non sono responsabile del form'', ``no, no, nemmeno io'').
\end{itemize}

\paragraph{La cartella condivisa.} Tutti gli studenti vengono aggiunti alla
cartella del corso su \textbf{Google Drive}, che si trova tra i file
``condivisi con me''. Gli studenti ne sono \textbf{gestori dei contenuti} e
possono caricarci quello che vogliono. Dentro ci sono:
\begin{itemize}
  \item le \textbf{attività in classe};
  \item le \textbf{dispense};
  \item i \textbf{materiali didattici}, con le slide (quelle della lezione sono
        già lì);
  \item le cartelle dei \textbf{project work} futuri;
  \item i \textbf{task}, che ``saranno pochi, ma ci saranno'';
  \item il \textbf{calendario studenti}, ``importantissimo'', con due schede: il
        calendario delle lezioni e quello del project work.
\end{itemize}

\paragraph{Discord.} Il corso ha anche un \textbf{server Discord}, a cui tutti
sono stati invitati; ci sono i docenti (Daniele, Ruggero, Diego e altri), che si
possono taggare. Di solito è lì che la segreteria pubblica le informazioni, per
esempio se una lezione ``dovesse per caso essere cancellata''. Nelle domande
finali Ruggero aggiungerà che c'è un canale apposito per il corso, che gli
studenti possono usare liberamente per organizzarsi tra loro.

\paragraph{Le reaction di Zoom.} Durante le lezioni su Zoom si possono usare le
\textbf{reaction} (il pollice in alto, il cuoricino, la mano alzata): un modo
rapido per far capire ai docenti se si è capito, se si vuole intervenire. Nel
frattempo arrivano i primi pollici.

% ---------------------------------------------------------------------
\section{Lo scopo e il metodo del corso}
\label{a1:sec:scopo}

La prima serie di slide (che Daniele, per un lapsus, presenta come ``cos'è il
corso di pulire il pavimento di casa'') spiega che cosa il corso vuole fare e
come.

\subsection{Diventare autori analogici}

Lo \textbf{scopo del corso} è fornire le competenze per \textbf{diventare autori
analogici} e stimolare la creatività e l'\textbf{autorialità} di ciascuno: ``voi
imparerete a diventare autori di giochi analogici''. La differenza tra
\emph{game designer} e \emph{game developer} verrà chiarita più avanti (sezione~\ref{a1:sec:design-development}),
ma la parola chiave qui è \emph{autorialità}: ``vogliamo voi, vogliamo le vostre
idee''. Magari non si potrà usare l'idea del gioco della propria vita, ma si
potranno portare nelle lezioni le proprie idee e le proprie sensibilità.

Gli obiettivi elencati nella slide sono questi:
\begin{itemize}
  \item dare \textbf{maggiore consapevolezza nel design e nello sviluppo}. In
        Italia, purtroppo, ``si fa design a volte in maniera non consapevole''.
        Consapevole vuol dire avere almeno un'introduzione alle basi del design e
        dello sviluppo, e fare le cose ``con metodo, con raziocinio, con
        consapevolezza'';
  \item fornire gli \textbf{strumenti essenziali} per un design consapevole;
  \item stimolare la \textbf{sperimentazione} e la \textbf{trasversalità
        mediatica}. Non si chiede di fare ``il \emph{Root} della situazione o il
        \emph{Carcassonne} della situazione''\footnote{La trascrizione dice ``il
        lumaven della situazione o il carcassone della situazione''; il primo
        titolo non è ricostruibile con certezza. Il senso è chiaro: non si
        chiede di rifare un grande classico del gioco da tavolo.}, ma un
        project work \textbf{calibrato sulla durata del corso}, e però
        \textbf{sperimentale} e ``assolutamente transmediale'';
  \item portare \textbf{numerosi esempi} di tipologie di gioco, dall'analogico al
        tavolo, al pervasivo ``e oltre''. Lo si vedrà nella lezione sui
        \textbf{casi studio} del venerdì successivo: ``voi non avete idea di
        quante cose esistono, bellissime, sconosciutissime e favolose''. Daniele
        confessa che preparando e aggiornando le slide i colleghi gli hanno
        mostrato cose che lui stesso non sapeva nemmeno che esistessero;
  \item dare gli strumenti per \textbf{analizzare} i giochi altrui. Ruggero lo
        sottolinea: se il corso riesce nel suo intento, gli studenti sapranno
        non solo fare i propri giochi ma anche analizzare giochi già fatti, che è
        ``una delle skill importanti anche nella parte di ricerca e sviluppo''.
        Si vedranno casi studio e giochi pubblicati, sempre con l'obiettivo di
        capire \textbf{che cosa fanno e perché lo fanno}, per poi fare meglio il
        proprio lavoro;
  \item imparare a \textbf{fare giochi insieme}. Esistono gli autori che lavorano
        da soli (i \emph{solo developer}), ``ma sappiate che anche l'autore che
        lavora da solo poi si deve interfacciare con il pubblico, con un editore,
        magari con un editor, con un publisher''. ``Non lavorerete mai da soli'':
        bisogna imparare a lavorare con gli altri.
\end{itemize}

\subsection{Un corso intersezionale}

Il corso è \textbf{intersezionale}: \textbf{non ci sono lezioni dedicate a una
singola tipologia di gioco}. Non ci sarà ``la lezione sul gioco da tavolo'',
``la lezione sul gioco di ruolo'', ``la lezione sul gioco pervasivo''. Ci saranno
lezioni su temi diversi, ciascuna pensata per essere utile a \textbf{tutte e tre
le sfaccettature} del corso; alcune saranno più importanti per il gioco da
tavolo, altre per il pervasivo, altre per il gioco di ruolo, ma nessuna sarà
esclusiva. I docenti vogliono ``contaminazioni e applicazioni assolutamente
trasversali''.

\subsection{Un corso cross-generazionale}

È anche un corso \textbf{cross-generazionale}: i gradi di esperienza di gioco
quest'anno sono abbastanza allineati, ma le esperienze di vita, le età, le
estrazioni sociali sono varie. ``Ed è un pro, non è un contro. Più diversificata
è la classe, più imparate meglio anche voi''. La slide ha una piccola nota per
fasce d'età:
\begin{description}
  \item[Under 30] Concentratevi sulla vostra \textbf{crescita}. Evitate di
        confrontarvi con le persone più grandi che magari ``ce la fanno meglio''
        in alcune cose.
  \item[Over 30] Concentratevi sul \textbf{consolidamento}, ma siate
        \textbf{umili}: i docenti diranno cose che faranno ``storcere il naso''.
        ``Bisogna un po' reimparare per imparare'': bisogna sapersi adattare.
\end{description}
Ruggero contestualizza: lo dicono perché in passato hanno visto studenti che,
essendo più adulti o avendo giocato di più, pensavano che l'esperienza
\emph{di gioco} si trasformasse automaticamente in competenza \emph{di
creazione} di giochi. Non è così. E non bisogna sottovalutare chi gioca meno:
proprio perché gioca cose diverse, per risolvere un problema di design specifico
``è una ricchezza incredibile avere qualcun altro che riesce a uscire dalla
scatola''.

\subsection{Un corso pratico e adattabile}

Il corso è \textbf{pratico}: gli studenti si concentreranno su un progetto di
gruppo, il \textbf{project work}, lavorando in team. Ed è \textbf{adattabile}: il
programma è definito, ma ``si aggiorna sui vostri stimoli, priorità e
necessità''. Se durante il corso i docenti vedono che serve più una cosa
dell'altra, sono ``prontissimi'' a dare più spazio a un argomento invece che a
un altro.

\subsection{Inquiry-based learning: domandate}

Il metodo è quello dell'\textbf{inquiry-based learning}, l'apprendimento basato
sulle domande. Di solito le persone dicono ``ho una domanda, magari sarà
stupida''. ``Non esistono domande stupide. Pensate che ci sia una domanda
stupida? Fatela, alzate la mano''. Vanno bene anche le
\textbf{provocazioni}: se qualcosa non torna, è meglio chiedere. La condivisione
è fondamentale.

I docenti danno \textbf{metodi, dati e conoscenze}, ma lasciano trovare le
risposte agli studenti: ``vi lasciamo un pochettino annegare nel mare prima di
darvi la soluzione, in modo tale che la possiate trovare da soli''. Di
conseguenza \textbf{non ci saranno interrogazioni sulle definizioni}: interessa
vedere \textbf{come} gli studenti applicano ciò che imparano. Se una
definizione non viene in mente, non è un problema; l'importante è saperla usare.

\subsection{Feedback, non valutazione}

``Sì al feedback, no alla valutazione''. I docenti danno \textbf{feedback
verbali}, anche sulle attività da fare a casa (i task, che non saranno tanti).
Il corso è calibrato per \textbf{studenti lavoratori}: c'è un po' da lavorare
fuori dalle lezioni, ma in modo calibrato.

Il principio è \textbf{fallire oggi per imparare per il futuro}: \emph{fail
forward}, \emph{fail to learn}. I docenti lasciano provare e fallire almeno una
volta prima di dare la soluzione: ``il classico metodo Montessori, quello
originale, non quello travisato''.

\subsection{Le soft skill e i criteri di valutazione}

Grande attenzione va alle \textbf{soft skill}: i lavori di gruppo, lo scambio
propositivo, l'atteggiamento positivo, la propositività, la proattività.
``Lavoro di gruppo'' non significa stare ogni giorno in sala studio con gli
altri: va benissimo che un gruppo si organizzi dandosi priorità e tempi in base
alla disponibilità di ciascuno. L'importante è \textbf{collaborare}, essere
\textbf{professionali}, rispettarsi a vicenda, rispettare le idee proprie e
altrui e lasciare spazio agli altri, perché i progetti ``non saranno solo
vostri, saranno di tutte le persone del vostro gruppo''.

La valutazione del project work non riguarda solo la \textbf{documentazione}
prodotta, ma anche:
\begin{itemize}
  \item la \textbf{presentazione} finale;
  \item l'\textbf{autovalutazione};
  \item la \textbf{crescita personale};
  \item l'\textbf{atteggiamento} tenuto nel corso.
\end{itemize}
Daniele cita: ``da ognuno secondo le sue capacità, a ognuno secondo i suoi
bisogni''. Lavorare in gruppo non vuol dire che, in tre, ciascuno lavori
esattamente il 33\% del tempo. Ai docenti serve piuttosto capire \textbf{chi fa
che cosa} nel gruppo, e che il processo interno sia \textbf{trasparente}.

\paragraph{Lavorare in cloud.} Per lo stesso motivo il consiglio è di
\textbf{non lavorare in locale}, ma direttamente nelle cartelle condivise del
project work e dei task. Si è liberi di fare altrimenti, ma lavorando in cloud
si dà ai docenti la possibilità di vedere \emph{come} si lavora, il proprio
``processo mentale''.

% ---------------------------------------------------------------------
\section{Le regole di comportamento e il carico di lavoro}
\label{a1:sec:regole}

Alcune slide sono di quelle che si ripetono in ogni corso (``\emph{repetita
iuvant}'', storpiato per scherzo in ``ripetita Juventus'').

\paragraph{Comportamento.} Prestare attenzione durante le lezioni; interagire;
portare rispetto; essere proattivi; alzare la mano e fare domande.
\textbf{Tolleranza zero} per pregiudizi, maleducazione, mancanza di rispetto,
\emph{cheat} e scorciatoie. Daniele ammette che ``cheat e scorciatoie'' viene da
altri corsi dove il tema è più importante, ma Ruggero lo ferma: ``non
sottovalutate questa cosa'', perché implica anche \textbf{non far fare a ChatGPT
i task}. Daniele ha una slide apposta, ``perché è un discorso molto complesso nel
2026'' (sezione~\ref{a1:sec:ia}).

\paragraph{Le lezioni sono intense.} Le lezioni possono essere \textbf{dense},
``soprattutto le mie, sappiatelo: io sono logorroico'', avverte Daniele. Bisogna
chiedere spiegazioni se serve, prendere appunti, rivedere le registrazioni. Le
lezioni sono ``delle scintille'': modi per far esplorare i temi oltre le slide.
``Andate oltre, se riuscite''.

\paragraph{Il carico di lavoro.} Le regole sui task sono queste:
\begin{itemize}
  \item i task \textbf{sono per voi, non per noi};
  \item se si ha poco tempo, è meglio farli rapidamente piuttosto che non farli:
        ``provateci'';
  \item i docenti cercano di non chiedere più di \textbf{un'ora o due a
        settimana} oltre alle ore di lezione;
  \item non tutte le lezioni avranno un task, e il \textbf{project work} è il più
        importante di tutti.
\end{itemize}

\paragraph{I feedback dei docenti.} I feedback ``si guadagnano'' man mano che si
reagisce bene: non bisogna aspettarsi ``dei papiri'' sui task. Se i docenti
hanno voglia di dare feedback più ampi lo fanno, soprattutto se vedono che lo
studente se li merita. Ad Analog di solito non ci sono task \emph{rejected}; ma
se un docente rifiuta un task, vuol dire che si è fatto \textbf{meno del livello
minimo} richiesto dal corso. I feedback possono non essere immediati: ``siate un
pochettino impietositi dal fatto che siamo dei docenti''. E se un feedback non
arriva, o se ne vogliono di più, bisogna scrivere, ``perché siamo tre esseri
umani e quindi possiamo anche dimenticarci''.

% ---------------------------------------------------------------------
\section{Il project work e il calendario}
\label{a1:sec:project}

\subsection{Come funziona il project work}

Una prima slide con le date del project work è da aggiornare (``perdonatemi il
refuso''); Daniele la corregge al volo e rimanda alla slide successiva, più
precisa. Il funzionamento è questo:
\begin{itemize}
  \item i \textbf{gruppi li formano i docenti}. Gli studenti sono nove, quindi si
        faranno \textbf{tre gruppi da tre}: una dimensione pensata anche per il
        fatto che molti sono studenti lavoratori, così da non ritrovarsi ``in due,
        a esplodere'';
  \item ogni gruppo deve prima preparare un \textbf{pitch} del proprio gioco, da
        presentare in anticipo: un'idea, una bozza del gioco;
  \item poi si lavora al gioco \textbf{nelle pause del corso}: da qui a dicembre
        ci saranno due pause, quindi non ci sarà lezione tutti i lunedì e tutti i
        venerdì;
  \item ci saranno \textbf{review intermedie} con i \textbf{tutoraggi}, per lo più
        il \textbf{sabato mattina}, dalle 11 alle 12, con Ruggero su Google Meet.
        Servono ad aiutare i gruppi a portare avanti il project work.
\end{itemize}

\subsection{Il calendario delle lezioni}

Il \textbf{calendario} è nel calendario studenti della cartella condivisa.
Attenzione: le date \textbf{possono variare}. ``Il 90\% rimarranno così'', ma
può capitare che, per motivi di lavoro, qualcuno stia male o non possa, o che
``scompaia l'elettricità da tutte le città d'Italia''. E l'ordine delle lezioni
può cambiare se i docenti vedono che agli studenti serve prima una lezione
invece di un'altra.

Ci sono \textbf{due pause importantissime}:
\begin{itemize}
  \item dal \textbf{13 ottobre al 5 novembre}: è il periodo delle grandi fiere
        (\textbf{Essen} e \textbf{Lucca}\footnote{La trascrizione dice ``Essen,
        Luca, Cose'': si tratta con ogni probabilità della fiera Spiel di Essen,
        in Germania, e di Lucca Comics \& Games, che si tengono entrambe in
        ottobre.}), ``un periodo un po' frenetico'', e cominciano anche i corsi
        accademici, quindi la pausa permette ai docenti di respirare. Per gli
        studenti è la pausa più importante: quasi un mese in cui sviluppare la
        parte importante del gioco;
  \item dal \textbf{24 novembre all'11 dicembre}, dopo qualche lezione a novembre.
\end{itemize}

E ci sono \textbf{due date da evidenziare}:
\begin{itemize}
  \item il \textbf{25 settembre}: nei primi 5--10 minuti della lezione (che per il
        resto sarà tenuta da Beatrice) i docenti annunceranno i tre gruppi;
  \item il \textbf{12 ottobre}: il \textbf{pitch} dei giochi, con l'idea
        embrionale. Tra il 25 settembre e il 12 ottobre ogni gruppo prepara un
        breve pitch in formato slide; quel giorno i docenti lo rivedono e danno
        feedback: se il gioco va bene, se non va bene, come svilupparlo.
\end{itemize}
Dal 12 ottobre fino all'inizio di dicembre si sviluppa il gioco, con i
tutoraggi: ``sfruttate questi momenti di pausa''.

\subsection{Il calendario del project work}

La slide successiva, ``reale e aggiornata'', riporta il calendario del project
work (tabella~\ref{a1:tab:calendario} e figura~\ref{a1:fig:calendario}).
``Per ora non pensate al project work: ci penserete dal 25 settembre in poi''.

\begin{table}[htbp]
\centering
\begin{tabular}{@{}lp{9.2cm}@{}}
\toprule
\textbf{Data} & \textbf{Che cosa succede} \\
\midrule
ven 25 settembre\footnotemark & lezione: annuncio dei gruppi e del
\emph{mandate}, cioè di che cosa si dovrà fare \\
sab 3 ottobre, ore 11 & prima review e tutoraggio con Ruggero (dopo una
decina di giorni di lavoro) \\
lun 12 ottobre & pitch del gioco durante la lezione \\
13 ottobre -- 5 novembre & prima pausa: sviluppo del gioco \\
sab 7 novembre & tutoraggio \\
sab 14 novembre & tutoraggio \\
24 novembre -- 11 dicembre & seconda pausa: sviluppo del gioco \\
ven 4 dicembre, entro mezzanotte & consegna del gioco sul Drive (non c'è
lezione) \\
ven 11 dicembre & i giochi vengono visti a lezione \\
14 dicembre & giornata finale, con la tavola rotonda con gli editori \\
\bottomrule
\end{tabular}
\caption{Il calendario del project work come lo presenta Daniele. Le date
possono variare.}
\label{a1:tab:calendario}
\end{table}
\footnotetext{Daniele dice ``lunedì 25 settembre'', ma nel 2026 il 25 settembre
è un venerdì; il giorno della settimana nella tabella è corretto di
conseguenza.}

\paragraph{Perché la consegna è il 4 dicembre.} Il 4 dicembre non c'è lezione:
si deve solo consegnare il gioco, entro mezzanotte, sul Drive. La presentazione
a lezione è l'11 dicembre, ma i tre docenti hanno bisogno di \textbf{una
settimana} per rivedere il gioco in privato e valutare anche la
\textbf{documentazione}; per questo la consegna non coincide con la
presentazione. Che cosa voglia dire esattamente ``consegnare il gioco'' verrà
spiegato più avanti.

\paragraph{La giornata finale.} Il 14 dicembre c'è la \textbf{giornata finale},
con una \textbf{tavola rotonda con gli editori}: di solito viene invitato Flavio,
che è anche docente di game production, più qualche altro editore, ``vediamo
quest'anno chi riusciamo a portarvi''.

\begin{figure}[htbp]
\centering
\begin{tikzpicture}[
    font=\footnotesize,
    x=0.1cm,
    evento/.style={draw, thick, rounded corners=2pt, fill=white,
                   text width=1.8cm, align=center, inner sep=2pt},
    pausa/.style={fill=gray!20, draw=gray!60},
    tacca/.style={thick}
  ]
  % asse: giorno 0 = 25 settembre; 1 unità = 1 giorno
  \draw[pausa] (18,-0.35) rectangle (41,0.35);
  \draw[pausa] (60,-0.35) rectangle (77,0.35);
  \node[font=\scriptsize\itshape] at (29.5,0.18) {pausa};
  \node[font=\scriptsize\itshape] at (65,0.18) {pausa};
  \draw[-{Stealth[length=2.5mm]}, thick] (-3,0) -- (85,0);
  % eventi sopra
  \draw[tacca] (0,-0.15) -- (0,0.9);
  \node[evento, above] at (0,0.9) {25 set\\ gruppi e \emph{mandate}};
  \draw[tacca] (17,-0.15) -- (17,2.2);
  \node[evento, above, fill=gray!8] at (17,2.2) {12 ott\\ \textbf{pitch}};
  \draw[tacca] (70,-0.15) -- (70,0.9);
  \node[evento, above, fill=gray!8] at (70,0.9) {4 dic\\ \textbf{consegna}};
  % eventi sotto
  \draw[tacca] (8,0.15) -- (8,-0.9);
  \node[evento, below] at (8,-0.9) {3 ott\\ tutoraggio};
  \draw[tacca] (43,0.15) -- (43,-0.7);
  \draw[tacca] (50,0.15) -- (50,-0.7);
  \draw (43,-0.7) -- (50,-0.7);
  \draw[tacca] (46.5,-0.7) -- (46.5,-0.9);
  \node[evento, below] at (46.5,-0.9) {7 e 14 nov\\ tutoraggi};
  \draw[tacca] (77,0.15) -- (77,-0.9);
  \node[evento, below] at (77,-0.9) {11 dic\\ giochi a lezione};
  \draw[tacca] (80,-0.15) -- (80,2.2);
  \node[evento, above] at (80,2.2) {14 dic\\ giornata finale};
\end{tikzpicture}
\caption{Una linea del tempo del project work, ricostruita dalle slide (scala
approssimativa: dal 25 settembre al 14 dicembre). Le bande grigie sono le due
pause senza lezioni, in cui i gruppi sviluppano il gioco.}
\label{a1:fig:calendario}
\end{figure}

Finita questa parte, Daniele apre alle domande: ``Domande sul calendario? Sulle
slide che ho fatto? Considerazioni? Non abbiate timore, io non ci credo''. E la
domanda arriva, sull'unico tema che aveva lasciato da parte.

% ---------------------------------------------------------------------
\section{L'intelligenza artificiale nel corso}
\label{a1:sec:ia}

\subsection{La slide di Daniele}

Daniele aveva già presentato una slide che definisce ``importantissima''. ``Lo so
che è paradossale, me ne rendo conto. Lo so che magari mi darete del luddista,
dell'anti-tecnologia: non è così. Uso l'intelligenza artificiale per lavoro. Ma
in questo corso, vi prego, \textbf{non usate l'intelligenza artificiale}''. La
richiesta vale per le attività, per i task e \textbf{anche per la
prototipazione}. Le ragioni:
\begin{itemize}
  \item ``voglio che siate voi a fare il gioco, voglio che siate voi a imparare a
        farlo''. Daniele parla da utilizzatore: prima di ChatGPT, che ha smesso
        di usare, oggi soprattutto di Claude, di cui ``si sottovaluta quanto
        possa essere incerto, insicuro, impreciso, non all'altezza'';
  \item nella game industry, \textbf{i publisher rifiutano automaticamente} i
        pitch e i prototipi fatti con l'intelligenza artificiale. Daniele lo ha
        visto di persona in ``tantissime pitch review'' quell'estate, anche al
        \textbf{First Playable} di Firenze\footnote{La trascrizione dice ``First
        Playboy a Firenze'': si tratta del First Playable, l'evento di business
        dell'industria videoludica che si tiene a Firenze.}: anche nel mondo dei
        videogiochi i publisher che vedono pitch o prototipi fatti con l'IA
        ``rigettano automaticamente la build'', perché hanno ancora timore
        dell'IA e delle sue possibili conseguenze;
  \item ``non so come sarà il futuro'': non sa se bisognerà abituarsi a usarla
        tutto il giorno. Non è ``totalmente contrario'', purché sia usata ``con
        raziocinio''; secondo lui oggi, negli Stati Uniti, non lo è;
  \item ma in un corso che si chiama \emph{Analog Game Design} gli studenti devono
        saper fare i giochi da soli. In futuro, se vorranno farsi aiutare
        dall'IA per essere più efficienti, saranno liberissimi di farlo; durante
        il corso no, perché i docenti vogliono valutare \textbf{le loro
        competenze e capacità}.
\end{itemize}

\subsection{La domanda: che cosa vuol dire usare l'IA in un gioco analogico?}

La domanda di uno studente è precisa: che cosa si intende per ``task fatti con
l'IA''? Qual è l'uso dell'IA in un gioco analogico?

\paragraph{Ruggero: l'esperienza di chi valuta i prototipi.} Ruggero risponde
dalla sua esperienza nel settore, dove riceve prototipi da valutare per
un'eventuale pubblicazione. Distingue due usi:
\begin{enumerate}
  \item le \textbf{immagini generate}: illustrazioni, graphic design e così via.
        Per molti publisher è ``un \emph{reject} immediato, proprio immediato'';
  \item l'IA usata per il \textbf{regolamento}, per scrivere le regole e le
        meccaniche. È più difficile da individuare, ma molti publisher se ne
        accorgono comunque, perché le IA ``non sono ancora brave abbastanza da
        creare un regolamento coeso meccanicamente e scritto con un senso per il
        vostro gioco''. Quello che fanno è prendere parti di altri regolamenti,
        elaborare le strutture che vedono nei regolamenti pubblicati e
        riapplicarle. Ma molti giochi si prestano poco a un regolamento
        strutturato come quello di altri, perché \textbf{ogni gioco è un
        ecosistema meccanico a sé}. Il risultato sono regolamenti con
        riferimenti non immediati, e soprattutto con \textbf{la stessa meccanica
        descritta in modo diverso in pagine diverse}: due difetti che fanno
        perdere tempo all'editore, e in questo caso ai docenti.
\end{enumerate}
Poi aggiunge un punto ``personale, soggettivo'': i docenti sono interessati a
vedere il lavoro, la creatività, ``il frutto della vostra creazione''. E
viceversa gli studenti vogliono i commenti dei docenti: ``immaginatevi se voi
fate un gioco e noi vi diamo i commenti scritti con ChatGPT''. Altrimenti, dice
ridendo, gli studenti fanno il gioco con ChatGPT, i docenti lo danno a ChatGPT
per i commenti, ``e facciamo molto alla svelta''.

\paragraph{Daniele: perché gli studenti usano l'IA.} Daniele racconta la sua
esperienza con gli studenti, che secondo lui ricorrono all'IA in due momenti:
\begin{enumerate}
  \item quando vanno in \textbf{crisi}: in ansia da prestazione o per i tempi.
        Cercano di risolvere il problema della propria disorganizzazione
        facendolo risolvere all'IA, senza accorgersi che l'output è ``molto,
        molto scarso'';
  \item quando hanno l'impressione che l'intelligenza dell'IA sia
        \textbf{migliore} di quello che scrivono loro, o di quello che c'è in
        giro. Per Daniele è una ``mancanza di palato critico'': l'output dell'IA
        sembra buono in superficie, ma quando si va nel dettaglio e ci si fa
        delle domande ``di solito fa un po' di \emph{crumbling}'', si sbriciola.
\end{enumerate}
Lui usa Claude ogni giorno e ne conosce i limiti: ``non c'è nemmeno una frase
che io copio direttamente''. Gli fa fare il lavoro più noioso, più ripetitivo e
soprattutto \textbf{strutturato}, ``replicare delle cose come se fosse una
macro''. Lo usa per il testo e per il codice, mai per le immagini: per motivi di
\textbf{diritto d'autore}, e perché ha molti amici artisti ed è stufo di vedere
il loro lavoro ``dilaniato'' dall'IA.

Ammette che i docenti non possono controllare gli studenti; per questo
``ve lo consiglio''. E riprende quella che definisce la cosa più bella detta da
Ruggero (``ti ruberò la frase per il corso di narrative''): i docenti i feedback
li danno di persona. Daniele a volte manda 30 minuti di audio a uno studente, o
fa intere call per rivedere script, pitch e giochi: ``è il mio lavoro, mi piace
farlo e voglio farlo io, quindi mi aspetto che lo studente lo faccia lui''. Come
freelance, poi, è ``stufo marcio'' dei clienti che arrivano con qualcosa fatto
con l'IA (``guarda che bello, pubblica questo''), perché il più delle volte il
risultato è pessimo: non sanno come usarla. ``È anche una questione di
competenza''.

\paragraph{Il prototipo non deve essere bello.} Questo vale anche per il
\textbf{prototipo}, su cui si tornerà nelle lezioni di prototipazione: in un
prototipo non servono ``graficoni allucinanti''. Le grafiche bellissime fatte con
l'IA spesso \textbf{peccano di user experience}: sembrano bellissime, ma quando
il giocatore deve leggere la carta non capisce qual è il costo, qual è il
potere, qual è l'effetto.

\paragraph{Beatrice: si perde l'occasione di imparare.} Beatrice aggiunge due
motivi che valgono ``al di là dei nostri credo''. Il primo: chi usa l'IA in
sostituzione (per esempio per scrivere un testo, per accelerare i tempi o per
trovare strategie) \textbf{perde la parte in cui può imparare}. Nel game design
bisogna ``tenere insieme i pezzi'' di un regolamento, ma anche dell'esperienza
dell'utente e di come trasmetterla; i passaggi necessari (creare il regolamento,
creare le meccaniche, sperimentare, raccogliere i feedback) ``si capiscono solo
vivendoli''. E si sta pagando un corso per imparare. Il secondo, sul piano
editoriale: chi presenta a un editore un prodotto che ha come base l'IA rischia
da un lato di non avere imparato abbastanza (e di ritrovarsi ``a piedi'' sul
campo senza IA), dall'altro che l'editore se ne accorga, ed è un tema
\textbf{reputazionale}. Ci sono editori ``più illuminati di altri'', quindi
anche qui c'è un mercato da capire; ma i docenti preferiscono dire che cosa è
più sensato per il mercato come lo conoscono.

\paragraph{Ruggero: non abbiate paura del giudizio.} Ruggero vuole dare uno
``spotlight'' a un punto detto da Daniele e rimasto ``semi-ignorato'': a volte
l'uso dell'IA nasce dalla \textbf{paura del giudizio}, dal voler consegnare
qualcosa di più pulito, più strutturato possibile. ``Vi giuro, non ve ne
preoccupate''. Ruggero detesta i voti e non ragiona per voti; non valuterà
nessuno per grafiche fatte a matita. Insegna da anni, ha avuto studenti con
dislessia e discalculia: refusi ed errori non gli interessano, ``non è un corso
di editing o di proofreading''. Gli interessa che le \textbf{meccaniche siano
coerenti con lo scopo di design}. Immagini fatte a matita o con Paint ``non sono
un problema, anzi, le preferisco'' rispetto all'``ennesima immagine fatta con
ChatGPT o Midjourney''.

\paragraph{Daniele: il pensiero critico, costo e qualità.} L'ultimo intervento
è di Daniele, sul \textbf{pensiero critico}. Ragionando per assurdo (``un
pensiero di Belzebù''), immagina che un domani l'IA venga usata ovunque. I suoi
limiti, secondo lui, resterebbero due:
\begin{description}
  \item[Il costo] Tutti dicono che sta diminuendo e che diminuirà, ma non se ne
        può essere sicuri al 100\%.
  \item[La qualità] Cioè la capacità di leggere l'output. Se si dà un regolamento
        a Claude o a ChatGPT chiedendo quali problemi ha, il modello risponde con
        i problemi che, in percentuale, è più probabile che gli esseri umani
        abbiano già individuato nei dati su cui è stato addestrato. Ne derivano
        due problemi: se il problema del proprio gioco non è mai stato
        individuato prima, il modello non lo vedrà; e viceversa può segnalare
        problemi che il gioco non ha.
\end{description}
Fa l'esempio di Grammarly, che usa nella versione premium per rivedere i testi
in inglese: spesso non accetta i suoi consigli, perché se vuole scrivere un
dialogo di videogioco in \textbf{dialetto Cockney} londinese, lo deve scrivere
in Cockney. Se Grammarly dice che è sbagliato, ``sei tu che non hai capito il
motivo per cui voglio scrivere questo testo in Cockney''. La conclusione: prima
di usare un'IA per un'attività bisogna sapere \textbf{che cosa si vuole
ottenere}, qual è l'obiettivo, e quali \textbf{competenze} servono per valutare
in modo critico se l'output di un qualunque LLM è davvero valido. Daniele non
vuole dilungarsi oltre: è un argomento ``molto caldo'' per lui, che è anche
scrittore e si sente dire da tutti ``ti sostituiranno''.

Subito dopo si presenta finalmente Gianmarco (lo si è visto nella
sezione~\ref{a1:sec:partecipanti}), e poi si fa una pausa di una decina di
minuti. Daniele avverte che nella terza parte lui non ci sarà (alle sette e un
quarto ha un impegno di lavoro), e lo stesso vale per Beatrice: la parte finale
di domande e risposte sarà con il solo Ruggero, ``in ottime mani''.

% ---------------------------------------------------------------------
\section{Le figure professionali dell'industria analogica}
\label{a1:sec:figure}

Dopo la pausa (in cui si scopre che Angelo lavora in un piccolo birrificio in
Brianza e si scherza sul rischio di un \emph{drinking game}, tra la birra e la
vendemmia) la lezione riprende con le slide di Ruggero. ``Oggi è una lezione
introduttiva'': gli argomenti si vedranno ``molto dall'alto'' e verranno
approfonditi nelle lezioni successive. Ma servono già ora, per creare
\textbf{l'inizio di un linguaggio condiviso}: ci sono termini che torneranno
spesso, a cominciare da quelli del mercato e dell'industria del gioco
analogico, un'industria ``molto varia, molto caotica, anche molto mutevole'',
con tanti tipi diversi di giochi e altrettanti tipi diversi di entità che ci
lavorano.

\subsection{Una decina di etichette}

I docenti hanno raccolto ``una decina di etichette'', cioè di \textbf{ruoli} che
nell'industria sono o possono essere presenti. Due avvertenze:
\begin{itemize}
  \item nelle aziende più piccole \textbf{la stessa persona copre spesso più
        ruoli}, per forza di cose. Sarà anche il caso degli studenti, che
        lavoreranno in gruppi di tre e non di venti: è facile che ciascuno finisca
        a fare più cose. I docenti ne terranno conto, e qui entreranno in gioco
        anche competenze diverse dal game design (programmi di grafica,
        illustrazione, impaginazione, layout, \emph{proofreading}\ldots);
  \item Daniele precisa che \textbf{non si chiede agli studenti di diventare}
        illustratori o altro. Ruggero concorda, ma aggiunge che è utile sapere
        anche \emph{che cosa non è il proprio lavoro}. Se in un gruppo nessuno sa
        illustrare, ai docenti va benissimo un prototipo che mostri dove vanno
        le immagini: \textbf{placeholder}, concept, anche semplici box di testo
        con scritto ``qui vorremmo un'immagine di\ldots''. ``È importante anche
        sapere che cosa sapete fare e che cosa non sapete fare''.
\end{itemize}
La tabella~\ref{a1:tab:figure} riassume le figure; le sottosezioni seguenti le
riprendono una per una, con i commenti dei docenti.

\begin{table}[htbp]
\centering
\small
\begin{tabular}{@{}p{3.3cm}p{9.6cm}@{}}
\toprule
\textbf{Figura} & \textbf{Di che cosa si occupa} \\
\midrule
Game designer & crea le regole alla base dell'esperienza: le meccaniche, come
interagiscono tra loro e come chi gioca interagisce con esse \\
Game developer & fa da ponte tra chi fa game design e il publisher: bilancia il
prototipo, cura i modelli matematici, adatta il gioco (e il paragame) alle
necessità del publisher \\
Editor (proofreader) & revisiona e adatta i testi a partire dalla prima stesura
del designer (correzione, linguaggio inclusivo\ldots) \\
Publisher (editore) & pubblica e produce fisicamente il gioco; il suo marchio è
sulla scatola o sulla copertina \\
Sales e marketing & il lato commerciale: vendite (anche B2B) e comunicazione del
prodotto \\
Graphic designer & icone meccaniche, wireframe, UX: il lato funzionale degli
elementi visivi \\
Illustratore & le illustrazioni: il \emph{look and feel}, il world building, il
``vibe'' del gioco \\
Sensitivity reader (cultural consultant) & verifica l'accuratezza e la
rappresentazione di un tema sensibile che il team non conosce \\
Producer & coordina il progetto (quasi project management), parla con le
tipografie, verifica che i materiali siano sostenibili economicamente \\
\bottomrule
\end{tabular}
\caption{Le figure professionali dell'industria del gioco analogico presentate
nella slide di Ruggero, in sintesi.}
\label{a1:tab:figure}
\end{table}

\subsection{Game designer}

Il \textbf{game designer} (o la game designer) è chi si occupa, ``sembra banale
a dirlo'', di \textbf{creare le regole alla base dell'esperienza} che si vuole
far arrivare. Vale per tutti i tipi di giochi, non solo per quelli analogici, e
che cosa faccia concretamente in un progetto o in una serie specifica ``può
mutare tantissimo''. Nell'analogico non verrà chiesto di programmare o di
prototipare in digitale; verrà chiesto di \textbf{creare le meccaniche di
gioco}, di capire \textbf{come le meccaniche si interfacciano tra loro} e
\textbf{come le persone che giocano si interfacceranno con le meccaniche}.

\subsection{Game developer}

Il \textbf{game developer} (o la game developer) sarà l'argomento delle due
lezioni di Marta. La parola \emph{development} è nota soprattutto
dall'ambito videoludico, ma nell'analogico ha un significato molto specifico
``che non c'entra niente con la programmazione'': il game developer è
\textbf{un ponte tra chi fa game design e il publisher}. Di solito si occupa di
\textbf{bilanciare il prototipo}, di curare i \textbf{modelli matematici}, di
adattare il gioco e il \textbf{paragame}\footnote{La trascrizione dice ``il
Wix, il paragame''. Il \emph{paragame} è, in termini generali, tutto ciò che
accompagna il gioco senza esserne il sistema di regole (regolamento, scatola,
materiali, comunicazione); i docenti lo tratteranno in una lezione dedicata, e
qui il termine è solo citato.} alle necessità del publisher. La distinzione
viene ripresa e approfondita nella sezione~\ref{a1:sec:design-development}.

\subsection{Editor}

L'\textbf{editor}, a volte chiamato \textbf{proofreader} a seconda dell'azienda,
si occupa della revisione dei \textbf{testi}: non della stesura iniziale, perché
di norma, data la scala dell'industria, è chi fa game design a scrivere un
\textbf{primo draft}, e l'editor ci passa sopra. Non solo per correggere i
refusi: spesso il suo lavoro è anche \textbf{adattare} i testi, per esempio
renderli \textbf{gender neutral}. Se ne vedranno casi studio più avanti.

\subsection{Publisher}

Il \textbf{publisher}, o \textbf{editore}, è l'entità che nell'ambito
commerciale \textbf{pubblica} giochi, giochi di ruolo o giochi da tavolo. È una
figura ``più solidificata'' nel gioco da tavolo e nel gioco di ruolo da tavolo;
\textbf{nell'ambito pervasivo non esiste proprio in questa forma}, come
mostrerà Beatrice parlando dei casi studio. Se si compra un gioco da tavolo o un
gioco di ruolo in un negozio, il marchio del publisher è di solito sul retro
della copertina o sul fronte della scatola: è l'entità che, di solito,
\textbf{produce fisicamente} il gioco.

\subsection{Sales e marketing}

L'ambito \textbf{sales e marketing}, a volte affidato a un'unica persona e a
volte separato, riguarda il lato commerciale. In alcune aziende è interno al
publisher, e si occupa di parlare con le fabbriche e le tipografie, di ``far
quadrare i conti''. In altre realtà è completamente diverso, specialmente nel
\textbf{B2B} o nei giochi \textbf{didattici} e \textbf{formativi}, di cui
Daniele e Beatrice ``sapranno parlare meglio''. Il tema viene ripreso più avanti
in dettaglio (sezione~\ref{a1:sec:marketing}).

\subsection{Graphic designer e illustratore}

Ruggero presenta insieme \textbf{graphic designer} e \textbf{illustratore}
proprio perché sono ruoli molto diversi che si tende a far coincidere:
\begin{itemize}
  \item chi fa \textbf{graphic design} si occupa delle \textbf{icone meccaniche},
        del \textbf{wireframing}, della \textbf{UX}: del lato \textbf{meccanico e
        funzionale} degli elementi visivi, cioè di tutto ciò che serve
        meccanicamente nel gioco, o anche alla sua ``qualità della vita'';
  \item chi fa \textbf{illustrazione} si occupa dell'\emph{art}, delle
        illustrazioni: copre il \textbf{look and feel}, la parte di
        \textbf{world building} e il ``vibe'' che si dà al gioco.
\end{itemize}

\paragraph{Daniele: che cos'è la user experience.} Daniele interviene sulla
\textbf{user experience}, ``importantissima per me'': ha un background da
\textbf{layout artist} (impaginava magazine), un mestiere in cui la UX conta
moltissimo. Per user experience non si intende solo l'\textbf{interfaccia
grafica}, la UI, che anche nei videogiochi è disegnata da un artista. La UX è
qualcosa che può fare anche il game designer: \textbf{disegnare l'esperienza del
giocatore}, cioè ``il gesto, l'intenzione, l'intuito'' che stanno dietro
l'aspetto visivo dell'interfaccia. Per fare UX non serve saper disegnare.
L'esempio è quello della \textbf{porta}: se si progetta una porta, bisogna
indicare all'utente come si apre, come si prende la maniglia; ``quella è user
experience''. Il design della maniglia in sé, invece, è
\textbf{interfaccia}.

\paragraph{Un aneddoto sulle icone.} Ruggero aggiunge che molte persone che
fanno \textbf{graphic design ben fatto non sanno illustrare}, e spesso vale il
viceversa. Lo si nota quando le icone di un gioco sono bellissime ma non
funzionali: andando a guardare i crediti, di solito si scopre che le icone le ha
fatte la persona che ha fatto le illustrazioni. ``Comunque non andiamo adesso in
questo, perché è un altro vaso di Pandora''.

\subsection{Sensitivity reader}

Il \textbf{sensitivity reader}, o \textbf{cultural consultant} a seconda del
progetto, fa un lavoro di \emph{sensitivity reading} del gioco. È una figura
molto usata nell'ambito \textbf{formativo e didattico}. L'esempio di Ruggero: se
si fa un gioco per un ospedale su una malattia specifica, chi lo progetta non è
un medico, ma il gioco deve essere \textbf{accurato} e \textbf{rappresentativo}
delle dinamiche che quella malattia comporta. Serve una persona che conosca quel
tema, anche se non sa che cosa sia il game design. Sono figure chiamate di
solito \textbf{a progetto}, per giochi specifici, che compensano competenze che
il team non ha ``per forza di cose, perché noi non sappiamo tutto''; e a volte,
per lavoro, si devono fare giochi su temi che non si conoscono personalmente.

\paragraph{Daniele: non è un censore.} Daniele segnala un problema italiano:
alcune persone, anche online, hanno confuso il sensitivity reader con un
\textbf{censore}, ``una persona che fa censura, soprattutto in materia di
\emph{woke}''. ``Assolutamente non è così''. La parola \emph{sensitive} non
significa che il pubblico sia sensibile a un tema, ma che \textbf{l'argomento di
cui tratta il gioco è sensibile}. In un gioco sulla medicina, realistico o no,
non si possono inserire imprecisioni. Per questo esistono i sensitivity reader
``più disparati'': da quelli per la cultura indiana o afroamericana fino a quello
su ``come si allevano gli gnu in Africa''. È una persona competente su un
argomento che aiuta a essere \textbf{veritieri}. Certo, uno degli ambiti
principali in cui si usano è quello culturale, ma ``non fate l'errore di pensare
che siano dei censori''. Daniele ha lavorato con ``centinaia'' di sensitivity
reader e nessuno gli ha mai detto ``non puoi scrivere questa cosa''. Di solito
sono come degli \textbf{editor}: danno opzioni, mostrano i punti problematici,
anche quelli di cui non ci si era nemmeno accorti.

Ruggero conferma con un esempio di quando lavorava in \textbf{Danimarca} con
l'\textbf{OMS}: giochi sul \textbf{diabete di tipo 2} nei bambini, per un target
di ragazzi delle scuole medie. C'era un team dell'OMS specializzato nella
comunicazione sul diabete, per garantire che il gioco fosse accurato e ne desse
una rappresentazione adatta al target. Il sensitivity reader aiuta a rendere il
gioco più accurato e ``più tematicamente coeso''.

\subsection{Producer}

Il \textbf{producer} è una figura a volte interna al publisher. In un progetto
grande è la figura di riferimento, ``quasi'' un project manager. In un progetto
piccolo si occupa spesso anche di parlare con le \textbf{tipografie} e di fare la
\textbf{lista dei materiali}, per assicurarsi che i materiali non siano solo
funzionali al gioco ma anche \textbf{economicamente sostenibili} da produrre:
``magari è un gioco bellissimo, ma poi costa più di quanto potete permettervi, e
nessuno vi comprerà mai un card game a 200 euro''.

\subsection{Perché conoscere queste figure}

Ruggero chiude la slide ``tornando a ombrello'': perché parlarne, se gli studenti
sono lì per fare game design? Non solo per saper dare un \textbf{nome} alle
figure. Non c'è niente di male a creare un gioco solo per sé; ma chi vuole fare
giochi a livello lavorativo, perché li giochino altre persone, ``magari più
persone possibile'', \textbf{si interfaccerà con queste figure}, e sapere che
cosa fanno e come parlarci è parte importante del percorso. Nel corso se ne
toccheranno diverse: il development con Marta, la produzione con Flavio, e
lezioni sulla UX, sul paragame, sul graphic design. ``Abituatevi mentalmente:
non siete in solitaria contro il mondo''.

% ---------------------------------------------------------------------
\section{Marketing e sales}
\label{a1:sec:marketing}

Arrivato all'angolo della slide con marketing e sales, Ruggero passa la parola
ai colleghi: ``io è meglio che non mi fate parlare di marketing e sales'', anche
per esperienze personali.

\subsection{Beatrice: i giochi trattati come prodotti}

Beatrice premette che le due voci sono in parte autoesplicative. Nel mondo
capitalista in cui viviamo, i giochi sono ``strumenti molto potenti per fare
\emph{n} cose'', ma sono trattati come \textbf{prodotti}, e il corso ragionerà
anche in quest'ottica. Nel \textbf{marketing} delle industrie ludiche e
videoludiche ci sono varie professionalità che gli studenti probabilmente già
conoscono: dal \emph{copy} alla grafica, tutto ciò che riguarda il
\textbf{vendere} il gioco o l'attività di gioco.

Il lato \textbf{sales} è forse la distinzione meno nota. Come in tante aziende, è
la parte che si occupa delle vendite \textbf{B2B}, cioè tra aziende: fare
accordi per vendere prodotti ludici a diverse realtà, per esempio ai
\textbf{distributori}, o a una casa editrice che vende alle librerie, che poi
vendono al consumatore finale. Chi si occupa di sales nel mondo ludico capisce
come vendere \textbf{alle aziende}. I due mondi si parlano spesso, e soprattutto
nelle realtà piccole si trovano figure che fanno ``un po' sales, un po'
marketing''.

\subsection{Daniele: che cos'è (e che cosa non è) il marketing}

Daniele il marketing si è trovato ``purtroppo'' a farlo, e lo fa anche per
l'accademia. ``Purtroppo'' perché in Italia è mal compreso: molti lo
identificano con l'\emph{advertisement} online. Non è così: il marketing è una
\textbf{disciplina molto più antica}, e in Italia ``abbiamo molto dimenticato''
come si fa il marketing tradizionale. Alcune cose da ricordare:
\begin{enumerate}
  \item \textbf{Il marketing non dice che cosa pubblicare, ma come pubblicarlo}.
        Esistono, nei videogiochi e nei giochi da tavolo, realtà che fanno
        prodotti \textbf{marketing driven}: l'esempio sono gli
        \textbf{hyper-casual game} per smartphone, quelli con il personaggio che
        va a destra e a sinistra e l'esercito che si allarga. Possono anche
        guadagnare, ma ``sono un po' la morte della creatività''. Il marketing
        ideale \textbf{prende un prodotto pensato da altri e gli dà le maggiori
        chance di essere venduto}: cerca il pubblico ideale per il prodotto che
        si ha, senza doverlo cambiare. Anzi, il mondo videoludico, e ludico in
        generale, ``è lastricato'' di errori di dipartimenti marketing che hanno
        cambiato completamente non il gioco, ma il modo in cui l'hanno
        pubblicizzato, e il gioco è fallito.
  \item \textbf{Usare strumenti del marketing non significa fare marketing}. Game
        design e game development a volte usano strumenti che vengono dal
        marketing (gli studenti ne vedranno con Beatrice, con Daniele e con
        Marta).
  \item \textbf{Il marketing è un gerundio}: si fa. È un'azione che si compie per
        comunicare; la comunicazione è attiva, e il marketing è l'insieme delle
        attività con cui si comunica il proprio gioco, dal \emph{copy} ai social.
        Può includere il digital marketing, ``ma non è solo quello''.
  \item \textbf{Il marketing esiste anche se si crede che non esista}. In Italia
        molti non vogliono ``fare marketing'' perché lo associano al lucro. Ma
        il marketing è \textbf{lo studio di come gli esseri umani comprano e
        vendono}, ed è innato in noi: ci sono logiche di cui siamo vittime, e
        ``anche se sei il bastian contrario, l'anticapitalista'', magari si
        compra comunque su Amazon perché fa comodo. Sono meccanismi che
        esistono, ``affascinanti anche, a volte''.
\end{enumerate}
Gli studenti non dovranno fare marketing, ma dovranno ricordare alcuni di questi
concetti, soprattutto nelle lezioni con Marta, che tratta argomenti legati al
marketing. E, chiude Daniele, come game designer e game developer ``per lo più
speriamo che il marketing non ci rovini i giochi'', soprattutto nel videoludico.

% ---------------------------------------------------------------------
\section{Game design e game development}
\label{a1:sec:design-development}

Ruggero torna sulla distinzione più importante della slide, ``proseguendo con
questa strada'':
\begin{description}
  \item[Game design] è chi \textbf{progetta l'esperienza} tramite la
        \textbf{creazione delle meccaniche} di gioco.
  \item[Game development] è chi fa \textbf{mediazione} tra la persona che scrive
        il game design e la \textbf{casa editrice}.
\end{description}

\subsection{I red flag dei designer}

Ruggero lo ripete perché lo vede ogni giorno nel suo lavoro: tra giochi
pubblicati e prototipi, ne valuta ``quasi uno al giorno, se faccio un conto
annuale'', e solo una minima percentuale arriva alla pubblicazione. Molte volte
i \textbf{red flag} vengono proprio da designer che non hanno presente questa
definizione, e che cominciano a mettere \textbf{clausole sulla non modifica del
gioco}: il gioco deve restare com'è. ``Questa cosa è impossibile con un
publisher, spesso''. Il publisher, per forza di cose, potrà chiedere di cambiare
la grafica, di reimpaginare il regolamento, di cambiare i materiali perché
troppo costosi: chiederà dei \textbf{compromessi}. E la persona che, all'interno
del publisher, fa questo lavoro è il \textbf{developer}.

Sta al designer capire \textbf{su che cosa del proprio gioco si può trovare un
compromesso e su che cosa no}. Publisher diversi lavorano in modo diverso, e il
designer ha sempre l'autorità di dire ``no, non voglio pubblicare questo gioco
se queste sono le modifiche''. Ma deve sapere chi fa il lavoro di development,
perché se il suo obiettivo è la \textbf{pubblicazione tradizionale} si
interfaccerà con questa figura.

\subsection{Daniele: una distinzione che manca al videogioco}

Daniele sottolinea che questa distinzione è \textbf{tipica del mondo analogico}
ed è ``molto interessante''. In Italia la parola \emph{developer} è diventata
negli anni quasi un sinonimo di programmatore di videogiochi, o di chiunque
lavori nel videoludico. Secondo lui c'entra anche il fatto che qualche anno fa
il sito \textbf{Gamasutra} ha cambiato nome in \textbf{Game Developer}, e la
parola è diventata un ombrello per tutte le professioni del videogioco. ``Mi fa
un po' ribrezzo'': il game developer dovrebbe essere un ruolo di design, un
designer che, come nell'analogico, fa \textbf{mediazione}. Nel videogioco i
publisher hanno in effetti i loro developer interni, che adattano il gioco, per
esempio alle console e ai loro kit di sviluppo. Ma la distinzione dovrebbe
esistere anche nel videogioco, perché è ``difficilissimo'' avere la mentalità
per fare le due cose insieme: chi fa design ha un focus che raramente si può
sovrapporre a quello di chi fa mediazione.

A Daniele piace la parola \emph{develop}, \textbf{sviluppare}, nel senso di
prendere qualcosa che è stato progettato e \textbf{renderlo possibile nel mondo
reale}. Ci sono progetti che hanno bisogno di uno sviluppo che ``polishi''
l'esperienza: la semplifichi, la migliori, la renda comprensibile e la concluda,
ma soprattutto la renda \textbf{stampabile} e \textbf{fruibile} per l'utente. La
metafora crossmediale è quella della \textbf{fotografia}: chi sviluppa le
fotografie non è chi scatta la foto, ma chi la rende visibile, anche
fisicamente. ``È un po' retro come metafora, ma ci capiamo''. ``È vero, però è
giusto'', conferma Ruggero.

% ---------------------------------------------------------------------
\section{Il mercato: Italia ed estero}
\label{a1:sec:mercato}

Il corso è in italiano e si svolge in Italia, ma bisogna conoscere anche il
\textbf{mercato estero}: molti giochi progettati da autori italiani vengono
pubblicati all'estero, e conoscere il mercato estero aiuta a capire come il
proprio gioco può arrivarci, ``o viceversa, come non può arrivarci''. La
tabella~\ref{a1:tab:mercato} riassume il confronto della slide.

\begin{table}[htbp]
\centering
\small
\begin{tabular}{@{}p{6.4cm}p{6.4cm}@{}}
\toprule
\textbf{Italia} & \textbf{Estero (soprattutto Europa)} \\
\midrule
il gioco è ancora visto come un'attività \textbf{infantile}, anche se la
normalizzazione è in crescita &
forte tradizione di gioco anche in un contesto \textbf{familiare adulto},
soprattutto ``dalla Germania in su'' \\
risposta \textbf{conservativa} sul tema ludico &
\textbf{premi di prestigio} riconosciuti in tutto il mondo (Spiel des Jahres,
As d'Or) \\
forte focus sulla \textbf{localizzazione} di giochi stranieri &
si cercano soprattutto giochi \textbf{originali} \\
investimento \textbf{molto basso} sulla creazione di giochi originali interni;
pochissimi designer assunti &
\textbf{developer interni} pagati con un contratto (Francia, Germania,
Danimarca\ldots) \\
\bottomrule
\end{tabular}
\caption{Il mercato del gioco analogico in Italia e all'estero, secondo la
slide di Ruggero e i commenti di Daniele.}
\label{a1:tab:mercato}
\end{table}

\subsection{L'Italia}

\begin{itemize}
  \item In Italia il gioco è ancora visto come un'\textbf{attività infantile},
        anche se la \textbf{normalizzazione} è in crescita. C'è ancora una forte
        risposta \textbf{conservativa} sul tema ludico.
  \item C'è un forte focus sulla \textbf{localizzazione} dei giochi stranieri: in
        una qualunque ludoteca o negozio, molti dei giochi da tavolo e di ruolo
        sono di publisher o designer esteri, pubblicati in italiano da publisher
        italiani. (La localizzazione, cioè portare un titolo dall'estero in
        Italia, esiste come funzione all'interno del publisher; Ruggero non l'ha
        messa nella slide delle figure perché non è il focus del corso.)
  \item C'è, ``purtroppo, preparatevi psicologicamente'', un \textbf{investimento
        molto basso} sulla creazione di giochi originali interni alla casa
        editrice. Pochissime case editrici italiane hanno designer interni pagati
        a tempo pieno con contratto a tempo indeterminato. Daniele lo dice senza
        mezzi termini: ``il mercato è molto piccolo, conosciamo praticamente
        tutto il settore'', e le persone che in Italia fanno game design
        analogico a tempo pieno con un contratto a tempo indeterminato ``sono
        meno delle dita delle mie due mani''. Di solito la formula è quella del
        \textbf{freelance}, del contratto, delle \textbf{royalty}. Chi va da un
        publisher aspettandosi di essere assunto a tempo pieno come game
        designer deve sapere che ``non siamo purtroppo ancora in quel tipo di
        situazione''.
  \item C'è però un contraltare: in Italia, soprattutto nel gioco di ruolo ma
        anche nel gioco da tavolo, c'è stata molta \textbf{imprenditoria}, anche
        grazie a \textbf{Kickstarter}, che ha aperto le porte ad altri sistemi
        di \textbf{crowdfunding}. È un altro argomento, perché richiede un team,
        un investimento e così via.
  \item E gli autori non mancano, anzi: le occasioni di incontro tra autori
        vedono ogni anno un aumento dei possibili autori\footnote{La trascrizione
        dice ``i DAG a Parma''; il riferimento è probabilmente a IDEAG, il
        circuito di incontri tra autori di giochi italiani, ma non è
        ricostruibile con certezza.}. Bisogna ``sgomitare'' tra tanti autori,
        che magari arrivano da esperienze molto diverse.
\end{itemize}
Ruggero precisa: l'investimento basso è ``dall'alto''. Le persone che fanno e
vogliono fare giochi sono tantissime e in aumento; scarseggia l'investimento,
sia dei publisher disposti a pagare stipendi sia dei \textbf{bandi} pubblici. I
bandi esistono (Beatrice ne parlerà per i giochi pervasivi, che spesso si
muovono in quella direzione, e Ruggero cita qualcuno che sta facendo un card
game con un bando italiano, ``una cosa assurda, ma lo sta facendo''), ma sono
una rarità, soprattutto rispetto all'estero.

\subsection{L'estero}

\begin{itemize}
  \item L'estero (dipende dalla nazione, ma in generale il resto d'Europa) ha una
        forte \textbf{tradizione di gioco} anche in un contesto \textbf{familiare
        adulto}, soprattutto ``dalla Germania in su''.
  \item Ha \textbf{premi di prestigio} riconosciuti in tutto il mondo, come lo
        \textbf{Spiel des Jahres} tedesco e l'\textbf{As d'Or} francese\footnote{La
        trascrizione dice ``lo Spield AR, la SDOR''.}.
  \item Cerca soprattutto \textbf{giochi originali}. Lo si vede dal fatto che
        molti autori italiani di giochi da tavolo pubblicano in Germania o in
        Francia e \textbf{non vengono pubblicati in Italia}, ``che è tutto
        dire''.
  \item Ha \textbf{figure di development interne}: non è raro vedere publisher
        francesi, tedeschi, danesi con developer pagati internamente, a
        contratto.
\end{itemize}
``Poi la situazione è ancora più drammatica'', ammette Ruggero, ``ma per ora
vogliamo tenere il morale un po' alto. È l'introduzione. Ci sarà tempo di andare
nella tristezza''.

Prima di proseguire, Ruggero insiste sulle domande: con lui si può interrompere
anche mentre parla, alzando la mano con le reaction o semplicemente togliendo il
microfono e chiedendo. ``Non mi offenderò mai, anzi, sarò felice di vedervelo
fare e mi ricorderò positivamente di voi''.

% ---------------------------------------------------------------------
\section{Perché serve una terminologia comune}
\label{a1:sec:terminologia}

La serie di slide successiva è ``una piccolissima parte di fondamenti di
design'' (``non spaventatevi, vi vedo che vi spaventate un po'\,''). Anche qui
molti temi verranno ripresi in dettaglio, ma l'idea è arrivare alla prossima
lezione con almeno un \textbf{glossario}, un \textbf{linguaggio condiviso}: se
qualcuno dice \emph{game developer}, tutti sanno di che cosa sta parlando.

Ruggero parte da una domanda, e precisa che non è retorica: \textbf{perché ci
serve una terminologia comune?} Gli studenti rispondono:
\begin{itemize}
  \item Gianmarco, ``molto pragmatico'': perché bisogna lavorare in team e
        capirsi. Vero, nel corso e anche fuori;
  \item un'altra studentessa: per la \textbf{comunicazione}, tra di noi e con
        tutti quelli che sono interessati nel settore. Vero anche questo:
        conoscere i termini comuni permette di parlare con chi fa game design o
        development nel dettaglio, nel tecnico, nello specifico.
\end{itemize}
Ruggero aggiunge altre due ragioni.

\paragraph{Saper comunicare il proprio gioco.} La prima riguarda la
\textbf{comunicazione del gioco} verso l'esterno. ``E lo dico non tirando acqua
al mio mulino, ma anzi togliendola'': \textbf{non è spesso il buon game design
a far diffondere un gioco}. ``Vorrei che lo fosse, sarebbe bellissimo, ma
purtroppo non è così''. Ci sono giochi bellissimi che non sanno comunicare, che
non sanno dire che cosa fanno, e che quindi è difficile, da fuori, trovare
invitanti; e ci sono giochi mediocri dal punto di vista del design che si
comunicano molto bene, e che quindi conoscono in molti. Saper parlare di quello
che si fa è utile \textbf{a sé stessi}, in primis.

\paragraph{Documentare per sé stessi.} La seconda è la \textbf{documentazione}:
documenti di design, di brainstorming e così via. ``Sembrerà una battuta, ma
provate a recuperare un documento di brainstorming una settimana dopo, se non
avete scritto bene le cose che avete scritto''. Chi l'ha fatto sa già di che
cosa si parla; ``chi non l'ha ancora vissuto, forse lo vivrà''.

% ---------------------------------------------------------------------
\section{Game e play}
\label{a1:sec:game-play}

La \textbf{prima, importantissima divisione terminologica} è quella tra
\emph{game} e \emph{play}. In italiano c'è un problema: \textbf{gioco} e
\textbf{giocare} hanno la stessa radice, mentre in altre lingue, tra cui
l'inglese, sono due parole strutturalmente diverse:
\begin{description}
  \item[Game] il \textbf{gioco come artefatto}: il gioco come ecosistema di
        meccaniche, come struttura di regole, come prodotto (``tra
        virgolette''), come oggetto fisico.
  \item[Play] l'\textbf{atto}, il verbo: le persone che giocano, l'approccio
        fisico all'attività del gioco.
\end{description}
È importante perché gli studenti faranno \textbf{game design}; esistono anche
ambiti di \emph{play design}, ma qui ``creerete giochi, non giocattoli''.

\subsection{Giochi e giocattoli}

Qual è la differenza? Uno studente risponde: \textbf{la presenza di regole}.
Esatto: un gioco ha delle regole, un'attività di \emph{play} non
necessariamente. Con un giocattolo si può fare \emph{play}: c'è chi gioca con i
Lego, chi gioca a guardia e ladri o a nascondino al parco, e magari alcune di
quelle regole sono ``più o meno semi-solidificate''. Ma il \emph{play} permette
spesso un approccio non strettamente legato alle regole.

Viceversa, \textbf{un game può esistere anche se nessuno lo sta giocando}. Se si
prepara su un tavolo una scacchiera con i pezzi degli scacchi, chi passa e la
vede riconosce il gioco degli scacchi, anche se nessuno sta giocando. Se invece
si vedono dei bambini correre in un parco giochi, quei bambini \emph{stanno
giocando}, ma non stanno giocando \emph{a un gioco} specifico. ``È importante
separare questa cosa''; ci si tornerà, ``ma per ora abituatevi a questa
differenza''.

\paragraph{Approfondimento: giocattoli, giochi e regole.} La stessa distinzione
è il punto di partenza del capitolo che Jesse Schell, in \emph{The Art of Game
Design}, dedica alla definizione di gioco (\emph{The Experience Rises Out of a
Game}). Schell nota che usiamo perfino un linguaggio diverso: ``un gioco è
qualcosa a cui si gioca'' (\emph{a game is something you play}), ``un giocattolo
è qualcosa con cui si gioca'' (\emph{a toy is something you play with}); e
arriva a definire il giocattolo come ``un oggetto con cui si gioca''. Analizzando
poi la definizione di gioco di Avedon e Sutton-Smith, ne ricava che ``i giochi
hanno regole, i giocattoli no'', e che le regole sono uno degli aspetti
\emph{definitori} del gioco, insieme alla volontarietà, agli obiettivi, al
conflitto e alla possibilità di vincere o perdere. Schell osserva anche che
molti giochi sono costruiti \textbf{sopra} un giocattolo (``una palla è un
giocattolo, ma il baseball è un gioco''), e tra i suoi consigli di
prototipazione c'è proprio \emph{Build the Toy First}: costruire prima il
giocattolo, cioè qualcosa che sia già piacevole da manipolare, e poi il gioco
intorno. È una lettura utile della distinzione di Ruggero: il \emph{play} può
fare a meno delle regole, il \emph{game} no.

% ---------------------------------------------------------------------
\section{Il framework MDA (e il DPE)}
\label{a1:sec:mda}

\subsection{Meccaniche, dinamiche, estetiche}

Un'altra terminologia fondamentale è quella del framework \textbf{MDA}
(\emph{Mechanics, Dynamics, Aesthetics}; in italiano a volte ``MDE''). Viene da
un \textbf{paper} nato nell'ambito del videogioco\footnote{È il paper di Robin
Hunicke, Marc LeBlanc e Robert Zubek, \emph{MDA: A Formal Approach to Game
Design and Game Research}, del 2004. Daniele aggiunge il paper alla cartella
condivisa durante la lezione.}, che ha ``solidificato'' queste terminologie, e
divide un gioco in tre aree:
\begin{description}
  \item[Meccaniche] le \textbf{regole per come sono scritte}. ``Ogni volta che
        pescate un 2 fate dei punti'': quella è una meccanica.
  \item[Dinamiche] l'\textbf{applicazione in tempo reale} delle meccaniche. La
        dinamica è il giocatore che, pescando una carta, \textbf{spera} che esca
        un 2. Quel processo non è una regola: ``il gioco non vi dice che dovete
        sperare che esca un 2''.
  \item[Estetiche] l'insieme di quello che il gioco vuole comunicare,
        l'\textbf{esperienziale}.
\end{description}

L'esempio di Ruggero mette insieme le tre aree. In un gioco in cui pescare un 2
fa fare punti (meccanica) e in cui quindi si spera che esca un 2 (dinamica),
\textbf{che estetica si ottiene se nel mazzo di 100 carte c'è un solo 2? E se ce
ne sono 50?} L'emozione di aver pescato \emph{l'unico} 2 del mazzo è
completamente diversa da quella di pescarne uno sapendo che ce ne sono
cinquanta e che è molto probabile trovarlo. Le regole sono quasi identiche; cambia
l'esperienza.

\subsection{Il falso amico: ``estetica''}

Daniele avverte di non sottovalutare la parola \textbf{estetica}, che in italiano
è un \textbf{falso amico}: sembra indicare la parte grafica, e invece non lo è.
Si incontreranno designer che hanno studiato MDA e altri che non l'hanno
studiato, e lo si capirà proprio da questo: i secondi diranno che le
\emph{aesthetics} sono la grafica. Le \emph{aesthetics} sono la parte
\textbf{esperienziale}: come nell'esempio di Ruggero, quello che lascia al
giocatore aver pescato l'unico 2 del mazzo.

\subsection{Il DPE: design, play, experience}

Per gli italiani, aggiunge Daniele, è più facile capire l'MDA attraverso un
altro paper, quello del \textbf{DPE}\footnote{È il framework \emph{Design, Play,
Experience} di Brian Winn, proposto come estensione dell'MDA per i giochi
educativi (\emph{serious games}). Anche questo paper viene aggiunto da Daniele
alla cartella condivisa.}, che di fatto fa un parallelismo assegnando
\textbf{design, play, experience}:
\begin{itemize}
  \item le \textbf{meccaniche} sono il lato \textbf{design};
  \item le \textbf{dinamiche} sono il lato \textbf{play};
  \item le \textbf{estetiche} sono l'\textbf{experience}, l'esperienza che riceve
        chi gioca.
\end{itemize}
``Se vivete le \emph{aesthetics} come l'esperienza che vive chi sta giocando,
state a postissimo''. In questo momento Beatrice lascia la lezione (``ciao,
Bea''). I paper di MDA e DPE sono in inglese; chi vuole può dargli una lettura.

\begin{figure}[htbp]
\centering
\begin{tikzpicture}[
    font=\small,
    blocco/.style={draw, thick, rounded corners=3pt, fill=gray!12,
                   minimum width=3.2cm, minimum height=1.0cm, align=center},
    dpe/.style={font=\footnotesize\itshape, align=center},
    esempio/.style={draw, rounded corners=2pt, text width=3.3cm,
                    align=left, font=\footnotesize, inner sep=4pt},
    freccia/.style={-{Stealth[length=2.5mm]}, thick}
  ]
  \node[blocco] (m) at (0,0)   {\textbf{Meccaniche}\\ \footnotesize le regole scritte};
  \node[blocco] (d) at (4.3,0) {\textbf{Dinamiche}\\ \footnotesize le regole in azione};
  \node[blocco] (a) at (8.6,0) {\textbf{Estetiche}\\ \footnotesize l'esperienza};
  \draw[freccia] (m) -- (d);
  \draw[freccia] (d) -- (a);
  \node[dpe, above=0.15cm of m] {DPE: \textbf{design}};
  \node[dpe, above=0.15cm of d] {DPE: \textbf{play}};
  \node[dpe, above=0.15cm of a] {DPE: \textbf{experience}};
  \node[esempio, below=0.45cm of m] (em)
    {``Ogni volta che peschi un 2 fai punti''};
  \node[esempio, below=0.45cm of d] (ed)
    {pescando, il giocatore \emph{spera} che esca un 2 (nessuna regola glielo
     chiede)};
  \node[esempio, below=0.45cm of a] (ea)
    {un solo 2 su 100 carte: pescarlo è un evento;\\
     50 su 100: è quasi scontato};
  \draw[dashed] (m) -- (em);
  \draw[dashed] (d) -- (ed);
  \draw[dashed] (a) -- (ea);
  % prospettive
  \draw[freccia, gray] ([yshift=-0.3cm]em.south west) --
    node[below, font=\footnotesize\itshape, text=black]
    {il designer agisce sulle meccaniche}
    ([yshift=-0.3cm]ea.south east);
  \draw[freccia, gray] ([yshift=-1.2cm]ea.south east) --
    node[below, font=\footnotesize\itshape, text=black]
    {chi gioca incontra prima l'esperienza}
    ([yshift=-1.2cm]em.south west);
\end{tikzpicture}
\caption{Il framework MDA con l'esempio del ``2'' usato da Ruggero, e la
corrispondenza con il DPE suggerita da Daniele. Le due frecce in basso
riassumono la lettura classica del paper: il designer costruisce le meccaniche
per ottenere un'esperienza, mentre il giocatore parte dall'esperienza.}
\label{a1:fig:mda}
\end{figure}

\paragraph{Approfondimento: le due direzioni dell'MDA e le otto estetiche.} Le
frecce della figura~\ref{a1:fig:mda} riprendono un'idea centrale del paper
originale, che la lezione non esplicita: designer e giocatore guardano il gioco
da lati opposti. In \emph{Challenges for Game Designers}, Brenda Brathwaite e
Ian Schreiber la riassumono così: i designer si chiedono quale estetica
vogliono ottenere, definiscono le dinamiche che porterebbero a quella
sensazione, e poi creano le meccaniche che producono quelle dinamiche. Lo stesso
libro conferma l'avvertenza di Daniele sul falso amico: ``le \emph{aesthetics}
non si riferiscono all'aspetto del gioco, ma alla risposta emotiva'' che il
designer vuole suscitare. E offre esempi di dinamiche nel senso di Ruggero: in
uno sparatutto i punti di \emph{spawn} sono una meccanica, mentre appostarsi
accanto a uno di essi per colpire chi riappare (lo \emph{spawn camping}) è una
dinamica; negli scacchi i movimenti dei pezzi sono meccaniche, mentre le
aperture codificate sono dinamiche. Per dare un nome alle estetiche, il paper
MDA propone una tassonomia, poi detta degli ``otto tipi di \emph{fun}'', che
Schell riprende in \emph{The Art of Game Design} come \emph{LeBlanc's Taxonomy of
Game Pleasures}: \textbf{sensazione}, \textbf{fantasia}, \textbf{narrativa},
\textbf{sfida}, \textbf{compagnia} (\emph{fellowship}), \textbf{scoperta},
\textbf{espressione} e \textbf{sottomissione} (il piacere di entrare nel
``cerchio magico'' e lasciarsi andare a un passatempo). È un primo esempio di
quello che i docenti chiedono con le parole vietate della prossima sezione:
invece di dire che un gioco è ``divertente'', dire \emph{quale} esperienza
offre.

% ---------------------------------------------------------------------
\section{Le parole vietate}
\label{a1:sec:parole-vietate}

``Attenzione, adesso qua triggeriamo delle persone''. Ruggero annuncia una regola
che ``tornerà'': nel corso ci sono delle \textbf{parole vietate}. ``Non perché
siamo dei censori''; il perché si vedrà in una lezione apposita, perché
spiegarlo adesso richiederebbe un'ora. Ma Daniele vuole spiegarne ``un
minimo''.

\subsection{La regola}

In \textbf{nessun materiale scritto} consegnato ai docenti, task compresi, sono
accettate queste parole, né in inglese né in italiano:
\begin{center}
\begin{tabular}{@{}ll@{}}
\toprule
\textbf{Inglese} & \textbf{Italiano} \\
\midrule
\emph{fun} & divertimento \\
\emph{immersion} & immersione \\
\emph{entertainment} & intrattenimento \\
\bottomrule
\end{tabular}
\end{center}
``Non li voglio neanche vedere''. Daniele lo ripete con un paragone: ``se mi fate
i disegni con Paint, benissimo, top, fantastico, siamo ancora best friend. Se mi
scrivete \emph{divertimento}, io sbrocco''. È una \textbf{regola dell'accademia}:
negli stessi corsi triennali di narrative design e game design le parole sono
vietate agli studenti per tre anni; ``per voi solo per tre mesi''. Daniele
aggiungerebbe volentieri un quarto termine, \textbf{gamification}.

\subsection{Perché}

Il motivo approfondito lo darà Ruggero nella lezione apposita. Daniele dà il
suo. Ruggero e Diego Ricchiuti sono molto contrari a \emph{fun} e a
\emph{immersion}; lui ha aggiunto \emph{intrattenimento}. Sono \textbf{parole
troppo passe-partout e troppo ambigue}:
\begin{itemize}
  \item \textbf{intrattenimento} in Italia è stata usata per anni come un modo
        per \textbf{sminuire}: ``ma sì, tanto è un gioco solo di
        intrattenimento, è un romanzo solo di intrattenimento''. Ma anche se un
        prodotto è ``solo di intrattenimento'', non vuol dire che non ci si debba
        lavorare con attenzione. E poi intrattenimento ``non significa niente'':
        ci sono tanti modi in cui una cosa può intrattenere;
  \item lo stesso vale per \textbf{fun}, e per \textbf{immersion}, che ha un
        problema in più: è una parola passe-partout per cui si può dire che un
        gioco è ``immersivo'' nel senso generale del termine, ma poi non si sa
        davvero che cosa voglia dire, \emph{immersivo in che senso}.
\end{itemize}

L'obiettivo didattico è che gli studenti \textbf{vadano oltre queste parole}:
``un po' come quando alle medie vi dicono smettetela di dire che tutto è bello,
perché ci sono vari gradi di bello, ci sono altri modi di dire bello''. Bisogna
\textbf{decostruirle}: quando si sta per usarle, chiedersi \textbf{che cosa si
intende} con quella parola e cercarne una \textbf{più precisa}.

Ruggero chiede, in uno dei ``rari casi'', di \textbf{non fare domande adesso} su
questa regola: ogni anno arrivano domande e commenti, e ci saranno lezioni
apposite di tre ore, con strumenti per capire il perché. Ma da subito gli
studenti devono abituarsi a evitare questi termini, sia perché si spera che
abbiano termini più precisi, sia perché queste parole ``portano con sé dei
problemi agganciati'' che si vedranno. ``Non spaventatevi, non verrete sparati in
un vicolo cieco di notte''.

\paragraph{Una piccola gamification.} Il divieto va visto come uno
\textbf{stimolo}, ``un po' come se stessimo gamificando la vostra esperienza'',
una sorta di user experience delle lezioni. Imponendosi di non usare queste
parole, si noterà come migliora la capacità di \textbf{descrivere i giochi} e le
proprie esperienze ludiche, e di essere più precisi. È un invito a ``obbligarsi a
essere più precisi e ad avere una mentalità più critica''. La lezione di
Ruggero sull'\textbf{involvement}, promette Daniele, sarà ``fenomenale'' e
``vi aprirà la testa''.

\paragraph{Approfondimento: anche i libri faticano a definire il ``fun''.} La
difficoltà che il divieto vuole far emergere è visibile anche nei testi di game
design. Schell, in \emph{The Art of Game Design}, prova a definire il \emph{fun}
e propone ``piacere con sorprese'' (\emph{pleasure with surprises}), ammettendo
subito che è strano poter usare una parola per tutta la vita, sapendo con
certezza che cosa significa, senza riuscire a esprimerla chiaramente quando
qualcuno lo chiede; e invita a mettere alla prova la definizione con
controesempi. La sua lente sul \emph{fun} (\emph{The Lens of Fun}) si limita,
non a caso, a chiedere quali parti del gioco siano \emph{fun} e \emph{perché}.
Anche \emph{immersion} ha molte letture: in \emph{21st Century Game Design},
Bateman e Boon la descrivono, tra l'altro, come la capacità di ``perdersi''
nel mondo artificiale presentato dal gioco, legata alla sospensione
dell'incredulità e ai giochi di mimesi. Sono significati diversi sotto la stessa
parola: proprio il motivo per cui, in un documento di design, conviene dire
quale si intende.

% ---------------------------------------------------------------------
\section{Chi crea, chi gioca, chi facilita}
\label{a1:sec:ruoli-esperienza}

Un'altra distinzione importante è quella tra \textbf{chi crea}, \textbf{chi
gioca} e \textbf{chi facilita}. Spesso, quando si fa un gioco, lo si fa provare
di persona a chi gioca. Ma ``la norma dei giochi pubblicati non funziona così'':
quando si pubblica un gioco non si vede chi lo compra e chi lo gioca, e non si
sa se qualcuno, a casa propria con i propri amici, segue le regole alla lettera e
interpreta correttamente quello che il designer aveva descritto. Quindi già in
fase di progettazione bisogna tenere a mente \textbf{chi sta facendo game design
e chi sta giocando}.

E nei giochi a evento, nei giochi pervasivi, nei giochi di ruolo da tavolo e dal
vivo (LARP) c'è un terzo pezzo del puzzle, ``importantissimo'': \textbf{chi
facilita}. La domanda da farsi è: \textbf{il regolamento è scritto per chi gioca
o per chi deve facilitare il gioco?} Bisogna abituarsi a pensare a tutte le
persone coinvolte nell'esperienza.

% ---------------------------------------------------------------------
\section{Playtest e iterazione}
\label{a1:sec:playtest}

\subsection{Nessun gioco nasce perfetto}

L'ultima voce del glossario è ``un'altra cosa importantissima'':
\textbf{playtest} e \textbf{iterazione}. Il corso è \textbf{pragmatico}:
``specialmente da parte mia'', dice Ruggero, non si vedranno teorie ``messe lì
per fare accademia sul game design nelle nostre torri d'avorio''. Ci saranno
teorie e framework, ma tutti estremamente pragmatici, perché l'obiettivo non è
solo capire i giochi, ma \textbf{saperli fare e iterare}.

Playtest e iterazione sono \textbf{inevitabili} nel game design, perché
\textbf{nessun gioco nasce perfetto}, soprattutto perché le persone sono
diverse: un gioco perfetto per una persona può essere incomprensibile per
un'altra. Chi progetta giochi per altre persone deve capire, dopo averli fatti
provare, \textbf{che cosa cambiare, che cosa sistemare, che cosa aggiungere, che
cosa togliere}. Ci saranno lezioni su playtest, prototipazione, raccolta dei
feedback, metodi etnografici. Ma bisogna abituarsi già ora all'idea che
\textbf{non esiste la consegna definitiva}. A un certo punto i task vanno
consegnati per forza di cose, ma, come diceva Daniele all'inizio, alla fine del
corso gli studenti avranno un \textbf{prototipo}, ``che sarà un prototipo'', e
gli strumenti per portare avanti il lavoro anche dopo.

\subsection{L'aiuto dopo il corso}

Non solo gli strumenti: anche l'aiuto dei docenti. Daniele racconta di aver
portato a un editore, a Play\footnote{Probabilmente \emph{Play -- Festival del
Gioco} di Modena; la trascrizione non permette di esserne certi.}, uno dei
progetti realizzati l'anno precedente da due studenti, un gioco di ruolo in
solitario (``un \emph{solo game} RPG molto, molto bello''). Se gli studenti
vorranno trasformare il proprio lavoro in un prototipo vero e proprio, possono
contattarlo: lo fa ``per lavoro, ma anche per passione''. Gestisce anche una
community per aiutare i designer ``un po' depressi'' a rimettersi a progettare.

Quindi: ``non sentitevi persi se durante la fase di iterazione vi sentite
persi, è normale''.

\subsection{Il dizionario comune e il production dictionary}

Daniele riprende un punto detto ``benissimo'' da Ruggero e che secondo lui non è
stato sottolineato abbastanza: il \textbf{dizionario comune} serve anche a
\textbf{non perdersi} mentre si itera. Se da un'iterazione all'altra i
documenti usano termini completamente diversi, e non c'è un dizionario di base,
ci si perde.

Nel videogioco, una delle primissime cose che un designer fa quando entra in
un'azienda è il \textbf{production dictionary}: si annota, in un file privato,
tutte le parole che usano i colleghi, per allinearsi al dizionario comune
dell'azienda, perché lì su un prototipo possono lavorare anche cento persone.
Nel gioco da tavolo non si fa molto, perché il mondo dell'analogico è
``talmente più piccolo''; ma serve comunque imparare un \textbf{dizionario
minimo}, dai docenti e dai professionisti. E il corso ne darà una parte,
entrando nel dettaglio di alcuni termini e framework.

\paragraph{L'esempio del feedback.} Perché serva, lo mostra un esempio. Si
testa un gioco e ci si annota come feedback: ``questo gioco deve essere più
divertente''. Che cosa vuol dire? Come lo si migliora? Quella frase non dà
elementi utili. Se invece si scrive che deve essere \textbf{più adrenalinico},
\textbf{più frenetico}, \textbf{più rilassato}, \textbf{più cozy}, o
\textbf{più bilanciato matematicamente}, si sa dove andare a lavorare.

\paragraph{Approfondimento: la regola del ciclo.} Il principio ``nessun gioco
nasce perfetto'' ha un nome preciso in \emph{The Art of Game Design}: è la
\textbf{Rule of the Loop}, la regola del ciclo. ``Più volte provi e migliori il
tuo design, migliore sarà il tuo gioco''. Schell precisa che non è una lente,
cioè un punto di vista, ma ``una verità assoluta'', senza eccezioni: ci saranno
momenti in cui si cercherà di convincersi che ``questa volta il design è così
buono che non serve testarlo'', e ogni volta se ne pagherà il prezzo. Ed è
interessante la ragione per cui la regola pesa di più sul digitale: secondo
Schell il tempo e il denaro necessari per provare e correggere un sistema sono
molto maggiori nei videogiochi che nei giochi tradizionali. Per un gioco di
carte o da tavolo, con il tempo per provarlo e cambiarlo più e più volte, il
consiglio è semplicemente di farlo. Il gioco analogico è, da questo punto di
vista, la palestra ideale dell'iterazione.

% ---------------------------------------------------------------------
\section{Tabletop Simulator}
\label{a1:sec:tts}

\subsection{Il software di settore}

``Last but not least'': il corso usa un software, \textbf{Tabletop Simulator},
che si compra su \textbf{Steam} e costa tra i \textbf{10 e i 20 euro} (``aspettate
che vada in sconto, va in sconto ogni due per tre''). Daniele spiega perché:
\begin{itemize}
  \item è \textbf{il software di settore} del gioco da tavolo, quello che i
        publisher usano per i prodotti digitali, interni ed esterni;
  \item essendo su Steam, permette di fare \textbf{playtest} e di giocare con gli
        amici di Steam. Daniele stesso sta testando lì il proprio gioco;
  \item funziona anche su Mac: Windows o Mac, nessun problema.
\end{itemize}
Ci sarà una lezione apposita, di Ruggero, sulla creazione del proprio
\textbf{modulo} di gioco.

\paragraph{Che cosa non è.} Daniele chiarisce subito un ``fraintendimento
gigantesco'', soprattutto per chi ha già esperienza: Tabletop Simulator
\textbf{non richiede programmazione} e \textbf{non è un software per
trasformare i giochi da tavolo in videogiochi}. È un software che
\textbf{simula dei componenti fisici su un tavolo}: serve per giocare i board
game in digitale. Simula anche la fisica: c'è perfino il comando per
\textbf{ribaltare il tavolo} (``non usatelo ora'').

\paragraph{La partnership mancata.} Da anni Daniele prova a fare una
partnership con gli sviluppatori: due anni prima ha contattato l'azienda per
acquistare delle chiavi per gli studenti, dato che il software si usa nei corsi.
L'azienda ha risposto che non fa accordi commerciali. ``Credo sia la prima
azienda nella mia vita che dice che non vuole che io le dia dei soldi''. Gli
studenti dovranno quindi provare ad acquistarlo da sé; se i docenti riusciranno a
darglielo, bene.

\paragraph{La domanda di Gianmarco.} Gianmarco segue le lezioni da un computer
aziendale ``super chiuso'', su cui non può installare Steam: deve usare Tabletop
Simulator durante le lezioni? Ruggero risponde che per la lezione dedicata è
meglio poter seguire e rifare in diretta. Ma il software non sarà richiesto
necessariamente per il project work: un gruppo potrebbe fare un gioco pervasivo
o un gioco di ruolo, non un gioco da tavolo. Essendo un software di settore, è
comunque bene conoscerlo. Ruggero aggiunge che le versioni non ufficiali che
circolano non permettono di giocare in molti, e quindi non servono per
playtestare o giocare insieme.

\subsection{Una dimostrazione}

Daniele mostra il software aperto da Steam:
\begin{itemize}
  \item si può creare una partita \textbf{single player} o \textbf{multiplayer};
        per giocare tra compagni si crea un server multiplayer e ci si aggiunge
        come amici su Steam;
  \item si dà al server un nome, un numero di giocatori, eventualmente una
        password; creato il server, si apre la scelta dei giochi;
  \item ci sono giochi gratuiti, giochi da acquistare (DLC, non obbligatori) e
        soprattutto un \textbf{Workshop} ``ampissimo'', con decine di migliaia di
        giochi caricati dagli utenti: ``se c'è un board game che conoscete,
        quel board game è su Tabletop Simulator, giocabile dal Workshop''.
        L'esempio mostrato è l'open beta ufficiale del gioco da tavolo di
        \emph{Slay the Spire}\footnote{La trascrizione dice ``Slade Spire''.}:
        molti publisher, per farsi pubblicità, mettono il proprio gioco in demo o
        giocabile in digitale;
  \item c'è un \textbf{tutorial} iniziale che spiega come muoversi sul tavolo;
        con \emph{Esc} e \emph{Help} si trovano i comandi da tastiera.
\end{itemize}

\subsection{Il primo task}

Il primo (e per ora unico) task è questo: \textbf{entro il 10 ottobre}, giorno
della lezione di Ruggero su Tabletop Simulator, giocare \textbf{tre o quattro
volte in totale}, online, tra compagni di corso o con amici. Non si tratta di
imparare a creare moduli, ma di prendere la mano \textbf{lato giocatore}: aprire
il software, impararne i comandi, giocare. ``Per me potete anche andare a
prendere domino o mahjong'' (``mahjong è bellissimo tra l'altro'', aggiunge
Daniele, che sta facendo un videogioco sul mahjong), o il pachisi o il go. Gli
anni precedenti si chiedeva a ciascuno di proporre ai compagni un gioco da
tavolo che apprezzava. È anche un modo per conoscersi. Ruggero spiega il
perché: usarlo per giocare è facile e intuitivo, usarlo per creare moduli è più
complesso; ma chi arriva al 10 ottobre avendo già aperto il software e giocato,
troverà più intuitivo anche creare un modulo. E l'avvertimento: ``non riducetevi
all'ultimo, non fate una partita il 9 sera''.

\subsection{Tabletopia e altri vantaggi}

Tancredi racconta di aver provato a caricare un suo gioco su
\textbf{Tabletopia}. ``Questo è un vaso di Pandora'', risponde Daniele.
Tabletopia è molto usato, ma secondo lui il suo strumento di creazione dei moduli
è pessimo, ed è il motivo per cui la maggioranza dei moduli ufficiali dei
publisher è su Tabletop Simulator. Tabletopia ha un vantaggio vero, è
\textbf{gratuito}, ma ha scelte di ``qualità della vita'' poco felici.

Tabletop Simulator ha invece un vantaggio ``fenomenale'' che si tende a
sottovalutare: essendo su Steam, \textbf{il modulo si pubblica direttamente sul
Workshop di Steam}; si possono invitare gli amici a iscriversi e giocare con
loro, purché abbiano un account Steam. Rispetto ad altri software da imparare,
costa ``tra i 10 e i 20 euro, ragazzi: si parla proprio di spiccioli''. ``Io con
Tabletop Simulator ci ho playtestato la vita''.

Altre caratteristiche utili:
\begin{itemize}
  \item si possono creare \textbf{moduli non elencati} (\emph{unlisted}),
        perfetti da passare a un publisher prima che il gioco sia pubblicato:
        non sono pubblici, e nessuno può rubare i materiali;
  \item il software ha una grafica datata (``\emph{World Art 1998}'', scherza
        uno dei docenti), ma lo stanno portando su \textbf{Unity} con molti investimenti
        e aggiornamenti; ascolta la comunità e introduce funzioni utili anche per
        chi fa design;
  \item il suo server Discord è molto frequentato, e c'è un canale YouTube
        dedicato solo al design di giochi di carte\footnote{La trascrizione ne
        riporta il nome come ``Distraction Maker''; non è verificabile.} con un
        server Discord dove si fanno playtest: il 99\% è su Tabletop Simulator;
  \item molti editori, quando si propone loro un gioco da tavolo,
        \textbf{chiedono se c'è un modulo per Tabletop Simulator}.
\end{itemize}

% ---------------------------------------------------------------------
\section{Le domande a Daniele: il tema del project work}
\label{a1:sec:domande-daniele}

Daniele ha ``un minuto'' per rispondere alle domande sul corso prima di dover
andare. Uno studente chiede, ``giusto per curiosità'', se il project work avrà
un tema o sarà libero.

\begin{itemize}
  \item C'è un \textbf{mandate}, che i docenti propongono ogni anno e
        aggiornano leggermente; verrà presentato il \textbf{25 settembre}.
        ``Non facciamo spoiler''. Conterrà quello che si chiede di fare e quello
        che non si chiede: dei \textbf{must}, dei \textbf{must not} e dei
        \textbf{nice to have}.
  \item Verrà chiesto di lavorare su un'\textbf{idea originale} che non si è
        portata al corso: originale nel senso di \textbf{creata apposta}, non
        nel senso di ``mai vista prima''.
  \item L'idea va \textbf{partorita insieme}. Altri anni alcuni studenti hanno
        provato a imporre la propria idea al gruppo, ``e l'esperienza è stata
        pessima per tutti''. Quindi: \textbf{non imponete le vostre idee, fate
        brainstorming con gli altri}. Daniele, che fa ``molto spesso
        brainstorming'', darà strumenti per farlo: nessuno arriverà impreparato
        a quella lezione.
\end{itemize}

Ruggero aggiunge che il mandate è \textbf{abbastanza flessibile} perché tutti
possano fare qualcosa che li interessi: ``l'ultima cosa che vogliamo è forzarvi a
lavorare su qualcosa che non vi piace''. Bisogna lavorare su ciò che piace e per
cui si ha uno stimolo personale, anche nella \textbf{tipologia} di gioco: nessuno
sarà obbligato a fare un gioco di ruolo o un board game, e la lezione sui casi
studio mostrerà quante forme diverse di gioco analogico il corso accetta.

L'unica cosa da \textbf{non fare}, ``perché è contraddittorio per questo corso'',
sono i \textbf{giochi grandi}: ``lo chiediamo per voi''. Niente progetti
mastodontici\footnote{La trascrizione dice ``non fate delle cose da gnome'':
probabilmente il nome di un gioco da tavolo molto grande e complesso, non
ricostruibile con certezza.}, niente ``cose alla \emph{Dungeons \& Dragons}''.
Si può fare un gioco di ruolo, ma con un'esperienza ``meno corposa'': ci sono
tanti giochi di ruolo da 15 o 20 pagine, ``bellissimi, fichissimi'', che si
possono leggere online. Il consiglio è farsi una cultura anche su
\textbf{itch.io}, ``uno dei portali più importanti del gioco di ruolo
internazionale'', di cui si parlerà anche nella parte di ricerca e sviluppo.

Poi Daniele saluta (``noi ci vediamo venerdì a lezione'') e lascia la
conduzione a Ruggero.

% ---------------------------------------------------------------------
\section{Le domande finali con Ruggero}
\label{a1:sec:domande}

\subsection{Le regole delle domande}

Ruggero spiega che cosa vuol dire ``stare con lui''. Alla fine di ogni sua
lezione ci saranno \textbf{almeno 20--30 minuti di domande} libere, \emph{free
form}. Vuole che fare domande sia \textbf{normalizzato}: ``io gestisco malissimo
l'\emph{horror vacui} e divento super cringe molto alla svelta se non ci sono
domande in giro''. Sono benvenute le domande sulla lezione appena fatta, ma non
solo: il corso è crossmediale, e anche le domande che vengono da
\textbf{divagazioni}, da parti laterali, ``fatele''. Quello che sembra scontato o
banale a uno studente può non esserlo per il docente, e viceversa; e può essere
utile anche agli altri. ``Se vi sembrano stupide, fatele il doppio''. ``Magari
sembro incazzato, ma vi giuro sono rilassatissimo. Non mordo, magari abbaio un
po'\,''.

\subsection{Come approcciarsi a Steam e Tabletop Simulator}

Rita non ha mai usato né Steam né Tabletop Simulator e chiede come
approcciarsi, e se c'è modo di mettersi d'accordo tra compagni su Discord, magari
con un sondaggio sui giochi, per formare gruppi di gioco. Ruggero conferma: sul
Discord dell'accademia c'è un \textbf{canale apposito per il corso}, ed è degli
studenti. Lo possono usare per organizzarsi a giocare, per chiedere a un
compagno di prestare un gioco, per domande tecniche (come si scarica, come si
installa). ``Sarebbe proprio ideale che vi parlaste e vi organizzaste''.

Anche perché, aggiunge con un ``piccolo spoiler'' (``nessuno vi ha detto
niente''), se fossero i docenti a organizzare i gruppi di gioco, il task lo
starebbero facendo loro. Bisogna \textbf{abituarsi a organizzare le sessioni di
gioco}, perché arriveranno i task sul \textbf{playtest} e su come organizzarlo,
e a quel punto si sarà ``meno disperati''. Sembra scontato, ma saper
\textbf{organizzarsi per giocare in autonomia} ``vi sarà utilissimo come skill in
questo settore, ve lo giuro''.

\subsection{Il bilanciamento}

Angelo confessa ``una paranoia grossissima'': ci saranno lezioni sul
\textbf{bilanciamento}? ``Ma certo''. Sono proprio le lezioni di Marta sul
development: il bilanciamento è di solito un tema che tocca \textbf{i developer
più dei designer}. L'anno precedente c'era una sola lezione con Marta, e i
docenti si sono resi conto che serviva andare ancora più a fondo: per questo
quest'anno le lezioni sono due. ``Non è una paranoia, è una legittima domanda'':
Ruggero insegna da anni, vede altri corsi, parla con altri docenti e studenti, e
``mi incazzo spesso'' perché molti corsi di game design \textbf{non fanno
bilanciamento}. I docenti volevano invece un corso che coprisse tutte le parti
del game design: non solo l'idea del gioco, ma il bilanciamento, la produzione,
i playtest, l'iterazione, la scrittura dei regolamenti.

\subsection{Si può giocare con i docenti?}

Una studentessa chiede se si può giocare con i docenti. Per motivi etici,
durante il corso questo è limitato alle parti di \textbf{valutazione} dei giochi
e di \textbf{tutoraggio}. Dopo il corso, ``nessun problema, almeno con me''.

\subsection{Dove si trovano i giochi su Tabletop Simulator}

Ruggero mostra dove si trovano i giochi: nella propria libreria di Steam, nella
pagina di Tabletop Simulator, alla voce \textbf{Workshop}, dove ci sono ``più o
meno 36.000'' giochi. Una studentessa scopre intanto che l'app di Steam per
telefono è solo una \emph{companion app}: per registrarsi e giocare serve il
PC.

Per provare, chiede agli studenti un titolo:
\begin{itemize}
  \item \emph{Keep the Heroes Out}: non compare, ``non tutti i giochi ci sono,
        perché molti publisher non li mettono''; forse, essendoci già una
        versione del gioco su Steam, il modulo è stato tolto;
  \item \emph{Magic: The Gathering}: ``almeno 50 moduli'', anzi ``123 pagine'';
  \item \emph{Bang!}: tre moduli.
\end{itemize}
Bisogna fare attenzione: alcuni moduli sono in lingue specifiche (``magari il
primo è in russo''), alcuni hanno già l'espansione, altri sono solo il gioco
base. Conviene sfogliare e non fermarsi al primo risultato. E nel Workshop si
trovano anche cose molto particolari: giochi fuori produzione, giochi vecchi,
giochi ancora in prototipo.

\subsection{I ruoli nei team e le dinamiche tra le persone}

Uno studente torna sulla slide delle figure: nelle realtà piccole molti ruoli
hanno ``il doppio o triplo cappello''. Quali ruoli restano distinti, e come si
creano le sinergie tra persone con obiettivi diversi?

Ruggero premette che l'ideale ``platonico'' è un ambiente di lavoro costruttivo,
basato sulla fiducia e ``il più \emph{safe} possibile''; ma nell'industria non
sempre è così, e le strutture variano moltissimo: alcuni publisher sono molto
verticali, altri hanno ruoli già solidificati (se un publisher ha già i suoi
developer e non sta crescendo, entrarci come developer è praticamente
impossibile). ``Se prendiamo due aziende in Italia, vediamo due strutture
interne completamente diverse''. Non esiste una scorciatoia per farsi assumere
da un editore.

Poi prende ``l'esempio più comune di tutti'', la media dei giochi da tavolo
pubblicati (figura~\ref{a1:fig:relazioni}):
\begin{itemize}
  \item il \textbf{game designer} è una figura \textbf{esterna}, con un
        \textbf{contratto a royalty} con il publisher;
  \item la sua figura di riferimento è il \textbf{publisher} stesso, e la
        conversazione designer--publisher è ``molto calda'';
  \item se il publisher ha dei \textbf{developer}, il developer entra in questa
        conversazione: designer e developer collaborano molto, e internamente
        developer e publisher collaborano molto. Spesso, nei publisher piccoli,
        è il publisher a fare le veci del developer;
  \item \textbf{graphic design} e \textbf{illustrazione} tendono a parlare molto
        con il publisher e \textbf{molto poco con il designer}. A meno che non si
        tratti di uno studio indie molto piccolo, è difficile che il designer
        sappia anche solo chi è l'illustratore prima che il gioco esca.
\end{itemize}

\begin{figure}[htbp]
\centering
\begin{tikzpicture}[
    font=\small,
    ruolo/.style={draw, thick, rounded corners=3pt, fill=gray!12,
                  minimum width=2.8cm, minimum height=0.9cm, align=center},
    esterno/.style={draw, thick, dashed, rounded corners=3pt,
                    minimum width=2.8cm, minimum height=0.9cm, align=center},
    calda/.style={{Stealth[length=2.5mm]}-{Stealth[length=2.5mm]}, very thick},
    tiepida/.style={{Stealth[length=2mm]}-{Stealth[length=2mm]}, thin, dashed}
  ]
  \node[esterno] (gd) at (0,0) {Game designer\\ \footnotesize esterno, a royalty};
  \node[ruolo] (dev) at (5,1.4) {Developer\\ \footnotesize se c'è};
  \node[ruolo] (pub) at (5,-1.4) {Publisher};
  \node[ruolo] (art) at (10,-1.4) {Graphic designer\\ e illustratore};
  \draw[calda] (gd) -- (dev);
  \draw[calda] (gd) -- (pub);
  \draw[calda] (dev) -- (pub);
  \draw[calda] (pub) -- (art);
  \draw[tiepida] (gd.south) to[bend right=25]
     node[below, font=\footnotesize\itshape] {quasi nessun contatto} (art.south);
  \node[draw, dotted, rounded corners=4pt, fit=(dev)(pub)(art), inner sep=8pt,
        label={[font=\footnotesize\itshape]above right:interno al publisher (di solito)}] {};
\end{tikzpicture}
\caption{Le relazioni tipiche intorno a un gioco da tavolo pubblicato, secondo la
risposta di Ruggero. Le frecce spesse indicano una collaborazione stretta e
frequente; quella tratteggiata, un contatto quasi assente. Nei publisher piccoli
il publisher fa spesso anche da developer.}
\label{a1:fig:relazioni}
\end{figure}

\subsection{Fare da developer di sé stessi?}

Rita chiede in chat se un designer può fare anche da developer del proprio
gioco, o se Ruggero lo sconsiglia. In generale \textbf{lo sconsiglia}, per lo
stesso motivo per cui si sconsiglia in altri media: i \textbf{darling}. Il
designer ha delle meccaniche a cui è affezionato, e il developer a volte dovrà
\textbf{tagliarle}; è molto difficile, se si è ``gestori di sé stessi'', tagliare
cose che si vogliono nel gioco. ``È un po' come lasciare fare il lavoro di editing
allo scrittore che ha già scritto il libro'': può rivedersi, certo, ma difficilmente
cambierà cose che farebbe bene cambiare. Nelle lezioni con Marta si vedrà perché,
e alla fine del corso ``non avrete dubbi nemmeno voi'' nel separare designer e
developer.

Ma tutto dipende dallo \textbf{scopo del progetto}: se si fa un gioco in tre,
qualcuno quel lavoro deve farlo per forza, e magari il designer fa anche il
developer, compensando con \textbf{più playtest}. ``Il mondo attuale e il mondo
ideale sono due cose molto diverse, e lo vediamo non solo dai giochi ma dalla
politica'', ma ``non entriamo in questo vaso di Pandora''.

\subsection{Il ruolo indispensabile e la pubblicazione}

Lo studente di prima ne trae una conclusione: senza game designer il gioco non
esiste. ``Quello sicuramente''. Ma Ruggero coglie l'occasione per approfondire
un punto che forse non era passato: \textbf{fare design di un gioco per
pubblicarlo non è l'unica opzione possibile}. Un gioco non è ``meno gioco''
se non è pubblicato da un editore. È legittimo mettere un gioco su itch.io in
\textbf{print and play}, fare un gioco a evento senza venderlo, fare un
\emph{murder party}, una escape room, un board game per una persona amica:
``quell'esperienza lì non è meno gioco''. Dipende dall'obiettivo del progetto.

Ruggero ha lavorato molto anche con le \textbf{scuole}, dove il publisher non
esiste: è un gioco per le classi, e la figura più vicina è l'editoria scolastica,
``che però non fa il publisher''. Ci sono compromessi diversi sulla materialità
e su come il gioco si presenta, che si vedranno nella prototipazione e nei casi
studio. Quindi l'unica figura che deve esserci sempre è \textbf{chi fa game
design}, anche se magari non sa di chiamarsi game designer e non ha mai
sentito parlare di game design: ``questo non la rende meno game designer''.

\paragraph{Un libro consigliato.} Ruggero fa anche mentoring e gestisce gruppi
di settore, e spesso gli chiedono che cosa leggere per imparare il game design.
Non è il tipo che obbliga a leggere libri (``non vi obbligherò a leggere
niente''), ma ce n'è uno che consiglia sempre, ed è il primo che ha scelto di
leggere: \textbf{\emph{Rise of the Videogame Zinesters}} di \textbf{Anna
Anthropy}\footnote{La trascrizione dice ``Anna Antrofi'' e ``Rise of the
Videogames in Esters''.}. È un libro un po' datato, ma ancora reperibile; parla
soprattutto di videogiochi, ma ``per l'80\%'' il contenuto si applica al gioco e
al game design in generale. È un libro sulla \textbf{democratizzazione del gioco}:
il gioco come strumento per parlare di sé e come \textbf{medium legittimo}. Come
esistono film drammatici, tristi, intensi, possono esistere giochi su tematiche
delicate, e se ne vedranno nei casi studio. Molti giochi fatti con questa
direttiva non hanno un publisher, che tende a pubblicare ciò che dà un ritorno
commerciale.

\paragraph{Commerciale e personale.} ``Io spero che voi non farete giochi solo
per scopo commerciale''. Certo, servono un tetto sulla testa, il cibo in
tavola, un posto in cui dormire, e così si finisce per fare giochi che non sono
il proprio gioco dei sogni. Spesso è un mix: Ruggero ha fatto giochi commerciali
``per pagarmi l'affitto'' e giochi che voleva fare e su cui ha perso soldi,
``sapendolo proprio che sarebbero stati un fallimento''. I giochi, come i libri e
le canzoni, si fanno per molti motivi, e quelli commerciali di successo non sono
più importanti degli altri. Tra i giochi della sua slide c'erano il gioco da
tavolo ufficiale di \emph{Cyberpunk}\footnote{La trascrizione dice ``Cyberpunk
del 1977'': quasi certamente un gioco legato a \emph{Cyberpunk 2077}.} e
\emph{Marvel United}, su IP grosse e famose, ma anche giochi fatti di ``una
candela e una busta su un tavolo''. E se gli si chiede quale preferisce in
termini di contenuto, non sono quelli sui supereroi.

\subsection{Il gioco di ruolo in Italia}

Nicola, che è interessato soprattutto al gioco di ruolo e conosce già persone
del settore, ha ``una leggera paura'': che chi studia game design finisca per
fare il \textbf{freelance} a vita, ``perché non ci sarà mai un'occupazione per
bene'', o il \textbf{copywriter}. Conosce un ragazzo che ha lavorato in quel
mondo e che ora fa altro.

Ruggero risponde che la maggioranza dei giochi nel suo portfolio sono giochi di
ruolo, quindi il tema lo tocca da vicino, e che in Italia quello che dice Nicola
``è semplicemente vero''. Molti publisher piccoli hanno una persona di
riferimento che copre più ruoli: impagina, fa i playtest, fa un po' di
development. È ``un settore economicamente completamente disastrato'', per
motivi culturali che non apre: oggi in Italia è praticamente impossibile
\textbf{vivere di giochi di ruolo}. Chi ci lavora lo fa nel tempo libero, al di
fuori del lavoro principale, oppure coprendo molte altre cose: molti publisher
italiani di giochi di ruolo pubblicano anche titoli esteri e localizzano, non
fanno solo cose originali. ``È un grosso calderone''. Poi si ferma: ``questo è
un \emph{rant}'', e ``i \emph{rant} tornano sempre'', ma non vuole che gli
studenti escano dalla prima lezione ``con la disperazione nel cuore''. Ci sarà
spazio per parlarne quando avranno gli strumenti per capirlo, ``per non farvi
abbattere''.

\subsection{Che cosa bisogna sapere già}

Uno studente chiede se ci sono nozioni che dovrebbero già avere e che possono
recuperare: per esempio, sente parlare di giochi ``German'' o ``American'' e non
sa che cosa siano. Ruggero spiega l'organizzazione del corso:
\begin{itemize}
  \item le \textbf{prime due lezioni} sono di fondamenta. Questa era ``dall'alto'':
        come è strutturato il corso, le tempistiche;
  \item la \textbf{prossima lezione} è di tre ore di soli \textbf{casi studio},
        ``a manetta, a raffica'', divisi per categorie e tipologie, con
        ragionamenti su categorie e generi di gioco: lì si vedranno German,
        American ``e anche altre più utili'';
  \item nel resto del corso le terminologie specifiche sono inserite nelle
        lezioni dedicate: il \textbf{paragame} nelle lezioni sul paragame, il
        \textbf{modello matematico} nelle lezioni di Marta sul development.
\end{itemize}
Per ora agli studenti si chiede di sapere quello che si è visto oggi:
meccaniche, dinamiche ed estetiche; le figure professionali, designer e
developer; e soprattutto \textbf{game e play}.

\subsection{Game e play: i quattro quadranti}

``Specialmente \emph{game} e \emph{play}'', insiste Ruggero: è più difficile di
quello che sembra, perché l'italiano non abitua a separare le due cose, ed è
possibile averle usate per anni come sinonimi. Nella sua esperienza con chi si
avvicina al design è ``uno dei punti più critici''. Spesso i commenti che si
fanno sui giochi giocati sono commenti sul \emph{play} applicati al
\emph{game}: ``il gioco non mi è piaciuto perché, giocando con Tizio, è successo
questo''; ``il gioco mi sembra sbilanciato perché Tizio al tavolo ha fatto
questo''.

Separare \emph{game} e \emph{play} serve a capire quali giochi sono progettati
bene e quali male, e soprattutto a riconoscere i \textbf{quattro possibili
incroci} (figura~\ref{a1:fig:quadranti}): un gioco progettato bene che piace
giocare; uno progettato male che non piace; uno progettato male che piace; uno
progettato bene che non piace. ``Non è importantissimo in generale nella vita,
ci sono cose più importanti'', ma per chi fa game design lo è: ``a volte vi
chiederanno di lavorare su un gioco che non vi piace giocare, e quel gioco deve
essere fatto bene lo stesso''.

\begin{figure}[htbp]
\centering
\begin{tikzpicture}[
    font=\small,
    cella/.style={draw, thick, minimum width=4.4cm, minimum height=1.9cm,
                  text width=4.0cm, align=center}
  ]
  \node[cella, fill=gray!20] (a) at (0,0)     {\textbf{progettato bene}\\ e \textbf{mi piace} giocarlo};
  \node[cella]               (b) at (4.4,0)   {\textbf{progettato bene}\\ ma \textbf{non mi piace} giocarlo};
  \node[cella]               (c) at (0,-1.9)  {\textbf{progettato male}\\ ma \textbf{mi piace} giocarlo};
  \node[cella, fill=gray!8]  (d) at (4.4,-1.9){\textbf{progettato male}\\ e \textbf{non mi piace} giocarlo};
  \node[font=\footnotesize\itshape, above=0.15cm of a] {mi piace (\emph{play})};
  \node[font=\footnotesize\itshape, above=0.15cm of b] {non mi piace (\emph{play})};
  \node[font=\footnotesize\itshape, rotate=90, left=0.15cm of a, anchor=south]
       {bene (\emph{game})};
  \node[font=\footnotesize\itshape, rotate=90, left=0.15cm of c, anchor=south]
       {male (\emph{game})};
\end{tikzpicture}
\caption{I quattro incroci tra la qualità del \emph{game} (come è progettato) e
l'esperienza di \emph{play} (se piace giocarlo), descritti da Ruggero. Per un
designer i quadranti ``scomodi'' sono i due fuori diagonale: un gioco può
piacere pur essendo progettato male, e va saputo progettare bene anche un gioco
che non piace.}
\label{a1:fig:quadranti}
\end{figure}

\paragraph{La domanda di Tancredi.} Tancredi chiede conferma: il \emph{play}
non è il ``giocare senza regole'', ma può essere anche il giocare a un gioco? Sì:
il \emph{play} è l'\textbf{atto}, e può essere fatto con un \emph{game} o senza.
E un \emph{game} può esistere senza che nessuno lo giochi: gli scacchi non sono
meno scacchi perché in quel momento nessuno ci gioca. E se, giocando, una
persona applica le regole in modo diverso? Allora \textbf{non si sta giocando
quel gioco}.

\subsection{House rule: il caso di Monopoly}

Qui Ruggero fa un po' di \emph{foreshadowing} su una lezione di Daniele: il
discorso degli \textbf{hack}, delle \textbf{house rule}, delle \textbf{mod},
cioè di come modificare un gioco esistente. Il punto è che \textbf{un gioco
modificato da una house rule è un altro gioco}: cambiando le regole cambiano
necessariamente anche le \textbf{dinamiche} che ne derivano, e quindi
l'esperienza. Si immagini \emph{Monopoly} in cui ogni turno si possa, volendo,
prendere gratis dei soldi dalla banca: tutto il resto è identico, eppure il gioco
cambia.

E questo è talmente interiorizzato che ci sono regole che si applicano ``senza
neanche rendercene conto''. Ruggero chiede quanto duri in media una partita a
\emph{Monopoly}. Gianmarco: ``quattro ore''; altri: ``un pomeriggio''; un altro:
``finché uno dei giocatori non aggiunge la regola: ora vince chi\ldots''. In
realtà, \textbf{seguendo le regole scritte}, una partita non dura mai più di
\textbf{un'ora e venti, un'ora e mezza}. La domanda che fa sempre, e la risposta
è sempre la stessa: perché si pensa che \emph{Monopoly} duri tantissimo? Qual è
la regola che nessuno rispetta?

Uno studente indovina: \textbf{le aste}. Secondo il regolamento, se chi finisce
su un terreno non lo compra, il terreno \textbf{va all'asta} tra i giocatori, e
chiunque può aggiudicarselo. Molti invece giocano come se finire su un terreno
desse solo la \emph{facoltà} di comprarlo o di metterlo all'asta\footnote{Più
precisamente, nel regolamento ufficiale, se il giocatore che finisce su una
proprietà libera decide di non comprarla, la banca la mette all'asta e tutti i
giocatori possono fare offerte, compreso chi l'ha rifiutata. L'asta quindi non è
una scelta del giocatore, ma la conseguenza obbligata del suo rifiuto.}. Con la
regola applicata, di aste se ne fanno moltissime: all'ottavo turno ci sono già
otto terreni assegnati, al ventesimo venti, la gente va in bancarotta
velocissimamente e la partita dura un'ora e mezza. Anche le partite competitive
di \emph{Monopoly}, nota Ruggero, durano un'ora e venti, un'ora e mezza.

Uno studente obietta che giocava con le aste e durava comunque quattro ore.
``Avete dei nemici con l'analisi paralisi'', ride Ruggero; e cita un'altra house
rule tipica dell'infanzia: genitori che danno soldi per non far andare in
bancarotta. Ma \textbf{la bancarotta è parte del gioco}: se la si toglie,
\emph{Monopoly} non regge. Sul tema che \emph{Monopoly} vuole trattare, e sul
suo antenato (\emph{The Landlord's Game}), ``non apriamo adesso un altro vaso di
Pandora''.

\paragraph{Approfondimento: Monopoly, la durata e le house rule nei libri.}
L'esempio di Ruggero si ritrova quasi identico in \emph{The Art of Game Design}.
Parlando del bilanciamento tra partite brevi e lunghe (\emph{Short vs. Long}),
Schell scrive che \emph{Monopoly}, giocato secondo le regole ufficiali,
``finisce spesso in circa novanta minuti'', ma molti giocatori trovano quelle
regole troppo dure e le modificano, distribuendo premi in denaro e allentando i
vincoli sull'acquisto delle proprietà: con l'effetto collaterale di far durare la
partita molto di più, ``tipicamente tre ore, o anche di più''. Schell classifica
le \textbf{house rule} come regole create dai giocatori in risposta a un
difetto percepito dopo qualche partita. Anche \emph{Challenges for Game
Designers}, nel capitolo sull'aggiungere e togliere meccaniche, parte dallo
stesso gioco: il \emph{Parcheggio gratuito}, che per le regole originali non fa
succedere nulla, mentre quasi tutti giocano con una house rule che dà un bonus a
chi ci finisce sopra. Le due fonti confermano la lezione: una piccola regola
cambiata cambia le dinamiche, e le dinamiche cambiano l'esperienza.

\subsection{Il fine del game è il play?}

Un altro studente commenta: va bene la distinzione, ma il fine del \emph{game} è
il \emph{play}; quello che si vede è sempre l'effetto. Ruggero è d'accordo: di
solito, se si fa un gioco, lo si fa per farlo giocare, come si scrive un libro
per farlo leggere (anche se esistono, ed ``è già stato esplorato'', libri scritti
per non essere letti e giochi per zero giocatori). Ma la distinzione serve a una
cosa precisa: \textbf{separare, nel playtest, i feedback sul gioco dai feedback
sul gruppo di gioco} e sulle sue dinamiche. Se una persona esce frustratissima da
un playtest, non è necessariamente colpa del gioco; se una persona dice ``questo
è il gioco migliore del mondo'', potrebbe non essere merito del gioco ma del
gruppo. Capire a che cosa si riferisce un feedback serve a capire \textbf{dove
iterare} il prototipo.

Lo studente osserva che l'area è sfumata: alcuni giochi si basano proprio
sull'avere le persone giuste al tavolo. Un gioco di \textbf{negoziazione} con
giocatori abituati ai giochi ``multi-solitario'' diventa ingiocabile; e nei giochi
a \textbf{economia chiusa} (cita \emph{Container}\footnote{La trascrizione cita
anche un secondo titolo non ricostruibile (``Tokyo su Vigimarket'').}) un gruppo
economicamente conservativo, che spende poco e non fa girare l'economia, fa
durare la partita il doppio. Ruggero rilancia: e se il gioco dura il doppio,
questa dinamica cambia la \textbf{percezione di peso} del gioco, le energie spese,
quanto si è stanchi a fine partita. ``Tutto è collegato. \textbf{I giochi sono
ecosistemi di meccaniche}, e cambiandone una, le altre cambiano di
conseguenza''.

\paragraph{La parola vietata, alla prova.} Lo studente aggiunge che ci sono giochi
che ha trovato ``molto sbagliati, ma che mi divertono''. ``Attenzione: non si può
dire \emph{divertimento}!'', lo ferma Ruggero; nelle presentazioni aveva lasciato
correre, ``ma lì non vi avevo ancora letto la regola''. Lo studente ci prova:
``giochi che poi voglio \textbf{rigiocare}'', anche se ``sono evidentemente
sbagliati, rotti''. Ruggero apprezza, ma non vuole anticipare cose che meritano il
loro tempo: se ne parlerà nelle lezioni sui framework e soprattutto in quella di
\textbf{emotion design}, ``forse la mia lezione preferita del corso''. (Lo
studente ammette che ``forse alla fine si gioca per quello, per l'emozione''.)

\subsection{I giochi ``comunicati male''}

L'ultima domanda corposa riprende un'affermazione di inizio lezione: ci sono
giochi belli ma comunicati male. Qualche esempio? ``Mi basta girarmi per
vederne 48'', risponde Ruggero indicando gli scaffali alle sue spalle. Più che
titoli, elenca le \textbf{ragioni}:
\begin{itemize}
  \item \textbf{l'Italia localizza lentamente}. Basta guardare i premi del gioco
        dell'anno in Italia: spesso vanno a giochi usciti due o tre anni prima in
        Germania, Francia, Stati Uniti o Regno Unito, perché localizzare richiede i
        suoi tempi. Il mercato italiano è quindi molto dipendente dall'estero,
        ``molto impaurito'', e spesso aspetta che un gioco vada bene all'estero
        prima di portarlo in Italia; e alcuni giochi ``perdono il treno'';
  \item \textbf{il publisher originale chiude}: come si pubblica un gioco se non
        c'è più l'azienda con cui fare un contratto? A Ruggero è successo più
        volte;
  \item \textbf{il dipartimento commerciale italiano non sa comunicare il gioco}:
        si fa male un'edizione, la stampa vende pochissimo e non viene più
        ristampata;
  \item \textbf{la scala del mercato}: in Italia un gioco da tavolo di peso medio
        (un \emph{eurogame}, un \emph{German}) che vende 1.500 copie ``è tanto'';
        all'estero se ne vendono magari 10.000. Con 10.000 copie in inglese ci
        sono i soldi per una campagna di comunicazione: un trailer, i
        recensori, i tutorial online. In Italia tutto questo è ``ridotto all'osso'',
        e la cosa si riflette su stampa e divulgazione: anche i canali per
        informarsi su un gioco prima di comprarlo sono pochi (se ne parlerà nella
        lezione su ricerca e sviluppo). Un gioco spesso non arriva al mercato non
        perché sia brutto, ma perché l'editore non lo sa comunicare, ne fa poche
        copie e va fuori produzione;
  \item \textbf{il tempismo}: party game progettati nel 2019 e usciti nel 2020, in
        pieno Covid, quando non si poteva giocare al tavolo, sono falliti tutti, e
        gli editori ci hanno perso. ``Non potevano immaginarselo'';
  \item \textbf{i temi caldi}: dopo l'invasione russa dell'Ucraina e la guerra in
        Palestina, alcuni editori non pubblicano o non localizzano giochi
        strategici e di guerra. Ruggero ha lavorato a \textbf{\emph{Status}}, un
        gioco sulla geopolitica che tocca, tra gli altri temi, l'invasione russa
        dell'Ucraina e il genocidio in Palestina, e alcuni editori che di
        solito localizzano i loro giochi (per esempio il partner di riferimento in
        Russia) non lo hanno localizzato per motivi politici.
\end{itemize}
Lo studente voleva capire che cosa voglia dire comunicare male: non metterci
abbastanza energia, o comunicare una cosa per un'altra, ``come il trailer
sbagliato, che poi vai al cinema e ti trovi un altro film''. ``È un mix'', e la
domanda ``era complessa e meritava più spazio''.

\subsection{I giochi di carte collezionabili e i saluti}

L'ultima domanda (``posso rompere le balle un'ultima volta?'') riguarda i
\textbf{giochi di carte collezionabili}: funzionano con le stesse regole dei
giochi da tavolo? Se ne parlerà alla prossima lezione: è proprio il motivo per
cui la lezione sui casi studio dura tre ore, perché ogni genere di gioco ha le
sue dinamiche e le sue regole, e anche il modo di pubblicare cambia molto tra un
TCG, un gioco di ruolo da tavolo, un LARP e un gioco pervasivo.

Ruggero saluta: le email dei tre docenti sono all'inizio delle slide, la
registrazione arriverà, ``ci si vede questo venerdì''.

% ---------------------------------------------------------------------
\section*{Punti chiave del capitolo}
\addcontentsline{toc}{section}{Punti chiave del capitolo}

\begin{itemize}
  \item Il corso di \textbf{Analog Game Design} (seconda edizione) dura circa tre
        mesi, fino a dicembre, ed è tenuto da tre docenti principali:
        \textbf{Daniele} (narrative design, gioco di ruolo, fiction),
        \textbf{Ruggero} (game design, ricerca e sviluppo) e \textbf{Beatrice}
        (giochi pervasivi), con tre ospiti su RNG, development e game production.
  \item Lo scopo è formare \textbf{autori analogici} consapevoli: strumenti da
        usare anche fuori dal corso, non il ``gioco dei sogni''. Il corso è
        \textbf{intersezionale} (nessuna lezione dedicata a un solo tipo di
        gioco), \textbf{cross-generazionale}, \textbf{pratico e adattabile},
        basato sulle \textbf{domande} (nessuna domanda è stupida), sul
        \textbf{feedback} invece che sul voto e sul \emph{fail forward}.
  \item Il \textbf{project work} si fa in tre gruppi da tre, formati dai docenti;
        pitch il 12 ottobre, sviluppo nelle due pause, tutoraggi il sabato,
        consegna il 4 dicembre, presentazione l'11 e giornata finale con gli
        editori il 14. Si valutano documentazione, presentazione,
        autovalutazione, crescita e atteggiamento.
  \item \textbf{Niente intelligenza artificiale} nel corso: i publisher
        rifiutano pitch e prototipi con immagini generate, i regolamenti scritti
        dall'IA descrivono la stessa meccanica in modi diversi, e soprattutto chi
        la usa perde l'occasione di imparare. Un prototipo può essere fatto a
        matita.
  \item Le \textbf{figure dell'industria}: game designer, game developer,
        editor, publisher, sales e marketing, graphic designer, illustratore,
        sensitivity reader (che non è un censore) e producer. Nelle realtà
        piccole una persona copre più ruoli.
  \item Nell'analogico il \textbf{game designer} progetta l'esperienza creando
        le meccaniche; il \textbf{game developer} fa da \textbf{mediatore} con il
        publisher (bilanciamento, modelli matematici, adattamenti). Pretendere
        che il proprio gioco non venga modificato è un \emph{red flag}.
  \item Il \textbf{marketing} non decide che cosa pubblicare ma come; è
        un'azione, ed esiste anche se lo si ignora.
  \item In \textbf{Italia} il gioco è ancora visto come infantile, il mercato
        punta sulla localizzazione e investe poco su giochi e designer interni;
        all'estero (soprattutto dalla Germania in su) ci sono tradizione adulta,
        premi di prestigio e developer interni.
  \item \textbf{Game} è l'artefatto (le regole, il sistema, l'oggetto);
        \textbf{play} è l'atto di giocare. Si fanno giochi, non giocattoli.
        Separarli serve a distinguere i feedback sul gioco da quelli sul gruppo
        e a riconoscere i quattro incroci tra ``progettato bene/male'' e ``mi
        piace/non mi piace''.
  \item \textbf{MDA}: meccaniche (regole scritte), dinamiche (regole in azione),
        estetiche (esperienza, non grafica); nel \textbf{DPE}: design, play,
        experience.
  \item \textbf{Parole vietate} in ogni materiale scritto: \emph{fun},
        \emph{immersion}, \emph{entertainment} (divertimento, immersione,
        intrattenimento). Bisogna trovare parole più precise.
  \item Distinguere \textbf{chi crea, chi gioca e chi facilita}; nessun gioco
        nasce perfetto: \textbf{playtest e iterazione} sono inevitabili, e un
        \textbf{dizionario comune} aiuta a non perdersi.
  \item \textbf{Tabletop Simulator} è il software di settore per giocare (non
        creare videogiochi) i giochi da tavolo in digitale. Primo task: giocarci
        tre o quattro volte entro il 10 ottobre.
  \item Una \textbf{house rule} crea un altro gioco: \emph{Monopoly} con le aste
        applicate come da regolamento dura circa un'ora e mezza.
\end{itemize}

%    % =====================================================================
%  Analog Game Design, lezione 2 — I casi studio: una panoramica a volo
%  d'aquila dei giochi pervasivi, dei giochi di ruolo dal vivo, dei
%  giochi di ruolo da tavolo e dei giochi da tavolo
%  Fonte: trascrizione della seconda lezione (AGD2) del corso breve di
%  Analog Game Design, seconda edizione: una lezione online su Zoom di
%  circa tre ore, di venerdì sera, tenuta a più voci dai tre docenti
%  "cuore" del corso (Ruggero, Daniele e Beatrice), con una pausa di
%  pochi minuti e moltissime slide di esempi.
%  Pacchetti richiesti nel preambolo del libro (già presenti nel
%  main.tex): \usepackage{booktabs}, \usepackage{tikz} con
%    \usetikzlibrary{shapes.geometric,shapes.symbols,shapes.misc,
%      arrows.meta,positioning,calc,fit}
%  (le librerie vanno caricate nel preambolo, non nel capitolo).
%  Consigliati, come per gli altri capitoli: \usepackage[italian]{babel}
%  e \usepackage[T1]{fontenc} (il capitolo compila anche senza).
%
%  Collocazione nel libro: parte "Analog Game Design", secondo capitolo
%  (numerazione propria 2.2). Nel file principale del libro, subito dopo
%  il capitolo 2.1:
%
%    % =====================================================================
%  Analog Game Design, lezione 1 — Introduzione al corso: docenti,
%  metodo, project work, figure dell'industria e primo glossario
%  Fonte: trascrizione della prima lezione (AGD1) del corso breve di
%  Analog Game Design, seconda edizione: una lezione online su Zoom di
%  circa tre ore, di pomeriggio, tenuta insieme dai tre docenti "cuore"
%  del corso (Daniele, Ruggero e Beatrice), con nove iscritti, una
%  pausa di una decina di minuti, molte slide e una breve dimostrazione
%  dal vivo di Tabletop Simulator.
%  Pacchetti richiesti nel preambolo del libro (già presenti nel
%  main.tex): \usepackage{booktabs}, \usepackage{tikz} con
%    \usetikzlibrary{shapes.geometric,shapes.symbols,shapes.misc,
%      arrows.meta,positioning,calc,fit}
%  (le librerie vanno caricate nel preambolo, non nel capitolo).
%  Consigliati, come per gli altri capitoli: \usepackage[italian]{babel}
%  e \usepackage[T1]{fontenc} (il capitolo compila anche senza).
%
%  Collocazione nel libro: parte "Analog Game Design", primo capitolo
%  (numerazione propria 2.1). Nel file principale del libro:
%
%    \part{Analog Game Design}
%    \renewcommand{\thechapter}{2.\arabic{chapter}}
%    \setcounter{chapter}{0}
%    \input{AGD/AGD1_introduzione_all_analog_game_design}
%
%  (Con hyperref conviene aggiungere anche
%  \renewcommand{\theHchapter}{agd.\arabic{chapter}} nella parte, per
%  evitare ancore duplicate con i capitoli 1, 2, ... di Game Design.)
%  Tutte le etichette di questo capitolo hanno il prefisso "a1:"
%  (a1:cap:..., a1:sec:..., a1:fig:..., a1:tab:...). Il capitolo è
%  autonomo: non contiene \ref verso altri capitoli.
%  Le figure sono ricostruzioni schematiche di quanto i docenti spiegano
%  a voce e con le slide, non copie delle slide.
% =====================================================================

\chapter[Introduzione all'Analog Game Design]{Introduzione all'Analog Game
Design: il corso, le figure dell'industria e un primo linguaggio comune}
\label{a1:cap:intro}

Questa è la \textbf{prima lezione} del corso di \textbf{Analog Game Design},
cioè di progettazione di giochi \emph{analogici}: giochi da tavolo, giochi di
carte, giochi di ruolo da tavolo, giochi di ruolo dal vivo (LARP), giochi
\emph{pervasivi} ed esperienze a evento. È un corso \textbf{breve}, di circa tre
mesi (``da qui a dicembre''), separato dal corso di Game Design, e siamo alla sua
\textbf{seconda edizione}. Come avvertono subito i docenti, è ``una giornata di
introduzione'': si spiega come funziona il corso, chi sono i docenti, con quali
metodi e con quale calendario si lavorerà, e ci si presenta a vicenda. ``C'è poco
gioco analogico'', si scusano. Ma nella seconda parte la lezione comincia
comunque a costruire un \textbf{linguaggio comune}: le figure professionali
dell'industria del gioco analogico, la differenza tra \emph{game design} e
\emph{game development}, la situazione del mercato italiano rispetto a quello
estero, e un primo glossario di design (\emph{game} e \emph{play}, il framework
MDA, le parole ``vietate'', il playtest e l'iterazione). Si chiude con lo
strumento di settore per i prototipi digitali dei giochi da tavolo,
\textbf{Tabletop Simulator}, e con una lunga sessione di domande.

La lezione si svolge \textbf{online}, su Zoom, ed è condotta a più voci dai tre
docenti principali del corso: \textbf{Daniele}, narrative designer, che guida la
parte organizzativa; \textbf{Ruggero}, game designer, che guida la parte sui
fondamenti e le domande finali; \textbf{Beatrice}, progettista di giochi
pervasivi, che interviene soprattutto nella prima metà. Gli studenti iscritti
sono \textbf{nove}; uno è assente e un altro arriva a lezione iniziata. Ci sono
una pausa di una decina di minuti (si riprende alle 18:15) e, dopo le 19:15,
un'ultima mezz'ora di domande con il solo Ruggero.

Tre avvertenze di lettura. La prima: la trascrizione è automatica e storpia
molti nomi propri, compresi quelli dei docenti e dell'accademia (``Dive'',
``IVE'', ``Accademia Live'', ``Ai''); dove la ricostruzione è sicura, o molto
probabile, il capitolo usa la forma corretta e lo segnala in nota; dove non lo
è, riporta la forma della trascrizione e lo dice. La seconda: durante la pausa la
trascrizione registra, ripetuta in loop per decine di righe, una frase senza
senso che non ha nulla a che fare con la lezione, e che il capitolo ignora. La
terza: in molti passaggi più docenti parlano uno dopo l'altro senza che la
trascrizione segni il cambio di voce; il capitolo attribuisce le parole a chi le
pronuncia quando il contesto lo rende chiaro (chi parla di sé, chi viene
chiamato per nome, chi ``passa la parola''), e resta generico altrimenti.

Quanto alla data, la lezione si colloca con ogni probabilità \textbf{nella prima
metà di settembre del 2026}\footnote{La data non è mai detta. La si deduce da
alcuni indizi: il form iniziale va compilato ``entro e non oltre il 15
settembre''; i gruppi verranno annunciati il 25 settembre; le date del
calendario (sabato 3 ottobre, lunedì 12 ottobre, sabato 7 e 14 novembre,
venerdì 4 dicembre) cadono nei giorni della settimana indicati soltanto nel
2026; un docente parla di un argomento ``molto complesso nel 2026''; e fa
ancora molto caldo (``le temperature sono ancora indecenti''). È una
deduzione, non un'informazione esplicita.}: i riferimenti a ``oggi'',
``quest'estate'' e ``l'anno scorso'' vanno letti rispetto a quel momento.

Il percorso del capitolo segue quello della lezione:
\begin{enumerate}
  \item l'avvio, i \textbf{docenti} e le \textbf{presentazioni} dei partecipanti
        (sezioni da~\ref{a1:sec:avvio} a~\ref{a1:sec:partecipanti});
  \item i \textbf{materiali} e gli strumenti del corso: il form, la cartella
        condivisa, Discord (sezione~\ref{a1:sec:materiali});
  \item lo \textbf{scopo} e il \textbf{metodo} del corso, le regole di
        comportamento e il carico di lavoro (sezioni~\ref{a1:sec:scopo}
        e~\ref{a1:sec:regole});
  \item il \textbf{project work} e il \textbf{calendario} (sezione~\ref{a1:sec:project});
  \item la regola sull'\textbf{intelligenza artificiale} e la lunga discussione
        che ne segue (sezione~\ref{a1:sec:ia});
  \item le \textbf{figure professionali} dell'industria del gioco analogico,
        con un approfondimento su marketing e sales
        (sezioni~\ref{a1:sec:figure} e~\ref{a1:sec:marketing});
  \item la differenza tra \textbf{game design} e \textbf{game development}
        (sezione~\ref{a1:sec:design-development});
  \item il \textbf{mercato}: Italia ed estero (sezione~\ref{a1:sec:mercato});
  \item i \textbf{fondamenti di design}: perché serve una terminologia comune,
        \emph{game} e \emph{play}, il framework \textbf{MDA} (e il DPE), le
        \textbf{parole vietate}, chi crea, chi gioca e chi facilita, il
        \textbf{playtest} e l'\textbf{iterazione} (sezioni
        da~\ref{a1:sec:terminologia} a~\ref{a1:sec:playtest});
  \item \textbf{Tabletop Simulator} e il primo task (sezione~\ref{a1:sec:tts});
  \item le \textbf{domande finali}: il tema del project work, il bilanciamento,
        i ruoli nei team, la pubblicazione, il gioco di ruolo in Italia, le
        \emph{house rule} di Monopoly, i giochi ``comunicati male''
        (sezioni~\ref{a1:sec:domande-daniele} e~\ref{a1:sec:domande}).
\end{enumerate}

% ---------------------------------------------------------------------
\section{L'avvio della lezione}
\label{a1:sec:avvio}

La lezione comincia con qualche minuto di attesa. ``Siamo in nove in classe'', e
si aspetta che arrivino tutti; uno dei docenti scherza sul fatto che l'unica
persona che arriverà ``sempre in ritardo di venerdì'' è Gianmarco. Mancano
Giordano e Tiziana; Tiziana arriva, mentre la segreteria fa sapere poco dopo che
Giordano non potrà esserci e vedrà la registrazione. Si aspetta anche Beatrice,
la terza docente, che arriva a sua volta dopo pochi minuti.

Nell'attesa, Daniele anticipa il programma della giornata e qualche
informazione pratica:
\begin{itemize}
  \item oggi è una giornata di \textbf{introduzione}: come funziona il corso, chi
        sono i docenti, le presentazioni degli studenti, le \textbf{metodologie}
        e il \textbf{calendario}. Gli studenti avranno ``un fracco di domande'', e
        alla fine della lezione c'è ``proprio uno slot apposta'' per le domande;
  \item le lezioni sono \textbf{registrate}: di solito la registrazione viene
        consegnata dopo circa \textbf{12--24 ore} sul profilo personale della
        piattaforma dell'accademia, che va attivato con la segreteria (Tiziana,
        già studentessa, probabilmente ce l'ha già);
  \item i docenti ``cuore'' del corso sono tre (Daniele, Ruggero e Beatrice),
        ``soprattutto Ruggero''; a loro si aggiungeranno tre docenti ospiti, per
        meno ore, probabilmente una lezione a testa, ``tranne Marta, che avrà due
        lezioni quest'anno''.
\end{itemize}
E una battuta sul caldo: ``spero che voi siate carichi, anche se le temperature
sono ancora indecenti''. ``Carichi moderatamente, per quello che potete''.

Poi Daniele condivide lo schermo (``non solo le slide, ma anche un po' di altre
cose'') e apre ufficialmente il corso: ``benvenuti e benvenute nella seconda
edizione del corso di Analog Game Design. Questa è la prima lezione''. Il
percorso è di circa \textbf{tre mesi}, fino a dicembre, e durante il corso i
docenti cercheranno di far provare agli studenti ``cosa significa fare Analog
Game Design'' con un \textbf{project work}, un progetto di gruppo di cui si
spiegheranno più avanti la struttura e le richieste.

% ---------------------------------------------------------------------
\section{I docenti}
\label{a1:sec:docenti}

``Anzitutto partiamo da delle presentazioni doverose''. Ognuno dei tre docenti
principali si presenta con una slide che riporta i contatti e una selezione dei
lavori fatti.

\subsection{Daniele: narrativa, fiction e gioco di ruolo}

\textbf{Daniele}\footnote{La trascrizione riporta ``Daniel Cusetto''; nel
resto della lezione tutti lo chiamano ``Daniele''. Il cognome non è
ricostruibile con certezza.} ha già conosciuto quasi tutti gli studenti, perché
ha fatto loro il colloquio di ammissione. È \textbf{narrative designer}, o meglio
\emph{game narrative designer}; insegna nell'accademia che ospita il
corso\footnote{La trascrizione storpia il nome dell'accademia in molti modi
(``Dive'', ``IVE'', ``Accademia Live'', ``Ai'', ``AIVA''). Con ogni probabilità
si tratta dell'\textbf{AIV}, l'Accademia Italiana Videogiochi, che ha sede a
Roma: lo fanno pensare i corsi triennali di game design e narrative design citati
dai docenti e il fatto che una studentessa racconti di essersi trasferita a Roma.
È una ricostruzione, non un'informazione esplicita; nel seguito il capitolo
parla semplicemente dell'``accademia''.} ma lavora soprattutto come \textbf{brand
advisor}. Nel corso si occuperà della parte di \textbf{gioco di ruolo},
\textbf{narrazione} e \textbf{fiction}.

I suoi contatti, ``in quest'ordine'', sono email, Telegram e Discord. Ci tiene a
dare alcune indicazioni precise:
\begin{itemize}
  \item se si vuole avere la certezza di una risposta, bisogna scrivergli su
        \textbf{Telegram}, perché lì ha tutti i suoi clienti;
  \item Discord lo guarda di meno; l'\textbf{email} è per le cose ufficiali;
  \item non bisogna contattarlo altrove: non su Facebook, non su LinkedIn, non su
        Instagram, ``perché li vedo molto di meno''.
\end{itemize}
La slide mostra anche alcuni giochi, attività e brand su cui ha lavorato negli
anni, ma non si dilunga: ``nel caso ne parleremo''.

\subsection{Ruggero: game design e ricerca e sviluppo}

\textbf{Ruggero} (spesso chiamato ``Ruger'') ha fatto il \textbf{game designer}
per varie aziende e oggi è game designer e si occupa di \textbf{ricerca e
sviluppo} per un'azienda del settore\footnote{La trascrizione dice ``per MS'';
il nome dell'azienda non è ricostruibile con certezza. Più avanti Ruggero cita
tra i giochi su cui ha lavorato il gioco da tavolo ufficiale di
\emph{Cyberpunk}, \emph{Marvel United} e \emph{Status}, un gioco sulla
geopolitica.}. Ha un background ``misto'': nel suo portfolio ci sono giochi da
tavolo, giochi di ruolo, giochi analogici, LARP, videogiochi, giochi in realtà
virtuale. Si occupa anche di \textbf{divulgazione} sul gioco e di analisi di
\textbf{game studies} e game design. Più avanti, nelle domande finali, dirà di
avere ``50 e più giochi prodotti e progettati'' e di insegnare da vari anni.

Lo si può contattare liberamente via Gmail e Telegram, e cerca di rispondere
\textbf{entro 24 ore}, salvo imprevisti o eventi: ``sentitevi in libertà di
scrivermi se avete dubbi, senza problemi''. Quando arriva a parlare dei propri
giochi, la slide mostra ``un mix di varie cose, dalle più famose alle meno
famose'', ma Ruggero si imbarazza a presentarli, tra le risate dei colleghi.

\subsection{Beatrice: i giochi pervasivi}

\textbf{Beatrice} (``Bea'') si occupa di \textbf{giochi pervasivi}
(\emph{pervasive games}). ``Se non li avete mai sentiti nominare, non
preoccupatevi'': questo corso è ``uno dei pochi, forse l'unico in Italia'' che
li racconta esplicitamente. Li si vedrà nelle prossime lezioni, ma la definizione
di base è questa: i giochi pervasivi \textbf{utilizzano gli spazi e gli oggetti
della vita quotidiana e li trasformano in spazi e oggetti di gioco}.

Beatrice lavora come \textbf{freelance} e si occupa di diversi tipi di
progettazione: lavora sia con le \textbf{aziende} sia con gli \textbf{eventi dal
vivo}. La sua attività diretta si trova sotto l'etichetta \textbf{Minime
Differenze}\footnote{La trascrizione riporta il nome dell'etichetta come
``Minime Differenze'' e l'indirizzo del sito come ``mininedifferenze.com''; la
grafia esatta del sito non è verificabile dalla trascrizione.}, mentre i lavori
per le aziende sono ``sparpagliati'' sui siti delle varie aziende. Molte delle
cose mostrate nella sua slide sono attività proposte al pubblico ``con un
biglietto, con un gioco in formato evento'': un tema che si toccherà nel corso,
per capire come \textbf{comunicare} i propri giochi in modo che siano
comprensibili, accessibili per il pubblico e \textbf{vendibili}.

Ammette un ``grosso limite'' con i messaggi: è un po' lenta a rispondere, ``ma
rispondo sempre a tutto''. La si può contattare per email (cognome punto nome) o
su Telegram. È ``molto curiosa di vedere cosa produrrete''.

\subsection{I docenti ospiti}

Daniele presenta poi i tre docenti che il corso avrà ``solo per alcune
lezioni'', tutte verso \textbf{novembre}:
\begin{itemize}
  \item \textbf{Diego Ricchiuti}, direttore del corso di Game Design
        dell'accademia, farà una lezione sull'\textbf{RNG} (\emph{random number
        generation}, la casualità nei giochi): ``vi parleremo di RNG quando sarà
        necessario. Comunque avrete una lezione con lui, sappiatelo'';
  \item \textbf{Marta}\footnote{La trascrizione riporta il cognome come
        ``Cercasassi''; la grafia non è verificabile.} farà \textbf{due lezioni
        sul development}: quest'anno ``abbiamo deciso di fare doppietta'',
        perché l'anno scorso il suo intervento è stato ``molto, molto utile'';
  \item \textbf{Flavio}\footnote{La trascrizione riporta il cognome come
        ``Mortarino'' e, più avanti, ``Mortalino'', e il nome della sua casa
        editrice come ``Grampibert''; nessuno dei due è ricostruibile con
        certezza.} farà la lezione di \textbf{game production}; è un editore, e
        di solito partecipa anche alla tavola rotonda con gli editori della
        giornata finale.
\end{itemize}
Le slide sono già online e gli studenti avranno accesso alla cartella condivisa:
``non vi preoccupate di queste cose, ne parleremo a breve''. Ma prima di parlare
del corso, ``vogliamo conoscere voi''.

% ---------------------------------------------------------------------
\section{Le presentazioni dei partecipanti}
\label{a1:sec:partecipanti}

Le domande per la presentazione sono semplici, da affrontare ``a turno, senza
fretta, con calma'':
\begin{itemize}
  \item nome e occupazione (``se volete dirla'');
  \item perché si vogliono fare giochi analogici, e a che livello: che cosa ha
        portato lo studente al corso;
  \item che tipo di esperienze o conoscenze si hanno già nell'ambito, se si è già
        fatto qualcosa oppure no.
\end{itemize}
``Chi vuole partire senza nessun tipo di ritegno e di vergogne?''. Per far
parlare gli studenti guardandoli, Daniele interrompe la condivisione dello
schermo e ricorda lui le domande quando serve. Le presentazioni sono interessanti
anche per i commenti dei docenti, che anticipano diversi principi del corso.

\subsection{Tiziana}

Tiziana fa l'\textbf{infermiera}, ma nella vita ha fatto molti lavori:
impiegata, insegnante e ``tanti altri lavoretti''. Gioca ``da sempre'': il gioco
fa parte anche della sua vita familiare, perché suo fratello fa il programmatore
di videogiochi. Ha già \textbf{organizzato} giochi: una caccia al tesoro per far
conoscere il territorio, due scene con un delitto (giochi con omicidio da
risolvere), giochi per intrattenere le persone. Ha riscoperto di recente i giochi
da tavolo, arrivando a \emph{Carcassonne} con famiglia e amici; poiché è
difficile trovare persone con voglia di giocare nella propria cerchia ristretta,
ha scoperto il mondo delle \textbf{associazioni} ludiche, ``e lì è stato
un'esplosione'': prima in Veneto (è di Vicenza), poi a Roma, dove si è
trasferita.

Ha in mente un progetto ambizioso: far conoscere una certa realtà con la modalità
del \textbf{gioco di ruolo}, o comunque proporre un gioco suo.

\paragraph{Il commento: gli strumenti, non il gioco dei sogni.} Ruggero coglie
l'occasione per un'avvertenza che vale per tutti, visto che molti studenti
arrivano con idee già pronte. Nel corso si lavorerà a un progetto, ma
l'obiettivo dei docenti è dare \textbf{strumenti} da applicare anche
\emph{fuori} dal progetto, ai giochi che gli studenti vorranno fare un domani:
``che non siano esercizi per il corso, ma per voi in primis''. Quindi: ``non
rimaneteci male se non riuscirete a lavorare al vostro gioco dei sogni in questo
corso''. Anche perché ``sarebbe impossibile, in due mesi, fare il gioco dei
sogni''. Il corso darà invece la possibilità di fare un'esperienza ludica
insieme e di imparare a replicarla sulle proprie idee: ``le idee non buttatele,
anzi, il contrario: tenetele''.

\subsection{Nicola}

Nicola è uno \textbf{studente}. Sei mesi prima ha fatto un colloquio con Daniele
perché stava valutando di passare all'accademia, e ha deciso, ``nel dubbio'', di
cominciare da questo corso breve per capire come funziona la formazione;
eventualmente si iscriverà al corso di game design vero e proprio. Gli interessa
in particolare il lato dell'esperienza del giocatore, su cui è ``molto critico''
(qui usa la parola ``divertimento'', e Daniele lo ferma con un ``ne parleremo'':
è una delle parole che il corso vieta, come si vedrà nella
sezione~\ref{a1:sec:parole-vietate}).

Ha 21 anni ed è \textbf{giocatore di ruolo da circa otto}. Prima ancora di
conoscere i giochi di ruolo si era creato un sistema tutto suo, ``pessimo,
disgustoso, terribile'', un sistema a percentuale con il d100 ``terrificante'';
poi ha scoperto i giochi di ruolo da tavolo, i classici (\emph{Dungeons \&
Dragons}, \emph{Daggerheart}, \emph{Vampire}\footnote{La trascrizione riporta
``DD'', ``Daggerart'' e ``Vampire'': quasi certamente \emph{Dungeons \&
Dragons}, \emph{Daggerheart} e \emph{Vampire: The Masquerade}.}). Gioca anche a
\textbf{card game} e ad altri giochi sociali, pur avendo sempre avuto difficoltà
a trovare gruppi, ``avendo vissuto in un buco dimenticato dal mondo''. Come
giocatore è ``l'eterno narratore'' del gruppo, quello che prepara la storia e si
costruisce le regole. Come autore ha fatto soprattutto giochi su carta, ``perché
sono incapace a programmare, per adesso''.

Daniele osserva che il tema del giocatore e del suo rapporto con il gioco è ``un
argomento importante'', anzi ``un vaso di Pandora enorme'' (espressione che
tornerà molte volte nella lezione), e che il corso tratterà anche che cos'è
l'analog game design in generale e quanti tipi di prodotti esistono.

\subsection{Rita}

Rita\footnote{Daniele la chiama prima ``Nino'' (probabilmente un errore della
trascrizione) e poi la ringrazia come Rita.} ha 30 anni e si collega da una
regione del Sud, dove ``ancora si muore di caldo''. Ha studiato
all'\textbf{Accademia di Belle Arti di Napoli}, \textbf{design della
comunicazione visiva}: si occupa dell'aspetto più imprenditoriale della
progettazione artistica. Non volendo progettare per un'impresa (``fare loghi,
robe del genere''), vorrebbe applicare la propria creatività alla realizzazione
di giochi analogici.

Non è una giocatrice di ruolo e non lo è mai stata: è una giocatrice
\textbf{tradizionale}, cresciuta con Monopoly, la dama e giochi tradizionali
cinesi. Ha già \textbf{idee, prototipi e progetti} realizzati, e vuole capire
come renderli \textbf{prodotti per la vendita} e se questo è il campo adatto a
lei; magari anche confrontarsi con i colleghi. La sua progettazione è orientata
al \textbf{sociale} e all'\textbf{educativo}: ``responsabilizzare un po' le nuove
generazioni su come sta andando il mondo e su come potrebbe andare, visto che è
tutto nelle nostre mani''.

\paragraph{Il commento: siate curiosi.} ``Pane per i tuoi denti'', le risponde
Daniele, soprattutto per i LARP e i giochi pervasivi: ``vi apriremo un po' la
testa da questo punto di vista''. E aggiunge un consiglio per tutti: ``siate
curiosi'', perché imparare da altri ambiti dà moltissimo. Ognuno è lì per fare
qualcosa di preciso, ``però siate aperti anche ad altre esperienze''.

\subsection{Tancredi}

Tancredi ha 21 anni e non ha ancora grande esperienza lavorativa: qualche
``piccolo progettino'' che non è andato come voleva. Gli interessa il mondo dei
giochi, digitali e non, con un approccio \textbf{educativo} e \textbf{sociale}.
È arrivato soprattutto dai videogiochi, a cui gioca da quando era piccolo; ai
giochi da tavolo si è avvicinato solo nell'ultimo anno, grazie a una realtà
ludica, perché in famiglia ``non siamo molto da giochi''. Ha realizzato qualche
piccolo videogioco con vari engine: prima quelli senza un vero linguaggio di
programmazione, ``più a nodi'', poi qualcosa di più complicato, e ora sta
lavorando a un piccolo prototipo.

\paragraph{Il commento: niente di quello che vedrete è sprecato.} ``Non butterai
via niente di analog game design'', gli risponde Daniele, ``perché il 99 per cento
delle cose che vedremo si ritrovano quasi uno a uno nel mondo dei videogiochi''.
Tanto che nei corsi triennali dell'accademia (anche in quello di narrative
design, dove insegna Ruggero) il primo anno si parte dai \textbf{giochi da tavolo
e di ruolo}, proprio per spiegare meccaniche più complesse che servono poi per i
videogiochi. Ruggero aggiunge che vale anche il contrario: alcuni dei
\textbf{framework} di game design che si vedranno nelle lezioni tecniche vengono
dal videogioco e sono stati riapplicati all'analogico. ``Evitare la
crossmedialità qui è impossibile''. Anzi, tutte le conoscenze ``mediatiche
adiacenti'' (musica, letteratura, cinema) saranno utili.

\subsection{Lorenzo Vergata}

I Lorenzo in classe sono due; il primo a presentarsi è \textbf{Lorenzo
Vergata}, 26 anni, \textbf{dottorando al Politecnico di Milano}. Gioca ai
videogiochi da quando aveva due anni e mezzo, ma ha giocato al suo primo gioco da
tavolo ``credo nove mesi fa''. Gli fa piacere sentire che molti in classe sono
interessati ai giochi \textbf{sociali} ed \textbf{educativi}, un settore che
interessa anche a lui e che ``può essere esplorato molto'', anche con qualche
aiuto dell'intelligenza artificiale per esperienze più ricercate.

Di esperienza nell'analogico dice di avere fatto ``il minimo sindacale''. Ma
racconta che diciotto anni prima, da bambino, era ``più game designer di
adesso'': prendeva un gioco e ne cambiava le regole, anche al biliardino (``no,
ora introduciamo questa regola''). Gli piace sperimentare, creare regole e
vedere che cosa accade in sé e negli altri. Soprattutto vorrebbe ``una chiave di
lettura per poter rimanipolare la realtà'': tre settimane prima è stato a fare la
\textbf{vendemmia in Francia} e ha visto tantissime dinamiche che potrebbero
diventare un gioco da tavolo o un videogioco. Per esempio: finire contro un rovo
pieno di semi che si attaccano ai vestiti lo rallentava nella raccolta
dell'uva, mentre in un altro filare era più rapido. ``Si può fare pure una
dinamica'', e vorrebbe gli strumenti per farne giochi.

Arriva subito la battuta: una ``vendemmia gamificata'' da proporre in Oltrepò
Pavese (``io sono dell'Oltrepò Pavese, quindi devo per forza bere vino'')\footnote{La
trascrizione dice ``nell'altro papavese'' e non chiarisce chi pronunci la
battuta.}. Daniele conferma che quest'anno il corso avrà ``tante cose sul
sociale'': ``ci piace parecchio''.

\subsection{L'altro Lorenzo}

Il secondo Lorenzo (``secondo, ma non meno importante'') è uno \textbf{scrittore
di racconti brevi} e sta scrivendo il suo primo romanzo. Gli è sempre piaciuto
raccontare storie, e da più grande scriverle; gioca ai videogiochi ed è ``un
amante delle storie e della cultura pop''. Negli ultimi anni è entrato nel mondo
del \textbf{gioco di ruolo da tavolo} e vorrebbe unire le due passioni, le
storie e il gioco di ruolo, magari facendone un lavoro; gli interessa anche lo
\emph{strategic writing}. Daniele gli ricorda che ci sarà una lezione proprio
sulla \textbf{fiction}.

\subsection{Angelo}

Angelo si collega senza videocamera dal lavoro: fa il \textbf{birraio} (in un
piccolo birrificio in Brianza, a Bernareggio, come racconterà dopo la pausa) e,
con il caldo, ``vi risparmio l'orrore''. Per molto tempo era convinto che i
giochi da tavolo fossero solo i classici (Monopoly, Risiko) o i party game; poi,
quattro o cinque anni fa, ha ``sbattuto la faccia'' contro giochi davvero
appassionanti, con una specie di narrazione dietro che li rendeva ancora più
coinvolgenti, e non ha più smesso. È uno che ``ha sempre un botto di idee, ma poi
non mette in pratica assolutamente nulla'': cercava un corso che gli desse le
basi e soprattutto lo \textbf{motivasse} a provare a fare qualcosa.

\paragraph{Il commento: bisogna buttarsi.} Daniele racconta che anche l'anno
prima un suo caro amico si era iscritto proprio per questo motivo. Il corso è
breve, ma cerca di portare gli studenti ``da 0 a 100'' e di far realizzare un
\textbf{prodotto} da presentare alla fine. Il punto è che i giochi ``bisogna
farli, bisogna buttarsi''. E, venendo dalla narrazione, Daniele lo paragona allo
scrivere: si scrive e si riscrive, e allo stesso modo si fa \textbf{playtest}, si
rifà, si aggiorna. ``Non è che si fanno così''.

\subsection{Gianmarco (dopo le slide)}

Gianmarco arriva a lezione iniziata e si presenta più tardi, alla fine della
prima parte, dopo che i docenti hanno chiuso le slide organizzative e la
discussione sull'intelligenza artificiale (``Gianmarco, ti giuro, ti facciamo
presentare''). Lavora nel \textbf{team di innovazione} di un'azienda italiana.
Vuole fare giochi perché è una sua passione, che lo gratifica, e in passato ci ha
provato ``in maniera scomposta''. Cerca un \textbf{metodo} e un confronto con
persone più esperte: ``studiare solo dei manuali è molto sterile'', ``stare
chiusi nel pertugio non aiuta'', ed ``essere autodidatta ha estremamente i suoi
limiti''. Come esperienza è un giocatore e ha letto ``un po' di materiale, di
blog, di manuali'', ``niente di rilevante''.

Manca solo Giordano, che si presenterà alla lezione successiva.

\paragraph{Uno sguardo d'insieme.} La classe è molto varia: un'infermiera,
una designer della comunicazione visiva, un birraio, un dottorando, uno
scrittore, un impiegato nell'innovazione, due studenti ventunenni. Alcuni
vengono dal gioco di ruolo, altri dai giochi tradizionali, altri dai
videogiochi; quasi tutti hanno scoperto i giochi da tavolo moderni da pochi anni
o da pochi mesi. E almeno tre studenti dichiarano un interesse per i giochi
\textbf{sociali} ed \textbf{educativi}. Come si vedrà nella
sezione~\ref{a1:sec:scopo}, i docenti considerano questa varietà una risorsa.

% ---------------------------------------------------------------------
\section{I materiali e gli strumenti del corso}
\label{a1:sec:materiali}

Prima delle slide sul corso, Daniele sbriga un po' di ``clean up casalingo''.

\paragraph{Il form iniziale.} C'è un \textbf{form} da compilare ``entro e non
oltre il 15 settembre'': il QR code è nelle slide, e la segreteria lo ricorderà
su Discord. Richiede circa 10--15 minuti; si può fare subito o in giornata, ``e
ve lo togliete di torno''. È un form standardizzato per tutti i corsi (``non
siamo io e Ruggero a farvi passare la tortura del test''), diviso in due parti:
\begin{itemize}
  \item una parte di \textbf{conoscenza del candidato}: quanto tempo si ha a
        disposizione per il corso, come ci si organizza con il lavoro, che cosa
        piace giocare. È la parte che interessa ai docenti, perché li aiuta a
        organizzare le lezioni;
  \item una parte di \textbf{test logico}, che i docenti non hanno scritto (``io
        non sono responsabile del form'', ``no, no, nemmeno io'').
\end{itemize}

\paragraph{La cartella condivisa.} Tutti gli studenti vengono aggiunti alla
cartella del corso su \textbf{Google Drive}, che si trova tra i file
``condivisi con me''. Gli studenti ne sono \textbf{gestori dei contenuti} e
possono caricarci quello che vogliono. Dentro ci sono:
\begin{itemize}
  \item le \textbf{attività in classe};
  \item le \textbf{dispense};
  \item i \textbf{materiali didattici}, con le slide (quelle della lezione sono
        già lì);
  \item le cartelle dei \textbf{project work} futuri;
  \item i \textbf{task}, che ``saranno pochi, ma ci saranno'';
  \item il \textbf{calendario studenti}, ``importantissimo'', con due schede: il
        calendario delle lezioni e quello del project work.
\end{itemize}

\paragraph{Discord.} Il corso ha anche un \textbf{server Discord}, a cui tutti
sono stati invitati; ci sono i docenti (Daniele, Ruggero, Diego e altri), che si
possono taggare. Di solito è lì che la segreteria pubblica le informazioni, per
esempio se una lezione ``dovesse per caso essere cancellata''. Nelle domande
finali Ruggero aggiungerà che c'è un canale apposito per il corso, che gli
studenti possono usare liberamente per organizzarsi tra loro.

\paragraph{Le reaction di Zoom.} Durante le lezioni su Zoom si possono usare le
\textbf{reaction} (il pollice in alto, il cuoricino, la mano alzata): un modo
rapido per far capire ai docenti se si è capito, se si vuole intervenire. Nel
frattempo arrivano i primi pollici.

% ---------------------------------------------------------------------
\section{Lo scopo e il metodo del corso}
\label{a1:sec:scopo}

La prima serie di slide (che Daniele, per un lapsus, presenta come ``cos'è il
corso di pulire il pavimento di casa'') spiega che cosa il corso vuole fare e
come.

\subsection{Diventare autori analogici}

Lo \textbf{scopo del corso} è fornire le competenze per \textbf{diventare autori
analogici} e stimolare la creatività e l'\textbf{autorialità} di ciascuno: ``voi
imparerete a diventare autori di giochi analogici''. La differenza tra
\emph{game designer} e \emph{game developer} verrà chiarita più avanti (sezione~\ref{a1:sec:design-development}),
ma la parola chiave qui è \emph{autorialità}: ``vogliamo voi, vogliamo le vostre
idee''. Magari non si potrà usare l'idea del gioco della propria vita, ma si
potranno portare nelle lezioni le proprie idee e le proprie sensibilità.

Gli obiettivi elencati nella slide sono questi:
\begin{itemize}
  \item dare \textbf{maggiore consapevolezza nel design e nello sviluppo}. In
        Italia, purtroppo, ``si fa design a volte in maniera non consapevole''.
        Consapevole vuol dire avere almeno un'introduzione alle basi del design e
        dello sviluppo, e fare le cose ``con metodo, con raziocinio, con
        consapevolezza'';
  \item fornire gli \textbf{strumenti essenziali} per un design consapevole;
  \item stimolare la \textbf{sperimentazione} e la \textbf{trasversalità
        mediatica}. Non si chiede di fare ``il \emph{Root} della situazione o il
        \emph{Carcassonne} della situazione''\footnote{La trascrizione dice ``il
        lumaven della situazione o il carcassone della situazione''; il primo
        titolo non è ricostruibile con certezza. Il senso è chiaro: non si
        chiede di rifare un grande classico del gioco da tavolo.}, ma un
        project work \textbf{calibrato sulla durata del corso}, e però
        \textbf{sperimentale} e ``assolutamente transmediale'';
  \item portare \textbf{numerosi esempi} di tipologie di gioco, dall'analogico al
        tavolo, al pervasivo ``e oltre''. Lo si vedrà nella lezione sui
        \textbf{casi studio} del venerdì successivo: ``voi non avete idea di
        quante cose esistono, bellissime, sconosciutissime e favolose''. Daniele
        confessa che preparando e aggiornando le slide i colleghi gli hanno
        mostrato cose che lui stesso non sapeva nemmeno che esistessero;
  \item dare gli strumenti per \textbf{analizzare} i giochi altrui. Ruggero lo
        sottolinea: se il corso riesce nel suo intento, gli studenti sapranno
        non solo fare i propri giochi ma anche analizzare giochi già fatti, che è
        ``una delle skill importanti anche nella parte di ricerca e sviluppo''.
        Si vedranno casi studio e giochi pubblicati, sempre con l'obiettivo di
        capire \textbf{che cosa fanno e perché lo fanno}, per poi fare meglio il
        proprio lavoro;
  \item imparare a \textbf{fare giochi insieme}. Esistono gli autori che lavorano
        da soli (i \emph{solo developer}), ``ma sappiate che anche l'autore che
        lavora da solo poi si deve interfacciare con il pubblico, con un editore,
        magari con un editor, con un publisher''. ``Non lavorerete mai da soli'':
        bisogna imparare a lavorare con gli altri.
\end{itemize}

\subsection{Un corso intersezionale}

Il corso è \textbf{intersezionale}: \textbf{non ci sono lezioni dedicate a una
singola tipologia di gioco}. Non ci sarà ``la lezione sul gioco da tavolo'',
``la lezione sul gioco di ruolo'', ``la lezione sul gioco pervasivo''. Ci saranno
lezioni su temi diversi, ciascuna pensata per essere utile a \textbf{tutte e tre
le sfaccettature} del corso; alcune saranno più importanti per il gioco da
tavolo, altre per il pervasivo, altre per il gioco di ruolo, ma nessuna sarà
esclusiva. I docenti vogliono ``contaminazioni e applicazioni assolutamente
trasversali''.

\subsection{Un corso cross-generazionale}

È anche un corso \textbf{cross-generazionale}: i gradi di esperienza di gioco
quest'anno sono abbastanza allineati, ma le esperienze di vita, le età, le
estrazioni sociali sono varie. ``Ed è un pro, non è un contro. Più diversificata
è la classe, più imparate meglio anche voi''. La slide ha una piccola nota per
fasce d'età:
\begin{description}
  \item[Under 30] Concentratevi sulla vostra \textbf{crescita}. Evitate di
        confrontarvi con le persone più grandi che magari ``ce la fanno meglio''
        in alcune cose.
  \item[Over 30] Concentratevi sul \textbf{consolidamento}, ma siate
        \textbf{umili}: i docenti diranno cose che faranno ``storcere il naso''.
        ``Bisogna un po' reimparare per imparare'': bisogna sapersi adattare.
\end{description}
Ruggero contestualizza: lo dicono perché in passato hanno visto studenti che,
essendo più adulti o avendo giocato di più, pensavano che l'esperienza
\emph{di gioco} si trasformasse automaticamente in competenza \emph{di
creazione} di giochi. Non è così. E non bisogna sottovalutare chi gioca meno:
proprio perché gioca cose diverse, per risolvere un problema di design specifico
``è una ricchezza incredibile avere qualcun altro che riesce a uscire dalla
scatola''.

\subsection{Un corso pratico e adattabile}

Il corso è \textbf{pratico}: gli studenti si concentreranno su un progetto di
gruppo, il \textbf{project work}, lavorando in team. Ed è \textbf{adattabile}: il
programma è definito, ma ``si aggiorna sui vostri stimoli, priorità e
necessità''. Se durante il corso i docenti vedono che serve più una cosa
dell'altra, sono ``prontissimi'' a dare più spazio a un argomento invece che a
un altro.

\subsection{Inquiry-based learning: domandate}

Il metodo è quello dell'\textbf{inquiry-based learning}, l'apprendimento basato
sulle domande. Di solito le persone dicono ``ho una domanda, magari sarà
stupida''. ``Non esistono domande stupide. Pensate che ci sia una domanda
stupida? Fatela, alzate la mano''. Vanno bene anche le
\textbf{provocazioni}: se qualcosa non torna, è meglio chiedere. La condivisione
è fondamentale.

I docenti danno \textbf{metodi, dati e conoscenze}, ma lasciano trovare le
risposte agli studenti: ``vi lasciamo un pochettino annegare nel mare prima di
darvi la soluzione, in modo tale che la possiate trovare da soli''. Di
conseguenza \textbf{non ci saranno interrogazioni sulle definizioni}: interessa
vedere \textbf{come} gli studenti applicano ciò che imparano. Se una
definizione non viene in mente, non è un problema; l'importante è saperla usare.

\subsection{Feedback, non valutazione}

``Sì al feedback, no alla valutazione''. I docenti danno \textbf{feedback
verbali}, anche sulle attività da fare a casa (i task, che non saranno tanti).
Il corso è calibrato per \textbf{studenti lavoratori}: c'è un po' da lavorare
fuori dalle lezioni, ma in modo calibrato.

Il principio è \textbf{fallire oggi per imparare per il futuro}: \emph{fail
forward}, \emph{fail to learn}. I docenti lasciano provare e fallire almeno una
volta prima di dare la soluzione: ``il classico metodo Montessori, quello
originale, non quello travisato''.

\subsection{Le soft skill e i criteri di valutazione}

Grande attenzione va alle \textbf{soft skill}: i lavori di gruppo, lo scambio
propositivo, l'atteggiamento positivo, la propositività, la proattività.
``Lavoro di gruppo'' non significa stare ogni giorno in sala studio con gli
altri: va benissimo che un gruppo si organizzi dandosi priorità e tempi in base
alla disponibilità di ciascuno. L'importante è \textbf{collaborare}, essere
\textbf{professionali}, rispettarsi a vicenda, rispettare le idee proprie e
altrui e lasciare spazio agli altri, perché i progetti ``non saranno solo
vostri, saranno di tutte le persone del vostro gruppo''.

La valutazione del project work non riguarda solo la \textbf{documentazione}
prodotta, ma anche:
\begin{itemize}
  \item la \textbf{presentazione} finale;
  \item l'\textbf{autovalutazione};
  \item la \textbf{crescita personale};
  \item l'\textbf{atteggiamento} tenuto nel corso.
\end{itemize}
Daniele cita: ``da ognuno secondo le sue capacità, a ognuno secondo i suoi
bisogni''. Lavorare in gruppo non vuol dire che, in tre, ciascuno lavori
esattamente il 33\% del tempo. Ai docenti serve piuttosto capire \textbf{chi fa
che cosa} nel gruppo, e che il processo interno sia \textbf{trasparente}.

\paragraph{Lavorare in cloud.} Per lo stesso motivo il consiglio è di
\textbf{non lavorare in locale}, ma direttamente nelle cartelle condivise del
project work e dei task. Si è liberi di fare altrimenti, ma lavorando in cloud
si dà ai docenti la possibilità di vedere \emph{come} si lavora, il proprio
``processo mentale''.

% ---------------------------------------------------------------------
\section{Le regole di comportamento e il carico di lavoro}
\label{a1:sec:regole}

Alcune slide sono di quelle che si ripetono in ogni corso (``\emph{repetita
iuvant}'', storpiato per scherzo in ``ripetita Juventus'').

\paragraph{Comportamento.} Prestare attenzione durante le lezioni; interagire;
portare rispetto; essere proattivi; alzare la mano e fare domande.
\textbf{Tolleranza zero} per pregiudizi, maleducazione, mancanza di rispetto,
\emph{cheat} e scorciatoie. Daniele ammette che ``cheat e scorciatoie'' viene da
altri corsi dove il tema è più importante, ma Ruggero lo ferma: ``non
sottovalutate questa cosa'', perché implica anche \textbf{non far fare a ChatGPT
i task}. Daniele ha una slide apposta, ``perché è un discorso molto complesso nel
2026'' (sezione~\ref{a1:sec:ia}).

\paragraph{Le lezioni sono intense.} Le lezioni possono essere \textbf{dense},
``soprattutto le mie, sappiatelo: io sono logorroico'', avverte Daniele. Bisogna
chiedere spiegazioni se serve, prendere appunti, rivedere le registrazioni. Le
lezioni sono ``delle scintille'': modi per far esplorare i temi oltre le slide.
``Andate oltre, se riuscite''.

\paragraph{Il carico di lavoro.} Le regole sui task sono queste:
\begin{itemize}
  \item i task \textbf{sono per voi, non per noi};
  \item se si ha poco tempo, è meglio farli rapidamente piuttosto che non farli:
        ``provateci'';
  \item i docenti cercano di non chiedere più di \textbf{un'ora o due a
        settimana} oltre alle ore di lezione;
  \item non tutte le lezioni avranno un task, e il \textbf{project work} è il più
        importante di tutti.
\end{itemize}

\paragraph{I feedback dei docenti.} I feedback ``si guadagnano'' man mano che si
reagisce bene: non bisogna aspettarsi ``dei papiri'' sui task. Se i docenti
hanno voglia di dare feedback più ampi lo fanno, soprattutto se vedono che lo
studente se li merita. Ad Analog di solito non ci sono task \emph{rejected}; ma
se un docente rifiuta un task, vuol dire che si è fatto \textbf{meno del livello
minimo} richiesto dal corso. I feedback possono non essere immediati: ``siate un
pochettino impietositi dal fatto che siamo dei docenti''. E se un feedback non
arriva, o se ne vogliono di più, bisogna scrivere, ``perché siamo tre esseri
umani e quindi possiamo anche dimenticarci''.

% ---------------------------------------------------------------------
\section{Il project work e il calendario}
\label{a1:sec:project}

\subsection{Come funziona il project work}

Una prima slide con le date del project work è da aggiornare (``perdonatemi il
refuso''); Daniele la corregge al volo e rimanda alla slide successiva, più
precisa. Il funzionamento è questo:
\begin{itemize}
  \item i \textbf{gruppi li formano i docenti}. Gli studenti sono nove, quindi si
        faranno \textbf{tre gruppi da tre}: una dimensione pensata anche per il
        fatto che molti sono studenti lavoratori, così da non ritrovarsi ``in due,
        a esplodere'';
  \item ogni gruppo deve prima preparare un \textbf{pitch} del proprio gioco, da
        presentare in anticipo: un'idea, una bozza del gioco;
  \item poi si lavora al gioco \textbf{nelle pause del corso}: da qui a dicembre
        ci saranno due pause, quindi non ci sarà lezione tutti i lunedì e tutti i
        venerdì;
  \item ci saranno \textbf{review intermedie} con i \textbf{tutoraggi}, per lo più
        il \textbf{sabato mattina}, dalle 11 alle 12, con Ruggero su Google Meet.
        Servono ad aiutare i gruppi a portare avanti il project work.
\end{itemize}

\subsection{Il calendario delle lezioni}

Il \textbf{calendario} è nel calendario studenti della cartella condivisa.
Attenzione: le date \textbf{possono variare}. ``Il 90\% rimarranno così'', ma
può capitare che, per motivi di lavoro, qualcuno stia male o non possa, o che
``scompaia l'elettricità da tutte le città d'Italia''. E l'ordine delle lezioni
può cambiare se i docenti vedono che agli studenti serve prima una lezione
invece di un'altra.

Ci sono \textbf{due pause importantissime}:
\begin{itemize}
  \item dal \textbf{13 ottobre al 5 novembre}: è il periodo delle grandi fiere
        (\textbf{Essen} e \textbf{Lucca}\footnote{La trascrizione dice ``Essen,
        Luca, Cose'': si tratta con ogni probabilità della fiera Spiel di Essen,
        in Germania, e di Lucca Comics \& Games, che si tengono entrambe in
        ottobre.}), ``un periodo un po' frenetico'', e cominciano anche i corsi
        accademici, quindi la pausa permette ai docenti di respirare. Per gli
        studenti è la pausa più importante: quasi un mese in cui sviluppare la
        parte importante del gioco;
  \item dal \textbf{24 novembre all'11 dicembre}, dopo qualche lezione a novembre.
\end{itemize}

E ci sono \textbf{due date da evidenziare}:
\begin{itemize}
  \item il \textbf{25 settembre}: nei primi 5--10 minuti della lezione (che per il
        resto sarà tenuta da Beatrice) i docenti annunceranno i tre gruppi;
  \item il \textbf{12 ottobre}: il \textbf{pitch} dei giochi, con l'idea
        embrionale. Tra il 25 settembre e il 12 ottobre ogni gruppo prepara un
        breve pitch in formato slide; quel giorno i docenti lo rivedono e danno
        feedback: se il gioco va bene, se non va bene, come svilupparlo.
\end{itemize}
Dal 12 ottobre fino all'inizio di dicembre si sviluppa il gioco, con i
tutoraggi: ``sfruttate questi momenti di pausa''.

\subsection{Il calendario del project work}

La slide successiva, ``reale e aggiornata'', riporta il calendario del project
work (tabella~\ref{a1:tab:calendario} e figura~\ref{a1:fig:calendario}).
``Per ora non pensate al project work: ci penserete dal 25 settembre in poi''.

\begin{table}[htbp]
\centering
\begin{tabular}{@{}lp{9.2cm}@{}}
\toprule
\textbf{Data} & \textbf{Che cosa succede} \\
\midrule
ven 25 settembre\footnotemark & lezione: annuncio dei gruppi e del
\emph{mandate}, cioè di che cosa si dovrà fare \\
sab 3 ottobre, ore 11 & prima review e tutoraggio con Ruggero (dopo una
decina di giorni di lavoro) \\
lun 12 ottobre & pitch del gioco durante la lezione \\
13 ottobre -- 5 novembre & prima pausa: sviluppo del gioco \\
sab 7 novembre & tutoraggio \\
sab 14 novembre & tutoraggio \\
24 novembre -- 11 dicembre & seconda pausa: sviluppo del gioco \\
ven 4 dicembre, entro mezzanotte & consegna del gioco sul Drive (non c'è
lezione) \\
ven 11 dicembre & i giochi vengono visti a lezione \\
14 dicembre & giornata finale, con la tavola rotonda con gli editori \\
\bottomrule
\end{tabular}
\caption{Il calendario del project work come lo presenta Daniele. Le date
possono variare.}
\label{a1:tab:calendario}
\end{table}
\footnotetext{Daniele dice ``lunedì 25 settembre'', ma nel 2026 il 25 settembre
è un venerdì; il giorno della settimana nella tabella è corretto di
conseguenza.}

\paragraph{Perché la consegna è il 4 dicembre.} Il 4 dicembre non c'è lezione:
si deve solo consegnare il gioco, entro mezzanotte, sul Drive. La presentazione
a lezione è l'11 dicembre, ma i tre docenti hanno bisogno di \textbf{una
settimana} per rivedere il gioco in privato e valutare anche la
\textbf{documentazione}; per questo la consegna non coincide con la
presentazione. Che cosa voglia dire esattamente ``consegnare il gioco'' verrà
spiegato più avanti.

\paragraph{La giornata finale.} Il 14 dicembre c'è la \textbf{giornata finale},
con una \textbf{tavola rotonda con gli editori}: di solito viene invitato Flavio,
che è anche docente di game production, più qualche altro editore, ``vediamo
quest'anno chi riusciamo a portarvi''.

\begin{figure}[htbp]
\centering
\begin{tikzpicture}[
    font=\footnotesize,
    x=0.1cm,
    evento/.style={draw, thick, rounded corners=2pt, fill=white,
                   text width=1.8cm, align=center, inner sep=2pt},
    pausa/.style={fill=gray!20, draw=gray!60},
    tacca/.style={thick}
  ]
  % asse: giorno 0 = 25 settembre; 1 unità = 1 giorno
  \draw[pausa] (18,-0.35) rectangle (41,0.35);
  \draw[pausa] (60,-0.35) rectangle (77,0.35);
  \node[font=\scriptsize\itshape] at (29.5,0.18) {pausa};
  \node[font=\scriptsize\itshape] at (65,0.18) {pausa};
  \draw[-{Stealth[length=2.5mm]}, thick] (-3,0) -- (85,0);
  % eventi sopra
  \draw[tacca] (0,-0.15) -- (0,0.9);
  \node[evento, above] at (0,0.9) {25 set\\ gruppi e \emph{mandate}};
  \draw[tacca] (17,-0.15) -- (17,2.2);
  \node[evento, above, fill=gray!8] at (17,2.2) {12 ott\\ \textbf{pitch}};
  \draw[tacca] (70,-0.15) -- (70,0.9);
  \node[evento, above, fill=gray!8] at (70,0.9) {4 dic\\ \textbf{consegna}};
  % eventi sotto
  \draw[tacca] (8,0.15) -- (8,-0.9);
  \node[evento, below] at (8,-0.9) {3 ott\\ tutoraggio};
  \draw[tacca] (43,0.15) -- (43,-0.7);
  \draw[tacca] (50,0.15) -- (50,-0.7);
  \draw (43,-0.7) -- (50,-0.7);
  \draw[tacca] (46.5,-0.7) -- (46.5,-0.9);
  \node[evento, below] at (46.5,-0.9) {7 e 14 nov\\ tutoraggi};
  \draw[tacca] (77,0.15) -- (77,-0.9);
  \node[evento, below] at (77,-0.9) {11 dic\\ giochi a lezione};
  \draw[tacca] (80,-0.15) -- (80,2.2);
  \node[evento, above] at (80,2.2) {14 dic\\ giornata finale};
\end{tikzpicture}
\caption{Una linea del tempo del project work, ricostruita dalle slide (scala
approssimativa: dal 25 settembre al 14 dicembre). Le bande grigie sono le due
pause senza lezioni, in cui i gruppi sviluppano il gioco.}
\label{a1:fig:calendario}
\end{figure}

Finita questa parte, Daniele apre alle domande: ``Domande sul calendario? Sulle
slide che ho fatto? Considerazioni? Non abbiate timore, io non ci credo''. E la
domanda arriva, sull'unico tema che aveva lasciato da parte.

% ---------------------------------------------------------------------
\section{L'intelligenza artificiale nel corso}
\label{a1:sec:ia}

\subsection{La slide di Daniele}

Daniele aveva già presentato una slide che definisce ``importantissima''. ``Lo so
che è paradossale, me ne rendo conto. Lo so che magari mi darete del luddista,
dell'anti-tecnologia: non è così. Uso l'intelligenza artificiale per lavoro. Ma
in questo corso, vi prego, \textbf{non usate l'intelligenza artificiale}''. La
richiesta vale per le attività, per i task e \textbf{anche per la
prototipazione}. Le ragioni:
\begin{itemize}
  \item ``voglio che siate voi a fare il gioco, voglio che siate voi a imparare a
        farlo''. Daniele parla da utilizzatore: prima di ChatGPT, che ha smesso
        di usare, oggi soprattutto di Claude, di cui ``si sottovaluta quanto
        possa essere incerto, insicuro, impreciso, non all'altezza'';
  \item nella game industry, \textbf{i publisher rifiutano automaticamente} i
        pitch e i prototipi fatti con l'intelligenza artificiale. Daniele lo ha
        visto di persona in ``tantissime pitch review'' quell'estate, anche al
        \textbf{First Playable} di Firenze\footnote{La trascrizione dice ``First
        Playboy a Firenze'': si tratta del First Playable, l'evento di business
        dell'industria videoludica che si tiene a Firenze.}: anche nel mondo dei
        videogiochi i publisher che vedono pitch o prototipi fatti con l'IA
        ``rigettano automaticamente la build'', perché hanno ancora timore
        dell'IA e delle sue possibili conseguenze;
  \item ``non so come sarà il futuro'': non sa se bisognerà abituarsi a usarla
        tutto il giorno. Non è ``totalmente contrario'', purché sia usata ``con
        raziocinio''; secondo lui oggi, negli Stati Uniti, non lo è;
  \item ma in un corso che si chiama \emph{Analog Game Design} gli studenti devono
        saper fare i giochi da soli. In futuro, se vorranno farsi aiutare
        dall'IA per essere più efficienti, saranno liberissimi di farlo; durante
        il corso no, perché i docenti vogliono valutare \textbf{le loro
        competenze e capacità}.
\end{itemize}

\subsection{La domanda: che cosa vuol dire usare l'IA in un gioco analogico?}

La domanda di uno studente è precisa: che cosa si intende per ``task fatti con
l'IA''? Qual è l'uso dell'IA in un gioco analogico?

\paragraph{Ruggero: l'esperienza di chi valuta i prototipi.} Ruggero risponde
dalla sua esperienza nel settore, dove riceve prototipi da valutare per
un'eventuale pubblicazione. Distingue due usi:
\begin{enumerate}
  \item le \textbf{immagini generate}: illustrazioni, graphic design e così via.
        Per molti publisher è ``un \emph{reject} immediato, proprio immediato'';
  \item l'IA usata per il \textbf{regolamento}, per scrivere le regole e le
        meccaniche. È più difficile da individuare, ma molti publisher se ne
        accorgono comunque, perché le IA ``non sono ancora brave abbastanza da
        creare un regolamento coeso meccanicamente e scritto con un senso per il
        vostro gioco''. Quello che fanno è prendere parti di altri regolamenti,
        elaborare le strutture che vedono nei regolamenti pubblicati e
        riapplicarle. Ma molti giochi si prestano poco a un regolamento
        strutturato come quello di altri, perché \textbf{ogni gioco è un
        ecosistema meccanico a sé}. Il risultato sono regolamenti con
        riferimenti non immediati, e soprattutto con \textbf{la stessa meccanica
        descritta in modo diverso in pagine diverse}: due difetti che fanno
        perdere tempo all'editore, e in questo caso ai docenti.
\end{enumerate}
Poi aggiunge un punto ``personale, soggettivo'': i docenti sono interessati a
vedere il lavoro, la creatività, ``il frutto della vostra creazione''. E
viceversa gli studenti vogliono i commenti dei docenti: ``immaginatevi se voi
fate un gioco e noi vi diamo i commenti scritti con ChatGPT''. Altrimenti, dice
ridendo, gli studenti fanno il gioco con ChatGPT, i docenti lo danno a ChatGPT
per i commenti, ``e facciamo molto alla svelta''.

\paragraph{Daniele: perché gli studenti usano l'IA.} Daniele racconta la sua
esperienza con gli studenti, che secondo lui ricorrono all'IA in due momenti:
\begin{enumerate}
  \item quando vanno in \textbf{crisi}: in ansia da prestazione o per i tempi.
        Cercano di risolvere il problema della propria disorganizzazione
        facendolo risolvere all'IA, senza accorgersi che l'output è ``molto,
        molto scarso'';
  \item quando hanno l'impressione che l'intelligenza dell'IA sia
        \textbf{migliore} di quello che scrivono loro, o di quello che c'è in
        giro. Per Daniele è una ``mancanza di palato critico'': l'output dell'IA
        sembra buono in superficie, ma quando si va nel dettaglio e ci si fa
        delle domande ``di solito fa un po' di \emph{crumbling}'', si sbriciola.
\end{enumerate}
Lui usa Claude ogni giorno e ne conosce i limiti: ``non c'è nemmeno una frase
che io copio direttamente''. Gli fa fare il lavoro più noioso, più ripetitivo e
soprattutto \textbf{strutturato}, ``replicare delle cose come se fosse una
macro''. Lo usa per il testo e per il codice, mai per le immagini: per motivi di
\textbf{diritto d'autore}, e perché ha molti amici artisti ed è stufo di vedere
il loro lavoro ``dilaniato'' dall'IA.

Ammette che i docenti non possono controllare gli studenti; per questo
``ve lo consiglio''. E riprende quella che definisce la cosa più bella detta da
Ruggero (``ti ruberò la frase per il corso di narrative''): i docenti i feedback
li danno di persona. Daniele a volte manda 30 minuti di audio a uno studente, o
fa intere call per rivedere script, pitch e giochi: ``è il mio lavoro, mi piace
farlo e voglio farlo io, quindi mi aspetto che lo studente lo faccia lui''. Come
freelance, poi, è ``stufo marcio'' dei clienti che arrivano con qualcosa fatto
con l'IA (``guarda che bello, pubblica questo''), perché il più delle volte il
risultato è pessimo: non sanno come usarla. ``È anche una questione di
competenza''.

\paragraph{Il prototipo non deve essere bello.} Questo vale anche per il
\textbf{prototipo}, su cui si tornerà nelle lezioni di prototipazione: in un
prototipo non servono ``graficoni allucinanti''. Le grafiche bellissime fatte con
l'IA spesso \textbf{peccano di user experience}: sembrano bellissime, ma quando
il giocatore deve leggere la carta non capisce qual è il costo, qual è il
potere, qual è l'effetto.

\paragraph{Beatrice: si perde l'occasione di imparare.} Beatrice aggiunge due
motivi che valgono ``al di là dei nostri credo''. Il primo: chi usa l'IA in
sostituzione (per esempio per scrivere un testo, per accelerare i tempi o per
trovare strategie) \textbf{perde la parte in cui può imparare}. Nel game design
bisogna ``tenere insieme i pezzi'' di un regolamento, ma anche dell'esperienza
dell'utente e di come trasmetterla; i passaggi necessari (creare il regolamento,
creare le meccaniche, sperimentare, raccogliere i feedback) ``si capiscono solo
vivendoli''. E si sta pagando un corso per imparare. Il secondo, sul piano
editoriale: chi presenta a un editore un prodotto che ha come base l'IA rischia
da un lato di non avere imparato abbastanza (e di ritrovarsi ``a piedi'' sul
campo senza IA), dall'altro che l'editore se ne accorga, ed è un tema
\textbf{reputazionale}. Ci sono editori ``più illuminati di altri'', quindi
anche qui c'è un mercato da capire; ma i docenti preferiscono dire che cosa è
più sensato per il mercato come lo conoscono.

\paragraph{Ruggero: non abbiate paura del giudizio.} Ruggero vuole dare uno
``spotlight'' a un punto detto da Daniele e rimasto ``semi-ignorato'': a volte
l'uso dell'IA nasce dalla \textbf{paura del giudizio}, dal voler consegnare
qualcosa di più pulito, più strutturato possibile. ``Vi giuro, non ve ne
preoccupate''. Ruggero detesta i voti e non ragiona per voti; non valuterà
nessuno per grafiche fatte a matita. Insegna da anni, ha avuto studenti con
dislessia e discalculia: refusi ed errori non gli interessano, ``non è un corso
di editing o di proofreading''. Gli interessa che le \textbf{meccaniche siano
coerenti con lo scopo di design}. Immagini fatte a matita o con Paint ``non sono
un problema, anzi, le preferisco'' rispetto all'``ennesima immagine fatta con
ChatGPT o Midjourney''.

\paragraph{Daniele: il pensiero critico, costo e qualità.} L'ultimo intervento
è di Daniele, sul \textbf{pensiero critico}. Ragionando per assurdo (``un
pensiero di Belzebù''), immagina che un domani l'IA venga usata ovunque. I suoi
limiti, secondo lui, resterebbero due:
\begin{description}
  \item[Il costo] Tutti dicono che sta diminuendo e che diminuirà, ma non se ne
        può essere sicuri al 100\%.
  \item[La qualità] Cioè la capacità di leggere l'output. Se si dà un regolamento
        a Claude o a ChatGPT chiedendo quali problemi ha, il modello risponde con
        i problemi che, in percentuale, è più probabile che gli esseri umani
        abbiano già individuato nei dati su cui è stato addestrato. Ne derivano
        due problemi: se il problema del proprio gioco non è mai stato
        individuato prima, il modello non lo vedrà; e viceversa può segnalare
        problemi che il gioco non ha.
\end{description}
Fa l'esempio di Grammarly, che usa nella versione premium per rivedere i testi
in inglese: spesso non accetta i suoi consigli, perché se vuole scrivere un
dialogo di videogioco in \textbf{dialetto Cockney} londinese, lo deve scrivere
in Cockney. Se Grammarly dice che è sbagliato, ``sei tu che non hai capito il
motivo per cui voglio scrivere questo testo in Cockney''. La conclusione: prima
di usare un'IA per un'attività bisogna sapere \textbf{che cosa si vuole
ottenere}, qual è l'obiettivo, e quali \textbf{competenze} servono per valutare
in modo critico se l'output di un qualunque LLM è davvero valido. Daniele non
vuole dilungarsi oltre: è un argomento ``molto caldo'' per lui, che è anche
scrittore e si sente dire da tutti ``ti sostituiranno''.

Subito dopo si presenta finalmente Gianmarco (lo si è visto nella
sezione~\ref{a1:sec:partecipanti}), e poi si fa una pausa di una decina di
minuti. Daniele avverte che nella terza parte lui non ci sarà (alle sette e un
quarto ha un impegno di lavoro), e lo stesso vale per Beatrice: la parte finale
di domande e risposte sarà con il solo Ruggero, ``in ottime mani''.

% ---------------------------------------------------------------------
\section{Le figure professionali dell'industria analogica}
\label{a1:sec:figure}

Dopo la pausa (in cui si scopre che Angelo lavora in un piccolo birrificio in
Brianza e si scherza sul rischio di un \emph{drinking game}, tra la birra e la
vendemmia) la lezione riprende con le slide di Ruggero. ``Oggi è una lezione
introduttiva'': gli argomenti si vedranno ``molto dall'alto'' e verranno
approfonditi nelle lezioni successive. Ma servono già ora, per creare
\textbf{l'inizio di un linguaggio condiviso}: ci sono termini che torneranno
spesso, a cominciare da quelli del mercato e dell'industria del gioco
analogico, un'industria ``molto varia, molto caotica, anche molto mutevole'',
con tanti tipi diversi di giochi e altrettanti tipi diversi di entità che ci
lavorano.

\subsection{Una decina di etichette}

I docenti hanno raccolto ``una decina di etichette'', cioè di \textbf{ruoli} che
nell'industria sono o possono essere presenti. Due avvertenze:
\begin{itemize}
  \item nelle aziende più piccole \textbf{la stessa persona copre spesso più
        ruoli}, per forza di cose. Sarà anche il caso degli studenti, che
        lavoreranno in gruppi di tre e non di venti: è facile che ciascuno finisca
        a fare più cose. I docenti ne terranno conto, e qui entreranno in gioco
        anche competenze diverse dal game design (programmi di grafica,
        illustrazione, impaginazione, layout, \emph{proofreading}\ldots);
  \item Daniele precisa che \textbf{non si chiede agli studenti di diventare}
        illustratori o altro. Ruggero concorda, ma aggiunge che è utile sapere
        anche \emph{che cosa non è il proprio lavoro}. Se in un gruppo nessuno sa
        illustrare, ai docenti va benissimo un prototipo che mostri dove vanno
        le immagini: \textbf{placeholder}, concept, anche semplici box di testo
        con scritto ``qui vorremmo un'immagine di\ldots''. ``È importante anche
        sapere che cosa sapete fare e che cosa non sapete fare''.
\end{itemize}
La tabella~\ref{a1:tab:figure} riassume le figure; le sottosezioni seguenti le
riprendono una per una, con i commenti dei docenti.

\begin{table}[htbp]
\centering
\small
\begin{tabular}{@{}p{3.3cm}p{9.6cm}@{}}
\toprule
\textbf{Figura} & \textbf{Di che cosa si occupa} \\
\midrule
Game designer & crea le regole alla base dell'esperienza: le meccaniche, come
interagiscono tra loro e come chi gioca interagisce con esse \\
Game developer & fa da ponte tra chi fa game design e il publisher: bilancia il
prototipo, cura i modelli matematici, adatta il gioco (e il paragame) alle
necessità del publisher \\
Editor (proofreader) & revisiona e adatta i testi a partire dalla prima stesura
del designer (correzione, linguaggio inclusivo\ldots) \\
Publisher (editore) & pubblica e produce fisicamente il gioco; il suo marchio è
sulla scatola o sulla copertina \\
Sales e marketing & il lato commerciale: vendite (anche B2B) e comunicazione del
prodotto \\
Graphic designer & icone meccaniche, wireframe, UX: il lato funzionale degli
elementi visivi \\
Illustratore & le illustrazioni: il \emph{look and feel}, il world building, il
``vibe'' del gioco \\
Sensitivity reader (cultural consultant) & verifica l'accuratezza e la
rappresentazione di un tema sensibile che il team non conosce \\
Producer & coordina il progetto (quasi project management), parla con le
tipografie, verifica che i materiali siano sostenibili economicamente \\
\bottomrule
\end{tabular}
\caption{Le figure professionali dell'industria del gioco analogico presentate
nella slide di Ruggero, in sintesi.}
\label{a1:tab:figure}
\end{table}

\subsection{Game designer}

Il \textbf{game designer} (o la game designer) è chi si occupa, ``sembra banale
a dirlo'', di \textbf{creare le regole alla base dell'esperienza} che si vuole
far arrivare. Vale per tutti i tipi di giochi, non solo per quelli analogici, e
che cosa faccia concretamente in un progetto o in una serie specifica ``può
mutare tantissimo''. Nell'analogico non verrà chiesto di programmare o di
prototipare in digitale; verrà chiesto di \textbf{creare le meccaniche di
gioco}, di capire \textbf{come le meccaniche si interfacciano tra loro} e
\textbf{come le persone che giocano si interfacceranno con le meccaniche}.

\subsection{Game developer}

Il \textbf{game developer} (o la game developer) sarà l'argomento delle due
lezioni di Marta. La parola \emph{development} è nota soprattutto
dall'ambito videoludico, ma nell'analogico ha un significato molto specifico
``che non c'entra niente con la programmazione'': il game developer è
\textbf{un ponte tra chi fa game design e il publisher}. Di solito si occupa di
\textbf{bilanciare il prototipo}, di curare i \textbf{modelli matematici}, di
adattare il gioco e il \textbf{paragame}\footnote{La trascrizione dice ``il
Wix, il paragame''. Il \emph{paragame} è, in termini generali, tutto ciò che
accompagna il gioco senza esserne il sistema di regole (regolamento, scatola,
materiali, comunicazione); i docenti lo tratteranno in una lezione dedicata, e
qui il termine è solo citato.} alle necessità del publisher. La distinzione
viene ripresa e approfondita nella sezione~\ref{a1:sec:design-development}.

\subsection{Editor}

L'\textbf{editor}, a volte chiamato \textbf{proofreader} a seconda dell'azienda,
si occupa della revisione dei \textbf{testi}: non della stesura iniziale, perché
di norma, data la scala dell'industria, è chi fa game design a scrivere un
\textbf{primo draft}, e l'editor ci passa sopra. Non solo per correggere i
refusi: spesso il suo lavoro è anche \textbf{adattare} i testi, per esempio
renderli \textbf{gender neutral}. Se ne vedranno casi studio più avanti.

\subsection{Publisher}

Il \textbf{publisher}, o \textbf{editore}, è l'entità che nell'ambito
commerciale \textbf{pubblica} giochi, giochi di ruolo o giochi da tavolo. È una
figura ``più solidificata'' nel gioco da tavolo e nel gioco di ruolo da tavolo;
\textbf{nell'ambito pervasivo non esiste proprio in questa forma}, come
mostrerà Beatrice parlando dei casi studio. Se si compra un gioco da tavolo o un
gioco di ruolo in un negozio, il marchio del publisher è di solito sul retro
della copertina o sul fronte della scatola: è l'entità che, di solito,
\textbf{produce fisicamente} il gioco.

\subsection{Sales e marketing}

L'ambito \textbf{sales e marketing}, a volte affidato a un'unica persona e a
volte separato, riguarda il lato commerciale. In alcune aziende è interno al
publisher, e si occupa di parlare con le fabbriche e le tipografie, di ``far
quadrare i conti''. In altre realtà è completamente diverso, specialmente nel
\textbf{B2B} o nei giochi \textbf{didattici} e \textbf{formativi}, di cui
Daniele e Beatrice ``sapranno parlare meglio''. Il tema viene ripreso più avanti
in dettaglio (sezione~\ref{a1:sec:marketing}).

\subsection{Graphic designer e illustratore}

Ruggero presenta insieme \textbf{graphic designer} e \textbf{illustratore}
proprio perché sono ruoli molto diversi che si tende a far coincidere:
\begin{itemize}
  \item chi fa \textbf{graphic design} si occupa delle \textbf{icone meccaniche},
        del \textbf{wireframing}, della \textbf{UX}: del lato \textbf{meccanico e
        funzionale} degli elementi visivi, cioè di tutto ciò che serve
        meccanicamente nel gioco, o anche alla sua ``qualità della vita'';
  \item chi fa \textbf{illustrazione} si occupa dell'\emph{art}, delle
        illustrazioni: copre il \textbf{look and feel}, la parte di
        \textbf{world building} e il ``vibe'' che si dà al gioco.
\end{itemize}

\paragraph{Daniele: che cos'è la user experience.} Daniele interviene sulla
\textbf{user experience}, ``importantissima per me'': ha un background da
\textbf{layout artist} (impaginava magazine), un mestiere in cui la UX conta
moltissimo. Per user experience non si intende solo l'\textbf{interfaccia
grafica}, la UI, che anche nei videogiochi è disegnata da un artista. La UX è
qualcosa che può fare anche il game designer: \textbf{disegnare l'esperienza del
giocatore}, cioè ``il gesto, l'intenzione, l'intuito'' che stanno dietro
l'aspetto visivo dell'interfaccia. Per fare UX non serve saper disegnare.
L'esempio è quello della \textbf{porta}: se si progetta una porta, bisogna
indicare all'utente come si apre, come si prende la maniglia; ``quella è user
experience''. Il design della maniglia in sé, invece, è
\textbf{interfaccia}.

\paragraph{Un aneddoto sulle icone.} Ruggero aggiunge che molte persone che
fanno \textbf{graphic design ben fatto non sanno illustrare}, e spesso vale il
viceversa. Lo si nota quando le icone di un gioco sono bellissime ma non
funzionali: andando a guardare i crediti, di solito si scopre che le icone le ha
fatte la persona che ha fatto le illustrazioni. ``Comunque non andiamo adesso in
questo, perché è un altro vaso di Pandora''.

\subsection{Sensitivity reader}

Il \textbf{sensitivity reader}, o \textbf{cultural consultant} a seconda del
progetto, fa un lavoro di \emph{sensitivity reading} del gioco. È una figura
molto usata nell'ambito \textbf{formativo e didattico}. L'esempio di Ruggero: se
si fa un gioco per un ospedale su una malattia specifica, chi lo progetta non è
un medico, ma il gioco deve essere \textbf{accurato} e \textbf{rappresentativo}
delle dinamiche che quella malattia comporta. Serve una persona che conosca quel
tema, anche se non sa che cosa sia il game design. Sono figure chiamate di
solito \textbf{a progetto}, per giochi specifici, che compensano competenze che
il team non ha ``per forza di cose, perché noi non sappiamo tutto''; e a volte,
per lavoro, si devono fare giochi su temi che non si conoscono personalmente.

\paragraph{Daniele: non è un censore.} Daniele segnala un problema italiano:
alcune persone, anche online, hanno confuso il sensitivity reader con un
\textbf{censore}, ``una persona che fa censura, soprattutto in materia di
\emph{woke}''. ``Assolutamente non è così''. La parola \emph{sensitive} non
significa che il pubblico sia sensibile a un tema, ma che \textbf{l'argomento di
cui tratta il gioco è sensibile}. In un gioco sulla medicina, realistico o no,
non si possono inserire imprecisioni. Per questo esistono i sensitivity reader
``più disparati'': da quelli per la cultura indiana o afroamericana fino a quello
su ``come si allevano gli gnu in Africa''. È una persona competente su un
argomento che aiuta a essere \textbf{veritieri}. Certo, uno degli ambiti
principali in cui si usano è quello culturale, ma ``non fate l'errore di pensare
che siano dei censori''. Daniele ha lavorato con ``centinaia'' di sensitivity
reader e nessuno gli ha mai detto ``non puoi scrivere questa cosa''. Di solito
sono come degli \textbf{editor}: danno opzioni, mostrano i punti problematici,
anche quelli di cui non ci si era nemmeno accorti.

Ruggero conferma con un esempio di quando lavorava in \textbf{Danimarca} con
l'\textbf{OMS}: giochi sul \textbf{diabete di tipo 2} nei bambini, per un target
di ragazzi delle scuole medie. C'era un team dell'OMS specializzato nella
comunicazione sul diabete, per garantire che il gioco fosse accurato e ne desse
una rappresentazione adatta al target. Il sensitivity reader aiuta a rendere il
gioco più accurato e ``più tematicamente coeso''.

\subsection{Producer}

Il \textbf{producer} è una figura a volte interna al publisher. In un progetto
grande è la figura di riferimento, ``quasi'' un project manager. In un progetto
piccolo si occupa spesso anche di parlare con le \textbf{tipografie} e di fare la
\textbf{lista dei materiali}, per assicurarsi che i materiali non siano solo
funzionali al gioco ma anche \textbf{economicamente sostenibili} da produrre:
``magari è un gioco bellissimo, ma poi costa più di quanto potete permettervi, e
nessuno vi comprerà mai un card game a 200 euro''.

\subsection{Perché conoscere queste figure}

Ruggero chiude la slide ``tornando a ombrello'': perché parlarne, se gli studenti
sono lì per fare game design? Non solo per saper dare un \textbf{nome} alle
figure. Non c'è niente di male a creare un gioco solo per sé; ma chi vuole fare
giochi a livello lavorativo, perché li giochino altre persone, ``magari più
persone possibile'', \textbf{si interfaccerà con queste figure}, e sapere che
cosa fanno e come parlarci è parte importante del percorso. Nel corso se ne
toccheranno diverse: il development con Marta, la produzione con Flavio, e
lezioni sulla UX, sul paragame, sul graphic design. ``Abituatevi mentalmente:
non siete in solitaria contro il mondo''.

% ---------------------------------------------------------------------
\section{Marketing e sales}
\label{a1:sec:marketing}

Arrivato all'angolo della slide con marketing e sales, Ruggero passa la parola
ai colleghi: ``io è meglio che non mi fate parlare di marketing e sales'', anche
per esperienze personali.

\subsection{Beatrice: i giochi trattati come prodotti}

Beatrice premette che le due voci sono in parte autoesplicative. Nel mondo
capitalista in cui viviamo, i giochi sono ``strumenti molto potenti per fare
\emph{n} cose'', ma sono trattati come \textbf{prodotti}, e il corso ragionerà
anche in quest'ottica. Nel \textbf{marketing} delle industrie ludiche e
videoludiche ci sono varie professionalità che gli studenti probabilmente già
conoscono: dal \emph{copy} alla grafica, tutto ciò che riguarda il
\textbf{vendere} il gioco o l'attività di gioco.

Il lato \textbf{sales} è forse la distinzione meno nota. Come in tante aziende, è
la parte che si occupa delle vendite \textbf{B2B}, cioè tra aziende: fare
accordi per vendere prodotti ludici a diverse realtà, per esempio ai
\textbf{distributori}, o a una casa editrice che vende alle librerie, che poi
vendono al consumatore finale. Chi si occupa di sales nel mondo ludico capisce
come vendere \textbf{alle aziende}. I due mondi si parlano spesso, e soprattutto
nelle realtà piccole si trovano figure che fanno ``un po' sales, un po'
marketing''.

\subsection{Daniele: che cos'è (e che cosa non è) il marketing}

Daniele il marketing si è trovato ``purtroppo'' a farlo, e lo fa anche per
l'accademia. ``Purtroppo'' perché in Italia è mal compreso: molti lo
identificano con l'\emph{advertisement} online. Non è così: il marketing è una
\textbf{disciplina molto più antica}, e in Italia ``abbiamo molto dimenticato''
come si fa il marketing tradizionale. Alcune cose da ricordare:
\begin{enumerate}
  \item \textbf{Il marketing non dice che cosa pubblicare, ma come pubblicarlo}.
        Esistono, nei videogiochi e nei giochi da tavolo, realtà che fanno
        prodotti \textbf{marketing driven}: l'esempio sono gli
        \textbf{hyper-casual game} per smartphone, quelli con il personaggio che
        va a destra e a sinistra e l'esercito che si allarga. Possono anche
        guadagnare, ma ``sono un po' la morte della creatività''. Il marketing
        ideale \textbf{prende un prodotto pensato da altri e gli dà le maggiori
        chance di essere venduto}: cerca il pubblico ideale per il prodotto che
        si ha, senza doverlo cambiare. Anzi, il mondo videoludico, e ludico in
        generale, ``è lastricato'' di errori di dipartimenti marketing che hanno
        cambiato completamente non il gioco, ma il modo in cui l'hanno
        pubblicizzato, e il gioco è fallito.
  \item \textbf{Usare strumenti del marketing non significa fare marketing}. Game
        design e game development a volte usano strumenti che vengono dal
        marketing (gli studenti ne vedranno con Beatrice, con Daniele e con
        Marta).
  \item \textbf{Il marketing è un gerundio}: si fa. È un'azione che si compie per
        comunicare; la comunicazione è attiva, e il marketing è l'insieme delle
        attività con cui si comunica il proprio gioco, dal \emph{copy} ai social.
        Può includere il digital marketing, ``ma non è solo quello''.
  \item \textbf{Il marketing esiste anche se si crede che non esista}. In Italia
        molti non vogliono ``fare marketing'' perché lo associano al lucro. Ma
        il marketing è \textbf{lo studio di come gli esseri umani comprano e
        vendono}, ed è innato in noi: ci sono logiche di cui siamo vittime, e
        ``anche se sei il bastian contrario, l'anticapitalista'', magari si
        compra comunque su Amazon perché fa comodo. Sono meccanismi che
        esistono, ``affascinanti anche, a volte''.
\end{enumerate}
Gli studenti non dovranno fare marketing, ma dovranno ricordare alcuni di questi
concetti, soprattutto nelle lezioni con Marta, che tratta argomenti legati al
marketing. E, chiude Daniele, come game designer e game developer ``per lo più
speriamo che il marketing non ci rovini i giochi'', soprattutto nel videoludico.

% ---------------------------------------------------------------------
\section{Game design e game development}
\label{a1:sec:design-development}

Ruggero torna sulla distinzione più importante della slide, ``proseguendo con
questa strada'':
\begin{description}
  \item[Game design] è chi \textbf{progetta l'esperienza} tramite la
        \textbf{creazione delle meccaniche} di gioco.
  \item[Game development] è chi fa \textbf{mediazione} tra la persona che scrive
        il game design e la \textbf{casa editrice}.
\end{description}

\subsection{I red flag dei designer}

Ruggero lo ripete perché lo vede ogni giorno nel suo lavoro: tra giochi
pubblicati e prototipi, ne valuta ``quasi uno al giorno, se faccio un conto
annuale'', e solo una minima percentuale arriva alla pubblicazione. Molte volte
i \textbf{red flag} vengono proprio da designer che non hanno presente questa
definizione, e che cominciano a mettere \textbf{clausole sulla non modifica del
gioco}: il gioco deve restare com'è. ``Questa cosa è impossibile con un
publisher, spesso''. Il publisher, per forza di cose, potrà chiedere di cambiare
la grafica, di reimpaginare il regolamento, di cambiare i materiali perché
troppo costosi: chiederà dei \textbf{compromessi}. E la persona che, all'interno
del publisher, fa questo lavoro è il \textbf{developer}.

Sta al designer capire \textbf{su che cosa del proprio gioco si può trovare un
compromesso e su che cosa no}. Publisher diversi lavorano in modo diverso, e il
designer ha sempre l'autorità di dire ``no, non voglio pubblicare questo gioco
se queste sono le modifiche''. Ma deve sapere chi fa il lavoro di development,
perché se il suo obiettivo è la \textbf{pubblicazione tradizionale} si
interfaccerà con questa figura.

\subsection{Daniele: una distinzione che manca al videogioco}

Daniele sottolinea che questa distinzione è \textbf{tipica del mondo analogico}
ed è ``molto interessante''. In Italia la parola \emph{developer} è diventata
negli anni quasi un sinonimo di programmatore di videogiochi, o di chiunque
lavori nel videoludico. Secondo lui c'entra anche il fatto che qualche anno fa
il sito \textbf{Gamasutra} ha cambiato nome in \textbf{Game Developer}, e la
parola è diventata un ombrello per tutte le professioni del videogioco. ``Mi fa
un po' ribrezzo'': il game developer dovrebbe essere un ruolo di design, un
designer che, come nell'analogico, fa \textbf{mediazione}. Nel videogioco i
publisher hanno in effetti i loro developer interni, che adattano il gioco, per
esempio alle console e ai loro kit di sviluppo. Ma la distinzione dovrebbe
esistere anche nel videogioco, perché è ``difficilissimo'' avere la mentalità
per fare le due cose insieme: chi fa design ha un focus che raramente si può
sovrapporre a quello di chi fa mediazione.

A Daniele piace la parola \emph{develop}, \textbf{sviluppare}, nel senso di
prendere qualcosa che è stato progettato e \textbf{renderlo possibile nel mondo
reale}. Ci sono progetti che hanno bisogno di uno sviluppo che ``polishi''
l'esperienza: la semplifichi, la migliori, la renda comprensibile e la concluda,
ma soprattutto la renda \textbf{stampabile} e \textbf{fruibile} per l'utente. La
metafora crossmediale è quella della \textbf{fotografia}: chi sviluppa le
fotografie non è chi scatta la foto, ma chi la rende visibile, anche
fisicamente. ``È un po' retro come metafora, ma ci capiamo''. ``È vero, però è
giusto'', conferma Ruggero.

% ---------------------------------------------------------------------
\section{Il mercato: Italia ed estero}
\label{a1:sec:mercato}

Il corso è in italiano e si svolge in Italia, ma bisogna conoscere anche il
\textbf{mercato estero}: molti giochi progettati da autori italiani vengono
pubblicati all'estero, e conoscere il mercato estero aiuta a capire come il
proprio gioco può arrivarci, ``o viceversa, come non può arrivarci''. La
tabella~\ref{a1:tab:mercato} riassume il confronto della slide.

\begin{table}[htbp]
\centering
\small
\begin{tabular}{@{}p{6.4cm}p{6.4cm}@{}}
\toprule
\textbf{Italia} & \textbf{Estero (soprattutto Europa)} \\
\midrule
il gioco è ancora visto come un'attività \textbf{infantile}, anche se la
normalizzazione è in crescita &
forte tradizione di gioco anche in un contesto \textbf{familiare adulto},
soprattutto ``dalla Germania in su'' \\
risposta \textbf{conservativa} sul tema ludico &
\textbf{premi di prestigio} riconosciuti in tutto il mondo (Spiel des Jahres,
As d'Or) \\
forte focus sulla \textbf{localizzazione} di giochi stranieri &
si cercano soprattutto giochi \textbf{originali} \\
investimento \textbf{molto basso} sulla creazione di giochi originali interni;
pochissimi designer assunti &
\textbf{developer interni} pagati con un contratto (Francia, Germania,
Danimarca\ldots) \\
\bottomrule
\end{tabular}
\caption{Il mercato del gioco analogico in Italia e all'estero, secondo la
slide di Ruggero e i commenti di Daniele.}
\label{a1:tab:mercato}
\end{table}

\subsection{L'Italia}

\begin{itemize}
  \item In Italia il gioco è ancora visto come un'\textbf{attività infantile},
        anche se la \textbf{normalizzazione} è in crescita. C'è ancora una forte
        risposta \textbf{conservativa} sul tema ludico.
  \item C'è un forte focus sulla \textbf{localizzazione} dei giochi stranieri: in
        una qualunque ludoteca o negozio, molti dei giochi da tavolo e di ruolo
        sono di publisher o designer esteri, pubblicati in italiano da publisher
        italiani. (La localizzazione, cioè portare un titolo dall'estero in
        Italia, esiste come funzione all'interno del publisher; Ruggero non l'ha
        messa nella slide delle figure perché non è il focus del corso.)
  \item C'è, ``purtroppo, preparatevi psicologicamente'', un \textbf{investimento
        molto basso} sulla creazione di giochi originali interni alla casa
        editrice. Pochissime case editrici italiane hanno designer interni pagati
        a tempo pieno con contratto a tempo indeterminato. Daniele lo dice senza
        mezzi termini: ``il mercato è molto piccolo, conosciamo praticamente
        tutto il settore'', e le persone che in Italia fanno game design
        analogico a tempo pieno con un contratto a tempo indeterminato ``sono
        meno delle dita delle mie due mani''. Di solito la formula è quella del
        \textbf{freelance}, del contratto, delle \textbf{royalty}. Chi va da un
        publisher aspettandosi di essere assunto a tempo pieno come game
        designer deve sapere che ``non siamo purtroppo ancora in quel tipo di
        situazione''.
  \item C'è però un contraltare: in Italia, soprattutto nel gioco di ruolo ma
        anche nel gioco da tavolo, c'è stata molta \textbf{imprenditoria}, anche
        grazie a \textbf{Kickstarter}, che ha aperto le porte ad altri sistemi
        di \textbf{crowdfunding}. È un altro argomento, perché richiede un team,
        un investimento e così via.
  \item E gli autori non mancano, anzi: le occasioni di incontro tra autori
        vedono ogni anno un aumento dei possibili autori\footnote{La trascrizione
        dice ``i DAG a Parma''; il riferimento è probabilmente a IDEAG, il
        circuito di incontri tra autori di giochi italiani, ma non è
        ricostruibile con certezza.}. Bisogna ``sgomitare'' tra tanti autori,
        che magari arrivano da esperienze molto diverse.
\end{itemize}
Ruggero precisa: l'investimento basso è ``dall'alto''. Le persone che fanno e
vogliono fare giochi sono tantissime e in aumento; scarseggia l'investimento,
sia dei publisher disposti a pagare stipendi sia dei \textbf{bandi} pubblici. I
bandi esistono (Beatrice ne parlerà per i giochi pervasivi, che spesso si
muovono in quella direzione, e Ruggero cita qualcuno che sta facendo un card
game con un bando italiano, ``una cosa assurda, ma lo sta facendo''), ma sono
una rarità, soprattutto rispetto all'estero.

\subsection{L'estero}

\begin{itemize}
  \item L'estero (dipende dalla nazione, ma in generale il resto d'Europa) ha una
        forte \textbf{tradizione di gioco} anche in un contesto \textbf{familiare
        adulto}, soprattutto ``dalla Germania in su''.
  \item Ha \textbf{premi di prestigio} riconosciuti in tutto il mondo, come lo
        \textbf{Spiel des Jahres} tedesco e l'\textbf{As d'Or} francese\footnote{La
        trascrizione dice ``lo Spield AR, la SDOR''.}.
  \item Cerca soprattutto \textbf{giochi originali}. Lo si vede dal fatto che
        molti autori italiani di giochi da tavolo pubblicano in Germania o in
        Francia e \textbf{non vengono pubblicati in Italia}, ``che è tutto
        dire''.
  \item Ha \textbf{figure di development interne}: non è raro vedere publisher
        francesi, tedeschi, danesi con developer pagati internamente, a
        contratto.
\end{itemize}
``Poi la situazione è ancora più drammatica'', ammette Ruggero, ``ma per ora
vogliamo tenere il morale un po' alto. È l'introduzione. Ci sarà tempo di andare
nella tristezza''.

Prima di proseguire, Ruggero insiste sulle domande: con lui si può interrompere
anche mentre parla, alzando la mano con le reaction o semplicemente togliendo il
microfono e chiedendo. ``Non mi offenderò mai, anzi, sarò felice di vedervelo
fare e mi ricorderò positivamente di voi''.

% ---------------------------------------------------------------------
\section{Perché serve una terminologia comune}
\label{a1:sec:terminologia}

La serie di slide successiva è ``una piccolissima parte di fondamenti di
design'' (``non spaventatevi, vi vedo che vi spaventate un po'\,''). Anche qui
molti temi verranno ripresi in dettaglio, ma l'idea è arrivare alla prossima
lezione con almeno un \textbf{glossario}, un \textbf{linguaggio condiviso}: se
qualcuno dice \emph{game developer}, tutti sanno di che cosa sta parlando.

Ruggero parte da una domanda, e precisa che non è retorica: \textbf{perché ci
serve una terminologia comune?} Gli studenti rispondono:
\begin{itemize}
  \item Gianmarco, ``molto pragmatico'': perché bisogna lavorare in team e
        capirsi. Vero, nel corso e anche fuori;
  \item un'altra studentessa: per la \textbf{comunicazione}, tra di noi e con
        tutti quelli che sono interessati nel settore. Vero anche questo:
        conoscere i termini comuni permette di parlare con chi fa game design o
        development nel dettaglio, nel tecnico, nello specifico.
\end{itemize}
Ruggero aggiunge altre due ragioni.

\paragraph{Saper comunicare il proprio gioco.} La prima riguarda la
\textbf{comunicazione del gioco} verso l'esterno. ``E lo dico non tirando acqua
al mio mulino, ma anzi togliendola'': \textbf{non è spesso il buon game design
a far diffondere un gioco}. ``Vorrei che lo fosse, sarebbe bellissimo, ma
purtroppo non è così''. Ci sono giochi bellissimi che non sanno comunicare, che
non sanno dire che cosa fanno, e che quindi è difficile, da fuori, trovare
invitanti; e ci sono giochi mediocri dal punto di vista del design che si
comunicano molto bene, e che quindi conoscono in molti. Saper parlare di quello
che si fa è utile \textbf{a sé stessi}, in primis.

\paragraph{Documentare per sé stessi.} La seconda è la \textbf{documentazione}:
documenti di design, di brainstorming e così via. ``Sembrerà una battuta, ma
provate a recuperare un documento di brainstorming una settimana dopo, se non
avete scritto bene le cose che avete scritto''. Chi l'ha fatto sa già di che
cosa si parla; ``chi non l'ha ancora vissuto, forse lo vivrà''.

% ---------------------------------------------------------------------
\section{Game e play}
\label{a1:sec:game-play}

La \textbf{prima, importantissima divisione terminologica} è quella tra
\emph{game} e \emph{play}. In italiano c'è un problema: \textbf{gioco} e
\textbf{giocare} hanno la stessa radice, mentre in altre lingue, tra cui
l'inglese, sono due parole strutturalmente diverse:
\begin{description}
  \item[Game] il \textbf{gioco come artefatto}: il gioco come ecosistema di
        meccaniche, come struttura di regole, come prodotto (``tra
        virgolette''), come oggetto fisico.
  \item[Play] l'\textbf{atto}, il verbo: le persone che giocano, l'approccio
        fisico all'attività del gioco.
\end{description}
È importante perché gli studenti faranno \textbf{game design}; esistono anche
ambiti di \emph{play design}, ma qui ``creerete giochi, non giocattoli''.

\subsection{Giochi e giocattoli}

Qual è la differenza? Uno studente risponde: \textbf{la presenza di regole}.
Esatto: un gioco ha delle regole, un'attività di \emph{play} non
necessariamente. Con un giocattolo si può fare \emph{play}: c'è chi gioca con i
Lego, chi gioca a guardia e ladri o a nascondino al parco, e magari alcune di
quelle regole sono ``più o meno semi-solidificate''. Ma il \emph{play} permette
spesso un approccio non strettamente legato alle regole.

Viceversa, \textbf{un game può esistere anche se nessuno lo sta giocando}. Se si
prepara su un tavolo una scacchiera con i pezzi degli scacchi, chi passa e la
vede riconosce il gioco degli scacchi, anche se nessuno sta giocando. Se invece
si vedono dei bambini correre in un parco giochi, quei bambini \emph{stanno
giocando}, ma non stanno giocando \emph{a un gioco} specifico. ``È importante
separare questa cosa''; ci si tornerà, ``ma per ora abituatevi a questa
differenza''.

\paragraph{Approfondimento: giocattoli, giochi e regole.} La stessa distinzione
è il punto di partenza del capitolo che Jesse Schell, in \emph{The Art of Game
Design}, dedica alla definizione di gioco (\emph{The Experience Rises Out of a
Game}). Schell nota che usiamo perfino un linguaggio diverso: ``un gioco è
qualcosa a cui si gioca'' (\emph{a game is something you play}), ``un giocattolo
è qualcosa con cui si gioca'' (\emph{a toy is something you play with}); e
arriva a definire il giocattolo come ``un oggetto con cui si gioca''. Analizzando
poi la definizione di gioco di Avedon e Sutton-Smith, ne ricava che ``i giochi
hanno regole, i giocattoli no'', e che le regole sono uno degli aspetti
\emph{definitori} del gioco, insieme alla volontarietà, agli obiettivi, al
conflitto e alla possibilità di vincere o perdere. Schell osserva anche che
molti giochi sono costruiti \textbf{sopra} un giocattolo (``una palla è un
giocattolo, ma il baseball è un gioco''), e tra i suoi consigli di
prototipazione c'è proprio \emph{Build the Toy First}: costruire prima il
giocattolo, cioè qualcosa che sia già piacevole da manipolare, e poi il gioco
intorno. È una lettura utile della distinzione di Ruggero: il \emph{play} può
fare a meno delle regole, il \emph{game} no.

% ---------------------------------------------------------------------
\section{Il framework MDA (e il DPE)}
\label{a1:sec:mda}

\subsection{Meccaniche, dinamiche, estetiche}

Un'altra terminologia fondamentale è quella del framework \textbf{MDA}
(\emph{Mechanics, Dynamics, Aesthetics}; in italiano a volte ``MDE''). Viene da
un \textbf{paper} nato nell'ambito del videogioco\footnote{È il paper di Robin
Hunicke, Marc LeBlanc e Robert Zubek, \emph{MDA: A Formal Approach to Game
Design and Game Research}, del 2004. Daniele aggiunge il paper alla cartella
condivisa durante la lezione.}, che ha ``solidificato'' queste terminologie, e
divide un gioco in tre aree:
\begin{description}
  \item[Meccaniche] le \textbf{regole per come sono scritte}. ``Ogni volta che
        pescate un 2 fate dei punti'': quella è una meccanica.
  \item[Dinamiche] l'\textbf{applicazione in tempo reale} delle meccaniche. La
        dinamica è il giocatore che, pescando una carta, \textbf{spera} che esca
        un 2. Quel processo non è una regola: ``il gioco non vi dice che dovete
        sperare che esca un 2''.
  \item[Estetiche] l'insieme di quello che il gioco vuole comunicare,
        l'\textbf{esperienziale}.
\end{description}

L'esempio di Ruggero mette insieme le tre aree. In un gioco in cui pescare un 2
fa fare punti (meccanica) e in cui quindi si spera che esca un 2 (dinamica),
\textbf{che estetica si ottiene se nel mazzo di 100 carte c'è un solo 2? E se ce
ne sono 50?} L'emozione di aver pescato \emph{l'unico} 2 del mazzo è
completamente diversa da quella di pescarne uno sapendo che ce ne sono
cinquanta e che è molto probabile trovarlo. Le regole sono quasi identiche; cambia
l'esperienza.

\subsection{Il falso amico: ``estetica''}

Daniele avverte di non sottovalutare la parola \textbf{estetica}, che in italiano
è un \textbf{falso amico}: sembra indicare la parte grafica, e invece non lo è.
Si incontreranno designer che hanno studiato MDA e altri che non l'hanno
studiato, e lo si capirà proprio da questo: i secondi diranno che le
\emph{aesthetics} sono la grafica. Le \emph{aesthetics} sono la parte
\textbf{esperienziale}: come nell'esempio di Ruggero, quello che lascia al
giocatore aver pescato l'unico 2 del mazzo.

\subsection{Il DPE: design, play, experience}

Per gli italiani, aggiunge Daniele, è più facile capire l'MDA attraverso un
altro paper, quello del \textbf{DPE}\footnote{È il framework \emph{Design, Play,
Experience} di Brian Winn, proposto come estensione dell'MDA per i giochi
educativi (\emph{serious games}). Anche questo paper viene aggiunto da Daniele
alla cartella condivisa.}, che di fatto fa un parallelismo assegnando
\textbf{design, play, experience}:
\begin{itemize}
  \item le \textbf{meccaniche} sono il lato \textbf{design};
  \item le \textbf{dinamiche} sono il lato \textbf{play};
  \item le \textbf{estetiche} sono l'\textbf{experience}, l'esperienza che riceve
        chi gioca.
\end{itemize}
``Se vivete le \emph{aesthetics} come l'esperienza che vive chi sta giocando,
state a postissimo''. In questo momento Beatrice lascia la lezione (``ciao,
Bea''). I paper di MDA e DPE sono in inglese; chi vuole può dargli una lettura.

\begin{figure}[htbp]
\centering
\begin{tikzpicture}[
    font=\small,
    blocco/.style={draw, thick, rounded corners=3pt, fill=gray!12,
                   minimum width=3.2cm, minimum height=1.0cm, align=center},
    dpe/.style={font=\footnotesize\itshape, align=center},
    esempio/.style={draw, rounded corners=2pt, text width=3.3cm,
                    align=left, font=\footnotesize, inner sep=4pt},
    freccia/.style={-{Stealth[length=2.5mm]}, thick}
  ]
  \node[blocco] (m) at (0,0)   {\textbf{Meccaniche}\\ \footnotesize le regole scritte};
  \node[blocco] (d) at (4.3,0) {\textbf{Dinamiche}\\ \footnotesize le regole in azione};
  \node[blocco] (a) at (8.6,0) {\textbf{Estetiche}\\ \footnotesize l'esperienza};
  \draw[freccia] (m) -- (d);
  \draw[freccia] (d) -- (a);
  \node[dpe, above=0.15cm of m] {DPE: \textbf{design}};
  \node[dpe, above=0.15cm of d] {DPE: \textbf{play}};
  \node[dpe, above=0.15cm of a] {DPE: \textbf{experience}};
  \node[esempio, below=0.45cm of m] (em)
    {``Ogni volta che peschi un 2 fai punti''};
  \node[esempio, below=0.45cm of d] (ed)
    {pescando, il giocatore \emph{spera} che esca un 2 (nessuna regola glielo
     chiede)};
  \node[esempio, below=0.45cm of a] (ea)
    {un solo 2 su 100 carte: pescarlo è un evento;\\
     50 su 100: è quasi scontato};
  \draw[dashed] (m) -- (em);
  \draw[dashed] (d) -- (ed);
  \draw[dashed] (a) -- (ea);
  % prospettive
  \draw[freccia, gray] ([yshift=-0.3cm]em.south west) --
    node[below, font=\footnotesize\itshape, text=black]
    {il designer agisce sulle meccaniche}
    ([yshift=-0.3cm]ea.south east);
  \draw[freccia, gray] ([yshift=-1.2cm]ea.south east) --
    node[below, font=\footnotesize\itshape, text=black]
    {chi gioca incontra prima l'esperienza}
    ([yshift=-1.2cm]em.south west);
\end{tikzpicture}
\caption{Il framework MDA con l'esempio del ``2'' usato da Ruggero, e la
corrispondenza con il DPE suggerita da Daniele. Le due frecce in basso
riassumono la lettura classica del paper: il designer costruisce le meccaniche
per ottenere un'esperienza, mentre il giocatore parte dall'esperienza.}
\label{a1:fig:mda}
\end{figure}

\paragraph{Approfondimento: le due direzioni dell'MDA e le otto estetiche.} Le
frecce della figura~\ref{a1:fig:mda} riprendono un'idea centrale del paper
originale, che la lezione non esplicita: designer e giocatore guardano il gioco
da lati opposti. In \emph{Challenges for Game Designers}, Brenda Brathwaite e
Ian Schreiber la riassumono così: i designer si chiedono quale estetica
vogliono ottenere, definiscono le dinamiche che porterebbero a quella
sensazione, e poi creano le meccaniche che producono quelle dinamiche. Lo stesso
libro conferma l'avvertenza di Daniele sul falso amico: ``le \emph{aesthetics}
non si riferiscono all'aspetto del gioco, ma alla risposta emotiva'' che il
designer vuole suscitare. E offre esempi di dinamiche nel senso di Ruggero: in
uno sparatutto i punti di \emph{spawn} sono una meccanica, mentre appostarsi
accanto a uno di essi per colpire chi riappare (lo \emph{spawn camping}) è una
dinamica; negli scacchi i movimenti dei pezzi sono meccaniche, mentre le
aperture codificate sono dinamiche. Per dare un nome alle estetiche, il paper
MDA propone una tassonomia, poi detta degli ``otto tipi di \emph{fun}'', che
Schell riprende in \emph{The Art of Game Design} come \emph{LeBlanc's Taxonomy of
Game Pleasures}: \textbf{sensazione}, \textbf{fantasia}, \textbf{narrativa},
\textbf{sfida}, \textbf{compagnia} (\emph{fellowship}), \textbf{scoperta},
\textbf{espressione} e \textbf{sottomissione} (il piacere di entrare nel
``cerchio magico'' e lasciarsi andare a un passatempo). È un primo esempio di
quello che i docenti chiedono con le parole vietate della prossima sezione:
invece di dire che un gioco è ``divertente'', dire \emph{quale} esperienza
offre.

% ---------------------------------------------------------------------
\section{Le parole vietate}
\label{a1:sec:parole-vietate}

``Attenzione, adesso qua triggeriamo delle persone''. Ruggero annuncia una regola
che ``tornerà'': nel corso ci sono delle \textbf{parole vietate}. ``Non perché
siamo dei censori''; il perché si vedrà in una lezione apposita, perché
spiegarlo adesso richiederebbe un'ora. Ma Daniele vuole spiegarne ``un
minimo''.

\subsection{La regola}

In \textbf{nessun materiale scritto} consegnato ai docenti, task compresi, sono
accettate queste parole, né in inglese né in italiano:
\begin{center}
\begin{tabular}{@{}ll@{}}
\toprule
\textbf{Inglese} & \textbf{Italiano} \\
\midrule
\emph{fun} & divertimento \\
\emph{immersion} & immersione \\
\emph{entertainment} & intrattenimento \\
\bottomrule
\end{tabular}
\end{center}
``Non li voglio neanche vedere''. Daniele lo ripete con un paragone: ``se mi fate
i disegni con Paint, benissimo, top, fantastico, siamo ancora best friend. Se mi
scrivete \emph{divertimento}, io sbrocco''. È una \textbf{regola dell'accademia}:
negli stessi corsi triennali di narrative design e game design le parole sono
vietate agli studenti per tre anni; ``per voi solo per tre mesi''. Daniele
aggiungerebbe volentieri un quarto termine, \textbf{gamification}.

\subsection{Perché}

Il motivo approfondito lo darà Ruggero nella lezione apposita. Daniele dà il
suo. Ruggero e Diego Ricchiuti sono molto contrari a \emph{fun} e a
\emph{immersion}; lui ha aggiunto \emph{intrattenimento}. Sono \textbf{parole
troppo passe-partout e troppo ambigue}:
\begin{itemize}
  \item \textbf{intrattenimento} in Italia è stata usata per anni come un modo
        per \textbf{sminuire}: ``ma sì, tanto è un gioco solo di
        intrattenimento, è un romanzo solo di intrattenimento''. Ma anche se un
        prodotto è ``solo di intrattenimento'', non vuol dire che non ci si debba
        lavorare con attenzione. E poi intrattenimento ``non significa niente'':
        ci sono tanti modi in cui una cosa può intrattenere;
  \item lo stesso vale per \textbf{fun}, e per \textbf{immersion}, che ha un
        problema in più: è una parola passe-partout per cui si può dire che un
        gioco è ``immersivo'' nel senso generale del termine, ma poi non si sa
        davvero che cosa voglia dire, \emph{immersivo in che senso}.
\end{itemize}

L'obiettivo didattico è che gli studenti \textbf{vadano oltre queste parole}:
``un po' come quando alle medie vi dicono smettetela di dire che tutto è bello,
perché ci sono vari gradi di bello, ci sono altri modi di dire bello''. Bisogna
\textbf{decostruirle}: quando si sta per usarle, chiedersi \textbf{che cosa si
intende} con quella parola e cercarne una \textbf{più precisa}.

Ruggero chiede, in uno dei ``rari casi'', di \textbf{non fare domande adesso} su
questa regola: ogni anno arrivano domande e commenti, e ci saranno lezioni
apposite di tre ore, con strumenti per capire il perché. Ma da subito gli
studenti devono abituarsi a evitare questi termini, sia perché si spera che
abbiano termini più precisi, sia perché queste parole ``portano con sé dei
problemi agganciati'' che si vedranno. ``Non spaventatevi, non verrete sparati in
un vicolo cieco di notte''.

\paragraph{Una piccola gamification.} Il divieto va visto come uno
\textbf{stimolo}, ``un po' come se stessimo gamificando la vostra esperienza'',
una sorta di user experience delle lezioni. Imponendosi di non usare queste
parole, si noterà come migliora la capacità di \textbf{descrivere i giochi} e le
proprie esperienze ludiche, e di essere più precisi. È un invito a ``obbligarsi a
essere più precisi e ad avere una mentalità più critica''. La lezione di
Ruggero sull'\textbf{involvement}, promette Daniele, sarà ``fenomenale'' e
``vi aprirà la testa''.

\paragraph{Approfondimento: anche i libri faticano a definire il ``fun''.} La
difficoltà che il divieto vuole far emergere è visibile anche nei testi di game
design. Schell, in \emph{The Art of Game Design}, prova a definire il \emph{fun}
e propone ``piacere con sorprese'' (\emph{pleasure with surprises}), ammettendo
subito che è strano poter usare una parola per tutta la vita, sapendo con
certezza che cosa significa, senza riuscire a esprimerla chiaramente quando
qualcuno lo chiede; e invita a mettere alla prova la definizione con
controesempi. La sua lente sul \emph{fun} (\emph{The Lens of Fun}) si limita,
non a caso, a chiedere quali parti del gioco siano \emph{fun} e \emph{perché}.
Anche \emph{immersion} ha molte letture: in \emph{21st Century Game Design},
Bateman e Boon la descrivono, tra l'altro, come la capacità di ``perdersi''
nel mondo artificiale presentato dal gioco, legata alla sospensione
dell'incredulità e ai giochi di mimesi. Sono significati diversi sotto la stessa
parola: proprio il motivo per cui, in un documento di design, conviene dire
quale si intende.

% ---------------------------------------------------------------------
\section{Chi crea, chi gioca, chi facilita}
\label{a1:sec:ruoli-esperienza}

Un'altra distinzione importante è quella tra \textbf{chi crea}, \textbf{chi
gioca} e \textbf{chi facilita}. Spesso, quando si fa un gioco, lo si fa provare
di persona a chi gioca. Ma ``la norma dei giochi pubblicati non funziona così'':
quando si pubblica un gioco non si vede chi lo compra e chi lo gioca, e non si
sa se qualcuno, a casa propria con i propri amici, segue le regole alla lettera e
interpreta correttamente quello che il designer aveva descritto. Quindi già in
fase di progettazione bisogna tenere a mente \textbf{chi sta facendo game design
e chi sta giocando}.

E nei giochi a evento, nei giochi pervasivi, nei giochi di ruolo da tavolo e dal
vivo (LARP) c'è un terzo pezzo del puzzle, ``importantissimo'': \textbf{chi
facilita}. La domanda da farsi è: \textbf{il regolamento è scritto per chi gioca
o per chi deve facilitare il gioco?} Bisogna abituarsi a pensare a tutte le
persone coinvolte nell'esperienza.

% ---------------------------------------------------------------------
\section{Playtest e iterazione}
\label{a1:sec:playtest}

\subsection{Nessun gioco nasce perfetto}

L'ultima voce del glossario è ``un'altra cosa importantissima'':
\textbf{playtest} e \textbf{iterazione}. Il corso è \textbf{pragmatico}:
``specialmente da parte mia'', dice Ruggero, non si vedranno teorie ``messe lì
per fare accademia sul game design nelle nostre torri d'avorio''. Ci saranno
teorie e framework, ma tutti estremamente pragmatici, perché l'obiettivo non è
solo capire i giochi, ma \textbf{saperli fare e iterare}.

Playtest e iterazione sono \textbf{inevitabili} nel game design, perché
\textbf{nessun gioco nasce perfetto}, soprattutto perché le persone sono
diverse: un gioco perfetto per una persona può essere incomprensibile per
un'altra. Chi progetta giochi per altre persone deve capire, dopo averli fatti
provare, \textbf{che cosa cambiare, che cosa sistemare, che cosa aggiungere, che
cosa togliere}. Ci saranno lezioni su playtest, prototipazione, raccolta dei
feedback, metodi etnografici. Ma bisogna abituarsi già ora all'idea che
\textbf{non esiste la consegna definitiva}. A un certo punto i task vanno
consegnati per forza di cose, ma, come diceva Daniele all'inizio, alla fine del
corso gli studenti avranno un \textbf{prototipo}, ``che sarà un prototipo'', e
gli strumenti per portare avanti il lavoro anche dopo.

\subsection{L'aiuto dopo il corso}

Non solo gli strumenti: anche l'aiuto dei docenti. Daniele racconta di aver
portato a un editore, a Play\footnote{Probabilmente \emph{Play -- Festival del
Gioco} di Modena; la trascrizione non permette di esserne certi.}, uno dei
progetti realizzati l'anno precedente da due studenti, un gioco di ruolo in
solitario (``un \emph{solo game} RPG molto, molto bello''). Se gli studenti
vorranno trasformare il proprio lavoro in un prototipo vero e proprio, possono
contattarlo: lo fa ``per lavoro, ma anche per passione''. Gestisce anche una
community per aiutare i designer ``un po' depressi'' a rimettersi a progettare.

Quindi: ``non sentitevi persi se durante la fase di iterazione vi sentite
persi, è normale''.

\subsection{Il dizionario comune e il production dictionary}

Daniele riprende un punto detto ``benissimo'' da Ruggero e che secondo lui non è
stato sottolineato abbastanza: il \textbf{dizionario comune} serve anche a
\textbf{non perdersi} mentre si itera. Se da un'iterazione all'altra i
documenti usano termini completamente diversi, e non c'è un dizionario di base,
ci si perde.

Nel videogioco, una delle primissime cose che un designer fa quando entra in
un'azienda è il \textbf{production dictionary}: si annota, in un file privato,
tutte le parole che usano i colleghi, per allinearsi al dizionario comune
dell'azienda, perché lì su un prototipo possono lavorare anche cento persone.
Nel gioco da tavolo non si fa molto, perché il mondo dell'analogico è
``talmente più piccolo''; ma serve comunque imparare un \textbf{dizionario
minimo}, dai docenti e dai professionisti. E il corso ne darà una parte,
entrando nel dettaglio di alcuni termini e framework.

\paragraph{L'esempio del feedback.} Perché serva, lo mostra un esempio. Si
testa un gioco e ci si annota come feedback: ``questo gioco deve essere più
divertente''. Che cosa vuol dire? Come lo si migliora? Quella frase non dà
elementi utili. Se invece si scrive che deve essere \textbf{più adrenalinico},
\textbf{più frenetico}, \textbf{più rilassato}, \textbf{più cozy}, o
\textbf{più bilanciato matematicamente}, si sa dove andare a lavorare.

\paragraph{Approfondimento: la regola del ciclo.} Il principio ``nessun gioco
nasce perfetto'' ha un nome preciso in \emph{The Art of Game Design}: è la
\textbf{Rule of the Loop}, la regola del ciclo. ``Più volte provi e migliori il
tuo design, migliore sarà il tuo gioco''. Schell precisa che non è una lente,
cioè un punto di vista, ma ``una verità assoluta'', senza eccezioni: ci saranno
momenti in cui si cercherà di convincersi che ``questa volta il design è così
buono che non serve testarlo'', e ogni volta se ne pagherà il prezzo. Ed è
interessante la ragione per cui la regola pesa di più sul digitale: secondo
Schell il tempo e il denaro necessari per provare e correggere un sistema sono
molto maggiori nei videogiochi che nei giochi tradizionali. Per un gioco di
carte o da tavolo, con il tempo per provarlo e cambiarlo più e più volte, il
consiglio è semplicemente di farlo. Il gioco analogico è, da questo punto di
vista, la palestra ideale dell'iterazione.

% ---------------------------------------------------------------------
\section{Tabletop Simulator}
\label{a1:sec:tts}

\subsection{Il software di settore}

``Last but not least'': il corso usa un software, \textbf{Tabletop Simulator},
che si compra su \textbf{Steam} e costa tra i \textbf{10 e i 20 euro} (``aspettate
che vada in sconto, va in sconto ogni due per tre''). Daniele spiega perché:
\begin{itemize}
  \item è \textbf{il software di settore} del gioco da tavolo, quello che i
        publisher usano per i prodotti digitali, interni ed esterni;
  \item essendo su Steam, permette di fare \textbf{playtest} e di giocare con gli
        amici di Steam. Daniele stesso sta testando lì il proprio gioco;
  \item funziona anche su Mac: Windows o Mac, nessun problema.
\end{itemize}
Ci sarà una lezione apposita, di Ruggero, sulla creazione del proprio
\textbf{modulo} di gioco.

\paragraph{Che cosa non è.} Daniele chiarisce subito un ``fraintendimento
gigantesco'', soprattutto per chi ha già esperienza: Tabletop Simulator
\textbf{non richiede programmazione} e \textbf{non è un software per
trasformare i giochi da tavolo in videogiochi}. È un software che
\textbf{simula dei componenti fisici su un tavolo}: serve per giocare i board
game in digitale. Simula anche la fisica: c'è perfino il comando per
\textbf{ribaltare il tavolo} (``non usatelo ora'').

\paragraph{La partnership mancata.} Da anni Daniele prova a fare una
partnership con gli sviluppatori: due anni prima ha contattato l'azienda per
acquistare delle chiavi per gli studenti, dato che il software si usa nei corsi.
L'azienda ha risposto che non fa accordi commerciali. ``Credo sia la prima
azienda nella mia vita che dice che non vuole che io le dia dei soldi''. Gli
studenti dovranno quindi provare ad acquistarlo da sé; se i docenti riusciranno a
darglielo, bene.

\paragraph{La domanda di Gianmarco.} Gianmarco segue le lezioni da un computer
aziendale ``super chiuso'', su cui non può installare Steam: deve usare Tabletop
Simulator durante le lezioni? Ruggero risponde che per la lezione dedicata è
meglio poter seguire e rifare in diretta. Ma il software non sarà richiesto
necessariamente per il project work: un gruppo potrebbe fare un gioco pervasivo
o un gioco di ruolo, non un gioco da tavolo. Essendo un software di settore, è
comunque bene conoscerlo. Ruggero aggiunge che le versioni non ufficiali che
circolano non permettono di giocare in molti, e quindi non servono per
playtestare o giocare insieme.

\subsection{Una dimostrazione}

Daniele mostra il software aperto da Steam:
\begin{itemize}
  \item si può creare una partita \textbf{single player} o \textbf{multiplayer};
        per giocare tra compagni si crea un server multiplayer e ci si aggiunge
        come amici su Steam;
  \item si dà al server un nome, un numero di giocatori, eventualmente una
        password; creato il server, si apre la scelta dei giochi;
  \item ci sono giochi gratuiti, giochi da acquistare (DLC, non obbligatori) e
        soprattutto un \textbf{Workshop} ``ampissimo'', con decine di migliaia di
        giochi caricati dagli utenti: ``se c'è un board game che conoscete,
        quel board game è su Tabletop Simulator, giocabile dal Workshop''.
        L'esempio mostrato è l'open beta ufficiale del gioco da tavolo di
        \emph{Slay the Spire}\footnote{La trascrizione dice ``Slade Spire''.}:
        molti publisher, per farsi pubblicità, mettono il proprio gioco in demo o
        giocabile in digitale;
  \item c'è un \textbf{tutorial} iniziale che spiega come muoversi sul tavolo;
        con \emph{Esc} e \emph{Help} si trovano i comandi da tastiera.
\end{itemize}

\subsection{Il primo task}

Il primo (e per ora unico) task è questo: \textbf{entro il 10 ottobre}, giorno
della lezione di Ruggero su Tabletop Simulator, giocare \textbf{tre o quattro
volte in totale}, online, tra compagni di corso o con amici. Non si tratta di
imparare a creare moduli, ma di prendere la mano \textbf{lato giocatore}: aprire
il software, impararne i comandi, giocare. ``Per me potete anche andare a
prendere domino o mahjong'' (``mahjong è bellissimo tra l'altro'', aggiunge
Daniele, che sta facendo un videogioco sul mahjong), o il pachisi o il go. Gli
anni precedenti si chiedeva a ciascuno di proporre ai compagni un gioco da
tavolo che apprezzava. È anche un modo per conoscersi. Ruggero spiega il
perché: usarlo per giocare è facile e intuitivo, usarlo per creare moduli è più
complesso; ma chi arriva al 10 ottobre avendo già aperto il software e giocato,
troverà più intuitivo anche creare un modulo. E l'avvertimento: ``non riducetevi
all'ultimo, non fate una partita il 9 sera''.

\subsection{Tabletopia e altri vantaggi}

Tancredi racconta di aver provato a caricare un suo gioco su
\textbf{Tabletopia}. ``Questo è un vaso di Pandora'', risponde Daniele.
Tabletopia è molto usato, ma secondo lui il suo strumento di creazione dei moduli
è pessimo, ed è il motivo per cui la maggioranza dei moduli ufficiali dei
publisher è su Tabletop Simulator. Tabletopia ha un vantaggio vero, è
\textbf{gratuito}, ma ha scelte di ``qualità della vita'' poco felici.

Tabletop Simulator ha invece un vantaggio ``fenomenale'' che si tende a
sottovalutare: essendo su Steam, \textbf{il modulo si pubblica direttamente sul
Workshop di Steam}; si possono invitare gli amici a iscriversi e giocare con
loro, purché abbiano un account Steam. Rispetto ad altri software da imparare,
costa ``tra i 10 e i 20 euro, ragazzi: si parla proprio di spiccioli''. ``Io con
Tabletop Simulator ci ho playtestato la vita''.

Altre caratteristiche utili:
\begin{itemize}
  \item si possono creare \textbf{moduli non elencati} (\emph{unlisted}),
        perfetti da passare a un publisher prima che il gioco sia pubblicato:
        non sono pubblici, e nessuno può rubare i materiali;
  \item il software ha una grafica datata (``\emph{World Art 1998}'', scherza
        uno dei docenti), ma lo stanno portando su \textbf{Unity} con molti investimenti
        e aggiornamenti; ascolta la comunità e introduce funzioni utili anche per
        chi fa design;
  \item il suo server Discord è molto frequentato, e c'è un canale YouTube
        dedicato solo al design di giochi di carte\footnote{La trascrizione ne
        riporta il nome come ``Distraction Maker''; non è verificabile.} con un
        server Discord dove si fanno playtest: il 99\% è su Tabletop Simulator;
  \item molti editori, quando si propone loro un gioco da tavolo,
        \textbf{chiedono se c'è un modulo per Tabletop Simulator}.
\end{itemize}

% ---------------------------------------------------------------------
\section{Le domande a Daniele: il tema del project work}
\label{a1:sec:domande-daniele}

Daniele ha ``un minuto'' per rispondere alle domande sul corso prima di dover
andare. Uno studente chiede, ``giusto per curiosità'', se il project work avrà
un tema o sarà libero.

\begin{itemize}
  \item C'è un \textbf{mandate}, che i docenti propongono ogni anno e
        aggiornano leggermente; verrà presentato il \textbf{25 settembre}.
        ``Non facciamo spoiler''. Conterrà quello che si chiede di fare e quello
        che non si chiede: dei \textbf{must}, dei \textbf{must not} e dei
        \textbf{nice to have}.
  \item Verrà chiesto di lavorare su un'\textbf{idea originale} che non si è
        portata al corso: originale nel senso di \textbf{creata apposta}, non
        nel senso di ``mai vista prima''.
  \item L'idea va \textbf{partorita insieme}. Altri anni alcuni studenti hanno
        provato a imporre la propria idea al gruppo, ``e l'esperienza è stata
        pessima per tutti''. Quindi: \textbf{non imponete le vostre idee, fate
        brainstorming con gli altri}. Daniele, che fa ``molto spesso
        brainstorming'', darà strumenti per farlo: nessuno arriverà impreparato
        a quella lezione.
\end{itemize}

Ruggero aggiunge che il mandate è \textbf{abbastanza flessibile} perché tutti
possano fare qualcosa che li interessi: ``l'ultima cosa che vogliamo è forzarvi a
lavorare su qualcosa che non vi piace''. Bisogna lavorare su ciò che piace e per
cui si ha uno stimolo personale, anche nella \textbf{tipologia} di gioco: nessuno
sarà obbligato a fare un gioco di ruolo o un board game, e la lezione sui casi
studio mostrerà quante forme diverse di gioco analogico il corso accetta.

L'unica cosa da \textbf{non fare}, ``perché è contraddittorio per questo corso'',
sono i \textbf{giochi grandi}: ``lo chiediamo per voi''. Niente progetti
mastodontici\footnote{La trascrizione dice ``non fate delle cose da gnome'':
probabilmente il nome di un gioco da tavolo molto grande e complesso, non
ricostruibile con certezza.}, niente ``cose alla \emph{Dungeons \& Dragons}''.
Si può fare un gioco di ruolo, ma con un'esperienza ``meno corposa'': ci sono
tanti giochi di ruolo da 15 o 20 pagine, ``bellissimi, fichissimi'', che si
possono leggere online. Il consiglio è farsi una cultura anche su
\textbf{itch.io}, ``uno dei portali più importanti del gioco di ruolo
internazionale'', di cui si parlerà anche nella parte di ricerca e sviluppo.

Poi Daniele saluta (``noi ci vediamo venerdì a lezione'') e lascia la
conduzione a Ruggero.

% ---------------------------------------------------------------------
\section{Le domande finali con Ruggero}
\label{a1:sec:domande}

\subsection{Le regole delle domande}

Ruggero spiega che cosa vuol dire ``stare con lui''. Alla fine di ogni sua
lezione ci saranno \textbf{almeno 20--30 minuti di domande} libere, \emph{free
form}. Vuole che fare domande sia \textbf{normalizzato}: ``io gestisco malissimo
l'\emph{horror vacui} e divento super cringe molto alla svelta se non ci sono
domande in giro''. Sono benvenute le domande sulla lezione appena fatta, ma non
solo: il corso è crossmediale, e anche le domande che vengono da
\textbf{divagazioni}, da parti laterali, ``fatele''. Quello che sembra scontato o
banale a uno studente può non esserlo per il docente, e viceversa; e può essere
utile anche agli altri. ``Se vi sembrano stupide, fatele il doppio''. ``Magari
sembro incazzato, ma vi giuro sono rilassatissimo. Non mordo, magari abbaio un
po'\,''.

\subsection{Come approcciarsi a Steam e Tabletop Simulator}

Rita non ha mai usato né Steam né Tabletop Simulator e chiede come
approcciarsi, e se c'è modo di mettersi d'accordo tra compagni su Discord, magari
con un sondaggio sui giochi, per formare gruppi di gioco. Ruggero conferma: sul
Discord dell'accademia c'è un \textbf{canale apposito per il corso}, ed è degli
studenti. Lo possono usare per organizzarsi a giocare, per chiedere a un
compagno di prestare un gioco, per domande tecniche (come si scarica, come si
installa). ``Sarebbe proprio ideale che vi parlaste e vi organizzaste''.

Anche perché, aggiunge con un ``piccolo spoiler'' (``nessuno vi ha detto
niente''), se fossero i docenti a organizzare i gruppi di gioco, il task lo
starebbero facendo loro. Bisogna \textbf{abituarsi a organizzare le sessioni di
gioco}, perché arriveranno i task sul \textbf{playtest} e su come organizzarlo,
e a quel punto si sarà ``meno disperati''. Sembra scontato, ma saper
\textbf{organizzarsi per giocare in autonomia} ``vi sarà utilissimo come skill in
questo settore, ve lo giuro''.

\subsection{Il bilanciamento}

Angelo confessa ``una paranoia grossissima'': ci saranno lezioni sul
\textbf{bilanciamento}? ``Ma certo''. Sono proprio le lezioni di Marta sul
development: il bilanciamento è di solito un tema che tocca \textbf{i developer
più dei designer}. L'anno precedente c'era una sola lezione con Marta, e i
docenti si sono resi conto che serviva andare ancora più a fondo: per questo
quest'anno le lezioni sono due. ``Non è una paranoia, è una legittima domanda'':
Ruggero insegna da anni, vede altri corsi, parla con altri docenti e studenti, e
``mi incazzo spesso'' perché molti corsi di game design \textbf{non fanno
bilanciamento}. I docenti volevano invece un corso che coprisse tutte le parti
del game design: non solo l'idea del gioco, ma il bilanciamento, la produzione,
i playtest, l'iterazione, la scrittura dei regolamenti.

\subsection{Si può giocare con i docenti?}

Una studentessa chiede se si può giocare con i docenti. Per motivi etici,
durante il corso questo è limitato alle parti di \textbf{valutazione} dei giochi
e di \textbf{tutoraggio}. Dopo il corso, ``nessun problema, almeno con me''.

\subsection{Dove si trovano i giochi su Tabletop Simulator}

Ruggero mostra dove si trovano i giochi: nella propria libreria di Steam, nella
pagina di Tabletop Simulator, alla voce \textbf{Workshop}, dove ci sono ``più o
meno 36.000'' giochi. Una studentessa scopre intanto che l'app di Steam per
telefono è solo una \emph{companion app}: per registrarsi e giocare serve il
PC.

Per provare, chiede agli studenti un titolo:
\begin{itemize}
  \item \emph{Keep the Heroes Out}: non compare, ``non tutti i giochi ci sono,
        perché molti publisher non li mettono''; forse, essendoci già una
        versione del gioco su Steam, il modulo è stato tolto;
  \item \emph{Magic: The Gathering}: ``almeno 50 moduli'', anzi ``123 pagine'';
  \item \emph{Bang!}: tre moduli.
\end{itemize}
Bisogna fare attenzione: alcuni moduli sono in lingue specifiche (``magari il
primo è in russo''), alcuni hanno già l'espansione, altri sono solo il gioco
base. Conviene sfogliare e non fermarsi al primo risultato. E nel Workshop si
trovano anche cose molto particolari: giochi fuori produzione, giochi vecchi,
giochi ancora in prototipo.

\subsection{I ruoli nei team e le dinamiche tra le persone}

Uno studente torna sulla slide delle figure: nelle realtà piccole molti ruoli
hanno ``il doppio o triplo cappello''. Quali ruoli restano distinti, e come si
creano le sinergie tra persone con obiettivi diversi?

Ruggero premette che l'ideale ``platonico'' è un ambiente di lavoro costruttivo,
basato sulla fiducia e ``il più \emph{safe} possibile''; ma nell'industria non
sempre è così, e le strutture variano moltissimo: alcuni publisher sono molto
verticali, altri hanno ruoli già solidificati (se un publisher ha già i suoi
developer e non sta crescendo, entrarci come developer è praticamente
impossibile). ``Se prendiamo due aziende in Italia, vediamo due strutture
interne completamente diverse''. Non esiste una scorciatoia per farsi assumere
da un editore.

Poi prende ``l'esempio più comune di tutti'', la media dei giochi da tavolo
pubblicati (figura~\ref{a1:fig:relazioni}):
\begin{itemize}
  \item il \textbf{game designer} è una figura \textbf{esterna}, con un
        \textbf{contratto a royalty} con il publisher;
  \item la sua figura di riferimento è il \textbf{publisher} stesso, e la
        conversazione designer--publisher è ``molto calda'';
  \item se il publisher ha dei \textbf{developer}, il developer entra in questa
        conversazione: designer e developer collaborano molto, e internamente
        developer e publisher collaborano molto. Spesso, nei publisher piccoli,
        è il publisher a fare le veci del developer;
  \item \textbf{graphic design} e \textbf{illustrazione} tendono a parlare molto
        con il publisher e \textbf{molto poco con il designer}. A meno che non si
        tratti di uno studio indie molto piccolo, è difficile che il designer
        sappia anche solo chi è l'illustratore prima che il gioco esca.
\end{itemize}

\begin{figure}[htbp]
\centering
\begin{tikzpicture}[
    font=\small,
    ruolo/.style={draw, thick, rounded corners=3pt, fill=gray!12,
                  minimum width=2.8cm, minimum height=0.9cm, align=center},
    esterno/.style={draw, thick, dashed, rounded corners=3pt,
                    minimum width=2.8cm, minimum height=0.9cm, align=center},
    calda/.style={{Stealth[length=2.5mm]}-{Stealth[length=2.5mm]}, very thick},
    tiepida/.style={{Stealth[length=2mm]}-{Stealth[length=2mm]}, thin, dashed}
  ]
  \node[esterno] (gd) at (0,0) {Game designer\\ \footnotesize esterno, a royalty};
  \node[ruolo] (dev) at (5,1.4) {Developer\\ \footnotesize se c'è};
  \node[ruolo] (pub) at (5,-1.4) {Publisher};
  \node[ruolo] (art) at (10,-1.4) {Graphic designer\\ e illustratore};
  \draw[calda] (gd) -- (dev);
  \draw[calda] (gd) -- (pub);
  \draw[calda] (dev) -- (pub);
  \draw[calda] (pub) -- (art);
  \draw[tiepida] (gd.south) to[bend right=25]
     node[below, font=\footnotesize\itshape] {quasi nessun contatto} (art.south);
  \node[draw, dotted, rounded corners=4pt, fit=(dev)(pub)(art), inner sep=8pt,
        label={[font=\footnotesize\itshape]above right:interno al publisher (di solito)}] {};
\end{tikzpicture}
\caption{Le relazioni tipiche intorno a un gioco da tavolo pubblicato, secondo la
risposta di Ruggero. Le frecce spesse indicano una collaborazione stretta e
frequente; quella tratteggiata, un contatto quasi assente. Nei publisher piccoli
il publisher fa spesso anche da developer.}
\label{a1:fig:relazioni}
\end{figure}

\subsection{Fare da developer di sé stessi?}

Rita chiede in chat se un designer può fare anche da developer del proprio
gioco, o se Ruggero lo sconsiglia. In generale \textbf{lo sconsiglia}, per lo
stesso motivo per cui si sconsiglia in altri media: i \textbf{darling}. Il
designer ha delle meccaniche a cui è affezionato, e il developer a volte dovrà
\textbf{tagliarle}; è molto difficile, se si è ``gestori di sé stessi'', tagliare
cose che si vogliono nel gioco. ``È un po' come lasciare fare il lavoro di editing
allo scrittore che ha già scritto il libro'': può rivedersi, certo, ma difficilmente
cambierà cose che farebbe bene cambiare. Nelle lezioni con Marta si vedrà perché,
e alla fine del corso ``non avrete dubbi nemmeno voi'' nel separare designer e
developer.

Ma tutto dipende dallo \textbf{scopo del progetto}: se si fa un gioco in tre,
qualcuno quel lavoro deve farlo per forza, e magari il designer fa anche il
developer, compensando con \textbf{più playtest}. ``Il mondo attuale e il mondo
ideale sono due cose molto diverse, e lo vediamo non solo dai giochi ma dalla
politica'', ma ``non entriamo in questo vaso di Pandora''.

\subsection{Il ruolo indispensabile e la pubblicazione}

Lo studente di prima ne trae una conclusione: senza game designer il gioco non
esiste. ``Quello sicuramente''. Ma Ruggero coglie l'occasione per approfondire
un punto che forse non era passato: \textbf{fare design di un gioco per
pubblicarlo non è l'unica opzione possibile}. Un gioco non è ``meno gioco''
se non è pubblicato da un editore. È legittimo mettere un gioco su itch.io in
\textbf{print and play}, fare un gioco a evento senza venderlo, fare un
\emph{murder party}, una escape room, un board game per una persona amica:
``quell'esperienza lì non è meno gioco''. Dipende dall'obiettivo del progetto.

Ruggero ha lavorato molto anche con le \textbf{scuole}, dove il publisher non
esiste: è un gioco per le classi, e la figura più vicina è l'editoria scolastica,
``che però non fa il publisher''. Ci sono compromessi diversi sulla materialità
e su come il gioco si presenta, che si vedranno nella prototipazione e nei casi
studio. Quindi l'unica figura che deve esserci sempre è \textbf{chi fa game
design}, anche se magari non sa di chiamarsi game designer e non ha mai
sentito parlare di game design: ``questo non la rende meno game designer''.

\paragraph{Un libro consigliato.} Ruggero fa anche mentoring e gestisce gruppi
di settore, e spesso gli chiedono che cosa leggere per imparare il game design.
Non è il tipo che obbliga a leggere libri (``non vi obbligherò a leggere
niente''), ma ce n'è uno che consiglia sempre, ed è il primo che ha scelto di
leggere: \textbf{\emph{Rise of the Videogame Zinesters}} di \textbf{Anna
Anthropy}\footnote{La trascrizione dice ``Anna Antrofi'' e ``Rise of the
Videogames in Esters''.}. È un libro un po' datato, ma ancora reperibile; parla
soprattutto di videogiochi, ma ``per l'80\%'' il contenuto si applica al gioco e
al game design in generale. È un libro sulla \textbf{democratizzazione del gioco}:
il gioco come strumento per parlare di sé e come \textbf{medium legittimo}. Come
esistono film drammatici, tristi, intensi, possono esistere giochi su tematiche
delicate, e se ne vedranno nei casi studio. Molti giochi fatti con questa
direttiva non hanno un publisher, che tende a pubblicare ciò che dà un ritorno
commerciale.

\paragraph{Commerciale e personale.} ``Io spero che voi non farete giochi solo
per scopo commerciale''. Certo, servono un tetto sulla testa, il cibo in
tavola, un posto in cui dormire, e così si finisce per fare giochi che non sono
il proprio gioco dei sogni. Spesso è un mix: Ruggero ha fatto giochi commerciali
``per pagarmi l'affitto'' e giochi che voleva fare e su cui ha perso soldi,
``sapendolo proprio che sarebbero stati un fallimento''. I giochi, come i libri e
le canzoni, si fanno per molti motivi, e quelli commerciali di successo non sono
più importanti degli altri. Tra i giochi della sua slide c'erano il gioco da
tavolo ufficiale di \emph{Cyberpunk}\footnote{La trascrizione dice ``Cyberpunk
del 1977'': quasi certamente un gioco legato a \emph{Cyberpunk 2077}.} e
\emph{Marvel United}, su IP grosse e famose, ma anche giochi fatti di ``una
candela e una busta su un tavolo''. E se gli si chiede quale preferisce in
termini di contenuto, non sono quelli sui supereroi.

\subsection{Il gioco di ruolo in Italia}

Nicola, che è interessato soprattutto al gioco di ruolo e conosce già persone
del settore, ha ``una leggera paura'': che chi studia game design finisca per
fare il \textbf{freelance} a vita, ``perché non ci sarà mai un'occupazione per
bene'', o il \textbf{copywriter}. Conosce un ragazzo che ha lavorato in quel
mondo e che ora fa altro.

Ruggero risponde che la maggioranza dei giochi nel suo portfolio sono giochi di
ruolo, quindi il tema lo tocca da vicino, e che in Italia quello che dice Nicola
``è semplicemente vero''. Molti publisher piccoli hanno una persona di
riferimento che copre più ruoli: impagina, fa i playtest, fa un po' di
development. È ``un settore economicamente completamente disastrato'', per
motivi culturali che non apre: oggi in Italia è praticamente impossibile
\textbf{vivere di giochi di ruolo}. Chi ci lavora lo fa nel tempo libero, al di
fuori del lavoro principale, oppure coprendo molte altre cose: molti publisher
italiani di giochi di ruolo pubblicano anche titoli esteri e localizzano, non
fanno solo cose originali. ``È un grosso calderone''. Poi si ferma: ``questo è
un \emph{rant}'', e ``i \emph{rant} tornano sempre'', ma non vuole che gli
studenti escano dalla prima lezione ``con la disperazione nel cuore''. Ci sarà
spazio per parlarne quando avranno gli strumenti per capirlo, ``per non farvi
abbattere''.

\subsection{Che cosa bisogna sapere già}

Uno studente chiede se ci sono nozioni che dovrebbero già avere e che possono
recuperare: per esempio, sente parlare di giochi ``German'' o ``American'' e non
sa che cosa siano. Ruggero spiega l'organizzazione del corso:
\begin{itemize}
  \item le \textbf{prime due lezioni} sono di fondamenta. Questa era ``dall'alto'':
        come è strutturato il corso, le tempistiche;
  \item la \textbf{prossima lezione} è di tre ore di soli \textbf{casi studio},
        ``a manetta, a raffica'', divisi per categorie e tipologie, con
        ragionamenti su categorie e generi di gioco: lì si vedranno German,
        American ``e anche altre più utili'';
  \item nel resto del corso le terminologie specifiche sono inserite nelle
        lezioni dedicate: il \textbf{paragame} nelle lezioni sul paragame, il
        \textbf{modello matematico} nelle lezioni di Marta sul development.
\end{itemize}
Per ora agli studenti si chiede di sapere quello che si è visto oggi:
meccaniche, dinamiche ed estetiche; le figure professionali, designer e
developer; e soprattutto \textbf{game e play}.

\subsection{Game e play: i quattro quadranti}

``Specialmente \emph{game} e \emph{play}'', insiste Ruggero: è più difficile di
quello che sembra, perché l'italiano non abitua a separare le due cose, ed è
possibile averle usate per anni come sinonimi. Nella sua esperienza con chi si
avvicina al design è ``uno dei punti più critici''. Spesso i commenti che si
fanno sui giochi giocati sono commenti sul \emph{play} applicati al
\emph{game}: ``il gioco non mi è piaciuto perché, giocando con Tizio, è successo
questo''; ``il gioco mi sembra sbilanciato perché Tizio al tavolo ha fatto
questo''.

Separare \emph{game} e \emph{play} serve a capire quali giochi sono progettati
bene e quali male, e soprattutto a riconoscere i \textbf{quattro possibili
incroci} (figura~\ref{a1:fig:quadranti}): un gioco progettato bene che piace
giocare; uno progettato male che non piace; uno progettato male che piace; uno
progettato bene che non piace. ``Non è importantissimo in generale nella vita,
ci sono cose più importanti'', ma per chi fa game design lo è: ``a volte vi
chiederanno di lavorare su un gioco che non vi piace giocare, e quel gioco deve
essere fatto bene lo stesso''.

\begin{figure}[htbp]
\centering
\begin{tikzpicture}[
    font=\small,
    cella/.style={draw, thick, minimum width=4.4cm, minimum height=1.9cm,
                  text width=4.0cm, align=center}
  ]
  \node[cella, fill=gray!20] (a) at (0,0)     {\textbf{progettato bene}\\ e \textbf{mi piace} giocarlo};
  \node[cella]               (b) at (4.4,0)   {\textbf{progettato bene}\\ ma \textbf{non mi piace} giocarlo};
  \node[cella]               (c) at (0,-1.9)  {\textbf{progettato male}\\ ma \textbf{mi piace} giocarlo};
  \node[cella, fill=gray!8]  (d) at (4.4,-1.9){\textbf{progettato male}\\ e \textbf{non mi piace} giocarlo};
  \node[font=\footnotesize\itshape, above=0.15cm of a] {mi piace (\emph{play})};
  \node[font=\footnotesize\itshape, above=0.15cm of b] {non mi piace (\emph{play})};
  \node[font=\footnotesize\itshape, rotate=90, left=0.15cm of a, anchor=south]
       {bene (\emph{game})};
  \node[font=\footnotesize\itshape, rotate=90, left=0.15cm of c, anchor=south]
       {male (\emph{game})};
\end{tikzpicture}
\caption{I quattro incroci tra la qualità del \emph{game} (come è progettato) e
l'esperienza di \emph{play} (se piace giocarlo), descritti da Ruggero. Per un
designer i quadranti ``scomodi'' sono i due fuori diagonale: un gioco può
piacere pur essendo progettato male, e va saputo progettare bene anche un gioco
che non piace.}
\label{a1:fig:quadranti}
\end{figure}

\paragraph{La domanda di Tancredi.} Tancredi chiede conferma: il \emph{play}
non è il ``giocare senza regole'', ma può essere anche il giocare a un gioco? Sì:
il \emph{play} è l'\textbf{atto}, e può essere fatto con un \emph{game} o senza.
E un \emph{game} può esistere senza che nessuno lo giochi: gli scacchi non sono
meno scacchi perché in quel momento nessuno ci gioca. E se, giocando, una
persona applica le regole in modo diverso? Allora \textbf{non si sta giocando
quel gioco}.

\subsection{House rule: il caso di Monopoly}

Qui Ruggero fa un po' di \emph{foreshadowing} su una lezione di Daniele: il
discorso degli \textbf{hack}, delle \textbf{house rule}, delle \textbf{mod},
cioè di come modificare un gioco esistente. Il punto è che \textbf{un gioco
modificato da una house rule è un altro gioco}: cambiando le regole cambiano
necessariamente anche le \textbf{dinamiche} che ne derivano, e quindi
l'esperienza. Si immagini \emph{Monopoly} in cui ogni turno si possa, volendo,
prendere gratis dei soldi dalla banca: tutto il resto è identico, eppure il gioco
cambia.

E questo è talmente interiorizzato che ci sono regole che si applicano ``senza
neanche rendercene conto''. Ruggero chiede quanto duri in media una partita a
\emph{Monopoly}. Gianmarco: ``quattro ore''; altri: ``un pomeriggio''; un altro:
``finché uno dei giocatori non aggiunge la regola: ora vince chi\ldots''. In
realtà, \textbf{seguendo le regole scritte}, una partita non dura mai più di
\textbf{un'ora e venti, un'ora e mezza}. La domanda che fa sempre, e la risposta
è sempre la stessa: perché si pensa che \emph{Monopoly} duri tantissimo? Qual è
la regola che nessuno rispetta?

Uno studente indovina: \textbf{le aste}. Secondo il regolamento, se chi finisce
su un terreno non lo compra, il terreno \textbf{va all'asta} tra i giocatori, e
chiunque può aggiudicarselo. Molti invece giocano come se finire su un terreno
desse solo la \emph{facoltà} di comprarlo o di metterlo all'asta\footnote{Più
precisamente, nel regolamento ufficiale, se il giocatore che finisce su una
proprietà libera decide di non comprarla, la banca la mette all'asta e tutti i
giocatori possono fare offerte, compreso chi l'ha rifiutata. L'asta quindi non è
una scelta del giocatore, ma la conseguenza obbligata del suo rifiuto.}. Con la
regola applicata, di aste se ne fanno moltissime: all'ottavo turno ci sono già
otto terreni assegnati, al ventesimo venti, la gente va in bancarotta
velocissimamente e la partita dura un'ora e mezza. Anche le partite competitive
di \emph{Monopoly}, nota Ruggero, durano un'ora e venti, un'ora e mezza.

Uno studente obietta che giocava con le aste e durava comunque quattro ore.
``Avete dei nemici con l'analisi paralisi'', ride Ruggero; e cita un'altra house
rule tipica dell'infanzia: genitori che danno soldi per non far andare in
bancarotta. Ma \textbf{la bancarotta è parte del gioco}: se la si toglie,
\emph{Monopoly} non regge. Sul tema che \emph{Monopoly} vuole trattare, e sul
suo antenato (\emph{The Landlord's Game}), ``non apriamo adesso un altro vaso di
Pandora''.

\paragraph{Approfondimento: Monopoly, la durata e le house rule nei libri.}
L'esempio di Ruggero si ritrova quasi identico in \emph{The Art of Game Design}.
Parlando del bilanciamento tra partite brevi e lunghe (\emph{Short vs. Long}),
Schell scrive che \emph{Monopoly}, giocato secondo le regole ufficiali,
``finisce spesso in circa novanta minuti'', ma molti giocatori trovano quelle
regole troppo dure e le modificano, distribuendo premi in denaro e allentando i
vincoli sull'acquisto delle proprietà: con l'effetto collaterale di far durare la
partita molto di più, ``tipicamente tre ore, o anche di più''. Schell classifica
le \textbf{house rule} come regole create dai giocatori in risposta a un
difetto percepito dopo qualche partita. Anche \emph{Challenges for Game
Designers}, nel capitolo sull'aggiungere e togliere meccaniche, parte dallo
stesso gioco: il \emph{Parcheggio gratuito}, che per le regole originali non fa
succedere nulla, mentre quasi tutti giocano con una house rule che dà un bonus a
chi ci finisce sopra. Le due fonti confermano la lezione: una piccola regola
cambiata cambia le dinamiche, e le dinamiche cambiano l'esperienza.

\subsection{Il fine del game è il play?}

Un altro studente commenta: va bene la distinzione, ma il fine del \emph{game} è
il \emph{play}; quello che si vede è sempre l'effetto. Ruggero è d'accordo: di
solito, se si fa un gioco, lo si fa per farlo giocare, come si scrive un libro
per farlo leggere (anche se esistono, ed ``è già stato esplorato'', libri scritti
per non essere letti e giochi per zero giocatori). Ma la distinzione serve a una
cosa precisa: \textbf{separare, nel playtest, i feedback sul gioco dai feedback
sul gruppo di gioco} e sulle sue dinamiche. Se una persona esce frustratissima da
un playtest, non è necessariamente colpa del gioco; se una persona dice ``questo
è il gioco migliore del mondo'', potrebbe non essere merito del gioco ma del
gruppo. Capire a che cosa si riferisce un feedback serve a capire \textbf{dove
iterare} il prototipo.

Lo studente osserva che l'area è sfumata: alcuni giochi si basano proprio
sull'avere le persone giuste al tavolo. Un gioco di \textbf{negoziazione} con
giocatori abituati ai giochi ``multi-solitario'' diventa ingiocabile; e nei giochi
a \textbf{economia chiusa} (cita \emph{Container}\footnote{La trascrizione cita
anche un secondo titolo non ricostruibile (``Tokyo su Vigimarket'').}) un gruppo
economicamente conservativo, che spende poco e non fa girare l'economia, fa
durare la partita il doppio. Ruggero rilancia: e se il gioco dura il doppio,
questa dinamica cambia la \textbf{percezione di peso} del gioco, le energie spese,
quanto si è stanchi a fine partita. ``Tutto è collegato. \textbf{I giochi sono
ecosistemi di meccaniche}, e cambiandone una, le altre cambiano di
conseguenza''.

\paragraph{La parola vietata, alla prova.} Lo studente aggiunge che ci sono giochi
che ha trovato ``molto sbagliati, ma che mi divertono''. ``Attenzione: non si può
dire \emph{divertimento}!'', lo ferma Ruggero; nelle presentazioni aveva lasciato
correre, ``ma lì non vi avevo ancora letto la regola''. Lo studente ci prova:
``giochi che poi voglio \textbf{rigiocare}'', anche se ``sono evidentemente
sbagliati, rotti''. Ruggero apprezza, ma non vuole anticipare cose che meritano il
loro tempo: se ne parlerà nelle lezioni sui framework e soprattutto in quella di
\textbf{emotion design}, ``forse la mia lezione preferita del corso''. (Lo
studente ammette che ``forse alla fine si gioca per quello, per l'emozione''.)

\subsection{I giochi ``comunicati male''}

L'ultima domanda corposa riprende un'affermazione di inizio lezione: ci sono
giochi belli ma comunicati male. Qualche esempio? ``Mi basta girarmi per
vederne 48'', risponde Ruggero indicando gli scaffali alle sue spalle. Più che
titoli, elenca le \textbf{ragioni}:
\begin{itemize}
  \item \textbf{l'Italia localizza lentamente}. Basta guardare i premi del gioco
        dell'anno in Italia: spesso vanno a giochi usciti due o tre anni prima in
        Germania, Francia, Stati Uniti o Regno Unito, perché localizzare richiede i
        suoi tempi. Il mercato italiano è quindi molto dipendente dall'estero,
        ``molto impaurito'', e spesso aspetta che un gioco vada bene all'estero
        prima di portarlo in Italia; e alcuni giochi ``perdono il treno'';
  \item \textbf{il publisher originale chiude}: come si pubblica un gioco se non
        c'è più l'azienda con cui fare un contratto? A Ruggero è successo più
        volte;
  \item \textbf{il dipartimento commerciale italiano non sa comunicare il gioco}:
        si fa male un'edizione, la stampa vende pochissimo e non viene più
        ristampata;
  \item \textbf{la scala del mercato}: in Italia un gioco da tavolo di peso medio
        (un \emph{eurogame}, un \emph{German}) che vende 1.500 copie ``è tanto'';
        all'estero se ne vendono magari 10.000. Con 10.000 copie in inglese ci
        sono i soldi per una campagna di comunicazione: un trailer, i
        recensori, i tutorial online. In Italia tutto questo è ``ridotto all'osso'',
        e la cosa si riflette su stampa e divulgazione: anche i canali per
        informarsi su un gioco prima di comprarlo sono pochi (se ne parlerà nella
        lezione su ricerca e sviluppo). Un gioco spesso non arriva al mercato non
        perché sia brutto, ma perché l'editore non lo sa comunicare, ne fa poche
        copie e va fuori produzione;
  \item \textbf{il tempismo}: party game progettati nel 2019 e usciti nel 2020, in
        pieno Covid, quando non si poteva giocare al tavolo, sono falliti tutti, e
        gli editori ci hanno perso. ``Non potevano immaginarselo'';
  \item \textbf{i temi caldi}: dopo l'invasione russa dell'Ucraina e la guerra in
        Palestina, alcuni editori non pubblicano o non localizzano giochi
        strategici e di guerra. Ruggero ha lavorato a \textbf{\emph{Status}}, un
        gioco sulla geopolitica che tocca, tra gli altri temi, l'invasione russa
        dell'Ucraina e il genocidio in Palestina, e alcuni editori che di
        solito localizzano i loro giochi (per esempio il partner di riferimento in
        Russia) non lo hanno localizzato per motivi politici.
\end{itemize}
Lo studente voleva capire che cosa voglia dire comunicare male: non metterci
abbastanza energia, o comunicare una cosa per un'altra, ``come il trailer
sbagliato, che poi vai al cinema e ti trovi un altro film''. ``È un mix'', e la
domanda ``era complessa e meritava più spazio''.

\subsection{I giochi di carte collezionabili e i saluti}

L'ultima domanda (``posso rompere le balle un'ultima volta?'') riguarda i
\textbf{giochi di carte collezionabili}: funzionano con le stesse regole dei
giochi da tavolo? Se ne parlerà alla prossima lezione: è proprio il motivo per
cui la lezione sui casi studio dura tre ore, perché ogni genere di gioco ha le
sue dinamiche e le sue regole, e anche il modo di pubblicare cambia molto tra un
TCG, un gioco di ruolo da tavolo, un LARP e un gioco pervasivo.

Ruggero saluta: le email dei tre docenti sono all'inizio delle slide, la
registrazione arriverà, ``ci si vede questo venerdì''.

% ---------------------------------------------------------------------
\section*{Punti chiave del capitolo}
\addcontentsline{toc}{section}{Punti chiave del capitolo}

\begin{itemize}
  \item Il corso di \textbf{Analog Game Design} (seconda edizione) dura circa tre
        mesi, fino a dicembre, ed è tenuto da tre docenti principali:
        \textbf{Daniele} (narrative design, gioco di ruolo, fiction),
        \textbf{Ruggero} (game design, ricerca e sviluppo) e \textbf{Beatrice}
        (giochi pervasivi), con tre ospiti su RNG, development e game production.
  \item Lo scopo è formare \textbf{autori analogici} consapevoli: strumenti da
        usare anche fuori dal corso, non il ``gioco dei sogni''. Il corso è
        \textbf{intersezionale} (nessuna lezione dedicata a un solo tipo di
        gioco), \textbf{cross-generazionale}, \textbf{pratico e adattabile},
        basato sulle \textbf{domande} (nessuna domanda è stupida), sul
        \textbf{feedback} invece che sul voto e sul \emph{fail forward}.
  \item Il \textbf{project work} si fa in tre gruppi da tre, formati dai docenti;
        pitch il 12 ottobre, sviluppo nelle due pause, tutoraggi il sabato,
        consegna il 4 dicembre, presentazione l'11 e giornata finale con gli
        editori il 14. Si valutano documentazione, presentazione,
        autovalutazione, crescita e atteggiamento.
  \item \textbf{Niente intelligenza artificiale} nel corso: i publisher
        rifiutano pitch e prototipi con immagini generate, i regolamenti scritti
        dall'IA descrivono la stessa meccanica in modi diversi, e soprattutto chi
        la usa perde l'occasione di imparare. Un prototipo può essere fatto a
        matita.
  \item Le \textbf{figure dell'industria}: game designer, game developer,
        editor, publisher, sales e marketing, graphic designer, illustratore,
        sensitivity reader (che non è un censore) e producer. Nelle realtà
        piccole una persona copre più ruoli.
  \item Nell'analogico il \textbf{game designer} progetta l'esperienza creando
        le meccaniche; il \textbf{game developer} fa da \textbf{mediatore} con il
        publisher (bilanciamento, modelli matematici, adattamenti). Pretendere
        che il proprio gioco non venga modificato è un \emph{red flag}.
  \item Il \textbf{marketing} non decide che cosa pubblicare ma come; è
        un'azione, ed esiste anche se lo si ignora.
  \item In \textbf{Italia} il gioco è ancora visto come infantile, il mercato
        punta sulla localizzazione e investe poco su giochi e designer interni;
        all'estero (soprattutto dalla Germania in su) ci sono tradizione adulta,
        premi di prestigio e developer interni.
  \item \textbf{Game} è l'artefatto (le regole, il sistema, l'oggetto);
        \textbf{play} è l'atto di giocare. Si fanno giochi, non giocattoli.
        Separarli serve a distinguere i feedback sul gioco da quelli sul gruppo
        e a riconoscere i quattro incroci tra ``progettato bene/male'' e ``mi
        piace/non mi piace''.
  \item \textbf{MDA}: meccaniche (regole scritte), dinamiche (regole in azione),
        estetiche (esperienza, non grafica); nel \textbf{DPE}: design, play,
        experience.
  \item \textbf{Parole vietate} in ogni materiale scritto: \emph{fun},
        \emph{immersion}, \emph{entertainment} (divertimento, immersione,
        intrattenimento). Bisogna trovare parole più precise.
  \item Distinguere \textbf{chi crea, chi gioca e chi facilita}; nessun gioco
        nasce perfetto: \textbf{playtest e iterazione} sono inevitabili, e un
        \textbf{dizionario comune} aiuta a non perdersi.
  \item \textbf{Tabletop Simulator} è il software di settore per giocare (non
        creare videogiochi) i giochi da tavolo in digitale. Primo task: giocarci
        tre o quattro volte entro il 10 ottobre.
  \item Una \textbf{house rule} crea un altro gioco: \emph{Monopoly} con le aste
        applicate come da regolamento dura circa un'ora e mezza.
\end{itemize}

%    % =====================================================================
%  Analog Game Design, lezione 2 — I casi studio: una panoramica a volo
%  d'aquila dei giochi pervasivi, dei giochi di ruolo dal vivo, dei
%  giochi di ruolo da tavolo e dei giochi da tavolo
%  Fonte: trascrizione della seconda lezione (AGD2) del corso breve di
%  Analog Game Design, seconda edizione: una lezione online su Zoom di
%  circa tre ore, di venerdì sera, tenuta a più voci dai tre docenti
%  "cuore" del corso (Ruggero, Daniele e Beatrice), con una pausa di
%  pochi minuti e moltissime slide di esempi.
%  Pacchetti richiesti nel preambolo del libro (già presenti nel
%  main.tex): \usepackage{booktabs}, \usepackage{tikz} con
%    \usetikzlibrary{shapes.geometric,shapes.symbols,shapes.misc,
%      arrows.meta,positioning,calc,fit}
%  (le librerie vanno caricate nel preambolo, non nel capitolo).
%  Consigliati, come per gli altri capitoli: \usepackage[italian]{babel}
%  e \usepackage[T1]{fontenc} (il capitolo compila anche senza).
%
%  Collocazione nel libro: parte "Analog Game Design", secondo capitolo
%  (numerazione propria 2.2). Nel file principale del libro, subito dopo
%  il capitolo 2.1:
%
%    \input{AGD/AGD1_introduzione_all_analog_game_design}
%    \input{AGD/AGD2_casi_studio}
%
%  Tutte le etichette di questo capitolo hanno il prefisso "a2:"
%  (a2:cap:..., a2:sec:..., a2:fig:..., a2:tab:...). Il capitolo è
%  autonomo: non contiene \ref verso altri capitoli.
%  Le figure e le tabelle sono ricostruzioni schematiche di quanto i
%  docenti spiegano a voce e con le slide, non copie delle slide.
% =====================================================================

\chapter[I casi studio dell'Analog Game Design]{I casi studio: giochi
pervasivi, giochi di ruolo dal vivo, giochi di ruolo da tavolo e giochi da
tavolo}
\label{a2:cap:casi}

Questa è la \textbf{seconda lezione} del corso di \textbf{Analog Game Design}, e
mantiene la promessa fatta alla fine della prima: è una lezione di tre ore
dedicata quasi interamente ai \textbf{casi studio} (\emph{case studies}). I
docenti la definiscono ``un overview'', una panoramica ``a volo d'aquila'':
``voleremo in tutti i settori ludici analogici e guarderemo che cosa c'è''.
Non si scende nei dettagli tecnici dell'analisi (ci saranno lezioni apposite,
soprattutto quella di \textbf{ricerca e sviluppo}), ma si apre la ``diga'' degli
esempi: decine di giochi, eventi, esperienze, movimenti e manifesti, scelti per
mostrare quante cose esistono, quanto sono diverse tra loro e quanto i confini
tra i generi siano porosi.

La lezione si svolge \textbf{online}, su Zoom, un \textbf{venerdì} sera (``che
bella la lezione del venerdì! Sembra quasi un aperitivo''), ed è divisa tra i
tre docenti principali del corso, ciascuno sul proprio territorio:
\textbf{Ruggero} introduce l'argomento e presenta i \textbf{giochi di ruolo dal
vivo} (LARP) e i \textbf{giochi da tavolo}; \textbf{Beatrice} presenta i
\textbf{giochi pervasivi}; \textbf{Daniele} presenta i \textbf{giochi di ruolo
da tavolo}. ``Ce lo smezziamo un po' tutti'', perché l'argomento è diviso ``in
tante categorie, sottocategorie, sotto-sottocategorie''. Dopo i pervasivi c'è
una pausa di pochi minuti; gli ultimi dieci minuti sono lasciati alle domande.

Tre avvertenze di lettura, le stesse del capitolo precedente. La prima: la
trascrizione è automatica e storpia moltissimi nomi propri, soprattutto i titoli
dei giochi, quasi sempre in inglese e pronunciati velocemente (``Blaze in the
Dark'', ``Cobo Bobby'', ``Sifol'', ``Jeep Form''\ldots). Dove la ricostruzione è
sicura, o molto probabile, il capitolo usa la forma corretta e lo segnala in
nota; dove non lo è, riporta la forma della trascrizione e lo dice. La seconda:
durante la pausa la trascrizione registra, ripetute in loop per decine di righe,
alcune frasi senza senso (anche all'inizio della registrazione succede qualcosa
di simile), che il capitolo ignora. La terza: i docenti si passano spesso la
parola senza che la trascrizione segni il cambio di voce; il capitolo attribuisce
le parole a chi le pronuncia quando il contesto lo rende chiaro (chi parla di sé,
chi viene chiamato per nome, chi ``passa la parola''), e resta generico
altrimenti.

Quanto alla data, la lezione si colloca con ogni probabilità \textbf{venerdì 11
settembre 2026}\footnote{La data non è mai detta. La si deduce da alcuni
indizi: la prima lezione si era chiusa con ``ci si vede questo venerdì''; uno
studente assente ha recuperato ``la lezione di lunedì''; a fine lezione si
ricorda che il questionario iniziale va compilato ``entro il 15'' (settembre) e
che la lezione successiva, sull'\emph{emotional design}, sarà di lunedì. Nel
2026 il primo venerdì utile prima del 15 settembre, dopo il lunedì della prima
lezione, è l'11. È una deduzione, non un'informazione esplicita.}: i riferimenti
a ``l'altra volta'' e ``l'anno scorso'' vanno letti rispetto a quel momento.

\paragraph{Come leggere questo capitolo.} Più ancora della lezione precedente,
questa è un \textbf{catalogo}. I docenti lo dicono esplicitamente: non ci sarà
nessuna interrogazione sui singoli giochi, e lo scopo non è ricordarli tutti ma
sapere che esistono, avere dei punti di riferimento e, soprattutto, abituarsi a
guardare oltre il proprio genere preferito. Per questo il capitolo riporta
\emph{tutti} i casi citati, con quello che i docenti ne dicono, ma li organizza
in tabelle e schemi per renderli consultabili ``in fase di ricerca e sviluppo'',
esattamente come fa Ruggero con le sue slide (``non sapete quante volte mi sono
trovato a riaprire proprio esattamente queste slide''). Qua e là, dei paragrafi
di \textbf{approfondimento} collegano i casi ai libri di game design di
riferimento (\emph{The Art of Game Design} di Jesse Schell, \emph{Challenges for
Game Designers} di Brenda Brathwaite e Ian Schreiber, \emph{21st Century Game
Design} di Chris Bateman e Richard Boon).

Il percorso del capitolo segue quello della lezione:
\begin{enumerate}
  \item l'avvio, Discord, i materiali e i \textbf{libri consigliati} (e
        sconsigliati) per cominciare (sezioni~\ref{a2:sec:avvio}
        e~\ref{a2:sec:libri});
  \item \textbf{perché} analizzare i casi studio e \textbf{che cosa} sono
        (sezione~\ref{a2:sec:perche});
  \item le \textbf{quattro macrocategorie} del gioco analogico, l'ibridazione e
        l'invito a non chiudersi nel proprio genere (sezioni~\ref{a2:sec:macro}
        e~\ref{a2:sec:aperti});
  \item i \textbf{giochi pervasivi}, con Beatrice: la definizione, il cerchio
        magico e diciotto casi studio, dalle app agli \emph{alternate reality
        game}, dal teatro interattivo alle cacce al tesoro
        (sezioni da~\ref{a2:sec:pervasivi} a~\ref{a2:sec:pervasivi-edu});
  \item i \textbf{giochi di ruolo dal vivo}, con Ruggero: boffer, Mind's Eye
        Theatre, chamber e parlor LARP, American freeform, murder party, il
        \textbf{Nordic LARP} con i suoi manifesti, lo Jeepform, il 360°, i tape
        e i black box LARP (sezioni da~\ref{a2:sec:larp} a~\ref{a2:sec:oppresso});
  \item i \textbf{giochi di ruolo da tavolo}, con Daniele: le quattro parole
        (conflitto, azione, risoluzione, effetto), la storia dal primo
        \emph{Dungeons \& Dragons} alla New Wave di The Forge, l'OSR, i giochi
        focalizzati, gli story game, i giochi per due, in solitario e i micro
        giochi di ruolo (sezioni da~\ref{a2:sec:gdr} a~\ref{a2:sec:gdr-formati});
  \item i \textbf{giochi da tavolo}, di nuovo con Ruggero: diciassette generi,
        dal \emph{roll and move} ai Legacy, e gli ibridi
        (sezioni da~\ref{a2:sec:tavolo} a~\ref{a2:sec:tavolo-ibridi});
  \item la chiusura, i materiali caricati, le \textbf{domande finali} e due
        avvertimenti importanti: non fare un gioco di carte collezionabili, non
        fare un gioco per ``aggiustare'' un altro gioco
        (sezioni da~\ref{a2:sec:chiusura} a~\ref{a2:sec:cimiteri}).
\end{enumerate}

% ---------------------------------------------------------------------
\section{L'avvio della lezione}
\label{a2:sec:avvio}

La registrazione parte mentre i docenti si stanno ancora collegando (``mamma
mia, la burocrazia italiana'') e si controlla che il simbolo REC lampeggi.
Mentre Daniele prova a condividere le slide (``aspettate, sono un attimo
bloccato'') e chiede di essere nominato co-organizzatore della riunione, gli
studenti entrano uno alla volta: Bea, Lorenzo, Nicola, Rita, Gianmarco,
Tiziana, Giordano. L'atmosfera è rilassata. Si scherza sul fatto che la lezione
del venerdì ``sembra quasi un aperitivo'': un docente si mette ``bello comodo'',
un altro ha a casa ``le cose per fare l'aperitivo'' e chiede il permesso di
prepararlo nel frattempo, ``purché siate sul pezzo''. Daniele scrive anche di
dover ancora rispondere a un messaggio di Beatrice; ``non preoccuparti'', gli
risponde lei: settembre è un mese di ripartenza, ``in realtà è l'inizio
dell'anno reale''.

\paragraph{Giordano.} Arriva \textbf{Giordano}, che era assente alla prima
lezione. Ha recuperato la registrazione di lunedì, e sa che gli si chiederà di
presentarsi. (Alla fine, per mancanza di tempo, la presentazione verrà rinviata
alla lezione successiva: si veda la sezione~\ref{a2:sec:domande}.) All'inizio
mancano ancora Lorenzo e Tancredi; Lorenzo arriva poco dopo, Tancredi dopo la
pausa.

\paragraph{Discord.} Uno dei docenti nota che gli studenti sono ``belli attivi
su Discord'' e se ne rallegra: ``sfruttatelo, perché l'anno scorso lo si
sfruttava poco e invece è uno spazio molto importante''. I docenti leggono il
server ogni tanto, e le domande sono benvenute. L'unica richiesta è
\textbf{taggarli}: ``tra tanti server Discord non riusciamo a leggere'' tutto,
quindi una domanda rivolta a loro va segnalata con un tag. Lo spazio, ricorda,
``è fatto per voi'': per confrontarsi, condividere cose e materiali.

\paragraph{Le dispense.} Gli studenti hanno visto i materiali caricati: le
slide, le dispense ``e quant'altro''. Qualcuno chiede se le dispense siano in
inglese e se siano articoli. Sono i \textbf{saggi citati nella prima lezione}
(per esempio quelli sull'MDA e sul DPE), e \textbf{non è obbligatorio leggerli}:
sono lì ``per completezza'' e per mostrare che dietro gli argomenti delle
lezioni ``c'è una base teorica''. L'importante è che le lezioni facciano capire
gli argomenti. Spesso i docenti citeranno \textbf{paper accademici} di game
design e di \emph{game studies}, e ``purtroppo c'è poco da fare'': molti non sono
mai stati tradotti, o lo sono stati ``in maniera un po' approssimativa''. Per
questo di solito si linka la fonte originale; nulla vieta di cercare traduzioni
o di tradurseli per conto proprio. In ogni caso i docenti cercano di scegliere
paper \textbf{accessibili}: ``non vi linkeremo delle tesi di master di 300
pagine, a meno che non ce le chiediate voi''.

% ---------------------------------------------------------------------
\section{Due libri per cominciare (e uno da rimandare)}
\label{a2:sec:libri}

Prima di entrare nel tema, Daniele torna su due libri citati la volta
precedente, mostrandoli in video.

\subsection{Un volume sull'edu-larp}

A Tiziana, che ha qualche difficoltà con l'inglese, Daniele consiglia un libro
italiano in cui si trovano molti degli argomenti del corso: un volume
sull'\textbf{edu-larp}, il gioco di ruolo dal vivo in ambito educativo, a cura
di \textbf{Andrea Maragliano} ed edito da \textbf{FrancoAngeli}\footnote{La
trascrizione riporta il titolo come ``EduLARP'', ``Edu LARP'' e ``DualArp'', e
il nome del curatore come ``Andrea Maragliano''; il titolo completo del volume
non è detto a lezione.}. Con una frecciata affettuosa, Daniele aggiunge che
contiene ``un'introduzione molto interessante scritta da un certo Ruggero. Non
so chi sia''. Ruggero protesta (``non fare queste cose, che poi sembra che
consiglio i libri in cui ho partecipato''): il libro non è suo, ci hanno scritto
``tutti i miei amici'', e lui per principio non consiglia i libri su cui scrive
(``possiamo anche comprarlo e non leggere l'introduzione'', ride un collega).

Daniele però lo consiglia convintamente:
\begin{itemize}
  \item molti concetti che si trovano nei siti in inglese ``qua dentro sono
        molto spiegati''; lui stesso lo ha usato per studiare;
  \item c'è una parte, ``bellissima'', sull'\textbf{involvement} (il
        coinvolgimento), che spiega anche ``perché non usiamo la parola
        immersione'' (una delle parole vietate della prima lezione); se ne
        parlerà anche con Ruggero;
  \item finalmente si può consigliare ``roba genuina, sana, italiana di game
        design'';
  \item copre l'ambito \textbf{didattico e formativo}, cioè il gioco nelle
        scuole e il gioco come strumento di formazione, che interessa ad alcuni
        studenti del corso.
\end{itemize}
Non bisogna leggerlo tutto: è utile ``leggerlo pian pianino mentre fate il
corso'', se si vuole approfondire.

\subsection{La progettazione dei giochi da tavolo di Walter Nuccio}

Il secondo libro mostrato la volta precedente è \emph{La progettazione dei
giochi da tavolo} di \textbf{Walter Nuccio}. Qui Daniele fa un ``piccolo
cesellamento'': ``se fossi in voi non lo comprerei per il momento'', e lascia la
parola a Ruggero, che ne ha una copia alle spalle. Il giudizio di Ruggero è
preciso:
\begin{itemize}
  \item è più utile \textbf{dopo} aver interiorizzato alcune basi su come si
        progettano i giochi; come primo libro è un rischio, ``perché sembra che
        ti dia delle scorciatoie'';
  \item più che un saggio di game design, è \textbf{``un grosso archivio di
        meccaniche di gioco''}, una sorta di dizionario (``che cos'è un malus,
        per esempio?''), con una parte di \emph{pattern};
  \item il pericolo è applicare quelle meccaniche ai propri giochi
        \textbf{senza capire che cosa fanno}: ``come se apriamo una cassetta
        degli attrezzi e non sappiamo che cosa dobbiamo fare. Gli attrezzi
        sembrano tutti bellissimi, però magari li usiamo male'';
  \item il consiglio è aspettare a comprarlo, o almeno a leggerlo, fino a dopo
        la lezione sull'\textbf{emotional design}, quando si avrà già un'idea di
        come partire con il proprio gioco.
\end{itemize}
Daniele è ``assolutamente d'accordo''. (Nel frattempo si è spostato in cucina:
``benvenuti nella mia cucina, davanti al mio frigorifero, per parlare di user
experience''.) Del libro di Nuccio, e di un'alternativa in inglese, si
riparlerà a proposito dei giochi da tavolo (sezione~\ref{a2:sec:tavolo-libri}).

% ---------------------------------------------------------------------
\section{Perché analizzare i casi studio}
\label{a2:sec:perche}

Ruggero apre l'argomento del giorno, i \textbf{casi studio}, con due domande
alla classe: \emph{perché} ci teniamo ad analizzarli? E prima ancora: che cosa
ci si aspetta che siano? ``Non è una domanda provocatoria''.

\paragraph{Le risposte degli studenti.} Una prima risposta: sono \textbf{altri
giochi già realizzati}, da analizzare da molti punti di vista, ``anche
insoliti'', non solo da giocatori: dal punto di vista della produzione, o
dell'esperienza utente che se ne è avuta, per capire come migliorarli. Una
seconda risposta: i casi studio sono delle \textbf{milestone}, giochi o
situazioni ``lapidarie'' da cui trarre ``figure di merito'' da sviscerare per
poi generare altre cose. Tiziana chiede invece una cosa ``più basica'': che cosa
si intende esattamente? Modelli, esempi, ma di che cosa?

\paragraph{La definizione.} ``Ottima domanda''. Con \textbf{caso studio} si
intende un \textbf{gioco specifico}, o una \textbf{categoria di giochi}, che,
analizzato, ci dà delle \textbf{informazioni utili}: osservando da vicino che
cosa un gioco fa, possiamo estrapolare informazioni da riproporre, o metodologie
che ci servono per creare i nostri giochi. L'analogia è con la scrittura: chi
vuole scrivere un libro legge tanti libri, e capisce che cosa fanno di positivo
quelli che gli piacciono e che cosa non fanno quelli che non gli piacciono, per
evitare le cose che non gli piacciono e ``iterare sulle cose che ci piacciono''.
Con i giochi è lo stesso: si guarda che cosa fanno di positivo e di negativo,
\textbf{tenendo conto anche del contesto} in cui sono nati.

\paragraph{Oggi la panoramica, poi gli strumenti.} Gli strumenti più tecnici
per analizzare nel dettaglio arriveranno con le lezioni di \textbf{ricerca e
sviluppo}. Oggi si fa una ricognizione di tutti i settori.

% ---------------------------------------------------------------------
\section{Il gioco analogico e le sue quattro macrocategorie}
\label{a2:sec:macro}

\subsection{Un termine ombrello}

Prima una precisazione. Nel corso, \textbf{gioco analogico} è un
\textbf{termine ombrello} che copre i giochi da tavolo (\emph{board game}), i
giochi di carte (\emph{card game}), i giochi di ruolo da tavolo, i giochi di
ruolo dal vivo, i giochi pervasivi, le cacce al tesoro, le \emph{escape room},
i \emph{murder party}: tutti quei giochi che ``generalmente non sono focalizzati
e incentrati sulla programmazione, sull'uso di un terminale, sull'uso di uno
schermo''; in breve, che \textbf{non sono videogiochi}. Esistono cose ibride,
che si vedranno, ma l'area di interesse principale resta il gioco analogico.

\subsection{Le quattro macrocategorie}

\slidefig{img/AGD/agd2-categorie}{Le macrocategorie del gioco analogico: giochi da tavolo, giochi di ruolo da tavolo, LARP e, dietro, l'ibridazione.}{slide \emph{Analog -- Case Studies}}{fig:agd2-categorie}


La lezione parte dal ``frame attuale'': quattro \textbf{macrocategorie}, a loro
volta divise in sottocategorie. Ruggero avverte subito che ``non sono separate da
muri di acciaio non comunicanti'': i casi ibridi, difficili da mettere in una
sola categoria, sono frequenti. I tre docenti presentano ciascuno la propria.

\begin{description}
  \item[Gioco da tavolo (Ruggero).] La categoria dei \emph{board game}. ``Per
        motivi di industria'' vi si includono anche i giochi di carte e i
        \emph{party game}: tutti i giochi \textbf{pubblicati e prodotti
        all'interno dell'industria del gioco da tavolo}, che abbiano o no un
        tabellone.
  \item[Gioco di ruolo da tavolo (Daniele).] GDR in italiano, TTRPG
        (\emph{TableTop Role-Playing Game}) in inglese. È una categoria ``molto
        particolare'', perché anzitutto \textbf{non c'è una definizione
        univoca}: sul mercato si trovano cose molto diverse, perfino giochi che
        sono quasi ``dei LARP da tavolo'', vicini ai \emph{chamber LARP}. Anche
        i \textbf{formati} sono diversissimi: la slide mostra \emph{Dungeons \&
        Dragons}; \emph{Neverland}, un supplemento\footnote{Daniele lo chiama
        ``un'espansione per D\&D''; più avanti lo descrive come una
        reinterpretazione di Peter Pan in chiave OSR, compatibile con le regole
        di D\&D.}; \emph{Dune}, in volumi cartonati ``importanti''; ma anche
        \emph{The Quiet Year}\footnote{La trascrizione dice ``Quietier'' e, più
        avanti, ``The Quiet Ear''.}, che sta in un sacchettino; e \emph{Lady
        Blackbird}\footnote{La trascrizione dice ``Lady Blackboard''. \emph{Lady
        Blackbird} è un gioco di ruolo distribuito gratuitamente in PDF.}, che è
        solo digitale: ``tu lo prendi, lo stampi in bianco e nero e te lo
        giochi''. L'edizione italiana di \emph{Montsegur 1244}\footnote{La
        trascrizione dice ``Monsegur''.} è una sorta di scatola che si apre;
        \emph{CBR+PNK}\footnote{La trascrizione dice ``CBR PNK''.} ``si apre in
        un modo stranissimo'' e contiene dei pieghevoli (\emph{leaflet}); ci sono
        giochi di ruolo in formato brochure. Daniele li mostrerà meglio nella
        lezione sulla \emph{fiction}. La conclusione: nei giochi di ruolo
        \textbf{non si può usare il metro del formato}.
  \item[Gioco di ruolo dal vivo (Ruggero).] Il LARP, \emph{Live Action
        Role-Playing game}: ``anche questa è una categoria gigantesca''.
  \item[Gioco pervasivo (Beatrice).] La categoria meno conosciuta in Italia, e
        anche per questo inserita nel corso.
\end{description}

\paragraph{Che cosa distingue il gioco di ruolo.} Daniele chiede alla classe
che cosa distingua i \emph{role-playing game}. La risposta è nel nome: il
\textbf{role playing}, cioè un \textbf{elemento di interpretazione}, che non è
detto sia interpretazione di un \emph{personaggio} (in \emph{The Quiet Year},
per esempio, non c'è un personaggio, ``abbiamo qualcosa di diverso''). E c'è un
elemento ``importantissimo'': la \textbf{fiction}, cioè la narrazione. Non la
narrazione intesa come ``scrivere una trama e giocarla'', ma quella che nasce al
tavolo, la \textbf{narrazione emergente}: non si è previsto che l'avventura
andasse in quel modo, ma mentre si gioca gli altri dicono cose che la portano in
una certa direzione. Se ne parlerà nella lezione sulla fiction emergente.
Categorizzare i GDR è ``veramente molto arduo'', e anche il pubblico è
particolare: ``ognuno ha il suo palato'', c'è chi gioca tutto e chi gioca solo un
sottogenere. A una domanda dalla classe (il gioco \emph{Lovecraftesque}, citato
da qualcuno, rientra tra i giochi di ruolo da tavolo?) la risposta è sì: e di
che \emph{tipo} di gioco di ruolo sia si capirà con le sottocategorie che Daniele
presenterà più tardi.

\paragraph{Perché il LARP è separato dal GDR.} Ruggero spiega che i LARP sono
separati dal GDR da tavolo ``semplicemente per motivi di settore e di
industria''. Il gioco di ruolo da tavolo viene venduto con una produzione più
vicina al gioco da tavolo: si produce un \textbf{oggetto fisico} che poi si
vende. Il LARP invece è ancora spesso organizzato \textbf{a eventi}: si partecipa
\textbf{comprando un biglietto}, e si va a giocare in un posto specifico, a
un'ora specifica, con persone che sono lì per giocare quel gioco. Da questo
punto di vista è più vicino ad alcuni giochi pervasivi. Le sottocategorie sono
``una valanga'': la più famosa è il \textbf{boffer} (il cliché del LARP: persone
in costume medievale ``che si picchiano con spade di plastazote o di gommapiuma''),
ma ce ne sono di molto più astratte e di molto più accessibili, che si possono
organizzare anche in casa.

\subsection{I giochi pervasivi: una prima definizione}

Beatrice chiede se qualcuno abbia mai sentito parlare di \emph{pervasive games}.
Quasi nessuno; una studentessa ne ha sentito parlare solo leggendo la
descrizione del corso, e si è incuriosita. ``Ci sta tantissimo'': in Italia non
sono ancora molto conosciuti, ed è uno dei motivi per cui sono nel percorso.

Tiziana chiede se le \textbf{cacce al tesoro} organizzate nelle città rientrino
negli \emph{urban game}. ``Assolutamente sì'': gli \textbf{urban game} sono una
categoria dei giochi pervasivi, e nella slide si vedono proprio come un
sottoinsieme del ``grande cappello''. Beatrice dà due definizioni:
\begin{enumerate}
  \item la definizione ``standard'', dai \emph{game studies}: un gioco
        pervasivo è un gioco che \textbf{espande il proprio cerchio magico
        spazialmente, temporalmente e socialmente};
  \item la definizione più intuitiva, da cui partire per capire poi meglio la
        prima: un gioco pervasivo è un gioco che \textbf{utilizza i luoghi e gli
        oggetti della vita quotidiana e li trasforma in luoghi e oggetti di
        gioco}.
\end{enumerate}
Se la prima non è chiara, ``non è un problema, la capiremo con i casi studio''.
Negli urban game ``non giochiamo seduti al tavolo, ma giochiamo in città, e
stiamo comunque giocando''. La parola \emph{pervasivo} indica la pervasività del
gioco, che entra in situazioni in cui non siamo abituati a immaginarlo, anche se
magari siamo abituatissimi a viverlo senza pensarci.

\paragraph{Le cene con delitto.} Tiziana fa un'altra domanda (Beatrice ricorda
che si può sempre interrompere, anche con la mano alzata delle reazioni di
Zoom): le \textbf{cene con delitto} rientrano tra i pervasivi? E il
\textbf{teatro interattivo}, visto che in certe cene con delitto chi le propone
recita una vera rappresentazione? Per Beatrice, ``a mio dire, sì, entrambe''.
Ma le cene con delitto, o \emph{murder party}, esistono in molte versioni:
\begin{itemize}
  \item in una, gli \textbf{attori} vengono al tavolo con te, magari
        insospettabili, e tu, come persona giocante, li interroghi: questa
        versione assomiglia di più al teatro interattivo, che è incluso nei
        giochi pervasivi;
  \item in un'altra, \textbf{ognuno riceve un ruolo}, con una sorta di copione
        (``stasera sei la sorella della vittima''): questa va più nella
        direzione dei \textbf{giochi di ruolo dal vivo}.
\end{itemize}
Sono distinzioni ``che facciamo per allenarci a comprendere meglio'', ma che
spesso si sovrappongono: anche quando si è investigatori e si interrogano degli
attori, si è dentro una dinamica di gioco di ruolo, ``anche se magari non come
la immaginiamo in D\&D o in giochi più blasonati''. Le definizioni ``sono
malleabili'': si imparano le parole e poi si vede come usarle nel caso
specifico.

\subsection{Ibridazione e sperimentazione}

Nell'angolo della slide, dietro tutte le categorie, c'è ``una piccola cosina'':
\textbf{ibridazione e sperimentazione}. ``In realtà le categorie non esistono'',
dice Ruggero: esistono come divisioni tra industrie, ma escono continuamente
giochi che fanno da ponte. La slide mostra alcuni esempi:
\begin{itemize}
  \item \emph{Fateforge}\footnote{La trascrizione dice ``FitForge'' e, più
        avanti, ``Fateforge''.}, un gioco da tavolo che usa un'\textbf{app}
        digitale integrata nel gioco: ``un po' videogioco, un po' board game'';
  \item un gioco di aste che usa anch'esso un cellulare (il titolo non è
        detto);
  \item dei \emph{mock trial} (processi simulati) ibridi, a metà tra il murder
        party e il roleplay, ``un po' roleplay, ma anche un po' giochi
        pervasivi'';
  \item \emph{Alice is Missing}, un gioco di ruolo da tavolo che usa una
        \textbf{chat} come Telegram tra le meccaniche di gioco.
\end{itemize}
Le ibridazioni rendono spesso difficile mettere un gioco ``al 100\%'' in uno di
questi ``cesti''.

\begin{figure}[htbp]
\centering
\begin{tikzpicture}[
    font=\small,
    macro/.style={draw, thick, rounded corners=4pt, fill=gray!12,
                  text width=5.2cm, align=left, inner sep=6pt},
    titolo/.style={font=\bfseries},
    ponte/.style={draw, dashed, rounded corners=3pt, fill=white,
                  text width=3.4cm, align=center, font=\footnotesize,
                  inner sep=4pt}
  ]
  \node[macro] (tav) at (0,0) {{\bfseries Gioco da tavolo} \hfill
    \footnotesize Ruggero\\[2pt]
    \footnotesize board game, card game, party game: tutto ciò che è
    prodotto dall'industria del gioco da tavolo};
  \node[macro] (gdr) at (6.6,0) {{\bfseries Gioco di ruolo da tavolo}
    \hfill \footnotesize Daniele\\[2pt]
    \footnotesize GDR / TTRPG: interpretazione e fiction emergente; nessuna
    definizione univoca, formati diversissimi};
  \node[macro] (larp) at (0,-3.6) {{\bfseries Gioco di ruolo dal vivo}
    \hfill \footnotesize Ruggero\\[2pt]
    \footnotesize LARP: spesso a evento, con biglietto; dal boffer al
    black box};
  \node[macro] (perv) at (6.6,-3.6) {{\bfseries Gioco pervasivo}
    \hfill \footnotesize Beatrice\\[2pt]
    \footnotesize luoghi e oggetti quotidiani trasformati in luoghi e oggetti
    di gioco; urban game, ARG, teatro interattivo};
  \node[ponte] (ibr) at (3.3,-1.8) {\textbf{ibridazione e
    sperimentazione}\\ app, chat, murder party, mock trial, eventi\ldots};
  \draw[dashed, thick] (ibr) -- (tav);
  \draw[dashed, thick] (ibr) -- (gdr);
  \draw[dashed, thick] (ibr) -- (larp);
  \draw[dashed, thick] (ibr) -- (perv);
  \draw[{Stealth[length=2mm]}-{Stealth[length=2mm]}, gray, thick]
    (gdr.south) -- node[right, font=\footnotesize\itshape, text=black,
    text width=2.4cm] {oggetto venduto vs evento con biglietto} (perv.north);
  \draw[{Stealth[length=2mm]}-{Stealth[length=2mm]}, gray, thick]
    (tav.south) -- node[left, font=\footnotesize\itshape, text=black,
    text width=2.4cm, align=right] {separati per motivi di industria}
    (larp.north);
\end{tikzpicture}
\caption{Le quattro macrocategorie del gioco analogico usate nella lezione, con
il docente che le presenta. Al centro, la ``piccola cosina'' della slide:
ibridazione e sperimentazione, che rendono i confini porosi. Le frecce grigie
riassumono due osservazioni dei docenti: il GDR da tavolo e il LARP sono
separati soprattutto per ragioni di industria (oggetto venduto o evento), e per
lo stesso motivo il LARP è vicino a molti giochi pervasivi.}
\label{a2:fig:macro}
\end{figure}

% ---------------------------------------------------------------------
\section{Non chiudersi nel proprio genere}
\label{a2:sec:aperti}

Prima di aprire la ``diga dei casi studio'', i docenti insistono su alcuni punti
di metodo.

\paragraph{Le categorie servono a ragionare, non a recintare.} Ruggero: le
categorie non vengono date ``per dirvi: le cose sono A, B, C e basta'', ma per
far ragionare su \textbf{quanto sono flessibili i confini}, e per dare esempi di
cose \textbf{adiacenti} a quelle che interessano, che ``possono sbloccarvi delle
cose''. Ognuno è lì per interessi diversi (i giochi di ruolo, le carte, i giochi
formativi e didattici\ldots), ma l'invito è a \textbf{non mettersi i paraocchi}
sul proprio genere: ``è grazie alla \textbf{cross-modalità} e al vedere qualcosa
di risolto in un altro gioco che vi arriva poi un'idea per un gioco vostro che
non c'entra apparentemente niente'', magari nemmeno meccanicamente, con l'altro.

\paragraph{Un corso agnostico di genere.} Daniele aggiunge che il corso è
\textbf{agnostico di genere}: alcune lezioni saranno più importanti per una
categoria, ma i \textbf{concetti di base del design si applicano a qualsiasi
sviluppo ludico}. È fondamentale, perché altrimenti sembra che chi fa un certo
tipo di prodotto abbia bisogno di ``un game design particolare'', e non è così.
Anzi: quando è successo, storicamente anche in Italia, si è finito per
\textbf{ostracizzare e separare le community}, mentre i docenti vogliono
\textbf{metterle insieme}. ``Voi non sapete quante volte i videogiochi vengono
ispirati dai giochi da tavolo, e viceversa''.

\paragraph{Il project work è libero.} C'è anche un motivo pratico. Il progetto
di gruppo finale è \textbf{libero}: i docenti non assegneranno una categoria di
gioco, saranno i gruppi a sceglierla, e può essere anche
un'\textbf{ibridazione}, una \textbf{sperimentazione}, un gioco che non si
colloca perfettamente in un genere (``anzi, sarebbe fantastico''). Ma i gruppi
saranno formati da persone con interessi diversi, ed è proprio per questo che
bisogna restare aperti. L'unica cosa che i docenti chiederanno di non fare è
\textbf{un gioco enorme}: il progetto deve essere \textbf{calibrato}.

\paragraph{``Per me, questo è design''.} Infine, Ruggero chiede di \textbf{non
sottovalutare} la lezione: ad alcuni può sembrare ``meramente un elenco di cose
che vi facciamo vedere per nessun motivo'', ma ``per me questo è design''.
Quando deve trovare una soluzione, va a cercare altri giochi; e gli è capitato
moltissime volte di riaprire proprio quelle slide per ricordarsi com'era fatto un
gioco. E come hanno fatto Tiziana e Rita, bisogna \textbf{fare domande}:
``fermateci, chiedeteci'', i docenti sono lì per approfondire.

% ---------------------------------------------------------------------
\section{I giochi pervasivi: la premessa di Beatrice}
\label{a2:sec:pervasivi}

Si comincia ``in scivolata'' con i giochi pervasivi. Beatrice anticipa che
diversi casi studio hanno \textbf{accenni di digitale}, per tre motivi:
\begin{enumerate}
  \item alcuni casi sono \textbf{molto famosi} e aiutano a capire il concetto di
        pervasive game, anche se non sono prettamente analogici;
  \item un gioco analogico \textbf{non esclude completamente la componente
        digitale}: basti pensare ai giochi da tavolo che usano un'app;
  \item l'\textbf{ibridazione}: si vedranno casi studio ``molto strani''.
\end{enumerate}
Aggiunge una considerazione personale e professionale: quello dei pervasivi è
un campo in cui c'è ancora \textbf{ampio spazio}, perché sono poco conosciuti e
la loro \textbf{definizione è ancora in costruzione}. C'è spazio ``per mettere
anche la vostra voce'', per esprimersi e sperimentare, proprio perché le
definizioni non sono state date in modo rigido come altrove. E anticipa un tema
ricorrente del corso: ``parleremo molto di \textbf{emozioni}, anche di come
accompagnarle all'interno di un'esperienza di gioco''.

\paragraph{Approfondimento: il cerchio magico.} La definizione ``da game
studies'' presuppone il concetto di \textbf{cerchio magico}, che a lezione non
viene spiegato. In \emph{The Art of Game Design}, Schell lo introduce
discutendo le definizioni di gioco: quasi tutte implicano che un gioco sia un
sistema \emph{chiuso}, con un confine, e ``Johan Huizinga lo ha chiamato il
\emph{cerchio magico}'', che ha davvero qualcosa di magico: quando siamo
mentalmente ``dentro il gioco'' abbiamo pensieri, sentimenti e valori diversi da
quando ne siamo fuori. Schell lo reinterpreta come il nostro \textbf{sistema
interno di risoluzione dei problemi}: una ``microrealtà'' semplificata in cui gli
elementi acquistano un'importanza che non devono avere fuori dal gioco. Nel
capitolo sulle meccaniche, lo stesso Schell identifica il cerchio magico con lo
\textbf{spazio} del gioco, la prima delle sue sette categorie di meccaniche. Il
gioco pervasivo è dunque quello che sposta, allarga o rende sfumato questo
confine \textbf{nello spazio} (si gioca nella città, non sul tavolo), \textbf{nel
tempo} (non si sa bene quando la partita inizia e finisce) e \textbf{nella
società} (vi entrano persone che non sanno di giocare, o che giocano senza
chiamarlo gioco)\footnote{La formula ``espandere il cerchio magico
spazialmente, temporalmente e socialmente'' è quella proposta da Markus Montola
nei suoi studi sui giochi pervasivi, poi ripresa nel volume \emph{Pervasive
Games: Theory and Design} di Markus Montola, Jaakko Stenros e Annika Waern
(2009). Beatrice non cita la fonte; la nota è un'integrazione.}. È esattamente
quello che Beatrice dirà di \emph{GuruShot}: ``non sappiamo bene dove inizia e
dove finisce l'attività di gioco, ed è per questo che viene utilizzata la parola
pervasivo''.

\begin{figure}[htbp]
\centering
\begin{tikzpicture}[
    font=\small,
    freccia/.style={-{Stealth[length=2.5mm]}, thick},
    etichetta/.style={font=\footnotesize, align=center, text width=3.6cm}
  ]
  % cerchio classico
  \draw[thick, fill=gray!15] (0,0) circle (1.0cm);
  \node[font=\footnotesize, align=center] at (0,0) {cerchio\\magico};
  \node[font=\footnotesize\itshape, align=center] at (0,-1.45)
    {gioco ``classico'':\\tavolo, durata, giocatori};
  % cerchio espanso
  \draw[thick, dashed, fill=gray!6] (7,0) circle (2.1cm);
  \draw[thick, fill=gray!15] (7,0) circle (1.0cm);
  \node[font=\footnotesize, align=center] at (7,0) {cerchio\\magico};
  \draw[freccia] (1.3,0) -- node[above, font=\footnotesize]{gioco pervasivo}
    (4.6,0);
  \draw[freccia, gray] (7,1.0) -- (7,2.1);
  \draw[freccia, gray] (5.95,-0.6) -- (5.18,-1.05);
  \draw[freccia, gray] (8.05,-0.6) -- (8.82,-1.05);
  \node[etichetta, above=0.05cm] at (7,2.15) {\textbf{spazialmente}: la città,
    l'hotel, la stazione dell'autobus};
  \node[etichetta, anchor=north east] at (5.4,-1.5) {\textbf{temporalmente}:
    quando inizia e quando finisce la partita?};
  \node[etichetta, anchor=north west] at (8.6,-1.5) {\textbf{socialmente}:
    passanti, pasticcieri, vicini di quartiere};
\end{tikzpicture}
\caption{La definizione ``da game studies'' citata da Beatrice: il gioco
pervasivo espande il proprio cerchio magico nello spazio, nel tempo e nella
società. La definizione più intuitiva dice la stessa cosa dal lato degli
oggetti: luoghi e oggetti della vita quotidiana diventano luoghi e oggetti di
gioco.}
\label{a2:fig:cerchio}
\end{figure}

% ---------------------------------------------------------------------
\section{Pervasivi con un'app: Pokémon Go, geocaching, GuruShot}
\label{a2:sec:pervasivi-app}

\subsection{Pokémon Go}

\slidefig{img/AGD/agd2-pokemon}{Pervasivi con un'app: Pokémon Go, il geocaching e GuruShots.}{slide \emph{Analog -- Case Studies}}{fig:agd2-pokemon}


Il primo esempio è scelto ``proprio per farvi comprendere'': \textbf{Pokémon
Go}, che conoscono tutti, o di cui almeno tutti hanno sentito il nome. Beatrice
lo spiega comunque, ``così ci allineiamo'' (scusandosi con i grandi fan per la
descrizione imprecisa). È un gioco in formato \textbf{app}. I Pokémon, ``questi
animaletti bizzarri'', nella narrativa dei Pokémon popolano il loro mondo; in
\emph{Pokémon Go} popolano \textbf{il nostro mondo reale}. Il gioco consiste nel
cercarli, catturarli, andare nelle palestre, e tutto questo si fa
\textbf{camminando nel mondo}: la stessa via in cui si va a fare colazione al
bar o che si percorre per andare in ufficio è la via in cui si può trovare un
Pokémon o una palestra. Ed è ``tutto costruito dagli utenti''.

Su \emph{Pokémon Go} ci sarebbe molto da dire, ``anche su come è stato
utilizzato per fare altre cose meno belle ed eleganti dell'esperienza di
gioco'', ma Beatrice non ne parla per ragioni di tempo: ci tiene solo a far
notare che ``tante cose possono avere tante sfumature'', e che i casi più
spinosi si possono discutere a parte. Qui serve per una ragione semplice:
\textbf{è pervasivo perché usa gli stessi spazi della vita quotidiana come spazi
di gioco}.

\subsection{Il geocaching}

Il \textbf{geocaching} è anch'esso abbastanza conosciuto. ``Di base è una
gigantesca caccia al tesoro'', tenuta in piedi da alcuni portali, e nata
proprio con il \textbf{GPS} e la geolocalizzazione. Il principio:
\begin{itemize}
  \item le persone giocanti \textbf{nascondono} delle scatole, piccole o
        grandi, con contenuti vari (dei ``tesori'') e con una sorta di
        \textbf{diario di bordo} (un \emph{log}: fogli o quaderni) su cui chi le
        trova può scrivere ``sono stato qui, ho trovato questa scatoletta'';
  \item la scatola viene nascosta nel posto che l'utente ritiene più idoneo, e
        sulla piattaforma l'utente dà le \textbf{indicazioni} per trovarla;
  \item la grande regola è che \textbf{le scatole non vengono mai tolte}: chi le
        trova può scambiare gli oggetti al loro interno (piccoli, di qualunque
        valore, con piena libertà) e segnare la propria presenza.
\end{itemize}
Non è ``un arraffare i tesori'', ma un modo di \textbf{viverli e dar loro
spazio collettivamente}. In Italia, ``per la cultura che abbiamo'', funziona
così così, perché le scatole spesso vengono rubate o spostate. Ma vale la pena
guardare online i video dei modi in cui vengono nascoste: alcune sono ``piccole
opere di ingegneria artistica'', quasi piccole escape room, con un gioco di pesi
o una chiave ``nascosta in bella vista'' che si vede solo con grande
attenzione; alcune non sono scatole ma \textbf{cabine telefoniche} in cui si
entra e succedono cose, o intere automobili.

\subsection{GuruShot}

\textbf{GuruShot}\footnote{Più avanti la trascrizione lo chiama anche ``blue
shot''.} è di nuovo in formato app, ma Beatrice lo inserisce perché usa uno
strumento che lei usa ``moltissimo'' e che trova ``molto potente per fare molte
cose'': \textbf{la fotografia}. È un gioco dedicato a persone molto dentro al
tema, fotografi anche di un certo livello di passione e di competenza: una
\textbf{gara} o maratona fotografica in cui vengono date delle \textbf{challenge
fotografiche}. Chi vuole giocare apre l'app, sceglie la sfida, la svolge, carica
la foto; e poi ``succedono cose'': valutazioni, momenti di incontro con la
community, premi veri e propri offerti dagli \textbf{sponsor}. È un'attività
\textbf{continua}, online da molti anni (``almeno 7--8'', corregge Beatrice), con
una nicchia internazionale molto nutrita.

Ed è ``molto pervasivo'': \textbf{non si può dire quando una partita è iniziata
o finita}, né quando si sta giocando e quando no. Rispetto alla definizione del
cerchio magico, non si sa bene dove inizi e dove finisca l'attività di gioco:
per questo si parla di pervasività.

% ---------------------------------------------------------------------
\section{La città come campo di gioco: Journey, SF0 e Critical City}
\label{a2:sec:pervasivi-citta}

\subsection{Journey to the End of the Night}

\slidefig{img/AGD/agd2-journey}{La città come campo di gioco: Journey to the End of the Night e SF0.}{slide \emph{Analog -- Case Studies}}{fig:agd2-journey}


Si torna sotto il cappello dei \textbf{giochi urbani}, con un gioco che si
gioca in città, sempre \textbf{in formato evento}: un'ora di inizio, un'ora di
fine e una zona di gioco all'interno della città. È ``una sorta di \textbf{guardie
e ladri gigante}''. Lo scopo dei creatori era \textbf{riappropriarsi della vita
notturna} e viverla in modo ludico, esplorando la città ``in maniera
\emph{safe} da un lato e un po' \emph{noir} dall'altro''. Un piccolo gruppo di
giocatori fa i \textbf{chaser} (gli inseguitori, le guardie) e gli altri i
\textbf{runner} (i ladri); ci sono delle \textbf{tappe} da raggiungere, e il
gioco dura svariate ore. Non è un gioco prettamente sportivo (ci sono altri
pervasivi molto più sportivi): ci sono molte dinamiche \textbf{strategiche} su
come arrivare alla fine del percorso. Beatrice lo inserisce perché, rispetto ai
casi precedenti, è \textbf{completamente analogico}: non usa nessuna app, nulla.
È un evento.

Il gioco è \textbf{Journey to the End of the Night}, nato in America da un
collettivo con base a \textbf{San Francisco}. È arrivato in Italia nel
\textbf{2008}, durante un'ondata di giochi pervasivi arrivata ``in modo molto
forte'', e poi in Giappone e in tantissimi altri posti, perché i creatori lo
hanno reso \textbf{libero}: chi vuole può organizzarlo, purché \textbf{non metta
un biglietto a pagamento}. Ha avuto per anni un grande seguito: cercando
``Journey to the End'' Italia si trovano tantissimi eventi, e a \textbf{Milano},
tra il 2008 e il 2013, se ne sono fatte svariate edizioni con
\textbf{2\,000--3\,000 partecipanti} ciascuna. Beatrice lo racconta ``non tanto
per le meccaniche'', che più o meno si conoscono, ma per \textbf{come è
arrivato}: un'onda lunga che dura ancora, anche se si sta affievolendo.

\subsection{SF0}

Degli stessi creatori è \textbf{SF0}\footnote{La trascrizione dice ``SF0'',
``SFZero'' e, una volta, ``essere al vero''.}, che era una sorta di grande
\textbf{piattaforma online} con lo stesso principio di \emph{GuruShot}: delle
\textbf{missioni} da svolgere e documentare con video o foto. Ma invece di
puntare sull'estetica della fotografia, le missioni puntavano sulle
\textbf{interazioni cittadine e urbane}.

\subsection{Critical City Upload}

SF0 ha avuto ricadute anche in Italia, con ``un figlio, chiamiamolo così'':
\textbf{Critical City Upload}, un gioco che è andato molto tra il 2009 e il
2013. Lo scopo è vincere, accumulando punti con le missioni, ma il gioco è
molto legato agli \textbf{scambi umani e al territorio}. Le missioni andavano da
``organizza una cena con una persona del tuo quartiere con cui non hai mai
parlato'' a una delle preferite di Beatrice: ``crea un'opera d'arte interattiva
che intrattenga le persone che aspettano l'autobus e posizionala in una delle
pensiline del quartiere o della città''.

In entrambi i casi \textbf{è la community che valuta la community}: non ci sono
giudici. E dopo i primi anni, sia SF0 sia Critical City hanno permesso ai
partecipanti che avevano raggiunto un certo livello di \textbf{creare essi
stessi nuove missioni} per i nuovi giocatori: la comunità non si limitava a
valutarsi a vicenda, ma creava il gioco per gli altri. Oggi nessuno dei due è
più attivo, ma i siti sono ancora online, con molte delle missioni visibili.
Quanto a ``stranezza'' rispetto all'idea comune di gioco analogico, ``questo
spinge molto forte''.

% ---------------------------------------------------------------------
\section{Teatro, ARG e marketing: da Sleep No More a The Institute}
\label{a2:sec:pervasivi-arg}

\subsection{Sleep No More}

\slidefig{img/AGD/agd2-sleep}{Teatro e ARG: Sleep No More e Why So Serious?}{slide \emph{Analog -- Case Studies}}{fig:agd2-sleep}


Chi conosce il \textbf{teatro interattivo} conosce di sicuro \textbf{Sleep No
More}\footnote{La trascrizione dice ``Slip no more'' e ``SlipMod''. Lo
spettacolo è della compagnia britannica Punchdrunk; la versione più nota è
quella di New York, allestita in un edificio trasformato nel fittizio McKittrick
Hotel. Il nome della compagnia non è detto a lezione.}: un'opera messa in scena
all'interno di un \textbf{hotel}. È interattiva nel senso che gli attori recitano
la loro parte in interazione con il pubblico, e il pubblico, con regole precise
sui contatti fisici e su come interagire con gli spazi, è \textbf{libero di
esplorare l'hotel} come vuole: in ogni zona ci sono attori che fanno cose. Resta
un'opera teatrale, con un inizio, una fine e degli eventi nel mezzo; c'è molta
documentazione online (interviste, spezzoni). E di nuovo si gioca su
\textbf{qual è lo spazio di gioco e quale non lo è}.

\subsection{Gli alternate reality game e Why So Serious}

Qui Beatrice introduce un nuovo termine: gli \textbf{alternate reality game}
(ARG), i giochi di realtà alternativa, ``giochi in cui c'è un intero mondo
alternativo al nostro, in cui succedono cose''. Il caso scelto, ``uno dei miei
preferiti'', porta anche nel mondo dei \textbf{giochi per il marketing e la
comunicazione}: \textbf{Why So Serious}, nato come gioco promozionale per
\emph{Il cavaliere oscuro} di Christopher Nolan.
\begin{itemize}
  \item Iniziava con una \textbf{mail} ad alcune persone iscritte a mailing
        list di fan di Batman, contattate dagli organizzatori tramite le loro
        associazioni.
  \item La mail era un \textbf{invito strano a entrare nella banda del Joker},
        con un enigma; da lì partivano attività di socializzazione organizzate
        attraverso un sito, che era una serie di indirizzi ``in cui succedevano
        cose''.
  \item Il gioco raccontava che \textbf{nel nostro mondo il Joker esisteva
        davvero}, e gli indirizzi erano indirizzi reali.
  \item Una delle prime cose: andare a questi indirizzi, tutti
        \textbf{pasticcerie}, e dire la parola in codice ricevuta. Si riceveva
        una \textbf{torta} con sopra scritto un numero di telefono; chiamandolo,
        squillava la torta stessa, perché dentro c'era un \textbf{cellulare}:
        lo strumento con cui il Joker gestiva le comunicazioni con i giocatori.
\end{itemize}
Molte strutture di questo tipo, molti colpi di scena ``anche proprio pratici'',
che sfruttano la pervasività: cose che nel nostro mondo non hanno motivo di
succedere (``non esiste che ci sia un cellulare in una torta'') ma che in quel
mondo hanno un senso. Gli ARG giocano moltissimo su ciò che \textbf{potrebbe
essere reale ma non lo è}, e che il gioco dice che lo è: ``in quel mondo voi
siete gli aiutanti del Joker''. \emph{Why So Serious} è opera di un'azienda che
si occupa da molti anni di ARG\footnote{Beatrice non fa il nome dell'azienda,
che si troverà nell'estratto distribuito. Si tratta di 42 Entertainment
(integrazione).}; i dettagli sono nell'estratto che Beatrice caricherà, tratto
da \emph{Immersi nelle storie} di \textbf{Frank Rose}\footnote{Edizione
italiana di \emph{The Art of Immersion} (2011).}, che racconta proprio l'inizio
di questo gioco.

Perché raccontarlo? Perché mostra ``come possiamo giocare sulle
\textbf{emozioni}'' e come \textbf{connettere le persone alle cose che vogliamo
comunicare}, anche costruendo un mondo alternativo. Non è un obbligo, ``ma è bene
sapere'' che esistono questi giochi, ``giganteschi''.

\subsection{Un gioco da tavolo che è anche un ARG}

Tornando alla malleabilità delle definizioni, Beatrice presenta un gioco che è
insieme \textbf{gioco da tavolo}, \textbf{alternate reality game} e gioco con
\textbf{componenti digitali}, edito da \textbf{MS Edizioni}\footnote{La
trascrizione chiama il gioco ``The Ice''; il titolo non è ricostruibile con
certezza (potrebbe trattarsi di un titolo come \emph{The Heist}). MS Edizioni
è un editore italiano di giochi; è forse la stessa ``MS'' per cui, nella prima
lezione, Ruggero dice di lavorare, ma la trascrizione non permette di
confermarlo.}. Si riceve un \textbf{fascicolo} di informazioni e poi dei
\textbf{messaggi}, da cui si capisce di far parte di un gruppo che \textbf{sta
rapinando una banca}, già nel pieno della rapina. I giocatori sono ``l'entità che
ha accesso a internet'', la voce fuori dall'azione che coordina: parlano con dei
personaggi, scambiano informazioni, guardano video, analizzano dati e
documenti. È un gioco ``con una profondità notevole''. Due cose lo rendono
interessante:
\begin{itemize}
  \item mescola cose che \textbf{esistono davvero} (la cartina di una città
        reale, gli scambi della metropolitana di una città reale) con cose del
        \textbf{mondo narrativo} (i personaggi, i fatti, la rapina: ``spero che
        sia chiaro che non state rapinando la banca davvero'');
  \item a differenza degli altri casi, ha un \textbf{inizio e una fine ben
        definiti} ed è \textbf{completamente accessibile}: lo si può provare.
\end{itemize}
Per Beatrice è un caso ``unico sul panorama italiano'' (``poi mi correggeranno i
colleghi a livello internazionale'').

\subsection{Collettivo Cinetico: Cinetico 4.4 e Amleto}

\slidefig{img/AGD/agd2-teatro}{Cinetico 4.4, Lettere anonime per un camminatore, The MP3 Experiment, The Institute e l'ARG di Stranger Things.}{slide \emph{Analog -- Case Studies}}{fig:agd2-teatro}


\textbf{Cinetico 4.4} è ``un altro bell'esperimento'': un \textbf{gioco da
tavolo creato da un collettivo di danzatori e danzatrici}, il \textbf{Collettivo
Cinetico}, per generare una \textbf{performance artistica dal vivo}, teatrale o
di danza. Non è un gioco da prendere dallo scaffale di una ludoteca: è pensato
per una performance. Chi ci gioca partecipa a un gioco da tavolo che in realtà è
un \textbf{evento}: si paga un biglietto, si entra, ci si siede, succedono cose.
Il collettivo, di cui Beatrice si dichiara ``di parte'', fa moltissime attività
``né danza, né teatro, né gioco, o meglio tutte e tre le cose insieme'', difficili
da mettere sotto un cappello. E qui Beatrice aggancia un tema delle prossime
lezioni: la \textbf{transmedialità} e la \textbf{ricerca e sviluppo}, cioè
attingere anche a cose lontane dal gioco che si vuole creare, perché ``più
conosciamo, più possiamo aprirci la testa''.

Un caso studio in più: l'\textbf{Amleto} del Collettivo Cinetico. È l'Amleto di
Shakespeare, ma sul palco non ci sono i danzatori del collettivo. Tre giorni
prima dello spettacolo fanno una vera e propria \textbf{audizione} aperta alla
popolazione; tra chi si presenta scelgono \textbf{sei potenziali Amleto}, che
portano sul palco a danzare e recitare, e creano momenti di
\textbf{interazione con il pubblico} per scegliere, tra i sei, l'Amleto di
quell'edizione. Ci sono giochi veri e propri sia nell'audizione sia nello
spettacolo. Loro, da danzatori, lo chiamano spettacolo di danza, ``ma se lo
vedete c'è tantissimo''.

\subsection{Lettere anonime per un camminatore}

``Rispetto a quello che vi ho detto, questo è semplice''. \textbf{Lettere
anonime per un camminatore} è del collettivo \textbf{Amigdala}, anche loro una
compagnia teatrale (Gianmarco, che ha vissuto a Modena, conferma che sono di
lì). È una \textbf{passeggiata} con piccole interazioni ludiche, tutta basata su
un \textbf{audio}:
\begin{itemize}
  \item è un gioco \textbf{in solitario}: una persona sola segue un percorso
        segnato da frecce, per terra o sui muri di una città, guidata da una
        \textbf{voce in cuffia} che le racconta delle cose;
  \item la struttura del racconto è fissa, ma la \textbf{parte centrale} cambia
        ogni volta: prima di ogni edizione il collettivo \textbf{intervista le
        persone che hanno abitato o abitano il percorso}, e mentre si cammina si
        ascolta, oltre alla narrazione, il racconto di chi quei luoghi li vive o
        li ha vissuti;
  \item la \textbf{conclusione} è sempre ``col botto'': il collettivo sceglie di
        finire in modo scenografico in luoghi normalmente inaccessibili.
        Beatrice lo ha giocato a Ferrara, e alla fine del percorso si è
        ritrovata seduta in poltrona, con un tavolino e un drink, al
        \textbf{centro dello stadio comunale}, con lo stadio vuoto e le cuffie
        nelle orecchie.
\end{itemize}
Torna il tema del \textbf{percorso emotivo} che si fa fare alla persona giocante.
E lo strumento dell'\textbf{audio che guida} è interessante anche per il corso:
un'app non si riesce a fare nel tempo del project work, ``ma per utilizzare
l'audio invece possiamo''.

\subsection{MP3 Experiment}

Un altro esempio basato sull'audio, più simile a un \textbf{flash mob}:
\textbf{MP3 Experiment}, del collettivo newyorkese \textbf{Improv
Everywhere}\footnote{La trascrizione dice ``Improve Everywhere''.}, che fa
moltissimi esperimenti di questo tipo. Ci si ritrova tutti nel giorno, nell'ora
e nel posto concordati, ci si collega con le cuffie a una trasmissione in
diretta, e la voce dice che cosa fare. Beatrice non svela di più (``guardatevi i
video su YouTube, ci sono delle perle''), ma racconta un dettaglio: gli
organizzatori \textbf{elencano gli oggetti da portare}, e in un'edizione chiesero
di portare \textbf{due matite}; a un certo punto tutto il gruppo si mise a
suonare la batteria con le matite sugli oggetti della città. Quella
\textbf{scintilla di curiosità} per oggetti di cui non si sa ancora l'uso è uno
strumento del percorso emotivo, ``non difficile'' da usare: ``basta solo
ragionare di creatività''. Come le istruzioni audio.

\subsection{The Institute e Dispatches from Elsewhere}

\textbf{The Institute}\footnote{Il gioco è noto anche come \emph{The Jejune
Institute}, di San Francisco (integrazione).} è, nel modo in cui Beatrice ha
imparato questo mondo, uno dei casi più blasonati e famosi di alternate reality
game, uno dei primi raccontati come tale. È un'esperienza così complessa che
Beatrice preferisce indicare due fonti per approfondirla:
\begin{itemize}
  \item un \textbf{documentario}, intitolato anch'esso \emph{The Institute}, che
        racconta tutto l'ARG, durato circa un anno e fatto di molti passaggi, con
        tantissimi colpi di scena del tipo ``torta con il cellulare'';
  \item una \textbf{serie TV}, \emph{Dispatches from Elsewhere} (Daniele
        conferma il titolo), che ha preso l'esperienza di The Institute, ci ha
        aggiunto una storia dei personaggi e ne ha fatto una serie. Tolta la
        storia dei personaggi (``comunque molto carina''), resta più o meno
        l'idea dell'ARG.
\end{itemize}

\paragraph{Con budget zero.} Qui Beatrice fa una precisazione importante. Questi
esempi sono \textbf{difficili da replicare}: ``con che soldi compriamo la torta e
ci mettiamo il cellulare? Con che soldi affittiamo un intero palazzo come in The
Institute?''. Serve un intero lato commerciale, di ricerca di sponsor, che non è
parte del design del gioco e quindi non è il motivo per cui si è nel corso.
Quello che vuole che gli studenti si portino a casa è che \textbf{con molto meno}
si possono ottenere esperienze ``estremamente significative a livello
emozionale'' e creazioni di game design molto profonde \textbf{con budget zero}.
I casi grandi servono come \textbf{punti di riferimento} e per sapere che
esistono cose ``molto distanti da quelle che immaginiamo tutti i giorni''.

\subsection{L'ARG di Stranger Things a Milano}

Nello stesso filone, un ARG che ``è stato un'operazione di marketing e basta''
(``per me sono molto affascinanti, però era un'operazione di marketing''), ma
fatto \textbf{a Milano}: quello per una stagione di \emph{Stranger Things} (la
terza o la quarta, Beatrice non ricorda). Era una sorta di \textbf{caccia al
tesoro} in città, con punti di gioco preparati in modo ``altamente scenografico''
(nella slide si vede un autobus), mescolata con una \textbf{telefonata} a un
numero in cui uno dei personaggi diceva che il Sottosopra era arrivato a Milano e
che serviva aiuto. Il messaggio: anche da noi questi giochi ``sono arrivati,
arrivano, e credo che col tempo avranno sempre più spazio'', e per chi pensa in
termini di lavoro ``è un tema''.

% ---------------------------------------------------------------------
\section{Cacce al tesoro, corse e indagini}
\label{a2:sec:pervasivi-cacce}

\subsection{Le guide-caccia al tesoro su carta}

Beatrice presenta poi delle \textbf{guide per esplorare la città} in forma di
caccia al tesoro\footnote{La trascrizione riporta il nome come ``Waiwai''; la
grafia non è verificabile.}, uscite negli stessi anni dell'ondata dei pervasivi
in Italia (la prima nel 2007--2008), disponibili solo per alcune città e
vendute come \textbf{libretti} sul loro sito o da rivenditori autorizzati.
Tiziana le conosce: aveva comprato quelle di \textbf{Venezia} e di
\textbf{Verona}, ha giocato in parte quella di Venezia, nella versione breve
(c'erano due possibilità: un gioco di un paio d'ore o uno da una giornata
intera), e ricorda il \textbf{numero di cellulare} con cui interagire per avere
le indicazioni. Beatrice le aveva giocate con gli SMS, ``era un'altra vita'';
oggi pare si siano aggiornate a WhatsApp (e chiede aggiornamenti a chi le
giocherà).

Sono un altro ibrido di \textbf{analogico e digitale}: un'entità parla con i
giocatori tramite messaggi, mentre il libretto cartaceo resta in mano a loro. La
particolarità è la meccanica per passare da una tappa all'altra: \textbf{ogni
pagina del libro è divisa in strisce}, e quando si risolve una tappa la guida
comunica dei \textbf{codici} (``il foglietto 3, il 32C e il 18F'') con cui
ricomporre la pagina che porta alla tappa successiva.

\subsection{PacManhattan}

\slidefig{img/AGD/agd2-cacce}{PacManhattan, Zombies, Run! e Reptilarium.}{slide \emph{Analog -- Case Studies}}{fig:agd2-cacce}


\textbf{PacManhattan}\footnote{La trascrizione dice ``Pac-Man-Hatton''.}, degli
stessi anni, è \textbf{Pac-Man per le strade di Manhattan}: una persona fa
letteralmente Pac-Man, guidata da un'altra persona (``la mente'') tramite
telefono o walkie-talkie; altri giocatori fanno i fantasmi. È di nuovo ``una
sorta di guardie e ladri'', ed è uno dei primi urban game che si trovano nei
manuali di game studies. Beatrice lo mostra per un'idea precisa: si può
\textbf{prendere qualcosa che esiste già in una forma e cambiarle forma
completamente}. Non è lo stesso gioco di Pac-Man, ma si ispira alla sua
narrativa, alle sue meccaniche, alla sua estetica.

\subsection{Zombies, Run!}

\textbf{Zombies, Run!}\footnote{La trascrizione dice ``Zombie Run''.} è
un'\textbf{app per fare jogging}: invece di ascoltare musica, o i messaggi
vocali degli amici, si ascoltano \textbf{gli zombie che stanno per prenderti}, e
c'è tutta una narrativa su come si sta scappando mentre si corre. Si trova sugli
store.

\subsection{Reptilarium}

Si torna al teatro, ma in versione analogico-digitale. \textbf{Reptilarium} è un
caso studio di cui purtroppo c'è pochissima documentazione. È uscito durante il
\textbf{COVID}, da una compagnia teatrale di Roma, \textbf{Project
XX1}\footnote{Grafia come nella trascrizione.}. Era a modalità evento (un certo
giorno, a una certa ora) e le persone accedevano da \textbf{casa propria}, davanti
al computer, a un'esperienza con degli \textbf{attori} che si mettevano in
comunicazione con loro: ricevevano mail e telefonate, potevano vedere video
``recuperati da altre persone'', potevano \textbf{coordinare i movimenti degli
attori} e vedere che cosa succedeva mentre li eseguivano. La trama partiva da una
\textbf{persona scomparsa} da ritrovare, e poi si espandeva molto. Come per le
cene con delitto (dove si può interrogare l'attore o essere l'indiziato), anche
questo esperimento ``può diventare un LARP'', un ibrido tra teatro e gioco di
ruolo, o tra gioco investigativo e alternate reality game.

Ruggero aggiunge una nota: il \textbf{lockdown} ha colpito moltissimo i giochi
pervasivi, ma anche altri generi come i LARP, e in quel periodo sono nate
esperienze che hanno spostato \textbf{online e a distanza} giochi nati per essere
giocati di persona. Se ne riparlerà.

% ---------------------------------------------------------------------
\section{Pervasivi per la formazione: Echi resistenti}
\label{a2:sec:pervasivi-edu}

\slidefig{img/AGD/agd2-echi}{Un gioco pervasivo per la formazione storica.}{slide \emph{Analog -- Case Studies}}{fig:agd2-echi}


``Tutte le cose che vedete hanno enorme spazio per essere esplorate in tema
\textbf{educativo e formativo}'', e Beatrice chiude con un accenno a questo
ambito. \textbf{Pop History} è un'\textbf{associazione di promozione sociale}:
non si presenta come collettivo teatrale né come gruppo di game designer, ma come
associazione (``credo siano anche loro emiliani''), e fa tante piccole
sperimentazioni. Una di queste è \textbf{Echi resistenti}\footnote{La
trascrizione dice ``Ecchi Resistenti'' e ``Eh, che resistenti''.}, un gioco
sulla \textbf{Resistenza e sulla Liberazione}, che ibridava:
\begin{itemize}
  \item un'\textbf{indagine} con dei fascicoli;
  \item un'\textbf{attività teatrale}, con personaggi storici con cui si poteva
        interagire direttamente;
  \item una \textbf{caccia al tesoro}, in cui si andava a cercare questi
        personaggi per la città.
\end{itemize}
Altri casi simili sono sul loro sito.

\paragraph{Approfondimento: giochi con migliaia di giocatori.} In
\emph{Challenges for Game Designers}, nel capitolo sui giochi multigiocatore,
Brathwaite e Schreiber osservano che nei giochi da tavolo il problema
dell'interfaccia è evidentissimo (``come si mettono venti persone attorno a un
solo tabellone? E mille?''), eppure ``esistono giochi non digitali con migliaia
e perfino centinaia di migliaia di giocatori'': gli \textbf{alternate reality
game} e i grandi \textbf{urban game}, che ``portano i giocatori via dai loro
computer e dentro il mondo reale''. È la stessa scala dei casi di Beatrice:
\emph{Journey to the End of the Night} con migliaia di partecipanti a Milano,
\emph{Why So Serious} con le sue mailing list di fan, \emph{Pokémon Go} con il
mondo intero come mappa.

\begin{table}[htbp]
\centering
\small
\begin{tabular}{@{}p{3.6cm}p{2.3cm}p{7.0cm}@{}}
\toprule
\textbf{Caso studio} & \textbf{Formato} & \textbf{Perché è nella lezione} \\
\midrule
Pokémon Go & app & gli spazi quotidiani diventano spazi di gioco \\
Geocaching & piattaforma + scatole & caccia al tesoro collettiva: le scatole
non si tolgono mai \\
GuruShot & app & la fotografia come strumento; non si sa quando la partita
inizi e finisca \\
Journey to the End of the Night & evento urbano & guardie e ladri notturno,
completamente analogico, libero e non a pagamento \\
SF0 & piattaforma & missioni urbane documentate; la community valuta e crea
missioni \\
Critical City Upload & piattaforma & il ``figlio'' italiano di SF0
(2009--2013): missioni di relazione sul territorio \\
Sleep No More & teatro interattivo & il pubblico esplora liberamente un hotel
\\
Why So Serious & ARG di marketing & la torta con il cellulare: il Joker esiste
davvero \\
Gioco di MS Edizioni & gioco da tavolo + ARG & rapina in corso tra documenti
reali e finzione; ha inizio e fine \\
Cinetico 4.4, Amleto & performance & un gioco da tavolo che è uno spettacolo; il
pubblico sceglie Amleto \\
Lettere anonime per un camminatore & passeggiata audio & gioco in solitario,
voci degli abitanti, finale scenografico \\
MP3 Experiment & flash mob audio & la curiosità per gli oggetti da portare \\
The Institute & ARG & il caso blasonato, con documentario e serie TV \\
ARG di Stranger Things & ARG di marketing & un ARG fatto a Milano \\
Guide-caccia al tesoro & libretto + SMS & pagine in strisce da ricomporre con
codici \\
PacManhattan & urban game & un videogioco cambiato completamente di forma \\
Zombies, Run! & app & narrativa audio durante la corsa \\
Reptilarium & teatro a distanza & attori guidati da casa durante il COVID \\
Echi resistenti & gioco formativo & indagine, teatro e caccia al tesoro sulla
Resistenza \\
\bottomrule
\end{tabular}
\caption{I casi studio di giochi pervasivi presentati da Beatrice, nell'ordine
della lezione, con il motivo per cui sono stati scelti.}
\label{a2:tab:pervasivi}
\end{table}

% ---------------------------------------------------------------------
\section{I giochi di ruolo dal vivo: le macrocategorie}
\label{a2:sec:larp}

\slidefig{img/AGD/agd2-larp}{Le macrocategorie del LARP: boffer, Mind's Eye Theatre, parlor/chamber LARP, American freeform e murder party.}{slide \emph{Analog -- Case Studies}}{fig:agd2-larp}


Finita la parte sui pervasivi, si fa una pausa di pochi minuti (``facciamo
sette minuti, così andiamo ai 20 tondi''). Al rientro ci sono tutti, compreso
Tancredi (``ci sono tutti e due'' i Lorenzo), e Ruggero passa ai
\textbf{giochi di ruolo dal vivo}.

\paragraph{Una struttura diversa.} Questa parte è organizzata diversamente da
quella sui pervasivi. Il gioco di ruolo dal vivo, come il gioco di ruolo da
tavolo e il gioco da tavolo, ha alle spalle \textbf{anni di industria
solidificata}, ed è quindi più facile identificare \textbf{macrocategorie} con
nomi e formati specifici. Prima, però, Ruggero chiede chi abbia mai giocato un
LARP. Una persona della classe sì: un LARP fantasy, al \textbf{festival del gioco
di ruolo di Bracciano}, ``con le spade di gommapiuma, insomma, quelle cose lì''.
È il primo impatto di molti.

\subsection{Il boffer LARP}

Chi ha giocato un LARP probabilmente ha giocato un \textbf{boffer}, e chi ne ha
solo sentito parlare ha in mente quello. È uno dei primi giochi di ruolo dal
vivo diffusi, nato all'epoca proprio come \textbf{tentativo di adattare
\emph{Dungeons \& Dragons} dal vivo}, poi evoluto nel tempo. (Daniele mostra
una foto e Ruggero protesta: ``non evidenziarmi, che già è cringe''; ne ha
giocati anche lui.) Il nome viene dalle armi: le armi \emph{boffer} sono armi in
plastazote che devono avere caratteristiche precise per non essere pericolose.
Sono giochi \textbf{incentrati sul combattimento}, su trame epiche e così via.
Sono anche quelli che possono essere più grandi: un boffer anche grosso può
ospitare \textbf{400--500 persone}, ``per non dire quelli in Germania, che ne
tengono anche migliaia''.

\subsection{Il Mind's Eye Theatre}

Un altro tipo abbastanza diffuso è il \textbf{Mind's Eye Theatre}: se il boffer
deriva da \emph{Dungeons \& Dragons}, questo deriva da \textbf{\emph{Vampire: The
Masquerade}}. Le trame sono più dark, gotiche, horror, legate alla proprietà
intellettuale di Vampire, e chi partecipa di solito interpreta uno dei vampiri
delle loro società segrete. Sono giochi \textbf{più politici che di
combattimento}: l'intrigo, i segreti, il tradirsi a vicenda, ``il pugnalare
alle spalle qualcuno'' sono parte delle dinamiche. Come i boffer, sono giochi
\textbf{in costume}: il costume è richiesto, e si chiede a chi gioca di vestirsi
come il proprio personaggio.

\subsection{Chamber LARP e parlor LARP}

Il costume sembra scontato, ``ma in realtà non tutti i LARP usano il costume
come strumento necessario''. I \textbf{parlor LARP} o \textbf{chamber LARP}
sono giochi di ruolo dal vivo in cui il costume \textbf{non è richiesto}: si
gioca con i propri vestiti di tutti i giorni. Lo studente di prima interviene: a
Bracciano, in effetti, giocavano con i propri vestiti (``di solito gli eventi
così sono degli ibridi''); e ricorda di aver giocato un altro LARP al
\textbf{Play}\footnote{La trascrizione dice ``Play di Bologna''; il festival
Play, dedicato al gioco, si tiene a Modena.}, attorno a un tavolo, ognuno con un
ruolo e con i propri vestiti: \textbf{\emph{A cena col nonno}}. ``Perfetto, quello
è un chamber LARP''.

Ruggero precisa un punto sottile: alcuni chamber LARP si giocano \textbf{a un
tavolo}, ma questo non li rende giochi di ruolo da tavolo, perché si giocano
personaggi \emph{che sono seduti a un tavolo}, mentre in \emph{Dungeons \&
Dragons} il personaggio non è necessariamente seduto a un tavolo in quel momento.
Con gli altri LARP hanno in comune la \textbf{sovrapposizione tra la fisicità di
chi gioca e quella del personaggio}, l'\textbf{embodiment}. Rispetto ad altri
LARP, i chamber LARP hanno di solito:
\begin{itemize}
  \item un livello di \textbf{accessibilità} più alto, perché spesso non
        richiedono costumi;
  \item un \textbf{budget ridotto};
  \item \textbf{non sono a evento}: molti si scaricano in \emph{print and play},
        si stampano a casa e si giocano con gli amici;
  \item per questo, un \textbf{numero di giocatori} più ridotto, di solito da
        3 a 20 persone.
\end{itemize}

\subsection{L'American freeform}

Una categoria dei chamber LARP è l'\textbf{American freeform}: LARP progettati da
designer che si occupano soprattutto di giochi di ruolo da tavolo, e infatti è
un \textbf{genere ibrido}. Visti da fuori sono chamber LARP (gente che fa
\emph{embodiment} del personaggio in delle stanze, non necessariamente in
costume), ma hanno ``un sacco di meccaniche'' prese dal gioco da tavolo e dal
gioco di ruolo da tavolo: le cosiddette \textbf{metatecniche} o metameccaniche,
di cui si parlerà in un'altra lezione. Anche questi sono spesso in \emph{print
and play}; alcuni sono venduti in un formato più vicino al gioco da tavolo, come
mazzi di carte o libretti. Quindi non sono a evento, a differenza dei boffer,
del Mind's Eye Theatre e dei murder party.

\subsection{I murder party}

Ruggero riprende l'ottima distinzione fatta da Beatrice. ``Murder party'', da noi,
indica culturalmente due cose:
\begin{itemize}
  \item le \textbf{cene con delitto}, in cui si va al ristorante, degli attori
        mettono in scena un omicidio e la successiva indagine, e i commensali
        mangiano e intanto risolvono il caso. \textbf{Nessuno dei commensali è
        l'assassino}: è uno degli attori;
  \item la categoria ``più british'', i \textbf{murder party all'inglese}, o i
        loro equivalenti che si stanno diffondendo moltissimo in Asia in questo
        periodo\footnote{La trascrizione dice ``Jubel Shock'' e ``Juvenile
        Show''. Si tratta quasi certamente dei \emph{jubensha}, i giochi di
        investigazione con ruoli molto diffusi in Cina.}, in cui \textbf{si
        interpretano i personaggi della storia} su cui si indaga: tra i giocatori
        può esserci l'assassino.
\end{itemize}
La seconda forma rientra, per il design dell'esperienza e delle meccaniche,
sotto l'ombrello dei LARP; gli aspetti ambientali e le decorazioni la
avvicinano poi al mondo del pervasivo.

% ---------------------------------------------------------------------
\section{Il Nordic LARP}
\label{a2:sec:nordic}

\slidefig{img/AGD/agd2-nordic}{Il Nordic LARP e i Knutebooks, i volumi collettivi del Knutepunkt.}{slide \emph{Analog -- Case Studies}}{fig:agd2-nordic}


C'è poi ``una grossissima categoria'', che si rivedrà più avanti perché ``è stata
importantissima'': il \textbf{Nordic LARP}, il LARP progettato in
\textbf{Scandinavia} (Danimarca, Norvegia, Svezia e Finlandia).

\subsection{Perché è così importante}

Ha un'importanza gigantesca nel settore \textbf{perché ha documentato tutto}. La
maggior parte dei libri accademici di game design e di analisi di questo tipo di
giochi, gli eventi e le conferenze, si fanno in Scandinavia, e c'è una lunga
storia di analisi e documentazione: la prima conferenza è del \textbf{1997}, e da
allora è annuale. È ``gente che si incontra da parecchi anni a parlare di come si
progetta e come si gioca''. Ed è una scena molto \textbf{cross-mediale}: prende
dagli alternate reality game, dai LARP americani, dai parchi a tema, dalle
simulazioni economiche, da tante teorie parallele a quelle del gioco, cercando
di interiorizzare, iterare e integrare tecniche di altri settori. Anche il Nordic
LARP ha le sue sottocategorie: alcune esperienze sono più vicine al chamber
LARP, altre al boffer.

\subsection{Come si riconosce}

Da fuori, un Nordic LARP si riconosce da \textbf{come è progettato il design
dell'esperienza}:
\begin{itemize}
  \item un focus sul \textbf{coinvolgimento}, sull'\textbf{emozione} e sulla
        \textbf{coesione tematica};
  \item pochissime cose \textbf{vestigiali}: per esempio, anche LARP molto
        grandi possono non avere il combattimento, se il combattimento non fa
        parte del tema;
  \item una cura del \textbf{post-esperienza}: molti giochi hanno una
        struttura per ``che cosa succede dopo il gioco''.
\end{itemize}
Molti concetti nati da questa scena si sono ``solidificati nel gergo del game
design'' e torneranno nella lezione sull'emotional design.

\paragraph{Origini diverse da D\&D.} Storicamente, il LARP nordico ha
un'origine \textbf{precedente a \emph{Dungeons \& Dragons}}. Si tende a pensare
che il LARP nasca da ``prendiamo D\&D e lo facciamo dal vivo'', e per alcune
nazioni, come l'Inghilterra e l'America, è proprio così; per altre no. Il Nordic
LARP prende strategie e strumenti dal \textbf{teatro}, dalle avanguardie, dallo
\textbf{psicodramma}, dalla \textbf{Gestalt therapy}; ``è ovvio che quello che
esce fuori è un gioco molto particolare''.

\paragraph{I manifesti.} È anche un ambiente ``pieno, carico, zeppo di
manifesti''. La slide ne mostra alcuni, nati negli anni nella scena; non si
approfondiscono ora, ma sono tutti pubblici e si possono leggere a parte.

\subsection{Knutpunkt e i Knutepunkt books}

La cosa che interessa di più è la conferenza citata prima, che cambia nome a
seconda del paese che la ospita: \textbf{Knutpunkt}, \textbf{Knudepunkt},
\textbf{Knutepunkt} o \textbf{Solmukohta}\footnote{Rispettivamente in svedese,
danese, norvegese e finlandese: la conferenza ruota tra i quattro paesi. La
trascrizione dice ``Knutpunkt o Knutpunkt o Knudpunkt o Solmukott'' e, più
avanti, ``Knut Punk''.}. È il principale punto di ritrovo scandinavo sul gioco
di ruolo dal vivo: \textbf{quattro giorni di conferenza non stop}, con un
programma di centinaia di slot, spesso \textbf{otto track in parallelo} di talk,
panel e workshop. Si tiene dal 1997, e \textbf{ogni anno produce dei libri}
che riassumono lo stato dell'arte in forma di articoli: la documentazione è
``costantemente in aggiornamento''. Ruggero ne ha diversi, ``dei bei mattoni'':
ne mostra uno di \textbf{600 pagine} su come si progetta un LARP, uscito da una
delle edizioni (``nel 2018, se non ricordo male''). È forse ``una delle
comunità ludiche che più documenta quello che fa'', e soprattutto lo documenta
in modo \textbf{pragmatico}, lasciando framework e strumenti a chi vuole creare
giochi di questo tipo. Molte di queste cose si applicano anche ad altri generi, e
torneranno.

\subsection{Dogma 99}

In chat qualcuno chiede se c'entri il \textbf{Dogma} di Lars von Trier.
``Sì e no'': il \textbf{Dogma 99} è un manifesto di LARP che riprende il
manifesto cinematografico del \textbf{Dogma 95} per fare un parallelo con il
LARP. Il Dogma 95 aveva alcune direttive, per esempio che tutto debba essere
\textbf{diegetico}, senza nulla di extradiegetico; il Dogma 99 fa lo stesso: in
un LARP storico non ci devono essere cose extradiegetiche rispetto alla storia.
Una conseguenza: il Dogma 99 \textbf{non fa LARP fantasy}, perché la magia non
esiste e non si può riprodurre in modo perfetto; si fanno giochi ambientati
tendenzialmente ai giorni nostri, che si possono riprodurre facilmente. Da
questa catena di ragionamenti sono nati ``giochi molto interessanti''.

% ---------------------------------------------------------------------
\section{Jeepform: il dramma quotidiano}
\label{a2:sec:jeepform}

\slidefig{img/AGD/agd2-jeepform}{Jeepform tra tabletop, freeform e LARP; la locandina di Fat Man Down.}{slide \emph{Analog -- Case Studies}}{fig:agd2-jeepform}


Tra i manifesti c'era lo \textbf{Jeepform}\footnote{La trascrizione dice
``G-Form'' e ``Jeep Form''.}, nato all'interno del \textbf{freeform}, la zona ``a
metà tra il LARP e il tabletop''. Il suo manifesto, che si trova online, ha tra i
punti più importanti questo: \textbf{il dramma quotidiano è più interessante del
dramma epico}. Tutti i giochi Jeepform sono giochi di ruolo dal vivo su
\textbf{situazioni quotidiane}, ``drammaticamente quotidiane''.

Alcuni sono emotivamente molto pesanti, e storicamente importanti per vari
motivi; Ruggero ha lasciato fuori i più pesanti (se ne riparlerà nella lezione
sull'emotional design) e ne presenta due ``più leggeri'', che fanno comunque
capire in che toni si va.

\subsection{Il gioco della coppia perfetta}

Il primo\footnote{La trascrizione riporta il titolo come ``Lady 8'' e i nomi dei
personaggi come ``Lady'' e ``8''; né il titolo né i nomi sono verificabili.} è
un gioco per \textbf{quattro persone} e \textbf{due personaggi}, che per premessa
\textbf{si amano tantissimo}. Due persone alla volta interpretano i personaggi,
in un gioco \textbf{diviso in scene}; ogni scena è descritta dalle due persone che
non sono in gioco. L'unico obiettivo di chi gioca la coppia è \textbf{essere la
coppia perfetta}, mostrare quanto si amano. Se la scena si sposta su qualcos'altro,
\textbf{si rifà}, all'infinito, ``come nel giorno della marmotta'', finché non
parla solo di quanto si amano. Esempio: lei entra in bagno, la porta si rompe e
lei rimane chiusa dentro, lui è dall'altra parte; se lui comincia a riparare la
porta, si rifà la scena. ``Una paradossalità esasperata, quasi''.

\subsection{Fat Man Down}

Il secondo è introdotto con un \emph{content warning}: \textbf{Fat Man Down} ``è
un gioco molto pesante''. Non è il più pesante dei Jeepform, ma è il più pesante
di cui si parla nella lezione. È un LARP dello stile Jeepform, quindi un chamber
LARP tendenzialmente freeform.
\begin{itemize}
  \item Il personaggio principale è un uomo di successo, con un buon lavoro e
        una buona posizione economica, ma \textbf{obeso}.
  \item Tutti i partecipanti \textbf{si pesano su una bilancia}, e chi pesa di
        più interpreta questo personaggio.
  \item Il gioco è diviso in scene, ognuna con un \emph{framing} che è uno
        \emph{slice of life}: un'ora in ufficio, la strada da casa al
        supermercato, dieci minuti di pausa al parco.
  \item Tutte le altre persone in scena hanno la \textbf{direttiva meccanica} di
        farlo stare malissimo.
  \item Il protagonista ha una meccanica ``importantissima'': \textbf{tre
        cioccolatini} in tasca. In qualunque momento può scartarne uno e
        mangiarlo, e quando lo fa tutti i personaggi in scena devono trattarlo
        come la persona migliore del mondo.
\end{itemize}
Il personaggio è messo sotto \textbf{costante stress narrativo} finché non è
spinto a pensare ``devo mangiare per stare meglio''. ``La roba è molto pesante'',
e ci si ferma qui.

\subsection{Le meccaniche di sicurezza}

Avere un manifesto così noto nella scena nordica, con giochi molto pesanti, molto
intimi e soprattutto \textbf{senza un filtro fantasy o fantascientifico}, cioè
su drammi quotidiani, ha generato e normalizzato cose poi rivelatesi molto utili,
come le \textbf{meccaniche di sicurezza}. Se qualcuno ne ha sentito parlare, la
scena nordica è ``la principale responsabile'' della normalizzazione, della
ricerca e dell'iterazione delle meccaniche di sicurezza nel gioco di ruolo. Se ne
parlerà nella lezione sull'emotional design.

% ---------------------------------------------------------------------
\section{Il 360° e l'alta risoluzione}
\label{a2:sec:360}

\slidefig{img/AGD/agd2-360}{Il 360° e l'alta risoluzione: Ground Zero e Monitor Celestra.}{slide \emph{Analog -- Case Studies}}{fig:agd2-360}


Il Nordic LARP non è solo chamber, freeform e Jeepform. C'è anche ``quasi
l'opposto'': i LARP \textbf{360°} o ad \textbf{alta risoluzione} (\emph{high
resolution})\footnote{La trascrizione dice ``360 o High Resolution'' e, più
avanti, ``IRES 360'', ``iRes'', ``Ires''.}. Lo scopo è essere \textbf{il più
fisicamente coinvolgenti possibile}: tutto quello che si ha attorno deve essere
\textbf{diegetico}. Ma la soluzione non è quella del Dogma 99 (ambientare nel
quotidiano, così non serve la magia): ``la soluzione è buttare al muro un budget
talmente alto che la realtà a quel punto sembra magia''. Niente oggetti di scena
finti o armi finte: anche la scenografia è curata al punto che sembra di stare
davvero nel gioco. I casi studio sono ``anche abbastanza storici''.

\begin{description}
  \item[Ground Zero (1998--2001).] La premessa è la \textbf{Guerra Fredda}.
        Una cittadina americana riceve l'allarme per una bomba nucleare, e i
        partecipanti vivono \textbf{48 ore di fila}, non stop, chiusi in uno
        \textbf{scantinato} al buio, senza poter uscire. Il gioco si conclude
        con la bomba che cade: dietro le finte pareti di cartone dello
        scantinato c'erano dei \textbf{subwoofer giganti}, e alla fine i muri
        cadono addosso ai giocatori con un rumore devastante.
  \item[Panopticorp (2003, 2013).] Due \emph{run}, anche questo di 48 ore non
        stop: in teoria si dorme in gioco, ``ma in questo caso non si dormiva''.
        È una \textbf{distopia corporate} spinta al massimo: si è dipendenti di
        un'agenzia pubblicitaria sull'orlo della chiusura, che deve licenziare
        ``a manetta'' e tenere solo i migliori. Una spinta ultra competitiva e
        ``ultra tossica'' per vedere fin dove si è disposti a spingersi eticamente
        per restare in azienda, dormendo sotto le scrivanie.
  \item[Monitor Celestra.] Un 360° \textbf{fantascientifico} basato su
        \emph{Battlestar Galactica}: per giocarlo è stato noleggiato un
        \textbf{incrociatore svedese} in disarmo. Una nave, non un'astronave,
        ``ma se non usciamo all'esterno, ci mettiamo un paio di LED, ed ecco
        qua'': per i giocatori è un'astronave. Il gioco durava circa
        \textbf{60 ore}, con cabine allestite in modo che si dormisse in
        personaggio.
  \item[Demeter.] Un 360° sul viaggio della bara di \textbf{Dracula} dalla
        Transilvania a Londra, come nel romanzo di Stoker, in cui la nave
        \emph{Demeter} arriva in porto con tutti morti a bordo e il capitano
        legato al timone. I giocatori sono i passeggeri, e \textbf{da giocatori
        sanno che moriranno tutti}: per design è impossibile sopravvivere,
        ``altrimenti Dracula non esisterebbe come romanzo''. Si gioca in
        costume su una \textbf{nave vera} che, naturalmente, non va dalla
        Transilvania a Londra ma gira nel mare della Germania del nord;
        narrativamente, però, si gioca quel viaggio.
  \item[Bunker 101.] Un caso \textbf{italiano}, ispirato alla proprietà
        intellettuale di \emph{Fallout}: si giocava \textbf{72 ore} in un bunker.
        Qui il dormire non era diegetico: nell'ultima run hanno separato la parte
        di gioco dal sonno. Per tutto il resto, però, valeva il principio 360°:
        tutto ciò che si vede deve essere diegetico.
\end{description}

% ---------------------------------------------------------------------
\section{Tape LARP e black box}
\label{a2:sec:blackbox}

\slidefig{img/AGD/agd2-tape}{Tape LARP e black box LARP.}{slide \emph{Analog -- Case Studies}}{fig:agd2-tape}


Accanto al Jeepform (``in cui giochiamo in casa con i nostri costumi'') e al
360° (``l'ultra-diegesi''), ci sono tipologie scandinave ``molto più particolari
e quasi astratte''.

\subsection{Il tape LARP}

Visto che qualcuno conosce Lars von Trier: chi ha visto \emph{Dogville}, il film
con lo scotch per terra, sa da dove nascono i \textbf{tape LARP}. ``Abbiamo visto
Dogville, ci ispira un casino'': giochi in cui \textbf{lo scotch per terra indica
gli ambienti}, con tutte le meccaniche che ne derivano (il personaggio vede e
sente cose che normalmente non potrebbe sentire; come si agisce e non si agisce
da giocatori\ldots). Ne sono stati fatti parecchi, in particolare in Danimarca.

\subsection{Il black box LARP}

Ruggero è ``un particolare amante'' del \textbf{black box LARP}, che prende il
nome dalle \emph{black box}: per chi ha fatto teatro ``dalla Germania in su'',
spazi teatrali di solito senza spalti né sedie, con \textbf{pareti, soffitto e
pavimento neri}, molto diffusi in Germania e nel Nord Europa. A
\textbf{Copenaghen} c'è un festival, il \textbf{Black Box LARP Festival}, che usa
questi spazi per giochi creati apposta: giochi di ruolo che usano \textbf{luci e
colori come meccaniche di gioco}. Gli esempi della slide:

\begin{description}
  \item[Borderline.] Un gioco ``estremamente astratto''. Si gioca \textbf{in
        coppie}, e ogni coppia gioca \textbf{un solo personaggio}, con dei
        conflitti interni: tutta la conversazione è un monologo interiore del
        personaggio, e i due giocatori, con una candela, lottano per portare ciò
        che il personaggio fa dalla propria parte o da quella dell'altro.
  \item[Waiting for Flight.] Ruggero lo ha giocato. Si aspetta un aereo con a
        bordo persone care, amici, colleghi; \textbf{l'aereo precipita}, e
        muoiono quasi tutti. Si gioca la \textbf{sala d'attesa} di chi li
        aspettava: è un gioco sull'attesa, ``sul logorante non avere
        informazioni'' finché vengono centellinate. Si conclude con un momento
        molto catartico: arrivano le \textbf{valigie} delle persone attese, e in
        alcune (3 su 20) è nascosto un piccolo \textbf{cuore di vetro}: se lo
        trovi, la persona si è salvata; se no, non si è salvata. C'è una fase in
        cui speri che ci sia, ma potrebbe essere nascosto: ``è una cosa
        lacerante. Abbiamo pianto in 19 su 20''.
  \item[Stranger.] Di una designer danese\footnote{La trascrizione dice ``Nino
        Ressendrop'': con ogni probabilità Nina Runa Essendrop, designer danese
        nota per i LARP fisici e non verbali. Anche il titolo, ``Stranger'', è
        riportato come nella trascrizione.}. È un black box \textbf{non verbale}:
        è vietato parlare. I partecipanti sono divisi in due gruppi, i
        \textbf{blu} e i \textbf{verdi}, e prima del gioco ciascun gruppo crea un
        \textbf{linguaggio gestuale condiviso}: come la propria nazione esprime
        felicità, rabbia, paura, tristezza, amicizia, odio. Il twist: i verdi
        devono avere gesti squadrati, rigidi, geometrici; i blu gesti curvi,
        lenti, morbidi, sinuosi. Il gioco è sull'\textbf{immigrazione forzata}:
        ci sono fasi di giorno e di notte, e di notte chi facilita porta alcuni
        verdi nella nazione blu, dove si risveglieranno, con la regola che il
        confine non si può riattraversare. Le dinamiche: il primo che si
        risveglia tra i blu impara facilmente la loro lingua, e i blu lo aiutano
        con un \emph{onboarding} non verbale; ma quando arrivano molti verdi,
        tendono a muoversi \textbf{in gruppi}, perché a quel punto c'è già chi
        parla la propria lingua, anche se non è la lingua della nazione. ``Un
        microscopio'' di dinamiche della vita reale: emigrazione, immigrazione,
        inclusione, e quanto è più facile essere inclusi se il sistema ti
        supporta.
  \item[Il gioco di Ruggero.] Uno dei giochi della slide è suo, ``quindi non ne
        parliamo'', nonostante le proteste dei colleghi.
  \item[Un gioco su Cuba.] Di una designer che è andata a vivere a
        \textbf{Cuba} e ha partecipato da vicino alla vita di una delle storiche
        comunità queer cubane\footnote{La trascrizione riporta il titolo come
        ``Socialism and Word''; non è verificabile. Anche le date sono incerte
        (``nel 97, 89 e 91''): il contesto, la caduta del Muro di Berlino, fa
        pensare al 1987, 1989 e 1991.}. Il gioco racconta che cosa è successo
        alle comunità \textbf{punk, rock e queer}, molto sovrapposte, prima,
        durante e dopo la caduta del Muro di Berlino: lo Stato, rimasto senza i
        fondi sovietici, decide chi è cittadino di serie A e di serie B, e le
        persone queer e punk (la musica occidentale) vengono picchiate e
        imprigionate. Erano anche gli anni dello scoppio dell'\textbf{AIDS}. Una
        comunità famosa, quella dei \emph{frikis}\footnote{La trascrizione dice
        ``Os Fricchios''.} cubani (di cui sono ancora vive due persone), decise:
        ``tanto ci ammazza la polizia, tanto vale iniettarci sangue infetto da
        HIV'', per andare a morire non per strada ma nei \textbf{centri di
        recupero}, gestiti da medici stranieri, uno dei pochissimi casi in cui a
        Cuba poteva esserci personale non cubano. Lì si moriva, ma ``almeno si
        poteva ascoltare musica punk'', magari con il dottore che arrivava dalla
        Florida con le cassette rock, e tra simili. ``Una roba drammaticissima'':
        anche questo è un gioco devastante.
\end{description}

\begin{table}[htbp]
\centering
\small
\begin{tabular}{@{}p{3.1cm}p{9.8cm}@{}}
\toprule
\textbf{Tipo di LARP} & \textbf{Caratteristiche ed esempi citati} \\
\midrule
Boffer & nato per adattare D\&D dal vivo; costume, armi in plastazote,
combattimento e trame epiche; fino a centinaia o migliaia di persone \\
Mind's Eye Theatre & da \emph{Vampire: The Masquerade}; costume, dark/gotico,
intrigo politico più che combattimento \\
Chamber / parlor LARP & senza costume, accessibile, budget ridotto, spesso
\emph{print and play}, 3--20 persone (\emph{A cena col nonno}) \\
American freeform & chamber LARP con metatecniche prese dal gioco da tavolo e
dal GDR; anche in carte o libretti \\
Murder party & cena con delitto (gli attori indagano, l'assassino è un attore)
o murder party ``all'inglese'' (si interpretano i sospettati) \\
Nordic LARP & coinvolgimento, emozione, coesione tematica, cura del dopo-gioco;
manifesti; Knutpunkt dal 1997; Dogma 99 \\
Jeepform & freeform nordico sul dramma quotidiano (il gioco della coppia
perfetta, \emph{Fat Man Down}); meccaniche di sicurezza \\
360° / alta risoluzione & tutto diegetico grazie a un budget altissimo
(\emph{Ground Zero}, \emph{Panopticorp}, \emph{Monitor Celestra},
\emph{Demeter}, \emph{Bunker 101}) \\
Tape LARP & lo scotch per terra indica gli ambienti, come in \emph{Dogville} \\
Black box & spazi teatrali neri; luci e colori come meccaniche
(\emph{Borderline}, \emph{Waiting for Flight}, \emph{Stranger}) \\
\bottomrule
\end{tabular}
\caption{Le tipologie di gioco di ruolo dal vivo presentate da Ruggero, in
sintesi.}
\label{a2:tab:larp}
\end{table}

% ---------------------------------------------------------------------
\section{Chi osserva? Il LARP e il teatro dell'oppresso}
\label{a2:sec:oppresso}

\paragraph{La domanda di Tiziana.} Tiziana confessa un'impressione: tutti questi
giochi ``sembrano cose da studi'', come se ci fosse qualcuno che sta sopra e
osserva le dinamiche per ricavarne informazioni. Ruggero la rassicura: \textbf{non
è così}. I designer di questa scena, specialmente quella black box, vengono
dall'ambiente ibrido del teatro interattivo, degli \emph{happening}, dello
psicodramma, e \textbf{non documentano nulla di quel tipo}: non registrano le
giocate, non estrapolano informazioni da chi gioca. Spesso sono giochi per poche
persone (è difficile che un black box superi le 20), e molti hanno avuto
\textbf{run singole}, o al massimo due o tre. ``Si tratta solo di fare un certo
tipo di esperienza''. Come in altri media esistono generi con scopi diversissimi
(un film e un documentario sono entrambi operazioni visive), anche qui gli scopi
sono diversi. E chi non ha mai incontrato giochi del genere si chiederà
``perché dovrei giocare una roba del genere?'': la risposta arriverà con la
lezione sull'emotional design (``lo so, ve l'ho detto quattro volte'').

\paragraph{Il teatro dell'oppresso.} Prima di passare la parola, uno studente
interviene: in queste cose ritrova molto il metodo del \textbf{teatro
dell'oppresso}\footnote{Il metodo teatrale ideato dal regista brasiliano Augusto
Boal (integrazione).}. ``Ah, certo''. Ruggero torna indietro di una slide, a
quella sulle \textbf{influenze fondanti del Nordic LARP}: ci sono le tecniche di
\textbf{Moreno} (sociodramma e psicodramma), il \emph{sensitivity training}, il
\textbf{roleplay come metodo di ricerca}, e tutti quei metodi che usano il teatro
e l'interazione come strumento per far emergere il personale e l'intimo. Lo
studente precisa che il teatro dell'oppresso nasceva per \textbf{trovare
soluzioni} a situazioni che apparentemente non ne hanno: a turno il pubblico
proponeva una soluzione, si rifaceva sempre la stessa scena, e si studiava come
reagivano le persone; quindi, in un certo senso, il comportamento delle persone
\emph{si osserva}, ma dall'interno della scena.

Per Ruggero, nella scena LARP questa dimensione ``viene un po' meno'', perché lo
scopo resta \textbf{ludico} (``e scopo ludico può voler dire molte cose, come
vedremo''): è meno pragmatico del ``prendiamo quello che abbiamo giocato e il
gioco ci dà già gli strumenti per applicarlo''. Alcuni \textbf{debriefing}, i
momenti strutturati del dopo-gioco, possono dare strumenti a chi ha giocato; ma
questi giochi nascono come esperienze che devono essere ``ludicamente
interessanti per chi gioca'', \textbf{non come esperimenti scientifici ``con dei
topi dentro un labirinto''}.

\paragraph{Approfondimento: il LARP in classe e il debriefing.} Che il LARP
possa essere uno strumento formativo, come nel volume sull'edu-larp citato
all'inizio, lo mostra anche \emph{Challenges for Game Designers}. Nel capitolo
sui giochi come strumento di insegnamento, Brathwaite e Schreiber propongono una
sfida, \emph{Living History}: progettare un LARP che rappresenti un evento storico
(una convenzione costituzionale, un protocollo internazionale, una guerra)
giocabile da una classe di 20--30 studenti in una o due ore. I passi suggeriti
sono: scegliere un evento in cui molte parti, ciascuna con la propria agenda, si
incontrano in una negoziazione o in un conflitto; scegliere i giocatori (o le
squadre); definire gli \textbf{obiettivi} di ciascuno; creare regole semplici e
rapide per risolvere i conflitti, perché i giocatori si muovono e parlano di
continuo (``alcuni LARP usano un singolo tiro di dado o una partita a
sasso-carta-forbici''); e preparare, oltre alle schede dei personaggi, delle
\textbf{note di debriefing}: che cosa è successo davvero e quali domande
discutere confrontando la storia con l'esito della partita. Lo scopo dichiarato è
che gli studenti capiscano che ``l'evento era complicato'' e come l'esito sia
\emph{emerso} dalla situazione iniziale. È la versione didattica della cura del
dopo-gioco che Ruggero attribuisce alla scena nordica.

% ---------------------------------------------------------------------
\section{I giochi di ruolo da tavolo: quattro parole per orientarsi}
\label{a2:sec:gdr}

``Daniele, ti lascio la parola''. Daniele (``mi giri tu le slide quando te le
chiedo'') avverte che fare una panoramica dei giochi di ruolo in poco tempo è
``veramente difficile'': ce ne sono tantissimi, e negli anni si sono evoluti e
diversificati sempre di più, fino a diventare molto diversi tra loro.

\subsection{Il tribalismo e le regole tramandate a voce}

C'è anche un altro problema, che si vedrà nella lezione sulla fiction e che
esiste in tutto il gioco analogico: il \textbf{tribalismo}, cioè
l'\textbf{ereditarietà orale delle regole}. Spesso i giochi di ruolo venivano
insegnati da altre persone, anche per motivi economici. Daniele ha cominciato
così: da bambino, giocando per la prima volta a D\&D, l'unico ad avere il
manuale era il master; ``ergo io non avevo accesso alle regole'', e il master
poteva inventarsi quello che voleva.

\subsection{Conflitto, azione, risoluzione, effetto}

\slidefig{img/AGD/agd2-car}{Conflitto, azione, risoluzione, effetto: le quattro parole per orientarsi nei giochi di ruolo da tavolo.}{slide \emph{Analog -- Case Studies}}{fig:agd2-car}


Per capire un po' i giochi di ruolo bisogna ricordare \textbf{quattro parole},
che torneranno nella lezione sulla fiction: \textbf{conflitto}, \textbf{azione},
\textbf{risoluzione} ed \textbf{effetto}. Nella slide \emph{conflitto} e
\emph{risoluzione} hanno un asterisco: per ora basta capirli quanto serve a
leggere i casi studio, senza analizzarli troppo.
\begin{description}
  \item[Conflitto.] L'ostacolo, l'incertezza, l'evento che \textbf{stimola
        un'azione o la contrasta}. ``Tornerà molto''.
  \item[Azione.] Quello che si descrive di voler fare nei panni del personaggio
        per risolvere un conflitto.
  \item[Risoluzione.] Come si risolve l'azione dentro il gameplay, con le regole:
        che cosa succede?
  \item[Effetto.] La conseguenza che si applica.
\end{description}
L'esempio: ``voglio colpire il goblin con la mia spada''. Il goblin ha il suo
scudo e la sua armatura. Che azione fai? ``Lo becco con la spada''. Si tirano i
dadi, si risolve, si applica un effetto. Ma non è così semplice, perché in
mezzo c'è la \textbf{fiction}: nel gioco di ruolo si tende a riscrivere tutto
quello che succede nel gameplay \emph{dentro} la narrazione. Se faccio tre danni
agli HP del goblin, che cosa vuol dire narrativamente? Come li descrivo? Da qui
nasce ``una serie di problematiche''.

\begin{figure}[htbp]
\centering
\begin{tikzpicture}[
    font=\small,
    passo/.style={draw, thick, rounded corners=3pt, fill=gray!12,
                  minimum width=2.6cm, minimum height=0.95cm, align=center},
    esempio/.style={font=\footnotesize\itshape, align=center, text width=2.8cm},
    freccia/.style={-{Stealth[length=2.5mm]}, thick}
  ]
  \node[passo] (c) at (0,0)    {\textbf{Conflitto}$^{*}$};
  \node[passo] (a) at (3.4,0)  {\textbf{Azione}};
  \node[passo] (r) at (6.8,0)  {\textbf{Risoluzione}$^{*}$};
  \node[passo] (e) at (10.2,0) {\textbf{Effetto}};
  \draw[freccia] (c) -- (a);
  \draw[freccia] (a) -- (r);
  \draw[freccia] (r) -- (e);
  \node[esempio, below=0.25cm of c] {il goblin, con scudo e armatura};
  \node[esempio, below=0.25cm of a] {``lo colpisco con la spada''};
  \node[esempio, below=0.25cm of r] {tiro dei dadi secondo le regole};
  \node[esempio, below=0.25cm of e] {tre danni agli HP};
  \node[draw, dashed, rounded corners=3pt, fit=(c)(e), inner sep=0.5cm,
        yshift=-0.35cm, minimum height=2.6cm] (fic) {};
  \node[font=\footnotesize, anchor=south] at (fic.north)
    {\textbf{fiction}: tutto viene ridescritto nella narrazione (che cosa
    \emph{sono}, nella storia, tre danni?)};
\end{tikzpicture}
\caption{Le quattro parole di Daniele per leggere i giochi di ruolo da tavolo,
con l'esempio del goblin. Gli asterischi sono quelli della slide: conflitto e
risoluzione verranno analizzati meglio nella lezione sulla fiction. Nel gioco di
ruolo ``classico'' il conflitto precede l'azione, la risoluzione è un tiro di
dado e l'effetto è descritto dal master e dal regolamento; molti giochi
successivi cambiano proprio questo schema.}
\label{a2:fig:car}
\end{figure}

% ---------------------------------------------------------------------
\section{Le origini: da Dungeons \& Dragons agli anni Novanta}
\label{a2:sec:gdr-storia}

Le quattro parole servono a Daniele per una breve panoramica, anche storica,
dei giochi di ruolo.

\subsection{Dungeons \& Dragons e il gioco di guerra}

Il primo gioco di ruolo da tavolo ``così come viene concepito'' è sicuramente
\textbf{\emph{Dungeons \& Dragons}}, di \textbf{Gygax e Arneson}, nel
\textbf{1974}. C'è in realtà un precedente, \textbf{\emph{Chainmail}}\footnote{La
trascrizione dice ``C'è in Mail'' e data \emph{Chainmail} al '72; il gioco è
del 1971.}, e come per tutte le origini (``anche \emph{Cyberpunk} ha lo stesso
problema'') c'è un po' di ambiguità su che cosa venga prima. Ci sono libri che
raccontano la nascita di D\&D, c'è la storia della \textbf{TSR}, la casa
editrice che lo pubblicò per prima, ci sono versioni discordanti e, negli anni,
parecchi contrasti tra le famiglie di Gygax, di Arneson e di altri.

D\&D viene dal \textbf{gioco di guerra}, il gioco di \textbf{miniature}. ``La
leggenda vuole'' che Gygax e Arneson, giocando, si siano detti: ho battuto il
tuo esercito; che cosa succede se il mio generale entra nel tuo castello e va a
trattare con il tuo generale? Abbiamo regole per questo? No: inventiamocele.
``Ripeto, una leggenda'': non si sa se sia andata così, e in fondo non interessa
molto.

\subsection{I generi e la struttura del gioco di ruolo ``early''}

\slidefig{img/AGD/agd2-dnd}{Il flusso di Dungeons \& Dragons letto con conflitto, azione, risoluzione ed effetto.}{slide \emph{Analog -- Case Studies}}{fig:agd2-dnd}


Quello che interessa è che dopo D\&D, nell'\textbf{RPG delle origini}, vengono
coperti più o meno \textbf{tutti i generi letterari}: \emph{Call of Cthulhu}
l'orrore lovecraftiano, \emph{Vampire: The Masquerade} l'orrore gotico,
\emph{Traveller} la fantascienza, \emph{Cyberpunk} il cyberpunk, e \emph{GURPS},
un \textbf{sistema generico} con cui ``teoricamente'' si può giocare qualunque
genere.

Nel gioco di ruolo classico:
\begin{itemize}
  \item il \textbf{conflitto} poteva venire prima o dopo l'\textbf{azione}:
        ``voglio colpire quel goblin'' (e il goblin ha la sua difesa), oppure
        ``corri, ti sta inseguendo una palla'', come in \emph{Indiana Jones};
  \item la \textbf{risoluzione} era il tiro di dado, e l'\textbf{effetto} veniva
        descritto dal master e dal regolamento;
  \item in pratica, \textbf{tutto tranne l'azione}, che il giocatore gestiva e
        narrava, era gestito dal \textbf{master}.
\end{itemize}

\paragraph{Il master come arbitro.} All'inizio il master non era un ``maestro''
né una figura che insegna: era interpretato soprattutto come un
\textbf{arbitro}. Nella prima versione di D\&D non c'è ancora l'idea, venuta
dopo, del master come ``dio del tavolo'': quando deve decidere, dovrebbe farlo
secondo criteri di realismo e oggettività (``puoi essere oggettivo a un tavolo
di gioco? La voglio vedere''). Questa idea è importantissima perché verrà ripresa
dall'\textbf{OSR}, l'\emph{Old School Revival}, il movimento che cerca di
riportare nella modernità il modo in cui, a suo dire, si giocava D\&D agli
inizi.

\paragraph{Il master è un giocatore.} Nei giochi di ruolo ci sono quindi
\textbf{ruoli asimmetrici}. Ma Daniele raccomanda di ricordare una cosa, che
tornerà nella lezione sulla fiction: \textbf{il master è un giocatore}, magari un
giocatore con più regole e più autorità al tavolo, ma ``un giocatore tanto quanto
i giocatori che interpretano un personaggio''. Nei GDR delle origini: si
interpreta un personaggio, con una \textbf{scheda} e le sue statistiche, dentro
una fiction e un \textbf{setting} stabilito dal manuale o dal master; quando si
dice che cosa si vuole fare, il master reagisce. E lentamente, con gli anni, il
master assume anche il ruolo di \textbf{regista} e \textbf{sceneggiatore}: scrive
l'avventura, decide che cosa succede, decide i tagli di scena.

\paragraph{Approfondimento: dal wargame al contratto sociale.} La genealogia
raccontata da Daniele si ritrova in \emph{21st Century Game Design} di Bateman e
Boon. Le radici del gioco di strategia, scrivono, precedono i videogiochi di
decenni: alla fine degli anni Cinquanta Charles Roberts fonda Avalon Hill e
pubblica i primi \textbf{wargame da tavolo} (1958); in un'epoca senza
videogiochi, giochi di ruolo e giochi di carte collezionabili, i wargame
``fondano l'industria dei \emph{hobby games}'', e \emph{Chainmail} (1971),
ispirato a quei giochi, è il ``trampolino'' verso \emph{Dungeons \& Dragons}. Gli
stessi autori descrivono la forma minima del gioco di ruolo da tavolo come un
\textbf{contratto sociale} tra l'arbitro (il \emph{games master}) e i
giocatori: i giocatori accettano che sia l'arbitro a decidere come le loro azioni
sviluppano la storia, e l'arbitro si impegna a garantire intrattenimento o
imparzialità, come mediatore e narratore principale. Citano anche un'osservazione
attribuita a Gygax: se i giocatori di ruolo si accorgessero di poter fare tutto
da soli, i game designer resterebbero senza lavoro. Il fatto, notano, è che la
maggior parte dei master ha bisogno di meccaniche per gestire il mondo, e i
giocatori si fidano di più dell'arbitro quando una struttura di meccaniche
casuali, i dadi, fa da cornice. Anche Schell, in \emph{The Art of Game Design},
cita il master dei giochi di ruolo da tavolo come esempio di stato del gioco
\emph{nascosto}: è lui a conoscere privatamente gran parte dello stato del gioco,
perché ne è ``il meccanismo operativo''.

\subsection{Gli anni Ottanta e Novanta: adattare un regolamento a tutto}

Qui Daniele fa un ``piccolo inciso editoriale''. Nella grande espansione degli
anni Ottanta e Novanta \textbf{non c'era un professionismo} del settore. Gli
editori compravano delle proprietà intellettuali, per esempio \textbf{007},
prendevano un master appassionato e competente e gli chiedevano di
\textbf{adattare} il regolamento di D\&D, o un altro regolamento
esistente\footnote{La trascrizione dice ``Call of Duty 2'', che non può essere
corretto in questo contesto; il senso è ``il regolamento di un altro gioco
già esistente''.}, a quella IP. A volte nascevano giochi innovativi, a volte
giochi ``non molto interessanti da giocare'': immaginate di giocare con le
regole di D\&D qualcosa che quelle regole non supportano. Verso la fine degli
anni Novanta il problema della qualità ``ha iniziato a farsi sentire''.

\paragraph{Approfondimento: James Bond 007 e l'esperienza essenziale.} Proprio
il gioco di ruolo di 007 è il caso con cui Schell illustra, in \emph{The Art of
Game Design}, la \textbf{lente dell'esperienza essenziale} (\emph{Lens of
Essential Experience}), e mostra l'altra faccia della storia di Daniele. Il
designer Chris Klug era frustrato dai precedenti giochi di ruolo di spie, come
\emph{Top Secret} della TSR, perché ``si giocavano troppo come dei wargame'':
l'essenza di ciò che rende eccitanti i film di spionaggio non c'era. Per il
gioco di James Bond, Klug progettò le meccaniche in modo che restituissero in
ogni modo possibile la sensazione dei film. L'esempio è l'invenzione degli
\textbf{Hero Points}: nei GDR tradizionali, davanti a un'azione rischiosa (saltare
da una finestra su un elicottero in volo), il master calcola una probabilità, il
giocatore tira i dadi ``e questo è quanto''; ma se la probabilità è bassa
nessuno rischia, e se è alta tutti si comportano da supereroi. I punti eroe
risolvono il dilemma in modo coerente con il genere. È esattamente quello che
la New Wave rivendicherà: \textbf{il sistema conta}, e un regolamento adattato ``a
nastro'' non restituisce l'esperienza di un genere.

% ---------------------------------------------------------------------
\section{La New Wave e The Forge}
\label{a2:sec:gdr-newwave}

\subsection{Il sistema conta}

La slide successiva è sulla \textbf{New Wave}: un movimento nato tra la fine
degli anni Novanta e l'inizio dei Duemila. Daniele avverte che non vuole fare
un'analisi geografica ma ``un riassunto della storiografia''. La New Wave cerca
un \textbf{metodo diverso di interpretare il gioco di ruolo}, e la prima cosa
che stabilisce, in un forum chiamato \textbf{The Forge} (c'è anche un libro in
inglese su questo movimento), è che \textbf{il sistema dei giochi di ruolo è
importante}: le regole, il gameplay, la parte meccanica \textbf{non sono un
orpello}. Non si può tenere lo stesso sistema, cambiare il setting e sperare che
funzioni da sé. Tanto è vero che perfino \emph{GURPS}, il sistema generalista,
aveva tantissimi manuali di espansione delle regole per ogni genere:
\emph{GURPS} base, \emph{GURPS Fantasy}, \emph{GURPS} fantascienza\ldots

\subsection{I giochi della New Wave}

\slidefig{img/AGD/agd2-newwave}{Alcuni giochi della New Wave.}{slide \emph{Analog -- Case Studies}}{fig:agd2-newwave}


Dentro la New Wave nascono esperienze diversissime, che spostano l'attenzione
su due cose: la \textbf{scrittura del manuale}, diversa da quella tradizionale,
e la \textbf{rottura degli schemi}.
\begin{itemize}
  \item \textbf{\emph{Apocalypse World}}\footnote{La trascrizione dice ``A
        Polly's Ward'' e ``Minster Baker''.}, di \textbf{Vincent Baker}, cerca di
        scardinare il modo classico di strutturare un manuale di gioco di ruolo
        (Daniele accenna a una nuova edizione recente).
  \item \textbf{\emph{Nobilis}}, uno dei primi giochi abbastanza famosi
        \textbf{diceless}, senza dadi: azioni e conflitti si risolvono in un
        altro modo.
  \item \textbf{\emph{Don't Rest Your Head}} (in italiano \emph{Non cedere al
        sonno}) e \textbf{\emph{Polaris}}, che lavorano sul concetto di chi è e
        chi non è il personaggio. \emph{Non cedere al sonno} è un gioco
        ``surrealista'' in cui ci si sveglia a \textbf{Mad City}; la parte
        importante sono i \textbf{poteri}. All'epoca \emph{Vampire} aveva
        stabilito che i poteri fossero una lista di possibilità con statistiche
        a cui applicare le regole; in \emph{Don't Rest Your Head} ci sono poteri
        ``totalmente assurdi, fenomenali'', che non richiedono meccaniche in più
        ma usano le meccaniche base ``in maniera molto innovativa''.
  \item \textbf{\emph{Fate}}, un sistema anch'esso generalista, ``più o meno
        come \emph{GURPS}'', ma con delle sue specifiche: ci sono
        \textbf{procedure per creare i setting} e per stabilire quali regole un
        setting richiede.
  \item \textbf{\emph{Dread}}, uno dei primi giochi a usare per l'incertezza
        \textbf{un elemento che non è un dado}: una \textbf{torre Jenga}. È un
        gioco dell'orrore in cui ogni volta che si fa un'azione si toglie un
        blocchetto, e se la torre cade il personaggio muore (``non è proprio
        così, è più complicato, però per semplificare'').
\end{itemize}

\paragraph{Le guerre di religione.} La New Wave ha avuto ``una serie di
problematiche'' in giro per il mondo, e soprattutto in Italia ha creato una
sorta di \textbf{guerra teologico-filosofica} tra i ``tradizionalisti'' di D\&D e
i ``forgiti'' della New Wave. ``Per vostra sanità mentale non ci entriamo''.

\paragraph{Una reazione a due cose.} La New Wave è stata una reazione:
\begin{enumerate}
  \item ai manuali in cui ``c'è un sistema e quello è il sistema'', con il
        setting applicato sopra ``a nastro'', senza intreccio tra la narrazione
        e le regole che la raccontano: si voleva \textbf{coerenza};
  \item all'assenza di studio: è stato il \textbf{primo movimento a voler
        studiare i giochi di ruolo}, a capire come sono fatti. Da qui viene il
        \textbf{Big Model}, un modello di analisi dei giochi di ruolo di cui si
        vedranno alcune cose nella lezione sulla fiction.
\end{enumerate}
Ma il Big Model ``non è fondamentale'', perché per Daniele uno dei problemi della
New Wave è aver creato \textbf{un dizionario collaterale} a quello del game design
del gioco da tavolo e del videogioco: termini diversi che dicono la stessa cosa.

\subsection{The Forge e la scena nordica}

Ruggero fa da ponte: \textbf{The Forge ha fatto per il gioco di ruolo da tavolo
quello che la scena scandinava ha fatto per il LARP}. Inventavano e scrivevano
moltissimo, hanno lasciato framework e strumenti. La scena scandinava è nata un
po' prima, The Forge un po' dopo, ma hanno agito in parallelo. La differenza è
che The Forge era \textbf{meno strutturato}: essendo un forum, era più caotico,
con discussioni anche accese, soprattutto sulle interpretazioni. Anche la
\textbf{preservazione storica} è diversa: della scena nordica ci sono i libri
fisici, mentre per il forum bisogna passare dalla Wayback Machine. Ma in termini
di normalizzazione di termini, meccaniche e framework, entrambi hanno fatto
moltissimo.

\subsection{Fiasco}

\slidefig{img/AGD/agd2-fiasco}{Il flusso di Fiasco: chi stabilisce la scena e chi la risolve.}{slide \emph{Analog -- Case Studies}}{fig:agd2-fiasco}


Daniele riprende: a un certo punto The Forge chiude, e nascono altri giochi
dagli stimoli e dalle discussioni che quell'ambiente aveva generato, non
necessariamente da persone che ne facevano parte. Per esempio quelli di
\textbf{Jason Morningstar}, noto anche per dei chamber LARP (il migliore amico
di Daniele, Edoardo, ``più esperto di me'', ne è talmente fan che il suo gatto si
chiama Giasone).

\textbf{\emph{Fiasco}} è per Daniele uno dei primi giochi che \textbf{rompe il
canone}. Racconta, appunto, un fiasco, come nei film dei \textbf{fratelli Coen}
o di Tarantino: c'è uno scenario, ci sono dei personaggi, e ogni volta che si fa
una scena \textbf{il giocatore di turno sceglie se aprirla} (e gli altri la
risolvono) \textbf{o se lasciare che gli altri la impostino} (e risolverla lui).
Già solo questo ``micro elemento'' (non l'unica meccanica: c'è anche quella,
interessantissima, dei dadi) rompe l'idea che in un gioco di ruolo si debba
descrivere l'azione, tirare i dadi e vedere che cosa succede, oppure reagire a
ciò che descrive il master. Per chi veniva da \emph{AD\&D} seconda edizione o da
\emph{Vampire} era ``assolutamente interessante'': più leggero e più legato alla
fiction. Anche i \textbf{dadi} sono usati in modo nuovo: vengono tirati
\textbf{prima}, messi sul tavolo, selezionati, usati quasi come
\textbf{token}: un'innovazione anche nel modo di usare i materiali.

% ---------------------------------------------------------------------
\section{Il gioco di ruolo contemporaneo}
\label{a2:sec:gdr-oggi}

\subsection{Prima del COVID: i giochi focalizzati}

\slidefig{img/AGD/agd2-contemporary}{Giochi di ruolo contemporanei.}{slide \emph{Analog -- Case Studies}}{fig:agd2-contemporary}


Daniele salta alla contemporaneità. Prima del COVID c'è un grande aumento di
giochi che, mentre le persone si appassionano ai giochi di ruolo, vanno o nella
direzione di D\&D o in una \textbf{direzione più ibrida}:
\begin{itemize}
  \item \textbf{\emph{Blades in the Dark}}\footnote{La trascrizione dice
        ``Blaze in the Dark''. Daniele lo data ``2016--18 più o meno''.};
  \item più recenti, il gioco di ruolo di \textbf{\emph{Cowboy
        Bebop}}\footnote{La trascrizione dice ``Cabobby Bop'' e ``Cobo
        Bobby''.} e \textbf{\emph{The Clay That Woke}} di \textbf{Paul
        Czege}\footnote{La trascrizione dice ``Paul Ksege, The Clayed Walk''.},
        autore di ``uno dei principi di design più interessanti'';
  \item \textbf{\emph{Fabula Ultima}}, di \textbf{Emanuele Galletto}, italiano
        (``che salutiamo'').
\end{itemize}
Sono prodotti che \textbf{non cercano di fare la stessa cosa di D\&D}. Hanno un
impianto simile (scheda del personaggio, dadi, azioni descritte, tiri), ma si
pongono \textbf{obiettivi di design più coerenti}: ``si chiedono che cosa vogliamo
che il giocatore senta''. Il sistema di \emph{Cowboy Bebop}, per esempio, usa
gli \textbf{orologi} (\emph{clock}) che si trovano anche in \emph{Blades in the
Dark} e nei giochi \emph{Powered by the Apocalypse}, cioè la stessa meccanica,
ma è un sistema completamente diverso, perché deve raccontare una fiction
diversa:
\begin{itemize}
  \item in \emph{Blades in the Dark} si giocano dei \textbf{criminali} a
        \textbf{Doskvol}\footnote{La trascrizione dice ``Oswald''.}, città di un
        mondo alternativo dove è sempre notte, dove i fantasmi non passano
        nell'aldilà e vivono con i vivi, e dove si è costretti a restare perché
        fuori dalle mura i fantasmi uccidono. Si fa parte di una banda che cerca
        di scalare la gerarchia della città, ``tipo \emph{Peaky Blinders}'';
  \item in \emph{Cowboy Bebop} si è ``letteralmente una ciurma di persone nello
        spazio, nel futuro, che fanno cose'', perché ``l'anime si racconta con i
        personaggi che fanno cose'' (Daniele, ``molto fan dell'anime'', lo dice
        come un complimento).
\end{itemize}

\subsection{Dopo il COVID: il ritorno all'origine}

Soprattutto dopo il COVID si è notato un ``ritorno fortissimo'' al
\textbf{giocare come una volta}. Per Daniele entrambe le tendenze rispondono a
una necessità: \textbf{semplificare}. \emph{D\&D} 3 e 3.5, le nuove edizioni di
\emph{Vampire} e altri sono diventati ``pieni di regole, pieni di contenuti'', e
chiedono a master e giocatori di imparare moltissimo, senza contare il tempo
necessario per una campagna. Allora, da un lato, si cercano sistemi più semplici
``da intavolare'' che facciano bene una cosa; dall'altro si cerca la stessa cosa
per altri motivi, ricreando le atmosfere del D\&D \emph{old school}: \textbf{alta
mortalità}, parte narrativa e interpretativa meno importante, immediatezza.
\textbf{\emph{Mörk Borg}}\footnote{La trascrizione dice ``Morg Borg''.}, per
esempio, sono circa \textbf{cento pagine}, con poche regole. C'è più attenzione
al ``realismo'' (termine un po' sbagliato, si vedrà nella lezione sulla
fiction), e il vantaggio è che questi giochi possono usare sia le proprie regole
sia quelle di \textbf{D\&D quinta edizione}, che è stata importante soprattutto
``per un motivo editoriale'' (anche questo si vedrà nella lezione sulla
fiction).

\paragraph{La domanda di Gianmarco.} Gianmarco chiede se questa tendenza a
manuali e partite ``asciutti'' segua un pubblico più adulto, con meno tempo, o
sia una tendenza generale del mercato. È soprattutto la prima. Il gioco di ruolo
ha avuto parecchie difficoltà ad attirare nuove leve, e i giocatori sono
cresciuti. Daniele gioca da molti anni, ma il suo ``vero inizio'' è stato con un
gioco della New Wave\footnote{La trascrizione dice ``Paura dell'Epoclips'':
probabilmente \emph{Il mondo dell'apocalisse}, l'edizione italiana di
\emph{Apocalypse World}, ma non è certo.}: per lui, all'università, avere in un
solo manuale da 200 pagine tutte le regole e tutti i personaggi è stato ``un
toccasana''. Chi inizia con D\&D, magari alle superiori o all'università, man
mano che cresce tende a preferire:
\begin{itemize}
  \item o giochi più focalizzati sull'aspetto \textbf{narrativo}, che hanno
        comunque scheda, tiro di dado e struttura tradizionale, ma una scheda e
        un sistema più semplici;
  \item o un'esperienza simile al D\&D canonico, ma con meno ``regoline e
        regolette'' e meno \textbf{simulazionismo}.
\end{itemize}
Due differenze pratiche. I manuali \textbf{OSR}, che un po' di simulazionismo ce
l'hanno, sono pieni di \textbf{soluzioni per l'arbitro} (``che non è il master,
è l'arbitro'') per risolvere i problemi sul momento; e l'OSR è \textbf{sandbox}:
c'è una mappa, una zona in cui si va dove si vuole, e il master non ha scritto
un'avventura. I \textbf{giochi narrativi}, dal canto loro, non hanno una campagna
scritta ma uno \textbf{scenario}. In entrambi i casi il tempo di preparazione del
master si riduce ``di gran lunga''.

\subsection{Focalizzati, non semplificati}

Ruggero interviene sulle parole: non vorrebbe che passasse l'idea che questi
giochi siano ``più semplici'' o \emph{dumbed down} rispetto a D\&D. Si parla di
\textbf{giochi focalizzati}: invece di semplificare, \textbf{focalizzano}, sono
molto più \textbf{coesi tematicamente}, hanno ``un cuore tematico molto forte'' e
si giocano per giocare quello. D\&D invece, a livello di cuore meccanico, è
giocato in modi diversissimi da dieci gruppi diversi.

Daniele fa un esempio: nei manuali di D\&D (dalla terza edizione e 3.5 alla
quinta) ci sono le regole per il \textbf{danno da caduta}. In \emph{Blades in
the Dark} non servono: c'è un altro metodo, ``non è così generalista, è
focalizzato'', e la narrazione si risolve in altro modo. Anche nell'OSR:
\emph{Shadowdark} ha il danno da caduta, perché magari è importante;
\emph{Neverland}, una reinterpretazione di \textbf{Peter Pan} in chiave OSR, no,
perché lì conta di più che il manuale offra tutti i personaggi del Peter Pan
tradizionale in una versione \emph{grimdark}. ``Anche l'OSR diventa un po'
focalizzato quando ci prova''.

\subsection{Not the End}

\slidefig{img/AGD/agd2-nottheend}{Il flusso di Not the End: il sacchetto di token come motore della risoluzione.}{slide \emph{Analog -- Case Studies}}{fig:agd2-nottheend}


Nicola, per cui questo è ``il settore che più adoro del gioco analogico'',
chiede dove si collochi \textbf{\emph{Not the End}}, un gioco italiano. La
risposta è nella slide successiva. \emph{Not the End} continua nella direzione
più narrativa e focalizzata (``semplicistica è un termine non corretto''): ha
procedure, meccaniche, un gameplay focalizzato a fare quello che il gioco
chiede. Si possono creare i propri setting, ma il gioco resta quello: il
personaggio principale è un \textbf{eroe}, ed è il giocatore a decidere quando
sarà la sua fine, quando giocherà il suo momento di eroismo stoico (da cui il
titolo). È un gioco \textbf{senza dadi}: c'è un conflitto, si fa e si descrive
un'azione, e la \textbf{risoluzione} consiste nell'\textbf{estrarre dei token}
bianchi e neri (o arancioni e viola, per le persone con difficoltà nella visione
dei colori); l'effetto viene descritto secondo le procedure del gioco. Il
panorama, insomma, si è diviso: c'è chi cerca un'esperienza più simile
all'originale, e chi un'esperienza più focalizzata sulla narrazione, magari
senza un \emph{world building} preciso.

\subsection{Che cos'è l'OSR, e com'è nato davvero il gioco di ruolo}

Un'altra studentessa chiede che cosa significhi OSR: \textbf{Old School
Renaissance}, ``rinascimento della vecchia scuola''\footnote{Poco prima Daniele
aveva sciolto la sigla come \emph{Old School Revival}; entrambe le forme sono in
uso.}. È un movimento ``molto, molto particolare'': ogni persona che ne fa parte
ha il proprio principio di OSR, gioca un OSR diverso dagli altri, e gli
appassionati ``adorano scontrarsi'' su che cosa intendono per OSR. Per questo
nella lezione sulla fiction Daniele parlerà di \textbf{culture} del GDR, più che
di community.

La stessa studentessa (iscritta a un altro corso dell'accademia, di cui sta
recuperando le lezioni del primo anno) riporta quello che ha capito dell'origine
di D\&D: i due autori avrebbero avuto visioni opposte, uno più libera e meno
regolamentata, partendo dai giochi di guerra, l'altro, che ha portato il gioco a
livello commerciale, più strutturata, dettagliata e pesante. ``È giusto, è un
riassunto corretto, ma è un po' più complicato di così'', e lo si vedrà nella
lezione sulla fiction. Ruggero aggiunge che D\&D ``non è l'unica origine del
gioco di ruolo'': il gioco di ruolo diventa \emph{famoso} con D\&D, ma gli stessi
autori si ispirarono alle \textbf{simulazioni militari}, molto diffuse
nell'America di quegli anni, e al \textbf{reenactment storico} (``ci sono due
episodi della \emph{Signora in giallo} sul reenactment della guerra civile, un
argomento allucinante''). Esistono esempi storici, dell'Ottocento e perfino del
Seicento, di \textbf{giochi di ruolo} veri e propri, con quest, avventure e
personaggi. Il \emph{mainstream} ha normalizzato l'idea che D\&D sia il primo
gioco di ruolo, ``ma non è così'': nel manuale originale l'espressione
\emph{role-playing game} non compare da nessuna parte, anche perché Gygax
odiava il termine. ``Un vaso di Pandora che non apriamo''.

Daniele però raccoglie l'osservazione. Nei primi anni editoriali di D\&D il
problema maggiore è stato il \textbf{simulazionismo}: interpretare
``simulazione'' come la necessità di avere una regola per ogni cosa che esiste nel
realismo del gioco. Ne è venuto un gioco molto strutturato, in cui ``ogni cosa
deve avere una regola'', perché ciò che non è coperto dalle regole il giocatore
può usarlo come \emph{loophole}. E ne è venuta anche l'idea che il master sia il
\textbf{padrone del tavolo} e che il giocatore debba in qualche modo
\textbf{vincere contro il master}: un altro ``vaso di Pandora enorme''.

\paragraph{Non esiste solo D\&D.} È per questo che Daniele non si concentra solo
su D\&D: ``l'idea con cui dovete uscire da qui è che non esiste solo D\&D''.
Anzi, oggi D\&D è \textbf{in controtendenza}, non per motivi di design ma per
\textbf{motivi editoriali}: il problema principale di D\&D come proprietà
intellettuale ``è sempre stato editoriale'', e di recente la Wizards of the Coast
ha fatto un po' di confusione con il \emph{branding}.

% ---------------------------------------------------------------------
\section{Story game e giochi di ruolo sperimentali}
\label{a2:sec:gdr-story}

\subsection{Gli story game}

\slidefig{img/AGD/agd2-storygames}{Alcuni story game.}{slide \emph{Analog -- Case Studies}}{fig:agd2-storygames}


La slide successiva è sugli \textbf{story game}, che in realtà nascono anche
prima di The Forge. ``Story game è un termine che io odio, perché vuol dire tutto
e non vuol dire niente'': sono giochi di ruolo da tavolo che si focalizzano
\textbf{ancora di più su una storia precisa}.
\begin{itemize}
  \item \textbf{\emph{Alice is Missing}}: la storia di Alice, scomparsa, e dei
        suoi amici che cercano di capire che cosa le è successo. Si gioca con
        una \textbf{chat} sul telefono, con un video a schermo, e con
        \textbf{due ore di tempo}.
  \item \textbf{\emph{Microscope}}: ``letteralmente \emph{world building}
        narrativo''.
  \item \textbf{\emph{Bleak Spirit}}: giocare un \emph{Dark Souls} in cui tutti
        controllano a turno il viandante protagonista.
  \item \textbf{\emph{The Zone}}\footnote{La trascrizione attribuisce il gioco a
        ``Ruff D'Amico''; il nome non è verificabile.}, giocabile anche con un
        supporto online: si narra, di base, \emph{Annientamento} (``se avete
        visto il film''). Si entra in una zona in cui una presenza aliena sta
        cambiando il mondo, e \textbf{moriranno tutti tranne una persona}, che
        non si sa in anticipo chi sarà. È un'influenza del LARP: come in
        \emph{Demeter}, si sa che il proprio personaggio morirà.
  \item \textbf{\emph{Archipelago}} e \textbf{\emph{For the Queen}}, e altri
        giochi che \textbf{tolgono i dadi} e mettono le carte, o tolgono
        perfino la \textbf{scheda del personaggio}: in \emph{Microscope} ci sono
        solo dei post-it.
\end{itemize}

\subsection{The Quiet Year e i map drawing game}

\slidefig{img/AGD/agd2-quietyear}{Il flusso di The Quiet Year.}{slide \emph{Analog -- Case Studies}}{fig:agd2-quietyear}


``C'è un minimo di suspense'': si arriva a esperienze ancora più lontane.
In \textbf{\emph{The Quiet Year}} non si interpreta nemmeno un personaggio: si è
parte di una \textbf{comunità post-apocalittica} e si gioca \textbf{un anno} della
sua vita. Non si può discutere di ciò che si gioca. A ogni turno si gira una
\textbf{carta} che dice che cosa succede, e durante il proprio turno si può fare
una tra due o tre azioni (disegnare, iniziare un progetto\ldots). Tutto si gioca
su una \textbf{mappa} su cui si disegna ciò che si vuole intendere: c'è chi gioca
in modo più rigoroso, senza poter nemmeno scrivere parole sulla mappa, e chi lo
permette. È ``un gioco fenomenale'', perché fa capire che il gioco di ruolo può
essere anche altro; e ha normalizzato un genere, quello dei \textbf{map drawing
role-playing game}, i giochi di ruolo in cui si disegnano mappe.

\subsection{Dialect}

\slidefig{img/AGD/agd2-dialect}{Il flusso di Dialect.}{slide \emph{Analog -- Case Studies}}{fig:agd2-dialect}


E poi \textbf{\emph{Dialect}}. Qui si ha un personaggio, ma è solo un nome e un
ruolo, un archetipo. Si gioca \textbf{la morte di una comunità attraverso la
morte di una lingua}: una comunità umana isolata dalle altre, che ha creato il
proprio linguaggio, e che giocherà la fine di quell'isolamento attraverso la
morte del linguaggio. A ogni turno si gioca una carta che è un \textbf{concetto
della lingua} (per esempio: come si dice ``amico'' nella nostra lingua?); si
stabilisce la parola e si gioca una breve scena in personaggio, senza tiri di
dado né altro; quando la parola viene usata in personaggio, chiunque al tavolo
può dire stop: ``abbiamo usato la parola, va bene così'', e si va avanti.

Non è un gioco ``non complesso'', come diceva Ruggero: \emph{Dialect} è
``chiarissimo'' da leggere, ma la complessità sta nel \textbf{design}, un design
elegante che crea una \emph{one shot} con ``una profondità'' enorme. Daniele con
\emph{Dialect} ha pianto; e quando incontra gli amici con cui lo ha giocato anni
fa, si parlano ancora nella lingua inventata allora: ``ehi tu, oceano'', e chi
c'era sa che cosa vuol dire. Anche qui si vede come le esperienze nate dal Nordic
LARP siano \textbf{``pervase''} (``rubo una parola a Beatrice'') in altri ambiti.

\subsection{Servono davvero le categorie? Il problema editoriale}

Ci sono ancora esperienze più tradizionali; ``il mondo è molto variegato'', e si
parlerà di \textbf{culture}: spesso chi si avvicina al gioco di ruolo vuole
giocare una certa cosa, una cultura di gioco con titoli che ``vanno bene'' e
titoli che no. L'ultima slide è ``una piccola critica'': è davvero necessario
mettere le categorie ai giochi di ruolo? Da un lato sì, dall'altro no, perché i
problemi principali sono due, e il primo è \textbf{il mondo editoriale}: i giochi
di ruolo sono stampati come \textbf{libri}, e l'editore ha ``una vita travagliata''
con il gioco. Esempi:
\begin{itemize}
  \item \textbf{\emph{In the Name of God}}\footnote{Titolo e autore (``Alessandro
        Pirodi'') come nella trascrizione; non verificabili.}, stampato solo in
        \textbf{self-publishing} e assente dal mercato, ``ma bellissimo''; la sua
        meccanica di risoluzione chiede di fare qualcosa fisicamente al tavolo
        (la trascrizione dice ``fare home'');
  \item un'\textbf{espansione per qualsiasi gioco di ruolo}\footnote{La
        trascrizione dice ``Mew Wild in the Subway''.};
  \item un \textbf{GDR stampato su una maglietta}, di due autori che Daniele
        saluta;
  \item un gioco di \textbf{cultisti contro investigatori} che prende da
        esperienze di chamber LARP\footnote{La trascrizione lo chiama ``SOT''.}.
\end{itemize}

\begin{figure}[htbp]
\centering
\begin{tikzpicture}[
    font=\footnotesize,
    epoca/.style={draw, thick, rounded corners=3pt, fill=gray!12,
                  text width=2.55cm, align=center, minimum height=1.0cm},
    nota/.style={align=left, text width=2.7cm, font=\scriptsize},
    freccia/.style={-{Stealth[length=2.5mm]}, thick}
  ]
  \draw[freccia] (-0.2,0) -- (15.4,0);
  \foreach \x/\a in {0.9/{anni '50--'70}, 4.0/{'74--anni '90},
                     7.1/{fine '90--2000}, 10.2/{anni 2010}, 13.3/{post-COVID}}{
    \draw[thick] (\x,0.12) -- (\x,-0.12);
    \node[below=0.15cm, font=\scriptsize\itshape] at (\x,0) {\a};
  }
  \node[epoca] (e1) at (0.9,1.2)  {\textbf{Wargame e miniature}};
  \node[epoca] (e2) at (4.0,1.2)  {\textbf{D\&D e GDR ``early''}};
  \node[epoca] (e3) at (7.1,1.2)  {\textbf{New Wave / The Forge}};
  \node[epoca] (e4) at (10.2,1.2) {\textbf{Giochi focalizzati, story game}};
  \node[epoca] (e5) at (13.3,1.2) {\textbf{OSR e ritorno all'origine}};
  \node[nota, anchor=north] at (0.9,-0.65) {simulazioni militari,
    reenactment; \emph{Chainmail} (1971)};
  \node[nota, anchor=north] at (4.0,-0.65) {Gygax e Arneson;
    master arbitro; tutti i generi (\emph{Cthulhu}, \emph{Vampire},
    \emph{Traveller}, \emph{GURPS}); IP adattate};
  \node[nota, anchor=north] at (7.1,-0.65) {il sistema conta;
    \emph{Apocalypse World}, \emph{Nobilis}, \emph{Dread}, \emph{Fate}; Big
    Model; \emph{Fiasco}};
  \node[nota, anchor=north] at (10.2,-0.65) {\emph{Blades in the
    Dark}, \emph{Fabula Ultima}, \emph{Not the End}; \emph{The Quiet Year},
    \emph{Dialect}; due giocatori, solo, micro};
  \node[nota, anchor=north] at (13.3,-0.65) {meno regole, alta
    mortalità, sandbox; \emph{Mörk Borg}, \emph{Shadowdark}};
\end{tikzpicture}
\caption{Una ricostruzione schematica della ``piccola storiografia'' dei giochi
di ruolo da tavolo raccontata da Daniele. Le epoche sono indicative e si
sovrappongono (gli story game, per esempio, nascono anche prima di The Forge);
la storia completa è rimandata alla lezione sulla fiction.}
\label{a2:fig:storia-gdr}
\end{figure}

% ---------------------------------------------------------------------
\section{Giochi di ruolo per due, in solitario e micro}
\label{a2:sec:gdr-formati}

\slidefig{img/AGD/agd2-micro}{Giochi di ruolo per due, in solitario e nano-RPG.}{slide \emph{Analog -- Case Studies}}{fig:agd2-micro}


L'ultima parte sui giochi di ruolo torna a Ruggero, che presenta tre formati.

\begin{description}
  \item[Two player only.] Storicamente i giochi di ruolo nascono in gruppi di
        gioco ampi, ma c'è stata ``una grossissima wave'' di giochi di ruolo
        \textbf{per due persone}, e non nel senso che una fa il master e l'altra
        un personaggio: spesso l'\textbf{agency narrativa}, il potere di decidere
        che cosa succede, viene \textbf{passata meccanicamente} al tavolo o
        divisa in responsabilità diverse. Ce ne sono tantissimi, e per Ruggero
        sono molto interessanti anche perché, come si vedrà nella lezione
        sull'emotional design, permettono di fare cose molto \textbf{intime}: non
        serve un gruppo di gioco, si possono giocare con il partner o con
        un'altra persona.
  \item[Solo RPG.] Dal 2013 circa si sono normalizzati i \textbf{giochi di
        ruolo in solitario}: si gioca da soli, senza un'altra persona e senza
        fare da master a sé stessi. Spesso non c'è un master: ci sono
        \textbf{procedure} che guidano il gioco, e il manuale è il sistema di
        gioco, strutturato ``quasi come un board game'' in termini di regole. Ma
        resta tutto ciò che è tipico del gioco di ruolo: l'agency emergente,
        l'interpretazione di un personaggio, la creazione di una narrazione. Di
        solito si giocano in formato \textbf{diario} (\emph{journaling}) o in
        formato \textbf{audio}, e sono in grande tendenza.
  \item[Micro RPG.] Infine, i \textbf{micro giochi di ruolo}\footnote{La
        trascrizione dice ``handheld RPG''; dal contesto si tratta dei
        \emph{micro RPG}.}: giochi di ruolo su \textbf{biglietti da visita}, o
        raccolte in cui un libro ne contiene cento. Il \textbf{200 Word RPG}
        (\emph{200 Word RPG Challenge}) era un contest che chiedeva di scrivere un
        gioco di ruolo in \textbf{massimo 200 parole}, e ogni edizione ne
        raccoglieva migliaia: nel 2017, ricorda Ruggero, circa 900 giochi, su
        circa mille proposte.
\end{description}

% ---------------------------------------------------------------------
\section{I giochi da tavolo: i generi storici}
\label{a2:sec:tavolo}

\slidefig{img/AGD/agd2-rollmove}{Roll and move, astratti e party game.}{slide \emph{Analog -- Case Studies}}{fig:agd2-rollmove}


``Passerei ai giochi da tavolo, correndo velocemente e senza fare pause, se ce la
fate''. Come per i LARP e i giochi di ruolo, Ruggero parla di
\textbf{categorie dall'alto}, e per ognuna mostra come il genere si sia
\textbf{evoluto} nel tempo.

\subsection{Roll and move}

I \textbf{roll and move} si basano, come dice il nome, sul \textbf{tirare un
dado e muoversi}. Sono forse i più comuni nella storia dell'umanità, con quello
che da noi è il \textbf{gioco dell'oca}: un tracciato su cui si tirano i dadi e
si avanza. L'\textbf{agency} in questi giochi è ``praticamente inesistente'', e
il genere si è presto evoluto integrando altre meccaniche: con
\textbf{elementi economici} (come \emph{Monopoly}) o con \textbf{elementi
deduttivi}, come \emph{Cluedo}, che è un roll and move.

\subsection{Astratti}

Un altro genere storicamente gigantesco è quello degli \textbf{astratti}:
scacchi, dama, go, tutti quei giochi la cui componente tematica è ridotta
``quasi all'osso''. Gli scacchi mantengono ancora qualche forma di tema; il go, di
fatto, è un rettangolo di legno con delle caselle e delle pietre bianche e nere.
Anche qui il genere non si è fermato: ci sono iterazioni moderne come
\emph{ZÈRTZ} e i giochi del \textbf{progetto GIPF}\footnote{La trascrizione dice
``Zerz o il Gift Project''.}, e altri astratti contemporanei (uno è mostrato a
destra nella slide\footnote{La trascrizione non permette di ricostruirne il
titolo (``i borsa astratti, per esempio, di Supernova'').}), che mantengono
dinamiche ``molto asciutte, molto distillate'': \textbf{informazione aperta},
trasparente, niente carte imprevisti, niente mosse segrete. Diventano quasi una
\textbf{sfida d'intelletto} con l'avversario.

\subsection{Party game}

I \textbf{party game}: \emph{Taboo}, \emph{Lupus in Tabula} o \emph{Werewolf} a
seconda dell'edizione. Sono pensati per \textbf{molte persone} (di solito dalle
quattro in su), con regole ``molto asciutte, molto leggere'', perché si giocano
in contesti di festa, in cui non è detto che chi è nella stanza ami giocare o
abbia mai giocato. Poche regole, leggerezza, durata breve (spesso non più di
mezz'ora): sono \textbf{un possibile punto di accesso}. Anche questi si sono
evoluti, con giochi come \emph{Hitster}\footnote{La trascrizione dice ``Hit'';
la ricostruzione è probabile ma non certa.} e \emph{Wavelength}, che portano
sensibilità più moderne e meccaniche più interessanti del ``tirare alla cieca
per capire chi tra di noi è un lupo mannaro''.

\subsection{Giochi di parole}

I \textbf{word game}\footnote{La trascrizione dice ``war games'', ma gli esempi
(\emph{Boggle}, \emph{Scarabeo}) chiariscono che si tratta dei giochi di
parole.}: \emph{Boggle}, \emph{Scarabeo}, giochi in cui la meccanica centrale è
\textbf{creare parole}, e più la parola è lunga, complessa o usa lettere rare,
più punti si fanno. Anche qui le iterazioni sono interessanti: \emph{Cross
Clues}\footnote{La trascrizione dice ``Cross Cruise''.} trasforma il genere in
un gioco \textbf{cooperativo}, in cui non si compete per la parola più lunga ma si
gioca con le parole per capire ``quanta intesa c'è tra di noi''; e
\emph{Illiterati} è praticamente un gioco con le \emph{boss fight}, in cui si
sconfiggono i boss formando parole. Ibridazioni anche qui, perché chi crea giochi
li crea in base a ciò che ha giocato prima.

\subsection{Giochi di carte}

\slidefig{img/AGD/agd2-carte}{Giochi di parole, giochi di carte, TCG e deckbuilding.}{slide \emph{Analog -- Case Studies}}{fig:agd2-carte}


I \textbf{card game}, che in Occidente si diffondono storicamente con i mazzi
\textbf{francesi} e, in Italia, con quelli ``da briscola'' (ma ci sono anche altri
casi storici\footnote{La trascrizione cita i ``Mazda 1''; il riferimento non è
ricostruibile.}). Le evoluzioni:
\begin{itemize}
  \item i \textbf{micro card game} come \emph{Love Letter}, che in una decina o
        dozzina di carte racchiude un gioco complesso;
  \item giochi come \emph{Munchkin}, in una direzione più ampia, con il
        \textbf{take that} (``voglio mettere i bastoni tra le ruote'' agli altri);
  \item la ``gigantesca categoria'' dei \textbf{trading card game}, già
        menzionata nella lezione precedente: \emph{Magic: The Gathering},
        \emph{Yu-Gi-Oh!}, \emph{Pokémon}. Il \emph{trading} non è solo una
        dinamica tra persone: è nel nome della categoria. Scambiare carte
        \textbf{fuori dalla partita} è parte della dinamica comunitaria;
  \item i \textbf{Living Card Game} (LCG), che provano a risolvere uno dei
        principali problemi editoriali dei TCG: ``se io ho più soldi di te, tu
        contro di me a Magic non vinci'', perché comprerò carte da 800 o 900
        euro l'una. Gli LCG \textbf{non hanno booster casuali}: ogni espansione ha
        l'elenco preciso del contenuto. Restano giochi di carte con il
        \textbf{deck building pre-partita}, con mazzi ``fortemente identitari''
        (si possono scegliere le carte anche solo perché sono più simpatiche),
        ma senza l'aspetto economico della \textbf{rarità indotta};
  \item i \textbf{deck-building game}, in cui la costruzione del mazzo
        \textbf{non è più pre-partita ma avviene durante la partita}: spesso si
        parte con mazzi identici, e durante il gioco si acquisiscono carte che
        cambiano \textbf{permanentemente} il proprio mazzo, con lo scopo di
        costruire il mazzo migliore possibile.
\end{itemize}

\subsection{German game ed Eurogame}

\slidefig{img/AGD/agd2-euro}{German game, Eurogame e American.}{slide \emph{Analog -- Case Studies}}{fig:agd2-euro}


Altre categorie, che Ruggero mette tra virgolette, sono i \textbf{German game}:
un genere sviluppato in \textbf{Germania} dopo la metà degli anni '50--'60, con
giochi a \textbf{interazione indiretta}, competitivi ma ``senza troppo dar
fastidio direttamente all'altra persona'': non ci si può attaccare
direttamente. Sono giochi poco bellici, anche per \textbf{rigetto della Seconda
guerra mondiale}, e si è sviluppato moltissimo il filone \textbf{comunitario}, di
\emph{community building} e di scambio: costruire la propria fattoria, costruire
un villaggio, scambiarsi le risorse (``io ho delle risorse, tu altre,
scambiamocele''), senza distruggere ciò che l'altro ha costruito.

Il genere si è poi reso più complesso e si è espanso oltre la Germania con gli
\textbf{Eurogame}, che includono Spagna, Portogallo, Francia e anche l'Italia
(Ruggero ne ha fatti parecchi). L'interazione resta indiretta, ma i giochi
diventano molto più complessi: spesso durano \textbf{due o tre ore} e provano a
rendere \textbf{una simulazione focalizzata di qualcosa del mondo reale}. Nella
slide, a sinistra, un gioco in cui si gestisce una \textbf{galleria
d'arte}\footnote{La trascrizione dice ``The Gallery''; probabilmente \emph{The
Gallerist}.}, con tutte le dinamiche di compravendita di quadri; e
\emph{Endeavor}, che Ruggero ha alle spalle, sul commercio mondiale in epoca
moderna. Dinamiche che simulano un tipo di mercanteggiare: ``a volte compro e
vendo quadri, a volte compro e vendo vasi''.

\paragraph{Approfondimento: l'Eurogame secondo Brathwaite e Schreiber.} In
\emph{Challenges for Game Designers}, nel capitolo sul \emph{targeting} di un
mercato, gli Eurogame (``noti anche come \emph{German games} o \emph{designer
games}'') sono un caso studio di gioco per il mercato di massa: giochi per la
famiglia, godibili a tutte le età; \emph{I coloni di Catan} è il primo ad aver
reso popolare il genere in Nord America, seguito da \emph{Carcassonne} e
\emph{Puerto Rico}. Il libro ne dà una definizione operativa, utile per
confrontare le osservazioni di Ruggero:
\begin{itemize}
  \item \textbf{durata tra 20 e 90 minuti} (troppo brevi annoiano gli adulti,
        troppo lunghi i bambini);
  \item \textbf{setup breve}, perché il setup non contiene decisioni;
  \item \textbf{regole semplici}, curva di apprendimento breve;
  \item \textbf{cooperazione più che scontro}: scambi e negoziazioni invece di
        meccaniche distruttive in cui un giocatore danneggia direttamente un
        altro (il che non impedisce una competizione intensa);
  \item \textbf{decisioni strategiche con casualità misurata};
  \item \textbf{tempi morti minimi}: turni brevi, o azioni anche nei turni
        altrui.
\end{itemize}
La ``cooperazione invece dello scontro'' è l'interazione indiretta del German
game; la durata, invece, mostra bene l'evoluzione che racconta Ruggero: gli
Eurogame più complessi, da due o tre ore, sono ormai un'altra cosa rispetto al
gioco per famiglie descritto dal libro.

\subsection{American}

In contrapposizione, oltreoceano c'erano gli \textbf{American} (o
\emph{Ameritrash}): giochi molto focalizzati sulla \textbf{produzione} (le
miniature di plastica), con il \textbf{dungeon crawling}, tematicamente molto
vicini a \emph{Dungeons \& Dragons}, quasi versioni da tavolo del combattimento a
griglia normalizzato da D\&D. Continuano la tradizione di \emph{Chainmail} e dei
giochi strategici e tattici: sono spesso giochi di combattimento, a
\textbf{interazione diretta} (si attacca e ci si difende), e buona parte dei loro
regolamenti è dedicata alla simulazione del combattimento.

\subsection{Strategici e wargame}

\slidefig{img/AGD/agd2-strategici}{Strategici, movimento nascosto, destrezza, 18XX, tempo reale e legacy.}{slide \emph{Analog -- Case Studies}}{fig:agd2-strategici}


Un altro genere gigantesco è quello dei \textbf{giochi di strategia}, usati
storicamente anche nella \textbf{formazione militare}, ``non a caso'': la
\textbf{RAND Corporation} ne ha creati diversi, e ci sono editori specializzati.
Lo scopo era ricreare battaglie storiche. Anche qui, evoluzioni:
\begin{itemize}
  \item i giochi \textbf{a zone} o ad \textbf{aree} (\emph{area control}) con
        unità su una mappa, di cui \emph{Risiko} è il più famoso: io ho un
        esercito, tu ne hai un altro, dobbiamo conquistare territori, e il
        conflitto diretto è inevitabile perché gli obiettivi si sovrappongono;
  \item i \textbf{card-driven strategy game}: strategici con mappe e zone di
        controllo, ma con \textbf{mazzi di eventi}. I designer prendono eventi
        reali di conflitti reali e cercano di renderli coesi tematicamente con
        le meccaniche, dando a ogni evento una risoluzione precisa: ``l'importante
        non è tanto quante unità ho, ma quali carte evento ho in mano''. Sono,
        se ci si pensa, più vicini ai card game che agli strategici.
\end{itemize}
Nicola chiede di \textbf{\emph{Campaign for North Africa}}\footnote{La
trascrizione dice ``Campi for Nord Africa''.}: è uno \textbf{strategico
classico}, con esagoni e pedine di cartoncino, ``un'esperienza devastante in
termini di lunghezza'': forse non è il caso di cominciare da lì.

\subsection{Movimento nascosto}

Il genere degli \textbf{hidden movement}\footnote{La trascrizione dice ``Eat the
Movement''.}: \emph{Scotland Yard}, \emph{Letters from Whitechapel}, ma anche
altri titoli più recenti\footnote{La trascrizione cita ``Sniper Elite, Mind Dame
e GMT''; il secondo è probabilmente \emph{Mind MGMT}, GMT è un editore di
wargame.}. Sono giochi \textbf{asimmetrici}: di solito una persona si nasconde e
scappa, e le altre vincono collettivamente se la catturano. Uno contro tutti.

\subsection{Destrezza}

I \textbf{dexterity game} richiedono \textbf{destrezza fisica}: colpire
componenti con le schicchere, impilarli senza farli cadere. Non basta la
strategia che si ha in testa: ``posso anche avere la strategia migliore, ma se poi
le mie mani non ce la fanno a mettere in piedi questi cilindretti, crolla
tutto''.

\subsection{18XX}

Un genere che Ruggero ama particolarmente, e di cui è ``abbastanza sicuro''
che nessuno abbia mai sentito parlare, anche se conta centinaia di titoli (``più
di 400 giochi'' sotto l'etichetta): i \textbf{18XX}. Il nome viene dal secolo
delle \textbf{ferrovie}, l'Ottocento: sono tutti giochi ferroviari, ognuno
centrato su un anno e un luogo specifici (\emph{1849}, le ferrovie della
Sicilia nel 1849; \emph{1889}, quelle dell'isola di Shikoku). Hanno una forte
componente di \textbf{costruzione delle ferrovie} sulla mappa e una componente
\textbf{economico-azionaria}: comprare e vendere azioni, speculazione
finanziaria. Sono ``molto belli, ma anche molto intensi'': \emph{1889}, che
Ruggero usa di solito come introduttivo, ``dura tre ore''.

\subsection{Tempo reale}

Culturalmente si pensa ai giochi da tavolo come a qualcosa che avviene
\textbf{a turni}, ma non è sempre così. I giochi \textbf{in tempo reale} usano
clessidre, timer o altre meccaniche per cui non si ha tutto il tempo per pensare:
bisogna reagire in fretta. \emph{Kitchen Rush}, per esempio, è ``praticamente
\emph{Overcooked} da tavolo'': si gestisce un ristorante, e \textbf{la clessidra
è il lavoratore}, che non si può spostare finché non si è esaurita. Chi prende
gli ingredienti, chi li prepara sul fornello, chi va in sala a prendere gli
ordini, chi porta i piatti, chi li lava: ne nasce un certo caos nell'intreccio.

\subsection{Legacy}

I \textbf{Legacy} sono ``un po' una rivoluzione nel genere''. Di solito un gioco
da tavolo resta sempre quello: lo si prepara, lo si gioca, la settimana dopo lo
si gioca con un altro gruppo. I giochi Legacy invece chiedono di avere un
\textbf{impatto fisico sul gioco}, per cui \textbf{le partite successive sono
influenzate dalle precedenti}:
\begin{itemize}
  \item \emph{Pandemic Legacy} chiede di \textbf{strappare carte}; alcune città
        vengono distrutte dalla pandemia, coperte con degli \textbf{adesivi}, e
        non esistono più sulla mappa;
  \item \emph{SeaFall}\footnote{La trascrizione dice ``Sifol''.}, sulla scoperta
        di isole in un Ottocento alternativo: le isole hanno il nome vuoto, e chi
        le scopre le battezza; da lì in poi l'isola porterà quel nome per tutto
        il resto del gioco.
\end{itemize}
``So che qualcuno sta già sudando freddo all'idea di strappare carte o scrivere
sul tabellone'', ma sono ``esperienze emotivamente incredibili'': ``giocatevi dei
Legacy''.

\paragraph{Approfondimento: Risk Legacy e la rottura di un'assunzione.} Schell
dedica ai Legacy due passaggi di \emph{The Art of Game Design}. Nel capitolo sul
\textbf{tema}, presenta \emph{Risk Legacy} di \textbf{Rob Daviau} come esempio di
tema fondato su una verità: Daviau ha fatto ``qualcosa di mai tentato prima nel
mondo dei giochi da tavolo'', creando meccaniche per cui le scelte dei giocatori
\textbf{cambiano il gioco per sempre}. Le regole impongono di modificare il
tabellone con adesivi non rimovibili, di scriverci sopra con un pennarello
indelebile per rivendicare i territori, di strappare e scartare carte e perfino
di modificare permanentemente il regolamento. Queste meccaniche, nota Schell,
sono nuove e intriganti, ma soprattutto \textbf{rafforzano il tema centrale}:
``la guerra cambia il mondo''. Nel capitolo sulle tecniche di brainstorming,
Schell racconta come Daviau ci sia arrivato: elencando le \textbf{assunzioni}
date per scontate sui giochi da tavolo e provando a ``farne saltare'' una,
``una partita non influenza la successiva''. È un buon esempio di quello che
Ruggero chiedeva all'inizio: guardare un genere che si conosce e chiedersi che
cosa succederebbe se una delle sue regole implicite non valesse più.

% ---------------------------------------------------------------------
\section{Due libri sul gioco da tavolo}
\label{a2:sec:tavolo-libri}

\slidefig{img/AGD/agd2-libri}{I due libri consigliati sul gioco da tavolo.}{slide \emph{Analog -- Case Studies}}{fig:agd2-libri}


A proposito di risorse, Ruggero torna sul libro di Nuccio, \emph{La
progettazione dei giochi da tavolo}, citato a inizio lezione, e indica
un'alternativa in inglese che a suo parere ``fa quello che fa Nuccio, meglio'':
\textbf{\emph{Building Blocks of Tabletop Game Design}} di \textbf{Geoffrey
Engelstein e Isaac Shalev}\footnote{La trascrizione dice ``Engelstein e
Schaler'' e ``Building Blocks of Tablet of Game Design''.}. È più
\textbf{contemporaneo} negli esempi e ``più funzionale''. Per i giochi da tavolo,
se si ha già un framework di riferimento, sono entrambi ``una risorsa
incredibile''. (Resta valido l'avvertimento di inizio lezione: sono archivi di
meccaniche, da usare dopo aver capito che cosa le meccaniche fanno.)

% ---------------------------------------------------------------------
\section{Libri-gioco, megagame e altri ibridi}
\label{a2:sec:tavolo-ibridi}

\slidefig{img/AGD/agd2-ibridi}{Libri-gioco, megagame e ibridi tra gioco da tavolo, gioco di ruolo e videogioco.}{slide \emph{Analog -- Case Studies}}{fig:agd2-ibridi}


``Manca ancora qualcosa?'' Mancano diverse cose.
\begin{itemize}
  \item I \textbf{libri-gioco}, nel formato libro tradizionale e nel formato
        ``più contemporaneo'' dei giochi da tavolo che sono libri-gioco, come
        \emph{Sherlock Holmes Consulente Investigativo}.
  \item I \textbf{megagame}: giochi da tavolo \textbf{per cento persone}. Non
        hanno un'organizzazione editoriale come \emph{Monopoly} o gli altri
        board game: sono organizzati in modo più vicino ai giochi pervasivi e al
        LARP, spesso \textbf{a evento}, spesso con biglietto, anche perché
        gestire un gioco da tavolo con cento persone in parallelo ``può diventare
        molto intenso''.
  \item I \textbf{LARP urban game}, di cui ha parlato Beatrice, con esempi
        storici anche italiani, come \emph{Pathos}\footnote{Grafia come nella
        trascrizione.}: giochi di ruolo dal vivo che sono anche urban game, e
        un po' ARG.
  \item Produzioni ispirate al freeform ma vendute come giochi da tavolo, come
        \textbf{\emph{Inhuman Conditions}}\footnote{La trascrizione dice ``In
        Human Condition''.}: praticamente il \textbf{test di \emph{Blade Runner}}
        per scoprire chi è un androide, per due persone. È un gioco di ruolo, ma
        prodotto esattamente come un gioco da tavolo (carte, scatola,
        regolamento) e non si presenta come gioco di ruolo, ``ma questo non lo
        rende meno gioco di ruolo''.
  \item Ibridi con il party game, come \textbf{\emph{Sì, Oscuro Signore!}}: un
        party game, ma anche un \emph{role-playing game}, perché si fa il role
        play del proprio personaggio.
  \item Ibridi con il videogioco, come \emph{Fateforge} e gli altri giochi che
        usano un'\textbf{app} non solo come segnapunti ma come meccanica durante
        la partita.
\end{itemize}

\begin{table}[htbp]
\centering
\small
\begin{tabular}{@{}p{3.0cm}p{9.9cm}@{}}
\toprule
\textbf{Genere} & \textbf{Che cosa lo caratterizza ed evoluzioni citate} \\
\midrule
Roll and move & tira e muovi (gioco dell'oca); agency quasi nulla; poi
elementi economici (\emph{Monopoly}) o deduttivi (\emph{Cluedo}) \\
Astratti & tema ridotto all'osso, informazione aperta (scacchi, dama, go,
\emph{ZÈRTZ}, progetto GIPF) \\
Party game & molte persone, regole leggere, partite brevi, punto di accesso
(\emph{Taboo}, \emph{Lupus in Tabula}, \emph{Wavelength}) \\
Giochi di parole & creare parole per fare punti (\emph{Boggle},
\emph{Scarabeo}); cooperativi (\emph{Cross Clues}), boss fight
(\emph{Illiterati}) \\
Giochi di carte & mazzi tradizionali; \emph{Love Letter}; take that
(\emph{Munchkin}); TCG, LCG e deck-building \\
German game ed Eurogame & interazione indiretta, costruzione e scambio;
simulazioni focalizzate di due o tre ore \\
American & produzione in plastica, dungeon crawling, combattimento e
interazione diretta \\
Strategici & formazione militare; aree e zone (\emph{Risiko}); card-driven;
esagoni (\emph{Campaign for North Africa}) \\
Movimento nascosto & uno contro tutti, asimmetrici (\emph{Scotland Yard},
\emph{Letters from Whitechapel}) \\
Destrezza & schicchere, impilare: conta anche la mano \\
18XX & ferrovie e azioni dell'Ottocento; più di 400 titoli (\emph{1849},
\emph{1889}) \\
Tempo reale & clessidre e timer (\emph{Kitchen Rush}) \\
Legacy & le partite cambiano fisicamente il gioco (\emph{Pandemic Legacy},
\emph{SeaFall}) \\
Libri-gioco, megagame, ibridi & \emph{Sherlock Holmes Consulente
Investigativo}; cento giocatori a evento; \emph{Inhuman Conditions},
\emph{Sì, Oscuro Signore!}, giochi con app \\
\bottomrule
\end{tabular}
\caption{I generi di gioco da tavolo presentati da Ruggero, in sintesi.}
\label{a2:tab:tavolo}
\end{table}

% ---------------------------------------------------------------------
\section{La chiusura: niente interrogazioni, e i materiali}
\label{a2:sec:chiusura}

``E ho corso un po', ma ce l'abbiamo fatta''. Mancano gli ultimi dieci minuti, e
Ruggero si complimenta con chi è ``arrivato fino a qui''. Ripete il punto
fondamentale: \textbf{non ci sarà nessuna interrogazione} su questi giochi
(``Tiziana, ti do subito il sollievo''). Non si chiede di sapere che cosa sia ogni
caso studio: lo scopo era dare una \textbf{panoramica}, degli spunti, e far
sapere che ``questa roba esiste''. Le slide restano, e ci si potrà tornare in
fase di ricerca e sviluppo. Con la lezione di \textbf{ricerca e sviluppo}, che
darà gli strumenti per farlo in autonomia, ognuno troverà \textbf{i casi studio
che servono al proprio gioco}, ``e io ve li chiederò, noi ve li chiederemo''.
Questa lezione serviva ad \textbf{aprire le possibili direzioni}.

\paragraph{I materiali caricati.} Mentre Ruggero parla, le slide vengono
caricate nella cartella degli studenti. Beatrice aggiunge di aver caricato:
\begin{itemize}
  \item il paragrafo di \textbf{Frank Rose} su \emph{Why So Serious};
  \item un estratto sulla \textbf{definizione dei pervasive game} tratto da
        \emph{Game Design}, un libro di \textbf{Maresa Bertolo e Ilaria
        Mariani}\footnote{La trascrizione dice ``Marianne Bertolo'' e
        ``Bertolo Mariani''. Il libro è con ogni probabilità \emph{Game
        Design. Gioco e giocare tra teoria e progetto} di Maresa Bertolo e Ilaria
        Mariani.}, ``un ottimo libro''.
\end{itemize}
Non sono letture obbligatorie.

\paragraph{Un avvertimento sui libri italiani di gioco di ruolo.} Daniele
aggiunge una cautela sulla parte dedicata al gioco di ruolo del libro di Bertolo
e Mariani: alcuni paragrafi ``generalizzano molto verso il mainstream''. E poi un
\emph{warning} più generale: \textbf{la maggior parte della letteratura italiana
sul gioco di ruolo è ``terribile''}, al punto che alcuni libri i docenti non solo
non li consigliano, ma li \textbf{sconsigliano attivamente}. Con la premessa che
non è nulla di personale contro l'autore, Daniele sconsiglia il manuale
\emph{Giochi di ruolo} di \textbf{Mauro Longo}\footnote{La trascrizione riporta
l'editore come ``Dino Dino Editori'': probabilmente Dino Audino Editore.},
perché è molto soggettivo, molto legato alla visione del gioco di ruolo
dell'autore e a un solo tipo di gioco, quello di \emph{Lex Arcana}\footnote{La
trascrizione dice ``l'ex arcana''.}, con teorie proprie che non corrispondono al
design generale dei giochi di ruolo. ``Va bene se volete fare \emph{Lex Arcana}
esattamente come \emph{Lex Arcana}'', non se si vuole fare altro.

Meglio il volume sull'edu-larp, e soprattutto \textbf{giocare}: ``giocate i giochi,
perché in Italia abbiamo una creatività sfrenata''. Si trovano ``robe favolose''
online: tantissimi \textbf{one page game} di ruolo, il mondo del \textbf{print and
play} anche nei giochi da tavolo, chamber LARP e freeform da stampare
direttamente da internet. ``Fatevi una ricerca e vedete quello che vi garba''.

% ---------------------------------------------------------------------
\section{Le domande finali}
\label{a2:sec:domande}

``Siamo stati più veloci dell'anno scorso'' (quando si erano accumulati trenta
minuti di ritardo), e restano sette minuti per le domande.

\subsection{Che cosa fa scegliere una forma?}

Gianmarco, ``un po' per rompere il ghiaccio'', fa una domanda ``molto generica'',
che sa non avere una risposta semplice: si sono visti giochi dalle forme
diversissime, ma i \textbf{framework} per crearli sono più o meno gli stessi. Che
cosa determina la scelta di una forma rispetto a un'altra? Ruggero rimanda la
risposta alla lezione sull'\textbf{emotional design}, ``che è proprio una risposta
a questa cosa'': la risposta è così lunga che ``abbiamo una lezione di tre ore per
rispondere alla tua domanda''.

\subsection{Che cosa va per la maggiore in Italia?}

Angelo chiede quale genere ``vada per la maggiore'' in Italia. Ma che cosa
significa ``andare per la maggiore''? Il successo commerciale non è uguale per
tutti i generi, ``nemmeno nei videogiochi''. Uno dei docenti (con ogni probabilità Daniele) risponde con due domande
``provocatorie'', che valgono anche per i libri (ne parlava quel giorno con
un'amica narrative designer di videogiochi):
\begin{enumerate}
  \item quante copie deve vendere un gioco per essere un successo?
  \item il numero di copie vendute coincide con il successo
        commerciale ed economico?
\end{enumerate}
La risposta a entrambe è \textbf{no}, ma si capirà con la lezione di
\textbf{Flavio} sull'aspetto produttivo: quando si ragiona in logica commerciale
si è ``in logica di prodotto'', e ci sono varie \textbf{metriche} e varie scale di
successo. A questo si aggiungono altri livelli:
\begin{itemize}
  \item \textbf{quale mercato}? In Italia \textbf{i giochi che non sono in
        italiano non vendono}, e molti editori italiani sono innanzitutto
        \textbf{editori di localizzazione}: prendono giochi dall'estero, li
        traducono, e questo è buona parte della loro sopravvivenza, nel gioco
        da tavolo come nel gioco di ruolo. Un gioco in inglese si vende in
        Inghilterra, in America e in moltissimi paesi dove l'inglese è la seconda
        lingua; un gioco in italiano si vende solo in Italia, e un editore
        italiano non ha facilmente la struttura per arrivare all'estero. Per
        questo anche autori italiani di giochi da tavolo pubblicano con
        \textbf{editori esteri};
  \item \textbf{perché si fa il gioco}? Non è detto che lo scopo sia
        commerciale. Un gioco può essere \textbf{formativo o didattico}, fatto
        con un'azienda o una scuola, o \textbf{commissionato} da un cliente per
        uno scopo specifico: allora non interessa che venda mille copie. Il
        designer viene pagato, ma il pagamento non corrisponde alle vendite.
\end{itemize}

Beatrice chiede ad Angelo perché faccia la domanda, perché i \textbf{settori}
appena visti spesso ``non si parlano minimamente'', e alcuni non sanno nemmeno
dell'esistenza degli altri: la cross-medialità presentata nella lezione non viene
presentata spesso così. Se la domanda nasce da prospettive lavorative, come per
alcuni nel corso, la risposta è che ci sono ``immense possibilità'' a seconda
della direzione: fare un gioco da tavolo è una cosa, lavorare nella
\textbf{formazione} un'altra, e c'è un intero ambito non trattato, quello dei
\textbf{giochi per le aziende}.

Angelo precisa: voleva solo sapere \textbf{che cosa giocano di solito gli
italiani}. È ``un altro volo pindarico'': in Italia è difficile trovare dati
statistici affidabili, perché i questionari e le survey sono spesso invalidati
dal modo in cui sono formulate le domande. Da quello che sa il docente che risponde\footnote{La trascrizione dice ``poi Lucio mi correggerà''; il nome non è ricostruibile con certezza.}, ci sono sicuramente \textbf{più giocatori di giochi da
tavolo che di giochi di ruolo}: il gioco da tavolo è più conosciuto e
frequentato, anche nelle ludoteche, mentre il gioco di ruolo e il LARP richiedono
un investimento economico diverso e sono soprattutto pratiche
\textbf{comunitarie}, ``tribali'', che nascono dal basso e si trasmettono
\textbf{oralmente}. Chi gioca ai giochi da tavolo li compra; chi gioca di ruolo
spesso lo fa perché qualcuno lo fa giocare, e non è detto che abbia il gioco in
casa.

\subsection{Avvisi}

Intanto Daniele ricorda che il \textbf{questionario} della prima slide va
compilato \textbf{entro il 15}, e che Giordano deve ancora presentarsi (``siamo
delle pessime persone''). Visto che è ora di cena, Giordano si presenterà come
\textbf{prima cosa} alla lezione successiva, \textbf{lunedì}, sull'emotional
design (``ricordacelo, perché se no ci dimentichiamo'').

% ---------------------------------------------------------------------
\section{Due cimiteri: i giochi di carte collezionabili e i ``fix''}
\label{a2:sec:cimiteri}

Prima di chiudere, a un docente torna in mente ``una cosa importantissima''
che voleva dire.

\subsection{Non fate Magic: The Gathering 2}

``Vi prego, non cadete nella trappola'' di vedere un gioco di successo commerciale
e farne uno simile pensando che venderà moltissimo. ``Ci possiamo fare un
cimitero, una fossa comune'' con migliaia di \textbf{trading card game} e
\textbf{collectible card game}. ``Non fate un trading card game. Ve lo dico per
voi: non fate \emph{Magic: The Gathering 2}''. Chi vede \emph{Magic} o
\emph{Pokémon} pensa di poter fare un sacco di soldi, ma \textbf{non vede la
struttura commerciale e produttiva gigantesca} che c'è dietro, quasi impossibile
da emulare dal basso. Quasi tutti, anche editori grossi, quelli che hanno provato
a competere con \emph{Pokémon}, \emph{Yu-Gi-Oh!} e \emph{Magic} ``sono andati a
picco''. ``Se volete perdere soldi, quello è il modo perfetto''.

Un collega aggiunge che c'è un \textbf{revival dei TCG}, ma ogni titolo dura un
anno o due e poi muore. Si sperimentano anche dei \textbf{micro TCG} annunciati
fin dall'inizio come limitati (``questo gioco avrà tre espansioni e poi
chiuderà''), una sorta di ibrido con gli LCG. Ma attenzione: spesso sono fatti
da ex designer di LCG, di \emph{Magic} o di \emph{Hearthstone}, persone con
contatti, entourage e un'azienda alle spalle; da fuori si vede una persona, ma
dietro c'è una struttura che non si può replicare.

\paragraph{Approfondimento: perché Pokémon non si replica.} \emph{The Art of
Game Design} offre un buon esempio della ``struttura gigantesca'' di cui parlano
i docenti. Nel capitolo sui mondi transmediali, Schell racconta che il gioco di
carte collezionabili di \emph{Pokémon} nacque nel 1999 da una collaborazione tra
Nintendo e Wizards of the Coast, la casa di \emph{Magic}, e che fu progettato per
restare il più possibile fedele alle meccaniche del videogioco per Game Boy. Il
gioco di carte era la \textbf{terza porta d'ingresso} in un unico mondo, accanto
al videogioco e alla serie TV: portatile, sociale, economico, perfetto per lo
slogan ``acchiappali tutti''. Le tre porte si rafforzavano a vicenda (voler
vincere spingeva a guardare la serie, e la serie rendeva il gioco più vivido), e
Schell ricorda un dirigente dell'intrattenimento che nel 1999 profetizzava la
fine della ``moda'' con l'uscita del film. Non andò così, appunto, perché il
successo non stava in un singolo prodotto ma nell'ecosistema. Un gioco di carte
collezionabili nato dal basso, senza un ecosistema del genere, compete con questo.

\subsection{Non fate un gioco per ``aggiustarne'' un altro}

Un altro docente (probabilmente Daniele, che cita ``le lezioni di framework con Ruggero'') aggiunge un altro cimitero: quello di chi vuole fare un gioco per
\textbf{``fixare'' un gioco esistente}, cioè farne una versione migliore. Come
esperimento può essere interessante, ma, come si vedrà nella lezione
sull'emotional design e nelle lezioni sui framework con Ruggero, l'intenzione deve
essere ``\textbf{eticamente onesta}''. Non si vuole aggiustare D\&D per il piacere
di aggiustare D\&D: bisogna chiedersi sempre \textbf{che cosa non mi piace
dell'esperienza} di D\&D, che cosa voglio aggiustare di quell'esperienza. E non
creare ``D\&D dove tolgo una regola e ne metto un'altra, ed è sistemato'', ma
\textbf{partire dall'esperienza che si vuole raggiungere}, quella che non si è
avuta da D\&D, e costruirla. Con la lezione di ricerca e sviluppo si vedrà anche
come prendere spunto dai giochi esistenti, non per copiarli ma per capire che
cosa fanno, perché \textbf{prenderne pezzi senza capire che cosa fanno} è proprio
l'errore da evitare (lo stesso del libro di Nuccio letto troppo presto).

\paragraph{I saluti.} Un docente deve scappare e saluta; gli altri ricordano il
questionario, che le loro email sono nelle slide per eventuali domande, e
l'appuntamento a lunedì.

% ---------------------------------------------------------------------
\section*{Punti chiave del capitolo}
\addcontentsline{toc}{section}{Punti chiave del capitolo}

\begin{itemize}
  \item Un \textbf{caso studio} è un gioco, o una categoria di giochi, che,
        analizzato tenendo conto del suo contesto, dà informazioni e metodi
        utili per i propri giochi: come leggere libri per imparare a scrivere.
        Questa lezione è una panoramica; gli strumenti di analisi arriveranno con
        la ricerca e sviluppo. ``Per me, questo è design'' (Ruggero).
  \item \textbf{Gioco analogico} è un termine ombrello per tutto ciò che non è
        videogioco. Le quattro \textbf{macrocategorie} sono gioco da tavolo,
        gioco di ruolo da tavolo, gioco di ruolo dal vivo e gioco pervasivo, ma
        i confini sono porosi: \textbf{ibridazione e sperimentazione}. Il corso è
        \textbf{agnostico di genere}, e il project work è libero, ma non deve
        essere un gioco enorme.
  \item Il \textbf{gioco pervasivo} usa luoghi e oggetti della vita quotidiana
        come luoghi e oggetti di gioco; in termini di game studies,
        \textbf{espande il cerchio magico} nello spazio, nel tempo e nella
        società. Casi: \emph{Pokémon Go}, geocaching, \emph{GuruShot},
        \emph{Journey to the End of the Night}, SF0, \emph{Critical City Upload},
        \emph{Sleep No More}, gli ARG (\emph{Why So Serious}, \emph{The
        Institute}, \emph{Stranger Things} a Milano), Collettivo Cinetico,
        Amigdala, \emph{MP3 Experiment}, \emph{PacManhattan}, \emph{Zombies,
        Run!}, \emph{Reptilarium}, \emph{Echi resistenti}.
  \item Esperienze emotivamente significative si possono fare \textbf{con budget
        zero}: l'audio che guida, gli oggetti da portare, un finale
        scenografico.
  \item Il \textbf{LARP} è separato dal GDR per ragioni di industria (evento con
        biglietto contro oggetto venduto). Tipi: boffer, Mind's Eye Theatre,
        chamber/parlor, American freeform, murder party, Nordic LARP (Jeepform,
        360°, tape, black box).
  \item Il \textbf{Nordic LARP} ha documentato tutto (Knutpunkt dal 1997, libri
        annuali, manifesti come il Dogma 99), punta su coinvolgimento, emozione,
        coesione tematica e dopo-gioco, e ha normalizzato le \textbf{meccaniche di
        sicurezza}. Non è un esperimento scientifico: lo scopo resta ludico.
  \item Nel \textbf{GDR da tavolo} contano la fiction emergente e quattro parole:
        \textbf{conflitto, azione, risoluzione, effetto}. Il master è un
        giocatore. Dalla nascita di D\&D (1974) dal wargame, attraverso gli
        adattamenti a nastro degli anni '80--'90, si arriva alla \textbf{New
        Wave} di The Forge (``il sistema conta''), ai \textbf{giochi
        focalizzati} (non semplificati), all'OSR, agli story game, ai giochi per
        due, in solitario e micro.
  \item Nel \textbf{gioco da tavolo} ogni genere si è evoluto: roll and move,
        astratti, party game, giochi di parole, giochi di carte (TCG, LCG,
        deck-building), German game ed Eurogame, American, strategici,
        movimento nascosto, destrezza, 18XX, tempo reale, \textbf{Legacy},
        libri-gioco, megagame, ibridi.
  \item Libri: il volume sull'\textbf{edu-larp} sì; Nuccio e Engelstein--Shalev
        dopo l'emotional design; Bertolo--Mariani con cautela sulla parte di GDR;
        gran parte della letteratura italiana sul GDR no. Soprattutto,
        \textbf{giocare}.
  \item Il \textbf{successo} non si misura solo in copie, e dipende dal mercato
        (in Italia vende solo ciò che è in italiano) e dallo scopo del gioco.
  \item Due cimiteri: \textbf{non fare un gioco di carte collezionabili} e non
        fare un gioco per \textbf{``aggiustarne'' un altro} senza partire
        dall'esperienza che si vuole creare.
\end{itemize}

%
%  Tutte le etichette di questo capitolo hanno il prefisso "a2:"
%  (a2:cap:..., a2:sec:..., a2:fig:..., a2:tab:...). Il capitolo è
%  autonomo: non contiene \ref verso altri capitoli.
%  Le figure e le tabelle sono ricostruzioni schematiche di quanto i
%  docenti spiegano a voce e con le slide, non copie delle slide.
% =====================================================================

\chapter[I casi studio dell'Analog Game Design]{I casi studio: giochi
pervasivi, giochi di ruolo dal vivo, giochi di ruolo da tavolo e giochi da
tavolo}
\label{a2:cap:casi}

Questa è la \textbf{seconda lezione} del corso di \textbf{Analog Game Design}, e
mantiene la promessa fatta alla fine della prima: è una lezione di tre ore
dedicata quasi interamente ai \textbf{casi studio} (\emph{case studies}). I
docenti la definiscono ``un overview'', una panoramica ``a volo d'aquila'':
``voleremo in tutti i settori ludici analogici e guarderemo che cosa c'è''.
Non si scende nei dettagli tecnici dell'analisi (ci saranno lezioni apposite,
soprattutto quella di \textbf{ricerca e sviluppo}), ma si apre la ``diga'' degli
esempi: decine di giochi, eventi, esperienze, movimenti e manifesti, scelti per
mostrare quante cose esistono, quanto sono diverse tra loro e quanto i confini
tra i generi siano porosi.

La lezione si svolge \textbf{online}, su Zoom, un \textbf{venerdì} sera (``che
bella la lezione del venerdì! Sembra quasi un aperitivo''), ed è divisa tra i
tre docenti principali del corso, ciascuno sul proprio territorio:
\textbf{Ruggero} introduce l'argomento e presenta i \textbf{giochi di ruolo dal
vivo} (LARP) e i \textbf{giochi da tavolo}; \textbf{Beatrice} presenta i
\textbf{giochi pervasivi}; \textbf{Daniele} presenta i \textbf{giochi di ruolo
da tavolo}. ``Ce lo smezziamo un po' tutti'', perché l'argomento è diviso ``in
tante categorie, sottocategorie, sotto-sottocategorie''. Dopo i pervasivi c'è
una pausa di pochi minuti; gli ultimi dieci minuti sono lasciati alle domande.

Tre avvertenze di lettura, le stesse del capitolo precedente. La prima: la
trascrizione è automatica e storpia moltissimi nomi propri, soprattutto i titoli
dei giochi, quasi sempre in inglese e pronunciati velocemente (``Blaze in the
Dark'', ``Cobo Bobby'', ``Sifol'', ``Jeep Form''\ldots). Dove la ricostruzione è
sicura, o molto probabile, il capitolo usa la forma corretta e lo segnala in
nota; dove non lo è, riporta la forma della trascrizione e lo dice. La seconda:
durante la pausa la trascrizione registra, ripetute in loop per decine di righe,
alcune frasi senza senso (anche all'inizio della registrazione succede qualcosa
di simile), che il capitolo ignora. La terza: i docenti si passano spesso la
parola senza che la trascrizione segni il cambio di voce; il capitolo attribuisce
le parole a chi le pronuncia quando il contesto lo rende chiaro (chi parla di sé,
chi viene chiamato per nome, chi ``passa la parola''), e resta generico
altrimenti.

Quanto alla data, la lezione si colloca con ogni probabilità \textbf{venerdì 11
settembre 2026}\footnote{La data non è mai detta. La si deduce da alcuni
indizi: la prima lezione si era chiusa con ``ci si vede questo venerdì''; uno
studente assente ha recuperato ``la lezione di lunedì''; a fine lezione si
ricorda che il questionario iniziale va compilato ``entro il 15'' (settembre) e
che la lezione successiva, sull'\emph{emotional design}, sarà di lunedì. Nel
2026 il primo venerdì utile prima del 15 settembre, dopo il lunedì della prima
lezione, è l'11. È una deduzione, non un'informazione esplicita.}: i riferimenti
a ``l'altra volta'' e ``l'anno scorso'' vanno letti rispetto a quel momento.

\paragraph{Come leggere questo capitolo.} Più ancora della lezione precedente,
questa è un \textbf{catalogo}. I docenti lo dicono esplicitamente: non ci sarà
nessuna interrogazione sui singoli giochi, e lo scopo non è ricordarli tutti ma
sapere che esistono, avere dei punti di riferimento e, soprattutto, abituarsi a
guardare oltre il proprio genere preferito. Per questo il capitolo riporta
\emph{tutti} i casi citati, con quello che i docenti ne dicono, ma li organizza
in tabelle e schemi per renderli consultabili ``in fase di ricerca e sviluppo'',
esattamente come fa Ruggero con le sue slide (``non sapete quante volte mi sono
trovato a riaprire proprio esattamente queste slide''). Qua e là, dei paragrafi
di \textbf{approfondimento} collegano i casi ai libri di game design di
riferimento (\emph{The Art of Game Design} di Jesse Schell, \emph{Challenges for
Game Designers} di Brenda Brathwaite e Ian Schreiber, \emph{21st Century Game
Design} di Chris Bateman e Richard Boon).

Il percorso del capitolo segue quello della lezione:
\begin{enumerate}
  \item l'avvio, Discord, i materiali e i \textbf{libri consigliati} (e
        sconsigliati) per cominciare (sezioni~\ref{a2:sec:avvio}
        e~\ref{a2:sec:libri});
  \item \textbf{perché} analizzare i casi studio e \textbf{che cosa} sono
        (sezione~\ref{a2:sec:perche});
  \item le \textbf{quattro macrocategorie} del gioco analogico, l'ibridazione e
        l'invito a non chiudersi nel proprio genere (sezioni~\ref{a2:sec:macro}
        e~\ref{a2:sec:aperti});
  \item i \textbf{giochi pervasivi}, con Beatrice: la definizione, il cerchio
        magico e diciotto casi studio, dalle app agli \emph{alternate reality
        game}, dal teatro interattivo alle cacce al tesoro
        (sezioni da~\ref{a2:sec:pervasivi} a~\ref{a2:sec:pervasivi-edu});
  \item i \textbf{giochi di ruolo dal vivo}, con Ruggero: boffer, Mind's Eye
        Theatre, chamber e parlor LARP, American freeform, murder party, il
        \textbf{Nordic LARP} con i suoi manifesti, lo Jeepform, il 360°, i tape
        e i black box LARP (sezioni da~\ref{a2:sec:larp} a~\ref{a2:sec:oppresso});
  \item i \textbf{giochi di ruolo da tavolo}, con Daniele: le quattro parole
        (conflitto, azione, risoluzione, effetto), la storia dal primo
        \emph{Dungeons \& Dragons} alla New Wave di The Forge, l'OSR, i giochi
        focalizzati, gli story game, i giochi per due, in solitario e i micro
        giochi di ruolo (sezioni da~\ref{a2:sec:gdr} a~\ref{a2:sec:gdr-formati});
  \item i \textbf{giochi da tavolo}, di nuovo con Ruggero: diciassette generi,
        dal \emph{roll and move} ai Legacy, e gli ibridi
        (sezioni da~\ref{a2:sec:tavolo} a~\ref{a2:sec:tavolo-ibridi});
  \item la chiusura, i materiali caricati, le \textbf{domande finali} e due
        avvertimenti importanti: non fare un gioco di carte collezionabili, non
        fare un gioco per ``aggiustare'' un altro gioco
        (sezioni da~\ref{a2:sec:chiusura} a~\ref{a2:sec:cimiteri}).
\end{enumerate}

% ---------------------------------------------------------------------
\section{L'avvio della lezione}
\label{a2:sec:avvio}

La registrazione parte mentre i docenti si stanno ancora collegando (``mamma
mia, la burocrazia italiana'') e si controlla che il simbolo REC lampeggi.
Mentre Daniele prova a condividere le slide (``aspettate, sono un attimo
bloccato'') e chiede di essere nominato co-organizzatore della riunione, gli
studenti entrano uno alla volta: Bea, Lorenzo, Nicola, Rita, Gianmarco,
Tiziana, Giordano. L'atmosfera è rilassata. Si scherza sul fatto che la lezione
del venerdì ``sembra quasi un aperitivo'': un docente si mette ``bello comodo'',
un altro ha a casa ``le cose per fare l'aperitivo'' e chiede il permesso di
prepararlo nel frattempo, ``purché siate sul pezzo''. Daniele scrive anche di
dover ancora rispondere a un messaggio di Beatrice; ``non preoccuparti'', gli
risponde lei: settembre è un mese di ripartenza, ``in realtà è l'inizio
dell'anno reale''.

\paragraph{Giordano.} Arriva \textbf{Giordano}, che era assente alla prima
lezione. Ha recuperato la registrazione di lunedì, e sa che gli si chiederà di
presentarsi. (Alla fine, per mancanza di tempo, la presentazione verrà rinviata
alla lezione successiva: si veda la sezione~\ref{a2:sec:domande}.) All'inizio
mancano ancora Lorenzo e Tancredi; Lorenzo arriva poco dopo, Tancredi dopo la
pausa.

\paragraph{Discord.} Uno dei docenti nota che gli studenti sono ``belli attivi
su Discord'' e se ne rallegra: ``sfruttatelo, perché l'anno scorso lo si
sfruttava poco e invece è uno spazio molto importante''. I docenti leggono il
server ogni tanto, e le domande sono benvenute. L'unica richiesta è
\textbf{taggarli}: ``tra tanti server Discord non riusciamo a leggere'' tutto,
quindi una domanda rivolta a loro va segnalata con un tag. Lo spazio, ricorda,
``è fatto per voi'': per confrontarsi, condividere cose e materiali.

\paragraph{Le dispense.} Gli studenti hanno visto i materiali caricati: le
slide, le dispense ``e quant'altro''. Qualcuno chiede se le dispense siano in
inglese e se siano articoli. Sono i \textbf{saggi citati nella prima lezione}
(per esempio quelli sull'MDA e sul DPE), e \textbf{non è obbligatorio leggerli}:
sono lì ``per completezza'' e per mostrare che dietro gli argomenti delle
lezioni ``c'è una base teorica''. L'importante è che le lezioni facciano capire
gli argomenti. Spesso i docenti citeranno \textbf{paper accademici} di game
design e di \emph{game studies}, e ``purtroppo c'è poco da fare'': molti non sono
mai stati tradotti, o lo sono stati ``in maniera un po' approssimativa''. Per
questo di solito si linka la fonte originale; nulla vieta di cercare traduzioni
o di tradurseli per conto proprio. In ogni caso i docenti cercano di scegliere
paper \textbf{accessibili}: ``non vi linkeremo delle tesi di master di 300
pagine, a meno che non ce le chiediate voi''.

% ---------------------------------------------------------------------
\section{Due libri per cominciare (e uno da rimandare)}
\label{a2:sec:libri}

Prima di entrare nel tema, Daniele torna su due libri citati la volta
precedente, mostrandoli in video.

\subsection{Un volume sull'edu-larp}

A Tiziana, che ha qualche difficoltà con l'inglese, Daniele consiglia un libro
italiano in cui si trovano molti degli argomenti del corso: un volume
sull'\textbf{edu-larp}, il gioco di ruolo dal vivo in ambito educativo, a cura
di \textbf{Andrea Maragliano} ed edito da \textbf{FrancoAngeli}\footnote{La
trascrizione riporta il titolo come ``EduLARP'', ``Edu LARP'' e ``DualArp'', e
il nome del curatore come ``Andrea Maragliano''; il titolo completo del volume
non è detto a lezione.}. Con una frecciata affettuosa, Daniele aggiunge che
contiene ``un'introduzione molto interessante scritta da un certo Ruggero. Non
so chi sia''. Ruggero protesta (``non fare queste cose, che poi sembra che
consiglio i libri in cui ho partecipato''): il libro non è suo, ci hanno scritto
``tutti i miei amici'', e lui per principio non consiglia i libri su cui scrive
(``possiamo anche comprarlo e non leggere l'introduzione'', ride un collega).

Daniele però lo consiglia convintamente:
\begin{itemize}
  \item molti concetti che si trovano nei siti in inglese ``qua dentro sono
        molto spiegati''; lui stesso lo ha usato per studiare;
  \item c'è una parte, ``bellissima'', sull'\textbf{involvement} (il
        coinvolgimento), che spiega anche ``perché non usiamo la parola
        immersione'' (una delle parole vietate della prima lezione); se ne
        parlerà anche con Ruggero;
  \item finalmente si può consigliare ``roba genuina, sana, italiana di game
        design'';
  \item copre l'ambito \textbf{didattico e formativo}, cioè il gioco nelle
        scuole e il gioco come strumento di formazione, che interessa ad alcuni
        studenti del corso.
\end{itemize}
Non bisogna leggerlo tutto: è utile ``leggerlo pian pianino mentre fate il
corso'', se si vuole approfondire.

\subsection{La progettazione dei giochi da tavolo di Walter Nuccio}

Il secondo libro mostrato la volta precedente è \emph{La progettazione dei
giochi da tavolo} di \textbf{Walter Nuccio}. Qui Daniele fa un ``piccolo
cesellamento'': ``se fossi in voi non lo comprerei per il momento'', e lascia la
parola a Ruggero, che ne ha una copia alle spalle. Il giudizio di Ruggero è
preciso:
\begin{itemize}
  \item è più utile \textbf{dopo} aver interiorizzato alcune basi su come si
        progettano i giochi; come primo libro è un rischio, ``perché sembra che
        ti dia delle scorciatoie'';
  \item più che un saggio di game design, è \textbf{``un grosso archivio di
        meccaniche di gioco''}, una sorta di dizionario (``che cos'è un malus,
        per esempio?''), con una parte di \emph{pattern};
  \item il pericolo è applicare quelle meccaniche ai propri giochi
        \textbf{senza capire che cosa fanno}: ``come se apriamo una cassetta
        degli attrezzi e non sappiamo che cosa dobbiamo fare. Gli attrezzi
        sembrano tutti bellissimi, però magari li usiamo male'';
  \item il consiglio è aspettare a comprarlo, o almeno a leggerlo, fino a dopo
        la lezione sull'\textbf{emotional design}, quando si avrà già un'idea di
        come partire con il proprio gioco.
\end{itemize}
Daniele è ``assolutamente d'accordo''. (Nel frattempo si è spostato in cucina:
``benvenuti nella mia cucina, davanti al mio frigorifero, per parlare di user
experience''.) Del libro di Nuccio, e di un'alternativa in inglese, si
riparlerà a proposito dei giochi da tavolo (sezione~\ref{a2:sec:tavolo-libri}).

% ---------------------------------------------------------------------
\section{Perché analizzare i casi studio}
\label{a2:sec:perche}

Ruggero apre l'argomento del giorno, i \textbf{casi studio}, con due domande
alla classe: \emph{perché} ci teniamo ad analizzarli? E prima ancora: che cosa
ci si aspetta che siano? ``Non è una domanda provocatoria''.

\paragraph{Le risposte degli studenti.} Una prima risposta: sono \textbf{altri
giochi già realizzati}, da analizzare da molti punti di vista, ``anche
insoliti'', non solo da giocatori: dal punto di vista della produzione, o
dell'esperienza utente che se ne è avuta, per capire come migliorarli. Una
seconda risposta: i casi studio sono delle \textbf{milestone}, giochi o
situazioni ``lapidarie'' da cui trarre ``figure di merito'' da sviscerare per
poi generare altre cose. Tiziana chiede invece una cosa ``più basica'': che cosa
si intende esattamente? Modelli, esempi, ma di che cosa?

\paragraph{La definizione.} ``Ottima domanda''. Con \textbf{caso studio} si
intende un \textbf{gioco specifico}, o una \textbf{categoria di giochi}, che,
analizzato, ci dà delle \textbf{informazioni utili}: osservando da vicino che
cosa un gioco fa, possiamo estrapolare informazioni da riproporre, o metodologie
che ci servono per creare i nostri giochi. L'analogia è con la scrittura: chi
vuole scrivere un libro legge tanti libri, e capisce che cosa fanno di positivo
quelli che gli piacciono e che cosa non fanno quelli che non gli piacciono, per
evitare le cose che non gli piacciono e ``iterare sulle cose che ci piacciono''.
Con i giochi è lo stesso: si guarda che cosa fanno di positivo e di negativo,
\textbf{tenendo conto anche del contesto} in cui sono nati.

\paragraph{Oggi la panoramica, poi gli strumenti.} Gli strumenti più tecnici
per analizzare nel dettaglio arriveranno con le lezioni di \textbf{ricerca e
sviluppo}. Oggi si fa una ricognizione di tutti i settori.

% ---------------------------------------------------------------------
\section{Il gioco analogico e le sue quattro macrocategorie}
\label{a2:sec:macro}

\subsection{Un termine ombrello}

Prima una precisazione. Nel corso, \textbf{gioco analogico} è un
\textbf{termine ombrello} che copre i giochi da tavolo (\emph{board game}), i
giochi di carte (\emph{card game}), i giochi di ruolo da tavolo, i giochi di
ruolo dal vivo, i giochi pervasivi, le cacce al tesoro, le \emph{escape room},
i \emph{murder party}: tutti quei giochi che ``generalmente non sono focalizzati
e incentrati sulla programmazione, sull'uso di un terminale, sull'uso di uno
schermo''; in breve, che \textbf{non sono videogiochi}. Esistono cose ibride,
che si vedranno, ma l'area di interesse principale resta il gioco analogico.

\subsection{Le quattro macrocategorie}

\slidefig{img/AGD/agd2-categorie}{Le macrocategorie del gioco analogico: giochi da tavolo, giochi di ruolo da tavolo, LARP e, dietro, l'ibridazione.}{slide \emph{Analog -- Case Studies}}{fig:agd2-categorie}


La lezione parte dal ``frame attuale'': quattro \textbf{macrocategorie}, a loro
volta divise in sottocategorie. Ruggero avverte subito che ``non sono separate da
muri di acciaio non comunicanti'': i casi ibridi, difficili da mettere in una
sola categoria, sono frequenti. I tre docenti presentano ciascuno la propria.

\begin{description}
  \item[Gioco da tavolo (Ruggero).] La categoria dei \emph{board game}. ``Per
        motivi di industria'' vi si includono anche i giochi di carte e i
        \emph{party game}: tutti i giochi \textbf{pubblicati e prodotti
        all'interno dell'industria del gioco da tavolo}, che abbiano o no un
        tabellone.
  \item[Gioco di ruolo da tavolo (Daniele).] GDR in italiano, TTRPG
        (\emph{TableTop Role-Playing Game}) in inglese. È una categoria ``molto
        particolare'', perché anzitutto \textbf{non c'è una definizione
        univoca}: sul mercato si trovano cose molto diverse, perfino giochi che
        sono quasi ``dei LARP da tavolo'', vicini ai \emph{chamber LARP}. Anche
        i \textbf{formati} sono diversissimi: la slide mostra \emph{Dungeons \&
        Dragons}; \emph{Neverland}, un supplemento\footnote{Daniele lo chiama
        ``un'espansione per D\&D''; più avanti lo descrive come una
        reinterpretazione di Peter Pan in chiave OSR, compatibile con le regole
        di D\&D.}; \emph{Dune}, in volumi cartonati ``importanti''; ma anche
        \emph{The Quiet Year}\footnote{La trascrizione dice ``Quietier'' e, più
        avanti, ``The Quiet Ear''.}, che sta in un sacchettino; e \emph{Lady
        Blackbird}\footnote{La trascrizione dice ``Lady Blackboard''. \emph{Lady
        Blackbird} è un gioco di ruolo distribuito gratuitamente in PDF.}, che è
        solo digitale: ``tu lo prendi, lo stampi in bianco e nero e te lo
        giochi''. L'edizione italiana di \emph{Montsegur 1244}\footnote{La
        trascrizione dice ``Monsegur''.} è una sorta di scatola che si apre;
        \emph{CBR+PNK}\footnote{La trascrizione dice ``CBR PNK''.} ``si apre in
        un modo stranissimo'' e contiene dei pieghevoli (\emph{leaflet}); ci sono
        giochi di ruolo in formato brochure. Daniele li mostrerà meglio nella
        lezione sulla \emph{fiction}. La conclusione: nei giochi di ruolo
        \textbf{non si può usare il metro del formato}.
  \item[Gioco di ruolo dal vivo (Ruggero).] Il LARP, \emph{Live Action
        Role-Playing game}: ``anche questa è una categoria gigantesca''.
  \item[Gioco pervasivo (Beatrice).] La categoria meno conosciuta in Italia, e
        anche per questo inserita nel corso.
\end{description}

\paragraph{Che cosa distingue il gioco di ruolo.} Daniele chiede alla classe
che cosa distingua i \emph{role-playing game}. La risposta è nel nome: il
\textbf{role playing}, cioè un \textbf{elemento di interpretazione}, che non è
detto sia interpretazione di un \emph{personaggio} (in \emph{The Quiet Year},
per esempio, non c'è un personaggio, ``abbiamo qualcosa di diverso''). E c'è un
elemento ``importantissimo'': la \textbf{fiction}, cioè la narrazione. Non la
narrazione intesa come ``scrivere una trama e giocarla'', ma quella che nasce al
tavolo, la \textbf{narrazione emergente}: non si è previsto che l'avventura
andasse in quel modo, ma mentre si gioca gli altri dicono cose che la portano in
una certa direzione. Se ne parlerà nella lezione sulla fiction emergente.
Categorizzare i GDR è ``veramente molto arduo'', e anche il pubblico è
particolare: ``ognuno ha il suo palato'', c'è chi gioca tutto e chi gioca solo un
sottogenere. A una domanda dalla classe (il gioco \emph{Lovecraftesque}, citato
da qualcuno, rientra tra i giochi di ruolo da tavolo?) la risposta è sì: e di
che \emph{tipo} di gioco di ruolo sia si capirà con le sottocategorie che Daniele
presenterà più tardi.

\paragraph{Perché il LARP è separato dal GDR.} Ruggero spiega che i LARP sono
separati dal GDR da tavolo ``semplicemente per motivi di settore e di
industria''. Il gioco di ruolo da tavolo viene venduto con una produzione più
vicina al gioco da tavolo: si produce un \textbf{oggetto fisico} che poi si
vende. Il LARP invece è ancora spesso organizzato \textbf{a eventi}: si partecipa
\textbf{comprando un biglietto}, e si va a giocare in un posto specifico, a
un'ora specifica, con persone che sono lì per giocare quel gioco. Da questo
punto di vista è più vicino ad alcuni giochi pervasivi. Le sottocategorie sono
``una valanga'': la più famosa è il \textbf{boffer} (il cliché del LARP: persone
in costume medievale ``che si picchiano con spade di plastazote o di gommapiuma''),
ma ce ne sono di molto più astratte e di molto più accessibili, che si possono
organizzare anche in casa.

\subsection{I giochi pervasivi: una prima definizione}

Beatrice chiede se qualcuno abbia mai sentito parlare di \emph{pervasive games}.
Quasi nessuno; una studentessa ne ha sentito parlare solo leggendo la
descrizione del corso, e si è incuriosita. ``Ci sta tantissimo'': in Italia non
sono ancora molto conosciuti, ed è uno dei motivi per cui sono nel percorso.

Tiziana chiede se le \textbf{cacce al tesoro} organizzate nelle città rientrino
negli \emph{urban game}. ``Assolutamente sì'': gli \textbf{urban game} sono una
categoria dei giochi pervasivi, e nella slide si vedono proprio come un
sottoinsieme del ``grande cappello''. Beatrice dà due definizioni:
\begin{enumerate}
  \item la definizione ``standard'', dai \emph{game studies}: un gioco
        pervasivo è un gioco che \textbf{espande il proprio cerchio magico
        spazialmente, temporalmente e socialmente};
  \item la definizione più intuitiva, da cui partire per capire poi meglio la
        prima: un gioco pervasivo è un gioco che \textbf{utilizza i luoghi e gli
        oggetti della vita quotidiana e li trasforma in luoghi e oggetti di
        gioco}.
\end{enumerate}
Se la prima non è chiara, ``non è un problema, la capiremo con i casi studio''.
Negli urban game ``non giochiamo seduti al tavolo, ma giochiamo in città, e
stiamo comunque giocando''. La parola \emph{pervasivo} indica la pervasività del
gioco, che entra in situazioni in cui non siamo abituati a immaginarlo, anche se
magari siamo abituatissimi a viverlo senza pensarci.

\paragraph{Le cene con delitto.} Tiziana fa un'altra domanda (Beatrice ricorda
che si può sempre interrompere, anche con la mano alzata delle reazioni di
Zoom): le \textbf{cene con delitto} rientrano tra i pervasivi? E il
\textbf{teatro interattivo}, visto che in certe cene con delitto chi le propone
recita una vera rappresentazione? Per Beatrice, ``a mio dire, sì, entrambe''.
Ma le cene con delitto, o \emph{murder party}, esistono in molte versioni:
\begin{itemize}
  \item in una, gli \textbf{attori} vengono al tavolo con te, magari
        insospettabili, e tu, come persona giocante, li interroghi: questa
        versione assomiglia di più al teatro interattivo, che è incluso nei
        giochi pervasivi;
  \item in un'altra, \textbf{ognuno riceve un ruolo}, con una sorta di copione
        (``stasera sei la sorella della vittima''): questa va più nella
        direzione dei \textbf{giochi di ruolo dal vivo}.
\end{itemize}
Sono distinzioni ``che facciamo per allenarci a comprendere meglio'', ma che
spesso si sovrappongono: anche quando si è investigatori e si interrogano degli
attori, si è dentro una dinamica di gioco di ruolo, ``anche se magari non come
la immaginiamo in D\&D o in giochi più blasonati''. Le definizioni ``sono
malleabili'': si imparano le parole e poi si vede come usarle nel caso
specifico.

\subsection{Ibridazione e sperimentazione}

Nell'angolo della slide, dietro tutte le categorie, c'è ``una piccola cosina'':
\textbf{ibridazione e sperimentazione}. ``In realtà le categorie non esistono'',
dice Ruggero: esistono come divisioni tra industrie, ma escono continuamente
giochi che fanno da ponte. La slide mostra alcuni esempi:
\begin{itemize}
  \item \emph{Fateforge}\footnote{La trascrizione dice ``FitForge'' e, più
        avanti, ``Fateforge''.}, un gioco da tavolo che usa un'\textbf{app}
        digitale integrata nel gioco: ``un po' videogioco, un po' board game'';
  \item un gioco di aste che usa anch'esso un cellulare (il titolo non è
        detto);
  \item dei \emph{mock trial} (processi simulati) ibridi, a metà tra il murder
        party e il roleplay, ``un po' roleplay, ma anche un po' giochi
        pervasivi'';
  \item \emph{Alice is Missing}, un gioco di ruolo da tavolo che usa una
        \textbf{chat} come Telegram tra le meccaniche di gioco.
\end{itemize}
Le ibridazioni rendono spesso difficile mettere un gioco ``al 100\%'' in uno di
questi ``cesti''.

\begin{figure}[htbp]
\centering
\begin{tikzpicture}[
    font=\small,
    macro/.style={draw, thick, rounded corners=4pt, fill=gray!12,
                  text width=5.2cm, align=left, inner sep=6pt},
    titolo/.style={font=\bfseries},
    ponte/.style={draw, dashed, rounded corners=3pt, fill=white,
                  text width=3.4cm, align=center, font=\footnotesize,
                  inner sep=4pt}
  ]
  \node[macro] (tav) at (0,0) {{\bfseries Gioco da tavolo} \hfill
    \footnotesize Ruggero\\[2pt]
    \footnotesize board game, card game, party game: tutto ciò che è
    prodotto dall'industria del gioco da tavolo};
  \node[macro] (gdr) at (6.6,0) {{\bfseries Gioco di ruolo da tavolo}
    \hfill \footnotesize Daniele\\[2pt]
    \footnotesize GDR / TTRPG: interpretazione e fiction emergente; nessuna
    definizione univoca, formati diversissimi};
  \node[macro] (larp) at (0,-3.6) {{\bfseries Gioco di ruolo dal vivo}
    \hfill \footnotesize Ruggero\\[2pt]
    \footnotesize LARP: spesso a evento, con biglietto; dal boffer al
    black box};
  \node[macro] (perv) at (6.6,-3.6) {{\bfseries Gioco pervasivo}
    \hfill \footnotesize Beatrice\\[2pt]
    \footnotesize luoghi e oggetti quotidiani trasformati in luoghi e oggetti
    di gioco; urban game, ARG, teatro interattivo};
  \node[ponte] (ibr) at (3.3,-1.8) {\textbf{ibridazione e
    sperimentazione}\\ app, chat, murder party, mock trial, eventi\ldots};
  \draw[dashed, thick] (ibr) -- (tav);
  \draw[dashed, thick] (ibr) -- (gdr);
  \draw[dashed, thick] (ibr) -- (larp);
  \draw[dashed, thick] (ibr) -- (perv);
  \draw[{Stealth[length=2mm]}-{Stealth[length=2mm]}, gray, thick]
    (gdr.south) -- node[right, font=\footnotesize\itshape, text=black,
    text width=2.4cm] {oggetto venduto vs evento con biglietto} (perv.north);
  \draw[{Stealth[length=2mm]}-{Stealth[length=2mm]}, gray, thick]
    (tav.south) -- node[left, font=\footnotesize\itshape, text=black,
    text width=2.4cm, align=right] {separati per motivi di industria}
    (larp.north);
\end{tikzpicture}
\caption{Le quattro macrocategorie del gioco analogico usate nella lezione, con
il docente che le presenta. Al centro, la ``piccola cosina'' della slide:
ibridazione e sperimentazione, che rendono i confini porosi. Le frecce grigie
riassumono due osservazioni dei docenti: il GDR da tavolo e il LARP sono
separati soprattutto per ragioni di industria (oggetto venduto o evento), e per
lo stesso motivo il LARP è vicino a molti giochi pervasivi.}
\label{a2:fig:macro}
\end{figure}

% ---------------------------------------------------------------------
\section{Non chiudersi nel proprio genere}
\label{a2:sec:aperti}

Prima di aprire la ``diga dei casi studio'', i docenti insistono su alcuni punti
di metodo.

\paragraph{Le categorie servono a ragionare, non a recintare.} Ruggero: le
categorie non vengono date ``per dirvi: le cose sono A, B, C e basta'', ma per
far ragionare su \textbf{quanto sono flessibili i confini}, e per dare esempi di
cose \textbf{adiacenti} a quelle che interessano, che ``possono sbloccarvi delle
cose''. Ognuno è lì per interessi diversi (i giochi di ruolo, le carte, i giochi
formativi e didattici\ldots), ma l'invito è a \textbf{non mettersi i paraocchi}
sul proprio genere: ``è grazie alla \textbf{cross-modalità} e al vedere qualcosa
di risolto in un altro gioco che vi arriva poi un'idea per un gioco vostro che
non c'entra apparentemente niente'', magari nemmeno meccanicamente, con l'altro.

\paragraph{Un corso agnostico di genere.} Daniele aggiunge che il corso è
\textbf{agnostico di genere}: alcune lezioni saranno più importanti per una
categoria, ma i \textbf{concetti di base del design si applicano a qualsiasi
sviluppo ludico}. È fondamentale, perché altrimenti sembra che chi fa un certo
tipo di prodotto abbia bisogno di ``un game design particolare'', e non è così.
Anzi: quando è successo, storicamente anche in Italia, si è finito per
\textbf{ostracizzare e separare le community}, mentre i docenti vogliono
\textbf{metterle insieme}. ``Voi non sapete quante volte i videogiochi vengono
ispirati dai giochi da tavolo, e viceversa''.

\paragraph{Il project work è libero.} C'è anche un motivo pratico. Il progetto
di gruppo finale è \textbf{libero}: i docenti non assegneranno una categoria di
gioco, saranno i gruppi a sceglierla, e può essere anche
un'\textbf{ibridazione}, una \textbf{sperimentazione}, un gioco che non si
colloca perfettamente in un genere (``anzi, sarebbe fantastico''). Ma i gruppi
saranno formati da persone con interessi diversi, ed è proprio per questo che
bisogna restare aperti. L'unica cosa che i docenti chiederanno di non fare è
\textbf{un gioco enorme}: il progetto deve essere \textbf{calibrato}.

\paragraph{``Per me, questo è design''.} Infine, Ruggero chiede di \textbf{non
sottovalutare} la lezione: ad alcuni può sembrare ``meramente un elenco di cose
che vi facciamo vedere per nessun motivo'', ma ``per me questo è design''.
Quando deve trovare una soluzione, va a cercare altri giochi; e gli è capitato
moltissime volte di riaprire proprio quelle slide per ricordarsi com'era fatto un
gioco. E come hanno fatto Tiziana e Rita, bisogna \textbf{fare domande}:
``fermateci, chiedeteci'', i docenti sono lì per approfondire.

% ---------------------------------------------------------------------
\section{I giochi pervasivi: la premessa di Beatrice}
\label{a2:sec:pervasivi}

Si comincia ``in scivolata'' con i giochi pervasivi. Beatrice anticipa che
diversi casi studio hanno \textbf{accenni di digitale}, per tre motivi:
\begin{enumerate}
  \item alcuni casi sono \textbf{molto famosi} e aiutano a capire il concetto di
        pervasive game, anche se non sono prettamente analogici;
  \item un gioco analogico \textbf{non esclude completamente la componente
        digitale}: basti pensare ai giochi da tavolo che usano un'app;
  \item l'\textbf{ibridazione}: si vedranno casi studio ``molto strani''.
\end{enumerate}
Aggiunge una considerazione personale e professionale: quello dei pervasivi è
un campo in cui c'è ancora \textbf{ampio spazio}, perché sono poco conosciuti e
la loro \textbf{definizione è ancora in costruzione}. C'è spazio ``per mettere
anche la vostra voce'', per esprimersi e sperimentare, proprio perché le
definizioni non sono state date in modo rigido come altrove. E anticipa un tema
ricorrente del corso: ``parleremo molto di \textbf{emozioni}, anche di come
accompagnarle all'interno di un'esperienza di gioco''.

\paragraph{Approfondimento: il cerchio magico.} La definizione ``da game
studies'' presuppone il concetto di \textbf{cerchio magico}, che a lezione non
viene spiegato. In \emph{The Art of Game Design}, Schell lo introduce
discutendo le definizioni di gioco: quasi tutte implicano che un gioco sia un
sistema \emph{chiuso}, con un confine, e ``Johan Huizinga lo ha chiamato il
\emph{cerchio magico}'', che ha davvero qualcosa di magico: quando siamo
mentalmente ``dentro il gioco'' abbiamo pensieri, sentimenti e valori diversi da
quando ne siamo fuori. Schell lo reinterpreta come il nostro \textbf{sistema
interno di risoluzione dei problemi}: una ``microrealtà'' semplificata in cui gli
elementi acquistano un'importanza che non devono avere fuori dal gioco. Nel
capitolo sulle meccaniche, lo stesso Schell identifica il cerchio magico con lo
\textbf{spazio} del gioco, la prima delle sue sette categorie di meccaniche. Il
gioco pervasivo è dunque quello che sposta, allarga o rende sfumato questo
confine \textbf{nello spazio} (si gioca nella città, non sul tavolo), \textbf{nel
tempo} (non si sa bene quando la partita inizia e finisce) e \textbf{nella
società} (vi entrano persone che non sanno di giocare, o che giocano senza
chiamarlo gioco)\footnote{La formula ``espandere il cerchio magico
spazialmente, temporalmente e socialmente'' è quella proposta da Markus Montola
nei suoi studi sui giochi pervasivi, poi ripresa nel volume \emph{Pervasive
Games: Theory and Design} di Markus Montola, Jaakko Stenros e Annika Waern
(2009). Beatrice non cita la fonte; la nota è un'integrazione.}. È esattamente
quello che Beatrice dirà di \emph{GuruShot}: ``non sappiamo bene dove inizia e
dove finisce l'attività di gioco, ed è per questo che viene utilizzata la parola
pervasivo''.

\begin{figure}[htbp]
\centering
\begin{tikzpicture}[
    font=\small,
    freccia/.style={-{Stealth[length=2.5mm]}, thick},
    etichetta/.style={font=\footnotesize, align=center, text width=3.6cm}
  ]
  % cerchio classico
  \draw[thick, fill=gray!15] (0,0) circle (1.0cm);
  \node[font=\footnotesize, align=center] at (0,0) {cerchio\\magico};
  \node[font=\footnotesize\itshape, align=center] at (0,-1.45)
    {gioco ``classico'':\\tavolo, durata, giocatori};
  % cerchio espanso
  \draw[thick, dashed, fill=gray!6] (7,0) circle (2.1cm);
  \draw[thick, fill=gray!15] (7,0) circle (1.0cm);
  \node[font=\footnotesize, align=center] at (7,0) {cerchio\\magico};
  \draw[freccia] (1.3,0) -- node[above, font=\footnotesize]{gioco pervasivo}
    (4.6,0);
  \draw[freccia, gray] (7,1.0) -- (7,2.1);
  \draw[freccia, gray] (5.95,-0.6) -- (5.18,-1.05);
  \draw[freccia, gray] (8.05,-0.6) -- (8.82,-1.05);
  \node[etichetta, above=0.05cm] at (7,2.15) {\textbf{spazialmente}: la città,
    l'hotel, la stazione dell'autobus};
  \node[etichetta, anchor=north east] at (5.4,-1.5) {\textbf{temporalmente}:
    quando inizia e quando finisce la partita?};
  \node[etichetta, anchor=north west] at (8.6,-1.5) {\textbf{socialmente}:
    passanti, pasticcieri, vicini di quartiere};
\end{tikzpicture}
\caption{La definizione ``da game studies'' citata da Beatrice: il gioco
pervasivo espande il proprio cerchio magico nello spazio, nel tempo e nella
società. La definizione più intuitiva dice la stessa cosa dal lato degli
oggetti: luoghi e oggetti della vita quotidiana diventano luoghi e oggetti di
gioco.}
\label{a2:fig:cerchio}
\end{figure}

% ---------------------------------------------------------------------
\section{Pervasivi con un'app: Pokémon Go, geocaching, GuruShot}
\label{a2:sec:pervasivi-app}

\subsection{Pokémon Go}

\slidefig{img/AGD/agd2-pokemon}{Pervasivi con un'app: Pokémon Go, il geocaching e GuruShots.}{slide \emph{Analog -- Case Studies}}{fig:agd2-pokemon}


Il primo esempio è scelto ``proprio per farvi comprendere'': \textbf{Pokémon
Go}, che conoscono tutti, o di cui almeno tutti hanno sentito il nome. Beatrice
lo spiega comunque, ``così ci allineiamo'' (scusandosi con i grandi fan per la
descrizione imprecisa). È un gioco in formato \textbf{app}. I Pokémon, ``questi
animaletti bizzarri'', nella narrativa dei Pokémon popolano il loro mondo; in
\emph{Pokémon Go} popolano \textbf{il nostro mondo reale}. Il gioco consiste nel
cercarli, catturarli, andare nelle palestre, e tutto questo si fa
\textbf{camminando nel mondo}: la stessa via in cui si va a fare colazione al
bar o che si percorre per andare in ufficio è la via in cui si può trovare un
Pokémon o una palestra. Ed è ``tutto costruito dagli utenti''.

Su \emph{Pokémon Go} ci sarebbe molto da dire, ``anche su come è stato
utilizzato per fare altre cose meno belle ed eleganti dell'esperienza di
gioco'', ma Beatrice non ne parla per ragioni di tempo: ci tiene solo a far
notare che ``tante cose possono avere tante sfumature'', e che i casi più
spinosi si possono discutere a parte. Qui serve per una ragione semplice:
\textbf{è pervasivo perché usa gli stessi spazi della vita quotidiana come spazi
di gioco}.

\subsection{Il geocaching}

Il \textbf{geocaching} è anch'esso abbastanza conosciuto. ``Di base è una
gigantesca caccia al tesoro'', tenuta in piedi da alcuni portali, e nata
proprio con il \textbf{GPS} e la geolocalizzazione. Il principio:
\begin{itemize}
  \item le persone giocanti \textbf{nascondono} delle scatole, piccole o
        grandi, con contenuti vari (dei ``tesori'') e con una sorta di
        \textbf{diario di bordo} (un \emph{log}: fogli o quaderni) su cui chi le
        trova può scrivere ``sono stato qui, ho trovato questa scatoletta'';
  \item la scatola viene nascosta nel posto che l'utente ritiene più idoneo, e
        sulla piattaforma l'utente dà le \textbf{indicazioni} per trovarla;
  \item la grande regola è che \textbf{le scatole non vengono mai tolte}: chi le
        trova può scambiare gli oggetti al loro interno (piccoli, di qualunque
        valore, con piena libertà) e segnare la propria presenza.
\end{itemize}
Non è ``un arraffare i tesori'', ma un modo di \textbf{viverli e dar loro
spazio collettivamente}. In Italia, ``per la cultura che abbiamo'', funziona
così così, perché le scatole spesso vengono rubate o spostate. Ma vale la pena
guardare online i video dei modi in cui vengono nascoste: alcune sono ``piccole
opere di ingegneria artistica'', quasi piccole escape room, con un gioco di pesi
o una chiave ``nascosta in bella vista'' che si vede solo con grande
attenzione; alcune non sono scatole ma \textbf{cabine telefoniche} in cui si
entra e succedono cose, o intere automobili.

\subsection{GuruShot}

\textbf{GuruShot}\footnote{Più avanti la trascrizione lo chiama anche ``blue
shot''.} è di nuovo in formato app, ma Beatrice lo inserisce perché usa uno
strumento che lei usa ``moltissimo'' e che trova ``molto potente per fare molte
cose'': \textbf{la fotografia}. È un gioco dedicato a persone molto dentro al
tema, fotografi anche di un certo livello di passione e di competenza: una
\textbf{gara} o maratona fotografica in cui vengono date delle \textbf{challenge
fotografiche}. Chi vuole giocare apre l'app, sceglie la sfida, la svolge, carica
la foto; e poi ``succedono cose'': valutazioni, momenti di incontro con la
community, premi veri e propri offerti dagli \textbf{sponsor}. È un'attività
\textbf{continua}, online da molti anni (``almeno 7--8'', corregge Beatrice), con
una nicchia internazionale molto nutrita.

Ed è ``molto pervasivo'': \textbf{non si può dire quando una partita è iniziata
o finita}, né quando si sta giocando e quando no. Rispetto alla definizione del
cerchio magico, non si sa bene dove inizi e dove finisca l'attività di gioco:
per questo si parla di pervasività.

% ---------------------------------------------------------------------
\section{La città come campo di gioco: Journey, SF0 e Critical City}
\label{a2:sec:pervasivi-citta}

\subsection{Journey to the End of the Night}

\slidefig{img/AGD/agd2-journey}{La città come campo di gioco: Journey to the End of the Night e SF0.}{slide \emph{Analog -- Case Studies}}{fig:agd2-journey}


Si torna sotto il cappello dei \textbf{giochi urbani}, con un gioco che si
gioca in città, sempre \textbf{in formato evento}: un'ora di inizio, un'ora di
fine e una zona di gioco all'interno della città. È ``una sorta di \textbf{guardie
e ladri gigante}''. Lo scopo dei creatori era \textbf{riappropriarsi della vita
notturna} e viverla in modo ludico, esplorando la città ``in maniera
\emph{safe} da un lato e un po' \emph{noir} dall'altro''. Un piccolo gruppo di
giocatori fa i \textbf{chaser} (gli inseguitori, le guardie) e gli altri i
\textbf{runner} (i ladri); ci sono delle \textbf{tappe} da raggiungere, e il
gioco dura svariate ore. Non è un gioco prettamente sportivo (ci sono altri
pervasivi molto più sportivi): ci sono molte dinamiche \textbf{strategiche} su
come arrivare alla fine del percorso. Beatrice lo inserisce perché, rispetto ai
casi precedenti, è \textbf{completamente analogico}: non usa nessuna app, nulla.
È un evento.

Il gioco è \textbf{Journey to the End of the Night}, nato in America da un
collettivo con base a \textbf{San Francisco}. È arrivato in Italia nel
\textbf{2008}, durante un'ondata di giochi pervasivi arrivata ``in modo molto
forte'', e poi in Giappone e in tantissimi altri posti, perché i creatori lo
hanno reso \textbf{libero}: chi vuole può organizzarlo, purché \textbf{non metta
un biglietto a pagamento}. Ha avuto per anni un grande seguito: cercando
``Journey to the End'' Italia si trovano tantissimi eventi, e a \textbf{Milano},
tra il 2008 e il 2013, se ne sono fatte svariate edizioni con
\textbf{2\,000--3\,000 partecipanti} ciascuna. Beatrice lo racconta ``non tanto
per le meccaniche'', che più o meno si conoscono, ma per \textbf{come è
arrivato}: un'onda lunga che dura ancora, anche se si sta affievolendo.

\subsection{SF0}

Degli stessi creatori è \textbf{SF0}\footnote{La trascrizione dice ``SF0'',
``SFZero'' e, una volta, ``essere al vero''.}, che era una sorta di grande
\textbf{piattaforma online} con lo stesso principio di \emph{GuruShot}: delle
\textbf{missioni} da svolgere e documentare con video o foto. Ma invece di
puntare sull'estetica della fotografia, le missioni puntavano sulle
\textbf{interazioni cittadine e urbane}.

\subsection{Critical City Upload}

SF0 ha avuto ricadute anche in Italia, con ``un figlio, chiamiamolo così'':
\textbf{Critical City Upload}, un gioco che è andato molto tra il 2009 e il
2013. Lo scopo è vincere, accumulando punti con le missioni, ma il gioco è
molto legato agli \textbf{scambi umani e al territorio}. Le missioni andavano da
``organizza una cena con una persona del tuo quartiere con cui non hai mai
parlato'' a una delle preferite di Beatrice: ``crea un'opera d'arte interattiva
che intrattenga le persone che aspettano l'autobus e posizionala in una delle
pensiline del quartiere o della città''.

In entrambi i casi \textbf{è la community che valuta la community}: non ci sono
giudici. E dopo i primi anni, sia SF0 sia Critical City hanno permesso ai
partecipanti che avevano raggiunto un certo livello di \textbf{creare essi
stessi nuove missioni} per i nuovi giocatori: la comunità non si limitava a
valutarsi a vicenda, ma creava il gioco per gli altri. Oggi nessuno dei due è
più attivo, ma i siti sono ancora online, con molte delle missioni visibili.
Quanto a ``stranezza'' rispetto all'idea comune di gioco analogico, ``questo
spinge molto forte''.

% ---------------------------------------------------------------------
\section{Teatro, ARG e marketing: da Sleep No More a The Institute}
\label{a2:sec:pervasivi-arg}

\subsection{Sleep No More}

\slidefig{img/AGD/agd2-sleep}{Teatro e ARG: Sleep No More e Why So Serious?}{slide \emph{Analog -- Case Studies}}{fig:agd2-sleep}


Chi conosce il \textbf{teatro interattivo} conosce di sicuro \textbf{Sleep No
More}\footnote{La trascrizione dice ``Slip no more'' e ``SlipMod''. Lo
spettacolo è della compagnia britannica Punchdrunk; la versione più nota è
quella di New York, allestita in un edificio trasformato nel fittizio McKittrick
Hotel. Il nome della compagnia non è detto a lezione.}: un'opera messa in scena
all'interno di un \textbf{hotel}. È interattiva nel senso che gli attori recitano
la loro parte in interazione con il pubblico, e il pubblico, con regole precise
sui contatti fisici e su come interagire con gli spazi, è \textbf{libero di
esplorare l'hotel} come vuole: in ogni zona ci sono attori che fanno cose. Resta
un'opera teatrale, con un inizio, una fine e degli eventi nel mezzo; c'è molta
documentazione online (interviste, spezzoni). E di nuovo si gioca su
\textbf{qual è lo spazio di gioco e quale non lo è}.

\subsection{Gli alternate reality game e Why So Serious}

Qui Beatrice introduce un nuovo termine: gli \textbf{alternate reality game}
(ARG), i giochi di realtà alternativa, ``giochi in cui c'è un intero mondo
alternativo al nostro, in cui succedono cose''. Il caso scelto, ``uno dei miei
preferiti'', porta anche nel mondo dei \textbf{giochi per il marketing e la
comunicazione}: \textbf{Why So Serious}, nato come gioco promozionale per
\emph{Il cavaliere oscuro} di Christopher Nolan.
\begin{itemize}
  \item Iniziava con una \textbf{mail} ad alcune persone iscritte a mailing
        list di fan di Batman, contattate dagli organizzatori tramite le loro
        associazioni.
  \item La mail era un \textbf{invito strano a entrare nella banda del Joker},
        con un enigma; da lì partivano attività di socializzazione organizzate
        attraverso un sito, che era una serie di indirizzi ``in cui succedevano
        cose''.
  \item Il gioco raccontava che \textbf{nel nostro mondo il Joker esisteva
        davvero}, e gli indirizzi erano indirizzi reali.
  \item Una delle prime cose: andare a questi indirizzi, tutti
        \textbf{pasticcerie}, e dire la parola in codice ricevuta. Si riceveva
        una \textbf{torta} con sopra scritto un numero di telefono; chiamandolo,
        squillava la torta stessa, perché dentro c'era un \textbf{cellulare}:
        lo strumento con cui il Joker gestiva le comunicazioni con i giocatori.
\end{itemize}
Molte strutture di questo tipo, molti colpi di scena ``anche proprio pratici'',
che sfruttano la pervasività: cose che nel nostro mondo non hanno motivo di
succedere (``non esiste che ci sia un cellulare in una torta'') ma che in quel
mondo hanno un senso. Gli ARG giocano moltissimo su ciò che \textbf{potrebbe
essere reale ma non lo è}, e che il gioco dice che lo è: ``in quel mondo voi
siete gli aiutanti del Joker''. \emph{Why So Serious} è opera di un'azienda che
si occupa da molti anni di ARG\footnote{Beatrice non fa il nome dell'azienda,
che si troverà nell'estratto distribuito. Si tratta di 42 Entertainment
(integrazione).}; i dettagli sono nell'estratto che Beatrice caricherà, tratto
da \emph{Immersi nelle storie} di \textbf{Frank Rose}\footnote{Edizione
italiana di \emph{The Art of Immersion} (2011).}, che racconta proprio l'inizio
di questo gioco.

Perché raccontarlo? Perché mostra ``come possiamo giocare sulle
\textbf{emozioni}'' e come \textbf{connettere le persone alle cose che vogliamo
comunicare}, anche costruendo un mondo alternativo. Non è un obbligo, ``ma è bene
sapere'' che esistono questi giochi, ``giganteschi''.

\subsection{Un gioco da tavolo che è anche un ARG}

Tornando alla malleabilità delle definizioni, Beatrice presenta un gioco che è
insieme \textbf{gioco da tavolo}, \textbf{alternate reality game} e gioco con
\textbf{componenti digitali}, edito da \textbf{MS Edizioni}\footnote{La
trascrizione chiama il gioco ``The Ice''; il titolo non è ricostruibile con
certezza (potrebbe trattarsi di un titolo come \emph{The Heist}). MS Edizioni
è un editore italiano di giochi; è forse la stessa ``MS'' per cui, nella prima
lezione, Ruggero dice di lavorare, ma la trascrizione non permette di
confermarlo.}. Si riceve un \textbf{fascicolo} di informazioni e poi dei
\textbf{messaggi}, da cui si capisce di far parte di un gruppo che \textbf{sta
rapinando una banca}, già nel pieno della rapina. I giocatori sono ``l'entità che
ha accesso a internet'', la voce fuori dall'azione che coordina: parlano con dei
personaggi, scambiano informazioni, guardano video, analizzano dati e
documenti. È un gioco ``con una profondità notevole''. Due cose lo rendono
interessante:
\begin{itemize}
  \item mescola cose che \textbf{esistono davvero} (la cartina di una città
        reale, gli scambi della metropolitana di una città reale) con cose del
        \textbf{mondo narrativo} (i personaggi, i fatti, la rapina: ``spero che
        sia chiaro che non state rapinando la banca davvero'');
  \item a differenza degli altri casi, ha un \textbf{inizio e una fine ben
        definiti} ed è \textbf{completamente accessibile}: lo si può provare.
\end{itemize}
Per Beatrice è un caso ``unico sul panorama italiano'' (``poi mi correggeranno i
colleghi a livello internazionale'').

\subsection{Collettivo Cinetico: Cinetico 4.4 e Amleto}

\slidefig{img/AGD/agd2-teatro}{Cinetico 4.4, Lettere anonime per un camminatore, The MP3 Experiment, The Institute e l'ARG di Stranger Things.}{slide \emph{Analog -- Case Studies}}{fig:agd2-teatro}


\textbf{Cinetico 4.4} è ``un altro bell'esperimento'': un \textbf{gioco da
tavolo creato da un collettivo di danzatori e danzatrici}, il \textbf{Collettivo
Cinetico}, per generare una \textbf{performance artistica dal vivo}, teatrale o
di danza. Non è un gioco da prendere dallo scaffale di una ludoteca: è pensato
per una performance. Chi ci gioca partecipa a un gioco da tavolo che in realtà è
un \textbf{evento}: si paga un biglietto, si entra, ci si siede, succedono cose.
Il collettivo, di cui Beatrice si dichiara ``di parte'', fa moltissime attività
``né danza, né teatro, né gioco, o meglio tutte e tre le cose insieme'', difficili
da mettere sotto un cappello. E qui Beatrice aggancia un tema delle prossime
lezioni: la \textbf{transmedialità} e la \textbf{ricerca e sviluppo}, cioè
attingere anche a cose lontane dal gioco che si vuole creare, perché ``più
conosciamo, più possiamo aprirci la testa''.

Un caso studio in più: l'\textbf{Amleto} del Collettivo Cinetico. È l'Amleto di
Shakespeare, ma sul palco non ci sono i danzatori del collettivo. Tre giorni
prima dello spettacolo fanno una vera e propria \textbf{audizione} aperta alla
popolazione; tra chi si presenta scelgono \textbf{sei potenziali Amleto}, che
portano sul palco a danzare e recitare, e creano momenti di
\textbf{interazione con il pubblico} per scegliere, tra i sei, l'Amleto di
quell'edizione. Ci sono giochi veri e propri sia nell'audizione sia nello
spettacolo. Loro, da danzatori, lo chiamano spettacolo di danza, ``ma se lo
vedete c'è tantissimo''.

\subsection{Lettere anonime per un camminatore}

``Rispetto a quello che vi ho detto, questo è semplice''. \textbf{Lettere
anonime per un camminatore} è del collettivo \textbf{Amigdala}, anche loro una
compagnia teatrale (Gianmarco, che ha vissuto a Modena, conferma che sono di
lì). È una \textbf{passeggiata} con piccole interazioni ludiche, tutta basata su
un \textbf{audio}:
\begin{itemize}
  \item è un gioco \textbf{in solitario}: una persona sola segue un percorso
        segnato da frecce, per terra o sui muri di una città, guidata da una
        \textbf{voce in cuffia} che le racconta delle cose;
  \item la struttura del racconto è fissa, ma la \textbf{parte centrale} cambia
        ogni volta: prima di ogni edizione il collettivo \textbf{intervista le
        persone che hanno abitato o abitano il percorso}, e mentre si cammina si
        ascolta, oltre alla narrazione, il racconto di chi quei luoghi li vive o
        li ha vissuti;
  \item la \textbf{conclusione} è sempre ``col botto'': il collettivo sceglie di
        finire in modo scenografico in luoghi normalmente inaccessibili.
        Beatrice lo ha giocato a Ferrara, e alla fine del percorso si è
        ritrovata seduta in poltrona, con un tavolino e un drink, al
        \textbf{centro dello stadio comunale}, con lo stadio vuoto e le cuffie
        nelle orecchie.
\end{itemize}
Torna il tema del \textbf{percorso emotivo} che si fa fare alla persona giocante.
E lo strumento dell'\textbf{audio che guida} è interessante anche per il corso:
un'app non si riesce a fare nel tempo del project work, ``ma per utilizzare
l'audio invece possiamo''.

\subsection{MP3 Experiment}

Un altro esempio basato sull'audio, più simile a un \textbf{flash mob}:
\textbf{MP3 Experiment}, del collettivo newyorkese \textbf{Improv
Everywhere}\footnote{La trascrizione dice ``Improve Everywhere''.}, che fa
moltissimi esperimenti di questo tipo. Ci si ritrova tutti nel giorno, nell'ora
e nel posto concordati, ci si collega con le cuffie a una trasmissione in
diretta, e la voce dice che cosa fare. Beatrice non svela di più (``guardatevi i
video su YouTube, ci sono delle perle''), ma racconta un dettaglio: gli
organizzatori \textbf{elencano gli oggetti da portare}, e in un'edizione chiesero
di portare \textbf{due matite}; a un certo punto tutto il gruppo si mise a
suonare la batteria con le matite sugli oggetti della città. Quella
\textbf{scintilla di curiosità} per oggetti di cui non si sa ancora l'uso è uno
strumento del percorso emotivo, ``non difficile'' da usare: ``basta solo
ragionare di creatività''. Come le istruzioni audio.

\subsection{The Institute e Dispatches from Elsewhere}

\textbf{The Institute}\footnote{Il gioco è noto anche come \emph{The Jejune
Institute}, di San Francisco (integrazione).} è, nel modo in cui Beatrice ha
imparato questo mondo, uno dei casi più blasonati e famosi di alternate reality
game, uno dei primi raccontati come tale. È un'esperienza così complessa che
Beatrice preferisce indicare due fonti per approfondirla:
\begin{itemize}
  \item un \textbf{documentario}, intitolato anch'esso \emph{The Institute}, che
        racconta tutto l'ARG, durato circa un anno e fatto di molti passaggi, con
        tantissimi colpi di scena del tipo ``torta con il cellulare'';
  \item una \textbf{serie TV}, \emph{Dispatches from Elsewhere} (Daniele
        conferma il titolo), che ha preso l'esperienza di The Institute, ci ha
        aggiunto una storia dei personaggi e ne ha fatto una serie. Tolta la
        storia dei personaggi (``comunque molto carina''), resta più o meno
        l'idea dell'ARG.
\end{itemize}

\paragraph{Con budget zero.} Qui Beatrice fa una precisazione importante. Questi
esempi sono \textbf{difficili da replicare}: ``con che soldi compriamo la torta e
ci mettiamo il cellulare? Con che soldi affittiamo un intero palazzo come in The
Institute?''. Serve un intero lato commerciale, di ricerca di sponsor, che non è
parte del design del gioco e quindi non è il motivo per cui si è nel corso.
Quello che vuole che gli studenti si portino a casa è che \textbf{con molto meno}
si possono ottenere esperienze ``estremamente significative a livello
emozionale'' e creazioni di game design molto profonde \textbf{con budget zero}.
I casi grandi servono come \textbf{punti di riferimento} e per sapere che
esistono cose ``molto distanti da quelle che immaginiamo tutti i giorni''.

\subsection{L'ARG di Stranger Things a Milano}

Nello stesso filone, un ARG che ``è stato un'operazione di marketing e basta''
(``per me sono molto affascinanti, però era un'operazione di marketing''), ma
fatto \textbf{a Milano}: quello per una stagione di \emph{Stranger Things} (la
terza o la quarta, Beatrice non ricorda). Era una sorta di \textbf{caccia al
tesoro} in città, con punti di gioco preparati in modo ``altamente scenografico''
(nella slide si vede un autobus), mescolata con una \textbf{telefonata} a un
numero in cui uno dei personaggi diceva che il Sottosopra era arrivato a Milano e
che serviva aiuto. Il messaggio: anche da noi questi giochi ``sono arrivati,
arrivano, e credo che col tempo avranno sempre più spazio'', e per chi pensa in
termini di lavoro ``è un tema''.

% ---------------------------------------------------------------------
\section{Cacce al tesoro, corse e indagini}
\label{a2:sec:pervasivi-cacce}

\subsection{Le guide-caccia al tesoro su carta}

Beatrice presenta poi delle \textbf{guide per esplorare la città} in forma di
caccia al tesoro\footnote{La trascrizione riporta il nome come ``Waiwai''; la
grafia non è verificabile.}, uscite negli stessi anni dell'ondata dei pervasivi
in Italia (la prima nel 2007--2008), disponibili solo per alcune città e
vendute come \textbf{libretti} sul loro sito o da rivenditori autorizzati.
Tiziana le conosce: aveva comprato quelle di \textbf{Venezia} e di
\textbf{Verona}, ha giocato in parte quella di Venezia, nella versione breve
(c'erano due possibilità: un gioco di un paio d'ore o uno da una giornata
intera), e ricorda il \textbf{numero di cellulare} con cui interagire per avere
le indicazioni. Beatrice le aveva giocate con gli SMS, ``era un'altra vita'';
oggi pare si siano aggiornate a WhatsApp (e chiede aggiornamenti a chi le
giocherà).

Sono un altro ibrido di \textbf{analogico e digitale}: un'entità parla con i
giocatori tramite messaggi, mentre il libretto cartaceo resta in mano a loro. La
particolarità è la meccanica per passare da una tappa all'altra: \textbf{ogni
pagina del libro è divisa in strisce}, e quando si risolve una tappa la guida
comunica dei \textbf{codici} (``il foglietto 3, il 32C e il 18F'') con cui
ricomporre la pagina che porta alla tappa successiva.

\subsection{PacManhattan}

\slidefig{img/AGD/agd2-cacce}{PacManhattan, Zombies, Run! e Reptilarium.}{slide \emph{Analog -- Case Studies}}{fig:agd2-cacce}


\textbf{PacManhattan}\footnote{La trascrizione dice ``Pac-Man-Hatton''.}, degli
stessi anni, è \textbf{Pac-Man per le strade di Manhattan}: una persona fa
letteralmente Pac-Man, guidata da un'altra persona (``la mente'') tramite
telefono o walkie-talkie; altri giocatori fanno i fantasmi. È di nuovo ``una
sorta di guardie e ladri'', ed è uno dei primi urban game che si trovano nei
manuali di game studies. Beatrice lo mostra per un'idea precisa: si può
\textbf{prendere qualcosa che esiste già in una forma e cambiarle forma
completamente}. Non è lo stesso gioco di Pac-Man, ma si ispira alla sua
narrativa, alle sue meccaniche, alla sua estetica.

\subsection{Zombies, Run!}

\textbf{Zombies, Run!}\footnote{La trascrizione dice ``Zombie Run''.} è
un'\textbf{app per fare jogging}: invece di ascoltare musica, o i messaggi
vocali degli amici, si ascoltano \textbf{gli zombie che stanno per prenderti}, e
c'è tutta una narrativa su come si sta scappando mentre si corre. Si trova sugli
store.

\subsection{Reptilarium}

Si torna al teatro, ma in versione analogico-digitale. \textbf{Reptilarium} è un
caso studio di cui purtroppo c'è pochissima documentazione. È uscito durante il
\textbf{COVID}, da una compagnia teatrale di Roma, \textbf{Project
XX1}\footnote{Grafia come nella trascrizione.}. Era a modalità evento (un certo
giorno, a una certa ora) e le persone accedevano da \textbf{casa propria}, davanti
al computer, a un'esperienza con degli \textbf{attori} che si mettevano in
comunicazione con loro: ricevevano mail e telefonate, potevano vedere video
``recuperati da altre persone'', potevano \textbf{coordinare i movimenti degli
attori} e vedere che cosa succedeva mentre li eseguivano. La trama partiva da una
\textbf{persona scomparsa} da ritrovare, e poi si espandeva molto. Come per le
cene con delitto (dove si può interrogare l'attore o essere l'indiziato), anche
questo esperimento ``può diventare un LARP'', un ibrido tra teatro e gioco di
ruolo, o tra gioco investigativo e alternate reality game.

Ruggero aggiunge una nota: il \textbf{lockdown} ha colpito moltissimo i giochi
pervasivi, ma anche altri generi come i LARP, e in quel periodo sono nate
esperienze che hanno spostato \textbf{online e a distanza} giochi nati per essere
giocati di persona. Se ne riparlerà.

% ---------------------------------------------------------------------
\section{Pervasivi per la formazione: Echi resistenti}
\label{a2:sec:pervasivi-edu}

\slidefig{img/AGD/agd2-echi}{Un gioco pervasivo per la formazione storica.}{slide \emph{Analog -- Case Studies}}{fig:agd2-echi}


``Tutte le cose che vedete hanno enorme spazio per essere esplorate in tema
\textbf{educativo e formativo}'', e Beatrice chiude con un accenno a questo
ambito. \textbf{Pop History} è un'\textbf{associazione di promozione sociale}:
non si presenta come collettivo teatrale né come gruppo di game designer, ma come
associazione (``credo siano anche loro emiliani''), e fa tante piccole
sperimentazioni. Una di queste è \textbf{Echi resistenti}\footnote{La
trascrizione dice ``Ecchi Resistenti'' e ``Eh, che resistenti''.}, un gioco
sulla \textbf{Resistenza e sulla Liberazione}, che ibridava:
\begin{itemize}
  \item un'\textbf{indagine} con dei fascicoli;
  \item un'\textbf{attività teatrale}, con personaggi storici con cui si poteva
        interagire direttamente;
  \item una \textbf{caccia al tesoro}, in cui si andava a cercare questi
        personaggi per la città.
\end{itemize}
Altri casi simili sono sul loro sito.

\paragraph{Approfondimento: giochi con migliaia di giocatori.} In
\emph{Challenges for Game Designers}, nel capitolo sui giochi multigiocatore,
Brathwaite e Schreiber osservano che nei giochi da tavolo il problema
dell'interfaccia è evidentissimo (``come si mettono venti persone attorno a un
solo tabellone? E mille?''), eppure ``esistono giochi non digitali con migliaia
e perfino centinaia di migliaia di giocatori'': gli \textbf{alternate reality
game} e i grandi \textbf{urban game}, che ``portano i giocatori via dai loro
computer e dentro il mondo reale''. È la stessa scala dei casi di Beatrice:
\emph{Journey to the End of the Night} con migliaia di partecipanti a Milano,
\emph{Why So Serious} con le sue mailing list di fan, \emph{Pokémon Go} con il
mondo intero come mappa.

\begin{table}[htbp]
\centering
\small
\begin{tabular}{@{}p{3.6cm}p{2.3cm}p{7.0cm}@{}}
\toprule
\textbf{Caso studio} & \textbf{Formato} & \textbf{Perché è nella lezione} \\
\midrule
Pokémon Go & app & gli spazi quotidiani diventano spazi di gioco \\
Geocaching & piattaforma + scatole & caccia al tesoro collettiva: le scatole
non si tolgono mai \\
GuruShot & app & la fotografia come strumento; non si sa quando la partita
inizi e finisca \\
Journey to the End of the Night & evento urbano & guardie e ladri notturno,
completamente analogico, libero e non a pagamento \\
SF0 & piattaforma & missioni urbane documentate; la community valuta e crea
missioni \\
Critical City Upload & piattaforma & il ``figlio'' italiano di SF0
(2009--2013): missioni di relazione sul territorio \\
Sleep No More & teatro interattivo & il pubblico esplora liberamente un hotel
\\
Why So Serious & ARG di marketing & la torta con il cellulare: il Joker esiste
davvero \\
Gioco di MS Edizioni & gioco da tavolo + ARG & rapina in corso tra documenti
reali e finzione; ha inizio e fine \\
Cinetico 4.4, Amleto & performance & un gioco da tavolo che è uno spettacolo; il
pubblico sceglie Amleto \\
Lettere anonime per un camminatore & passeggiata audio & gioco in solitario,
voci degli abitanti, finale scenografico \\
MP3 Experiment & flash mob audio & la curiosità per gli oggetti da portare \\
The Institute & ARG & il caso blasonato, con documentario e serie TV \\
ARG di Stranger Things & ARG di marketing & un ARG fatto a Milano \\
Guide-caccia al tesoro & libretto + SMS & pagine in strisce da ricomporre con
codici \\
PacManhattan & urban game & un videogioco cambiato completamente di forma \\
Zombies, Run! & app & narrativa audio durante la corsa \\
Reptilarium & teatro a distanza & attori guidati da casa durante il COVID \\
Echi resistenti & gioco formativo & indagine, teatro e caccia al tesoro sulla
Resistenza \\
\bottomrule
\end{tabular}
\caption{I casi studio di giochi pervasivi presentati da Beatrice, nell'ordine
della lezione, con il motivo per cui sono stati scelti.}
\label{a2:tab:pervasivi}
\end{table}

% ---------------------------------------------------------------------
\section{I giochi di ruolo dal vivo: le macrocategorie}
\label{a2:sec:larp}

\slidefig{img/AGD/agd2-larp}{Le macrocategorie del LARP: boffer, Mind's Eye Theatre, parlor/chamber LARP, American freeform e murder party.}{slide \emph{Analog -- Case Studies}}{fig:agd2-larp}


Finita la parte sui pervasivi, si fa una pausa di pochi minuti (``facciamo
sette minuti, così andiamo ai 20 tondi''). Al rientro ci sono tutti, compreso
Tancredi (``ci sono tutti e due'' i Lorenzo), e Ruggero passa ai
\textbf{giochi di ruolo dal vivo}.

\paragraph{Una struttura diversa.} Questa parte è organizzata diversamente da
quella sui pervasivi. Il gioco di ruolo dal vivo, come il gioco di ruolo da
tavolo e il gioco da tavolo, ha alle spalle \textbf{anni di industria
solidificata}, ed è quindi più facile identificare \textbf{macrocategorie} con
nomi e formati specifici. Prima, però, Ruggero chiede chi abbia mai giocato un
LARP. Una persona della classe sì: un LARP fantasy, al \textbf{festival del gioco
di ruolo di Bracciano}, ``con le spade di gommapiuma, insomma, quelle cose lì''.
È il primo impatto di molti.

\subsection{Il boffer LARP}

Chi ha giocato un LARP probabilmente ha giocato un \textbf{boffer}, e chi ne ha
solo sentito parlare ha in mente quello. È uno dei primi giochi di ruolo dal
vivo diffusi, nato all'epoca proprio come \textbf{tentativo di adattare
\emph{Dungeons \& Dragons} dal vivo}, poi evoluto nel tempo. (Daniele mostra
una foto e Ruggero protesta: ``non evidenziarmi, che già è cringe''; ne ha
giocati anche lui.) Il nome viene dalle armi: le armi \emph{boffer} sono armi in
plastazote che devono avere caratteristiche precise per non essere pericolose.
Sono giochi \textbf{incentrati sul combattimento}, su trame epiche e così via.
Sono anche quelli che possono essere più grandi: un boffer anche grosso può
ospitare \textbf{400--500 persone}, ``per non dire quelli in Germania, che ne
tengono anche migliaia''.

\subsection{Il Mind's Eye Theatre}

Un altro tipo abbastanza diffuso è il \textbf{Mind's Eye Theatre}: se il boffer
deriva da \emph{Dungeons \& Dragons}, questo deriva da \textbf{\emph{Vampire: The
Masquerade}}. Le trame sono più dark, gotiche, horror, legate alla proprietà
intellettuale di Vampire, e chi partecipa di solito interpreta uno dei vampiri
delle loro società segrete. Sono giochi \textbf{più politici che di
combattimento}: l'intrigo, i segreti, il tradirsi a vicenda, ``il pugnalare
alle spalle qualcuno'' sono parte delle dinamiche. Come i boffer, sono giochi
\textbf{in costume}: il costume è richiesto, e si chiede a chi gioca di vestirsi
come il proprio personaggio.

\subsection{Chamber LARP e parlor LARP}

Il costume sembra scontato, ``ma in realtà non tutti i LARP usano il costume
come strumento necessario''. I \textbf{parlor LARP} o \textbf{chamber LARP}
sono giochi di ruolo dal vivo in cui il costume \textbf{non è richiesto}: si
gioca con i propri vestiti di tutti i giorni. Lo studente di prima interviene: a
Bracciano, in effetti, giocavano con i propri vestiti (``di solito gli eventi
così sono degli ibridi''); e ricorda di aver giocato un altro LARP al
\textbf{Play}\footnote{La trascrizione dice ``Play di Bologna''; il festival
Play, dedicato al gioco, si tiene a Modena.}, attorno a un tavolo, ognuno con un
ruolo e con i propri vestiti: \textbf{\emph{A cena col nonno}}. ``Perfetto, quello
è un chamber LARP''.

Ruggero precisa un punto sottile: alcuni chamber LARP si giocano \textbf{a un
tavolo}, ma questo non li rende giochi di ruolo da tavolo, perché si giocano
personaggi \emph{che sono seduti a un tavolo}, mentre in \emph{Dungeons \&
Dragons} il personaggio non è necessariamente seduto a un tavolo in quel momento.
Con gli altri LARP hanno in comune la \textbf{sovrapposizione tra la fisicità di
chi gioca e quella del personaggio}, l'\textbf{embodiment}. Rispetto ad altri
LARP, i chamber LARP hanno di solito:
\begin{itemize}
  \item un livello di \textbf{accessibilità} più alto, perché spesso non
        richiedono costumi;
  \item un \textbf{budget ridotto};
  \item \textbf{non sono a evento}: molti si scaricano in \emph{print and play},
        si stampano a casa e si giocano con gli amici;
  \item per questo, un \textbf{numero di giocatori} più ridotto, di solito da
        3 a 20 persone.
\end{itemize}

\subsection{L'American freeform}

Una categoria dei chamber LARP è l'\textbf{American freeform}: LARP progettati da
designer che si occupano soprattutto di giochi di ruolo da tavolo, e infatti è
un \textbf{genere ibrido}. Visti da fuori sono chamber LARP (gente che fa
\emph{embodiment} del personaggio in delle stanze, non necessariamente in
costume), ma hanno ``un sacco di meccaniche'' prese dal gioco da tavolo e dal
gioco di ruolo da tavolo: le cosiddette \textbf{metatecniche} o metameccaniche,
di cui si parlerà in un'altra lezione. Anche questi sono spesso in \emph{print
and play}; alcuni sono venduti in un formato più vicino al gioco da tavolo, come
mazzi di carte o libretti. Quindi non sono a evento, a differenza dei boffer,
del Mind's Eye Theatre e dei murder party.

\subsection{I murder party}

Ruggero riprende l'ottima distinzione fatta da Beatrice. ``Murder party'', da noi,
indica culturalmente due cose:
\begin{itemize}
  \item le \textbf{cene con delitto}, in cui si va al ristorante, degli attori
        mettono in scena un omicidio e la successiva indagine, e i commensali
        mangiano e intanto risolvono il caso. \textbf{Nessuno dei commensali è
        l'assassino}: è uno degli attori;
  \item la categoria ``più british'', i \textbf{murder party all'inglese}, o i
        loro equivalenti che si stanno diffondendo moltissimo in Asia in questo
        periodo\footnote{La trascrizione dice ``Jubel Shock'' e ``Juvenile
        Show''. Si tratta quasi certamente dei \emph{jubensha}, i giochi di
        investigazione con ruoli molto diffusi in Cina.}, in cui \textbf{si
        interpretano i personaggi della storia} su cui si indaga: tra i giocatori
        può esserci l'assassino.
\end{itemize}
La seconda forma rientra, per il design dell'esperienza e delle meccaniche,
sotto l'ombrello dei LARP; gli aspetti ambientali e le decorazioni la
avvicinano poi al mondo del pervasivo.

% ---------------------------------------------------------------------
\section{Il Nordic LARP}
\label{a2:sec:nordic}

\slidefig{img/AGD/agd2-nordic}{Il Nordic LARP e i Knutebooks, i volumi collettivi del Knutepunkt.}{slide \emph{Analog -- Case Studies}}{fig:agd2-nordic}


C'è poi ``una grossissima categoria'', che si rivedrà più avanti perché ``è stata
importantissima'': il \textbf{Nordic LARP}, il LARP progettato in
\textbf{Scandinavia} (Danimarca, Norvegia, Svezia e Finlandia).

\subsection{Perché è così importante}

Ha un'importanza gigantesca nel settore \textbf{perché ha documentato tutto}. La
maggior parte dei libri accademici di game design e di analisi di questo tipo di
giochi, gli eventi e le conferenze, si fanno in Scandinavia, e c'è una lunga
storia di analisi e documentazione: la prima conferenza è del \textbf{1997}, e da
allora è annuale. È ``gente che si incontra da parecchi anni a parlare di come si
progetta e come si gioca''. Ed è una scena molto \textbf{cross-mediale}: prende
dagli alternate reality game, dai LARP americani, dai parchi a tema, dalle
simulazioni economiche, da tante teorie parallele a quelle del gioco, cercando
di interiorizzare, iterare e integrare tecniche di altri settori. Anche il Nordic
LARP ha le sue sottocategorie: alcune esperienze sono più vicine al chamber
LARP, altre al boffer.

\subsection{Come si riconosce}

Da fuori, un Nordic LARP si riconosce da \textbf{come è progettato il design
dell'esperienza}:
\begin{itemize}
  \item un focus sul \textbf{coinvolgimento}, sull'\textbf{emozione} e sulla
        \textbf{coesione tematica};
  \item pochissime cose \textbf{vestigiali}: per esempio, anche LARP molto
        grandi possono non avere il combattimento, se il combattimento non fa
        parte del tema;
  \item una cura del \textbf{post-esperienza}: molti giochi hanno una
        struttura per ``che cosa succede dopo il gioco''.
\end{itemize}
Molti concetti nati da questa scena si sono ``solidificati nel gergo del game
design'' e torneranno nella lezione sull'emotional design.

\paragraph{Origini diverse da D\&D.} Storicamente, il LARP nordico ha
un'origine \textbf{precedente a \emph{Dungeons \& Dragons}}. Si tende a pensare
che il LARP nasca da ``prendiamo D\&D e lo facciamo dal vivo'', e per alcune
nazioni, come l'Inghilterra e l'America, è proprio così; per altre no. Il Nordic
LARP prende strategie e strumenti dal \textbf{teatro}, dalle avanguardie, dallo
\textbf{psicodramma}, dalla \textbf{Gestalt therapy}; ``è ovvio che quello che
esce fuori è un gioco molto particolare''.

\paragraph{I manifesti.} È anche un ambiente ``pieno, carico, zeppo di
manifesti''. La slide ne mostra alcuni, nati negli anni nella scena; non si
approfondiscono ora, ma sono tutti pubblici e si possono leggere a parte.

\subsection{Knutpunkt e i Knutepunkt books}

La cosa che interessa di più è la conferenza citata prima, che cambia nome a
seconda del paese che la ospita: \textbf{Knutpunkt}, \textbf{Knudepunkt},
\textbf{Knutepunkt} o \textbf{Solmukohta}\footnote{Rispettivamente in svedese,
danese, norvegese e finlandese: la conferenza ruota tra i quattro paesi. La
trascrizione dice ``Knutpunkt o Knutpunkt o Knudpunkt o Solmukott'' e, più
avanti, ``Knut Punk''.}. È il principale punto di ritrovo scandinavo sul gioco
di ruolo dal vivo: \textbf{quattro giorni di conferenza non stop}, con un
programma di centinaia di slot, spesso \textbf{otto track in parallelo} di talk,
panel e workshop. Si tiene dal 1997, e \textbf{ogni anno produce dei libri}
che riassumono lo stato dell'arte in forma di articoli: la documentazione è
``costantemente in aggiornamento''. Ruggero ne ha diversi, ``dei bei mattoni'':
ne mostra uno di \textbf{600 pagine} su come si progetta un LARP, uscito da una
delle edizioni (``nel 2018, se non ricordo male''). È forse ``una delle
comunità ludiche che più documenta quello che fa'', e soprattutto lo documenta
in modo \textbf{pragmatico}, lasciando framework e strumenti a chi vuole creare
giochi di questo tipo. Molte di queste cose si applicano anche ad altri generi, e
torneranno.

\subsection{Dogma 99}

In chat qualcuno chiede se c'entri il \textbf{Dogma} di Lars von Trier.
``Sì e no'': il \textbf{Dogma 99} è un manifesto di LARP che riprende il
manifesto cinematografico del \textbf{Dogma 95} per fare un parallelo con il
LARP. Il Dogma 95 aveva alcune direttive, per esempio che tutto debba essere
\textbf{diegetico}, senza nulla di extradiegetico; il Dogma 99 fa lo stesso: in
un LARP storico non ci devono essere cose extradiegetiche rispetto alla storia.
Una conseguenza: il Dogma 99 \textbf{non fa LARP fantasy}, perché la magia non
esiste e non si può riprodurre in modo perfetto; si fanno giochi ambientati
tendenzialmente ai giorni nostri, che si possono riprodurre facilmente. Da
questa catena di ragionamenti sono nati ``giochi molto interessanti''.

% ---------------------------------------------------------------------
\section{Jeepform: il dramma quotidiano}
\label{a2:sec:jeepform}

\slidefig{img/AGD/agd2-jeepform}{Jeepform tra tabletop, freeform e LARP; la locandina di Fat Man Down.}{slide \emph{Analog -- Case Studies}}{fig:agd2-jeepform}


Tra i manifesti c'era lo \textbf{Jeepform}\footnote{La trascrizione dice
``G-Form'' e ``Jeep Form''.}, nato all'interno del \textbf{freeform}, la zona ``a
metà tra il LARP e il tabletop''. Il suo manifesto, che si trova online, ha tra i
punti più importanti questo: \textbf{il dramma quotidiano è più interessante del
dramma epico}. Tutti i giochi Jeepform sono giochi di ruolo dal vivo su
\textbf{situazioni quotidiane}, ``drammaticamente quotidiane''.

Alcuni sono emotivamente molto pesanti, e storicamente importanti per vari
motivi; Ruggero ha lasciato fuori i più pesanti (se ne riparlerà nella lezione
sull'emotional design) e ne presenta due ``più leggeri'', che fanno comunque
capire in che toni si va.

\subsection{Il gioco della coppia perfetta}

Il primo\footnote{La trascrizione riporta il titolo come ``Lady 8'' e i nomi dei
personaggi come ``Lady'' e ``8''; né il titolo né i nomi sono verificabili.} è
un gioco per \textbf{quattro persone} e \textbf{due personaggi}, che per premessa
\textbf{si amano tantissimo}. Due persone alla volta interpretano i personaggi,
in un gioco \textbf{diviso in scene}; ogni scena è descritta dalle due persone che
non sono in gioco. L'unico obiettivo di chi gioca la coppia è \textbf{essere la
coppia perfetta}, mostrare quanto si amano. Se la scena si sposta su qualcos'altro,
\textbf{si rifà}, all'infinito, ``come nel giorno della marmotta'', finché non
parla solo di quanto si amano. Esempio: lei entra in bagno, la porta si rompe e
lei rimane chiusa dentro, lui è dall'altra parte; se lui comincia a riparare la
porta, si rifà la scena. ``Una paradossalità esasperata, quasi''.

\subsection{Fat Man Down}

Il secondo è introdotto con un \emph{content warning}: \textbf{Fat Man Down} ``è
un gioco molto pesante''. Non è il più pesante dei Jeepform, ma è il più pesante
di cui si parla nella lezione. È un LARP dello stile Jeepform, quindi un chamber
LARP tendenzialmente freeform.
\begin{itemize}
  \item Il personaggio principale è un uomo di successo, con un buon lavoro e
        una buona posizione economica, ma \textbf{obeso}.
  \item Tutti i partecipanti \textbf{si pesano su una bilancia}, e chi pesa di
        più interpreta questo personaggio.
  \item Il gioco è diviso in scene, ognuna con un \emph{framing} che è uno
        \emph{slice of life}: un'ora in ufficio, la strada da casa al
        supermercato, dieci minuti di pausa al parco.
  \item Tutte le altre persone in scena hanno la \textbf{direttiva meccanica} di
        farlo stare malissimo.
  \item Il protagonista ha una meccanica ``importantissima'': \textbf{tre
        cioccolatini} in tasca. In qualunque momento può scartarne uno e
        mangiarlo, e quando lo fa tutti i personaggi in scena devono trattarlo
        come la persona migliore del mondo.
\end{itemize}
Il personaggio è messo sotto \textbf{costante stress narrativo} finché non è
spinto a pensare ``devo mangiare per stare meglio''. ``La roba è molto pesante'',
e ci si ferma qui.

\subsection{Le meccaniche di sicurezza}

Avere un manifesto così noto nella scena nordica, con giochi molto pesanti, molto
intimi e soprattutto \textbf{senza un filtro fantasy o fantascientifico}, cioè
su drammi quotidiani, ha generato e normalizzato cose poi rivelatesi molto utili,
come le \textbf{meccaniche di sicurezza}. Se qualcuno ne ha sentito parlare, la
scena nordica è ``la principale responsabile'' della normalizzazione, della
ricerca e dell'iterazione delle meccaniche di sicurezza nel gioco di ruolo. Se ne
parlerà nella lezione sull'emotional design.

% ---------------------------------------------------------------------
\section{Il 360° e l'alta risoluzione}
\label{a2:sec:360}

\slidefig{img/AGD/agd2-360}{Il 360° e l'alta risoluzione: Ground Zero e Monitor Celestra.}{slide \emph{Analog -- Case Studies}}{fig:agd2-360}


Il Nordic LARP non è solo chamber, freeform e Jeepform. C'è anche ``quasi
l'opposto'': i LARP \textbf{360°} o ad \textbf{alta risoluzione} (\emph{high
resolution})\footnote{La trascrizione dice ``360 o High Resolution'' e, più
avanti, ``IRES 360'', ``iRes'', ``Ires''.}. Lo scopo è essere \textbf{il più
fisicamente coinvolgenti possibile}: tutto quello che si ha attorno deve essere
\textbf{diegetico}. Ma la soluzione non è quella del Dogma 99 (ambientare nel
quotidiano, così non serve la magia): ``la soluzione è buttare al muro un budget
talmente alto che la realtà a quel punto sembra magia''. Niente oggetti di scena
finti o armi finte: anche la scenografia è curata al punto che sembra di stare
davvero nel gioco. I casi studio sono ``anche abbastanza storici''.

\begin{description}
  \item[Ground Zero (1998--2001).] La premessa è la \textbf{Guerra Fredda}.
        Una cittadina americana riceve l'allarme per una bomba nucleare, e i
        partecipanti vivono \textbf{48 ore di fila}, non stop, chiusi in uno
        \textbf{scantinato} al buio, senza poter uscire. Il gioco si conclude
        con la bomba che cade: dietro le finte pareti di cartone dello
        scantinato c'erano dei \textbf{subwoofer giganti}, e alla fine i muri
        cadono addosso ai giocatori con un rumore devastante.
  \item[Panopticorp (2003, 2013).] Due \emph{run}, anche questo di 48 ore non
        stop: in teoria si dorme in gioco, ``ma in questo caso non si dormiva''.
        È una \textbf{distopia corporate} spinta al massimo: si è dipendenti di
        un'agenzia pubblicitaria sull'orlo della chiusura, che deve licenziare
        ``a manetta'' e tenere solo i migliori. Una spinta ultra competitiva e
        ``ultra tossica'' per vedere fin dove si è disposti a spingersi eticamente
        per restare in azienda, dormendo sotto le scrivanie.
  \item[Monitor Celestra.] Un 360° \textbf{fantascientifico} basato su
        \emph{Battlestar Galactica}: per giocarlo è stato noleggiato un
        \textbf{incrociatore svedese} in disarmo. Una nave, non un'astronave,
        ``ma se non usciamo all'esterno, ci mettiamo un paio di LED, ed ecco
        qua'': per i giocatori è un'astronave. Il gioco durava circa
        \textbf{60 ore}, con cabine allestite in modo che si dormisse in
        personaggio.
  \item[Demeter.] Un 360° sul viaggio della bara di \textbf{Dracula} dalla
        Transilvania a Londra, come nel romanzo di Stoker, in cui la nave
        \emph{Demeter} arriva in porto con tutti morti a bordo e il capitano
        legato al timone. I giocatori sono i passeggeri, e \textbf{da giocatori
        sanno che moriranno tutti}: per design è impossibile sopravvivere,
        ``altrimenti Dracula non esisterebbe come romanzo''. Si gioca in
        costume su una \textbf{nave vera} che, naturalmente, non va dalla
        Transilvania a Londra ma gira nel mare della Germania del nord;
        narrativamente, però, si gioca quel viaggio.
  \item[Bunker 101.] Un caso \textbf{italiano}, ispirato alla proprietà
        intellettuale di \emph{Fallout}: si giocava \textbf{72 ore} in un bunker.
        Qui il dormire non era diegetico: nell'ultima run hanno separato la parte
        di gioco dal sonno. Per tutto il resto, però, valeva il principio 360°:
        tutto ciò che si vede deve essere diegetico.
\end{description}

% ---------------------------------------------------------------------
\section{Tape LARP e black box}
\label{a2:sec:blackbox}

\slidefig{img/AGD/agd2-tape}{Tape LARP e black box LARP.}{slide \emph{Analog -- Case Studies}}{fig:agd2-tape}


Accanto al Jeepform (``in cui giochiamo in casa con i nostri costumi'') e al
360° (``l'ultra-diegesi''), ci sono tipologie scandinave ``molto più particolari
e quasi astratte''.

\subsection{Il tape LARP}

Visto che qualcuno conosce Lars von Trier: chi ha visto \emph{Dogville}, il film
con lo scotch per terra, sa da dove nascono i \textbf{tape LARP}. ``Abbiamo visto
Dogville, ci ispira un casino'': giochi in cui \textbf{lo scotch per terra indica
gli ambienti}, con tutte le meccaniche che ne derivano (il personaggio vede e
sente cose che normalmente non potrebbe sentire; come si agisce e non si agisce
da giocatori\ldots). Ne sono stati fatti parecchi, in particolare in Danimarca.

\subsection{Il black box LARP}

Ruggero è ``un particolare amante'' del \textbf{black box LARP}, che prende il
nome dalle \emph{black box}: per chi ha fatto teatro ``dalla Germania in su'',
spazi teatrali di solito senza spalti né sedie, con \textbf{pareti, soffitto e
pavimento neri}, molto diffusi in Germania e nel Nord Europa. A
\textbf{Copenaghen} c'è un festival, il \textbf{Black Box LARP Festival}, che usa
questi spazi per giochi creati apposta: giochi di ruolo che usano \textbf{luci e
colori come meccaniche di gioco}. Gli esempi della slide:

\begin{description}
  \item[Borderline.] Un gioco ``estremamente astratto''. Si gioca \textbf{in
        coppie}, e ogni coppia gioca \textbf{un solo personaggio}, con dei
        conflitti interni: tutta la conversazione è un monologo interiore del
        personaggio, e i due giocatori, con una candela, lottano per portare ciò
        che il personaggio fa dalla propria parte o da quella dell'altro.
  \item[Waiting for Flight.] Ruggero lo ha giocato. Si aspetta un aereo con a
        bordo persone care, amici, colleghi; \textbf{l'aereo precipita}, e
        muoiono quasi tutti. Si gioca la \textbf{sala d'attesa} di chi li
        aspettava: è un gioco sull'attesa, ``sul logorante non avere
        informazioni'' finché vengono centellinate. Si conclude con un momento
        molto catartico: arrivano le \textbf{valigie} delle persone attese, e in
        alcune (3 su 20) è nascosto un piccolo \textbf{cuore di vetro}: se lo
        trovi, la persona si è salvata; se no, non si è salvata. C'è una fase in
        cui speri che ci sia, ma potrebbe essere nascosto: ``è una cosa
        lacerante. Abbiamo pianto in 19 su 20''.
  \item[Stranger.] Di una designer danese\footnote{La trascrizione dice ``Nino
        Ressendrop'': con ogni probabilità Nina Runa Essendrop, designer danese
        nota per i LARP fisici e non verbali. Anche il titolo, ``Stranger'', è
        riportato come nella trascrizione.}. È un black box \textbf{non verbale}:
        è vietato parlare. I partecipanti sono divisi in due gruppi, i
        \textbf{blu} e i \textbf{verdi}, e prima del gioco ciascun gruppo crea un
        \textbf{linguaggio gestuale condiviso}: come la propria nazione esprime
        felicità, rabbia, paura, tristezza, amicizia, odio. Il twist: i verdi
        devono avere gesti squadrati, rigidi, geometrici; i blu gesti curvi,
        lenti, morbidi, sinuosi. Il gioco è sull'\textbf{immigrazione forzata}:
        ci sono fasi di giorno e di notte, e di notte chi facilita porta alcuni
        verdi nella nazione blu, dove si risveglieranno, con la regola che il
        confine non si può riattraversare. Le dinamiche: il primo che si
        risveglia tra i blu impara facilmente la loro lingua, e i blu lo aiutano
        con un \emph{onboarding} non verbale; ma quando arrivano molti verdi,
        tendono a muoversi \textbf{in gruppi}, perché a quel punto c'è già chi
        parla la propria lingua, anche se non è la lingua della nazione. ``Un
        microscopio'' di dinamiche della vita reale: emigrazione, immigrazione,
        inclusione, e quanto è più facile essere inclusi se il sistema ti
        supporta.
  \item[Il gioco di Ruggero.] Uno dei giochi della slide è suo, ``quindi non ne
        parliamo'', nonostante le proteste dei colleghi.
  \item[Un gioco su Cuba.] Di una designer che è andata a vivere a
        \textbf{Cuba} e ha partecipato da vicino alla vita di una delle storiche
        comunità queer cubane\footnote{La trascrizione riporta il titolo come
        ``Socialism and Word''; non è verificabile. Anche le date sono incerte
        (``nel 97, 89 e 91''): il contesto, la caduta del Muro di Berlino, fa
        pensare al 1987, 1989 e 1991.}. Il gioco racconta che cosa è successo
        alle comunità \textbf{punk, rock e queer}, molto sovrapposte, prima,
        durante e dopo la caduta del Muro di Berlino: lo Stato, rimasto senza i
        fondi sovietici, decide chi è cittadino di serie A e di serie B, e le
        persone queer e punk (la musica occidentale) vengono picchiate e
        imprigionate. Erano anche gli anni dello scoppio dell'\textbf{AIDS}. Una
        comunità famosa, quella dei \emph{frikis}\footnote{La trascrizione dice
        ``Os Fricchios''.} cubani (di cui sono ancora vive due persone), decise:
        ``tanto ci ammazza la polizia, tanto vale iniettarci sangue infetto da
        HIV'', per andare a morire non per strada ma nei \textbf{centri di
        recupero}, gestiti da medici stranieri, uno dei pochissimi casi in cui a
        Cuba poteva esserci personale non cubano. Lì si moriva, ma ``almeno si
        poteva ascoltare musica punk'', magari con il dottore che arrivava dalla
        Florida con le cassette rock, e tra simili. ``Una roba drammaticissima'':
        anche questo è un gioco devastante.
\end{description}

\begin{table}[htbp]
\centering
\small
\begin{tabular}{@{}p{3.1cm}p{9.8cm}@{}}
\toprule
\textbf{Tipo di LARP} & \textbf{Caratteristiche ed esempi citati} \\
\midrule
Boffer & nato per adattare D\&D dal vivo; costume, armi in plastazote,
combattimento e trame epiche; fino a centinaia o migliaia di persone \\
Mind's Eye Theatre & da \emph{Vampire: The Masquerade}; costume, dark/gotico,
intrigo politico più che combattimento \\
Chamber / parlor LARP & senza costume, accessibile, budget ridotto, spesso
\emph{print and play}, 3--20 persone (\emph{A cena col nonno}) \\
American freeform & chamber LARP con metatecniche prese dal gioco da tavolo e
dal GDR; anche in carte o libretti \\
Murder party & cena con delitto (gli attori indagano, l'assassino è un attore)
o murder party ``all'inglese'' (si interpretano i sospettati) \\
Nordic LARP & coinvolgimento, emozione, coesione tematica, cura del dopo-gioco;
manifesti; Knutpunkt dal 1997; Dogma 99 \\
Jeepform & freeform nordico sul dramma quotidiano (il gioco della coppia
perfetta, \emph{Fat Man Down}); meccaniche di sicurezza \\
360° / alta risoluzione & tutto diegetico grazie a un budget altissimo
(\emph{Ground Zero}, \emph{Panopticorp}, \emph{Monitor Celestra},
\emph{Demeter}, \emph{Bunker 101}) \\
Tape LARP & lo scotch per terra indica gli ambienti, come in \emph{Dogville} \\
Black box & spazi teatrali neri; luci e colori come meccaniche
(\emph{Borderline}, \emph{Waiting for Flight}, \emph{Stranger}) \\
\bottomrule
\end{tabular}
\caption{Le tipologie di gioco di ruolo dal vivo presentate da Ruggero, in
sintesi.}
\label{a2:tab:larp}
\end{table}

% ---------------------------------------------------------------------
\section{Chi osserva? Il LARP e il teatro dell'oppresso}
\label{a2:sec:oppresso}

\paragraph{La domanda di Tiziana.} Tiziana confessa un'impressione: tutti questi
giochi ``sembrano cose da studi'', come se ci fosse qualcuno che sta sopra e
osserva le dinamiche per ricavarne informazioni. Ruggero la rassicura: \textbf{non
è così}. I designer di questa scena, specialmente quella black box, vengono
dall'ambiente ibrido del teatro interattivo, degli \emph{happening}, dello
psicodramma, e \textbf{non documentano nulla di quel tipo}: non registrano le
giocate, non estrapolano informazioni da chi gioca. Spesso sono giochi per poche
persone (è difficile che un black box superi le 20), e molti hanno avuto
\textbf{run singole}, o al massimo due o tre. ``Si tratta solo di fare un certo
tipo di esperienza''. Come in altri media esistono generi con scopi diversissimi
(un film e un documentario sono entrambi operazioni visive), anche qui gli scopi
sono diversi. E chi non ha mai incontrato giochi del genere si chiederà
``perché dovrei giocare una roba del genere?'': la risposta arriverà con la
lezione sull'emotional design (``lo so, ve l'ho detto quattro volte'').

\paragraph{Il teatro dell'oppresso.} Prima di passare la parola, uno studente
interviene: in queste cose ritrova molto il metodo del \textbf{teatro
dell'oppresso}\footnote{Il metodo teatrale ideato dal regista brasiliano Augusto
Boal (integrazione).}. ``Ah, certo''. Ruggero torna indietro di una slide, a
quella sulle \textbf{influenze fondanti del Nordic LARP}: ci sono le tecniche di
\textbf{Moreno} (sociodramma e psicodramma), il \emph{sensitivity training}, il
\textbf{roleplay come metodo di ricerca}, e tutti quei metodi che usano il teatro
e l'interazione come strumento per far emergere il personale e l'intimo. Lo
studente precisa che il teatro dell'oppresso nasceva per \textbf{trovare
soluzioni} a situazioni che apparentemente non ne hanno: a turno il pubblico
proponeva una soluzione, si rifaceva sempre la stessa scena, e si studiava come
reagivano le persone; quindi, in un certo senso, il comportamento delle persone
\emph{si osserva}, ma dall'interno della scena.

Per Ruggero, nella scena LARP questa dimensione ``viene un po' meno'', perché lo
scopo resta \textbf{ludico} (``e scopo ludico può voler dire molte cose, come
vedremo''): è meno pragmatico del ``prendiamo quello che abbiamo giocato e il
gioco ci dà già gli strumenti per applicarlo''. Alcuni \textbf{debriefing}, i
momenti strutturati del dopo-gioco, possono dare strumenti a chi ha giocato; ma
questi giochi nascono come esperienze che devono essere ``ludicamente
interessanti per chi gioca'', \textbf{non come esperimenti scientifici ``con dei
topi dentro un labirinto''}.

\paragraph{Approfondimento: il LARP in classe e il debriefing.} Che il LARP
possa essere uno strumento formativo, come nel volume sull'edu-larp citato
all'inizio, lo mostra anche \emph{Challenges for Game Designers}. Nel capitolo
sui giochi come strumento di insegnamento, Brathwaite e Schreiber propongono una
sfida, \emph{Living History}: progettare un LARP che rappresenti un evento storico
(una convenzione costituzionale, un protocollo internazionale, una guerra)
giocabile da una classe di 20--30 studenti in una o due ore. I passi suggeriti
sono: scegliere un evento in cui molte parti, ciascuna con la propria agenda, si
incontrano in una negoziazione o in un conflitto; scegliere i giocatori (o le
squadre); definire gli \textbf{obiettivi} di ciascuno; creare regole semplici e
rapide per risolvere i conflitti, perché i giocatori si muovono e parlano di
continuo (``alcuni LARP usano un singolo tiro di dado o una partita a
sasso-carta-forbici''); e preparare, oltre alle schede dei personaggi, delle
\textbf{note di debriefing}: che cosa è successo davvero e quali domande
discutere confrontando la storia con l'esito della partita. Lo scopo dichiarato è
che gli studenti capiscano che ``l'evento era complicato'' e come l'esito sia
\emph{emerso} dalla situazione iniziale. È la versione didattica della cura del
dopo-gioco che Ruggero attribuisce alla scena nordica.

% ---------------------------------------------------------------------
\section{I giochi di ruolo da tavolo: quattro parole per orientarsi}
\label{a2:sec:gdr}

``Daniele, ti lascio la parola''. Daniele (``mi giri tu le slide quando te le
chiedo'') avverte che fare una panoramica dei giochi di ruolo in poco tempo è
``veramente difficile'': ce ne sono tantissimi, e negli anni si sono evoluti e
diversificati sempre di più, fino a diventare molto diversi tra loro.

\subsection{Il tribalismo e le regole tramandate a voce}

C'è anche un altro problema, che si vedrà nella lezione sulla fiction e che
esiste in tutto il gioco analogico: il \textbf{tribalismo}, cioè
l'\textbf{ereditarietà orale delle regole}. Spesso i giochi di ruolo venivano
insegnati da altre persone, anche per motivi economici. Daniele ha cominciato
così: da bambino, giocando per la prima volta a D\&D, l'unico ad avere il
manuale era il master; ``ergo io non avevo accesso alle regole'', e il master
poteva inventarsi quello che voleva.

\subsection{Conflitto, azione, risoluzione, effetto}

\slidefig{img/AGD/agd2-car}{Conflitto, azione, risoluzione, effetto: le quattro parole per orientarsi nei giochi di ruolo da tavolo.}{slide \emph{Analog -- Case Studies}}{fig:agd2-car}


Per capire un po' i giochi di ruolo bisogna ricordare \textbf{quattro parole},
che torneranno nella lezione sulla fiction: \textbf{conflitto}, \textbf{azione},
\textbf{risoluzione} ed \textbf{effetto}. Nella slide \emph{conflitto} e
\emph{risoluzione} hanno un asterisco: per ora basta capirli quanto serve a
leggere i casi studio, senza analizzarli troppo.
\begin{description}
  \item[Conflitto.] L'ostacolo, l'incertezza, l'evento che \textbf{stimola
        un'azione o la contrasta}. ``Tornerà molto''.
  \item[Azione.] Quello che si descrive di voler fare nei panni del personaggio
        per risolvere un conflitto.
  \item[Risoluzione.] Come si risolve l'azione dentro il gameplay, con le regole:
        che cosa succede?
  \item[Effetto.] La conseguenza che si applica.
\end{description}
L'esempio: ``voglio colpire il goblin con la mia spada''. Il goblin ha il suo
scudo e la sua armatura. Che azione fai? ``Lo becco con la spada''. Si tirano i
dadi, si risolve, si applica un effetto. Ma non è così semplice, perché in
mezzo c'è la \textbf{fiction}: nel gioco di ruolo si tende a riscrivere tutto
quello che succede nel gameplay \emph{dentro} la narrazione. Se faccio tre danni
agli HP del goblin, che cosa vuol dire narrativamente? Come li descrivo? Da qui
nasce ``una serie di problematiche''.

\begin{figure}[htbp]
\centering
\begin{tikzpicture}[
    font=\small,
    passo/.style={draw, thick, rounded corners=3pt, fill=gray!12,
                  minimum width=2.6cm, minimum height=0.95cm, align=center},
    esempio/.style={font=\footnotesize\itshape, align=center, text width=2.8cm},
    freccia/.style={-{Stealth[length=2.5mm]}, thick}
  ]
  \node[passo] (c) at (0,0)    {\textbf{Conflitto}$^{*}$};
  \node[passo] (a) at (3.4,0)  {\textbf{Azione}};
  \node[passo] (r) at (6.8,0)  {\textbf{Risoluzione}$^{*}$};
  \node[passo] (e) at (10.2,0) {\textbf{Effetto}};
  \draw[freccia] (c) -- (a);
  \draw[freccia] (a) -- (r);
  \draw[freccia] (r) -- (e);
  \node[esempio, below=0.25cm of c] {il goblin, con scudo e armatura};
  \node[esempio, below=0.25cm of a] {``lo colpisco con la spada''};
  \node[esempio, below=0.25cm of r] {tiro dei dadi secondo le regole};
  \node[esempio, below=0.25cm of e] {tre danni agli HP};
  \node[draw, dashed, rounded corners=3pt, fit=(c)(e), inner sep=0.5cm,
        yshift=-0.35cm, minimum height=2.6cm] (fic) {};
  \node[font=\footnotesize, anchor=south] at (fic.north)
    {\textbf{fiction}: tutto viene ridescritto nella narrazione (che cosa
    \emph{sono}, nella storia, tre danni?)};
\end{tikzpicture}
\caption{Le quattro parole di Daniele per leggere i giochi di ruolo da tavolo,
con l'esempio del goblin. Gli asterischi sono quelli della slide: conflitto e
risoluzione verranno analizzati meglio nella lezione sulla fiction. Nel gioco di
ruolo ``classico'' il conflitto precede l'azione, la risoluzione è un tiro di
dado e l'effetto è descritto dal master e dal regolamento; molti giochi
successivi cambiano proprio questo schema.}
\label{a2:fig:car}
\end{figure}

% ---------------------------------------------------------------------
\section{Le origini: da Dungeons \& Dragons agli anni Novanta}
\label{a2:sec:gdr-storia}

Le quattro parole servono a Daniele per una breve panoramica, anche storica,
dei giochi di ruolo.

\subsection{Dungeons \& Dragons e il gioco di guerra}

Il primo gioco di ruolo da tavolo ``così come viene concepito'' è sicuramente
\textbf{\emph{Dungeons \& Dragons}}, di \textbf{Gygax e Arneson}, nel
\textbf{1974}. C'è in realtà un precedente, \textbf{\emph{Chainmail}}\footnote{La
trascrizione dice ``C'è in Mail'' e data \emph{Chainmail} al '72; il gioco è
del 1971.}, e come per tutte le origini (``anche \emph{Cyberpunk} ha lo stesso
problema'') c'è un po' di ambiguità su che cosa venga prima. Ci sono libri che
raccontano la nascita di D\&D, c'è la storia della \textbf{TSR}, la casa
editrice che lo pubblicò per prima, ci sono versioni discordanti e, negli anni,
parecchi contrasti tra le famiglie di Gygax, di Arneson e di altri.

D\&D viene dal \textbf{gioco di guerra}, il gioco di \textbf{miniature}. ``La
leggenda vuole'' che Gygax e Arneson, giocando, si siano detti: ho battuto il
tuo esercito; che cosa succede se il mio generale entra nel tuo castello e va a
trattare con il tuo generale? Abbiamo regole per questo? No: inventiamocele.
``Ripeto, una leggenda'': non si sa se sia andata così, e in fondo non interessa
molto.

\subsection{I generi e la struttura del gioco di ruolo ``early''}

\slidefig{img/AGD/agd2-dnd}{Il flusso di Dungeons \& Dragons letto con conflitto, azione, risoluzione ed effetto.}{slide \emph{Analog -- Case Studies}}{fig:agd2-dnd}


Quello che interessa è che dopo D\&D, nell'\textbf{RPG delle origini}, vengono
coperti più o meno \textbf{tutti i generi letterari}: \emph{Call of Cthulhu}
l'orrore lovecraftiano, \emph{Vampire: The Masquerade} l'orrore gotico,
\emph{Traveller} la fantascienza, \emph{Cyberpunk} il cyberpunk, e \emph{GURPS},
un \textbf{sistema generico} con cui ``teoricamente'' si può giocare qualunque
genere.

Nel gioco di ruolo classico:
\begin{itemize}
  \item il \textbf{conflitto} poteva venire prima o dopo l'\textbf{azione}:
        ``voglio colpire quel goblin'' (e il goblin ha la sua difesa), oppure
        ``corri, ti sta inseguendo una palla'', come in \emph{Indiana Jones};
  \item la \textbf{risoluzione} era il tiro di dado, e l'\textbf{effetto} veniva
        descritto dal master e dal regolamento;
  \item in pratica, \textbf{tutto tranne l'azione}, che il giocatore gestiva e
        narrava, era gestito dal \textbf{master}.
\end{itemize}

\paragraph{Il master come arbitro.} All'inizio il master non era un ``maestro''
né una figura che insegna: era interpretato soprattutto come un
\textbf{arbitro}. Nella prima versione di D\&D non c'è ancora l'idea, venuta
dopo, del master come ``dio del tavolo'': quando deve decidere, dovrebbe farlo
secondo criteri di realismo e oggettività (``puoi essere oggettivo a un tavolo
di gioco? La voglio vedere''). Questa idea è importantissima perché verrà ripresa
dall'\textbf{OSR}, l'\emph{Old School Revival}, il movimento che cerca di
riportare nella modernità il modo in cui, a suo dire, si giocava D\&D agli
inizi.

\paragraph{Il master è un giocatore.} Nei giochi di ruolo ci sono quindi
\textbf{ruoli asimmetrici}. Ma Daniele raccomanda di ricordare una cosa, che
tornerà nella lezione sulla fiction: \textbf{il master è un giocatore}, magari un
giocatore con più regole e più autorità al tavolo, ma ``un giocatore tanto quanto
i giocatori che interpretano un personaggio''. Nei GDR delle origini: si
interpreta un personaggio, con una \textbf{scheda} e le sue statistiche, dentro
una fiction e un \textbf{setting} stabilito dal manuale o dal master; quando si
dice che cosa si vuole fare, il master reagisce. E lentamente, con gli anni, il
master assume anche il ruolo di \textbf{regista} e \textbf{sceneggiatore}: scrive
l'avventura, decide che cosa succede, decide i tagli di scena.

\paragraph{Approfondimento: dal wargame al contratto sociale.} La genealogia
raccontata da Daniele si ritrova in \emph{21st Century Game Design} di Bateman e
Boon. Le radici del gioco di strategia, scrivono, precedono i videogiochi di
decenni: alla fine degli anni Cinquanta Charles Roberts fonda Avalon Hill e
pubblica i primi \textbf{wargame da tavolo} (1958); in un'epoca senza
videogiochi, giochi di ruolo e giochi di carte collezionabili, i wargame
``fondano l'industria dei \emph{hobby games}'', e \emph{Chainmail} (1971),
ispirato a quei giochi, è il ``trampolino'' verso \emph{Dungeons \& Dragons}. Gli
stessi autori descrivono la forma minima del gioco di ruolo da tavolo come un
\textbf{contratto sociale} tra l'arbitro (il \emph{games master}) e i
giocatori: i giocatori accettano che sia l'arbitro a decidere come le loro azioni
sviluppano la storia, e l'arbitro si impegna a garantire intrattenimento o
imparzialità, come mediatore e narratore principale. Citano anche un'osservazione
attribuita a Gygax: se i giocatori di ruolo si accorgessero di poter fare tutto
da soli, i game designer resterebbero senza lavoro. Il fatto, notano, è che la
maggior parte dei master ha bisogno di meccaniche per gestire il mondo, e i
giocatori si fidano di più dell'arbitro quando una struttura di meccaniche
casuali, i dadi, fa da cornice. Anche Schell, in \emph{The Art of Game Design},
cita il master dei giochi di ruolo da tavolo come esempio di stato del gioco
\emph{nascosto}: è lui a conoscere privatamente gran parte dello stato del gioco,
perché ne è ``il meccanismo operativo''.

\subsection{Gli anni Ottanta e Novanta: adattare un regolamento a tutto}

Qui Daniele fa un ``piccolo inciso editoriale''. Nella grande espansione degli
anni Ottanta e Novanta \textbf{non c'era un professionismo} del settore. Gli
editori compravano delle proprietà intellettuali, per esempio \textbf{007},
prendevano un master appassionato e competente e gli chiedevano di
\textbf{adattare} il regolamento di D\&D, o un altro regolamento
esistente\footnote{La trascrizione dice ``Call of Duty 2'', che non può essere
corretto in questo contesto; il senso è ``il regolamento di un altro gioco
già esistente''.}, a quella IP. A volte nascevano giochi innovativi, a volte
giochi ``non molto interessanti da giocare'': immaginate di giocare con le
regole di D\&D qualcosa che quelle regole non supportano. Verso la fine degli
anni Novanta il problema della qualità ``ha iniziato a farsi sentire''.

\paragraph{Approfondimento: James Bond 007 e l'esperienza essenziale.} Proprio
il gioco di ruolo di 007 è il caso con cui Schell illustra, in \emph{The Art of
Game Design}, la \textbf{lente dell'esperienza essenziale} (\emph{Lens of
Essential Experience}), e mostra l'altra faccia della storia di Daniele. Il
designer Chris Klug era frustrato dai precedenti giochi di ruolo di spie, come
\emph{Top Secret} della TSR, perché ``si giocavano troppo come dei wargame'':
l'essenza di ciò che rende eccitanti i film di spionaggio non c'era. Per il
gioco di James Bond, Klug progettò le meccaniche in modo che restituissero in
ogni modo possibile la sensazione dei film. L'esempio è l'invenzione degli
\textbf{Hero Points}: nei GDR tradizionali, davanti a un'azione rischiosa (saltare
da una finestra su un elicottero in volo), il master calcola una probabilità, il
giocatore tira i dadi ``e questo è quanto''; ma se la probabilità è bassa
nessuno rischia, e se è alta tutti si comportano da supereroi. I punti eroe
risolvono il dilemma in modo coerente con il genere. È esattamente quello che
la New Wave rivendicherà: \textbf{il sistema conta}, e un regolamento adattato ``a
nastro'' non restituisce l'esperienza di un genere.

% ---------------------------------------------------------------------
\section{La New Wave e The Forge}
\label{a2:sec:gdr-newwave}

\subsection{Il sistema conta}

La slide successiva è sulla \textbf{New Wave}: un movimento nato tra la fine
degli anni Novanta e l'inizio dei Duemila. Daniele avverte che non vuole fare
un'analisi geografica ma ``un riassunto della storiografia''. La New Wave cerca
un \textbf{metodo diverso di interpretare il gioco di ruolo}, e la prima cosa
che stabilisce, in un forum chiamato \textbf{The Forge} (c'è anche un libro in
inglese su questo movimento), è che \textbf{il sistema dei giochi di ruolo è
importante}: le regole, il gameplay, la parte meccanica \textbf{non sono un
orpello}. Non si può tenere lo stesso sistema, cambiare il setting e sperare che
funzioni da sé. Tanto è vero che perfino \emph{GURPS}, il sistema generalista,
aveva tantissimi manuali di espansione delle regole per ogni genere:
\emph{GURPS} base, \emph{GURPS Fantasy}, \emph{GURPS} fantascienza\ldots

\subsection{I giochi della New Wave}

\slidefig{img/AGD/agd2-newwave}{Alcuni giochi della New Wave.}{slide \emph{Analog -- Case Studies}}{fig:agd2-newwave}


Dentro la New Wave nascono esperienze diversissime, che spostano l'attenzione
su due cose: la \textbf{scrittura del manuale}, diversa da quella tradizionale,
e la \textbf{rottura degli schemi}.
\begin{itemize}
  \item \textbf{\emph{Apocalypse World}}\footnote{La trascrizione dice ``A
        Polly's Ward'' e ``Minster Baker''.}, di \textbf{Vincent Baker}, cerca di
        scardinare il modo classico di strutturare un manuale di gioco di ruolo
        (Daniele accenna a una nuova edizione recente).
  \item \textbf{\emph{Nobilis}}, uno dei primi giochi abbastanza famosi
        \textbf{diceless}, senza dadi: azioni e conflitti si risolvono in un
        altro modo.
  \item \textbf{\emph{Don't Rest Your Head}} (in italiano \emph{Non cedere al
        sonno}) e \textbf{\emph{Polaris}}, che lavorano sul concetto di chi è e
        chi non è il personaggio. \emph{Non cedere al sonno} è un gioco
        ``surrealista'' in cui ci si sveglia a \textbf{Mad City}; la parte
        importante sono i \textbf{poteri}. All'epoca \emph{Vampire} aveva
        stabilito che i poteri fossero una lista di possibilità con statistiche
        a cui applicare le regole; in \emph{Don't Rest Your Head} ci sono poteri
        ``totalmente assurdi, fenomenali'', che non richiedono meccaniche in più
        ma usano le meccaniche base ``in maniera molto innovativa''.
  \item \textbf{\emph{Fate}}, un sistema anch'esso generalista, ``più o meno
        come \emph{GURPS}'', ma con delle sue specifiche: ci sono
        \textbf{procedure per creare i setting} e per stabilire quali regole un
        setting richiede.
  \item \textbf{\emph{Dread}}, uno dei primi giochi a usare per l'incertezza
        \textbf{un elemento che non è un dado}: una \textbf{torre Jenga}. È un
        gioco dell'orrore in cui ogni volta che si fa un'azione si toglie un
        blocchetto, e se la torre cade il personaggio muore (``non è proprio
        così, è più complicato, però per semplificare'').
\end{itemize}

\paragraph{Le guerre di religione.} La New Wave ha avuto ``una serie di
problematiche'' in giro per il mondo, e soprattutto in Italia ha creato una
sorta di \textbf{guerra teologico-filosofica} tra i ``tradizionalisti'' di D\&D e
i ``forgiti'' della New Wave. ``Per vostra sanità mentale non ci entriamo''.

\paragraph{Una reazione a due cose.} La New Wave è stata una reazione:
\begin{enumerate}
  \item ai manuali in cui ``c'è un sistema e quello è il sistema'', con il
        setting applicato sopra ``a nastro'', senza intreccio tra la narrazione
        e le regole che la raccontano: si voleva \textbf{coerenza};
  \item all'assenza di studio: è stato il \textbf{primo movimento a voler
        studiare i giochi di ruolo}, a capire come sono fatti. Da qui viene il
        \textbf{Big Model}, un modello di analisi dei giochi di ruolo di cui si
        vedranno alcune cose nella lezione sulla fiction.
\end{enumerate}
Ma il Big Model ``non è fondamentale'', perché per Daniele uno dei problemi della
New Wave è aver creato \textbf{un dizionario collaterale} a quello del game design
del gioco da tavolo e del videogioco: termini diversi che dicono la stessa cosa.

\subsection{The Forge e la scena nordica}

Ruggero fa da ponte: \textbf{The Forge ha fatto per il gioco di ruolo da tavolo
quello che la scena scandinava ha fatto per il LARP}. Inventavano e scrivevano
moltissimo, hanno lasciato framework e strumenti. La scena scandinava è nata un
po' prima, The Forge un po' dopo, ma hanno agito in parallelo. La differenza è
che The Forge era \textbf{meno strutturato}: essendo un forum, era più caotico,
con discussioni anche accese, soprattutto sulle interpretazioni. Anche la
\textbf{preservazione storica} è diversa: della scena nordica ci sono i libri
fisici, mentre per il forum bisogna passare dalla Wayback Machine. Ma in termini
di normalizzazione di termini, meccaniche e framework, entrambi hanno fatto
moltissimo.

\subsection{Fiasco}

\slidefig{img/AGD/agd2-fiasco}{Il flusso di Fiasco: chi stabilisce la scena e chi la risolve.}{slide \emph{Analog -- Case Studies}}{fig:agd2-fiasco}


Daniele riprende: a un certo punto The Forge chiude, e nascono altri giochi
dagli stimoli e dalle discussioni che quell'ambiente aveva generato, non
necessariamente da persone che ne facevano parte. Per esempio quelli di
\textbf{Jason Morningstar}, noto anche per dei chamber LARP (il migliore amico
di Daniele, Edoardo, ``più esperto di me'', ne è talmente fan che il suo gatto si
chiama Giasone).

\textbf{\emph{Fiasco}} è per Daniele uno dei primi giochi che \textbf{rompe il
canone}. Racconta, appunto, un fiasco, come nei film dei \textbf{fratelli Coen}
o di Tarantino: c'è uno scenario, ci sono dei personaggi, e ogni volta che si fa
una scena \textbf{il giocatore di turno sceglie se aprirla} (e gli altri la
risolvono) \textbf{o se lasciare che gli altri la impostino} (e risolverla lui).
Già solo questo ``micro elemento'' (non l'unica meccanica: c'è anche quella,
interessantissima, dei dadi) rompe l'idea che in un gioco di ruolo si debba
descrivere l'azione, tirare i dadi e vedere che cosa succede, oppure reagire a
ciò che descrive il master. Per chi veniva da \emph{AD\&D} seconda edizione o da
\emph{Vampire} era ``assolutamente interessante'': più leggero e più legato alla
fiction. Anche i \textbf{dadi} sono usati in modo nuovo: vengono tirati
\textbf{prima}, messi sul tavolo, selezionati, usati quasi come
\textbf{token}: un'innovazione anche nel modo di usare i materiali.

% ---------------------------------------------------------------------
\section{Il gioco di ruolo contemporaneo}
\label{a2:sec:gdr-oggi}

\subsection{Prima del COVID: i giochi focalizzati}

\slidefig{img/AGD/agd2-contemporary}{Giochi di ruolo contemporanei.}{slide \emph{Analog -- Case Studies}}{fig:agd2-contemporary}


Daniele salta alla contemporaneità. Prima del COVID c'è un grande aumento di
giochi che, mentre le persone si appassionano ai giochi di ruolo, vanno o nella
direzione di D\&D o in una \textbf{direzione più ibrida}:
\begin{itemize}
  \item \textbf{\emph{Blades in the Dark}}\footnote{La trascrizione dice
        ``Blaze in the Dark''. Daniele lo data ``2016--18 più o meno''.};
  \item più recenti, il gioco di ruolo di \textbf{\emph{Cowboy
        Bebop}}\footnote{La trascrizione dice ``Cabobby Bop'' e ``Cobo
        Bobby''.} e \textbf{\emph{The Clay That Woke}} di \textbf{Paul
        Czege}\footnote{La trascrizione dice ``Paul Ksege, The Clayed Walk''.},
        autore di ``uno dei principi di design più interessanti'';
  \item \textbf{\emph{Fabula Ultima}}, di \textbf{Emanuele Galletto}, italiano
        (``che salutiamo'').
\end{itemize}
Sono prodotti che \textbf{non cercano di fare la stessa cosa di D\&D}. Hanno un
impianto simile (scheda del personaggio, dadi, azioni descritte, tiri), ma si
pongono \textbf{obiettivi di design più coerenti}: ``si chiedono che cosa vogliamo
che il giocatore senta''. Il sistema di \emph{Cowboy Bebop}, per esempio, usa
gli \textbf{orologi} (\emph{clock}) che si trovano anche in \emph{Blades in the
Dark} e nei giochi \emph{Powered by the Apocalypse}, cioè la stessa meccanica,
ma è un sistema completamente diverso, perché deve raccontare una fiction
diversa:
\begin{itemize}
  \item in \emph{Blades in the Dark} si giocano dei \textbf{criminali} a
        \textbf{Doskvol}\footnote{La trascrizione dice ``Oswald''.}, città di un
        mondo alternativo dove è sempre notte, dove i fantasmi non passano
        nell'aldilà e vivono con i vivi, e dove si è costretti a restare perché
        fuori dalle mura i fantasmi uccidono. Si fa parte di una banda che cerca
        di scalare la gerarchia della città, ``tipo \emph{Peaky Blinders}'';
  \item in \emph{Cowboy Bebop} si è ``letteralmente una ciurma di persone nello
        spazio, nel futuro, che fanno cose'', perché ``l'anime si racconta con i
        personaggi che fanno cose'' (Daniele, ``molto fan dell'anime'', lo dice
        come un complimento).
\end{itemize}

\subsection{Dopo il COVID: il ritorno all'origine}

Soprattutto dopo il COVID si è notato un ``ritorno fortissimo'' al
\textbf{giocare come una volta}. Per Daniele entrambe le tendenze rispondono a
una necessità: \textbf{semplificare}. \emph{D\&D} 3 e 3.5, le nuove edizioni di
\emph{Vampire} e altri sono diventati ``pieni di regole, pieni di contenuti'', e
chiedono a master e giocatori di imparare moltissimo, senza contare il tempo
necessario per una campagna. Allora, da un lato, si cercano sistemi più semplici
``da intavolare'' che facciano bene una cosa; dall'altro si cerca la stessa cosa
per altri motivi, ricreando le atmosfere del D\&D \emph{old school}: \textbf{alta
mortalità}, parte narrativa e interpretativa meno importante, immediatezza.
\textbf{\emph{Mörk Borg}}\footnote{La trascrizione dice ``Morg Borg''.}, per
esempio, sono circa \textbf{cento pagine}, con poche regole. C'è più attenzione
al ``realismo'' (termine un po' sbagliato, si vedrà nella lezione sulla
fiction), e il vantaggio è che questi giochi possono usare sia le proprie regole
sia quelle di \textbf{D\&D quinta edizione}, che è stata importante soprattutto
``per un motivo editoriale'' (anche questo si vedrà nella lezione sulla
fiction).

\paragraph{La domanda di Gianmarco.} Gianmarco chiede se questa tendenza a
manuali e partite ``asciutti'' segua un pubblico più adulto, con meno tempo, o
sia una tendenza generale del mercato. È soprattutto la prima. Il gioco di ruolo
ha avuto parecchie difficoltà ad attirare nuove leve, e i giocatori sono
cresciuti. Daniele gioca da molti anni, ma il suo ``vero inizio'' è stato con un
gioco della New Wave\footnote{La trascrizione dice ``Paura dell'Epoclips'':
probabilmente \emph{Il mondo dell'apocalisse}, l'edizione italiana di
\emph{Apocalypse World}, ma non è certo.}: per lui, all'università, avere in un
solo manuale da 200 pagine tutte le regole e tutti i personaggi è stato ``un
toccasana''. Chi inizia con D\&D, magari alle superiori o all'università, man
mano che cresce tende a preferire:
\begin{itemize}
  \item o giochi più focalizzati sull'aspetto \textbf{narrativo}, che hanno
        comunque scheda, tiro di dado e struttura tradizionale, ma una scheda e
        un sistema più semplici;
  \item o un'esperienza simile al D\&D canonico, ma con meno ``regoline e
        regolette'' e meno \textbf{simulazionismo}.
\end{itemize}
Due differenze pratiche. I manuali \textbf{OSR}, che un po' di simulazionismo ce
l'hanno, sono pieni di \textbf{soluzioni per l'arbitro} (``che non è il master,
è l'arbitro'') per risolvere i problemi sul momento; e l'OSR è \textbf{sandbox}:
c'è una mappa, una zona in cui si va dove si vuole, e il master non ha scritto
un'avventura. I \textbf{giochi narrativi}, dal canto loro, non hanno una campagna
scritta ma uno \textbf{scenario}. In entrambi i casi il tempo di preparazione del
master si riduce ``di gran lunga''.

\subsection{Focalizzati, non semplificati}

Ruggero interviene sulle parole: non vorrebbe che passasse l'idea che questi
giochi siano ``più semplici'' o \emph{dumbed down} rispetto a D\&D. Si parla di
\textbf{giochi focalizzati}: invece di semplificare, \textbf{focalizzano}, sono
molto più \textbf{coesi tematicamente}, hanno ``un cuore tematico molto forte'' e
si giocano per giocare quello. D\&D invece, a livello di cuore meccanico, è
giocato in modi diversissimi da dieci gruppi diversi.

Daniele fa un esempio: nei manuali di D\&D (dalla terza edizione e 3.5 alla
quinta) ci sono le regole per il \textbf{danno da caduta}. In \emph{Blades in
the Dark} non servono: c'è un altro metodo, ``non è così generalista, è
focalizzato'', e la narrazione si risolve in altro modo. Anche nell'OSR:
\emph{Shadowdark} ha il danno da caduta, perché magari è importante;
\emph{Neverland}, una reinterpretazione di \textbf{Peter Pan} in chiave OSR, no,
perché lì conta di più che il manuale offra tutti i personaggi del Peter Pan
tradizionale in una versione \emph{grimdark}. ``Anche l'OSR diventa un po'
focalizzato quando ci prova''.

\subsection{Not the End}

\slidefig{img/AGD/agd2-nottheend}{Il flusso di Not the End: il sacchetto di token come motore della risoluzione.}{slide \emph{Analog -- Case Studies}}{fig:agd2-nottheend}


Nicola, per cui questo è ``il settore che più adoro del gioco analogico'',
chiede dove si collochi \textbf{\emph{Not the End}}, un gioco italiano. La
risposta è nella slide successiva. \emph{Not the End} continua nella direzione
più narrativa e focalizzata (``semplicistica è un termine non corretto''): ha
procedure, meccaniche, un gameplay focalizzato a fare quello che il gioco
chiede. Si possono creare i propri setting, ma il gioco resta quello: il
personaggio principale è un \textbf{eroe}, ed è il giocatore a decidere quando
sarà la sua fine, quando giocherà il suo momento di eroismo stoico (da cui il
titolo). È un gioco \textbf{senza dadi}: c'è un conflitto, si fa e si descrive
un'azione, e la \textbf{risoluzione} consiste nell'\textbf{estrarre dei token}
bianchi e neri (o arancioni e viola, per le persone con difficoltà nella visione
dei colori); l'effetto viene descritto secondo le procedure del gioco. Il
panorama, insomma, si è diviso: c'è chi cerca un'esperienza più simile
all'originale, e chi un'esperienza più focalizzata sulla narrazione, magari
senza un \emph{world building} preciso.

\subsection{Che cos'è l'OSR, e com'è nato davvero il gioco di ruolo}

Un'altra studentessa chiede che cosa significhi OSR: \textbf{Old School
Renaissance}, ``rinascimento della vecchia scuola''\footnote{Poco prima Daniele
aveva sciolto la sigla come \emph{Old School Revival}; entrambe le forme sono in
uso.}. È un movimento ``molto, molto particolare'': ogni persona che ne fa parte
ha il proprio principio di OSR, gioca un OSR diverso dagli altri, e gli
appassionati ``adorano scontrarsi'' su che cosa intendono per OSR. Per questo
nella lezione sulla fiction Daniele parlerà di \textbf{culture} del GDR, più che
di community.

La stessa studentessa (iscritta a un altro corso dell'accademia, di cui sta
recuperando le lezioni del primo anno) riporta quello che ha capito dell'origine
di D\&D: i due autori avrebbero avuto visioni opposte, uno più libera e meno
regolamentata, partendo dai giochi di guerra, l'altro, che ha portato il gioco a
livello commerciale, più strutturata, dettagliata e pesante. ``È giusto, è un
riassunto corretto, ma è un po' più complicato di così'', e lo si vedrà nella
lezione sulla fiction. Ruggero aggiunge che D\&D ``non è l'unica origine del
gioco di ruolo'': il gioco di ruolo diventa \emph{famoso} con D\&D, ma gli stessi
autori si ispirarono alle \textbf{simulazioni militari}, molto diffuse
nell'America di quegli anni, e al \textbf{reenactment storico} (``ci sono due
episodi della \emph{Signora in giallo} sul reenactment della guerra civile, un
argomento allucinante''). Esistono esempi storici, dell'Ottocento e perfino del
Seicento, di \textbf{giochi di ruolo} veri e propri, con quest, avventure e
personaggi. Il \emph{mainstream} ha normalizzato l'idea che D\&D sia il primo
gioco di ruolo, ``ma non è così'': nel manuale originale l'espressione
\emph{role-playing game} non compare da nessuna parte, anche perché Gygax
odiava il termine. ``Un vaso di Pandora che non apriamo''.

Daniele però raccoglie l'osservazione. Nei primi anni editoriali di D\&D il
problema maggiore è stato il \textbf{simulazionismo}: interpretare
``simulazione'' come la necessità di avere una regola per ogni cosa che esiste nel
realismo del gioco. Ne è venuto un gioco molto strutturato, in cui ``ogni cosa
deve avere una regola'', perché ciò che non è coperto dalle regole il giocatore
può usarlo come \emph{loophole}. E ne è venuta anche l'idea che il master sia il
\textbf{padrone del tavolo} e che il giocatore debba in qualche modo
\textbf{vincere contro il master}: un altro ``vaso di Pandora enorme''.

\paragraph{Non esiste solo D\&D.} È per questo che Daniele non si concentra solo
su D\&D: ``l'idea con cui dovete uscire da qui è che non esiste solo D\&D''.
Anzi, oggi D\&D è \textbf{in controtendenza}, non per motivi di design ma per
\textbf{motivi editoriali}: il problema principale di D\&D come proprietà
intellettuale ``è sempre stato editoriale'', e di recente la Wizards of the Coast
ha fatto un po' di confusione con il \emph{branding}.

% ---------------------------------------------------------------------
\section{Story game e giochi di ruolo sperimentali}
\label{a2:sec:gdr-story}

\subsection{Gli story game}

\slidefig{img/AGD/agd2-storygames}{Alcuni story game.}{slide \emph{Analog -- Case Studies}}{fig:agd2-storygames}


La slide successiva è sugli \textbf{story game}, che in realtà nascono anche
prima di The Forge. ``Story game è un termine che io odio, perché vuol dire tutto
e non vuol dire niente'': sono giochi di ruolo da tavolo che si focalizzano
\textbf{ancora di più su una storia precisa}.
\begin{itemize}
  \item \textbf{\emph{Alice is Missing}}: la storia di Alice, scomparsa, e dei
        suoi amici che cercano di capire che cosa le è successo. Si gioca con
        una \textbf{chat} sul telefono, con un video a schermo, e con
        \textbf{due ore di tempo}.
  \item \textbf{\emph{Microscope}}: ``letteralmente \emph{world building}
        narrativo''.
  \item \textbf{\emph{Bleak Spirit}}: giocare un \emph{Dark Souls} in cui tutti
        controllano a turno il viandante protagonista.
  \item \textbf{\emph{The Zone}}\footnote{La trascrizione attribuisce il gioco a
        ``Ruff D'Amico''; il nome non è verificabile.}, giocabile anche con un
        supporto online: si narra, di base, \emph{Annientamento} (``se avete
        visto il film''). Si entra in una zona in cui una presenza aliena sta
        cambiando il mondo, e \textbf{moriranno tutti tranne una persona}, che
        non si sa in anticipo chi sarà. È un'influenza del LARP: come in
        \emph{Demeter}, si sa che il proprio personaggio morirà.
  \item \textbf{\emph{Archipelago}} e \textbf{\emph{For the Queen}}, e altri
        giochi che \textbf{tolgono i dadi} e mettono le carte, o tolgono
        perfino la \textbf{scheda del personaggio}: in \emph{Microscope} ci sono
        solo dei post-it.
\end{itemize}

\subsection{The Quiet Year e i map drawing game}

\slidefig{img/AGD/agd2-quietyear}{Il flusso di The Quiet Year.}{slide \emph{Analog -- Case Studies}}{fig:agd2-quietyear}


``C'è un minimo di suspense'': si arriva a esperienze ancora più lontane.
In \textbf{\emph{The Quiet Year}} non si interpreta nemmeno un personaggio: si è
parte di una \textbf{comunità post-apocalittica} e si gioca \textbf{un anno} della
sua vita. Non si può discutere di ciò che si gioca. A ogni turno si gira una
\textbf{carta} che dice che cosa succede, e durante il proprio turno si può fare
una tra due o tre azioni (disegnare, iniziare un progetto\ldots). Tutto si gioca
su una \textbf{mappa} su cui si disegna ciò che si vuole intendere: c'è chi gioca
in modo più rigoroso, senza poter nemmeno scrivere parole sulla mappa, e chi lo
permette. È ``un gioco fenomenale'', perché fa capire che il gioco di ruolo può
essere anche altro; e ha normalizzato un genere, quello dei \textbf{map drawing
role-playing game}, i giochi di ruolo in cui si disegnano mappe.

\subsection{Dialect}

\slidefig{img/AGD/agd2-dialect}{Il flusso di Dialect.}{slide \emph{Analog -- Case Studies}}{fig:agd2-dialect}


E poi \textbf{\emph{Dialect}}. Qui si ha un personaggio, ma è solo un nome e un
ruolo, un archetipo. Si gioca \textbf{la morte di una comunità attraverso la
morte di una lingua}: una comunità umana isolata dalle altre, che ha creato il
proprio linguaggio, e che giocherà la fine di quell'isolamento attraverso la
morte del linguaggio. A ogni turno si gioca una carta che è un \textbf{concetto
della lingua} (per esempio: come si dice ``amico'' nella nostra lingua?); si
stabilisce la parola e si gioca una breve scena in personaggio, senza tiri di
dado né altro; quando la parola viene usata in personaggio, chiunque al tavolo
può dire stop: ``abbiamo usato la parola, va bene così'', e si va avanti.

Non è un gioco ``non complesso'', come diceva Ruggero: \emph{Dialect} è
``chiarissimo'' da leggere, ma la complessità sta nel \textbf{design}, un design
elegante che crea una \emph{one shot} con ``una profondità'' enorme. Daniele con
\emph{Dialect} ha pianto; e quando incontra gli amici con cui lo ha giocato anni
fa, si parlano ancora nella lingua inventata allora: ``ehi tu, oceano'', e chi
c'era sa che cosa vuol dire. Anche qui si vede come le esperienze nate dal Nordic
LARP siano \textbf{``pervase''} (``rubo una parola a Beatrice'') in altri ambiti.

\subsection{Servono davvero le categorie? Il problema editoriale}

Ci sono ancora esperienze più tradizionali; ``il mondo è molto variegato'', e si
parlerà di \textbf{culture}: spesso chi si avvicina al gioco di ruolo vuole
giocare una certa cosa, una cultura di gioco con titoli che ``vanno bene'' e
titoli che no. L'ultima slide è ``una piccola critica'': è davvero necessario
mettere le categorie ai giochi di ruolo? Da un lato sì, dall'altro no, perché i
problemi principali sono due, e il primo è \textbf{il mondo editoriale}: i giochi
di ruolo sono stampati come \textbf{libri}, e l'editore ha ``una vita travagliata''
con il gioco. Esempi:
\begin{itemize}
  \item \textbf{\emph{In the Name of God}}\footnote{Titolo e autore (``Alessandro
        Pirodi'') come nella trascrizione; non verificabili.}, stampato solo in
        \textbf{self-publishing} e assente dal mercato, ``ma bellissimo''; la sua
        meccanica di risoluzione chiede di fare qualcosa fisicamente al tavolo
        (la trascrizione dice ``fare home'');
  \item un'\textbf{espansione per qualsiasi gioco di ruolo}\footnote{La
        trascrizione dice ``Mew Wild in the Subway''.};
  \item un \textbf{GDR stampato su una maglietta}, di due autori che Daniele
        saluta;
  \item un gioco di \textbf{cultisti contro investigatori} che prende da
        esperienze di chamber LARP\footnote{La trascrizione lo chiama ``SOT''.}.
\end{itemize}

\begin{figure}[htbp]
\centering
\begin{tikzpicture}[
    font=\footnotesize,
    epoca/.style={draw, thick, rounded corners=3pt, fill=gray!12,
                  text width=2.55cm, align=center, minimum height=1.0cm},
    nota/.style={align=left, text width=2.7cm, font=\scriptsize},
    freccia/.style={-{Stealth[length=2.5mm]}, thick}
  ]
  \draw[freccia] (-0.2,0) -- (15.4,0);
  \foreach \x/\a in {0.9/{anni '50--'70}, 4.0/{'74--anni '90},
                     7.1/{fine '90--2000}, 10.2/{anni 2010}, 13.3/{post-COVID}}{
    \draw[thick] (\x,0.12) -- (\x,-0.12);
    \node[below=0.15cm, font=\scriptsize\itshape] at (\x,0) {\a};
  }
  \node[epoca] (e1) at (0.9,1.2)  {\textbf{Wargame e miniature}};
  \node[epoca] (e2) at (4.0,1.2)  {\textbf{D\&D e GDR ``early''}};
  \node[epoca] (e3) at (7.1,1.2)  {\textbf{New Wave / The Forge}};
  \node[epoca] (e4) at (10.2,1.2) {\textbf{Giochi focalizzati, story game}};
  \node[epoca] (e5) at (13.3,1.2) {\textbf{OSR e ritorno all'origine}};
  \node[nota, anchor=north] at (0.9,-0.65) {simulazioni militari,
    reenactment; \emph{Chainmail} (1971)};
  \node[nota, anchor=north] at (4.0,-0.65) {Gygax e Arneson;
    master arbitro; tutti i generi (\emph{Cthulhu}, \emph{Vampire},
    \emph{Traveller}, \emph{GURPS}); IP adattate};
  \node[nota, anchor=north] at (7.1,-0.65) {il sistema conta;
    \emph{Apocalypse World}, \emph{Nobilis}, \emph{Dread}, \emph{Fate}; Big
    Model; \emph{Fiasco}};
  \node[nota, anchor=north] at (10.2,-0.65) {\emph{Blades in the
    Dark}, \emph{Fabula Ultima}, \emph{Not the End}; \emph{The Quiet Year},
    \emph{Dialect}; due giocatori, solo, micro};
  \node[nota, anchor=north] at (13.3,-0.65) {meno regole, alta
    mortalità, sandbox; \emph{Mörk Borg}, \emph{Shadowdark}};
\end{tikzpicture}
\caption{Una ricostruzione schematica della ``piccola storiografia'' dei giochi
di ruolo da tavolo raccontata da Daniele. Le epoche sono indicative e si
sovrappongono (gli story game, per esempio, nascono anche prima di The Forge);
la storia completa è rimandata alla lezione sulla fiction.}
\label{a2:fig:storia-gdr}
\end{figure}

% ---------------------------------------------------------------------
\section{Giochi di ruolo per due, in solitario e micro}
\label{a2:sec:gdr-formati}

\slidefig{img/AGD/agd2-micro}{Giochi di ruolo per due, in solitario e nano-RPG.}{slide \emph{Analog -- Case Studies}}{fig:agd2-micro}


L'ultima parte sui giochi di ruolo torna a Ruggero, che presenta tre formati.

\begin{description}
  \item[Two player only.] Storicamente i giochi di ruolo nascono in gruppi di
        gioco ampi, ma c'è stata ``una grossissima wave'' di giochi di ruolo
        \textbf{per due persone}, e non nel senso che una fa il master e l'altra
        un personaggio: spesso l'\textbf{agency narrativa}, il potere di decidere
        che cosa succede, viene \textbf{passata meccanicamente} al tavolo o
        divisa in responsabilità diverse. Ce ne sono tantissimi, e per Ruggero
        sono molto interessanti anche perché, come si vedrà nella lezione
        sull'emotional design, permettono di fare cose molto \textbf{intime}: non
        serve un gruppo di gioco, si possono giocare con il partner o con
        un'altra persona.
  \item[Solo RPG.] Dal 2013 circa si sono normalizzati i \textbf{giochi di
        ruolo in solitario}: si gioca da soli, senza un'altra persona e senza
        fare da master a sé stessi. Spesso non c'è un master: ci sono
        \textbf{procedure} che guidano il gioco, e il manuale è il sistema di
        gioco, strutturato ``quasi come un board game'' in termini di regole. Ma
        resta tutto ciò che è tipico del gioco di ruolo: l'agency emergente,
        l'interpretazione di un personaggio, la creazione di una narrazione. Di
        solito si giocano in formato \textbf{diario} (\emph{journaling}) o in
        formato \textbf{audio}, e sono in grande tendenza.
  \item[Micro RPG.] Infine, i \textbf{micro giochi di ruolo}\footnote{La
        trascrizione dice ``handheld RPG''; dal contesto si tratta dei
        \emph{micro RPG}.}: giochi di ruolo su \textbf{biglietti da visita}, o
        raccolte in cui un libro ne contiene cento. Il \textbf{200 Word RPG}
        (\emph{200 Word RPG Challenge}) era un contest che chiedeva di scrivere un
        gioco di ruolo in \textbf{massimo 200 parole}, e ogni edizione ne
        raccoglieva migliaia: nel 2017, ricorda Ruggero, circa 900 giochi, su
        circa mille proposte.
\end{description}

% ---------------------------------------------------------------------
\section{I giochi da tavolo: i generi storici}
\label{a2:sec:tavolo}

\slidefig{img/AGD/agd2-rollmove}{Roll and move, astratti e party game.}{slide \emph{Analog -- Case Studies}}{fig:agd2-rollmove}


``Passerei ai giochi da tavolo, correndo velocemente e senza fare pause, se ce la
fate''. Come per i LARP e i giochi di ruolo, Ruggero parla di
\textbf{categorie dall'alto}, e per ognuna mostra come il genere si sia
\textbf{evoluto} nel tempo.

\subsection{Roll and move}

I \textbf{roll and move} si basano, come dice il nome, sul \textbf{tirare un
dado e muoversi}. Sono forse i più comuni nella storia dell'umanità, con quello
che da noi è il \textbf{gioco dell'oca}: un tracciato su cui si tirano i dadi e
si avanza. L'\textbf{agency} in questi giochi è ``praticamente inesistente'', e
il genere si è presto evoluto integrando altre meccaniche: con
\textbf{elementi economici} (come \emph{Monopoly}) o con \textbf{elementi
deduttivi}, come \emph{Cluedo}, che è un roll and move.

\subsection{Astratti}

Un altro genere storicamente gigantesco è quello degli \textbf{astratti}:
scacchi, dama, go, tutti quei giochi la cui componente tematica è ridotta
``quasi all'osso''. Gli scacchi mantengono ancora qualche forma di tema; il go, di
fatto, è un rettangolo di legno con delle caselle e delle pietre bianche e nere.
Anche qui il genere non si è fermato: ci sono iterazioni moderne come
\emph{ZÈRTZ} e i giochi del \textbf{progetto GIPF}\footnote{La trascrizione dice
``Zerz o il Gift Project''.}, e altri astratti contemporanei (uno è mostrato a
destra nella slide\footnote{La trascrizione non permette di ricostruirne il
titolo (``i borsa astratti, per esempio, di Supernova'').}), che mantengono
dinamiche ``molto asciutte, molto distillate'': \textbf{informazione aperta},
trasparente, niente carte imprevisti, niente mosse segrete. Diventano quasi una
\textbf{sfida d'intelletto} con l'avversario.

\subsection{Party game}

I \textbf{party game}: \emph{Taboo}, \emph{Lupus in Tabula} o \emph{Werewolf} a
seconda dell'edizione. Sono pensati per \textbf{molte persone} (di solito dalle
quattro in su), con regole ``molto asciutte, molto leggere'', perché si giocano
in contesti di festa, in cui non è detto che chi è nella stanza ami giocare o
abbia mai giocato. Poche regole, leggerezza, durata breve (spesso non più di
mezz'ora): sono \textbf{un possibile punto di accesso}. Anche questi si sono
evoluti, con giochi come \emph{Hitster}\footnote{La trascrizione dice ``Hit'';
la ricostruzione è probabile ma non certa.} e \emph{Wavelength}, che portano
sensibilità più moderne e meccaniche più interessanti del ``tirare alla cieca
per capire chi tra di noi è un lupo mannaro''.

\subsection{Giochi di parole}

I \textbf{word game}\footnote{La trascrizione dice ``war games'', ma gli esempi
(\emph{Boggle}, \emph{Scarabeo}) chiariscono che si tratta dei giochi di
parole.}: \emph{Boggle}, \emph{Scarabeo}, giochi in cui la meccanica centrale è
\textbf{creare parole}, e più la parola è lunga, complessa o usa lettere rare,
più punti si fanno. Anche qui le iterazioni sono interessanti: \emph{Cross
Clues}\footnote{La trascrizione dice ``Cross Cruise''.} trasforma il genere in
un gioco \textbf{cooperativo}, in cui non si compete per la parola più lunga ma si
gioca con le parole per capire ``quanta intesa c'è tra di noi''; e
\emph{Illiterati} è praticamente un gioco con le \emph{boss fight}, in cui si
sconfiggono i boss formando parole. Ibridazioni anche qui, perché chi crea giochi
li crea in base a ciò che ha giocato prima.

\subsection{Giochi di carte}

\slidefig{img/AGD/agd2-carte}{Giochi di parole, giochi di carte, TCG e deckbuilding.}{slide \emph{Analog -- Case Studies}}{fig:agd2-carte}


I \textbf{card game}, che in Occidente si diffondono storicamente con i mazzi
\textbf{francesi} e, in Italia, con quelli ``da briscola'' (ma ci sono anche altri
casi storici\footnote{La trascrizione cita i ``Mazda 1''; il riferimento non è
ricostruibile.}). Le evoluzioni:
\begin{itemize}
  \item i \textbf{micro card game} come \emph{Love Letter}, che in una decina o
        dozzina di carte racchiude un gioco complesso;
  \item giochi come \emph{Munchkin}, in una direzione più ampia, con il
        \textbf{take that} (``voglio mettere i bastoni tra le ruote'' agli altri);
  \item la ``gigantesca categoria'' dei \textbf{trading card game}, già
        menzionata nella lezione precedente: \emph{Magic: The Gathering},
        \emph{Yu-Gi-Oh!}, \emph{Pokémon}. Il \emph{trading} non è solo una
        dinamica tra persone: è nel nome della categoria. Scambiare carte
        \textbf{fuori dalla partita} è parte della dinamica comunitaria;
  \item i \textbf{Living Card Game} (LCG), che provano a risolvere uno dei
        principali problemi editoriali dei TCG: ``se io ho più soldi di te, tu
        contro di me a Magic non vinci'', perché comprerò carte da 800 o 900
        euro l'una. Gli LCG \textbf{non hanno booster casuali}: ogni espansione ha
        l'elenco preciso del contenuto. Restano giochi di carte con il
        \textbf{deck building pre-partita}, con mazzi ``fortemente identitari''
        (si possono scegliere le carte anche solo perché sono più simpatiche),
        ma senza l'aspetto economico della \textbf{rarità indotta};
  \item i \textbf{deck-building game}, in cui la costruzione del mazzo
        \textbf{non è più pre-partita ma avviene durante la partita}: spesso si
        parte con mazzi identici, e durante il gioco si acquisiscono carte che
        cambiano \textbf{permanentemente} il proprio mazzo, con lo scopo di
        costruire il mazzo migliore possibile.
\end{itemize}

\subsection{German game ed Eurogame}

\slidefig{img/AGD/agd2-euro}{German game, Eurogame e American.}{slide \emph{Analog -- Case Studies}}{fig:agd2-euro}


Altre categorie, che Ruggero mette tra virgolette, sono i \textbf{German game}:
un genere sviluppato in \textbf{Germania} dopo la metà degli anni '50--'60, con
giochi a \textbf{interazione indiretta}, competitivi ma ``senza troppo dar
fastidio direttamente all'altra persona'': non ci si può attaccare
direttamente. Sono giochi poco bellici, anche per \textbf{rigetto della Seconda
guerra mondiale}, e si è sviluppato moltissimo il filone \textbf{comunitario}, di
\emph{community building} e di scambio: costruire la propria fattoria, costruire
un villaggio, scambiarsi le risorse (``io ho delle risorse, tu altre,
scambiamocele''), senza distruggere ciò che l'altro ha costruito.

Il genere si è poi reso più complesso e si è espanso oltre la Germania con gli
\textbf{Eurogame}, che includono Spagna, Portogallo, Francia e anche l'Italia
(Ruggero ne ha fatti parecchi). L'interazione resta indiretta, ma i giochi
diventano molto più complessi: spesso durano \textbf{due o tre ore} e provano a
rendere \textbf{una simulazione focalizzata di qualcosa del mondo reale}. Nella
slide, a sinistra, un gioco in cui si gestisce una \textbf{galleria
d'arte}\footnote{La trascrizione dice ``The Gallery''; probabilmente \emph{The
Gallerist}.}, con tutte le dinamiche di compravendita di quadri; e
\emph{Endeavor}, che Ruggero ha alle spalle, sul commercio mondiale in epoca
moderna. Dinamiche che simulano un tipo di mercanteggiare: ``a volte compro e
vendo quadri, a volte compro e vendo vasi''.

\paragraph{Approfondimento: l'Eurogame secondo Brathwaite e Schreiber.} In
\emph{Challenges for Game Designers}, nel capitolo sul \emph{targeting} di un
mercato, gli Eurogame (``noti anche come \emph{German games} o \emph{designer
games}'') sono un caso studio di gioco per il mercato di massa: giochi per la
famiglia, godibili a tutte le età; \emph{I coloni di Catan} è il primo ad aver
reso popolare il genere in Nord America, seguito da \emph{Carcassonne} e
\emph{Puerto Rico}. Il libro ne dà una definizione operativa, utile per
confrontare le osservazioni di Ruggero:
\begin{itemize}
  \item \textbf{durata tra 20 e 90 minuti} (troppo brevi annoiano gli adulti,
        troppo lunghi i bambini);
  \item \textbf{setup breve}, perché il setup non contiene decisioni;
  \item \textbf{regole semplici}, curva di apprendimento breve;
  \item \textbf{cooperazione più che scontro}: scambi e negoziazioni invece di
        meccaniche distruttive in cui un giocatore danneggia direttamente un
        altro (il che non impedisce una competizione intensa);
  \item \textbf{decisioni strategiche con casualità misurata};
  \item \textbf{tempi morti minimi}: turni brevi, o azioni anche nei turni
        altrui.
\end{itemize}
La ``cooperazione invece dello scontro'' è l'interazione indiretta del German
game; la durata, invece, mostra bene l'evoluzione che racconta Ruggero: gli
Eurogame più complessi, da due o tre ore, sono ormai un'altra cosa rispetto al
gioco per famiglie descritto dal libro.

\subsection{American}

In contrapposizione, oltreoceano c'erano gli \textbf{American} (o
\emph{Ameritrash}): giochi molto focalizzati sulla \textbf{produzione} (le
miniature di plastica), con il \textbf{dungeon crawling}, tematicamente molto
vicini a \emph{Dungeons \& Dragons}, quasi versioni da tavolo del combattimento a
griglia normalizzato da D\&D. Continuano la tradizione di \emph{Chainmail} e dei
giochi strategici e tattici: sono spesso giochi di combattimento, a
\textbf{interazione diretta} (si attacca e ci si difende), e buona parte dei loro
regolamenti è dedicata alla simulazione del combattimento.

\subsection{Strategici e wargame}

\slidefig{img/AGD/agd2-strategici}{Strategici, movimento nascosto, destrezza, 18XX, tempo reale e legacy.}{slide \emph{Analog -- Case Studies}}{fig:agd2-strategici}


Un altro genere gigantesco è quello dei \textbf{giochi di strategia}, usati
storicamente anche nella \textbf{formazione militare}, ``non a caso'': la
\textbf{RAND Corporation} ne ha creati diversi, e ci sono editori specializzati.
Lo scopo era ricreare battaglie storiche. Anche qui, evoluzioni:
\begin{itemize}
  \item i giochi \textbf{a zone} o ad \textbf{aree} (\emph{area control}) con
        unità su una mappa, di cui \emph{Risiko} è il più famoso: io ho un
        esercito, tu ne hai un altro, dobbiamo conquistare territori, e il
        conflitto diretto è inevitabile perché gli obiettivi si sovrappongono;
  \item i \textbf{card-driven strategy game}: strategici con mappe e zone di
        controllo, ma con \textbf{mazzi di eventi}. I designer prendono eventi
        reali di conflitti reali e cercano di renderli coesi tematicamente con
        le meccaniche, dando a ogni evento una risoluzione precisa: ``l'importante
        non è tanto quante unità ho, ma quali carte evento ho in mano''. Sono,
        se ci si pensa, più vicini ai card game che agli strategici.
\end{itemize}
Nicola chiede di \textbf{\emph{Campaign for North Africa}}\footnote{La
trascrizione dice ``Campi for Nord Africa''.}: è uno \textbf{strategico
classico}, con esagoni e pedine di cartoncino, ``un'esperienza devastante in
termini di lunghezza'': forse non è il caso di cominciare da lì.

\subsection{Movimento nascosto}

Il genere degli \textbf{hidden movement}\footnote{La trascrizione dice ``Eat the
Movement''.}: \emph{Scotland Yard}, \emph{Letters from Whitechapel}, ma anche
altri titoli più recenti\footnote{La trascrizione cita ``Sniper Elite, Mind Dame
e GMT''; il secondo è probabilmente \emph{Mind MGMT}, GMT è un editore di
wargame.}. Sono giochi \textbf{asimmetrici}: di solito una persona si nasconde e
scappa, e le altre vincono collettivamente se la catturano. Uno contro tutti.

\subsection{Destrezza}

I \textbf{dexterity game} richiedono \textbf{destrezza fisica}: colpire
componenti con le schicchere, impilarli senza farli cadere. Non basta la
strategia che si ha in testa: ``posso anche avere la strategia migliore, ma se poi
le mie mani non ce la fanno a mettere in piedi questi cilindretti, crolla
tutto''.

\subsection{18XX}

Un genere che Ruggero ama particolarmente, e di cui è ``abbastanza sicuro''
che nessuno abbia mai sentito parlare, anche se conta centinaia di titoli (``più
di 400 giochi'' sotto l'etichetta): i \textbf{18XX}. Il nome viene dal secolo
delle \textbf{ferrovie}, l'Ottocento: sono tutti giochi ferroviari, ognuno
centrato su un anno e un luogo specifici (\emph{1849}, le ferrovie della
Sicilia nel 1849; \emph{1889}, quelle dell'isola di Shikoku). Hanno una forte
componente di \textbf{costruzione delle ferrovie} sulla mappa e una componente
\textbf{economico-azionaria}: comprare e vendere azioni, speculazione
finanziaria. Sono ``molto belli, ma anche molto intensi'': \emph{1889}, che
Ruggero usa di solito come introduttivo, ``dura tre ore''.

\subsection{Tempo reale}

Culturalmente si pensa ai giochi da tavolo come a qualcosa che avviene
\textbf{a turni}, ma non è sempre così. I giochi \textbf{in tempo reale} usano
clessidre, timer o altre meccaniche per cui non si ha tutto il tempo per pensare:
bisogna reagire in fretta. \emph{Kitchen Rush}, per esempio, è ``praticamente
\emph{Overcooked} da tavolo'': si gestisce un ristorante, e \textbf{la clessidra
è il lavoratore}, che non si può spostare finché non si è esaurita. Chi prende
gli ingredienti, chi li prepara sul fornello, chi va in sala a prendere gli
ordini, chi porta i piatti, chi li lava: ne nasce un certo caos nell'intreccio.

\subsection{Legacy}

I \textbf{Legacy} sono ``un po' una rivoluzione nel genere''. Di solito un gioco
da tavolo resta sempre quello: lo si prepara, lo si gioca, la settimana dopo lo
si gioca con un altro gruppo. I giochi Legacy invece chiedono di avere un
\textbf{impatto fisico sul gioco}, per cui \textbf{le partite successive sono
influenzate dalle precedenti}:
\begin{itemize}
  \item \emph{Pandemic Legacy} chiede di \textbf{strappare carte}; alcune città
        vengono distrutte dalla pandemia, coperte con degli \textbf{adesivi}, e
        non esistono più sulla mappa;
  \item \emph{SeaFall}\footnote{La trascrizione dice ``Sifol''.}, sulla scoperta
        di isole in un Ottocento alternativo: le isole hanno il nome vuoto, e chi
        le scopre le battezza; da lì in poi l'isola porterà quel nome per tutto
        il resto del gioco.
\end{itemize}
``So che qualcuno sta già sudando freddo all'idea di strappare carte o scrivere
sul tabellone'', ma sono ``esperienze emotivamente incredibili'': ``giocatevi dei
Legacy''.

\paragraph{Approfondimento: Risk Legacy e la rottura di un'assunzione.} Schell
dedica ai Legacy due passaggi di \emph{The Art of Game Design}. Nel capitolo sul
\textbf{tema}, presenta \emph{Risk Legacy} di \textbf{Rob Daviau} come esempio di
tema fondato su una verità: Daviau ha fatto ``qualcosa di mai tentato prima nel
mondo dei giochi da tavolo'', creando meccaniche per cui le scelte dei giocatori
\textbf{cambiano il gioco per sempre}. Le regole impongono di modificare il
tabellone con adesivi non rimovibili, di scriverci sopra con un pennarello
indelebile per rivendicare i territori, di strappare e scartare carte e perfino
di modificare permanentemente il regolamento. Queste meccaniche, nota Schell,
sono nuove e intriganti, ma soprattutto \textbf{rafforzano il tema centrale}:
``la guerra cambia il mondo''. Nel capitolo sulle tecniche di brainstorming,
Schell racconta come Daviau ci sia arrivato: elencando le \textbf{assunzioni}
date per scontate sui giochi da tavolo e provando a ``farne saltare'' una,
``una partita non influenza la successiva''. È un buon esempio di quello che
Ruggero chiedeva all'inizio: guardare un genere che si conosce e chiedersi che
cosa succederebbe se una delle sue regole implicite non valesse più.

% ---------------------------------------------------------------------
\section{Due libri sul gioco da tavolo}
\label{a2:sec:tavolo-libri}

\slidefig{img/AGD/agd2-libri}{I due libri consigliati sul gioco da tavolo.}{slide \emph{Analog -- Case Studies}}{fig:agd2-libri}


A proposito di risorse, Ruggero torna sul libro di Nuccio, \emph{La
progettazione dei giochi da tavolo}, citato a inizio lezione, e indica
un'alternativa in inglese che a suo parere ``fa quello che fa Nuccio, meglio'':
\textbf{\emph{Building Blocks of Tabletop Game Design}} di \textbf{Geoffrey
Engelstein e Isaac Shalev}\footnote{La trascrizione dice ``Engelstein e
Schaler'' e ``Building Blocks of Tablet of Game Design''.}. È più
\textbf{contemporaneo} negli esempi e ``più funzionale''. Per i giochi da tavolo,
se si ha già un framework di riferimento, sono entrambi ``una risorsa
incredibile''. (Resta valido l'avvertimento di inizio lezione: sono archivi di
meccaniche, da usare dopo aver capito che cosa le meccaniche fanno.)

% ---------------------------------------------------------------------
\section{Libri-gioco, megagame e altri ibridi}
\label{a2:sec:tavolo-ibridi}

\slidefig{img/AGD/agd2-ibridi}{Libri-gioco, megagame e ibridi tra gioco da tavolo, gioco di ruolo e videogioco.}{slide \emph{Analog -- Case Studies}}{fig:agd2-ibridi}


``Manca ancora qualcosa?'' Mancano diverse cose.
\begin{itemize}
  \item I \textbf{libri-gioco}, nel formato libro tradizionale e nel formato
        ``più contemporaneo'' dei giochi da tavolo che sono libri-gioco, come
        \emph{Sherlock Holmes Consulente Investigativo}.
  \item I \textbf{megagame}: giochi da tavolo \textbf{per cento persone}. Non
        hanno un'organizzazione editoriale come \emph{Monopoly} o gli altri
        board game: sono organizzati in modo più vicino ai giochi pervasivi e al
        LARP, spesso \textbf{a evento}, spesso con biglietto, anche perché
        gestire un gioco da tavolo con cento persone in parallelo ``può diventare
        molto intenso''.
  \item I \textbf{LARP urban game}, di cui ha parlato Beatrice, con esempi
        storici anche italiani, come \emph{Pathos}\footnote{Grafia come nella
        trascrizione.}: giochi di ruolo dal vivo che sono anche urban game, e
        un po' ARG.
  \item Produzioni ispirate al freeform ma vendute come giochi da tavolo, come
        \textbf{\emph{Inhuman Conditions}}\footnote{La trascrizione dice ``In
        Human Condition''.}: praticamente il \textbf{test di \emph{Blade Runner}}
        per scoprire chi è un androide, per due persone. È un gioco di ruolo, ma
        prodotto esattamente come un gioco da tavolo (carte, scatola,
        regolamento) e non si presenta come gioco di ruolo, ``ma questo non lo
        rende meno gioco di ruolo''.
  \item Ibridi con il party game, come \textbf{\emph{Sì, Oscuro Signore!}}: un
        party game, ma anche un \emph{role-playing game}, perché si fa il role
        play del proprio personaggio.
  \item Ibridi con il videogioco, come \emph{Fateforge} e gli altri giochi che
        usano un'\textbf{app} non solo come segnapunti ma come meccanica durante
        la partita.
\end{itemize}

\begin{table}[htbp]
\centering
\small
\begin{tabular}{@{}p{3.0cm}p{9.9cm}@{}}
\toprule
\textbf{Genere} & \textbf{Che cosa lo caratterizza ed evoluzioni citate} \\
\midrule
Roll and move & tira e muovi (gioco dell'oca); agency quasi nulla; poi
elementi economici (\emph{Monopoly}) o deduttivi (\emph{Cluedo}) \\
Astratti & tema ridotto all'osso, informazione aperta (scacchi, dama, go,
\emph{ZÈRTZ}, progetto GIPF) \\
Party game & molte persone, regole leggere, partite brevi, punto di accesso
(\emph{Taboo}, \emph{Lupus in Tabula}, \emph{Wavelength}) \\
Giochi di parole & creare parole per fare punti (\emph{Boggle},
\emph{Scarabeo}); cooperativi (\emph{Cross Clues}), boss fight
(\emph{Illiterati}) \\
Giochi di carte & mazzi tradizionali; \emph{Love Letter}; take that
(\emph{Munchkin}); TCG, LCG e deck-building \\
German game ed Eurogame & interazione indiretta, costruzione e scambio;
simulazioni focalizzate di due o tre ore \\
American & produzione in plastica, dungeon crawling, combattimento e
interazione diretta \\
Strategici & formazione militare; aree e zone (\emph{Risiko}); card-driven;
esagoni (\emph{Campaign for North Africa}) \\
Movimento nascosto & uno contro tutti, asimmetrici (\emph{Scotland Yard},
\emph{Letters from Whitechapel}) \\
Destrezza & schicchere, impilare: conta anche la mano \\
18XX & ferrovie e azioni dell'Ottocento; più di 400 titoli (\emph{1849},
\emph{1889}) \\
Tempo reale & clessidre e timer (\emph{Kitchen Rush}) \\
Legacy & le partite cambiano fisicamente il gioco (\emph{Pandemic Legacy},
\emph{SeaFall}) \\
Libri-gioco, megagame, ibridi & \emph{Sherlock Holmes Consulente
Investigativo}; cento giocatori a evento; \emph{Inhuman Conditions},
\emph{Sì, Oscuro Signore!}, giochi con app \\
\bottomrule
\end{tabular}
\caption{I generi di gioco da tavolo presentati da Ruggero, in sintesi.}
\label{a2:tab:tavolo}
\end{table}

% ---------------------------------------------------------------------
\section{La chiusura: niente interrogazioni, e i materiali}
\label{a2:sec:chiusura}

``E ho corso un po', ma ce l'abbiamo fatta''. Mancano gli ultimi dieci minuti, e
Ruggero si complimenta con chi è ``arrivato fino a qui''. Ripete il punto
fondamentale: \textbf{non ci sarà nessuna interrogazione} su questi giochi
(``Tiziana, ti do subito il sollievo''). Non si chiede di sapere che cosa sia ogni
caso studio: lo scopo era dare una \textbf{panoramica}, degli spunti, e far
sapere che ``questa roba esiste''. Le slide restano, e ci si potrà tornare in
fase di ricerca e sviluppo. Con la lezione di \textbf{ricerca e sviluppo}, che
darà gli strumenti per farlo in autonomia, ognuno troverà \textbf{i casi studio
che servono al proprio gioco}, ``e io ve li chiederò, noi ve li chiederemo''.
Questa lezione serviva ad \textbf{aprire le possibili direzioni}.

\paragraph{I materiali caricati.} Mentre Ruggero parla, le slide vengono
caricate nella cartella degli studenti. Beatrice aggiunge di aver caricato:
\begin{itemize}
  \item il paragrafo di \textbf{Frank Rose} su \emph{Why So Serious};
  \item un estratto sulla \textbf{definizione dei pervasive game} tratto da
        \emph{Game Design}, un libro di \textbf{Maresa Bertolo e Ilaria
        Mariani}\footnote{La trascrizione dice ``Marianne Bertolo'' e
        ``Bertolo Mariani''. Il libro è con ogni probabilità \emph{Game
        Design. Gioco e giocare tra teoria e progetto} di Maresa Bertolo e Ilaria
        Mariani.}, ``un ottimo libro''.
\end{itemize}
Non sono letture obbligatorie.

\paragraph{Un avvertimento sui libri italiani di gioco di ruolo.} Daniele
aggiunge una cautela sulla parte dedicata al gioco di ruolo del libro di Bertolo
e Mariani: alcuni paragrafi ``generalizzano molto verso il mainstream''. E poi un
\emph{warning} più generale: \textbf{la maggior parte della letteratura italiana
sul gioco di ruolo è ``terribile''}, al punto che alcuni libri i docenti non solo
non li consigliano, ma li \textbf{sconsigliano attivamente}. Con la premessa che
non è nulla di personale contro l'autore, Daniele sconsiglia il manuale
\emph{Giochi di ruolo} di \textbf{Mauro Longo}\footnote{La trascrizione riporta
l'editore come ``Dino Dino Editori'': probabilmente Dino Audino Editore.},
perché è molto soggettivo, molto legato alla visione del gioco di ruolo
dell'autore e a un solo tipo di gioco, quello di \emph{Lex Arcana}\footnote{La
trascrizione dice ``l'ex arcana''.}, con teorie proprie che non corrispondono al
design generale dei giochi di ruolo. ``Va bene se volete fare \emph{Lex Arcana}
esattamente come \emph{Lex Arcana}'', non se si vuole fare altro.

Meglio il volume sull'edu-larp, e soprattutto \textbf{giocare}: ``giocate i giochi,
perché in Italia abbiamo una creatività sfrenata''. Si trovano ``robe favolose''
online: tantissimi \textbf{one page game} di ruolo, il mondo del \textbf{print and
play} anche nei giochi da tavolo, chamber LARP e freeform da stampare
direttamente da internet. ``Fatevi una ricerca e vedete quello che vi garba''.

% ---------------------------------------------------------------------
\section{Le domande finali}
\label{a2:sec:domande}

``Siamo stati più veloci dell'anno scorso'' (quando si erano accumulati trenta
minuti di ritardo), e restano sette minuti per le domande.

\subsection{Che cosa fa scegliere una forma?}

Gianmarco, ``un po' per rompere il ghiaccio'', fa una domanda ``molto generica'',
che sa non avere una risposta semplice: si sono visti giochi dalle forme
diversissime, ma i \textbf{framework} per crearli sono più o meno gli stessi. Che
cosa determina la scelta di una forma rispetto a un'altra? Ruggero rimanda la
risposta alla lezione sull'\textbf{emotional design}, ``che è proprio una risposta
a questa cosa'': la risposta è così lunga che ``abbiamo una lezione di tre ore per
rispondere alla tua domanda''.

\subsection{Che cosa va per la maggiore in Italia?}

Angelo chiede quale genere ``vada per la maggiore'' in Italia. Ma che cosa
significa ``andare per la maggiore''? Il successo commerciale non è uguale per
tutti i generi, ``nemmeno nei videogiochi''. Uno dei docenti (con ogni probabilità Daniele) risponde con due domande
``provocatorie'', che valgono anche per i libri (ne parlava quel giorno con
un'amica narrative designer di videogiochi):
\begin{enumerate}
  \item quante copie deve vendere un gioco per essere un successo?
  \item il numero di copie vendute coincide con il successo
        commerciale ed economico?
\end{enumerate}
La risposta a entrambe è \textbf{no}, ma si capirà con la lezione di
\textbf{Flavio} sull'aspetto produttivo: quando si ragiona in logica commerciale
si è ``in logica di prodotto'', e ci sono varie \textbf{metriche} e varie scale di
successo. A questo si aggiungono altri livelli:
\begin{itemize}
  \item \textbf{quale mercato}? In Italia \textbf{i giochi che non sono in
        italiano non vendono}, e molti editori italiani sono innanzitutto
        \textbf{editori di localizzazione}: prendono giochi dall'estero, li
        traducono, e questo è buona parte della loro sopravvivenza, nel gioco
        da tavolo come nel gioco di ruolo. Un gioco in inglese si vende in
        Inghilterra, in America e in moltissimi paesi dove l'inglese è la seconda
        lingua; un gioco in italiano si vende solo in Italia, e un editore
        italiano non ha facilmente la struttura per arrivare all'estero. Per
        questo anche autori italiani di giochi da tavolo pubblicano con
        \textbf{editori esteri};
  \item \textbf{perché si fa il gioco}? Non è detto che lo scopo sia
        commerciale. Un gioco può essere \textbf{formativo o didattico}, fatto
        con un'azienda o una scuola, o \textbf{commissionato} da un cliente per
        uno scopo specifico: allora non interessa che venda mille copie. Il
        designer viene pagato, ma il pagamento non corrisponde alle vendite.
\end{itemize}

Beatrice chiede ad Angelo perché faccia la domanda, perché i \textbf{settori}
appena visti spesso ``non si parlano minimamente'', e alcuni non sanno nemmeno
dell'esistenza degli altri: la cross-medialità presentata nella lezione non viene
presentata spesso così. Se la domanda nasce da prospettive lavorative, come per
alcuni nel corso, la risposta è che ci sono ``immense possibilità'' a seconda
della direzione: fare un gioco da tavolo è una cosa, lavorare nella
\textbf{formazione} un'altra, e c'è un intero ambito non trattato, quello dei
\textbf{giochi per le aziende}.

Angelo precisa: voleva solo sapere \textbf{che cosa giocano di solito gli
italiani}. È ``un altro volo pindarico'': in Italia è difficile trovare dati
statistici affidabili, perché i questionari e le survey sono spesso invalidati
dal modo in cui sono formulate le domande. Da quello che sa il docente che risponde\footnote{La trascrizione dice ``poi Lucio mi correggerà''; il nome non è ricostruibile con certezza.}, ci sono sicuramente \textbf{più giocatori di giochi da
tavolo che di giochi di ruolo}: il gioco da tavolo è più conosciuto e
frequentato, anche nelle ludoteche, mentre il gioco di ruolo e il LARP richiedono
un investimento economico diverso e sono soprattutto pratiche
\textbf{comunitarie}, ``tribali'', che nascono dal basso e si trasmettono
\textbf{oralmente}. Chi gioca ai giochi da tavolo li compra; chi gioca di ruolo
spesso lo fa perché qualcuno lo fa giocare, e non è detto che abbia il gioco in
casa.

\subsection{Avvisi}

Intanto Daniele ricorda che il \textbf{questionario} della prima slide va
compilato \textbf{entro il 15}, e che Giordano deve ancora presentarsi (``siamo
delle pessime persone''). Visto che è ora di cena, Giordano si presenterà come
\textbf{prima cosa} alla lezione successiva, \textbf{lunedì}, sull'emotional
design (``ricordacelo, perché se no ci dimentichiamo'').

% ---------------------------------------------------------------------
\section{Due cimiteri: i giochi di carte collezionabili e i ``fix''}
\label{a2:sec:cimiteri}

Prima di chiudere, a un docente torna in mente ``una cosa importantissima''
che voleva dire.

\subsection{Non fate Magic: The Gathering 2}

``Vi prego, non cadete nella trappola'' di vedere un gioco di successo commerciale
e farne uno simile pensando che venderà moltissimo. ``Ci possiamo fare un
cimitero, una fossa comune'' con migliaia di \textbf{trading card game} e
\textbf{collectible card game}. ``Non fate un trading card game. Ve lo dico per
voi: non fate \emph{Magic: The Gathering 2}''. Chi vede \emph{Magic} o
\emph{Pokémon} pensa di poter fare un sacco di soldi, ma \textbf{non vede la
struttura commerciale e produttiva gigantesca} che c'è dietro, quasi impossibile
da emulare dal basso. Quasi tutti, anche editori grossi, quelli che hanno provato
a competere con \emph{Pokémon}, \emph{Yu-Gi-Oh!} e \emph{Magic} ``sono andati a
picco''. ``Se volete perdere soldi, quello è il modo perfetto''.

Un collega aggiunge che c'è un \textbf{revival dei TCG}, ma ogni titolo dura un
anno o due e poi muore. Si sperimentano anche dei \textbf{micro TCG} annunciati
fin dall'inizio come limitati (``questo gioco avrà tre espansioni e poi
chiuderà''), una sorta di ibrido con gli LCG. Ma attenzione: spesso sono fatti
da ex designer di LCG, di \emph{Magic} o di \emph{Hearthstone}, persone con
contatti, entourage e un'azienda alle spalle; da fuori si vede una persona, ma
dietro c'è una struttura che non si può replicare.

\paragraph{Approfondimento: perché Pokémon non si replica.} \emph{The Art of
Game Design} offre un buon esempio della ``struttura gigantesca'' di cui parlano
i docenti. Nel capitolo sui mondi transmediali, Schell racconta che il gioco di
carte collezionabili di \emph{Pokémon} nacque nel 1999 da una collaborazione tra
Nintendo e Wizards of the Coast, la casa di \emph{Magic}, e che fu progettato per
restare il più possibile fedele alle meccaniche del videogioco per Game Boy. Il
gioco di carte era la \textbf{terza porta d'ingresso} in un unico mondo, accanto
al videogioco e alla serie TV: portatile, sociale, economico, perfetto per lo
slogan ``acchiappali tutti''. Le tre porte si rafforzavano a vicenda (voler
vincere spingeva a guardare la serie, e la serie rendeva il gioco più vivido), e
Schell ricorda un dirigente dell'intrattenimento che nel 1999 profetizzava la
fine della ``moda'' con l'uscita del film. Non andò così, appunto, perché il
successo non stava in un singolo prodotto ma nell'ecosistema. Un gioco di carte
collezionabili nato dal basso, senza un ecosistema del genere, compete con questo.

\subsection{Non fate un gioco per ``aggiustarne'' un altro}

Un altro docente (probabilmente Daniele, che cita ``le lezioni di framework con Ruggero'') aggiunge un altro cimitero: quello di chi vuole fare un gioco per
\textbf{``fixare'' un gioco esistente}, cioè farne una versione migliore. Come
esperimento può essere interessante, ma, come si vedrà nella lezione
sull'emotional design e nelle lezioni sui framework con Ruggero, l'intenzione deve
essere ``\textbf{eticamente onesta}''. Non si vuole aggiustare D\&D per il piacere
di aggiustare D\&D: bisogna chiedersi sempre \textbf{che cosa non mi piace
dell'esperienza} di D\&D, che cosa voglio aggiustare di quell'esperienza. E non
creare ``D\&D dove tolgo una regola e ne metto un'altra, ed è sistemato'', ma
\textbf{partire dall'esperienza che si vuole raggiungere}, quella che non si è
avuta da D\&D, e costruirla. Con la lezione di ricerca e sviluppo si vedrà anche
come prendere spunto dai giochi esistenti, non per copiarli ma per capire che
cosa fanno, perché \textbf{prenderne pezzi senza capire che cosa fanno} è proprio
l'errore da evitare (lo stesso del libro di Nuccio letto troppo presto).

\paragraph{I saluti.} Un docente deve scappare e saluta; gli altri ricordano il
questionario, che le loro email sono nelle slide per eventuali domande, e
l'appuntamento a lunedì.

% ---------------------------------------------------------------------
\section*{Punti chiave del capitolo}
\addcontentsline{toc}{section}{Punti chiave del capitolo}

\begin{itemize}
  \item Un \textbf{caso studio} è un gioco, o una categoria di giochi, che,
        analizzato tenendo conto del suo contesto, dà informazioni e metodi
        utili per i propri giochi: come leggere libri per imparare a scrivere.
        Questa lezione è una panoramica; gli strumenti di analisi arriveranno con
        la ricerca e sviluppo. ``Per me, questo è design'' (Ruggero).
  \item \textbf{Gioco analogico} è un termine ombrello per tutto ciò che non è
        videogioco. Le quattro \textbf{macrocategorie} sono gioco da tavolo,
        gioco di ruolo da tavolo, gioco di ruolo dal vivo e gioco pervasivo, ma
        i confini sono porosi: \textbf{ibridazione e sperimentazione}. Il corso è
        \textbf{agnostico di genere}, e il project work è libero, ma non deve
        essere un gioco enorme.
  \item Il \textbf{gioco pervasivo} usa luoghi e oggetti della vita quotidiana
        come luoghi e oggetti di gioco; in termini di game studies,
        \textbf{espande il cerchio magico} nello spazio, nel tempo e nella
        società. Casi: \emph{Pokémon Go}, geocaching, \emph{GuruShot},
        \emph{Journey to the End of the Night}, SF0, \emph{Critical City Upload},
        \emph{Sleep No More}, gli ARG (\emph{Why So Serious}, \emph{The
        Institute}, \emph{Stranger Things} a Milano), Collettivo Cinetico,
        Amigdala, \emph{MP3 Experiment}, \emph{PacManhattan}, \emph{Zombies,
        Run!}, \emph{Reptilarium}, \emph{Echi resistenti}.
  \item Esperienze emotivamente significative si possono fare \textbf{con budget
        zero}: l'audio che guida, gli oggetti da portare, un finale
        scenografico.
  \item Il \textbf{LARP} è separato dal GDR per ragioni di industria (evento con
        biglietto contro oggetto venduto). Tipi: boffer, Mind's Eye Theatre,
        chamber/parlor, American freeform, murder party, Nordic LARP (Jeepform,
        360°, tape, black box).
  \item Il \textbf{Nordic LARP} ha documentato tutto (Knutpunkt dal 1997, libri
        annuali, manifesti come il Dogma 99), punta su coinvolgimento, emozione,
        coesione tematica e dopo-gioco, e ha normalizzato le \textbf{meccaniche di
        sicurezza}. Non è un esperimento scientifico: lo scopo resta ludico.
  \item Nel \textbf{GDR da tavolo} contano la fiction emergente e quattro parole:
        \textbf{conflitto, azione, risoluzione, effetto}. Il master è un
        giocatore. Dalla nascita di D\&D (1974) dal wargame, attraverso gli
        adattamenti a nastro degli anni '80--'90, si arriva alla \textbf{New
        Wave} di The Forge (``il sistema conta''), ai \textbf{giochi
        focalizzati} (non semplificati), all'OSR, agli story game, ai giochi per
        due, in solitario e micro.
  \item Nel \textbf{gioco da tavolo} ogni genere si è evoluto: roll and move,
        astratti, party game, giochi di parole, giochi di carte (TCG, LCG,
        deck-building), German game ed Eurogame, American, strategici,
        movimento nascosto, destrezza, 18XX, tempo reale, \textbf{Legacy},
        libri-gioco, megagame, ibridi.
  \item Libri: il volume sull'\textbf{edu-larp} sì; Nuccio e Engelstein--Shalev
        dopo l'emotional design; Bertolo--Mariani con cautela sulla parte di GDR;
        gran parte della letteratura italiana sul GDR no. Soprattutto,
        \textbf{giocare}.
  \item Il \textbf{successo} non si misura solo in copie, e dipende dal mercato
        (in Italia vende solo ciò che è in italiano) e dallo scopo del gioco.
  \item Due cimiteri: \textbf{non fare un gioco di carte collezionabili} e non
        fare un gioco per \textbf{``aggiustarne'' un altro} senza partire
        dall'esperienza che si vuole creare.
\end{itemize}

%
%  Tutte le etichette di questo capitolo hanno il prefisso "a2:"
%  (a2:cap:..., a2:sec:..., a2:fig:..., a2:tab:...). Il capitolo è
%  autonomo: non contiene \ref verso altri capitoli.
%  Le figure e le tabelle sono ricostruzioni schematiche di quanto i
%  docenti spiegano a voce e con le slide, non copie delle slide.
% =====================================================================

\chapter[I casi studio dell'Analog Game Design]{I casi studio: giochi
pervasivi, giochi di ruolo dal vivo, giochi di ruolo da tavolo e giochi da
tavolo}
\label{a2:cap:casi}

Questa è la \textbf{seconda lezione} del corso di \textbf{Analog Game Design}, e
mantiene la promessa fatta alla fine della prima: è una lezione di tre ore
dedicata quasi interamente ai \textbf{casi studio} (\emph{case studies}). I
docenti la definiscono ``un overview'', una panoramica ``a volo d'aquila'':
``voleremo in tutti i settori ludici analogici e guarderemo che cosa c'è''.
Non si scende nei dettagli tecnici dell'analisi (ci saranno lezioni apposite,
soprattutto quella di \textbf{ricerca e sviluppo}), ma si apre la ``diga'' degli
esempi: decine di giochi, eventi, esperienze, movimenti e manifesti, scelti per
mostrare quante cose esistono, quanto sono diverse tra loro e quanto i confini
tra i generi siano porosi.

La lezione si svolge \textbf{online}, su Zoom, un \textbf{venerdì} sera (``che
bella la lezione del venerdì! Sembra quasi un aperitivo''), ed è divisa tra i
tre docenti principali del corso, ciascuno sul proprio territorio:
\textbf{Ruggero} introduce l'argomento e presenta i \textbf{giochi di ruolo dal
vivo} (LARP) e i \textbf{giochi da tavolo}; \textbf{Beatrice} presenta i
\textbf{giochi pervasivi}; \textbf{Daniele} presenta i \textbf{giochi di ruolo
da tavolo}. ``Ce lo smezziamo un po' tutti'', perché l'argomento è diviso ``in
tante categorie, sottocategorie, sotto-sottocategorie''. Dopo i pervasivi c'è
una pausa di pochi minuti; gli ultimi dieci minuti sono lasciati alle domande.

Tre avvertenze di lettura, le stesse del capitolo precedente. La prima: la
trascrizione è automatica e storpia moltissimi nomi propri, soprattutto i titoli
dei giochi, quasi sempre in inglese e pronunciati velocemente (``Blaze in the
Dark'', ``Cobo Bobby'', ``Sifol'', ``Jeep Form''\ldots). Dove la ricostruzione è
sicura, o molto probabile, il capitolo usa la forma corretta e lo segnala in
nota; dove non lo è, riporta la forma della trascrizione e lo dice. La seconda:
durante la pausa la trascrizione registra, ripetute in loop per decine di righe,
alcune frasi senza senso (anche all'inizio della registrazione succede qualcosa
di simile), che il capitolo ignora. La terza: i docenti si passano spesso la
parola senza che la trascrizione segni il cambio di voce; il capitolo attribuisce
le parole a chi le pronuncia quando il contesto lo rende chiaro (chi parla di sé,
chi viene chiamato per nome, chi ``passa la parola''), e resta generico
altrimenti.

Quanto alla data, la lezione si colloca con ogni probabilità \textbf{venerdì 11
settembre 2026}\footnote{La data non è mai detta. La si deduce da alcuni
indizi: la prima lezione si era chiusa con ``ci si vede questo venerdì''; uno
studente assente ha recuperato ``la lezione di lunedì''; a fine lezione si
ricorda che il questionario iniziale va compilato ``entro il 15'' (settembre) e
che la lezione successiva, sull'\emph{emotional design}, sarà di lunedì. Nel
2026 il primo venerdì utile prima del 15 settembre, dopo il lunedì della prima
lezione, è l'11. È una deduzione, non un'informazione esplicita.}: i riferimenti
a ``l'altra volta'' e ``l'anno scorso'' vanno letti rispetto a quel momento.

\paragraph{Come leggere questo capitolo.} Più ancora della lezione precedente,
questa è un \textbf{catalogo}. I docenti lo dicono esplicitamente: non ci sarà
nessuna interrogazione sui singoli giochi, e lo scopo non è ricordarli tutti ma
sapere che esistono, avere dei punti di riferimento e, soprattutto, abituarsi a
guardare oltre il proprio genere preferito. Per questo il capitolo riporta
\emph{tutti} i casi citati, con quello che i docenti ne dicono, ma li organizza
in tabelle e schemi per renderli consultabili ``in fase di ricerca e sviluppo'',
esattamente come fa Ruggero con le sue slide (``non sapete quante volte mi sono
trovato a riaprire proprio esattamente queste slide''). Qua e là, dei paragrafi
di \textbf{approfondimento} collegano i casi ai libri di game design di
riferimento (\emph{The Art of Game Design} di Jesse Schell, \emph{Challenges for
Game Designers} di Brenda Brathwaite e Ian Schreiber, \emph{21st Century Game
Design} di Chris Bateman e Richard Boon).

Il percorso del capitolo segue quello della lezione:
\begin{enumerate}
  \item l'avvio, Discord, i materiali e i \textbf{libri consigliati} (e
        sconsigliati) per cominciare (sezioni~\ref{a2:sec:avvio}
        e~\ref{a2:sec:libri});
  \item \textbf{perché} analizzare i casi studio e \textbf{che cosa} sono
        (sezione~\ref{a2:sec:perche});
  \item le \textbf{quattro macrocategorie} del gioco analogico, l'ibridazione e
        l'invito a non chiudersi nel proprio genere (sezioni~\ref{a2:sec:macro}
        e~\ref{a2:sec:aperti});
  \item i \textbf{giochi pervasivi}, con Beatrice: la definizione, il cerchio
        magico e diciotto casi studio, dalle app agli \emph{alternate reality
        game}, dal teatro interattivo alle cacce al tesoro
        (sezioni da~\ref{a2:sec:pervasivi} a~\ref{a2:sec:pervasivi-edu});
  \item i \textbf{giochi di ruolo dal vivo}, con Ruggero: boffer, Mind's Eye
        Theatre, chamber e parlor LARP, American freeform, murder party, il
        \textbf{Nordic LARP} con i suoi manifesti, lo Jeepform, il 360°, i tape
        e i black box LARP (sezioni da~\ref{a2:sec:larp} a~\ref{a2:sec:oppresso});
  \item i \textbf{giochi di ruolo da tavolo}, con Daniele: le quattro parole
        (conflitto, azione, risoluzione, effetto), la storia dal primo
        \emph{Dungeons \& Dragons} alla New Wave di The Forge, l'OSR, i giochi
        focalizzati, gli story game, i giochi per due, in solitario e i micro
        giochi di ruolo (sezioni da~\ref{a2:sec:gdr} a~\ref{a2:sec:gdr-formati});
  \item i \textbf{giochi da tavolo}, di nuovo con Ruggero: diciassette generi,
        dal \emph{roll and move} ai Legacy, e gli ibridi
        (sezioni da~\ref{a2:sec:tavolo} a~\ref{a2:sec:tavolo-ibridi});
  \item la chiusura, i materiali caricati, le \textbf{domande finali} e due
        avvertimenti importanti: non fare un gioco di carte collezionabili, non
        fare un gioco per ``aggiustare'' un altro gioco
        (sezioni da~\ref{a2:sec:chiusura} a~\ref{a2:sec:cimiteri}).
\end{enumerate}

% ---------------------------------------------------------------------
\section{L'avvio della lezione}
\label{a2:sec:avvio}

La registrazione parte mentre i docenti si stanno ancora collegando (``mamma
mia, la burocrazia italiana'') e si controlla che il simbolo REC lampeggi.
Mentre Daniele prova a condividere le slide (``aspettate, sono un attimo
bloccato'') e chiede di essere nominato co-organizzatore della riunione, gli
studenti entrano uno alla volta: Bea, Lorenzo, Nicola, Rita, Gianmarco,
Tiziana, Giordano. L'atmosfera è rilassata. Si scherza sul fatto che la lezione
del venerdì ``sembra quasi un aperitivo'': un docente si mette ``bello comodo'',
un altro ha a casa ``le cose per fare l'aperitivo'' e chiede il permesso di
prepararlo nel frattempo, ``purché siate sul pezzo''. Daniele scrive anche di
dover ancora rispondere a un messaggio di Beatrice; ``non preoccuparti'', gli
risponde lei: settembre è un mese di ripartenza, ``in realtà è l'inizio
dell'anno reale''.

\paragraph{Giordano.} Arriva \textbf{Giordano}, che era assente alla prima
lezione. Ha recuperato la registrazione di lunedì, e sa che gli si chiederà di
presentarsi. (Alla fine, per mancanza di tempo, la presentazione verrà rinviata
alla lezione successiva: si veda la sezione~\ref{a2:sec:domande}.) All'inizio
mancano ancora Lorenzo e Tancredi; Lorenzo arriva poco dopo, Tancredi dopo la
pausa.

\paragraph{Discord.} Uno dei docenti nota che gli studenti sono ``belli attivi
su Discord'' e se ne rallegra: ``sfruttatelo, perché l'anno scorso lo si
sfruttava poco e invece è uno spazio molto importante''. I docenti leggono il
server ogni tanto, e le domande sono benvenute. L'unica richiesta è
\textbf{taggarli}: ``tra tanti server Discord non riusciamo a leggere'' tutto,
quindi una domanda rivolta a loro va segnalata con un tag. Lo spazio, ricorda,
``è fatto per voi'': per confrontarsi, condividere cose e materiali.

\paragraph{Le dispense.} Gli studenti hanno visto i materiali caricati: le
slide, le dispense ``e quant'altro''. Qualcuno chiede se le dispense siano in
inglese e se siano articoli. Sono i \textbf{saggi citati nella prima lezione}
(per esempio quelli sull'MDA e sul DPE), e \textbf{non è obbligatorio leggerli}:
sono lì ``per completezza'' e per mostrare che dietro gli argomenti delle
lezioni ``c'è una base teorica''. L'importante è che le lezioni facciano capire
gli argomenti. Spesso i docenti citeranno \textbf{paper accademici} di game
design e di \emph{game studies}, e ``purtroppo c'è poco da fare'': molti non sono
mai stati tradotti, o lo sono stati ``in maniera un po' approssimativa''. Per
questo di solito si linka la fonte originale; nulla vieta di cercare traduzioni
o di tradurseli per conto proprio. In ogni caso i docenti cercano di scegliere
paper \textbf{accessibili}: ``non vi linkeremo delle tesi di master di 300
pagine, a meno che non ce le chiediate voi''.

% ---------------------------------------------------------------------
\section{Due libri per cominciare (e uno da rimandare)}
\label{a2:sec:libri}

Prima di entrare nel tema, Daniele torna su due libri citati la volta
precedente, mostrandoli in video.

\subsection{Un volume sull'edu-larp}

A Tiziana, che ha qualche difficoltà con l'inglese, Daniele consiglia un libro
italiano in cui si trovano molti degli argomenti del corso: un volume
sull'\textbf{edu-larp}, il gioco di ruolo dal vivo in ambito educativo, a cura
di \textbf{Andrea Maragliano} ed edito da \textbf{FrancoAngeli}\footnote{La
trascrizione riporta il titolo come ``EduLARP'', ``Edu LARP'' e ``DualArp'', e
il nome del curatore come ``Andrea Maragliano''; il titolo completo del volume
non è detto a lezione.}. Con una frecciata affettuosa, Daniele aggiunge che
contiene ``un'introduzione molto interessante scritta da un certo Ruggero. Non
so chi sia''. Ruggero protesta (``non fare queste cose, che poi sembra che
consiglio i libri in cui ho partecipato''): il libro non è suo, ci hanno scritto
``tutti i miei amici'', e lui per principio non consiglia i libri su cui scrive
(``possiamo anche comprarlo e non leggere l'introduzione'', ride un collega).

Daniele però lo consiglia convintamente:
\begin{itemize}
  \item molti concetti che si trovano nei siti in inglese ``qua dentro sono
        molto spiegati''; lui stesso lo ha usato per studiare;
  \item c'è una parte, ``bellissima'', sull'\textbf{involvement} (il
        coinvolgimento), che spiega anche ``perché non usiamo la parola
        immersione'' (una delle parole vietate della prima lezione); se ne
        parlerà anche con Ruggero;
  \item finalmente si può consigliare ``roba genuina, sana, italiana di game
        design'';
  \item copre l'ambito \textbf{didattico e formativo}, cioè il gioco nelle
        scuole e il gioco come strumento di formazione, che interessa ad alcuni
        studenti del corso.
\end{itemize}
Non bisogna leggerlo tutto: è utile ``leggerlo pian pianino mentre fate il
corso'', se si vuole approfondire.

\subsection{La progettazione dei giochi da tavolo di Walter Nuccio}

Il secondo libro mostrato la volta precedente è \emph{La progettazione dei
giochi da tavolo} di \textbf{Walter Nuccio}. Qui Daniele fa un ``piccolo
cesellamento'': ``se fossi in voi non lo comprerei per il momento'', e lascia la
parola a Ruggero, che ne ha una copia alle spalle. Il giudizio di Ruggero è
preciso:
\begin{itemize}
  \item è più utile \textbf{dopo} aver interiorizzato alcune basi su come si
        progettano i giochi; come primo libro è un rischio, ``perché sembra che
        ti dia delle scorciatoie'';
  \item più che un saggio di game design, è \textbf{``un grosso archivio di
        meccaniche di gioco''}, una sorta di dizionario (``che cos'è un malus,
        per esempio?''), con una parte di \emph{pattern};
  \item il pericolo è applicare quelle meccaniche ai propri giochi
        \textbf{senza capire che cosa fanno}: ``come se apriamo una cassetta
        degli attrezzi e non sappiamo che cosa dobbiamo fare. Gli attrezzi
        sembrano tutti bellissimi, però magari li usiamo male'';
  \item il consiglio è aspettare a comprarlo, o almeno a leggerlo, fino a dopo
        la lezione sull'\textbf{emotional design}, quando si avrà già un'idea di
        come partire con il proprio gioco.
\end{itemize}
Daniele è ``assolutamente d'accordo''. (Nel frattempo si è spostato in cucina:
``benvenuti nella mia cucina, davanti al mio frigorifero, per parlare di user
experience''.) Del libro di Nuccio, e di un'alternativa in inglese, si
riparlerà a proposito dei giochi da tavolo (sezione~\ref{a2:sec:tavolo-libri}).

% ---------------------------------------------------------------------
\section{Perché analizzare i casi studio}
\label{a2:sec:perche}

Ruggero apre l'argomento del giorno, i \textbf{casi studio}, con due domande
alla classe: \emph{perché} ci teniamo ad analizzarli? E prima ancora: che cosa
ci si aspetta che siano? ``Non è una domanda provocatoria''.

\paragraph{Le risposte degli studenti.} Una prima risposta: sono \textbf{altri
giochi già realizzati}, da analizzare da molti punti di vista, ``anche
insoliti'', non solo da giocatori: dal punto di vista della produzione, o
dell'esperienza utente che se ne è avuta, per capire come migliorarli. Una
seconda risposta: i casi studio sono delle \textbf{milestone}, giochi o
situazioni ``lapidarie'' da cui trarre ``figure di merito'' da sviscerare per
poi generare altre cose. Tiziana chiede invece una cosa ``più basica'': che cosa
si intende esattamente? Modelli, esempi, ma di che cosa?

\paragraph{La definizione.} ``Ottima domanda''. Con \textbf{caso studio} si
intende un \textbf{gioco specifico}, o una \textbf{categoria di giochi}, che,
analizzato, ci dà delle \textbf{informazioni utili}: osservando da vicino che
cosa un gioco fa, possiamo estrapolare informazioni da riproporre, o metodologie
che ci servono per creare i nostri giochi. L'analogia è con la scrittura: chi
vuole scrivere un libro legge tanti libri, e capisce che cosa fanno di positivo
quelli che gli piacciono e che cosa non fanno quelli che non gli piacciono, per
evitare le cose che non gli piacciono e ``iterare sulle cose che ci piacciono''.
Con i giochi è lo stesso: si guarda che cosa fanno di positivo e di negativo,
\textbf{tenendo conto anche del contesto} in cui sono nati.

\paragraph{Oggi la panoramica, poi gli strumenti.} Gli strumenti più tecnici
per analizzare nel dettaglio arriveranno con le lezioni di \textbf{ricerca e
sviluppo}. Oggi si fa una ricognizione di tutti i settori.

% ---------------------------------------------------------------------
\section{Il gioco analogico e le sue quattro macrocategorie}
\label{a2:sec:macro}

\subsection{Un termine ombrello}

Prima una precisazione. Nel corso, \textbf{gioco analogico} è un
\textbf{termine ombrello} che copre i giochi da tavolo (\emph{board game}), i
giochi di carte (\emph{card game}), i giochi di ruolo da tavolo, i giochi di
ruolo dal vivo, i giochi pervasivi, le cacce al tesoro, le \emph{escape room},
i \emph{murder party}: tutti quei giochi che ``generalmente non sono focalizzati
e incentrati sulla programmazione, sull'uso di un terminale, sull'uso di uno
schermo''; in breve, che \textbf{non sono videogiochi}. Esistono cose ibride,
che si vedranno, ma l'area di interesse principale resta il gioco analogico.

\subsection{Le quattro macrocategorie}

\slidefig{img/AGD/agd2-categorie}{Le macrocategorie del gioco analogico: giochi da tavolo, giochi di ruolo da tavolo, LARP e, dietro, l'ibridazione.}{slide \emph{Analog -- Case Studies}}{fig:agd2-categorie}


La lezione parte dal ``frame attuale'': quattro \textbf{macrocategorie}, a loro
volta divise in sottocategorie. Ruggero avverte subito che ``non sono separate da
muri di acciaio non comunicanti'': i casi ibridi, difficili da mettere in una
sola categoria, sono frequenti. I tre docenti presentano ciascuno la propria.

\begin{description}
  \item[Gioco da tavolo (Ruggero).] La categoria dei \emph{board game}. ``Per
        motivi di industria'' vi si includono anche i giochi di carte e i
        \emph{party game}: tutti i giochi \textbf{pubblicati e prodotti
        all'interno dell'industria del gioco da tavolo}, che abbiano o no un
        tabellone.
  \item[Gioco di ruolo da tavolo (Daniele).] GDR in italiano, TTRPG
        (\emph{TableTop Role-Playing Game}) in inglese. È una categoria ``molto
        particolare'', perché anzitutto \textbf{non c'è una definizione
        univoca}: sul mercato si trovano cose molto diverse, perfino giochi che
        sono quasi ``dei LARP da tavolo'', vicini ai \emph{chamber LARP}. Anche
        i \textbf{formati} sono diversissimi: la slide mostra \emph{Dungeons \&
        Dragons}; \emph{Neverland}, un supplemento\footnote{Daniele lo chiama
        ``un'espansione per D\&D''; più avanti lo descrive come una
        reinterpretazione di Peter Pan in chiave OSR, compatibile con le regole
        di D\&D.}; \emph{Dune}, in volumi cartonati ``importanti''; ma anche
        \emph{The Quiet Year}\footnote{La trascrizione dice ``Quietier'' e, più
        avanti, ``The Quiet Ear''.}, che sta in un sacchettino; e \emph{Lady
        Blackbird}\footnote{La trascrizione dice ``Lady Blackboard''. \emph{Lady
        Blackbird} è un gioco di ruolo distribuito gratuitamente in PDF.}, che è
        solo digitale: ``tu lo prendi, lo stampi in bianco e nero e te lo
        giochi''. L'edizione italiana di \emph{Montsegur 1244}\footnote{La
        trascrizione dice ``Monsegur''.} è una sorta di scatola che si apre;
        \emph{CBR+PNK}\footnote{La trascrizione dice ``CBR PNK''.} ``si apre in
        un modo stranissimo'' e contiene dei pieghevoli (\emph{leaflet}); ci sono
        giochi di ruolo in formato brochure. Daniele li mostrerà meglio nella
        lezione sulla \emph{fiction}. La conclusione: nei giochi di ruolo
        \textbf{non si può usare il metro del formato}.
  \item[Gioco di ruolo dal vivo (Ruggero).] Il LARP, \emph{Live Action
        Role-Playing game}: ``anche questa è una categoria gigantesca''.
  \item[Gioco pervasivo (Beatrice).] La categoria meno conosciuta in Italia, e
        anche per questo inserita nel corso.
\end{description}

\paragraph{Che cosa distingue il gioco di ruolo.} Daniele chiede alla classe
che cosa distingua i \emph{role-playing game}. La risposta è nel nome: il
\textbf{role playing}, cioè un \textbf{elemento di interpretazione}, che non è
detto sia interpretazione di un \emph{personaggio} (in \emph{The Quiet Year},
per esempio, non c'è un personaggio, ``abbiamo qualcosa di diverso''). E c'è un
elemento ``importantissimo'': la \textbf{fiction}, cioè la narrazione. Non la
narrazione intesa come ``scrivere una trama e giocarla'', ma quella che nasce al
tavolo, la \textbf{narrazione emergente}: non si è previsto che l'avventura
andasse in quel modo, ma mentre si gioca gli altri dicono cose che la portano in
una certa direzione. Se ne parlerà nella lezione sulla fiction emergente.
Categorizzare i GDR è ``veramente molto arduo'', e anche il pubblico è
particolare: ``ognuno ha il suo palato'', c'è chi gioca tutto e chi gioca solo un
sottogenere. A una domanda dalla classe (il gioco \emph{Lovecraftesque}, citato
da qualcuno, rientra tra i giochi di ruolo da tavolo?) la risposta è sì: e di
che \emph{tipo} di gioco di ruolo sia si capirà con le sottocategorie che Daniele
presenterà più tardi.

\paragraph{Perché il LARP è separato dal GDR.} Ruggero spiega che i LARP sono
separati dal GDR da tavolo ``semplicemente per motivi di settore e di
industria''. Il gioco di ruolo da tavolo viene venduto con una produzione più
vicina al gioco da tavolo: si produce un \textbf{oggetto fisico} che poi si
vende. Il LARP invece è ancora spesso organizzato \textbf{a eventi}: si partecipa
\textbf{comprando un biglietto}, e si va a giocare in un posto specifico, a
un'ora specifica, con persone che sono lì per giocare quel gioco. Da questo
punto di vista è più vicino ad alcuni giochi pervasivi. Le sottocategorie sono
``una valanga'': la più famosa è il \textbf{boffer} (il cliché del LARP: persone
in costume medievale ``che si picchiano con spade di plastazote o di gommapiuma''),
ma ce ne sono di molto più astratte e di molto più accessibili, che si possono
organizzare anche in casa.

\subsection{I giochi pervasivi: una prima definizione}

Beatrice chiede se qualcuno abbia mai sentito parlare di \emph{pervasive games}.
Quasi nessuno; una studentessa ne ha sentito parlare solo leggendo la
descrizione del corso, e si è incuriosita. ``Ci sta tantissimo'': in Italia non
sono ancora molto conosciuti, ed è uno dei motivi per cui sono nel percorso.

Tiziana chiede se le \textbf{cacce al tesoro} organizzate nelle città rientrino
negli \emph{urban game}. ``Assolutamente sì'': gli \textbf{urban game} sono una
categoria dei giochi pervasivi, e nella slide si vedono proprio come un
sottoinsieme del ``grande cappello''. Beatrice dà due definizioni:
\begin{enumerate}
  \item la definizione ``standard'', dai \emph{game studies}: un gioco
        pervasivo è un gioco che \textbf{espande il proprio cerchio magico
        spazialmente, temporalmente e socialmente};
  \item la definizione più intuitiva, da cui partire per capire poi meglio la
        prima: un gioco pervasivo è un gioco che \textbf{utilizza i luoghi e gli
        oggetti della vita quotidiana e li trasforma in luoghi e oggetti di
        gioco}.
\end{enumerate}
Se la prima non è chiara, ``non è un problema, la capiremo con i casi studio''.
Negli urban game ``non giochiamo seduti al tavolo, ma giochiamo in città, e
stiamo comunque giocando''. La parola \emph{pervasivo} indica la pervasività del
gioco, che entra in situazioni in cui non siamo abituati a immaginarlo, anche se
magari siamo abituatissimi a viverlo senza pensarci.

\paragraph{Le cene con delitto.} Tiziana fa un'altra domanda (Beatrice ricorda
che si può sempre interrompere, anche con la mano alzata delle reazioni di
Zoom): le \textbf{cene con delitto} rientrano tra i pervasivi? E il
\textbf{teatro interattivo}, visto che in certe cene con delitto chi le propone
recita una vera rappresentazione? Per Beatrice, ``a mio dire, sì, entrambe''.
Ma le cene con delitto, o \emph{murder party}, esistono in molte versioni:
\begin{itemize}
  \item in una, gli \textbf{attori} vengono al tavolo con te, magari
        insospettabili, e tu, come persona giocante, li interroghi: questa
        versione assomiglia di più al teatro interattivo, che è incluso nei
        giochi pervasivi;
  \item in un'altra, \textbf{ognuno riceve un ruolo}, con una sorta di copione
        (``stasera sei la sorella della vittima''): questa va più nella
        direzione dei \textbf{giochi di ruolo dal vivo}.
\end{itemize}
Sono distinzioni ``che facciamo per allenarci a comprendere meglio'', ma che
spesso si sovrappongono: anche quando si è investigatori e si interrogano degli
attori, si è dentro una dinamica di gioco di ruolo, ``anche se magari non come
la immaginiamo in D\&D o in giochi più blasonati''. Le definizioni ``sono
malleabili'': si imparano le parole e poi si vede come usarle nel caso
specifico.

\subsection{Ibridazione e sperimentazione}

Nell'angolo della slide, dietro tutte le categorie, c'è ``una piccola cosina'':
\textbf{ibridazione e sperimentazione}. ``In realtà le categorie non esistono'',
dice Ruggero: esistono come divisioni tra industrie, ma escono continuamente
giochi che fanno da ponte. La slide mostra alcuni esempi:
\begin{itemize}
  \item \emph{Fateforge}\footnote{La trascrizione dice ``FitForge'' e, più
        avanti, ``Fateforge''.}, un gioco da tavolo che usa un'\textbf{app}
        digitale integrata nel gioco: ``un po' videogioco, un po' board game'';
  \item un gioco di aste che usa anch'esso un cellulare (il titolo non è
        detto);
  \item dei \emph{mock trial} (processi simulati) ibridi, a metà tra il murder
        party e il roleplay, ``un po' roleplay, ma anche un po' giochi
        pervasivi'';
  \item \emph{Alice is Missing}, un gioco di ruolo da tavolo che usa una
        \textbf{chat} come Telegram tra le meccaniche di gioco.
\end{itemize}
Le ibridazioni rendono spesso difficile mettere un gioco ``al 100\%'' in uno di
questi ``cesti''.

\begin{figure}[htbp]
\centering
\begin{tikzpicture}[
    font=\small,
    macro/.style={draw, thick, rounded corners=4pt, fill=gray!12,
                  text width=5.2cm, align=left, inner sep=6pt},
    titolo/.style={font=\bfseries},
    ponte/.style={draw, dashed, rounded corners=3pt, fill=white,
                  text width=3.4cm, align=center, font=\footnotesize,
                  inner sep=4pt}
  ]
  \node[macro] (tav) at (0,0) {{\bfseries Gioco da tavolo} \hfill
    \footnotesize Ruggero\\[2pt]
    \footnotesize board game, card game, party game: tutto ciò che è
    prodotto dall'industria del gioco da tavolo};
  \node[macro] (gdr) at (6.6,0) {{\bfseries Gioco di ruolo da tavolo}
    \hfill \footnotesize Daniele\\[2pt]
    \footnotesize GDR / TTRPG: interpretazione e fiction emergente; nessuna
    definizione univoca, formati diversissimi};
  \node[macro] (larp) at (0,-3.6) {{\bfseries Gioco di ruolo dal vivo}
    \hfill \footnotesize Ruggero\\[2pt]
    \footnotesize LARP: spesso a evento, con biglietto; dal boffer al
    black box};
  \node[macro] (perv) at (6.6,-3.6) {{\bfseries Gioco pervasivo}
    \hfill \footnotesize Beatrice\\[2pt]
    \footnotesize luoghi e oggetti quotidiani trasformati in luoghi e oggetti
    di gioco; urban game, ARG, teatro interattivo};
  \node[ponte] (ibr) at (3.3,-1.8) {\textbf{ibridazione e
    sperimentazione}\\ app, chat, murder party, mock trial, eventi\ldots};
  \draw[dashed, thick] (ibr) -- (tav);
  \draw[dashed, thick] (ibr) -- (gdr);
  \draw[dashed, thick] (ibr) -- (larp);
  \draw[dashed, thick] (ibr) -- (perv);
  \draw[{Stealth[length=2mm]}-{Stealth[length=2mm]}, gray, thick]
    (gdr.south) -- node[right, font=\footnotesize\itshape, text=black,
    text width=2.4cm] {oggetto venduto vs evento con biglietto} (perv.north);
  \draw[{Stealth[length=2mm]}-{Stealth[length=2mm]}, gray, thick]
    (tav.south) -- node[left, font=\footnotesize\itshape, text=black,
    text width=2.4cm, align=right] {separati per motivi di industria}
    (larp.north);
\end{tikzpicture}
\caption{Le quattro macrocategorie del gioco analogico usate nella lezione, con
il docente che le presenta. Al centro, la ``piccola cosina'' della slide:
ibridazione e sperimentazione, che rendono i confini porosi. Le frecce grigie
riassumono due osservazioni dei docenti: il GDR da tavolo e il LARP sono
separati soprattutto per ragioni di industria (oggetto venduto o evento), e per
lo stesso motivo il LARP è vicino a molti giochi pervasivi.}
\label{a2:fig:macro}
\end{figure}

% ---------------------------------------------------------------------
\section{Non chiudersi nel proprio genere}
\label{a2:sec:aperti}

Prima di aprire la ``diga dei casi studio'', i docenti insistono su alcuni punti
di metodo.

\paragraph{Le categorie servono a ragionare, non a recintare.} Ruggero: le
categorie non vengono date ``per dirvi: le cose sono A, B, C e basta'', ma per
far ragionare su \textbf{quanto sono flessibili i confini}, e per dare esempi di
cose \textbf{adiacenti} a quelle che interessano, che ``possono sbloccarvi delle
cose''. Ognuno è lì per interessi diversi (i giochi di ruolo, le carte, i giochi
formativi e didattici\ldots), ma l'invito è a \textbf{non mettersi i paraocchi}
sul proprio genere: ``è grazie alla \textbf{cross-modalità} e al vedere qualcosa
di risolto in un altro gioco che vi arriva poi un'idea per un gioco vostro che
non c'entra apparentemente niente'', magari nemmeno meccanicamente, con l'altro.

\paragraph{Un corso agnostico di genere.} Daniele aggiunge che il corso è
\textbf{agnostico di genere}: alcune lezioni saranno più importanti per una
categoria, ma i \textbf{concetti di base del design si applicano a qualsiasi
sviluppo ludico}. È fondamentale, perché altrimenti sembra che chi fa un certo
tipo di prodotto abbia bisogno di ``un game design particolare'', e non è così.
Anzi: quando è successo, storicamente anche in Italia, si è finito per
\textbf{ostracizzare e separare le community}, mentre i docenti vogliono
\textbf{metterle insieme}. ``Voi non sapete quante volte i videogiochi vengono
ispirati dai giochi da tavolo, e viceversa''.

\paragraph{Il project work è libero.} C'è anche un motivo pratico. Il progetto
di gruppo finale è \textbf{libero}: i docenti non assegneranno una categoria di
gioco, saranno i gruppi a sceglierla, e può essere anche
un'\textbf{ibridazione}, una \textbf{sperimentazione}, un gioco che non si
colloca perfettamente in un genere (``anzi, sarebbe fantastico''). Ma i gruppi
saranno formati da persone con interessi diversi, ed è proprio per questo che
bisogna restare aperti. L'unica cosa che i docenti chiederanno di non fare è
\textbf{un gioco enorme}: il progetto deve essere \textbf{calibrato}.

\paragraph{``Per me, questo è design''.} Infine, Ruggero chiede di \textbf{non
sottovalutare} la lezione: ad alcuni può sembrare ``meramente un elenco di cose
che vi facciamo vedere per nessun motivo'', ma ``per me questo è design''.
Quando deve trovare una soluzione, va a cercare altri giochi; e gli è capitato
moltissime volte di riaprire proprio quelle slide per ricordarsi com'era fatto un
gioco. E come hanno fatto Tiziana e Rita, bisogna \textbf{fare domande}:
``fermateci, chiedeteci'', i docenti sono lì per approfondire.

% ---------------------------------------------------------------------
\section{I giochi pervasivi: la premessa di Beatrice}
\label{a2:sec:pervasivi}

Si comincia ``in scivolata'' con i giochi pervasivi. Beatrice anticipa che
diversi casi studio hanno \textbf{accenni di digitale}, per tre motivi:
\begin{enumerate}
  \item alcuni casi sono \textbf{molto famosi} e aiutano a capire il concetto di
        pervasive game, anche se non sono prettamente analogici;
  \item un gioco analogico \textbf{non esclude completamente la componente
        digitale}: basti pensare ai giochi da tavolo che usano un'app;
  \item l'\textbf{ibridazione}: si vedranno casi studio ``molto strani''.
\end{enumerate}
Aggiunge una considerazione personale e professionale: quello dei pervasivi è
un campo in cui c'è ancora \textbf{ampio spazio}, perché sono poco conosciuti e
la loro \textbf{definizione è ancora in costruzione}. C'è spazio ``per mettere
anche la vostra voce'', per esprimersi e sperimentare, proprio perché le
definizioni non sono state date in modo rigido come altrove. E anticipa un tema
ricorrente del corso: ``parleremo molto di \textbf{emozioni}, anche di come
accompagnarle all'interno di un'esperienza di gioco''.

\paragraph{Approfondimento: il cerchio magico.} La definizione ``da game
studies'' presuppone il concetto di \textbf{cerchio magico}, che a lezione non
viene spiegato. In \emph{The Art of Game Design}, Schell lo introduce
discutendo le definizioni di gioco: quasi tutte implicano che un gioco sia un
sistema \emph{chiuso}, con un confine, e ``Johan Huizinga lo ha chiamato il
\emph{cerchio magico}'', che ha davvero qualcosa di magico: quando siamo
mentalmente ``dentro il gioco'' abbiamo pensieri, sentimenti e valori diversi da
quando ne siamo fuori. Schell lo reinterpreta come il nostro \textbf{sistema
interno di risoluzione dei problemi}: una ``microrealtà'' semplificata in cui gli
elementi acquistano un'importanza che non devono avere fuori dal gioco. Nel
capitolo sulle meccaniche, lo stesso Schell identifica il cerchio magico con lo
\textbf{spazio} del gioco, la prima delle sue sette categorie di meccaniche. Il
gioco pervasivo è dunque quello che sposta, allarga o rende sfumato questo
confine \textbf{nello spazio} (si gioca nella città, non sul tavolo), \textbf{nel
tempo} (non si sa bene quando la partita inizia e finisce) e \textbf{nella
società} (vi entrano persone che non sanno di giocare, o che giocano senza
chiamarlo gioco)\footnote{La formula ``espandere il cerchio magico
spazialmente, temporalmente e socialmente'' è quella proposta da Markus Montola
nei suoi studi sui giochi pervasivi, poi ripresa nel volume \emph{Pervasive
Games: Theory and Design} di Markus Montola, Jaakko Stenros e Annika Waern
(2009). Beatrice non cita la fonte; la nota è un'integrazione.}. È esattamente
quello che Beatrice dirà di \emph{GuruShot}: ``non sappiamo bene dove inizia e
dove finisce l'attività di gioco, ed è per questo che viene utilizzata la parola
pervasivo''.

\begin{figure}[htbp]
\centering
\begin{tikzpicture}[
    font=\small,
    freccia/.style={-{Stealth[length=2.5mm]}, thick},
    etichetta/.style={font=\footnotesize, align=center, text width=3.6cm}
  ]
  % cerchio classico
  \draw[thick, fill=gray!15] (0,0) circle (1.0cm);
  \node[font=\footnotesize, align=center] at (0,0) {cerchio\\magico};
  \node[font=\footnotesize\itshape, align=center] at (0,-1.45)
    {gioco ``classico'':\\tavolo, durata, giocatori};
  % cerchio espanso
  \draw[thick, dashed, fill=gray!6] (7,0) circle (2.1cm);
  \draw[thick, fill=gray!15] (7,0) circle (1.0cm);
  \node[font=\footnotesize, align=center] at (7,0) {cerchio\\magico};
  \draw[freccia] (1.3,0) -- node[above, font=\footnotesize]{gioco pervasivo}
    (4.6,0);
  \draw[freccia, gray] (7,1.0) -- (7,2.1);
  \draw[freccia, gray] (5.95,-0.6) -- (5.18,-1.05);
  \draw[freccia, gray] (8.05,-0.6) -- (8.82,-1.05);
  \node[etichetta, above=0.05cm] at (7,2.15) {\textbf{spazialmente}: la città,
    l'hotel, la stazione dell'autobus};
  \node[etichetta, anchor=north east] at (5.4,-1.5) {\textbf{temporalmente}:
    quando inizia e quando finisce la partita?};
  \node[etichetta, anchor=north west] at (8.6,-1.5) {\textbf{socialmente}:
    passanti, pasticcieri, vicini di quartiere};
\end{tikzpicture}
\caption{La definizione ``da game studies'' citata da Beatrice: il gioco
pervasivo espande il proprio cerchio magico nello spazio, nel tempo e nella
società. La definizione più intuitiva dice la stessa cosa dal lato degli
oggetti: luoghi e oggetti della vita quotidiana diventano luoghi e oggetti di
gioco.}
\label{a2:fig:cerchio}
\end{figure}

% ---------------------------------------------------------------------
\section{Pervasivi con un'app: Pokémon Go, geocaching, GuruShot}
\label{a2:sec:pervasivi-app}

\subsection{Pokémon Go}

\slidefig{img/AGD/agd2-pokemon}{Pervasivi con un'app: Pokémon Go, il geocaching e GuruShots.}{slide \emph{Analog -- Case Studies}}{fig:agd2-pokemon}


Il primo esempio è scelto ``proprio per farvi comprendere'': \textbf{Pokémon
Go}, che conoscono tutti, o di cui almeno tutti hanno sentito il nome. Beatrice
lo spiega comunque, ``così ci allineiamo'' (scusandosi con i grandi fan per la
descrizione imprecisa). È un gioco in formato \textbf{app}. I Pokémon, ``questi
animaletti bizzarri'', nella narrativa dei Pokémon popolano il loro mondo; in
\emph{Pokémon Go} popolano \textbf{il nostro mondo reale}. Il gioco consiste nel
cercarli, catturarli, andare nelle palestre, e tutto questo si fa
\textbf{camminando nel mondo}: la stessa via in cui si va a fare colazione al
bar o che si percorre per andare in ufficio è la via in cui si può trovare un
Pokémon o una palestra. Ed è ``tutto costruito dagli utenti''.

Su \emph{Pokémon Go} ci sarebbe molto da dire, ``anche su come è stato
utilizzato per fare altre cose meno belle ed eleganti dell'esperienza di
gioco'', ma Beatrice non ne parla per ragioni di tempo: ci tiene solo a far
notare che ``tante cose possono avere tante sfumature'', e che i casi più
spinosi si possono discutere a parte. Qui serve per una ragione semplice:
\textbf{è pervasivo perché usa gli stessi spazi della vita quotidiana come spazi
di gioco}.

\subsection{Il geocaching}

Il \textbf{geocaching} è anch'esso abbastanza conosciuto. ``Di base è una
gigantesca caccia al tesoro'', tenuta in piedi da alcuni portali, e nata
proprio con il \textbf{GPS} e la geolocalizzazione. Il principio:
\begin{itemize}
  \item le persone giocanti \textbf{nascondono} delle scatole, piccole o
        grandi, con contenuti vari (dei ``tesori'') e con una sorta di
        \textbf{diario di bordo} (un \emph{log}: fogli o quaderni) su cui chi le
        trova può scrivere ``sono stato qui, ho trovato questa scatoletta'';
  \item la scatola viene nascosta nel posto che l'utente ritiene più idoneo, e
        sulla piattaforma l'utente dà le \textbf{indicazioni} per trovarla;
  \item la grande regola è che \textbf{le scatole non vengono mai tolte}: chi le
        trova può scambiare gli oggetti al loro interno (piccoli, di qualunque
        valore, con piena libertà) e segnare la propria presenza.
\end{itemize}
Non è ``un arraffare i tesori'', ma un modo di \textbf{viverli e dar loro
spazio collettivamente}. In Italia, ``per la cultura che abbiamo'', funziona
così così, perché le scatole spesso vengono rubate o spostate. Ma vale la pena
guardare online i video dei modi in cui vengono nascoste: alcune sono ``piccole
opere di ingegneria artistica'', quasi piccole escape room, con un gioco di pesi
o una chiave ``nascosta in bella vista'' che si vede solo con grande
attenzione; alcune non sono scatole ma \textbf{cabine telefoniche} in cui si
entra e succedono cose, o intere automobili.

\subsection{GuruShot}

\textbf{GuruShot}\footnote{Più avanti la trascrizione lo chiama anche ``blue
shot''.} è di nuovo in formato app, ma Beatrice lo inserisce perché usa uno
strumento che lei usa ``moltissimo'' e che trova ``molto potente per fare molte
cose'': \textbf{la fotografia}. È un gioco dedicato a persone molto dentro al
tema, fotografi anche di un certo livello di passione e di competenza: una
\textbf{gara} o maratona fotografica in cui vengono date delle \textbf{challenge
fotografiche}. Chi vuole giocare apre l'app, sceglie la sfida, la svolge, carica
la foto; e poi ``succedono cose'': valutazioni, momenti di incontro con la
community, premi veri e propri offerti dagli \textbf{sponsor}. È un'attività
\textbf{continua}, online da molti anni (``almeno 7--8'', corregge Beatrice), con
una nicchia internazionale molto nutrita.

Ed è ``molto pervasivo'': \textbf{non si può dire quando una partita è iniziata
o finita}, né quando si sta giocando e quando no. Rispetto alla definizione del
cerchio magico, non si sa bene dove inizi e dove finisca l'attività di gioco:
per questo si parla di pervasività.

% ---------------------------------------------------------------------
\section{La città come campo di gioco: Journey, SF0 e Critical City}
\label{a2:sec:pervasivi-citta}

\subsection{Journey to the End of the Night}

\slidefig{img/AGD/agd2-journey}{La città come campo di gioco: Journey to the End of the Night e SF0.}{slide \emph{Analog -- Case Studies}}{fig:agd2-journey}


Si torna sotto il cappello dei \textbf{giochi urbani}, con un gioco che si
gioca in città, sempre \textbf{in formato evento}: un'ora di inizio, un'ora di
fine e una zona di gioco all'interno della città. È ``una sorta di \textbf{guardie
e ladri gigante}''. Lo scopo dei creatori era \textbf{riappropriarsi della vita
notturna} e viverla in modo ludico, esplorando la città ``in maniera
\emph{safe} da un lato e un po' \emph{noir} dall'altro''. Un piccolo gruppo di
giocatori fa i \textbf{chaser} (gli inseguitori, le guardie) e gli altri i
\textbf{runner} (i ladri); ci sono delle \textbf{tappe} da raggiungere, e il
gioco dura svariate ore. Non è un gioco prettamente sportivo (ci sono altri
pervasivi molto più sportivi): ci sono molte dinamiche \textbf{strategiche} su
come arrivare alla fine del percorso. Beatrice lo inserisce perché, rispetto ai
casi precedenti, è \textbf{completamente analogico}: non usa nessuna app, nulla.
È un evento.

Il gioco è \textbf{Journey to the End of the Night}, nato in America da un
collettivo con base a \textbf{San Francisco}. È arrivato in Italia nel
\textbf{2008}, durante un'ondata di giochi pervasivi arrivata ``in modo molto
forte'', e poi in Giappone e in tantissimi altri posti, perché i creatori lo
hanno reso \textbf{libero}: chi vuole può organizzarlo, purché \textbf{non metta
un biglietto a pagamento}. Ha avuto per anni un grande seguito: cercando
``Journey to the End'' Italia si trovano tantissimi eventi, e a \textbf{Milano},
tra il 2008 e il 2013, se ne sono fatte svariate edizioni con
\textbf{2\,000--3\,000 partecipanti} ciascuna. Beatrice lo racconta ``non tanto
per le meccaniche'', che più o meno si conoscono, ma per \textbf{come è
arrivato}: un'onda lunga che dura ancora, anche se si sta affievolendo.

\subsection{SF0}

Degli stessi creatori è \textbf{SF0}\footnote{La trascrizione dice ``SF0'',
``SFZero'' e, una volta, ``essere al vero''.}, che era una sorta di grande
\textbf{piattaforma online} con lo stesso principio di \emph{GuruShot}: delle
\textbf{missioni} da svolgere e documentare con video o foto. Ma invece di
puntare sull'estetica della fotografia, le missioni puntavano sulle
\textbf{interazioni cittadine e urbane}.

\subsection{Critical City Upload}

SF0 ha avuto ricadute anche in Italia, con ``un figlio, chiamiamolo così'':
\textbf{Critical City Upload}, un gioco che è andato molto tra il 2009 e il
2013. Lo scopo è vincere, accumulando punti con le missioni, ma il gioco è
molto legato agli \textbf{scambi umani e al territorio}. Le missioni andavano da
``organizza una cena con una persona del tuo quartiere con cui non hai mai
parlato'' a una delle preferite di Beatrice: ``crea un'opera d'arte interattiva
che intrattenga le persone che aspettano l'autobus e posizionala in una delle
pensiline del quartiere o della città''.

In entrambi i casi \textbf{è la community che valuta la community}: non ci sono
giudici. E dopo i primi anni, sia SF0 sia Critical City hanno permesso ai
partecipanti che avevano raggiunto un certo livello di \textbf{creare essi
stessi nuove missioni} per i nuovi giocatori: la comunità non si limitava a
valutarsi a vicenda, ma creava il gioco per gli altri. Oggi nessuno dei due è
più attivo, ma i siti sono ancora online, con molte delle missioni visibili.
Quanto a ``stranezza'' rispetto all'idea comune di gioco analogico, ``questo
spinge molto forte''.

% ---------------------------------------------------------------------
\section{Teatro, ARG e marketing: da Sleep No More a The Institute}
\label{a2:sec:pervasivi-arg}

\subsection{Sleep No More}

\slidefig{img/AGD/agd2-sleep}{Teatro e ARG: Sleep No More e Why So Serious?}{slide \emph{Analog -- Case Studies}}{fig:agd2-sleep}


Chi conosce il \textbf{teatro interattivo} conosce di sicuro \textbf{Sleep No
More}\footnote{La trascrizione dice ``Slip no more'' e ``SlipMod''. Lo
spettacolo è della compagnia britannica Punchdrunk; la versione più nota è
quella di New York, allestita in un edificio trasformato nel fittizio McKittrick
Hotel. Il nome della compagnia non è detto a lezione.}: un'opera messa in scena
all'interno di un \textbf{hotel}. È interattiva nel senso che gli attori recitano
la loro parte in interazione con il pubblico, e il pubblico, con regole precise
sui contatti fisici e su come interagire con gli spazi, è \textbf{libero di
esplorare l'hotel} come vuole: in ogni zona ci sono attori che fanno cose. Resta
un'opera teatrale, con un inizio, una fine e degli eventi nel mezzo; c'è molta
documentazione online (interviste, spezzoni). E di nuovo si gioca su
\textbf{qual è lo spazio di gioco e quale non lo è}.

\subsection{Gli alternate reality game e Why So Serious}

Qui Beatrice introduce un nuovo termine: gli \textbf{alternate reality game}
(ARG), i giochi di realtà alternativa, ``giochi in cui c'è un intero mondo
alternativo al nostro, in cui succedono cose''. Il caso scelto, ``uno dei miei
preferiti'', porta anche nel mondo dei \textbf{giochi per il marketing e la
comunicazione}: \textbf{Why So Serious}, nato come gioco promozionale per
\emph{Il cavaliere oscuro} di Christopher Nolan.
\begin{itemize}
  \item Iniziava con una \textbf{mail} ad alcune persone iscritte a mailing
        list di fan di Batman, contattate dagli organizzatori tramite le loro
        associazioni.
  \item La mail era un \textbf{invito strano a entrare nella banda del Joker},
        con un enigma; da lì partivano attività di socializzazione organizzate
        attraverso un sito, che era una serie di indirizzi ``in cui succedevano
        cose''.
  \item Il gioco raccontava che \textbf{nel nostro mondo il Joker esisteva
        davvero}, e gli indirizzi erano indirizzi reali.
  \item Una delle prime cose: andare a questi indirizzi, tutti
        \textbf{pasticcerie}, e dire la parola in codice ricevuta. Si riceveva
        una \textbf{torta} con sopra scritto un numero di telefono; chiamandolo,
        squillava la torta stessa, perché dentro c'era un \textbf{cellulare}:
        lo strumento con cui il Joker gestiva le comunicazioni con i giocatori.
\end{itemize}
Molte strutture di questo tipo, molti colpi di scena ``anche proprio pratici'',
che sfruttano la pervasività: cose che nel nostro mondo non hanno motivo di
succedere (``non esiste che ci sia un cellulare in una torta'') ma che in quel
mondo hanno un senso. Gli ARG giocano moltissimo su ciò che \textbf{potrebbe
essere reale ma non lo è}, e che il gioco dice che lo è: ``in quel mondo voi
siete gli aiutanti del Joker''. \emph{Why So Serious} è opera di un'azienda che
si occupa da molti anni di ARG\footnote{Beatrice non fa il nome dell'azienda,
che si troverà nell'estratto distribuito. Si tratta di 42 Entertainment
(integrazione).}; i dettagli sono nell'estratto che Beatrice caricherà, tratto
da \emph{Immersi nelle storie} di \textbf{Frank Rose}\footnote{Edizione
italiana di \emph{The Art of Immersion} (2011).}, che racconta proprio l'inizio
di questo gioco.

Perché raccontarlo? Perché mostra ``come possiamo giocare sulle
\textbf{emozioni}'' e come \textbf{connettere le persone alle cose che vogliamo
comunicare}, anche costruendo un mondo alternativo. Non è un obbligo, ``ma è bene
sapere'' che esistono questi giochi, ``giganteschi''.

\subsection{Un gioco da tavolo che è anche un ARG}

Tornando alla malleabilità delle definizioni, Beatrice presenta un gioco che è
insieme \textbf{gioco da tavolo}, \textbf{alternate reality game} e gioco con
\textbf{componenti digitali}, edito da \textbf{MS Edizioni}\footnote{La
trascrizione chiama il gioco ``The Ice''; il titolo non è ricostruibile con
certezza (potrebbe trattarsi di un titolo come \emph{The Heist}). MS Edizioni
è un editore italiano di giochi; è forse la stessa ``MS'' per cui, nella prima
lezione, Ruggero dice di lavorare, ma la trascrizione non permette di
confermarlo.}. Si riceve un \textbf{fascicolo} di informazioni e poi dei
\textbf{messaggi}, da cui si capisce di far parte di un gruppo che \textbf{sta
rapinando una banca}, già nel pieno della rapina. I giocatori sono ``l'entità che
ha accesso a internet'', la voce fuori dall'azione che coordina: parlano con dei
personaggi, scambiano informazioni, guardano video, analizzano dati e
documenti. È un gioco ``con una profondità notevole''. Due cose lo rendono
interessante:
\begin{itemize}
  \item mescola cose che \textbf{esistono davvero} (la cartina di una città
        reale, gli scambi della metropolitana di una città reale) con cose del
        \textbf{mondo narrativo} (i personaggi, i fatti, la rapina: ``spero che
        sia chiaro che non state rapinando la banca davvero'');
  \item a differenza degli altri casi, ha un \textbf{inizio e una fine ben
        definiti} ed è \textbf{completamente accessibile}: lo si può provare.
\end{itemize}
Per Beatrice è un caso ``unico sul panorama italiano'' (``poi mi correggeranno i
colleghi a livello internazionale'').

\subsection{Collettivo Cinetico: Cinetico 4.4 e Amleto}

\slidefig{img/AGD/agd2-teatro}{Cinetico 4.4, Lettere anonime per un camminatore, The MP3 Experiment, The Institute e l'ARG di Stranger Things.}{slide \emph{Analog -- Case Studies}}{fig:agd2-teatro}


\textbf{Cinetico 4.4} è ``un altro bell'esperimento'': un \textbf{gioco da
tavolo creato da un collettivo di danzatori e danzatrici}, il \textbf{Collettivo
Cinetico}, per generare una \textbf{performance artistica dal vivo}, teatrale o
di danza. Non è un gioco da prendere dallo scaffale di una ludoteca: è pensato
per una performance. Chi ci gioca partecipa a un gioco da tavolo che in realtà è
un \textbf{evento}: si paga un biglietto, si entra, ci si siede, succedono cose.
Il collettivo, di cui Beatrice si dichiara ``di parte'', fa moltissime attività
``né danza, né teatro, né gioco, o meglio tutte e tre le cose insieme'', difficili
da mettere sotto un cappello. E qui Beatrice aggancia un tema delle prossime
lezioni: la \textbf{transmedialità} e la \textbf{ricerca e sviluppo}, cioè
attingere anche a cose lontane dal gioco che si vuole creare, perché ``più
conosciamo, più possiamo aprirci la testa''.

Un caso studio in più: l'\textbf{Amleto} del Collettivo Cinetico. È l'Amleto di
Shakespeare, ma sul palco non ci sono i danzatori del collettivo. Tre giorni
prima dello spettacolo fanno una vera e propria \textbf{audizione} aperta alla
popolazione; tra chi si presenta scelgono \textbf{sei potenziali Amleto}, che
portano sul palco a danzare e recitare, e creano momenti di
\textbf{interazione con il pubblico} per scegliere, tra i sei, l'Amleto di
quell'edizione. Ci sono giochi veri e propri sia nell'audizione sia nello
spettacolo. Loro, da danzatori, lo chiamano spettacolo di danza, ``ma se lo
vedete c'è tantissimo''.

\subsection{Lettere anonime per un camminatore}

``Rispetto a quello che vi ho detto, questo è semplice''. \textbf{Lettere
anonime per un camminatore} è del collettivo \textbf{Amigdala}, anche loro una
compagnia teatrale (Gianmarco, che ha vissuto a Modena, conferma che sono di
lì). È una \textbf{passeggiata} con piccole interazioni ludiche, tutta basata su
un \textbf{audio}:
\begin{itemize}
  \item è un gioco \textbf{in solitario}: una persona sola segue un percorso
        segnato da frecce, per terra o sui muri di una città, guidata da una
        \textbf{voce in cuffia} che le racconta delle cose;
  \item la struttura del racconto è fissa, ma la \textbf{parte centrale} cambia
        ogni volta: prima di ogni edizione il collettivo \textbf{intervista le
        persone che hanno abitato o abitano il percorso}, e mentre si cammina si
        ascolta, oltre alla narrazione, il racconto di chi quei luoghi li vive o
        li ha vissuti;
  \item la \textbf{conclusione} è sempre ``col botto'': il collettivo sceglie di
        finire in modo scenografico in luoghi normalmente inaccessibili.
        Beatrice lo ha giocato a Ferrara, e alla fine del percorso si è
        ritrovata seduta in poltrona, con un tavolino e un drink, al
        \textbf{centro dello stadio comunale}, con lo stadio vuoto e le cuffie
        nelle orecchie.
\end{itemize}
Torna il tema del \textbf{percorso emotivo} che si fa fare alla persona giocante.
E lo strumento dell'\textbf{audio che guida} è interessante anche per il corso:
un'app non si riesce a fare nel tempo del project work, ``ma per utilizzare
l'audio invece possiamo''.

\subsection{MP3 Experiment}

Un altro esempio basato sull'audio, più simile a un \textbf{flash mob}:
\textbf{MP3 Experiment}, del collettivo newyorkese \textbf{Improv
Everywhere}\footnote{La trascrizione dice ``Improve Everywhere''.}, che fa
moltissimi esperimenti di questo tipo. Ci si ritrova tutti nel giorno, nell'ora
e nel posto concordati, ci si collega con le cuffie a una trasmissione in
diretta, e la voce dice che cosa fare. Beatrice non svela di più (``guardatevi i
video su YouTube, ci sono delle perle''), ma racconta un dettaglio: gli
organizzatori \textbf{elencano gli oggetti da portare}, e in un'edizione chiesero
di portare \textbf{due matite}; a un certo punto tutto il gruppo si mise a
suonare la batteria con le matite sugli oggetti della città. Quella
\textbf{scintilla di curiosità} per oggetti di cui non si sa ancora l'uso è uno
strumento del percorso emotivo, ``non difficile'' da usare: ``basta solo
ragionare di creatività''. Come le istruzioni audio.

\subsection{The Institute e Dispatches from Elsewhere}

\textbf{The Institute}\footnote{Il gioco è noto anche come \emph{The Jejune
Institute}, di San Francisco (integrazione).} è, nel modo in cui Beatrice ha
imparato questo mondo, uno dei casi più blasonati e famosi di alternate reality
game, uno dei primi raccontati come tale. È un'esperienza così complessa che
Beatrice preferisce indicare due fonti per approfondirla:
\begin{itemize}
  \item un \textbf{documentario}, intitolato anch'esso \emph{The Institute}, che
        racconta tutto l'ARG, durato circa un anno e fatto di molti passaggi, con
        tantissimi colpi di scena del tipo ``torta con il cellulare'';
  \item una \textbf{serie TV}, \emph{Dispatches from Elsewhere} (Daniele
        conferma il titolo), che ha preso l'esperienza di The Institute, ci ha
        aggiunto una storia dei personaggi e ne ha fatto una serie. Tolta la
        storia dei personaggi (``comunque molto carina''), resta più o meno
        l'idea dell'ARG.
\end{itemize}

\paragraph{Con budget zero.} Qui Beatrice fa una precisazione importante. Questi
esempi sono \textbf{difficili da replicare}: ``con che soldi compriamo la torta e
ci mettiamo il cellulare? Con che soldi affittiamo un intero palazzo come in The
Institute?''. Serve un intero lato commerciale, di ricerca di sponsor, che non è
parte del design del gioco e quindi non è il motivo per cui si è nel corso.
Quello che vuole che gli studenti si portino a casa è che \textbf{con molto meno}
si possono ottenere esperienze ``estremamente significative a livello
emozionale'' e creazioni di game design molto profonde \textbf{con budget zero}.
I casi grandi servono come \textbf{punti di riferimento} e per sapere che
esistono cose ``molto distanti da quelle che immaginiamo tutti i giorni''.

\subsection{L'ARG di Stranger Things a Milano}

Nello stesso filone, un ARG che ``è stato un'operazione di marketing e basta''
(``per me sono molto affascinanti, però era un'operazione di marketing''), ma
fatto \textbf{a Milano}: quello per una stagione di \emph{Stranger Things} (la
terza o la quarta, Beatrice non ricorda). Era una sorta di \textbf{caccia al
tesoro} in città, con punti di gioco preparati in modo ``altamente scenografico''
(nella slide si vede un autobus), mescolata con una \textbf{telefonata} a un
numero in cui uno dei personaggi diceva che il Sottosopra era arrivato a Milano e
che serviva aiuto. Il messaggio: anche da noi questi giochi ``sono arrivati,
arrivano, e credo che col tempo avranno sempre più spazio'', e per chi pensa in
termini di lavoro ``è un tema''.

% ---------------------------------------------------------------------
\section{Cacce al tesoro, corse e indagini}
\label{a2:sec:pervasivi-cacce}

\subsection{Le guide-caccia al tesoro su carta}

Beatrice presenta poi delle \textbf{guide per esplorare la città} in forma di
caccia al tesoro\footnote{La trascrizione riporta il nome come ``Waiwai''; la
grafia non è verificabile.}, uscite negli stessi anni dell'ondata dei pervasivi
in Italia (la prima nel 2007--2008), disponibili solo per alcune città e
vendute come \textbf{libretti} sul loro sito o da rivenditori autorizzati.
Tiziana le conosce: aveva comprato quelle di \textbf{Venezia} e di
\textbf{Verona}, ha giocato in parte quella di Venezia, nella versione breve
(c'erano due possibilità: un gioco di un paio d'ore o uno da una giornata
intera), e ricorda il \textbf{numero di cellulare} con cui interagire per avere
le indicazioni. Beatrice le aveva giocate con gli SMS, ``era un'altra vita'';
oggi pare si siano aggiornate a WhatsApp (e chiede aggiornamenti a chi le
giocherà).

Sono un altro ibrido di \textbf{analogico e digitale}: un'entità parla con i
giocatori tramite messaggi, mentre il libretto cartaceo resta in mano a loro. La
particolarità è la meccanica per passare da una tappa all'altra: \textbf{ogni
pagina del libro è divisa in strisce}, e quando si risolve una tappa la guida
comunica dei \textbf{codici} (``il foglietto 3, il 32C e il 18F'') con cui
ricomporre la pagina che porta alla tappa successiva.

\subsection{PacManhattan}

\slidefig{img/AGD/agd2-cacce}{PacManhattan, Zombies, Run! e Reptilarium.}{slide \emph{Analog -- Case Studies}}{fig:agd2-cacce}


\textbf{PacManhattan}\footnote{La trascrizione dice ``Pac-Man-Hatton''.}, degli
stessi anni, è \textbf{Pac-Man per le strade di Manhattan}: una persona fa
letteralmente Pac-Man, guidata da un'altra persona (``la mente'') tramite
telefono o walkie-talkie; altri giocatori fanno i fantasmi. È di nuovo ``una
sorta di guardie e ladri'', ed è uno dei primi urban game che si trovano nei
manuali di game studies. Beatrice lo mostra per un'idea precisa: si può
\textbf{prendere qualcosa che esiste già in una forma e cambiarle forma
completamente}. Non è lo stesso gioco di Pac-Man, ma si ispira alla sua
narrativa, alle sue meccaniche, alla sua estetica.

\subsection{Zombies, Run!}

\textbf{Zombies, Run!}\footnote{La trascrizione dice ``Zombie Run''.} è
un'\textbf{app per fare jogging}: invece di ascoltare musica, o i messaggi
vocali degli amici, si ascoltano \textbf{gli zombie che stanno per prenderti}, e
c'è tutta una narrativa su come si sta scappando mentre si corre. Si trova sugli
store.

\subsection{Reptilarium}

Si torna al teatro, ma in versione analogico-digitale. \textbf{Reptilarium} è un
caso studio di cui purtroppo c'è pochissima documentazione. È uscito durante il
\textbf{COVID}, da una compagnia teatrale di Roma, \textbf{Project
XX1}\footnote{Grafia come nella trascrizione.}. Era a modalità evento (un certo
giorno, a una certa ora) e le persone accedevano da \textbf{casa propria}, davanti
al computer, a un'esperienza con degli \textbf{attori} che si mettevano in
comunicazione con loro: ricevevano mail e telefonate, potevano vedere video
``recuperati da altre persone'', potevano \textbf{coordinare i movimenti degli
attori} e vedere che cosa succedeva mentre li eseguivano. La trama partiva da una
\textbf{persona scomparsa} da ritrovare, e poi si espandeva molto. Come per le
cene con delitto (dove si può interrogare l'attore o essere l'indiziato), anche
questo esperimento ``può diventare un LARP'', un ibrido tra teatro e gioco di
ruolo, o tra gioco investigativo e alternate reality game.

Ruggero aggiunge una nota: il \textbf{lockdown} ha colpito moltissimo i giochi
pervasivi, ma anche altri generi come i LARP, e in quel periodo sono nate
esperienze che hanno spostato \textbf{online e a distanza} giochi nati per essere
giocati di persona. Se ne riparlerà.

% ---------------------------------------------------------------------
\section{Pervasivi per la formazione: Echi resistenti}
\label{a2:sec:pervasivi-edu}

\slidefig{img/AGD/agd2-echi}{Un gioco pervasivo per la formazione storica.}{slide \emph{Analog -- Case Studies}}{fig:agd2-echi}


``Tutte le cose che vedete hanno enorme spazio per essere esplorate in tema
\textbf{educativo e formativo}'', e Beatrice chiude con un accenno a questo
ambito. \textbf{Pop History} è un'\textbf{associazione di promozione sociale}:
non si presenta come collettivo teatrale né come gruppo di game designer, ma come
associazione (``credo siano anche loro emiliani''), e fa tante piccole
sperimentazioni. Una di queste è \textbf{Echi resistenti}\footnote{La
trascrizione dice ``Ecchi Resistenti'' e ``Eh, che resistenti''.}, un gioco
sulla \textbf{Resistenza e sulla Liberazione}, che ibridava:
\begin{itemize}
  \item un'\textbf{indagine} con dei fascicoli;
  \item un'\textbf{attività teatrale}, con personaggi storici con cui si poteva
        interagire direttamente;
  \item una \textbf{caccia al tesoro}, in cui si andava a cercare questi
        personaggi per la città.
\end{itemize}
Altri casi simili sono sul loro sito.

\paragraph{Approfondimento: giochi con migliaia di giocatori.} In
\emph{Challenges for Game Designers}, nel capitolo sui giochi multigiocatore,
Brathwaite e Schreiber osservano che nei giochi da tavolo il problema
dell'interfaccia è evidentissimo (``come si mettono venti persone attorno a un
solo tabellone? E mille?''), eppure ``esistono giochi non digitali con migliaia
e perfino centinaia di migliaia di giocatori'': gli \textbf{alternate reality
game} e i grandi \textbf{urban game}, che ``portano i giocatori via dai loro
computer e dentro il mondo reale''. È la stessa scala dei casi di Beatrice:
\emph{Journey to the End of the Night} con migliaia di partecipanti a Milano,
\emph{Why So Serious} con le sue mailing list di fan, \emph{Pokémon Go} con il
mondo intero come mappa.

\begin{table}[htbp]
\centering
\small
\begin{tabular}{@{}p{3.6cm}p{2.3cm}p{7.0cm}@{}}
\toprule
\textbf{Caso studio} & \textbf{Formato} & \textbf{Perché è nella lezione} \\
\midrule
Pokémon Go & app & gli spazi quotidiani diventano spazi di gioco \\
Geocaching & piattaforma + scatole & caccia al tesoro collettiva: le scatole
non si tolgono mai \\
GuruShot & app & la fotografia come strumento; non si sa quando la partita
inizi e finisca \\
Journey to the End of the Night & evento urbano & guardie e ladri notturno,
completamente analogico, libero e non a pagamento \\
SF0 & piattaforma & missioni urbane documentate; la community valuta e crea
missioni \\
Critical City Upload & piattaforma & il ``figlio'' italiano di SF0
(2009--2013): missioni di relazione sul territorio \\
Sleep No More & teatro interattivo & il pubblico esplora liberamente un hotel
\\
Why So Serious & ARG di marketing & la torta con il cellulare: il Joker esiste
davvero \\
Gioco di MS Edizioni & gioco da tavolo + ARG & rapina in corso tra documenti
reali e finzione; ha inizio e fine \\
Cinetico 4.4, Amleto & performance & un gioco da tavolo che è uno spettacolo; il
pubblico sceglie Amleto \\
Lettere anonime per un camminatore & passeggiata audio & gioco in solitario,
voci degli abitanti, finale scenografico \\
MP3 Experiment & flash mob audio & la curiosità per gli oggetti da portare \\
The Institute & ARG & il caso blasonato, con documentario e serie TV \\
ARG di Stranger Things & ARG di marketing & un ARG fatto a Milano \\
Guide-caccia al tesoro & libretto + SMS & pagine in strisce da ricomporre con
codici \\
PacManhattan & urban game & un videogioco cambiato completamente di forma \\
Zombies, Run! & app & narrativa audio durante la corsa \\
Reptilarium & teatro a distanza & attori guidati da casa durante il COVID \\
Echi resistenti & gioco formativo & indagine, teatro e caccia al tesoro sulla
Resistenza \\
\bottomrule
\end{tabular}
\caption{I casi studio di giochi pervasivi presentati da Beatrice, nell'ordine
della lezione, con il motivo per cui sono stati scelti.}
\label{a2:tab:pervasivi}
\end{table}

% ---------------------------------------------------------------------
\section{I giochi di ruolo dal vivo: le macrocategorie}
\label{a2:sec:larp}

\slidefig{img/AGD/agd2-larp}{Le macrocategorie del LARP: boffer, Mind's Eye Theatre, parlor/chamber LARP, American freeform e murder party.}{slide \emph{Analog -- Case Studies}}{fig:agd2-larp}


Finita la parte sui pervasivi, si fa una pausa di pochi minuti (``facciamo
sette minuti, così andiamo ai 20 tondi''). Al rientro ci sono tutti, compreso
Tancredi (``ci sono tutti e due'' i Lorenzo), e Ruggero passa ai
\textbf{giochi di ruolo dal vivo}.

\paragraph{Una struttura diversa.} Questa parte è organizzata diversamente da
quella sui pervasivi. Il gioco di ruolo dal vivo, come il gioco di ruolo da
tavolo e il gioco da tavolo, ha alle spalle \textbf{anni di industria
solidificata}, ed è quindi più facile identificare \textbf{macrocategorie} con
nomi e formati specifici. Prima, però, Ruggero chiede chi abbia mai giocato un
LARP. Una persona della classe sì: un LARP fantasy, al \textbf{festival del gioco
di ruolo di Bracciano}, ``con le spade di gommapiuma, insomma, quelle cose lì''.
È il primo impatto di molti.

\subsection{Il boffer LARP}

Chi ha giocato un LARP probabilmente ha giocato un \textbf{boffer}, e chi ne ha
solo sentito parlare ha in mente quello. È uno dei primi giochi di ruolo dal
vivo diffusi, nato all'epoca proprio come \textbf{tentativo di adattare
\emph{Dungeons \& Dragons} dal vivo}, poi evoluto nel tempo. (Daniele mostra
una foto e Ruggero protesta: ``non evidenziarmi, che già è cringe''; ne ha
giocati anche lui.) Il nome viene dalle armi: le armi \emph{boffer} sono armi in
plastazote che devono avere caratteristiche precise per non essere pericolose.
Sono giochi \textbf{incentrati sul combattimento}, su trame epiche e così via.
Sono anche quelli che possono essere più grandi: un boffer anche grosso può
ospitare \textbf{400--500 persone}, ``per non dire quelli in Germania, che ne
tengono anche migliaia''.

\subsection{Il Mind's Eye Theatre}

Un altro tipo abbastanza diffuso è il \textbf{Mind's Eye Theatre}: se il boffer
deriva da \emph{Dungeons \& Dragons}, questo deriva da \textbf{\emph{Vampire: The
Masquerade}}. Le trame sono più dark, gotiche, horror, legate alla proprietà
intellettuale di Vampire, e chi partecipa di solito interpreta uno dei vampiri
delle loro società segrete. Sono giochi \textbf{più politici che di
combattimento}: l'intrigo, i segreti, il tradirsi a vicenda, ``il pugnalare
alle spalle qualcuno'' sono parte delle dinamiche. Come i boffer, sono giochi
\textbf{in costume}: il costume è richiesto, e si chiede a chi gioca di vestirsi
come il proprio personaggio.

\subsection{Chamber LARP e parlor LARP}

Il costume sembra scontato, ``ma in realtà non tutti i LARP usano il costume
come strumento necessario''. I \textbf{parlor LARP} o \textbf{chamber LARP}
sono giochi di ruolo dal vivo in cui il costume \textbf{non è richiesto}: si
gioca con i propri vestiti di tutti i giorni. Lo studente di prima interviene: a
Bracciano, in effetti, giocavano con i propri vestiti (``di solito gli eventi
così sono degli ibridi''); e ricorda di aver giocato un altro LARP al
\textbf{Play}\footnote{La trascrizione dice ``Play di Bologna''; il festival
Play, dedicato al gioco, si tiene a Modena.}, attorno a un tavolo, ognuno con un
ruolo e con i propri vestiti: \textbf{\emph{A cena col nonno}}. ``Perfetto, quello
è un chamber LARP''.

Ruggero precisa un punto sottile: alcuni chamber LARP si giocano \textbf{a un
tavolo}, ma questo non li rende giochi di ruolo da tavolo, perché si giocano
personaggi \emph{che sono seduti a un tavolo}, mentre in \emph{Dungeons \&
Dragons} il personaggio non è necessariamente seduto a un tavolo in quel momento.
Con gli altri LARP hanno in comune la \textbf{sovrapposizione tra la fisicità di
chi gioca e quella del personaggio}, l'\textbf{embodiment}. Rispetto ad altri
LARP, i chamber LARP hanno di solito:
\begin{itemize}
  \item un livello di \textbf{accessibilità} più alto, perché spesso non
        richiedono costumi;
  \item un \textbf{budget ridotto};
  \item \textbf{non sono a evento}: molti si scaricano in \emph{print and play},
        si stampano a casa e si giocano con gli amici;
  \item per questo, un \textbf{numero di giocatori} più ridotto, di solito da
        3 a 20 persone.
\end{itemize}

\subsection{L'American freeform}

Una categoria dei chamber LARP è l'\textbf{American freeform}: LARP progettati da
designer che si occupano soprattutto di giochi di ruolo da tavolo, e infatti è
un \textbf{genere ibrido}. Visti da fuori sono chamber LARP (gente che fa
\emph{embodiment} del personaggio in delle stanze, non necessariamente in
costume), ma hanno ``un sacco di meccaniche'' prese dal gioco da tavolo e dal
gioco di ruolo da tavolo: le cosiddette \textbf{metatecniche} o metameccaniche,
di cui si parlerà in un'altra lezione. Anche questi sono spesso in \emph{print
and play}; alcuni sono venduti in un formato più vicino al gioco da tavolo, come
mazzi di carte o libretti. Quindi non sono a evento, a differenza dei boffer,
del Mind's Eye Theatre e dei murder party.

\subsection{I murder party}

Ruggero riprende l'ottima distinzione fatta da Beatrice. ``Murder party'', da noi,
indica culturalmente due cose:
\begin{itemize}
  \item le \textbf{cene con delitto}, in cui si va al ristorante, degli attori
        mettono in scena un omicidio e la successiva indagine, e i commensali
        mangiano e intanto risolvono il caso. \textbf{Nessuno dei commensali è
        l'assassino}: è uno degli attori;
  \item la categoria ``più british'', i \textbf{murder party all'inglese}, o i
        loro equivalenti che si stanno diffondendo moltissimo in Asia in questo
        periodo\footnote{La trascrizione dice ``Jubel Shock'' e ``Juvenile
        Show''. Si tratta quasi certamente dei \emph{jubensha}, i giochi di
        investigazione con ruoli molto diffusi in Cina.}, in cui \textbf{si
        interpretano i personaggi della storia} su cui si indaga: tra i giocatori
        può esserci l'assassino.
\end{itemize}
La seconda forma rientra, per il design dell'esperienza e delle meccaniche,
sotto l'ombrello dei LARP; gli aspetti ambientali e le decorazioni la
avvicinano poi al mondo del pervasivo.

% ---------------------------------------------------------------------
\section{Il Nordic LARP}
\label{a2:sec:nordic}

\slidefig{img/AGD/agd2-nordic}{Il Nordic LARP e i Knutebooks, i volumi collettivi del Knutepunkt.}{slide \emph{Analog -- Case Studies}}{fig:agd2-nordic}


C'è poi ``una grossissima categoria'', che si rivedrà più avanti perché ``è stata
importantissima'': il \textbf{Nordic LARP}, il LARP progettato in
\textbf{Scandinavia} (Danimarca, Norvegia, Svezia e Finlandia).

\subsection{Perché è così importante}

Ha un'importanza gigantesca nel settore \textbf{perché ha documentato tutto}. La
maggior parte dei libri accademici di game design e di analisi di questo tipo di
giochi, gli eventi e le conferenze, si fanno in Scandinavia, e c'è una lunga
storia di analisi e documentazione: la prima conferenza è del \textbf{1997}, e da
allora è annuale. È ``gente che si incontra da parecchi anni a parlare di come si
progetta e come si gioca''. Ed è una scena molto \textbf{cross-mediale}: prende
dagli alternate reality game, dai LARP americani, dai parchi a tema, dalle
simulazioni economiche, da tante teorie parallele a quelle del gioco, cercando
di interiorizzare, iterare e integrare tecniche di altri settori. Anche il Nordic
LARP ha le sue sottocategorie: alcune esperienze sono più vicine al chamber
LARP, altre al boffer.

\subsection{Come si riconosce}

Da fuori, un Nordic LARP si riconosce da \textbf{come è progettato il design
dell'esperienza}:
\begin{itemize}
  \item un focus sul \textbf{coinvolgimento}, sull'\textbf{emozione} e sulla
        \textbf{coesione tematica};
  \item pochissime cose \textbf{vestigiali}: per esempio, anche LARP molto
        grandi possono non avere il combattimento, se il combattimento non fa
        parte del tema;
  \item una cura del \textbf{post-esperienza}: molti giochi hanno una
        struttura per ``che cosa succede dopo il gioco''.
\end{itemize}
Molti concetti nati da questa scena si sono ``solidificati nel gergo del game
design'' e torneranno nella lezione sull'emotional design.

\paragraph{Origini diverse da D\&D.} Storicamente, il LARP nordico ha
un'origine \textbf{precedente a \emph{Dungeons \& Dragons}}. Si tende a pensare
che il LARP nasca da ``prendiamo D\&D e lo facciamo dal vivo'', e per alcune
nazioni, come l'Inghilterra e l'America, è proprio così; per altre no. Il Nordic
LARP prende strategie e strumenti dal \textbf{teatro}, dalle avanguardie, dallo
\textbf{psicodramma}, dalla \textbf{Gestalt therapy}; ``è ovvio che quello che
esce fuori è un gioco molto particolare''.

\paragraph{I manifesti.} È anche un ambiente ``pieno, carico, zeppo di
manifesti''. La slide ne mostra alcuni, nati negli anni nella scena; non si
approfondiscono ora, ma sono tutti pubblici e si possono leggere a parte.

\subsection{Knutpunkt e i Knutepunkt books}

La cosa che interessa di più è la conferenza citata prima, che cambia nome a
seconda del paese che la ospita: \textbf{Knutpunkt}, \textbf{Knudepunkt},
\textbf{Knutepunkt} o \textbf{Solmukohta}\footnote{Rispettivamente in svedese,
danese, norvegese e finlandese: la conferenza ruota tra i quattro paesi. La
trascrizione dice ``Knutpunkt o Knutpunkt o Knudpunkt o Solmukott'' e, più
avanti, ``Knut Punk''.}. È il principale punto di ritrovo scandinavo sul gioco
di ruolo dal vivo: \textbf{quattro giorni di conferenza non stop}, con un
programma di centinaia di slot, spesso \textbf{otto track in parallelo} di talk,
panel e workshop. Si tiene dal 1997, e \textbf{ogni anno produce dei libri}
che riassumono lo stato dell'arte in forma di articoli: la documentazione è
``costantemente in aggiornamento''. Ruggero ne ha diversi, ``dei bei mattoni'':
ne mostra uno di \textbf{600 pagine} su come si progetta un LARP, uscito da una
delle edizioni (``nel 2018, se non ricordo male''). È forse ``una delle
comunità ludiche che più documenta quello che fa'', e soprattutto lo documenta
in modo \textbf{pragmatico}, lasciando framework e strumenti a chi vuole creare
giochi di questo tipo. Molte di queste cose si applicano anche ad altri generi, e
torneranno.

\subsection{Dogma 99}

In chat qualcuno chiede se c'entri il \textbf{Dogma} di Lars von Trier.
``Sì e no'': il \textbf{Dogma 99} è un manifesto di LARP che riprende il
manifesto cinematografico del \textbf{Dogma 95} per fare un parallelo con il
LARP. Il Dogma 95 aveva alcune direttive, per esempio che tutto debba essere
\textbf{diegetico}, senza nulla di extradiegetico; il Dogma 99 fa lo stesso: in
un LARP storico non ci devono essere cose extradiegetiche rispetto alla storia.
Una conseguenza: il Dogma 99 \textbf{non fa LARP fantasy}, perché la magia non
esiste e non si può riprodurre in modo perfetto; si fanno giochi ambientati
tendenzialmente ai giorni nostri, che si possono riprodurre facilmente. Da
questa catena di ragionamenti sono nati ``giochi molto interessanti''.

% ---------------------------------------------------------------------
\section{Jeepform: il dramma quotidiano}
\label{a2:sec:jeepform}

\slidefig{img/AGD/agd2-jeepform}{Jeepform tra tabletop, freeform e LARP; la locandina di Fat Man Down.}{slide \emph{Analog -- Case Studies}}{fig:agd2-jeepform}


Tra i manifesti c'era lo \textbf{Jeepform}\footnote{La trascrizione dice
``G-Form'' e ``Jeep Form''.}, nato all'interno del \textbf{freeform}, la zona ``a
metà tra il LARP e il tabletop''. Il suo manifesto, che si trova online, ha tra i
punti più importanti questo: \textbf{il dramma quotidiano è più interessante del
dramma epico}. Tutti i giochi Jeepform sono giochi di ruolo dal vivo su
\textbf{situazioni quotidiane}, ``drammaticamente quotidiane''.

Alcuni sono emotivamente molto pesanti, e storicamente importanti per vari
motivi; Ruggero ha lasciato fuori i più pesanti (se ne riparlerà nella lezione
sull'emotional design) e ne presenta due ``più leggeri'', che fanno comunque
capire in che toni si va.

\subsection{Il gioco della coppia perfetta}

Il primo\footnote{La trascrizione riporta il titolo come ``Lady 8'' e i nomi dei
personaggi come ``Lady'' e ``8''; né il titolo né i nomi sono verificabili.} è
un gioco per \textbf{quattro persone} e \textbf{due personaggi}, che per premessa
\textbf{si amano tantissimo}. Due persone alla volta interpretano i personaggi,
in un gioco \textbf{diviso in scene}; ogni scena è descritta dalle due persone che
non sono in gioco. L'unico obiettivo di chi gioca la coppia è \textbf{essere la
coppia perfetta}, mostrare quanto si amano. Se la scena si sposta su qualcos'altro,
\textbf{si rifà}, all'infinito, ``come nel giorno della marmotta'', finché non
parla solo di quanto si amano. Esempio: lei entra in bagno, la porta si rompe e
lei rimane chiusa dentro, lui è dall'altra parte; se lui comincia a riparare la
porta, si rifà la scena. ``Una paradossalità esasperata, quasi''.

\subsection{Fat Man Down}

Il secondo è introdotto con un \emph{content warning}: \textbf{Fat Man Down} ``è
un gioco molto pesante''. Non è il più pesante dei Jeepform, ma è il più pesante
di cui si parla nella lezione. È un LARP dello stile Jeepform, quindi un chamber
LARP tendenzialmente freeform.
\begin{itemize}
  \item Il personaggio principale è un uomo di successo, con un buon lavoro e
        una buona posizione economica, ma \textbf{obeso}.
  \item Tutti i partecipanti \textbf{si pesano su una bilancia}, e chi pesa di
        più interpreta questo personaggio.
  \item Il gioco è diviso in scene, ognuna con un \emph{framing} che è uno
        \emph{slice of life}: un'ora in ufficio, la strada da casa al
        supermercato, dieci minuti di pausa al parco.
  \item Tutte le altre persone in scena hanno la \textbf{direttiva meccanica} di
        farlo stare malissimo.
  \item Il protagonista ha una meccanica ``importantissima'': \textbf{tre
        cioccolatini} in tasca. In qualunque momento può scartarne uno e
        mangiarlo, e quando lo fa tutti i personaggi in scena devono trattarlo
        come la persona migliore del mondo.
\end{itemize}
Il personaggio è messo sotto \textbf{costante stress narrativo} finché non è
spinto a pensare ``devo mangiare per stare meglio''. ``La roba è molto pesante'',
e ci si ferma qui.

\subsection{Le meccaniche di sicurezza}

Avere un manifesto così noto nella scena nordica, con giochi molto pesanti, molto
intimi e soprattutto \textbf{senza un filtro fantasy o fantascientifico}, cioè
su drammi quotidiani, ha generato e normalizzato cose poi rivelatesi molto utili,
come le \textbf{meccaniche di sicurezza}. Se qualcuno ne ha sentito parlare, la
scena nordica è ``la principale responsabile'' della normalizzazione, della
ricerca e dell'iterazione delle meccaniche di sicurezza nel gioco di ruolo. Se ne
parlerà nella lezione sull'emotional design.

% ---------------------------------------------------------------------
\section{Il 360° e l'alta risoluzione}
\label{a2:sec:360}

\slidefig{img/AGD/agd2-360}{Il 360° e l'alta risoluzione: Ground Zero e Monitor Celestra.}{slide \emph{Analog -- Case Studies}}{fig:agd2-360}


Il Nordic LARP non è solo chamber, freeform e Jeepform. C'è anche ``quasi
l'opposto'': i LARP \textbf{360°} o ad \textbf{alta risoluzione} (\emph{high
resolution})\footnote{La trascrizione dice ``360 o High Resolution'' e, più
avanti, ``IRES 360'', ``iRes'', ``Ires''.}. Lo scopo è essere \textbf{il più
fisicamente coinvolgenti possibile}: tutto quello che si ha attorno deve essere
\textbf{diegetico}. Ma la soluzione non è quella del Dogma 99 (ambientare nel
quotidiano, così non serve la magia): ``la soluzione è buttare al muro un budget
talmente alto che la realtà a quel punto sembra magia''. Niente oggetti di scena
finti o armi finte: anche la scenografia è curata al punto che sembra di stare
davvero nel gioco. I casi studio sono ``anche abbastanza storici''.

\begin{description}
  \item[Ground Zero (1998--2001).] La premessa è la \textbf{Guerra Fredda}.
        Una cittadina americana riceve l'allarme per una bomba nucleare, e i
        partecipanti vivono \textbf{48 ore di fila}, non stop, chiusi in uno
        \textbf{scantinato} al buio, senza poter uscire. Il gioco si conclude
        con la bomba che cade: dietro le finte pareti di cartone dello
        scantinato c'erano dei \textbf{subwoofer giganti}, e alla fine i muri
        cadono addosso ai giocatori con un rumore devastante.
  \item[Panopticorp (2003, 2013).] Due \emph{run}, anche questo di 48 ore non
        stop: in teoria si dorme in gioco, ``ma in questo caso non si dormiva''.
        È una \textbf{distopia corporate} spinta al massimo: si è dipendenti di
        un'agenzia pubblicitaria sull'orlo della chiusura, che deve licenziare
        ``a manetta'' e tenere solo i migliori. Una spinta ultra competitiva e
        ``ultra tossica'' per vedere fin dove si è disposti a spingersi eticamente
        per restare in azienda, dormendo sotto le scrivanie.
  \item[Monitor Celestra.] Un 360° \textbf{fantascientifico} basato su
        \emph{Battlestar Galactica}: per giocarlo è stato noleggiato un
        \textbf{incrociatore svedese} in disarmo. Una nave, non un'astronave,
        ``ma se non usciamo all'esterno, ci mettiamo un paio di LED, ed ecco
        qua'': per i giocatori è un'astronave. Il gioco durava circa
        \textbf{60 ore}, con cabine allestite in modo che si dormisse in
        personaggio.
  \item[Demeter.] Un 360° sul viaggio della bara di \textbf{Dracula} dalla
        Transilvania a Londra, come nel romanzo di Stoker, in cui la nave
        \emph{Demeter} arriva in porto con tutti morti a bordo e il capitano
        legato al timone. I giocatori sono i passeggeri, e \textbf{da giocatori
        sanno che moriranno tutti}: per design è impossibile sopravvivere,
        ``altrimenti Dracula non esisterebbe come romanzo''. Si gioca in
        costume su una \textbf{nave vera} che, naturalmente, non va dalla
        Transilvania a Londra ma gira nel mare della Germania del nord;
        narrativamente, però, si gioca quel viaggio.
  \item[Bunker 101.] Un caso \textbf{italiano}, ispirato alla proprietà
        intellettuale di \emph{Fallout}: si giocava \textbf{72 ore} in un bunker.
        Qui il dormire non era diegetico: nell'ultima run hanno separato la parte
        di gioco dal sonno. Per tutto il resto, però, valeva il principio 360°:
        tutto ciò che si vede deve essere diegetico.
\end{description}

% ---------------------------------------------------------------------
\section{Tape LARP e black box}
\label{a2:sec:blackbox}

\slidefig{img/AGD/agd2-tape}{Tape LARP e black box LARP.}{slide \emph{Analog -- Case Studies}}{fig:agd2-tape}


Accanto al Jeepform (``in cui giochiamo in casa con i nostri costumi'') e al
360° (``l'ultra-diegesi''), ci sono tipologie scandinave ``molto più particolari
e quasi astratte''.

\subsection{Il tape LARP}

Visto che qualcuno conosce Lars von Trier: chi ha visto \emph{Dogville}, il film
con lo scotch per terra, sa da dove nascono i \textbf{tape LARP}. ``Abbiamo visto
Dogville, ci ispira un casino'': giochi in cui \textbf{lo scotch per terra indica
gli ambienti}, con tutte le meccaniche che ne derivano (il personaggio vede e
sente cose che normalmente non potrebbe sentire; come si agisce e non si agisce
da giocatori\ldots). Ne sono stati fatti parecchi, in particolare in Danimarca.

\subsection{Il black box LARP}

Ruggero è ``un particolare amante'' del \textbf{black box LARP}, che prende il
nome dalle \emph{black box}: per chi ha fatto teatro ``dalla Germania in su'',
spazi teatrali di solito senza spalti né sedie, con \textbf{pareti, soffitto e
pavimento neri}, molto diffusi in Germania e nel Nord Europa. A
\textbf{Copenaghen} c'è un festival, il \textbf{Black Box LARP Festival}, che usa
questi spazi per giochi creati apposta: giochi di ruolo che usano \textbf{luci e
colori come meccaniche di gioco}. Gli esempi della slide:

\begin{description}
  \item[Borderline.] Un gioco ``estremamente astratto''. Si gioca \textbf{in
        coppie}, e ogni coppia gioca \textbf{un solo personaggio}, con dei
        conflitti interni: tutta la conversazione è un monologo interiore del
        personaggio, e i due giocatori, con una candela, lottano per portare ciò
        che il personaggio fa dalla propria parte o da quella dell'altro.
  \item[Waiting for Flight.] Ruggero lo ha giocato. Si aspetta un aereo con a
        bordo persone care, amici, colleghi; \textbf{l'aereo precipita}, e
        muoiono quasi tutti. Si gioca la \textbf{sala d'attesa} di chi li
        aspettava: è un gioco sull'attesa, ``sul logorante non avere
        informazioni'' finché vengono centellinate. Si conclude con un momento
        molto catartico: arrivano le \textbf{valigie} delle persone attese, e in
        alcune (3 su 20) è nascosto un piccolo \textbf{cuore di vetro}: se lo
        trovi, la persona si è salvata; se no, non si è salvata. C'è una fase in
        cui speri che ci sia, ma potrebbe essere nascosto: ``è una cosa
        lacerante. Abbiamo pianto in 19 su 20''.
  \item[Stranger.] Di una designer danese\footnote{La trascrizione dice ``Nino
        Ressendrop'': con ogni probabilità Nina Runa Essendrop, designer danese
        nota per i LARP fisici e non verbali. Anche il titolo, ``Stranger'', è
        riportato come nella trascrizione.}. È un black box \textbf{non verbale}:
        è vietato parlare. I partecipanti sono divisi in due gruppi, i
        \textbf{blu} e i \textbf{verdi}, e prima del gioco ciascun gruppo crea un
        \textbf{linguaggio gestuale condiviso}: come la propria nazione esprime
        felicità, rabbia, paura, tristezza, amicizia, odio. Il twist: i verdi
        devono avere gesti squadrati, rigidi, geometrici; i blu gesti curvi,
        lenti, morbidi, sinuosi. Il gioco è sull'\textbf{immigrazione forzata}:
        ci sono fasi di giorno e di notte, e di notte chi facilita porta alcuni
        verdi nella nazione blu, dove si risveglieranno, con la regola che il
        confine non si può riattraversare. Le dinamiche: il primo che si
        risveglia tra i blu impara facilmente la loro lingua, e i blu lo aiutano
        con un \emph{onboarding} non verbale; ma quando arrivano molti verdi,
        tendono a muoversi \textbf{in gruppi}, perché a quel punto c'è già chi
        parla la propria lingua, anche se non è la lingua della nazione. ``Un
        microscopio'' di dinamiche della vita reale: emigrazione, immigrazione,
        inclusione, e quanto è più facile essere inclusi se il sistema ti
        supporta.
  \item[Il gioco di Ruggero.] Uno dei giochi della slide è suo, ``quindi non ne
        parliamo'', nonostante le proteste dei colleghi.
  \item[Un gioco su Cuba.] Di una designer che è andata a vivere a
        \textbf{Cuba} e ha partecipato da vicino alla vita di una delle storiche
        comunità queer cubane\footnote{La trascrizione riporta il titolo come
        ``Socialism and Word''; non è verificabile. Anche le date sono incerte
        (``nel 97, 89 e 91''): il contesto, la caduta del Muro di Berlino, fa
        pensare al 1987, 1989 e 1991.}. Il gioco racconta che cosa è successo
        alle comunità \textbf{punk, rock e queer}, molto sovrapposte, prima,
        durante e dopo la caduta del Muro di Berlino: lo Stato, rimasto senza i
        fondi sovietici, decide chi è cittadino di serie A e di serie B, e le
        persone queer e punk (la musica occidentale) vengono picchiate e
        imprigionate. Erano anche gli anni dello scoppio dell'\textbf{AIDS}. Una
        comunità famosa, quella dei \emph{frikis}\footnote{La trascrizione dice
        ``Os Fricchios''.} cubani (di cui sono ancora vive due persone), decise:
        ``tanto ci ammazza la polizia, tanto vale iniettarci sangue infetto da
        HIV'', per andare a morire non per strada ma nei \textbf{centri di
        recupero}, gestiti da medici stranieri, uno dei pochissimi casi in cui a
        Cuba poteva esserci personale non cubano. Lì si moriva, ma ``almeno si
        poteva ascoltare musica punk'', magari con il dottore che arrivava dalla
        Florida con le cassette rock, e tra simili. ``Una roba drammaticissima'':
        anche questo è un gioco devastante.
\end{description}

\begin{table}[htbp]
\centering
\small
\begin{tabular}{@{}p{3.1cm}p{9.8cm}@{}}
\toprule
\textbf{Tipo di LARP} & \textbf{Caratteristiche ed esempi citati} \\
\midrule
Boffer & nato per adattare D\&D dal vivo; costume, armi in plastazote,
combattimento e trame epiche; fino a centinaia o migliaia di persone \\
Mind's Eye Theatre & da \emph{Vampire: The Masquerade}; costume, dark/gotico,
intrigo politico più che combattimento \\
Chamber / parlor LARP & senza costume, accessibile, budget ridotto, spesso
\emph{print and play}, 3--20 persone (\emph{A cena col nonno}) \\
American freeform & chamber LARP con metatecniche prese dal gioco da tavolo e
dal GDR; anche in carte o libretti \\
Murder party & cena con delitto (gli attori indagano, l'assassino è un attore)
o murder party ``all'inglese'' (si interpretano i sospettati) \\
Nordic LARP & coinvolgimento, emozione, coesione tematica, cura del dopo-gioco;
manifesti; Knutpunkt dal 1997; Dogma 99 \\
Jeepform & freeform nordico sul dramma quotidiano (il gioco della coppia
perfetta, \emph{Fat Man Down}); meccaniche di sicurezza \\
360° / alta risoluzione & tutto diegetico grazie a un budget altissimo
(\emph{Ground Zero}, \emph{Panopticorp}, \emph{Monitor Celestra},
\emph{Demeter}, \emph{Bunker 101}) \\
Tape LARP & lo scotch per terra indica gli ambienti, come in \emph{Dogville} \\
Black box & spazi teatrali neri; luci e colori come meccaniche
(\emph{Borderline}, \emph{Waiting for Flight}, \emph{Stranger}) \\
\bottomrule
\end{tabular}
\caption{Le tipologie di gioco di ruolo dal vivo presentate da Ruggero, in
sintesi.}
\label{a2:tab:larp}
\end{table}

% ---------------------------------------------------------------------
\section{Chi osserva? Il LARP e il teatro dell'oppresso}
\label{a2:sec:oppresso}

\paragraph{La domanda di Tiziana.} Tiziana confessa un'impressione: tutti questi
giochi ``sembrano cose da studi'', come se ci fosse qualcuno che sta sopra e
osserva le dinamiche per ricavarne informazioni. Ruggero la rassicura: \textbf{non
è così}. I designer di questa scena, specialmente quella black box, vengono
dall'ambiente ibrido del teatro interattivo, degli \emph{happening}, dello
psicodramma, e \textbf{non documentano nulla di quel tipo}: non registrano le
giocate, non estrapolano informazioni da chi gioca. Spesso sono giochi per poche
persone (è difficile che un black box superi le 20), e molti hanno avuto
\textbf{run singole}, o al massimo due o tre. ``Si tratta solo di fare un certo
tipo di esperienza''. Come in altri media esistono generi con scopi diversissimi
(un film e un documentario sono entrambi operazioni visive), anche qui gli scopi
sono diversi. E chi non ha mai incontrato giochi del genere si chiederà
``perché dovrei giocare una roba del genere?'': la risposta arriverà con la
lezione sull'emotional design (``lo so, ve l'ho detto quattro volte'').

\paragraph{Il teatro dell'oppresso.} Prima di passare la parola, uno studente
interviene: in queste cose ritrova molto il metodo del \textbf{teatro
dell'oppresso}\footnote{Il metodo teatrale ideato dal regista brasiliano Augusto
Boal (integrazione).}. ``Ah, certo''. Ruggero torna indietro di una slide, a
quella sulle \textbf{influenze fondanti del Nordic LARP}: ci sono le tecniche di
\textbf{Moreno} (sociodramma e psicodramma), il \emph{sensitivity training}, il
\textbf{roleplay come metodo di ricerca}, e tutti quei metodi che usano il teatro
e l'interazione come strumento per far emergere il personale e l'intimo. Lo
studente precisa che il teatro dell'oppresso nasceva per \textbf{trovare
soluzioni} a situazioni che apparentemente non ne hanno: a turno il pubblico
proponeva una soluzione, si rifaceva sempre la stessa scena, e si studiava come
reagivano le persone; quindi, in un certo senso, il comportamento delle persone
\emph{si osserva}, ma dall'interno della scena.

Per Ruggero, nella scena LARP questa dimensione ``viene un po' meno'', perché lo
scopo resta \textbf{ludico} (``e scopo ludico può voler dire molte cose, come
vedremo''): è meno pragmatico del ``prendiamo quello che abbiamo giocato e il
gioco ci dà già gli strumenti per applicarlo''. Alcuni \textbf{debriefing}, i
momenti strutturati del dopo-gioco, possono dare strumenti a chi ha giocato; ma
questi giochi nascono come esperienze che devono essere ``ludicamente
interessanti per chi gioca'', \textbf{non come esperimenti scientifici ``con dei
topi dentro un labirinto''}.

\paragraph{Approfondimento: il LARP in classe e il debriefing.} Che il LARP
possa essere uno strumento formativo, come nel volume sull'edu-larp citato
all'inizio, lo mostra anche \emph{Challenges for Game Designers}. Nel capitolo
sui giochi come strumento di insegnamento, Brathwaite e Schreiber propongono una
sfida, \emph{Living History}: progettare un LARP che rappresenti un evento storico
(una convenzione costituzionale, un protocollo internazionale, una guerra)
giocabile da una classe di 20--30 studenti in una o due ore. I passi suggeriti
sono: scegliere un evento in cui molte parti, ciascuna con la propria agenda, si
incontrano in una negoziazione o in un conflitto; scegliere i giocatori (o le
squadre); definire gli \textbf{obiettivi} di ciascuno; creare regole semplici e
rapide per risolvere i conflitti, perché i giocatori si muovono e parlano di
continuo (``alcuni LARP usano un singolo tiro di dado o una partita a
sasso-carta-forbici''); e preparare, oltre alle schede dei personaggi, delle
\textbf{note di debriefing}: che cosa è successo davvero e quali domande
discutere confrontando la storia con l'esito della partita. Lo scopo dichiarato è
che gli studenti capiscano che ``l'evento era complicato'' e come l'esito sia
\emph{emerso} dalla situazione iniziale. È la versione didattica della cura del
dopo-gioco che Ruggero attribuisce alla scena nordica.

% ---------------------------------------------------------------------
\section{I giochi di ruolo da tavolo: quattro parole per orientarsi}
\label{a2:sec:gdr}

``Daniele, ti lascio la parola''. Daniele (``mi giri tu le slide quando te le
chiedo'') avverte che fare una panoramica dei giochi di ruolo in poco tempo è
``veramente difficile'': ce ne sono tantissimi, e negli anni si sono evoluti e
diversificati sempre di più, fino a diventare molto diversi tra loro.

\subsection{Il tribalismo e le regole tramandate a voce}

C'è anche un altro problema, che si vedrà nella lezione sulla fiction e che
esiste in tutto il gioco analogico: il \textbf{tribalismo}, cioè
l'\textbf{ereditarietà orale delle regole}. Spesso i giochi di ruolo venivano
insegnati da altre persone, anche per motivi economici. Daniele ha cominciato
così: da bambino, giocando per la prima volta a D\&D, l'unico ad avere il
manuale era il master; ``ergo io non avevo accesso alle regole'', e il master
poteva inventarsi quello che voleva.

\subsection{Conflitto, azione, risoluzione, effetto}

\slidefig{img/AGD/agd2-car}{Conflitto, azione, risoluzione, effetto: le quattro parole per orientarsi nei giochi di ruolo da tavolo.}{slide \emph{Analog -- Case Studies}}{fig:agd2-car}


Per capire un po' i giochi di ruolo bisogna ricordare \textbf{quattro parole},
che torneranno nella lezione sulla fiction: \textbf{conflitto}, \textbf{azione},
\textbf{risoluzione} ed \textbf{effetto}. Nella slide \emph{conflitto} e
\emph{risoluzione} hanno un asterisco: per ora basta capirli quanto serve a
leggere i casi studio, senza analizzarli troppo.
\begin{description}
  \item[Conflitto.] L'ostacolo, l'incertezza, l'evento che \textbf{stimola
        un'azione o la contrasta}. ``Tornerà molto''.
  \item[Azione.] Quello che si descrive di voler fare nei panni del personaggio
        per risolvere un conflitto.
  \item[Risoluzione.] Come si risolve l'azione dentro il gameplay, con le regole:
        che cosa succede?
  \item[Effetto.] La conseguenza che si applica.
\end{description}
L'esempio: ``voglio colpire il goblin con la mia spada''. Il goblin ha il suo
scudo e la sua armatura. Che azione fai? ``Lo becco con la spada''. Si tirano i
dadi, si risolve, si applica un effetto. Ma non è così semplice, perché in
mezzo c'è la \textbf{fiction}: nel gioco di ruolo si tende a riscrivere tutto
quello che succede nel gameplay \emph{dentro} la narrazione. Se faccio tre danni
agli HP del goblin, che cosa vuol dire narrativamente? Come li descrivo? Da qui
nasce ``una serie di problematiche''.

\begin{figure}[htbp]
\centering
\begin{tikzpicture}[
    font=\small,
    passo/.style={draw, thick, rounded corners=3pt, fill=gray!12,
                  minimum width=2.6cm, minimum height=0.95cm, align=center},
    esempio/.style={font=\footnotesize\itshape, align=center, text width=2.8cm},
    freccia/.style={-{Stealth[length=2.5mm]}, thick}
  ]
  \node[passo] (c) at (0,0)    {\textbf{Conflitto}$^{*}$};
  \node[passo] (a) at (3.4,0)  {\textbf{Azione}};
  \node[passo] (r) at (6.8,0)  {\textbf{Risoluzione}$^{*}$};
  \node[passo] (e) at (10.2,0) {\textbf{Effetto}};
  \draw[freccia] (c) -- (a);
  \draw[freccia] (a) -- (r);
  \draw[freccia] (r) -- (e);
  \node[esempio, below=0.25cm of c] {il goblin, con scudo e armatura};
  \node[esempio, below=0.25cm of a] {``lo colpisco con la spada''};
  \node[esempio, below=0.25cm of r] {tiro dei dadi secondo le regole};
  \node[esempio, below=0.25cm of e] {tre danni agli HP};
  \node[draw, dashed, rounded corners=3pt, fit=(c)(e), inner sep=0.5cm,
        yshift=-0.35cm, minimum height=2.6cm] (fic) {};
  \node[font=\footnotesize, anchor=south] at (fic.north)
    {\textbf{fiction}: tutto viene ridescritto nella narrazione (che cosa
    \emph{sono}, nella storia, tre danni?)};
\end{tikzpicture}
\caption{Le quattro parole di Daniele per leggere i giochi di ruolo da tavolo,
con l'esempio del goblin. Gli asterischi sono quelli della slide: conflitto e
risoluzione verranno analizzati meglio nella lezione sulla fiction. Nel gioco di
ruolo ``classico'' il conflitto precede l'azione, la risoluzione è un tiro di
dado e l'effetto è descritto dal master e dal regolamento; molti giochi
successivi cambiano proprio questo schema.}
\label{a2:fig:car}
\end{figure}

% ---------------------------------------------------------------------
\section{Le origini: da Dungeons \& Dragons agli anni Novanta}
\label{a2:sec:gdr-storia}

Le quattro parole servono a Daniele per una breve panoramica, anche storica,
dei giochi di ruolo.

\subsection{Dungeons \& Dragons e il gioco di guerra}

Il primo gioco di ruolo da tavolo ``così come viene concepito'' è sicuramente
\textbf{\emph{Dungeons \& Dragons}}, di \textbf{Gygax e Arneson}, nel
\textbf{1974}. C'è in realtà un precedente, \textbf{\emph{Chainmail}}\footnote{La
trascrizione dice ``C'è in Mail'' e data \emph{Chainmail} al '72; il gioco è
del 1971.}, e come per tutte le origini (``anche \emph{Cyberpunk} ha lo stesso
problema'') c'è un po' di ambiguità su che cosa venga prima. Ci sono libri che
raccontano la nascita di D\&D, c'è la storia della \textbf{TSR}, la casa
editrice che lo pubblicò per prima, ci sono versioni discordanti e, negli anni,
parecchi contrasti tra le famiglie di Gygax, di Arneson e di altri.

D\&D viene dal \textbf{gioco di guerra}, il gioco di \textbf{miniature}. ``La
leggenda vuole'' che Gygax e Arneson, giocando, si siano detti: ho battuto il
tuo esercito; che cosa succede se il mio generale entra nel tuo castello e va a
trattare con il tuo generale? Abbiamo regole per questo? No: inventiamocele.
``Ripeto, una leggenda'': non si sa se sia andata così, e in fondo non interessa
molto.

\subsection{I generi e la struttura del gioco di ruolo ``early''}

\slidefig{img/AGD/agd2-dnd}{Il flusso di Dungeons \& Dragons letto con conflitto, azione, risoluzione ed effetto.}{slide \emph{Analog -- Case Studies}}{fig:agd2-dnd}


Quello che interessa è che dopo D\&D, nell'\textbf{RPG delle origini}, vengono
coperti più o meno \textbf{tutti i generi letterari}: \emph{Call of Cthulhu}
l'orrore lovecraftiano, \emph{Vampire: The Masquerade} l'orrore gotico,
\emph{Traveller} la fantascienza, \emph{Cyberpunk} il cyberpunk, e \emph{GURPS},
un \textbf{sistema generico} con cui ``teoricamente'' si può giocare qualunque
genere.

Nel gioco di ruolo classico:
\begin{itemize}
  \item il \textbf{conflitto} poteva venire prima o dopo l'\textbf{azione}:
        ``voglio colpire quel goblin'' (e il goblin ha la sua difesa), oppure
        ``corri, ti sta inseguendo una palla'', come in \emph{Indiana Jones};
  \item la \textbf{risoluzione} era il tiro di dado, e l'\textbf{effetto} veniva
        descritto dal master e dal regolamento;
  \item in pratica, \textbf{tutto tranne l'azione}, che il giocatore gestiva e
        narrava, era gestito dal \textbf{master}.
\end{itemize}

\paragraph{Il master come arbitro.} All'inizio il master non era un ``maestro''
né una figura che insegna: era interpretato soprattutto come un
\textbf{arbitro}. Nella prima versione di D\&D non c'è ancora l'idea, venuta
dopo, del master come ``dio del tavolo'': quando deve decidere, dovrebbe farlo
secondo criteri di realismo e oggettività (``puoi essere oggettivo a un tavolo
di gioco? La voglio vedere''). Questa idea è importantissima perché verrà ripresa
dall'\textbf{OSR}, l'\emph{Old School Revival}, il movimento che cerca di
riportare nella modernità il modo in cui, a suo dire, si giocava D\&D agli
inizi.

\paragraph{Il master è un giocatore.} Nei giochi di ruolo ci sono quindi
\textbf{ruoli asimmetrici}. Ma Daniele raccomanda di ricordare una cosa, che
tornerà nella lezione sulla fiction: \textbf{il master è un giocatore}, magari un
giocatore con più regole e più autorità al tavolo, ma ``un giocatore tanto quanto
i giocatori che interpretano un personaggio''. Nei GDR delle origini: si
interpreta un personaggio, con una \textbf{scheda} e le sue statistiche, dentro
una fiction e un \textbf{setting} stabilito dal manuale o dal master; quando si
dice che cosa si vuole fare, il master reagisce. E lentamente, con gli anni, il
master assume anche il ruolo di \textbf{regista} e \textbf{sceneggiatore}: scrive
l'avventura, decide che cosa succede, decide i tagli di scena.

\paragraph{Approfondimento: dal wargame al contratto sociale.} La genealogia
raccontata da Daniele si ritrova in \emph{21st Century Game Design} di Bateman e
Boon. Le radici del gioco di strategia, scrivono, precedono i videogiochi di
decenni: alla fine degli anni Cinquanta Charles Roberts fonda Avalon Hill e
pubblica i primi \textbf{wargame da tavolo} (1958); in un'epoca senza
videogiochi, giochi di ruolo e giochi di carte collezionabili, i wargame
``fondano l'industria dei \emph{hobby games}'', e \emph{Chainmail} (1971),
ispirato a quei giochi, è il ``trampolino'' verso \emph{Dungeons \& Dragons}. Gli
stessi autori descrivono la forma minima del gioco di ruolo da tavolo come un
\textbf{contratto sociale} tra l'arbitro (il \emph{games master}) e i
giocatori: i giocatori accettano che sia l'arbitro a decidere come le loro azioni
sviluppano la storia, e l'arbitro si impegna a garantire intrattenimento o
imparzialità, come mediatore e narratore principale. Citano anche un'osservazione
attribuita a Gygax: se i giocatori di ruolo si accorgessero di poter fare tutto
da soli, i game designer resterebbero senza lavoro. Il fatto, notano, è che la
maggior parte dei master ha bisogno di meccaniche per gestire il mondo, e i
giocatori si fidano di più dell'arbitro quando una struttura di meccaniche
casuali, i dadi, fa da cornice. Anche Schell, in \emph{The Art of Game Design},
cita il master dei giochi di ruolo da tavolo come esempio di stato del gioco
\emph{nascosto}: è lui a conoscere privatamente gran parte dello stato del gioco,
perché ne è ``il meccanismo operativo''.

\subsection{Gli anni Ottanta e Novanta: adattare un regolamento a tutto}

Qui Daniele fa un ``piccolo inciso editoriale''. Nella grande espansione degli
anni Ottanta e Novanta \textbf{non c'era un professionismo} del settore. Gli
editori compravano delle proprietà intellettuali, per esempio \textbf{007},
prendevano un master appassionato e competente e gli chiedevano di
\textbf{adattare} il regolamento di D\&D, o un altro regolamento
esistente\footnote{La trascrizione dice ``Call of Duty 2'', che non può essere
corretto in questo contesto; il senso è ``il regolamento di un altro gioco
già esistente''.}, a quella IP. A volte nascevano giochi innovativi, a volte
giochi ``non molto interessanti da giocare'': immaginate di giocare con le
regole di D\&D qualcosa che quelle regole non supportano. Verso la fine degli
anni Novanta il problema della qualità ``ha iniziato a farsi sentire''.

\paragraph{Approfondimento: James Bond 007 e l'esperienza essenziale.} Proprio
il gioco di ruolo di 007 è il caso con cui Schell illustra, in \emph{The Art of
Game Design}, la \textbf{lente dell'esperienza essenziale} (\emph{Lens of
Essential Experience}), e mostra l'altra faccia della storia di Daniele. Il
designer Chris Klug era frustrato dai precedenti giochi di ruolo di spie, come
\emph{Top Secret} della TSR, perché ``si giocavano troppo come dei wargame'':
l'essenza di ciò che rende eccitanti i film di spionaggio non c'era. Per il
gioco di James Bond, Klug progettò le meccaniche in modo che restituissero in
ogni modo possibile la sensazione dei film. L'esempio è l'invenzione degli
\textbf{Hero Points}: nei GDR tradizionali, davanti a un'azione rischiosa (saltare
da una finestra su un elicottero in volo), il master calcola una probabilità, il
giocatore tira i dadi ``e questo è quanto''; ma se la probabilità è bassa
nessuno rischia, e se è alta tutti si comportano da supereroi. I punti eroe
risolvono il dilemma in modo coerente con il genere. È esattamente quello che
la New Wave rivendicherà: \textbf{il sistema conta}, e un regolamento adattato ``a
nastro'' non restituisce l'esperienza di un genere.

% ---------------------------------------------------------------------
\section{La New Wave e The Forge}
\label{a2:sec:gdr-newwave}

\subsection{Il sistema conta}

La slide successiva è sulla \textbf{New Wave}: un movimento nato tra la fine
degli anni Novanta e l'inizio dei Duemila. Daniele avverte che non vuole fare
un'analisi geografica ma ``un riassunto della storiografia''. La New Wave cerca
un \textbf{metodo diverso di interpretare il gioco di ruolo}, e la prima cosa
che stabilisce, in un forum chiamato \textbf{The Forge} (c'è anche un libro in
inglese su questo movimento), è che \textbf{il sistema dei giochi di ruolo è
importante}: le regole, il gameplay, la parte meccanica \textbf{non sono un
orpello}. Non si può tenere lo stesso sistema, cambiare il setting e sperare che
funzioni da sé. Tanto è vero che perfino \emph{GURPS}, il sistema generalista,
aveva tantissimi manuali di espansione delle regole per ogni genere:
\emph{GURPS} base, \emph{GURPS Fantasy}, \emph{GURPS} fantascienza\ldots

\subsection{I giochi della New Wave}

\slidefig{img/AGD/agd2-newwave}{Alcuni giochi della New Wave.}{slide \emph{Analog -- Case Studies}}{fig:agd2-newwave}


Dentro la New Wave nascono esperienze diversissime, che spostano l'attenzione
su due cose: la \textbf{scrittura del manuale}, diversa da quella tradizionale,
e la \textbf{rottura degli schemi}.
\begin{itemize}
  \item \textbf{\emph{Apocalypse World}}\footnote{La trascrizione dice ``A
        Polly's Ward'' e ``Minster Baker''.}, di \textbf{Vincent Baker}, cerca di
        scardinare il modo classico di strutturare un manuale di gioco di ruolo
        (Daniele accenna a una nuova edizione recente).
  \item \textbf{\emph{Nobilis}}, uno dei primi giochi abbastanza famosi
        \textbf{diceless}, senza dadi: azioni e conflitti si risolvono in un
        altro modo.
  \item \textbf{\emph{Don't Rest Your Head}} (in italiano \emph{Non cedere al
        sonno}) e \textbf{\emph{Polaris}}, che lavorano sul concetto di chi è e
        chi non è il personaggio. \emph{Non cedere al sonno} è un gioco
        ``surrealista'' in cui ci si sveglia a \textbf{Mad City}; la parte
        importante sono i \textbf{poteri}. All'epoca \emph{Vampire} aveva
        stabilito che i poteri fossero una lista di possibilità con statistiche
        a cui applicare le regole; in \emph{Don't Rest Your Head} ci sono poteri
        ``totalmente assurdi, fenomenali'', che non richiedono meccaniche in più
        ma usano le meccaniche base ``in maniera molto innovativa''.
  \item \textbf{\emph{Fate}}, un sistema anch'esso generalista, ``più o meno
        come \emph{GURPS}'', ma con delle sue specifiche: ci sono
        \textbf{procedure per creare i setting} e per stabilire quali regole un
        setting richiede.
  \item \textbf{\emph{Dread}}, uno dei primi giochi a usare per l'incertezza
        \textbf{un elemento che non è un dado}: una \textbf{torre Jenga}. È un
        gioco dell'orrore in cui ogni volta che si fa un'azione si toglie un
        blocchetto, e se la torre cade il personaggio muore (``non è proprio
        così, è più complicato, però per semplificare'').
\end{itemize}

\paragraph{Le guerre di religione.} La New Wave ha avuto ``una serie di
problematiche'' in giro per il mondo, e soprattutto in Italia ha creato una
sorta di \textbf{guerra teologico-filosofica} tra i ``tradizionalisti'' di D\&D e
i ``forgiti'' della New Wave. ``Per vostra sanità mentale non ci entriamo''.

\paragraph{Una reazione a due cose.} La New Wave è stata una reazione:
\begin{enumerate}
  \item ai manuali in cui ``c'è un sistema e quello è il sistema'', con il
        setting applicato sopra ``a nastro'', senza intreccio tra la narrazione
        e le regole che la raccontano: si voleva \textbf{coerenza};
  \item all'assenza di studio: è stato il \textbf{primo movimento a voler
        studiare i giochi di ruolo}, a capire come sono fatti. Da qui viene il
        \textbf{Big Model}, un modello di analisi dei giochi di ruolo di cui si
        vedranno alcune cose nella lezione sulla fiction.
\end{enumerate}
Ma il Big Model ``non è fondamentale'', perché per Daniele uno dei problemi della
New Wave è aver creato \textbf{un dizionario collaterale} a quello del game design
del gioco da tavolo e del videogioco: termini diversi che dicono la stessa cosa.

\subsection{The Forge e la scena nordica}

Ruggero fa da ponte: \textbf{The Forge ha fatto per il gioco di ruolo da tavolo
quello che la scena scandinava ha fatto per il LARP}. Inventavano e scrivevano
moltissimo, hanno lasciato framework e strumenti. La scena scandinava è nata un
po' prima, The Forge un po' dopo, ma hanno agito in parallelo. La differenza è
che The Forge era \textbf{meno strutturato}: essendo un forum, era più caotico,
con discussioni anche accese, soprattutto sulle interpretazioni. Anche la
\textbf{preservazione storica} è diversa: della scena nordica ci sono i libri
fisici, mentre per il forum bisogna passare dalla Wayback Machine. Ma in termini
di normalizzazione di termini, meccaniche e framework, entrambi hanno fatto
moltissimo.

\subsection{Fiasco}

\slidefig{img/AGD/agd2-fiasco}{Il flusso di Fiasco: chi stabilisce la scena e chi la risolve.}{slide \emph{Analog -- Case Studies}}{fig:agd2-fiasco}


Daniele riprende: a un certo punto The Forge chiude, e nascono altri giochi
dagli stimoli e dalle discussioni che quell'ambiente aveva generato, non
necessariamente da persone che ne facevano parte. Per esempio quelli di
\textbf{Jason Morningstar}, noto anche per dei chamber LARP (il migliore amico
di Daniele, Edoardo, ``più esperto di me'', ne è talmente fan che il suo gatto si
chiama Giasone).

\textbf{\emph{Fiasco}} è per Daniele uno dei primi giochi che \textbf{rompe il
canone}. Racconta, appunto, un fiasco, come nei film dei \textbf{fratelli Coen}
o di Tarantino: c'è uno scenario, ci sono dei personaggi, e ogni volta che si fa
una scena \textbf{il giocatore di turno sceglie se aprirla} (e gli altri la
risolvono) \textbf{o se lasciare che gli altri la impostino} (e risolverla lui).
Già solo questo ``micro elemento'' (non l'unica meccanica: c'è anche quella,
interessantissima, dei dadi) rompe l'idea che in un gioco di ruolo si debba
descrivere l'azione, tirare i dadi e vedere che cosa succede, oppure reagire a
ciò che descrive il master. Per chi veniva da \emph{AD\&D} seconda edizione o da
\emph{Vampire} era ``assolutamente interessante'': più leggero e più legato alla
fiction. Anche i \textbf{dadi} sono usati in modo nuovo: vengono tirati
\textbf{prima}, messi sul tavolo, selezionati, usati quasi come
\textbf{token}: un'innovazione anche nel modo di usare i materiali.

% ---------------------------------------------------------------------
\section{Il gioco di ruolo contemporaneo}
\label{a2:sec:gdr-oggi}

\subsection{Prima del COVID: i giochi focalizzati}

\slidefig{img/AGD/agd2-contemporary}{Giochi di ruolo contemporanei.}{slide \emph{Analog -- Case Studies}}{fig:agd2-contemporary}


Daniele salta alla contemporaneità. Prima del COVID c'è un grande aumento di
giochi che, mentre le persone si appassionano ai giochi di ruolo, vanno o nella
direzione di D\&D o in una \textbf{direzione più ibrida}:
\begin{itemize}
  \item \textbf{\emph{Blades in the Dark}}\footnote{La trascrizione dice
        ``Blaze in the Dark''. Daniele lo data ``2016--18 più o meno''.};
  \item più recenti, il gioco di ruolo di \textbf{\emph{Cowboy
        Bebop}}\footnote{La trascrizione dice ``Cabobby Bop'' e ``Cobo
        Bobby''.} e \textbf{\emph{The Clay That Woke}} di \textbf{Paul
        Czege}\footnote{La trascrizione dice ``Paul Ksege, The Clayed Walk''.},
        autore di ``uno dei principi di design più interessanti'';
  \item \textbf{\emph{Fabula Ultima}}, di \textbf{Emanuele Galletto}, italiano
        (``che salutiamo'').
\end{itemize}
Sono prodotti che \textbf{non cercano di fare la stessa cosa di D\&D}. Hanno un
impianto simile (scheda del personaggio, dadi, azioni descritte, tiri), ma si
pongono \textbf{obiettivi di design più coerenti}: ``si chiedono che cosa vogliamo
che il giocatore senta''. Il sistema di \emph{Cowboy Bebop}, per esempio, usa
gli \textbf{orologi} (\emph{clock}) che si trovano anche in \emph{Blades in the
Dark} e nei giochi \emph{Powered by the Apocalypse}, cioè la stessa meccanica,
ma è un sistema completamente diverso, perché deve raccontare una fiction
diversa:
\begin{itemize}
  \item in \emph{Blades in the Dark} si giocano dei \textbf{criminali} a
        \textbf{Doskvol}\footnote{La trascrizione dice ``Oswald''.}, città di un
        mondo alternativo dove è sempre notte, dove i fantasmi non passano
        nell'aldilà e vivono con i vivi, e dove si è costretti a restare perché
        fuori dalle mura i fantasmi uccidono. Si fa parte di una banda che cerca
        di scalare la gerarchia della città, ``tipo \emph{Peaky Blinders}'';
  \item in \emph{Cowboy Bebop} si è ``letteralmente una ciurma di persone nello
        spazio, nel futuro, che fanno cose'', perché ``l'anime si racconta con i
        personaggi che fanno cose'' (Daniele, ``molto fan dell'anime'', lo dice
        come un complimento).
\end{itemize}

\subsection{Dopo il COVID: il ritorno all'origine}

Soprattutto dopo il COVID si è notato un ``ritorno fortissimo'' al
\textbf{giocare come una volta}. Per Daniele entrambe le tendenze rispondono a
una necessità: \textbf{semplificare}. \emph{D\&D} 3 e 3.5, le nuove edizioni di
\emph{Vampire} e altri sono diventati ``pieni di regole, pieni di contenuti'', e
chiedono a master e giocatori di imparare moltissimo, senza contare il tempo
necessario per una campagna. Allora, da un lato, si cercano sistemi più semplici
``da intavolare'' che facciano bene una cosa; dall'altro si cerca la stessa cosa
per altri motivi, ricreando le atmosfere del D\&D \emph{old school}: \textbf{alta
mortalità}, parte narrativa e interpretativa meno importante, immediatezza.
\textbf{\emph{Mörk Borg}}\footnote{La trascrizione dice ``Morg Borg''.}, per
esempio, sono circa \textbf{cento pagine}, con poche regole. C'è più attenzione
al ``realismo'' (termine un po' sbagliato, si vedrà nella lezione sulla
fiction), e il vantaggio è che questi giochi possono usare sia le proprie regole
sia quelle di \textbf{D\&D quinta edizione}, che è stata importante soprattutto
``per un motivo editoriale'' (anche questo si vedrà nella lezione sulla
fiction).

\paragraph{La domanda di Gianmarco.} Gianmarco chiede se questa tendenza a
manuali e partite ``asciutti'' segua un pubblico più adulto, con meno tempo, o
sia una tendenza generale del mercato. È soprattutto la prima. Il gioco di ruolo
ha avuto parecchie difficoltà ad attirare nuove leve, e i giocatori sono
cresciuti. Daniele gioca da molti anni, ma il suo ``vero inizio'' è stato con un
gioco della New Wave\footnote{La trascrizione dice ``Paura dell'Epoclips'':
probabilmente \emph{Il mondo dell'apocalisse}, l'edizione italiana di
\emph{Apocalypse World}, ma non è certo.}: per lui, all'università, avere in un
solo manuale da 200 pagine tutte le regole e tutti i personaggi è stato ``un
toccasana''. Chi inizia con D\&D, magari alle superiori o all'università, man
mano che cresce tende a preferire:
\begin{itemize}
  \item o giochi più focalizzati sull'aspetto \textbf{narrativo}, che hanno
        comunque scheda, tiro di dado e struttura tradizionale, ma una scheda e
        un sistema più semplici;
  \item o un'esperienza simile al D\&D canonico, ma con meno ``regoline e
        regolette'' e meno \textbf{simulazionismo}.
\end{itemize}
Due differenze pratiche. I manuali \textbf{OSR}, che un po' di simulazionismo ce
l'hanno, sono pieni di \textbf{soluzioni per l'arbitro} (``che non è il master,
è l'arbitro'') per risolvere i problemi sul momento; e l'OSR è \textbf{sandbox}:
c'è una mappa, una zona in cui si va dove si vuole, e il master non ha scritto
un'avventura. I \textbf{giochi narrativi}, dal canto loro, non hanno una campagna
scritta ma uno \textbf{scenario}. In entrambi i casi il tempo di preparazione del
master si riduce ``di gran lunga''.

\subsection{Focalizzati, non semplificati}

Ruggero interviene sulle parole: non vorrebbe che passasse l'idea che questi
giochi siano ``più semplici'' o \emph{dumbed down} rispetto a D\&D. Si parla di
\textbf{giochi focalizzati}: invece di semplificare, \textbf{focalizzano}, sono
molto più \textbf{coesi tematicamente}, hanno ``un cuore tematico molto forte'' e
si giocano per giocare quello. D\&D invece, a livello di cuore meccanico, è
giocato in modi diversissimi da dieci gruppi diversi.

Daniele fa un esempio: nei manuali di D\&D (dalla terza edizione e 3.5 alla
quinta) ci sono le regole per il \textbf{danno da caduta}. In \emph{Blades in
the Dark} non servono: c'è un altro metodo, ``non è così generalista, è
focalizzato'', e la narrazione si risolve in altro modo. Anche nell'OSR:
\emph{Shadowdark} ha il danno da caduta, perché magari è importante;
\emph{Neverland}, una reinterpretazione di \textbf{Peter Pan} in chiave OSR, no,
perché lì conta di più che il manuale offra tutti i personaggi del Peter Pan
tradizionale in una versione \emph{grimdark}. ``Anche l'OSR diventa un po'
focalizzato quando ci prova''.

\subsection{Not the End}

\slidefig{img/AGD/agd2-nottheend}{Il flusso di Not the End: il sacchetto di token come motore della risoluzione.}{slide \emph{Analog -- Case Studies}}{fig:agd2-nottheend}


Nicola, per cui questo è ``il settore che più adoro del gioco analogico'',
chiede dove si collochi \textbf{\emph{Not the End}}, un gioco italiano. La
risposta è nella slide successiva. \emph{Not the End} continua nella direzione
più narrativa e focalizzata (``semplicistica è un termine non corretto''): ha
procedure, meccaniche, un gameplay focalizzato a fare quello che il gioco
chiede. Si possono creare i propri setting, ma il gioco resta quello: il
personaggio principale è un \textbf{eroe}, ed è il giocatore a decidere quando
sarà la sua fine, quando giocherà il suo momento di eroismo stoico (da cui il
titolo). È un gioco \textbf{senza dadi}: c'è un conflitto, si fa e si descrive
un'azione, e la \textbf{risoluzione} consiste nell'\textbf{estrarre dei token}
bianchi e neri (o arancioni e viola, per le persone con difficoltà nella visione
dei colori); l'effetto viene descritto secondo le procedure del gioco. Il
panorama, insomma, si è diviso: c'è chi cerca un'esperienza più simile
all'originale, e chi un'esperienza più focalizzata sulla narrazione, magari
senza un \emph{world building} preciso.

\subsection{Che cos'è l'OSR, e com'è nato davvero il gioco di ruolo}

Un'altra studentessa chiede che cosa significhi OSR: \textbf{Old School
Renaissance}, ``rinascimento della vecchia scuola''\footnote{Poco prima Daniele
aveva sciolto la sigla come \emph{Old School Revival}; entrambe le forme sono in
uso.}. È un movimento ``molto, molto particolare'': ogni persona che ne fa parte
ha il proprio principio di OSR, gioca un OSR diverso dagli altri, e gli
appassionati ``adorano scontrarsi'' su che cosa intendono per OSR. Per questo
nella lezione sulla fiction Daniele parlerà di \textbf{culture} del GDR, più che
di community.

La stessa studentessa (iscritta a un altro corso dell'accademia, di cui sta
recuperando le lezioni del primo anno) riporta quello che ha capito dell'origine
di D\&D: i due autori avrebbero avuto visioni opposte, uno più libera e meno
regolamentata, partendo dai giochi di guerra, l'altro, che ha portato il gioco a
livello commerciale, più strutturata, dettagliata e pesante. ``È giusto, è un
riassunto corretto, ma è un po' più complicato di così'', e lo si vedrà nella
lezione sulla fiction. Ruggero aggiunge che D\&D ``non è l'unica origine del
gioco di ruolo'': il gioco di ruolo diventa \emph{famoso} con D\&D, ma gli stessi
autori si ispirarono alle \textbf{simulazioni militari}, molto diffuse
nell'America di quegli anni, e al \textbf{reenactment storico} (``ci sono due
episodi della \emph{Signora in giallo} sul reenactment della guerra civile, un
argomento allucinante''). Esistono esempi storici, dell'Ottocento e perfino del
Seicento, di \textbf{giochi di ruolo} veri e propri, con quest, avventure e
personaggi. Il \emph{mainstream} ha normalizzato l'idea che D\&D sia il primo
gioco di ruolo, ``ma non è così'': nel manuale originale l'espressione
\emph{role-playing game} non compare da nessuna parte, anche perché Gygax
odiava il termine. ``Un vaso di Pandora che non apriamo''.

Daniele però raccoglie l'osservazione. Nei primi anni editoriali di D\&D il
problema maggiore è stato il \textbf{simulazionismo}: interpretare
``simulazione'' come la necessità di avere una regola per ogni cosa che esiste nel
realismo del gioco. Ne è venuto un gioco molto strutturato, in cui ``ogni cosa
deve avere una regola'', perché ciò che non è coperto dalle regole il giocatore
può usarlo come \emph{loophole}. E ne è venuta anche l'idea che il master sia il
\textbf{padrone del tavolo} e che il giocatore debba in qualche modo
\textbf{vincere contro il master}: un altro ``vaso di Pandora enorme''.

\paragraph{Non esiste solo D\&D.} È per questo che Daniele non si concentra solo
su D\&D: ``l'idea con cui dovete uscire da qui è che non esiste solo D\&D''.
Anzi, oggi D\&D è \textbf{in controtendenza}, non per motivi di design ma per
\textbf{motivi editoriali}: il problema principale di D\&D come proprietà
intellettuale ``è sempre stato editoriale'', e di recente la Wizards of the Coast
ha fatto un po' di confusione con il \emph{branding}.

% ---------------------------------------------------------------------
\section{Story game e giochi di ruolo sperimentali}
\label{a2:sec:gdr-story}

\subsection{Gli story game}

\slidefig{img/AGD/agd2-storygames}{Alcuni story game.}{slide \emph{Analog -- Case Studies}}{fig:agd2-storygames}


La slide successiva è sugli \textbf{story game}, che in realtà nascono anche
prima di The Forge. ``Story game è un termine che io odio, perché vuol dire tutto
e non vuol dire niente'': sono giochi di ruolo da tavolo che si focalizzano
\textbf{ancora di più su una storia precisa}.
\begin{itemize}
  \item \textbf{\emph{Alice is Missing}}: la storia di Alice, scomparsa, e dei
        suoi amici che cercano di capire che cosa le è successo. Si gioca con
        una \textbf{chat} sul telefono, con un video a schermo, e con
        \textbf{due ore di tempo}.
  \item \textbf{\emph{Microscope}}: ``letteralmente \emph{world building}
        narrativo''.
  \item \textbf{\emph{Bleak Spirit}}: giocare un \emph{Dark Souls} in cui tutti
        controllano a turno il viandante protagonista.
  \item \textbf{\emph{The Zone}}\footnote{La trascrizione attribuisce il gioco a
        ``Ruff D'Amico''; il nome non è verificabile.}, giocabile anche con un
        supporto online: si narra, di base, \emph{Annientamento} (``se avete
        visto il film''). Si entra in una zona in cui una presenza aliena sta
        cambiando il mondo, e \textbf{moriranno tutti tranne una persona}, che
        non si sa in anticipo chi sarà. È un'influenza del LARP: come in
        \emph{Demeter}, si sa che il proprio personaggio morirà.
  \item \textbf{\emph{Archipelago}} e \textbf{\emph{For the Queen}}, e altri
        giochi che \textbf{tolgono i dadi} e mettono le carte, o tolgono
        perfino la \textbf{scheda del personaggio}: in \emph{Microscope} ci sono
        solo dei post-it.
\end{itemize}

\subsection{The Quiet Year e i map drawing game}

\slidefig{img/AGD/agd2-quietyear}{Il flusso di The Quiet Year.}{slide \emph{Analog -- Case Studies}}{fig:agd2-quietyear}


``C'è un minimo di suspense'': si arriva a esperienze ancora più lontane.
In \textbf{\emph{The Quiet Year}} non si interpreta nemmeno un personaggio: si è
parte di una \textbf{comunità post-apocalittica} e si gioca \textbf{un anno} della
sua vita. Non si può discutere di ciò che si gioca. A ogni turno si gira una
\textbf{carta} che dice che cosa succede, e durante il proprio turno si può fare
una tra due o tre azioni (disegnare, iniziare un progetto\ldots). Tutto si gioca
su una \textbf{mappa} su cui si disegna ciò che si vuole intendere: c'è chi gioca
in modo più rigoroso, senza poter nemmeno scrivere parole sulla mappa, e chi lo
permette. È ``un gioco fenomenale'', perché fa capire che il gioco di ruolo può
essere anche altro; e ha normalizzato un genere, quello dei \textbf{map drawing
role-playing game}, i giochi di ruolo in cui si disegnano mappe.

\subsection{Dialect}

\slidefig{img/AGD/agd2-dialect}{Il flusso di Dialect.}{slide \emph{Analog -- Case Studies}}{fig:agd2-dialect}


E poi \textbf{\emph{Dialect}}. Qui si ha un personaggio, ma è solo un nome e un
ruolo, un archetipo. Si gioca \textbf{la morte di una comunità attraverso la
morte di una lingua}: una comunità umana isolata dalle altre, che ha creato il
proprio linguaggio, e che giocherà la fine di quell'isolamento attraverso la
morte del linguaggio. A ogni turno si gioca una carta che è un \textbf{concetto
della lingua} (per esempio: come si dice ``amico'' nella nostra lingua?); si
stabilisce la parola e si gioca una breve scena in personaggio, senza tiri di
dado né altro; quando la parola viene usata in personaggio, chiunque al tavolo
può dire stop: ``abbiamo usato la parola, va bene così'', e si va avanti.

Non è un gioco ``non complesso'', come diceva Ruggero: \emph{Dialect} è
``chiarissimo'' da leggere, ma la complessità sta nel \textbf{design}, un design
elegante che crea una \emph{one shot} con ``una profondità'' enorme. Daniele con
\emph{Dialect} ha pianto; e quando incontra gli amici con cui lo ha giocato anni
fa, si parlano ancora nella lingua inventata allora: ``ehi tu, oceano'', e chi
c'era sa che cosa vuol dire. Anche qui si vede come le esperienze nate dal Nordic
LARP siano \textbf{``pervase''} (``rubo una parola a Beatrice'') in altri ambiti.

\subsection{Servono davvero le categorie? Il problema editoriale}

Ci sono ancora esperienze più tradizionali; ``il mondo è molto variegato'', e si
parlerà di \textbf{culture}: spesso chi si avvicina al gioco di ruolo vuole
giocare una certa cosa, una cultura di gioco con titoli che ``vanno bene'' e
titoli che no. L'ultima slide è ``una piccola critica'': è davvero necessario
mettere le categorie ai giochi di ruolo? Da un lato sì, dall'altro no, perché i
problemi principali sono due, e il primo è \textbf{il mondo editoriale}: i giochi
di ruolo sono stampati come \textbf{libri}, e l'editore ha ``una vita travagliata''
con il gioco. Esempi:
\begin{itemize}
  \item \textbf{\emph{In the Name of God}}\footnote{Titolo e autore (``Alessandro
        Pirodi'') come nella trascrizione; non verificabili.}, stampato solo in
        \textbf{self-publishing} e assente dal mercato, ``ma bellissimo''; la sua
        meccanica di risoluzione chiede di fare qualcosa fisicamente al tavolo
        (la trascrizione dice ``fare home'');
  \item un'\textbf{espansione per qualsiasi gioco di ruolo}\footnote{La
        trascrizione dice ``Mew Wild in the Subway''.};
  \item un \textbf{GDR stampato su una maglietta}, di due autori che Daniele
        saluta;
  \item un gioco di \textbf{cultisti contro investigatori} che prende da
        esperienze di chamber LARP\footnote{La trascrizione lo chiama ``SOT''.}.
\end{itemize}

\begin{figure}[htbp]
\centering
\begin{tikzpicture}[
    font=\footnotesize,
    epoca/.style={draw, thick, rounded corners=3pt, fill=gray!12,
                  text width=2.55cm, align=center, minimum height=1.0cm},
    nota/.style={align=left, text width=2.7cm, font=\scriptsize},
    freccia/.style={-{Stealth[length=2.5mm]}, thick}
  ]
  \draw[freccia] (-0.2,0) -- (15.4,0);
  \foreach \x/\a in {0.9/{anni '50--'70}, 4.0/{'74--anni '90},
                     7.1/{fine '90--2000}, 10.2/{anni 2010}, 13.3/{post-COVID}}{
    \draw[thick] (\x,0.12) -- (\x,-0.12);
    \node[below=0.15cm, font=\scriptsize\itshape] at (\x,0) {\a};
  }
  \node[epoca] (e1) at (0.9,1.2)  {\textbf{Wargame e miniature}};
  \node[epoca] (e2) at (4.0,1.2)  {\textbf{D\&D e GDR ``early''}};
  \node[epoca] (e3) at (7.1,1.2)  {\textbf{New Wave / The Forge}};
  \node[epoca] (e4) at (10.2,1.2) {\textbf{Giochi focalizzati, story game}};
  \node[epoca] (e5) at (13.3,1.2) {\textbf{OSR e ritorno all'origine}};
  \node[nota, anchor=north] at (0.9,-0.65) {simulazioni militari,
    reenactment; \emph{Chainmail} (1971)};
  \node[nota, anchor=north] at (4.0,-0.65) {Gygax e Arneson;
    master arbitro; tutti i generi (\emph{Cthulhu}, \emph{Vampire},
    \emph{Traveller}, \emph{GURPS}); IP adattate};
  \node[nota, anchor=north] at (7.1,-0.65) {il sistema conta;
    \emph{Apocalypse World}, \emph{Nobilis}, \emph{Dread}, \emph{Fate}; Big
    Model; \emph{Fiasco}};
  \node[nota, anchor=north] at (10.2,-0.65) {\emph{Blades in the
    Dark}, \emph{Fabula Ultima}, \emph{Not the End}; \emph{The Quiet Year},
    \emph{Dialect}; due giocatori, solo, micro};
  \node[nota, anchor=north] at (13.3,-0.65) {meno regole, alta
    mortalità, sandbox; \emph{Mörk Borg}, \emph{Shadowdark}};
\end{tikzpicture}
\caption{Una ricostruzione schematica della ``piccola storiografia'' dei giochi
di ruolo da tavolo raccontata da Daniele. Le epoche sono indicative e si
sovrappongono (gli story game, per esempio, nascono anche prima di The Forge);
la storia completa è rimandata alla lezione sulla fiction.}
\label{a2:fig:storia-gdr}
\end{figure}

% ---------------------------------------------------------------------
\section{Giochi di ruolo per due, in solitario e micro}
\label{a2:sec:gdr-formati}

\slidefig{img/AGD/agd2-micro}{Giochi di ruolo per due, in solitario e nano-RPG.}{slide \emph{Analog -- Case Studies}}{fig:agd2-micro}


L'ultima parte sui giochi di ruolo torna a Ruggero, che presenta tre formati.

\begin{description}
  \item[Two player only.] Storicamente i giochi di ruolo nascono in gruppi di
        gioco ampi, ma c'è stata ``una grossissima wave'' di giochi di ruolo
        \textbf{per due persone}, e non nel senso che una fa il master e l'altra
        un personaggio: spesso l'\textbf{agency narrativa}, il potere di decidere
        che cosa succede, viene \textbf{passata meccanicamente} al tavolo o
        divisa in responsabilità diverse. Ce ne sono tantissimi, e per Ruggero
        sono molto interessanti anche perché, come si vedrà nella lezione
        sull'emotional design, permettono di fare cose molto \textbf{intime}: non
        serve un gruppo di gioco, si possono giocare con il partner o con
        un'altra persona.
  \item[Solo RPG.] Dal 2013 circa si sono normalizzati i \textbf{giochi di
        ruolo in solitario}: si gioca da soli, senza un'altra persona e senza
        fare da master a sé stessi. Spesso non c'è un master: ci sono
        \textbf{procedure} che guidano il gioco, e il manuale è il sistema di
        gioco, strutturato ``quasi come un board game'' in termini di regole. Ma
        resta tutto ciò che è tipico del gioco di ruolo: l'agency emergente,
        l'interpretazione di un personaggio, la creazione di una narrazione. Di
        solito si giocano in formato \textbf{diario} (\emph{journaling}) o in
        formato \textbf{audio}, e sono in grande tendenza.
  \item[Micro RPG.] Infine, i \textbf{micro giochi di ruolo}\footnote{La
        trascrizione dice ``handheld RPG''; dal contesto si tratta dei
        \emph{micro RPG}.}: giochi di ruolo su \textbf{biglietti da visita}, o
        raccolte in cui un libro ne contiene cento. Il \textbf{200 Word RPG}
        (\emph{200 Word RPG Challenge}) era un contest che chiedeva di scrivere un
        gioco di ruolo in \textbf{massimo 200 parole}, e ogni edizione ne
        raccoglieva migliaia: nel 2017, ricorda Ruggero, circa 900 giochi, su
        circa mille proposte.
\end{description}

% ---------------------------------------------------------------------
\section{I giochi da tavolo: i generi storici}
\label{a2:sec:tavolo}

\slidefig{img/AGD/agd2-rollmove}{Roll and move, astratti e party game.}{slide \emph{Analog -- Case Studies}}{fig:agd2-rollmove}


``Passerei ai giochi da tavolo, correndo velocemente e senza fare pause, se ce la
fate''. Come per i LARP e i giochi di ruolo, Ruggero parla di
\textbf{categorie dall'alto}, e per ognuna mostra come il genere si sia
\textbf{evoluto} nel tempo.

\subsection{Roll and move}

I \textbf{roll and move} si basano, come dice il nome, sul \textbf{tirare un
dado e muoversi}. Sono forse i più comuni nella storia dell'umanità, con quello
che da noi è il \textbf{gioco dell'oca}: un tracciato su cui si tirano i dadi e
si avanza. L'\textbf{agency} in questi giochi è ``praticamente inesistente'', e
il genere si è presto evoluto integrando altre meccaniche: con
\textbf{elementi economici} (come \emph{Monopoly}) o con \textbf{elementi
deduttivi}, come \emph{Cluedo}, che è un roll and move.

\subsection{Astratti}

Un altro genere storicamente gigantesco è quello degli \textbf{astratti}:
scacchi, dama, go, tutti quei giochi la cui componente tematica è ridotta
``quasi all'osso''. Gli scacchi mantengono ancora qualche forma di tema; il go, di
fatto, è un rettangolo di legno con delle caselle e delle pietre bianche e nere.
Anche qui il genere non si è fermato: ci sono iterazioni moderne come
\emph{ZÈRTZ} e i giochi del \textbf{progetto GIPF}\footnote{La trascrizione dice
``Zerz o il Gift Project''.}, e altri astratti contemporanei (uno è mostrato a
destra nella slide\footnote{La trascrizione non permette di ricostruirne il
titolo (``i borsa astratti, per esempio, di Supernova'').}), che mantengono
dinamiche ``molto asciutte, molto distillate'': \textbf{informazione aperta},
trasparente, niente carte imprevisti, niente mosse segrete. Diventano quasi una
\textbf{sfida d'intelletto} con l'avversario.

\subsection{Party game}

I \textbf{party game}: \emph{Taboo}, \emph{Lupus in Tabula} o \emph{Werewolf} a
seconda dell'edizione. Sono pensati per \textbf{molte persone} (di solito dalle
quattro in su), con regole ``molto asciutte, molto leggere'', perché si giocano
in contesti di festa, in cui non è detto che chi è nella stanza ami giocare o
abbia mai giocato. Poche regole, leggerezza, durata breve (spesso non più di
mezz'ora): sono \textbf{un possibile punto di accesso}. Anche questi si sono
evoluti, con giochi come \emph{Hitster}\footnote{La trascrizione dice ``Hit'';
la ricostruzione è probabile ma non certa.} e \emph{Wavelength}, che portano
sensibilità più moderne e meccaniche più interessanti del ``tirare alla cieca
per capire chi tra di noi è un lupo mannaro''.

\subsection{Giochi di parole}

I \textbf{word game}\footnote{La trascrizione dice ``war games'', ma gli esempi
(\emph{Boggle}, \emph{Scarabeo}) chiariscono che si tratta dei giochi di
parole.}: \emph{Boggle}, \emph{Scarabeo}, giochi in cui la meccanica centrale è
\textbf{creare parole}, e più la parola è lunga, complessa o usa lettere rare,
più punti si fanno. Anche qui le iterazioni sono interessanti: \emph{Cross
Clues}\footnote{La trascrizione dice ``Cross Cruise''.} trasforma il genere in
un gioco \textbf{cooperativo}, in cui non si compete per la parola più lunga ma si
gioca con le parole per capire ``quanta intesa c'è tra di noi''; e
\emph{Illiterati} è praticamente un gioco con le \emph{boss fight}, in cui si
sconfiggono i boss formando parole. Ibridazioni anche qui, perché chi crea giochi
li crea in base a ciò che ha giocato prima.

\subsection{Giochi di carte}

\slidefig{img/AGD/agd2-carte}{Giochi di parole, giochi di carte, TCG e deckbuilding.}{slide \emph{Analog -- Case Studies}}{fig:agd2-carte}


I \textbf{card game}, che in Occidente si diffondono storicamente con i mazzi
\textbf{francesi} e, in Italia, con quelli ``da briscola'' (ma ci sono anche altri
casi storici\footnote{La trascrizione cita i ``Mazda 1''; il riferimento non è
ricostruibile.}). Le evoluzioni:
\begin{itemize}
  \item i \textbf{micro card game} come \emph{Love Letter}, che in una decina o
        dozzina di carte racchiude un gioco complesso;
  \item giochi come \emph{Munchkin}, in una direzione più ampia, con il
        \textbf{take that} (``voglio mettere i bastoni tra le ruote'' agli altri);
  \item la ``gigantesca categoria'' dei \textbf{trading card game}, già
        menzionata nella lezione precedente: \emph{Magic: The Gathering},
        \emph{Yu-Gi-Oh!}, \emph{Pokémon}. Il \emph{trading} non è solo una
        dinamica tra persone: è nel nome della categoria. Scambiare carte
        \textbf{fuori dalla partita} è parte della dinamica comunitaria;
  \item i \textbf{Living Card Game} (LCG), che provano a risolvere uno dei
        principali problemi editoriali dei TCG: ``se io ho più soldi di te, tu
        contro di me a Magic non vinci'', perché comprerò carte da 800 o 900
        euro l'una. Gli LCG \textbf{non hanno booster casuali}: ogni espansione ha
        l'elenco preciso del contenuto. Restano giochi di carte con il
        \textbf{deck building pre-partita}, con mazzi ``fortemente identitari''
        (si possono scegliere le carte anche solo perché sono più simpatiche),
        ma senza l'aspetto economico della \textbf{rarità indotta};
  \item i \textbf{deck-building game}, in cui la costruzione del mazzo
        \textbf{non è più pre-partita ma avviene durante la partita}: spesso si
        parte con mazzi identici, e durante il gioco si acquisiscono carte che
        cambiano \textbf{permanentemente} il proprio mazzo, con lo scopo di
        costruire il mazzo migliore possibile.
\end{itemize}

\subsection{German game ed Eurogame}

\slidefig{img/AGD/agd2-euro}{German game, Eurogame e American.}{slide \emph{Analog -- Case Studies}}{fig:agd2-euro}


Altre categorie, che Ruggero mette tra virgolette, sono i \textbf{German game}:
un genere sviluppato in \textbf{Germania} dopo la metà degli anni '50--'60, con
giochi a \textbf{interazione indiretta}, competitivi ma ``senza troppo dar
fastidio direttamente all'altra persona'': non ci si può attaccare
direttamente. Sono giochi poco bellici, anche per \textbf{rigetto della Seconda
guerra mondiale}, e si è sviluppato moltissimo il filone \textbf{comunitario}, di
\emph{community building} e di scambio: costruire la propria fattoria, costruire
un villaggio, scambiarsi le risorse (``io ho delle risorse, tu altre,
scambiamocele''), senza distruggere ciò che l'altro ha costruito.

Il genere si è poi reso più complesso e si è espanso oltre la Germania con gli
\textbf{Eurogame}, che includono Spagna, Portogallo, Francia e anche l'Italia
(Ruggero ne ha fatti parecchi). L'interazione resta indiretta, ma i giochi
diventano molto più complessi: spesso durano \textbf{due o tre ore} e provano a
rendere \textbf{una simulazione focalizzata di qualcosa del mondo reale}. Nella
slide, a sinistra, un gioco in cui si gestisce una \textbf{galleria
d'arte}\footnote{La trascrizione dice ``The Gallery''; probabilmente \emph{The
Gallerist}.}, con tutte le dinamiche di compravendita di quadri; e
\emph{Endeavor}, che Ruggero ha alle spalle, sul commercio mondiale in epoca
moderna. Dinamiche che simulano un tipo di mercanteggiare: ``a volte compro e
vendo quadri, a volte compro e vendo vasi''.

\paragraph{Approfondimento: l'Eurogame secondo Brathwaite e Schreiber.} In
\emph{Challenges for Game Designers}, nel capitolo sul \emph{targeting} di un
mercato, gli Eurogame (``noti anche come \emph{German games} o \emph{designer
games}'') sono un caso studio di gioco per il mercato di massa: giochi per la
famiglia, godibili a tutte le età; \emph{I coloni di Catan} è il primo ad aver
reso popolare il genere in Nord America, seguito da \emph{Carcassonne} e
\emph{Puerto Rico}. Il libro ne dà una definizione operativa, utile per
confrontare le osservazioni di Ruggero:
\begin{itemize}
  \item \textbf{durata tra 20 e 90 minuti} (troppo brevi annoiano gli adulti,
        troppo lunghi i bambini);
  \item \textbf{setup breve}, perché il setup non contiene decisioni;
  \item \textbf{regole semplici}, curva di apprendimento breve;
  \item \textbf{cooperazione più che scontro}: scambi e negoziazioni invece di
        meccaniche distruttive in cui un giocatore danneggia direttamente un
        altro (il che non impedisce una competizione intensa);
  \item \textbf{decisioni strategiche con casualità misurata};
  \item \textbf{tempi morti minimi}: turni brevi, o azioni anche nei turni
        altrui.
\end{itemize}
La ``cooperazione invece dello scontro'' è l'interazione indiretta del German
game; la durata, invece, mostra bene l'evoluzione che racconta Ruggero: gli
Eurogame più complessi, da due o tre ore, sono ormai un'altra cosa rispetto al
gioco per famiglie descritto dal libro.

\subsection{American}

In contrapposizione, oltreoceano c'erano gli \textbf{American} (o
\emph{Ameritrash}): giochi molto focalizzati sulla \textbf{produzione} (le
miniature di plastica), con il \textbf{dungeon crawling}, tematicamente molto
vicini a \emph{Dungeons \& Dragons}, quasi versioni da tavolo del combattimento a
griglia normalizzato da D\&D. Continuano la tradizione di \emph{Chainmail} e dei
giochi strategici e tattici: sono spesso giochi di combattimento, a
\textbf{interazione diretta} (si attacca e ci si difende), e buona parte dei loro
regolamenti è dedicata alla simulazione del combattimento.

\subsection{Strategici e wargame}

\slidefig{img/AGD/agd2-strategici}{Strategici, movimento nascosto, destrezza, 18XX, tempo reale e legacy.}{slide \emph{Analog -- Case Studies}}{fig:agd2-strategici}


Un altro genere gigantesco è quello dei \textbf{giochi di strategia}, usati
storicamente anche nella \textbf{formazione militare}, ``non a caso'': la
\textbf{RAND Corporation} ne ha creati diversi, e ci sono editori specializzati.
Lo scopo era ricreare battaglie storiche. Anche qui, evoluzioni:
\begin{itemize}
  \item i giochi \textbf{a zone} o ad \textbf{aree} (\emph{area control}) con
        unità su una mappa, di cui \emph{Risiko} è il più famoso: io ho un
        esercito, tu ne hai un altro, dobbiamo conquistare territori, e il
        conflitto diretto è inevitabile perché gli obiettivi si sovrappongono;
  \item i \textbf{card-driven strategy game}: strategici con mappe e zone di
        controllo, ma con \textbf{mazzi di eventi}. I designer prendono eventi
        reali di conflitti reali e cercano di renderli coesi tematicamente con
        le meccaniche, dando a ogni evento una risoluzione precisa: ``l'importante
        non è tanto quante unità ho, ma quali carte evento ho in mano''. Sono,
        se ci si pensa, più vicini ai card game che agli strategici.
\end{itemize}
Nicola chiede di \textbf{\emph{Campaign for North Africa}}\footnote{La
trascrizione dice ``Campi for Nord Africa''.}: è uno \textbf{strategico
classico}, con esagoni e pedine di cartoncino, ``un'esperienza devastante in
termini di lunghezza'': forse non è il caso di cominciare da lì.

\subsection{Movimento nascosto}

Il genere degli \textbf{hidden movement}\footnote{La trascrizione dice ``Eat the
Movement''.}: \emph{Scotland Yard}, \emph{Letters from Whitechapel}, ma anche
altri titoli più recenti\footnote{La trascrizione cita ``Sniper Elite, Mind Dame
e GMT''; il secondo è probabilmente \emph{Mind MGMT}, GMT è un editore di
wargame.}. Sono giochi \textbf{asimmetrici}: di solito una persona si nasconde e
scappa, e le altre vincono collettivamente se la catturano. Uno contro tutti.

\subsection{Destrezza}

I \textbf{dexterity game} richiedono \textbf{destrezza fisica}: colpire
componenti con le schicchere, impilarli senza farli cadere. Non basta la
strategia che si ha in testa: ``posso anche avere la strategia migliore, ma se poi
le mie mani non ce la fanno a mettere in piedi questi cilindretti, crolla
tutto''.

\subsection{18XX}

Un genere che Ruggero ama particolarmente, e di cui è ``abbastanza sicuro''
che nessuno abbia mai sentito parlare, anche se conta centinaia di titoli (``più
di 400 giochi'' sotto l'etichetta): i \textbf{18XX}. Il nome viene dal secolo
delle \textbf{ferrovie}, l'Ottocento: sono tutti giochi ferroviari, ognuno
centrato su un anno e un luogo specifici (\emph{1849}, le ferrovie della
Sicilia nel 1849; \emph{1889}, quelle dell'isola di Shikoku). Hanno una forte
componente di \textbf{costruzione delle ferrovie} sulla mappa e una componente
\textbf{economico-azionaria}: comprare e vendere azioni, speculazione
finanziaria. Sono ``molto belli, ma anche molto intensi'': \emph{1889}, che
Ruggero usa di solito come introduttivo, ``dura tre ore''.

\subsection{Tempo reale}

Culturalmente si pensa ai giochi da tavolo come a qualcosa che avviene
\textbf{a turni}, ma non è sempre così. I giochi \textbf{in tempo reale} usano
clessidre, timer o altre meccaniche per cui non si ha tutto il tempo per pensare:
bisogna reagire in fretta. \emph{Kitchen Rush}, per esempio, è ``praticamente
\emph{Overcooked} da tavolo'': si gestisce un ristorante, e \textbf{la clessidra
è il lavoratore}, che non si può spostare finché non si è esaurita. Chi prende
gli ingredienti, chi li prepara sul fornello, chi va in sala a prendere gli
ordini, chi porta i piatti, chi li lava: ne nasce un certo caos nell'intreccio.

\subsection{Legacy}

I \textbf{Legacy} sono ``un po' una rivoluzione nel genere''. Di solito un gioco
da tavolo resta sempre quello: lo si prepara, lo si gioca, la settimana dopo lo
si gioca con un altro gruppo. I giochi Legacy invece chiedono di avere un
\textbf{impatto fisico sul gioco}, per cui \textbf{le partite successive sono
influenzate dalle precedenti}:
\begin{itemize}
  \item \emph{Pandemic Legacy} chiede di \textbf{strappare carte}; alcune città
        vengono distrutte dalla pandemia, coperte con degli \textbf{adesivi}, e
        non esistono più sulla mappa;
  \item \emph{SeaFall}\footnote{La trascrizione dice ``Sifol''.}, sulla scoperta
        di isole in un Ottocento alternativo: le isole hanno il nome vuoto, e chi
        le scopre le battezza; da lì in poi l'isola porterà quel nome per tutto
        il resto del gioco.
\end{itemize}
``So che qualcuno sta già sudando freddo all'idea di strappare carte o scrivere
sul tabellone'', ma sono ``esperienze emotivamente incredibili'': ``giocatevi dei
Legacy''.

\paragraph{Approfondimento: Risk Legacy e la rottura di un'assunzione.} Schell
dedica ai Legacy due passaggi di \emph{The Art of Game Design}. Nel capitolo sul
\textbf{tema}, presenta \emph{Risk Legacy} di \textbf{Rob Daviau} come esempio di
tema fondato su una verità: Daviau ha fatto ``qualcosa di mai tentato prima nel
mondo dei giochi da tavolo'', creando meccaniche per cui le scelte dei giocatori
\textbf{cambiano il gioco per sempre}. Le regole impongono di modificare il
tabellone con adesivi non rimovibili, di scriverci sopra con un pennarello
indelebile per rivendicare i territori, di strappare e scartare carte e perfino
di modificare permanentemente il regolamento. Queste meccaniche, nota Schell,
sono nuove e intriganti, ma soprattutto \textbf{rafforzano il tema centrale}:
``la guerra cambia il mondo''. Nel capitolo sulle tecniche di brainstorming,
Schell racconta come Daviau ci sia arrivato: elencando le \textbf{assunzioni}
date per scontate sui giochi da tavolo e provando a ``farne saltare'' una,
``una partita non influenza la successiva''. È un buon esempio di quello che
Ruggero chiedeva all'inizio: guardare un genere che si conosce e chiedersi che
cosa succederebbe se una delle sue regole implicite non valesse più.

% ---------------------------------------------------------------------
\section{Due libri sul gioco da tavolo}
\label{a2:sec:tavolo-libri}

\slidefig{img/AGD/agd2-libri}{I due libri consigliati sul gioco da tavolo.}{slide \emph{Analog -- Case Studies}}{fig:agd2-libri}


A proposito di risorse, Ruggero torna sul libro di Nuccio, \emph{La
progettazione dei giochi da tavolo}, citato a inizio lezione, e indica
un'alternativa in inglese che a suo parere ``fa quello che fa Nuccio, meglio'':
\textbf{\emph{Building Blocks of Tabletop Game Design}} di \textbf{Geoffrey
Engelstein e Isaac Shalev}\footnote{La trascrizione dice ``Engelstein e
Schaler'' e ``Building Blocks of Tablet of Game Design''.}. È più
\textbf{contemporaneo} negli esempi e ``più funzionale''. Per i giochi da tavolo,
se si ha già un framework di riferimento, sono entrambi ``una risorsa
incredibile''. (Resta valido l'avvertimento di inizio lezione: sono archivi di
meccaniche, da usare dopo aver capito che cosa le meccaniche fanno.)

% ---------------------------------------------------------------------
\section{Libri-gioco, megagame e altri ibridi}
\label{a2:sec:tavolo-ibridi}

\slidefig{img/AGD/agd2-ibridi}{Libri-gioco, megagame e ibridi tra gioco da tavolo, gioco di ruolo e videogioco.}{slide \emph{Analog -- Case Studies}}{fig:agd2-ibridi}


``Manca ancora qualcosa?'' Mancano diverse cose.
\begin{itemize}
  \item I \textbf{libri-gioco}, nel formato libro tradizionale e nel formato
        ``più contemporaneo'' dei giochi da tavolo che sono libri-gioco, come
        \emph{Sherlock Holmes Consulente Investigativo}.
  \item I \textbf{megagame}: giochi da tavolo \textbf{per cento persone}. Non
        hanno un'organizzazione editoriale come \emph{Monopoly} o gli altri
        board game: sono organizzati in modo più vicino ai giochi pervasivi e al
        LARP, spesso \textbf{a evento}, spesso con biglietto, anche perché
        gestire un gioco da tavolo con cento persone in parallelo ``può diventare
        molto intenso''.
  \item I \textbf{LARP urban game}, di cui ha parlato Beatrice, con esempi
        storici anche italiani, come \emph{Pathos}\footnote{Grafia come nella
        trascrizione.}: giochi di ruolo dal vivo che sono anche urban game, e
        un po' ARG.
  \item Produzioni ispirate al freeform ma vendute come giochi da tavolo, come
        \textbf{\emph{Inhuman Conditions}}\footnote{La trascrizione dice ``In
        Human Condition''.}: praticamente il \textbf{test di \emph{Blade Runner}}
        per scoprire chi è un androide, per due persone. È un gioco di ruolo, ma
        prodotto esattamente come un gioco da tavolo (carte, scatola,
        regolamento) e non si presenta come gioco di ruolo, ``ma questo non lo
        rende meno gioco di ruolo''.
  \item Ibridi con il party game, come \textbf{\emph{Sì, Oscuro Signore!}}: un
        party game, ma anche un \emph{role-playing game}, perché si fa il role
        play del proprio personaggio.
  \item Ibridi con il videogioco, come \emph{Fateforge} e gli altri giochi che
        usano un'\textbf{app} non solo come segnapunti ma come meccanica durante
        la partita.
\end{itemize}

\begin{table}[htbp]
\centering
\small
\begin{tabular}{@{}p{3.0cm}p{9.9cm}@{}}
\toprule
\textbf{Genere} & \textbf{Che cosa lo caratterizza ed evoluzioni citate} \\
\midrule
Roll and move & tira e muovi (gioco dell'oca); agency quasi nulla; poi
elementi economici (\emph{Monopoly}) o deduttivi (\emph{Cluedo}) \\
Astratti & tema ridotto all'osso, informazione aperta (scacchi, dama, go,
\emph{ZÈRTZ}, progetto GIPF) \\
Party game & molte persone, regole leggere, partite brevi, punto di accesso
(\emph{Taboo}, \emph{Lupus in Tabula}, \emph{Wavelength}) \\
Giochi di parole & creare parole per fare punti (\emph{Boggle},
\emph{Scarabeo}); cooperativi (\emph{Cross Clues}), boss fight
(\emph{Illiterati}) \\
Giochi di carte & mazzi tradizionali; \emph{Love Letter}; take that
(\emph{Munchkin}); TCG, LCG e deck-building \\
German game ed Eurogame & interazione indiretta, costruzione e scambio;
simulazioni focalizzate di due o tre ore \\
American & produzione in plastica, dungeon crawling, combattimento e
interazione diretta \\
Strategici & formazione militare; aree e zone (\emph{Risiko}); card-driven;
esagoni (\emph{Campaign for North Africa}) \\
Movimento nascosto & uno contro tutti, asimmetrici (\emph{Scotland Yard},
\emph{Letters from Whitechapel}) \\
Destrezza & schicchere, impilare: conta anche la mano \\
18XX & ferrovie e azioni dell'Ottocento; più di 400 titoli (\emph{1849},
\emph{1889}) \\
Tempo reale & clessidre e timer (\emph{Kitchen Rush}) \\
Legacy & le partite cambiano fisicamente il gioco (\emph{Pandemic Legacy},
\emph{SeaFall}) \\
Libri-gioco, megagame, ibridi & \emph{Sherlock Holmes Consulente
Investigativo}; cento giocatori a evento; \emph{Inhuman Conditions},
\emph{Sì, Oscuro Signore!}, giochi con app \\
\bottomrule
\end{tabular}
\caption{I generi di gioco da tavolo presentati da Ruggero, in sintesi.}
\label{a2:tab:tavolo}
\end{table}

% ---------------------------------------------------------------------
\section{La chiusura: niente interrogazioni, e i materiali}
\label{a2:sec:chiusura}

``E ho corso un po', ma ce l'abbiamo fatta''. Mancano gli ultimi dieci minuti, e
Ruggero si complimenta con chi è ``arrivato fino a qui''. Ripete il punto
fondamentale: \textbf{non ci sarà nessuna interrogazione} su questi giochi
(``Tiziana, ti do subito il sollievo''). Non si chiede di sapere che cosa sia ogni
caso studio: lo scopo era dare una \textbf{panoramica}, degli spunti, e far
sapere che ``questa roba esiste''. Le slide restano, e ci si potrà tornare in
fase di ricerca e sviluppo. Con la lezione di \textbf{ricerca e sviluppo}, che
darà gli strumenti per farlo in autonomia, ognuno troverà \textbf{i casi studio
che servono al proprio gioco}, ``e io ve li chiederò, noi ve li chiederemo''.
Questa lezione serviva ad \textbf{aprire le possibili direzioni}.

\paragraph{I materiali caricati.} Mentre Ruggero parla, le slide vengono
caricate nella cartella degli studenti. Beatrice aggiunge di aver caricato:
\begin{itemize}
  \item il paragrafo di \textbf{Frank Rose} su \emph{Why So Serious};
  \item un estratto sulla \textbf{definizione dei pervasive game} tratto da
        \emph{Game Design}, un libro di \textbf{Maresa Bertolo e Ilaria
        Mariani}\footnote{La trascrizione dice ``Marianne Bertolo'' e
        ``Bertolo Mariani''. Il libro è con ogni probabilità \emph{Game
        Design. Gioco e giocare tra teoria e progetto} di Maresa Bertolo e Ilaria
        Mariani.}, ``un ottimo libro''.
\end{itemize}
Non sono letture obbligatorie.

\paragraph{Un avvertimento sui libri italiani di gioco di ruolo.} Daniele
aggiunge una cautela sulla parte dedicata al gioco di ruolo del libro di Bertolo
e Mariani: alcuni paragrafi ``generalizzano molto verso il mainstream''. E poi un
\emph{warning} più generale: \textbf{la maggior parte della letteratura italiana
sul gioco di ruolo è ``terribile''}, al punto che alcuni libri i docenti non solo
non li consigliano, ma li \textbf{sconsigliano attivamente}. Con la premessa che
non è nulla di personale contro l'autore, Daniele sconsiglia il manuale
\emph{Giochi di ruolo} di \textbf{Mauro Longo}\footnote{La trascrizione riporta
l'editore come ``Dino Dino Editori'': probabilmente Dino Audino Editore.},
perché è molto soggettivo, molto legato alla visione del gioco di ruolo
dell'autore e a un solo tipo di gioco, quello di \emph{Lex Arcana}\footnote{La
trascrizione dice ``l'ex arcana''.}, con teorie proprie che non corrispondono al
design generale dei giochi di ruolo. ``Va bene se volete fare \emph{Lex Arcana}
esattamente come \emph{Lex Arcana}'', non se si vuole fare altro.

Meglio il volume sull'edu-larp, e soprattutto \textbf{giocare}: ``giocate i giochi,
perché in Italia abbiamo una creatività sfrenata''. Si trovano ``robe favolose''
online: tantissimi \textbf{one page game} di ruolo, il mondo del \textbf{print and
play} anche nei giochi da tavolo, chamber LARP e freeform da stampare
direttamente da internet. ``Fatevi una ricerca e vedete quello che vi garba''.

% ---------------------------------------------------------------------
\section{Le domande finali}
\label{a2:sec:domande}

``Siamo stati più veloci dell'anno scorso'' (quando si erano accumulati trenta
minuti di ritardo), e restano sette minuti per le domande.

\subsection{Che cosa fa scegliere una forma?}

Gianmarco, ``un po' per rompere il ghiaccio'', fa una domanda ``molto generica'',
che sa non avere una risposta semplice: si sono visti giochi dalle forme
diversissime, ma i \textbf{framework} per crearli sono più o meno gli stessi. Che
cosa determina la scelta di una forma rispetto a un'altra? Ruggero rimanda la
risposta alla lezione sull'\textbf{emotional design}, ``che è proprio una risposta
a questa cosa'': la risposta è così lunga che ``abbiamo una lezione di tre ore per
rispondere alla tua domanda''.

\subsection{Che cosa va per la maggiore in Italia?}

Angelo chiede quale genere ``vada per la maggiore'' in Italia. Ma che cosa
significa ``andare per la maggiore''? Il successo commerciale non è uguale per
tutti i generi, ``nemmeno nei videogiochi''. Uno dei docenti (con ogni probabilità Daniele) risponde con due domande
``provocatorie'', che valgono anche per i libri (ne parlava quel giorno con
un'amica narrative designer di videogiochi):
\begin{enumerate}
  \item quante copie deve vendere un gioco per essere un successo?
  \item il numero di copie vendute coincide con il successo
        commerciale ed economico?
\end{enumerate}
La risposta a entrambe è \textbf{no}, ma si capirà con la lezione di
\textbf{Flavio} sull'aspetto produttivo: quando si ragiona in logica commerciale
si è ``in logica di prodotto'', e ci sono varie \textbf{metriche} e varie scale di
successo. A questo si aggiungono altri livelli:
\begin{itemize}
  \item \textbf{quale mercato}? In Italia \textbf{i giochi che non sono in
        italiano non vendono}, e molti editori italiani sono innanzitutto
        \textbf{editori di localizzazione}: prendono giochi dall'estero, li
        traducono, e questo è buona parte della loro sopravvivenza, nel gioco
        da tavolo come nel gioco di ruolo. Un gioco in inglese si vende in
        Inghilterra, in America e in moltissimi paesi dove l'inglese è la seconda
        lingua; un gioco in italiano si vende solo in Italia, e un editore
        italiano non ha facilmente la struttura per arrivare all'estero. Per
        questo anche autori italiani di giochi da tavolo pubblicano con
        \textbf{editori esteri};
  \item \textbf{perché si fa il gioco}? Non è detto che lo scopo sia
        commerciale. Un gioco può essere \textbf{formativo o didattico}, fatto
        con un'azienda o una scuola, o \textbf{commissionato} da un cliente per
        uno scopo specifico: allora non interessa che venda mille copie. Il
        designer viene pagato, ma il pagamento non corrisponde alle vendite.
\end{itemize}

Beatrice chiede ad Angelo perché faccia la domanda, perché i \textbf{settori}
appena visti spesso ``non si parlano minimamente'', e alcuni non sanno nemmeno
dell'esistenza degli altri: la cross-medialità presentata nella lezione non viene
presentata spesso così. Se la domanda nasce da prospettive lavorative, come per
alcuni nel corso, la risposta è che ci sono ``immense possibilità'' a seconda
della direzione: fare un gioco da tavolo è una cosa, lavorare nella
\textbf{formazione} un'altra, e c'è un intero ambito non trattato, quello dei
\textbf{giochi per le aziende}.

Angelo precisa: voleva solo sapere \textbf{che cosa giocano di solito gli
italiani}. È ``un altro volo pindarico'': in Italia è difficile trovare dati
statistici affidabili, perché i questionari e le survey sono spesso invalidati
dal modo in cui sono formulate le domande. Da quello che sa il docente che risponde\footnote{La trascrizione dice ``poi Lucio mi correggerà''; il nome non è ricostruibile con certezza.}, ci sono sicuramente \textbf{più giocatori di giochi da
tavolo che di giochi di ruolo}: il gioco da tavolo è più conosciuto e
frequentato, anche nelle ludoteche, mentre il gioco di ruolo e il LARP richiedono
un investimento economico diverso e sono soprattutto pratiche
\textbf{comunitarie}, ``tribali'', che nascono dal basso e si trasmettono
\textbf{oralmente}. Chi gioca ai giochi da tavolo li compra; chi gioca di ruolo
spesso lo fa perché qualcuno lo fa giocare, e non è detto che abbia il gioco in
casa.

\subsection{Avvisi}

Intanto Daniele ricorda che il \textbf{questionario} della prima slide va
compilato \textbf{entro il 15}, e che Giordano deve ancora presentarsi (``siamo
delle pessime persone''). Visto che è ora di cena, Giordano si presenterà come
\textbf{prima cosa} alla lezione successiva, \textbf{lunedì}, sull'emotional
design (``ricordacelo, perché se no ci dimentichiamo'').

% ---------------------------------------------------------------------
\section{Due cimiteri: i giochi di carte collezionabili e i ``fix''}
\label{a2:sec:cimiteri}

Prima di chiudere, a un docente torna in mente ``una cosa importantissima''
che voleva dire.

\subsection{Non fate Magic: The Gathering 2}

``Vi prego, non cadete nella trappola'' di vedere un gioco di successo commerciale
e farne uno simile pensando che venderà moltissimo. ``Ci possiamo fare un
cimitero, una fossa comune'' con migliaia di \textbf{trading card game} e
\textbf{collectible card game}. ``Non fate un trading card game. Ve lo dico per
voi: non fate \emph{Magic: The Gathering 2}''. Chi vede \emph{Magic} o
\emph{Pokémon} pensa di poter fare un sacco di soldi, ma \textbf{non vede la
struttura commerciale e produttiva gigantesca} che c'è dietro, quasi impossibile
da emulare dal basso. Quasi tutti, anche editori grossi, quelli che hanno provato
a competere con \emph{Pokémon}, \emph{Yu-Gi-Oh!} e \emph{Magic} ``sono andati a
picco''. ``Se volete perdere soldi, quello è il modo perfetto''.

Un collega aggiunge che c'è un \textbf{revival dei TCG}, ma ogni titolo dura un
anno o due e poi muore. Si sperimentano anche dei \textbf{micro TCG} annunciati
fin dall'inizio come limitati (``questo gioco avrà tre espansioni e poi
chiuderà''), una sorta di ibrido con gli LCG. Ma attenzione: spesso sono fatti
da ex designer di LCG, di \emph{Magic} o di \emph{Hearthstone}, persone con
contatti, entourage e un'azienda alle spalle; da fuori si vede una persona, ma
dietro c'è una struttura che non si può replicare.

\paragraph{Approfondimento: perché Pokémon non si replica.} \emph{The Art of
Game Design} offre un buon esempio della ``struttura gigantesca'' di cui parlano
i docenti. Nel capitolo sui mondi transmediali, Schell racconta che il gioco di
carte collezionabili di \emph{Pokémon} nacque nel 1999 da una collaborazione tra
Nintendo e Wizards of the Coast, la casa di \emph{Magic}, e che fu progettato per
restare il più possibile fedele alle meccaniche del videogioco per Game Boy. Il
gioco di carte era la \textbf{terza porta d'ingresso} in un unico mondo, accanto
al videogioco e alla serie TV: portatile, sociale, economico, perfetto per lo
slogan ``acchiappali tutti''. Le tre porte si rafforzavano a vicenda (voler
vincere spingeva a guardare la serie, e la serie rendeva il gioco più vivido), e
Schell ricorda un dirigente dell'intrattenimento che nel 1999 profetizzava la
fine della ``moda'' con l'uscita del film. Non andò così, appunto, perché il
successo non stava in un singolo prodotto ma nell'ecosistema. Un gioco di carte
collezionabili nato dal basso, senza un ecosistema del genere, compete con questo.

\subsection{Non fate un gioco per ``aggiustarne'' un altro}

Un altro docente (probabilmente Daniele, che cita ``le lezioni di framework con Ruggero'') aggiunge un altro cimitero: quello di chi vuole fare un gioco per
\textbf{``fixare'' un gioco esistente}, cioè farne una versione migliore. Come
esperimento può essere interessante, ma, come si vedrà nella lezione
sull'emotional design e nelle lezioni sui framework con Ruggero, l'intenzione deve
essere ``\textbf{eticamente onesta}''. Non si vuole aggiustare D\&D per il piacere
di aggiustare D\&D: bisogna chiedersi sempre \textbf{che cosa non mi piace
dell'esperienza} di D\&D, che cosa voglio aggiustare di quell'esperienza. E non
creare ``D\&D dove tolgo una regola e ne metto un'altra, ed è sistemato'', ma
\textbf{partire dall'esperienza che si vuole raggiungere}, quella che non si è
avuta da D\&D, e costruirla. Con la lezione di ricerca e sviluppo si vedrà anche
come prendere spunto dai giochi esistenti, non per copiarli ma per capire che
cosa fanno, perché \textbf{prenderne pezzi senza capire che cosa fanno} è proprio
l'errore da evitare (lo stesso del libro di Nuccio letto troppo presto).

\paragraph{I saluti.} Un docente deve scappare e saluta; gli altri ricordano il
questionario, che le loro email sono nelle slide per eventuali domande, e
l'appuntamento a lunedì.

% ---------------------------------------------------------------------
\section*{Punti chiave del capitolo}
\addcontentsline{toc}{section}{Punti chiave del capitolo}

\begin{itemize}
  \item Un \textbf{caso studio} è un gioco, o una categoria di giochi, che,
        analizzato tenendo conto del suo contesto, dà informazioni e metodi
        utili per i propri giochi: come leggere libri per imparare a scrivere.
        Questa lezione è una panoramica; gli strumenti di analisi arriveranno con
        la ricerca e sviluppo. ``Per me, questo è design'' (Ruggero).
  \item \textbf{Gioco analogico} è un termine ombrello per tutto ciò che non è
        videogioco. Le quattro \textbf{macrocategorie} sono gioco da tavolo,
        gioco di ruolo da tavolo, gioco di ruolo dal vivo e gioco pervasivo, ma
        i confini sono porosi: \textbf{ibridazione e sperimentazione}. Il corso è
        \textbf{agnostico di genere}, e il project work è libero, ma non deve
        essere un gioco enorme.
  \item Il \textbf{gioco pervasivo} usa luoghi e oggetti della vita quotidiana
        come luoghi e oggetti di gioco; in termini di game studies,
        \textbf{espande il cerchio magico} nello spazio, nel tempo e nella
        società. Casi: \emph{Pokémon Go}, geocaching, \emph{GuruShot},
        \emph{Journey to the End of the Night}, SF0, \emph{Critical City Upload},
        \emph{Sleep No More}, gli ARG (\emph{Why So Serious}, \emph{The
        Institute}, \emph{Stranger Things} a Milano), Collettivo Cinetico,
        Amigdala, \emph{MP3 Experiment}, \emph{PacManhattan}, \emph{Zombies,
        Run!}, \emph{Reptilarium}, \emph{Echi resistenti}.
  \item Esperienze emotivamente significative si possono fare \textbf{con budget
        zero}: l'audio che guida, gli oggetti da portare, un finale
        scenografico.
  \item Il \textbf{LARP} è separato dal GDR per ragioni di industria (evento con
        biglietto contro oggetto venduto). Tipi: boffer, Mind's Eye Theatre,
        chamber/parlor, American freeform, murder party, Nordic LARP (Jeepform,
        360°, tape, black box).
  \item Il \textbf{Nordic LARP} ha documentato tutto (Knutpunkt dal 1997, libri
        annuali, manifesti come il Dogma 99), punta su coinvolgimento, emozione,
        coesione tematica e dopo-gioco, e ha normalizzato le \textbf{meccaniche di
        sicurezza}. Non è un esperimento scientifico: lo scopo resta ludico.
  \item Nel \textbf{GDR da tavolo} contano la fiction emergente e quattro parole:
        \textbf{conflitto, azione, risoluzione, effetto}. Il master è un
        giocatore. Dalla nascita di D\&D (1974) dal wargame, attraverso gli
        adattamenti a nastro degli anni '80--'90, si arriva alla \textbf{New
        Wave} di The Forge (``il sistema conta''), ai \textbf{giochi
        focalizzati} (non semplificati), all'OSR, agli story game, ai giochi per
        due, in solitario e micro.
  \item Nel \textbf{gioco da tavolo} ogni genere si è evoluto: roll and move,
        astratti, party game, giochi di parole, giochi di carte (TCG, LCG,
        deck-building), German game ed Eurogame, American, strategici,
        movimento nascosto, destrezza, 18XX, tempo reale, \textbf{Legacy},
        libri-gioco, megagame, ibridi.
  \item Libri: il volume sull'\textbf{edu-larp} sì; Nuccio e Engelstein--Shalev
        dopo l'emotional design; Bertolo--Mariani con cautela sulla parte di GDR;
        gran parte della letteratura italiana sul GDR no. Soprattutto,
        \textbf{giocare}.
  \item Il \textbf{successo} non si misura solo in copie, e dipende dal mercato
        (in Italia vende solo ciò che è in italiano) e dallo scopo del gioco.
  \item Due cimiteri: \textbf{non fare un gioco di carte collezionabili} e non
        fare un gioco per \textbf{``aggiustarne'' un altro} senza partire
        dall'esperienza che si vuole creare.
\end{itemize}

%
%  Tutte le etichette di questo capitolo hanno il prefisso "a2:"
%  (a2:cap:..., a2:sec:..., a2:fig:..., a2:tab:...). Il capitolo è
%  autonomo: non contiene \ref verso altri capitoli.
%  Le figure e le tabelle sono ricostruzioni schematiche di quanto i
%  docenti spiegano a voce e con le slide, non copie delle slide.
% =====================================================================

\chapter[I casi studio dell'Analog Game Design]{I casi studio: giochi
pervasivi, giochi di ruolo dal vivo, giochi di ruolo da tavolo e giochi da
tavolo}
\label{a2:cap:casi}

Questa è la \textbf{seconda lezione} del corso di \textbf{Analog Game Design}, e
mantiene la promessa fatta alla fine della prima: è una lezione di tre ore
dedicata quasi interamente ai \textbf{casi studio} (\emph{case studies}). I
docenti la definiscono ``un overview'', una panoramica ``a volo d'aquila'':
``voleremo in tutti i settori ludici analogici e guarderemo che cosa c'è''.
Non si scende nei dettagli tecnici dell'analisi (ci saranno lezioni apposite,
soprattutto quella di \textbf{ricerca e sviluppo}), ma si apre la ``diga'' degli
esempi: decine di giochi, eventi, esperienze, movimenti e manifesti, scelti per
mostrare quante cose esistono, quanto sono diverse tra loro e quanto i confini
tra i generi siano porosi.

La lezione si svolge \textbf{online}, su Zoom, un \textbf{venerdì} sera (``che
bella la lezione del venerdì! Sembra quasi un aperitivo''), ed è divisa tra i
tre docenti principali del corso, ciascuno sul proprio territorio:
\textbf{Ruggero} introduce l'argomento e presenta i \textbf{giochi di ruolo dal
vivo} (LARP) e i \textbf{giochi da tavolo}; \textbf{Beatrice} presenta i
\textbf{giochi pervasivi}; \textbf{Daniele} presenta i \textbf{giochi di ruolo
da tavolo}. ``Ce lo smezziamo un po' tutti'', perché l'argomento è diviso ``in
tante categorie, sottocategorie, sotto-sottocategorie''. Dopo i pervasivi c'è
una pausa di pochi minuti; gli ultimi dieci minuti sono lasciati alle domande.

Tre avvertenze di lettura, le stesse del capitolo precedente. La prima: la
trascrizione è automatica e storpia moltissimi nomi propri, soprattutto i titoli
dei giochi, quasi sempre in inglese e pronunciati velocemente (``Blaze in the
Dark'', ``Cobo Bobby'', ``Sifol'', ``Jeep Form''\ldots). Dove la ricostruzione è
sicura, o molto probabile, il capitolo usa la forma corretta e lo segnala in
nota; dove non lo è, riporta la forma della trascrizione e lo dice. La seconda:
durante la pausa la trascrizione registra, ripetute in loop per decine di righe,
alcune frasi senza senso (anche all'inizio della registrazione succede qualcosa
di simile), che il capitolo ignora. La terza: i docenti si passano spesso la
parola senza che la trascrizione segni il cambio di voce; il capitolo attribuisce
le parole a chi le pronuncia quando il contesto lo rende chiaro (chi parla di sé,
chi viene chiamato per nome, chi ``passa la parola''), e resta generico
altrimenti.

Quanto alla data, la lezione si colloca con ogni probabilità \textbf{venerdì 11
settembre 2026}\footnote{La data non è mai detta. La si deduce da alcuni
indizi: la prima lezione si era chiusa con ``ci si vede questo venerdì''; uno
studente assente ha recuperato ``la lezione di lunedì''; a fine lezione si
ricorda che il questionario iniziale va compilato ``entro il 15'' (settembre) e
che la lezione successiva, sull'\emph{emotional design}, sarà di lunedì. Nel
2026 il primo venerdì utile prima del 15 settembre, dopo il lunedì della prima
lezione, è l'11. È una deduzione, non un'informazione esplicita.}: i riferimenti
a ``l'altra volta'' e ``l'anno scorso'' vanno letti rispetto a quel momento.

\paragraph{Come leggere questo capitolo.} Più ancora della lezione precedente,
questa è un \textbf{catalogo}. I docenti lo dicono esplicitamente: non ci sarà
nessuna interrogazione sui singoli giochi, e lo scopo non è ricordarli tutti ma
sapere che esistono, avere dei punti di riferimento e, soprattutto, abituarsi a
guardare oltre il proprio genere preferito. Per questo il capitolo riporta
\emph{tutti} i casi citati, con quello che i docenti ne dicono, ma li organizza
in tabelle e schemi per renderli consultabili ``in fase di ricerca e sviluppo'',
esattamente come fa Ruggero con le sue slide (``non sapete quante volte mi sono
trovato a riaprire proprio esattamente queste slide''). Qua e là, dei paragrafi
di \textbf{approfondimento} collegano i casi ai libri di game design di
riferimento (\emph{The Art of Game Design} di Jesse Schell, \emph{Challenges for
Game Designers} di Brenda Brathwaite e Ian Schreiber, \emph{21st Century Game
Design} di Chris Bateman e Richard Boon).

Il percorso del capitolo segue quello della lezione:
\begin{enumerate}
  \item l'avvio, Discord, i materiali e i \textbf{libri consigliati} (e
        sconsigliati) per cominciare (sezioni~\ref{a2:sec:avvio}
        e~\ref{a2:sec:libri});
  \item \textbf{perché} analizzare i casi studio e \textbf{che cosa} sono
        (sezione~\ref{a2:sec:perche});
  \item le \textbf{quattro macrocategorie} del gioco analogico, l'ibridazione e
        l'invito a non chiudersi nel proprio genere (sezioni~\ref{a2:sec:macro}
        e~\ref{a2:sec:aperti});
  \item i \textbf{giochi pervasivi}, con Beatrice: la definizione, il cerchio
        magico e diciotto casi studio, dalle app agli \emph{alternate reality
        game}, dal teatro interattivo alle cacce al tesoro
        (sezioni da~\ref{a2:sec:pervasivi} a~\ref{a2:sec:pervasivi-edu});
  \item i \textbf{giochi di ruolo dal vivo}, con Ruggero: boffer, Mind's Eye
        Theatre, chamber e parlor LARP, American freeform, murder party, il
        \textbf{Nordic LARP} con i suoi manifesti, lo Jeepform, il 360°, i tape
        e i black box LARP (sezioni da~\ref{a2:sec:larp} a~\ref{a2:sec:oppresso});
  \item i \textbf{giochi di ruolo da tavolo}, con Daniele: le quattro parole
        (conflitto, azione, risoluzione, effetto), la storia dal primo
        \emph{Dungeons \& Dragons} alla New Wave di The Forge, l'OSR, i giochi
        focalizzati, gli story game, i giochi per due, in solitario e i micro
        giochi di ruolo (sezioni da~\ref{a2:sec:gdr} a~\ref{a2:sec:gdr-formati});
  \item i \textbf{giochi da tavolo}, di nuovo con Ruggero: diciassette generi,
        dal \emph{roll and move} ai Legacy, e gli ibridi
        (sezioni da~\ref{a2:sec:tavolo} a~\ref{a2:sec:tavolo-ibridi});
  \item la chiusura, i materiali caricati, le \textbf{domande finali} e due
        avvertimenti importanti: non fare un gioco di carte collezionabili, non
        fare un gioco per ``aggiustare'' un altro gioco
        (sezioni da~\ref{a2:sec:chiusura} a~\ref{a2:sec:cimiteri}).
\end{enumerate}

% ---------------------------------------------------------------------
\section{L'avvio della lezione}
\label{a2:sec:avvio}

La registrazione parte mentre i docenti si stanno ancora collegando (``mamma
mia, la burocrazia italiana'') e si controlla che il simbolo REC lampeggi.
Mentre Daniele prova a condividere le slide (``aspettate, sono un attimo
bloccato'') e chiede di essere nominato co-organizzatore della riunione, gli
studenti entrano uno alla volta: Bea, Lorenzo, Nicola, Rita, Gianmarco,
Tiziana, Giordano. L'atmosfera è rilassata. Si scherza sul fatto che la lezione
del venerdì ``sembra quasi un aperitivo'': un docente si mette ``bello comodo'',
un altro ha a casa ``le cose per fare l'aperitivo'' e chiede il permesso di
prepararlo nel frattempo, ``purché siate sul pezzo''. Daniele scrive anche di
dover ancora rispondere a un messaggio di Beatrice; ``non preoccuparti'', gli
risponde lei: settembre è un mese di ripartenza, ``in realtà è l'inizio
dell'anno reale''.

\paragraph{Giordano.} Arriva \textbf{Giordano}, che era assente alla prima
lezione. Ha recuperato la registrazione di lunedì, e sa che gli si chiederà di
presentarsi. (Alla fine, per mancanza di tempo, la presentazione verrà rinviata
alla lezione successiva: si veda la sezione~\ref{a2:sec:domande}.) All'inizio
mancano ancora Lorenzo e Tancredi; Lorenzo arriva poco dopo, Tancredi dopo la
pausa.

\paragraph{Discord.} Uno dei docenti nota che gli studenti sono ``belli attivi
su Discord'' e se ne rallegra: ``sfruttatelo, perché l'anno scorso lo si
sfruttava poco e invece è uno spazio molto importante''. I docenti leggono il
server ogni tanto, e le domande sono benvenute. L'unica richiesta è
\textbf{taggarli}: ``tra tanti server Discord non riusciamo a leggere'' tutto,
quindi una domanda rivolta a loro va segnalata con un tag. Lo spazio, ricorda,
``è fatto per voi'': per confrontarsi, condividere cose e materiali.

\paragraph{Le dispense.} Gli studenti hanno visto i materiali caricati: le
slide, le dispense ``e quant'altro''. Qualcuno chiede se le dispense siano in
inglese e se siano articoli. Sono i \textbf{saggi citati nella prima lezione}
(per esempio quelli sull'MDA e sul DPE), e \textbf{non è obbligatorio leggerli}:
sono lì ``per completezza'' e per mostrare che dietro gli argomenti delle
lezioni ``c'è una base teorica''. L'importante è che le lezioni facciano capire
gli argomenti. Spesso i docenti citeranno \textbf{paper accademici} di game
design e di \emph{game studies}, e ``purtroppo c'è poco da fare'': molti non sono
mai stati tradotti, o lo sono stati ``in maniera un po' approssimativa''. Per
questo di solito si linka la fonte originale; nulla vieta di cercare traduzioni
o di tradurseli per conto proprio. In ogni caso i docenti cercano di scegliere
paper \textbf{accessibili}: ``non vi linkeremo delle tesi di master di 300
pagine, a meno che non ce le chiediate voi''.

% ---------------------------------------------------------------------
\section{Due libri per cominciare (e uno da rimandare)}
\label{a2:sec:libri}

Prima di entrare nel tema, Daniele torna su due libri citati la volta
precedente, mostrandoli in video.

\subsection{Un volume sull'edu-larp}

A Tiziana, che ha qualche difficoltà con l'inglese, Daniele consiglia un libro
italiano in cui si trovano molti degli argomenti del corso: un volume
sull'\textbf{edu-larp}, il gioco di ruolo dal vivo in ambito educativo, a cura
di \textbf{Andrea Maragliano} ed edito da \textbf{FrancoAngeli}\footnote{La
trascrizione riporta il titolo come ``EduLARP'', ``Edu LARP'' e ``DualArp'', e
il nome del curatore come ``Andrea Maragliano''; il titolo completo del volume
non è detto a lezione.}. Con una frecciata affettuosa, Daniele aggiunge che
contiene ``un'introduzione molto interessante scritta da un certo Ruggero. Non
so chi sia''. Ruggero protesta (``non fare queste cose, che poi sembra che
consiglio i libri in cui ho partecipato''): il libro non è suo, ci hanno scritto
``tutti i miei amici'', e lui per principio non consiglia i libri su cui scrive
(``possiamo anche comprarlo e non leggere l'introduzione'', ride un collega).

Daniele però lo consiglia convintamente:
\begin{itemize}
  \item molti concetti che si trovano nei siti in inglese ``qua dentro sono
        molto spiegati''; lui stesso lo ha usato per studiare;
  \item c'è una parte, ``bellissima'', sull'\textbf{involvement} (il
        coinvolgimento), che spiega anche ``perché non usiamo la parola
        immersione'' (una delle parole vietate della prima lezione); se ne
        parlerà anche con Ruggero;
  \item finalmente si può consigliare ``roba genuina, sana, italiana di game
        design'';
  \item copre l'ambito \textbf{didattico e formativo}, cioè il gioco nelle
        scuole e il gioco come strumento di formazione, che interessa ad alcuni
        studenti del corso.
\end{itemize}
Non bisogna leggerlo tutto: è utile ``leggerlo pian pianino mentre fate il
corso'', se si vuole approfondire.

\subsection{La progettazione dei giochi da tavolo di Walter Nuccio}

Il secondo libro mostrato la volta precedente è \emph{La progettazione dei
giochi da tavolo} di \textbf{Walter Nuccio}. Qui Daniele fa un ``piccolo
cesellamento'': ``se fossi in voi non lo comprerei per il momento'', e lascia la
parola a Ruggero, che ne ha una copia alle spalle. Il giudizio di Ruggero è
preciso:
\begin{itemize}
  \item è più utile \textbf{dopo} aver interiorizzato alcune basi su come si
        progettano i giochi; come primo libro è un rischio, ``perché sembra che
        ti dia delle scorciatoie'';
  \item più che un saggio di game design, è \textbf{``un grosso archivio di
        meccaniche di gioco''}, una sorta di dizionario (``che cos'è un malus,
        per esempio?''), con una parte di \emph{pattern};
  \item il pericolo è applicare quelle meccaniche ai propri giochi
        \textbf{senza capire che cosa fanno}: ``come se apriamo una cassetta
        degli attrezzi e non sappiamo che cosa dobbiamo fare. Gli attrezzi
        sembrano tutti bellissimi, però magari li usiamo male'';
  \item il consiglio è aspettare a comprarlo, o almeno a leggerlo, fino a dopo
        la lezione sull'\textbf{emotional design}, quando si avrà già un'idea di
        come partire con il proprio gioco.
\end{itemize}
Daniele è ``assolutamente d'accordo''. (Nel frattempo si è spostato in cucina:
``benvenuti nella mia cucina, davanti al mio frigorifero, per parlare di user
experience''.) Del libro di Nuccio, e di un'alternativa in inglese, si
riparlerà a proposito dei giochi da tavolo (sezione~\ref{a2:sec:tavolo-libri}).

% ---------------------------------------------------------------------
\section{Perché analizzare i casi studio}
\label{a2:sec:perche}

Ruggero apre l'argomento del giorno, i \textbf{casi studio}, con due domande
alla classe: \emph{perché} ci teniamo ad analizzarli? E prima ancora: che cosa
ci si aspetta che siano? ``Non è una domanda provocatoria''.

\paragraph{Le risposte degli studenti.} Una prima risposta: sono \textbf{altri
giochi già realizzati}, da analizzare da molti punti di vista, ``anche
insoliti'', non solo da giocatori: dal punto di vista della produzione, o
dell'esperienza utente che se ne è avuta, per capire come migliorarli. Una
seconda risposta: i casi studio sono delle \textbf{milestone}, giochi o
situazioni ``lapidarie'' da cui trarre ``figure di merito'' da sviscerare per
poi generare altre cose. Tiziana chiede invece una cosa ``più basica'': che cosa
si intende esattamente? Modelli, esempi, ma di che cosa?

\paragraph{La definizione.} ``Ottima domanda''. Con \textbf{caso studio} si
intende un \textbf{gioco specifico}, o una \textbf{categoria di giochi}, che,
analizzato, ci dà delle \textbf{informazioni utili}: osservando da vicino che
cosa un gioco fa, possiamo estrapolare informazioni da riproporre, o metodologie
che ci servono per creare i nostri giochi. L'analogia è con la scrittura: chi
vuole scrivere un libro legge tanti libri, e capisce che cosa fanno di positivo
quelli che gli piacciono e che cosa non fanno quelli che non gli piacciono, per
evitare le cose che non gli piacciono e ``iterare sulle cose che ci piacciono''.
Con i giochi è lo stesso: si guarda che cosa fanno di positivo e di negativo,
\textbf{tenendo conto anche del contesto} in cui sono nati.

\paragraph{Oggi la panoramica, poi gli strumenti.} Gli strumenti più tecnici
per analizzare nel dettaglio arriveranno con le lezioni di \textbf{ricerca e
sviluppo}. Oggi si fa una ricognizione di tutti i settori.

% ---------------------------------------------------------------------
\section{Il gioco analogico e le sue quattro macrocategorie}
\label{a2:sec:macro}

\subsection{Un termine ombrello}

Prima una precisazione. Nel corso, \textbf{gioco analogico} è un
\textbf{termine ombrello} che copre i giochi da tavolo (\emph{board game}), i
giochi di carte (\emph{card game}), i giochi di ruolo da tavolo, i giochi di
ruolo dal vivo, i giochi pervasivi, le cacce al tesoro, le \emph{escape room},
i \emph{murder party}: tutti quei giochi che ``generalmente non sono focalizzati
e incentrati sulla programmazione, sull'uso di un terminale, sull'uso di uno
schermo''; in breve, che \textbf{non sono videogiochi}. Esistono cose ibride,
che si vedranno, ma l'area di interesse principale resta il gioco analogico.

\subsection{Le quattro macrocategorie}

\slidefig{img/AGD/agd2-categorie}{Le macrocategorie del gioco analogico: giochi da tavolo, giochi di ruolo da tavolo, LARP e, dietro, l'ibridazione.}{slide \emph{Analog -- Case Studies}}{fig:agd2-categorie}


La lezione parte dal ``frame attuale'': quattro \textbf{macrocategorie}, a loro
volta divise in sottocategorie. Ruggero avverte subito che ``non sono separate da
muri di acciaio non comunicanti'': i casi ibridi, difficili da mettere in una
sola categoria, sono frequenti. I tre docenti presentano ciascuno la propria.

\begin{description}
  \item[Gioco da tavolo (Ruggero).] La categoria dei \emph{board game}. ``Per
        motivi di industria'' vi si includono anche i giochi di carte e i
        \emph{party game}: tutti i giochi \textbf{pubblicati e prodotti
        all'interno dell'industria del gioco da tavolo}, che abbiano o no un
        tabellone.
  \item[Gioco di ruolo da tavolo (Daniele).] GDR in italiano, TTRPG
        (\emph{TableTop Role-Playing Game}) in inglese. È una categoria ``molto
        particolare'', perché anzitutto \textbf{non c'è una definizione
        univoca}: sul mercato si trovano cose molto diverse, perfino giochi che
        sono quasi ``dei LARP da tavolo'', vicini ai \emph{chamber LARP}. Anche
        i \textbf{formati} sono diversissimi: la slide mostra \emph{Dungeons \&
        Dragons}; \emph{Neverland}, un supplemento\footnote{Daniele lo chiama
        ``un'espansione per D\&D''; più avanti lo descrive come una
        reinterpretazione di Peter Pan in chiave OSR, compatibile con le regole
        di D\&D.}; \emph{Dune}, in volumi cartonati ``importanti''; ma anche
        \emph{The Quiet Year}\footnote{La trascrizione dice ``Quietier'' e, più
        avanti, ``The Quiet Ear''.}, che sta in un sacchettino; e \emph{Lady
        Blackbird}\footnote{La trascrizione dice ``Lady Blackboard''. \emph{Lady
        Blackbird} è un gioco di ruolo distribuito gratuitamente in PDF.}, che è
        solo digitale: ``tu lo prendi, lo stampi in bianco e nero e te lo
        giochi''. L'edizione italiana di \emph{Montsegur 1244}\footnote{La
        trascrizione dice ``Monsegur''.} è una sorta di scatola che si apre;
        \emph{CBR+PNK}\footnote{La trascrizione dice ``CBR PNK''.} ``si apre in
        un modo stranissimo'' e contiene dei pieghevoli (\emph{leaflet}); ci sono
        giochi di ruolo in formato brochure. Daniele li mostrerà meglio nella
        lezione sulla \emph{fiction}. La conclusione: nei giochi di ruolo
        \textbf{non si può usare il metro del formato}.
  \item[Gioco di ruolo dal vivo (Ruggero).] Il LARP, \emph{Live Action
        Role-Playing game}: ``anche questa è una categoria gigantesca''.
  \item[Gioco pervasivo (Beatrice).] La categoria meno conosciuta in Italia, e
        anche per questo inserita nel corso.
\end{description}

\paragraph{Che cosa distingue il gioco di ruolo.} Daniele chiede alla classe
che cosa distingua i \emph{role-playing game}. La risposta è nel nome: il
\textbf{role playing}, cioè un \textbf{elemento di interpretazione}, che non è
detto sia interpretazione di un \emph{personaggio} (in \emph{The Quiet Year},
per esempio, non c'è un personaggio, ``abbiamo qualcosa di diverso''). E c'è un
elemento ``importantissimo'': la \textbf{fiction}, cioè la narrazione. Non la
narrazione intesa come ``scrivere una trama e giocarla'', ma quella che nasce al
tavolo, la \textbf{narrazione emergente}: non si è previsto che l'avventura
andasse in quel modo, ma mentre si gioca gli altri dicono cose che la portano in
una certa direzione. Se ne parlerà nella lezione sulla fiction emergente.
Categorizzare i GDR è ``veramente molto arduo'', e anche il pubblico è
particolare: ``ognuno ha il suo palato'', c'è chi gioca tutto e chi gioca solo un
sottogenere. A una domanda dalla classe (il gioco \emph{Lovecraftesque}, citato
da qualcuno, rientra tra i giochi di ruolo da tavolo?) la risposta è sì: e di
che \emph{tipo} di gioco di ruolo sia si capirà con le sottocategorie che Daniele
presenterà più tardi.

\paragraph{Perché il LARP è separato dal GDR.} Ruggero spiega che i LARP sono
separati dal GDR da tavolo ``semplicemente per motivi di settore e di
industria''. Il gioco di ruolo da tavolo viene venduto con una produzione più
vicina al gioco da tavolo: si produce un \textbf{oggetto fisico} che poi si
vende. Il LARP invece è ancora spesso organizzato \textbf{a eventi}: si partecipa
\textbf{comprando un biglietto}, e si va a giocare in un posto specifico, a
un'ora specifica, con persone che sono lì per giocare quel gioco. Da questo
punto di vista è più vicino ad alcuni giochi pervasivi. Le sottocategorie sono
``una valanga'': la più famosa è il \textbf{boffer} (il cliché del LARP: persone
in costume medievale ``che si picchiano con spade di plastazote o di gommapiuma''),
ma ce ne sono di molto più astratte e di molto più accessibili, che si possono
organizzare anche in casa.

\subsection{I giochi pervasivi: una prima definizione}

Beatrice chiede se qualcuno abbia mai sentito parlare di \emph{pervasive games}.
Quasi nessuno; una studentessa ne ha sentito parlare solo leggendo la
descrizione del corso, e si è incuriosita. ``Ci sta tantissimo'': in Italia non
sono ancora molto conosciuti, ed è uno dei motivi per cui sono nel percorso.

Tiziana chiede se le \textbf{cacce al tesoro} organizzate nelle città rientrino
negli \emph{urban game}. ``Assolutamente sì'': gli \textbf{urban game} sono una
categoria dei giochi pervasivi, e nella slide si vedono proprio come un
sottoinsieme del ``grande cappello''. Beatrice dà due definizioni:
\begin{enumerate}
  \item la definizione ``standard'', dai \emph{game studies}: un gioco
        pervasivo è un gioco che \textbf{espande il proprio cerchio magico
        spazialmente, temporalmente e socialmente};
  \item la definizione più intuitiva, da cui partire per capire poi meglio la
        prima: un gioco pervasivo è un gioco che \textbf{utilizza i luoghi e gli
        oggetti della vita quotidiana e li trasforma in luoghi e oggetti di
        gioco}.
\end{enumerate}
Se la prima non è chiara, ``non è un problema, la capiremo con i casi studio''.
Negli urban game ``non giochiamo seduti al tavolo, ma giochiamo in città, e
stiamo comunque giocando''. La parola \emph{pervasivo} indica la pervasività del
gioco, che entra in situazioni in cui non siamo abituati a immaginarlo, anche se
magari siamo abituatissimi a viverlo senza pensarci.

\paragraph{Le cene con delitto.} Tiziana fa un'altra domanda (Beatrice ricorda
che si può sempre interrompere, anche con la mano alzata delle reazioni di
Zoom): le \textbf{cene con delitto} rientrano tra i pervasivi? E il
\textbf{teatro interattivo}, visto che in certe cene con delitto chi le propone
recita una vera rappresentazione? Per Beatrice, ``a mio dire, sì, entrambe''.
Ma le cene con delitto, o \emph{murder party}, esistono in molte versioni:
\begin{itemize}
  \item in una, gli \textbf{attori} vengono al tavolo con te, magari
        insospettabili, e tu, come persona giocante, li interroghi: questa
        versione assomiglia di più al teatro interattivo, che è incluso nei
        giochi pervasivi;
  \item in un'altra, \textbf{ognuno riceve un ruolo}, con una sorta di copione
        (``stasera sei la sorella della vittima''): questa va più nella
        direzione dei \textbf{giochi di ruolo dal vivo}.
\end{itemize}
Sono distinzioni ``che facciamo per allenarci a comprendere meglio'', ma che
spesso si sovrappongono: anche quando si è investigatori e si interrogano degli
attori, si è dentro una dinamica di gioco di ruolo, ``anche se magari non come
la immaginiamo in D\&D o in giochi più blasonati''. Le definizioni ``sono
malleabili'': si imparano le parole e poi si vede come usarle nel caso
specifico.

\subsection{Ibridazione e sperimentazione}

Nell'angolo della slide, dietro tutte le categorie, c'è ``una piccola cosina'':
\textbf{ibridazione e sperimentazione}. ``In realtà le categorie non esistono'',
dice Ruggero: esistono come divisioni tra industrie, ma escono continuamente
giochi che fanno da ponte. La slide mostra alcuni esempi:
\begin{itemize}
  \item \emph{Fateforge}\footnote{La trascrizione dice ``FitForge'' e, più
        avanti, ``Fateforge''.}, un gioco da tavolo che usa un'\textbf{app}
        digitale integrata nel gioco: ``un po' videogioco, un po' board game'';
  \item un gioco di aste che usa anch'esso un cellulare (il titolo non è
        detto);
  \item dei \emph{mock trial} (processi simulati) ibridi, a metà tra il murder
        party e il roleplay, ``un po' roleplay, ma anche un po' giochi
        pervasivi'';
  \item \emph{Alice is Missing}, un gioco di ruolo da tavolo che usa una
        \textbf{chat} come Telegram tra le meccaniche di gioco.
\end{itemize}
Le ibridazioni rendono spesso difficile mettere un gioco ``al 100\%'' in uno di
questi ``cesti''.

\begin{figure}[htbp]
\centering
\begin{tikzpicture}[
    font=\small,
    macro/.style={draw, thick, rounded corners=4pt, fill=gray!12,
                  text width=5.2cm, align=left, inner sep=6pt},
    titolo/.style={font=\bfseries},
    ponte/.style={draw, dashed, rounded corners=3pt, fill=white,
                  text width=3.4cm, align=center, font=\footnotesize,
                  inner sep=4pt}
  ]
  \node[macro] (tav) at (0,0) {{\bfseries Gioco da tavolo} \hfill
    \footnotesize Ruggero\\[2pt]
    \footnotesize board game, card game, party game: tutto ciò che è
    prodotto dall'industria del gioco da tavolo};
  \node[macro] (gdr) at (6.6,0) {{\bfseries Gioco di ruolo da tavolo}
    \hfill \footnotesize Daniele\\[2pt]
    \footnotesize GDR / TTRPG: interpretazione e fiction emergente; nessuna
    definizione univoca, formati diversissimi};
  \node[macro] (larp) at (0,-3.6) {{\bfseries Gioco di ruolo dal vivo}
    \hfill \footnotesize Ruggero\\[2pt]
    \footnotesize LARP: spesso a evento, con biglietto; dal boffer al
    black box};
  \node[macro] (perv) at (6.6,-3.6) {{\bfseries Gioco pervasivo}
    \hfill \footnotesize Beatrice\\[2pt]
    \footnotesize luoghi e oggetti quotidiani trasformati in luoghi e oggetti
    di gioco; urban game, ARG, teatro interattivo};
  \node[ponte] (ibr) at (3.3,-1.8) {\textbf{ibridazione e
    sperimentazione}\\ app, chat, murder party, mock trial, eventi\ldots};
  \draw[dashed, thick] (ibr) -- (tav);
  \draw[dashed, thick] (ibr) -- (gdr);
  \draw[dashed, thick] (ibr) -- (larp);
  \draw[dashed, thick] (ibr) -- (perv);
  \draw[{Stealth[length=2mm]}-{Stealth[length=2mm]}, gray, thick]
    (gdr.south) -- node[right, font=\footnotesize\itshape, text=black,
    text width=2.4cm] {oggetto venduto vs evento con biglietto} (perv.north);
  \draw[{Stealth[length=2mm]}-{Stealth[length=2mm]}, gray, thick]
    (tav.south) -- node[left, font=\footnotesize\itshape, text=black,
    text width=2.4cm, align=right] {separati per motivi di industria}
    (larp.north);
\end{tikzpicture}
\caption{Le quattro macrocategorie del gioco analogico usate nella lezione, con
il docente che le presenta. Al centro, la ``piccola cosina'' della slide:
ibridazione e sperimentazione, che rendono i confini porosi. Le frecce grigie
riassumono due osservazioni dei docenti: il GDR da tavolo e il LARP sono
separati soprattutto per ragioni di industria (oggetto venduto o evento), e per
lo stesso motivo il LARP è vicino a molti giochi pervasivi.}
\label{a2:fig:macro}
\end{figure}

% ---------------------------------------------------------------------
\section{Non chiudersi nel proprio genere}
\label{a2:sec:aperti}

Prima di aprire la ``diga dei casi studio'', i docenti insistono su alcuni punti
di metodo.

\paragraph{Le categorie servono a ragionare, non a recintare.} Ruggero: le
categorie non vengono date ``per dirvi: le cose sono A, B, C e basta'', ma per
far ragionare su \textbf{quanto sono flessibili i confini}, e per dare esempi di
cose \textbf{adiacenti} a quelle che interessano, che ``possono sbloccarvi delle
cose''. Ognuno è lì per interessi diversi (i giochi di ruolo, le carte, i giochi
formativi e didattici\ldots), ma l'invito è a \textbf{non mettersi i paraocchi}
sul proprio genere: ``è grazie alla \textbf{cross-modalità} e al vedere qualcosa
di risolto in un altro gioco che vi arriva poi un'idea per un gioco vostro che
non c'entra apparentemente niente'', magari nemmeno meccanicamente, con l'altro.

\paragraph{Un corso agnostico di genere.} Daniele aggiunge che il corso è
\textbf{agnostico di genere}: alcune lezioni saranno più importanti per una
categoria, ma i \textbf{concetti di base del design si applicano a qualsiasi
sviluppo ludico}. È fondamentale, perché altrimenti sembra che chi fa un certo
tipo di prodotto abbia bisogno di ``un game design particolare'', e non è così.
Anzi: quando è successo, storicamente anche in Italia, si è finito per
\textbf{ostracizzare e separare le community}, mentre i docenti vogliono
\textbf{metterle insieme}. ``Voi non sapete quante volte i videogiochi vengono
ispirati dai giochi da tavolo, e viceversa''.

\paragraph{Il project work è libero.} C'è anche un motivo pratico. Il progetto
di gruppo finale è \textbf{libero}: i docenti non assegneranno una categoria di
gioco, saranno i gruppi a sceglierla, e può essere anche
un'\textbf{ibridazione}, una \textbf{sperimentazione}, un gioco che non si
colloca perfettamente in un genere (``anzi, sarebbe fantastico''). Ma i gruppi
saranno formati da persone con interessi diversi, ed è proprio per questo che
bisogna restare aperti. L'unica cosa che i docenti chiederanno di non fare è
\textbf{un gioco enorme}: il progetto deve essere \textbf{calibrato}.

\paragraph{``Per me, questo è design''.} Infine, Ruggero chiede di \textbf{non
sottovalutare} la lezione: ad alcuni può sembrare ``meramente un elenco di cose
che vi facciamo vedere per nessun motivo'', ma ``per me questo è design''.
Quando deve trovare una soluzione, va a cercare altri giochi; e gli è capitato
moltissime volte di riaprire proprio quelle slide per ricordarsi com'era fatto un
gioco. E come hanno fatto Tiziana e Rita, bisogna \textbf{fare domande}:
``fermateci, chiedeteci'', i docenti sono lì per approfondire.

% ---------------------------------------------------------------------
\section{I giochi pervasivi: la premessa di Beatrice}
\label{a2:sec:pervasivi}

Si comincia ``in scivolata'' con i giochi pervasivi. Beatrice anticipa che
diversi casi studio hanno \textbf{accenni di digitale}, per tre motivi:
\begin{enumerate}
  \item alcuni casi sono \textbf{molto famosi} e aiutano a capire il concetto di
        pervasive game, anche se non sono prettamente analogici;
  \item un gioco analogico \textbf{non esclude completamente la componente
        digitale}: basti pensare ai giochi da tavolo che usano un'app;
  \item l'\textbf{ibridazione}: si vedranno casi studio ``molto strani''.
\end{enumerate}
Aggiunge una considerazione personale e professionale: quello dei pervasivi è
un campo in cui c'è ancora \textbf{ampio spazio}, perché sono poco conosciuti e
la loro \textbf{definizione è ancora in costruzione}. C'è spazio ``per mettere
anche la vostra voce'', per esprimersi e sperimentare, proprio perché le
definizioni non sono state date in modo rigido come altrove. E anticipa un tema
ricorrente del corso: ``parleremo molto di \textbf{emozioni}, anche di come
accompagnarle all'interno di un'esperienza di gioco''.

\paragraph{Approfondimento: il cerchio magico.} La definizione ``da game
studies'' presuppone il concetto di \textbf{cerchio magico}, che a lezione non
viene spiegato. In \emph{The Art of Game Design}, Schell lo introduce
discutendo le definizioni di gioco: quasi tutte implicano che un gioco sia un
sistema \emph{chiuso}, con un confine, e ``Johan Huizinga lo ha chiamato il
\emph{cerchio magico}'', che ha davvero qualcosa di magico: quando siamo
mentalmente ``dentro il gioco'' abbiamo pensieri, sentimenti e valori diversi da
quando ne siamo fuori. Schell lo reinterpreta come il nostro \textbf{sistema
interno di risoluzione dei problemi}: una ``microrealtà'' semplificata in cui gli
elementi acquistano un'importanza che non devono avere fuori dal gioco. Nel
capitolo sulle meccaniche, lo stesso Schell identifica il cerchio magico con lo
\textbf{spazio} del gioco, la prima delle sue sette categorie di meccaniche. Il
gioco pervasivo è dunque quello che sposta, allarga o rende sfumato questo
confine \textbf{nello spazio} (si gioca nella città, non sul tavolo), \textbf{nel
tempo} (non si sa bene quando la partita inizia e finisce) e \textbf{nella
società} (vi entrano persone che non sanno di giocare, o che giocano senza
chiamarlo gioco)\footnote{La formula ``espandere il cerchio magico
spazialmente, temporalmente e socialmente'' è quella proposta da Markus Montola
nei suoi studi sui giochi pervasivi, poi ripresa nel volume \emph{Pervasive
Games: Theory and Design} di Markus Montola, Jaakko Stenros e Annika Waern
(2009). Beatrice non cita la fonte; la nota è un'integrazione.}. È esattamente
quello che Beatrice dirà di \emph{GuruShot}: ``non sappiamo bene dove inizia e
dove finisce l'attività di gioco, ed è per questo che viene utilizzata la parola
pervasivo''.

\begin{figure}[htbp]
\centering
\begin{tikzpicture}[
    font=\small,
    freccia/.style={-{Stealth[length=2.5mm]}, thick},
    etichetta/.style={font=\footnotesize, align=center, text width=3.6cm}
  ]
  % cerchio classico
  \draw[thick, fill=gray!15] (0,0) circle (1.0cm);
  \node[font=\footnotesize, align=center] at (0,0) {cerchio\\magico};
  \node[font=\footnotesize\itshape, align=center] at (0,-1.45)
    {gioco ``classico'':\\tavolo, durata, giocatori};
  % cerchio espanso
  \draw[thick, dashed, fill=gray!6] (7,0) circle (2.1cm);
  \draw[thick, fill=gray!15] (7,0) circle (1.0cm);
  \node[font=\footnotesize, align=center] at (7,0) {cerchio\\magico};
  \draw[freccia] (1.3,0) -- node[above, font=\footnotesize]{gioco pervasivo}
    (4.6,0);
  \draw[freccia, gray] (7,1.0) -- (7,2.1);
  \draw[freccia, gray] (5.95,-0.6) -- (5.18,-1.05);
  \draw[freccia, gray] (8.05,-0.6) -- (8.82,-1.05);
  \node[etichetta, above=0.05cm] at (7,2.15) {\textbf{spazialmente}: la città,
    l'hotel, la stazione dell'autobus};
  \node[etichetta, anchor=north east] at (5.4,-1.5) {\textbf{temporalmente}:
    quando inizia e quando finisce la partita?};
  \node[etichetta, anchor=north west] at (8.6,-1.5) {\textbf{socialmente}:
    passanti, pasticcieri, vicini di quartiere};
\end{tikzpicture}
\caption{La definizione ``da game studies'' citata da Beatrice: il gioco
pervasivo espande il proprio cerchio magico nello spazio, nel tempo e nella
società. La definizione più intuitiva dice la stessa cosa dal lato degli
oggetti: luoghi e oggetti della vita quotidiana diventano luoghi e oggetti di
gioco.}
\label{a2:fig:cerchio}
\end{figure}

% ---------------------------------------------------------------------
\section{Pervasivi con un'app: Pokémon Go, geocaching, GuruShot}
\label{a2:sec:pervasivi-app}

\subsection{Pokémon Go}

\slidefig{img/AGD/agd2-pokemon}{Pervasivi con un'app: Pokémon Go, il geocaching e GuruShots.}{slide \emph{Analog -- Case Studies}}{fig:agd2-pokemon}


Il primo esempio è scelto ``proprio per farvi comprendere'': \textbf{Pokémon
Go}, che conoscono tutti, o di cui almeno tutti hanno sentito il nome. Beatrice
lo spiega comunque, ``così ci allineiamo'' (scusandosi con i grandi fan per la
descrizione imprecisa). È un gioco in formato \textbf{app}. I Pokémon, ``questi
animaletti bizzarri'', nella narrativa dei Pokémon popolano il loro mondo; in
\emph{Pokémon Go} popolano \textbf{il nostro mondo reale}. Il gioco consiste nel
cercarli, catturarli, andare nelle palestre, e tutto questo si fa
\textbf{camminando nel mondo}: la stessa via in cui si va a fare colazione al
bar o che si percorre per andare in ufficio è la via in cui si può trovare un
Pokémon o una palestra. Ed è ``tutto costruito dagli utenti''.

Su \emph{Pokémon Go} ci sarebbe molto da dire, ``anche su come è stato
utilizzato per fare altre cose meno belle ed eleganti dell'esperienza di
gioco'', ma Beatrice non ne parla per ragioni di tempo: ci tiene solo a far
notare che ``tante cose possono avere tante sfumature'', e che i casi più
spinosi si possono discutere a parte. Qui serve per una ragione semplice:
\textbf{è pervasivo perché usa gli stessi spazi della vita quotidiana come spazi
di gioco}.

\subsection{Il geocaching}

Il \textbf{geocaching} è anch'esso abbastanza conosciuto. ``Di base è una
gigantesca caccia al tesoro'', tenuta in piedi da alcuni portali, e nata
proprio con il \textbf{GPS} e la geolocalizzazione. Il principio:
\begin{itemize}
  \item le persone giocanti \textbf{nascondono} delle scatole, piccole o
        grandi, con contenuti vari (dei ``tesori'') e con una sorta di
        \textbf{diario di bordo} (un \emph{log}: fogli o quaderni) su cui chi le
        trova può scrivere ``sono stato qui, ho trovato questa scatoletta'';
  \item la scatola viene nascosta nel posto che l'utente ritiene più idoneo, e
        sulla piattaforma l'utente dà le \textbf{indicazioni} per trovarla;
  \item la grande regola è che \textbf{le scatole non vengono mai tolte}: chi le
        trova può scambiare gli oggetti al loro interno (piccoli, di qualunque
        valore, con piena libertà) e segnare la propria presenza.
\end{itemize}
Non è ``un arraffare i tesori'', ma un modo di \textbf{viverli e dar loro
spazio collettivamente}. In Italia, ``per la cultura che abbiamo'', funziona
così così, perché le scatole spesso vengono rubate o spostate. Ma vale la pena
guardare online i video dei modi in cui vengono nascoste: alcune sono ``piccole
opere di ingegneria artistica'', quasi piccole escape room, con un gioco di pesi
o una chiave ``nascosta in bella vista'' che si vede solo con grande
attenzione; alcune non sono scatole ma \textbf{cabine telefoniche} in cui si
entra e succedono cose, o intere automobili.

\subsection{GuruShot}

\textbf{GuruShot}\footnote{Più avanti la trascrizione lo chiama anche ``blue
shot''.} è di nuovo in formato app, ma Beatrice lo inserisce perché usa uno
strumento che lei usa ``moltissimo'' e che trova ``molto potente per fare molte
cose'': \textbf{la fotografia}. È un gioco dedicato a persone molto dentro al
tema, fotografi anche di un certo livello di passione e di competenza: una
\textbf{gara} o maratona fotografica in cui vengono date delle \textbf{challenge
fotografiche}. Chi vuole giocare apre l'app, sceglie la sfida, la svolge, carica
la foto; e poi ``succedono cose'': valutazioni, momenti di incontro con la
community, premi veri e propri offerti dagli \textbf{sponsor}. È un'attività
\textbf{continua}, online da molti anni (``almeno 7--8'', corregge Beatrice), con
una nicchia internazionale molto nutrita.

Ed è ``molto pervasivo'': \textbf{non si può dire quando una partita è iniziata
o finita}, né quando si sta giocando e quando no. Rispetto alla definizione del
cerchio magico, non si sa bene dove inizi e dove finisca l'attività di gioco:
per questo si parla di pervasività.

% ---------------------------------------------------------------------
\section{La città come campo di gioco: Journey, SF0 e Critical City}
\label{a2:sec:pervasivi-citta}

\subsection{Journey to the End of the Night}

\slidefig{img/AGD/agd2-journey}{La città come campo di gioco: Journey to the End of the Night e SF0.}{slide \emph{Analog -- Case Studies}}{fig:agd2-journey}


Si torna sotto il cappello dei \textbf{giochi urbani}, con un gioco che si
gioca in città, sempre \textbf{in formato evento}: un'ora di inizio, un'ora di
fine e una zona di gioco all'interno della città. È ``una sorta di \textbf{guardie
e ladri gigante}''. Lo scopo dei creatori era \textbf{riappropriarsi della vita
notturna} e viverla in modo ludico, esplorando la città ``in maniera
\emph{safe} da un lato e un po' \emph{noir} dall'altro''. Un piccolo gruppo di
giocatori fa i \textbf{chaser} (gli inseguitori, le guardie) e gli altri i
\textbf{runner} (i ladri); ci sono delle \textbf{tappe} da raggiungere, e il
gioco dura svariate ore. Non è un gioco prettamente sportivo (ci sono altri
pervasivi molto più sportivi): ci sono molte dinamiche \textbf{strategiche} su
come arrivare alla fine del percorso. Beatrice lo inserisce perché, rispetto ai
casi precedenti, è \textbf{completamente analogico}: non usa nessuna app, nulla.
È un evento.

Il gioco è \textbf{Journey to the End of the Night}, nato in America da un
collettivo con base a \textbf{San Francisco}. È arrivato in Italia nel
\textbf{2008}, durante un'ondata di giochi pervasivi arrivata ``in modo molto
forte'', e poi in Giappone e in tantissimi altri posti, perché i creatori lo
hanno reso \textbf{libero}: chi vuole può organizzarlo, purché \textbf{non metta
un biglietto a pagamento}. Ha avuto per anni un grande seguito: cercando
``Journey to the End'' Italia si trovano tantissimi eventi, e a \textbf{Milano},
tra il 2008 e il 2013, se ne sono fatte svariate edizioni con
\textbf{2\,000--3\,000 partecipanti} ciascuna. Beatrice lo racconta ``non tanto
per le meccaniche'', che più o meno si conoscono, ma per \textbf{come è
arrivato}: un'onda lunga che dura ancora, anche se si sta affievolendo.

\subsection{SF0}

Degli stessi creatori è \textbf{SF0}\footnote{La trascrizione dice ``SF0'',
``SFZero'' e, una volta, ``essere al vero''.}, che era una sorta di grande
\textbf{piattaforma online} con lo stesso principio di \emph{GuruShot}: delle
\textbf{missioni} da svolgere e documentare con video o foto. Ma invece di
puntare sull'estetica della fotografia, le missioni puntavano sulle
\textbf{interazioni cittadine e urbane}.

\subsection{Critical City Upload}

SF0 ha avuto ricadute anche in Italia, con ``un figlio, chiamiamolo così'':
\textbf{Critical City Upload}, un gioco che è andato molto tra il 2009 e il
2013. Lo scopo è vincere, accumulando punti con le missioni, ma il gioco è
molto legato agli \textbf{scambi umani e al territorio}. Le missioni andavano da
``organizza una cena con una persona del tuo quartiere con cui non hai mai
parlato'' a una delle preferite di Beatrice: ``crea un'opera d'arte interattiva
che intrattenga le persone che aspettano l'autobus e posizionala in una delle
pensiline del quartiere o della città''.

In entrambi i casi \textbf{è la community che valuta la community}: non ci sono
giudici. E dopo i primi anni, sia SF0 sia Critical City hanno permesso ai
partecipanti che avevano raggiunto un certo livello di \textbf{creare essi
stessi nuove missioni} per i nuovi giocatori: la comunità non si limitava a
valutarsi a vicenda, ma creava il gioco per gli altri. Oggi nessuno dei due è
più attivo, ma i siti sono ancora online, con molte delle missioni visibili.
Quanto a ``stranezza'' rispetto all'idea comune di gioco analogico, ``questo
spinge molto forte''.

% ---------------------------------------------------------------------
\section{Teatro, ARG e marketing: da Sleep No More a The Institute}
\label{a2:sec:pervasivi-arg}

\subsection{Sleep No More}

\slidefig{img/AGD/agd2-sleep}{Teatro e ARG: Sleep No More e Why So Serious?}{slide \emph{Analog -- Case Studies}}{fig:agd2-sleep}


Chi conosce il \textbf{teatro interattivo} conosce di sicuro \textbf{Sleep No
More}\footnote{La trascrizione dice ``Slip no more'' e ``SlipMod''. Lo
spettacolo è della compagnia britannica Punchdrunk; la versione più nota è
quella di New York, allestita in un edificio trasformato nel fittizio McKittrick
Hotel. Il nome della compagnia non è detto a lezione.}: un'opera messa in scena
all'interno di un \textbf{hotel}. È interattiva nel senso che gli attori recitano
la loro parte in interazione con il pubblico, e il pubblico, con regole precise
sui contatti fisici e su come interagire con gli spazi, è \textbf{libero di
esplorare l'hotel} come vuole: in ogni zona ci sono attori che fanno cose. Resta
un'opera teatrale, con un inizio, una fine e degli eventi nel mezzo; c'è molta
documentazione online (interviste, spezzoni). E di nuovo si gioca su
\textbf{qual è lo spazio di gioco e quale non lo è}.

\subsection{Gli alternate reality game e Why So Serious}

Qui Beatrice introduce un nuovo termine: gli \textbf{alternate reality game}
(ARG), i giochi di realtà alternativa, ``giochi in cui c'è un intero mondo
alternativo al nostro, in cui succedono cose''. Il caso scelto, ``uno dei miei
preferiti'', porta anche nel mondo dei \textbf{giochi per il marketing e la
comunicazione}: \textbf{Why So Serious}, nato come gioco promozionale per
\emph{Il cavaliere oscuro} di Christopher Nolan.
\begin{itemize}
  \item Iniziava con una \textbf{mail} ad alcune persone iscritte a mailing
        list di fan di Batman, contattate dagli organizzatori tramite le loro
        associazioni.
  \item La mail era un \textbf{invito strano a entrare nella banda del Joker},
        con un enigma; da lì partivano attività di socializzazione organizzate
        attraverso un sito, che era una serie di indirizzi ``in cui succedevano
        cose''.
  \item Il gioco raccontava che \textbf{nel nostro mondo il Joker esisteva
        davvero}, e gli indirizzi erano indirizzi reali.
  \item Una delle prime cose: andare a questi indirizzi, tutti
        \textbf{pasticcerie}, e dire la parola in codice ricevuta. Si riceveva
        una \textbf{torta} con sopra scritto un numero di telefono; chiamandolo,
        squillava la torta stessa, perché dentro c'era un \textbf{cellulare}:
        lo strumento con cui il Joker gestiva le comunicazioni con i giocatori.
\end{itemize}
Molte strutture di questo tipo, molti colpi di scena ``anche proprio pratici'',
che sfruttano la pervasività: cose che nel nostro mondo non hanno motivo di
succedere (``non esiste che ci sia un cellulare in una torta'') ma che in quel
mondo hanno un senso. Gli ARG giocano moltissimo su ciò che \textbf{potrebbe
essere reale ma non lo è}, e che il gioco dice che lo è: ``in quel mondo voi
siete gli aiutanti del Joker''. \emph{Why So Serious} è opera di un'azienda che
si occupa da molti anni di ARG\footnote{Beatrice non fa il nome dell'azienda,
che si troverà nell'estratto distribuito. Si tratta di 42 Entertainment
(integrazione).}; i dettagli sono nell'estratto che Beatrice caricherà, tratto
da \emph{Immersi nelle storie} di \textbf{Frank Rose}\footnote{Edizione
italiana di \emph{The Art of Immersion} (2011).}, che racconta proprio l'inizio
di questo gioco.

Perché raccontarlo? Perché mostra ``come possiamo giocare sulle
\textbf{emozioni}'' e come \textbf{connettere le persone alle cose che vogliamo
comunicare}, anche costruendo un mondo alternativo. Non è un obbligo, ``ma è bene
sapere'' che esistono questi giochi, ``giganteschi''.

\subsection{Un gioco da tavolo che è anche un ARG}

Tornando alla malleabilità delle definizioni, Beatrice presenta un gioco che è
insieme \textbf{gioco da tavolo}, \textbf{alternate reality game} e gioco con
\textbf{componenti digitali}, edito da \textbf{MS Edizioni}\footnote{La
trascrizione chiama il gioco ``The Ice''; il titolo non è ricostruibile con
certezza (potrebbe trattarsi di un titolo come \emph{The Heist}). MS Edizioni
è un editore italiano di giochi; è forse la stessa ``MS'' per cui, nella prima
lezione, Ruggero dice di lavorare, ma la trascrizione non permette di
confermarlo.}. Si riceve un \textbf{fascicolo} di informazioni e poi dei
\textbf{messaggi}, da cui si capisce di far parte di un gruppo che \textbf{sta
rapinando una banca}, già nel pieno della rapina. I giocatori sono ``l'entità che
ha accesso a internet'', la voce fuori dall'azione che coordina: parlano con dei
personaggi, scambiano informazioni, guardano video, analizzano dati e
documenti. È un gioco ``con una profondità notevole''. Due cose lo rendono
interessante:
\begin{itemize}
  \item mescola cose che \textbf{esistono davvero} (la cartina di una città
        reale, gli scambi della metropolitana di una città reale) con cose del
        \textbf{mondo narrativo} (i personaggi, i fatti, la rapina: ``spero che
        sia chiaro che non state rapinando la banca davvero'');
  \item a differenza degli altri casi, ha un \textbf{inizio e una fine ben
        definiti} ed è \textbf{completamente accessibile}: lo si può provare.
\end{itemize}
Per Beatrice è un caso ``unico sul panorama italiano'' (``poi mi correggeranno i
colleghi a livello internazionale'').

\subsection{Collettivo Cinetico: Cinetico 4.4 e Amleto}

\slidefig{img/AGD/agd2-teatro}{Cinetico 4.4, Lettere anonime per un camminatore, The MP3 Experiment, The Institute e l'ARG di Stranger Things.}{slide \emph{Analog -- Case Studies}}{fig:agd2-teatro}


\textbf{Cinetico 4.4} è ``un altro bell'esperimento'': un \textbf{gioco da
tavolo creato da un collettivo di danzatori e danzatrici}, il \textbf{Collettivo
Cinetico}, per generare una \textbf{performance artistica dal vivo}, teatrale o
di danza. Non è un gioco da prendere dallo scaffale di una ludoteca: è pensato
per una performance. Chi ci gioca partecipa a un gioco da tavolo che in realtà è
un \textbf{evento}: si paga un biglietto, si entra, ci si siede, succedono cose.
Il collettivo, di cui Beatrice si dichiara ``di parte'', fa moltissime attività
``né danza, né teatro, né gioco, o meglio tutte e tre le cose insieme'', difficili
da mettere sotto un cappello. E qui Beatrice aggancia un tema delle prossime
lezioni: la \textbf{transmedialità} e la \textbf{ricerca e sviluppo}, cioè
attingere anche a cose lontane dal gioco che si vuole creare, perché ``più
conosciamo, più possiamo aprirci la testa''.

Un caso studio in più: l'\textbf{Amleto} del Collettivo Cinetico. È l'Amleto di
Shakespeare, ma sul palco non ci sono i danzatori del collettivo. Tre giorni
prima dello spettacolo fanno una vera e propria \textbf{audizione} aperta alla
popolazione; tra chi si presenta scelgono \textbf{sei potenziali Amleto}, che
portano sul palco a danzare e recitare, e creano momenti di
\textbf{interazione con il pubblico} per scegliere, tra i sei, l'Amleto di
quell'edizione. Ci sono giochi veri e propri sia nell'audizione sia nello
spettacolo. Loro, da danzatori, lo chiamano spettacolo di danza, ``ma se lo
vedete c'è tantissimo''.

\subsection{Lettere anonime per un camminatore}

``Rispetto a quello che vi ho detto, questo è semplice''. \textbf{Lettere
anonime per un camminatore} è del collettivo \textbf{Amigdala}, anche loro una
compagnia teatrale (Gianmarco, che ha vissuto a Modena, conferma che sono di
lì). È una \textbf{passeggiata} con piccole interazioni ludiche, tutta basata su
un \textbf{audio}:
\begin{itemize}
  \item è un gioco \textbf{in solitario}: una persona sola segue un percorso
        segnato da frecce, per terra o sui muri di una città, guidata da una
        \textbf{voce in cuffia} che le racconta delle cose;
  \item la struttura del racconto è fissa, ma la \textbf{parte centrale} cambia
        ogni volta: prima di ogni edizione il collettivo \textbf{intervista le
        persone che hanno abitato o abitano il percorso}, e mentre si cammina si
        ascolta, oltre alla narrazione, il racconto di chi quei luoghi li vive o
        li ha vissuti;
  \item la \textbf{conclusione} è sempre ``col botto'': il collettivo sceglie di
        finire in modo scenografico in luoghi normalmente inaccessibili.
        Beatrice lo ha giocato a Ferrara, e alla fine del percorso si è
        ritrovata seduta in poltrona, con un tavolino e un drink, al
        \textbf{centro dello stadio comunale}, con lo stadio vuoto e le cuffie
        nelle orecchie.
\end{itemize}
Torna il tema del \textbf{percorso emotivo} che si fa fare alla persona giocante.
E lo strumento dell'\textbf{audio che guida} è interessante anche per il corso:
un'app non si riesce a fare nel tempo del project work, ``ma per utilizzare
l'audio invece possiamo''.

\subsection{MP3 Experiment}

Un altro esempio basato sull'audio, più simile a un \textbf{flash mob}:
\textbf{MP3 Experiment}, del collettivo newyorkese \textbf{Improv
Everywhere}\footnote{La trascrizione dice ``Improve Everywhere''.}, che fa
moltissimi esperimenti di questo tipo. Ci si ritrova tutti nel giorno, nell'ora
e nel posto concordati, ci si collega con le cuffie a una trasmissione in
diretta, e la voce dice che cosa fare. Beatrice non svela di più (``guardatevi i
video su YouTube, ci sono delle perle''), ma racconta un dettaglio: gli
organizzatori \textbf{elencano gli oggetti da portare}, e in un'edizione chiesero
di portare \textbf{due matite}; a un certo punto tutto il gruppo si mise a
suonare la batteria con le matite sugli oggetti della città. Quella
\textbf{scintilla di curiosità} per oggetti di cui non si sa ancora l'uso è uno
strumento del percorso emotivo, ``non difficile'' da usare: ``basta solo
ragionare di creatività''. Come le istruzioni audio.

\subsection{The Institute e Dispatches from Elsewhere}

\textbf{The Institute}\footnote{Il gioco è noto anche come \emph{The Jejune
Institute}, di San Francisco (integrazione).} è, nel modo in cui Beatrice ha
imparato questo mondo, uno dei casi più blasonati e famosi di alternate reality
game, uno dei primi raccontati come tale. È un'esperienza così complessa che
Beatrice preferisce indicare due fonti per approfondirla:
\begin{itemize}
  \item un \textbf{documentario}, intitolato anch'esso \emph{The Institute}, che
        racconta tutto l'ARG, durato circa un anno e fatto di molti passaggi, con
        tantissimi colpi di scena del tipo ``torta con il cellulare'';
  \item una \textbf{serie TV}, \emph{Dispatches from Elsewhere} (Daniele
        conferma il titolo), che ha preso l'esperienza di The Institute, ci ha
        aggiunto una storia dei personaggi e ne ha fatto una serie. Tolta la
        storia dei personaggi (``comunque molto carina''), resta più o meno
        l'idea dell'ARG.
\end{itemize}

\paragraph{Con budget zero.} Qui Beatrice fa una precisazione importante. Questi
esempi sono \textbf{difficili da replicare}: ``con che soldi compriamo la torta e
ci mettiamo il cellulare? Con che soldi affittiamo un intero palazzo come in The
Institute?''. Serve un intero lato commerciale, di ricerca di sponsor, che non è
parte del design del gioco e quindi non è il motivo per cui si è nel corso.
Quello che vuole che gli studenti si portino a casa è che \textbf{con molto meno}
si possono ottenere esperienze ``estremamente significative a livello
emozionale'' e creazioni di game design molto profonde \textbf{con budget zero}.
I casi grandi servono come \textbf{punti di riferimento} e per sapere che
esistono cose ``molto distanti da quelle che immaginiamo tutti i giorni''.

\subsection{L'ARG di Stranger Things a Milano}

Nello stesso filone, un ARG che ``è stato un'operazione di marketing e basta''
(``per me sono molto affascinanti, però era un'operazione di marketing''), ma
fatto \textbf{a Milano}: quello per una stagione di \emph{Stranger Things} (la
terza o la quarta, Beatrice non ricorda). Era una sorta di \textbf{caccia al
tesoro} in città, con punti di gioco preparati in modo ``altamente scenografico''
(nella slide si vede un autobus), mescolata con una \textbf{telefonata} a un
numero in cui uno dei personaggi diceva che il Sottosopra era arrivato a Milano e
che serviva aiuto. Il messaggio: anche da noi questi giochi ``sono arrivati,
arrivano, e credo che col tempo avranno sempre più spazio'', e per chi pensa in
termini di lavoro ``è un tema''.

% ---------------------------------------------------------------------
\section{Cacce al tesoro, corse e indagini}
\label{a2:sec:pervasivi-cacce}

\subsection{Le guide-caccia al tesoro su carta}

Beatrice presenta poi delle \textbf{guide per esplorare la città} in forma di
caccia al tesoro\footnote{La trascrizione riporta il nome come ``Waiwai''; la
grafia non è verificabile.}, uscite negli stessi anni dell'ondata dei pervasivi
in Italia (la prima nel 2007--2008), disponibili solo per alcune città e
vendute come \textbf{libretti} sul loro sito o da rivenditori autorizzati.
Tiziana le conosce: aveva comprato quelle di \textbf{Venezia} e di
\textbf{Verona}, ha giocato in parte quella di Venezia, nella versione breve
(c'erano due possibilità: un gioco di un paio d'ore o uno da una giornata
intera), e ricorda il \textbf{numero di cellulare} con cui interagire per avere
le indicazioni. Beatrice le aveva giocate con gli SMS, ``era un'altra vita'';
oggi pare si siano aggiornate a WhatsApp (e chiede aggiornamenti a chi le
giocherà).

Sono un altro ibrido di \textbf{analogico e digitale}: un'entità parla con i
giocatori tramite messaggi, mentre il libretto cartaceo resta in mano a loro. La
particolarità è la meccanica per passare da una tappa all'altra: \textbf{ogni
pagina del libro è divisa in strisce}, e quando si risolve una tappa la guida
comunica dei \textbf{codici} (``il foglietto 3, il 32C e il 18F'') con cui
ricomporre la pagina che porta alla tappa successiva.

\subsection{PacManhattan}

\slidefig{img/AGD/agd2-cacce}{PacManhattan, Zombies, Run! e Reptilarium.}{slide \emph{Analog -- Case Studies}}{fig:agd2-cacce}


\textbf{PacManhattan}\footnote{La trascrizione dice ``Pac-Man-Hatton''.}, degli
stessi anni, è \textbf{Pac-Man per le strade di Manhattan}: una persona fa
letteralmente Pac-Man, guidata da un'altra persona (``la mente'') tramite
telefono o walkie-talkie; altri giocatori fanno i fantasmi. È di nuovo ``una
sorta di guardie e ladri'', ed è uno dei primi urban game che si trovano nei
manuali di game studies. Beatrice lo mostra per un'idea precisa: si può
\textbf{prendere qualcosa che esiste già in una forma e cambiarle forma
completamente}. Non è lo stesso gioco di Pac-Man, ma si ispira alla sua
narrativa, alle sue meccaniche, alla sua estetica.

\subsection{Zombies, Run!}

\textbf{Zombies, Run!}\footnote{La trascrizione dice ``Zombie Run''.} è
un'\textbf{app per fare jogging}: invece di ascoltare musica, o i messaggi
vocali degli amici, si ascoltano \textbf{gli zombie che stanno per prenderti}, e
c'è tutta una narrativa su come si sta scappando mentre si corre. Si trova sugli
store.

\subsection{Reptilarium}

Si torna al teatro, ma in versione analogico-digitale. \textbf{Reptilarium} è un
caso studio di cui purtroppo c'è pochissima documentazione. È uscito durante il
\textbf{COVID}, da una compagnia teatrale di Roma, \textbf{Project
XX1}\footnote{Grafia come nella trascrizione.}. Era a modalità evento (un certo
giorno, a una certa ora) e le persone accedevano da \textbf{casa propria}, davanti
al computer, a un'esperienza con degli \textbf{attori} che si mettevano in
comunicazione con loro: ricevevano mail e telefonate, potevano vedere video
``recuperati da altre persone'', potevano \textbf{coordinare i movimenti degli
attori} e vedere che cosa succedeva mentre li eseguivano. La trama partiva da una
\textbf{persona scomparsa} da ritrovare, e poi si espandeva molto. Come per le
cene con delitto (dove si può interrogare l'attore o essere l'indiziato), anche
questo esperimento ``può diventare un LARP'', un ibrido tra teatro e gioco di
ruolo, o tra gioco investigativo e alternate reality game.

Ruggero aggiunge una nota: il \textbf{lockdown} ha colpito moltissimo i giochi
pervasivi, ma anche altri generi come i LARP, e in quel periodo sono nate
esperienze che hanno spostato \textbf{online e a distanza} giochi nati per essere
giocati di persona. Se ne riparlerà.

% ---------------------------------------------------------------------
\section{Pervasivi per la formazione: Echi resistenti}
\label{a2:sec:pervasivi-edu}

\slidefig{img/AGD/agd2-echi}{Un gioco pervasivo per la formazione storica.}{slide \emph{Analog -- Case Studies}}{fig:agd2-echi}


``Tutte le cose che vedete hanno enorme spazio per essere esplorate in tema
\textbf{educativo e formativo}'', e Beatrice chiude con un accenno a questo
ambito. \textbf{Pop History} è un'\textbf{associazione di promozione sociale}:
non si presenta come collettivo teatrale né come gruppo di game designer, ma come
associazione (``credo siano anche loro emiliani''), e fa tante piccole
sperimentazioni. Una di queste è \textbf{Echi resistenti}\footnote{La
trascrizione dice ``Ecchi Resistenti'' e ``Eh, che resistenti''.}, un gioco
sulla \textbf{Resistenza e sulla Liberazione}, che ibridava:
\begin{itemize}
  \item un'\textbf{indagine} con dei fascicoli;
  \item un'\textbf{attività teatrale}, con personaggi storici con cui si poteva
        interagire direttamente;
  \item una \textbf{caccia al tesoro}, in cui si andava a cercare questi
        personaggi per la città.
\end{itemize}
Altri casi simili sono sul loro sito.

\paragraph{Approfondimento: giochi con migliaia di giocatori.} In
\emph{Challenges for Game Designers}, nel capitolo sui giochi multigiocatore,
Brathwaite e Schreiber osservano che nei giochi da tavolo il problema
dell'interfaccia è evidentissimo (``come si mettono venti persone attorno a un
solo tabellone? E mille?''), eppure ``esistono giochi non digitali con migliaia
e perfino centinaia di migliaia di giocatori'': gli \textbf{alternate reality
game} e i grandi \textbf{urban game}, che ``portano i giocatori via dai loro
computer e dentro il mondo reale''. È la stessa scala dei casi di Beatrice:
\emph{Journey to the End of the Night} con migliaia di partecipanti a Milano,
\emph{Why So Serious} con le sue mailing list di fan, \emph{Pokémon Go} con il
mondo intero come mappa.

\begin{table}[htbp]
\centering
\small
\begin{tabular}{@{}p{3.6cm}p{2.3cm}p{7.0cm}@{}}
\toprule
\textbf{Caso studio} & \textbf{Formato} & \textbf{Perché è nella lezione} \\
\midrule
Pokémon Go & app & gli spazi quotidiani diventano spazi di gioco \\
Geocaching & piattaforma + scatole & caccia al tesoro collettiva: le scatole
non si tolgono mai \\
GuruShot & app & la fotografia come strumento; non si sa quando la partita
inizi e finisca \\
Journey to the End of the Night & evento urbano & guardie e ladri notturno,
completamente analogico, libero e non a pagamento \\
SF0 & piattaforma & missioni urbane documentate; la community valuta e crea
missioni \\
Critical City Upload & piattaforma & il ``figlio'' italiano di SF0
(2009--2013): missioni di relazione sul territorio \\
Sleep No More & teatro interattivo & il pubblico esplora liberamente un hotel
\\
Why So Serious & ARG di marketing & la torta con il cellulare: il Joker esiste
davvero \\
Gioco di MS Edizioni & gioco da tavolo + ARG & rapina in corso tra documenti
reali e finzione; ha inizio e fine \\
Cinetico 4.4, Amleto & performance & un gioco da tavolo che è uno spettacolo; il
pubblico sceglie Amleto \\
Lettere anonime per un camminatore & passeggiata audio & gioco in solitario,
voci degli abitanti, finale scenografico \\
MP3 Experiment & flash mob audio & la curiosità per gli oggetti da portare \\
The Institute & ARG & il caso blasonato, con documentario e serie TV \\
ARG di Stranger Things & ARG di marketing & un ARG fatto a Milano \\
Guide-caccia al tesoro & libretto + SMS & pagine in strisce da ricomporre con
codici \\
PacManhattan & urban game & un videogioco cambiato completamente di forma \\
Zombies, Run! & app & narrativa audio durante la corsa \\
Reptilarium & teatro a distanza & attori guidati da casa durante il COVID \\
Echi resistenti & gioco formativo & indagine, teatro e caccia al tesoro sulla
Resistenza \\
\bottomrule
\end{tabular}
\caption{I casi studio di giochi pervasivi presentati da Beatrice, nell'ordine
della lezione, con il motivo per cui sono stati scelti.}
\label{a2:tab:pervasivi}
\end{table}

% ---------------------------------------------------------------------
\section{I giochi di ruolo dal vivo: le macrocategorie}
\label{a2:sec:larp}

\slidefig{img/AGD/agd2-larp}{Le macrocategorie del LARP: boffer, Mind's Eye Theatre, parlor/chamber LARP, American freeform e murder party.}{slide \emph{Analog -- Case Studies}}{fig:agd2-larp}


Finita la parte sui pervasivi, si fa una pausa di pochi minuti (``facciamo
sette minuti, così andiamo ai 20 tondi''). Al rientro ci sono tutti, compreso
Tancredi (``ci sono tutti e due'' i Lorenzo), e Ruggero passa ai
\textbf{giochi di ruolo dal vivo}.

\paragraph{Una struttura diversa.} Questa parte è organizzata diversamente da
quella sui pervasivi. Il gioco di ruolo dal vivo, come il gioco di ruolo da
tavolo e il gioco da tavolo, ha alle spalle \textbf{anni di industria
solidificata}, ed è quindi più facile identificare \textbf{macrocategorie} con
nomi e formati specifici. Prima, però, Ruggero chiede chi abbia mai giocato un
LARP. Una persona della classe sì: un LARP fantasy, al \textbf{festival del gioco
di ruolo di Bracciano}, ``con le spade di gommapiuma, insomma, quelle cose lì''.
È il primo impatto di molti.

\subsection{Il boffer LARP}

Chi ha giocato un LARP probabilmente ha giocato un \textbf{boffer}, e chi ne ha
solo sentito parlare ha in mente quello. È uno dei primi giochi di ruolo dal
vivo diffusi, nato all'epoca proprio come \textbf{tentativo di adattare
\emph{Dungeons \& Dragons} dal vivo}, poi evoluto nel tempo. (Daniele mostra
una foto e Ruggero protesta: ``non evidenziarmi, che già è cringe''; ne ha
giocati anche lui.) Il nome viene dalle armi: le armi \emph{boffer} sono armi in
plastazote che devono avere caratteristiche precise per non essere pericolose.
Sono giochi \textbf{incentrati sul combattimento}, su trame epiche e così via.
Sono anche quelli che possono essere più grandi: un boffer anche grosso può
ospitare \textbf{400--500 persone}, ``per non dire quelli in Germania, che ne
tengono anche migliaia''.

\subsection{Il Mind's Eye Theatre}

Un altro tipo abbastanza diffuso è il \textbf{Mind's Eye Theatre}: se il boffer
deriva da \emph{Dungeons \& Dragons}, questo deriva da \textbf{\emph{Vampire: The
Masquerade}}. Le trame sono più dark, gotiche, horror, legate alla proprietà
intellettuale di Vampire, e chi partecipa di solito interpreta uno dei vampiri
delle loro società segrete. Sono giochi \textbf{più politici che di
combattimento}: l'intrigo, i segreti, il tradirsi a vicenda, ``il pugnalare
alle spalle qualcuno'' sono parte delle dinamiche. Come i boffer, sono giochi
\textbf{in costume}: il costume è richiesto, e si chiede a chi gioca di vestirsi
come il proprio personaggio.

\subsection{Chamber LARP e parlor LARP}

Il costume sembra scontato, ``ma in realtà non tutti i LARP usano il costume
come strumento necessario''. I \textbf{parlor LARP} o \textbf{chamber LARP}
sono giochi di ruolo dal vivo in cui il costume \textbf{non è richiesto}: si
gioca con i propri vestiti di tutti i giorni. Lo studente di prima interviene: a
Bracciano, in effetti, giocavano con i propri vestiti (``di solito gli eventi
così sono degli ibridi''); e ricorda di aver giocato un altro LARP al
\textbf{Play}\footnote{La trascrizione dice ``Play di Bologna''; il festival
Play, dedicato al gioco, si tiene a Modena.}, attorno a un tavolo, ognuno con un
ruolo e con i propri vestiti: \textbf{\emph{A cena col nonno}}. ``Perfetto, quello
è un chamber LARP''.

Ruggero precisa un punto sottile: alcuni chamber LARP si giocano \textbf{a un
tavolo}, ma questo non li rende giochi di ruolo da tavolo, perché si giocano
personaggi \emph{che sono seduti a un tavolo}, mentre in \emph{Dungeons \&
Dragons} il personaggio non è necessariamente seduto a un tavolo in quel momento.
Con gli altri LARP hanno in comune la \textbf{sovrapposizione tra la fisicità di
chi gioca e quella del personaggio}, l'\textbf{embodiment}. Rispetto ad altri
LARP, i chamber LARP hanno di solito:
\begin{itemize}
  \item un livello di \textbf{accessibilità} più alto, perché spesso non
        richiedono costumi;
  \item un \textbf{budget ridotto};
  \item \textbf{non sono a evento}: molti si scaricano in \emph{print and play},
        si stampano a casa e si giocano con gli amici;
  \item per questo, un \textbf{numero di giocatori} più ridotto, di solito da
        3 a 20 persone.
\end{itemize}

\subsection{L'American freeform}

Una categoria dei chamber LARP è l'\textbf{American freeform}: LARP progettati da
designer che si occupano soprattutto di giochi di ruolo da tavolo, e infatti è
un \textbf{genere ibrido}. Visti da fuori sono chamber LARP (gente che fa
\emph{embodiment} del personaggio in delle stanze, non necessariamente in
costume), ma hanno ``un sacco di meccaniche'' prese dal gioco da tavolo e dal
gioco di ruolo da tavolo: le cosiddette \textbf{metatecniche} o metameccaniche,
di cui si parlerà in un'altra lezione. Anche questi sono spesso in \emph{print
and play}; alcuni sono venduti in un formato più vicino al gioco da tavolo, come
mazzi di carte o libretti. Quindi non sono a evento, a differenza dei boffer,
del Mind's Eye Theatre e dei murder party.

\subsection{I murder party}

Ruggero riprende l'ottima distinzione fatta da Beatrice. ``Murder party'', da noi,
indica culturalmente due cose:
\begin{itemize}
  \item le \textbf{cene con delitto}, in cui si va al ristorante, degli attori
        mettono in scena un omicidio e la successiva indagine, e i commensali
        mangiano e intanto risolvono il caso. \textbf{Nessuno dei commensali è
        l'assassino}: è uno degli attori;
  \item la categoria ``più british'', i \textbf{murder party all'inglese}, o i
        loro equivalenti che si stanno diffondendo moltissimo in Asia in questo
        periodo\footnote{La trascrizione dice ``Jubel Shock'' e ``Juvenile
        Show''. Si tratta quasi certamente dei \emph{jubensha}, i giochi di
        investigazione con ruoli molto diffusi in Cina.}, in cui \textbf{si
        interpretano i personaggi della storia} su cui si indaga: tra i giocatori
        può esserci l'assassino.
\end{itemize}
La seconda forma rientra, per il design dell'esperienza e delle meccaniche,
sotto l'ombrello dei LARP; gli aspetti ambientali e le decorazioni la
avvicinano poi al mondo del pervasivo.

% ---------------------------------------------------------------------
\section{Il Nordic LARP}
\label{a2:sec:nordic}

\slidefig{img/AGD/agd2-nordic}{Il Nordic LARP e i Knutebooks, i volumi collettivi del Knutepunkt.}{slide \emph{Analog -- Case Studies}}{fig:agd2-nordic}


C'è poi ``una grossissima categoria'', che si rivedrà più avanti perché ``è stata
importantissima'': il \textbf{Nordic LARP}, il LARP progettato in
\textbf{Scandinavia} (Danimarca, Norvegia, Svezia e Finlandia).

\subsection{Perché è così importante}

Ha un'importanza gigantesca nel settore \textbf{perché ha documentato tutto}. La
maggior parte dei libri accademici di game design e di analisi di questo tipo di
giochi, gli eventi e le conferenze, si fanno in Scandinavia, e c'è una lunga
storia di analisi e documentazione: la prima conferenza è del \textbf{1997}, e da
allora è annuale. È ``gente che si incontra da parecchi anni a parlare di come si
progetta e come si gioca''. Ed è una scena molto \textbf{cross-mediale}: prende
dagli alternate reality game, dai LARP americani, dai parchi a tema, dalle
simulazioni economiche, da tante teorie parallele a quelle del gioco, cercando
di interiorizzare, iterare e integrare tecniche di altri settori. Anche il Nordic
LARP ha le sue sottocategorie: alcune esperienze sono più vicine al chamber
LARP, altre al boffer.

\subsection{Come si riconosce}

Da fuori, un Nordic LARP si riconosce da \textbf{come è progettato il design
dell'esperienza}:
\begin{itemize}
  \item un focus sul \textbf{coinvolgimento}, sull'\textbf{emozione} e sulla
        \textbf{coesione tematica};
  \item pochissime cose \textbf{vestigiali}: per esempio, anche LARP molto
        grandi possono non avere il combattimento, se il combattimento non fa
        parte del tema;
  \item una cura del \textbf{post-esperienza}: molti giochi hanno una
        struttura per ``che cosa succede dopo il gioco''.
\end{itemize}
Molti concetti nati da questa scena si sono ``solidificati nel gergo del game
design'' e torneranno nella lezione sull'emotional design.

\paragraph{Origini diverse da D\&D.} Storicamente, il LARP nordico ha
un'origine \textbf{precedente a \emph{Dungeons \& Dragons}}. Si tende a pensare
che il LARP nasca da ``prendiamo D\&D e lo facciamo dal vivo'', e per alcune
nazioni, come l'Inghilterra e l'America, è proprio così; per altre no. Il Nordic
LARP prende strategie e strumenti dal \textbf{teatro}, dalle avanguardie, dallo
\textbf{psicodramma}, dalla \textbf{Gestalt therapy}; ``è ovvio che quello che
esce fuori è un gioco molto particolare''.

\paragraph{I manifesti.} È anche un ambiente ``pieno, carico, zeppo di
manifesti''. La slide ne mostra alcuni, nati negli anni nella scena; non si
approfondiscono ora, ma sono tutti pubblici e si possono leggere a parte.

\subsection{Knutpunkt e i Knutepunkt books}

La cosa che interessa di più è la conferenza citata prima, che cambia nome a
seconda del paese che la ospita: \textbf{Knutpunkt}, \textbf{Knudepunkt},
\textbf{Knutepunkt} o \textbf{Solmukohta}\footnote{Rispettivamente in svedese,
danese, norvegese e finlandese: la conferenza ruota tra i quattro paesi. La
trascrizione dice ``Knutpunkt o Knutpunkt o Knudpunkt o Solmukott'' e, più
avanti, ``Knut Punk''.}. È il principale punto di ritrovo scandinavo sul gioco
di ruolo dal vivo: \textbf{quattro giorni di conferenza non stop}, con un
programma di centinaia di slot, spesso \textbf{otto track in parallelo} di talk,
panel e workshop. Si tiene dal 1997, e \textbf{ogni anno produce dei libri}
che riassumono lo stato dell'arte in forma di articoli: la documentazione è
``costantemente in aggiornamento''. Ruggero ne ha diversi, ``dei bei mattoni'':
ne mostra uno di \textbf{600 pagine} su come si progetta un LARP, uscito da una
delle edizioni (``nel 2018, se non ricordo male''). È forse ``una delle
comunità ludiche che più documenta quello che fa'', e soprattutto lo documenta
in modo \textbf{pragmatico}, lasciando framework e strumenti a chi vuole creare
giochi di questo tipo. Molte di queste cose si applicano anche ad altri generi, e
torneranno.

\subsection{Dogma 99}

In chat qualcuno chiede se c'entri il \textbf{Dogma} di Lars von Trier.
``Sì e no'': il \textbf{Dogma 99} è un manifesto di LARP che riprende il
manifesto cinematografico del \textbf{Dogma 95} per fare un parallelo con il
LARP. Il Dogma 95 aveva alcune direttive, per esempio che tutto debba essere
\textbf{diegetico}, senza nulla di extradiegetico; il Dogma 99 fa lo stesso: in
un LARP storico non ci devono essere cose extradiegetiche rispetto alla storia.
Una conseguenza: il Dogma 99 \textbf{non fa LARP fantasy}, perché la magia non
esiste e non si può riprodurre in modo perfetto; si fanno giochi ambientati
tendenzialmente ai giorni nostri, che si possono riprodurre facilmente. Da
questa catena di ragionamenti sono nati ``giochi molto interessanti''.

% ---------------------------------------------------------------------
\section{Jeepform: il dramma quotidiano}
\label{a2:sec:jeepform}

\slidefig{img/AGD/agd2-jeepform}{Jeepform tra tabletop, freeform e LARP; la locandina di Fat Man Down.}{slide \emph{Analog -- Case Studies}}{fig:agd2-jeepform}


Tra i manifesti c'era lo \textbf{Jeepform}\footnote{La trascrizione dice
``G-Form'' e ``Jeep Form''.}, nato all'interno del \textbf{freeform}, la zona ``a
metà tra il LARP e il tabletop''. Il suo manifesto, che si trova online, ha tra i
punti più importanti questo: \textbf{il dramma quotidiano è più interessante del
dramma epico}. Tutti i giochi Jeepform sono giochi di ruolo dal vivo su
\textbf{situazioni quotidiane}, ``drammaticamente quotidiane''.

Alcuni sono emotivamente molto pesanti, e storicamente importanti per vari
motivi; Ruggero ha lasciato fuori i più pesanti (se ne riparlerà nella lezione
sull'emotional design) e ne presenta due ``più leggeri'', che fanno comunque
capire in che toni si va.

\subsection{Il gioco della coppia perfetta}

Il primo\footnote{La trascrizione riporta il titolo come ``Lady 8'' e i nomi dei
personaggi come ``Lady'' e ``8''; né il titolo né i nomi sono verificabili.} è
un gioco per \textbf{quattro persone} e \textbf{due personaggi}, che per premessa
\textbf{si amano tantissimo}. Due persone alla volta interpretano i personaggi,
in un gioco \textbf{diviso in scene}; ogni scena è descritta dalle due persone che
non sono in gioco. L'unico obiettivo di chi gioca la coppia è \textbf{essere la
coppia perfetta}, mostrare quanto si amano. Se la scena si sposta su qualcos'altro,
\textbf{si rifà}, all'infinito, ``come nel giorno della marmotta'', finché non
parla solo di quanto si amano. Esempio: lei entra in bagno, la porta si rompe e
lei rimane chiusa dentro, lui è dall'altra parte; se lui comincia a riparare la
porta, si rifà la scena. ``Una paradossalità esasperata, quasi''.

\subsection{Fat Man Down}

Il secondo è introdotto con un \emph{content warning}: \textbf{Fat Man Down} ``è
un gioco molto pesante''. Non è il più pesante dei Jeepform, ma è il più pesante
di cui si parla nella lezione. È un LARP dello stile Jeepform, quindi un chamber
LARP tendenzialmente freeform.
\begin{itemize}
  \item Il personaggio principale è un uomo di successo, con un buon lavoro e
        una buona posizione economica, ma \textbf{obeso}.
  \item Tutti i partecipanti \textbf{si pesano su una bilancia}, e chi pesa di
        più interpreta questo personaggio.
  \item Il gioco è diviso in scene, ognuna con un \emph{framing} che è uno
        \emph{slice of life}: un'ora in ufficio, la strada da casa al
        supermercato, dieci minuti di pausa al parco.
  \item Tutte le altre persone in scena hanno la \textbf{direttiva meccanica} di
        farlo stare malissimo.
  \item Il protagonista ha una meccanica ``importantissima'': \textbf{tre
        cioccolatini} in tasca. In qualunque momento può scartarne uno e
        mangiarlo, e quando lo fa tutti i personaggi in scena devono trattarlo
        come la persona migliore del mondo.
\end{itemize}
Il personaggio è messo sotto \textbf{costante stress narrativo} finché non è
spinto a pensare ``devo mangiare per stare meglio''. ``La roba è molto pesante'',
e ci si ferma qui.

\subsection{Le meccaniche di sicurezza}

Avere un manifesto così noto nella scena nordica, con giochi molto pesanti, molto
intimi e soprattutto \textbf{senza un filtro fantasy o fantascientifico}, cioè
su drammi quotidiani, ha generato e normalizzato cose poi rivelatesi molto utili,
come le \textbf{meccaniche di sicurezza}. Se qualcuno ne ha sentito parlare, la
scena nordica è ``la principale responsabile'' della normalizzazione, della
ricerca e dell'iterazione delle meccaniche di sicurezza nel gioco di ruolo. Se ne
parlerà nella lezione sull'emotional design.

% ---------------------------------------------------------------------
\section{Il 360° e l'alta risoluzione}
\label{a2:sec:360}

\slidefig{img/AGD/agd2-360}{Il 360° e l'alta risoluzione: Ground Zero e Monitor Celestra.}{slide \emph{Analog -- Case Studies}}{fig:agd2-360}


Il Nordic LARP non è solo chamber, freeform e Jeepform. C'è anche ``quasi
l'opposto'': i LARP \textbf{360°} o ad \textbf{alta risoluzione} (\emph{high
resolution})\footnote{La trascrizione dice ``360 o High Resolution'' e, più
avanti, ``IRES 360'', ``iRes'', ``Ires''.}. Lo scopo è essere \textbf{il più
fisicamente coinvolgenti possibile}: tutto quello che si ha attorno deve essere
\textbf{diegetico}. Ma la soluzione non è quella del Dogma 99 (ambientare nel
quotidiano, così non serve la magia): ``la soluzione è buttare al muro un budget
talmente alto che la realtà a quel punto sembra magia''. Niente oggetti di scena
finti o armi finte: anche la scenografia è curata al punto che sembra di stare
davvero nel gioco. I casi studio sono ``anche abbastanza storici''.

\begin{description}
  \item[Ground Zero (1998--2001).] La premessa è la \textbf{Guerra Fredda}.
        Una cittadina americana riceve l'allarme per una bomba nucleare, e i
        partecipanti vivono \textbf{48 ore di fila}, non stop, chiusi in uno
        \textbf{scantinato} al buio, senza poter uscire. Il gioco si conclude
        con la bomba che cade: dietro le finte pareti di cartone dello
        scantinato c'erano dei \textbf{subwoofer giganti}, e alla fine i muri
        cadono addosso ai giocatori con un rumore devastante.
  \item[Panopticorp (2003, 2013).] Due \emph{run}, anche questo di 48 ore non
        stop: in teoria si dorme in gioco, ``ma in questo caso non si dormiva''.
        È una \textbf{distopia corporate} spinta al massimo: si è dipendenti di
        un'agenzia pubblicitaria sull'orlo della chiusura, che deve licenziare
        ``a manetta'' e tenere solo i migliori. Una spinta ultra competitiva e
        ``ultra tossica'' per vedere fin dove si è disposti a spingersi eticamente
        per restare in azienda, dormendo sotto le scrivanie.
  \item[Monitor Celestra.] Un 360° \textbf{fantascientifico} basato su
        \emph{Battlestar Galactica}: per giocarlo è stato noleggiato un
        \textbf{incrociatore svedese} in disarmo. Una nave, non un'astronave,
        ``ma se non usciamo all'esterno, ci mettiamo un paio di LED, ed ecco
        qua'': per i giocatori è un'astronave. Il gioco durava circa
        \textbf{60 ore}, con cabine allestite in modo che si dormisse in
        personaggio.
  \item[Demeter.] Un 360° sul viaggio della bara di \textbf{Dracula} dalla
        Transilvania a Londra, come nel romanzo di Stoker, in cui la nave
        \emph{Demeter} arriva in porto con tutti morti a bordo e il capitano
        legato al timone. I giocatori sono i passeggeri, e \textbf{da giocatori
        sanno che moriranno tutti}: per design è impossibile sopravvivere,
        ``altrimenti Dracula non esisterebbe come romanzo''. Si gioca in
        costume su una \textbf{nave vera} che, naturalmente, non va dalla
        Transilvania a Londra ma gira nel mare della Germania del nord;
        narrativamente, però, si gioca quel viaggio.
  \item[Bunker 101.] Un caso \textbf{italiano}, ispirato alla proprietà
        intellettuale di \emph{Fallout}: si giocava \textbf{72 ore} in un bunker.
        Qui il dormire non era diegetico: nell'ultima run hanno separato la parte
        di gioco dal sonno. Per tutto il resto, però, valeva il principio 360°:
        tutto ciò che si vede deve essere diegetico.
\end{description}

% ---------------------------------------------------------------------
\section{Tape LARP e black box}
\label{a2:sec:blackbox}

\slidefig{img/AGD/agd2-tape}{Tape LARP e black box LARP.}{slide \emph{Analog -- Case Studies}}{fig:agd2-tape}


Accanto al Jeepform (``in cui giochiamo in casa con i nostri costumi'') e al
360° (``l'ultra-diegesi''), ci sono tipologie scandinave ``molto più particolari
e quasi astratte''.

\subsection{Il tape LARP}

Visto che qualcuno conosce Lars von Trier: chi ha visto \emph{Dogville}, il film
con lo scotch per terra, sa da dove nascono i \textbf{tape LARP}. ``Abbiamo visto
Dogville, ci ispira un casino'': giochi in cui \textbf{lo scotch per terra indica
gli ambienti}, con tutte le meccaniche che ne derivano (il personaggio vede e
sente cose che normalmente non potrebbe sentire; come si agisce e non si agisce
da giocatori\ldots). Ne sono stati fatti parecchi, in particolare in Danimarca.

\subsection{Il black box LARP}

Ruggero è ``un particolare amante'' del \textbf{black box LARP}, che prende il
nome dalle \emph{black box}: per chi ha fatto teatro ``dalla Germania in su'',
spazi teatrali di solito senza spalti né sedie, con \textbf{pareti, soffitto e
pavimento neri}, molto diffusi in Germania e nel Nord Europa. A
\textbf{Copenaghen} c'è un festival, il \textbf{Black Box LARP Festival}, che usa
questi spazi per giochi creati apposta: giochi di ruolo che usano \textbf{luci e
colori come meccaniche di gioco}. Gli esempi della slide:

\begin{description}
  \item[Borderline.] Un gioco ``estremamente astratto''. Si gioca \textbf{in
        coppie}, e ogni coppia gioca \textbf{un solo personaggio}, con dei
        conflitti interni: tutta la conversazione è un monologo interiore del
        personaggio, e i due giocatori, con una candela, lottano per portare ciò
        che il personaggio fa dalla propria parte o da quella dell'altro.
  \item[Waiting for Flight.] Ruggero lo ha giocato. Si aspetta un aereo con a
        bordo persone care, amici, colleghi; \textbf{l'aereo precipita}, e
        muoiono quasi tutti. Si gioca la \textbf{sala d'attesa} di chi li
        aspettava: è un gioco sull'attesa, ``sul logorante non avere
        informazioni'' finché vengono centellinate. Si conclude con un momento
        molto catartico: arrivano le \textbf{valigie} delle persone attese, e in
        alcune (3 su 20) è nascosto un piccolo \textbf{cuore di vetro}: se lo
        trovi, la persona si è salvata; se no, non si è salvata. C'è una fase in
        cui speri che ci sia, ma potrebbe essere nascosto: ``è una cosa
        lacerante. Abbiamo pianto in 19 su 20''.
  \item[Stranger.] Di una designer danese\footnote{La trascrizione dice ``Nino
        Ressendrop'': con ogni probabilità Nina Runa Essendrop, designer danese
        nota per i LARP fisici e non verbali. Anche il titolo, ``Stranger'', è
        riportato come nella trascrizione.}. È un black box \textbf{non verbale}:
        è vietato parlare. I partecipanti sono divisi in due gruppi, i
        \textbf{blu} e i \textbf{verdi}, e prima del gioco ciascun gruppo crea un
        \textbf{linguaggio gestuale condiviso}: come la propria nazione esprime
        felicità, rabbia, paura, tristezza, amicizia, odio. Il twist: i verdi
        devono avere gesti squadrati, rigidi, geometrici; i blu gesti curvi,
        lenti, morbidi, sinuosi. Il gioco è sull'\textbf{immigrazione forzata}:
        ci sono fasi di giorno e di notte, e di notte chi facilita porta alcuni
        verdi nella nazione blu, dove si risveglieranno, con la regola che il
        confine non si può riattraversare. Le dinamiche: il primo che si
        risveglia tra i blu impara facilmente la loro lingua, e i blu lo aiutano
        con un \emph{onboarding} non verbale; ma quando arrivano molti verdi,
        tendono a muoversi \textbf{in gruppi}, perché a quel punto c'è già chi
        parla la propria lingua, anche se non è la lingua della nazione. ``Un
        microscopio'' di dinamiche della vita reale: emigrazione, immigrazione,
        inclusione, e quanto è più facile essere inclusi se il sistema ti
        supporta.
  \item[Il gioco di Ruggero.] Uno dei giochi della slide è suo, ``quindi non ne
        parliamo'', nonostante le proteste dei colleghi.
  \item[Un gioco su Cuba.] Di una designer che è andata a vivere a
        \textbf{Cuba} e ha partecipato da vicino alla vita di una delle storiche
        comunità queer cubane\footnote{La trascrizione riporta il titolo come
        ``Socialism and Word''; non è verificabile. Anche le date sono incerte
        (``nel 97, 89 e 91''): il contesto, la caduta del Muro di Berlino, fa
        pensare al 1987, 1989 e 1991.}. Il gioco racconta che cosa è successo
        alle comunità \textbf{punk, rock e queer}, molto sovrapposte, prima,
        durante e dopo la caduta del Muro di Berlino: lo Stato, rimasto senza i
        fondi sovietici, decide chi è cittadino di serie A e di serie B, e le
        persone queer e punk (la musica occidentale) vengono picchiate e
        imprigionate. Erano anche gli anni dello scoppio dell'\textbf{AIDS}. Una
        comunità famosa, quella dei \emph{frikis}\footnote{La trascrizione dice
        ``Os Fricchios''.} cubani (di cui sono ancora vive due persone), decise:
        ``tanto ci ammazza la polizia, tanto vale iniettarci sangue infetto da
        HIV'', per andare a morire non per strada ma nei \textbf{centri di
        recupero}, gestiti da medici stranieri, uno dei pochissimi casi in cui a
        Cuba poteva esserci personale non cubano. Lì si moriva, ma ``almeno si
        poteva ascoltare musica punk'', magari con il dottore che arrivava dalla
        Florida con le cassette rock, e tra simili. ``Una roba drammaticissima'':
        anche questo è un gioco devastante.
\end{description}

\begin{table}[htbp]
\centering
\small
\begin{tabular}{@{}p{3.1cm}p{9.8cm}@{}}
\toprule
\textbf{Tipo di LARP} & \textbf{Caratteristiche ed esempi citati} \\
\midrule
Boffer & nato per adattare D\&D dal vivo; costume, armi in plastazote,
combattimento e trame epiche; fino a centinaia o migliaia di persone \\
Mind's Eye Theatre & da \emph{Vampire: The Masquerade}; costume, dark/gotico,
intrigo politico più che combattimento \\
Chamber / parlor LARP & senza costume, accessibile, budget ridotto, spesso
\emph{print and play}, 3--20 persone (\emph{A cena col nonno}) \\
American freeform & chamber LARP con metatecniche prese dal gioco da tavolo e
dal GDR; anche in carte o libretti \\
Murder party & cena con delitto (gli attori indagano, l'assassino è un attore)
o murder party ``all'inglese'' (si interpretano i sospettati) \\
Nordic LARP & coinvolgimento, emozione, coesione tematica, cura del dopo-gioco;
manifesti; Knutpunkt dal 1997; Dogma 99 \\
Jeepform & freeform nordico sul dramma quotidiano (il gioco della coppia
perfetta, \emph{Fat Man Down}); meccaniche di sicurezza \\
360° / alta risoluzione & tutto diegetico grazie a un budget altissimo
(\emph{Ground Zero}, \emph{Panopticorp}, \emph{Monitor Celestra},
\emph{Demeter}, \emph{Bunker 101}) \\
Tape LARP & lo scotch per terra indica gli ambienti, come in \emph{Dogville} \\
Black box & spazi teatrali neri; luci e colori come meccaniche
(\emph{Borderline}, \emph{Waiting for Flight}, \emph{Stranger}) \\
\bottomrule
\end{tabular}
\caption{Le tipologie di gioco di ruolo dal vivo presentate da Ruggero, in
sintesi.}
\label{a2:tab:larp}
\end{table}

% ---------------------------------------------------------------------
\section{Chi osserva? Il LARP e il teatro dell'oppresso}
\label{a2:sec:oppresso}

\paragraph{La domanda di Tiziana.} Tiziana confessa un'impressione: tutti questi
giochi ``sembrano cose da studi'', come se ci fosse qualcuno che sta sopra e
osserva le dinamiche per ricavarne informazioni. Ruggero la rassicura: \textbf{non
è così}. I designer di questa scena, specialmente quella black box, vengono
dall'ambiente ibrido del teatro interattivo, degli \emph{happening}, dello
psicodramma, e \textbf{non documentano nulla di quel tipo}: non registrano le
giocate, non estrapolano informazioni da chi gioca. Spesso sono giochi per poche
persone (è difficile che un black box superi le 20), e molti hanno avuto
\textbf{run singole}, o al massimo due o tre. ``Si tratta solo di fare un certo
tipo di esperienza''. Come in altri media esistono generi con scopi diversissimi
(un film e un documentario sono entrambi operazioni visive), anche qui gli scopi
sono diversi. E chi non ha mai incontrato giochi del genere si chiederà
``perché dovrei giocare una roba del genere?'': la risposta arriverà con la
lezione sull'emotional design (``lo so, ve l'ho detto quattro volte'').

\paragraph{Il teatro dell'oppresso.} Prima di passare la parola, uno studente
interviene: in queste cose ritrova molto il metodo del \textbf{teatro
dell'oppresso}\footnote{Il metodo teatrale ideato dal regista brasiliano Augusto
Boal (integrazione).}. ``Ah, certo''. Ruggero torna indietro di una slide, a
quella sulle \textbf{influenze fondanti del Nordic LARP}: ci sono le tecniche di
\textbf{Moreno} (sociodramma e psicodramma), il \emph{sensitivity training}, il
\textbf{roleplay come metodo di ricerca}, e tutti quei metodi che usano il teatro
e l'interazione come strumento per far emergere il personale e l'intimo. Lo
studente precisa che il teatro dell'oppresso nasceva per \textbf{trovare
soluzioni} a situazioni che apparentemente non ne hanno: a turno il pubblico
proponeva una soluzione, si rifaceva sempre la stessa scena, e si studiava come
reagivano le persone; quindi, in un certo senso, il comportamento delle persone
\emph{si osserva}, ma dall'interno della scena.

Per Ruggero, nella scena LARP questa dimensione ``viene un po' meno'', perché lo
scopo resta \textbf{ludico} (``e scopo ludico può voler dire molte cose, come
vedremo''): è meno pragmatico del ``prendiamo quello che abbiamo giocato e il
gioco ci dà già gli strumenti per applicarlo''. Alcuni \textbf{debriefing}, i
momenti strutturati del dopo-gioco, possono dare strumenti a chi ha giocato; ma
questi giochi nascono come esperienze che devono essere ``ludicamente
interessanti per chi gioca'', \textbf{non come esperimenti scientifici ``con dei
topi dentro un labirinto''}.

\paragraph{Approfondimento: il LARP in classe e il debriefing.} Che il LARP
possa essere uno strumento formativo, come nel volume sull'edu-larp citato
all'inizio, lo mostra anche \emph{Challenges for Game Designers}. Nel capitolo
sui giochi come strumento di insegnamento, Brathwaite e Schreiber propongono una
sfida, \emph{Living History}: progettare un LARP che rappresenti un evento storico
(una convenzione costituzionale, un protocollo internazionale, una guerra)
giocabile da una classe di 20--30 studenti in una o due ore. I passi suggeriti
sono: scegliere un evento in cui molte parti, ciascuna con la propria agenda, si
incontrano in una negoziazione o in un conflitto; scegliere i giocatori (o le
squadre); definire gli \textbf{obiettivi} di ciascuno; creare regole semplici e
rapide per risolvere i conflitti, perché i giocatori si muovono e parlano di
continuo (``alcuni LARP usano un singolo tiro di dado o una partita a
sasso-carta-forbici''); e preparare, oltre alle schede dei personaggi, delle
\textbf{note di debriefing}: che cosa è successo davvero e quali domande
discutere confrontando la storia con l'esito della partita. Lo scopo dichiarato è
che gli studenti capiscano che ``l'evento era complicato'' e come l'esito sia
\emph{emerso} dalla situazione iniziale. È la versione didattica della cura del
dopo-gioco che Ruggero attribuisce alla scena nordica.

% ---------------------------------------------------------------------
\section{I giochi di ruolo da tavolo: quattro parole per orientarsi}
\label{a2:sec:gdr}

``Daniele, ti lascio la parola''. Daniele (``mi giri tu le slide quando te le
chiedo'') avverte che fare una panoramica dei giochi di ruolo in poco tempo è
``veramente difficile'': ce ne sono tantissimi, e negli anni si sono evoluti e
diversificati sempre di più, fino a diventare molto diversi tra loro.

\subsection{Il tribalismo e le regole tramandate a voce}

C'è anche un altro problema, che si vedrà nella lezione sulla fiction e che
esiste in tutto il gioco analogico: il \textbf{tribalismo}, cioè
l'\textbf{ereditarietà orale delle regole}. Spesso i giochi di ruolo venivano
insegnati da altre persone, anche per motivi economici. Daniele ha cominciato
così: da bambino, giocando per la prima volta a D\&D, l'unico ad avere il
manuale era il master; ``ergo io non avevo accesso alle regole'', e il master
poteva inventarsi quello che voleva.

\subsection{Conflitto, azione, risoluzione, effetto}

\slidefig{img/AGD/agd2-car}{Conflitto, azione, risoluzione, effetto: le quattro parole per orientarsi nei giochi di ruolo da tavolo.}{slide \emph{Analog -- Case Studies}}{fig:agd2-car}


Per capire un po' i giochi di ruolo bisogna ricordare \textbf{quattro parole},
che torneranno nella lezione sulla fiction: \textbf{conflitto}, \textbf{azione},
\textbf{risoluzione} ed \textbf{effetto}. Nella slide \emph{conflitto} e
\emph{risoluzione} hanno un asterisco: per ora basta capirli quanto serve a
leggere i casi studio, senza analizzarli troppo.
\begin{description}
  \item[Conflitto.] L'ostacolo, l'incertezza, l'evento che \textbf{stimola
        un'azione o la contrasta}. ``Tornerà molto''.
  \item[Azione.] Quello che si descrive di voler fare nei panni del personaggio
        per risolvere un conflitto.
  \item[Risoluzione.] Come si risolve l'azione dentro il gameplay, con le regole:
        che cosa succede?
  \item[Effetto.] La conseguenza che si applica.
\end{description}
L'esempio: ``voglio colpire il goblin con la mia spada''. Il goblin ha il suo
scudo e la sua armatura. Che azione fai? ``Lo becco con la spada''. Si tirano i
dadi, si risolve, si applica un effetto. Ma non è così semplice, perché in
mezzo c'è la \textbf{fiction}: nel gioco di ruolo si tende a riscrivere tutto
quello che succede nel gameplay \emph{dentro} la narrazione. Se faccio tre danni
agli HP del goblin, che cosa vuol dire narrativamente? Come li descrivo? Da qui
nasce ``una serie di problematiche''.

\begin{figure}[htbp]
\centering
\begin{tikzpicture}[
    font=\small,
    passo/.style={draw, thick, rounded corners=3pt, fill=gray!12,
                  minimum width=2.6cm, minimum height=0.95cm, align=center},
    esempio/.style={font=\footnotesize\itshape, align=center, text width=2.8cm},
    freccia/.style={-{Stealth[length=2.5mm]}, thick}
  ]
  \node[passo] (c) at (0,0)    {\textbf{Conflitto}$^{*}$};
  \node[passo] (a) at (3.4,0)  {\textbf{Azione}};
  \node[passo] (r) at (6.8,0)  {\textbf{Risoluzione}$^{*}$};
  \node[passo] (e) at (10.2,0) {\textbf{Effetto}};
  \draw[freccia] (c) -- (a);
  \draw[freccia] (a) -- (r);
  \draw[freccia] (r) -- (e);
  \node[esempio, below=0.25cm of c] {il goblin, con scudo e armatura};
  \node[esempio, below=0.25cm of a] {``lo colpisco con la spada''};
  \node[esempio, below=0.25cm of r] {tiro dei dadi secondo le regole};
  \node[esempio, below=0.25cm of e] {tre danni agli HP};
  \node[draw, dashed, rounded corners=3pt, fit=(c)(e), inner sep=0.5cm,
        yshift=-0.35cm, minimum height=2.6cm] (fic) {};
  \node[font=\footnotesize, anchor=south] at (fic.north)
    {\textbf{fiction}: tutto viene ridescritto nella narrazione (che cosa
    \emph{sono}, nella storia, tre danni?)};
\end{tikzpicture}
\caption{Le quattro parole di Daniele per leggere i giochi di ruolo da tavolo,
con l'esempio del goblin. Gli asterischi sono quelli della slide: conflitto e
risoluzione verranno analizzati meglio nella lezione sulla fiction. Nel gioco di
ruolo ``classico'' il conflitto precede l'azione, la risoluzione è un tiro di
dado e l'effetto è descritto dal master e dal regolamento; molti giochi
successivi cambiano proprio questo schema.}
\label{a2:fig:car}
\end{figure}

% ---------------------------------------------------------------------
\section{Le origini: da Dungeons \& Dragons agli anni Novanta}
\label{a2:sec:gdr-storia}

Le quattro parole servono a Daniele per una breve panoramica, anche storica,
dei giochi di ruolo.

\subsection{Dungeons \& Dragons e il gioco di guerra}

Il primo gioco di ruolo da tavolo ``così come viene concepito'' è sicuramente
\textbf{\emph{Dungeons \& Dragons}}, di \textbf{Gygax e Arneson}, nel
\textbf{1974}. C'è in realtà un precedente, \textbf{\emph{Chainmail}}\footnote{La
trascrizione dice ``C'è in Mail'' e data \emph{Chainmail} al '72; il gioco è
del 1971.}, e come per tutte le origini (``anche \emph{Cyberpunk} ha lo stesso
problema'') c'è un po' di ambiguità su che cosa venga prima. Ci sono libri che
raccontano la nascita di D\&D, c'è la storia della \textbf{TSR}, la casa
editrice che lo pubblicò per prima, ci sono versioni discordanti e, negli anni,
parecchi contrasti tra le famiglie di Gygax, di Arneson e di altri.

D\&D viene dal \textbf{gioco di guerra}, il gioco di \textbf{miniature}. ``La
leggenda vuole'' che Gygax e Arneson, giocando, si siano detti: ho battuto il
tuo esercito; che cosa succede se il mio generale entra nel tuo castello e va a
trattare con il tuo generale? Abbiamo regole per questo? No: inventiamocele.
``Ripeto, una leggenda'': non si sa se sia andata così, e in fondo non interessa
molto.

\subsection{I generi e la struttura del gioco di ruolo ``early''}

\slidefig{img/AGD/agd2-dnd}{Il flusso di Dungeons \& Dragons letto con conflitto, azione, risoluzione ed effetto.}{slide \emph{Analog -- Case Studies}}{fig:agd2-dnd}


Quello che interessa è che dopo D\&D, nell'\textbf{RPG delle origini}, vengono
coperti più o meno \textbf{tutti i generi letterari}: \emph{Call of Cthulhu}
l'orrore lovecraftiano, \emph{Vampire: The Masquerade} l'orrore gotico,
\emph{Traveller} la fantascienza, \emph{Cyberpunk} il cyberpunk, e \emph{GURPS},
un \textbf{sistema generico} con cui ``teoricamente'' si può giocare qualunque
genere.

Nel gioco di ruolo classico:
\begin{itemize}
  \item il \textbf{conflitto} poteva venire prima o dopo l'\textbf{azione}:
        ``voglio colpire quel goblin'' (e il goblin ha la sua difesa), oppure
        ``corri, ti sta inseguendo una palla'', come in \emph{Indiana Jones};
  \item la \textbf{risoluzione} era il tiro di dado, e l'\textbf{effetto} veniva
        descritto dal master e dal regolamento;
  \item in pratica, \textbf{tutto tranne l'azione}, che il giocatore gestiva e
        narrava, era gestito dal \textbf{master}.
\end{itemize}

\paragraph{Il master come arbitro.} All'inizio il master non era un ``maestro''
né una figura che insegna: era interpretato soprattutto come un
\textbf{arbitro}. Nella prima versione di D\&D non c'è ancora l'idea, venuta
dopo, del master come ``dio del tavolo'': quando deve decidere, dovrebbe farlo
secondo criteri di realismo e oggettività (``puoi essere oggettivo a un tavolo
di gioco? La voglio vedere''). Questa idea è importantissima perché verrà ripresa
dall'\textbf{OSR}, l'\emph{Old School Revival}, il movimento che cerca di
riportare nella modernità il modo in cui, a suo dire, si giocava D\&D agli
inizi.

\paragraph{Il master è un giocatore.} Nei giochi di ruolo ci sono quindi
\textbf{ruoli asimmetrici}. Ma Daniele raccomanda di ricordare una cosa, che
tornerà nella lezione sulla fiction: \textbf{il master è un giocatore}, magari un
giocatore con più regole e più autorità al tavolo, ma ``un giocatore tanto quanto
i giocatori che interpretano un personaggio''. Nei GDR delle origini: si
interpreta un personaggio, con una \textbf{scheda} e le sue statistiche, dentro
una fiction e un \textbf{setting} stabilito dal manuale o dal master; quando si
dice che cosa si vuole fare, il master reagisce. E lentamente, con gli anni, il
master assume anche il ruolo di \textbf{regista} e \textbf{sceneggiatore}: scrive
l'avventura, decide che cosa succede, decide i tagli di scena.

\paragraph{Approfondimento: dal wargame al contratto sociale.} La genealogia
raccontata da Daniele si ritrova in \emph{21st Century Game Design} di Bateman e
Boon. Le radici del gioco di strategia, scrivono, precedono i videogiochi di
decenni: alla fine degli anni Cinquanta Charles Roberts fonda Avalon Hill e
pubblica i primi \textbf{wargame da tavolo} (1958); in un'epoca senza
videogiochi, giochi di ruolo e giochi di carte collezionabili, i wargame
``fondano l'industria dei \emph{hobby games}'', e \emph{Chainmail} (1971),
ispirato a quei giochi, è il ``trampolino'' verso \emph{Dungeons \& Dragons}. Gli
stessi autori descrivono la forma minima del gioco di ruolo da tavolo come un
\textbf{contratto sociale} tra l'arbitro (il \emph{games master}) e i
giocatori: i giocatori accettano che sia l'arbitro a decidere come le loro azioni
sviluppano la storia, e l'arbitro si impegna a garantire intrattenimento o
imparzialità, come mediatore e narratore principale. Citano anche un'osservazione
attribuita a Gygax: se i giocatori di ruolo si accorgessero di poter fare tutto
da soli, i game designer resterebbero senza lavoro. Il fatto, notano, è che la
maggior parte dei master ha bisogno di meccaniche per gestire il mondo, e i
giocatori si fidano di più dell'arbitro quando una struttura di meccaniche
casuali, i dadi, fa da cornice. Anche Schell, in \emph{The Art of Game Design},
cita il master dei giochi di ruolo da tavolo come esempio di stato del gioco
\emph{nascosto}: è lui a conoscere privatamente gran parte dello stato del gioco,
perché ne è ``il meccanismo operativo''.

\subsection{Gli anni Ottanta e Novanta: adattare un regolamento a tutto}

Qui Daniele fa un ``piccolo inciso editoriale''. Nella grande espansione degli
anni Ottanta e Novanta \textbf{non c'era un professionismo} del settore. Gli
editori compravano delle proprietà intellettuali, per esempio \textbf{007},
prendevano un master appassionato e competente e gli chiedevano di
\textbf{adattare} il regolamento di D\&D, o un altro regolamento
esistente\footnote{La trascrizione dice ``Call of Duty 2'', che non può essere
corretto in questo contesto; il senso è ``il regolamento di un altro gioco
già esistente''.}, a quella IP. A volte nascevano giochi innovativi, a volte
giochi ``non molto interessanti da giocare'': immaginate di giocare con le
regole di D\&D qualcosa che quelle regole non supportano. Verso la fine degli
anni Novanta il problema della qualità ``ha iniziato a farsi sentire''.

\paragraph{Approfondimento: James Bond 007 e l'esperienza essenziale.} Proprio
il gioco di ruolo di 007 è il caso con cui Schell illustra, in \emph{The Art of
Game Design}, la \textbf{lente dell'esperienza essenziale} (\emph{Lens of
Essential Experience}), e mostra l'altra faccia della storia di Daniele. Il
designer Chris Klug era frustrato dai precedenti giochi di ruolo di spie, come
\emph{Top Secret} della TSR, perché ``si giocavano troppo come dei wargame'':
l'essenza di ciò che rende eccitanti i film di spionaggio non c'era. Per il
gioco di James Bond, Klug progettò le meccaniche in modo che restituissero in
ogni modo possibile la sensazione dei film. L'esempio è l'invenzione degli
\textbf{Hero Points}: nei GDR tradizionali, davanti a un'azione rischiosa (saltare
da una finestra su un elicottero in volo), il master calcola una probabilità, il
giocatore tira i dadi ``e questo è quanto''; ma se la probabilità è bassa
nessuno rischia, e se è alta tutti si comportano da supereroi. I punti eroe
risolvono il dilemma in modo coerente con il genere. È esattamente quello che
la New Wave rivendicherà: \textbf{il sistema conta}, e un regolamento adattato ``a
nastro'' non restituisce l'esperienza di un genere.

% ---------------------------------------------------------------------
\section{La New Wave e The Forge}
\label{a2:sec:gdr-newwave}

\subsection{Il sistema conta}

La slide successiva è sulla \textbf{New Wave}: un movimento nato tra la fine
degli anni Novanta e l'inizio dei Duemila. Daniele avverte che non vuole fare
un'analisi geografica ma ``un riassunto della storiografia''. La New Wave cerca
un \textbf{metodo diverso di interpretare il gioco di ruolo}, e la prima cosa
che stabilisce, in un forum chiamato \textbf{The Forge} (c'è anche un libro in
inglese su questo movimento), è che \textbf{il sistema dei giochi di ruolo è
importante}: le regole, il gameplay, la parte meccanica \textbf{non sono un
orpello}. Non si può tenere lo stesso sistema, cambiare il setting e sperare che
funzioni da sé. Tanto è vero che perfino \emph{GURPS}, il sistema generalista,
aveva tantissimi manuali di espansione delle regole per ogni genere:
\emph{GURPS} base, \emph{GURPS Fantasy}, \emph{GURPS} fantascienza\ldots

\subsection{I giochi della New Wave}

\slidefig{img/AGD/agd2-newwave}{Alcuni giochi della New Wave.}{slide \emph{Analog -- Case Studies}}{fig:agd2-newwave}


Dentro la New Wave nascono esperienze diversissime, che spostano l'attenzione
su due cose: la \textbf{scrittura del manuale}, diversa da quella tradizionale,
e la \textbf{rottura degli schemi}.
\begin{itemize}
  \item \textbf{\emph{Apocalypse World}}\footnote{La trascrizione dice ``A
        Polly's Ward'' e ``Minster Baker''.}, di \textbf{Vincent Baker}, cerca di
        scardinare il modo classico di strutturare un manuale di gioco di ruolo
        (Daniele accenna a una nuova edizione recente).
  \item \textbf{\emph{Nobilis}}, uno dei primi giochi abbastanza famosi
        \textbf{diceless}, senza dadi: azioni e conflitti si risolvono in un
        altro modo.
  \item \textbf{\emph{Don't Rest Your Head}} (in italiano \emph{Non cedere al
        sonno}) e \textbf{\emph{Polaris}}, che lavorano sul concetto di chi è e
        chi non è il personaggio. \emph{Non cedere al sonno} è un gioco
        ``surrealista'' in cui ci si sveglia a \textbf{Mad City}; la parte
        importante sono i \textbf{poteri}. All'epoca \emph{Vampire} aveva
        stabilito che i poteri fossero una lista di possibilità con statistiche
        a cui applicare le regole; in \emph{Don't Rest Your Head} ci sono poteri
        ``totalmente assurdi, fenomenali'', che non richiedono meccaniche in più
        ma usano le meccaniche base ``in maniera molto innovativa''.
  \item \textbf{\emph{Fate}}, un sistema anch'esso generalista, ``più o meno
        come \emph{GURPS}'', ma con delle sue specifiche: ci sono
        \textbf{procedure per creare i setting} e per stabilire quali regole un
        setting richiede.
  \item \textbf{\emph{Dread}}, uno dei primi giochi a usare per l'incertezza
        \textbf{un elemento che non è un dado}: una \textbf{torre Jenga}. È un
        gioco dell'orrore in cui ogni volta che si fa un'azione si toglie un
        blocchetto, e se la torre cade il personaggio muore (``non è proprio
        così, è più complicato, però per semplificare'').
\end{itemize}

\paragraph{Le guerre di religione.} La New Wave ha avuto ``una serie di
problematiche'' in giro per il mondo, e soprattutto in Italia ha creato una
sorta di \textbf{guerra teologico-filosofica} tra i ``tradizionalisti'' di D\&D e
i ``forgiti'' della New Wave. ``Per vostra sanità mentale non ci entriamo''.

\paragraph{Una reazione a due cose.} La New Wave è stata una reazione:
\begin{enumerate}
  \item ai manuali in cui ``c'è un sistema e quello è il sistema'', con il
        setting applicato sopra ``a nastro'', senza intreccio tra la narrazione
        e le regole che la raccontano: si voleva \textbf{coerenza};
  \item all'assenza di studio: è stato il \textbf{primo movimento a voler
        studiare i giochi di ruolo}, a capire come sono fatti. Da qui viene il
        \textbf{Big Model}, un modello di analisi dei giochi di ruolo di cui si
        vedranno alcune cose nella lezione sulla fiction.
\end{enumerate}
Ma il Big Model ``non è fondamentale'', perché per Daniele uno dei problemi della
New Wave è aver creato \textbf{un dizionario collaterale} a quello del game design
del gioco da tavolo e del videogioco: termini diversi che dicono la stessa cosa.

\subsection{The Forge e la scena nordica}

Ruggero fa da ponte: \textbf{The Forge ha fatto per il gioco di ruolo da tavolo
quello che la scena scandinava ha fatto per il LARP}. Inventavano e scrivevano
moltissimo, hanno lasciato framework e strumenti. La scena scandinava è nata un
po' prima, The Forge un po' dopo, ma hanno agito in parallelo. La differenza è
che The Forge era \textbf{meno strutturato}: essendo un forum, era più caotico,
con discussioni anche accese, soprattutto sulle interpretazioni. Anche la
\textbf{preservazione storica} è diversa: della scena nordica ci sono i libri
fisici, mentre per il forum bisogna passare dalla Wayback Machine. Ma in termini
di normalizzazione di termini, meccaniche e framework, entrambi hanno fatto
moltissimo.

\subsection{Fiasco}

\slidefig{img/AGD/agd2-fiasco}{Il flusso di Fiasco: chi stabilisce la scena e chi la risolve.}{slide \emph{Analog -- Case Studies}}{fig:agd2-fiasco}


Daniele riprende: a un certo punto The Forge chiude, e nascono altri giochi
dagli stimoli e dalle discussioni che quell'ambiente aveva generato, non
necessariamente da persone che ne facevano parte. Per esempio quelli di
\textbf{Jason Morningstar}, noto anche per dei chamber LARP (il migliore amico
di Daniele, Edoardo, ``più esperto di me'', ne è talmente fan che il suo gatto si
chiama Giasone).

\textbf{\emph{Fiasco}} è per Daniele uno dei primi giochi che \textbf{rompe il
canone}. Racconta, appunto, un fiasco, come nei film dei \textbf{fratelli Coen}
o di Tarantino: c'è uno scenario, ci sono dei personaggi, e ogni volta che si fa
una scena \textbf{il giocatore di turno sceglie se aprirla} (e gli altri la
risolvono) \textbf{o se lasciare che gli altri la impostino} (e risolverla lui).
Già solo questo ``micro elemento'' (non l'unica meccanica: c'è anche quella,
interessantissima, dei dadi) rompe l'idea che in un gioco di ruolo si debba
descrivere l'azione, tirare i dadi e vedere che cosa succede, oppure reagire a
ciò che descrive il master. Per chi veniva da \emph{AD\&D} seconda edizione o da
\emph{Vampire} era ``assolutamente interessante'': più leggero e più legato alla
fiction. Anche i \textbf{dadi} sono usati in modo nuovo: vengono tirati
\textbf{prima}, messi sul tavolo, selezionati, usati quasi come
\textbf{token}: un'innovazione anche nel modo di usare i materiali.

% ---------------------------------------------------------------------
\section{Il gioco di ruolo contemporaneo}
\label{a2:sec:gdr-oggi}

\subsection{Prima del COVID: i giochi focalizzati}

\slidefig{img/AGD/agd2-contemporary}{Giochi di ruolo contemporanei.}{slide \emph{Analog -- Case Studies}}{fig:agd2-contemporary}


Daniele salta alla contemporaneità. Prima del COVID c'è un grande aumento di
giochi che, mentre le persone si appassionano ai giochi di ruolo, vanno o nella
direzione di D\&D o in una \textbf{direzione più ibrida}:
\begin{itemize}
  \item \textbf{\emph{Blades in the Dark}}\footnote{La trascrizione dice
        ``Blaze in the Dark''. Daniele lo data ``2016--18 più o meno''.};
  \item più recenti, il gioco di ruolo di \textbf{\emph{Cowboy
        Bebop}}\footnote{La trascrizione dice ``Cabobby Bop'' e ``Cobo
        Bobby''.} e \textbf{\emph{The Clay That Woke}} di \textbf{Paul
        Czege}\footnote{La trascrizione dice ``Paul Ksege, The Clayed Walk''.},
        autore di ``uno dei principi di design più interessanti'';
  \item \textbf{\emph{Fabula Ultima}}, di \textbf{Emanuele Galletto}, italiano
        (``che salutiamo'').
\end{itemize}
Sono prodotti che \textbf{non cercano di fare la stessa cosa di D\&D}. Hanno un
impianto simile (scheda del personaggio, dadi, azioni descritte, tiri), ma si
pongono \textbf{obiettivi di design più coerenti}: ``si chiedono che cosa vogliamo
che il giocatore senta''. Il sistema di \emph{Cowboy Bebop}, per esempio, usa
gli \textbf{orologi} (\emph{clock}) che si trovano anche in \emph{Blades in the
Dark} e nei giochi \emph{Powered by the Apocalypse}, cioè la stessa meccanica,
ma è un sistema completamente diverso, perché deve raccontare una fiction
diversa:
\begin{itemize}
  \item in \emph{Blades in the Dark} si giocano dei \textbf{criminali} a
        \textbf{Doskvol}\footnote{La trascrizione dice ``Oswald''.}, città di un
        mondo alternativo dove è sempre notte, dove i fantasmi non passano
        nell'aldilà e vivono con i vivi, e dove si è costretti a restare perché
        fuori dalle mura i fantasmi uccidono. Si fa parte di una banda che cerca
        di scalare la gerarchia della città, ``tipo \emph{Peaky Blinders}'';
  \item in \emph{Cowboy Bebop} si è ``letteralmente una ciurma di persone nello
        spazio, nel futuro, che fanno cose'', perché ``l'anime si racconta con i
        personaggi che fanno cose'' (Daniele, ``molto fan dell'anime'', lo dice
        come un complimento).
\end{itemize}

\subsection{Dopo il COVID: il ritorno all'origine}

Soprattutto dopo il COVID si è notato un ``ritorno fortissimo'' al
\textbf{giocare come una volta}. Per Daniele entrambe le tendenze rispondono a
una necessità: \textbf{semplificare}. \emph{D\&D} 3 e 3.5, le nuove edizioni di
\emph{Vampire} e altri sono diventati ``pieni di regole, pieni di contenuti'', e
chiedono a master e giocatori di imparare moltissimo, senza contare il tempo
necessario per una campagna. Allora, da un lato, si cercano sistemi più semplici
``da intavolare'' che facciano bene una cosa; dall'altro si cerca la stessa cosa
per altri motivi, ricreando le atmosfere del D\&D \emph{old school}: \textbf{alta
mortalità}, parte narrativa e interpretativa meno importante, immediatezza.
\textbf{\emph{Mörk Borg}}\footnote{La trascrizione dice ``Morg Borg''.}, per
esempio, sono circa \textbf{cento pagine}, con poche regole. C'è più attenzione
al ``realismo'' (termine un po' sbagliato, si vedrà nella lezione sulla
fiction), e il vantaggio è che questi giochi possono usare sia le proprie regole
sia quelle di \textbf{D\&D quinta edizione}, che è stata importante soprattutto
``per un motivo editoriale'' (anche questo si vedrà nella lezione sulla
fiction).

\paragraph{La domanda di Gianmarco.} Gianmarco chiede se questa tendenza a
manuali e partite ``asciutti'' segua un pubblico più adulto, con meno tempo, o
sia una tendenza generale del mercato. È soprattutto la prima. Il gioco di ruolo
ha avuto parecchie difficoltà ad attirare nuove leve, e i giocatori sono
cresciuti. Daniele gioca da molti anni, ma il suo ``vero inizio'' è stato con un
gioco della New Wave\footnote{La trascrizione dice ``Paura dell'Epoclips'':
probabilmente \emph{Il mondo dell'apocalisse}, l'edizione italiana di
\emph{Apocalypse World}, ma non è certo.}: per lui, all'università, avere in un
solo manuale da 200 pagine tutte le regole e tutti i personaggi è stato ``un
toccasana''. Chi inizia con D\&D, magari alle superiori o all'università, man
mano che cresce tende a preferire:
\begin{itemize}
  \item o giochi più focalizzati sull'aspetto \textbf{narrativo}, che hanno
        comunque scheda, tiro di dado e struttura tradizionale, ma una scheda e
        un sistema più semplici;
  \item o un'esperienza simile al D\&D canonico, ma con meno ``regoline e
        regolette'' e meno \textbf{simulazionismo}.
\end{itemize}
Due differenze pratiche. I manuali \textbf{OSR}, che un po' di simulazionismo ce
l'hanno, sono pieni di \textbf{soluzioni per l'arbitro} (``che non è il master,
è l'arbitro'') per risolvere i problemi sul momento; e l'OSR è \textbf{sandbox}:
c'è una mappa, una zona in cui si va dove si vuole, e il master non ha scritto
un'avventura. I \textbf{giochi narrativi}, dal canto loro, non hanno una campagna
scritta ma uno \textbf{scenario}. In entrambi i casi il tempo di preparazione del
master si riduce ``di gran lunga''.

\subsection{Focalizzati, non semplificati}

Ruggero interviene sulle parole: non vorrebbe che passasse l'idea che questi
giochi siano ``più semplici'' o \emph{dumbed down} rispetto a D\&D. Si parla di
\textbf{giochi focalizzati}: invece di semplificare, \textbf{focalizzano}, sono
molto più \textbf{coesi tematicamente}, hanno ``un cuore tematico molto forte'' e
si giocano per giocare quello. D\&D invece, a livello di cuore meccanico, è
giocato in modi diversissimi da dieci gruppi diversi.

Daniele fa un esempio: nei manuali di D\&D (dalla terza edizione e 3.5 alla
quinta) ci sono le regole per il \textbf{danno da caduta}. In \emph{Blades in
the Dark} non servono: c'è un altro metodo, ``non è così generalista, è
focalizzato'', e la narrazione si risolve in altro modo. Anche nell'OSR:
\emph{Shadowdark} ha il danno da caduta, perché magari è importante;
\emph{Neverland}, una reinterpretazione di \textbf{Peter Pan} in chiave OSR, no,
perché lì conta di più che il manuale offra tutti i personaggi del Peter Pan
tradizionale in una versione \emph{grimdark}. ``Anche l'OSR diventa un po'
focalizzato quando ci prova''.

\subsection{Not the End}

\slidefig{img/AGD/agd2-nottheend}{Il flusso di Not the End: il sacchetto di token come motore della risoluzione.}{slide \emph{Analog -- Case Studies}}{fig:agd2-nottheend}


Nicola, per cui questo è ``il settore che più adoro del gioco analogico'',
chiede dove si collochi \textbf{\emph{Not the End}}, un gioco italiano. La
risposta è nella slide successiva. \emph{Not the End} continua nella direzione
più narrativa e focalizzata (``semplicistica è un termine non corretto''): ha
procedure, meccaniche, un gameplay focalizzato a fare quello che il gioco
chiede. Si possono creare i propri setting, ma il gioco resta quello: il
personaggio principale è un \textbf{eroe}, ed è il giocatore a decidere quando
sarà la sua fine, quando giocherà il suo momento di eroismo stoico (da cui il
titolo). È un gioco \textbf{senza dadi}: c'è un conflitto, si fa e si descrive
un'azione, e la \textbf{risoluzione} consiste nell'\textbf{estrarre dei token}
bianchi e neri (o arancioni e viola, per le persone con difficoltà nella visione
dei colori); l'effetto viene descritto secondo le procedure del gioco. Il
panorama, insomma, si è diviso: c'è chi cerca un'esperienza più simile
all'originale, e chi un'esperienza più focalizzata sulla narrazione, magari
senza un \emph{world building} preciso.

\subsection{Che cos'è l'OSR, e com'è nato davvero il gioco di ruolo}

Un'altra studentessa chiede che cosa significhi OSR: \textbf{Old School
Renaissance}, ``rinascimento della vecchia scuola''\footnote{Poco prima Daniele
aveva sciolto la sigla come \emph{Old School Revival}; entrambe le forme sono in
uso.}. È un movimento ``molto, molto particolare'': ogni persona che ne fa parte
ha il proprio principio di OSR, gioca un OSR diverso dagli altri, e gli
appassionati ``adorano scontrarsi'' su che cosa intendono per OSR. Per questo
nella lezione sulla fiction Daniele parlerà di \textbf{culture} del GDR, più che
di community.

La stessa studentessa (iscritta a un altro corso dell'accademia, di cui sta
recuperando le lezioni del primo anno) riporta quello che ha capito dell'origine
di D\&D: i due autori avrebbero avuto visioni opposte, uno più libera e meno
regolamentata, partendo dai giochi di guerra, l'altro, che ha portato il gioco a
livello commerciale, più strutturata, dettagliata e pesante. ``È giusto, è un
riassunto corretto, ma è un po' più complicato di così'', e lo si vedrà nella
lezione sulla fiction. Ruggero aggiunge che D\&D ``non è l'unica origine del
gioco di ruolo'': il gioco di ruolo diventa \emph{famoso} con D\&D, ma gli stessi
autori si ispirarono alle \textbf{simulazioni militari}, molto diffuse
nell'America di quegli anni, e al \textbf{reenactment storico} (``ci sono due
episodi della \emph{Signora in giallo} sul reenactment della guerra civile, un
argomento allucinante''). Esistono esempi storici, dell'Ottocento e perfino del
Seicento, di \textbf{giochi di ruolo} veri e propri, con quest, avventure e
personaggi. Il \emph{mainstream} ha normalizzato l'idea che D\&D sia il primo
gioco di ruolo, ``ma non è così'': nel manuale originale l'espressione
\emph{role-playing game} non compare da nessuna parte, anche perché Gygax
odiava il termine. ``Un vaso di Pandora che non apriamo''.

Daniele però raccoglie l'osservazione. Nei primi anni editoriali di D\&D il
problema maggiore è stato il \textbf{simulazionismo}: interpretare
``simulazione'' come la necessità di avere una regola per ogni cosa che esiste nel
realismo del gioco. Ne è venuto un gioco molto strutturato, in cui ``ogni cosa
deve avere una regola'', perché ciò che non è coperto dalle regole il giocatore
può usarlo come \emph{loophole}. E ne è venuta anche l'idea che il master sia il
\textbf{padrone del tavolo} e che il giocatore debba in qualche modo
\textbf{vincere contro il master}: un altro ``vaso di Pandora enorme''.

\paragraph{Non esiste solo D\&D.} È per questo che Daniele non si concentra solo
su D\&D: ``l'idea con cui dovete uscire da qui è che non esiste solo D\&D''.
Anzi, oggi D\&D è \textbf{in controtendenza}, non per motivi di design ma per
\textbf{motivi editoriali}: il problema principale di D\&D come proprietà
intellettuale ``è sempre stato editoriale'', e di recente la Wizards of the Coast
ha fatto un po' di confusione con il \emph{branding}.

% ---------------------------------------------------------------------
\section{Story game e giochi di ruolo sperimentali}
\label{a2:sec:gdr-story}

\subsection{Gli story game}

\slidefig{img/AGD/agd2-storygames}{Alcuni story game.}{slide \emph{Analog -- Case Studies}}{fig:agd2-storygames}


La slide successiva è sugli \textbf{story game}, che in realtà nascono anche
prima di The Forge. ``Story game è un termine che io odio, perché vuol dire tutto
e non vuol dire niente'': sono giochi di ruolo da tavolo che si focalizzano
\textbf{ancora di più su una storia precisa}.
\begin{itemize}
  \item \textbf{\emph{Alice is Missing}}: la storia di Alice, scomparsa, e dei
        suoi amici che cercano di capire che cosa le è successo. Si gioca con
        una \textbf{chat} sul telefono, con un video a schermo, e con
        \textbf{due ore di tempo}.
  \item \textbf{\emph{Microscope}}: ``letteralmente \emph{world building}
        narrativo''.
  \item \textbf{\emph{Bleak Spirit}}: giocare un \emph{Dark Souls} in cui tutti
        controllano a turno il viandante protagonista.
  \item \textbf{\emph{The Zone}}\footnote{La trascrizione attribuisce il gioco a
        ``Ruff D'Amico''; il nome non è verificabile.}, giocabile anche con un
        supporto online: si narra, di base, \emph{Annientamento} (``se avete
        visto il film''). Si entra in una zona in cui una presenza aliena sta
        cambiando il mondo, e \textbf{moriranno tutti tranne una persona}, che
        non si sa in anticipo chi sarà. È un'influenza del LARP: come in
        \emph{Demeter}, si sa che il proprio personaggio morirà.
  \item \textbf{\emph{Archipelago}} e \textbf{\emph{For the Queen}}, e altri
        giochi che \textbf{tolgono i dadi} e mettono le carte, o tolgono
        perfino la \textbf{scheda del personaggio}: in \emph{Microscope} ci sono
        solo dei post-it.
\end{itemize}

\subsection{The Quiet Year e i map drawing game}

\slidefig{img/AGD/agd2-quietyear}{Il flusso di The Quiet Year.}{slide \emph{Analog -- Case Studies}}{fig:agd2-quietyear}


``C'è un minimo di suspense'': si arriva a esperienze ancora più lontane.
In \textbf{\emph{The Quiet Year}} non si interpreta nemmeno un personaggio: si è
parte di una \textbf{comunità post-apocalittica} e si gioca \textbf{un anno} della
sua vita. Non si può discutere di ciò che si gioca. A ogni turno si gira una
\textbf{carta} che dice che cosa succede, e durante il proprio turno si può fare
una tra due o tre azioni (disegnare, iniziare un progetto\ldots). Tutto si gioca
su una \textbf{mappa} su cui si disegna ciò che si vuole intendere: c'è chi gioca
in modo più rigoroso, senza poter nemmeno scrivere parole sulla mappa, e chi lo
permette. È ``un gioco fenomenale'', perché fa capire che il gioco di ruolo può
essere anche altro; e ha normalizzato un genere, quello dei \textbf{map drawing
role-playing game}, i giochi di ruolo in cui si disegnano mappe.

\subsection{Dialect}

\slidefig{img/AGD/agd2-dialect}{Il flusso di Dialect.}{slide \emph{Analog -- Case Studies}}{fig:agd2-dialect}


E poi \textbf{\emph{Dialect}}. Qui si ha un personaggio, ma è solo un nome e un
ruolo, un archetipo. Si gioca \textbf{la morte di una comunità attraverso la
morte di una lingua}: una comunità umana isolata dalle altre, che ha creato il
proprio linguaggio, e che giocherà la fine di quell'isolamento attraverso la
morte del linguaggio. A ogni turno si gioca una carta che è un \textbf{concetto
della lingua} (per esempio: come si dice ``amico'' nella nostra lingua?); si
stabilisce la parola e si gioca una breve scena in personaggio, senza tiri di
dado né altro; quando la parola viene usata in personaggio, chiunque al tavolo
può dire stop: ``abbiamo usato la parola, va bene così'', e si va avanti.

Non è un gioco ``non complesso'', come diceva Ruggero: \emph{Dialect} è
``chiarissimo'' da leggere, ma la complessità sta nel \textbf{design}, un design
elegante che crea una \emph{one shot} con ``una profondità'' enorme. Daniele con
\emph{Dialect} ha pianto; e quando incontra gli amici con cui lo ha giocato anni
fa, si parlano ancora nella lingua inventata allora: ``ehi tu, oceano'', e chi
c'era sa che cosa vuol dire. Anche qui si vede come le esperienze nate dal Nordic
LARP siano \textbf{``pervase''} (``rubo una parola a Beatrice'') in altri ambiti.

\subsection{Servono davvero le categorie? Il problema editoriale}

Ci sono ancora esperienze più tradizionali; ``il mondo è molto variegato'', e si
parlerà di \textbf{culture}: spesso chi si avvicina al gioco di ruolo vuole
giocare una certa cosa, una cultura di gioco con titoli che ``vanno bene'' e
titoli che no. L'ultima slide è ``una piccola critica'': è davvero necessario
mettere le categorie ai giochi di ruolo? Da un lato sì, dall'altro no, perché i
problemi principali sono due, e il primo è \textbf{il mondo editoriale}: i giochi
di ruolo sono stampati come \textbf{libri}, e l'editore ha ``una vita travagliata''
con il gioco. Esempi:
\begin{itemize}
  \item \textbf{\emph{In the Name of God}}\footnote{Titolo e autore (``Alessandro
        Pirodi'') come nella trascrizione; non verificabili.}, stampato solo in
        \textbf{self-publishing} e assente dal mercato, ``ma bellissimo''; la sua
        meccanica di risoluzione chiede di fare qualcosa fisicamente al tavolo
        (la trascrizione dice ``fare home'');
  \item un'\textbf{espansione per qualsiasi gioco di ruolo}\footnote{La
        trascrizione dice ``Mew Wild in the Subway''.};
  \item un \textbf{GDR stampato su una maglietta}, di due autori che Daniele
        saluta;
  \item un gioco di \textbf{cultisti contro investigatori} che prende da
        esperienze di chamber LARP\footnote{La trascrizione lo chiama ``SOT''.}.
\end{itemize}

\begin{figure}[htbp]
\centering
\begin{tikzpicture}[
    font=\footnotesize,
    epoca/.style={draw, thick, rounded corners=3pt, fill=gray!12,
                  text width=2.55cm, align=center, minimum height=1.0cm},
    nota/.style={align=left, text width=2.7cm, font=\scriptsize},
    freccia/.style={-{Stealth[length=2.5mm]}, thick}
  ]
  \draw[freccia] (-0.2,0) -- (15.4,0);
  \foreach \x/\a in {0.9/{anni '50--'70}, 4.0/{'74--anni '90},
                     7.1/{fine '90--2000}, 10.2/{anni 2010}, 13.3/{post-COVID}}{
    \draw[thick] (\x,0.12) -- (\x,-0.12);
    \node[below=0.15cm, font=\scriptsize\itshape] at (\x,0) {\a};
  }
  \node[epoca] (e1) at (0.9,1.2)  {\textbf{Wargame e miniature}};
  \node[epoca] (e2) at (4.0,1.2)  {\textbf{D\&D e GDR ``early''}};
  \node[epoca] (e3) at (7.1,1.2)  {\textbf{New Wave / The Forge}};
  \node[epoca] (e4) at (10.2,1.2) {\textbf{Giochi focalizzati, story game}};
  \node[epoca] (e5) at (13.3,1.2) {\textbf{OSR e ritorno all'origine}};
  \node[nota, anchor=north] at (0.9,-0.65) {simulazioni militari,
    reenactment; \emph{Chainmail} (1971)};
  \node[nota, anchor=north] at (4.0,-0.65) {Gygax e Arneson;
    master arbitro; tutti i generi (\emph{Cthulhu}, \emph{Vampire},
    \emph{Traveller}, \emph{GURPS}); IP adattate};
  \node[nota, anchor=north] at (7.1,-0.65) {il sistema conta;
    \emph{Apocalypse World}, \emph{Nobilis}, \emph{Dread}, \emph{Fate}; Big
    Model; \emph{Fiasco}};
  \node[nota, anchor=north] at (10.2,-0.65) {\emph{Blades in the
    Dark}, \emph{Fabula Ultima}, \emph{Not the End}; \emph{The Quiet Year},
    \emph{Dialect}; due giocatori, solo, micro};
  \node[nota, anchor=north] at (13.3,-0.65) {meno regole, alta
    mortalità, sandbox; \emph{Mörk Borg}, \emph{Shadowdark}};
\end{tikzpicture}
\caption{Una ricostruzione schematica della ``piccola storiografia'' dei giochi
di ruolo da tavolo raccontata da Daniele. Le epoche sono indicative e si
sovrappongono (gli story game, per esempio, nascono anche prima di The Forge);
la storia completa è rimandata alla lezione sulla fiction.}
\label{a2:fig:storia-gdr}
\end{figure}

% ---------------------------------------------------------------------
\section{Giochi di ruolo per due, in solitario e micro}
\label{a2:sec:gdr-formati}

\slidefig{img/AGD/agd2-micro}{Giochi di ruolo per due, in solitario e nano-RPG.}{slide \emph{Analog -- Case Studies}}{fig:agd2-micro}


L'ultima parte sui giochi di ruolo torna a Ruggero, che presenta tre formati.

\begin{description}
  \item[Two player only.] Storicamente i giochi di ruolo nascono in gruppi di
        gioco ampi, ma c'è stata ``una grossissima wave'' di giochi di ruolo
        \textbf{per due persone}, e non nel senso che una fa il master e l'altra
        un personaggio: spesso l'\textbf{agency narrativa}, il potere di decidere
        che cosa succede, viene \textbf{passata meccanicamente} al tavolo o
        divisa in responsabilità diverse. Ce ne sono tantissimi, e per Ruggero
        sono molto interessanti anche perché, come si vedrà nella lezione
        sull'emotional design, permettono di fare cose molto \textbf{intime}: non
        serve un gruppo di gioco, si possono giocare con il partner o con
        un'altra persona.
  \item[Solo RPG.] Dal 2013 circa si sono normalizzati i \textbf{giochi di
        ruolo in solitario}: si gioca da soli, senza un'altra persona e senza
        fare da master a sé stessi. Spesso non c'è un master: ci sono
        \textbf{procedure} che guidano il gioco, e il manuale è il sistema di
        gioco, strutturato ``quasi come un board game'' in termini di regole. Ma
        resta tutto ciò che è tipico del gioco di ruolo: l'agency emergente,
        l'interpretazione di un personaggio, la creazione di una narrazione. Di
        solito si giocano in formato \textbf{diario} (\emph{journaling}) o in
        formato \textbf{audio}, e sono in grande tendenza.
  \item[Micro RPG.] Infine, i \textbf{micro giochi di ruolo}\footnote{La
        trascrizione dice ``handheld RPG''; dal contesto si tratta dei
        \emph{micro RPG}.}: giochi di ruolo su \textbf{biglietti da visita}, o
        raccolte in cui un libro ne contiene cento. Il \textbf{200 Word RPG}
        (\emph{200 Word RPG Challenge}) era un contest che chiedeva di scrivere un
        gioco di ruolo in \textbf{massimo 200 parole}, e ogni edizione ne
        raccoglieva migliaia: nel 2017, ricorda Ruggero, circa 900 giochi, su
        circa mille proposte.
\end{description}

% ---------------------------------------------------------------------
\section{I giochi da tavolo: i generi storici}
\label{a2:sec:tavolo}

\slidefig{img/AGD/agd2-rollmove}{Roll and move, astratti e party game.}{slide \emph{Analog -- Case Studies}}{fig:agd2-rollmove}


``Passerei ai giochi da tavolo, correndo velocemente e senza fare pause, se ce la
fate''. Come per i LARP e i giochi di ruolo, Ruggero parla di
\textbf{categorie dall'alto}, e per ognuna mostra come il genere si sia
\textbf{evoluto} nel tempo.

\subsection{Roll and move}

I \textbf{roll and move} si basano, come dice il nome, sul \textbf{tirare un
dado e muoversi}. Sono forse i più comuni nella storia dell'umanità, con quello
che da noi è il \textbf{gioco dell'oca}: un tracciato su cui si tirano i dadi e
si avanza. L'\textbf{agency} in questi giochi è ``praticamente inesistente'', e
il genere si è presto evoluto integrando altre meccaniche: con
\textbf{elementi economici} (come \emph{Monopoly}) o con \textbf{elementi
deduttivi}, come \emph{Cluedo}, che è un roll and move.

\subsection{Astratti}

Un altro genere storicamente gigantesco è quello degli \textbf{astratti}:
scacchi, dama, go, tutti quei giochi la cui componente tematica è ridotta
``quasi all'osso''. Gli scacchi mantengono ancora qualche forma di tema; il go, di
fatto, è un rettangolo di legno con delle caselle e delle pietre bianche e nere.
Anche qui il genere non si è fermato: ci sono iterazioni moderne come
\emph{ZÈRTZ} e i giochi del \textbf{progetto GIPF}\footnote{La trascrizione dice
``Zerz o il Gift Project''.}, e altri astratti contemporanei (uno è mostrato a
destra nella slide\footnote{La trascrizione non permette di ricostruirne il
titolo (``i borsa astratti, per esempio, di Supernova'').}), che mantengono
dinamiche ``molto asciutte, molto distillate'': \textbf{informazione aperta},
trasparente, niente carte imprevisti, niente mosse segrete. Diventano quasi una
\textbf{sfida d'intelletto} con l'avversario.

\subsection{Party game}

I \textbf{party game}: \emph{Taboo}, \emph{Lupus in Tabula} o \emph{Werewolf} a
seconda dell'edizione. Sono pensati per \textbf{molte persone} (di solito dalle
quattro in su), con regole ``molto asciutte, molto leggere'', perché si giocano
in contesti di festa, in cui non è detto che chi è nella stanza ami giocare o
abbia mai giocato. Poche regole, leggerezza, durata breve (spesso non più di
mezz'ora): sono \textbf{un possibile punto di accesso}. Anche questi si sono
evoluti, con giochi come \emph{Hitster}\footnote{La trascrizione dice ``Hit'';
la ricostruzione è probabile ma non certa.} e \emph{Wavelength}, che portano
sensibilità più moderne e meccaniche più interessanti del ``tirare alla cieca
per capire chi tra di noi è un lupo mannaro''.

\subsection{Giochi di parole}

I \textbf{word game}\footnote{La trascrizione dice ``war games'', ma gli esempi
(\emph{Boggle}, \emph{Scarabeo}) chiariscono che si tratta dei giochi di
parole.}: \emph{Boggle}, \emph{Scarabeo}, giochi in cui la meccanica centrale è
\textbf{creare parole}, e più la parola è lunga, complessa o usa lettere rare,
più punti si fanno. Anche qui le iterazioni sono interessanti: \emph{Cross
Clues}\footnote{La trascrizione dice ``Cross Cruise''.} trasforma il genere in
un gioco \textbf{cooperativo}, in cui non si compete per la parola più lunga ma si
gioca con le parole per capire ``quanta intesa c'è tra di noi''; e
\emph{Illiterati} è praticamente un gioco con le \emph{boss fight}, in cui si
sconfiggono i boss formando parole. Ibridazioni anche qui, perché chi crea giochi
li crea in base a ciò che ha giocato prima.

\subsection{Giochi di carte}

\slidefig{img/AGD/agd2-carte}{Giochi di parole, giochi di carte, TCG e deckbuilding.}{slide \emph{Analog -- Case Studies}}{fig:agd2-carte}


I \textbf{card game}, che in Occidente si diffondono storicamente con i mazzi
\textbf{francesi} e, in Italia, con quelli ``da briscola'' (ma ci sono anche altri
casi storici\footnote{La trascrizione cita i ``Mazda 1''; il riferimento non è
ricostruibile.}). Le evoluzioni:
\begin{itemize}
  \item i \textbf{micro card game} come \emph{Love Letter}, che in una decina o
        dozzina di carte racchiude un gioco complesso;
  \item giochi come \emph{Munchkin}, in una direzione più ampia, con il
        \textbf{take that} (``voglio mettere i bastoni tra le ruote'' agli altri);
  \item la ``gigantesca categoria'' dei \textbf{trading card game}, già
        menzionata nella lezione precedente: \emph{Magic: The Gathering},
        \emph{Yu-Gi-Oh!}, \emph{Pokémon}. Il \emph{trading} non è solo una
        dinamica tra persone: è nel nome della categoria. Scambiare carte
        \textbf{fuori dalla partita} è parte della dinamica comunitaria;
  \item i \textbf{Living Card Game} (LCG), che provano a risolvere uno dei
        principali problemi editoriali dei TCG: ``se io ho più soldi di te, tu
        contro di me a Magic non vinci'', perché comprerò carte da 800 o 900
        euro l'una. Gli LCG \textbf{non hanno booster casuali}: ogni espansione ha
        l'elenco preciso del contenuto. Restano giochi di carte con il
        \textbf{deck building pre-partita}, con mazzi ``fortemente identitari''
        (si possono scegliere le carte anche solo perché sono più simpatiche),
        ma senza l'aspetto economico della \textbf{rarità indotta};
  \item i \textbf{deck-building game}, in cui la costruzione del mazzo
        \textbf{non è più pre-partita ma avviene durante la partita}: spesso si
        parte con mazzi identici, e durante il gioco si acquisiscono carte che
        cambiano \textbf{permanentemente} il proprio mazzo, con lo scopo di
        costruire il mazzo migliore possibile.
\end{itemize}

\subsection{German game ed Eurogame}

\slidefig{img/AGD/agd2-euro}{German game, Eurogame e American.}{slide \emph{Analog -- Case Studies}}{fig:agd2-euro}


Altre categorie, che Ruggero mette tra virgolette, sono i \textbf{German game}:
un genere sviluppato in \textbf{Germania} dopo la metà degli anni '50--'60, con
giochi a \textbf{interazione indiretta}, competitivi ma ``senza troppo dar
fastidio direttamente all'altra persona'': non ci si può attaccare
direttamente. Sono giochi poco bellici, anche per \textbf{rigetto della Seconda
guerra mondiale}, e si è sviluppato moltissimo il filone \textbf{comunitario}, di
\emph{community building} e di scambio: costruire la propria fattoria, costruire
un villaggio, scambiarsi le risorse (``io ho delle risorse, tu altre,
scambiamocele''), senza distruggere ciò che l'altro ha costruito.

Il genere si è poi reso più complesso e si è espanso oltre la Germania con gli
\textbf{Eurogame}, che includono Spagna, Portogallo, Francia e anche l'Italia
(Ruggero ne ha fatti parecchi). L'interazione resta indiretta, ma i giochi
diventano molto più complessi: spesso durano \textbf{due o tre ore} e provano a
rendere \textbf{una simulazione focalizzata di qualcosa del mondo reale}. Nella
slide, a sinistra, un gioco in cui si gestisce una \textbf{galleria
d'arte}\footnote{La trascrizione dice ``The Gallery''; probabilmente \emph{The
Gallerist}.}, con tutte le dinamiche di compravendita di quadri; e
\emph{Endeavor}, che Ruggero ha alle spalle, sul commercio mondiale in epoca
moderna. Dinamiche che simulano un tipo di mercanteggiare: ``a volte compro e
vendo quadri, a volte compro e vendo vasi''.

\paragraph{Approfondimento: l'Eurogame secondo Brathwaite e Schreiber.} In
\emph{Challenges for Game Designers}, nel capitolo sul \emph{targeting} di un
mercato, gli Eurogame (``noti anche come \emph{German games} o \emph{designer
games}'') sono un caso studio di gioco per il mercato di massa: giochi per la
famiglia, godibili a tutte le età; \emph{I coloni di Catan} è il primo ad aver
reso popolare il genere in Nord America, seguito da \emph{Carcassonne} e
\emph{Puerto Rico}. Il libro ne dà una definizione operativa, utile per
confrontare le osservazioni di Ruggero:
\begin{itemize}
  \item \textbf{durata tra 20 e 90 minuti} (troppo brevi annoiano gli adulti,
        troppo lunghi i bambini);
  \item \textbf{setup breve}, perché il setup non contiene decisioni;
  \item \textbf{regole semplici}, curva di apprendimento breve;
  \item \textbf{cooperazione più che scontro}: scambi e negoziazioni invece di
        meccaniche distruttive in cui un giocatore danneggia direttamente un
        altro (il che non impedisce una competizione intensa);
  \item \textbf{decisioni strategiche con casualità misurata};
  \item \textbf{tempi morti minimi}: turni brevi, o azioni anche nei turni
        altrui.
\end{itemize}
La ``cooperazione invece dello scontro'' è l'interazione indiretta del German
game; la durata, invece, mostra bene l'evoluzione che racconta Ruggero: gli
Eurogame più complessi, da due o tre ore, sono ormai un'altra cosa rispetto al
gioco per famiglie descritto dal libro.

\subsection{American}

In contrapposizione, oltreoceano c'erano gli \textbf{American} (o
\emph{Ameritrash}): giochi molto focalizzati sulla \textbf{produzione} (le
miniature di plastica), con il \textbf{dungeon crawling}, tematicamente molto
vicini a \emph{Dungeons \& Dragons}, quasi versioni da tavolo del combattimento a
griglia normalizzato da D\&D. Continuano la tradizione di \emph{Chainmail} e dei
giochi strategici e tattici: sono spesso giochi di combattimento, a
\textbf{interazione diretta} (si attacca e ci si difende), e buona parte dei loro
regolamenti è dedicata alla simulazione del combattimento.

\subsection{Strategici e wargame}

\slidefig{img/AGD/agd2-strategici}{Strategici, movimento nascosto, destrezza, 18XX, tempo reale e legacy.}{slide \emph{Analog -- Case Studies}}{fig:agd2-strategici}


Un altro genere gigantesco è quello dei \textbf{giochi di strategia}, usati
storicamente anche nella \textbf{formazione militare}, ``non a caso'': la
\textbf{RAND Corporation} ne ha creati diversi, e ci sono editori specializzati.
Lo scopo era ricreare battaglie storiche. Anche qui, evoluzioni:
\begin{itemize}
  \item i giochi \textbf{a zone} o ad \textbf{aree} (\emph{area control}) con
        unità su una mappa, di cui \emph{Risiko} è il più famoso: io ho un
        esercito, tu ne hai un altro, dobbiamo conquistare territori, e il
        conflitto diretto è inevitabile perché gli obiettivi si sovrappongono;
  \item i \textbf{card-driven strategy game}: strategici con mappe e zone di
        controllo, ma con \textbf{mazzi di eventi}. I designer prendono eventi
        reali di conflitti reali e cercano di renderli coesi tematicamente con
        le meccaniche, dando a ogni evento una risoluzione precisa: ``l'importante
        non è tanto quante unità ho, ma quali carte evento ho in mano''. Sono,
        se ci si pensa, più vicini ai card game che agli strategici.
\end{itemize}
Nicola chiede di \textbf{\emph{Campaign for North Africa}}\footnote{La
trascrizione dice ``Campi for Nord Africa''.}: è uno \textbf{strategico
classico}, con esagoni e pedine di cartoncino, ``un'esperienza devastante in
termini di lunghezza'': forse non è il caso di cominciare da lì.

\subsection{Movimento nascosto}

Il genere degli \textbf{hidden movement}\footnote{La trascrizione dice ``Eat the
Movement''.}: \emph{Scotland Yard}, \emph{Letters from Whitechapel}, ma anche
altri titoli più recenti\footnote{La trascrizione cita ``Sniper Elite, Mind Dame
e GMT''; il secondo è probabilmente \emph{Mind MGMT}, GMT è un editore di
wargame.}. Sono giochi \textbf{asimmetrici}: di solito una persona si nasconde e
scappa, e le altre vincono collettivamente se la catturano. Uno contro tutti.

\subsection{Destrezza}

I \textbf{dexterity game} richiedono \textbf{destrezza fisica}: colpire
componenti con le schicchere, impilarli senza farli cadere. Non basta la
strategia che si ha in testa: ``posso anche avere la strategia migliore, ma se poi
le mie mani non ce la fanno a mettere in piedi questi cilindretti, crolla
tutto''.

\subsection{18XX}

Un genere che Ruggero ama particolarmente, e di cui è ``abbastanza sicuro''
che nessuno abbia mai sentito parlare, anche se conta centinaia di titoli (``più
di 400 giochi'' sotto l'etichetta): i \textbf{18XX}. Il nome viene dal secolo
delle \textbf{ferrovie}, l'Ottocento: sono tutti giochi ferroviari, ognuno
centrato su un anno e un luogo specifici (\emph{1849}, le ferrovie della
Sicilia nel 1849; \emph{1889}, quelle dell'isola di Shikoku). Hanno una forte
componente di \textbf{costruzione delle ferrovie} sulla mappa e una componente
\textbf{economico-azionaria}: comprare e vendere azioni, speculazione
finanziaria. Sono ``molto belli, ma anche molto intensi'': \emph{1889}, che
Ruggero usa di solito come introduttivo, ``dura tre ore''.

\subsection{Tempo reale}

Culturalmente si pensa ai giochi da tavolo come a qualcosa che avviene
\textbf{a turni}, ma non è sempre così. I giochi \textbf{in tempo reale} usano
clessidre, timer o altre meccaniche per cui non si ha tutto il tempo per pensare:
bisogna reagire in fretta. \emph{Kitchen Rush}, per esempio, è ``praticamente
\emph{Overcooked} da tavolo'': si gestisce un ristorante, e \textbf{la clessidra
è il lavoratore}, che non si può spostare finché non si è esaurita. Chi prende
gli ingredienti, chi li prepara sul fornello, chi va in sala a prendere gli
ordini, chi porta i piatti, chi li lava: ne nasce un certo caos nell'intreccio.

\subsection{Legacy}

I \textbf{Legacy} sono ``un po' una rivoluzione nel genere''. Di solito un gioco
da tavolo resta sempre quello: lo si prepara, lo si gioca, la settimana dopo lo
si gioca con un altro gruppo. I giochi Legacy invece chiedono di avere un
\textbf{impatto fisico sul gioco}, per cui \textbf{le partite successive sono
influenzate dalle precedenti}:
\begin{itemize}
  \item \emph{Pandemic Legacy} chiede di \textbf{strappare carte}; alcune città
        vengono distrutte dalla pandemia, coperte con degli \textbf{adesivi}, e
        non esistono più sulla mappa;
  \item \emph{SeaFall}\footnote{La trascrizione dice ``Sifol''.}, sulla scoperta
        di isole in un Ottocento alternativo: le isole hanno il nome vuoto, e chi
        le scopre le battezza; da lì in poi l'isola porterà quel nome per tutto
        il resto del gioco.
\end{itemize}
``So che qualcuno sta già sudando freddo all'idea di strappare carte o scrivere
sul tabellone'', ma sono ``esperienze emotivamente incredibili'': ``giocatevi dei
Legacy''.

\paragraph{Approfondimento: Risk Legacy e la rottura di un'assunzione.} Schell
dedica ai Legacy due passaggi di \emph{The Art of Game Design}. Nel capitolo sul
\textbf{tema}, presenta \emph{Risk Legacy} di \textbf{Rob Daviau} come esempio di
tema fondato su una verità: Daviau ha fatto ``qualcosa di mai tentato prima nel
mondo dei giochi da tavolo'', creando meccaniche per cui le scelte dei giocatori
\textbf{cambiano il gioco per sempre}. Le regole impongono di modificare il
tabellone con adesivi non rimovibili, di scriverci sopra con un pennarello
indelebile per rivendicare i territori, di strappare e scartare carte e perfino
di modificare permanentemente il regolamento. Queste meccaniche, nota Schell,
sono nuove e intriganti, ma soprattutto \textbf{rafforzano il tema centrale}:
``la guerra cambia il mondo''. Nel capitolo sulle tecniche di brainstorming,
Schell racconta come Daviau ci sia arrivato: elencando le \textbf{assunzioni}
date per scontate sui giochi da tavolo e provando a ``farne saltare'' una,
``una partita non influenza la successiva''. È un buon esempio di quello che
Ruggero chiedeva all'inizio: guardare un genere che si conosce e chiedersi che
cosa succederebbe se una delle sue regole implicite non valesse più.

% ---------------------------------------------------------------------
\section{Due libri sul gioco da tavolo}
\label{a2:sec:tavolo-libri}

\slidefig{img/AGD/agd2-libri}{I due libri consigliati sul gioco da tavolo.}{slide \emph{Analog -- Case Studies}}{fig:agd2-libri}


A proposito di risorse, Ruggero torna sul libro di Nuccio, \emph{La
progettazione dei giochi da tavolo}, citato a inizio lezione, e indica
un'alternativa in inglese che a suo parere ``fa quello che fa Nuccio, meglio'':
\textbf{\emph{Building Blocks of Tabletop Game Design}} di \textbf{Geoffrey
Engelstein e Isaac Shalev}\footnote{La trascrizione dice ``Engelstein e
Schaler'' e ``Building Blocks of Tablet of Game Design''.}. È più
\textbf{contemporaneo} negli esempi e ``più funzionale''. Per i giochi da tavolo,
se si ha già un framework di riferimento, sono entrambi ``una risorsa
incredibile''. (Resta valido l'avvertimento di inizio lezione: sono archivi di
meccaniche, da usare dopo aver capito che cosa le meccaniche fanno.)

% ---------------------------------------------------------------------
\section{Libri-gioco, megagame e altri ibridi}
\label{a2:sec:tavolo-ibridi}

\slidefig{img/AGD/agd2-ibridi}{Libri-gioco, megagame e ibridi tra gioco da tavolo, gioco di ruolo e videogioco.}{slide \emph{Analog -- Case Studies}}{fig:agd2-ibridi}


``Manca ancora qualcosa?'' Mancano diverse cose.
\begin{itemize}
  \item I \textbf{libri-gioco}, nel formato libro tradizionale e nel formato
        ``più contemporaneo'' dei giochi da tavolo che sono libri-gioco, come
        \emph{Sherlock Holmes Consulente Investigativo}.
  \item I \textbf{megagame}: giochi da tavolo \textbf{per cento persone}. Non
        hanno un'organizzazione editoriale come \emph{Monopoly} o gli altri
        board game: sono organizzati in modo più vicino ai giochi pervasivi e al
        LARP, spesso \textbf{a evento}, spesso con biglietto, anche perché
        gestire un gioco da tavolo con cento persone in parallelo ``può diventare
        molto intenso''.
  \item I \textbf{LARP urban game}, di cui ha parlato Beatrice, con esempi
        storici anche italiani, come \emph{Pathos}\footnote{Grafia come nella
        trascrizione.}: giochi di ruolo dal vivo che sono anche urban game, e
        un po' ARG.
  \item Produzioni ispirate al freeform ma vendute come giochi da tavolo, come
        \textbf{\emph{Inhuman Conditions}}\footnote{La trascrizione dice ``In
        Human Condition''.}: praticamente il \textbf{test di \emph{Blade Runner}}
        per scoprire chi è un androide, per due persone. È un gioco di ruolo, ma
        prodotto esattamente come un gioco da tavolo (carte, scatola,
        regolamento) e non si presenta come gioco di ruolo, ``ma questo non lo
        rende meno gioco di ruolo''.
  \item Ibridi con il party game, come \textbf{\emph{Sì, Oscuro Signore!}}: un
        party game, ma anche un \emph{role-playing game}, perché si fa il role
        play del proprio personaggio.
  \item Ibridi con il videogioco, come \emph{Fateforge} e gli altri giochi che
        usano un'\textbf{app} non solo come segnapunti ma come meccanica durante
        la partita.
\end{itemize}

\begin{table}[htbp]
\centering
\small
\begin{tabular}{@{}p{3.0cm}p{9.9cm}@{}}
\toprule
\textbf{Genere} & \textbf{Che cosa lo caratterizza ed evoluzioni citate} \\
\midrule
Roll and move & tira e muovi (gioco dell'oca); agency quasi nulla; poi
elementi economici (\emph{Monopoly}) o deduttivi (\emph{Cluedo}) \\
Astratti & tema ridotto all'osso, informazione aperta (scacchi, dama, go,
\emph{ZÈRTZ}, progetto GIPF) \\
Party game & molte persone, regole leggere, partite brevi, punto di accesso
(\emph{Taboo}, \emph{Lupus in Tabula}, \emph{Wavelength}) \\
Giochi di parole & creare parole per fare punti (\emph{Boggle},
\emph{Scarabeo}); cooperativi (\emph{Cross Clues}), boss fight
(\emph{Illiterati}) \\
Giochi di carte & mazzi tradizionali; \emph{Love Letter}; take that
(\emph{Munchkin}); TCG, LCG e deck-building \\
German game ed Eurogame & interazione indiretta, costruzione e scambio;
simulazioni focalizzate di due o tre ore \\
American & produzione in plastica, dungeon crawling, combattimento e
interazione diretta \\
Strategici & formazione militare; aree e zone (\emph{Risiko}); card-driven;
esagoni (\emph{Campaign for North Africa}) \\
Movimento nascosto & uno contro tutti, asimmetrici (\emph{Scotland Yard},
\emph{Letters from Whitechapel}) \\
Destrezza & schicchere, impilare: conta anche la mano \\
18XX & ferrovie e azioni dell'Ottocento; più di 400 titoli (\emph{1849},
\emph{1889}) \\
Tempo reale & clessidre e timer (\emph{Kitchen Rush}) \\
Legacy & le partite cambiano fisicamente il gioco (\emph{Pandemic Legacy},
\emph{SeaFall}) \\
Libri-gioco, megagame, ibridi & \emph{Sherlock Holmes Consulente
Investigativo}; cento giocatori a evento; \emph{Inhuman Conditions},
\emph{Sì, Oscuro Signore!}, giochi con app \\
\bottomrule
\end{tabular}
\caption{I generi di gioco da tavolo presentati da Ruggero, in sintesi.}
\label{a2:tab:tavolo}
\end{table}

% ---------------------------------------------------------------------
\section{La chiusura: niente interrogazioni, e i materiali}
\label{a2:sec:chiusura}

``E ho corso un po', ma ce l'abbiamo fatta''. Mancano gli ultimi dieci minuti, e
Ruggero si complimenta con chi è ``arrivato fino a qui''. Ripete il punto
fondamentale: \textbf{non ci sarà nessuna interrogazione} su questi giochi
(``Tiziana, ti do subito il sollievo''). Non si chiede di sapere che cosa sia ogni
caso studio: lo scopo era dare una \textbf{panoramica}, degli spunti, e far
sapere che ``questa roba esiste''. Le slide restano, e ci si potrà tornare in
fase di ricerca e sviluppo. Con la lezione di \textbf{ricerca e sviluppo}, che
darà gli strumenti per farlo in autonomia, ognuno troverà \textbf{i casi studio
che servono al proprio gioco}, ``e io ve li chiederò, noi ve li chiederemo''.
Questa lezione serviva ad \textbf{aprire le possibili direzioni}.

\paragraph{I materiali caricati.} Mentre Ruggero parla, le slide vengono
caricate nella cartella degli studenti. Beatrice aggiunge di aver caricato:
\begin{itemize}
  \item il paragrafo di \textbf{Frank Rose} su \emph{Why So Serious};
  \item un estratto sulla \textbf{definizione dei pervasive game} tratto da
        \emph{Game Design}, un libro di \textbf{Maresa Bertolo e Ilaria
        Mariani}\footnote{La trascrizione dice ``Marianne Bertolo'' e
        ``Bertolo Mariani''. Il libro è con ogni probabilità \emph{Game
        Design. Gioco e giocare tra teoria e progetto} di Maresa Bertolo e Ilaria
        Mariani.}, ``un ottimo libro''.
\end{itemize}
Non sono letture obbligatorie.

\paragraph{Un avvertimento sui libri italiani di gioco di ruolo.} Daniele
aggiunge una cautela sulla parte dedicata al gioco di ruolo del libro di Bertolo
e Mariani: alcuni paragrafi ``generalizzano molto verso il mainstream''. E poi un
\emph{warning} più generale: \textbf{la maggior parte della letteratura italiana
sul gioco di ruolo è ``terribile''}, al punto che alcuni libri i docenti non solo
non li consigliano, ma li \textbf{sconsigliano attivamente}. Con la premessa che
non è nulla di personale contro l'autore, Daniele sconsiglia il manuale
\emph{Giochi di ruolo} di \textbf{Mauro Longo}\footnote{La trascrizione riporta
l'editore come ``Dino Dino Editori'': probabilmente Dino Audino Editore.},
perché è molto soggettivo, molto legato alla visione del gioco di ruolo
dell'autore e a un solo tipo di gioco, quello di \emph{Lex Arcana}\footnote{La
trascrizione dice ``l'ex arcana''.}, con teorie proprie che non corrispondono al
design generale dei giochi di ruolo. ``Va bene se volete fare \emph{Lex Arcana}
esattamente come \emph{Lex Arcana}'', non se si vuole fare altro.

Meglio il volume sull'edu-larp, e soprattutto \textbf{giocare}: ``giocate i giochi,
perché in Italia abbiamo una creatività sfrenata''. Si trovano ``robe favolose''
online: tantissimi \textbf{one page game} di ruolo, il mondo del \textbf{print and
play} anche nei giochi da tavolo, chamber LARP e freeform da stampare
direttamente da internet. ``Fatevi una ricerca e vedete quello che vi garba''.

% ---------------------------------------------------------------------
\section{Le domande finali}
\label{a2:sec:domande}

``Siamo stati più veloci dell'anno scorso'' (quando si erano accumulati trenta
minuti di ritardo), e restano sette minuti per le domande.

\subsection{Che cosa fa scegliere una forma?}

Gianmarco, ``un po' per rompere il ghiaccio'', fa una domanda ``molto generica'',
che sa non avere una risposta semplice: si sono visti giochi dalle forme
diversissime, ma i \textbf{framework} per crearli sono più o meno gli stessi. Che
cosa determina la scelta di una forma rispetto a un'altra? Ruggero rimanda la
risposta alla lezione sull'\textbf{emotional design}, ``che è proprio una risposta
a questa cosa'': la risposta è così lunga che ``abbiamo una lezione di tre ore per
rispondere alla tua domanda''.

\subsection{Che cosa va per la maggiore in Italia?}

Angelo chiede quale genere ``vada per la maggiore'' in Italia. Ma che cosa
significa ``andare per la maggiore''? Il successo commerciale non è uguale per
tutti i generi, ``nemmeno nei videogiochi''. Uno dei docenti (con ogni probabilità Daniele) risponde con due domande
``provocatorie'', che valgono anche per i libri (ne parlava quel giorno con
un'amica narrative designer di videogiochi):
\begin{enumerate}
  \item quante copie deve vendere un gioco per essere un successo?
  \item il numero di copie vendute coincide con il successo
        commerciale ed economico?
\end{enumerate}
La risposta a entrambe è \textbf{no}, ma si capirà con la lezione di
\textbf{Flavio} sull'aspetto produttivo: quando si ragiona in logica commerciale
si è ``in logica di prodotto'', e ci sono varie \textbf{metriche} e varie scale di
successo. A questo si aggiungono altri livelli:
\begin{itemize}
  \item \textbf{quale mercato}? In Italia \textbf{i giochi che non sono in
        italiano non vendono}, e molti editori italiani sono innanzitutto
        \textbf{editori di localizzazione}: prendono giochi dall'estero, li
        traducono, e questo è buona parte della loro sopravvivenza, nel gioco
        da tavolo come nel gioco di ruolo. Un gioco in inglese si vende in
        Inghilterra, in America e in moltissimi paesi dove l'inglese è la seconda
        lingua; un gioco in italiano si vende solo in Italia, e un editore
        italiano non ha facilmente la struttura per arrivare all'estero. Per
        questo anche autori italiani di giochi da tavolo pubblicano con
        \textbf{editori esteri};
  \item \textbf{perché si fa il gioco}? Non è detto che lo scopo sia
        commerciale. Un gioco può essere \textbf{formativo o didattico}, fatto
        con un'azienda o una scuola, o \textbf{commissionato} da un cliente per
        uno scopo specifico: allora non interessa che venda mille copie. Il
        designer viene pagato, ma il pagamento non corrisponde alle vendite.
\end{itemize}

Beatrice chiede ad Angelo perché faccia la domanda, perché i \textbf{settori}
appena visti spesso ``non si parlano minimamente'', e alcuni non sanno nemmeno
dell'esistenza degli altri: la cross-medialità presentata nella lezione non viene
presentata spesso così. Se la domanda nasce da prospettive lavorative, come per
alcuni nel corso, la risposta è che ci sono ``immense possibilità'' a seconda
della direzione: fare un gioco da tavolo è una cosa, lavorare nella
\textbf{formazione} un'altra, e c'è un intero ambito non trattato, quello dei
\textbf{giochi per le aziende}.

Angelo precisa: voleva solo sapere \textbf{che cosa giocano di solito gli
italiani}. È ``un altro volo pindarico'': in Italia è difficile trovare dati
statistici affidabili, perché i questionari e le survey sono spesso invalidati
dal modo in cui sono formulate le domande. Da quello che sa il docente che risponde\footnote{La trascrizione dice ``poi Lucio mi correggerà''; il nome non è ricostruibile con certezza.}, ci sono sicuramente \textbf{più giocatori di giochi da
tavolo che di giochi di ruolo}: il gioco da tavolo è più conosciuto e
frequentato, anche nelle ludoteche, mentre il gioco di ruolo e il LARP richiedono
un investimento economico diverso e sono soprattutto pratiche
\textbf{comunitarie}, ``tribali'', che nascono dal basso e si trasmettono
\textbf{oralmente}. Chi gioca ai giochi da tavolo li compra; chi gioca di ruolo
spesso lo fa perché qualcuno lo fa giocare, e non è detto che abbia il gioco in
casa.

\subsection{Avvisi}

Intanto Daniele ricorda che il \textbf{questionario} della prima slide va
compilato \textbf{entro il 15}, e che Giordano deve ancora presentarsi (``siamo
delle pessime persone''). Visto che è ora di cena, Giordano si presenterà come
\textbf{prima cosa} alla lezione successiva, \textbf{lunedì}, sull'emotional
design (``ricordacelo, perché se no ci dimentichiamo'').

% ---------------------------------------------------------------------
\section{Due cimiteri: i giochi di carte collezionabili e i ``fix''}
\label{a2:sec:cimiteri}

Prima di chiudere, a un docente torna in mente ``una cosa importantissima''
che voleva dire.

\subsection{Non fate Magic: The Gathering 2}

``Vi prego, non cadete nella trappola'' di vedere un gioco di successo commerciale
e farne uno simile pensando che venderà moltissimo. ``Ci possiamo fare un
cimitero, una fossa comune'' con migliaia di \textbf{trading card game} e
\textbf{collectible card game}. ``Non fate un trading card game. Ve lo dico per
voi: non fate \emph{Magic: The Gathering 2}''. Chi vede \emph{Magic} o
\emph{Pokémon} pensa di poter fare un sacco di soldi, ma \textbf{non vede la
struttura commerciale e produttiva gigantesca} che c'è dietro, quasi impossibile
da emulare dal basso. Quasi tutti, anche editori grossi, quelli che hanno provato
a competere con \emph{Pokémon}, \emph{Yu-Gi-Oh!} e \emph{Magic} ``sono andati a
picco''. ``Se volete perdere soldi, quello è il modo perfetto''.

Un collega aggiunge che c'è un \textbf{revival dei TCG}, ma ogni titolo dura un
anno o due e poi muore. Si sperimentano anche dei \textbf{micro TCG} annunciati
fin dall'inizio come limitati (``questo gioco avrà tre espansioni e poi
chiuderà''), una sorta di ibrido con gli LCG. Ma attenzione: spesso sono fatti
da ex designer di LCG, di \emph{Magic} o di \emph{Hearthstone}, persone con
contatti, entourage e un'azienda alle spalle; da fuori si vede una persona, ma
dietro c'è una struttura che non si può replicare.

\paragraph{Approfondimento: perché Pokémon non si replica.} \emph{The Art of
Game Design} offre un buon esempio della ``struttura gigantesca'' di cui parlano
i docenti. Nel capitolo sui mondi transmediali, Schell racconta che il gioco di
carte collezionabili di \emph{Pokémon} nacque nel 1999 da una collaborazione tra
Nintendo e Wizards of the Coast, la casa di \emph{Magic}, e che fu progettato per
restare il più possibile fedele alle meccaniche del videogioco per Game Boy. Il
gioco di carte era la \textbf{terza porta d'ingresso} in un unico mondo, accanto
al videogioco e alla serie TV: portatile, sociale, economico, perfetto per lo
slogan ``acchiappali tutti''. Le tre porte si rafforzavano a vicenda (voler
vincere spingeva a guardare la serie, e la serie rendeva il gioco più vivido), e
Schell ricorda un dirigente dell'intrattenimento che nel 1999 profetizzava la
fine della ``moda'' con l'uscita del film. Non andò così, appunto, perché il
successo non stava in un singolo prodotto ma nell'ecosistema. Un gioco di carte
collezionabili nato dal basso, senza un ecosistema del genere, compete con questo.

\subsection{Non fate un gioco per ``aggiustarne'' un altro}

Un altro docente (probabilmente Daniele, che cita ``le lezioni di framework con Ruggero'') aggiunge un altro cimitero: quello di chi vuole fare un gioco per
\textbf{``fixare'' un gioco esistente}, cioè farne una versione migliore. Come
esperimento può essere interessante, ma, come si vedrà nella lezione
sull'emotional design e nelle lezioni sui framework con Ruggero, l'intenzione deve
essere ``\textbf{eticamente onesta}''. Non si vuole aggiustare D\&D per il piacere
di aggiustare D\&D: bisogna chiedersi sempre \textbf{che cosa non mi piace
dell'esperienza} di D\&D, che cosa voglio aggiustare di quell'esperienza. E non
creare ``D\&D dove tolgo una regola e ne metto un'altra, ed è sistemato'', ma
\textbf{partire dall'esperienza che si vuole raggiungere}, quella che non si è
avuta da D\&D, e costruirla. Con la lezione di ricerca e sviluppo si vedrà anche
come prendere spunto dai giochi esistenti, non per copiarli ma per capire che
cosa fanno, perché \textbf{prenderne pezzi senza capire che cosa fanno} è proprio
l'errore da evitare (lo stesso del libro di Nuccio letto troppo presto).

\paragraph{I saluti.} Un docente deve scappare e saluta; gli altri ricordano il
questionario, che le loro email sono nelle slide per eventuali domande, e
l'appuntamento a lunedì.

% ---------------------------------------------------------------------
\section*{Punti chiave del capitolo}
\addcontentsline{toc}{section}{Punti chiave del capitolo}

\begin{itemize}
  \item Un \textbf{caso studio} è un gioco, o una categoria di giochi, che,
        analizzato tenendo conto del suo contesto, dà informazioni e metodi
        utili per i propri giochi: come leggere libri per imparare a scrivere.
        Questa lezione è una panoramica; gli strumenti di analisi arriveranno con
        la ricerca e sviluppo. ``Per me, questo è design'' (Ruggero).
  \item \textbf{Gioco analogico} è un termine ombrello per tutto ciò che non è
        videogioco. Le quattro \textbf{macrocategorie} sono gioco da tavolo,
        gioco di ruolo da tavolo, gioco di ruolo dal vivo e gioco pervasivo, ma
        i confini sono porosi: \textbf{ibridazione e sperimentazione}. Il corso è
        \textbf{agnostico di genere}, e il project work è libero, ma non deve
        essere un gioco enorme.
  \item Il \textbf{gioco pervasivo} usa luoghi e oggetti della vita quotidiana
        come luoghi e oggetti di gioco; in termini di game studies,
        \textbf{espande il cerchio magico} nello spazio, nel tempo e nella
        società. Casi: \emph{Pokémon Go}, geocaching, \emph{GuruShot},
        \emph{Journey to the End of the Night}, SF0, \emph{Critical City Upload},
        \emph{Sleep No More}, gli ARG (\emph{Why So Serious}, \emph{The
        Institute}, \emph{Stranger Things} a Milano), Collettivo Cinetico,
        Amigdala, \emph{MP3 Experiment}, \emph{PacManhattan}, \emph{Zombies,
        Run!}, \emph{Reptilarium}, \emph{Echi resistenti}.
  \item Esperienze emotivamente significative si possono fare \textbf{con budget
        zero}: l'audio che guida, gli oggetti da portare, un finale
        scenografico.
  \item Il \textbf{LARP} è separato dal GDR per ragioni di industria (evento con
        biglietto contro oggetto venduto). Tipi: boffer, Mind's Eye Theatre,
        chamber/parlor, American freeform, murder party, Nordic LARP (Jeepform,
        360°, tape, black box).
  \item Il \textbf{Nordic LARP} ha documentato tutto (Knutpunkt dal 1997, libri
        annuali, manifesti come il Dogma 99), punta su coinvolgimento, emozione,
        coesione tematica e dopo-gioco, e ha normalizzato le \textbf{meccaniche di
        sicurezza}. Non è un esperimento scientifico: lo scopo resta ludico.
  \item Nel \textbf{GDR da tavolo} contano la fiction emergente e quattro parole:
        \textbf{conflitto, azione, risoluzione, effetto}. Il master è un
        giocatore. Dalla nascita di D\&D (1974) dal wargame, attraverso gli
        adattamenti a nastro degli anni '80--'90, si arriva alla \textbf{New
        Wave} di The Forge (``il sistema conta''), ai \textbf{giochi
        focalizzati} (non semplificati), all'OSR, agli story game, ai giochi per
        due, in solitario e micro.
  \item Nel \textbf{gioco da tavolo} ogni genere si è evoluto: roll and move,
        astratti, party game, giochi di parole, giochi di carte (TCG, LCG,
        deck-building), German game ed Eurogame, American, strategici,
        movimento nascosto, destrezza, 18XX, tempo reale, \textbf{Legacy},
        libri-gioco, megagame, ibridi.
  \item Libri: il volume sull'\textbf{edu-larp} sì; Nuccio e Engelstein--Shalev
        dopo l'emotional design; Bertolo--Mariani con cautela sulla parte di GDR;
        gran parte della letteratura italiana sul GDR no. Soprattutto,
        \textbf{giocare}.
  \item Il \textbf{successo} non si misura solo in copie, e dipende dal mercato
        (in Italia vende solo ciò che è in italiano) e dallo scopo del gioco.
  \item Due cimiteri: \textbf{non fare un gioco di carte collezionabili} e non
        fare un gioco per \textbf{``aggiustarne'' un altro} senza partire
        dall'esperienza che si vuole creare.
\end{itemize}
